%% gapil.tex
%%
%% Copyright (C) 2000-2018 Simone Piccardi.  Permission is granted to
%% copy, distribute and/or modify this document under the terms of the GNU Free
%% Documentation License, Version 1.3 or any later version published by the
%% Free Software Foundation; with the Invariant Sections being "Un preambolo",
%% with no Front-Cover Texts, and with no Back-Cover Texts.  A copy of the
%% license is included in the section entitled "GNU Free Documentation
%% License".
%%
%% GaPiL : Guida alla Programmazione in Linux
%%
%% S. Piccardi Oct. 2000
%% 
%% gapil.tex: file principale, gli altri vanno inclusi da questo.
%%
\documentclass[10pt,twoside]{book}

%\usepackage[paperwidth=18.91cm,paperheight=24.589cm,
%inner=2.5cm,outer=1.8cm,bottom=3.3cm,top=2.3cm]{geometry}

\usepackage[paperwidth=18.91cm,paperheight=24.589cm,
vmargin=1.9cm,inner=1.9cm,outer=1.9cm,bindingoffset=0.89cm]{geometry}




% encodings
\usepackage[T1]{fontenc}
\usepackage[utf8]{inputenc}
\usepackage{ae,aecompl}

% more fonts
\usepackage{inconsolata}


%
% New header style
%
\usepackage{fancyhdr}
\pagestyle{fancy}
\fancyhf{}
\fancyhead[RE]{\slshape \rightmark}
\fancyhead[RO,LE]{\thepage}
\fancyhead[LO]{\slshape \leftmark}

\renewcommand{\chaptermark}[1]{% 
\markboth{\thechapter.% 
\ #1}{}} 

\renewcommand{\sectionmark}[1]{% 
\markright{\ \thesection% 
\ #1}{}}


% fancy verbatim
\usepackage{fancyvrb}
\DefineVerbatimEnvironment{Example}{Verbatim}
{xleftmargin=1cm,xrightmargin=1cm}

\DefineVerbatimEnvironment{FileExample}{Verbatim}
{frame=lines,framerule=0.5mm,framesep=2mm,xleftmargin=1cm,xrightmargin=1cm,fontsize=\footnotesize}

\DefineVerbatimEnvironment{Terminal}{Verbatim}
{xleftmargin=\parindent,xrightmargin=\parindent,fontsize=\footnotesize}

\DefineVerbatimEnvironment{Command}{Verbatim}
{xleftmargin=\parindent,xrightmargin=\parindent,fontseries=b,
fontsize=\footnotesize}

\DefineVerbatimEnvironment{Console}{Verbatim}
{commandchars=\\\{\},xleftmargin=\parindent,xrightmargin=\parindent,fontsize=\footnotesize}


\usepackage[bookmarks=true,plainpages=false,pdfpagelabels,
hyperfootnotes=false]{hyperref}

\usepackage{makeidx}
\usepackage{pst-node}
\usepackage{eso-pic}
\usepackage{graphicx}
\usepackage[italian]{babel}
\usepackage{cite}
\usepackage{amsmath}
\usepackage{amsfonts}
\usepackage[hang,bf,footnotesize,it]{caption}
\usepackage{listings}
\lstloadlanguages{C++}
\usepackage{color} 

\providecommand{\defaultdesc}{%
\desclabelstyle{\pushlabel}%
\renewcommand{\makelabel}[1]{##1}%
\setlength{\labelwidth}{0pt}%
}
\usepackage{mdwlist}

\usepackage{boxedminipage}
\usepackage{multirow}
\usepackage{longtable}
\usepackage{varwidth}
%\usepackage{footnote} 
%\usepackage{mdwtab} 

\usepackage{tikz}
\usetikzlibrary{shapes}

\frenchspacing

%\includeonly{macro,pref,intro,fdl}

\title{\Huge \textbf{\textsl{G}a\textsl{P}i\textsl{L}}\\
\huge Guida alla Programmazione in Linux} 

\author{Simone Piccardi}

\makeindex

% Solo prima parte, scommentare 
% \includeonly{macro,preambolo,pref,intro,process,prochand,filedir,
%   fileunix,filestd,system,signal,session,fileadv,ipc,errors,
%   ringraziamenti,fdl}

% Solo seconda parte e appendici, scommentare
%\includeonly{macro,preambolo,network,socket,tcpsock,sockctrl,othersock,
%  sockadv,netlayer,trasplayer,errors,ringraziamenti,fdl}

\begin{document}

\frontmatter

\maketitle

\newcommand{\clearemptydoublepage}{\newpage{\pagestyle{empty}
\cleardoublepage}}

\begin{quote}
  
  Copyright \copyright\ 2000-2018 Simone Piccardi.  Permission is granted to
  copy, distribute and/or modify this document under the terms of the GNU Free
  Documentation License, Version 1.3 or any later version published by the
  Free Software Foundation; with the Invariant Sections being ``Un preambolo''
  in ``Prefazione'', with no Front-Cover Texts, and with no Back-Cover Texts.
  A copy of the license is included in the section entitled ``GNU Free
  Documentation License''.

\end{quote}
\clearemptydoublepage

%\ifx\hcode\undefined\else
\tableofcontents
%\fi

\clearemptydoublepage
%% macro.tex
%%
%% Copyright (C) 2000-2018 Simone Piccardi.  Permission is granted to
%% copy, distribute and/or modify this document under the terms of the GNU Free
%% Documentation License, Version 1.1 or any later version published by the
%% Free Software Foundation; with the Invariant Sections being "Un preambolo",
%% with no Front-Cover Texts, and with no Back-Cover Texts.  A copy of the
%% license is included in the section entitled "GNU Free Documentation
%% License".
%%
%
%
% Defining some special character for use inside typewriter 
% text without using the verbatim environment
%
\def\tild{\char'176}
\def\bslash{\char'134}
\def\circonf{\char'136}
\def\invap{\char'140}

\newcommand{\includecodesnip}[1]{\lstinputlisting[stepnumber=0,xleftmargin=\parindent,frame=]{#1}}{}
\newcommand{\includestruct}[1]{\lstinputlisting[stepnumber=0]{#1}}{}
\newcommand{\includecodesample}[1]{\lstinputlisting{#1}}{}

%
% Defining some commands to manipulate counter to avoid ude of 
% \label and \ref commands (and related problem to remeber the 
% used labels) to refer nearest objects
%
%
\newcounter{usercount}       % define a new counter for internal use
%
% equations commands
%
\newcommand{\cureq}{(\theequation)}
\newcommand{\nxeq}{%
\setcounter{usercount}{\value{equation}}%
\addtocounter{usercount}{1}%
(\thechapter.\theusercount)}
\newcommand{\preeq}{%
\setcounter{usercount}{\value{equation}}%
\addtocounter{usercount}{-1}%
(\thechapter.\theusercount)}
%
% Macro to put picture (in format PICT) inside a figure
%
\newcommand{\pictfig}[3]{
        \begin{minipage}[t][#1][b]{#2}
        \mbox{\special{pict=#3}}
%    \vspace{3cm}
    \end{minipage}
}
%
% Macro to create a special environment for function prototypes
% boxed description
%
\newenvironment{prototype}[2]
{% defining what is done by \begin
  \nobreak
  \center
  \begin{boxedminipage}[c]{14cm}
  \footnotesize
    \begin{description*}{}{} 
    \item \texttt{\#include <#1>}
    \item \texttt{#2} \par
%   \begin{lstlisting}{}
% #1
% #2
%   \end{lstlisting}
}
{% defining what is done by \end
%  \end{list} 
%  \par 
\end{description*}
\end{boxedminipage}
\vspace{6pt}
\par
\normalsize 
}
\newenvironment{errlist}{\begin{basedescript}{\desclabelwidth{1.0cm}}}
{\end{basedescript}}
%
% Slighty different environment to be used for multi-header, 
% multi-functions boxed description
%
\newcommand{\headdecl}[1]{\item\texttt{\#include <#1>}}
\newcommand{\funcdecl}[1]{\item\texttt{#1}\par}
\newcommand{\bodydesc}[1]{\par \end{description*} #1
 \begin{description*}{}{} \baselineskip=0pt
 \item \vspace{-4pt}
} 
% 
\newenvironment{functions}
{% defining what is done by \begin
  \nobreak
  \center
    \begin{boxedminipage}[c]{14cm}
      \footnotesize
    \begin{description*}{}{} 
}
{% defining what is done by \end
      \end{description*}
    \end{boxedminipage}
    \vspace{6pt}
     \par
    \normalsize 
%\break
}

%
% Wrapper for shell command, functions, filenames, links,
% variables, macros,  and everything can be useul, 
%
\newcommand{\cmd}[1]{\texttt{#1}}     % shell command
\newcommand{\code}[1]{\texttt{#1}}    % for simple code

\newcommand{\myfunc}[1]{\texttt{#1}}    % for my functions

\newcommand{\func}[1]{%
\texttt{#1}%
%\index{funzione!{#1}@{{\tt {#1}}}}\texttt{#1}%
%\index{#1@{{\tt {#1}} (funzione)}}\texttt{#1}%
}

\newcommand{\funcd}[1]{%
\index{funzione!{#1}@{{\tt {#1}}}}\texttt{#1}%
%\index{funzione!{#1}@{{\tt {#1}}}!definizione di}\texttt{#1}%
%\index{#1@{{\tt {#1}} (funzione)}!definizione di}\texttt{#1}%
}
\newcommand{\funcm}[1]{%
\index{funzione!{#1}@{{\tt {#1}}}}\texttt{#1}%
%\index{funzione!{#1}@{{\tt {#1}}}!menzione di}\texttt{#1}%
%\index{#1@{{\tt {#1}} (funzione)}!definizione di}\texttt{#1}%
}

\newcommand{\macro}[1]{%
\texttt{#1}%
%\index{macro!{#1}@{{\tt {#1}}}}\texttt{#1}%
%\index{#1@{{\tt {#1}} (macro)}}\texttt{#1}%
}

\newcommand{\macrod}[1]{%
\index{macro!{#1}@{{\tt {#1}}}}\texttt{#1}%
%\index{#1@{{\tt {#1}} (macro)}}\texttt{#1}%
}

\newcommand{\macrobeg}[1]{%
\index{macro!{#1}@{{\tt {#1}}}|(}%
%\index{#1@{{\tt {#1}} (macro)}}\texttt{#1}%
}
\newcommand{\macroend}[1]{%
\index{macro!{#1}@{{\tt {#1}}}|)}%
%\index{#1@{{\tt {#1}} (macro)}}\texttt{#1}%
}

\newcommand{\errcode}[1]{%
\texttt{#1}%
%\index{errore!{#1}@{{\tt {#1}}}}\texttt{#1}%
%\index{#1@{{\tt {#1}} (errore)}}\texttt{#1}%
}

\newcommand{\errval}[1]{\texttt{#1}}     % value 
\newcommand{\var}[1]{\texttt{#1}}     % variable 
\newcommand{\val}[1]{\texttt{#1}}     % value 

\newcommand{\signal}[1]{%
\texttt{#1}%
%\index{segnale!{#1}@{{\tt {#1}}}}\texttt{#1}%
%\index{#1@{{\tt {#1}} (costante)}}\texttt{#1}%
}                                     % constant name

\newcommand{\signald}[1]{%
\index{segnale!{#1}@{{\tt {#1}}}}\texttt{#1}%
%\index{#1@{{\tt {#1}} (costante)}}\texttt{#1}%
}                                     % constant name

\newcommand{\constd}[1]{%
\index{costante!{#1}@{{\tt {#1}}}}\texttt{#1}%
%\index{#1@{{\tt {#1}} (costante)}}\texttt{#1}%
}                                     % constant name

\newcommand{\constbeg}[1]{%
\index{costante!{#1}@{{\tt {#1}}}|(}%
%\index{#1@{{\tt {#1}} (costante)}}\texttt{#1}%
}                                     % constant name
\newcommand{\constend}[1]{%
\index{costante!{#1}@{{\tt {#1}}}|)}%
%\index{#1@{{\tt {#1}} (costante)}}\texttt{#1}%
}                                     % constant name

\newcommand{\const}[1]{%
\texttt{#1}%
%\index{costante!{#1}@{{\tt {#1}}}}\texttt{#1}%
%\index{#1@{{\tt {#1}} (costante)}}\texttt{#1}%
}                                     % constant name

\newcommand{\instruction}[1]{%
\index{istruzione linguaggio C!{#1}@{{\tt {#1}}}}\texttt{#1}%
%\index{#1@{{\tt {#1}} (direttiva)}}\texttt{#1}%
}                                      % instruction name

\newcommand{\instr}[1]{\texttt{#1}}    % instruction

\newcommand{\direct}[1]{%
\index{direttiva linguaggio C!{#1}@{{\tt {#1}}}}\texttt{#1}%
%\index{#1@{{\tt {#1}} (direttiva)}}\texttt{#1}%
}                                     % directive name

\newcommand{\dirct}[1]{\texttt{#1}}   % directive


\newcommand{\file}[1]{\texttt{#1}}    % file name
\newcommand{\link}[1]{\texttt{#1}}    % html link
\newcommand{\ctyp}[1]{\texttt{#1}}    % C standard type

\newcommand{\headfile}[1]{%
\texttt{#1}%
}                                     % header file name
\newcommand{\headfiled}[1]{%
\index{file!include!{#1}@{{\tt {#1}}}}\texttt{#1}%
}                                     % header file name

\newcommand{\procfile}[1]{%
\index{file!filesystem~\texttt{/proc}!{#1}@{{\tt {#1}}}}\texttt{#1}%
}                                     % /proc file name
\newcommand{\procfilem}[1]{%
\texttt{#1}%
%\index{file!filesystem~\texttt{/proc}!{#1}@{{\tt {#1}}}}\texttt{#1}%
%\index{#1@{{\tt {#1}} (direttiva)}}\texttt{#1}%
}                                     % /proc file name

\newcommand{\sysfile}[1]{%
\texttt{#1}%
%\index{file!di~sistema!{#1}@{{\tt {#1}}}}\texttt{#1}%
}                                     % system file name
\newcommand{\sysfiled}[1]{%
\index{file!di~sistema!{#1}@{{\tt {#1}}}}\texttt{#1}%
}                                     % system file name

\newcommand{\conffile}[1]{%
\texttt{#1}%
%\index{file!di~configurazione!{#1}@{{\tt {#1}}}}\texttt{#1}%
}                                     % configuration file name
\newcommand{\conffiled}[1]{%
\index{file!di~configurazione!{#1}@{{\tt {#1}}}}\texttt{#1}%
}                                     % configuration file name
\newcommand{\conffilebeg}[1]{%
\index{file!di~configurazione!{#1}@{{\tt {#1}}}}%
}                                     % configuration file name
\newcommand{\conffileend}[1]{%
\index{file!di~configurazione!{#1}@{{\tt {#1}}}}%
}                                     % configuration file name

\newcommand{\procrelfile}[2]{%
\index{file!filesystem~\texttt{/proc}!{#1/#2}@{{\tt {#1/#2}}}}\texttt{#2}%
}   

\newcommand{\sysctlfiled}[1]{%
\index{file!file di controllo (sotto \texttt{/proc/sys})!{#1}@{{\tt {#1}}}}\texttt{/proc/sys/#1}%
}                                     % /proc/sys file 

\newcommand{\sysctlfile}[1]{%
\texttt{/proc/sys/#1}%
%\index{file!file di controllo (sotto \texttt{/proc/sys})!{#1}@{{\tt {#1}}}}\texttt{/proc/sys/#1}%
}                                     % /proc/sys file 

\newcommand{\sysctlrelfile}[2]{%
\texttt{#2}%
%\index{file!file di controllo (sotto \texttt{/proc/sys})!{#1/#2}@{{\tt {#1/#2}}}}\texttt{#2}%
}                                     % /proc/sys file name

\newcommand{\sysctlrelfiled}[2]{%
\index{file!file di controllo (sotto \texttt{/proc/sys})!{#1/#2}@{{\tt {#1/#2}}}}\texttt{#2}%
}                                     % /proc/sys file name

\newcommand{\kstruct}[1]{%
\texttt{#1}%
%\index{struttura dati del kernel!{#1}@{{\tt {#1}}}}\texttt{#1}%
}                                     % struttura dati del kernel
\newcommand{\kstructd}[1]{%
\index{struttura dati del kernel!{#1}@{{\tt {#1}}}}\texttt{#1}%
%\index{struttura dati del kernel!{#1}@{{\tt {#1}}}!definizione di}\texttt{#1}%
}                                     % struttura dati

\newcommand{\typed}[1]{%
\index{tipo di dato!{#1}@{{\tt {#1}}}}\texttt{#1}%
%\index{#1@{{\tt {#1}} (tipo)}}\texttt{#1}%
}                                     % system type
\newcommand{\type}[1]{%
\texttt{#1}%
%\index{tipo di dato!{#1}@{{\tt {#1}}}}\texttt{#1}%
%\index{#1@{{\tt {#1}} (tipo)}}\texttt{#1}%
}                                     % system type
\newcommand{\struct}[1]{%
\texttt{#1}%
%\index{struttura dati!{#1}@{{\tt {#1}}}}\texttt{#1}%
%\index{#1@{{\tt {#1}} (struttura dati)}}\texttt{#1}%
}                                     % struttura dati
\newcommand{\structd}[1]{%
\index{struttura dati!{#1}@{{\tt {#1}}}}\texttt{#1}%
%\index{struttura dati!{#1}@{{\tt {#1}}}!definizione di}\texttt{#1}%
}                                     % struttura dati
\newcommand{\param}[1]{\texttt{#1}}   % function parameter
\newcommand{\acr}[1]{\textsl{#1}}     % acrostic (for glibc, ext2, ecc.)

\newcommand{\ids}[1]{\textsl{#1}}     % Identifier (PID, GID, UID, TID, ecc.)
\newcommand{\envvar}[1]{\texttt{#1}}   % environment variable
\newcommand{\samplefunc}[1]{\texttt{#1}}   % funzione degli esempi

\newcommand{\itindex}[1]{%
\index{#1@{\textit{#1}}}%
}

\newcommand{\itindbeg}[1]{%
\index{#1@{\textit{#1}}|(}%
}
\newcommand{\itindend}[1]{%
\index{#1@{\textit{#1}}|)}%
}
\newcommand{\itindsub}[2]{%
\index{#1@{\textit{#1}}!\textit{#2}}%
}
\newcommand{\itindsubbeg}[2]{%
\index{#1@{\textit{#1}}!\textit{#2}}%
}
\newcommand{\itindsubend}[2]{%
\index{#1@{\textit{#1}}!\textit{#2}}%
}

% Aggiunte di Mirko per la gestione delle tabelle complicate come immagini
% nella traslazione in HTML
\newenvironment{usepicture}{}{}{}{}

%
% Macro di definizione di alcune lunghezze usate per le formattazioni
%
\newlength{\codesamplewidth}
\setlength{\codesamplewidth}{0.95\textwidth}

\newlength{\funcboxwidth}
\setlength{\funcboxwidth}{0.85\textwidth}


%
% Nuove macro per diversa formattazione delle definizioni delle funzioni
%
\newcommand{\fhead}[1]{\texttt{\#include <#1>}\par}
\newcommand{\fdecl}[1]{\texttt{#1}\par}
\newcommand{\fdesc}[1]{\hfill{#1}\par}

\newenvironment{funcproto}[2]
{% defining what is done by \begin
\centering
\vspace{3pt}
\begin{funcbox}
#1
\end{funcbox}
\begin{funcbox}
#2
}
{% defining what is done by \end
\end{funcbox}
%\vspace{6pt}
%\break
}

\newenvironment{funcbox}
{% defining what is done by \begin
  \nobreak
  \begin{boxedminipage}[c]{\funcboxwidth}
    \footnotesize
}
{% defining what is done by \end
\end{boxedminipage}
\break
\normalsize
\par
}

\newcommand{\unavref}[1]{\ifx\incomplete\undefined\relax\else #1\fi}


%%% Local Variables: 
%%% mode: latex
%%% TeX-master: "gapil"
%%% End: 

%% preambolo.tex
%%
%% Copyright (C) 2000-2018 Simone Piccardi.  Permission is granted to
%% copy, distribute and/or modify this document under the terms of the GNU Free
%% Documentation License, Version 1.1 or any later version published by the
%% Free Software Foundation; with the Invariant Sections being "Un preambolo",
%% with no Front-Cover Texts, and with no Back-Cover Texts.  A copy of the
%% license is included in the section entitled "GNU Free Documentation
%% License".
%%

\chapter{Un preambolo}
\label{cha:preamble}

Questa guida nasce dalla mia profonda convinzione che le istanze di libertà e
di condivisione della conoscenza che hanno dato vita a quello straordinario
movimento di persone ed intelligenza che va sotto il nome di \textsl{software
  libero} hanno la stessa rilevanza anche quando applicate alla produzione
culturale in genere.

L'ambito più comune in cui questa filosofia viene applicata è quello della
documentazione perché il software, per quanto possa essere libero, se non
accompagnato da una buona documentazione che aiuti a comprenderne il
funzionamento, rischia di essere fortemente deficitario riguardo ad una delle
libertà fondamentali, quella di essere studiato e migliorato.

Ritengo inoltre che in campo tecnico ed educativo sia importante poter
disporre di testi didattici (come manuali, enciclopedie, dizionari, ecc.) in
grado di crescere, essere adattati alle diverse esigenze, modificati e
ampliati, o anche ridotti per usi specifici, nello stesso modo in cui si fa
per il software libero.

Questa guida è il mio tentativo di restituire indietro, nei limiti di quelle
che sono le mie capacità, un po' della conoscenza che ho ricevuto, mettendo a
disposizione un testo che possa fare da riferimento a chi si avvicina alla
programmazione su Linux, nella speranza anche di trasmettergli non solo delle
conoscenze tecniche, ma anche un po' di quella passione per la libertà e la
condivisione della conoscenza che sono la ricchezza maggiore che ho ricevuto.

E, come per il software libero, anche in questo caso è importante la
possibilità di accedere ai sorgenti (e non solo al risultato finale, sia
questo una stampa o un file formattato) e la libertà di modificarli per
apportarvi migliorie, aggiornamenti, ecc.

Per questo motivo la Free Software Foundation ha creato una apposita licenza
che potesse giocare lo stesso ruolo fondamentale che la GPL ha avuto per il
software libero nel garantire la permanenza delle libertà date, ma potesse
anche tenere conto delle differenze che comunque ci sono fra un testo ed un
programma.

Una di queste differenze è che in un testo, come in questa sezione, possono
venire espresse quelle che sono le idee ed i punti di vista dell'autore, e
mentre trovo che sia necessario permettere cambiamenti nei contenuti tecnici,
che devono essere aggiornati e corretti, non vale lo stesso per l'espressione
delle mie idee contenuta in questa sezione, che ho richiesto resti invariata.

Il progetto pertanto prevede il rilascio della guida con licenza GNU FDL, ed
una modalità di realizzazione aperta che permetta di accogliere i contributi
di chiunque sia interessato. Tutti i programmi di esempio sono rilasciati con
licenza GNU GPL.




%%% Local Variables: 
%%% mode: latex
%%% TeX-master: "gapil"
%%% End: 

% LocalWords:  Foundation FDL

%% pref.tex
%%
%% Copyright (C) 2000-2018 Simone Piccardi.  Permission is granted to
%% copy, distribute and/or modify this document under the terms of the GNU Free
%% Documentation License, Version 1.1 or any later version published by the
%% Free Software Foundation; with the Invariant Sections being "Un preambolo",
%% with no Front-Cover Texts, and with no Back-Cover Texts.  A copy of the
%% license is included in the section entitled "GNU Free Documentation
%% License".
%%

\chapter{Prefazione}
\label{cha:preface}

Questo progetto mira alla stesura di un testo il più completo e chiaro
possibile sulla programmazione di sistema su un kernel Linux.  Essendo i
concetti in gran parte gli stessi, il testo dovrebbe restare valido anche per
la programmazione in ambito di sistemi Unix generici, ma resta una attenzione
specifica alle caratteristiche peculiari del kernel Linux e delle versioni
delle librerie del C in uso con esso; in particolare si darà ampio spazio alla
versione realizzata dal progetto GNU, le cosiddette \textit{GNU C Library} o
\textsl{glibc}, che ormai sono usate nella stragrande maggioranza dei casi,
senza tralasciare, là dove note, le differenze con altre implementazioni come
le \textsl{libc5} o le \textsl{uclib}.

L'obiettivo finale di questo progetto è quello di riuscire a ottenere un testo
utilizzabile per apprendere i concetti fondamentali della programmazione di
sistema della stessa qualità dei libri del compianto R. W. Stevens (è un
progetto molto ambizioso ...).

Infatti benché le pagine di manuale del sistema (quelle che si accedono con il
comando \cmd{man}) e il manuale delle librerie del C GNU siano una fonte
inesauribile di informazioni (da cui si è costantemente attinto nella stesura
di tutto il testo) la loro struttura li rende totalmente inadatti ad una
trattazione che vada oltre la descrizione delle caratteristiche particolari
dello specifico argomento in esame (ed in particolare lo \textit{GNU C Library
  Reference Manual} non brilla per chiarezza espositiva).

Per questo motivo si è cercato di fare tesoro di quanto appreso dai testi di
R. W. Stevens (in particolare \cite{APUE} e \cite{UNP1}) per rendere la
trattazione dei vari argomenti in una sequenza logica il più esplicativa
possibile, corredando il tutto, quando possibile, con programmi di esempio.

Dato che sia il kernel che tutte le librerie fondamentali di GNU/Linux sono
scritte in C, questo sarà il linguaggio di riferimento del testo. In
particolare il compilatore usato per provare tutti i programmi e gli esempi
descritti nel testo è lo GNU GCC. Il testo presuppone una conoscenza media del
linguaggio, e di quanto necessario per scrivere, compilare ed eseguire un
programma.

Infine, dato che lo scopo del progetto è la produzione di un libro, si è
scelto di usare \LaTeX\ come "ambiente di sviluppo" del medesimo, sia per
l'impareggiabile qualità tipografica ottenibile, che per la congruenza dello
strumento con il fine, tanto sul piano pratico, quanto su quello filosofico.

Il testo sarà, almeno inizialmente, in italiano. Per il momento lo si è
suddiviso in due parti, la prima sulla programmazione di sistema, in cui si
trattano le varie funzionalità disponibili per i programmi che devono essere
eseguiti su una singola macchina, la seconda sulla programmazione di rete, in
cui si trattano le funzionalità per eseguire programmi che mettono in
comunicazione macchine diverse.



%%% Local Variables: 
%%% mode: latex
%%% TeX-master: "gapil"
%%% End: 

% LocalWords:  kernel Library glibc Stevens Reference Manual GCC

\clearemptydoublepage
\mainmatter

\lstset{language=C++} 

\lstset{basicstyle=\small,
  numbers=left,
  numberstyle=\tiny,
  stepnumber=1,
  numbersep=2pt, 
  frame=TB,
  framesep=5pt,
%  xleftmargin=0.3cm,
  basicstyle=\ttfamily,
  keywordstyle=\color{blue}\ttfamily,
  ndkeywordstyle=\color{yellow}\ttfamily,
%  rdkeywordstyle=\ttfamily,
  identifierstyle=\ttfamily,
  commentstyle=\color{red}\ttfamily,
  stringstyle=\color{green}\ttfamily,
%  texcsststyle=\ttfamily,
  directivestyle=\color{magenta}\ttfamily
} 

% distance from margins for boxedminipage
%\fboxsep=6pt

% less hspace for tables
\setlength{\tabcolsep}{0.5em}

% comment following to have reference to incomplete section using \unavref
% command removed
\def\incomplete{\relax}

\part{Programmazione di sistema}
\label{part:progr-di-sist}
%% intro.tex
%%
%% Copyright (C) 2000-2018 Simone Piccardi.  Permission is granted to
%% copy, distribute and/or modify this document under the terms of the GNU Free
%% Documentation License, Version 1.1 or any later version published by the
%% Free Software Foundation; with the Invariant Sections being "Un preambolo",
%% with no Front-Cover Texts, and with no Back-Cover Texts.  A copy of the
%% license is included in the section entitled "GNU Free Documentation
%% License".
%%

\chapter{L'architettura del sistema}
\label{cha:intro_unix}

In questo primo capitolo sarà fatta un'introduzione ai concetti generali su
cui è basato un sistema operativo di tipo Unix come GNU/Linux, in questo modo
potremo fornire una base di comprensione mirata a sottolineare le peculiarità
del sistema che sono più rilevanti per quello che riguarda la programmazione.

Dopo un'introduzione sulle caratteristiche principali di un sistema di tipo
Unix passeremo ad illustrare alcuni dei concetti base dell'architettura di
GNU/Linux (che sono comunque comuni a tutti i sistemi \textit{unix-like}) ed
introdurremo alcuni degli standard principali a cui viene fatto riferimento.


\section{Una panoramica}
\label{sec:intro_unix_struct}

In questa prima sezione faremo una breve panoramica sull'architettura di un
sistema operativo di tipo Unix, come GNU/Linux, e della relazione fra le varie
parti che lo compongono. Chi avesse già una conoscenza di questa materia può
tranquillamente saltare questa sezione.

\subsection{Concetti base}
\label{sec:intro_base_concept}

Il concetto principale su cui è basata l'architettura di un sistema unix-like
è quello di un nucleo del sistema, il cosiddetto \textit{kernel} (nel nostro
caso Linux) a cui si demanda la gestione delle risorse della propria macchina
(la CPU, la memoria, le periferiche) mentre tutto il resto, quindi anche la
parte che prevede l'interazione con l'utente, dev'essere realizzato tramite
programmi eseguiti dal kernel, che accedano alle risorse tramite opportune
richieste a quest'ultimo.

Fin dai suoi albori Unix nasce come sistema operativo \textit{multitasking},
cioè in grado di eseguire contemporaneamente più programmi, e multiutente, in
cui è possibile che più utenti siano connessi ad una macchina eseguendo più
programmi ``\textsl{in contemporanea}''. In realtà, almeno per le macchine a
processore singolo, i programmi vengono semplicemente eseguiti uno alla volta
in una opportuna \textsl{rotazione}.\footnote{anche se oggi, con la presenza
  di sistemi multiprocessore, si possono avere più processi eseguiti in
  contemporanea, il concetto di ``\textsl{rotazione}'' resta comunque valido,
  dato che in genere il numero di processi da eseguire eccede il numero dei
  precessori disponibili. }

% Questa e` una distinzione essenziale da capire,
%specie nei confronti dei sistemi operativi successivi, nati per i personal
%computer (e quindi per un uso personale), sui quali l'hardware (allora
%limitato) non consentiva la realizzazione di un sistema evoluto come uno unix.

I kernel Unix più recenti, come Linux, sono realizzati sfruttando alcune
caratteristiche dei processori moderni come la gestione hardware della memoria
e la modalità protetta. In sostanza con i processori moderni si può
disabilitare temporaneamente l'uso di certe istruzioni e l'accesso a certe
zone di memoria fisica.  Quello che succede è che il kernel è il solo
programma ad essere eseguito in modalità privilegiata, con il completo accesso
a tutte le risorse della macchina, mentre i programmi normali vengono eseguiti
in modalità protetta senza accesso diretto alle risorse.  Uno schema
elementare della struttura del sistema è riportato in
fig.~\ref{fig:intro_sys_struct}.

\begin{figure}[htb]
  \centering
  \includegraphics[width=11cm]{img/struct_sys}
  % \begin{tikzpicture}
  %   \filldraw[fill=black!20] (0,0) rectangle (7.5,1);
  %   \draw (3.75,0.5) node {\textsl{System Call Interface}};
  %   \filldraw[fill=black!35] (0,1) rectangle (7.5,4);
  %   \draw (3.75,2.5) node {\huge{\textsf{kernel}}};
  %   \filldraw[fill=black!20] (0,4) rectangle (2.5,5);
  %   \draw (1.25,4.5) node {\textsf{scheduler}};
  %   \filldraw[fill=black!20] (2.5,4) rectangle (5,5);
  %   \draw (3.75,4.5) node {\textsf{VM}};
  %   \filldraw[fill=black!20] (5,4) rectangle (7.5,5);
  %   \draw (6.25,4.5) node {\textsf{driver}};

  %   \draw (1.25,7) node(cpu) [ellipse,draw] {\textsf{CPU}};
  %   \draw (3.75,7) node(mem) [ellipse,draw] {\textsf{memoria}};
  %   \draw (6.25,7) node(disk) [ellipse,draw] {\textsf{disco}};

  %   \draw[<->] (cpu) -- (1.25,5);
  %   \draw[<->] (mem) -- (3.75,5);
  %   \draw[<->] (disk) -- (6.25,5);

  %   \draw (7.5,0) node [anchor=base west] {\textit{kernel space}};
  %   \draw (7.5,-1) node [anchor=west] {\textit{user space}};

  %   \draw (-1,-0.5) -- (8.5, -0.5);

  %   \draw (0,-2) rectangle (7.5,-1);
  %   \draw (3.75, -1.5) node {\textsl{GNU C Library}};
  %   \draw[->] (1.25,-1) -- (1.25,0);
  %   \draw[->] (3.75,-1) -- (3.75,0);
  %   \draw[->] (6.25,-1) -- (6.25,0);

  %   \draw (1.25,-3) node(proc1) [rectangle,draw] {\textsf{processo}};
  %   \draw (3.75,-3) node(proc2) [rectangle,draw] {\textsf{processo}};
  %   \draw (6.25,-3) node(proc3) [rectangle,draw] {\textsf{processo}};

  %   \draw[->] (1.25,-2) -- (proc1);
  %   \draw[->] (3.75,-2) -- (proc2);
  %   \draw[->] (6.25,-2) -- (proc3);
  % \end{tikzpicture}
  \caption{Schema di massima della struttura di interazione fra processi,
    kernel e dispositivi in Linux.}
  \label{fig:intro_sys_struct}
\end{figure}

\itindbeg{scheduler}

Una parte del kernel, lo \textit{scheduler}, si occupa di stabilire, sulla
base di un opportuno calcolo delle priorità e con una suddivisione appropriata
del tempo di processore, quali fra i vari ``\textsl{processi}'' presenti nel
sistema deve essere eseguito, realizzando il cosiddetto
\itindex{preemptive~multitasking} \textit{preemptive
  multitasking}.\footnote{si chiama così quella gestione del
  \textit{multitasking} in cui è il kernel a decidere a chi assegnare l'uso
  della CPU, potendo interrompere l'esecuzione di un processo in qualunque
  momento.}  Ogni processo verrà comunque eseguito in modalità protetta;
quando necessario esso potrà accedere alle risorse della macchina soltanto
attraverso delle ``\textsl{chiamate al sistema}'' (vedi
sez.~\ref{sec:intro_syscall}) che restituiranno il controllo al kernel per
eseguire le operazioni necessarie.

\itindend{scheduler}

La memoria viene sempre gestita dal kernel attraverso il meccanismo della
\textsl{memoria virtuale}, che consente di assegnare a ciascun processo uno
spazio di indirizzi ``\textsl{virtuale}'' (vedi sez.~\ref{sec:proc_memory})
che il kernel stesso, con l'ausilio della unità di gestione della memoria, si
incaricherà di rimappare automaticamente sulla memoria fisica disponibile, con
la possibilità ulteriore di spostare temporaneamente su disco (nella
cosiddetta area di \textit{swap}) parte di detta memoria qualora ci si trovi
nella necessità di liberare risorse.

Le periferiche infine vengono normalmente viste attraverso un'interfaccia
astratta che permette di trattarle come se fossero dei file, secondo uno dei
concetti base della architettura di Unix, per cui ``\textsl{tutto è in file}''
(\textit{everything is a file}) su cui torneremo in
sez.~\ref{sec:intro_file_dir}. In realtà questo non è sempre vero (ad esempio
non lo è per le interfacce di rete) dato che ci sono periferiche che non
rispondendo bene a questa astrazione richiedono un'interfaccia diversa.  Anche
in questo caso però resta valido il concetto generale che tutto il lavoro di
accesso e gestione delle periferiche a basso livello viene effettuato dal
kernel tramite l'opportuno codice di gestione delle stesse, che in
fig.~\ref{fig:intro_sys_struct} si è indicato come \textit{driver}.


\subsection{Il kernel e il sistema}
\label{sec:intro_kern_and_sys}

Uno dei concetti fondamentali su cui si basa l'architettura dei sistemi Unix è
quello della distinzione fra il cosiddetto \textit{user space}, che
contraddistingue l'ambiente in cui vengono eseguiti i programmi, e il
\textit{kernel space}, che è l'ambiente in cui viene eseguito il kernel. Ogni
programma vede sé stesso come se avesse la piena disponibilità della CPU e
della memoria ed è, salvo i meccanismi di comunicazione previsti dal sistema,
completamente ignaro del fatto che altri programmi possono essere messi in
esecuzione dal kernel.

Per questa separazione non è possibile ad un singolo programma disturbare
l'azione di un altro programma o del kernel stesso, e questo è il principale
motivo della stabilità di un sistema unix-like nei confronti di altri sistemi
in cui i processi non hanno di questi limiti o in cui essi vengono eseguiti
allo stesso livello del kernel. Pertanto deve essere chiaro a chi programma in
un sistema unix-like che l'accesso diretto all'hardware non può avvenire se
non all'interno del kernel; al di fuori dal kernel il programmatore deve usare
le opportune interfacce che quest'ultimo fornisce per i programmi in
\textit{user space}.

Per capire meglio la distinzione fra \textit{kernel space} e \textit{user
  space} si può prendere in esame la procedura di avvio di un sistema
unix-like. All'accensione il \textit{firmware} presente nella EPROM della
propria macchina (per i PC compatibili il BIOS), eseguirà la procedura di
avvio del sistema, il cosiddetto \textit{bootstrap},\footnote{il nome deriva
  da un'espressione gergale che significa ``sollevarsi da terra tirandosi per
  le stringhe delle scarpe'', per indicare il compito, almeno apparentemente
  impossibile, di far eseguire un programma a partire da un computer appena
  acceso che appunto non ne contiene nessuno; non è impossibile perché in
  realtà c'è un programma iniziale, che è il BIOS.} incaricandosi di caricare
il kernel in memoria e di farne partire l'esecuzione. 

A questo punto il controllo passerà al kernel, il quale però da parte sua, una
volta inizializzato opportunamente l'hardware, si limiterà a due sole
operazioni, montare il filesystem radice (torneremo su questo in
sez.~\ref{sec:file_arch_overview}) e lanciare il primo processo che eseguirà
il programma di inizializzazione del sistema, che in genere, visto il suo
scopo, si chiama \cmd{init}.

Una volta lanciato \cmd{init} tutto il lavoro successivo verrà eseguito
\textit{user space} da questo programma, che sua volta si incaricherà di
lanciare tutti gli altri programmi, fra cui ci sarà quello che si occupa di
dialogare con la tastiera e lo schermo della console, quello che mette a
disposizione un terminale e la \textit{shell} da cui inviare i comandi
all'utente che si vuole collegare, ed in generale tutto quanto necessario ad
avere un sistema utilizzabile.

E' da rimarcare come tutto ciò che riguarda l'interazione con l'utente, che
usualmente viene visto come parte del sistema, non abbia in realtà niente a
che fare con il kernel, ma sia effettuato da opportuni programmi che vengono
eseguiti, allo stesso modo di un qualunque programma di scrittura o di disegno
e della stessa interfaccia grafica, in \textit{user space}.

Questo significa ad esempio che il sistema di per sé non dispone di primitive
per tutta una serie di operazioni (ad esempio come la copia di un file) che
altri sistemi (come Windows) hanno invece al loro interno. Questo perché tutte
le operazioni di normale amministrazione di un sistema, sono effettuata
attraverso dei normali programmi utilizzando le interfacce di programmazione
che il kernel mette a disposizione.

È per questo motivo che quando ci si riferisce al sistema nella sua interezza
viene spesso sottolineato come sia corretto parlare di ``GNU/Linux'' e non di
Linux; da solo infatti il kernel non è sufficiente, quello che costruisce un
sistema operativo utilizzabile è la presenza di tutta una serie di librerie e
programmi di utilità, ed i più comuni sono appunto quelli realizzati dal
progetto GNU della Free Software Foundation, grazie ai quali si possono
eseguire le normali operazioni che ci si aspetta da un sistema operativo.


\subsection{\textit{System call} e funzioni di libreria}
\label{sec:intro_syscall}

Come illustrato in fig.~\ref{fig:intro_sys_struct} i programmi possono
accedere ai servizi forniti dal kernel tramite opportune interfacce dette
\textit{system call} (\textsl{chiamate al sistema}, appunto). Si tratta di un
insieme di funzioni che un programma può invocare, per le quali viene generata
un'interruzione nell'esecuzione del codice del processo, passando il controllo
al kernel. Sarà quest'ultimo che eseguirà in le operazioni relative alla
funzione richiesta in \textit{kernel space}, restituendo poi i risultati al
chiamante.

Ogni versione di Unix ha storicamente sempre avuto un certo numero di
\textit{system call}, che sono documentate nella seconda sezione del
\textsl{Manuale di programmazione di Unix}, quella cui si accede con il
comando \cmd{man 2 <nome>}, ed anche Linux non fa eccezione. Queste
\textit{system call} sono poi state codificate da vari standard, che
esamineremo brevemente in sez.~\ref{sec:intro_standard}.

Normalmente ciascuna \textit{system call} fornita dal kernel viene associata
ad una funzione con lo stesso nome definita all'interno della libreria
fondamentale del sistema, quella che viene chiamata \textsl{Libreria Standard
  del C} (\textit{C Standard Library}) in ragione del fatto che il primo
kernel Unix e tutti i programmi eseguiti su di esso vennero scritti in questo
linguaggio, usando le sue librerie. In seguito faremo riferimento alle
funzioni di questa libreria che si interfacciano alle \textit{system call}
come ``\textsl{funzioni di sistema}''.

Questa libreria infatti, oltre alle chiamate alle \textit{system call},
contiene anche tutta una serie di ulteriori funzioni di utilità che vengono
comunemente usate nella programmazione e sono definite nei vari standard che
documentano le interfacce di programmazione di un sistema unix-like. Questo
concetto è importante da tener presente perché programmare in Linux significa
anche essere in grado di usare le funzioni fornite dalla libreria Standard del
C, in quanto né il kernel né il linguaggio C implementano direttamente
operazioni ordinarie come l'allocazione dinamica della memoria, l'input/output
bufferizzato sui file, la manipolazione delle stringhe, la matematica in
virgola mobile, che sono comunemente usate da qualunque programma.

Tutto ciò mette nuovamente in evidenza il fatto che nella stragrande
maggioranza dei casi si dovrebbe usare il nome GNU/Linux in quanto una parte
essenziale del sistema, senza la quale niente funzionerebbe, è appunto la
\textit{GNU Standard C Library} (a cui faremo da qui in avanti riferimento
come \acr{glibc}), la libreria standard del C realizzata dalla Free Software
Foundation nella quale sono state implementate tutte le funzioni essenziali
definite negli standard POSIX e ANSI C (e molte altre), che vengono utilizzate
da qualunque programma.

Si tenga comunque presente che questo non è sempre vero, dato che esistono
implementazioni alternative della libreria standard del C, come la
\textit{libc5} o la \textit{uClib}, che non derivano dal progetto GNU. La
\textit{libc5}, che era usata con le prime versioni del kernel Linux, è oggi
ormai completamente soppiantata dalla \acr{glibc}. La \textit{uClib} invece,
pur non essendo completa come la \acr{glibc}, resta molto diffusa nel mondo
dei dispositivi \textit{embedded} per le sue dimensioni estremamente ridotte,
e soprattutto per la possibilità di togliere le parti non necessarie. Pertanto
costituisce un valido rimpiazzo della \acr{glibc} in tutti quei sistemi
specializzati che richiedono una minima occupazione di memoria. Infine per lo
sviluppo del sistema Android è stata realizzata da Google un'altra libreria
standard del C, utilizzata principalmente per evitare l'uso della \acr{glibc}.

Tradizionalmente le funzioni specifiche della libreria standard del C sono
riportate nella terza sezione del \textsl{Manuale di Programmazione di Unix}
(cioè accessibili con il comando \cmd{man 3 <nome>}) e come accennato non sono
direttamente associate ad una \textit{system call} anche se, ad esempio per la
gestione dei file o della allocazione dinamica della memoria, possono farne
uso nella loro implementazione.  Nonostante questa distinzione, fondamentale
per capire il funzionamento del sistema, l'uso da parte dei programmi di una
di queste funzioni resta lo stesso, sia che si tratti di una funzione interna
della libreria che di una \textit{system call}.


\subsection{Un sistema multiutente}
\label{sec:intro_multiuser}

Linux, come gli altri kernel Unix, nasce fin dall'inizio come sistema
multiutente, cioè in grado di fare lavorare più persone in contemporanea. Per
questo esistono una serie di meccanismi di sicurezza, che non sono previsti in
sistemi operativi monoutente, e che occorre tenere presenti. In questa sezione
parleremo brevemente soltanto dei meccanismi di sicurezza tradizionali di un
sistema unix-like, oggi molti di questi sono stati notevolmente estesi
rispetto al modello tradizionale, ma per il momento ignoreremo queste
estensioni.

Il concetto base è quello di utente (\textit{user}) del sistema, le cui
capacità sono sottoposte a limiti precisi.  Sono così previsti una serie di
meccanismi per identificare i singoli utenti ed una serie di permessi e
protezioni per impedire che utenti diversi possano danneggiarsi a vicenda o
danneggiare il sistema. Questi meccanismi sono realizzati dal kernel stesso ed
attengono alle operazioni più varie, e torneremo su di essi in dettaglio più
avanti.

Normalmente l'utente è identificato da un nome (il cosiddetto
\textit{username}), che ad esempio è quello che viene richiesto all'ingresso
nel sistema dalla procedura di \textit{login} (torneremo su questo in
sez.~\ref{sec:sess_login}).  Questa procedura si incarica di verificare
l'identità dell'utente, in genere attraverso la richiesta di una parola
d'ordine (la \textit{password}), anche se sono possibili meccanismi
diversi.\footnote{ad esempio usando la libreria PAM (\textit{Pluggable
    Autentication Methods}) è possibile generalizzare i meccanismi di
  autenticazione e sostituire ad esempio l'uso delle password con meccanismi
  di identificazione biometrica; per un approfondimento dell'argomento si
  rimanda alla sez.~4.3.7 di \cite{AGL}.} Eseguita la procedura di
riconoscimento in genere il sistema manda in esecuzione un programma di
interfaccia (che può essere la \textit{shell} su terminale o un'interfaccia
grafica) che mette a disposizione dell'utente un meccanismo con cui questo può
impartire comandi o eseguire altri programmi.

Ogni utente appartiene anche ad almeno un gruppo (il cosiddetto
\textit{default group}), ma può essere associato ad altri gruppi (i
\textit{supplementary group}), questo permette di gestire i permessi di
accesso ai file e quindi anche alle periferiche, in maniera più flessibile,
definendo gruppi di lavoro, di accesso a determinate risorse, ecc. 

L'utente e il gruppo sono identificati dal kernel un identificativo numerico,
la cui corrispondenza ad un nome espresso in caratteri è inserita nei due file
\conffile{/etc/passwd} e \conffile{/etc/group}.\footnote{in realtà negli
  sistemi più moderni, come vedremo in sez.~\ref{sec:sys_user_group} queste
  informazioni possono essere mantenute, con l'uso del \textit{Name Service
    Switch}, su varie tipologie di supporti, compresi server centralizzati
  come LDAP.}  Questi identificativi sono l'\textit{user identifier}, detto in
breve \textsl{user-ID}, ed indicato dall'acronimo \ids{UID}, e il
\textit{group identifier}, detto in breve \textsl{group-ID}, ed identificato
dall'acronimo \ids{GID}, torneremo in dettaglio su questo argomento in
sez.~\ref{sec:proc_perms}.  Il kernel conosce ed utilizza soltanto questi
valori numerici, i nomi ad essi associati sono interamente gestiti in
\textit{user space} con opportune funzioni di libreria, torneremo su questo
argomento in sez.~\ref{sec:sys_user_group}.
 
Grazie a questi identificativi il sistema è in grado di tenere traccia
dell'utente a cui appartiene ciascun processo ed impedire ad altri utenti di
interferire con quest'ultimo.  Inoltre con questo sistema viene anche
garantita una forma base di sicurezza interna in quanto anche l'accesso ai
file (vedi sez.~\ref{sec:file_access_control}) è regolato da questo meccanismo
di identificazione.

Infine in ogni sistema unix-like è presente uno speciale utente privilegiato,
il cosiddetto \textit{superuser}, il cui username è di norma \textit{root}, ed
il cui \ids{UID} è zero. Esso identifica l'amministratore del sistema, che
deve essere in grado di fare qualunque operazione; per l'utente \textit{root}
infatti i meccanismi di controllo cui si è accennato in precedenza sono
disattivati.\footnote{i controlli infatti vengono eseguiti da uno pseudo-codice
  del tipo: ``\code{if (uid) \{ \textellipsis\ \}}''.}


\section{L'architettura di file e directory}
\label{sec:intro_file_dir}

Come accennato in sez.~\ref{sec:intro_base_concept} uno dei concetti
fondamentali dell'architettura di un sistema Unix è il cosiddetto
\textit{everything is a file} (\textsl{tutto è un file}), cioè il fatto che
l'accesso ai vari dispositivi di input/output del computer viene effettuato
attraverso un'interfaccia astratta che tratta le periferiche allo stesso modo
dei normali file di dati.

In questa sezione forniremo una descrizione a grandi linee dell'architettura
della gestione dei file in Linux, partendo da una introduzione ai concetti di
base, per poi illustrare la struttura dell'albero dei file ed il significato
dei tipi di file, concludendo con una panoramica sulle caratteristiche
principali delle due interfacce con cui i processi possono effettuare l'I/O su
file.


\subsection{Una panoramica generale}
\label{sec:file_arch_overview}

Per poter accedere ai file, il kernel deve mettere a disposizione dei
programmi delle opportune \textit{system call} che consentano di leggere e
scrivere il contenuto. Tutto ciò ha due aspetti: il primo è che il kernel, per
il concetto dell'\textit{everything is a file}, deve fornire una interfaccia
che consenta di operare sui file, sia che questi corrispondano ai normali file
di dati, o ai cosiddetti ``\textsl{file speciali}'', come i file di
dispositivo (o \textit{device file}) che permettono di accedere alle
periferiche o le \textit{fifo} ed i socket che forniscono funzionalità di
comunicazione fra processi (torneremo su questo in sez.~\ref{sec:file_mknod}).

Il secondo aspetto è che per poter utilizzare dei normali file di dati il
kernel deve provvedere ad organizzare e rendere accessibile in maniera
opportuna l'informazione in essi contenuta memorizzandola sullo spazio grezzo
disponibile sui dischi.  Questo viene fatto strutturando l'informazione sul
disco attraverso quello che si chiama un
``\textit{filesystem}''. L'informazione così strutturata poi viene resa
disponibile ai processi attraverso quello che viene chiamato il
``\textsl{montaggio}'' del filesystem nell'albero dei file, dove il contenuto
sarà accessibile nella forma ordinaria di file e directory.

\itindbeg{Virtual~File~System~(VFS)}

In Linux il concetto di \textit{everything is a file} è stato implementato
attraverso il cosiddetto \textit{Virtual File System} (da qui in avanti VFS)
che è uno strato intermedio che il kernel usa per accedere ai più svariati
filesystem mantenendo la stessa interfaccia per i programmi in \textit{user
  space}.  Il VFS fornisce quel livello di astrazione che permette di
collegare le operazioni interne del kernel per la manipolazione sui file con
le \textit{system call} relative alle operazioni di I/O, e gestisce poi
l'organizzazione di dette operazioni nei vari modi in cui i diversi filesystem
le effettuano, permettendo la coesistenza di filesystem differenti all'interno
dello stesso albero delle directory. Approfondiremo il funzionamento di
interfaccia generica fornita dal VFS in sez.~\ref{sec:file_vfs_work}.

In sostanza quello che accade è che quando un processo esegue una
\textit{system call} che opera su un file, il kernel chiama sempre una
funzione implementata nel VFS. La funzione eseguirà le manipolazioni sulle
strutture generiche e utilizzerà poi delle chiamate alle opportune funzioni
del filesystem specifico a cui si fa riferimento. Saranno queste infine a
chiamare le funzioni di più basso livello che eseguono le operazioni di I/O
sul dispositivo fisico, secondo lo schema riportato in
fig.~\ref{fig:file_VFS_scheme}.

\begin{figure}[!htb]
  \centering
  \includegraphics[width=8cm]{img/vfs}
  \caption{Schema delle operazioni del VFS.}
  \label{fig:file_VFS_scheme}
\end{figure}

Questa interfaccia resta la stessa anche quando, invece che a dei normali
file, si accede alle periferiche coi citati file di dispositivo, solo che in
questo caso invece di usare il codice del filesystem che accede al disco, il
\textit{Virtual File System} eseguirà direttamente il codice del kernel che
permette di accedere alla periferica.

\itindend{Virtual~File~System~(VFS)}

Come accennato in precedenza una delle funzioni essenziali per il
funzionamento dell'interfaccia dei file è quella che consente di montare un
filesystem nell'albero dei file, e rendere così visibili i suoi contenuti. In
un sistema unix-like infatti, a differenza di quanto avviene in altri sistemi
operativi, tutti i file vengono mantenuti all'interno di un unico albero la
cui radice (quella che viene chiamata \textit{root directory}) viene montata
all'avvio direttamente dal kernel.

Come accennato in sez.~\ref{sec:intro_kern_and_sys}) montare la radice è,
insieme al lancio di \cmd{init},\footnote{l'operazione è ovviamente anche
  preliminare al lancio di \cmd{init}, dato il kernel deve poter accedere al
  file che contiene detto programma.} l'unica operazione che viene effettuata
direttamente dal kernel in fase di avvio. Il kernel infatti, completata la
fase di inizializzazione, utilizza l'indicazione passatagli dal
\textit{bootloader} su quale sia il dispositivo che contiene il filesystem da
usare come punto di partenza, lo monta come radice dell'albero dei file.

Tutti gli ulteriori filesystem che possono essere disponibili su altri
dispositivi dovranno a loro volta essere inseriti nell'albero, montandoli in
altrettante directory del filesystem radice, su quelli che vengono chiamati
\textit{mount point}.  Questo comunque avverrà sempre in un secondo tempo, a
cura dei programmi eseguiti nella procedura di inizializzazione del sistema,
grazie alle funzioni che tratteremo in sez.~\ref{sec:filesystem_mounting}.


\subsection{La risoluzione del nome di file e directory}
\label{sec:file_pathname}

\itindbeg{pathname}

Come appena illustrato sez.~\ref{sec:file_arch_overview} una delle
caratteristiche distintive di un sistema unix-like è quella di avere un unico
albero dei file. Un file deve essere identificato dall'utente usando quello
che viene chiamato il suo \textit{pathname},\footnote{il manuale della
  \acr{glibc} depreca questa nomenclatura, che genererebbe confusione poiché
  \textit{path} indica anche un insieme di directory su cui effettuare una
  ricerca (come quello in cui la shell cerca i comandi). Al suo posto viene
  proposto l'uso di \textit{filename} e di componente per il nome del file
  all'interno della directory. Non seguiremo questa scelta dato che l'uso
  della parola \textit{pathname} è ormai così comune che mantenerne l'uso è
  senz'altro più chiaro dell'alternativa proposta.} vale a dire tramite il
``\textsl{percorso}'' (nome che talvolta viene usato come traduzione di
\textit{pathname}) che si deve fare per accedere al file a partire da una
certa ``\textit{directory}''.

Una directory in realtà è anch'essa un file, nel senso che è anch'essa un
oggetto di un filesystem, solo che è un file particolare che il kernel
riconosce appositamente come tale per poterlo utilizzare come directory. Il
suo scopo è quello di contenere una lista di nomi di file e le informazioni
che associano ciascuno di questi nomi al relativo contenuto (torneremo su
questo in sez.~\ref{sec:file_arch_func}).

Dato che questi nomi possono corrispondere ad un qualunque altro oggetto del
filesystem, compresa un'altra directory, si ottiene naturalmente
un'organizzazione ad albero inserendo nomi di directory dentro altre
directory.  All'interno dello stesso albero si potranno poi inserire anche
tutti gli altri oggetti previsti l'interfaccia del VFS (su cui torneremo in
sez.~\ref{sec:file_file_types}), come le \textit{fifo}, i collegamenti
simbolici, i socket e gli stessi file di dispositivo.

La convenzione usata nei sistemi unix-like per indicare i \textit{pathname}
dei file è quella di usare il carattere ``\texttt{/}'' come separatore fra i
nomi che indicano le directory che lo compongono. Dato che la directory radice
sta in cima all'albero, essa viene indicata semplicemente con il
\textit{pathname} ``\file{/}''.

\itindbeg{pathname~resolution}

Un file può essere indicato rispetto ad una directory semplicemente
specificandone il nome, il manuale della \acr{glibc} chiama i nomi contenuti
nelle directory ``componenti'' (in inglese \textit{file name components}), noi
li chiameremo più semplicemente \textsl{nomi} o \textsl{voci}, riservando la
parola \textsl{componenti} ai nomi che, separati da una ``\texttt{/}'',
costituiscono il \textit{pathname}. Questi poi dovranno corrispondere, perché
il \textit{pathname} sia valido, a voci effettivamente presenti nelle
directory, ma non è detto che un \textit{pathname} debba per forza risultare
valido.  

Il procedimento con cui dato un \textit{pathname} si individua il file a cui
esso fa riferimento, è chiamato \textsl{risoluzione del percorso}
(\textit{filename resolution} o \textit{pathname resolution}). Lo stesso
procedimento ci può anche dire che il \textit{pathname} usato non è valido.
La risoluzione viene eseguita esaminando il \textit{pathname} da sinistra a
destra e localizzando ogni componente dello stesso come nome in una directory
a partire dalla directory iniziale, usando il carattere ``\texttt{/}'' come
separatore per scendere dall'una all'altra. Nel caso si indichi un componente
vuoto il costrutto ``\texttt{//}'' viene considerato equivalente a
``\texttt{/}''.

Ovviamente perché la risoluzione abbia successo occorre che i componenti
intermedi esistano e siano effettivamente directory, e che il file o la
directory indicata dall'ultimo componente esista.  Inoltre i permessi relativi
alle directory indicate nel \textit{pathname} (torneremo su questo
sez.~\ref{sec:file_access_control}) dovranno consentire l'accesso all'intero
\textit{pathname}.

\itindsubbeg{pathname}{assoluto}
\itindsubbeg{pathname}{relativo}

Se il \textit{pathname} comincia con il carattere ``\texttt{/}'' la ricerca
parte dalla directory radice del processo. Questa, a meno di non avere
eseguito una \func{chroot} (funzione su cui torneremo in
sez.~\ref{sec:file_chroot}) è la stessa per tutti i processi ed equivale alla
directory radice dell'albero dei file montata dal kernel all'avvio del
sistema; in questo caso si parla di un \textsl{pathname assoluto}. Altrimenti
la ricerca parte dalla directory di lavoro corrente del processo (su cui
torneremo in sez.~\ref{sec:file_work_dir}) ed il \textit{pathname} è detto
\textsl{pathname relativo}.

\itindsubend{pathname}{assoluto}
\itindsubend{pathname}{relativo}

Infine i nomi di directory ``\file{.}'' e ``\file{..}'' hanno un significato
speciale e vengono inseriti in ogni directory quando questa viene creata (vedi
sez.~\ref{sec:file_dir_creat_rem}). Il primo fa riferimento alla directory
corrente e il secondo alla directory \textsl{genitrice} (o \textit{parent
  directory}) cioè la directory che contiene il riferimento alla directory
corrente.

In questo modo con ``\file{..}'' si può usare un \textit{pathname} relativo
per indicare un file posto al di sopra della directory corrente, tornando
all'indietro nell'albero dei file.  Questa retromarcia però su fermerà una
volta raggiunta la directory radice, perché non esistendo in questo caso una
directory superiore, il nome ``\file{..}''  farà riferimento alla radice
stessa.

\itindend{pathname}
\itindend{pathname~resolution}


\subsection{I tipi di file}
\label{sec:file_file_types}

Parlare dei tipi di file su Linux, come per qualunque sistema unix-like,
significa anzitutto chiarire il proprio vocabolario e sottolineare le
differenze che ci sono rispetto ad altri sistemi operativi.

\index{file!di~dispositivo|(}
\index{file!speciali|(} 

Come accennato in sez.~\ref{sec:file_arch_overview} su Linux l'uso del
\textit{Virtual File System} consente di trattare come file oggetti molto
diversi fra loro. Oltre ai normali file di dati abbiamo già accennato ad altri
due di questi oggetti, i file di dispositivo e le directory, ma ne esistono
altri. In genere quando si parla di tipo di file su Linux si fa riferimento a
questi diversi tipi, di cui si riportato l'elenco completo in
tab.~\ref{tab:file_file_types}.

\begin{table}[htb]
  \footnotesize
  \centering
    \begin{tabular}[c]{|l|l|p{6cm}|}
    \hline
    \multicolumn{2}{|c|}{\textbf{Tipo di file}} & \textbf{Descrizione} \\
    \hline
    \hline
      \textit{regular file}& \textsl{file regolare}
                           & Un file che contiene dei dati (l'accezione
                             normale di file).\\
      \textit{directory}   &\textsl{cartella o direttorio}
                           & Un file che contiene una lista di nomi associati
                             a degli \textit{inode} (vedi
                             sez.~\ref{sec:file_vfs_work}).\\   
      \textit{symbolic link}&\textsl{collegamento simbolico}
                           & Un file che contiene un riferimento ad un altro 
                             file/directory.\\ 
      \textit{char device} &\textsl{dispositivo a caratteri} 
                           & Un file \textsl{speciale} che identifica una
                             periferica ad accesso a caratteri.\\
      \textit{block device}& \textsl{dispositivo a blocchi} 
                           & Un file \textsl{speciale} che identifica una
                             periferica ad accesso a blocchi.\\
      \textit{fifo}        & ``\textsl{coda}'' 
                           & Un file \textsl{speciale} che identifica una
                             linea di comunicazione unidirezionale (vedi
                             sez.~\ref{sec:ipc_named_pipe}).\\
      \textit{socket}      & ``\textsl{presa}''
                           & Un file \textsl{speciale} che identifica una
                             linea di comunicazione bidirezionale (vedi
                             cap.~\ref{cha:socket_intro}).\\
    \hline
    \end{tabular}
    \caption{Tipologia dei file definiti nel VFS}
    \label{tab:file_file_types}
\end{table}

Si tenga ben presente che questa classificazione non ha nulla a che fare con
una classificazione dei file in base al loro contenuto, dato che in tal caso
si avrebbe a che fare sempre e solo con dei file di dati. E non ha niente a
che fare neanche con le eventuali diverse modalità con cui si possa accedere
al contenuto dei file di dati.  La classificazione di
tab.~\ref{tab:file_file_types} riguarda il tipo di oggetti gestiti dal
\textit{Virtual File System}, ed è da notare la presenza dei cosiddetti file
``\textsl{speciali}''.

Alcuni di essi, come le \textit{fifo} (che tratteremo in
sez.~\ref{sec:ipc_named_pipe}) ed i \textit{socket} (che tratteremo in
cap.~\ref{cha:socket_intro}) non sono altro che dei riferimenti per utilizzare
alcune funzionalità di comunicazione fornite dal kernel. Gli altri sono
proprio quei \textsl{file di dispositivo} che costituiscono una interfaccia
diretta per leggere e scrivere sui dispositivi fisici. Anche se finora li
abbiamo chiamati genericamente così, essi sono tradizionalmente suddivisi in
due grandi categorie, \textsl{a blocchi} e \textsl{a caratteri}, a seconda
delle modalità in cui il dispositivo sottostante effettua le operazioni di
I/O.

I dispositivi a blocchi (ad esempio i dischi) sono quelli corrispondono a
periferiche per le quali è richiesto che l'I/O venga effettuato per blocchi di
dati di dimensioni fissate (nel caso dei dischi le dimensioni di un settore),
mentre i dispositivi a caratteri sono quelli per cui l'I/O può essere
effettuato senza nessuna particolare struttura, ed in generale anche un byte
alla volta, da cui il nome.

Una delle differenze principali con altri sistemi operativi come il VMS o
Windows è che per Unix tutti i file di dati sono identici e contengono un
flusso continuo di byte. Non esiste cioè differenza per come vengono visti dal
sistema file di diverso contenuto o formato, come nel caso di quella fra file
di testo e binari che c'è in Windows. Non c'è neanche una strutturazione a
record per il cosiddetto ``\textsl{accesso diretto}'' come nel caso del
VMS.\footnote{questo vale anche per i dispositivi a blocchi: la strutturazione
  dell'I/O in blocchi di dimensione fissa avviene solo all'interno del kernel,
  ed è completamente trasparente all'utente; inoltre talvolta si parla di
  \textsl{accesso diretto} riferendosi alla capacità, che non ha niente a che
  fare con tutto ciò, di effettuare attraverso degli appositi file di
  dispositivo delle operazioni di I/O direttamente sui dischi senza passare
  attraverso un filesystem, il cosiddetto \textit{raw access}, introdotto coi
  kernel della serie 2.4.x ma ormai in sostanziale disuso.}

\index{file!di~dispositivo|)}
\index{file!speciali|)} 

Una differenza che attiene ai contenuti di un file però esiste, ed è relativa
al formato dei file di testo. Nei sistemi unix-like la fine riga è codificata
in maniera diversa da Windows o dal vecchio MacOS, in particolare il fine riga
è il carattere \texttt{LF} (\verb|\n|) al posto del \texttt{CR} (\verb|\r|)
del vecchio MacOS e del \texttt{CR LF} (\verb|\r\n|) di Windows. Questo può
causare alcuni problemi qualora nei programmi si facciano assunzioni sul
terminatore della riga e per questo esistono dei programmi come \cmd{unix2dos}
e \cmd{dos2unix} che effettuano una conversione fra questi due formati di
testo.

Si ricordi comunque che un kernel unix-like non fornisce nessun supporto per
la tipizzazione dei file di dati in base al loro contenuto e che non c'è
nessun supporto per una qualche interpretazione delle estensioni (nel nome del
file) da parte del kernel,\footnote{non è così ad esempio nel filesystem HFS
  dei Mac, che supporta delle risorse associate ad ogni file, che specificano
  fra l'altro il contenuto ed il programma da usare per leggerlo; in realtà
  per alcuni filesystem esiste la possibilità di associare delle risorse ai
  file con gli \textit{extended attributes} (vedi sez.~\ref{sec:file_xattr}),
  ma è una caratteristica tutt'ora poco utilizzata, dato che non corrisponde
  al modello classico dei file in un sistema Unix.} ogni classificazione di
questo tipo avviene sempre in \textit{user space}. Gli unici file di cui il
kernel deve essere in grado di capire il contenuto sono i binari dei
programmi, per i quali sono supportati solo alcuni formati, anche se oggi
viene usato quasi esclusivamente
\itindex{Executable~and~Linkable~Format~(ELF)} l'ELF.\footnote{il nome è
  l'acronimo di \textit{Executable and Linkable Format}, un formato per
  eseguibili binari molto flessibile ed estendibile definito nel 1995 dal
  \textit{Tool Interface Standard} che per le sue caratteristiche di non
  essere legato a nessun tipo di processore o architettura è stato adottato da
  molti sistemi unix-like e non solo.}

\itindbeg{magic~number}

Nonostante l'assenza di supporto da parte del kernel per la classificazione
del contenuto dei file di dati, molti programmi adottano comunque delle
convenzioni per i nomi dei file, ad esempio il codice C normalmente si mette
in file con l'estensione ``\file{.c}''. Inoltre una tecnica molto usata per
classificare i contenuti da parte dei programmi è quella di utilizzare i primi
byte del file per memorizzare un ``\textit{magic number}'' che ne classifichi
il contenuto. Il concetto è quello di un numero intero, solitamente fra 2 e 10
byte, che identifichi il contenuto seguente, dato che questi sono anche
caratteri è comune trovare espresso tale numero con stringhe come
``\texttt{\%PDF}'' per i PDF o ``\texttt{\#!}'' per gli script. Entrambe
queste tecniche, per quanto usate ed accettate in maniera diffusa, restano
solo delle convenzioni il cui rispetto è demandato alle applicazioni stesse.

\itindend{magic~number}


\subsection{Le due interfacce per l'accesso ai file}
\label{sec:file_io_api}


\itindbeg{file~descriptor}

In Linux le interfacce di programmazione per l'I/O su file due.  La prima è
l'interfaccia nativa del sistema, quella che il manuale della \textsl{glibc}
chiama interfaccia dei ``\textit{file descriptor}'' (in italiano
\textsl{descrittori di file}). Si tratta di un'interfaccia specifica dei
sistemi unix-like che fornisce un accesso non bufferizzato.

L'interfaccia è essenziale, l'accesso viene detto non bufferizzato in quanto
la lettura e la scrittura vengono eseguite chiamando direttamente le
\textit{system call} del kernel, anche se in realtà il kernel effettua al suo
interno alcune bufferizzazioni per aumentare l'efficienza nell'accesso ai
dispositivi. L'accesso viene gestito attraverso i \textit{file descriptor} che
sono rappresentati da numeri interi (cioè semplici variabili di tipo
\ctyp{int}).  L'interfaccia è definita nell'\textit{header file}
\headfile{unistd.h} e la tratteremo in dettaglio in
sez.~\ref{sec:file_unix_interface}.

\itindbeg{file~stream}

La seconda interfaccia è quella che il manuale della \acr{glibc} chiama dei
\textit{file stream} o più semplicemente degli \textit{stream}.\footnote{in
  realtà una interfaccia con lo stesso nome è stata introdotta a livello di
  kernel negli Unix derivati da \textit{System V}, come strato di astrazione
  per file e socket; in Linux questa interfaccia, che comunque ha avuto poco
  successo, non esiste, per cui facendo riferimento agli \textit{stream}
  useremo il significato adottato dal manuale della \acr{glibc}.} Essa
fornisce funzioni più evolute e un accesso bufferizzato, controllato dalla
implementazione fatta nella \acr{glibc}.  Questa è l'interfaccia specificata
dallo standard ANSI C e perciò si trova anche su tutti i sistemi non Unix. Gli
\textit{stream} sono oggetti complessi e sono rappresentati da puntatori ad un
opportuna struttura definita dalle librerie del C, a cui si accede sempre in
maniera indiretta utilizzando il tipo \code{FILE *}; l'interfaccia è definita
nell'\textit{header file} \headfile{stdio.h} e la tratteremo in dettaglio in
sez.~\ref{sec:files_std_interface}.

Entrambe le interfacce possono essere usate per l'accesso ai file come agli
altri oggetti del VFS, ma per poter accedere alle operazioni di controllo
(descritte in sez.~\ref{sec:file_fcntl_ioctl}) su un qualunque tipo di oggetto
del VFS occorre usare l'interfaccia standard di Unix con i file
descriptor. Allo stesso modo devono essere usati i file descriptor se si vuole
ricorrere a modalità speciali di I/O come il \textit{file locking} o l'I/O
non-bloccante (vedi cap.~\ref{cha:file_advanced}).

Gli \textit{stream} forniscono un'interfaccia di alto livello costruita sopra
quella dei \textit{file descriptor}, che permette di poter scegliere tra
diversi stili di bufferizzazione.  Il maggior vantaggio degli \textit{stream}
è che l'interfaccia per le operazioni di input/output è molto più ricca di
quella dei \textit{file descriptor}, che forniscono solo funzioni elementari
per la lettura/scrittura diretta di blocchi di byte.  In particolare gli
\textit{stream} dispongono di tutte le funzioni di formattazione per l'input e
l'output adatte per manipolare anche i dati in forma di linee o singoli
caratteri.

In ogni caso, dato che gli \textit{stream} sono implementati sopra
l'interfaccia standard di Unix, è sempre possibile estrarre il \textit{file
  descriptor} da uno \textit{stream} ed eseguirvi sopra operazioni di basso
livello, o associare in un secondo tempo uno \textit{stream} ad un
\textit{file descriptor} per usare l'interfaccia più sofisticata.

In generale, se non necessitano specificatamente le funzionalità di basso
livello, è opportuno usare sempre gli \textit{stream} per la loro maggiore
portabilità, essendo questi ultimi definiti nello standard ANSI C;
l'interfaccia con i \textit{file descriptor} infatti segue solo lo standard
POSIX.1 dei sistemi Unix, ed è pertanto di portabilità più limitata.

\itindend{file~descriptor}
\itindend{file~stream}

\section{Gli standard}
\label{sec:intro_standard}

In questa sezione faremo una breve panoramica relativa ai vari standard che
nel tempo sono stati formalizzati da enti, associazioni, consorzi e
organizzazioni varie al riguardo ai sistemi operativi di tipo Unix o alle
caratteristiche che si sono stabilite come standard di fatto in quanto facenti
parte di alcune implementazioni molto diffuse come BSD o System V.

Ovviamente prenderemo in considerazione solo gli standard riguardanti
interfacce di programmazione e le altre caratteristiche di un sistema
unix-like (alcuni standardizzano pure i comandi base del sistema e la shell)
ed in particolare ci concentreremo sul come ed in che modo essi sono
supportati sia per quanto riguarda il kernel che la libreria standard del C,
con una particolare attenzione alla \acr{glibc}.


\subsection{Lo standard ANSI C}
\label{sec:intro_ansiC}

\itindbeg{ANSI~C}
Lo standard ANSI C è stato definito nel 1989 dall'\textit{American National
  Standard Institute} come prima standardizzazione del linguaggio C e per
questo si fa riferimento ad esso anche come C89. L'anno successivo è stato
adottato dalla ISO (\textit{International Standard Organisation}) come
standard internazionale con la sigla ISO/IEC 9899:1990, e per questo è noto
anche sotto il nome di standard ISO C, o ISO C90.  Nel 1999 è stata pubblicata
una revisione dello standard C89, che viene usualmente indicata come C99,
anche questa è stata ratificata dalla ISO con la sigla ISO/IEC 9899:1990, per
cui vi si fa riferimento anche come ISO C99.

Scopo dello standard è quello di garantire la portabilità dei programmi C fra
sistemi operativi diversi, ma oltre alla sintassi ed alla semantica del
linguaggio C (operatori, parole chiave, tipi di dati) lo standard prevede
anche una libreria di funzioni che devono poter essere implementate su
qualunque sistema operativo.

Per questo motivo, anche se lo standard non ha alcun riferimento ad un sistema
di tipo Unix, GNU/Linux (per essere precisi la \acr{glibc}), come molti Unix
moderni, provvede la compatibilità con questo standard, fornendo le funzioni
di libreria da esso previste. Queste sono dichiarate in una serie di
\textit{header file} anch'essi forniti dalla \acr{glibc} (tratteremo
l'argomento in sez.~\ref{sec:proc_syscall}).

In realtà la \acr{glibc} ed i relativi \textit{header file} definiscono un
insieme di funzionalità in cui sono incluse come sottoinsieme anche quelle
previste dallo standard ANSI C. È possibile ottenere una conformità stretta
allo standard (scartando le funzionalità addizionali) usando il \cmd{gcc} con
l'opzione \cmd{-ansi}. Questa opzione istruisce il compilatore a definire nei
vari \textit{header file} soltanto le funzionalità previste dallo standard
ANSI C e a non usare le varie estensioni al linguaggio e al preprocessore da
esso supportate.
\itindend{ANSI~C}


\subsection{I tipi di dati primitivi}
\label{sec:intro_data_types}

Uno dei problemi di portabilità del codice più comune è quello dei tipi di
dati utilizzati nei programmi, che spesso variano da sistema a sistema, o
anche da una architettura ad un'altra (ad esempio passando da macchine con
processori 32 bit a 64). In particolare questo è vero nell'uso dei cosiddetti
\index{tipo!elementare} \textit{tipi elementari} del linguaggio C (come
\ctyp{int}) la cui dimensione varia a seconda dell'architettura hardware.

Storicamente alcuni tipi nativi dello standard ANSI C sono sempre stati
associati ad alcune variabili nei sistemi Unix, dando per scontata la
dimensione. Ad esempio la posizione corrente all'interno di un file è stata
associata ad un intero a 32 bit, mentre il numero di dispositivo è stato
associato ad un intero a 16 bit. Storicamente questi erano definiti
rispettivamente come \ctyp{int} e \ctyp{short}, ma tutte le volte che, con
l'evolversi ed il mutare delle piattaforme hardware, alcuni di questi tipi si
sono rivelati inadeguati o sono cambiati, ci si è trovati di fronte ad una
infinita serie di problemi di portabilità.

\begin{table}[htb]
  \footnotesize
  \centering
  \begin{tabular}[c]{|l|l|}
    \hline
    \textbf{Tipo} & \textbf{Contenuto} \\
    \hline
    \hline
    \typed{caddr\_t} & Core address.\\
    \typed{clock\_t} & Contatore del \textit{process time} (vedi
                      sez.~\ref{sec:sys_cpu_times}.\\ 
    \typed{dev\_t}   & Numero di dispositivo (vedi sez.~\ref{sec:file_mknod}).\\
    \typed{gid\_t}   & Identificatore di un gruppo (vedi
                      sez.~\ref{sec:proc_access_id}).\\
    \typed{ino\_t}   & Numero di \textit{inode} 
                      (vedi sez.~\ref{sec:file_vfs_work}).\\ 
    \typed{key\_t}   & Chiave per il System V IPC (vedi
                      sez.~\ref{sec:ipc_sysv_generic}).\\
    \typed{loff\_t}  & Posizione corrente in un file.\\
    \typed{mode\_t}  & Attributi di un file.\\
    \typed{nlink\_t} & Contatore dei collegamenti su un file.\\
    \typed{off\_t}   & Posizione corrente in un file.\\
    \typed{pid\_t}   & Identificatore di un processo (vedi
                      sez.~\ref{sec:proc_pid}).\\
    \typed{rlim\_t}  & Limite sulle risorse.\\
    \typed{sigset\_t}& Insieme di segnali (vedi sez.~\ref{sec:sig_sigset}).\\
    \typed{size\_t}  & Dimensione di un oggetto.\\
    \typed{ssize\_t} & Dimensione in numero di byte ritornata dalle funzioni.\\
    \typed{ptrdiff\_t}& Differenza fra due puntatori.\\
    \typed{time\_t}  & Numero di secondi (in \textit{calendar time}, vedi 
                      sez.~\ref{sec:sys_time}).\\
    \typed{uid\_t}   & Identificatore di un utente (vedi
                      sez.~\ref{sec:proc_access_id}).\\
    \hline
  \end{tabular}
  \caption{Elenco dei tipi primitivi, definiti in \headfile{sys/types.h}.}
  \label{tab:intro_primitive_types}
\end{table}

Per questo motivo tutte le funzioni di libreria di solito non fanno
riferimento ai tipi elementari dello standard del linguaggio C, ma ad una
serie di \index{tipo!primitivo} \textsl{tipi primitivi} del sistema, riportati
in tab.~\ref{tab:intro_primitive_types}, e definiti nell'\textit{header file}
\headfiled{sys/types.h}, in modo da mantenere completamente indipendenti i tipi
utilizzati dalle funzioni di sistema dai tipi elementari supportati dal
compilatore C.


\subsection{Lo standard System V}
\label{sec:intro_sysv}

Come noto Unix nasce nei laboratori della AT\&T, che ne registrò il nome come
marchio depositato, sviluppandone una serie di versioni diverse; nel 1983 la
versione supportata ufficialmente venne rilasciata al pubblico con il nome di
Unix System V, e si fa rifermento a questa implementazione con la sigla SysV o
SV.

Negli anni successivi l'AT\&T proseguì lo sviluppo del sistema rilasciando
varie versioni con aggiunte e integrazioni, ed in particolare la
\textit{release 2} nel 1985, a cui si fa riferimento con il nome SVr2 e la
\textit{release 3} nel 1986 (denominata SVr3). Le interfacce di programmazione
di queste due versioni vennero descritte formalmente in due documenti
denominati \itindex{System~V~Interface~Definition~(SVID)} \textit{System V
  Interface Definition} (o SVID), pertanto nel 1985 venne rilasciata la
specifica SVID 1 e nel 1986 la specifica SVID 2.

Nel 1989 un accordo fra vari venditori (AT\&T, Sun, HP, ed altri) portò ad una
versione di System V che provvedeva un'unificazione delle interfacce
comprendente anche Xenix e BSD, questa venne denominata \textit{release 4} o
SVr4. Anche le relative interfacce vennero descritte in un documento dal
titolo \textit{System V Interface Description}, venendo a costituire lo
standard SVID 3, che viene considerato la specifica finale di System V, ed a
cui spesso si fa riferimento semplicemente con SVID. Anche SVID costituisce un
sovrainsieme delle interfacce definite dallo standard POSIX.  

Nel 1992 venne rilasciata una seconda versione del sistema, la SVr4.2; l'anno
successivo la divisione della AT\&T (già a suo tempo rinominata in Unix System
Laboratories) venne acquistata dalla Novell, che poi trasferì il marchio Unix
al consorzio X/Open. L'ultima versione di System V fu la SVr4.2MP rilasciata
nel dicembre 1993. Infine nel 1995 è stata rilasciata da SCO, che aveva
acquisito alcuni diritti sul codice di System V, una ulteriore versione delle
\textit{System V Interface Description}, che va sotto la denominazione di SVID
4.

Linux e la \acr{glibc} implementano le principali funzionalità richieste dalle
specifiche SVID che non sono già incluse negli standard POSIX ed ANSI C, per
compatibilità con lo Unix System V e con altri Unix (come SunOS) che le
includono. Tuttavia le funzionalità più oscure e meno utilizzate (che non sono
presenti neanche in System V) sono state tralasciate.

Le funzionalità implementate sono principalmente il meccanismo di
intercomunicazione fra i processi e la memoria condivisa (il cosiddetto System
V IPC, che vedremo in sez.~\ref{sec:ipc_sysv}) le funzioni della famiglia
\funcm{hsearch} e \funcm{drand48}, \funcm{fmtmsg} e svariate funzioni
matematiche.


\subsection{Lo ``\textsl{standard}'' BSD}
\label{sec:intro_bsd}

Lo sviluppo di BSD iniziò quando la fine della collaborazione fra l'Università
di Berkeley e la AT\&T generò una delle prime e più importanti fratture del
mondo Unix.  L'università di Berkeley proseguì nello sviluppo della base di
codice di cui disponeva, e che presentava parecchie migliorie rispetto alle
versioni allora disponibili, fino ad arrivare al rilascio di una versione
completa di Unix, chiamata appunto BSD, del tutto indipendente dal codice
della AT\&T.

Benché BSD non sia mai stato uno standard formalizzato, l'implementazione
dello Unix dell'Università di Berkeley nella sua storia ha introdotto una
serie di estensioni e interfacce di grandissima rilevanza, come i collegamenti
simbolici, la funzione \code{select} ed i socket di rete. Per questo motivo si
fa spesso riferimento esplicito alle interfacce presenti nelle varie versioni
dello Unix di Berkeley con una apposita sigla.

Nel 1983, con il rilascio della versione 4.2 di BSD, venne definita una
implementazione delle funzioni di interfaccia a cui si fa riferimento con la
sigla 4.2BSD. Per fare riferimento alle precedenti versioni si usano poi le
sigle 3BSD e 4BSD (per le due versioni pubblicate nel 1980), e 4.1BSD per
quella pubblicata nel 1981.

Le varie estensioni ideate a Berkeley sono state via via aggiunte al sistema
nelle varie versioni succedutesi negli anni, che vanno sotto il nome di
4.3BSD, per la versione rilasciata nel 1986 e 4.4BSD, per la versione
rilasciata nel 1993, che costituisce l'ultima release ufficiale
dell'università di Berkeley. Si tenga presente che molte di queste interfacce
sono presenti in derivati commerciali di BSD come SunOS. Il kernel Linux e la
\acr{glibc} forniscono tutte queste estensioni che sono state in gran parte
incorporate negli standard successivi.


\subsection{Gli standard IEEE -- POSIX}
\label{sec:intro_posix}

Lo standard ufficiale creato da un organismo indipendente più attinente alle
interfacce di un sistema unix-like nel suo complesso (e che concerne sia il
kernel che le librerie che i comandi) è stato lo standard POSIX. Esso prende
origine dallo standard ANSI C, che contiene come sottoinsieme, prevedendo
ulteriori capacità per le funzioni in esso definite, ed aggiungendone di
nuove.

In realtà POSIX è una famiglia di standard diversi, il cui nome, suggerito da
Richard Stallman, sta per \textit{Portable Operating System Interface}, ma la
X finale denuncia la sua stretta relazione con i sistemi Unix. Esso nasce dal
lavoro dell'IEEE (\textit{Institute of Electrical and Electronics Engeneers})
che ne produsse una prima versione, nota come \textsl{IEEE 1003.1-1988},
mirante a standardizzare l'interfaccia con il sistema operativo.

Ma gli standard POSIX non si limitano alla standardizzazione delle funzioni di
libreria, e in seguito sono stati prodotti anche altri standard per la shell e
i comandi di sistema (1003.2), per le estensioni \textit{real-time} e per i
\textit{thread} (rispettivamente 1003.1d e 1003.1c) per i socket (1003.1g) e
vari altri.  In tab.~\ref{tab:intro_posix_std} è riportata una classificazione
sommaria dei principali documenti prodotti, e di come sono identificati fra
IEEE ed ISO; si tenga conto inoltre che molto spesso si usa l'estensione IEEE
anche come aggiunta al nome POSIX; ad esempio è più comune parlare di POSIX.4
come di POSIX.1b.

Si tenga presente inoltre che nuove specifiche e proposte di standardizzazione
si aggiungono continuamente, mentre le versioni precedenti vengono riviste;
talvolta poi i riferimenti cambiano nome, per cui anche solo seguire le
denominazioni usate diventa particolarmente faticoso.

\begin{table}[htb]
  \footnotesize
  \centering
  \begin{tabular}[c]{|l|l|l|l|}
    \hline
    \textbf{Standard} & \textbf{IEEE} & \textbf{ISO} & \textbf{Contenuto} \\
    \hline
    \hline
    POSIX.1 & 1003.1 & 9945-1& Interfacce di base.                          \\
    POSIX.1a& 1003.1a& 9945-1& Estensioni a POSIX.1.                        \\
    POSIX.2 & 1003.2 & 9945-2& Comandi.                                     \\
    POSIX.3 & 2003   &TR13210& Metodi di test.                              \\
    POSIX.4 & 1003.1b &  --- & Estensioni real-time.                        \\
    POSIX.4a& 1003.1c &  --- & Thread.                                      \\
    POSIX.4b& 1003.1d &9945-1& Ulteriori estensioni real-time.              \\
    POSIX.5 & 1003.5  & 14519& Interfaccia per il linguaggio ADA.           \\
    POSIX.6 & 1003.2c,1e& 9945-2& Sicurezza.                                \\
    POSIX.8 & 1003.1f& 9945-1& Accesso ai file via rete.                    \\
    POSIX.9 & 1003.9  &  --- & Interfaccia per il Fortran-77.               \\
    POSIX.12& 1003.1g& 9945-1& Socket.                                      \\
    \hline
  \end{tabular}
  \caption{Elenco dei vari standard POSIX e relative denominazioni.}
  \label{tab:intro_posix_std}
\end{table}

Benché l'insieme degli standard POSIX siano basati sui sistemi Unix, essi
definiscono comunque un'interfaccia di programmazione generica e non fanno
riferimento ad una implementazione specifica (ad esempio esiste
un'implementazione di POSIX.1 anche sotto Windows NT).  

Linux e la \acr{glibc} implementano tutte le funzioni definite nello standard
POSIX.1, queste ultime forniscono in più alcune ulteriori capacità (per
funzioni di \textit{pattern matching} e per la manipolazione delle
\textit{regular expression}), che vengono usate dalla shell e dai comandi di
sistema e che sono definite nello standard POSIX.2.

Nelle versioni più recenti del kernel e delle librerie sono inoltre supportate
ulteriori funzionalità aggiunte dallo standard POSIX.1c per quanto riguarda i
\textit{thread} (vedi cap.~\ref{cha:threads}), e dallo standard POSIX.1b per
quanto riguarda i segnali e lo scheduling real-time
(sez.~\ref{sec:sig_real_time} e sez.~\ref{sec:proc_real_time}), la misura del
tempo, i meccanismi di intercomunicazione (sez.~\ref{sec:ipc_posix}) e l'I/O
asincrono (sez.~\ref{sec:file_asyncronous_io}).

Lo standard principale resta comunque POSIX.1, che continua ad evolversi; la
versione più nota, cui gran parte delle implementazioni fanno riferimento, e
che costituisce una base per molti altri tentativi di standardizzazione, è
stata rilasciata anche come standard internazionale con la sigla
\textsl{ISO/IEC 9945-1:1996} ed include i precedenti POSIX.1b e POSIX.1c. In
genere si fa riferimento ad essa come POSIX.1-1996.

Nel 2001 è stata poi eseguita una sintesi degli standard POSIX.1, POSIX.2 e
SUSv3 (vedi sez.~\ref{sec:intro_xopen}) in un unico documento, redatto sotto
gli auspici del cosiddetto gruppo Austin che va sotto il nome di POSIX.1-2001.
Questo standard definisce due livelli di conformità, quello POSIX, in cui sono
presenti solo le interfacce di base, e quello XSI che richiede la presenza di
una serie di estensioni opzionali per lo standard POSIX, riprese da SUSv3.
Inoltre lo standard è stato allineato allo standard C99, e segue lo stesso
nella definizione delle interfacce.

A questo standard sono stati aggiunti due documenti di correzione e
perfezionamento denominati \textit{Technical Corrigenda}, il TC1 del 2003 ed
il TC2 del 2004, e talvolta si fa riferimento agli stessi con le sigle
POSIX.1-2003 e POSIX.1-2004. 

Una ulteriore revisione degli standard POSIX e SUS è stata completata e
ratificata nel 2008, cosa che ha portato al rilascio di una nuova versione
sotto il nome di POSIX.1-2008 (e SUSv4), con l'incorporazione di alcune nuove
interfacce, la obsolescenza di altre, la trasformazione da opzionali a
richieste di alcune specifiche di base, oltre alle solite precisazioni ed
aggiornamenti. Anche in questo caso è prevista la suddivisione in una
conformità di base, e delle interfacce aggiuntive. Una ulteriore revisione è
stata pubblicata nel 2017 come POSIX.1-2017.

Le procedure di aggiornamento dello standard POSIX prevedono comunque un
percorso continuo, che prevede la possibilità di introduzione di nuove
interfacce e la definizione di precisazioni ed aggiornamenti, per questo in
futuro verranno rilasciate nuove versioni. Alla stesura di queste note
l'ultima revisione approvata resta POSIX.1-2017, uno stato della situazione
corrente del supporto degli standard è allegato alla documentazione della
\acr{glibc} e si può ottenere con il comando \texttt{man standards}.


\subsection{Gli standard X/Open -- Opengroup -- Unix}
\label{sec:intro_xopen}

Il consorzio X/Open nacque nel 1984 come consorzio di venditori di sistemi
Unix per giungere ad un'armonizzazione delle varie implementazioni.  Per far
questo iniziò a pubblicare una serie di documentazioni e specifiche sotto il
nome di \textit{X/Open Portability Guide} a cui di norma si fa riferimento con
l'abbreviazione XPG$n$, con $n$ che indica la versione.

Nel 1989 il consorzio produsse una terza versione di questa guida
particolarmente voluminosa (la \textit{X/Open Portability Guide, Issue 3}),
contenente una dettagliata standardizzazione dell'interfaccia di sistema di
Unix, che venne presa come riferimento da vari produttori. Questo standard,
detto anche XPG3 dal nome della suddetta guida, è sempre basato sullo standard
POSIX.1, ma prevede una serie di funzionalità aggiuntive fra cui le specifiche
delle API (sigla che sta per \textit{Application Programmable Interface}, in
italiano interfaccia di programmazione) per l'interfaccia grafica (X11).

Nel 1992 lo standard venne rivisto con una nuova versione della guida, la
Issue 4, da cui la sigla XPG4, che aggiungeva l'interfaccia XTI (\textit{X
  Transport Interface}) mirante a soppiantare (senza molto successo)
l'interfaccia dei socket derivata da BSD. Una seconda versione della guida fu
rilasciata nel 1994; questa è nota con il nome di Spec 1170 (dal numero delle
interfacce, intestazioni e comandi definiti) ma si fa riferimento ad essa
anche come XPG4v2.

Nel 1993 il marchio Unix passò di proprietà dalla Novell (che a sua volta lo
aveva comprato dalla AT\&T) al consorzio X/Open che iniziò a pubblicare le sue
specifiche sotto il nome di \textit{Single UNIX Specification} o SUS, l'ultima
versione di Spec 1170 diventò così la prima versione delle \textit{Single UNIX
  Specification}, detta SUS o SUSv1, ma più comunemente nota anche come
\textit{Unix 95}.

Nel 1996 la fusione del consorzio X/Open con la Open Software Foundation (nata
da un gruppo di aziende concorrenti rispetto ai fondatori di X/Open) portò
alla costituzione dell'\textit{Open Group}, un consorzio internazionale che
raccoglie produttori, utenti industriali, entità accademiche e governative.
Attualmente il consorzio è detentore del marchio depositato Unix, e prosegue
il lavoro di standardizzazione delle varie implementazioni, rilasciando
periodicamente nuove specifiche e strumenti per la verifica della conformità
alle stesse.

Nel 1997 fu annunciata la seconda versione delle \textit{Single UNIX
  Specification}, nota con la sigla SUSv2, in questa versione le interfacce
specificate salgono a 1434, e addirittura a 3030 se si considerano le stazioni
di lavoro grafiche, per le quali sono inserite pure le interfacce usate da CDE
che richiede sia X11 che Motif. La conformità a questa versione permette l'uso
del nome \textit{Unix 98}, usato spesso anche per riferirsi allo standard. Un
altro nome alternativo di queste specifiche, date le origini, è XPG5.

Come accennato nel 2001, con il rilascio dello standard POSIX.1-2001, è stato
effettuato uno sforzo di sintesi in cui sono state comprese, nella parte di
interfacce estese, anche le interfacce definite nelle \textit{Single UNIX
  Specification}, pertanto si può fare riferimento a detto standard, quando
comprensivo del rispetto delle estensioni XSI, come SUSv3, e fregiarsi del
marchio UNIX 03 se conformi ad esso. 

Infine, come avvenuto per POSIX.1-2001, anche con la successiva revisione
dello standard POSIX.1 (la POSIX.1-2008) è stato stabilito che la conformità
completa a tutte quelle che sono le nuove estensioni XSI previste
dall'aggiornamento vada a definire la quarta versione delle \textit{Single
  UNIX Specification}, chiamata appunto SUSv4.


\subsection{Il controllo di aderenza agli standard}
\label{sec:intro_gcc_glibc_std}

In Linux, se si usa la \acr{glibc}, la conformità agli standard appena
descritti può essere richiesta sia attraverso l'uso di alcune opzioni del
compilatore (il \texttt{gcc}) che con la definizione di opportune costanti
prima dell'inclusione dei file di intestazione (gli \textit{header file}, vedi
sez.~\ref{sec:proc_syscall}) in cui le varie funzioni di libreria vengono
definite.

Ad esempio se si vuole che i programmi seguano una stretta attinenza allo
standard ANSI C si può usare l'opzione \texttt{-ansi} del compilatore, e non
potrà essere utilizzata nessuna funzione non riconosciuta dalle specifiche
standard ISO per il C.  Il \texttt{gcc} possiede inoltre una specifica opzione
per richiedere la conformità ad uno standard, nella forma \texttt{-std=nome},
dove \texttt{nome} può essere \texttt{c89} per indicare lo standard ANSI C
(vedi sez.~\ref{sec:intro_ansiC}) o \texttt{c99} per indicare la conformità
allo standard C99 o \texttt{c11} per indicare la conformità allo standard C11
(revisione del 2011).\footnote{esistono anche le possibilità di usare i valori
  \texttt{gnu89}, che indica l'uso delle estensioni GNU al C89, riprese poi
  dal C99, \texttt{gnu99} che indica il dialetto GNU del C99, o \texttt{gnu11}
  che indica le estensioni GNU al C11, lo standard adottato di default dipende
  dalla versione del \texttt{gcc}, ed all'agosto 2018 con la versione 8.2 è
  \texttt{gnu11}.}

Per attivare le varie opzioni di controllo di aderenza agli standard è poi
possibile definire delle macro di preprocessore che controllano le
funzionalità che la \acr{glibc} può mettere a disposizione:\footnote{le macro
  sono definite nel file di dichiarazione \file{<features.h>}, ma non è
  necessario includerlo nei propri programmi in quanto viene automaticamente
  incluso da tutti gli altri file di dichiarazione che utilizzano le macro in
  esso definite; si tenga conto inoltre che il file definisce anche delle
  ulteriori macro interne, in genere con un doppio prefisso di \texttt{\_},
  che non devono assolutamente mai essere usate direttamente. } questo può
essere fatto attraverso l'opzione \texttt{-D} del compilatore, ma è buona
norma farlo inserendo gli opportuni \code{\#define} prima della inclusione dei
propri \textit{header file} (vedi sez.~\ref{sec:proc_syscall}).

Le macro disponibili per controllare l'aderenza ai vari standard messi a
disposizione della \acr{glibc}, che rendono disponibili soltanto le funzioni
in essi definite, sono illustrate nel seguente elenco:
\begin{basedescript}{\desclabelwidth{2.7cm}\desclabelstyle{\nextlinelabel}}
\item[\macrod{\_\_STRICT\_ANSI\_\_}] richiede l'aderenza stretta allo standard
  C ISO; viene automaticamente predefinita qualora si invochi il \texttt{gcc}
  con le opzione \texttt{-ansi} o \texttt{-std=c99}.

\item[\macrod{\_POSIX\_SOURCE}] definendo questa macro (considerata obsoleta)
  si rendono disponibili tutte le funzionalità dello standard POSIX.1 (la
  versione IEEE Standard 1003.1) insieme a tutte le funzionalità dello
  standard ISO C. Se viene anche definita con un intero positivo la macro
  \macro{\_POSIX\_C\_SOURCE} lo stato di questa non viene preso in
  considerazione.

\item[\macrod{\_POSIX\_C\_SOURCE}] definendo questa macro ad un valore intero
  positivo si controlla quale livello delle funzionalità specificate da POSIX
  viene messa a disposizione; più alto è il valore maggiori sono le
  funzionalità:
  \begin{itemize}
  \item un valore uguale a ``\texttt{1}'' rende disponibili le funzionalità
    specificate nella edizione del 1990 (IEEE Standard 1003.1-1990);
  \item valori maggiori o uguali a ``\texttt{2}'' rendono disponibili le
    funzionalità previste dallo standard POSIX.2 specificate nell'edizione del
    1992 (IEEE Standard 1003.2-1992),
  \item un valore maggiore o uguale a ``\texttt{199309L}'' rende disponibili
    le funzionalità previste dallo standard POSIX.1b specificate nell'edizione
    del 1993 (IEEE Standard 1003.1b-1993);
  \item un valore maggiore o uguale a ``\texttt{199506L}'' rende disponibili
    le funzionalità previste dallo standard POSIX.1 specificate nell'edizione
    del 1996 (\textit{ISO/IEC 9945-1:1996}), ed in particolare le definizioni
    dello standard POSIX.1c per i \textit{thread};
  \item a partire dalla versione 2.3.3 della \acr{glibc} un valore maggiore o
    uguale a ``\texttt{200112L}'' rende disponibili le funzionalità di base
    previste dallo standard POSIX.1-2001, escludendo le estensioni XSI;
  \item a partire dalla versione 2.10 della \acr{glibc} un valore maggiore o
    uguale a ``\texttt{200809L}'' rende disponibili le funzionalità di base
    previste dallo standard POSIX.1-2008, escludendo le estensioni XSI;
  \item in futuro valori superiori potranno abilitare ulteriori estensioni.
  \end{itemize}

\item[\macrod{\_BSD\_SOURCE}] definendo questa macro si rendono disponibili le
  funzionalità derivate da BSD4.3, insieme a quelle previste dagli standard
  ISO C, POSIX.1 e POSIX.2; alcune delle funzionalità previste da BSD sono
  però in conflitto con le corrispondenti definite nello standard POSIX.1, in
  questo caso se la macro è definita le definizioni previste da BSD4.3 avranno
  la precedenza rispetto a POSIX.

  A causa della natura dei conflitti con POSIX per ottenere una piena
  compatibilità con BSD4.3 può essere necessario anche usare una libreria di
  compatibilità, dato che alcune funzioni sono definite in modo diverso. In
  questo caso occorrerà anche usare l'opzione \cmd{-lbsd-compat} con il
  compilatore per indicargli di utilizzare le versioni nella libreria di
  compatibilità prima di quelle normali.

  Si tenga inoltre presente che la preferenza verso le versioni delle funzioni
  usate da BSD viene mantenuta soltanto se nessuna delle ulteriori macro di
  specificazione di standard successivi (vale a dire una fra
  \macro{\_POSIX\_C\_SOURCE}, \macro{\_POSIX\_SOURCE}, \macro{\_SVID\_SOURCE},
  \macro{\_XOPEN\_SOURCE}, \macro{\_XOPEN\_SOURCE\_EXTENDED} o
  \macro{\_GNU\_SOURCE}) è stata a sua volta attivata, nel qual caso queste
  hanno la precedenza. Se però si definisce \macro{\_BSD\_SOURCE} dopo aver
  definito una di queste macro, l'effetto sarà quello di dare la precedenza
  alle funzioni in forma BSD. Questa macro, essendo ricompresa in
  \macro{\_DEFAULT\_SOURCE} che è definita di default, è stata deprecata a
  partire dalla \acr{glibc} 2.20.

\item[\macrod{\_SVID\_SOURCE}] definendo questa macro si rendono disponibili le
  funzionalità derivate da SVID. Esse comprendono anche quelle definite negli
  standard ISO C, POSIX.1, POSIX.2, e X/Open (XPG$n$) illustrati in
  precedenza. Questa macro, essendo ricompresa in \macro{\_DEFAULT\_SOURCE}
  che è definita di default, è stata deprecata a partire dalla \acr{glibc}
  2.20.

\item[\macrod{\_DEFAULT\_SOURCE}] questa macro abilita le definizioni
  considerate il \textit{default}, comprese quelle richieste dallo standard
  POSIX.1-2008, ed è sostanzialmente equivalente all'insieme di
  \macro{\_SVID\_SOURCE}, \macro{\_BSD\_SOURCE} e
  \macro{\_POSIX\_C\_SOURCE}. Essendo predefinita non è necessario usarla a
  meno di non aver richiesto delle definizioni più restrittive sia con altre
  macro che con i flag del compilatore, nel qual caso abilita le funzioni che
  altrimenti sarebbero disabilitate. Questa macro è stata introdotta a partire
  dalla \acr{glibc} 2.19 e consente di deprecare \macro{\_SVID\_SOURCE} e
  \macro{\_BSD\_SOURCE}.

\item[\macrod{\_XOPEN\_SOURCE}] definendo questa macro si rendono disponibili
  le funzionalità descritte nella \textit{X/Open Portability Guide}. Anche
  queste sono un sovrainsieme di quelle definite negli standard POSIX.1 e
  POSIX.2 ed in effetti sia \macro{\_POSIX\_SOURCE} che
  \macro{\_POSIX\_C\_SOURCE} vengono automaticamente definite. Sono incluse
  anche ulteriori funzionalità disponibili in BSD e SVID, più una serie di
  estensioni a secondo dei seguenti valori:
  \begin{itemize}
  \item la definizione della macro ad un valore qualunque attiva le
    funzionalità specificate negli standard POSIX.1, POSIX.2 e XPG4;
  \item un valore di ``\texttt{500}'' o superiore rende disponibili anche le
    funzionalità introdotte con SUSv2, vale a dire la conformità ad Unix98;
  \item a partire dalla versione 2.2 della \acr{glibc} un valore uguale a
    ``\texttt{600}'' o superiore rende disponibili anche le funzionalità
    introdotte con SUSv3, corrispondenti allo standard POSIX.1-2001 più le
    estensioni XSI.
  \item a partire dalla versione 2.10 della \acr{glibc} un valore uguale a
    ``\texttt{700}'' o superiore rende disponibili anche le funzionalità
    introdotte con SUSv4, corrispondenti allo standard POSIX.1-2008 più le
    estensioni XSI.
  \end{itemize}

\item[\macrod{\_XOPEN\_SOURCE\_EXTENDED}] definendo questa macro si rendono
  disponibili le ulteriori funzionalità necessarie la conformità al
  rilascio del marchio \textit{X/Open Unix} corrispondenti allo standard
  Unix95, vale a dire quelle specificate da SUSv1/XPG4v2. Questa macro viene
  definita implicitamente tutte le volte che si imposta
  \macro{\_XOPEN\_SOURCE} ad un valore maggiore o uguale a 500.

\item[\macrod{\_ISOC99\_SOURCE}] definendo questa macro si rendono disponibili
  le funzionalità previste per la revisione delle librerie standard del C
  introdotte con lo standard ISO C99. La macro è definita a partire dalla
  versione 2.1.3 della \acr{glibc}. 

  Le versioni precedenti la serie 2.1.x riconoscevano le stesse estensioni con
  la macro \macro{\_ISOC9X\_SOURCE}, dato che lo standard non era stato
  finalizzato, ma la \acr{glibc} aveva già un'implementazione completa che
  poteva essere attivata definendo questa macro. Benché questa sia obsoleta
  viene tuttora riconosciuta come equivalente di \macro{\_ISOC99\_SOURCE} per
  compatibilità.

\item[\macrod{\_ISOC11\_SOURCE}] definendo questa macro si rendono disponibili
  le funzionalità previste per la revisione delle librerie standard del C
  introdotte con lo standard ISO C11, e abilita anche quelle previste dagli
  standard C99 e C95. La macro è definita a partire dalla versione 2.16 della
  \acr{glibc}.

\item[\macrod{\_GNU\_SOURCE}] definendo questa macro si rendono disponibili
  tutte le funzionalità disponibili nei vari standard oltre a varie estensioni
  specifiche presenti solo nella \acr{glibc} ed in Linux. Gli standard coperti
  sono: ISO C89, ISO C99, POSIX.1, POSIX.2, BSD, SVID, X/Open, SUS.

  L'uso di \macro{\_GNU\_SOURCE} è equivalente alla definizione contemporanea
  delle macro: \macro{\_BSD\_SOURCE}, \macro{\_SVID\_SOURCE},
  \macro{\_POSIX\_SOURCE}, \macro{\_ISOC99\_SOURCE}, e inoltre di
  \macro{\_POSIX\_C\_SOURCE} con valore ``\texttt{200112L}'' (o
  ``\texttt{199506L}'' per le versioni della \acr{glibc} precedenti la 2.5),
  \macro{\_XOPEN\_SOURCE\_EXTENDED} e \macro{\_XOPEN\_SOURCE} con valore 600
  (o 500 per le versioni della \acr{glibc} precedenti la 2.2); oltre a queste
  vengono pure attivate le ulteriori due macro \macro{\_ATFILE\_SOURCE} e
  \macrod{\_LARGEFILE64\_SOURCE} che definiscono funzioni previste
  esclusivamente dalla \acr{glibc}.
 
\end{basedescript}

Benché Linux supporti in maniera estensiva gli standard più diffusi, esistono
comunque delle estensioni e funzionalità specifiche, non presenti in altri
standard e lo stesso vale per la \acr{glibc}, che definisce anche delle
ulteriori funzioni di libreria. Ovviamente l'uso di queste funzionalità deve
essere evitato se si ha a cuore la portabilità, ma qualora questo non sia un
requisito esse possono rivelarsi molto utili.

Come per l'aderenza ai vari standard, le funzionalità aggiuntive possono
essere rese esplicitamente disponibili tramite la definizione di opportune
macro di preprocessore, alcune di queste vengono attivate con la definizione
di \macro{\_GNU\_SOURCE}, mentre altre devono essere attivate esplicitamente,
inoltre alcune estensioni possono essere attivate indipendentemente tramite
una opportuna macro; queste estensioni sono illustrate nel seguente elenco:

\begin{basedescript}{\desclabelwidth{2.7cm}\desclabelstyle{\nextlinelabel}}

\item[\macrod{\_LARGEFILE\_SOURCE}] definendo questa macro si rendono
  disponibili alcune funzioni che consentono di superare una inconsistenza
  presente negli standard con i file di grandi dimensioni, ed in particolare
  definire le due funzioni \func{fseeko} e \func{ftello} che al contrario
  delle corrispettive \func{fseek} e \func{ftell} usano il tipo di dato
  specifico \type{off\_t} (vedi sez.~\ref{sec:file_io}).

\item[\macrod{\_LARGEFILE64\_SOURCE}] definendo questa macro si rendono
  disponibili le funzioni di una interfaccia alternativa al supporto di valori
  a 64 bit nelle funzioni di gestione dei file (non supportati in certi
  sistemi), caratterizzate dal suffisso \texttt{64} aggiunto ai vari nomi di
  tipi di dato e funzioni (come \typed{off64\_t} al posto di \type{off\_t} o
  \funcm{lseek64} al posto di \func{lseek}).

  Le funzioni di questa interfaccia alternativa sono state proposte come una
  estensione ad uso di transizione per le \textit{Single UNIX Specification},
  per consentire la gestione di file di grandi dimensioni anche nei sistemi a
  32 bit, in cui la dimensione massima, espressa con un intero, non poteva
  superare i 2Gb.  Nei nuovi programmi queste funzioni devono essere evitate,
  a favore dell'uso macro \macro{\_FILE\_OFFSET\_BITS}, che definita al valore
  di \texttt{64} consente di usare in maniera trasparente le funzioni
  dell'interfaccia classica.

\item[\macrod{\_FILE\_OFFSET\_BITS}] la definizione di questa macro al valore
  di \texttt{64} consente di attivare la conversione automatica di tutti i
  riferimenti a dati e funzioni a 32 bit nelle funzioni di interfaccia ai file
  con le equivalenti a 64 bit, senza dover utilizzare esplicitamente
  l'interfaccia alternativa appena illustrata. In questo modo diventa
  possibile usare le ordinarie funzioni per effettuare operazioni a 64 bit sui
  file anche su sistemi a 32 bit.\footnote{basterà ricompilare il programma
    dopo averla definita, e saranno usate in modo trasparente le funzioni a 64
    bit.}

  Se la macro non è definita o è definita con valore \texttt{32} questo
  comportamento viene disabilitato, e sui sistemi a 32 bit verranno usate le
  ordinarie funzioni a 32 bit, non avendo più il supporto per file di grandi
  dimensioni. Su sistemi a 64 bit invece, dove il problema non sussiste, la
  macro non ha nessun effetto.

\item[\macrod{\_ATFILE\_SOURCE}] definendo questa macro si rendono disponibili
  le estensioni delle funzioni di creazione, accesso e modifica di file e
  directory che risolvono i problemi di sicurezza insiti nell'uso di
  \textit{pathname} relativi con programmi \textit{multi-thread} illustrate in
  sez.~\ref{sec:file_openat}. Dalla \acr{glibc} 2.10 questa macro viene
  definita implicitamente se si definisce \macro{\_POSIX\_C\_SOURCE} ad un
  valore maggiore o uguale di ``\texttt{200809L}''.

\item[\macrod{\_REENTRANT}] definendo questa macro, o la equivalente
  \macrod{\_THREAD\_SAFE} (fornita per compatibilità) si rendono disponibili
  le versioni rientranti (vedi sez.~\ref{sec:proc_reentrant}) di alcune
  funzioni, necessarie quando si usano i \textit{thread}.  Alcune di queste
  funzioni sono anche previste nello standard POSIX.1c, ma ve ne sono altre
  che sono disponibili soltanto su alcuni sistemi, o specifiche della
  \acr{glibc}, e possono essere utilizzate una volta definita la macro. Oggi
  la macro è obsoleta, già dalla \acr{glibc} 2.3 le funzioni erano già
  completamente rientranti e dalla \acr{glibc} 2.25 la macro è equivalente ad
  definire \macro{\_POSIX\_C\_SOURCE} con un valore di ``\texttt{199606L}''
  mentre se una qualunque delle altre macro che richiede un valore di
  conformità più alto è definita, la sua definizione non ha alcun effetto.

\item[\macrod{\_FORTIFY\_SOURCE}] definendo questa macro viene abilitata
  l'inserimento di alcuni controlli per alcune funzioni di allocazione e
  manipolazione di memoria e stringhe che consentono di rilevare
  automaticamente alcuni errori di \textit{buffer overflow} nell'uso delle
  stesse. La funzionalità è stata introdotta a partire dalla versione 2.3.4
  della \acr{glibc} e richiede anche il supporto da parte del compilatore, che
  è disponibile solo a partire dalla versione 4.0 del \texttt{gcc}.

  Le funzioni di libreria che vengono messe sotto controllo quando questa
  funzionalità viene attivata sono, al momento della stesura di queste note,
  le seguenti: \funcm{memcpy}, \funcm{mempcpy}, \funcm{memmove},
  \funcm{memset}, \funcm{stpcpy}, \funcm{strcpy}, \funcm{strncpy},
  \funcm{strcat}, \funcm{strncat}, \func{sprintf}, \func{snprintf},
  \func{vsprintf}, \func{vsnprintf}, e \func{gets}.

  La macro prevede due valori, con \texttt{1} vengono eseguiti dei controlli
  di base che non cambiano il comportamento dei programmi se si richiede una
  ottimizzazione di livello uno o superiore,\footnote{vale a dire se si usa
    l'opzione \texttt{-O1} o superiore del \texttt{gcc}.}  mentre con il
  valore \texttt{2} vengono aggiunti maggiori controlli. Dato che alcuni dei
  controlli vengono effettuati in fase di compilazione l'uso di questa macro
  richiede anche la collaborazione del compilatore, disponibile dalla
  versione 4.0 del \texttt{gcc}.

\end{basedescript}

Se non è stata specificata esplicitamente nessuna di queste macro il default
assunto è che siano definite \macro{\_BSD\_SOURCE}, \macro{\_SVID\_SOURCE},
\macro{\_POSIX\_SOURCE} e, con le versioni della \acr{glibc} più recenti, che
la macro \macro{\_POSIX\_C\_SOURCE} abbia il valore ``\texttt{200809L}'', per
versioni precedenti della \acr{glibc} il valore assegnato a
\macro{\_POSIX\_C\_SOURCE} era di ``\texttt{200112L}'' prima delle 2.10, di
``\texttt{199506L}'' prima delle 2.4, di ``\texttt{199506L}'' prima delle
2.1. Si ricordi infine che perché queste macro abbiano effetto devono essere
sempre definite prima dell'inclusione dei file di dichiarazione.


% vedi anche man feature_test_macros

% TODO: forse può essere interessante trattare i tunables delle glibc, vedi:
% https://siddhesh.in/posts/the-story-of-tunables.html e
% https://sourceware.org/git/gitweb.cgi?p=glibc.git;a=blob;f=README.tunables;h=0e9b0d7a470aabc2344ad997fca03c7663de0419;hb=HEAD 


% LocalWords:  like kernel multitasking scheduler preemptive sez swap is cap VM
% LocalWords:  everything bootstrap init shell Windows Foundation system call
% LocalWords:  fig libc uClib glibc embedded Library POSIX username PAM Methods
% LocalWords:  Pluggable Autentication group supplementary Name Service Switch
% LocalWords:  LDAP identifier uid gid superuser root if BSD SVr dall' American
% LocalWords:  National Institute International Organisation IEC header tab gcc
% LocalWords:  assert ctype dirent errno fcntl limits malloc setjmp signal utmp
% LocalWords:  stdarg stdio stdlib string times unistd library int short caddr
% LocalWords:  address clock dev ino inode key IPC loff nlink off pid rlim size
% LocalWords:  sigset ssize ptrdiff sys IEEE Richard Portable of TR filesystem
% LocalWords:  Operating Interface dell'IEEE Electrical and Electronics thread
% LocalWords:  Engeneers Socket NT matching regular expression scheduling l'I
% LocalWords:  XPG Portability Issue Application Programmable XTI Transport AT
% LocalWords:  socket Spec Novell Specification SUSv CDE Motif Berkley select
% LocalWords:  SunOS l'AT Sun HP Xenix Description SVID Laboratories MP hsearch
% LocalWords:  drand fmtmsg define SOURCE lbsd compat XOPEN version ISOC Large
% LocalWords:  LARGEFILE Support LFS dell' black rectangle node fill cpu draw
% LocalWords:  ellipse mem anchor west proc SysV SV Definition SCO Austin XSI
% LocalWords:  Technical TC SUS Opengroup features STRICT std ATFILE fseeko VFS
% LocalWords:  ftello fseek ftell lseek FORTIFY REENTRANT SAFE overflow memcpy
% LocalWords:  mempcpy memmove memset stpcpy strcpy strncpy strcat strncat gets
% LocalWords:  sprintf snprintf vsprintf vsnprintf syscall number calendar BITS
% LocalWords:  pathname Google Android standards device Virtual bootloader path
% LocalWords:  filename fifo name components resolution chroot parent symbolic
% LocalWords:  char block VMS raw access MacOS LF CR dos HFS Mac attributes
% LocalWords:  Executable Linkable Format Tool magic descriptor stream locking
% LocalWords:  process mount point ELF

%%% Local Variables: 
%%% mode: latex
%%% TeX-master: "gapil"
%%% End: 

%% process.tex
%%
%% Copyright (C) 2000-2018 by Simone Piccardi.  Permission is granted to
%% copy, distribute and/or modify this document under the terms of the GNU Free
%% Documentation License, Version 1.1 or any later version published by the
%% Free Software Foundation; with the Invariant Sections being "Un preambolo",
%% with no Front-Cover Texts, and with no Back-Cover Texts.  A copy of the
%% license is included in the section entitled "GNU Free Documentation
%% License".
%%

\chapter{L'interfaccia base con i processi}
\label{cha:process_interface}

Come accennato nell'introduzione il \textsl{processo} è l'unità di base con
cui un sistema unix-like alloca ed utilizza le risorse.  Questo capitolo
tratterà l'interfaccia base fra il sistema e i processi, come vengono passati
gli argomenti, come viene gestita e allocata la memoria, come un processo può
richiedere servizi al sistema e cosa deve fare quando ha finito la sua
esecuzione. Nella sezione finale accenneremo ad alcune problematiche generiche
di programmazione.

In genere un programma viene eseguito quando un processo lo fa partire
eseguendo una funzione della famiglia \func{exec}; torneremo su questo e sulla
creazione e gestione dei processi nel prossimo capitolo. In questo
affronteremo l'avvio e il funzionamento di un singolo processo partendo dal
punto di vista del programma che viene messo in esecuzione.


\section{Esecuzione e conclusione di un programma}

Uno dei concetti base di Unix è che un processo esegue sempre uno ed un solo
programma: si possono avere più processi che eseguono lo stesso programma ma
ciascun processo vedrà la sua copia del codice (in realtà il kernel fa sì che
tutte le parti uguali siano condivise), avrà un suo spazio di indirizzi,
variabili proprie e sarà eseguito in maniera completamente indipendente da
tutti gli altri. Questo non è del tutto vero nel caso di un programma
\textit{multi-thread}, ma la gestione dei \textit{thread} in Linux sarà
trattata a parte\unavref{in cap.~\ref{cha:threads}}.


\subsection{L'avvio e l'esecuzione di un programma}
\label{sec:proc_main}

\itindbeg{link-loader}
\itindbeg{shared~objects}
Quando un programma viene messo in esecuzione, cosa che può essere fatta solo
con una funzione della famiglia \func{exec} (vedi sez.~\ref{sec:proc_exec}),
il kernel esegue un opportuno codice di avvio, il cosiddetto
\textit{link-loader}, costituito dal programma \cmd{ld-linux.so}. Questo
programma è una parte fondamentale del sistema il cui compito è quello della
gestione delle cosiddette \textsl{librerie condivise}, quelle che nel mondo
Windows sono chiamate DLL (\textit{Dinamic Link Library}), e che invece in un
sistema unix-like vengono chiamate \textit{shared objects}.

Infatti, a meno di non aver specificato il flag \texttt{-static} durante la
compilazione, tutti i programmi in Linux sono compilati facendo riferimento a
librerie condivise, in modo da evitare di duplicare lo stesso codice nei
relativi eseguibili e consentire un uso più efficiente della memoria, dato che
il codice di uno \textit{shared objects} viene caricato in memoria dal kernel
una sola volta per tutti i programmi che lo usano.
\itindend{shared~objects}

Questo significa però che normalmente il codice di un programma è incompleto,
contenendo solo i riferimenti alle funzioni di libreria che vuole utilizzare e
non il relativo codice. Per questo motivo all'avvio del programma è necessario
l'intervento del \textit{link-loader} il cui compito è caricare in memoria le
librerie condivise eventualmente assenti, ed effettuare poi il collegamento
dinamico del codice del programma alle funzioni di libreria da esso utilizzate
prima di metterlo in esecuzione.

Il funzionamento di \cmd{ld-linux.so} è controllato da alcune variabili di
ambiente e dal contenuto del file \conffile{/etc/ld.so.conf} che consentono di
elencare le directory un cui cercare le librerie e determinare quali verranno
utilizzate.  In particolare con la variabile di ambiente
\envvar{LD\_LIBRARY\_PATH} si possono indicare ulteriori directory rispetto a
quelle di sistema in cui inserire versioni personali delle librerie che hanno
la precedenza su quelle di sistema, mentre con la variabile di ambiente
\envvar{LD\_PRELOAD} si può passare direttamente una lista di file di librerie
condivise da usare al posto di quelli di sistema. In questo modo è possibile
effettuare lo sviluppo o il test di nuove librerie senza dover sostituire
quelle di sistema. Ulteriori dettagli sono riportati nella pagina di manuale
di \cmd{ld.so} e per un approfondimento dell'argomento si può consultare
sez.~3.1.2 di \cite{AGL}.

Una volta completate le operazioni di inizializzazione di \cmd{ld-linux.so}, il
sistema fa partire qualunque programma chiamando la funzione \code{main}. Sta
al programmatore chiamare così la funzione principale del programma da cui si
suppone che inizi l'esecuzione. In ogni caso senza questa funzione lo stesso
\textit{link-loader} darebbe luogo ad errori.  Lo standard ISO C specifica che
la funzione \code{main} può non avere argomenti o prendere due argomenti che
rappresentano gli argomenti passati da linea di comando (su cui torneremo in
sez.~\ref{sec:proc_par_format}), in sostanza un prototipo che va sempre bene è
il seguente:
\includecodesnip{listati/main_def.c}

\itindend{link-loader}

In realtà nei sistemi Unix esiste un altro modo per definire la funzione
\code{main}, che prevede la presenza di un terzo argomento, \code{char
  *envp[]}, che fornisce l'\textsl{ambiente} del programma; questa forma però
non è prevista dallo standard POSIX.1 per cui se si vogliono scrivere
programmi portabili è meglio evitarla. Per accedere all'ambiente, come vedremo
in sez.~\ref{sec:proc_environ} si usa in genere una variabile globale che
viene sempre definita automaticamente.

Ogni programma viene fatto partire mettendo in esecuzione il codice contenuto
nella funzione \code{main}, ogni altra funzione usata dal programma, che sia
ottenuta da una libreria condivisa, o che sia direttamente definita nel
codice, dovrà essere invocata a partire dal codice di \code{main}. Nel caso di
funzioni definite nel programma occorre tenere conto che, nel momento stesso
in cui si usano le librerie di sistema (vale a dire la \acr{glibc}) alcuni
nomi sono riservati e non possono essere utilizzati. 

In particolare sono riservati a priori e non possono essere mai ridefiniti in
nessun caso i nomi di tutte le funzioni, le variabili, le macro di
preprocessore, ed i tipi di dati previsti dallo standard ISO C. Lo stesso
varrà per tutti i nomi definiti negli \textit{header file} che si sono
esplicitamente inclusi nel programma (vedi sez.~\ref{sec:proc_syscall}), ma
anche se è possibile riutilizzare nomi definiti in altri \textit{header file}
la pratica è da evitare nella maniera più assoluta per non generare ambiguità.

Oltre ai nomi delle funzioni di libreria sono poi riservati in maniera
generica tutti i nomi di variabili o funzioni globali che iniziano con il
carattere di sottolineato (``\texttt{\_}''), e qualunque nome che inizi con il
doppio sottolineato (``\texttt{\_\_}'') o con il sottolineato seguito da
lettera maiuscola. Questi identificativi infatti sono utilizzati per i nomi
usati internamente in forma privata dalle librerie, ed evitandone l'uso si
elimina il rischio di conflitti.

Infine esiste una serie di classi di nomi che sono riservati per un loro
eventuale uso futuro da parte degli standard ISO C e POSIX.1, questi in teoria
possono essere usati senza problemi oggi, ma potrebbero dare un conflitto con
una revisione futura di questi standard, per cui è comunque opportuno
evitarli, in particolare questi sono:
\begin{itemize*}
\item i nomi che iniziano per ``\texttt{E}'' costituiti da lettere maiuscole e
  numeri, che potrebbero essere utilizzati per nuovi codici di errore (vedi
  sez.~\ref{sec:sys_errors}),
\item i nomi che iniziano con ``\texttt{is}'' o ``\texttt{to}'' e costituiti
  da lettere minuscole che potrebbero essere utilizzati da nuove funzioni per
  il controllo e la conversione del tipo di caratteri,
\item i nomi che iniziano con ``\texttt{LC\_}'' e costituiti
  da lettere maiuscole che possono essere usato per macro attinenti la
  localizzazione,% mettere in seguito (vedi sez.~\ref{sec:proc_localization}),
\item nomi che iniziano con ``\texttt{SIG}'' o ``\texttt{SIG\_}'' e costituiti
  da lettere maiuscole che potrebbero essere usati per nuovi nomi di segnale
  (vedi sez.~\ref{sec:sig_classification}),
\item nomi che iniziano con ``\texttt{str}'', ``\texttt{mem}'', o
  ``\texttt{wcs}'' e costituiti da lettere minuscole che possono essere
  utilizzati per funzioni attinenti la manipolazione delle stringhe e delle
  aree di memoria,
\item nomi che terminano in ``\texttt{\_t}'' che potrebbero essere utilizzati
  per la definizione di nuovi tipi di dati di sistema oltre quelli di
  tab.~\ref{tab:intro_primitive_types}).
\end{itemize*}


\subsection{Chiamate a funzioni e \textit{system call}}
\label{sec:proc_syscall}

Come accennato in sez.~\ref{sec:intro_syscall} un programma può utilizzare le
risorse che il sistema gli mette a disposizione attraverso l'uso delle
opportune \textit{system call}. Abbiamo inoltre appena visto come all'avvio un
programma venga messo in grado di chiamare le funzioni fornite da eventuali
librerie condivise da esso utilizzate. 

Vedremo nel resto della guida quali sono le risorse del sistema accessibili
attraverso le \textit{system call} e tratteremo buona parte delle funzioni
messe a disposizione dalla libreria standard del C, in questa sezione però si
forniranno alcune indicazioni generali sul come fare perché un programma possa
utilizzare queste funzioni.

\itindbeg{header~file}

In sez.~\ref{sec:intro_standard} abbiamo accennato come le funzioni definite
nei vari standard siano definite in una serie di \textit{header file} (in
italiano \textsl{file di intestazione}).  Vengono chiamati in questo modo quei
file, forniti insieme al codice delle librerie, che contengono le
dichiarazioni delle variabili, dei tipi di dati, delle macro di preprocessore
e soprattutto delle funzioni che fanno parte di una libreria.

Questi file sono necessari al compilatore del linguaggio C per ottenere i
riferimenti ai nomi delle funzioni (e alle altre risorse) definite in una
libreria, per questo quando si vogliono usare le funzioni di una libreria
occorre includere nel proprio codice gli \textit{header file} che le
definiscono con la direttiva \code{\#include}. Dato che le funzioni devono
essere definite prima di poterle usare in genere gli \textit{header file}
vengono inclusi all'inizio del programma. Se inoltre si vogliono utilizzare le
macro di controllo delle funzionalità fornite dai vari standard illustrate in
sez.~\ref{sec:intro_gcc_glibc_std} queste, come accennato, dovranno a loro
volta essere definite prima delle varie inclusioni.

Ogni libreria fornisce i propri file di intestazione per i quali si deve
consultare la documentazione, ma in tab.~\ref{tab:intro_posix_header} si sono
riportati i principali \textit{header file} definiti nella libreria standard
del C (nel caso la \acr{glibc}) che contengono le varie funzioni previste
negli standard POSIX ed ANSI C, e che prevedono la definizione sia delle
funzioni di utilità generica che delle interfacce alle \textit{system call}. In
seguito per ciascuna funzione o \textit{system call} che tratteremo
indicheremo anche quali sono gli \textit{header file} contenenti le necessarie
definizioni.

\begin{table}[htb]
  \footnotesize
  \centering
  \begin{tabular}[c]{|l|c|c|l|}
    \hline
    \multirow{2}{*}{\textbf{Header}}&
    \multicolumn{2}{|c|}{\textbf{Standard}}&
    \multirow{2}{*}{\textbf{Contenuto}} \\
    \cline{2-3}
    & ANSI C& POSIX& \\
    \hline
    \hline
    \headfiled{assert.h}&$\bullet$&    --   & Verifica le asserzioni fatte in un
                                              programma.\\ 
    \headfiled{ctype.h} &$\bullet$&    --   & Tipi standard.\\
    \headfiled{dirent.h}&   --    &$\bullet$& Manipolazione delle directory.\\
    \headfiled{errno.h} &   --    &$\bullet$& Errori di sistema.\\
    \headfiled{fcntl.h} &   --    &$\bullet$& Controllo sulle opzioni dei
                                              file.\\ 
    \headfiled{limits.h}&   --    &$\bullet$& Limiti e parametri del sistema.\\
    \headfiled{malloc.h}&$\bullet$&    --   & Allocazione della memoria.\\
    \headfiled{setjmp.h}&$\bullet$&    --   & Salti non locali.\\
    \headfiled{signal.h}&   --    &$\bullet$& Gestione dei segnali.\\
    \headfiled{stdarg.h}&$\bullet$&    --   & Gestione di funzioni a argomenti
                                             variabili.\\ 
    \headfiled{stdio.h} &$\bullet$&    --   & I/O bufferizzato in standard ANSI
                                              C.\\ 
    \headfiled{stdlib.h}&$\bullet$&    --   & Definizioni della libreria
                                              standard.\\ 
    \headfiled{string.h}&$\bullet$&    --   & Manipolazione delle stringhe.\\
    \headfiled{time.h}  &   --    &$\bullet$& Gestione dei tempi.\\
    \headfiled{times.h} &$\bullet$&    --   & Gestione dei tempi.\\
    \headfiled{unistd.h}&   --    &$\bullet$& Unix standard library.\\
    \headfiled{utmp.h}  &   --    &$\bullet$& Registro connessioni utenti.\\
    \hline
  \end{tabular}
  \caption{Elenco dei principali \textit{header file} definiti dagli standard
    POSIX e ANSI C.}
  \label{tab:intro_posix_header}
\end{table}

Un esempio di inclusione di questi file, preso da uno dei programmi di
esempio, è il seguente, e si noti come gli \textit{header file} possano essere
referenziati con il nome fra parentesi angolari, nel qual caso si indica l'uso
di quelli installati con il sistema,\footnote{in un sistema GNU/Linux che
  segue le specifiche del \textit{Filesystem Hierarchy Standard} (per maggiori
  informazioni si consulti sez.~1.2.3 di \cite{AGL}) si trovano sotto
  \texttt{/usr/include}.}  o fra virgolette, nel qual caso si fa riferimento
ad una versione locale, da indicare con un \textit{pathname} relativo:
\includecodesnip{listati/main_include.c}

Si tenga presente che oltre ai nomi riservati a livello generale di cui si è
parlato in sez.~\ref{sec:proc_main}, alcuni di questi \textit{header file}
riservano degli ulteriori identificativi, il cui uso sarà da evitare, ad
esempio si avrà che:
\begin{itemize*}
\item in \headfile{dirent.h} vengono riservati i nomi che iniziano con
  ``\texttt{d\_}'' e costituiti da lettere minuscole,
\item in \headfile{fcntl.h} vengono riservati i nomi che iniziano con
  ``\texttt{l\_}'', ``\texttt{F\_}'',``\texttt{O\_}'' e ``\texttt{S\_}'',
\item in \headfile{limits.h} vengono riservati i nomi che finiscono in
  ``\texttt{\_MAX}'',
\item in \headfile{signal.h} vengono riservati i nomi che iniziano con
  ``\texttt{sa\_}'' e ``\texttt{SA\_}'',
\item in \headfile{sys/stat.h} vengono riservati i nomi che iniziano con
  ``\texttt{st\_}'' e ``\texttt{S\_}'',
\item in \headfile{sys/times.h} vengono riservati i nomi che iniziano con
  ``\texttt{tms\_}'',
\item in \headfile{termios.h} vengono riservati i nomi che iniziano con
  ``\texttt{c\_}'', ``\texttt{V}'', ``\texttt{I}'', ``\texttt{O}'' e
  ``\texttt{TC}'' e con ``\texttt{B}'' seguito da un numero,
\item in \headfile{grp.h} vengono riservati i nomi che iniziano con
  ``\texttt{gr\_}'',
\item in \headfile{pwd.h} vengono riservati i nomi che iniziano con
  ``\texttt{pw\_}'',
\end{itemize*}

\itindend{header~file}

Una volta inclusi gli \textit{header file} necessari un programma potrà
richiamare le funzioni di libreria direttamente nel proprio codice ed accedere
ai servizi del kernel; come accennato infatti normalmente ogni \textit{system
  call} è associata ad una omonima funzione di libreria, che è quella che si
usa normalmente per invocarla.

Occorre però tenere presente che anche se dal punto di vista della scrittura
del codice la chiamata di una \textit{system call} non è diversa da quella di
una qualunque funzione ordinaria, la situazione è totalmente diversa
nell'esecuzione del programma. Una funzione ordinaria infatti viene eseguita,
esattamente come il codice che si è scritto nel corpo del programma, in
\textit{user space}. Quando invece si esegue una \textit{system call}
l'esecuzione ordinaria del programma viene interrotta con quello che viene
usualmente chiamato un \itindex{context~switch} \textit{context
  switch};\footnote{in realtà si parla più comunemente di \textit{context
    switch} quando l'esecuzione di un processo viene interrotta dal kernel
  (tramite lo \textit{scheduler}) per metterne in esecuzione un altro, ma il
  concetto generale resta lo stesso: l'esecuzione del proprio codice in
  \textit{user space} viene interrotta e lo stato del processo deve essere
  salvato per poterne riprendere l'esecuzione in un secondo tempo.}  il
contesto di esecuzione del processo viene salvato in modo da poterne
riprendere in seguito l'esecuzione ed i dati forniti (come argomenti della
chiamata) vengono trasferiti al kernel che esegue il codice della
\textit{system call} (che è codice del kernel) in \textit{kernel space}; al
completamento della \textit{system call} i dati salvati nel \textit{context
  switch} saranno usati per riprendere l'esecuzione ordinaria del programma.

Dato che il passaggio dei dati ed il salvataggio del contesto di esecuzione
sono operazioni critiche per le prestazioni del sistema, per rendere il più
veloce possibile questa operazione sono state sviluppate una serie di
ottimizzazioni che richiedono alcune preparazioni abbastanza complesse dei
dati, che in genere dipendono dall'architettura del processore e sono scritte
direttamente in \textit{assembler}.


%
% TODO:trattare qui, quando sarà il momento vsyscall e vDSO, vedi:
% http://davisdoesdownunder.blogspot.com/2011/02/linux-syscall-vsyscall-and-vdso-oh-my.html 
% http://www.win.tue.nl/~aeb/linux/lk/lk-4.html
%
% Altro materiale al riguardo http://lwn.net/Articles/615809/
% http://man7.org/linux/man-pages/man7/vdso.7.html 

Inoltre alcune \textit{system call} sono state modificate nel corso degli anni
con lo sviluppo del kernel per aggiungere ad esempio funzionalità in forma di
nuovi argomenti, o per consolidare diverse varianti in una interfaccia
generica.  Per questo motivo dovendo utilizzare una \textit{system call} è
sempre preferibile usare l'interfaccia fornita dalla \textsl{glibc}, che si
cura di mantenere una uniformità chiamando le versioni più aggiornate.

Ci sono alcuni casi però in cui può essere necessario evitare questa
associazione, e lavorare a basso livello con una specifica versione, oppure si
può voler utilizzare una \textit{system call} che non è stata ancora associata
ad una funzione di libreria.  In tal caso, per evitare di dover effettuare
esplicitamente le operazioni di preparazione citate, all'interno della
\textsl{glibc} è fornita una specifica funzione,
\funcd{syscall},\footnote{fino a prima del kernel 2.6.18 per l'esecuzione
  diretta delle \textit{system call} erano disponibili anche una serie di
  macro \texttt{\_syscall\textsl{N}} (con $N$ pari al numero di argomenti
  della \textit{system call}); queste sono deprecate e pertanto non ne
  parleremo ulteriormente.} che consente eseguire direttamente una
\textit{system call}; il suo prototipo, accessibile se si è definita la macro
\macro{\_GNU\_SOURCE}, è:

\begin{funcproto}{
  \fhead{unistd.h}
  \fhead{sys/syscall.h}
  \fdecl{long syscall(int number, ...)}
  \fdesc{Esegue la \textit{system call} indicata da \param{number}.}
}
{La funzione ritorna un intero dipendente dalla \textit{system call} invocata,
 in generale $0$ indica il successo ed un valore negativo un errore.}
\end{funcproto}

La funzione richiede come primo argomento il numero della \textit{system call}
da invocare, seguita dagli argomenti da passare alla stessa, che ovviamente
dipendono da quest'ultima, e restituisce il codice di ritorno della
\textit{system call} invocata. In generale un valore nullo indica il successo
ed un valore negativo è un codice di errore che poi viene memorizzato nella
variabile \var{errno} (sulla gestione degli errori torneremo in dettaglio in
sez.~\ref{sec:sys_errors}).

Il valore di \param{number} dipende sia dalla versione di kernel che
dall'architettura,\footnote{in genere le vecchie \textit{system call} non
  vengono eliminate e se ne aggiungono di nuove con nuovi numeri.}  ma
ciascuna \textit{system call} viene in genere identificata da una costante
nella forma \texttt{SYS\_*} dove al prefisso viene aggiunto il nome che spesso
corrisponde anche alla omonima funzione di libreria. Queste costanti sono
definite nel file \headfiled{sys/syscall.h}, ma si possono anche usare
direttamente valori numerici.


\subsection{La terminazione di un programma}
\label{sec:proc_conclusion}

Normalmente un programma conclude la sua esecuzione quando si fa ritornare la
funzione \code{main}, si usa cioè l'istruzione \instruction{return} del
linguaggio C all'interno della stessa, o se si richiede esplicitamente la
chiusura invocando direttamente la funzione \func{exit}. Queste due modalità
sono assolutamente equivalenti, dato che \func{exit} viene chiamata in maniera
trasparente anche quando \code{main} ritorna, passandogli come argomento il
valore indicato da \instruction{return}.

La funzione \funcd{exit}, che è completamente generale, essendo definita dallo
standard ANSI C, è quella che deve essere invocata per una terminazione
``\textit{normale}'', il suo prototipo è:

\begin{funcproto}{
  \fhead{unistd.h}
  \fdecl{void exit(int status)}
  \fdesc{Causa la conclusione ordinaria del programma.}
}
{La funzione non ritorna, il processo viene terminato.}
\end{funcproto}

La funzione è pensata per eseguire una conclusione pulita di un programma che
usi la libreria standard del C; essa esegue tutte le funzioni che sono state
registrate con \func{atexit} e \func{on\_exit} (vedi
sez.~\ref{sec:proc_atexit}), chiude tutti gli \textit{stream} (vedi
sez.~\ref{sec:file_stream}) effettuando il salvataggio dei dati sospesi
(chiamando \func{fclose}, vedi sez.~\ref{sec:file_fopen}), infine passa il
controllo al kernel chiamando la \textit{system call} \func{\_exit} (che
vedremo a breve) che completa la terminazione del processo.

\itindbeg{exit~status}

Il valore dell'argomento \param{status} o il valore di ritorno di \code{main}
costituisce quello che viene chiamato lo \textsl{stato di uscita}
(l'\textit{exit status}) del processo. In generale si usa questo valore per
fornire al processo padre (come vedremo in sez.~\ref{sec:proc_wait}) delle
informazioni generiche sulla riuscita o il fallimento del programma appena
terminato.

Anche se l'argomento \param{status} (ed il valore di ritorno di \code{main})
sono numeri interi di tipo \ctyp{int}, si deve tener presente che il valore
dello stato di uscita viene comunque troncato ad 8 bit, per cui deve essere
sempre compreso fra 0 e 255. Si tenga presente che se si raggiunge la fine
della funzione \code{main} senza ritornare esplicitamente si ha un valore di
uscita indefinito, è pertanto consigliabile di concludere sempre in maniera
esplicita detta funzione.

Non esiste un significato intrinseco della stato di uscita, ma una convenzione
in uso pressoché universale è quella di restituire 0 in caso di successo e 1
in caso di fallimento. Una eccezione a questa convenzione è per i programmi
che effettuano dei confronti (come \cmd{diff}), che usano 0 per indicare la
corrispondenza, 1 per indicare la non corrispondenza e 2 per indicare
l'incapacità di effettuare il confronto. Un'altra convenzione riserva i valori
da 128 a 256 per usi speciali: ad esempio 128 viene usato per indicare
l'incapacità di eseguire un altro programma in un sottoprocesso. Benché le
convenzioni citate non siano seguite universalmente è una buona idea tenerle
presenti ed adottarle a seconda dei casi.

Si tenga presente inoltre che non è una buona idea usare eventuali codici di
errore restituiti nella variabile \var{errno} (vedi sez.~\ref{sec:sys_errors})
come \textit{exit status}. In generale infatti non ci si cura del valore dello
stato di uscita di un processo se non per vedere se è diverso da zero, come
indicazione di un qualche errore.  Dato che viene troncato ad 8 bit utilizzare
un intero di valore generico può comportare il rischio, qualora si vada ad
usare un multiplo di 256, di avere uno stato di uscita uguale a zero, che
verrebbe interpretato come un successo.

Per questo motivo in \headfile{stdlib.h} sono definite, seguendo lo standard
POSIX, le due costanti \constd{EXIT\_SUCCESS} e \constd{EXIT\_FAILURE}, da
usare sempre per specificare lo stato di uscita di un processo. Su Linux, ed
in generale in qualunque sistema POSIX, ad esse sono assegnati rispettivamente
i valori 0 e 1.

\itindend{exit~status}

Una forma alternativa per effettuare una terminazione esplicita di un
programma è quella di chiamare direttamente la \textit{system call}
\funcd{\_exit},\footnote{la stessa è definita anche come \funcd{\_Exit} in
  \headfile{stdlib.h}, inoltre a partire dalla \acr{glibc} 2.3 usando questa
  funzione viene invocata \func{exit\_group} che termina tutti i
  \textit{thread} del processo e non solo quello corrente (fintanto che non si
  usano i \textit{thread}\unavref{, vedi sez.~\ref{cha:threads},} questo non
  fa nessuna differenza).} che restituisce il controllo direttamente al
kernel, concludendo immediatamente il processo, il suo prototipo è:

\begin{funcproto}{ \fhead{unistd.h} \fdecl{void \_exit(int status)}
    \fdesc{Causa la conclusione immediata del programma.}  } {La funzione non
    ritorna, il processo viene terminato.}
\end{funcproto}

La funzione termina immediatamente il processo e le eventuali funzioni
registrate con \func{atexit} e \func{on\_exit} non vengono eseguite. La
funzione chiude tutti i file descriptor appartenenti al processo, cosa che
però non comporta il salvataggio dei dati eventualmente presenti nei buffer
degli \textit{stream}, (torneremo sulle due interfacce dei file in
sez.~\ref{sec:file_unix_interface} e
sez.~\ref{sec:files_std_interface}). Infine fa sì che ogni figlio del processo
sia adottato da \cmd{init} (vedi sez.~\ref{sec:proc_termination}), manda un
segnale \signal{SIGCHLD} al processo padre (vedi
sez.~\ref{sec:sig_job_control}) e salva lo stato di uscita specificato in
\param{status} che può essere raccolto usando la funzione \func{wait} (vedi
sez.~\ref{sec:proc_wait}).

Si tenga presente infine che oltre alla conclusione ``\textsl{normale}''
appena illustrata esiste anche la possibilità di una conclusione
``\textsl{anomala}'' del programma a causa della ricezione di un segnale
(tratteremo i segnali in cap.~\ref{cha:signals}) o della chiamata alla
funzione \func{abort}; torneremo su questo in sez.~\ref{sec:proc_termination}.


\subsection{Esecuzione di funzioni preliminari all'uscita}
\label{sec:proc_atexit}

Un'esigenza comune che si incontra è quella di dover effettuare una serie di
operazioni di pulizia prima della conclusione di un programma, ad esempio
salvare dei dati, ripristinare delle impostazioni, eliminare dei file
temporanei, ecc. In genere queste operazioni vengono fatte in un'apposita
sezione del programma, ma quando si realizza una libreria diventa antipatico
dover richiedere una chiamata esplicita ad una funzione di pulizia al
programmatore che la utilizza.

È invece molto meno soggetto ad errori, e completamente trasparente
all'utente, avere la possibilità di fare effettuare automaticamente la
chiamata ad una funzione che effettui tali operazioni all'uscita dal
programma. A questo scopo lo standard ANSI C prevede la possibilità di
registrare un certo numero di funzioni che verranno eseguite all'uscita dal
programma,\footnote{nel caso di \func{atexit} lo standard POSIX.1-2001
  richiede che siano registrabili almeno \constd{ATEXIT\_MAX} funzioni (il
  valore può essere ottenuto con \func{sysconf}, vedi
  sez.~\ref{sec:sys_limits}).} sia per la chiamata ad \func{exit} che per il
ritorno di \code{main}. La prima funzione che si può utilizzare a tal fine è
\funcd{atexit}, il cui prototipo è:

\begin{funcproto}{ 
\fhead{stdlib.h} 
\fdecl{int atexit(void (*function)(void))}
\fdesc{Registra la funzione \param{function} per la chiamata all'uscita
      dal programma.}  
} 
{La funzione ritorna $0$ in caso di successo e $-1$ per un errore, \var{errno}
  non viene modificata.}
\end{funcproto}

La funzione richiede come argomento \param{function} l'indirizzo di una
opportuna funzione di pulizia da chiamare all'uscita del programma, che non
deve prendere argomenti e non deve ritornare niente. In sostanza deve la
funzione di pulizia dovrà essere definita come \code{void function(void)}.

Un'estensione di \func{atexit} è la funzione \funcd{on\_exit}, che la
\acr{glibc} include per compatibilità con SunOS ma che non è detto sia
definita su altri sistemi,\footnote{la funzione è disponibile dalla
  \acr{glibc} 2.19 definendo la macro \macro{\_DEFAULT\_SOURCE}, mentre in
  precedenza erano necessarie \macro{\_BSD\_SOURCE} o \macro{\_SVID\_SOURCE};
  non essendo prevista dallo standard POSIX è in generale preferibile evitarne
  l'uso.} il suo prototipo è:

\begin{funcproto}{ 
\fhead{stdlib.h} 
\fdecl{int on\_exit(void (*function)(int, void *), void *arg))}
\fdesc{Registra la funzione \param{function} per la chiamata all'uscita dal
  programma.} 
}
{La funzione ritorna $0$ in caso di successo e $-1$ per un errore, \var{errno}
  non viene modificata.} 
\end{funcproto}

In questo caso la funzione da chiamare all'uscita prende i due argomenti
specificati nel prototipo, un intero ed un puntatore; dovrà cioè essere
definita come \code{void function(int status, void *argp)}. Il primo argomento
sarà inizializzato allo stato di uscita con cui è stata chiamata \func{exit}
ed il secondo al puntatore \param{arg} passato come secondo argomento di
\func{on\_exit}.  Così diventa possibile passare dei dati alla funzione di
chiusura.

Nella sequenza di chiusura tutte le funzioni registrate verranno chiamate in
ordine inverso rispetto a quello di registrazione, ed una stessa funzione
registrata più volte sarà chiamata più volte. Siccome entrambe le funzioni
\func{atexit} e \func{on\_exit} fanno riferimento alla stessa lista, l'ordine
di esecuzione sarà riferito alla registrazione in quanto tale,
indipendentemente dalla funzione usata per farla.

Una volta completata l'esecuzione di tutte le funzioni registrate verranno
chiusi tutti gli \textit{stream} aperti ed infine verrà chiamata \func{\_exit}
per la terminazione del programma. Questa è la sequenza ordinaria, eseguita a
meno che una delle funzioni registrate non esegua al suo interno
\func{\_exit}, nel qual caso la terminazione del programma sarà immediata ed
anche le successive funzioni registrate non saranno invocate.

Se invece all'interno di una delle funzioni registrate si chiama un'altra
volta \func{exit} lo standard POSIX.1-2001 prescrive un comportamento
indefinito, con la possibilità (che su Linux comunque non c'è) di una
ripetizione infinita. Pertanto questa eventualità è da evitare nel modo più
assoluto. Una altro comportamento indefinito si può avere se si termina
l'esecuzione di una delle funzioni registrate con \func{longjmp} (vedi
sez.~\ref{sec:proc_longjmp}).

Si tenga presente infine che in caso di terminazione anomala di un processo
(ad esempio a causa di un segnale) nessuna delle funzioni registrate verrà
eseguita e che se invece si crea un nuovo processo con \func{fork} (vedi
sez.~\ref{sec:proc_fork}) questo manterrà tutte le funzioni già registrate.


\subsection{Un riepilogo}
\label{sec:proc_term_conclusion}

Data l'importanza dell'argomento è opportuno un piccolo riepilogo dei fatti
essenziali relativi alla esecuzione di un programma. Il primo punto da
sottolineare è che in un sistema unix-like l'unico modo in cui un programma
può essere eseguito dal kernel è attraverso la chiamata alla \textit{system
  call} \func{execve}, sia direttamente che attraverso una delle funzioni
della famiglia \func{exec} che ne semplificano l'uso (vedi
sez.~\ref{sec:proc_exec}).

Allo stesso modo l'unico modo in cui un programma può concludere
volontariamente la propria esecuzione è attraverso una chiamata alla
\textit{system call} \func{\_exit}, sia che questa venga fatta esplicitamente,
o in maniera indiretta attraverso l'uso di \func{exit} o il ritorno di
\code{main}. 

Uno schema riassuntivo che illustra le modalità con cui si avvia e conclude
normalmente un programma è riportato in fig.~\ref{fig:proc_prog_start_stop}.

\begin{figure}[htb]
  \centering
  \includegraphics[width=9cm]{img/proc_beginend}
  % \begin{tikzpicture}[>=stealth]
  %   \filldraw[fill=black!35] (-0.3,0) rectangle (12,1);
  %   \draw(5.5,0.5) node {\large{\textsf{kernel}}};

  %   \filldraw[fill=black!15] (1.5,2) rectangle (4,3);
  %   \draw (2.75,2.5) node {\texttt{ld-linux.so}};
  %   \draw [->] (2.75,1) -- (2.75,2);
  %   \draw (2.75,1.5) node [anchor=west]{\texttt{execve}};

  %   \filldraw[fill=black!15,rounded corners] (1.5,4) rectangle (4,5);
  %   \draw (2.75,4.5) node {\texttt{main}};

  %   \draw [<->, dashed] (2.75,3) -- (2.75,4);
  %   \draw [->] (1.5,4.5) -- (0.3,4.5) -- (0.3,1);
  %   \draw (0.9,4.5) node [anchor=south] {\texttt{\_exit}};

  %   \filldraw[fill=black!15,rounded corners] (1.5,6) rectangle (4,7);
  %   \draw (2.75,6.5) node {\texttt{funzione}};

  %   \draw [<->, dashed] (2.75,5) -- (2.75,6);
  %   \draw [->] (1.5,6.5) -- (0.05,6.5) -- (0.05,1);
  %   \draw (0.9,6.5) node [anchor=south] {\texttt{\_exit}};

  %   \draw (6.75,4.5) node (exit) [rectangle,fill=black!15,minimum width=2.5cm,minimum height=1cm,rounded corners, draw]{\texttt{exit}};

  %   \draw[->] (4,6.5) -- node[anchor=south west]{\texttt{exit}} (exit);
  %   \draw[->] (4,4.5) -- node[anchor=south]{\texttt{exit}} (exit);
  %   \draw[->] (exit) -- node[anchor=east]{\texttt{\_exit}}(6.75,1);

  %   \draw (10,4.5) node (exithandler1) [rectangle,fill=black!15,rounded corners, draw]{\textsf{exit handler}};
  %   \draw (10,5.5) node (exithandler2) [rectangle,fill=black!15,rounded corners, draw]{\textsf{exit handler}};
  %   \draw (10,3.5) node (stream) [rectangle,fill=black!15,rounded corners, draw]{\textsf{chiusura stream}};

  %   \draw[<->, dashed] (exithandler1) -- (exit);
  %   \draw[<->, dashed] (exithandler2) -- (exit);
  %   \draw[<->, dashed] (stream) -- (exit);
  % \end{tikzpicture}
  \caption{Schema dell'avvio e della conclusione di un programma.}
  \label{fig:proc_prog_start_stop}
\end{figure}

Si ricordi infine che un programma può anche essere interrotto dall'esterno
attraverso l'uso di un segnale (modalità di conclusione non mostrata in
fig.~\ref{fig:proc_prog_start_stop}); tratteremo nei dettagli i segnali e la
loro gestione nel capitolo \ref{cha:signals}.



\section{I processi e l'uso della memoria}
\label{sec:proc_memory}

Una delle risorse più importanti che ciascun processo ha a disposizione è la
memoria, e la gestione della memoria è appunto uno degli aspetti più complessi
di un sistema unix-like. In questa sezione, dopo una breve introduzione ai
concetti di base, esamineremo come la memoria viene vista da parte di un
programma in esecuzione, e le varie funzioni utilizzabili per la sua gestione.


\subsection{I concetti generali}
\label{sec:proc_mem_gen}

\index{memoria~virtuale|(}

Ci sono vari modi in cui i sistemi operativi organizzano la memoria, ed i
dettagli di basso livello dipendono spesso in maniera diretta
dall'architettura dell'hardware, ma quello più tipico, usato dai sistemi
unix-like come Linux è la cosiddetta \textsl{memoria virtuale} che consiste
nell'assegnare ad ogni processo uno spazio virtuale di indirizzamento lineare,
in cui gli indirizzi vanno da zero ad un qualche valore massimo.\footnote{nel
  caso di Linux fino al kernel 2.2 detto massimo era, per macchine a 32bit, di
  2Gb. Con il kernel 2.4 ed il supporto per la \textit{high-memory} il limite
  è stato esteso anche per macchine a 32 bit.}  Come accennato nel
cap.~\ref{cha:intro_unix} questo spazio di indirizzi è virtuale e non
corrisponde all'effettiva posizione dei dati nella RAM del computer. In
generale detto spazio non è neppure continuo, cioè non tutti gli indirizzi
possibili sono utilizzabili, e quelli usabili non sono necessariamente
adiacenti.

\itindbeg{huge~page}

Per la gestione da parte del kernel la memoria viene divisa in pagine di
dimensione fissa. Inizialmente queste pagine erano di 4kb sulle macchine a 32
bit e di 8kb sulle alpha. Con le versioni più recenti del kernel è possibile
anche utilizzare pagine di dimensioni maggiori (di 4Mb, dette \textit{huge
  page}), per sistemi con grandi quantitativi di memoria in cui l'uso di
pagine troppo piccole comporta una perdita di prestazioni. In alcuni sistemi
la costante \constd{PAGE\_SIZE}, definita in \headfile{limits.h}, indica la
dimensione di una pagina in byte, con Linux questo non avviene e per ottenere
questa dimensione si deve ricorrere alla funzione \func{getpagesize} (vedi
sez.~\ref{sec:sys_memory_res}).

\itindend{huge~page}
\itindbeg{page~table}

Ciascuna pagina di memoria nello spazio di indirizzi virtuale è associata ad
un supporto che può essere una pagina di memoria reale o ad un dispositivo di
stoccaggio secondario (come lo spazio disco riservato alla \textit{swap}, o i
file che contengono il codice). Per ciascun processo il kernel si cura di
mantenere un mappa di queste corrispondenze nella cosiddetta \textit{page
  table}.\footnote{questa è una semplificazione brutale, il meccanismo è molto
  più complesso; una buona trattazione di come Linux gestisce la memoria
  virtuale si trova su \cite{LinVM}.}

\itindend{page~table}

Una stessa pagina di memoria reale può fare da supporto a diverse pagine di
memoria virtuale appartenenti a processi diversi, come accade in genere per le
pagine che contengono il codice delle librerie condivise. Ad esempio il codice
della funzione \func{printf} starà su una sola pagina di memoria reale che
farà da supporto a tutte le pagine di memoria virtuale di tutti i processi che
hanno detta funzione nel loro codice.

\index{paginazione|(}

La corrispondenza fra le pagine della memoria virtuale di un processo e quelle
della memoria fisica della macchina viene gestita in maniera trasparente dal
kernel.\footnote{in genere con l'ausilio dell'hardware di gestione della
  memoria (la \textit{Memory Management Unit} del processore), con i kernel
  della serie 2.6 è comunque diventato possibile utilizzare Linux anche su
  architetture che non dispongono di una MMU.}  Poiché in genere la memoria
fisica è solo una piccola frazione della memoria virtuale, è necessario un
meccanismo che permetta di trasferire le pagine che servono dal supporto su
cui si trovano in memoria, eliminando quelle che non servono.  Questo
meccanismo è detto \textsl{paginazione} (o \textit{paging}), ed è uno dei
compiti principali del kernel.

\itindbeg{page~fault} 

Quando un processo cerca di accedere ad una pagina che non è nella memoria
reale, avviene quello che viene chiamato un \textit{page fault}; la gestione
della memoria genera un'interruzione e passa il controllo al kernel il quale
sospende il processo e si incarica di mettere in RAM la pagina richiesta,
effettuando tutte le operazioni necessarie per reperire lo spazio necessario,
per poi restituire il controllo al processo.

Dal punto di vista di un processo questo meccanismo è completamente
trasparente, e tutto avviene come se tutte le pagine fossero sempre
disponibili in memoria.  L'unica differenza avvertibile è quella dei tempi di
esecuzione, che passano dai pochi nanosecondi necessari per l'accesso in RAM
se la pagina è direttamente disponibile, a tempi estremamente più lunghi,
dovuti all'intervento del kernel, qualora sia necessario reperire pagine
riposte nella \textit{swap}.

\itindend{page~fault} 

Normalmente questo è il prezzo da pagare per avere un \textit{multitasking}
reale, ed in genere il sistema è molto efficiente in questo lavoro; quando
però ci siano esigenze specifiche di prestazioni è possibile usare delle
funzioni che permettono di bloccare il meccanismo della paginazione e
mantenere fisse delle pagine in memoria (vedi sez.~\ref{sec:proc_mem_lock}).

\index{paginazione|)}
\index{memoria~virtuale|)}


\subsection{La struttura della memoria di un processo}
\label{sec:proc_mem_layout}

\itindbeg{segment~violation}

Benché lo spazio di indirizzi virtuali copra un intervallo molto ampio, solo
una parte di essi è effettivamente allocato ed utilizzabile dal processo; il
tentativo di accedere ad un indirizzo non allocato è un tipico errore che si
commette quando si è manipolato male un puntatore e genera quella che viene
chiamata una \textit{segment violation}. Se si tenta cioè di leggere o
scrivere con un indirizzo per il quale non esiste un'associazione nella
memoria virtuale, il kernel risponde al relativo \textit{page fault} mandando
un segnale \signal{SIGSEGV} al processo, che normalmente ne causa la
terminazione immediata.

\itindend{segment~violation}

È pertanto importante capire come viene strutturata la memoria virtuale di un
processo. Essa viene divisa in \textsl{segmenti}, cioè un insieme contiguo di
indirizzi virtuali ai quali il processo può accedere.  Solitamente un
programma C viene suddiviso nei seguenti segmenti:
\index{segmento!testo|(}
\index{segmento!dati|(}
\itindbeg{heap} 
\itindbeg{stack}
\begin{enumerate}
\item Il \textsl{segmento di testo} o \textit{text segment}.  Contiene il
  codice del programma, delle funzioni di librerie da esso utilizzate, e le
  costanti.  Normalmente viene condiviso fra tutti i processi che eseguono lo
  stesso programma e nel caso delle librerie anche da processi che eseguono
  altri programmi.

  Quando l'architettura hardware lo supporta viene marcato in sola lettura per
  evitare sovrascritture accidentali (o maliziose) che ne modifichino le
  istruzioni.  Viene allocato da \func{execve} all'avvio del programma e resta
  invariato per tutto il tempo dell'esecuzione.
\index{variabili!globali|(}
\index{variabili!statiche|(}
\item Il \textsl{segmento dei dati} o \textit{data segment}. Contiene tutti i
  dati del programma, come le \textsl{variabili globali}, cioè quelle definite
  al di fuori di tutte le funzioni che compongono il programma, e le
  \textsl{variabili statiche}, cioè quelle dichiarate con l'attributo
  \direct{static},\footnote{la direttiva \direct{static} indica al compilatore
    C che una variabile così dichiarata all'interno di una funzione deve
    essere mantenuta staticamente in memoria (nel segmento dati appunto);
    questo significa che la variabile verrà inizializzata una sola volta alla
    prima invocazione della funzione e che il suo valore sarà mantenuto fra
    diverse esecuzioni della funzione stessa, la differenza con una variabile
    globale è che essa può essere vista solo all'interno della funzione in cui
    è dichiarata.} e la memoria allocata dinamicamente. Di norma è diviso in
  tre parti:
  \begin{itemize}
  \item Il segmento dei dati inizializzati, che contiene le variabili il cui
    valore è stato assegnato esplicitamente. Ad esempio se si definisce:
    \includecodesnip{listati/pi.c}
    questo valore sarà immagazzinato in questo segmento. La memoria di questo
    segmento viene preallocata all'avvio del programma e inizializzata ai valori
    specificati.
  \item Il segmento dei dati non inizializzati, che contiene le variabili il
    cui valore non è stato assegnato esplicitamente. Ad esempio se si
    definisce:
    \includecodesnip{listati/vect.c}
    questo vettore sarà immagazzinato in questo segmento. Anch'esso viene
    allocato all'avvio, e tutte le variabili vengono inizializzate a zero (ed
    i puntatori a \val{NULL}).\footnote{si ricordi che questo vale solo per le
      variabili che vanno nel segmento dati, e non è affatto vero in
      generale.}  Storicamente questa seconda parte del segmento dati viene
    chiamata \itindex{Block~Started~by~Symbol~(BSS)} BSS (da \textit{Block
      Started by Symbol}). La sua dimensione è fissa.
    \index{variabili!globali|)} \index{variabili!statiche|)}
  \item Lo \textit{heap}, detto anche \textit{free store}. Tecnicamente lo si
    può considerare l'estensione del segmento dei dati non inizializzati, a
    cui di solito è posto giusto di seguito. Questo è il segmento che viene
    utilizzato per l'allocazione dinamica della memoria.  Lo \textit{heap} può
    essere ridimensionato allargandolo e restringendolo per allocare e
    disallocare la memoria dinamica con le apposite funzioni (vedi
    sez.~\ref{sec:proc_mem_alloc}), ma il suo limite inferiore, quello
    adiacente al segmento dei dati non inizializzati, ha una posizione fissa.
  \end{itemize}
\item Il segmento di \textit{stack}, che contiene quello che viene chiamato lo
  ``\textit{stack}'' del programma.  Tutte le volte che si effettua una
  chiamata ad una funzione è qui che viene salvato l'indirizzo di ritorno e le
  informazioni dello stato del chiamante (come il contenuto di alcuni registri
  della CPU), poi la funzione chiamata alloca qui lo spazio per le sue
  variabili locali. Tutti questi dati vengono \textit{impilati} (da questo
  viene il nome \textit{stack}) in sequenza uno sull'altro; in questo modo le
  funzioni possono essere chiamate ricorsivamente. Al ritorno della funzione
  lo spazio è automaticamente rilasciato e ``\textsl{ripulito}''.\footnote{il
    compilatore si incarica di generare automaticamente il codice necessario,
    seguendo quella che viene chiamata una \textit{calling convention}; quella
    standard usata con il C ed il C++ è detta \textit{cdecl} e prevede che gli
    argomenti siano caricati nello \textit{stack} dal chiamante da destra a
    sinistra, e che sia il chiamante stesso ad eseguire la ripulitura dello
    \textit{stack} al ritorno della funzione, se ne possono però utilizzare di
    alternative (ad esempio nel Pascal gli argomenti sono inseriti da sinistra
    a destra ed è compito del chiamato ripulire lo \textit{stack}), in genere
    non ci si deve preoccupare di questo fintanto che non si mescolano
    funzioni scritte con linguaggi diversi.}

  La dimensione di questo segmento aumenta seguendo la crescita dello
  \textit{stack} del programma, ma non viene ridotta quando quest'ultimo si
  restringe.
\end{enumerate}

\begin{figure}[htb]
  \centering
  \includegraphics[height=10cm]{img/memory_layout}
  % \begin{tikzpicture}
  % \draw (0,0) rectangle (4,1);
  % \draw (2,0.5) node {\textit{text}};
  % \draw (0,1) rectangle (4,2.5);
  % \draw (2,1.75) node {dati inizializzati};
  % \draw (0,2.5) rectangle (4,5);
  % \draw (2,3.75) node {dati non inizializzati};
  % \draw (0,5) rectangle (4,9);
  % \draw[dashed] (0,6) -- (4,6);
  % \draw[dashed] (0,8) -- (4,8);
  % \draw (2,5.5) node {\textit{heap}};
  % \draw (2,8.5) node {\textit{stack}};
  % \draw [->] (2,6) -- (2,6.5);
  % \draw [->] (2,8) -- (2,7.5);
  % \draw (0,9) rectangle (4,10);
  % \draw (2,9.5) node {\textit{environment}};
  % \draw (4,0) node [anchor=west] {\texttt{0x08000000}};
  % \draw (4,5) node [anchor=west] {\texttt{0x08xxxxxx}};
  % \draw (4,9) node [anchor=west] {\texttt{0xC0000000}};
  % \end{tikzpicture} 
  \caption{Disposizione tipica dei segmenti di memoria di un processo.}
  \label{fig:proc_mem_layout}
\end{figure}

Una disposizione tipica dei vari segmenti (testo, dati inizializzati e non
inizializzati, \textit{heap}, \textit{stack}, ecc.) è riportata in
fig.~\ref{fig:proc_mem_layout}. Si noti come in figura sia indicata una
ulteriore regione, marcata \textit{environment}, che è quella che contiene i
dati relativi alle variabili di ambiente passate al programma al suo avvio
(torneremo su questo argomento in sez.~\ref{sec:proc_environ}).

Usando il comando \cmd{size} su un programma se ne può stampare le dimensioni
dei segmenti di testo e di dati (solo però per i dati inizializzati ed il BSS,
dato che lo \textit{heap} ha una dimensione dinamica). Si tenga presente
comunque che il BSS, contrariamente al segmento dei dati inizializzati, non è
mai salvato sul file che contiene l'eseguibile, dato che viene sempre
inizializzato a zero al caricamento del programma.

\index{segmento!testo|)}
\index{segmento!dati|)}
\itindend{heap} 
\itindend{stack}


\subsection{Allocazione della memoria per i programmi C}
\label{sec:proc_mem_alloc}

Il C supporta direttamente, come linguaggio di programmazione, soltanto due
modalità di allocazione della memoria: l'\textsl{allocazione statica} e
l'\textsl{allocazione automatica}.

L'\textsl{allocazione statica} è quella con cui sono memorizzate le variabili
globali e le variabili statiche, cioè le variabili il cui valore deve essere
mantenuto per tutta la durata del programma. Come accennato queste variabili
vengono allocate nel segmento dei dati all'avvio del programma come parte
delle operazioni svolte da \func{exec}, e lo spazio da loro occupato non viene
liberato fino alla sua conclusione.

\index{variabili!automatiche|(}

L'\textsl{allocazione automatica} è quella che avviene per gli argomenti di
una funzione e per le sue variabili locali, quelle che vengono definite
all'interno della funzione che esistono solo per la durata della sua
esecuzione e che per questo vengono anche dette \textsl{variabili
  automatiche}.  Lo spazio per queste variabili viene allocato nello
\textit{stack} quando viene eseguita la funzione e liberato quando si esce
dalla medesima.

\index{variabili!automatiche|)}

Esiste però un terzo tipo di allocazione, l'\textsl{allocazione dinamica}
della memoria, che non è prevista direttamente all'interno del linguaggio C,
ma che è necessaria quando il quantitativo di memoria che serve è
determinabile solo durante il corso dell'esecuzione del programma. Il C non
consente di usare variabili allocate dinamicamente, non è possibile cioè
definire in fase di programmazione una variabile le cui dimensioni possano
essere modificate durante l'esecuzione del programma. Per questo la libreria
standard del C fornisce una opportuna serie di funzioni per eseguire
l'allocazione dinamica di memoria, che come accennato avviene nello
\textit{heap}.

Le variabili il cui contenuto è allocato in questo modo non potranno essere
usate direttamente come le altre (quelle nello \textit{stack}), ma l'accesso
sarà possibile solo in maniera indiretta, attraverso i puntatori alla memoria
loro riservata che si sono ottenuti dalle funzioni di allocazione.

Le funzioni previste dallo standard ANSI C per la gestione della memoria sono
quattro: \func{malloc}, \func{calloc}, \func{realloc} e \func{free}. Le prime
due, \funcd{malloc} e \funcd{calloc} allocano nuovo spazio di memoria; i
rispettivi prototipi sono:

\begin{funcproto}{ 
\fhead{stdlib.h} 
\fdecl{void *calloc(size\_t nmemb, size\_t size)}
\fdesc{Alloca un'area di memoria inizializzata a 0.}  
\fdecl{void *malloc(size\_t size)}
\fdesc{Alloca un'area di memoria non inizializzata.}  
}
{Entrambe le funzioni restituiscono il puntatore alla zona di memoria allocata
in caso di successo e \val{NULL} in caso di fallimento, nel qual caso
  \var{errno} assumerà il valore \errcode{ENOMEM}.}
\end{funcproto}

In genere si usano \func{malloc} e \func{calloc} per allocare dinamicamente
un'area di memoria.\footnote{queste funzioni presentano un comportamento
  diverso fra la \acr{glibc} e la \acr{uClib} quando il valore di \param{size}
  è nullo.  Nel primo caso viene comunque restituito un puntatore valido,
  anche se non è chiaro a cosa esso possa fare riferimento, nel secondo caso
  viene restituito \val{NULL}. Il comportamento è analogo con
  \code{realloc(NULL, 0)}.}  Dato che i puntatori ritornati sono di tipo
generico non è necessario effettuare un cast per assegnarli a puntatori al
tipo di variabile per la quale si effettua l'allocazione, inoltre le funzioni
garantiscono che i puntatori siano allineati correttamente per tutti i tipi di
dati; ad esempio sulle macchine a 32 bit in genere sono allineati a multipli
di 4 byte e sulle macchine a 64 bit a multipli di 8 byte.

Nel caso di \func{calloc} l'area di memoria viene allocata nello \textit{heap}
come un vettore di \param{nmemb} membri di \param{size} byte di dimensione, e
preventivamente inizializzata a zero, nel caso di \func{malloc} invece vengono
semplicemente allocati \param{size} byte e l'area di memoria non viene
inizializzata.

Una volta che non sia più necessaria la memoria allocata dinamicamente deve
essere esplicitamente rilasciata usando la funzione \funcd{free},\footnote{le
  glibc provvedono anche una funzione \funcm{cfree} definita per compatibilità
  con SunOS, che è deprecata.} il suo prototipo è:

\begin{funcproto}{ 
\fhead{stdlib.h} 
\fdecl{void free(void *ptr)}
\fdesc{Disalloca un'area di memoria precedentemente allocata.}  
}
{La funzione non ritorna nulla e non riporta errori.}
\end{funcproto}

Questa funzione vuole come argomento \var{ptr} il puntatore restituito da una
precedente chiamata ad una qualunque delle funzioni di allocazione che non sia
già stato liberato da un'altra chiamata a \func{free}. Se il valore di
\param{ptr} è \val{NULL} la funzione non fa niente, mentre se l'area di
memoria era già stata liberata da una precedente chiamata il comportamento
della funzione è dichiarato indefinito, ma in genere comporta la corruzione
dei dati di gestione dell'allocazione, che può dar luogo a problemi gravi, ad
esempio un \textit{segmentation fault} in una successiva chiamata di una di
queste funzioni.

\itindbeg{double~free}

Dato che questo errore, chiamato in gergo \textit{double free}, è abbastanza
frequente, specie quando si manipolano vettori di puntatori, e dato che le
conseguenze possono essere pesanti ed inaspettate, si suggerisce come
soluzione precauzionale di assegnare sempre a \val{NULL} ogni puntatore su cui
sia stata eseguita \func{free} immediatamente dopo l'esecuzione della
funzione. In questo modo, dato che con un puntatore nullo \func{free} non
esegue nessuna operazione, si evitano i problemi del \textit{double free}.

\itindend{double~free}

Infine la funzione \funcd{realloc} consente di modificare, in genere di
aumentare, la dimensione di un'area di memoria precedentemente allocata; il
suo prototipo è:

\begin{funcproto}{ 
\fhead{stdlib.h} 
\fdecl{void *realloc(void *ptr, size\_t size)}
\fdesc{Cambia la dimensione di un'area di memoria precedentemente allocata.}
}  {La funzione ritorna il puntatore alla zona di memoria allocata in caso
  di successo e \val{NULL} per un errore, nel qual caso \var{errno}
  assumerà il valore \errcode{ENOMEM}.}
\end{funcproto}

La funzione vuole come primo argomento il puntatore restituito da una
precedente chiamata a \func{malloc} o \func{calloc} e come secondo argomento
la nuova dimensione (in byte) che si intende ottenere. Se si passa
per \param{ptr} il valore \val{NULL} allora la funzione si comporta come
\func{malloc}.\footnote{questo è vero per Linux e l'implementazione secondo lo
  standard ANSI C, ma non è vero per alcune vecchie implementazioni, inoltre
  alcune versioni delle librerie del C consentivano di usare \func{realloc}
  anche per un puntatore liberato con \func{free} purché non ci fossero state
  nel frattempo altre chiamate a funzioni di allocazione, questa funzionalità
  è totalmente deprecata e non è consentita sotto Linux.}

La funzione si usa ad esempio quando si deve far crescere la dimensione di un
vettore. In questo caso se è disponibile dello spazio adiacente al precedente
la funzione lo utilizza, altrimenti rialloca altrove un blocco della
dimensione voluta, copiandoci automaticamente il contenuto; lo spazio aggiunto
non viene inizializzato. Se la funzione fallisce l'area di memoria originale
non viene assolutamente toccata.

Si deve sempre avere ben presente il fatto che il blocco di memoria restituito
da \func{realloc} può non essere un'estensione di quello che gli si è passato
in ingresso; per questo si dovrà \emph{sempre} eseguire la riassegnazione di
\param{ptr} al valore di ritorno della funzione, e reinizializzare o provvedere
ad un adeguato aggiornamento di tutti gli altri puntatori all'interno del
blocco di dati ridimensionato.

La \acr{glibc} ha un'implementazione delle funzioni di allocazione che è
controllabile dall'utente attraverso alcune variabili di ambiente (vedi
sez.~\ref{sec:proc_environ}), in particolare diventa possibile tracciare
questo tipo di errori usando la variabile di ambiente \envvar{MALLOC\_CHECK\_}
che quando viene definita mette in uso una versione meno efficiente delle
funzioni suddette, che però è più tollerante nei confronti di piccoli errori
come quello dei \textit{double free} o i \textit{buffer overrun} di un
byte.\footnote{uno degli errori più comuni, causato ad esempio dalla scrittura
  di una stringa di dimensione pari a quella del buffer, in cui ci si
  dimentica dello zero di terminazione finale.}  In particolare:
\begin{itemize*}
\item se la variabile è posta a $0$ gli errori vengono ignorati;
\item se la variabile è posta a $1$ viene stampato un avviso sullo
  \textit{standard error} (vedi sez.~\ref{sec:file_fd});
\item se la variabile è posta a $2$ viene chiamata la funzione \func{abort}
  (vedi sez.~\ref{sec:sig_alarm_abort}), che in genere causa l'immediata
  terminazione del programma;
\item se la variabile è posta a $3$ viene stampato l'avviso e chiamata
  \func{abort}. 
\end{itemize*}

\itindbeg{memory~leak}

L'errore di programmazione più comune e più difficile da risolvere che si
incontra con le funzioni di allocazione è quando non viene opportunamente
liberata la memoria non più utilizzata, quello che in inglese viene chiamato
\textit{memory leak}, cioè una \textsl{perdita di memoria}.

Un caso tipico che illustra il problema è quello in cui in una propria
funzione si alloca della memoria per uso locale senza liberarla prima di
uscire. La memoria resta così allocata fino alla terminazione del processo.
Chiamate ripetute alla stessa funzione continueranno ad effettuare altre
allocazioni, che si accumuleranno causando a lungo andare un esaurimento della
memoria disponibile e la probabile impossibilità di proseguire l'esecuzione
del programma.

Il problema è che l'esaurimento della memoria può avvenire in qualunque
momento, in corrispondenza ad una qualunque chiamata di \func{malloc} che può
essere in una sezione del codice che non ha alcuna relazione con la funzione
che contiene l'errore. Per questo motivo è sempre molto difficile trovare un
\textit{memory leak}.  In C e C++ il problema è particolarmente sentito. In
C++, per mezzo della programmazione ad oggetti, il problema dei \textit{memory
  leak} si può notevolmente ridimensionare attraverso l'uso accurato di
appositi oggetti come gli \textit{smartpointers}.  Questo però in genere va a
scapito delle prestazioni dell'applicazione in esecuzione.

% TODO decidere cosa fare di questo che segue In altri linguaggi come il java
% e recentemente il C\# il problema non si pone nemmeno perché la gestione
% della memoria viene fatta totalmente in maniera automatica, ovvero il
% programmatore non deve minimamente preoccuparsi di liberare la memoria
% allocata precedentemente quando non serve più, poiché l'infrastruttura del
% linguaggio gestisce automaticamente la cosiddetta
% \itindex{garbage~collection} \textit{garbage collection}. In tal caso,
% attraverso meccanismi simili a quelli del \textit{reference counting},
% quando una zona di memoria precedentemente allocata non è più riferita da
% nessuna parte del codice in esecuzione, può essere deallocata
% automaticamente in qualunque momento dall'infrastruttura.

% Anche questo va a scapito delle prestazioni dell'applicazione in esecuzione
% (inoltre le applicazioni sviluppate con tali linguaggi di solito non sono
% eseguibili compilati, come avviene invece per il C ed il C++, ed è necessaria
% la presenza di una infrastruttura per la loro interpretazione e pertanto hanno
% di per sé delle prestazioni più scadenti rispetto alle stesse applicazioni
% compilate direttamente).  Questo comporta però il problema della non
% predicibilità del momento in cui viene deallocata la memoria precedentemente
% allocata da un oggetto.

Per limitare l'impatto di questi problemi, e semplificare la ricerca di
eventuali errori, l'implementazione delle funzioni di allocazione nella
\acr{glibc} mette a disposizione una serie di funzionalità che permettono di
tracciare le allocazioni e le disallocazioni, e definisce anche una serie di
possibili \textit{hook} (\textsl{ganci}) che permettono di sostituire alle
funzioni di libreria una propria versione (che può essere più o meno
specializzata per il debugging). Esistono varie librerie che forniscono dei
sostituti opportuni delle funzioni di allocazione in grado, senza neanche
ricompilare il programma,\footnote{esempi sono \textit{Dmalloc}
  \url{http://dmalloc.com/} di Gray Watson ed \textit{Electric Fence} di Bruce
  Perens.} di eseguire diagnostiche anche molto complesse riguardo
l'allocazione della memoria. Vedremo alcune delle funzionalità di ausilio
presenti nella \acr{glibc} in sez.~\ref{sec:proc_memory_adv_management}.

\itindend{memory~leak}

Una possibile alternativa all'uso di \func{malloc}, per evitare di soffrire
dei problemi di \textit{memory leak} descritti in precedenza, è di allocare la
memoria nel segmento di \textit{stack} della funzione corrente invece che
nello \textit{heap}. Per farlo si può usare la funzione \funcd{alloca}, la cui
sintassi è identica a quella di \func{malloc}; il suo prototipo è:

\begin{funcproto}{ 
\fhead{stdlib.h} 
\fdecl{void *alloca(size\_t size)}
\fdesc{Alloca un'area di memoria nello \textit{stack}.} 
}
{La funzione ritorna il puntatore alla zona di memoria allocata, in caso
  di errore il comportamento è indefinito.}
\end{funcproto}

La funzione alloca la quantità di memoria (non inizializzata) richiesta
dall'argomento \param{size} nel segmento di \textit{stack} della funzione
chiamante. Con questa funzione non è più necessario liberare la memoria
allocata, e quindi non esiste un analogo della \func{free}, in quanto essa
viene rilasciata automaticamente al ritorno della funzione.

Come è evidente questa funzione ha alcuni vantaggi interessanti, anzitutto
permette di evitare alla radice i problemi di \textit{memory leak}, dato che
non serve più la deallocazione esplicita; inoltre la deallocazione automatica
funziona anche quando si usa \func{longjmp} per uscire da una subroutine con
un salto non locale da una funzione (vedi sez.~\ref{sec:proc_longjmp}).  Un
altro vantaggio è che in Linux la funzione è molto più veloce di \func{malloc}
e non viene sprecato spazio, infatti non è necessario gestire un pool di
memoria da riservare e si evitano così anche i problemi di frammentazione di
quest'ultimo, che comportano inefficienze sia nell'allocazione della memoria
che nell'esecuzione dell'allocazione.

Gli svantaggi sono che questa funzione non è disponibile su tutti gli Unix, e
non è inserita né nello standard POSIX né in SUSv3 (ma è presente in BSD), il
suo utilizzo quindi limita la portabilità dei programmi. Inoltre la funzione
non può essere usata nella lista degli argomenti di una funzione, perché lo
spazio verrebbe allocato nel mezzo degli stessi. Inoltre non è chiaramente
possibile usare \func{alloca} per allocare memoria che deve poi essere usata
anche al di fuori della funzione in cui essa viene chiamata, dato che
all'uscita dalla funzione lo spazio allocato diventerebbe libero, e potrebbe
essere sovrascritto all'invocazione di nuove funzioni.  Questo è lo stesso
problema che si può avere con le variabili automatiche, su cui torneremo in
sez.~\ref{sec:proc_var_passing}.

Infine non esiste un modo di sapere se l'allocazione ha avuto successo, la
funzione infatti viene realizzata inserendo del codice \textit{inline} nel
programma\footnote{questo comporta anche il fatto che non è possibile
  sostituirla con una propria versione o modificarne il comportamento
  collegando il proprio programma con un'altra libreria.} che si limita a
modificare il puntatore nello \textit{stack} e non c'è modo di sapere se se ne
sono superate le dimensioni, per cui in caso di fallimento nell'allocazione il
comportamento del programma può risultare indefinito, dando luogo ad una
\textit{segment violation} la prima volta che si cerchi di accedere alla
memoria non effettivamente disponibile.

\index{segmento!dati|(}
\itindbeg{heap} 

Le due funzioni seguenti vengono utilizzate soltanto quando è necessario
effettuare direttamente la gestione della memoria associata allo spazio dati
di un processo,\footnote{le due funzioni sono state definite con BSD 4.3, sono
  marcate obsolete in SUSv2 e non fanno parte delle librerie standard del C e
  mentre sono state esplicitamente rimosse dallo standard POSIX.1-2001.} per
poterle utilizzare è necessario definire una della macro di funzionalità (vedi
sez.~\ref{sec:intro_gcc_glibc_std}) fra \macro{\_BSD\_SOURCE},
\macro{\_SVID\_SOURCE} e \macro{\_XOPEN\_SOURCE} (ad un valore maggiore o
uguale di 500). La prima funzione è \funcd{brk}, ed il suo prototipo è:

\begin{funcproto}{ 
\fhead{unistd.h} 
\fdecl{int brk(void *addr)}
\fdesc{Sposta la fine del segmento dati del processo.} 
}
{La funzione ritorna $0$ in caso di successo e $-1$ per un errore,
  nel qual caso \var{errno} assumerà il valore \errcode{ENOMEM}.}
\end{funcproto}

La funzione è un'interfaccia all'omonima \textit{system call} ed imposta
l'indirizzo finale del segmento dati di un processo (più precisamente dello
\textit{heap}) all'indirizzo specificato da \param{addr}. Quest'ultimo deve
essere un valore ragionevole e la dimensione totale non deve comunque eccedere
un eventuale limite (vedi sez.~\ref{sec:sys_resource_limit}) sulle dimensioni
massime del segmento dati del processo.

Il valore di ritorno della funzione fa riferimento alla versione fornita dalla
\acr{glibc}, in realtà in Linux la \textit{system call} corrispondente
restituisce come valore di ritorno il nuovo valore della fine del segmento
dati in caso di successo e quello corrente in caso di fallimento, è la
funzione di interfaccia usata dalla \acr{glibc} che fornisce i valori di
ritorno appena descritti; se si usano librerie diverse questo potrebbe non
accadere.

Una seconda funzione per la manipolazione diretta delle dimensioni del
segmento dati\footnote{in questo caso si tratta soltanto di una funzione di
  libreria, anche se basata sulla stessa \textit{system call}.} è
\funcd{sbrk}, ed il suo prototipo è:

\begin{funcproto}{ 
\fhead{unistd.h} 
\fdecl{void *sbrk(intptr\_t increment)}
\fdesc{Incrementa la dimensione del segmento dati del processo.} 
}
{La funzione ritorna il puntatore all'inizio della nuova zona di memoria
  allocata in caso di successo e \val{NULL} per un errore, nel qual
  caso \var{errno} assumerà il valore \errcode{ENOMEM}.}
\end{funcproto}

La funzione incrementa la dimensione dello \textit{heap} di un programma del
valore indicato dall'argomento \param{increment}, restituendo il nuovo
indirizzo finale dello stesso.  L'argomento è definito come di tipo
\typed{intptr\_t}, ma a seconda della versione delle librerie e del sistema
può essere indicato con una serie di tipi equivalenti come \type{ptrdiff\_t},
\type{ssize\_t}, \ctyp{int}. Se invocata con un valore nullo la funzione
permette di ottenere l'attuale posizione della fine del segmento dati.

Queste due funzioni sono state deliberatamente escluse dallo standard POSIX.1
dato che per i normali programmi è sempre opportuno usare le funzioni di
allocazione standard descritte in precedenza, a meno di non voler realizzare
per proprio conto un diverso meccanismo di gestione della memoria del segmento
dati.
\itindend{heap} 
\index{segmento!dati|)}


\subsection{Il controllo della memoria virtuale}  
\label{sec:proc_mem_lock}

\index{memoria~virtuale|(}

Come spiegato in sez.~\ref{sec:proc_mem_gen} il kernel gestisce la memoria
virtuale in maniera trasparente ai processi, decidendo quando rimuovere pagine
dalla memoria per metterle nell'area di \textit{swap}, sulla base
dell'utilizzo corrente da parte dei vari processi.

Nell'uso comune un processo non deve preoccuparsi di tutto ciò, in quanto il
meccanismo della paginazione riporta in RAM, ed in maniera trasparente, tutte
le pagine che gli occorrono; esistono però esigenze particolari in cui non si
vuole che questo meccanismo si attivi. In generale i motivi per cui si possono
avere di queste necessità sono due:
\begin{itemize*}
\item \textsl{La velocità}. Il processo della paginazione è trasparente solo
  se il programma in esecuzione non è sensibile al tempo che occorre a
  riportare la pagina in memoria; per questo motivo processi critici che hanno
  esigenze di tempo reale o tolleranze critiche nelle risposte (ad esempio
  processi che trattano campionamenti sonori) possono non essere in grado di
  sopportare le variazioni della velocità di accesso dovuta alla paginazione.
  
  In certi casi poi un programmatore può conoscere meglio dell'algoritmo di
  allocazione delle pagine le esigenze specifiche del suo programma e decidere
  quali pagine di memoria è opportuno che restino in memoria per un aumento
  delle prestazioni. In genere queste sono esigenze particolari e richiedono
  anche un aumento delle priorità in esecuzione del processo (vedi
  sez.~\ref{sec:proc_real_time}).
  
\item \textsl{La sicurezza}. Se si hanno password o chiavi segrete in chiaro
  in memoria queste possono essere portate su disco dal meccanismo della
  paginazione. Questo rende più lungo il periodo di tempo in cui detti segreti
  sono presenti in chiaro e più complessa la loro cancellazione: un processo
  infatti può cancellare la memoria su cui scrive le sue variabili, ma non può
  toccare lo spazio disco su cui una pagina di memoria può essere stata
  salvata. Per questo motivo di solito i programmi di crittografia richiedono
  il blocco di alcune pagine di memoria.
\end{itemize*}

Per ottenere informazioni sulle modalità in cui un programma sta usando la
memoria virtuale è disponibile una apposita funzione di sistema,
\funcd{mincore}, che però non è standardizzata da POSIX e pertanto non è
disponibile su tutte le versioni di kernel unix-like;\footnote{nel caso di
  Linux devono essere comunque definite le macro \macro{\_BSD\_SOURCE} e
  \macro{\_SVID\_SOURCE} o \macro{\_DEFAULT\_SOURCE}.}  il suo prototipo è:

\begin{funcproto}{
\fhead{unistd.h}
\fhead{sys/mman.h}
\fdecl{int mincore(void *addr, size\_t length, unsigned char *vec)}
\fdesc{Ritorna lo stato delle pagine di memoria occupate da un processo.}
}
{La funzione ritorna $0$ in caso di successo e $-1$ per un errore, nel qual
caso \var{errno} assumerà uno dei valori:
\begin{errlist}
   \item[\errcode{EAGAIN}] il kernel è temporaneamente non in grado di fornire
     una risposta.
   \item[\errcode{EFAULT}] \param{vec} punta ad un indirizzo non valido.
   \item[\errcode{EINVAL}] \param{addr} non è un multiplo delle dimensioni di
     una pagina.
   \item[\errcode{ENOMEM}] o \param{addr}$+$\param{length} eccede la dimensione
     della memoria usata dal processo o l'intervallo di indirizzi specificato
     non è mappato.
\end{errlist}}
\end{funcproto}

La funzione permette di ottenere le informazioni sullo stato della mappatura
della memoria per il processo chiamante, specificando l'intervallo da
esaminare con l'indirizzo iniziale, indicato con l'argomento \param{addr}, e
la lunghezza, indicata con l'argomento \param{length}. L'indirizzo iniziale
deve essere un multiplo delle dimensioni di una pagina, mentre la lunghezza
può essere qualunque, fintanto che si resta nello spazio di indirizzi del
processo,\footnote{in caso contrario si avrà un errore di \errcode{ENOMEM};
  fino al kernel 2.6.11 in questo caso veniva invece restituito
  \errcode{EINVAL}, in considerazione che il caso più comune in cui si
  verifica questo errore è quando si usa per sbaglio un valore negativo
  di \param{length}, che nel caso verrebbe interpretato come un intero
  positivo di grandi dimensioni.}  ma il risultato verrà comunque fornito per
l'intervallo compreso fino al multiplo successivo.

I risultati della funzione vengono forniti nel vettore puntato da \param{vec},
che deve essere allocato preventivamente e deve essere di dimensione
sufficiente a contenere tanti byte quante sono le pagine contenute
nell'intervallo di indirizzi specificato, la dimensione cioè deve essere
almeno pari a \code{(length+PAGE\_SIZE-1)/PAGE\_SIZE}.  Al ritorno della
funzione il bit meno significativo di ciascun byte del vettore sarà acceso se
la pagina di memoria corrispondente è al momento residente in memoria, o
cancellato altrimenti. Il comportamento sugli altri bit è indefinito, essendo
questi al momento riservati per usi futuri. Per questo motivo in genere è
comunque opportuno inizializzare a zero il contenuto del vettore, così che le
pagine attualmente residenti in memoria saranno indicate da un valore non
nullo del byte corrispondente.

Dato che lo stato della memoria di un processo può cambiare continuamente, il
risultato di \func{mincore} è assolutamente provvisorio e lo stato delle
pagine potrebbe essere già cambiato al ritorno stesso della funzione, a meno
che, come vedremo ora, non si sia attivato il meccanismo che forza il
mantenimento di una pagina sulla memoria.  

\itindbeg{memory~locking}

Il meccanismo che previene la paginazione di parte della memoria virtuale di
un processo è chiamato \textit{memory locking} (o \textsl{blocco della
  memoria}). Il blocco è sempre associato alle pagine della memoria virtuale
del processo, e non al segmento reale di RAM su cui essa viene mantenuta.  La
regola è che se un segmento di RAM fa da supporto ad almeno una pagina
bloccata allora esso viene escluso dal meccanismo della paginazione. I blocchi
non si accumulano, se si blocca due volte la stessa pagina non è necessario
sbloccarla due volte, una pagina o è bloccata oppure no.

Il \textit{memory lock} persiste fintanto che il processo che detiene la
memoria bloccata non la sblocca. Chiaramente la terminazione del processo
comporta anche la fine dell'uso della sua memoria virtuale, e quindi anche di
tutti i suoi \textit{memory lock}.  Inoltre i \textit{memory lock} non sono
ereditati dai processi figli, ma siccome Linux usa il \textit{copy on write}
(vedi sez.~\ref{sec:proc_fork}) gli indirizzi virtuali del figlio sono
mantenuti sullo stesso segmento di RAM del padre, e quindi fintanto che un
figlio non scrive su un segmento bloccato, può usufruire del \textit{memory
  lock} del padre. Infine i \textit{memory lock} vengono automaticamente
rimossi se si pone in esecuzione un altro programma con \func{exec} (vedi
sez.~\ref{sec:proc_exec}).

Il sistema pone dei limiti all'ammontare di memoria di un processo che può
essere bloccata e al totale di memoria fisica che si può dedicare a questo, lo
standard POSIX.1 richiede che sia definita in \headfile{unistd.h} la macro
\macrod{\_POSIX\_MEMLOCK\_RANGE} per indicare la capacità di eseguire il
\textit{memory locking}.

Siccome la richiesta di un \textit{memory lock} da parte di un processo riduce
la memoria fisica disponibile nel sistema per gli altri processi, questo ha un
evidente impatto su tutti gli altri processi, per cui fino al kernel 2.6.9
solo un processo dotato di privilegi amministrativi (la \textit{capability}
\const{CAP\_IPC\_LOCK}, vedi sez.~\ref{sec:proc_capabilities}) aveva la
capacità di bloccare una pagina di memoria.

A partire dal kernel 2.6.9 anche un processo normale può bloccare la propria
memoria\footnote{la funzionalità è stata introdotta per non essere costretti a
  dare privilegi eccessivi a programmi di crittografia, che necessitano di
  questa funzionalità, ma che devono essere usati da utenti normali.} ma
mentre un processo privilegiato non ha limiti sulla quantità di memoria che
può bloccare, un processo normale è soggetto al limite della risorsa
\const{RLIMIT\_MEMLOCK} (vedi sez.~\ref{sec:sys_resource_limit}). In generale
poi ogni processo può sbloccare le pagine relative alla propria memoria, se
però diversi processi bloccano la stessa pagina questa resterà bloccata
fintanto che ci sarà almeno un processo che la blocca.

Le funzioni di sistema per bloccare e sbloccare la paginazione di singole
sezioni di memoria sono rispettivamente \funcd{mlock} e \funcd{munlock}; i
loro prototipi sono:

\begin{funcproto}{
  \fhead{sys/mman.h} 
  \fdecl{int mlock(const void *addr, size\_t len)}
  \fdesc{Blocca la paginazione su un intervallo di memoria.}

  \fdecl{int munlock(const void *addr, size\_t len)}
  \fdesc{Rimuove il blocco della paginazione su un intervallo di memoria.}
  }
{Entrambe le funzioni ritornano $0$ in caso di successo e $-1$ in caso di
  errore, nel qual caso \var{errno} assumerà uno dei valori:
  \begin{errlist}
  \item[\errcode{EAGAIN}] una parte o tutto l'intervallo richiesto non può
    essere bloccato per una mancanza temporanea di risorse.
  \item[\errcode{EINVAL}] \param{len} non è un valore positivo o la somma con
    \param{addr} causa un overflow.
  \item[\errcode{ENOMEM}] alcuni indirizzi dell’intervallo specificato non
    corrispondono allo spazio di indirizzi del processo o con \func{mlock} si
    è superato il limite di \const{RLIMIT\_MEMLOCK} per un processo non
    privilegiato (solo per kernel a partire dal 2.6.9) o si è superato il
    limite di regioni di memoria con attributi diversi.
  \item[\errcode{EPERM}] il processo non è privilegiato (per kernel precedenti
    il 2.6.9) o si ha un limite nullo per \const{RLIMIT\_MEMLOCK} e
    il processo non è privilegiato (per kernel a partire dal 2.6.9).
  \end{errlist}}
\end{funcproto}

Le due funzioni permettono rispettivamente di bloccare e sbloccare la
paginazione per l'intervallo di memoria iniziante all'indirizzo \param{addr} e
lungo \param{len} byte.  Al ritorno di \func{mlock} tutte le pagine che
contengono una parte dell'intervallo bloccato sono garantite essere in RAM e
vi verranno mantenute per tutta la durata del blocco. Con kernel diversi da
Linux si può ottenere un errore di \errcode{EINVAL} se \param{addr} non è un
multiplo della dimensione delle pagine di memoria, pertanto se si ha a cuore
la portabilità si deve avere cura di allinearne correttamente il valore. Il
blocco viene rimosso chiamando \func{munlock}.

Altre due funzioni di sistema, \funcd{mlockall} e \funcd{munlockall},
consentono di bloccare genericamente la paginazione per l'intero spazio di
indirizzi di un processo.  I prototipi di queste funzioni sono:

\begin{funcproto}{ 
\fhead{sys/mman.h} 
\fdecl{int mlockall(int flags)}
\fdesc{Blocca la paginazione per lo spazio di indirizzi del processo corrente.} 
\fdecl{int munlockall(void)}
\fdesc{Sblocca la paginazione per lo spazio di indirizzi del processo corrente.}
}
{Codici di ritorno ed errori sono gli stessi di \func{mlock} e \func{munlock},
  tranne per \errcode{EINVAL} che viene restituito solo se si è specificato
  con \func{mlockall} un valore sconosciuto per \param{flags}.}
\end{funcproto}

L'argomento \param{flags} di \func{mlockall} permette di controllarne il
comportamento; esso deve essere specificato come maschera binaria dei valori
espressi dalle costanti riportate in tab.~\ref{tab:mlockall_flags}. 

\begin{table}[htb]
  \footnotesize
  \centering
  \begin{tabular}[c]{|l|p{8cm}|}
    \hline
    \textbf{Valore} & \textbf{Significato} \\
    \hline
    \hline
    \constd{MCL\_CURRENT}& blocca tutte le pagine correntemente mappate nello
                           spazio di indirizzi del processo.\\
    \constd{MCL\_FUTURE} & blocca tutte le pagine che verranno mappate nello
                           spazio di indirizzi del processo.\\
    \constd{MCL\_ONFAULT}& esegue il blocco delle pagine selezionate solo
                           quando vengono utilizzate (dal kernel 4.4).\\
   \hline
  \end{tabular}
  \caption{Valori e significato dell'argomento \param{flags} della funzione
    \func{mlockall}.}
  \label{tab:mlockall_flags}
\end{table}

Con \func{mlockall} si possono bloccare tutte le pagine mappate nello spazio
di indirizzi del processo, sia che comprendano il segmento di testo, di dati,
lo \textit{stack}, lo \textit{heap} e pure le funzioni di libreria chiamate, i
file mappati in memoria, i dati del kernel mappati in \textit{user space}, la
memoria condivisa.  L'uso dell'argomento \param{flags} permette di selezionare
con maggior finezza le pagine da bloccare, ad esempio usando
\const{MCL\_FUTURE} ci si può limitare a tutte le pagine allocate a partire
dalla chiamata della funzione, mentre \const{MCL\_CURRENT} blocca tutte quelle
correntemente mappate. L'uso di \func{munlockall} invece sblocca sempre tutte
le pagine di memoria correntemente mappate nello spazio di indirizzi del
programma.

A partire dal kernel 4.4 alla funzione \func{mlockall} è stato aggiunto un
altro flag, \const{MCL\_ONFAULT}, che può essere abbinato a entrambi gli altri
due flag, e consente di modificare il comportamento della funzione per
ottenere migliori prestazioni.

Il problema che si presenta infatti è che eseguire un \textit{memory lock} per
un intervallo ampio di memoria richiede che questa venga comunque allocata in
RAM, con altrettanti \textit{page fault} che ne assicurino la presenza; questo
vale per tutto l'intervallo e può avere un notevole costo in termini di
prestazioni, anche quando poi, nell'esecuzione del programma, venisse usata
solo una piccola parte dello stesso. L'uso di \const{MCL\_ONFAULT} previene il
\textit{page faulting} immediato di tutto l'intervallo, le pagine
dell'intervallo verranno bloccate, ma solo quando un \textit{page fault}
dovuto all'accesso ne richiede l'allocazione effettiva in RAM.

Questo stesso comportamento non è ottenibile con \func{mlock}, che non dispone
di un argomento \param{flag} che consenta di richiederlo, per questo sempre
con il kernel 4.4 è stata aggiunta una ulteriore funzione di sistema,
\funcd{mlock2}, il cui prototipo è:

\begin{funcproto}{
  \fhead{sys/mman.h} 
  \fdecl{int mlock2(const void *addr, size\_t len, int flags)}
  \fdesc{Blocca la paginazione su un intervallo di memoria.}
}
{Le funzione ritornano $0$ in caso di successo e $-1$ in caso di errore, nel
  qual caso \var{errno} assume gli stessi valori di \func{mlock} con
  l'aggiunta id un possibile \errcode{EINVAL} anche se si è indicato un valore
  errato di \param{flags}.}
\end{funcproto}

% NOTA: per mlock2, introdotta con il kernel 4.4 (vedi
% http://lwn.net/Articles/650538/)

Indicando un valore nullo per \param{flags} il comportamento della funzione è
identico a quello di \func{mlock}, l'unico altro valore possibile è
\constd{MLOCK\_ONFAULT} che ha lo stesso effetto sull'allocazione delle pagine
in RAM già descritto per \const{MCL\_ONFAULT}.

Si tenga presente che un processo \textit{real-time} che intende usare il
\textit{memory locking} con \func{mlockall} per prevenire l'avvenire di un
eventuale \textit{page fault} ed il conseguente rallentamento (probabilmente
inaccettabile) dei tempi di esecuzione, deve comunque avere delle accortezze.
In particolare si deve assicurare di aver preventivamente bloccato una
quantità di spazio nello \textit{stack} sufficiente a garantire l'esecuzione
di tutte le funzioni che hanno i requisiti di criticità sui tempi. Infatti,
anche usando \const{MCL\_FUTURE}, in caso di allocazione di una nuova pagina
nello \textit{stack} durante l'esecuzione di una funzione (precedentemente non
usata e quindi non bloccata) si potrebbe avere un \textit{page fault}.

In genere si ovvia a questa problematica chiamando inizialmente una funzione
che definisca una quantità sufficientemente ampia di variabili automatiche
(che si ricordi vengono allocate nello \textit{stack}) e ci scriva, in modo da
esser sicuri che le corrispondenti pagine vengano mappate nello spazio di
indirizzi del processo, per poi bloccarle. La scrittura è necessaria perché il
kernel usa il meccanismo di \textit{copy on write} (vedi
sez.~\ref{sec:proc_fork}) e le pagine potrebbero non essere allocate
immediatamente.

\itindend{memory~locking}
\index{memoria~virtuale|)} 


\subsection{Gestione avanzata dell'allocazione della memoria} 
\label{sec:proc_memory_adv_management}

La trattazione delle funzioni di allocazione di sez.~\ref{sec:proc_mem_alloc}
si è limitata a coprire le esigenze generiche di un programma, in cui non si
hanno dei requisiti specifici e si lascia il controllo delle modalità di
allocazione alle funzioni di libreria.  Tuttavia esistono una serie di casi in
cui può essere necessario avere un controllo più dettagliato delle modalità
con cui la memoria viene allocata; nel qual caso potranno venire in aiuto le
funzioni trattate in questa sezione.

Le prime funzioni che tratteremo sono quelle che consentono di richiedere di
allocare un blocco di memoria ``\textsl{allineato}'' ad un multiplo una certa
dimensione. Questo tipo di esigenza emerge usualmente quando si devono
allocare dei buffer da utilizzare per eseguire dell'I/O diretto su dispositivi
a blocchi. In questo caso infatti il trasferimento di dati viene eseguito per
blocchi di dimensione fissa, ed è richiesto che l'indirizzo di partenza del
buffer sia un multiplo intero di questa dimensione, usualmente 512 byte. In
tal caso l'uso di \func{malloc} non è sufficiente, ed occorre utilizzare una
funzione specifica.

Tradizionalmente per rispondere a questa esigenza sono state create due
funzioni diverse, \funcd{memalign} e \funcd{valloc}, oggi obsolete, cui si
aggiunge \funcd{pvalloc} come estensione GNU, anch'essa obsoleta; i rispettivi
prototipi sono:

\begin{funcproto}{ 
\fhead{malloc.h} 
\fdecl{void *valloc(size\_t size)}
\fdesc{Alloca un blocco di memoria allineato alla dimensione di una pagina di
  memoria.}  
\fdecl{void *memalign(size\_t boundary, size\_t size)}
\fdesc{Alloca un blocco di memoria allineato ad un multiplo
  di \param{boundary}.} 
\fdecl{void *pvalloc(size\_t size)}
\fdesc{Alloca un blocco di memoria allineato alla dimensione di una pagina di
  memoria.}  
}
{Entrambe le funzioni ritornano un puntatore al blocco di memoria allocato in
  caso di successo e \val{NULL} in caso di errore, nel qual caso \var{errno}
  assumerà uno dei valori:
  \begin{errlist}
  \item[\errcode{EINVAL}] \param{boundary} non è una potenza di due.
  \item[\errcode{ENOMEM}] non c'è memoria sufficiente per l'allocazione.
  \end{errlist}}
\end{funcproto}

Le funzioni restituiscono il puntatore al buffer di memoria allocata di
dimensioni pari a \param{size}, che per \func{memalign} sarà un multiplo di
\param{boundary} mentre per \func{valloc} un multiplo della dimensione di una
pagina di memoria; lo stesso vale per \func{pvalloc} che però arrotonda
automaticamente la dimensione dell'allocazione al primo multiplo di una
pagina. Nel caso della versione fornita dalla \acr{glibc} la memoria allocata
con queste funzioni deve essere liberata con \func{free}, cosa che non è detto
accada con altre implementazioni.

Nessuna delle due funzioni ha una chiara standardizzazione e nessuna delle due
compare in POSIX.1, inoltre ci sono indicazioni discordi sui file che ne
contengono la definizione;\footnote{secondo SUSv2 \func{valloc} è definita in
  \headfile{stdlib.h}, mentre sia la \acr{glibc} che le precedenti \acr{libc4}
  e \acr{libc5} la dichiarano in \headfile{malloc.h}, lo stesso vale per
  \func{memalign} che in alcuni sistemi è dichiarata in \headfile{stdlib.h}.}
per questo motivo il loro uso è sconsigliato, essendo state sostituite dalla
nuova \funcd{posix\_memalign}, che è stata standardizzata in POSIX.1d e
disponibile dalla \acr{glibc} 2.1.91; il suo prototipo è:

\begin{funcproto}{ 
\fhead{stdlib.h} 
\fdecl{posix\_memalign(void **memptr, size\_t alignment, size\_t size)}
\fdesc{Alloca un buffer di memoria allineato ad un multiplo
  di \param{alignment}.}   
}
{Entrambe le funzioni ritornano un puntatore al blocco di memoria allocato in
  caso di successo e \val{NULL} in caso di errore, nel qual caso \var{errno}
  assumerà uno dei valori:
  \begin{errlist}
  \item[\errcode{EINVAL}] \param{alignment} non è potenza di due o un multiplo
    di \code{sizeof(void *)}.
  \item[\errcode{ENOMEM}] non c'è memoria sufficiente per l'allocazione.
  \end{errlist}}
\end{funcproto}

La funzione restituisce il puntatore al buffer allocato di dimensioni pari
a \param{size} nella variabile (di tipo \texttt{void *}) posta all'indirizzo
indicato da \param{memptr}. La funzione fallisce nelle stesse condizioni delle
due funzioni precedenti, ma a loro differenza restituisce direttamente come
valore di ritorno il codice di errore.  Come per le precedenti la memoria
allocata con \func{posix\_memalign} deve essere disallocata con \func{free},
che in questo caso però è quanto richiesto dallo standard.

Dalla versione 2.16 della \acr{glibc} è stata aggiunta anche la funzione
\funcd{aligned\_alloc}, prevista dallo standard C11 (e disponibile definendo
\const{\_ISOC11\_SOURCE}), il cui prototipo è:

\begin{funcproto}{ 
\fhead{malloc.h} 
\fdecl{void *aligned\_alloc(size\_t alignment, size\_t size)}
\fdesc{Alloca un blocco di memoria allineato ad un multiplo
  di \param{alignment}.} 
}
{La funzione ha gli stessi valori di ritorno e codici di errore di
  \func{memalign}.}
\end{funcproto}

La funzione è identica a \func{memalign} ma richiede che \param{size} sia un
multiplo di \param{alignment}.  Infine si tenga presente infine che nessuna di
queste funzioni inizializza il buffer di memoria allocato, il loro
comportamento cioè è analogo, allineamento a parte, a quello di \func{malloc}.

Un secondo caso in cui risulta estremamente utile poter avere un maggior
controllo delle modalità di allocazione della memoria è quello in cui cercano
errori di programmazione. Esempi di questi errori sono i \textit{double free},
o i cosiddetti \itindex{buffer~overrun} \textit{buffer overrun}, cioè le
scritture su un buffer oltre le dimensioni della sua
allocazione,\footnote{entrambe queste operazioni causano in genere la
  corruzione dei dati di controllo delle funzioni di allocazione, che vengono
  anch'essi mantenuti nello \textit{heap} per tenere traccia delle zone di
  memoria allocata.} o i classici \textit{memory leak}.

Abbiamo visto in sez.~\ref{sec:proc_mem_lock} come una prima funzionalità di
ausilio nella ricerca di questi errori sia l'uso della variabile di ambiente
\envvar{MALLOC\_CHECK\_}.  Una modalità alternativa per effettuare dei
controlli di consistenza sullo stato delle allocazioni di memoria eseguite con
\func{malloc}, anche questa fornita come estensione specifica (e non standard)
della \acr{glibc}, è quella di utilizzare la funzione \funcd{mcheck}, che deve
essere chiamata prima di eseguire qualunque allocazione con \func{malloc}; il
suo prototipo è:

\begin{funcproto}{ 
\fhead{mcheck.h} 
\fdecl{int mcheck(void (*abortfn) (enum mcheck\_status status))}
\fdesc{Attiva i controlli di consistenza delle allocazioni di memoria.}   
}
{La funzione ritorna $0$ in caso di successo e $-1$ per un errore;
  \var{errno} non viene impostata.} 
\end{funcproto}

La funzione consente di registrare una funzione di emergenza che verrà
eseguita tutte le volte che, in una successiva esecuzione di \func{malloc},
venissero trovate delle inconsistenze, come delle operazioni di scrittura
oltre i limiti dei buffer allocati. Per questo motivo la funzione deve essere
chiamata prima di qualunque allocazione di memoria, altrimenti fallirà.

Se come primo argomento di \func{mcheck} si passa \val{NULL} verrà utilizzata
una funzione predefinita che stampa un messaggio di errore ed invoca la
funzione \func{abort} (vedi sez.~\ref{sec:sig_alarm_abort}), altrimenti si
dovrà creare una funzione personalizzata in grado di ricevere il tipo di
errore ed agire di conseguenza.

Nonostante la scarsa leggibilità del prototipo si tratta semplicemente di
definire una funzione di tipo \code{void abortfn(enum mcheck\_status status)},
che non deve restituire nulla e che deve avere un unico argomento di tipo
\code{mcheck\_status}. In caso di errore la funzione verrà eseguita ricevendo
un opportuno valore di \param{status} che è un tipo enumerato che può assumere
soltanto i valori di tab.~\ref{tab:mcheck_status_value} che indicano la
tipologia di errore riscontrata.

\begin{table}[htb]
  \centering
  \footnotesize
  \begin{tabular}[c]{|l|p{7cm}|}
    \hline
    \textbf{Valore} & \textbf{Significato} \\
    \hline
    \hline
    \constd{MCHECK\_OK}      & Riportato a \func{mprobe} se nessuna
                               inconsistenza è presente.\\
    \constd{MCHECK\_DISABLED}& Riportato a \func{mprobe} se si è chiamata
                               \func{mcheck} dopo aver già usato
                               \func{malloc}.\\
    \constd{MCHECK\_HEAD}    & I dati immediatamente precedenti il buffer sono
                               stati modificati, avviene in genere quando si
                               decrementa eccessivamente il valore di un
                               puntatore scrivendo poi prima dell'inizio del
                               buffer.\\
    \constd{MCHECK\_TAIL}    & I dati immediatamente seguenti il buffer sono
                               stati modificati, succede quando si va scrivere
                               oltre la dimensione corretta del buffer.\\
    \constd{MCHECK\_FREE}    & Il buffer è già stato disallocato.\\
    \hline
  \end{tabular}
  \caption{Valori dello stato dell'allocazione di memoria ottenibili dalla
    funzione di terminazione installata con \func{mcheck}.} 
  \label{tab:mcheck_status_value}
\end{table}

Una volta che si sia chiamata \func{mcheck} con successo si può anche
controllare esplicitamente lo stato delle allocazioni senza aspettare un
errore nelle relative funzioni utilizzando la funzione \funcd{mprobe}, il cui
prototipo è:

\begin{funcproto}{ 
\fhead{mcheck.h} 
\fdecl{enum mcheck\_status mprobe(ptr)}
\fdesc{Esegue un controllo di consistenza delle allocazioni.}   
}
{La funzione ritorna un codice fra quelli riportati in
   tab.~\ref{tab:mcheck_status_value} e non ha errori.} 
\end{funcproto}

La funzione richiede che si passi come argomento un puntatore ad un blocco di
memoria precedentemente allocato con \func{malloc} o \func{realloc}, e
restituisce lo stesso codice di errore che si avrebbe per la funzione di
emergenza ad una successiva chiamata di una funzione di allocazione, e poi i
primi due codici che indicano rispettivamente quando tutto è a posto o il
controllo non è possibile per non aver chiamato \func{mcheck} in tempo.

% TODO: trattare le altre funzionalità avanzate di \func{malloc}, mallopt,
% mtrace, muntrace, mallinfo e gli hook con le glibc 2.10 c'è pure malloc_info
% a sostituire mallinfo, vedi http://udrepper.livejournal.com/20948.html


\section{Argomenti, ambiente ed altre proprietà di un processo}
\label{sec:proc_options}

In questa sezione esamineremo le funzioni che permettono di gestire gli
argomenti e le opzioni, e quelle che consentono di manipolare ed utilizzare le
variabili di ambiente. Accenneremo infine alle modalità con cui si può gestire
la localizzazione di un programma modificandone il comportamento a seconda
della lingua o del paese a cui si vuole faccia riferimento nelle sue
operazioni. 

\subsection{Il formato degli argomenti}
\label{sec:proc_par_format}

Tutti i programmi hanno la possibilità di ricevere argomenti e opzioni quando
vengono lanciati e come accennato in sez.~\ref{sec:proc_main} questo viene
effettuato attraverso gli argomenti \param{argc} e \param{argv} ricevuti nella
funzione \code{main} all'avvio del programma. Questi argomenti vengono passati
al programma dalla shell o dal processo che esegue la \func{exec} (secondo le
modalità che vedremo in sez.~\ref{sec:proc_exec}) quando questo viene messo in
esecuzione.

Nel caso più comune il passaggio di argomenti ed opzioni viene effettuato
dalla shell, che si incarica di leggere la linea di comando con cui si lancia
il programma e di effettuarne la scansione (il cosiddetto \textit{parsing})
per individuare le parole che la compongono, ciascuna delle quali potrà essere
considerata un argomento o un'opzione. 

Di norma per individuare le parole che andranno a costituire la lista degli
argomenti viene usato come carattere di separazione lo spazio o il tabulatore,
ma la cosa dipende ovviamente dalle modalità con cui si effettua la scansione
e dalle convenzioni adottate dal programma che la esegue: ad esempio la shell
consente di proteggere con opportuni caratteri di controllo argomenti che
contengono degli spazi evitando di spezzarli in parole diverse.

\begin{figure}[htb]
  \centering
  \includegraphics[width=13cm]{img/argv_argc}
  % \begin{tikzpicture}[>=stealth]
  % \draw (0.5,2.5) rectangle (3.5,3);
  % \draw (2,2.75) node {\texttt{argc = 5}};
  % \draw (5,2.5) rectangle (8,3);
  % \draw (6.5,2.75) node {\texttt{argv[0]}};
  % \draw [->] (8,2.75) -- (9,2.75);
  % \draw (9,2.75) node [anchor=west] {\texttt{"touch"}};
  % \draw (5,2) rectangle (8,2.5);
  % \draw (6.5,2.25) node {\texttt{argv[1]}};
  % \draw [->] (8,2.25) -- (9,2.25);
  % \draw (9,2.25) node [anchor=west] {\texttt{"-r"}};
  % \draw (5,1.5) rectangle (8,2);
  % \draw (6.5,1.75) node {\texttt{argv[2]}};
  % \draw [->] (8,1.75) -- (9,1.75);
  % \draw (9,1.75) node [anchor=west] {\texttt{"riferimento.txt"}};
  % \draw (5,1.0) rectangle (8,1.5);
  % \draw (6.5,1.25) node {\texttt{argv[3]}};
  % \draw [->] (8,1.25) -- (9,1.25);
  % \draw (9,1.25) node [anchor=west] {\texttt{"-m"}};
  % \draw (5,0.5) rectangle (8,1.0);
  % \draw (6.5,0.75) node {\texttt{argv[4]}};
  % \draw [->] (8,0.75) -- (9,0.75);
  % \draw (9,0.75) node [anchor=west] {\texttt{"questofile.txt"}};
  % \draw (4.25,3.5) node{\texttt{"touch -r riferimento.txt -m questofile.txt"}};
  % \end{tikzpicture}
  \caption{Esempio dei valori di \param{argv} e \param{argc} generati nella 
    scansione di una riga di comando.}
  \label{fig:proc_argv_argc}
\end{figure}

Indipendentemente da come viene eseguita, il risultato finale della scansione
dovrà comunque essere la costruzione del vettore di puntatori \param{argv} in
cui si devono inserire in successione i puntatori alle stringhe costituenti i
vari argomenti ed opzioni da passare al programma, e della
variabile \param{argc} che deve essere inizializzata al numero di stringhe
contenute in \param{argv}. Nel caso della shell questo comporta ad esempio che
il primo argomento sia sempre il nome del programma. Un esempio di questo
meccanismo è mostrato in fig.~\ref{fig:proc_argv_argc}, che illustra il
risultato della scansione di una riga di comando.


\subsection{La gestione delle opzioni}
\label{sec:proc_opt_handling}

In generale un programma Unix riceve da linea di comando sia gli argomenti che
le opzioni, queste ultime sono standardizzate per essere riconosciute come
tali: un elemento di \param{argv} successivo al primo che inizia con il
carattere ``\texttt{-}'' e che non sia un singolo ``\texttt{-}'' o un
``\texttt{-{}-}'' viene considerato un'opzione.  In genere le opzioni sono
costituite da una lettera singola (preceduta dal carattere ``\texttt{-}'') e
possono avere o no un parametro associato.

Un esempio tipico può essere quello mostrato in
fig.~\ref{fig:proc_argv_argc}. In quel caso le opzioni sono \cmd{-r} e
\cmd{-m} e la prima vuole un parametro mentre la seconda no
(\cmd{questofile.txt} è un argomento del programma, non un parametro di
\cmd{-m}).

Per gestire le opzioni all'interno degli argomenti a linea di comando passati
in \param{argv} la libreria standard del C fornisce la funzione
\funcd{getopt}, che ha il seguente prototipo:

\begin{funcproto}{ 
\fhead{unistd.h} 
\fdecl{int getopt(int argc, char * const argv[], const char *optstring)}
\fdesc{Esegue la scansione delle opzioni negli argomenti della funzione
  \code{main}.} 
}
{Ritorna il carattere che segue l'opzione, ``\texttt{:}'' se manca un
  parametro all'opzione, ``\texttt{?}'' se l'opzione è sconosciuta, e $-1$ se
  non esistono altre opzioni.} 
\end{funcproto}

Questa funzione prende come argomenti le due variabili \param{argc} e
\param{argv} che devono essere quelle passate come argomenti di \code{main}
all'esecuzione del programma, ed una stringa \param{optstring} che indica
quali sono le opzioni valide. La funzione effettua la scansione della lista
degli argomenti ricercando ogni stringa che comincia con il carattere
``\texttt{-}'' e ritorna ogni volta che trova un'opzione valida.

La stringa \param{optstring} indica quali sono le opzioni riconosciute ed è
costituita da tutti i caratteri usati per identificare le singole opzioni, se
l'opzione ha un parametro al carattere deve essere fatto seguire il carattere
di due punti (``\texttt{:}''); nel caso di fig.~\ref{fig:proc_argv_argc} ad
esempio la stringa di opzioni avrebbe dovuto contenere \texttt{"r:m"}.

La modalità di uso di \func{getopt} è pertanto quella di chiamare più volte la
funzione all'interno di un ciclo, fintanto che essa non ritorna il valore $-1$
che indica che non ci sono più opzioni. Nel caso si incontri un'opzione non
dichiarata in \param{optstring} viene ritornato il carattere ``\texttt{?}''
mentre se un'opzione che lo richiede non è seguita da un parametro viene
ritornato il carattere ``\texttt{:}'', infine se viene incontrato il valore
``\texttt{-{}-}'' la scansione viene considerata conclusa, anche se vi sono
altri elementi di \param{argv} che cominciano con il carattere ``\texttt{-}''.

Quando \func{getopt} trova un'opzione fra quelle indicate in \param{optstring}
essa ritorna il valore numerico del carattere, in questo modo si possono
eseguire azioni specifiche usando uno \instruction{switch}; la funzione
inoltre inizializza alcune variabili globali:
\begin{itemize*}
\item \var{char *optarg} contiene il puntatore alla stringa parametro
  dell'opzione.
\item \var{int optind} alla fine della scansione restituisce l'indice del
  primo elemento di \param{argv} che non è un'opzione.
\item \var{int opterr} previene, se posto a zero, la stampa di un messaggio
  di errore in caso di riconoscimento di opzioni non definite.
\item \var{int optopt} contiene il carattere dell'opzione non riconosciuta.
\end{itemize*}

\begin{figure}[!htb]
  \footnotesize \centering
  \begin{minipage}[c]{\codesamplewidth}
  \includecodesample{listati/option_code.c}
  \end{minipage}
  \normalsize
  \caption{Esempio di codice per la gestione delle opzioni.}
  \label{fig:proc_options_code}
\end{figure}

In fig.~\ref{fig:proc_options_code} si è mostrata la sezione del programma
\file{fork\_test.c}, che useremo nel prossimo capitolo per effettuare dei test
sulla creazione dei processi, deputata alla decodifica delle opzioni a riga di
comando da esso supportate.

Si può notare che si è anzitutto (\texttt{\small 1}) disabilitata la stampa di
messaggi di errore per opzioni non riconosciute, per poi passare al ciclo per
la verifica delle opzioni (\texttt{\small 2-27}); per ciascuna delle opzioni
possibili si è poi provveduto ad un'azione opportuna, ad esempio per le tre
opzioni che prevedono un parametro si è effettuata la decodifica del medesimo,
il cui indirizzo è contenuto nella variabile \var{optarg}), avvalorando la
relativa variabile (\texttt{\small 12-14}, \texttt{\small 15-17} e
\texttt{\small 18-20}). Completato il ciclo troveremo in \var{optind}
l'indice in \code{argv[]} del primo degli argomenti rimanenti nella linea di
comando.

Normalmente \func{getopt} compie una permutazione degli elementi di
\param{argv} cosicché alla fine della scansione gli elementi che non sono
opzioni sono spostati in coda al vettore. Oltre a questa esistono altre due
modalità di gestire gli elementi di \param{argv}; se \param{optstring} inizia
con il carattere ``\texttt{+}'' (o è impostata la variabile di ambiente
\cmd{POSIXLY\_CORRECT}) la scansione viene fermata non appena si incontra un
elemento che non è un'opzione.

L'ultima modalità, usata quando un programma può gestire la mescolanza fra
opzioni e argomenti, ma se li aspetta in un ordine definito, si attiva
quando \param{optstring} inizia con il carattere ``\texttt{-}''. In questo caso
ogni elemento che non è un'opzione viene considerato comunque un'opzione e
associato ad un valore di ritorno pari ad 1, questo permette di identificare
gli elementi che non sono opzioni, ma non effettua il riordinamento del
vettore \param{argv}.


\subsection{Le variabili di ambiente}
\label{sec:proc_environ}

\index{variabili!di~ambiente|(}
Oltre agli argomenti passati a linea di comando esiste un'altra modalità che
permette di trasferire ad un processo delle informazioni in modo da
modificarne il comportamento.  Ogni processo infatti riceve dal sistema, oltre
alle variabili \param{argv} e \param{argc} anche un \textsl{ambiente} (in
inglese \textit{environment}); questo viene espresso nella forma di una lista
(chiamata \textit{environment list}) delle cosiddette \textsl{variabili di
  ambiente}, i valori di queste variabili possono essere poi usati dal
programma.

Anche in questo caso la lista delle \textsl{variabili di ambiente} deve essere
costruita ed utilizzata nella chiamata alla funzione \func{exec} (torneremo su
questo in sez.~\ref{sec:proc_exec}) quando questo viene lanciato. Come per la
lista degli argomenti anche questa lista è un vettore di puntatori a
caratteri, ciascuno dei quali punta ad una stringa, terminata da un
\val{NULL}. A differenza di \code{argv[]} in questo caso non si ha una
lunghezza del vettore data da un equivalente di \param{argc}, ma la lista è
terminata da un puntatore nullo.

L'indirizzo della lista delle variabili di ambiente è passato attraverso la
variabile globale \var{environ}, che viene definita automaticamente per
ciascun processo, e a cui si può accedere attraverso una semplice
dichiarazione del tipo:
\includecodesnip{listati/env_ptr.c}
un esempio della struttura di questa lista, contenente alcune delle variabili
più comuni che normalmente sono definite dal sistema, è riportato in
fig.~\ref{fig:proc_envirno_list}.
\begin{figure}[htb]
  \centering
  \includegraphics[width=13cm]{img/environ_var}
  % \begin{tikzpicture}[>=stealth]
  % \draw (2,3.5) node {\textsf{Environment pointer}};
  % \draw (6,3.5) node {\textsf{Environment list}};
  % \draw (10.5,3.5) node {\textsf{Environment string}};
  % \draw (0.5,2.5) rectangle (3.5,3);
  % \draw (2,2.75) node {\texttt{environ}};
  % \draw [->] (3.5,2.75) -- (4.5,2.75);
  % \draw (4.5,2.5) rectangle (7.5,3);
  % \draw (6,2.75) node {\texttt{environ[0]}};
  % \draw (4.5,2) rectangle (7.5,2.5);
  % \draw (6,2.25) node {\texttt{environ[1]}};
  % \draw (4.5,1.5) rectangle (7.5,2);
  % \draw (4.5,1) rectangle (7.5,1.5);
  % \draw (4.5,0.5) rectangle (7.5,1);
  % \draw (4.5,0) rectangle (7.5,0.5);
  % \draw (6,0.25) node {\texttt{NULL}};
  % \draw [->] (7.5,2.75) -- (8.5,2.75);
  % \draw (8.5,2.75) node[right] {\texttt{HOME=/home/piccardi}};
  % \draw [->] (7.5,2.25) -- (8.5,2.25);
  % \draw (8.5,2.25) node[right] {\texttt{PATH=:/bin:/usr/bin}};
  % \draw [->] (7.5,1.75) -- (8.5,1.75);
  % \draw (8.5,1.75) node[right] {\texttt{SHELL=/bin/bash}};
  % \draw [->] (7.5,1.25) -- (8.5,1.25);
  % \draw (8.5,1.25) node[right] {\texttt{EDITOR=emacs}};
  % \draw [->] (7.5,0.75) -- (8.5,0.75);
  % \draw (8.5,0.75) node[right] {\texttt{OSTYPE=linux-gnu}};
  % \end{tikzpicture}
  \caption{Esempio di lista delle variabili di ambiente.}
  \label{fig:proc_envirno_list}
\end{figure}

Per convenzione le stringhe che definiscono l'ambiente sono tutte del tipo
\textsl{\texttt{NOME=valore}} ed in questa forma che le funzioni di gestione
che vedremo a breve se le aspettano, se pertanto si dovesse costruire
manualmente un ambiente si abbia cura di rispettare questa convenzione.
Inoltre alcune variabili, come quelle elencate in
fig.~\ref{fig:proc_envirno_list}, sono definite dal sistema per essere usate
da diversi programmi e funzioni: per queste c'è l'ulteriore convenzione di
usare nomi espressi in caratteri maiuscoli.\footnote{ma si tratta solo di una
  convenzione, niente vieta di usare caratteri minuscoli, come avviene in vari
  casi.}

Il kernel non usa mai queste variabili, il loro uso e la loro interpretazione è
riservata alle applicazioni e ad alcune funzioni di libreria; in genere esse
costituiscono un modo comodo per definire un comportamento specifico senza
dover ricorrere all'uso di opzioni a linea di comando o di file di
configurazione. É di norma cura della shell, quando esegue un comando, passare
queste variabili al programma messo in esecuzione attraverso un uso opportuno
delle relative chiamate (si veda sez.~\ref{sec:proc_exec}).

La shell ad esempio ne usa molte per il suo funzionamento, come \envvar{PATH}
per indicare la lista delle directory in cui effettuare la ricerca dei comandi
o \envvar{PS1} per impostare il proprio \textit{prompt}. Alcune di esse, come
\envvar{HOME}, \envvar{USER}, ecc. sono invece definite al login (per i
dettagli si veda sez.~\ref{sec:sess_login}), ed in genere è cura della propria
distribuzione definire le opportune variabili di ambiente in uno script di
avvio. Alcune servono poi come riferimento generico per molti programmi, come
\envvar{EDITOR} che indica l'editor preferito da invocare in caso di
necessità. Una in particolare, \envvar{LANG}, serve a controllare la
localizzazione del programma 
%(su cui torneremo in sez.~\ref{sec:proc_localization}) 
per adattarlo alla lingua ed alle convezioni
dei vari paesi.

Gli standard POSIX e XPG3 definiscono alcune di queste variabili (le più
comuni), come riportato in tab.~\ref{tab:proc_env_var}. GNU/Linux le supporta
tutte e ne definisce anche altre, in particolare poi alcune funzioni di
libreria prevedono la presenza di specifiche variabili di ambiente che ne
modificano il comportamento, come quelle usate per indicare una localizzazione
e quelle per indicare un fuso orario; una lista più completa che comprende
queste ed ulteriori variabili si può ottenere con il comando \cmd{man 7
  environ}.

\begin{table}[htb]
  \centering
  \footnotesize
  \begin{tabular}[c]{|l|c|c|c|l|}
    \hline
    \textbf{Variabile} & \textbf{POSIX} & \textbf{XPG3} 
    & \textbf{Linux} & \textbf{Descrizione} \\
    \hline
    \hline
    \texttt{USER}   &$\bullet$&$\bullet$&$\bullet$& Nome utente.\\
    \texttt{LOGNAME}&$\bullet$&$\bullet$&$\bullet$& Nome di login.\\
    \texttt{HOME}   &$\bullet$&$\bullet$&$\bullet$& Directory base
                                                    dell'utente.\\
    \texttt{LANG}   &$\bullet$&$\bullet$&$\bullet$& Localizzazione.\\
    \texttt{PATH}   &$\bullet$&$\bullet$&$\bullet$& Elenco delle directory
                                                    dei programmi.\\
    \texttt{PWD}    &$\bullet$&$\bullet$&$\bullet$& Directory corrente.\\
    \texttt{SHELL}  &$\bullet$&$\bullet$&$\bullet$& Shell in uso.\\
    \texttt{TERM}   &$\bullet$&$\bullet$&$\bullet$& Tipo di terminale.\\
    \texttt{PAGER}  &$\bullet$&$\bullet$&$\bullet$& Programma per vedere i
                                                    testi.\\
    \texttt{EDITOR} &$\bullet$&$\bullet$&$\bullet$& Editor preferito.\\
    \texttt{BROWSER}&$\bullet$&$\bullet$&$\bullet$& Browser preferito.\\
    \texttt{TMPDIR} &$\bullet$&$\bullet$&$\bullet$& Directory dei file
                                                    temporanei.\\
    \hline
  \end{tabular}
  \caption{Esempi delle variabili di ambiente più comuni definite da vari
    standard.} 
  \label{tab:proc_env_var}
\end{table}

Lo standard ANSI C prevede l'esistenza di un ambiente, e pur non entrando
nelle specifiche di come sono strutturati i contenuti, definisce la funzione
\funcd{getenv} che permette di ottenere i valori delle variabili di ambiente;
il suo prototipo è:

\begin{funcproto}{ 
\fhead{stdlib.h}
\fdecl{char *getenv(const char *name)}
\fdesc{Cerca una variabile di ambiente del processo.} 
}
{La funzione ritorna il puntatore alla stringa contenente il valore della
  variabile di ambiente in caso di successo e \val{NULL} per un errore.} 
\end{funcproto}

La funzione effettua una ricerca nell'ambiente del processo cercando una
variabile il cui nome corrisponda a quanto indicato con
l'argomento \param{name}, ed in caso di successo ritorna il puntatore alla
stringa che ne contiene il valore, nella forma ``\texttt{NOME=valore}''.

\begin{table}[htb]
  \centering
  \footnotesize
  \begin{tabular}[c]{|l|c|c|c|c|c|c|}
    \hline
    \textbf{Funzione} & \textbf{ANSI C} & \textbf{POSIX.1} & \textbf{XPG3} & 
    \textbf{SVr4} & \textbf{BSD} & \textbf{Linux} \\
    \hline
    \hline
    \func{getenv}  & $\bullet$ & $\bullet$ & $\bullet$ 
                   & $\bullet$ & $\bullet$ & $\bullet$ \\
    \func{setenv}  &    --     &    --     &   --      
                   &    --     & $\bullet$ & $\bullet$ \\
    \func{unsetenv}&    --     &    --     &   --       
                   &    --     & $\bullet$ & $\bullet$ \\
    \func{putenv}  &    --     & opz.      & $\bullet$ 
                   &    --     & $\bullet$ & $\bullet$ \\
    \func{clearenv}&    --     & opz.      &   --
                   &    --     &    --     & $\bullet$ \\
    \hline
  \end{tabular}
  \caption{Funzioni per la gestione delle variabili di ambiente.}
  \label{tab:proc_env_func}
\end{table}

Oltre a questa funzione di lettura, che è l'unica definita dallo standard ANSI
C, nell'evoluzione dei sistemi Unix ne sono state proposte altre, da
utilizzare per impostare, modificare e cancellare le variabili di
ambiente. Uno schema delle funzioni previste nei vari standard e disponibili
in Linux è riportato in tab.~\ref{tab:proc_env_func}. Tutte le funzioni sono
state comunque inserite nello standard POSIX.1-2001, ad eccetto di
\func{clearenv} che è stata rigettata.

In Linux sono definite tutte le funzioni elencate in
tab.~\ref{tab:proc_env_func},\footnote{in realtà nelle libc4 e libc5 sono
  definite solo le prime quattro, \func{clearenv} è stata introdotta con la
  \acr{glibc} 2.0.} anche se parte delle funzionalità sono ridondanti. La
prima funzione di manipolazione che prenderemo in considerazione è
\funcd{putenv}, che consente di aggiungere, modificare e cancellare una
variabile di ambiente; il suo prototipo è:

\begin{funcproto}{ 
\fdecl{int putenv(char *string)}
\fdesc{Inserisce, modifica o rimuove una variabile d'ambiente.} 
}
{La funzione ritorna $0$ in caso di successo e $-1$ per un errore, che può
  essere solo \errval{ENOMEM}.}
\end{funcproto}

La funzione prende come argomento una stringa analoga a quella restituita da
\func{getenv} e sempre nella forma ``\texttt{NOME=valore}''. Se la variabile
specificata (nel caso \texttt{NOME}) non esiste la stringa sarà aggiunta
all'ambiente, se invece esiste il suo valore sarà impostato a quello
specificato dal contenuto di \param{string} (nel caso \texttt{valore}).  Se
invece si passa come argomento solo il nome di una variabile di ambiente
(cioè \param{string} è nella forma ``\texttt{NOME}'' e non contiene il
carattere ``\texttt{=}'') allora questa, se presente nell'ambiente, verrà
cancellata.

Si tenga presente che, seguendo lo standard SUSv2, le \acr{glibc} successive
alla versione 2.1.2 aggiungono direttamente \param{string} nella lista delle
variabili di ambiente illustrata in fig.~\ref{fig:proc_envirno_list}
sostituendo il relativo puntatore;\footnote{il comportamento è lo stesso delle
  vecchie \acr{libc4} e \acr{libc5}; nella \acr{glibc}, dalla versione 2.0
  alla 2.1.1, veniva invece fatta una copia, seguendo il comportamento di
  BSD4.4; dato che questo può dar luogo a perdite di memoria e non rispetta lo
  standard il comportamento è stato modificato a partire dalla 2.1.2,
  eliminando anche, sempre in conformità a SUSv2, l'attributo \direct{const}
  dal prototipo.}  pertanto ogni cambiamento alla stringa in questione si
riflette automaticamente sull'ambiente, e quindi si deve evitare di passare a
questa funzione una variabile automatica (per evitare i problemi esposti in
sez.~\ref{sec:proc_var_passing}). Benché non sia richiesto dallo standard,
nelle versioni della \acr{glibc} a partire dalla 2.1 la funzione è rientrante
(vedi sez.~\ref{sec:proc_reentrant}).

Infine quando una chiamata a \func{putenv} comporta la necessità di creare una
nuova versione del vettore \var{environ} questo sarà allocato automaticamente,
ma la versione corrente sarà deallocata solo se anch'essa è risultante da
un'allocazione fatta in precedenza da un'altra \func{putenv}. Questo avviene
perché il vettore delle variabili di ambiente iniziale, creato dalla chiamata
ad \func{exec} (vedi sez.~\ref{sec:proc_exec}) è piazzato nella memoria al di
sopra dello \textit{stack}, (vedi fig.~\ref{fig:proc_mem_layout}) e non nello
\textit{heap} e quindi non può essere deallocato.  Inoltre la memoria
associata alle variabili di ambiente eliminate non viene liberata.

Come alternativa a \func{putenv} si può usare la funzione \funcd{setenv} che
però consente solo di aggiungere o modificare una variabile di ambiente; il
suo prototipo è:

\begin{funcproto}{ 
\fhead{stdlib.h}
\fdecl{int setenv(const char *name, const char *value, int overwrite)}
\fdesc{Inserisce o modifica una variabile di ambiente.} 
}
{La funzione ritorna $0$ in caso di successo e $-1$ per un errore,
  nel qual caso \var{errno} assumerà uno dei valori:
  \begin{errlist}
  \item[\errcode{EINVAL}] \param{name} è \val{NULL} o una stringa di lunghezza
  nulla o che contiene il carattere ``\texttt{=}''.
  \item[\errcode{ENOMEM}] non c'è memoria sufficiente per aggiungere una nuova
    variabile all'ambiente.
\end{errlist}}
\end{funcproto}

La funzione consente di specificare separatamente nome e valore della
variabile di ambiente da aggiungere negli argomenti \param{name}
e \param{value}. Se la variabile è già presente nell'ambiente
l'argomento \param{overwrite} specifica il comportamento della funzione, se
diverso da zero sarà sovrascritta, se uguale a zero sarà lasciata immutata.  A
differenza di \func{putenv} la funzione esegue delle copie del contenuto degli
argomenti \param{name} e \param{value} e non è necessario preoccuparsi di
allocarli in maniera permanente.

La cancellazione di una variabile di ambiente viene invece gestita
esplicitamente con \funcd{unsetenv}, il cui prototipo è:

\begin{funcproto}{ 
\fhead{stdlib.h}
\fdecl{int unsetenv(const char *name)}
\fdesc{Rimuove una variabile di ambiente.} 
}
{La funzione ritorna $0$ in caso di successo e $-1$ per un errore,
  nel qual caso \var{errno} assumerà uno dei valori:
  \begin{errlist}
  \item[\errcode{EINVAL}] \param{name} è \val{NULL} o una stringa di lunghezza
  nulla o che contiene il carattere ``\texttt{=}''.
\end{errlist}}
\end{funcproto}

La funzione richiede soltanto il nome della variabile di ambiente
nell'argomento \param{name}, se la variabile non esiste la funzione ritorna
comunque con un valore di successo.\footnote{questo con le versioni della
  \acr{glibc} successive la 2.2.2, per le precedenti \func{unsetenv} era
  definita come \texttt{void} e non restituiva nessuna informazione.}

L'ultima funzione per la gestione dell'ambiente è
\funcd{clearenv},\footnote{che come accennato è l'unica non presente nello
  standard POSIX.1-2000, ed è disponibili solo per versioni della \acr{glibc}
  a partire dalla 2.0; per poterla utilizzare occorre aver definito le macro
  \macro{\_SVID\_SOURCE} e \macro{\_XOPEN\_SOURCE}.} che viene usata per
cancellare completamente tutto l'ambiente; il suo prototipo è:

\begin{funcproto}{ 
\fhead{stdlib.h}
\fdecl{int clearenv(void)}
\fdesc{Cancella tutto l'ambiente.} 
}
{La funzione ritorna $0$ in caso di successo e un valore diverso da zero per
  un errore.}
\end{funcproto}

In genere si usa questa funzione in maniera precauzionale per evitare i
problemi di sicurezza connessi nel trasmettere ai programmi che si invocano un
ambiente che può contenere dei dati non controllati, le cui variabili possono
causare effetti indesiderati. Con l'uso della funzione si provvede alla
cancellazione di tutto l'ambiente originale in modo da poterne costruirne una
versione ``\textsl{sicura}'' da zero.

\index{variabili!di~ambiente|)}


% \subsection{La localizzazione}
% \label{sec:proc_localization}

% Abbiamo accennato in sez.~\ref{sec:proc_environ} come la variabile di ambiente
% \envvar{LANG} sia usata per indicare ai processi il valore della cosiddetta
% \textsl{localizzazione}. Si tratta di una funzionalità fornita dalle librerie
% di sistema\footnote{prenderemo in esame soltanto il caso della \acr{glibc}.}
% che consente di gestire in maniera automatica sia la lingua in cui vengono
% stampati i vari messaggi (come i messaggi associati agli errori che vedremo in
% sez.~\ref{sec:sys_strerror}) che le convenzioni usate nei vari paesi per una
% serie di aspetti come il formato dell'ora, quello delle date, gli ordinamenti
% alfabetici, le espressioni della valute, ecc.

% Da finire.

% La localizzazione di un programma si può selezionare con la 

% In realtà perché un programma sia effettivamente localizzato non è sufficiente 

% TODO trattare, quando ci sarà tempo, setlocale ed il resto


%\subsection{Opzioni in formato esteso}
%\label{sec:proc_opt_extended}

%Oltre alla modalità ordinaria di gestione delle opzioni trattata in
%sez.~\ref{sec:proc_opt_handling} la \acr{glibc} fornisce una modalità
%alternativa costituita dalle cosiddette \textit{long-options}, che consente di
%esprimere le opzioni in una forma più descrittiva che nel caso più generale è
%qualcosa del tipo di ``\texttt{-{}-option-name=parameter}''.

%(NdA: questa parte verrà inserita in seguito).

% TODO opzioni in formato esteso

% TODO trattare il vettore ausiliario e getauxval (vedi
% http://lwn.net/Articles/519085/)


\section{Problematiche di programmazione generica}
\label{sec:proc_gen_prog}

Benché questo non sia un libro sul linguaggio C, è opportuno affrontare alcune
delle problematiche generali che possono emergere nella programmazione con
questo linguaggio e di quali precauzioni o accorgimenti occorre prendere per
risolverle. Queste problematiche non sono specifiche di sistemi unix-like o
\textit{multitasking}, ma avendo trattato in questo capitolo il comportamento
dei processi visti come entità a sé stanti, le riportiamo qui.


\subsection{Il passaggio di variabili e valori di ritorno nelle funzioni}
\label{sec:proc_var_passing}

Una delle caratteristiche standard del C è che le variabili vengono passate
alle funzioni che si invocano in un programma attraverso un meccanismo che
viene chiamato \textit{by value}, diverso ad esempio da quanto avviene con il
Fortran, dove le variabili sono passate, come suol dirsi, \textit{by
  reference}, o dal C++ dove la modalità del passaggio può essere controllata
con l'operatore \cmd{\&}.

Il passaggio di una variabile \textit{by value} significa che in realtà quello
che viene passato alla funzione è una copia del valore attuale di quella
variabile, copia che la funzione potrà modificare a piacere, senza che il
valore originale nella funzione chiamante venga toccato. In questo modo non
occorre preoccuparsi di eventuali effetti delle operazioni svolte nella
funzione stessa sulla variabile passata come argomento.

Questo però va inteso nella maniera corretta. Il passaggio \textit{by value}
vale per qualunque variabile, puntatori compresi; quando però in una funzione
si usano dei puntatori (ad esempio per scrivere in un buffer) in realtà si va
a modificare la zona di memoria a cui essi puntano, per cui anche se i
puntatori sono copie, i dati a cui essi puntano saranno sempre gli stessi, e
le eventuali modifiche avranno effetto e saranno visibili anche nella funzione
chiamante.

Nella maggior parte delle funzioni di libreria e delle \textit{system call} i
puntatori vengono usati per scambiare dati (attraverso i buffer o le strutture
a cui fanno riferimento) e le variabili normali vengono usate per specificare
argomenti; in genere le informazioni a riguardo dei risultati vengono passate
alla funzione chiamante attraverso il valore di ritorno.  È buona norma
seguire questa pratica anche nella programmazione normale.

\itindbeg{value~result~argument}

Talvolta però è necessario che la funzione possa restituire indietro alla
funzione chiamante un valore relativo ad uno dei suoi argomenti usato anche in
ingresso.  Per far questo si usa il cosiddetto \textit{value result argument},
si passa cioè, invece di una normale variabile, un puntatore alla stessa. Gli
esempi di questa modalità di passaggio sono moltissimi, ad esempio essa viene
usata nelle funzioni che gestiscono i socket (in
sez.~\ref{sec:TCP_functions}), in cui, per permettere al kernel di restituire
informazioni sulle dimensioni delle strutture degli indirizzi utilizzate,
viene usato proprio questo meccanismo.

Occorre tenere ben presente questa differenza, perché le variabili passate in
maniera ordinaria, che vengono inserite nello \textit{stack}, cessano di
esistere al ritorno di una funzione, ed ogni loro eventuale modifica
all'interno della stessa sparisce con la conclusione della stessa, per poter
passare delle informazioni occorre quindi usare un puntatore che faccia
riferimento ad un indirizzo accessibile alla funzione chiamante.

\itindend{value~result~argument}

Questo requisito di accessibilità è fondamentale, infatti dei possibili
problemi che si possono avere con il passaggio dei dati è quello di restituire
alla funzione chiamante dei dati che sono contenuti in una variabile
automatica.  Ovviamente quando la funzione ritorna la sezione dello
\textit{stack} che conteneva la variabile automatica (si ricordi quanto detto
in sez.~\ref{sec:proc_mem_alloc}) verrà liberata automaticamente e potrà
essere riutilizzata all'invocazione di un'altra funzione, con le immaginabili
conseguenze, quasi invariabilmente catastrofiche, di sovrapposizione e
sovrascrittura dei dati.

Per questo una delle regole fondamentali della programmazione in C è che
all'uscita di una funzione non deve restare nessun riferimento alle sue
variabili locali. Qualora sia necessario utilizzare delle variabili che devono
essere viste anche dalla funzione chiamante queste devono essere allocate
esplicitamente, o in maniera statica usando variabili globali o dichiarate
come \direct{extern},\footnote{la direttiva \direct{extern} informa il
  compilatore che la variabile che si è dichiarata in una funzione non è da
  considerarsi locale, ma globale, e per questo allocata staticamente e
  visibile da tutte le funzioni dello stesso programma.} o dinamicamente con
una delle funzioni della famiglia \func{malloc}, passando opportunamente il
relativo puntatore fra le funzioni.


\subsection{Il passaggio di un numero variabile di argomenti}
\label{sec:proc_variadic}

\index{funzioni!\textit{variadic}|(}

Come vedremo nei capitoli successivi, non sempre è possibile specificare un
numero fisso di argomenti per una funzione.  Lo standard ISO C prevede nella
sua sintassi la possibilità di definire delle \textit{variadic function} che
abbiano un numero variabile di argomenti, attraverso l'uso nella dichiarazione
della funzione dello speciale costrutto ``\texttt{...}'', che viene chiamato
\textit{ellipsis}.

Lo standard però non provvede a livello di linguaggio alcun meccanismo con cui
dette funzioni possono accedere ai loro argomenti.  L'accesso viene pertanto
realizzato a livello della libreria standard del C che fornisce gli strumenti
adeguati.  L'uso di una \textit{variadic function} prevede quindi tre punti:
\begin{itemize*}
\item \textsl{dichiarare} la funzione come \textit{variadic} usando un
  prototipo che contenga una \textit{ellipsis};
\item \textsl{definire} la funzione come \textit{variadic} usando la stessa
  \textit{ellipsis}, ed utilizzare le apposite macro che consentono la
  gestione di un numero variabile di argomenti;
\item \textsl{invocare} la funzione specificando prima gli argomenti fissi, ed
  a seguire quelli addizionali.
\end{itemize*}

Lo standard ISO C prevede che una \textit{variadic function} abbia sempre
almeno un argomento fisso. Prima di effettuare la dichiarazione deve essere
incluso l'apposito \textit{header file} \headfile{stdarg.h}; un esempio di
dichiarazione è il prototipo della funzione \func{execl} che vedremo in
sez.~\ref{sec:proc_exec}:
\includecodesnip{listati/exec_sample.c}
in questo caso la funzione prende due argomenti fissi ed un numero variabile
di altri argomenti, che andranno a costituire gli elementi successivi al primo
del vettore \param{argv} passato al nuovo processo. Lo standard ISO C richiede
inoltre che l'ultimo degli argomenti fissi sia di tipo
\textit{self-promoting}\footnote{il linguaggio C prevede che quando si
  mescolano vari tipi di dati, alcuni di essi possano essere \textsl{promossi}
  per compatibilità; ad esempio i tipi \ctyp{float} vengono convertiti
  automaticamente a \ctyp{double} ed i \ctyp{char} e gli \ctyp{short} ad
  \ctyp{int}. Un tipo \textit{self-promoting} è un tipo che verrebbe promosso
  a sé stesso.} il che esclude vettori, puntatori a funzioni e interi di tipo
\ctyp{char} o \ctyp{short} (con segno o meno). Una restrizione ulteriore di
alcuni compilatori è di non dichiarare l'ultimo argomento fisso come variabile
di tipo \direct{register}.\footnote{la direttiva \direct{register} del
  compilatore chiede che la variabile dichiarata tale sia mantenuta, nei
  limiti del possibile, all'interno di un registro del processore; questa
  direttiva è originaria dell'epoca dai primi compilatori, quando stava al
  programmatore scrivere codice ottimizzato, riservando esplicitamente alle
  variabili più usate l'uso dei registri del processore, oggi questa direttiva
  è in disuso pressoché completo dato che tutti i compilatori sono normalmente
  in grado di valutare con maggior efficacia degli stessi programmatori quando
  sia il caso di eseguire questa ottimizzazione.}

Una volta dichiarata la funzione il secondo passo è accedere ai vari argomenti
quando la si va a definire. Gli argomenti fissi infatti hanno un loro nome, ma
quelli variabili vengono indicati in maniera generica dalla
\textit{ellipsis}. L'unica modalità in cui essi possono essere recuperati è
pertanto quella sequenziale, in cui vengono estratti dallo \textit{stack}
secondo l'ordine in cui sono stati scritti nel prototipo della funzione.

\macrobeg{va\_start}

Per fare questo in \headfile{stdarg.h} sono definite delle macro specifiche,
previste dallo standard ISO C89, che consentono di eseguire questa operazione.
La prima di queste macro è \macro{va\_start}, che inizializza opportunamente
una lista degli argomenti, la sua definizione è:

{\centering
\begin{funcbox}{ 
\fhead{stdarg.h}
\fdecl{void va\_start(va\_list ap, last)}
\fdesc{Inizializza una lista degli argomenti di una funzione
  \textit{variadic}.} 
}
\end{funcbox}}

La macro inizializza il puntatore alla lista di argomenti \param{ap} che deve
essere una apposita variabile di tipo \type{va\_list}; il
parametro \param{last} deve indicare il nome dell'ultimo degli argomenti fissi
dichiarati nel prototipo della funzione \textit{variadic}.

\macrobeg{va\_arg}

La seconda macro di gestione delle liste di argomenti di una funzione
\textit{variadic} è \macro{va\_arg}, che restituisce in successione un
argomento della lista; la sua definizione è:

{\centering
\begin{funcbox}{ 
\fhead{stdarg.h}
\fdecl{type va\_arg(va\_list ap, type)}
\fdesc{Restituisce il valore del successivo argomento opzionale.} 
}
\end{funcbox}}
 
La macro restituisce il valore di un argomento, modificando opportunamente la
lista \param{ap} perché una chiamata successiva restituisca l'argomento
seguente. La macro richiede che si specifichi il tipo dell'argomento che si
andrà ad estrarre attraverso il parametro \param{type} che sarà anche il tipo
del valore da essa restituito. Si ricordi che il tipo deve essere
\textit{self-promoting}.

In generale è perfettamente legittimo richiedere meno argomenti di quelli che
potrebbero essere stati effettivamente forniti, per cui nella esecuzione delle
\macro{va\_arg} ci si può fermare in qualunque momento ed i restanti argomenti
saranno ignorati. Se invece si richiedono più argomenti di quelli
effettivamente forniti si otterranno dei valori indefiniti. Si avranno
risultati indefiniti anche quando si chiama \macro{va\_arg} specificando un
tipo che non corrisponde a quello usato per il corrispondente argomento.

\macrobeg{va\_end}

Infine una volta completata l'estrazione occorre indicare che si sono concluse
le operazioni con la macro \macrod{va\_end}, la cui definizione è:

{\centering
\begin{funcbox}{ 
\fhead{stdarg.h}
\fdecl{void va\_end(va\_list ap)}
\fdesc{Conclude l'estrazione degli argomenti di una funzione
  \textit{variadic}.} 
}
\end{funcbox}}
 
Dopo l'uso di \macro{va\_end} la variabile \param{ap} diventa indefinita e
successive chiamate a \macro{va\_arg} non funzioneranno.  Nel caso del
\cmd{gcc} l'uso di \macro{va\_end} può risultare inutile, ma è comunque
necessario usarla per chiarezza del codice, per compatibilità con diverse
implementazioni e per eventuali modifiche future a questo comportamento.

Riassumendo la procedura da seguire per effettuare l'estrazione degli
argomenti di una funzione \textit{variadic} è la seguente:
\begin{enumerate*}
\item inizializzare una lista degli argomenti attraverso la macro
  \macro{va\_start};
\item accedere agli argomenti con chiamate successive alla macro
  \macro{va\_arg}: la prima chiamata restituirà il primo argomento, la seconda
  il secondo e così via;
\item dichiarare la conclusione dell'estrazione degli argomenti invocando la
  macro \macro{va\_end}.
\end{enumerate*}

Si tenga presente che si possono usare anche più liste degli argomenti,
ciascuna di esse andrà inizializzata con \macro{va\_start} e letta con
\macro{va\_arg}, e ciascuna potrà essere usata per scandire la lista degli
argomenti in modo indipendente. Infine ciascuna scansione dovrà essere
terminata con \macro{va\_end}.

Un limite di queste macro è che i passi 1) e 3) devono essere eseguiti nel
corpo principale della funzione, il passo 2) invece può essere eseguito anche
in un'altra funzione, passandole lista degli argomenti \param{ap}. In questo
caso però al ritorno della funzione \macro{va\_arg} non può più essere usata
(anche se non si era completata l'estrazione) dato che il valore di \param{ap}
risulterebbe indefinito.

\macroend{va\_start}
\macroend{va\_arg}
\macroend{va\_end}

Esistono dei casi in cui è necessario eseguire più volte la scansione degli
argomenti e poter memorizzare una posizione durante la stessa. In questo caso
sembrerebbe naturale copiarsi la lista degli argomenti \param{ap} con una
semplice assegnazione ad un'altra variabile dello stesso tipo. Dato che una
delle realizzazioni più comuni di \type{va\_list} è quella di un puntatore
nello \textit{stack} all'indirizzo dove sono stati salvati gli argomenti, è
assolutamente normale pensare di poter effettuare questa operazione.

\index{tipo!opaco|(}

In generale però possono esistere anche realizzazioni diverse, ed è per questo
motivo che invece che un semplice puntatore, \typed{va\_list} è quello che
viene chiamato un \textsl{tipo opaco}. Si chiamano così quei tipi di dati, in
genere usati da una libreria, la cui struttura interna non deve essere vista
dal programma chiamante (da cui deriva il nome opaco) che li devono utilizzare
solo attraverso dalle opportune funzioni di gestione.

\index{tipo!opaco|)}

Per questo motivo una variabile di tipo \typed{va\_list} non può essere
assegnata direttamente ad un'altra variabile dello stesso tipo, ma lo standard
ISO C99\footnote{alcuni sistemi che non hanno questa macro provvedono al suo
  posto \macrod{\_\_va\_copy} che era il nome proposto in una bozza dello
  standard.}  ha previsto una macro ulteriore che permette di eseguire la
copia di una lista degli argomenti:

{\centering
\begin{funcbox}{ 
\fhead{stdarg.h}
\fdecl{void va\_copy(va\_list dest, va\_list src)}
\fdesc{Copia la lista degli argomenti di una funzione \textit{variadic}.} 
}
\end{funcbox}}

La macro copia l'attuale della lista degli argomenti \param{src} su una nuova
lista \param{dest}. Anche in questo caso è buona norma chiudere ogni
esecuzione di una \macrod{va\_copy} con una corrispondente \macro{va\_end} sul
nuovo puntatore alla lista degli argomenti.

La chiamata di una funzione con un numero variabile di argomenti, posto che la
si sia dichiarata e definita come tale, non prevede nulla di particolare;
l'invocazione è identica alle altre, con gli argomenti, sia quelli fissi che
quelli opzionali, separati da virgole. Quello che però è necessario tenere
presente è come verranno convertiti gli argomenti variabili.

In Linux gli argomenti dello stesso tipo sono passati allo stesso modo, sia
che siano fissi sia che siano opzionali (alcuni sistemi trattano diversamente
gli opzionali), ma dato che il prototipo non può specificare il tipo degli
argomenti opzionali, questi verranno sempre promossi, pertanto nella ricezione
dei medesimi occorrerà tenerne conto (ad esempio un \ctyp{char} verrà visto da
\macro{va\_arg} come \ctyp{int}).

Un altro dei problemi che si devono affrontare con le funzioni con un numero
variabile di argomenti è che non esiste un modo generico che permetta di
stabilire quanti sono gli argomenti effettivamente passati in una chiamata.

Esistono varie modalità per affrontare questo problema; una delle più
immediate è quella di specificare il numero degli argomenti opzionali come uno
degli argomenti fissi. Una variazione di questo metodo è l'uso di un argomento
fisso per specificare anche il tipo degli argomenti variabili, come fa la
stringa di formato per \func{printf} (vedi sez.~\ref{sec:file_formatted_io}).

Infine una ulteriore modalità diversa, che può essere applicata solo quando il
tipo degli argomenti lo rende possibile, è quella che prevede di usare un
valore speciale per l'ultimo argomento, come fa ad esempio \func{execl} che
usa un puntatore \val{NULL} per indicare la fine della lista degli argomenti
(vedi sez.~\ref{sec:proc_exec}).

\index{funzioni!\textit{variadic}|)}

\subsection{Il controllo di flusso non locale}
\label{sec:proc_longjmp}

Il controllo del flusso di un programma in genere viene effettuato con le
varie istruzioni del linguaggio C; fra queste la più bistrattata è il
\instruction{goto}, che viene deprecato in favore dei costrutti della
programmazione strutturata, che rendono il codice più leggibile e
mantenibile. Esiste però un caso in cui l'uso di questa istruzione porta
all'implementazione più efficiente e più chiara anche dal punto di vista della
struttura del programma: quello dell'uscita in caso di errore.

\index{salto~non-locale|(} 

Il C però non consente di effettuare un salto ad una etichetta definita in
un'altra funzione, per cui se l'errore avviene in una funzione, e la sua
gestione ordinaria è in un'altra, occorre usare quello che viene chiamato un
\textsl{salto non-locale}.  Il caso classico in cui si ha questa necessità,
citato sia in \cite{APUE} che in \cite{GlibcMan}, è quello di un programma nel
cui corpo principale vengono letti dei dati in ingresso sui quali viene
eseguita, tramite una serie di funzioni di analisi, una scansione dei
contenuti, da cui si ottengono le indicazioni per l'esecuzione di opportune
operazioni.

Dato che l'analisi può risultare molto complessa, ed opportunamente suddivisa
in fasi diverse, la rilevazione di un errore nei dati in ingresso può accadere
all'interno di funzioni profondamente annidate l'una nell'altra. In questo
caso si dovrebbe gestire, per ciascuna fase, tutta la casistica del passaggio
all'indietro di tutti gli errori rilevabili dalle funzioni usate nelle fasi
successive.  Questo comporterebbe una notevole complessità, mentre sarebbe
molto più comodo poter tornare direttamente al ciclo di lettura principale,
scartando l'input come errato.\footnote{a meno che, come precisa
  \cite{GlibcMan}, alla chiusura di ciascuna fase non siano associate
  operazioni di pulizia specifiche (come deallocazioni, chiusure di file,
  ecc.), che non potrebbero essere eseguite con un salto non-locale.}

Tutto ciò può essere realizzato proprio con un salto non-locale; questo di
norma viene realizzato salvando il contesto dello \textit{stack} nel punto in
cui si vuole tornare in caso di errore, e ripristinandolo, in modo da tornare
quando serve nella funzione da cui si era partiti.  La funzione che permette
di salvare il contesto dello \textit{stack} è \funcd{setjmp}, il cui prototipo
è:

\begin{funcproto}{ 
\fhead{setjmp.h}
\fdecl{int setjmp(jmp\_buf env)}
\fdesc{Salva il contesto dello \textit{stack}.} 
}
{La funzione ritorna $0$ quando è chiamata direttamente ed un valore diverso
  da zero quando ritorna da una chiamata di \func{longjmp} che usa il contesto
  salvato in precedenza.}
\end{funcproto}
  
Quando si esegue la funzione il contesto corrente dello \textit{stack} viene
salvato nell'argomento \param{env}, una variabile di tipo
\typed{jmp\_buf}\footnote{anche questo è un classico esempio di variabile di
  \textsl{tipo opaco}.}  che deve essere stata definita in precedenza. In
genere le variabili di tipo \type{jmp\_buf} vengono definite come variabili
globali in modo da poter essere viste in tutte le funzioni del programma.

Quando viene eseguita direttamente la funzione ritorna sempre zero, un valore
diverso da zero viene restituito solo quando il ritorno è dovuto ad una
chiamata di \func{longjmp} in un'altra parte del programma che ripristina lo
\textit{stack} effettuando il salto non-locale. Si tenga conto che il contesto
salvato in \param{env} viene invalidato se la funzione che ha chiamato
\func{setjmp} ritorna, nel qual caso un successivo uso di \func{longjmp} può
comportare conseguenze imprevedibili (e di norma fatali) per il processo.
  
Come accennato per effettuare un salto non-locale ad un punto precedentemente
stabilito con \func{setjmp} si usa la funzione \funcd{longjmp}; il suo
prototipo è:

\begin{funcproto}{ 
\fhead{setjmp.h}
\fdecl{void longjmp(jmp\_buf env, int val)}
\fdesc{Ripristina il contesto dello stack.} 
}
{La funzione non ritorna.}   
\end{funcproto}

La funzione ripristina il contesto dello \textit{stack} salvato da una
chiamata a \func{setjmp} nell'argomento \param{env}. Dopo l'esecuzione della
funzione il programma prosegue nel codice successivo alla chiamata della
\func{setjmp} con cui si era salvato \param{env}, che restituirà il valore
dell'argomento \param{val} invece di zero.  Il valore
dell'argomento \param{val} deve essere sempre diverso da zero, se si è
specificato 0 sarà comunque restituito 1 al suo posto.

In sostanza l'esecuzione di \func{longjmp} è analoga a quella di una
istruzione \instr{return}, solo che invece di ritornare alla riga
successiva della funzione chiamante, il programma in questo caso ritorna alla
posizione della relativa \func{setjmp}. L'altra differenza fondamentale con
\instr{return} è che il ritorno può essere effettuato anche attraverso
diversi livelli di funzioni annidate.

L'implementazione di queste funzioni comporta alcune restrizioni dato che esse
interagiscono direttamente con la gestione dello \textit{stack} ed il
funzionamento del compilatore stesso. In particolare \func{setjmp} è
implementata con una macro, pertanto non si può cercare di ottenerne
l'indirizzo, ed inoltre le chiamate a questa funzione sono sicure solo in uno
dei seguenti casi:
\begin{itemize*}
\item come espressione di controllo in un comando condizionale, di selezione o
  di iterazione (come \instruction{if}, \instruction{switch} o
  \instruction{while});
\item come operando per un operatore di uguaglianza o confronto in una
  espressione di controllo di un comando condizionale, di selezione o di
  iterazione;
\item come operando per l'operatore di negazione (\code{!}) in una espressione
  di controllo di un comando condizionale, di selezione o di iterazione;
\item come espressione a sé stante.
\end{itemize*}

In generale, dato che l'unica differenza fra il risultato di una chiamata
diretta di \func{setjmp} e quello ottenuto nell'uscita con un \func{longjmp} è
costituita dal valore di ritorno della funzione, quest'ultima viene usualmente
chiamata all'interno di un una istruzione \instr{if} che permetta di
distinguere i due casi.

Uno dei punti critici dei salti non-locali è quello del valore delle
variabili, ed in particolare quello delle variabili automatiche della funzione
a cui si ritorna. In generale le variabili globali e statiche mantengono i
valori che avevano al momento della chiamata di \func{longjmp}, ma quelli
delle variabili automatiche (o di quelle dichiarate \dirct{register}) sono in
genere indeterminati.

Quello che succede infatti è che i valori delle variabili che sono tenute in
memoria manterranno il valore avuto al momento della chiamata di
\func{longjmp}, mentre quelli tenuti nei registri del processore (che nella
chiamata ad un'altra funzione vengono salvati nel contesto nello
\textit{stack}) torneranno al valore avuto al momento della chiamata di
\func{setjmp}; per questo quando si vuole avere un comportamento coerente si
può bloccare l'ottimizzazione che porta le variabili nei registri
dichiarandole tutte come \direct{volatile}.\footnote{la direttiva
  \direct{volatile} informa il compilatore che la variabile che è dichiarata
  può essere modificata, durante l'esecuzione del nostro, da altri programmi.
  Per questo motivo occorre dire al compilatore che non deve essere mai
  utilizzata l'ottimizzazione per cui quanto opportuno essa viene mantenuta in
  un registro, poiché in questo modo si perderebbero le eventuali modifiche
  fatte dagli altri programmi (che avvengono solo in una copia posta in
  memoria).}

\index{salto~non-locale|)}


% TODO trattare qui le restartable sequences (vedi
% https://lwn.net/Articles/664645/ e https://lwn.net/Articles/650333/) se e
% quando saranno introdotte

\subsection{La \textit{endianness}}
\label{sec:endianness}

\itindbeg{endianness} 

Un altro dei problemi di programmazione che può dar luogo ad effetti
imprevisti è quello relativo alla cosiddetta \textit{endianness}.  Questa è
una caratteristica generale dell'architettura hardware di un computer che
dipende dal fatto che la rappresentazione di un numero binario può essere
fatta in due modi, chiamati rispettivamente \textit{big endian} e
\textit{little endian}, a seconda di come i singoli bit vengono aggregati per
formare le variabili intere (ed in genere in diretta corrispondenza a come
sono poi in realtà cablati sui bus interni del computer).

\begin{figure}[!htb]
  \centering \includegraphics[height=3cm]{img/endianness}
  \caption{Schema della disposizione dei dati in memoria a seconda della
    \textit{endianness}.}
  \label{fig:sock_endianness}
\end{figure}

Per capire meglio il problema si consideri un intero a 32 bit scritto in una
locazione di memoria posta ad un certo indirizzo. Come illustrato in
fig.~\ref{fig:sock_endianness} i singoli bit possono essere disposti in memoria
in due modi: a partire dal più significativo o a partire dal meno
significativo.  Così nel primo caso si troverà il byte che contiene i bit più
significativi all'indirizzo menzionato e il byte con i bit meno significativi
nell'indirizzo successivo; questo ordinamento è detto \textit{big endian},
dato che si trova per prima la parte più grande. Il caso opposto, in cui si
parte dal bit meno significativo è detto per lo stesso motivo \textit{little
  endian}.

Si può allora verificare quale tipo di \textit{endianness} usa il proprio
computer con un programma elementare che si limita ad assegnare un valore ad
una variabile per poi ristamparne il contenuto leggendolo un byte alla volta.
Il codice di detto programma, \file{endtest.c}, è nei sorgenti allegati,
allora se lo eseguiamo su un normale PC compatibile, che è \textit{little
  endian} otterremo qualcosa del tipo:
\begin{Console}
[piccardi@gont sources]$ \textbf{./endtest}
Using value ABCDEF01
val[0]= 1
val[1]=EF
val[2]=CD
val[3]=AB
\end{Console}
%$
mentre su un vecchio Macintosh con PowerPC, che è \textit{big endian} avremo
qualcosa del tipo:
\begin{Console}
piccardi@anarres:~/gapil/sources$ \textbf{./endtest}
Using value ABCDEF01
val[0]=AB
val[1]=CD
val[2]=EF
val[3]= 1
\end{Console}
%$

L'attenzione alla \textit{endianness} nella programmazione è importante, perché
se si fanno assunzioni relative alla propria architettura non è detto che
queste restino valide su un'altra architettura. Inoltre, come vedremo ad
esempio in sez.~\ref{sec:sock_addr_func}, si possono avere problemi quando ci
si trova a usare valori di un formato con una infrastruttura che ne usa
un altro. 

La \textit{endianness} di un computer dipende essenzialmente dalla architettura
hardware usata; Intel e Digital usano il \textit{little endian}, Motorola,
IBM, Sun (sostanzialmente tutti gli altri) usano il \textit{big endian}. Il
formato dei dati contenuti nelle intestazioni dei protocolli di rete (il
cosiddetto \textit{network order}) è anch'esso \textit{big endian}; altri
esempi di uso di questi due diversi formati sono quello del bus PCI, che è
\textit{little endian}, o quello del bus VME che è \textit{big endian}.

Esistono poi anche dei processori che possono scegliere il tipo di formato
all'avvio e alcuni che, come il PowerPC o l'Intel i860, possono pure passare
da un tipo di ordinamento all'altro con una specifica istruzione. In ogni caso
in Linux l'ordinamento è definito dall'architettura e dopo l'avvio del sistema
in genere resta sempre lo stesso,\footnote{su architettura PowerPC è possibile
  cambiarlo, si veda sez.~\ref{sec:process_prctl}.} anche quando il processore
permetterebbe di eseguire questi cambiamenti.

\begin{figure}[!htbp]
  \footnotesize \centering
  \begin{minipage}[c]{\codesamplewidth}
    \includecodesample{listati/endian.c}
  \end{minipage} 
  \normalsize
  \caption{La funzione \samplefunc{endian}, usata per controllare il tipo di
    architettura della macchina.}
  \label{fig:sock_endian_code}
\end{figure}

Per controllare quale tipo di ordinamento si ha sul proprio computer si è
scritta una piccola funzione di controllo, il cui codice è riportato
fig.~\ref{fig:sock_endian_code}, che restituisce un valore nullo (falso) se
l'architettura è \textit{big endian} ed uno non nullo (vero) se l'architettura
è \textit{little endian}.

Come si vede la funzione è molto semplice, e si limita, una volta assegnato
(\texttt{\small 9}) un valore di test pari a \texttt{0xABCD} ad una variabile
di tipo \ctyp{short} (cioè a 16 bit), a ricostruirne una copia byte a byte.
Per questo prima (\texttt{\small 10}) si definisce il puntatore \var{ptr} per
accedere al contenuto della prima variabile, ed infine calcola (\texttt{\small
  11}) il valore della seconda assumendo che il primo byte sia quello meno
significativo (cioè, per quanto visto in fig.~\ref{fig:sock_endianness}, che sia
\textit{little endian}). Infine la funzione restituisce (\texttt{\small 12})
il valore del confronto delle due variabili. 

In generale non ci si deve preoccupare della \textit{endianness} all'interno
di un programma fintanto che questo non deve generare o manipolare dei dati
che sono scambiati con altre macchine, ad esempio via rete o tramite dei file
binari. Nel primo caso la scelta è già stata fatta nella standardizzazione dei
protocolli, che hanno adottato il \textit{big endian} (che viene detto anche
per questo \textit{network order}); vedremo in sez.~\ref{sec:sock_func_ord} le
funzioni di conversione che devono essere usate.

Nel secondo caso occorre sapere quale \textit{endianness} è stata usata nei
dati memorizzati sul file e tenerne conto nella rilettura e nella
manipolazione e relativa modifica (e salvataggio). La gran parte dei formati
binari standardizzati specificano quale \textit{endianness} viene utilizzata e
basterà identificare qual'è, se se ne deve definire uno per i propri scopi
basterà scegliere una volta per tutte quale usare e attenersi alla scelta.

\itindend{endianness}


% LocalWords:  like exec kernel thread main ld linux static linker char envp Gb
% LocalWords:  sez POSIX exit system call cap abort shell diff errno stdlib int
% LocalWords:  SUCCESS FAILURE void atexit stream fclose unistd descriptor init
% LocalWords:  SIGCHLD wait function glibc SunOS arg argp execve fig high kb Mb
% LocalWords:  memory alpha swap table printf Unit MMU paging fault SIGSEGV BSS
% LocalWords:  multitasking text segment NULL Block Started Symbol fill black
% LocalWords:  heap stack calling convention size malloc calloc realloc nmemb
% LocalWords:  ENOMEM ptr uClib cfree error leak smartpointers hook Dmalloc brk
% LocalWords:  Gray Watson Electric Fence Bruce Perens sbrk longjmp SUSv BSD ap
% LocalWords:  ptrdiff increment locking lock copy write capabilities IPC mlock
% LocalWords:  capability MEMLOCK limits getpagesize RLIMIT munlock sys const
% LocalWords:  addr len EINVAL EPERM mlockall munlockall flags l'OR CURRENT IFS
% LocalWords:  argc argv parsing questofile txt getopt optstring switch optarg
% LocalWords:  optind opterr optopt POSIXLY CORRECT long options NdA group
% LocalWords:  option parameter list environ PATH HOME XPG tab LOGNAME LANG PWD
% LocalWords:  TERM PAGER TMPDIR getenv name SVr setenv unsetenv putenv opz gcc
% LocalWords:  clearenv libc value overwrite string reference result argument
% LocalWords:  socket variadic ellipsis header stdarg execl self promoting last
% LocalWords:  float double short register type dest src extern setjmp jmp buf
% LocalWords:  env return if while Di page cdecl  rectangle node anchor west PS
% LocalWords:  environment rounded corners dashed south width height draw east
% LocalWords:  exithandler handler violation inline SOURCE SVID XOPEN mincore
% LocalWords:  length unsigned vec EFAULT EAGAIN dell'I memalign valloc posix
% LocalWords:  boundary memptr alignment sizeof overrun mcheck abortfn enum big
% LocalWords:  mprobe DISABLED HEAD TAIL touch right emacs OSTYPE endianness IBM
% LocalWords:  endian little endtest Macintosh PowerPC Intel Digital Motorola
% LocalWords:  Sun order VME  loader Windows DLL shared objects PRELOAD termios
% LocalWords:  is to LC SIG str mem wcs assert ctype dirent fcntl signal stdio
% LocalWords:  times library utmp syscall number Filesystem Hierarchy pathname
% LocalWords:  context assembler sysconf fork Dinamic huge segmentation program
% LocalWords:  break store using intptr ssize overflow ONFAULT faulting alloc
%  LocalWords:  scheduler pvalloc aligned ISOC ABCDEF

%%% Local Variables: 
%%% mode: latex
%%% TeX-master: "gapil"
%%% End: 

%% prochand.tex
%%
%% Copyright (C) 2000-2018 by Simone Piccardi.  Permission is granted to
%% copy, distribute and/or modify this document under the terms of the GNU Free
%% Documentation License, Version 1.1 or any later version published by the
%% Free Software Foundation; with the Invariant Sections being "Un preambolo",
%% with no Front-Cover Texts, and with no Back-Cover Texts.  A copy of the
%% license is included in the section entitled "GNU Free Documentation
%% License".
%%

\chapter{La gestione dei processi}
\label{cha:process_handling}

Come accennato nell'introduzione in un sistema unix-like tutte le operazioni
vengono svolte tramite opportuni processi.  In sostanza questi ultimi vengono
a costituire l'unità base per l'allocazione e l'uso delle risorse del sistema.

Nel precedente capitolo abbiamo esaminato il funzionamento di un processo come
unità a se stante, in questo esamineremo il funzionamento dei processi
all'interno del sistema. Saranno cioè affrontati i dettagli della creazione e
della terminazione dei processi, della gestione dei loro attributi e
privilegi, e di tutte le funzioni a questo connesse. Infine nella sezione
finale introdurremo alcune problematiche generiche della programmazione in
ambiente \textit{multitasking}.


\section{Le funzioni di base della gestione dei processi}
\label{sec:proc_handling}

In questa sezione tratteremo le problematiche della gestione dei processi
all'interno del sistema, illustrandone tutti i dettagli.  Inizieremo con una
panoramica dell'architettura dei processi, tratteremo poi le funzioni
elementari che permettono di leggerne gli identificatori, per poi passare alla
spiegazione delle funzioni base che si usano per la creazione e la
terminazione dei processi, e per la messa in esecuzione degli altri programmi.


\subsection{L'architettura della gestione dei processi}
\label{sec:proc_hierarchy}

A differenza di quanto avviene in altri sistemi, ad esempio nel VMS, dove la
generazione di nuovi processi è un'operazione privilegiata, una delle
caratteristiche fondanti di Unix, che esamineremo in dettaglio più avanti, è
che qualunque processo può a sua volta generarne altri. Ogni processo è
identificato presso il sistema da un numero univoco, il cosiddetto
\textit{Process ID}, o più brevemente \ids{PID}, assegnato in forma
progressiva (vedi sez.~\ref{sec:proc_pid}) quando il processo viene creato.

Una seconda caratteristica di un sistema unix-like è che la generazione di un
processo è un'operazione separata rispetto al lancio di un programma. In
genere la sequenza è sempre quella di creare un nuovo processo, il quale
eseguirà, in un passo successivo, il programma desiderato: questo è ad esempio
quello che fa la shell quando mette in esecuzione il programma che gli
indichiamo nella linea di comando.

Una terza caratteristica del sistema è che ogni processo è sempre stato
generato da un altro processo, il processo generato viene chiamato
\textit{processo figlio} (\textit{child process}) mentre quello che lo ha
generato viene chiamato \textsl{processo padre} (\textit{parent
  process}). Questo vale per tutti i processi, con una sola eccezione; dato
che ci deve essere un punto di partenza esiste un processo iniziale (che
normalmente è \cmd{/sbin/init}), che come accennato in
sez.~\ref{sec:intro_kern_and_sys} viene lanciato dal kernel alla conclusione
della fase di avvio. Essendo questo il primo processo lanciato dal sistema ha
sempre \ids{PID} uguale a 1 e non è figlio di nessun altro processo.

Ovviamente \cmd{init} è un processo particolare che in genere si occupa di
lanciare tutti gli altri processi necessari al funzionamento del sistema,
inoltre \cmd{init} è essenziale per svolgere una serie di compiti
amministrativi nelle operazioni ordinarie del sistema (torneremo su alcuni di
essi in sez.~\ref{sec:proc_termination}) e non può mai essere terminato. La
struttura del sistema comunque consente di lanciare al posto di \cmd{init}
qualunque altro programma, e in casi di emergenza (ad esempio se il file di
\cmd{init} si fosse corrotto) è ad esempio possibile lanciare una shell al suo
posto.\footnote{la cosa si fa passando la riga \cmd{init=/bin/sh} come
  parametro di avvio del kernel, l'argomento è di natura sistemistica e
  trattato in sez.~5.3 di \cite{AGL}.}

\begin{figure}[!htb]
  \footnotesize
\begin{Console}
[piccardi@gont piccardi]$ \textbf{pstree -n} 
init-+-keventd
     |-kapm-idled
     |-kreiserfsd
     |-portmap
     |-syslogd
     |-klogd
     |-named
     |-rpc.statd
     |-gpm
     |-inetd
     |-junkbuster
     |-master-+-qmgr
     |        `-pickup
     |-sshd
     |-xfs
     |-cron
     |-bash---startx---xinit-+-XFree86
     |                       `-WindowMaker-+-ssh-agent
     |                                     |-wmtime
     |                                     |-wmmon
     |                                     |-wmmount
     |                                     |-wmppp
     |                                     |-wmcube
     |                                     |-wmmixer
     |                                     |-wmgtemp
     |                                     |-wterm---bash---pstree
     |                                     `-wterm---bash-+-emacs
     |                                                    `-man---pager
     |-5*[getty]
     |-snort
     `-wwwoffled
\end{Console}
%$
  \caption{L'albero dei processi, così come riportato dal comando
    \cmd{pstree}.}
  \label{fig:proc_tree}
\end{figure}

Dato che tutti i processi attivi nel sistema sono comunque generati da
\cmd{init} o da uno dei suoi figli si possono classificare i processi con la
relazione padre/figlio in un'organizzazione gerarchica ad albero. In
fig.~\ref{fig:proc_tree} si è mostrato il risultato del comando \cmd{pstree}
che permette di visualizzare questa struttura, alla cui base c'è \cmd{init}
che è progenitore di tutti gli altri processi.\footnote{in realtà questo non è
  del tutto vero, in Linux, specialmente nelle versioni più recenti del
  kernel, ci sono alcuni processi speciali (come \cmd{keventd}, \cmd{kswapd},
  ecc.) che pur comparendo nei comandi come figli di \cmd{init}, o con
  \ids{PID} successivi ad uno, sono in realtà processi interni al kernel e che
  non rientrano in questa classificazione.}

\itindbeg{process~table}

Il kernel mantiene una tabella dei processi attivi, la cosiddetta
\textit{process table}. Per ciascun processo viene mantenuta una voce in
questa tabella, costituita da una struttura \kstruct{task\_struct}, che
contiene tutte le informazioni rilevanti per quel processo. Tutte le strutture
usate a questo scopo sono dichiarate nell'\textit{header file}
\file{linux/sched.h}, ed in fig.~\ref{fig:proc_task_struct} si è riportato uno
schema semplificato che mostra la struttura delle principali informazioni
contenute nella \texttt{task\_struct}, che in seguito incontreremo a più
riprese.

\begin{figure}[!htb]
  \centering \includegraphics[width=14cm]{img/task_struct}
  \caption{Schema semplificato dell'architettura delle strutture (
    \kstructd{task\_struct}, \kstructd{fs\_struct}, \kstructd{file\_struct})
    usate dal kernel nella gestione dei processi.}
  \label{fig:proc_task_struct}
\end{figure}

\itindend{process~table}

% TODO la task_struct è cambiata per qualche dettaglio vedi anche
% http://www.ibm.com/developerworks/linux/library/l-linux-process-management/
% TODO completare la parte su quando viene chiamato lo scheduler.

\itindbeg{scheduler}

Come accennato in sez.~\ref{sec:intro_unix_struct} è lo \textit{scheduler} che
decide quale processo mettere in esecuzione; esso viene eseguito in occasione
di dell'invocazione di ogni \textit{system call} ed per ogni interrupt
dall'hardware oltre che in una serie di altre occasioni, e può essere anche
attivato esplicitamente. Il timer di sistema provvede comunque a che esso sia
invocato periodicamente, generando un interrupt periodico secondo una
frequenza predeterminata, specificata dalla costante \const{HZ} del kernel
(torneremo su questo argomento in sez.~\ref{sec:sys_unix_time}), che assicura
che lo \textit{scheduler} venga comunque eseguito ad intervalli regolari e
possa prendere le sue decisioni.

A partire dal kernel 2.6.21 è stato introdotto anche un meccanismo
completamente diverso, detto \textit{tickless}, in cui non c'è più una
interruzione periodica con frequenza prefissata, ma ad ogni chiamata del timer
viene programmata l'interruzione successiva sulla base di una stima; in questo
modo si evita di dover eseguire un migliaio di interruzioni al secondo anche
su macchine che non stanno facendo nulla, con un forte risparmio nell'uso
dell'energia da parte del processore che può essere messo in stato di
sospensione anche per lunghi periodi di tempo.

Ma, indipendentemente dalle motivazioni per cui questo avviene, ogni volta che
viene eseguito lo \textit{scheduler} effettua il calcolo delle priorità dei
vari processi attivi (torneremo su questo in sez.~\ref{sec:proc_priority}) e
stabilisce quale di essi debba essere posto in esecuzione fino alla successiva
invocazione.

\itindend{scheduler}

\subsection{Gli identificatori dei processi}
\label{sec:proc_pid}

\itindbeg{Process~ID~(PID)}

Come accennato nella sezione precedente ogni processo viene identificato dal
sistema da un numero identificativo univoco, il \textit{process ID} o
\ids{PID}. Questo è un tipo di dato standard, \type{pid\_t}, che in genere è un
intero con segno (nel caso di Linux e della \acr{glibc} il tipo usato è
\ctyp{int}).

Il \ids{PID} viene assegnato in forma progressiva ogni volta che un nuovo
processo viene creato,\footnote{in genere viene assegnato il numero successivo
  a quello usato per l'ultimo processo creato, a meno che questo numero non
  sia già utilizzato per un altro \ids{PID}, \acr{pgid} o \acr{sid} (vedi
  sez.~\ref{sec:sess_proc_group}).} fino ad un limite che, essendo il
tradizionalmente il \ids{PID} un numero positivo memorizzato in un intero a 16
bit, arriva ad un massimo di 32768.  Oltre questo valore l'assegnazione
riparte dal numero più basso disponibile a partire da un minimo di
300,\footnote{questi valori, fino al kernel 2.4.x, erano definiti dalla macro
  \constd{PID\_MAX} nei file \file{threads.h} e \file{fork.c} dei sorgenti del
  kernel, con il 2.6.x e la nuova interfaccia per i \textit{thread} anche il
  meccanismo di allocazione dei \ids{PID} è stato modificato ed il valore
  massimo è impostabile attraverso il file \sysctlfiled{kernel/pid\_max} e di
  default vale 32768.} che serve a riservare i \ids{PID} più bassi ai processi
eseguiti direttamente dal kernel.  Per questo motivo, come visto in
sez.~\ref{sec:proc_hierarchy}, il processo di avvio (\cmd{init}) ha sempre il
\ids{PID} uguale a uno.

\itindbeg{Parent~Process~ID~(PPID)} 

Tutti i processi inoltre memorizzano anche il \ids{PID} del genitore da cui
sono stati creati, questo viene chiamato in genere \ids{PPID} (da
\textit{Parent Process ID}).  Questi due identificativi possono essere
ottenuti usando le due funzioni di sistema \funcd{getpid} e \funcd{getppid}, i
cui prototipi sono:

\begin{funcproto}{ 
\fhead{sys/types.h}
\fhead{unistd.h}
\fdecl{pid\_t getpid(void)}
\fdesc{Restituisce il \ids{PID} del processo corrente..} 
\fdecl{pid\_t getppid(void)}
\fdesc{Restituisce il \ids{PID} del padre del processo corrente.} 
}
{Entrambe le funzioni non riportano condizioni di errore.}   
\end{funcproto}

\noindent esempi dell'uso di queste funzioni sono riportati in
fig.~\ref{fig:proc_fork_code}, nel programma \file{fork\_test.c}.

Il fatto che il \ids{PID} sia un numero univoco per il sistema lo rende un
candidato per generare ulteriori indicatori associati al processo di cui
diventa possibile garantire l'unicità: ad esempio in alcune implementazioni la
funzione \func{tempnam} (si veda sez.~\ref{sec:file_temp_file}) usa il
\ids{PID} per generare un \textit{pathname} univoco, che non potrà essere
replicato da un altro processo che usi la stessa funzione. Questo utilizzo
però può risultare pericoloso, un \ids{PID} infatti è univoco solo fintanto
che un processo è attivo, una volta terminato esso potrà essere riutilizzato
da un processo completamente diverso, e di questo bisogna essere ben
consapevoli.

Tutti i processi figli dello stesso processo padre sono detti
\textit{sibling}, questa è una delle relazioni usate nel \textsl{controllo di
  sessione}, in cui si raggruppano i processi creati su uno stesso terminale,
o relativi allo stesso login. Torneremo su questo argomento in dettaglio in
cap.~\ref{cha:session}, dove esamineremo gli altri identificativi associati ad
un processo e le varie relazioni fra processi utilizzate per definire una
sessione.

Oltre al \ids{PID} e al \ids{PPID}, e a quelli che vedremo in
sez.~\ref{sec:sess_proc_group}, relativi al controllo di sessione, ad ogni
processo vengono associati degli ulteriori identificatori ed in particolare
quelli che vengono usati per il controllo di accesso.  Questi servono per
determinare se un processo può eseguire o meno le operazioni richieste, a
seconda dei privilegi e dell'identità di chi lo ha posto in esecuzione;
l'argomento è complesso e sarà affrontato in dettaglio in
sez.~\ref{sec:proc_perms}.

\itindend{Process~ID~(PID)}
\itindend{Parent~Process~ID~(PPID)} 

\subsection{La funzione \func{fork} e le funzioni di creazione dei processi}
\label{sec:proc_fork}

La funzione di sistema \func{fork} è la funzione fondamentale della gestione
dei processi: come si è detto tradizionalmente l'unico modo di creare un nuovo
processo era attraverso l'uso di questa funzione,\footnote{in realtà oggi la
  \textit{system call} usata da Linux per creare nuovi processi è \func{clone}
  (vedi \ref{sec:process_clone}), anche perché a partire dalla \acr{glibc}
  2.3.3 non viene più usata la \textit{system call} originale, ma la stessa
  \func{fork} viene implementata tramite \func{clone}, cosa che consente una
  migliore interazione coi \textit{thread}.} essa quindi riveste un ruolo
centrale tutte le volte che si devono scrivere programmi che usano il
\textit{multitasking}.\footnote{oggi questa rilevanza, con la diffusione
  dell'uso dei \textit{thread}\unavref{ che tratteremo al
    cap.~\ref{cha:threads}}, è in parte minore, ma \func{fork} resta comunque
  la funzione principale per la creazione di processi.} Il prototipo di
\funcd{fork} è:

\begin{funcproto}{ 
\fhead{unistd.h}
\fdecl{pid\_t fork(void)}
\fdesc{Crea un nuovo processo.} 
}

{La funzione ritorna in caso di successo il \ids{PID} del figlio nel padre e
  $0$ nel figlio mentre ritorna $-1$ nel padre, senza creare il figlio, per un
  errore, al caso \var{errno} assumerà uno dei valori:
  \begin{errlist}
  \item[\errcode{EAGAIN}] non ci sono risorse sufficienti per creare un altro
    processo (per allocare la tabella delle pagine e le strutture del task) o
    si è esaurito il numero di processi disponibili.
  \item[\errcode{ENOMEM}] non è stato possibile allocare la memoria per le
    strutture necessarie al kernel per creare il nuovo processo.
  \end{errlist}}
\end{funcproto}

Dopo il successo dell'esecuzione di una \func{fork} sia il processo padre che
il processo figlio continuano ad essere eseguiti normalmente, a partire
dall'istruzione successiva alla \func{fork}. Il processo figlio è una copia
del padre, e riceve una copia dei segmenti di testo, dati e dello
\textit{stack} (vedi sez.~\ref{sec:proc_mem_layout}), ed esegue esattamente lo
stesso codice del padre. Si tenga presente però che la memoria è copiata e non
condivisa, pertanto padre e figlio vedranno variabili diverse e le eventuali
modifiche saranno totalmente indipendenti.

\itindbeg{copy~on~write}

Per quanto riguarda la gestione della memoria, in generale il segmento di
testo, che è identico per i due processi, è condiviso e tenuto in sola lettura
per il padre e per i figli. Per gli altri segmenti Linux utilizza la tecnica
del \textit{copy on write}. Questa tecnica comporta che una pagina di memoria
viene effettivamente copiata per il nuovo processo solo quando ci viene
effettuata sopra una scrittura, e si ha quindi una reale differenza fra padre
e figlio.  In questo modo si rende molto più efficiente il meccanismo della
creazione di un nuovo processo, non essendo più necessaria la copia di tutto
lo spazio degli indirizzi virtuali del padre, ma solo delle pagine di memoria
che sono state modificate, e solo al momento della modifica stessa.

\itindend{copy~on~write}

La differenza che si ha nei due processi è che nel processo padre il valore di
ritorno della funzione \func{fork} è il \ids{PID} del processo figlio, mentre
nel figlio è zero; in questo modo il programma può identificare se viene
eseguito dal padre o dal figlio.  Si noti come la funzione \func{fork} ritorni
due volte, una nel padre e una nel figlio.

La scelta di questi valori di ritorno non è casuale, un processo infatti può
avere più figli, ed il valore di ritorno di \func{fork} è l'unico che gli
permette di identificare qual è quello appena creato. Al contrario un figlio
ha sempre un solo padre il cui \ids{PID}, come spiegato in
sez.~\ref{sec:proc_pid}, può sempre essere ottenuto con \func{getppid}; per
questo si ritorna un valore nullo, che non è il \ids{PID} di nessun processo.

Normalmente la chiamata a \func{fork} può fallire solo per due ragioni: o ci
sono già troppi processi nel sistema, il che di solito è sintomo che
qualcos'altro non sta andando per il verso giusto, o si è ecceduto il limite
sul numero totale di processi permessi all'utente, argomento che tratteremo in
dettaglio in sez.~\ref{sec:sys_resource_limit}.

L'uso di \func{fork} avviene secondo due modalità principali; la prima è
quella in cui all'interno di un programma si creano processi figli cui viene
affidata l'esecuzione di una certa sezione di codice, mentre il processo padre
ne esegue un'altra. È il caso tipico dei programmi server (il modello
\textit{client-server} è illustrato in sez.~\ref{sec:net_cliserv}) in cui il
padre riceve ed accetta le richieste da parte dei programmi client, per
ciascuna delle quali pone in esecuzione un figlio che è incaricato di fornire
le risposte associate al servizio.

La seconda modalità è quella in cui il processo vuole eseguire un altro
programma; questo è ad esempio il caso della shell. In questo caso il processo
crea un figlio la cui unica operazione è quella di fare una \func{exec} (di
cui parleremo in sez.~\ref{sec:proc_exec}) subito dopo la \func{fork}.

Alcuni sistemi operativi (il VMS ad esempio) combinano le operazioni di questa
seconda modalità (una \func{fork} seguita da una \func{exec}) in un'unica
operazione che viene chiamata \textit{spawn}. Nei sistemi unix-like è stato
scelto di mantenere questa separazione, dato che, come per la prima modalità
d'uso, esistono numerosi scenari in cui si può usare una \func{fork} senza
aver bisogno di eseguire una \func{exec}. 

Inoltre, anche nel caso della seconda modalità d'uso, avere le due funzioni
separate permette al figlio di cambiare alcune caratteristiche del processo
(maschera dei segnali, redirezione dell'output, utente per conto del cui viene
eseguito, e molto altro su cui torneremo in seguito) prima della \func{exec},
rendendo così relativamente facile intervenire sulle le modalità di esecuzione
del nuovo programma.

\begin{figure}[!htb]
  \footnotesize \centering
  \begin{minipage}[c]{\codesamplewidth}
  \includecodesample{listati/fork_test.c}
  \end{minipage}
  \normalsize
  \caption{Esempio di codice per la creazione di nuovi processi (da
    \file{fork\_test.c}).}
  \label{fig:proc_fork_code}
\end{figure}

In fig.~\ref{fig:proc_fork_code} è riportato il corpo del codice del programma
di esempio \cmd{forktest}, che permette di illustrare molte caratteristiche
dell'uso della funzione \func{fork}. Il programma crea un numero di figli
specificato da linea di comando, e prende anche alcune opzioni per indicare
degli eventuali tempi di attesa in secondi (eseguiti tramite la funzione
\func{sleep}) per il padre ed il figlio (con \cmd{forktest -h} si ottiene la
descrizione delle opzioni). Il codice completo, compresa la parte che gestisce
le opzioni a riga di comando, è disponibile nel file \file{fork\_test.c},
distribuito insieme agli altri sorgenti degli esempi della guida su
\url{http://gapil.gnulinux.it}.

Decifrato il numero di figli da creare, il ciclo principale del programma
(\texttt{\small 24-40}) esegue in successione la creazione dei processi figli
controllando il successo della chiamata a \func{fork} (\texttt{\small
  25-29}); ciascun figlio (\texttt{\small 31-34}) si limita a stampare il
suo numero di successione, eventualmente attendere il numero di secondi
specificato e scrivere un messaggio prima di uscire. Il processo padre invece
(\texttt{\small 36-38}) stampa un messaggio di creazione, eventualmente
attende il numero di secondi specificato, e procede nell'esecuzione del ciclo;
alla conclusione del ciclo, prima di uscire, può essere specificato un altro
periodo di attesa.

Se eseguiamo il comando, che è preceduto dall'istruzione \code{export
  LD\_LIBRARY\_PATH=./} per permettere l'uso delle librerie dinamiche, senza
specificare attese (come si può notare in (\texttt{\small 17-19}) i valori
predefiniti specificano di non attendere), otterremo come risultato sul
terminale:
\begin{Console}
[piccardi@selidor sources]$ \textbf{export LD_LIBRARY_PATH=./; ./forktest 3}
Process 1963: forking 3 child
Spawned 1 child, pid 1964 
Child 1 successfully executing
Child 1, parent 1963, exiting
Go to next child 
Spawned 2 child, pid 1965 
Child 2 successfully executing
Child 2, parent 1963, exiting
Go to next child 
Child 3 successfully executing
Child 3, parent 1963, exiting
Spawned 3 child, pid 1966 
Go to next child 
\end{Console}
%$

Esaminiamo questo risultato: una prima conclusione che si può trarre è che non
si può dire quale processo fra il padre ed il figlio venga eseguito per primo
dopo la chiamata a \func{fork}; dall'esempio si può notare infatti come nei
primi due cicli sia stato eseguito per primo il padre (con la stampa del
\ids{PID} del nuovo processo) per poi passare all'esecuzione del figlio
(completata con i due avvisi di esecuzione ed uscita), e tornare
all'esecuzione del padre (con la stampa del passaggio al ciclo successivo),
mentre la terza volta è stato prima eseguito il figlio (fino alla conclusione)
e poi il padre.

In generale l'ordine di esecuzione dipenderà, oltre che dall'algoritmo di
\textit{scheduling} usato dal kernel, dalla particolare situazione in cui si
trova la macchina al momento della chiamata, risultando del tutto
impredicibile.  Eseguendo più volte il programma di prova e producendo un
numero diverso di figli, si sono ottenute situazioni completamente diverse,
compreso il caso in cui il processo padre ha eseguito più di una \func{fork}
prima che uno dei figli venisse messo in esecuzione.

Pertanto non si può fare nessuna assunzione sulla sequenza di esecuzione delle
istruzioni del codice fra padre e figli, né sull'ordine in cui questi potranno
essere messi in esecuzione. Se è necessaria una qualche forma di precedenza
occorrerà provvedere ad espliciti meccanismi di sincronizzazione, pena il
rischio di incorrere nelle cosiddette \textit{race condition} (vedi
sez.~\ref{sec:proc_race_cond}).

In realtà con l'introduzione dei kernel della serie 2.6 lo \textit{scheduler}
è stato modificato per eseguire sempre per primo il figlio.\footnote{i
  risultati precedenti infatti sono stati ottenuti usando un kernel della
  serie 2.4.}  Questa è una ottimizzazione adottata per evitare che il padre,
effettuando per primo una operazione di scrittura in memoria, attivasse il
meccanismo del \textit{copy on write}, operazione inutile quando il figlio
viene creato solo per eseguire una \func{exec} per lanciare un altro programma
che scarta completamente lo spazio degli indirizzi e rende superflua la copia
della memoria modificata dal padre. Eseguendo sempre per primo il figlio la
\func{exec} verrebbe effettuata subito, con la certezza di utilizzare il
\textit{copy on write} solo quando necessario.

Con il kernel 2.6.32 però il comportamento è stato nuovamente cambiato,
stavolta facendo eseguire per primo sempre il padre. Si è realizzato infatti
che l'eventualità prospettata per la scelta precedente era comunque poco
probabile, mentre l'esecuzione immediata del padre presenta sempre il
vantaggio di poter utilizzare immediatamente tutti i dati che sono nella cache
della CPU e nell'unità di gestione della memoria virtuale, senza doverli
invalidare, cosa che per i processori moderni, che hanno linee di cache
interne molto profonde, avrebbe un forte impatto sulle prestazioni.

Allora anche se quanto detto in precedenza si verifica nel comportamento
effettivo dei programmi soltanto per i kernel fino alla serie 2.4, per
mantenere la portabilità con altri kernel unix-like e con i diversi
comportamenti adottati dalle Linux nella sua evoluzione, è comunque opportuno
non fare nessuna assunzione sull'ordine di esecuzione di padre e figlio dopo
la chiamata a \func{fork}.

Si noti infine come dopo la \func{fork}, essendo i segmenti di memoria
utilizzati dai singoli processi completamente indipendenti, le modifiche delle
variabili nei processi figli, come l'incremento di \var{i} in (\texttt{\small
  31}), sono visibili solo a loro, (ogni processo vede solo la propria copia
della memoria), e non hanno alcun effetto sul valore che le stesse variabili
hanno nel processo padre ed in eventuali altri processi figli che eseguano lo
stesso codice.

Un secondo aspetto molto importante nella creazione dei processi figli è
quello dell'interazione dei vari processi con i file. Ne parleremo qui anche
se buona parte dei concetti relativi ai file verranno trattati più avanti
(principalmente in sez.~\ref{sec:file_unix_interface}). Per illustrare meglio
quello che avviene si può redirigere su un file l'output del programma di
test, quello che otterremo è:
\begin{Console}
[piccardi@selidor sources]$ \textbf{./forktest 3 > output}
[piccardi@selidor sources]$ \textbf{cat output}
Process 1967: forking 3 child
Child 1 successfully executing
Child 1, parent 1967, exiting
Test for forking 3 child
Spawned 1 child, pid 1968 
Go to next child 
Child 2 successfully executing
Child 2, parent 1967, exiting
Test for forking 3 child
Spawned 1 child, pid 1968 
Go to next child 
Spawned 2 child, pid 1969 
Go to next child 
Child 3 successfully executing
Child 3, parent 1967, exiting
Test for forking 3 child
Spawned 1 child, pid 1968 
Go to next child 
Spawned 2 child, pid 1969 
Go to next child 
Spawned 3 child, pid 1970 
Go to next child 
\end{Console}
che come si vede è completamente diverso da quanto ottenevamo sul terminale.

Il comportamento delle varie funzioni di interfaccia con i file è analizzato
in gran dettaglio in sez.~\ref{sec:file_unix_interface} per l'interfaccia
nativa Unix ed in sez.~\ref{sec:files_std_interface} per la standardizzazione
adottata nelle librerie del linguaggio C, valida per qualunque sistema
operativo.

Qui basta accennare che si sono usate le funzioni standard della libreria del
C che prevedono l'output bufferizzato. Il punto è che questa bufferizzazione
(che tratteremo in dettaglio in sez.~\ref{sec:file_buffering}) varia a seconda
che si tratti di un file su disco, in cui il buffer viene scaricato su disco
solo quando necessario, o di un terminale, in cui il buffer viene scaricato ad
ogni carattere di ``a capo''.

Nel primo esempio allora avevamo che, essendovi un a capo nella stringa
stampata, ad ogni chiamata a \func{printf} il buffer veniva scaricato, per cui
le singole righe comparivano a video subito dopo l'esecuzione della
\func{printf}. Ma con la redirezione su file la scrittura non avviene più alla
fine di ogni riga e l'output resta nel buffer. 

Dato che ogni figlio riceve una copia della memoria del padre, esso riceverà
anche quanto c'è nel buffer delle funzioni di I/O, comprese le linee scritte
dal padre fino allora. Così quando il buffer viene scritto su disco all'uscita
del figlio, troveremo nel file anche tutto quello che il processo padre aveva
scritto prima della sua creazione. E alla fine del file (dato che in questo
caso il padre esce per ultimo) troveremo anche l'output completo del padre.

L'esempio ci mostra un altro aspetto fondamentale dell'interazione con i file,
valido anche per l'esempio precedente, ma meno evidente: il fatto cioè che non
solo processi diversi possono scrivere in contemporanea sullo stesso file
(l'argomento dell'accesso concorrente ai file è trattato in dettaglio in
sez.~\ref{sec:file_shared_access}), ma anche che, a differenza di quanto
avviene per le variabili in memoria, la posizione corrente sul file è
condivisa fra il padre e tutti i processi figli.

Quello che succede è che quando lo \textit{standard output}\footnote{si chiama
  così il file su cui di default un programma scrive i suoi dati in uscita,
  tratteremo l'argomento in dettaglio in sez.~\ref{sec:file_fd}.} del padre
viene rediretto come si è fatto nell'esempio, lo stesso avviene anche per
tutti i figli. La funzione \func{fork} infatti ha la caratteristica di
duplicare nei processi figli tutti i \textit{file descriptor} (vedi
sez.~\ref{sec:file_fd}) dei file aperti nel processo padre (allo stesso modo
in cui lo fa la funzione \func{dup}, trattata in sez.~\ref{sec:file_dup}). Ciò
fa sì che padre e figli condividano le stesse voci della \textit{file table}
(tratteremo in dettaglio questi termini in sez.~\ref{sec:file_fd} e
sez.~\ref{sec:file_shared_access}) fra le quali c'è anche la posizione
corrente nel file.

Quando un processo scrive su un file la posizione corrente viene aggiornata
sulla \textit{file table}, e tutti gli altri processi, che vedono la stessa
\textit{file table}, vedranno il nuovo valore. In questo modo si evita, in
casi come quello appena mostrato in cui diversi figli scrivono sullo stesso
file usato dal padre, che una scrittura eseguita in un secondo tempo da un
processo vada a sovrapporsi a quelle precedenti: l'output potrà risultare
mescolato, ma non ci saranno parti perdute per via di una sovrascrittura.

Questo tipo di comportamento è essenziale in tutti quei casi in cui il padre
crea un figlio e attende la sua conclusione per proseguire, ed entrambi
scrivono sullo stesso file. Un caso tipico di questo comportamento è la shell
quando lancia un programma.  In questo modo, anche se lo standard output viene
rediretto, il padre potrà sempre continuare a scrivere in coda a quanto
scritto dal figlio in maniera automatica; se così non fosse ottenere questo
comportamento sarebbe estremamente complesso, necessitando di una qualche
forma di comunicazione fra i due processi per far riprendere al padre la
scrittura al punto giusto.

In generale comunque non è buona norma far scrivere più processi sullo stesso
file senza una qualche forma di sincronizzazione in quanto, come visto anche
con il nostro esempio, le varie scritture risulteranno mescolate fra loro in
una sequenza impredicibile. Per questo le modalità con cui in genere si usano
i file dopo una \func{fork} sono sostanzialmente due:
\begin{enumerate*}
\item Il processo padre aspetta la conclusione del figlio. In questo caso non
  è necessaria nessuna azione riguardo ai file, in quanto la sincronizzazione
  della posizione corrente dopo eventuali operazioni di lettura e scrittura
  effettuate dal figlio è automatica.
\item L'esecuzione di padre e figlio procede indipendentemente. In questo caso
  ciascuno dei due processi deve chiudere i file che non gli servono una volta
  che la \func{fork} è stata eseguita, per evitare ogni forma di interferenza.
\end{enumerate*}

Oltre ai file aperti i processi figli ereditano dal padre una serie di altre
proprietà; la lista dettagliata delle proprietà che padre e figlio hanno in
comune dopo l'esecuzione di una \func{fork} è la seguente:
\begin{itemize*}
\item i file aperti e gli eventuali flag di \textit{close-on-exec} impostati
  (vedi sez.~\ref{sec:proc_exec} e sez.~\ref{sec:file_fcntl_ioctl});
\item gli identificatori per il controllo di accesso: l'\textsl{user-ID
    reale}, il \textsl{group-ID reale}, l'\textsl{user-ID effettivo}, il
  \textsl{group-ID effettivo} ed i \textsl{group-ID supplementari} (vedi
  sez.~\ref{sec:proc_access_id});
\item gli identificatori per il controllo di sessione: il \textit{process
    group-ID} e il \textit{session id} ed il terminale di controllo (vedi
  sez.~\ref{sec:sess_proc_group});
\item la directory di lavoro (vedi sez.~\ref{sec:file_work_dir}) e la
  directory radice (vedi sez.~\ref{sec:file_chroot});
\item la maschera dei permessi di creazione dei file (vedi
  sez.~\ref{sec:file_perm_management});
\item la maschera dei segnali bloccati (vedi
  sez.~\ref{sec:sig_sigmask}) e le azioni installate (vedi
  sez.~\ref{sec:sig_gen_beha});
\item i segmenti di memoria condivisa agganciati al processo (vedi
  sez.~\ref{sec:ipc_sysv_shm});
\item i limiti sulle risorse (vedi sez.~\ref{sec:sys_resource_limit});
\item il valori di \textit{nice}, le priorità \textit{real-time} e le affinità
  di processore (vedi sez.~\ref{sec:proc_sched_stand},
  sez.~\ref{sec:proc_real_time} e sez.~\ref{sec:proc_sched_multiprocess});
\item le variabili di ambiente (vedi sez.~\ref{sec:proc_environ}).
\item l'insieme dei descrittori associati alle code di messaggi POSIX (vedi
  sez.~\ref{sec:ipc_posix_mq}) che vengono copiate come i \textit{file
    descriptor}, questo significa che entrambi condivideranno gli stessi flag.
\end{itemize*}

Oltre a quelle relative ad un diverso spazio degli indirizzi (e una memoria
totalmente indipendente) le differenze fra padre e figlio dopo l'esecuzione di
una \func{fork} invece sono:\footnote{a parte le ultime quattro, relative a
  funzionalità specifiche di Linux, le altre sono esplicitamente menzionate
  dallo standard POSIX.1-2001.}
\begin{itemize*}
\item il valore di ritorno di \func{fork};
\item il \ids{PID} (\textit{process id}), quello del figlio viene assegnato ad
  un nuovo valore univoco;
\item il \ids{PPID} (\textit{parent process id}), quello del figlio viene
  impostato al \ids{PID} del padre;
\item i valori dei tempi di esecuzione (vedi sez.~\ref{sec:sys_cpu_times}) e
  delle risorse usate (vedi sez.~\ref{sec:sys_resource_use}), che nel figlio
  sono posti a zero;
\item i \textit{lock} sui file (vedi sez.~\ref{sec:file_locking}) e sulla
  memoria (vedi sez.~\ref{sec:proc_mem_lock}), che non vengono ereditati dal
  figlio;
\item gli allarmi, i timer (vedi sez.~\ref{sec:sig_alarm_abort}) ed i segnali
  pendenti (vedi sez.~\ref{sec:sig_gen_beha}), che per il figlio vengono
  cancellati.
\item le operazioni di I/O asincrono in corso (vedi
  sez.~\ref{sec:file_asyncronous_io}) che non vengono ereditate dal figlio;
\item gli aggiustamenti fatti dal padre ai semafori con \func{semop} (vedi
  sez.~\ref{sec:ipc_sysv_sem}).
\item le notifiche sui cambiamenti delle directory con \textit{dnotify} (vedi
  sez.~\ref{sec:sig_notification}), che non vengono ereditate dal figlio;
\item le mappature di memoria marcate come \const{MADV\_DONTFORK} (vedi
  sez.~\ref{sec:file_memory_map}) che non vengono ereditate dal figlio;
\item l'impostazione con \func{prctl} (vedi sez.~\ref{sec:process_prctl}) che
  notifica al figlio la terminazione del padre viene cancellata se presente
  nel padre;
\item il segnale di terminazione del figlio è sempre \signal{SIGCHLD} anche
  qualora nel padre fosse stato modificato (vedi sez.~\ref{sec:process_clone}). 
\end{itemize*}

Una seconda funzione storica usata per la creazione di un nuovo processo è
\funcm{vfork}, che è esattamente identica a \func{fork} ed ha la stessa
semantica e gli stessi errori; la sola differenza è che non viene creata la
tabella delle pagine né la struttura dei task per il nuovo processo. Il
processo padre è posto in attesa fintanto che il figlio non ha eseguito una
\func{execve} o non è uscito con una \func{\_exit}. Il figlio condivide la
memoria del padre (e modifiche possono avere effetti imprevedibili) e non deve
ritornare o uscire con \func{exit} ma usare esplicitamente \func{\_exit}.

Questa funzione è un rimasuglio dei vecchi tempi in cui eseguire una
\func{fork} comportava anche la copia completa del segmento dati del processo
padre, che costituiva un inutile appesantimento in tutti quei casi in cui la
\func{fork} veniva fatta solo per poi eseguire una \func{exec}. La funzione
venne introdotta in BSD per migliorare le prestazioni.

Dato che Linux supporta il \textit{copy on write} la perdita di prestazioni è
assolutamente trascurabile, e l'uso di questa funzione, che resta un caso
speciale della \textit{system call} \func{clone} (che tratteremo in dettaglio
in sez.~\ref{sec:process_clone}) è deprecato; per questo eviteremo di
trattarla ulteriormente.


\subsection{La conclusione di un processo}
\label{sec:proc_termination}

In sez.~\ref{sec:proc_conclusion} abbiamo già affrontato le modalità con cui
chiudere un programma, ma dall'interno del programma stesso. Avendo a che fare
con un sistema \textit{multitasking} resta da affrontare l'argomento dal punto
di vista di come il sistema gestisce la conclusione dei processi.

Abbiamo visto in sez.~\ref{sec:proc_conclusion} le tre modalità con cui un
programma viene terminato in maniera normale: la chiamata di \func{exit}, che
esegue le funzioni registrate per l'uscita e chiude gli \textit{stream} e poi
esegue \func{\_exit}, il ritorno dalla funzione \code{main} equivalente alla
chiamata di \func{exit}, e la chiamata diretta a \func{\_exit}, che passa
direttamente alle operazioni di terminazione del processo da parte del kernel.

Ma abbiamo accennato che oltre alla conclusione normale esistono anche delle
modalità di conclusione anomala. Queste sono in sostanza due: il programma può
chiamare la funzione \func{abort} (vedi sez.~\ref{sec:sig_alarm_abort}) per
invocare una chiusura anomala, o essere terminato da un segnale (torneremo sui
segnali in cap.~\ref{cha:signals}).  In realtà anche la prima modalità si
riconduce alla seconda, dato che \func{abort} si limita a generare il segnale
\signal{SIGABRT}.

Qualunque sia la modalità di conclusione di un processo, il kernel esegue
comunque una serie di operazioni di terminazione: chiude tutti i file aperti,
rilascia la memoria che stava usando, e così via; l'elenco completo delle
operazioni eseguite alla chiusura di un processo è il seguente:
\begin{itemize*}
\item tutti i \textit{file descriptor} (vedi sez.~\ref{sec:file_fd}) sono
  chiusi;
\item viene memorizzato lo stato di terminazione del processo;
\item ad ogni processo figlio viene assegnato un nuovo padre (in genere
  \cmd{init});
\item viene inviato il segnale \signal{SIGCHLD} al processo padre (vedi
  sez.~\ref{sec:sig_sigchld});
\item se il processo è un leader di sessione ed il suo terminale di controllo
  è quello della sessione viene mandato un segnale di \signal{SIGHUP} a tutti i
  processi del gruppo di \textit{foreground} e il terminale di controllo viene
  disconnesso (vedi sez.~\ref{sec:sess_ctrl_term});
\item se la conclusione di un processo rende orfano un \textit{process
    group} ciascun membro del gruppo viene bloccato, e poi gli vengono
  inviati in successione i segnali \signal{SIGHUP} e \signal{SIGCONT}
  (vedi ancora sez.~\ref{sec:sess_ctrl_term}).
\end{itemize*}

\itindbeg{termination~status} 

Oltre queste operazioni è però necessario poter disporre di un meccanismo
ulteriore che consenta di sapere come la terminazione è avvenuta: dato che in
un sistema unix-like tutto viene gestito attraverso i processi, il meccanismo
scelto consiste nel riportare lo \textsl{stato di terminazione} (il cosiddetto
\textit{termination status}) al processo padre.

Nel caso di conclusione normale, abbiamo visto in
sez.~\ref{sec:proc_conclusion} che lo stato di uscita del processo viene
caratterizzato tramite il valore del cosiddetto \textit{exit status}, cioè il
valore passato come argomento alle funzioni \func{exit} o \func{\_exit} o il
valore di ritorno per \code{main}.  Ma se il processo viene concluso in
maniera anomala il programma non può specificare nessun \textit{exit status},
ed è il kernel che deve generare autonomamente il \textit{termination status}
per indicare le ragioni della conclusione anomala.

Si noti la distinzione fra \textit{exit status} e \textit{termination status}:
quello che contraddistingue lo stato di chiusura del processo e viene
riportato attraverso le funzioni \func{wait} o \func{waitpid} (vedi
sez.~\ref{sec:proc_wait}) è sempre quest'ultimo; in caso di conclusione
normale il kernel usa il primo (nel codice eseguito da \func{\_exit}) per
produrre il secondo.

La scelta di riportare al padre lo stato di terminazione dei figli, pur
essendo l'unica possibile, comporta comunque alcune complicazioni: infatti se
alla sua creazione è scontato che ogni nuovo processo abbia un padre, non è
detto che sia così alla sua conclusione, dato che il padre potrebbe essere già
terminato; si potrebbe avere cioè quello che si chiama un processo
\textsl{orfano}.

Questa complicazione viene superata facendo in modo che il processo orfano
venga \textsl{adottato} da \cmd{init}, o meglio dal processo con \ids{PID}
1,\footnote{anche se, come vedremo in sez.~\ref{sec:process_prctl}, a partire
  dal kernel 3.4 è diventato possibile delegare questo compito anche ad un
  altro processo.} cioè quello lanciato direttamente dal kernel all'avvio, che
sta alla base dell'albero dei processi visto in sez.~\ref{sec:proc_hierarchy}
e che anche per questo motivo ha un ruolo essenziale nel sistema e non può mai
terminare.\footnote{almeno non senza un blocco completo del sistema, in caso
  di terminazione o di non esecuzione di \cmd{init} infatti il kernel si
  blocca con un cosiddetto \textit{kernel panic}, dato che questo è un errore
  fatale.}

Come già accennato quando un processo termina, il kernel controlla se è il
padre di altri processi in esecuzione: in caso positivo allora il \ids{PPID}
di tutti questi processi verrà sostituito dal kernel con il \ids{PID} di
\cmd{init}, cioè con 1. In questo modo ogni processo avrà sempre un padre (nel
caso possiamo parlare di un padre \textsl{adottivo}) cui riportare il suo
stato di terminazione.  

Come verifica di questo comportamento possiamo eseguire il nostro programma
\cmd{forktest} imponendo a ciascun processo figlio due secondi di attesa prima
di uscire, il risultato è:
\begin{Console}
[piccardi@selidor sources]$ \textbf{./forktest -c2 3}
Process 1972: forking 3 child
Spawned 1 child, pid 1973 
Child 1 successfully executing
Go to next child 
Spawned 2 child, pid 1974 
Child 2 successfully executing
Go to next child 
Child 3 successfully executing
Spawned 3 child, pid 1975 
Go to next child 

[piccardi@selidor sources]$ Child 3, parent 1, exiting
Child 2, parent 1, exiting
Child 1, parent 1, exiting
\end{Console}
come si può notare in questo caso il processo padre si conclude prima dei
figli, tornando alla shell, che stampa il prompt sul terminale: circa due
secondi dopo viene stampato a video anche l'output dei tre figli che
terminano, e come si può notare in questo caso, al contrario di quanto visto
in precedenza, essi riportano 1 come \ids{PPID}.

Altrettanto rilevante è il caso in cui il figlio termina prima del padre,
perché non è detto che il padre sia in esecuzione e possa ricevere
immediatamente lo stato di terminazione, quindi il kernel deve comunque
conservare una certa quantità di informazioni riguardo ai processi che sta
terminando.

Questo viene fatto mantenendo attiva la voce nella tabella dei processi, e
memorizzando alcuni dati essenziali, come il \ids{PID}, i tempi di CPU usati
dal processo (vedi sez.~\ref{sec:sys_unix_time}) e lo stato di terminazione,
mentre la memoria in uso ed i file aperti vengono rilasciati immediatamente. 

\itindbeg{zombie}

I processi che sono terminati, ma il cui stato di terminazione non è stato
ancora ricevuto dal padre sono chiamati \textit{zombie}, essi restano presenti
nella tabella dei processi ed in genere possono essere identificati
dall'output di \cmd{ps} per la presenza di una \texttt{Z} nella colonna che ne
indica lo stato (vedi tab.~\ref{tab:proc_proc_states}). Quando il padre
effettuerà la lettura dello stato di terminazione anche questa informazione,
non più necessaria, verrà scartata ed il processo potrà considerarsi
completamente concluso.

Possiamo utilizzare il nostro programma di prova per analizzare anche questa
condizione: lanciamo il comando \cmd{forktest} in \textit{background} (vedi
sez.~\ref{sec:sess_job_control}), indicando al processo padre di aspettare 10
secondi prima di uscire. In questo caso, usando \cmd{ps} sullo stesso
terminale (prima dello scadere dei 10 secondi) otterremo:
\begin{Console}
[piccardi@selidor sources]$ \textbf{ps T}
  PID TTY      STAT   TIME COMMAND
  419 pts/0    S      0:00 bash
  568 pts/0    S      0:00 ./forktest -e10 3
  569 pts/0    Z      0:00 [forktest <defunct>]
  570 pts/0    Z      0:00 [forktest <defunct>]
  571 pts/0    Z      0:00 [forktest <defunct>]
  572 pts/0    R      0:00 ps T
\end{Console}
%$
e come si vede, dato che non si è fatto nulla per riceverne lo stato di
terminazione, i tre processi figli sono ancora presenti pur essendosi
conclusi, con lo stato di \textit{zombie} e l'indicazione che sono terminati
(la scritta \texttt{defunct}).

La possibilità di avere degli \textit{zombie} deve essere tenuta sempre
presente quando si scrive un programma che deve essere mantenuto in esecuzione
a lungo e creare molti processi figli. In questo caso si deve sempre avere
cura di far leggere al programma l'eventuale stato di uscita di tutti i
figli. Una modalità comune di farlo è attraverso l'utilizzo di un apposito
\textit{signal handler} che chiami la funzione \func{wait}, (vedi
sez.~\ref{sec:proc_wait}), ne esamineremo in dettaglio un esempio
(fig.~\ref{fig:sig_sigchld_handl}) in sez.~\ref{sec:sig_sigchld}.

La lettura degli stati di uscita è necessaria perché anche se gli
\textit{zombie} non consumano risorse di memoria o processore, occupano
comunque una voce nella tabella dei processi e se li si lasciano accumulare a
lungo quest'ultima potrebbe esaurirsi, con la conseguente impossibilità di
lanciare nuovi processi.

Si noti tuttavia che quando un processo adottato da \cmd{init} termina, non
diviene mai uno \textit{zombie}. Questo perché una delle funzioni di
\cmd{init} è appunto quella di chiamare la funzione \func{wait} per i processi
a cui fa da padre, completandone la terminazione. Questo è quanto avviene
anche quando, come nel caso del precedente esempio con \cmd{forktest}, il
padre termina con dei figli in stato di \textit{zombie}. Questi scompaiono
quando, alla terminazione del padre dopo i secondi programmati, tutti figli
che avevamo generato, e che erano diventati \textit{zombie}, vengono adottati
da \cmd{init}, il quale provvede a completarne la terminazione.

Si tenga presente infine che siccome gli \textit{zombie} sono processi già
terminati, non c'è modo di eliminarli con il comando \cmd{kill} o inviandogli
un qualunque segnale di terminazione (l'argomento è trattato in
sez.~\ref{sec:sig_termination}). Qualora ci si trovi in questa situazione
l'unica possibilità di cancellarli dalla tabella dei processi è quella di
terminare il processo che li ha generati e che non sta facendo il suo lavoro,
in modo che \cmd{init} possa adottarli e concluderne correttamente la
terminazione.

Si tenga anche presente che la presenza di \textit{zombie} nella tabella dei
processi non è sempre indice di un qualche malfunzionamento, in una macchina
con molto carico infatti può esservi una presenza temporanea dovuta al fatto
che il processo padre ancora non ha avuto il tempo di gestirli. 

\itindend{zombie}


\subsection{Le funzioni di attesa e ricezione degli stati di uscita}
\label{sec:proc_wait}

Uno degli usi più comuni delle capacità \textit{multitasking} di un sistema
unix-like consiste nella creazione di programmi di tipo server, in cui un
processo principale attende le richieste che vengono poi soddisfatte da una
serie di processi figli.

Si è già sottolineato al paragrafo precedente come in questo caso diventi
necessario gestire esplicitamente la conclusione dei figli onde evitare di
riempire di \textit{zombie} la tabella dei processi. Tratteremo in questa
sezione le funzioni di sistema deputate a questo compito; la prima è
\funcd{wait} ed il suo prototipo è:

\begin{funcproto}{ 
\fhead{sys/types.h}
\fhead{sys/wait.h}
\fdecl{pid\_t wait(int *status)}
\fdesc{Attende la terminazione di un processo.} 
}
{La funzione ritorna il \ids{PID} del figlio in caso di successo e $-1$ per un
  errore, nel qual caso \var{errno} assumerà uno dei valori:
  \begin{errlist}
  \item[\errcode{ECHILD}] il processo non ha nessun figlio di cui attendere
    uno stato di terminazione.
  \item[\errcode{EINTR}] la funzione è stata interrotta da un segnale.
  \end{errlist}}
\end{funcproto}

Questa funzione è presente fin dalle prime versioni di Unix ed è quella usata
tradizionalmente per attendere la terminazione dei figli. La funzione sospende
l'esecuzione del processo corrente e ritorna non appena un qualunque processo
figlio termina. Se un figlio è già terminato prima della sua chiamata la
funzione ritorna immediatamente, se più processi figli sono già terminati
occorrerà continuare a chiamare la funzione più volte fintanto che non si è
recuperato lo stato di terminazione di tutti quanti.

Al ritorno della funzione lo stato di terminazione del figlio viene salvato
(come \textit{value result argument}) nella variabile puntata
da \param{status} e tutte le risorse del kernel relative al processo (vedi
sez.~\ref{sec:proc_termination}) vengono rilasciate.  Nel caso un processo
abbia più figli il valore di ritorno della funzione sarà impostato al
\ids{PID} del processo di cui si è ricevuto lo stato di terminazione, cosa che
permette di identificare qual è il figlio che è terminato.

\itindend{termination~status}

Questa funzione ha il difetto di essere poco flessibile, in quanto ritorna
all'uscita di un qualunque processo figlio. Nelle occasioni in cui è
necessario attendere la conclusione di uno specifico processo fra tutti quelli
esistenti occorre predisporre un meccanismo che tenga conto di tutti processi
che sono terminati, e provveda a ripetere la chiamata alla funzione nel caso
il processo cercato non risulti fra questi. Se infatti il processo cercato è
già terminato e se è già ricevuto lo stato di uscita senza registrarlo, la
funzione non ha modo di accorgersene, e si continuerà a chiamarla senza
accorgersi che quanto interessava è già accaduto.

Per questo motivo lo standard POSIX.1 ha introdotto una seconda funzione che
effettua lo stesso servizio, ma dispone di una serie di funzionalità più
ampie, legate anche al controllo di sessione (si veda
sez.~\ref{sec:sess_job_control}).  Dato che è possibile ottenere lo stesso
comportamento di \func{wait}\footnote{in effetti il codice
  \code{wait(\&status)} è del tutto equivalente a \code{waitpid(WAIT\_ANY,
    \&status, 0)}.} si consiglia di utilizzare sempre questa nuova funzione di
sistema, \funcd{waitpid}, il cui prototipo è:

\begin{funcproto}{ 
\fhead{sys/types.h}
\fhead{sys/wait.h}
\fdecl{pid\_t waitpid(pid\_t pid, int *status, int options)}
\fdesc{Attende il cambiamento di stato di un processo figlio.} 
}
{La funzione ritorna il \ids{PID} del processo che ha cambiato stato in caso
  di successo, o 0 se è stata specificata l'opzione \const{WNOHANG} e il
  processo non è uscito e $-1$ per un errore, nel qual caso \var{errno}
  assumerà uno dei valori:
  \begin{errlist}
  \item[\errcode{ECHILD}] il processo specificato da \param{pid} non esiste o
    non è figlio del processo chiamante.
  \item[\errcode{EINTR}] non è stata specificata l'opzione \const{WNOHANG} e
    la funzione è stata interrotta da un segnale.
  \item[\errcode{EINVAL}] si è specificato un valore non valido per
    l'argomento \param{options}.
  \end{errlist}}
\end{funcproto}

La prima differenza fra le due funzioni è che con \func{waitpid} si può
specificare in maniera flessibile quale processo attendere, sulla base del
valore fornito dall'argomento \param{pid}. Questo può assumere diversi valori,
secondo lo specchietto riportato in tab.~\ref{tab:proc_waidpid_pid}, dove si
sono riportate anche le costanti definite per indicare alcuni di essi.

\begin{table}[!htb]
  \centering
  \footnotesize
  \begin{tabular}[c]{|c|c|p{8cm}|}
    \hline
    \textbf{Valore} & \textbf{Costante} &\textbf{Significato}\\
    \hline
    \hline
    $<-1$& --               & Attende per un figlio il cui \textit{process
                              group} (vedi sez.~\ref{sec:sess_proc_group}) è
                              uguale al valore assoluto di \param{pid}.\\ 
    $-1$&\constd{WAIT\_ANY}  & Attende per un figlio qualsiasi, usata in
                               questa maniera senza specificare nessuna opzione
                               è equivalente a \func{wait}.\\ 
    $ 0$&\constd{WAIT\_MYPGRP}&Attende per un figlio il cui \textit{process
                               group} (vedi sez.~\ref{sec:sess_proc_group}) è
                                uguale a quello del processo chiamante.\\ 
    $>0$& --                & Attende per un figlio il cui \ids{PID} è uguale
                              al valore di \param{pid}.\\
    \hline
  \end{tabular}
  \caption{Significato dei valori dell'argomento \param{pid} della funzione
    \func{waitpid}.}
  \label{tab:proc_waidpid_pid}
\end{table}

Il comportamento di \func{waitpid} può inoltre essere modificato passando alla
funzione delle opportune opzioni tramite l'argomento \param{options}; questo
deve essere specificato come maschera binaria delle costanti riportati nella
prima parte in tab.~\ref{tab:proc_waitpid_options} che possono essere
combinate fra loro con un OR aritmetico. Nella seconda parte della stessa
tabella si sono riportati anche alcune opzioni non standard specifiche di
Linux, che consentono un controllo più dettagliato per i processi creati con
la \textit{system call} generica \func{clone} (vedi
sez.~\ref{sec:process_clone}) e che vengono usati principalmente per la
gestione della terminazione dei \textit{thread}\unavref{ (vedi
sez.~\ref{sec:thread_xxx})}.

\begin{table}[!htb]
  \centering
  \footnotesize
  \begin{tabular}[c]{|l|p{8cm}|}
    \hline
    \textbf{Costante} & \textbf{Descrizione}\\
    \hline
    \hline
    \const{WNOHANG}   & La funzione ritorna immediatamente anche se non è
                        terminato nessun processo figlio.\\
    \const{WUNTRACED} & Ritorna anche quando un processo figlio è stato
                        fermato.\\ 
    \const{WCONTINUED}& Ritorna anche quando un processo figlio che era stato
                        fermato ha ripreso l'esecuzione (dal kernel 2.6.10).\\
    \hline
    \constd{\_\_WCLONE}& Attende solo per i figli creati con \func{clone} 
                        (vedi sez.~\ref{sec:process_clone}), vale a dire
                        processi che non emettono nessun segnale 
                        o emettono un segnale diverso da \signal{SIGCHLD} alla
                        terminazione, il default è attendere soltanto i
                        processi figli ordinari ignorando quelli creati da
                        \func{clone}.\\
    \constd{\_\_WALL}  & Attende per qualunque figlio, sia ordinario che creato
                        con \func{clone}, se specificata con
                        \const{\_\_WCLONE} quest'ultima viene ignorata. \\
    \constd{\_\_WNOTHREAD}& Non attende per i figli di altri \textit{thread}
                        dello stesso \textit{thread group}, questo era il
                        comportamento di default del kernel 2.4 che non
                        supportava la possibilità, divenuta il default a
                        partire dal 2.6, di attendere un qualunque figlio
                        appartenente allo stesso \textit{thread group}. \\
    \hline
  \end{tabular}
  \caption{Costanti che identificano i bit dell'argomento \param{options}
    della funzione \func{waitpid}.} 
  \label{tab:proc_waitpid_options}
\end{table}

\constbeg{WNOHANG}

L'uso dell'opzione \const{WNOHANG} consente di prevenire il blocco della
funzione qualora nessun figlio sia uscito o non si siano verificate le altre
condizioni per l'uscita della funzione. in tal caso. In tal caso la funzione,
invece di restituire il \ids{PID} del processo (che è sempre un intero
positivo) ritornerà un valore nullo.

\constend{WNOHANG}
\constbeg{WUNTRACED}
\constbeg{WCONTINUED}

Le altre due opzioni, \const{WUNTRACED} e \const{WCONTINUED}, consentono
rispettivamente di tracciare non la terminazione di un processo, ma il fatto
che esso sia stato fermato, o fatto ripartire, e sono utilizzate per la
gestione del controllo di sessione (vedi sez.~\ref{sec:sess_job_control}).

Nel caso di \const{WUNTRACED} la funzione ritorna, restituendone il \ids{PID},
quando un processo figlio entra nello stato \textit{stopped}\footnote{in
  realtà viene notificato soltanto il caso in cui il processo è stato fermato
  da un segnale di stop (vedi sez.~\ref{sec:sess_ctrl_term}), e non quello in
  cui lo stato \textit{stopped} è dovuto all'uso di \func{ptrace}\unavref{
    (vedi sez.~\ref{sec:process_ptrace})}.} (vedi
tab.~\ref{tab:proc_proc_states}), mentre con \const{WCONTINUED} la funzione
ritorna quando un processo in stato \textit{stopped} riprende l'esecuzione per
la ricezione del segnale \signal{SIGCONT} (l'uso di questi segnali per il
controllo di sessione è trattato in sez.~\ref{sec:sess_ctrl_term}).

\constend{WUNTRACED}
\constend{WCONTINUED}

La terminazione di un processo figlio (così come gli altri eventi osservabili
con \func{waitpid}) è chiaramente un evento asincrono rispetto all'esecuzione
di un programma e può avvenire in un qualunque momento. Per questo motivo,
come accennato nella sezione precedente, una delle azioni prese dal kernel
alla conclusione di un processo è quella di mandare un segnale di
\signal{SIGCHLD} al padre. L'azione predefinita per questo segnale (si veda
sez.~\ref{sec:sig_base}) è di essere ignorato, ma la sua generazione
costituisce il meccanismo di comunicazione asincrona con cui il kernel avverte
il processo padre che uno dei suoi figli è terminato.

Il comportamento delle funzioni è però cambiato nel passaggio dal kernel 2.4
al kernel 2.6, quest'ultimo infatti si è adeguato alle prescrizioni dello
standard POSIX.1-2001 e come da esso richiesto se \signal{SIGCHLD} viene
ignorato, o se si imposta il flag di \const{SA\_NOCLDSTOP} nella ricezione
dello stesso (si veda sez.~\ref{sec:sig_sigaction}) i processi figli che
terminano non diventano \textit{zombie} e sia \func{wait} che \func{waitpid}
si bloccano fintanto che tutti i processi figli non sono terminati, dopo di
che falliscono con un errore di \errcode{ENOCHLD}.\footnote{questo è anche il
  motivo per cui le opzioni \const{WUNTRACED} e \const{WCONTINUED} sono
  utilizzabili soltanto qualora non si sia impostato il flag di
  \const{SA\_NOCLDSTOP} per il segnale \signal{SIGCHLD}.}

Con i kernel della serie 2.4 e tutti i kernel delle serie precedenti entrambe
le funzioni di attesa ignorano questa prescrizione e si comportano sempre
nello stesso modo,\footnote{lo standard POSIX.1 originale infatti lascia
  indefinito il comportamento di queste funzioni quando \signal{SIGCHLD} viene
  ignorato.} indipendentemente dal fatto \signal{SIGCHLD} sia ignorato o meno:
attendono la terminazione di un processo figlio e ritornano il relativo
\ids{PID} e lo stato di terminazione nell'argomento \param{status}.

In generale in un programma non si vuole essere forzati ad attendere la
conclusione di un processo figlio per proseguire l'esecuzione, specie se tutto
questo serve solo per leggerne lo stato di chiusura (ed evitare eventualmente
la presenza di \textit{zombie}).  Per questo la modalità più comune di
chiamare queste funzioni è quella di utilizzarle all'interno di un
\textit{signal handler} (vedremo un esempio di come gestire \signal{SIGCHLD}
con i segnali in sez.~\ref{sec:sig_example}). In questo caso infatti, dato che
il segnale è generato dalla terminazione di un figlio, avremo la certezza che
la chiamata a \func{waitpid} non si bloccherà.

Come accennato sia \func{wait} che \func{waitpid} restituiscono lo stato di
terminazione del processo tramite il puntatore \param{status}, e se non
interessa memorizzarlo si può passare un puntatore nullo. Il valore restituito
da entrambe le funzioni dipende dall'implementazione, ma tradizionalmente gli
8 bit meno significativi sono riservati per memorizzare lo stato di uscita del
processo, e gli altri per indicare il segnale che ha causato la terminazione
(in caso di conclusione anomala), uno per indicare se è stato generato un
\textit{core dump} (vedi sez.~\ref{sec:sig_standard}), ecc.\footnote{le
  definizioni esatte si possono trovare in \file{<bits/waitstatus.h>} ma
  questo file non deve mai essere usato direttamente, esso viene incluso
  attraverso \file{<sys/wait.h>}.}

\begin{table}[!htb]
  \centering
  \footnotesize
  \begin{tabular}[c]{|l|p{10cm}|}
    \hline
    \textbf{Macro} & \textbf{Descrizione}\\
    \hline
    \hline
    \macrod{WIFEXITED}\texttt{(s)}   & Condizione vera (valore non nullo) per
                                      un processo figlio che sia terminato
                                      normalmente. \\ 
    \macrod{WEXITSTATUS}\texttt{(s)} & Restituisce gli otto bit meno
                                      significativi dello stato di uscita del
                                      processo (passato attraverso
                                      \func{\_exit}, \func{exit} o come valore
                                      di ritorno di \code{main}); può essere
                                      valutata solo se \val{WIFEXITED} ha
                                      restituito un valore non nullo.\\ 
    \macrod{WIFSIGNALED}\texttt{(s)} & Condizione vera se il processo figlio è
                                      terminato in maniera anomala a causa di
                                      un segnale che non è stato catturato
                                      (vedi sez.~\ref{sec:sig_notification}).\\ 
    \macrod{WTERMSIG}\texttt{(s)}    & Restituisce il numero del segnale che ha
                                      causato la terminazione anomala del
                                      processo; può essere valutata solo se
                                      \val{WIFSIGNALED} ha restituito un
                                      valore non nullo.\\
    \macrod{WCOREDUMP}\texttt{(s)}   & Vera se il processo terminato ha
                                      generato un file di 
                                      \textit{core dump}; può essere valutata
                                      solo se \val{WIFSIGNALED} ha restituito
                                      un valore non nullo.\footnotemark \\
    \macrod{WIFSTOPPED}\texttt{(s)}  & Vera se il processo che ha causato il
                                      ritorno di \func{waitpid} è bloccato;
                                      l'uso è possibile solo con
                                      \func{waitpid} avendo specificato
                                      l'opzione \const{WUNTRACED}.\\
    \macrod{WSTOPSIG}\texttt{(s)}    & Restituisce il numero del segnale che ha
                                      bloccato il processo; può essere
                                      valutata solo se \val{WIFSTOPPED} ha
                                      restituito un valore non nullo. \\ 
    \macrod{WIFCONTINUED}\texttt{(s)}& Vera se il processo che ha causato il
                                      ritorno è stato riavviato da un
                                      \signal{SIGCONT} (disponibile solo a
                                      partire dal kernel 2.6.10).\\
    \hline
  \end{tabular}
  \caption{Descrizione delle varie macro di preprocessore utilizzabili per 
    verificare lo stato di terminazione \var{s} di un processo.}
  \label{tab:proc_status_macro}
\end{table}

\footnotetext{questa macro non è definita dallo standard POSIX.1-2001, ma è
  presente come estensione sia in Linux che in altri Unix, deve essere
  pertanto utilizzata con attenzione (ad esempio è il caso di usarla in un
  blocco \texttt{\#ifdef WCOREDUMP ... \#endif}.}

Lo standard POSIX.1 definisce una serie di macro di preprocessore da usare per
analizzare lo stato di uscita. Esse sono definite sempre in
\file{<sys/wait.h>} ed elencate in tab.~\ref{tab:proc_status_macro}. Si tenga
presente che queste macro prevedono che gli si passi come parametro la
variabile di tipo \ctyp{int} puntata dall'argomento \param{status} restituito
da \func{wait} o \func{waitpid}.

Si tenga conto che nel caso di conclusione anomala il valore restituito da
\val{WTERMSIG} può essere confrontato con le costanti che identificano i
segnali definite in \headfile{signal.h} ed elencate in
tab.~\ref{tab:sig_signal_list}, e stampato usando le apposite funzioni
trattate in sez.~\ref{sec:sig_strsignal}.

A partire dal kernel 2.6.9, sempre in conformità allo standard POSIX.1-2001, è
stata introdotta una nuova funzione di attesa che consente di avere un
controllo molto più preciso sui possibili cambiamenti di stato dei processi
figli e più dettagli sullo stato di uscita; la funzione di sistema è
\funcd{waitid} ed il suo prototipo è:

\begin{funcproto}{ 
\fhead{sys/types.h}
\fhead{sys/wait.h}
\fdecl{int waitid(idtype\_t idtype, id\_t id, siginfo\_t *infop, int options)}
\fdesc{Attende il cambiamento di stato di un processo figlio.} 
}
{La funzione ritorna $0$ in caso di successo e $-1$ per un errore, nel qual
  caso \var{errno} assumerà uno dei valori:
  \begin{errlist}
  \item[\errcode{ECHILD}] il processo specificato da \param{pid} non esiste o
    non è figlio del processo chiamante.
  \item[\errcode{EINTR}] non è stata specificata l'opzione \const{WNOHANG} e
    la funzione è stata interrotta da un segnale.
  \item[\errcode{EINVAL}] si è specificato un valore non valido per
    l'argomento \param{options}.
  \end{errlist}}
\end{funcproto}

La funzione prevede che si specifichi quali processi si intendono osservare
usando i due argomenti \param{idtype} ed \param{id}; il primo indica se ci si
vuole porre in attesa su un singolo processo, un gruppo di processi o un
processo qualsiasi, e deve essere specificato secondo uno dei valori di
tab.~\ref{tab:proc_waitid_idtype}; il secondo indica, a seconda del valore del
primo, quale processo o quale gruppo di processi selezionare.

\begin{table}[!htb]
  \centering
  \footnotesize
  \begin{tabular}[c]{|l|p{8cm}|}
    \hline
    \textbf{Valore} & \textbf{Descrizione}\\
    \hline
    \hline
    \constd{P\_PID} & Indica la richiesta di attendere per un processo figlio
                      il cui \ids{PID} corrisponda al valore dell'argomento
                      \param{id}.\\
    \constd{P\_PGID}& Indica la richiesta di attendere per un processo figlio
                      appartenente al \textit{process group} (vedi
                      sez.~\ref{sec:sess_proc_group}) il cui \acr{pgid}
                      corrisponda al valore dell'argomento \param{id}.\\
    \constd{P\_ALL} & Indica la richiesta di attendere per un processo figlio
                      generico, il valore dell'argomento \param{id} viene
                      ignorato.\\
    \hline
  \end{tabular}
  \caption{Costanti per i valori dell'argomento \param{idtype} della funzione
    \func{waitid}.}
  \label{tab:proc_waitid_idtype}
\end{table}

Come per \func{waitpid} anche il comportamento di \func{waitid} è
controllato dall'argomento \param{options}, da specificare come maschera
binaria dei valori riportati in tab.~\ref{tab:proc_waitid_options}. Benché
alcuni di questi siano identici come significato ed effetto ai precedenti di
tab.~\ref{tab:proc_waitpid_options}, ci sono delle differenze significative:
in questo caso si dovrà specificare esplicitamente l'attesa della terminazione
di un processo impostando l'opzione \const{WEXITED}, mentre il precedente
\const{WUNTRACED} è sostituito da \const{WSTOPPED}.  Infine è stata aggiunta
l'opzione \const{WNOWAIT} che consente una lettura dello stato mantenendo il
processo in attesa di ricezione, così che una successiva chiamata possa di
nuovo riceverne lo stato.

\begin{table}[!htb]
  \centering
  \footnotesize
  \begin{tabular}[c]{|l|p{8cm}|}
    \hline
    \textbf{Valore} & \textbf{Descrizione}\\
    \hline
    \hline
    \constd{WEXITED}   & Ritorna quando un processo figlio è terminato.\\
    \constd{WNOHANG}   & Ritorna immediatamente anche se non c'è niente da
                         notificare.\\ 
    \constd{WSTOPPED}  & Ritorna quando un processo figlio è stato fermato.\\
    \const{WCONTINUED} & Ritorna quando un processo figlio che era stato
                         fermato ha ripreso l'esecuzione.\\
    \constd{WNOWAIT}   & Lascia il processo ancora in attesa di ricezione, così
                         che una successiva chiamata possa di nuovo riceverne
                         lo stato.\\
    \hline
  \end{tabular}
  \caption{Costanti che identificano i bit dell'argomento \param{options}
    della funzione \func{waitid}.} 
  \label{tab:proc_waitid_options}
\end{table}

La funzione \func{waitid} restituisce un valore nullo in caso di successo, e
$-1$ in caso di errore; viene restituito un valore nullo anche se è stata
specificata l'opzione \const{WNOHANG} e la funzione è ritornata immediatamente
senza che nessun figlio sia terminato. Pertanto per verificare il motivo del
ritorno della funzione occorre analizzare le informazioni che essa
restituisce; queste, al contrario delle precedenti \func{wait} e
\func{waitpid} che usavano un semplice valore numerico, sono ritornate in una
struttura di tipo \struct{siginfo\_t} (vedi fig.~\ref{fig:sig_siginfo_t})
all'indirizzo puntato dall'argomento \param{infop}.

Tratteremo nei dettagli la struttura \struct{siginfo\_t} ed il significato dei
suoi vari campi in sez.~\ref{sec:sig_sigaction}, per quanto ci interessa qui
basta dire che al ritorno di \func{waitid} verranno avvalorati i seguenti
campi:
\begin{basedescript}{\desclabelwidth{1.8cm}}
\item[\var{si\_pid}] con il \ids{PID} del figlio.
\item[\var{si\_uid}] con l'\textsl{user-ID reale} (vedi
  sez.~\ref{sec:proc_perms}) del figlio.
\item[\var{si\_signo}] con \signal{SIGCHLD}.
\item[\var{si\_status}] con lo stato di uscita del figlio o con il segnale che
  lo ha terminato, fermato o riavviato.
\item[\var{si\_code}] con uno fra \const{CLD\_EXITED}, \const{CLD\_KILLED},
  \const{CLD\_STOPPED}, \const{CLD\_CONTINUED}, \const{CLD\_TRAPPED} e
  \const{CLD\_DUMPED} a indicare la ragione del ritorno della funzione, il cui
  significato è, nell'ordine: uscita normale, terminazione da segnale,
  processo fermato, processo riavviato, processo terminato in
  \textit{core dump} (vedi sez.~\ref{sec:sig_standard}).
\end{basedescript}

Infine Linux, seguendo un'estensione di BSD, supporta altre due funzioni per
la lettura dello stato di terminazione di un processo, analoghe alle
precedenti ma che prevedono un ulteriore argomento attraverso il quale il
kernel può restituire al padre informazioni sulle risorse (vedi
sez.~\ref{sec:sys_res_limits}) usate dal processo terminato e dai vari figli.
Le due funzioni di sistema sono \funcd{wait3} e \funcd{wait4}, che diventano
accessibili definendo la macro \macro{\_USE\_BSD}, i loro prototipi sono:

\begin{funcproto}{ 
\fhead{sys/types.h}
\fhead{sys/times.h}
\fhead{sys/resource.h}
\fhead{sys/wait.h}
\fdecl{int wait3(int *status, int options, struct rusage *rusage)}
\fdecl{int wait4(pid\_t pid, int *status, int options, struct rusage *rusage)}
\fdesc{Attende il cambiamento di stato di un processo figlio, riportando l'uso
  delle risorse.} 
}
{La funzione ha gli stessi valori di ritorno e codici di errore di
  \func{waitpid}. }
\end{funcproto}

La funzione \func{wait4} è identica \func{waitpid} sia nel comportamento che
per i valori dei primi tre argomenti, ma in più restituisce nell'argomento
aggiuntivo \param{rusage} un sommario delle risorse usate dal processo. Questo
argomento è una struttura di tipo \struct{rusage} definita in
\headfile{sys/resource.h}, che viene utilizzata anche dalla funzione
\func{getrusage} per ottenere le risorse di sistema usate da un processo. La
sua definizione è riportata in fig.~\ref{fig:sys_rusage_struct} e ne
tratteremo in dettaglio il significato sez.~\ref{sec:sys_resource_use}. La
funzione \func{wait3} è semplicemente un caso particolare di (e con Linux
viene realizzata con la stessa \textit{system call}), ed è equivalente a
chiamare \code{wait4(-1, \&status, opt, rusage)}, per questo motivo è ormai
deprecata in favore di \func{wait4}.



\subsection{La famiglia delle funzioni \func{exec} per l'esecuzione dei
  programmi}
\label{sec:proc_exec}

Abbiamo già detto che una delle modalità principali con cui si utilizzano i
processi in Unix è quella di usarli per lanciare nuovi programmi: questo viene
fatto attraverso una delle funzioni della famiglia \func{exec}. Quando un
processo chiama una di queste funzioni esso viene completamente sostituito dal
nuovo programma, il \ids{PID} del processo non cambia, dato che non viene
creato un nuovo processo, la funzione semplicemente rimpiazza lo
\textit{stack}, i dati ed il testo del processo corrente con un nuovo
programma letto da disco, eseguendo il \textit{link-loader} con gli effetti
illustrati in sez.~\ref{sec:proc_main}.

\begin{figure}[!htb]
  \centering \includegraphics[width=8cm]{img/exec_rel}
  \caption{La interrelazione fra le sei funzioni della famiglia \func{exec}.}
  \label{fig:proc_exec_relat}
\end{figure}

Ci sono sei diverse versioni di \func{exec} (per questo la si è chiamata
famiglia di funzioni) che possono essere usate per questo compito, in realtà
(come mostrato in fig.~\ref{fig:proc_exec_relat}), tutte queste funzioni sono
tutte varianti che consentono di invocare in modi diversi, semplificando il
passaggio degli argomenti, la funzione di sistema \funcd{execve}, il cui
prototipo è:

\begin{funcproto}{
\fhead{unistd.h}
\fdecl{int execve(const char *filename, char *const argv[], char *const envp[])}
\fdesc{Esegue un programma.} 
}
{La funzione ritorna solo in caso di errore, restituendo $-1$, nel qual
 caso \var{errno} assumerà uno dei valori:
\begin{errlist}
  \item[\errcode{EACCES}] il file o l'interprete non file ordinari, o non sono
    eseguibili, o il file è su un filesystem montato con l'opzione
    \cmd{noexec}, o manca  il permesso di attraversamento di una delle
    directory del \textit{pathname}.
  \item[\errcode{EAGAIN}] dopo un cambio di \ids{UID} si è ancora  sopra il
    numero massimo di processi consentiti per l'utente (dal kernel 3.1, per i
    dettagli vedi sez.~\ref{sec:proc_setuid}).
  \item[\errcode{EINVAL}] l'eseguibile ELF ha più di un segmento
    \const{PT\_INTERP}, cioè chiede di essere eseguito da più di un
    interprete.
  \item[\errcode{ELIBBAD}] un interprete ELF non è in un formato
    riconoscibile.
  \item[\errcode{ENOENT}] il file o una delle librerie dinamiche o l'interprete
    necessari per eseguirlo non esistono.
  \item[\errcode{ENOEXEC}] il file è in un formato non eseguibile o non
    riconosciuto come tale, o compilato per un'altra architettura.
  \item[\errcode{EPERM}] il file ha i bit \acr{suid} o \acr{sgid} e l'utente
    non è root, ed il processo viene tracciato, oppure il filesystem è montato
    con l'opzione \cmd{nosuid}. 
  \item[\errcode{ETXTBSY}] l'eseguibile è aperto in scrittura da uno o più
    processi. 
  \item[\errcode{E2BIG}] la lista degli argomenti è troppo grande.
  \end{errlist}
  ed inoltre \errval{EFAULT}, \errval{EIO}, \errval{EISDIR}, \errval{ELOOP},
  \errval{EMFILE}, \errval{ENAMETOOLONG}, \errval{ENFILE}, \errval{ENOMEM},
  \errval{ENOTDIR} nel loro significato generico.  }
\end{funcproto}

La funzione \func{execve} esegue il programma o lo script indicato dal
\textit{pathname} \param{filename}, passandogli la lista di argomenti indicata
da \param{argv} e come ambiente la lista di stringhe indicata
da \param{envp}. Entrambe le liste devono essere terminate da un puntatore
nullo. I vettori degli argomenti e dell'ambiente possono essere acceduti dal
nuovo programma quando la sua funzione \code{main} è dichiarata nella forma
\code{main(int argc, char *argv[], char *envp[])}. Si tenga presente per il
passaggio degli argomenti e dell'ambiente esistono comunque dei limiti, su cui
torneremo in sez.~\ref{sec:sys_res_limits}).
% TODO aggiungere la parte sul numero massimo di argomenti, da man execve

In caso di successo la funzione non ritorna, in quanto al posto del programma
chiamante viene eseguito il nuovo programma indicato da \param{filename}. Se
il processo corrente è tracciato con \func{ptrace}\unavref{ (vedi
sez.~\ref{sec:process_ptrace})} in caso di successo viene emesso il segnale
\signal{SIGTRAP}.

Le altre funzioni della famiglia (\funcd{execl}, \funcd{execv},
\funcd{execle}, \funcd{execlp}, \funcd{execvp}) servono per fornire all'utente
una serie di possibili diverse interfacce nelle modalità di passaggio degli
argomenti all'esecuzione del nuovo programma. I loro prototipi sono:

\begin{funcproto}{ 
\fhead{unistd.h}
\fdecl{int execl(const char *path, const char *arg, ...)}
\fdecl{int execv(const char *path, char *const argv[])}
\fdecl{int execle(const char *path, const char *arg, ..., char * const envp[])}
\fdecl{int execlp(const char *file, const char *arg, ...)}
\fdecl{int execvp(const char *file, char *const argv[])}
\fdesc{Eseguono un programma.} 
}
{Le funzioni ritornano solo in caso di errore, restituendo $-1$, i codici di
  errore sono gli stessi di \func{execve}.
}
\end{funcproto}

Tutte le funzioni mettono in esecuzione nel processo corrente il programma
indicati nel primo argomento. Gli argomenti successivi consentono di
specificare gli argomenti e l'ambiente che saranno ricevuti dal nuovo
processo. Per capire meglio le differenze fra le funzioni della famiglia si può
fare riferimento allo specchietto riportato in
tab.~\ref{tab:proc_exec_scheme}. La relazione fra le funzioni è invece
illustrata in fig.~\ref{fig:proc_exec_relat}.

\begin{table}[!htb]
  \footnotesize
  \centering
  \begin{tabular}[c]{|l|c|c|c||c|c|c|}
    \hline
    \multicolumn{1}{|c|}{\textbf{Caratteristiche}} & 
    \multicolumn{6}{|c|}{\textbf{Funzioni}} \\
    \hline
    &\func{execl}\texttt{ }&\func{execlp}&\func{execle}
    &\func{execv}\texttt{ }& \func{execvp}& \func{execve} \\
    \hline
    \hline
    argomenti a lista    &$\bullet$&$\bullet$&$\bullet$&&& \\
    argomenti a vettore  &&&&$\bullet$&$\bullet$&$\bullet$\\
    \hline
    filename completo     &$\bullet$&&$\bullet$&$\bullet$&&$\bullet$\\ 
    ricerca su \var{PATH} &&$\bullet$&&&$\bullet$& \\
    \hline
    ambiente a vettore   &&&$\bullet$&&&$\bullet$ \\
    uso di \var{environ} &$\bullet$&$\bullet$&&$\bullet$&$\bullet$& \\
    \hline
  \end{tabular}
  \caption{Confronto delle caratteristiche delle varie funzioni della 
    famiglia \func{exec}.}
  \label{tab:proc_exec_scheme}
\end{table}

La prima differenza fra le funzioni riguarda le modalità di passaggio dei
valori che poi andranno a costituire gli argomenti a linea di comando (cioè i
valori di \param{argv} e \param{argc} visti dalla funzione \code{main} del
programma chiamato). Queste modalità sono due e sono riassunte dagli mnemonici
``\texttt{v}'' e ``\texttt{l}'' che stanno rispettivamente per \textit{vector}
e \textit{list}.

Nel primo caso gli argomenti sono passati tramite il vettore di puntatori
\var{argv[]} a stringhe terminate con zero che costituiranno gli argomenti a
riga di comando, questo vettore \emph{deve} essere terminato da un puntatore
nullo. Nel secondo caso le stringhe degli argomenti sono passate alla funzione
come lista di puntatori, nella forma:
\includecodesnip{listati/char_list.c}
che deve essere terminata da un puntatore nullo.  In entrambi i casi vale la
convenzione che il primo argomento (\var{arg0} o \var{argv[0]}) viene usato
per indicare il nome del file che contiene il programma che verrà eseguito.

La seconda differenza fra le funzioni riguarda le modalità con cui si
specifica il programma che si vuole eseguire. Con lo mnemonico ``\texttt{p}''
si indicano le due funzioni che replicano il comportamento della shell nello
specificare il comando da eseguire; quando l'argomento \param{file} non
contiene una ``\texttt{/}'' esso viene considerato come un nome di programma,
e viene eseguita automaticamente una ricerca fra i file presenti nella lista
di directory specificate dalla variabile di ambiente \envvar{PATH}. Il file
che viene posto in esecuzione è il primo che viene trovato. Se si ha un errore
relativo a permessi di accesso insufficienti (cioè l'esecuzione della
sottostante \func{execve} ritorna un \errcode{EACCES}), la ricerca viene
proseguita nelle eventuali ulteriori directory indicate in \envvar{PATH}; solo
se non viene trovato nessun altro file viene finalmente restituito
\errcode{EACCES}.  Le altre quattro funzioni si limitano invece a cercare di
eseguire il file indicato dall'argomento \param{path}, che viene interpretato
come il \textit{pathname} del programma.

La terza differenza è come viene passata la lista delle variabili di ambiente.
Con lo mnemonico ``\texttt{e}'' vengono indicate quelle funzioni che
necessitano di un vettore di parametri \var{envp[]} analogo a quello usato per
gli argomenti a riga di comando (terminato quindi da un \val{NULL}), le altre
usano il valore della variabile \var{environ} (vedi
sez.~\ref{sec:proc_environ}) del processo di partenza per costruire
l'ambiente.

Oltre a mantenere lo stesso \ids{PID}, il nuovo programma fatto partire da una
delle funzioni della famiglia \func{exec} mantiene la gran parte delle
proprietà del processo chiamante; una lista delle più significative è la
seguente:
\begin{itemize*}
\item il \textit{process id} (\ids{PID}) ed il \textit{parent process id}
  (\ids{PPID});
\item l'\textsl{user-ID reale}, il \textsl{group-ID reale} ed i
  \textsl{group-ID supplementari} (vedi sez.~\ref{sec:proc_access_id});
\item la directory radice (vedi sez.~\ref{sec:file_chroot}) e la directory di
  lavoro corrente (vedi sez.~\ref{sec:file_work_dir});
\item la maschera di creazione dei file (\textit{umask}, vedi
  sez.~\ref{sec:file_perm_management}) ed i \textit{lock} sui file (vedi
  sez.~\ref{sec:file_locking});
\item il valori di \textit{nice}, le priorità real-time e le affinità di
  processore (vedi sez.~\ref{sec:proc_sched_stand};
  sez.~\ref{sec:proc_real_time} e sez.~\ref{sec:proc_sched_multiprocess});
\item il \textit{session ID} (\acr{sid}) ed il \textit{process group ID}
  (\acr{pgid}), vedi sez.~\ref{sec:sess_proc_group};
\item il terminale di controllo (vedi sez.~\ref{sec:sess_ctrl_term});
\item il tempo restante ad un allarme (vedi sez.~\ref{sec:sig_alarm_abort});
\item i limiti sulle risorse (vedi sez.~\ref{sec:sys_resource_limit});
\item i valori delle variabili \var{tms\_utime}, \var{tms\_stime};
  \var{tms\_cutime}, \var{tms\_ustime} (vedi sez.~\ref{sec:sys_cpu_times});
\item la maschera dei segnali (si veda sez.~\ref{sec:sig_sigmask});
\item l'insieme dei segnali pendenti (vedi sez.~\ref{sec:sig_gen_beha}).
\end{itemize*}

Una serie di proprietà del processo originale, che non avrebbe senso mantenere
in un programma che esegue un codice completamente diverso in uno spazio di
indirizzi totalmente indipendente e ricreato da zero, vengono perse con
l'esecuzione di una \func{exec}. Lo standard POSIX.1-2001 prevede che le
seguenti proprietà non vengano preservate:
\begin{itemize*}
\item gli eventuali stack alternativi per i segnali (vedi
  sez.~\ref{sec:sig_specific_features});
\item i \textit{directory stream} (vedi sez.~\ref{sec:file_dir_read}), che
  vengono chiusi;
\item le mappature dei file in memoria (vedi sez.~\ref{sec:file_memory_map});
\item i segmenti di memoria condivisa SysV (vedi sez.~\ref{sec:ipc_sysv_shm})
  e POSIX (vedi sez.~\ref{sec:ipc_posix_shm});
\item i \textit{memory lock} (vedi sez.~\ref{sec:proc_mem_lock});
\item le funzioni registrate all'uscita (vedi sez.~\ref{sec:proc_atexit});
\item i semafori e le code di messaggi POSIX (vedi
  sez.~\ref{sec:ipc_posix_sem} e sez.~\ref{sec:ipc_posix_mq});
\item i timer POSIX (vedi sez.~\ref{sec:sig_timer_adv}).
\end{itemize*}

Inoltre i segnali che sono stati impostati per essere ignorati nel processo
chiamante mantengono la stessa impostazione pure nel nuovo programma, ma tutti
gli altri segnali, ed in particolare quelli per i quali è stato installato un
gestore vengono impostati alla loro azione predefinita (vedi
sez.~\ref{sec:sig_gen_beha}). Un caso speciale è il segnale \signal{SIGCHLD}
che, quando impostato a \const{SIG\_IGN}, potrebbe anche essere reimpostato a
\const{SIG\_DFL}. Lo standard POSIX.1-2001 prevede che questo comportamento
sia deciso dalla singola implementazione, quella di Linux è di non modificare
l'impostazione precedente.

Oltre alle precedenti, che sono completamente generali e disponibili anche su
altri sistemi unix-like, esistono altre proprietà dei processi, attinenti alle
caratteristiche specifiche di Linux, che non vengono preservate
nell'esecuzione della funzione \func{exec}, queste sono:
\begin{itemize*}
\item le operazioni di I/O asincrono (vedi sez.~\ref{sec:file_asyncronous_io})
  pendenti vengono cancellate;
\item le \textit{capabilities} vengono modificate come
  illustrato in sez.~\ref{sec:proc_capabilities};
\item tutti i \textit{thread} tranne il chiamante\unavref{ (vedi
    sez.~\ref{sec:thread_xxx})} vengono cancellati e tutti gli oggetti ad essi
  relativi\unavref{ (vedi sez.~\ref{sec:thread_xxx})} sono rimossi;
\item viene impostato il flag \const{PR\_SET\_DUMPABLE} di \func{prctl} (vedi
  sez.~\ref{sec:process_prctl}) a meno che il programma da eseguire non sia
  \acr{suid} o \acr{sgid} (vedi sez.~\ref{sec:proc_access_id} e
  sez.~\ref{sec:file_special_perm});
\item il flag \const{PR\_SET\_KEEPCAPS} di \func{prctl} (vedi
  sez.~\ref{sec:process_prctl}) viene cancellato;
\item il nome del processo viene impostato al nome del file contenente il
  programma messo in esecuzione;
\item il segnale di terminazione viene reimpostato a \signal{SIGCHLD};
\item l'ambiente viene reinizializzato impostando le variabili attinenti alla
  localizzazione al valore di default POSIX. 
\end{itemize*}

\itindbeg{close-on-exec}

La gestione dei file aperti nel passaggio al nuovo programma lanciato con
\func{exec} dipende dal valore che ha il flag di \textit{close-on-exec} (vedi
sez.~\ref{sec:file_fcntl_ioctl}) per ciascun \textit{file descriptor}. I file
per cui è impostato vengono chiusi, tutti gli altri file restano
aperti. Questo significa che il comportamento predefinito è che i file restano
aperti attraverso una \func{exec}, a meno di una chiamata esplicita a
\func{fcntl} che imposti il suddetto flag.  Per le directory, lo standard
POSIX.1 richiede che esse vengano chiuse attraverso una \func{exec}, in genere
questo è fatto dalla funzione \func{opendir} (vedi
sez.~\ref{sec:file_dir_read}) che effettua da sola l'impostazione del flag di
\textit{close-on-exec} sulle directory che apre, in maniera trasparente
all'utente.

\itindend{close-on-exec}

Il comportamento della funzione in relazione agli identificatori relativi al
controllo di accesso verrà trattato in dettaglio in sez.~\ref{sec:proc_perms},
qui è sufficiente anticipare (si faccia riferimento a
sez.~\ref{sec:proc_access_id} per la definizione di questi identificatori)
come l'\textsl{user-ID reale} ed il \textsl{group-ID reale} restano sempre gli
stessi, mentre l'\textsl{user-ID salvato} ed il \textsl{group-ID salvato}
vengono impostati rispettivamente all'\textsl{user-ID effettivo} ed il
\textsl{group-ID effettivo}. Questi ultimi normalmente non vengono modificati,
a meno che il file di cui viene chiesta l'esecuzione non abbia o il \acr{suid}
bit o lo \acr{sgid} bit impostato (vedi sez.~\ref{sec:file_special_perm}), in
questo caso l'\textsl{user-ID effettivo} ed il \textsl{group-ID effettivo}
vengono impostati rispettivamente all'utente o al gruppo cui il file
appartiene.

Se il file da eseguire è in formato \emph{a.out} e necessita di librerie
condivise, viene lanciato il \textit{linker} dinamico \cmd{/lib/ld.so} prima
del programma per caricare le librerie necessarie ed effettuare il link
dell'eseguibile; il formato è ormai in completo disuso, per cui è molto
probabile che non il relativo supporto non sia disponibile. Se il programma è
in formato ELF per caricare le librerie dinamiche viene usato l'interprete
indicato nel segmento \constd{PT\_INTERP} previsto dal formato stesso, in
genere questo è \sysfiled{/lib/ld-linux.so.1} per programmi collegati con la
\acr{libc5}, e \sysfiled{/lib/ld-linux.so.2} per programmi collegati con la
\acr{glibc}.

Infine nel caso il programma che si vuole eseguire sia uno script e non un
binario, questo deve essere un file di testo che deve iniziare con una linea
nella forma:
\begin{Example}
#!/path/to/interpreter [argomenti]
\end{Example}
dove l'interprete indicato deve essere un eseguibile binario e non un altro
script, che verrà chiamato come se si fosse eseguito il comando
\cmd{interpreter [argomenti] filename}. 

Si tenga presente che con Linux quanto viene scritto come \texttt{argomenti}
viene passato all'interprete come un unico argomento con una unica stringa di
lunghezza massima di 127 caratteri e se questa dimensione viene ecceduta la
stringa viene troncata; altri Unix hanno dimensioni massime diverse, e diversi
comportamenti, ad esempio FreeBSD esegue la scansione della riga e la divide
nei vari argomenti e se è troppo lunga restituisce un errore di
\errval{ENAMETOOLONG}; una comparazione dei vari comportamenti sui diversi
sistemi unix-like si trova su
\url{http://www.in-ulm.de/~mascheck/various/shebang/}.

Con la famiglia delle \func{exec} si chiude il novero delle funzioni su cui è
basata la gestione tradizionale dei processi in Unix: con \func{fork} si crea
un nuovo processo, con \func{exec} si lancia un nuovo programma, con
\func{exit} e \func{wait} si effettua e verifica la conclusione dei
processi. Tutte le altre funzioni sono ausiliarie e servono per la lettura e
l'impostazione dei vari parametri connessi ai processi.



\section{Il controllo di accesso}
\label{sec:proc_perms}

In questa sezione esamineremo le problematiche relative al controllo di
accesso dal punto di vista dei processi; vedremo quali sono gli identificatori
usati, come questi possono essere modificati nella creazione e nel lancio di
nuovi processi, le varie funzioni per la loro manipolazione diretta e tutte le
problematiche connesse ad una gestione accorta dei privilegi.


\subsection{Gli identificatori del controllo di accesso}
\label{sec:proc_access_id}

Come accennato in sez.~\ref{sec:intro_multiuser} il modello base\footnote{in
  realtà già esistono estensioni di questo modello base, che lo rendono più
  flessibile e controllabile, come le \textit{capabilities} illustrate in
  sez.~\ref{sec:proc_capabilities}, le ACL per i file (vedi
  sez.~\ref{sec:file_ACL}) o il \textit{Mandatory Access Control} di
  \textit{SELinux}; inoltre basandosi sul lavoro effettuato con
  \textit{SELinux}, a partire dal kernel 2.5.x, è iniziato lo sviluppo di una
  infrastruttura di sicurezza, i \textit{Linux Security Modules}, o LSM, in
  grado di fornire diversi agganci a livello del kernel per modularizzare
  tutti i possibili controlli di accesso, cosa che ha permesso di realizzare
  diverse alternative a \textit{SELinux}.} 
di sicurezza di un sistema unix-like è fondato sui concetti di utente e
gruppo, e sulla separazione fra l'amministratore (\textsl{root}, detto spesso
anche \textit{superuser}) che non è sottoposto a restrizioni, ed il resto
degli utenti, per i quali invece vengono effettuati i vari controlli di
accesso.

Abbiamo già accennato come il sistema associ ad ogni utente e gruppo due
identificatori univoci, lo \itindex{User~ID~(UID)} \textsl{User-ID}
(abbreviato in \ids{UID}) ed il \itindex{Group~ID~(GID)} \textsl{Group-ID}
(abbreviato in \ids{GID}). Questi servono al kernel per identificare uno
specifico utente o un gruppo di utenti, per poi poter controllare che essi
siano autorizzati a compiere le operazioni richieste.  Ad esempio in
sez.~\ref{sec:file_access_control} vedremo come ad ogni file vengano associati
un utente ed un gruppo (i suoi \textsl{proprietari}, indicati appunto tramite
un \ids{UID} ed un \ids{GID}) che vengono controllati dal kernel nella
gestione dei permessi di accesso.

Dato che tutte le operazioni del sistema vengono compiute dai processi, è
evidente che per poter implementare un controllo sulle operazioni occorre
anche poter identificare chi è che ha lanciato un certo programma, e pertanto
anche a ciascun processo dovrà essere associato un utente e un gruppo.

Un semplice controllo di una corrispondenza fra identificativi non garantisce
però sufficiente flessibilità per tutti quei casi in cui è necessario poter
disporre di privilegi diversi, o dover impersonare un altro utente per un
limitato insieme di operazioni. Per questo motivo in generale tutti i sistemi
unix-like prevedono che i processi abbiano almeno due gruppi di
identificatori, chiamati rispettivamente \textit{real} ed \textit{effective}
(cioè \textsl{reali} ed \textsl{effettivi}). Nel caso di Linux si aggiungono
poi altri due gruppi, il \textit{saved} (\textsl{salvati}) ed il
\textit{filesystem} (\textsl{di filesystem}), secondo la situazione illustrata
in tab.~\ref{tab:proc_uid_gid}.

\begin{table}[htb]
  \footnotesize
  \centering
  \begin{tabular}[c]{|c|c|l|p{7cm}|}
    \hline
    \textbf{Suffisso} & \textbf{Gruppo} & \textbf{Denominazione} 
                                        & \textbf{Significato} \\ 
    \hline
    \hline
    \texttt{uid} & \textit{real} & \textsl{user-ID reale} 
                 & Indica l'utente che ha lanciato il programma.\\ 
    \texttt{gid} & '' &\textsl{group-ID reale} 
                 & Indica il gruppo principale dell'utente che ha lanciato 
                   il programma.\\ 
    \hline
    \texttt{euid}& \textit{effective} &\textsl{user-ID effettivo} 
                 & Indica l'utente usato nel controllo di accesso.\\ 
    \texttt{egid}& '' & \textsl{group-ID effettivo} 
                 & Indica il gruppo usato nel controllo di accesso.\\ 
    --           & -- & \textsl{group-ID supplementari} 
                 & Indicano gli ulteriori gruppi cui l'utente appartiene.\\ 
    \hline
    --           & \textit{saved} & \textsl{user-ID salvato} 
                 & Mantiene una copia dell'\acr{euid} iniziale.\\ 
    --           & '' & \textsl{group-ID salvato} 
                 & Mantiene una copia dell'\acr{egid} iniziale.\\ 
    \hline
    \texttt{fsuid}& \textit{filesystem} &\textsl{user-ID di filesystem} 
                 & Indica l'utente effettivo per l'accesso al filesystem. \\ 
    \texttt{fsgid}& '' & \textsl{group-ID di filesystem} 
                 & Indica il gruppo effettivo per l'accesso al filesystem.\\ 
    \hline
  \end{tabular}
  \caption{Identificatori di utente e gruppo associati a ciascun processo con
    indicazione dei suffissi usati dalle varie funzioni di manipolazione.}
  \label{tab:proc_uid_gid}
\end{table}

Al primo gruppo appartengono l'\ids{UID} \textsl{reale} ed il \ids{GID}
\textsl{reale}: questi vengono impostati al login ai valori corrispondenti
all'utente con cui si accede al sistema (e relativo gruppo principale).
Servono per l'identificazione dell'utente e normalmente non vengono mai
cambiati. In realtà vedremo (in sez.~\ref{sec:proc_setuid}) che è possibile
modificarli, ma solo ad un processo che abbia i privilegi di amministratore;
questa possibilità è usata proprio dal programma \cmd{login} che, una volta
completata la procedura di autenticazione, lancia una shell per la quale
imposta questi identificatori ai valori corrispondenti all'utente che entra
nel sistema.

Al secondo gruppo appartengono l'\ids{UID} \textsl{effettivo} e il \ids{GID}
\textsl{effettivo}, a cui si aggiungono gli eventuali \ids{GID}
\textsl{supplementari} dei gruppi dei quali l'utente fa parte.  Questi sono
invece gli identificatori usati nelle verifiche dei permessi del processo e
per il controllo di accesso ai file (argomento affrontato in dettaglio in
sez.~\ref{sec:file_perm_overview}).

Questi identificatori normalmente sono identici ai corrispondenti del gruppo
\textit{real} tranne nel caso in cui, come accennato in
sez.~\ref{sec:proc_exec}, il programma che si è posto in esecuzione abbia i
bit \acr{suid} o \acr{sgid} impostati (il significato di questi bit è
affrontato in dettaglio in sez.~\ref{sec:file_special_perm}). In questo caso
essi saranno impostati all'utente e al gruppo proprietari del file. Questo
consente, per programmi in cui ci sia questa necessità, di dare a qualunque
utente i privilegi o i permessi di un altro, compreso l'amministratore.

Come nel caso del \ids{PID} e del \ids{PPID}, anche tutti questi
identificatori possono essere ottenuti da un programma attraverso altrettante
funzioni di sistema dedicate alla loro lettura, queste sono \funcd{getuid},
\funcd{geteuid}, \funcd{getgid} e \funcd{getegid}, ed i loro prototipi sono:

\begin{funcproto}{ 
\fhead{unistd.h}
\fhead{sys/types.h}
\fdecl{uid\_t getuid(void)}
\fdesc{Legge l'\ids{UID} reale del processo corrente.} 
\fdecl{uid\_t geteuid(void)}
\fdesc{Legge l'\ids{UID} effettivo del processo corrente.} 
\fdecl{gid\_t getgid(void)}
\fdesc{Legge il \ids{GID} reale del processo corrente.} 
\fdecl{gid\_t getegid(void)}
\fdesc{Legge il \ids{GID} effettivo del processo corrente.}
}
{Le funzioni ritornano i rispettivi identificativi del processo corrente, e
  non sono previste condizioni di errore.}
\end{funcproto}

In generale l'uso di privilegi superiori, ottenibile con un \ids{UID}
\textsl{effettivo} diverso da quello reale, deve essere limitato il più
possibile, per evitare abusi e problemi di sicurezza, per questo occorre anche
un meccanismo che consenta ad un programma di rilasciare gli eventuali
maggiori privilegi necessari, una volta che si siano effettuate le operazioni
per i quali erano richiesti, e a poterli eventualmente recuperare in caso
servano di nuovo.

Questo in Linux viene fatto usando altri due gruppi di identificatori, il
\textit{saved} ed il \textit{filesystem}. Il primo gruppo è lo stesso usato in
SVr4, e previsto dallo standard POSIX quando è definita
\macro{\_POSIX\_SAVED\_IDS},\footnote{in caso si abbia a cuore la portabilità
  del programma su altri Unix è buona norma controllare sempre la
  disponibilità di queste funzioni controllando se questa costante è
  definita.} il secondo gruppo è specifico di Linux e viene usato per
migliorare la sicurezza con NFS (il \textit{Network File System}, protocollo
che consente di accedere ai file via rete).

L'\ids{UID} \textsl{salvato} ed il \ids{GID} \textsl{salvato} sono copie
dell'\ids{UID} \textsl{effettivo} e del \ids{GID} \textsl{effettivo} del
processo padre, e vengono impostati dalla funzione \func{exec} all'avvio del
processo, come copie dell'\ids{UID} \textsl{effettivo} e del \ids{GID}
\textsl{effettivo} dopo che questi sono stati impostati tenendo conto di
eventuali permessi \acr{suid} o \acr{sgid} (su cui torneremo in
sez.~\ref{sec:file_special_perm}).  Essi quindi consentono di tenere traccia
di quale fossero utente e gruppo effettivi all'inizio dell'esecuzione di un
nuovo programma.

L'\ids{UID} \textsl{di filesystem} e il \ids{GID} \textsl{di filesystem} sono
un'estensione introdotta in Linux per rendere più sicuro l'uso di NFS
(torneremo sull'argomento in sez.~\ref{sec:proc_setuid}). Essi sono una
replica dei corrispondenti identificatori del gruppo \textit{effective}, ai
quali si sostituiscono per tutte le operazioni di verifica dei permessi
relativi ai file (trattate in sez.~\ref{sec:file_perm_overview}).  Ogni
cambiamento effettuato sugli identificatori effettivi viene automaticamente
riportato su di essi, per cui in condizioni normali si può tranquillamente
ignorarne l'esistenza, in quanto saranno del tutto equivalenti ai precedenti.


\subsection{Le funzioni di gestione degli identificatori dei processi}
\label{sec:proc_setuid}

Le funzioni di sistema più comuni che vengono usate per cambiare identità
(cioè utente e gruppo di appartenenza) ad un processo, e che come accennato in
sez.~\ref{sec:proc_access_id} seguono la semantica POSIX che prevede
l'esistenza dell'\ids{UID} salvato e del \ids{GID} salvato, sono
rispettivamente \funcd{setuid} e \funcd{setgid}; i loro prototipi sono:

\begin{funcproto}{
\fhead{unistd.h}
\fhead{sys/types.h}
\fdecl{int setuid(uid\_t uid)}
\fdesc{Imposta l'\ids{UID} del processo corrente.}
\fdecl{int setgid(gid\_t gid)}
\fdesc{Imposta il \ids{GID} del processo corrente.}
}
{Le funzioni ritornano $0$ in caso di successo e $-1$ per un errore, nel qual
caso \var{errno} uno dei valori: 
\begin{errlist}
\item[\errcode{EAGAIN}] (solo per \func{setuid}) la chiamata cambierebbe
  l'\ids{UID} reale ma il kernel non dispone temporaneamente delle risorse per
  farlo, oppure, per i kernel precedenti il 3.1, il cambiamento
  dell'\ids{UID} reale farebbe superare il limite per il numero dei processi
  \const{RLIMIT\_NPROC} (vedi sez.~\ref{sec:sys_resource_limit}).
\item[\errcode{EINVAL}] il valore di dell'argomento non è valido per il
    \textit{namespace} corrente (vedi sez.~\ref{sec:process_namespaces}).
\item[\errcode{EPERM}] non si hanno i permessi per l'operazione richiesta.
\end{errlist} 
}
\end{funcproto}

Il funzionamento di queste due funzioni è analogo, per cui considereremo solo
la prima, la seconda si comporta esattamente allo stesso modo facendo
riferimento al \ids{GID} invece che all'\ids{UID}. Gli eventuali \ids{GID}
supplementari non vengono modificati.

L'effetto della chiamata è diverso a seconda dei privilegi del processo; se
l'\ids{UID} effettivo è zero (cioè è quello dell'amministratore di sistema o
il processo ha la capacità \const{CAP\_SETUID}) allora tutti gli
identificatori (\textit{real}, \textit{effective} e \textit{saved}) vengono
impostati al valore specificato da \param{uid}, altrimenti viene impostato
solo l'\ids{UID} effettivo, e soltanto se il valore specificato corrisponde o
all'\ids{UID} reale o all'\ids{UID} salvato, ottenendo un errore di
\errcode{EPERM} negli altri casi.

E' importante notare che la funzione può fallire, con
\errval{EAGAIN},\footnote{non affronteremo qui l'altro caso di errore, che può
  avvenire solo quando si esegue la funzione all'interno di un diverso
  \textit{user namespace}, argomento su cui torneremo in
  sez.~\ref{sec:process_namespaces} ma la considerazione di controllare sempre
  lo stato di uscita si applica allo stesso modo.} anche quando viene invocata
da un processo con privilegi di amministratore per cambiare il proprio
l'\ids{UID} reale, sia per una temporanea indisponibilità di risorse del
kernel, sia perché l'utente di cui si vuole assumere l'\ids{UID} andrebbe a
superare un eventuale limite sul numero di processi (il limite
\const{RLIMIT\_NPROC}, che tratteremo in sez.~\ref{sec:sys_resource_limit}),
pertanto occorre sempre verificare lo stato di uscita della funzione.

Non controllare questo tipo di errori perché si presume che la funzione abbia
sempre successo quando si hanno i privilegi di amministratore può avere
conseguente devastanti per la sicurezza, in particolare quando la si usa per
cedere i suddetti privilegi ed eseguire un programma per conto di un utente
non privilegiato.

E' per diminuire l'impatto di questo tipo di disattenzioni che a partire dal
kernel 3.1 il comportamento di \func{setuid} e di tutte le analoghe funzioni
che tratteremo nel seguito di questa sezione è stato modificato e nel caso di
superamento del limite sulle risorse esse hanno comunque successo. Quando
questo avviene il processo assume comunque il nuovo \ids{UID} ed il controllo
sul superamento di \const{RLIMIT\_NPROC} viene posticipato ad una eventuale
successiva invocazione di \func{execve} (essendo questo poi il caso d'uso più
comune). In tal caso, se alla chiamata ancora sussiste la situazione di
superamento del limite, sarà \func{execve} a fallire con un errore di
\const{EAGAIN}.\footnote{che pertanto, a partire dal kernel 3.1, può
  restituire anche questo errore, non presente in altri sistemi
  \textit{unix-like}.}

Come accennato l'uso principale di queste funzioni è quello di poter
consentire ad un programma con i bit \acr{suid} o \acr{sgid} impostati (vedi
sez.~\ref{sec:file_special_perm}) di riportare l'\ids{UID} effettivo a quello
dell'utente che ha lanciato il programma, effettuare il lavoro che non
necessita di privilegi aggiuntivi, ed eventualmente tornare indietro.

Come esempio per chiarire l'uso di queste funzioni prendiamo quello con cui
viene gestito l'accesso al file \sysfiled{/var/run/utmp}.  In questo file viene
registrato chi sta usando il sistema al momento corrente; chiaramente non può
essere lasciato aperto in scrittura a qualunque utente, che potrebbe
falsificare la registrazione.

Per questo motivo questo file (e l'analogo \sysfiled{/var/log/wtmp} su cui
vengono registrati login e logout) appartengono ad un gruppo dedicato (in
genere \acr{utmp}) ed i programmi che devono accedervi (ad esempio tutti i
programmi di terminale in X, o il programma \cmd{screen} che crea terminali
multipli su una console) appartengono a questo gruppo ed hanno il bit
\acr{sgid} impostato.

Quando uno di questi programmi (ad esempio \cmd{xterm}) viene lanciato, la
situazione degli identificatori è la seguente:
\begin{eqnarray*}
  \label{eq:1}
  \textsl{group-ID reale}      &=& \textrm{\ids{GID} (del chiamante)} \\
  \textsl{group-ID effettivo}  &=& \textrm{\acr{utmp}} \\
  \textsl{group-ID salvato}    &=& \textrm{\acr{utmp}}
\end{eqnarray*}
in questo modo, dato che il \textsl{group-ID effettivo} è quello giusto, il
programma può accedere a \sysfile{/var/run/utmp} in scrittura ed aggiornarlo.
A questo punto il programma può eseguire una \code{setgid(getgid())} per
impostare il \textsl{group-ID effettivo} a quello dell'utente (e dato che il
\textsl{group-ID reale} corrisponde la funzione avrà successo), in questo modo
non sarà possibile lanciare dal terminale programmi che modificano detto file,
in tal caso infatti la situazione degli identificatori sarebbe:
\begin{eqnarray*}
  \label{eq:2}
  \textsl{group-ID reale}      &=& \textrm{\ids{GID} (invariato)}  \\
  \textsl{group-ID effettivo}  &=& \textrm{\ids{GID}} \\
  \textsl{group-ID salvato}    &=& \textrm{\acr{utmp} (invariato)}
\end{eqnarray*}
e ogni processo lanciato dal terminale avrebbe comunque \ids{GID} come
\textsl{group-ID effettivo}. All'uscita dal terminale, per poter di nuovo
aggiornare lo stato di \sysfile{/var/run/utmp} il programma eseguirà una
\code{setgid(utmp)} (dove \var{utmp} è il valore numerico associato al gruppo
\acr{utmp}, ottenuto ad esempio con una precedente \func{getegid}), dato che
in questo caso il valore richiesto corrisponde al \textsl{group-ID salvato} la
funzione avrà successo e riporterà la situazione a:
\begin{eqnarray*}
  \label{eq:3}
  \textsl{group-ID reale}      &=& \textrm{\ids{GID} (invariato)}  \\
  \textsl{group-ID effettivo}  &=& \textrm{\acr{utmp}} \\
  \textsl{group-ID salvato}    &=& \textrm{\acr{utmp} (invariato)}
\end{eqnarray*}
consentendo l'accesso a \sysfile{/var/run/utmp}.

Occorre però tenere conto che tutto questo non è possibile con un processo con
i privilegi di amministratore, in tal caso infatti l'esecuzione di una
\func{setuid} comporta il cambiamento di tutti gli identificatori associati al
processo, rendendo impossibile riguadagnare i privilegi di amministratore.
Questo comportamento è corretto per l'uso che ne fa un programma come
\cmd{login} una volta che crea una nuova shell per l'utente, ma quando si
vuole cambiare soltanto l'\ids{UID} effettivo del processo per cedere i
privilegi occorre ricorrere ad altre funzioni.

Le due funzioni di sistema \funcd{setreuid} e \funcd{setregid} derivano da BSD
che, non supportando (almeno fino alla versione 4.3+BSD) gli identificatori
del gruppo \textit{saved}, le usa per poter scambiare fra di loro
\textit{effective} e \textit{real}; i rispettivi prototipi sono:

\begin{funcproto}{ 
\fhead{unistd.h}
\fhead{sys/types.h}
\fdecl{int setreuid(uid\_t ruid, uid\_t euid)}
\fdesc{Imposta \ids{UID} reale e \ids{UID} effettivo del processo corrente.} 
\fdecl{int setregid(gid\_t rgid, gid\_t egid)}
\fdesc{Imposta \ids{GID} reale e \ids{GID} effettivo del processo corrente.} 
}
{Le funzioni ritornano $0$ in caso di successo e $-1$ per un errore, nel qual
caso \var{errno} assume i valori visti per \func{setuid}/\func{setgid}.
}
\end{funcproto}

Le due funzioni sono identiche, quanto diremo per la prima riguardo gli
\ids{UID} si applica alla seconda per i \ids{GID}.  La funzione
\func{setreuid} imposta rispettivamente l'\ids{UID} reale e l'\ids{UID}
effettivo del processo corrente ai valori specificati da \param{ruid}
e \param{euid}.

I processi non privilegiati possono impostare solo valori che corrispondano o
al loro \ids{UID} effettivo o a quello reale o a quello salvato, valori
diversi comportano il fallimento della chiamata.  L'amministratore invece può
specificare un valore qualunque.  Specificando un argomento di valore $-1$
l'identificatore corrispondente verrà lasciato inalterato.

Con queste funzioni si possono scambiare fra loro gli \ids{UID} reale ed
effettivo, e pertanto è possibile implementare un comportamento simile a
quello visto in precedenza per \func{setgid}, cedendo i privilegi con un primo
scambio, e recuperandoli, una volta eseguito il lavoro non privilegiato, con
un secondo scambio.

In questo caso però occorre porre molta attenzione quando si creano nuovi
processi nella fase intermedia in cui si sono scambiati gli identificatori, in
questo caso infatti essi avranno un \ids{UID} reale privilegiato, che dovrà
essere esplicitamente eliminato prima di porre in esecuzione un nuovo
programma, occorrerà cioè eseguire un'altra chiamata dopo la \func{fork} e
prima della \func{exec} per uniformare l'\ids{UID} reale a quello effettivo,
perché in caso contrario il nuovo programma potrebbe a sua volta effettuare
uno scambio e riottenere dei privilegi non previsti.

Lo stesso problema di propagazione dei privilegi ad eventuali processi figli
si pone anche per l'\ids{UID} salvato. Ma la funzione \func{setreuid} deriva
da un'implementazione di sistema che non ne prevede la presenza, e quindi non
è possibile usarla per correggere la situazione come nel caso precedente. Per
questo motivo in Linux tutte le volte che si imposta un qualunque valore
diverso da quello dall'\ids{UID} reale corrente, l'\ids{UID} salvato viene
automaticamente uniformato al valore dell'\ids{UID} effettivo.

Altre due funzioni di sistema, \funcd{seteuid} e \funcd{setegid}, sono
un'estensione dello standard POSIX.1, ma sono comunque supportate dalla
maggior parte degli Unix, esse vengono usate per cambiare gli identificatori
del gruppo \textit{effective} ed i loro prototipi sono:

\begin{funcproto}{ 
\fhead{unistd.h}
\fhead{sys/types.h}
\fdecl{int seteuid(uid\_t uid)}
\fdesc{Imposta l'\ids{UID} effettivo del processo corrente.} 
\fdecl{int setegid(gid\_t gid)}
\fdesc{Imposta il \ids{GID} effettivo del processo corrente.} 
}
{Le funzioni ritornano $0$ in caso di successo e $-1$ per un errore, nel qual
  caso \var{errno} assume i valori visti per \func{setuid}/\func{setgid}
  tranne \errval{EAGAIN}. 
}
\end{funcproto}

Ancora una volta le due funzioni sono identiche, e quanto diremo per la prima
riguardo gli \ids{UID} si applica allo stesso modo alla seconda per i
\ids{GID}. Con \func{seteuid} gli utenti normali possono impostare l'\ids{UID}
effettivo solo al valore dell'\ids{UID} reale o dell'\ids{UID} salvato,
l'amministratore può specificare qualunque valore. Queste funzioni sono usate
per permettere all'amministratore di impostare solo l'\ids{UID} effettivo,
dato che l'uso normale di \func{setuid} comporta l'impostazione di tutti gli
identificatori.
 
Le due funzioni di sistema \funcd{setresuid} e \funcd{setresgid} sono invece
un'estensione introdotta in Linux (a partire dal kernel 2.1.44) e permettono
un completo controllo su tutti e tre i gruppi di identificatori
(\textit{real}, \textit{effective} e \textit{saved}), i loro prototipi sono:

\begin{funcproto}{ 
\fhead{unistd.h}
\fhead{sys/types.h}
\fdecl{int setresuid(uid\_t ruid, uid\_t euid, uid\_t suid)}
\fdesc{Imposta l'\ids{UID} reale, effettivo e salvato del processo corrente.} 
\fdecl{int setresgid(gid\_t rgid, gid\_t egid, gid\_t sgid)}
\fdesc{Imposta il \ids{GID} reale, effettivo e salvato del processo corrente.} 
}
{Le funzioni ritornano $0$ in caso di successo e $-1$ per un errore, nel qual
caso \var{errno} assume i valori visti per \func{setuid}/\func{setgid}.
}
\end{funcproto}

Di nuovo le due funzioni sono identiche e quanto detto per la prima riguardo
gli \ids{UID} si applica alla seconda per i \ids{GID}.  La funzione
\func{setresuid} imposta l'\ids{UID} reale, l'\ids{UID} effettivo e
l'\ids{UID} salvato del processo corrente ai valori specificati
rispettivamente dagli argomenti \param{ruid}, \param{euid} e \param{suid}.  I
processi non privilegiati possono cambiare uno qualunque degli \ids{UID} solo
ad un valore corrispondente o all'\ids{UID} reale, o a quello effettivo o a
quello salvato, l'amministratore può specificare i valori che vuole. Un valore
di $-1$ per un qualunque argomento lascia inalterato l'identificatore
corrispondente.

Per queste funzioni di sistema esistono anche due controparti,
\funcd{getresuid} e \funcd{getresgid},\footnote{le funzioni non sono standard,
  anche se appaiono in altri kernel, su Linux sono presenti dal kernel 2.1.44
  e con le versioni della \acr{glibc} a partire dalla 2.3.2, definendo la
  macro \macro{\_GNU\_SOURCE}.} che permettono di leggere in blocco i vari
identificatori; i loro prototipi sono:

\begin{funcproto}{ 
\fhead{unistd.h}
\fhead{sys/types.h}
\fdecl{int getresuid(uid\_t *ruid, uid\_t *euid, uid\_t *suid)}
\fdesc{Legge l'\ids{UID} reale, effettivo e salvato del processo corrente.} 
\fdecl{int getresgid(gid\_t *rgid, gid\_t *egid, gid\_t *sgid)}
\fdesc{Legge il \ids{GID} reale, effettivo e salvato del processo corrente.} 
}
{Le funzioni ritornano $0$ in caso di successo e $-1$ per un errore, nel qual
  caso \var{errno} può assumere solo il valore \errcode{EFAULT} se gli
  indirizzi delle variabili di ritorno non sono validi.  }
\end{funcproto}

Anche queste funzioni sono un'estensione specifica di Linux, e non richiedono
nessun privilegio. I valori sono restituiti negli argomenti, che vanno
specificati come puntatori (è un altro esempio di \textit{value result
  argument}). Si noti che queste funzioni sono le uniche in grado di leggere
gli identificatori del gruppo \textit{saved}.

Infine le funzioni \func{setfsuid} e \func{setfsgid} servono per impostare gli
identificatori del gruppo \textit{filesystem} che sono usati da Linux per il
controllo dell'accesso ai file.  Come già accennato in
sez.~\ref{sec:proc_access_id} Linux definisce questo ulteriore gruppo di
identificatori, che in circostanze normali sono assolutamente equivalenti a
quelli del gruppo \textit{effective}, dato che ogni cambiamento di questi
ultimi viene immediatamente riportato su di essi.

C'è un solo caso in cui si ha necessità di introdurre una differenza fra gli
identificatori dei gruppi \textit{effective} e \textit{filesystem}, ed è per
ovviare ad un problema di sicurezza che si presenta quando si deve
implementare un server NFS. 

Il server NFS infatti deve poter cambiare l'identificatore con cui accede ai
file per assumere l'identità del singolo utente remoto, ma se questo viene
fatto cambiando l'\ids{UID} effettivo o l'\ids{UID} reale il server si espone
alla ricezione di eventuali segnali ostili da parte dell'utente di cui ha
temporaneamente assunto l'identità.  Cambiando solo l'\ids{UID} di filesystem
si ottengono i privilegi necessari per accedere ai file, mantenendo quelli
originari per quanto riguarda tutti gli altri controlli di accesso, così che
l'utente non possa inviare segnali al server NFS.

Le due funzioni di sistema usate appositamente per cambiare questi
identificatori sono \funcd{setfsuid} e \funcd{setfsgid} ovviamente sono
specifiche di Linux e non devono essere usate se si intendono scrivere
programmi portabili; i loro prototipi sono:

\begin{funcproto}{ 
\fhead{sys/fsuid.h}
\fdecl{int setfsuid(uid\_t fsuid)}
\fdesc{Imposta l'\ids{UID} di filesystem del processo corrente.} 
\fdecl{int setfsgid(gid\_t fsgid)}
\fdesc{Legge il \ids{GID} di filesystem del processo corrente.} 
}

{Le funzioni restituiscono sia in caso di successo che di errore il valore
  corrente dell'identificativo, e in caso di errore non viene impostato nessun
  codice in \var{errno}.}
\end{funcproto}

Le due funzioni sono analoghe ed usano il valore passato come argomento per
effettuare l'impostazione dell'identificativo.  Le funzioni hanno successo
solo se il processo chiamante ha i privilegi di amministratore o, per gli
altri utenti, se il valore specificato coincide con uno dei di quelli del
gruppo \textit{real}, \textit{effective} o \textit{saved}.

Il problema di queste funzioni è che non restituiscono un codice di errore e
non c'è modo di sapere (con una singola chiamata) di sapere se hanno avuto
successo o meno, per verificarlo occorre eseguire una chiamata aggiuntiva
passando come argomento $-1$ (un valore impossibile per un identificativo),
così fallendo si può di ottenere il valore corrente e verificare se è
cambiato.

\subsection{Le funzioni per la gestione dei gruppi associati a un processo}
\label{sec:proc_setgroups}

Le ultime funzioni che esamineremo sono quelle che permettono di operare sui
gruppi supplementari cui un utente può appartenere. Ogni processo può avere
almeno \const{NGROUPS\_MAX} gruppi supplementari\footnote{il numero massimo di
  gruppi secondari può essere ottenuto con \func{sysconf} (vedi
  sez.~\ref{sec:sys_limits}), leggendo il parametro
  \const{\_SC\_NGROUPS\_MAX}.} in aggiunta al gruppo primario; questi vengono
ereditati dal processo padre e possono essere cambiati con queste funzioni.

La funzione di sistema che permette di leggere i gruppi supplementari
associati ad un processo è \funcd{getgroups}; questa funzione è definita nello
standard POSIX.1, ed il suo prototipo è:

\begin{funcproto}{ 
\fhead{sys/types.h}
\fhead{unistd.h}
\fdecl{int getgroups(int size, gid\_t list[])}
\fdesc{Legge gli identificatori dei gruppi supplementari.} 
}
{La funzione ritorna il numero di gruppi letti in caso di successo e $-1$ per
  un errore, nel qual caso \var{errno} assumerà uno dei valori:
\begin{errlist}
\item[\errcode{EFAULT}] \param{list} non ha un indirizzo valido.
\item[\errcode{EINVAL}] il valore di \param{size} è diverso da zero ma
  minore del numero di gruppi supplementari del processo.
\end{errlist}}
\end{funcproto}

La funzione legge gli identificatori dei gruppi supplementari del processo sul
vettore \param{list} che deve essere di dimensione pari a \param{size}. Non è
specificato se la funzione inserisca o meno nella lista il \ids{GID} effettivo
del processo. Se si specifica un valore di \param{size} uguale a $0$ allora
l'argomento \param{list} non viene modificato, ma si ottiene dal valore di
ritorno il numero di gruppi supplementari.

Una seconda funzione, \funcd{getgrouplist}, può invece essere usata per
ottenere tutti i gruppi a cui appartiene utente identificato per nome; il suo
prototipo è:

\begin{funcproto}{ 
\fhead{grp.h}
\fdecl{int getgrouplist(const char *user, gid\_t group, gid\_t *groups, int
  *ngroups)} 
\fdesc{Legge i gruppi cui appartiene un utente.} 
}
{La funzione ritorna il numero di gruppi ottenuto in caso di successo e $-1$
  per un errore, che avviene solo quando il numero di gruppi è maggiore di
  quelli specificati con \param{ngroups}.}
\end{funcproto}

La funzione esegue una scansione del database dei gruppi (si veda
sez.~\ref{sec:sys_user_group}) per leggere i gruppi supplementari dell'utente
specificato per nome (e non con un \ids{UID}) nella stringa passata con
l'argomento \param{user}. Ritorna poi nel vettore \param{groups} la lista dei
\ids{GID} dei gruppi a cui l'utente appartiene. Si noti che \param{ngroups},
che in ingresso deve indicare la dimensione di \param{group}, è passato come
\textit{value result argument} perché, qualora il valore specificato sia
troppo piccolo, la funzione ritorna $-1$, passando comunque indietro il numero
dei gruppi trovati, in modo da poter ripetere la chiamata con un vettore di
dimensioni adeguate.

Infine per impostare i gruppi supplementari di un processo ci sono due
funzioni, che possono essere usate solo se si hanno i privilegi di
amministratore.\footnote{e più precisamente se si ha la \textit{capability}
  \const{CAP\_SETGID}.} La prima delle due è la funzione di sistema
\funcd{setgroups},\footnote{la funzione è definita in BSD e SRv4, ma a
  differenza di \func{getgroups} non è stata inclusa in POSIX.1-2001, per
  poterla utilizzare deve essere definita la macro \macro{\_BSD\_SOURCE}.} ed
il suo prototipo è:

\begin{funcproto}{ 
\fhead{grp.h}
\fdecl{int setgroups(size\_t size, gid\_t *list)}
\fdesc{Imposta i gruppi supplementari del processo.} 
}
{La funzione ritorna $0$ in caso di successo e $-1$ per un errore, nel qual
caso \var{errno} assumerà uno dei valori:
\begin{errlist}
\item[\errcode{EFAULT}] \param{list} non ha un indirizzo valido.
\item[\errcode{EINVAL}] il valore di \param{size} è maggiore del valore
    massimo consentito di gruppi supplementari.
\item[\errcode{EPERM}] il processo non ha i privilegi di amministratore.
\end{errlist}}
\end{funcproto}

La funzione imposta i gruppi supplementari del processo corrente ai valori
specificati nel vettore passato con l'argomento \param{list}, di dimensioni
date dall'argomento \param{size}. Il numero massimo di gruppi supplementari
che si possono impostare è un parametro di sistema, che può essere ricavato
con le modalità spiegate in sez.~\ref{sec:sys_characteristics}.

Se invece si vogliono impostare i gruppi supplementari del processo a quelli
di un utente specifico, si può usare la funzione \funcd{initgroups} il cui
prototipo è:

\begin{funcproto}{ 
\fhead{sys/types.h}
\fhead{grp.h}
\fdecl{int initgroups(const char *user, gid\_t group)}
\fdesc{Inizializza la lista dei gruppi supplementari.} 
}
{La funzione ritorna $0$ in caso di successo e $-1$ per un errore, nel qual
caso \var{errno} assumerà uno dei valori:
\begin{errlist}
\item[\errcode{ENOMEM}] non c'è memoria sufficiente per allocare lo spazio per
  informazioni dei gruppi.
\item[\errcode{EPERM}] il processo non ha i privilegi di amministratore.
\end{errlist}}
\end{funcproto}

La funzione esegue la scansione del database dei gruppi (usualmente
\conffile{/etc/group}) cercando i gruppi di cui è membro l'utente \param{user}
(di nuovo specificato per nome e non per \ids{UID}) con cui costruisce una
lista di gruppi supplementari, a cui aggiunge anche \param{group}, infine
imposta questa lista per il processo corrente usando \func{setgroups}.  Si
tenga presente che sia \func{setgroups} che \func{initgroups} non sono
definite nello standard POSIX.1 e che pertanto non è possibile utilizzarle
quando si definisce \macro{\_POSIX\_SOURCE} o si compila con il flag
\cmd{-ansi}, è pertanto meglio evitarle se si vuole scrivere codice portabile.

 
\section{La gestione della priorità dei processi}
\label{sec:proc_priority}

In questa sezione tratteremo più approfonditamente i meccanismi con il quale
lo \textit{scheduler} assegna la CPU ai vari processi attivi.  In particolare
prenderemo in esame i vari meccanismi con cui viene gestita l'assegnazione del
tempo di CPU, ed illustreremo le varie funzioni di gestione. Tratteremo infine
anche le altre priorità dei processi (come quelle per l'accesso a disco)
divenute disponibili con i kernel più recenti.


\subsection{I meccanismi di \textit{scheduling}}
\label{sec:proc_sched}

\itindbeg{scheduler}

La scelta di un meccanismo che sia in grado di distribuire in maniera efficace
il tempo di CPU per l'esecuzione dei processi è sempre una questione delicata,
ed oggetto di numerose ricerche; in generale essa dipende in maniera
essenziale anche dal tipo di utilizzo che deve essere fatto del sistema, per
cui non esiste un meccanismo che sia valido per tutti gli usi.

La caratteristica specifica di un sistema \textit{multitasking} come Linux è
quella del cosiddetto \itindex{preemptive~multitasking} \textit{preemptive
  multitasking}: questo significa che al contrario di altri sistemi (che usano
invece il cosiddetto \itindex{cooperative~multitasking} \textit{cooperative
  multitasking}) non sono i singoli processi, ma il kernel stesso a decidere
quando la CPU deve essere passata ad un altro processo. Come accennato in
sez.~\ref{sec:proc_hierarchy} questa scelta viene eseguita da una sezione
apposita del kernel, lo \textit{scheduler}, il cui scopo è quello di
distribuire al meglio il tempo di CPU fra i vari processi.

La cosa è resa ancora più complicata dal fatto che con le architetture
multi-processore si deve anche scegliere quale sia la CPU più opportuna da
utilizzare.\footnote{nei processori moderni la presenza di ampie cache può
  rendere poco efficiente trasferire l'esecuzione di un processo da una CPU ad
  un'altra, per cui effettuare la migliore scelta fra le diverse CPU non è
  banale.}  Tutto questo comunque appartiene alle sottigliezze
dell'implementazione del kernel; dal punto di vista dei programmi che girano
in \textit{user space}, anche quando si hanno più processori (e dei processi
che sono eseguiti davvero in contemporanea), le politiche di
\textit{scheduling} riguardano semplicemente l'allocazione della risorsa
\textsl{tempo di esecuzione}, la cui assegnazione sarà governata dai
meccanismi di scelta delle priorità che restano gli stessi indipendentemente
dal numero di processori.

Si tenga conto poi che i processi non devono solo eseguire del codice: ad
esempio molto spesso saranno impegnati in operazioni di I/O, o potranno
venire bloccati da un comando dal terminale, o sospesi per un certo periodo di
tempo.  In tutti questi casi la CPU diventa disponibile ed è compito dello
kernel provvedere a mettere in esecuzione un altro processo.

Tutte queste possibilità sono caratterizzate da un diverso \textsl{stato} del
processo; in Linux un processo può trovarsi in uno degli stati riportati in
tab.~\ref{tab:proc_proc_states}; ma soltanto i processi che sono nello stato
\textit{runnable} concorrono per l'esecuzione. Questo vuol dire che, qualunque
sia la sua priorità, un processo non potrà mai essere messo in esecuzione
fintanto che esso si trova in uno qualunque degli altri stati.

\begin{table}[htb]
  \footnotesize
  \centering
  \begin{tabular}[c]{|p{2.4cm}|c|p{9cm}|}
    \hline
    \textbf{Stato} & \texttt{STAT} & \textbf{Descrizione} \\
    \hline
    \hline
    \textit{runnable}& \texttt{R} & Il processo è in esecuzione o è pronto ad
                                    essere eseguito (in attesa che gli
                                    venga assegnata la CPU).\\
    \textit{sleep}   & \texttt{S} & Il processo  è in attesa di un
                                    risposta dal sistema, ma può essere 
                                    interrotto da un segnale.\\
    \textit{uninterrutible sleep}& \texttt{D} & Il  processo è in
                                    attesa di un risposta dal sistema (in 
                                    genere per I/O), e non può essere
                                    interrotto in nessuna circostanza.\\
    \textit{stopped} & \texttt{T} & Il processo è stato fermato con un
                                    \signal{SIGSTOP}, o è tracciato.\\
    \textit{zombie}  & \texttt{Z} & Il processo è terminato ma il
                                    suo stato di terminazione non è ancora
                                    stato letto dal padre.\\
    \textit{killable}& \texttt{D} & Un nuovo stato introdotto con il kernel
                                    2.6.25, sostanzialmente identico
                                    all'\textit{uninterrutible sleep} con la
                                    sola differenza che il processo può
                                    terminato con \signal{SIGKILL} (usato per
                                    lo più per NFS).\\ 
    \hline
  \end{tabular}
  \caption{Elenco dei possibili stati di un processo in Linux, nella colonna
    \texttt{STAT} si è riportata la corrispondente lettera usata dal comando 
    \cmd{ps} nell'omonimo campo.}
  \label{tab:proc_proc_states}
\end{table}

Si deve quindi tenere presente che l'utilizzo della CPU è soltanto una delle
risorse che sono necessarie per l'esecuzione di un programma, e a seconda
dello scopo del programma non è detto neanche che sia la più importante, dato
che molti programmi dipendono in maniera molto più critica dall'I/O. Per
questo motivo non è affatto detto che dare ad un programma la massima priorità
di esecuzione abbia risultati significativi in termini di prestazioni.

Il meccanismo tradizionale di \textit{scheduling} di Unix (che tratteremo in
sez.~\ref{sec:proc_sched_stand}) è sempre stato basato su delle
\textsl{priorità dinamiche}, in modo da assicurare che tutti i processi, anche
i meno importanti, potessero ricevere un po' di tempo di CPU. In sostanza
quando un processo ottiene la CPU la sua priorità viene diminuita. In questo
modo alla fine, anche un processo con priorità iniziale molto bassa, finisce
per avere una priorità sufficiente per essere eseguito.

Lo standard POSIX.1b però ha introdotto il concetto di \textsl{priorità
  assoluta}, (chiamata anche \textsl{priorità statica}, in contrapposizione
alla normale priorità dinamica), per tenere conto dei sistemi
\textit{real-time},\footnote{per sistema \textit{real-time} si intende un
  sistema in grado di eseguire operazioni in un tempo ben determinato; in
  genere si tende a distinguere fra l'\textit{hard real-time} in cui è
  necessario che i tempi di esecuzione di un programma siano determinabili con
  certezza assoluta (come nel caso di meccanismi di controllo di macchine,
  dove uno sforamento dei tempi avrebbe conseguenze disastrose), e
  \textit{soft-real-time} in cui un occasionale sforamento è ritenuto
  accettabile.} in cui è vitale che i processi che devono essere eseguiti in
un determinato momento non debbano aspettare la conclusione di altri che non
hanno questa necessità.

Il concetto di priorità assoluta dice che quando due processi si contendono
l'esecuzione, vince sempre quello con la priorità assoluta più alta.
Ovviamente questo avviene solo per i processi che sono pronti per essere
eseguiti (cioè nello stato \textit{runnable}).  La priorità assoluta viene in
genere indicata con un numero intero, ed un valore più alto comporta una
priorità maggiore. Su questa politica di \textit{scheduling} torneremo in
sez.~\ref{sec:proc_real_time}.

In generale quello che succede in tutti gli Unix moderni è che ai processi
normali viene sempre data una priorità assoluta pari a zero, e la decisione di
assegnazione della CPU è fatta solo con il meccanismo tradizionale della
priorità dinamica. In Linux tuttavia è possibile assegnare anche una priorità
assoluta, nel qual caso un processo avrà la precedenza su tutti gli altri di
priorità inferiore, che saranno eseguiti solo quando quest'ultimo non avrà
bisogno della CPU.


\subsection{Il meccanismo di \textit{scheduling} standard}
\label{sec:proc_sched_stand}

A meno che non si abbiano esigenze specifiche,\footnote{per alcune delle quali
  sono state introdotte delle varianti specifiche.} l'unico meccanismo di
\textit{scheduling} con il quale si avrà a che fare è quello tradizionale, che
prevede solo priorità dinamiche. È di questo che, di norma, ci si dovrà
preoccupare nella programmazione.  Come accennato in Linux i processi ordinari
hanno tutti una priorità assoluta nulla; quello che determina quale, fra tutti
i processi in attesa di esecuzione, sarà eseguito per primo, è la cosiddetta
\textsl{priorità dinamica}, quella che viene mostrata nella colonna
\texttt{PR} del comando \texttt{top}, che è chiamata così proprio perché varia
nel corso dell'esecuzione di un processo.

Il meccanismo usato da Linux è in realtà piuttosto complesso,\footnote{e
  dipende strettamente dalla versione di kernel; in particolare a partire
  dalla serie 2.6.x lo \textit{scheduler} è stato riscritto completamente, con
  molte modifiche susseguitesi per migliorarne le prestazioni, per un certo
  periodo ed è stata anche introdotta la possibilità di usare diversi
  algoritmi, selezionabili sia in fase di compilazione, che, nelle versioni
  più recenti, all'avvio (addirittura è stato ideato un sistema modulare che
  permette di cambiare lo \textit{scheduler} a sistema attivo).} ma a grandi
linee si può dire che ad ogni processo è assegnata una \textit{time-slice},
cioè un intervallo di tempo (letteralmente una fetta) per il quale, a meno di
eventi esterni, esso viene eseguito senza essere interrotto.  Inoltre la
priorità dinamica viene calcolata dallo \textit{scheduler} a partire da un
valore iniziale che viene \textsl{diminuito} tutte le volte che un processo è
in stato \textit{runnable} ma non viene posto in esecuzione.\footnote{in
  realtà il calcolo della priorità dinamica e la conseguente scelta di quale
  processo mettere in esecuzione avviene con un algoritmo molto più
  complicato, che tiene conto anche della \textsl{interattività} del processo,
  utilizzando diversi fattori, questa è una brutale semplificazione per
  rendere l'idea del funzionamento, per una trattazione più dettagliata dei
  meccanismi di funzionamento dello \textit{scheduler}, anche se non
  aggiornatissima, si legga il quarto capitolo di \cite{LinKernDev}.}

Lo \textit{scheduler} infatti mette sempre in esecuzione, fra tutti i processi
in stato \textit{runnable}, quello che ha il valore di priorità dinamica più
basso; con le priorità dinamiche il significato del valore numerico ad esse
associato è infatti invertito, un valore più basso significa una priorità
maggiore. Il fatto che questo valore venga diminuito quando un processo non
viene posto in esecuzione pur essendo pronto, significa che la priorità dei
processi che non ottengono l'uso del processore viene progressivamente
incrementata, così che anche questi alla fine hanno la possibilità di essere
eseguiti.

Sia la dimensione della \textit{time-slice} che il valore di partenza della
priorità dinamica sono determinate dalla cosiddetta \textit{nice} (o
\textit{niceness}) del processo.\footnote{questa è una delle tante proprietà
  che ciascun processo si porta dietro, essa viene ereditata dai processi
  figli e mantenuta attraverso una \func{exec}; fino alla serie 2.4 essa era
  mantenuta nell'omonimo campo \texttt{nice} della \texttt{task\_struct}, con
  la riscrittura dello \textit{scheduler} eseguita nel 2.6 viene mantenuta nel
  campo \texttt{static\_prio} come per le priorità statiche.} L'origine del
nome di questo parametro sta nel fatto che generalmente questo viene usato per
\textsl{diminuire} la priorità di un processo, come misura di cortesia nei
confronti degli altri.  I processi infatti vengono creati dal sistema con un
valore nullo e nessuno è privilegiato rispetto agli altri. Specificando un
valore di \textit{nice} positivo si avrà una \textit{time-slice} più breve ed
un valore di priorità dinamica iniziale più alto, mentre un valore negativo
darà una \textit{time-slice} più lunga ed un valore di priorità dinamica
iniziale più basso.

Esistono diverse funzioni che consentono di indicare un valore di
\textit{nice} di un processo; la più semplice è \funcd{nice}, che opera sul
processo corrente, il suo prototipo è:

\begin{funcproto}{ 
\fhead{unistd.h}
\fdecl{int nice(int inc)}
\fdesc{Aumenta il valore di \textit{nice} del processo corrente.} 
}
{La funzione ritorna il nuovo valore di \textit{nice} in caso di successo e
  $-1$ per un errore, nel qual caso \var{errno} assumerà uno dei valori:
\begin{errlist}
  \item[\errcode{EPERM}] non si ha il permesso di specificare un valore
    di \param{inc} negativo. 
\end{errlist}}
\end{funcproto}

\constbeg{PRIO\_MIN}
\constbeg{PRIO\_MAX}

L'argomento \param{inc} indica l'incremento da effettuare rispetto al valore
di \textit{nice} corrente, che può assumere valori compresi fra
\const{PRIO\_MIN} e \const{PRIO\_MAX}; nel caso di Linux sono fra $-20$ e
$19$,\footnote{in realtà l'intervallo varia a seconda delle versioni di
  kernel, ed è questo a partire dal kernel 1.3.43, anche se oggi si può avere
  anche l'intervallo fra $-20$ e $20$.} ma per \param{inc} si può specificare
un valore qualunque, positivo o negativo, ed il sistema provvederà a troncare
il risultato nell'intervallo consentito. Valori positivi comportano maggiore
\textit{cortesia} e cioè una diminuzione della priorità, valori negativi
comportano invece un aumento della priorità. Con i kernel precedenti il 2.6.12
solo l'amministratore\footnote{o un processo con la \textit{capability}
  \const{CAP\_SYS\_NICE}, vedi sez.~\ref{sec:proc_capabilities}.} può
specificare valori negativi di \param{inc} che permettono di aumentare la
priorità di un processo, a partire da questa versione è consentito anche agli
utenti normali alzare (entro certi limiti, che vedremo in
sez.~\ref{sec:sys_resource_limit}) la priorità dei propri processi.

\constend{PRIO\_MIN}
\constend{PRIO\_MAX}

Gli standard SUSv2 e POSIX.1 prevedono che la funzione ritorni il nuovo valore
di \textit{nice} del processo; tuttavia la \textit{system call} di Linux non
segue questa convenzione e restituisce sempre $0$ in caso di successo e $-1$
in caso di errore; questo perché $-1$ è anche un valore di \textit{nice}
legittimo e questo comporta una confusione con una eventuale condizione di
errore. La \textit{system call} originaria inoltre non consente, se non dotati
di adeguati privilegi, di diminuire un valore di \textit{nice} precedentemente
innalzato.
 
Fino alla \acr{glibc} 2.2.4 la funzione di libreria riportava direttamente il
risultato dalla \textit{system call}, violando lo standard, per cui per
ottenere il nuovo valore occorreva una successiva chiamata alla funzione
\func{getpriority}. A partire dalla \acr{glibc} 2.2.4 \func{nice} è stata
reimplementata e non viene più chiamata la omonima \textit{system call}, con
questa versione viene restituito come valore di ritorno il valore di
\textit{nice}, come richiesto dallo standard.\footnote{questo viene fatto
  chiamando al suo interno \func{setpriority}, che tratteremo a breve.}  In
questo caso l'unico modo per rilevare in maniera affidabile una condizione di
errore è quello di azzerare \var{errno} prima della chiamata della funzione e
verificarne il valore quando \func{nice} restituisce $-1$.

Per leggere il valore di \textit{nice} di un processo occorre usare la
funzione di sistema \funcd{getpriority}, derivata da BSD; il suo prototipo è:

\begin{funcproto}{ 
\fhead{sys/time.h}
\fhead{sys/resource.h}
\fdecl{int getpriority(int which, int who)}
\fdesc{Legge un valore di \textit{nice}.} 
}
{La funzione ritorna $0$ in caso di successo e $-1$ per un errore, nel qual
caso \var{errno} assumerà uno dei valori:
\begin{errlist}
\item[\errcode{EINVAL}] il valore di \param{which} non è uno di quelli
    elencati in tab.~\ref{tab:proc_getpriority}.
\item[\errcode{ESRCH}] non c'è nessun processo che corrisponda ai valori di
  \param{which} e \param{who}.
\end{errlist}}
\end{funcproto}

La funzione permette, a seconda di quanto specificato nell'argomento
\param{which}, di leggere il valore di \textit{nice} o di un processo, o di un
gruppo di processi (vedi sez.~\ref{sec:sess_proc_group}) o di un utente,
indicati con l'argomento \param{who}. Nelle vecchie versioni può essere
necessario includere anche \headfiled{sys/time.h}, questo non è più necessario
con versioni recenti delle librerie, ma è comunque utile per portabilità.

I valori possibili per \param{which}, ed il tipo di valore che occorre usare
in corrispondenza per \param{who}, solo elencati nella legenda di
tab.~\ref{tab:proc_getpriority} insieme ai relativi significati. Usare un
valore nullo per \param{who} indica, a seconda della corrispondente
indicazione usata per \param{which}, il processo, il gruppo di processi o
l'utente correnti.

\begin{table}[htb]
  \centering
  \footnotesize
  \begin{tabular}[c]{|c|c|l|}
    \hline
    \param{which} & \param{who} & \textbf{Significato} \\
    \hline
    \hline
    \constd{PRIO\_PROCESS} & \type{pid\_t} & processo  \\
    \constd{PRIO\_PRGR}    & \type{pid\_t} & \textit{process group} (vedi
                                             sez.~\ref{sec:sess_proc_group})\\
    \constd{PRIO\_USER}    & \type{uid\_t} & utente \\
    \hline
  \end{tabular}
  \caption{Legenda del valore dell'argomento \param{which} e del tipo
    dell'argomento \param{who} delle funzioni \func{getpriority} e
    \func{setpriority} per le tre possibili scelte.}
  \label{tab:proc_getpriority}
\end{table}

In caso di una indicazione che faccia riferimento a più processi, la funzione
restituisce la priorità più alta (cioè il valore più basso) fra quelle dei
processi corrispondenti. Come per \func{nice}, $-1$ è un possibile valore
corretto, per cui di nuovo per poter rilevare una condizione di errore è
necessario cancellare sempre \var{errno} prima della chiamata alla funzione e
quando si ottiene un valore di ritorno uguale a $-1$ per verificare che essa
resti uguale a zero.

Analoga a \func{getpriority} è la funzione di sistema \funcd{setpriority} che
permette di impostare la priorità di uno o più processi; il suo prototipo è:

\begin{funcproto}{ 
\fhead{sys/time.h}
\fhead{sys/resource.h}
\fdecl{int setpriority(int which, int who, int prio)}
\fdesc{Imposta un valore di \textit{nice}.} 
}
{La funzione ritorna $0$ in caso di successo e $-1$ per un errore, nel qual
caso \var{errno} assumerà uno dei valori:
\begin{errlist}
\item[\errcode{EACCES}] si è richiesto un aumento di priorità senza avere
  sufficienti privilegi.
\item[\errcode{EINVAL}] il valore di \param{which} non è uno di quelli
  elencati in tab.~\ref{tab:proc_getpriority}.
\item[\errcode{EPERM}] un processo senza i privilegi di amministratore ha
  cercato di modificare la priorità di un processo di un altro utente.
\item[\errcode{ESRCH}] non c'è nessun processo che corrisponda ai valori di
  \param{which} e \param{who}.
\end{errlist}}
\end{funcproto}

La funzione imposta la priorità dinamica al valore specificato da \param{prio}
per tutti i processi indicati dagli argomenti \param{which} e \param{who}, per
i quali valgono le stesse considerazioni fatte per \func{getpriority} e lo
specchietto di tab.~\ref{tab:proc_getpriority}. 

In questo caso come valore di \param{prio} deve essere specificato il valore
di \textit{nice} da assegnare nell'intervallo fra \const{PRIO\_MIN} ($-20$) e
\const{PRIO\_MAX} ($19$), e non un incremento (positivo o negativo) come nel
caso di \func{nice}. La funzione restituisce il valore di \textit{nice}
assegnato in caso di successo e $-1$ in caso di errore, e come per \func{nice}
anche in questo caso per rilevare un errore occorre sempre porre a zero
\var{errno} prima della chiamata della funzione, essendo $-1$ un valore di
\textit{nice} valido.

Si tenga presente che solo l'amministratore\footnote{o più precisamente un
  processo con la \textit{capability} \const{CAP\_SYS\_NICE}, vedi
  sez.~\ref{sec:proc_capabilities}.} ha la possibilità di modificare
arbitrariamente le priorità di qualunque processo. Un utente normale infatti
può modificare solo la priorità dei suoi processi ed in genere soltanto
diminuirla.  Fino alla versione di kernel 2.6.12 Linux ha seguito le
specifiche dello standard SUSv3, e come per tutti i sistemi derivati da SysV
veniva richiesto che l'\ids{UID} reale o quello effettivo del processo
chiamante corrispondessero all'\ids{UID} reale (e solo a quello) del processo
di cui si intendeva cambiare la priorità. A partire dalla versione 2.6.12 è
stata adottata la semantica in uso presso i sistemi derivati da BSD (SunOS,
Ultrix, *BSD), in cui la corrispondenza può essere anche con l'\ids{UID}
effettivo.

Sempre a partire dal kernel 2.6.12 è divenuto possibile anche per gli utenti
ordinari poter aumentare la priorità dei propri processi specificando un
valore di \param{prio} negativo. Questa operazione non è possibile però in
maniera indiscriminata, ed in particolare può essere effettuata solo
nell'intervallo consentito dal valore del limite \const{RLIMIT\_NICE}
(torneremo su questo in sez.~\ref{sec:sys_resource_limit}).

Infine nonostante i valori siano sempre rimasti gli stessi, il significato del
valore di \textit{nice} è cambiato parecchio nelle progressive riscritture
dello \textit{scheduler} di Linux, ed in particolare a partire dal kernel
2.6.23 l'uso di diversi valori di \textit{nice} ha un impatto molto più forte
nella distribuzione della CPU ai processi. Infatti se viene comunque calcolata
una priorità dinamica per i processi che non ricevono la CPU, così che anche
essi possano essere messi in esecuzione, un alto valore di \textit{nice}
corrisponde comunque ad una \textit{time-slice} molto piccola che non cresce
comunque, per cui un processo a bassa priorità avrà davvero scarse possibilità
di essere eseguito in presenza di processi attivi a priorità più alta.


\subsection{Il meccanismo di \textit{scheduling real-time}}
\label{sec:proc_real_time}

Come spiegato in sez.~\ref{sec:proc_sched} lo standard POSIX.1b ha introdotto
le priorità assolute per permettere la gestione di processi
\textit{real-time}. In realtà nel caso di Linux non si tratta di un vero
\textit{hard real-time}, in quanto in presenza di eventuali interrupt il
kernel interrompe l'esecuzione di un processo, qualsiasi sia la sua
priorità,\footnote{questo a meno che non si siano installate le patch di
  RTLinux, RTAI o Adeos, con i quali è possibile ottenere un sistema
  effettivamente \textit{hard real-time}. In tal caso infatti gli interrupt
  vengono intercettati dall'interfaccia \textit{real-time} (o nel caso di
  Adeos gestiti dalle code del nano-kernel), in modo da poterli controllare
  direttamente qualora ci sia la necessità di avere un processo con priorità
  più elevata di un \textit{interrupt handler}.} mentre con l'incorrere in un
\textit{page fault} si possono avere ritardi non previsti.  Se l'ultimo
problema può essere aggirato attraverso l'uso delle funzioni di controllo
della memoria virtuale (vedi sez.~\ref{sec:proc_mem_lock}), il primo non è
superabile e può comportare ritardi non prevedibili riguardo ai tempi di
esecuzione di qualunque processo.

Nonostante questo, ed in particolare con una serie di miglioramenti che sono
stati introdotti nello sviluppo del kernel,\footnote{in particolare a partire
  dalla versione 2.6.18 sono stati inserite nel kernel una serie di modifiche
  che consentono di avvicinarsi sempre di più ad un vero e proprio sistema
  \textit{real-time} estendendo il concetto di \textit{preemption} alle
  operazioni dello stesso kernel; esistono vari livelli a cui questo può
  essere fatto, ottenibili attivando in fase di compilazione una fra le
  opzioni \texttt{CONFIG\_PREEMPT\_NONE}, \texttt{CONFIG\_PREEMPT\_VOLUNTARY}
  e \texttt{CONFIG\_PREEMPT\_DESKTOP}.} si può arrivare ad una ottima
approssimazione di sistema \textit{real-time} usando le priorità assolute.
Occorre farlo però con molta attenzione: se si dà ad un processo una priorità
assoluta e questo finisce in un loop infinito, nessun altro processo potrà
essere eseguito, ed esso sarà mantenuto in esecuzione permanentemente
assorbendo tutta la CPU e senza nessuna possibilità di riottenere l'accesso al
sistema. Per questo motivo è sempre opportuno, quando si lavora con processi
che usano priorità assolute, tenere attiva una shell cui si sia assegnata la
massima priorità assoluta, in modo da poter essere comunque in grado di
rientrare nel sistema.

Quando c'è un processo con priorità assoluta lo \textit{scheduler} lo metterà
in esecuzione prima di ogni processo normale. In caso di più processi sarà
eseguito per primo quello con priorità assoluta più alta. Quando ci sono più
processi con la stessa priorità assoluta questi vengono tenuti in una coda e
tocca al kernel decidere quale deve essere eseguito.  Il meccanismo con cui
vengono gestiti questi processi dipende dalla politica di \textit{scheduling}
che si è scelta; lo standard ne prevede due:
\begin{basedescript}{\desclabelwidth{1.2cm}\desclabelstyle{\nextlinelabel}}
\item[\textit{First In First Out} (FIFO)] Il processo viene eseguito
  fintanto che non cede volontariamente la CPU (con la funzione
  \func{sched\_yield}), si blocca, finisce o viene interrotto da un processo a
  priorità più alta. Se il processo viene interrotto da uno a priorità più
  alta esso resterà in cima alla lista e sarà il primo ad essere eseguito
  quando i processi a priorità più alta diverranno inattivi. Se invece lo si
  blocca volontariamente sarà posto in coda alla lista (ed altri processi con
  la stessa priorità potranno essere eseguiti).
\item[\textit{Round Robin} (RR)] Il comportamento è del tutto analogo a quello
  precedente, con la sola differenza che ciascun processo viene eseguito al
  massimo per un certo periodo di tempo (la cosiddetta \textit{time-slice})
  dopo di che viene automaticamente posto in fondo alla coda dei processi con
  la stessa priorità. In questo modo si ha comunque una esecuzione a turno di
  tutti i processi, da cui il nome della politica. Solo i processi con la
  stessa priorità ed in stato \textit{runnable} entrano nel
  \textsl{girotondo}.
\end{basedescript}

Lo standard POSIX.1-2001 prevede una funzione che consenta sia di modificare
le politiche di \textit{scheduling}, passando da \textit{real-time} a
ordinarie o viceversa, che di specificare, in caso di politiche
\textit{real-time}, la eventuale priorità statica; la funzione di sistema è
\funcd{sched\_setscheduler} ed il suo prototipo è:

\begin{funcproto}{ 
\fhead{sched.h}
\fdecl{int sched\_setscheduler(pid\_t pid, int policy, const struct
  sched\_param *p)}
\fdesc{Imposta priorità e politica di \textit{scheduling}.} 
}
{La funzione ritorna $0$ in caso di successo e $-1$ per un errore, nel qual
caso \var{errno} assumerà uno dei valori:
\begin{errlist}
    \item[\errcode{EINVAL}] il valore di \param{policy} non esiste o il
      valore di \param{p} non è valido per la politica scelta.
    \item[\errcode{EPERM}] il processo non ha i privilegi per attivare la
      politica richiesta.
    \item[\errcode{ESRCH}] il processo \param{pid} non esiste.
    \end{errlist}}
\end{funcproto}

La funzione esegue l'impostazione per il processo specificato dall'argomento
\param{pid}; un valore nullo di questo argomento esegue l'impostazione per il
processo corrente.  La politica di \textit{scheduling} è specificata
dall'argomento \param{policy} i cui possibili valori sono riportati in
tab.~\ref{tab:proc_sched_policy}; la parte alta della tabella indica le
politiche \textit{real-time}, quella bassa le politiche ordinarie. Un valore
negativo per \param{policy} mantiene la politica di \textit{scheduling}
corrente.

\begin{table}[htb]
  \centering
  \footnotesize
  \begin{tabular}[c]{|l|p{6cm}|}
    \hline
    \textbf{Politica}  & \textbf{Significato} \\
    \hline
    \hline
    \constd{SCHED\_FIFO} & \textit{Scheduling real-time} con politica
                           \textit{FIFO}. \\
    \constd{SCHED\_RR}   & \textit{Scheduling real-time} con politica
                           \textit{Round Robin}. \\ 
    \hline
    \constd{SCHED\_OTHER}& \textit{Scheduling} ordinario.\\
    \constd{SCHED\_BATCH}& \textit{Scheduling} ordinario con l'assunzione
                           ulteriore di lavoro \textit{CPU
                           intensive} (dal kernel 2.6.16).\\ 
    \constd{SCHED\_IDLE} & \textit{Scheduling} di priorità estremamente
                           bassa (dal kernel 2.6.23).\\
    \hline
  \end{tabular}
  \caption{Valori dell'argomento \param{policy} per la funzione
    \func{sched\_setscheduler}.}
  \label{tab:proc_sched_policy}
\end{table}

% TODO Aggiungere SCHED_DEADLINE, sulla nuova politica di scheduling aggiunta
% con il kernel 3.14, vedi anche Documentation/scheduler/sched-deadline.txt e
% http://lwn.net/Articles/575497/
% vedi anche man 7 sched, man sched_setattr

Con le versioni più recenti del kernel sono state introdotte anche delle
varianti sulla politica di \textit{scheduling} tradizionale per alcuni carichi
di lavoro specifici, queste due nuove politiche sono specifiche di Linux e non
devono essere usate se si vogliono scrivere programmi portabili.

La politica \const{SCHED\_BATCH} è una variante della politica ordinaria con
la sola differenza che i processi ad essa soggetti non ottengono, nel calcolo
delle priorità dinamiche fatto dallo \textit{scheduler}, il cosiddetto bonus
di interattività che mira a favorire i processi che si svegliano dallo stato
di \textit{sleep}.\footnote{cosa che accade con grande frequenza per i
  processi interattivi, dato che essi sono per la maggior parte del tempo in
  attesa di dati in ingresso da parte dell'utente.} La si usa pertanto, come
indica il nome, per processi che usano molta CPU (come programmi di calcolo)
che in questo modo, pur non perdendo il loro valore di \textit{nice}, sono
leggermente sfavoriti rispetto ai processi interattivi che devono rispondere a
dei dati in ingresso.

La politica \const{SCHED\_IDLE} invece è una politica dedicata ai processi che
si desidera siano eseguiti con la più bassa priorità possibile, ancora più
bassa di un processo con il minimo valore di \textit{nice}. In sostanza la si
può utilizzare per processi che devono essere eseguiti se non c'è niente altro
da fare. Va comunque sottolineato che anche un processo \const{SCHED\_IDLE}
avrà comunque una sua possibilità di utilizzo della CPU, sia pure in
percentuale molto bassa.

Qualora si sia richiesta una politica \textit{real-time} il valore della
priorità statica viene impostato attraverso la struttura
\struct{sched\_param}, riportata in fig.~\ref{fig:sig_sched_param}, il cui
solo campo attualmente definito è \var{sched\_priority}. Il campo deve
contenere il valore della priorità statica da assegnare al processo; lo
standard prevede che questo debba essere assegnato all'interno di un
intervallo fra un massimo ed un minimo che nel caso di Linux sono
rispettivamente 1 e 99.

\begin{figure}[!htb]
  \footnotesize \centering
  \begin{minipage}[c]{0.5\textwidth}
    \includestruct{listati/sched_param.c}
  \end{minipage} 
  \normalsize 
  \caption{La struttura \structd{sched\_param}.} 
  \label{fig:sig_sched_param}
\end{figure}

I processi con politica di \textit{scheduling} ordinaria devono sempre
specificare un valore nullo di \var{sched\_priority} altrimenti si avrà un
errore \errcode{EINVAL}, questo valore infatti non ha niente a che vedere con
la priorità dinamica determinata dal valore di \textit{nice}, che deve essere
impostato con le funzioni viste in precedenza.

Lo standard POSIX.1b prevede che l'intervallo dei valori delle priorità
statiche possa essere ottenuto con le funzioni di sistema
\funcd{sched\_get\_priority\_max} e \funcd{sched\_get\_priority\_min}, i cui
prototipi sono:

\begin{funcproto}{ 
\fhead{sched.h}
\fdecl{int sched\_get\_priority\_max(int policy)}
\fdesc{Legge il valore massimo di una priorità statica.} 
\fdecl{int sched\_get\_priority\_min(int policy)}
\fdesc{Legge il valore minimo di una priorità statica.} 
}
{Le funzioni ritornano il valore della priorità in caso di successo e $-1$ per
  un errore, nel qual caso \var{errno} assumerà il valore:
\begin{errlist}
\item[\errcode{EINVAL}] il valore di \param{policy} non è valido.
\end{errlist}}
\end{funcproto}

Le funzioni ritornano rispettivamente il valore massimo e minimo usabile per
la priorità statica di una delle politiche di \textit{scheduling}
\textit{real-time} indicata dall'argomento \param{policy}.

Si tenga presente che quando si imposta una politica di \textit{scheduling}
real-time per un processo o se ne cambia la priorità statica questo viene
messo in cima alla lista dei processi con la stessa priorità; questo comporta
che verrà eseguito subito, interrompendo eventuali altri processi con la
stessa priorità in quel momento in esecuzione.

Il kernel mantiene i processi con la stessa priorità assoluta in una lista, ed
esegue sempre il primo della lista, mentre un nuovo processo che torna in
stato \textit{runnable} viene sempre inserito in coda alla lista. Se la
politica scelta è \const{SCHED\_FIFO} quando il processo viene eseguito viene
automaticamente rimesso in coda alla lista, e la sua esecuzione continua
fintanto che non viene bloccato da una richiesta di I/O, o non rilascia
volontariamente la CPU (in tal caso, tornando nello stato \textit{runnable}
sarà in coda alla lista); l'esecuzione viene ripresa subito solo nel caso che
esso sia stato interrotto da un processo a priorità più alta.

Solo un processo con i privilegi di amministratore\footnote{più precisamente
  con la capacità \const{CAP\_SYS\_NICE}, vedi
  sez.~\ref{sec:proc_capabilities}.} può impostare senza restrizioni priorità
assolute diverse da zero o politiche \const{SCHED\_FIFO} e
\const{SCHED\_RR}. Un utente normale può modificare solo le priorità di
processi che gli appartengono; è cioè richiesto che l'\ids{UID} effettivo del
processo chiamante corrisponda all'\ids{UID} reale o effettivo del processo
indicato con \param{pid}.

Fino al kernel 2.6.12 gli utenti normali non potevano impostare politiche
\textit{real-time} o modificare la eventuale priorità statica di un loro
processo. A partire da questa versione è divenuto possibile anche per gli
utenti normali usare politiche \textit{real-time} fintanto che la priorità
assoluta che si vuole impostare è inferiore al limite \const{RLIMIT\_RTPRIO}
(vedi sez.~\ref{sec:sys_resource_limit}) ad essi assegnato. 

Unica eccezione a questa possibilità sono i processi \const{SCHED\_IDLE}, che
non possono cambiare politica di \textit{scheduling} indipendentemente dal
valore di \const{RLIMIT\_RTPRIO}. Inoltre, in caso di processo già sottoposto
ad una politica \textit{real-time}, un utente può sempre, indipendentemente
dal valore di \const{RLIMIT\_RTPRIO}, diminuirne la priorità o portarlo ad una
politica ordinaria.

Se si intende operare solo sulla priorità statica di un processo si possono
usare le due funzioni di sistema \funcd{sched\_setparam} e
\funcd{sched\_getparam} che consentono rispettivamente di impostarne e
leggerne il valore, i loro prototipi sono:

\begin{funcproto}{
\fhead{sched.h}
\fdecl{int sched\_setparam(pid\_t pid, const struct sched\_param *param)}
\fdesc{Imposta la priorità statica di un processo.} 
\fdecl{int sched\_getparam(pid\_t pid, struct sched\_param *param)}
\fdesc{Legge la priorità statica di un processo.} 
}
{Le funzioni ritornano $0$ in caso di successo e $-1$ per un errore, nel qual
caso \var{errno} assumerà uno dei valori:
\begin{errlist}
\item[\errcode{EINVAL}] il valore di \param{param} non ha senso per la
  politica usata dal processo.
\item[\errcode{EPERM}] non si hanno privilegi sufficienti per eseguire
  l'operazione.
\item[\errcode{ESRCH}] il processo \param{pid} non esiste.
\end{errlist}}
\end{funcproto}

Le funzioni richiedono di indicare nell'argomento \param{pid} il processo su
cui operare e usano l'argomento \param{param} per mantenere il valore della
priorità dinamica. Questo è ancora una struttura \struct{sched\_param} ed
assume gli stessi valori già visti per \func{sched\_setscheduler}.

L'uso di \func{sched\_setparam}, compresi i controlli di accesso che vi si
applicano, è del tutto equivalente a quello di \func{sched\_setscheduler} con
argomento \param{policy} uguale a $-1$. Come per \func{sched\_setscheduler}
specificando $0$ come valore dell'argomento \param{pid} si opera sul processo
corrente. Benché la funzione sia utilizzabile anche con processi sottoposti a
politica ordinaria essa ha senso soltanto per quelli \textit{real-time}, dato
che per i primi la priorità statica può essere soltanto nulla.  La
disponibilità di entrambe le funzioni può essere verificata controllando la
macro \macro{\_POSIX\_PRIORITY\_SCHEDULING} che è definita nell'\textit{header
  file} \headfiled{sched.h}.

Se invece si vuole sapere quale è politica di \textit{scheduling} di un
processo si può usare la funzione di sistema \funcd{sched\_getscheduler}, il
cui prototipo è:

\begin{funcproto}{ 
\fhead{sched.h}
\fdecl{int sched\_getscheduler(pid\_t pid)}
\fdesc{Legge la politica di \textit{scheduling}.} 
}
{La funzione ritorna la politica di \textit{scheduling}  in caso di successo e
  $-1$ per un errore, nel qual caso \var{errno} assumerà uno dei valori:
\begin{errlist}
    \item[\errcode{EPERM}] non si hanno privilegi sufficienti per eseguire
      l'operazione.
    \item[\errcode{ESRCH}] il processo \param{pid} non esiste.
\end{errlist}}
\end{funcproto}

La funzione restituisce il valore, secondo quanto elencato in
tab.~\ref{tab:proc_sched_policy}, della politica di \textit{scheduling} per il
processo specificato dall'argomento \param{pid}, se questo è nullo viene
restituito il valore relativo al processo chiamante.

L'ultima funzione di sistema che permette di leggere le informazioni relative
ai processi real-time è \funcd{sched\_rr\_get\_interval}, che permette di
ottenere la lunghezza della \textit{time-slice} usata dalla politica
\textit{round robin}; il suo prototipo è:

\begin{funcproto}{ 
\fhead{sched.h}
\fdecl{int sched\_rr\_get\_interval(pid\_t pid, struct timespec *tp)}
\fdesc{Legge la durata della \textit{time-slice} per lo \textit{scheduling}
  \textit{round robin}.}  
}
{La funzione ritorna $0$ in caso di successo e $-1$ per un errore, nel qual
caso \var{errno} assumerà uno dei valori:
\begin{errlist}
\item[\errcode{EINVAL}] l'argomento \param{pid} non è valido. 
\item[\errcode{ENOSYS}] la \textit{system call} non è presente (solo per
  kernel arcaici).
\item[\errcode{ESRCH}] il processo \param{pid} non esiste.
\end{errlist}
ed inoltre anche \errval{EFAULT} nel suo significato generico.}
\end{funcproto}

La funzione restituisce nell'argomento \param{tp} come una struttura
\struct{timespec}, (la cui definizione si può trovare in
fig.~\ref{fig:sys_timeval_struct}) il valore dell'intervallo di tempo usato
per la politica \textit{round robin} dal processo indicato da \ids{PID}. Il
valore dipende dalla versione del kernel, a lungo infatti questo intervallo di
tempo era prefissato e non modificabile ad un valore di 150 millisecondi,
restituito indipendentemente dal \ids{PID} indicato. 

Con kernel recenti però è possibile ottenere una variazione della
\textit{time-slice}, modificando il valore di \textit{nice} del processo
(anche se questo non incide assolutamente sulla priorità statica) che come
accennato in precedenza modifica il valore assegnato alla \textit{time-slice}
di un processo ordinario, che però viene usato anche dai processi
\textit{real-time}.

Come accennato ogni processo può rilasciare volontariamente la CPU in modo da
consentire agli altri processi di essere eseguiti; la funzione di sistema che
consente di fare tutto questo è \funcd{sched\_yield}, il cui prototipo è:

\begin{funcproto}{ 
\fhead{sched.h}
\fdecl{int sched\_yield(void)}
\fdesc{Rilascia volontariamente l'esecuzione.} 
}
{La funzione ritorna $0$ in caso di successo e teoricamente $-1$ per un
  errore, ma su Linux ha sempre successo.}
\end{funcproto}


Questa funzione ha un utilizzo effettivo soltanto quando si usa lo
\textit{scheduling} \textit{real-time}, e serve a far sì che il processo
corrente rilasci la CPU, in modo da essere rimesso in coda alla lista dei
processi con la stessa priorità per permettere ad un altro di essere eseguito;
se però il processo è l'unico ad essere presente sulla coda l'esecuzione non
sarà interrotta. In genere usano questa funzione i processi con politica
\const{SCHED\_FIFO}, per permettere l'esecuzione degli altri processi con pari
priorità quando la sezione più urgente è finita.

La funzione può essere utilizzata anche con processi che usano lo
\textit{scheduling} ordinario, ma in questo caso il comportamento non è ben
definito, e dipende dall'implementazione. Fino al kernel 2.6.23 questo
comportava che i processi venissero messi in fondo alla coda di quelli attivi,
con la possibilità di essere rimessi in esecuzione entro breve tempo, con
l'introduzione del \textit{Completely Fair Scheduler} questo comportamento è
cambiato ed un processo che chiama la funzione viene inserito nella lista dei
processi inattivi, con un tempo molto maggiore.\footnote{è comunque possibile
  ripristinare un comportamento analogo al precedente scrivendo il valore 1
  nel file \sysctlfiled{kernel/sched\_compat\_yield}.}

L'uso delle funzione nella programmazione ordinaria può essere utile e
migliorare le prestazioni generali del sistema quando si è appena rilasciata
una risorsa contesa con altri processi, e si vuole dare agli altri una
possibilità di approfittarne mettendoli in esecuzione, ma chiamarla senza
necessità, specie se questo avviene ripetutamente all'interno di un qualche
ciclo, può avere invece un forte impatto negativo per la generazione di
\textit{context switch} inutili.


\subsection{Il controllo dello \textit{scheduler} per i sistemi
  multiprocessore}
\label{sec:proc_sched_multiprocess}

\index{effetto~ping-pong|(} 

Con il supporto dei sistemi multiprocessore sono state introdotte delle
funzioni che permettono di controllare in maniera più dettagliata la scelta di
quale processore utilizzare per eseguire un certo programma. Uno dei problemi
che si pongono nei sistemi multiprocessore è infatti quello del cosiddetto
\textsl{effetto ping-pong}. Può accadere cioè che lo \textit{scheduler},
quando riavvia un processo precedentemente interrotto scegliendo il primo
processore disponibile, lo faccia eseguire da un processore diverso rispetto a
quello su cui era stato eseguito in precedenza. Se il processo passa da un
processore all'altro in questo modo, cosa che avveniva abbastanza di frequente
con i kernel della seria 2.4.x, si ha l'effetto ping-pong.

Questo tipo di comportamento può generare dei seri problemi di prestazioni;
infatti tutti i processori moderni utilizzano una memoria interna (la
\textit{cache}) contenente i dati più usati, che permette di evitare di
eseguire un accesso (molto più lento) alla memoria principale sulla scheda
madre.  Chiaramente un processo sarà favorito se i suoi dati sono nella cache
del processore, ma è ovvio che questo può essere vero solo per un processore
alla volta, perché in presenza di più copie degli stessi dati su più
processori, non si potrebbe determinare quale di questi ha la versione dei
dati aggiornata rispetto alla memoria principale.

Questo comporta che quando un processore inserisce un dato nella sua cache,
tutti gli altri processori che hanno lo stesso dato devono invalidarlo, e
questa operazione è molto costosa in termini di prestazioni. Il problema
diventa serio quando si verifica l'effetto ping-pong, in tal caso infatti un
processo \textsl{rimbalza} continuamente da un processore all'altro e si ha
una continua invalidazione della cache, che non diventa mai disponibile.

\itindbeg{CPU~affinity}

Per ovviare a questo tipo di problemi è nato il concetto di \textsl{affinità
  di processore} (o \textit{CPU affinity}); la possibilità cioè di far sì che
un processo possa essere assegnato per l'esecuzione sempre allo stesso
processore. Lo \textit{scheduler} dei kernel della serie 2.4.x aveva una
scarsa \textit{CPU affinity}, e l'effetto ping-pong era comune; con il nuovo
\textit{scheduler} dei kernel della 2.6.x questo problema è stato risolto ed
esso cerca di mantenere il più possibile ciascun processo sullo stesso
processore.

\index{effetto~ping-pong|)}

In certi casi però resta l'esigenza di poter essere sicuri che un processo sia
sempre eseguito dallo stesso processore,\footnote{quella che viene detta
  \textit{hard CPU affinity}, in contrasto con quella fornita dallo
  \textit{scheduler}, detta \textit{soft CPU affinity}, che di norma indica
  solo una preferenza, non un requisito assoluto.} e per poter risolvere
questo tipo di problematiche nei nuovi kernel\footnote{le due \textit{system
    call} per la gestione della \textit{CPU affinity} sono state introdotte
  nel kernel 2.5.8, e le corrispondenti funzioni di sistema nella
  \textsl{glibc} 2.3.} è stata introdotta l'opportuna infrastruttura ed una
nuova \textit{system call} che permette di impostare su quali processori far
eseguire un determinato processo attraverso una \textsl{maschera di
  affinità}. La corrispondente funzione di sistema è
\funcd{sched\_setaffinity} ed il suo prototipo è:

\index{insieme~di~processori|(}

\begin{funcproto}{ 
\fhead{sched.h}
\fdecl{int sched\_setaffinity(pid\_t pid, size\_t setsize, 
  cpu\_set\_t *mask)}
\fdesc{Imposta la maschera di affinità di un processo.} 
}
{La funzione ritorna $0$ in caso di successo e $-1$ per un errore, nel qual
caso \var{errno} assumerà uno dei valori:
\begin{errlist}
\item[\errcode{EINVAL}] il valore di \param{mask} contiene riferimenti a
  processori non esistenti nel sistema o a cui non è consentito l'accesso.
\item[\errcode{EPERM}] il processo non ha i privilegi sufficienti per
  eseguire l'operazione.
\item[\errcode{ESRCH}] il processo \param{pid} non esiste.
\end{errlist}
ed inoltre anche \errval{EFAULT} nel suo significato generico.}
\end{funcproto}

Questa funzione e la corrispondente \func{sched\_getaffinity} hanno una storia
abbastanza complessa, la sottostante \textit{system call} infatti prevede
l'uso di due soli argomenti (per il pid e l'indicazione della maschera dei
processori), che corrispondono al fatto che l'implementazione effettiva usa
una semplice maschera binaria. Quando le funzioni vennero incluse nella
\acr{glibc} assunsero invece un prototipo simile a quello mostrato però con il
secondo argomento di tipo \ctyp{unsigned int}. A complicare la cosa si
aggiunge il fatto che nella versione 2.3.3 della \acr{glibc} detto argomento
venne stato eliminato, per poi essere ripristinato nella versione 2.3.4 nella
forma attuale.\footnote{pertanto se la vostra pagina di manuale non è
  aggiornata, o usate quella particolare versione della \acr{glibc}, potrete
  trovare indicazioni diverse, il prototipo illustrato è quello riportato
  nella versione corrente (maggio 2008) delle pagine di manuale e
  corrispondente alla definizione presente in \headfile{sched.h}.}

La funzione imposta, con l'uso del valore contenuto all'indirizzo
\param{mask}, l'insieme dei processori sui quali deve essere eseguito il
processo identificato tramite il valore passato in \param{pid}. Come in
precedenza il valore nullo di \param{pid} indica il processo corrente.  Per
poter utilizzare questa funzione sono richiesti i privilegi di amministratore
(è necessaria la capacità \const{CAP\_SYS\_NICE}) altrimenti essa fallirà con
un errore di \errcode{EPERM}. Una volta impostata una maschera di affinità,
questa viene ereditata attraverso una \func{fork}, in questo modo diventa
possibile legare automaticamente un gruppo di processi ad un singolo
processore.

Nell'uso comune, almeno con i kernel successivi alla serie 2.6.x, utilizzare
questa funzione non è necessario, in quanto è lo \textit{scheduler} stesso che
provvede a mantenere al meglio l'affinità di processore. Esistono però
esigenze particolari, ad esempio quando un processo (o un gruppo di processi)
è utilizzato per un compito importante (ad esempio per applicazioni
\textit{real-time} o la cui risposta è critica) e si vuole la massima
velocità; con questa interfaccia diventa possibile selezionare gruppi di
processori utilizzabili in maniera esclusiva.  Lo stesso dicasi quando
l'accesso a certe risorse (memoria o periferiche) può avere un costo diverso a
seconda del processore, come avviene nelle architetture NUMA
(\textit{Non-Uniform Memory Access}).

Infine se un gruppo di processi accede alle stesse risorse condivise (ad
esempio una applicazione con più \textit{thread}) può avere senso usare lo
stesso processore in modo da sfruttare meglio l'uso della sua cache; questo
ovviamente riduce i benefici di un sistema multiprocessore nell'esecuzione
contemporanea dei \textit{thread}, ma in certi casi (quando i \textit{thread}
sono inerentemente serializzati nell'accesso ad una risorsa) possono esserci
sufficienti vantaggi nell'evitare la perdita della cache da rendere
conveniente l'uso dell'affinità di processore.

Dato che il numero di processori può variare a seconda delle architetture, per
semplificare l'uso dell'argomento \param{mask} la \acr{glibc} ha introdotto un
apposito dato di tipo, \typed{cpu\_set\_t},\footnote{questa è una estensione
  specifica della \acr{glibc}, da attivare definendo la macro
  \macro{\_GNU\_SOURCE}, non esiste infatti una standardizzazione per questo
  tipo di interfaccia e POSIX al momento non prevede nulla al riguardo.} che
permette di identificare un insieme di processori. Il dato è normalmente una
maschera binaria: nei casi più comuni potrebbe bastare un intero a 32 bit, in
cui ogni bit corrisponde ad un processore, ma oggi esistono architetture in
cui questo numero può non essere sufficiente, e per questo è stato creato
questo tipo opaco e una interfaccia di gestione che permette di usare a basso
livello un tipo di dato qualunque rendendosi indipendenti dal numero di bit e
dalla loro disposizione.  Per questo le funzioni di libreria richiedono che
oltre all'insieme di processori si indichi anche la dimensione dello stesso
con l'argomento \param{setsize}, per il quale, se non si usa l'allocazione
dinamica che vedremo a breve, è in genere sufficiente passare il valore
\code{sizeof(cpu\_set\_t)}.

L'interfaccia di gestione degli insiemi di processori, oltre alla definizione
del tipo \type{cpu\_set\_t}, prevede una serie di macro di preprocessore per
la manipolazione degli stessi. Quelle di base, che consentono rispettivamente
di svuotare un insieme, di aggiungere o togliere un processore o di verificare
se esso è già presente in un insieme, sono le seguenti:

{\centering
\vspace{3pt}
\begin{funcbox}{ 
\fhead{sched.h}
\fdecl{void \macrod{CPU\_ZERO}(cpu\_set\_t *set)}
\fdesc{Inizializza un insieme di processori vuoto \param{set}.} 
\fdecl{void \macrod{CPU\_SET}(int cpu, cpu\_set\_t *set)}
\fdesc{Inserisce il processore \param{cpu} nell'insieme di
  processori \param{set}.}  
\fdecl{void \macrod{CPU\_CLR}(int cpu, cpu\_set\_t *set)}
\fdesc{Rimuove il processore \param{cpu} nell'insieme di
  processori \param{set}.}  
\fdecl{int \macrod{CPU\_ISSET}(int cpu, cpu\_set\_t *set)}
\fdesc{Controlla se il processore \param{cpu} è nell'insieme di
  processori \param{set}.}  
}
\end{funcbox}}

Queste macro che sono ispirate dalle analoghe usate per gli insiemi di
\textit{file descriptor} (vedi sez.~\ref{sec:file_select}) e sono state
introdotte con la versione 2.3.3 della \acr{glibc}. Tutte richiedono che si
specifichi il numero di una CPU nell'argomento \param{cpu}, ed un insieme su
cui operare. L'unica che ritorna un risultato è \macro{CPU\_ISSET}, che
restituisce un intero da usare come valore logico (zero se la CPU non è
presente, diverso da zero se è presente).

\itindbeg{side~effects}
Si tenga presente che trattandosi di macro l'argomento \param{cpu} può essere
valutato più volte. Questo significa ad esempio che non si può usare al suo
posto una funzione o un'altra macro, altrimenti queste verrebbero eseguite più
volte; l'argomento cioè non deve avere \textsl{effetti collaterali} (in gergo
 \textit{side effects}).\footnote{nel linguaggio C si
  parla appunto di \textit{side effects} quando si usano istruzioni la cui
  valutazione comporta effetti al di fuori dell'istruzione stessa, come il
  caso indicato in cui si passa una funzione ad una macro che usa l'argomento
  al suo interno più volte, o si scrivono espressioni come \code{a=a++} in cui
  non è chiaro se prima avvenga l'incremento e poi l'assegnazione, ed il cui
  risultato dipende dall'implementazione del compilatore.}
\itindend{side~effects}


Le CPU sono numerate da zero (che indica la prima disponibile) fino ad un
numero massimo che dipende dall'architettura hardware. La costante
\constd{CPU\_SETSIZE} indica il numero massimo di processori che possono far
parte di un insieme (al momento vale sempre 1024), e costituisce un limite
massimo al valore dell'argomento \param{cpu}.  Dalla versione 2.6 della
\acr{glibc} alle precedenti macro è stata aggiunta, per contare il numero di
processori in un insieme, l'ulteriore:

{\centering
\vspace{3pt}
\begin{funcbox}{ 
\fhead{sched.h}
\fdecl{int \macrod{CPU\_COUNT}(cpu\_set\_t *set)}
\fdesc{Conta il numero di processori presenti nell'insieme \param{set}.} 
}
\end{funcbox}}

A partire dalla versione 2.7 della \acr{glibc} sono state introdotte altre
macro che consentono ulteriori manipolazioni, in particolare si possono
compiere delle operazioni logiche sugli insiemi di processori con:

{\centering
\vspace{3pt}
\begin{funcbox}{ 
\fhead{sched.h}
\fdecl{void \macrod{CPU\_AND}(cpu\_set\_t *destset, cpu\_set\_t *srcset1, cpu\_set\_t *srcset2)}
\fdesc{Esegue l'AND logico di due insiemi di processori.} 
\fdecl{void \macrod{CPU\_OR}(cpu\_set\_t *destset, cpu\_set\_t *srcset1, cpu\_set\_t *srcset2)}
\fdesc{Esegue l'OR logico di due insiemi di processori.} 
\fdecl{void \macrod{CPU\_XOR}(cpu\_set\_t *destset, cpu\_set\_t *srcset1, cpu\_set\_t *srcset2)}
\fdesc{Esegue lo XOR logico di due insiemi di processori.} 
\fdecl{int \macrod{CPU\_EQUAL}(cpu\_set\_t *set1, cpu\_set\_t *set2)}
\fdesc{Verifica se due insiemi di processori sono uguali.} 
}
\end{funcbox}}

Le prime tre macro richiedono due insiemi di partenza, \param{srcset1} e
\param{srcset2} e forniscono in un terzo insieme \param{destset} (che può
essere anche lo stesso di uno dei precedenti) il risultato della rispettiva
operazione logica sui contenuti degli stessi. In sostanza con \macro{CPU\_AND}
si otterrà come risultato l'insieme che contiene le CPU presenti in entrambi
gli insiemi di partenza, con \macro{CPU\_OR} l'insieme che contiene le CPU
presenti in uno qualunque dei due insiemi di partenza, e con \macro{CPU\_XOR}
l'insieme che contiene le CPU presenti in uno solo dei due insiemi di
partenza. Infine \macro{CPU\_EQUAL} confronta due insiemi ed è l'unica che
restituisce un intero, da usare come valore logico che indica se sono identici
o meno.

Inoltre, sempre a partire dalla versione 2.7 della \acr{glibc}, è stata
introdotta la possibilità di una allocazione dinamica degli insiemi di
processori, per poterli avere di dimensioni corrispondenti al numero di CPU
effettivamente in gioco, senza dover fare riferimento necessariamente alla
precedente dimensione preimpostata di 1024. Per questo motivo sono state
definite tre ulteriori macro, che consentono rispettivamente di allocare,
disallocare ed ottenere la dimensione in byte di un insieme di processori:

{\centering
\vspace{3pt}
\begin{funcbox}{ 
\fhead{sched.h}
\fdecl{cpu\_set\_t * \macrod{CPU\_ALLOC}(num\_cpus)}
\fdesc{Alloca dinamicamente un insieme di processori di dimensione voluta.} 
\fdecl{void \macrod{CPU\_FREE}(cpu\_set\_t *set)}
\fdesc{Disalloca un insieme di processori allocato dinamicamente.} 
\fdecl{size\_t \macrod{CPU\_ALLOC\_SIZE}(num\_cpus)}
\fdesc{Ritorna la dimensione di un insieme di processori allocato dinamicamente.} 
}
\end{funcbox}}

La prima macro, \macro{CPU\_ALLOC}, restituisce il puntatore ad un insieme di
processori in grado di contenere almeno \param{num\_cpus} che viene allocato
dinamicamente. Ogni insieme così allocato dovrà essere disallocato con
\macro{CPU\_FREE} passandogli un puntatore ottenuto da una precedente
\macro{CPU\_ALLOC}. La terza macro, \macro{CPU\_ALLOC\_SIZE}, consente di
ottenere la dimensione in byte di un insieme allocato dinamicamente che
contenga \param{num\_cpus} processori.

Dato che le dimensioni effettive possono essere diverse le macro di gestione e
manipolazione che abbiamo trattato in precedenza non si applicano agli insiemi
allocati dinamicamente, per i quali dovranno sono state definite altrettante
macro equivalenti contraddistinte dal suffisso \texttt{\_S}, che effettuano le
stesse operazioni, ma richiedono in più un argomento
aggiuntivo \param{setsize} che deve essere assegnato al valore ottenuto con
\macro{CPU\_ALLOC\_SIZE}. Questo stesso valore deve essere usato per l'omonimo
argomento delle funzioni \func{sched\_setaffinity} o \func{sched\_getaffinity}
quando si vuole usare per l'argomento che indica la maschera di affinità un
insieme di processori allocato dinamicamente.

\index{insieme~di~processori|)}

A meno di non aver utilizzato \func{sched\_setaffinity}, in condizioni
ordinarie la maschera di affinità di un processo è preimpostata dal sistema in
modo che esso possa essere eseguito su qualunque processore. Se ne può
comunque ottenere il valore corrente usando la funzione di sistema
\funcd{sched\_getaffinity}, il cui prototipo è:

\begin{funcproto}{ 
\fhead{sched.h}
\fdecl{int sched\_getaffinity (pid\_t pid, size\_t setsize, 
  cpu\_set\_t *mask)}
\fdesc{Legge la maschera di affinità di un processo.} 
}
{La funzione ritorna $0$ in caso di successo e $-1$ per un errore, nel qual
caso \var{errno} assumerà uno dei valori:
\begin{errlist}
\item[\errcode{EINVAL}] \param{setsize} è più piccolo delle dimensioni
  della maschera di affinità usata dal kernel.
\item[\errcode{ESRCH}] il processo \param{pid} non esiste.
\end{errlist}
ed inoltre anche \errval{EFAULT} nel suo significato generico.}
\end{funcproto}

La funzione restituirà all'indirizzo specificato da \param{mask} il valore
della maschera di affinità del processo indicato dall'argomento \param{pid}
(al solito un valore nullo indica il processo corrente) così da poterla
riutilizzare per una successiva reimpostazione.

È chiaro che queste funzioni per la gestione dell'affinità hanno significato
soltanto su un sistema multiprocessore, esse possono comunque essere
utilizzate anche in un sistema con un processore singolo, nel qual caso però
non avranno alcun risultato effettivo.

\itindend{scheduler}
\itindend{CPU~affinity}


\subsection{Le priorità per le operazioni di I/O}
\label{sec:io_priority}

A lungo l'unica priorità usata per i processi è stata quella relativa
all'assegnazione dell'uso del processore. Ma il processore non è l'unica
risorsa che i processi devono contendersi, un'altra, altrettanto importante
per le prestazioni, è quella dell'accesso a disco. Per questo motivo nello
sviluppo del kernel sono stati introdotti diversi \textit{I/O scheduler} in
grado di distribuire in maniera opportuna questa risorsa ai vari processi.

Fino al kernel 2.6.17 era possibile soltanto differenziare le politiche
generali di gestione, scegliendo di usare un diverso \textit{I/O scheduler}. A
partire da questa versione, con l'introduzione dello \textit{scheduler} CFQ
(\textit{Completely Fair Queuing}) è divenuto possibile, qualora si usi questo
\textit{scheduler}, impostare anche delle diverse priorità di accesso per i
singoli processi.\footnote{al momento (kernel 2.6.31), le priorità di I/O sono
  disponibili soltanto per questo \textit{scheduler}.}

La scelta di uno \textit{scheduler} di I/O si può fare in maniera generica per
tutto il sistema all'avvio del kernel con il parametro di avvio
\texttt{elevator},\footnote{per la trattazione dei parametri di avvio del
  kernel si rimanda al solito alla sez.~5.3 di \cite{AGL}.} cui assegnare il
nome dello \textit{scheduler}, ma se ne può anche indicare uno specifico per
l'accesso al singolo disco scrivendo nel file
\texttt{/sys/block/\textit{<dev>}/queue/scheduler} (dove
\texttt{\textit{<dev>}} è il nome del dispositivo associato al disco).

Gli \textit{scheduler} disponibili sono mostrati dal contenuto dello stesso
file che riporta fra parentesi quadre quello attivo, il default in tutti i
kernel recenti è proprio il \texttt{cfq}, nome con cui si indica appunto lo
\textit{scheduler} CFQ, che supporta le priorità. Per i dettagli sulle
caratteristiche specifiche degli altri \textit{scheduler}, la cui discussione
attiene a problematiche di ambito sistemistico, si consulti la documentazione
nella directory \texttt{Documentation/block/} dei sorgenti del kernel.

Una volta che si sia impostato lo \textit{scheduler} CFQ ci sono due
specifiche \textit{system call}, specifiche di Linux, che consentono di
leggere ed impostare le priorità di I/O.\footnote{se usate in corrispondenza
  ad uno \textit{scheduler} diverso il loro utilizzo non avrà alcun effetto.}
Dato che non esiste una interfaccia diretta nella \acr{glibc} per queste due
funzioni\footnote{almeno al momento della scrittura di questa sezione, con la
  versione 2.11 della \acr{glibc}.} occorrerà invocarle tramite la funzione
\func{syscall} (come illustrato in sez.~\ref{sec:proc_syscall}). Le due
\textit{system call} sono \funcd{ioprio\_get} ed \funcd{ioprio\_set}; i
rispettivi prototipi sono:

\begin{funcproto}{ 
\fhead{linux/ioprio.h}
\fdecl{int ioprio\_get(int which, int who)}
\fdesc{Legge la priorità di I/O di un processo.} 
\fdecl{int ioprio\_set(int which, int who, int ioprio)}
\fdesc{Imposta la priorità di I/O di un processo.} 
}
{Le funzioni ritornano rispettivamente un intero positivo o 0 in caso di
  successo e $-1$ per un errore, nel qual caso \var{errno} assumerà uno dei
  valori:
\begin{errlist}
\item[\errcode{EINVAL}] i valori di \param{which} o di \param{ioprio} non
  sono validi. 
\item[\errcode{EPERM}] non si hanno i privilegi per eseguire
  l'impostazione (solo per \func{ioprio\_set}). 
\item[\errcode{ESRCH}] non esiste un processo corrispondente alle indicazioni.
\end{errlist}}
\end{funcproto}

Le funzioni leggono o impostano la priorità di I/O sulla base dell'indicazione
dei due argomenti \param{which} e \param{who} che hanno lo stesso significato
già visto per gli omonimi argomenti di \func{getpriority} e
\func{setpriority}. Anche in questo caso si deve specificare il valore di
\param{which} tramite le opportune costanti riportate in
tab.~\ref{tab:ioprio_args} che consentono di indicare un singolo processo, i
processi di un \textit{process group} (vedi sez.~\ref{sec:sess_proc_group}) o
tutti i processi di un utente.

\begin{table}[htb]
  \centering
  \footnotesize
  \begin{tabular}[c]{|c|c|l|}
    \hline
    \param{which} & \param{who} & \textbf{Significato} \\
    \hline
    \hline
    \constd{IPRIO\_WHO\_PROCESS} & \type{pid\_t} & processo\\
    \constd{IPRIO\_WHO\_PRGR}    & \type{pid\_t} & \textit{process group}\\ 
    \constd{IPRIO\_WHO\_USER}    & \type{uid\_t} & utente\\
    \hline
  \end{tabular}
  \caption{Legenda del valore dell'argomento \param{which} e del tipo
    dell'argomento \param{who} delle funzioni \func{ioprio\_get} e
    \func{ioprio\_set} per le tre possibili scelte.}
  \label{tab:ioprio_args}
\end{table}

In caso di successo \func{ioprio\_get} restituisce un intero positivo che
esprime il valore della priorità di I/O, questo valore è una maschera binaria
composta da due parti, una che esprime la \textsl{classe} di
\textit{scheduling} di I/O del processo, l'altra che esprime, quando la classe
di \textit{scheduling} lo prevede, la priorità del processo all'interno della
classe stessa. Questo stesso formato viene utilizzato per indicare il valore
della priorità da impostare con l'argomento \param{ioprio} di
\func{ioprio\_set}.

\begin{table}[htb]
  \centering
  \footnotesize
  \begin{tabular}[c]{|l|p{8cm}|}
    \hline
    \textbf{Macro} & \textbf{Significato}\\
    \hline
    \hline
    \macrod{IOPRIO\_PRIO\_CLASS}\texttt{(\textit{value})}
                                & Dato il valore di una priorità come
                                  restituito da \func{ioprio\_get} estrae il
                                  valore della classe.\\
    \macrod{IOPRIO\_PRIO\_DATA}\texttt{(\textit{value})}
                                & Dato il valore di una priorità come
                                  restituito da \func{ioprio\_get} estrae il
                                  valore della priorità.\\
    \macrod{IOPRIO\_PRIO\_VALUE}\texttt{(\textit{class},\textit{prio})}
                                & Dato un valore di priorità ed una classe
                                  ottiene il valore numerico da passare a
                                  \func{ioprio\_set}.\\
    \hline
  \end{tabular}
  \caption{Le macro per la gestione dei valori numerici .}
  \label{tab:IOsched_class_macro}
\end{table}


Per la gestione dei valori che esprimono le priorità di I/O sono state
definite delle opportune macro di preprocessore, riportate in
tab.~\ref{tab:IOsched_class_macro}. I valori delle priorità si ottengono o si
impostano usando queste macro.  Le prime due si usano con il valore restituito
da \func{ioprio\_get} e per ottenere rispettivamente la classe di
\textit{scheduling}\footnote{restituita dalla macro con i valori di
  tab.~\ref{tab:IOsched_class}.} e l'eventuale valore della priorità. La terza
macro viene invece usata per creare un valore di priorità da usare come
argomento di \func{ioprio\_set} per eseguire una impostazione.

Le classi di \textit{scheduling} previste dallo \textit{scheduler} CFQ sono
tre, e ricalcano tre diverse modalità di distribuzione delle risorse, analoghe
a quelle già adottate anche nel funzionamento dello \textit{scheduler} del
processore. Ciascuna di esse è identificata tramite una opportuna costante,
secondo quanto riportato in tab.~\ref{tab:IOsched_class}.

\begin{table}[htb]
  \centering
  \footnotesize
  \begin{tabular}[c]{|l|l|}
    \hline
    \textbf{Classe}  & \textbf{Significato} \\
    \hline
    \hline
    \const{IOPRIO\_CLASS\_RT}  & \textit{Scheduling} di I/O
                                 \textit{real-time}.\\  
    \const{IOPRIO\_CLASS\_BE}  & \textit{Scheduling} di I/O ordinario.\\ 
    \const{IOPRIO\_CLASS\_IDLE}& \textit{Scheduling} di I/O di priorità
                                  minima.\\
    \hline
  \end{tabular}
  \caption{Costanti che identificano le classi di \textit{scheduling} di I/O.}
  \label{tab:IOsched_class}
\end{table}

La classe di priorità più bassa è \constd{IOPRIO\_CLASS\_IDLE}; i processi in
questa classe riescono ad accedere a disco soltanto quando nessun altro
processo richiede l'accesso. Occorre pertanto usarla con molta attenzione,
perché un processo in questa classe può venire completamente bloccato quando
ci sono altri processi in una qualunque delle altre due classi che stanno
accedendo al disco. Quando si usa questa classe non ha senso indicare un
valore di priorità, dato che in questo caso non esiste nessuna gerarchia e la
priorità è identica, la minima possibile, per tutti i processi che la usano.

La seconda classe di priorità di I/O è \constd{IOPRIO\_CLASS\_BE} (il nome sta
per \textit{best-effort}), che è quella usata ordinariamente da tutti
processi. In questo caso esistono priorità diverse che consentono di
assegnazione di una maggiore banda passante nell'accesso a disco ad un
processo rispetto agli altri, con meccanismo simile a quello dei valori di
\textit{nice} in cui si evita che un processo a priorità più alta possa
bloccare indefinitamente quelli a priorità più bassa. In questo caso però le
diverse priorità sono soltanto otto, indicate da un valore numerico fra 0 e 7
e come per \textit{nice} anche in questo caso un valore più basso indica una
priorità maggiore.

Infine la classe di priorità di I/O \textit{real-time}
\constd{IOPRIO\_CLASS\_RT} ricalca le omonime priorità di processore: un
processo in questa classe ha sempre la precedenza nell'accesso a disco
rispetto a tutti i processi delle altre classi e di un processo nella stessa
classe ma con priorità inferiore, ed è pertanto in grado di bloccare
completamente tutti gli altri. Anche in questo caso ci sono 8 priorità diverse
con un valore numerico fra 0 e 7, con una priorità più elevata per valori più
bassi.

In generale nel funzionamento ordinario la priorità di I/O di un processo
viene impostata in maniera automatica nella classe \const{IOPRIO\_CLASS\_BE}
con un valore ottenuto a partire dal corrispondente valore di \textit{nice}
tramite la formula: $\mathtt{\mathit{prio}}=(\mathtt{\mathit{nice}}+20)/5$. Un
utente ordinario può modificare con \func{ioprio\_set} soltanto le priorità
dei processi che gli appartengono,\footnote{per la modifica delle priorità di
  altri processi occorrono privilegi amministrativi, ed in particolare la
  capacità \const{CAP\_SYS\_NICE} (vedi sez.~\ref{sec:proc_capabilities}).}
cioè quelli il cui \ids{UID} reale corrisponde all'\ids{UID} reale o effettivo
del chiamante. Data la possibilità di ottenere un blocco totale del sistema,
solo l'amministratore\footnote{o un processo con la capacità
  \const{CAP\_SYS\_ADMIN} (vedi sez.~\ref{sec:proc_capabilities}).} può
impostare un processo ad una priorità di I/O nella classe
\const{IOPRIO\_CLASS\_RT}, lo stesso privilegio era richiesto anche per la
classe \const{IOPRIO\_CLASS\_IDLE} fino al kernel 2.6.24, ma dato che in
questo caso non ci sono effetti sugli altri processi questo limite è stato
rimosso a partire dal kernel 2.6.25.

%TODO verificare http://lwn.net/Articles/355987/


\section{Problematiche di programmazione \textit{multitasking}}
\label{sec:proc_multi_prog}

Benché i processi siano strutturati in modo da apparire il più possibile come
indipendenti l'uno dall'altro, nella programmazione in un sistema
\textit{multitasking} occorre tenere conto di una serie di problematiche che
normalmente non esistono quando si ha a che fare con un sistema in cui viene
eseguito un solo programma alla volta.

Per questo motivo, essendo questo argomento di carattere generale, ci è parso
opportuno introdurre sinteticamente queste problematiche, che ritroveremo a
più riprese in capitoli successivi, in questa sezione conclusiva del capitolo
in cui abbiamo affrontato la gestione dei processi, sottolineando come esse
diventino cogenti quando invece si usano i \textit{thread}.


\subsection{Le operazioni atomiche}
\label{sec:proc_atom_oper}

La nozione di \textsl{operazione atomica} deriva dal significato greco della
parola atomo, cioè indivisibile; si dice infatti che un'operazione è atomica
quando si ha la certezza che, qualora essa venga effettuata, tutti i passaggi
che devono essere compiuti per realizzarla verranno eseguiti senza possibilità
di interruzione in una fase intermedia.

In un ambiente \textit{multitasking} il concetto è essenziale, dato che un
processo può essere interrotto in qualunque momento dal kernel che mette in
esecuzione un altro processo o dalla ricezione di un segnale. Occorre pertanto
essere accorti nei confronti delle possibili \textit{race condition} (vedi
sez.~\ref{sec:proc_race_cond}) derivanti da operazioni interrotte in una fase
in cui non erano ancora state completate.

Nel caso dell'interazione fra processi la situazione è molto più semplice, ed
occorre preoccuparsi della atomicità delle operazioni solo quando si ha a che
fare con meccanismi di intercomunicazione (che esamineremo in dettaglio in
cap.~\ref{cha:IPC}) o nelle operazioni con i file (vedremo alcuni esempi in
sez.~\ref{sec:file_shared_access}). In questi casi in genere l'uso delle
appropriate funzioni di libreria per compiere le operazioni necessarie è
garanzia sufficiente di atomicità in quanto le \textit{system call} con cui
esse sono realizzate non possono essere interrotte (o subire interferenze
pericolose) da altri processi.

Nel caso dei segnali invece la situazione è molto più delicata, in quanto lo
stesso processo, e pure alcune \textit{system call}, possono essere interrotti
in qualunque momento, e le operazioni di un eventuale \textit{signal handler}
sono compiute nello stesso spazio di indirizzi del processo. Per questo, anche
il solo accesso o l'assegnazione di una variabile possono non essere più
operazioni atomiche (torneremo su questi aspetti in
sez.~\ref{sec:sig_adv_control}).

Qualora invece si usino i \textit{thread}, in cui lo spazio degli indirizzi è
condiviso, il problema è sempre presente, perché qualunque \textit{thread} può
interromperne un altro in qualunque momento e l'atomicità di qualunque
operazione è messa in discussione, per cui l'assenza di eventuali \textit{race
  condition} (vedi sez.~\ref{sec:proc_race_cond}) deve essere sempre
verificata nei minimi dettagli.

In questo caso il sistema provvede un tipo di dato, il \typed{sig\_atomic\_t},
il cui accesso è assicurato essere atomico.  In pratica comunque si può
assumere che, in ogni piattaforma su cui è implementato Linux, il tipo
\ctyp{int}, gli altri interi di dimensione inferiore ed i puntatori sono
atomici. Non è affatto detto che lo stesso valga per interi di dimensioni
maggiori (in cui l'accesso può comportare più istruzioni in assembler) o per
le strutture di dati. In tutti questi casi è anche opportuno marcare come
\dirct{volatile} le variabili che possono essere interessate ad accesso
condiviso, onde evitare problemi con le ottimizzazioni del codice.



\subsection{Le \textit{race condition} ed i \textit{deadlock}}
\label{sec:proc_race_cond}

\itindbeg{race~condition}

Si definiscono \textit{race condition} tutte quelle situazioni in cui processi
diversi operano su una risorsa comune, ed in cui il risultato viene a
dipendere dall'ordine in cui essi effettuano le loro operazioni. Il caso
tipico è quello di un'operazione che viene eseguita da un processo in più
passi, e può essere compromessa dall'intervento di un altro processo che
accede alla stessa risorsa quando ancora non tutti i passi sono stati
completati.

Dato che in un sistema \textit{multitasking} ogni processo può essere
interrotto in qualunque momento per farne subentrare un altro in esecuzione,
niente può assicurare un preciso ordine di esecuzione fra processi diversi o
che una sezione di un programma possa essere eseguita senza interruzioni da
parte di altri. Queste situazioni comportano pertanto errori estremamente
subdoli e difficili da tracciare, in quanto nella maggior parte dei casi tutto
funzionerà regolarmente, e solo occasionalmente si avranno degli errori.

Per questo occorre essere ben consapevoli di queste problematiche, e del fatto
che l'unico modo per evitarle è quello di riconoscerle come tali e prendere
gli adeguati provvedimenti per far sì che non si verifichino. Casi tipici di
\textit{race condition} si hanno quando diversi processi accedono allo stesso
file, o nell'accesso a meccanismi di intercomunicazione come la memoria
condivisa. 

\index{sezione~critica|(}

In questi casi, se non si dispone della possibilità di eseguire atomicamente
le operazioni necessarie, occorre che quelle parti di codice in cui si
compiono le operazioni sulle risorse condivise, quelle che in genere vengono
denominate ``\textsl{sezioni critiche}'' del programma, siano opportunamente
protette da meccanismi di sincronizzazione (vedremo alcune problematiche di
questo tipo in cap.~\ref{cha:IPC}).

\index{sezione~critica|)}

Nel caso dei \textit{thread} invece la situazione è molto più delicata e
sostanzialmente qualunque accesso in memoria (a buffer, variabili o altro) può
essere soggetto a \textit{race condition} dato potrebbe essere interrotto in
qualunque momento da un altro \textit{thread}. In tal caso occorre pianificare
con estrema attenzione l'uso delle variabili ed utilizzare i vari meccanismi
di sincronizzazione che anche in questo caso sono disponibili (torneremo su
queste problematiche di questo tipo in cap.~\ref{sec:pthread_sync})

\itindbeg{deadlock} 

Un caso particolare di \textit{race condition} sono poi i cosiddetti
\textit{deadlock} (traducibile in \textsl{condizione di stallo}), che
particolarmente gravi in quanto comportano spesso il blocco completo di un
servizio, e non il fallimento di una singola operazione. Per definizione un
\textit{deadlock} è una situazione in cui due o più processi non sono più in
grado di proseguire perché ciascuno aspetta il risultato di una operazione che
dovrebbe essere eseguita dall'altro.

L'esempio tipico di una situazione che può condurre ad un
\textit{deadlock} è quello in cui un flag di
``\textsl{occupazione}'' viene rilasciato da un evento asincrono (come un
segnale o un altro processo) fra il momento in cui lo si è controllato
(trovandolo occupato) e la successiva operazione di attesa per lo sblocco. In
questo caso, dato che l'evento di sblocco del flag è avvenuto senza che ce ne
accorgessimo proprio fra il controllo e la messa in attesa, quest'ultima
diventerà perpetua (da cui il nome di \textit{deadlock}).

In tutti questi casi è di fondamentale importanza il concetto di atomicità
visto in sez.~\ref{sec:proc_atom_oper}; questi problemi infatti possono essere
risolti soltanto assicurandosi, quando essa sia richiesta, che sia possibile
eseguire in maniera atomica le operazioni necessarie.

\itindend{race~condition}
\itindend{deadlock}


\subsection{Le funzioni rientranti}
\label{sec:proc_reentrant}

\index{funzioni!rientranti|(}

Si dice \textsl{rientrante} una funzione che può essere interrotta in
qualunque punto della sua esecuzione ed essere chiamata una seconda volta da
un altro \textit{thread} di esecuzione senza che questo comporti nessun
problema nell'esecuzione della stessa. La problematica è comune nella
programmazione con i \textit{thread}, ma si hanno gli stessi problemi quando
si vogliono chiamare delle funzioni all'interno dei gestori dei segnali.

Fintanto che una funzione opera soltanto con le variabili locali è rientrante;
queste infatti vengono allocate nello \textit{stack}, ed un'altra invocazione
non fa altro che allocarne un'altra copia. Una funzione può non essere
rientrante quando opera su memoria che non è nello \textit{stack}.  Ad esempio
una funzione non è mai rientrante se usa una variabile globale o statica.

Nel caso invece la funzione operi su un oggetto allocato dinamicamente, la
cosa viene a dipendere da come avvengono le operazioni: se l'oggetto è creato
ogni volta e ritornato indietro la funzione può essere rientrante, se invece
esso viene individuato dalla funzione stessa due chiamate alla stessa funzione
potranno interferire quando entrambe faranno riferimento allo stesso oggetto.
Allo stesso modo una funzione può non essere rientrante se usa e modifica un
oggetto che le viene fornito dal chiamante: due chiamate possono interferire
se viene passato lo stesso oggetto; in tutti questi casi occorre molta cura da
parte del programmatore.

In genere le funzioni di libreria non sono rientranti, molte di esse ad
esempio utilizzano variabili statiche, la \acr{glibc} però mette a
disposizione due macro di compilatore, \macro{\_REENTRANT} e
\macro{\_THREAD\_SAFE}, la cui definizione attiva le versioni rientranti di
varie funzioni di libreria, che sono identificate aggiungendo il suffisso
\code{\_r} al nome della versione normale.

\index{funzioni!rientranti|)}


% LocalWords:  multitasking like VMS child process identifier pid sez shell fig
% LocalWords:  parent kernel init pstree keventd kswapd table struct linux call
% LocalWords:  nell'header scheduler system interrupt timer HZ asm Hertz clock
% LocalWords:  l'alpha tick fork wait waitpid exit exec image glibc int pgid ps
% LocalWords:  sid thread Ingo Molnar ppid getpid getppid sys unistd LD threads
% LocalWords:  void tempnam pathname sibling cap errno EAGAIN ENOMEM context
% LocalWords:  stack read only copy write tab client spawn forktest sleep PATH
% LocalWords:  source LIBRARY scheduling race condition printf descriptor dup
% LocalWords:  close group session tms lock vfork execve BSD stream main abort
% LocalWords:  SIGABRT SIGCHLD SIGHUP foreground SIGCONT termination signal ANY
% LocalWords:  handler kill EINTR POSIX options WNOHANG ECHILD option WUNTRACED
% LocalWords:  dump bits rusage getrusage heap const filename argv envp EACCES
% LocalWords:  filesystem noexec EPERM suid sgid root nosuid ENOEXEC ENOENT ELF
% LocalWords:  ETXTBSY EINVAL ELIBBAD BIG EFAULT EIO ENAMETOOLONG ELOOP ENOTDIR
% LocalWords:  ENFILE EMFILE argc execl path execv execle execlp execvp vector
% LocalWords:  list environ NULL umask utime cutime ustime fcntl linker Posix
% LocalWords:  opendir libc interpreter FreeBSD capabilities mandatory access
% LocalWords:  control MAC SELinux security modules LSM superuser uid gid saved
% LocalWords:  effective euid egid dell' fsuid fsgid getuid geteuid getgid SVr
% LocalWords:  getegid IDS NFS setuid setgid all' logout utmp screen xterm TODO
% LocalWords:  setreuid setregid FIXME ruid rgid seteuid setegid setresuid size
% LocalWords:  setresgid getresuid getresgid value result argument setfsuid DAC
% LocalWords:  setfsgid NGROUPS sysconf getgroups getgrouplist groups ngroups
% LocalWords:  setgroups initgroups patch LIDS CHOWN OVERRIDE Discrectionary 
% LocalWords:  SEARCH chattr sticky NOATIME socket domain immutable append mmap
% LocalWords:  broadcast multicast multicasting memory locking mlock mlockall
% LocalWords:  shmctl ioperm iopl chroot ptrace accounting swap reboot hangup
% LocalWords:  vhangup mknod lease permitted inherited inheritable bounding AND
% LocalWords:  capability capget capset header ESRCH undef version obj clear PT
% LocalWords:  pag ssize length proc capgetp preemptive cache runnable idled
% LocalWords:  SIGSTOP soft slice nice niceness counter which SC switch side
% LocalWords:  getpriority who setpriority RTLinux RTAI Adeos fault FIFO  COUNT
% LocalWords:  yield Robin setscheduler policy param OTHER priority setparam to
% LocalWords:  min getparam getscheduler interval robin ENOSYS fifo ping long
% LocalWords:  affinity setaffinity unsigned mask cpu NUMA CLR ISSET SETSIZE RR
% LocalWords:  getaffinity assembler deadlock REENTRANT SAFE tgz MYPGRP l'OR rr
% LocalWords:  WIFEXITED WEXITSTATUS WIFSIGNALED WTERMSIG WCOREDUMP WIFSTOPPED
% LocalWords:  WSTOPSIG opt char INTERP arg SIG IGN DFL mascheck grp FOWNER RAW
% LocalWords:  FSETID SETPCAP BIND SERVICE ADMIN PACKET IPC OWNER MODULE RAWIO
% LocalWords:  PACCT RESOURCE TTY CONFIG SETFCAP hdrp datap libcap lcap text tp
% LocalWords:  get ncap caps CapInh CapPrm fffffeff CapEff getcap STAT dall'I
% LocalWords:  inc PRIO SUSv PRGR prio SysV SunOS Ultrix sched timespec len sig
% LocalWords:  cpusetsize cpuset atomic tickless redirezione WCONTINUED stopped
% LocalWords:  waitid NOCLDSTOP ENOCHLD WIFCONTINUED ifdef endif idtype siginfo
% LocalWords:  infop ALL WEXITED WSTOPPED WNOWAIT signo CLD EXITED KILLED page
% LocalWords:  CONTINUED sources forking Spawned successfully executing exiting
% LocalWords:  next cat for COMMAND pts bash defunct TRAPPED DUMPED PR effects
% LocalWords:  SIGKILL static RLIMIT preemption PREEMPT VOLUNTARY IDLE RTPRIO
% LocalWords:  completely fair compat uniform CFQ queuing elevator dev cfq RT
% LocalWords:  documentation block syscall ioprio IPRIO CLASS class best effort
% LocalWords:  refresh semop dnotify MADV DONTFORK prctl WCLONE WALL big mount
% LocalWords:  WNOTHREAD DUMPABLE KEEPCAPS IRIX CAPBSET endianness endian flags
% LocalWords:  little PPC PowerPC FPEMU NOPRINT SIGFPE FPEXC point FP SW malloc
% LocalWords:  exception EXC ENABLE OVF overflow UND underflow RES INV DISABLED
% LocalWords:  NONRECOV ASYNC KEEP securebits NAME NUL PDEATHSIG SECCOMP VM FS
% LocalWords:  secure computing sigreturn TIMING STATISTICAL TSC MCE conditions
% LocalWords:  timestamp Stamp SIGSEGV UNALIGN SIGBUS MCEERR AO failure early
% LocalWords:  namespace vsyscall SETTID FILES NEWIPC NEWNET NEWNS NEWPID ptid
% LocalWords:  NEWUTS SETTLS SIGHAND SYSVSEM UNTRACED tls ctid CLEARTID panic
% LocalWords:  loader EISDIR SIGTRAP uninterrutible killable EQUAL sizeof XOR
% LocalWords:  destset srcset ALLOC num cpus setsize emacs pager getty TID
%  LocalWords:  reaper SUBREAPER Library futex klogd named rpc statd NPROC
%  LocalWords:  EACCESS EBADF EBUSY ENXIO EOPNOTSUPP DISABLE tracer Yama
 
%%% Local Variables: 
%%% mode: latex
%%% TeX-master: "gapil"
%%% End: 

%% filedir.tex
%%
%% Copyright (C) 2000-2018 Simone Piccardi.  Permission is granted to
%% copy, distribute and/or modify this document under the terms of the GNU Free
%% Documentation License, Version 1.1 or any later version published by the
%% Free Software Foundation; with the Invariant Sections being "Un preambolo",
%% with no Front-Cover Texts, and with no Back-Cover Texts.  A copy of the
%% license is included in the section entitled "GNU Free Documentation
%% License".
%%

\chapter{La gestione di file e directory}
\label{cha:files_and_dirs}

In questo capitolo tratteremo in dettaglio le modalità con cui si gestiscono
file e directory, iniziando da un approfondimento dell'architettura del
sistema illustrata a grandi linee in sez.~\ref{sec:file_arch_overview} ed
illustrando le principali caratteristiche di un filesystem e le interfacce
che consentono di controllarne il montaggio e lo smontaggio. 

Esamineremo poi le funzioni di libreria che si usano per copiare, spostare e
cambiare i nomi di file e directory e l'interfaccia che permette la
manipolazione dei loro attributi. Tratteremo inoltre la struttura di base del
sistema delle protezioni e del controllo dell'accesso ai file e le successive
estensioni (\textit{Extended Attributes}, ACL, quote disco). Tutto quello che
riguarda invece la gestione dell'I/O sui file è lasciato al capitolo
successivo.



\section{L'architettura della gestione dei file}
\label{sec:file_arch_func}

In questa sezione tratteremo con maggiori dettagli rispetto a quanto visto in
sez.~\ref{sec:file_arch_overview} il \textit{Virtual File System} di Linux e
come il kernel può gestire diversi tipi di filesystem, descrivendo prima le
caratteristiche generali di un filesystem di un sistema unix-like, per poi
fare una panoramica sul filesystem tradizionalmente più usato con Linux,
l'\acr{ext2} ed i suoi successori.


\subsection{Il funzionamento del \textit{Virtual File System} di Linux}
\label{sec:file_vfs_work}

% NOTE articolo interessante:
% http://www.ibm.com/developerworks/linux/library/l-virtual-filesystem-switch/index.html?ca=dgr-lnxw97Linux-VFSdth-LXdW&S_TACT=105AGX59&S_CMP=GRlnxw97

\itindbeg{Virtual~File~System~(VFS)}

Come illustrato brevemente in sez.~\ref{sec:file_arch_overview} in Linux il
concetto di \textit{everything is a file} è stato implementato attraverso il
\textit{Virtual File System}, la cui struttura generale è illustrata in
fig.~\ref{fig:file_VFS_scheme}.  Il VFS definisce un insieme di funzioni che
tutti i filesystem devono implementare per l'accesso ai file che contengono e
l'interfaccia che consente di eseguire l'I/O sui file, che questi siano di
dati o dispositivi. 

\itindbeg{inode}

L'interfaccia fornita dal VFS comprende in sostanza tutte le funzioni che
riguardano i file, le operazioni implementate dal VFS sono realizzate con una
astrazione che prevede quattro tipi di oggetti strettamente correlati: i
filesystem, le \textit{dentry}, gli \textit{inode} ed i file. A questi oggetti
corrispondono una serie di apposite strutture definite dal kernel che
contengono come campi le funzioni di gestione e realizzano l'infrastruttura
del VFS. L'interfaccia è molto complessa, ne faremo pertanto una trattazione
estremamente semplificata che consenta di comprenderne i principi
di funzionamento.

Il VFS usa una tabella mantenuta dal kernel che contiene il nome di ciascun
filesystem supportato, quando si vuole inserire il supporto di un nuovo
filesystem tutto quello che occorre è chiamare la funzione
\code{register\_filesystem} passando come argomento la struttura
\kstruct{file\_system\_type} (la cui definizione è riportata in
fig.~\ref{fig:kstruct_file_system_type}) relativa a quel filesystem. Questa
verrà inserita nella tabella, ed il nuovo filesystem comparirà in
\procfile{/proc/filesystems}.

\begin{figure}[!htb]
  \footnotesize \centering
  \begin{minipage}[c]{0.80\textwidth}
    \includestruct{listati/file_system_type.h}
  \end{minipage}
  \normalsize 
  \caption{Estratto della struttura \kstructd{file\_system\_type} usata dal
    VFS (da \texttt{include/linux/fs.h}).}
  \label{fig:kstruct_file_system_type}
\end{figure}

La struttura \kstruct{file\_system\_type}, oltre ad una serie di dati interni,
come il nome del tipo di filesystem nel campo \var{name},\footnote{quello che
  viene riportato in \procfile{/proc/filesystems} e che viene usato come
  valore del parametro dell'opzione \texttt{-t} del comando \texttt{mount} che
  indica il tipo di filesystem.}  contiene i riferimenti alle funzioni di base
che consentono l'utilizzo di quel filesystem. In particolare la funzione
\code{mount} del quarto campo è quella che verrà invocata tutte le volte che
si dovrà effettuare il montaggio di un filesystem di quel tipo. Per ogni nuovo
filesystem si dovrà allocare una di queste strutture ed inizializzare i
relativi campi con i dati specifici di quel filesystem, ed in particolare si
dovrà creare anche la relativa versione della funzione \code{mount}.

\itindbeg{pathname}
\itindbeg{pathname~resolution}

Come illustrato in fig.~\ref{fig:kstruct_file_system_type} questa funzione
restituisce una \textit{dentry}, abbreviazione che sta per \textit{directory
  entry}. Le \textit{dentry} sono gli oggetti che il kernel usa per eseguire
la \textit{pathname resolution}, ciascuna di esse corrisponde ad un
\textit{pathname} e contiene il riferimento ad un \textit{inode}, che come
vedremo a breve è l'oggetto usato dal kernel per identificare un
file.\footnote{in questo caso si parla di file come di un qualunque oggetto
  generico che sta sul filesystem e non dell'oggetto file del VFS cui
  accennavamo prima.} La \textit{dentry} ottenuta dalla chiamata alla funzione
\code{mount} sarà inserita in corrispondenza al \textit{pathname} della
directory in cui il filesystem è stato montato.

% NOTA: struct dentry è dichiarata in include/linux/dcache.h

Le \textit{dentry} sono oggetti del VFS che vivono esclusivamente in memoria,
nella cosiddetta \textit{directory entry cache} (spesso chiamata in breve
\textit{dcache}). Ogni volta che una \textit{system call} specifica un
\textit{pathname} viene effettuata una ricerca nella \textit{dcache} per
ottenere immediatamente la \textit{dentry} corrispondente,\footnote{il buon
  funzionamento della \textit{dcache} è in effetti di una delle parti più
  critiche per le prestazioni del sistema.} che a sua volta ci darà, tramite
l'\textit{inode}, il riferimento al file.

Dato che normalmente non è possibile mantenere nella \textit{dcache} le
informazioni relative a tutto l'albero dei file la procedura della
\textit{pathname resolution} richiede un meccanismo con cui riempire gli
eventuali vuoti. Il meccanismo prevede che tutte le volte che si arriva ad una
\textit{dentry} mancante venga invocata la funzione \texttt{lookup}
dell'\textit{inode} associato alla \textit{dentry} precedente nella
risoluzione del \textit{pathname},\footnote{che a questo punto è una
  directory, per cui si può cercare al suo interno il nome di un file.} il cui
scopo è risolvere il nome mancante e fornire la sua \textit{dentry} che a
questo punto verrà inserita nella cache.

Dato che tutte le volte che si monta un filesystem la funzione \texttt{mount}
(vedi sez.~\ref{sec:filesystem_mounting}) della corrispondente
\kstruct{file\_system\_type} inserisce la \textit{dentry} iniziale nel
\textit{mount point} dello stesso, si avrà comunque un punto di
partenza. Inoltre essendo questa \textit{dentry} relativa a quel tipo di
filesystem essa farà riferimento ad un \textit{inode} di quel filesystem, e
come vedremo questo farà sì che venga eseguita una \texttt{lookup} adatta per
effettuare la risoluzione dei nomi per quel filesystem.

\itindend{pathname}
\itindend{pathname~resolution}

% Un secondo effetto della chiamata funzione \texttt{mount} di
% \kstruct{file\_system\_type} è quello di allocare una struttura
% \kstruct{super\_block} per ciascuna istanza montata, che contiene le
% informazioni generali di un qualunque filesystem montato, come le opzioni di
% montaggio, le dimensioni dei blocchi, quando il filesystem è stato montato
% ecc. Fra queste però viene pure inserta, nel campo \var{s\_op}, una ulteriore
% struttura \kstruct{super\_operations}, il cui contenuto sono i puntatori
% alle funzioni di gestione di un filesystem, anche inizializzata in modo da
% utilizzare le versioni specifiche di quel filesystem.

L'oggetto più importante per il funzionamento del VFS è probabilmente
l'\textit{inode}, ma con questo nome si può fare riferimento a due cose
diverse.  La prima è la struttura su disco (su cui torneremo anche in
sez.~\ref{sec:file_filesystem}) che fa parte della organizzazione dei dati
realizzata dal filesystem e che contiene le informazioni relative alle
proprietà (i cosiddetti \textsl{metadati}) di ogni oggetto presente su di esso
(si intende al solito uno qualunque dei tipi di file di
tab.~\ref{tab:file_file_types}).

La seconda è la corrispondente struttura \kstruct{inode}, della cui
definizione si è riportato un estratto in
fig.~\ref{fig:kstruct_inode}.\footnote{l'estratto fa riferimento alla versione
  del kernel 2.6.37.} Questa struttura viene mantenuta in memoria ed è a
questa che facevamo riferimento quando parlavamo dell'\textit{inode} associato
a ciascuna \textit{dentry}. Nella struttura in memoria sono presenti gli
stessi \textsl{metadati} memorizzati su disco, che vengono letti quando questa
struttura viene allocata e trascritti all'indietro se modificati.

\begin{figure}[!htb]
  \footnotesize \centering
  \begin{minipage}[c]{0.8\textwidth}
    \includestruct{listati/inode.h}
  \end{minipage}
  \normalsize 
  \caption{Estratto della struttura \kstructd{inode} del kernel (da
    \texttt{include/linux/fs.h}).}
  \label{fig:kstruct_inode}
\end{figure}

Il fatto che la struttura \kstruct{inode} sia mantenuta in memoria,
direttamente associata ad una \textit{dentry}, rende sostanzialmente immediate
le operazioni che devono semplicemente effettuare un accesso ai dati in essa
contenuti: è così ad esempio che viene realizzata la \textit{system call}
\func{stat} che vedremo in sez.~\ref{sec:file_stat}. Rispetto ai dati salvati
sul disco questa struttura contiene però anche quanto necessario alla
implementazione del VFS, ed in particolare è importante il campo \var{i\_op}
che, come illustrato in fig.~\ref{fig:kstruct_inode}, contiene il puntatore ad
una struttura di tipo \kstruct{inode\_operation}, la cui definizione si può
trovare nel file \texttt{include/kernel/fs.h} dei sorgenti del kernel.

Questa struttura non è altro che una tabella di funzioni, ogni suo membro cioè
è un puntatore ad una funzione e, come suggerisce il nome della struttura
stessa, queste funzioni sono quelle che definiscono le operazioni che il VFS
può compiere su un \textit{inode}. Si sono riportate in
tab.~\ref{tab:file_inode_operations} le più rilevanti.

\begin{table}[htb]
  \centering
  \footnotesize
  \begin{tabular}[c]{|l|l|}
    \hline
    \textbf{Funzione} & \textbf{Operazione} \\
    \hline
    \hline
    \textsl{\code{create}} & Chiamata per creare un nuovo file (vedi
                             sez.~\ref{sec:file_open_close}).\\ 
    \textsl{\code{link}}   & Crea un \textit{hard link} (vedi
                             sez.~\ref{sec:link_symlink_rename}).\\
    \textsl{\code{unlink}} & Cancella un \textit{hard link} (vedi
                             sez.~\ref{sec:link_symlink_rename}).\\
    \textsl{\code{symlink}}& Crea un collegamento simbolico (vedi
                             sez.~\ref{sec:link_symlink_rename}).\\
    \textsl{\code{mkdir}}  & Crea una directory (vedi
                             sez.~\ref{sec:file_dir_creat_rem}).\\
    \textsl{\code{rmdir}}  & Rimuove una directory (vedi
                             sez.~\ref{sec:file_dir_creat_rem}).\\
    \textsl{\code{mknod}}  & Crea un file speciale (vedi
                             sez.~\ref{sec:file_mknod}).\\
    \textsl{\code{rename}} & Cambia il nome di un file (vedi
                             sez.~\ref{sec:link_symlink_rename}).\\
    \textsl{\code{lookup}}&  Risolve il nome di un file.\\
    \hline
  \end{tabular}
  \caption{Le principali operazioni sugli \textit{inode} definite tramite
    \kstructd{inode\_operation}.} 
  \label{tab:file_inode_operations}
\end{table}

Possiamo notare come molte di queste funzioni abbiano nomi sostanzialmente
identici alle varie \textit{system call} con le quali si gestiscono file e
directory, che tratteremo nel resto del capitolo. Quello che succede è che
tutte le volte che deve essere eseguita una \textit{system call}, o una
qualunque altra operazione su un \textit{inode} (come \texttt{lookup}) il VFS
andrà ad utilizzare la funzione corrispondente attraverso il puntatore
\var{i\_op}.

Sarà allora sufficiente che nella realizzazione di un filesystem si crei una
implementazione di queste funzioni per quel filesystem e si allochi una
opportuna istanza di \kstruct{inode\_operation} contenente i puntatori a dette
funzioni. A quel punto le strutture \kstruct{inode} usate per gli oggetti di
quel filesystem otterranno il puntatore alla relativa istanza di
\kstruct{inode\_operation} e verranno automaticamente usate le funzioni
corrette.

Si noti però come in tab.~\ref{tab:file_inode_operations} non sia presente la
funzione \texttt{open} che invece è citata in
tab.~\ref{tab:file_file_operations}.\footnote{essa può essere comunque
  invocata dato che nella struttura \kstruct{inode} è presente anche il
  puntatore \var{i\_fop} alla struttura \kstruct{file\_operation} che fornisce
  detta funzione.} Questo avviene perché su Linux l'apertura di un file
richiede comunque un'altra operazione che mette in gioco l'omonimo oggetto del
VFS: l'allocazione di una struttura di tipo \kstruct{file} che viene associata
ad ogni file aperto nel sistema.  I motivi per cui viene usata una struttura a
parte sono diversi, anzitutto, come illustrato in sez.~\ref{sec:file_fd},
questa è necessaria per le operazioni eseguite dai processi con l'interfaccia
dei file descriptor. Ogni processo infatti mantiene il riferimento ad una
struttura \kstruct{file} per ogni file che ha aperto, ed è tramite essa che
esegue le operazioni di I/O. Inoltre il kernel mantiene un elenco di tutti i
file aperti nella \textit{file table} (torneremo su questo in
sez.~\ref{sec:file_fd}).

Inoltre se le operazioni relative agli \textit{inode} fanno riferimento ad
oggetti posti all'interno di un filesystem e vi si applicano quindi le
funzioni fornite nell'implementazione di quest'ultimo, quando si apre un file
questo può essere anche un file di dispositivo, ed in questo caso il VFS
invece di usare le operazioni fornite dal filesystem (come farebbe per un file
di dati) dovrà invece ricorrere a quelle fornite dal driver del dispositivo.

\itindend{inode}

\begin{figure}[!htb]
  \footnotesize \centering
  \begin{minipage}[c]{0.8\textwidth}
    \includestruct{listati/file.h}
  \end{minipage}
  \normalsize 
  \caption{Estratto della struttura \kstructd{file} del kernel (da
    \texttt{include/linux/fs.h}).}
  \label{fig:kstruct_file}
\end{figure}

Come si può notare dall'estratto di fig.~\ref{fig:kstruct_file}, la struttura
\kstruct{file} contiene, oltre ad alcune informazioni usate dall'interfaccia
dei file descriptor il cui significato emergerà più avanti (vedi
sez.~\ref{sec:file_unix_interface}), il puntatore \var{f\_op} ad una struttura
\kstruct{file\_operation}. Questa è l'analoga per i file di
\kstruct{inode\_operation}, e definisce le operazioni generiche fornite dal
VFS per i file. Si sono riportate in tab.~\ref{tab:file_file_operations} le
più significative.

\begin{table}[htb]
  \centering
  \footnotesize
  \begin{tabular}[c]{|l|p{8cm}|}
    \hline
    \textbf{Funzione} & \textbf{Operazione} \\
    \hline
    \hline
    \textsl{\code{open}}   & Apre il file (vedi
                             sez.~\ref{sec:file_open_close}).\\ 
    \textsl{\code{read}}   & Legge dal file (vedi sez.~\ref{sec:file_read}).\\
    \textsl{\code{write}}  & Scrive sul file (vedi 
                             sez.~\ref{sec:file_write}).\\
    \textsl{\code{llseek}} & Sposta la posizione corrente sul file (vedi
                             sez.~\ref{sec:file_lseek}).\\
    \textsl{\code{ioctl}}  & Accede alle operazioni di controllo 
                             (vedi sez.~\ref{sec:file_fcntl_ioctl}).\\
    \textsl{\code{readdir}}& Legge il contenuto di una directory (vedi 
                             sez.~\ref{sec:file_dir_read}).\\
    \textsl{\code{poll}}   & Usata nell'I/O multiplexing (vedi
                             sez.~\ref{sec:file_multiplexing}).\\
    \textsl{\code{mmap}}   & Mappa il file in memoria (vedi 
                             sez.~\ref{sec:file_memory_map}).\\
    \textsl{\code{release}}& Chiamata quando l'ultimo riferimento a un file 
                             aperto è chiuso.\\
    \textsl{\code{fsync}}  & Sincronizza il contenuto del file (vedi
                             sez.~\ref{sec:file_sync}).\\
    \textsl{\code{fasync}} & Abilita l'I/O asincrono (vedi
                             sez.~\ref{sec:file_asyncronous_io}) sul file.\\
    \hline
  \end{tabular}
  \caption{Operazioni sui file definite tramite \kstructd{file\_operation}.}
  \label{tab:file_file_operations}
\end{table}

Anche in questo caso tutte le volte che deve essere eseguita una
\textit{system call} o una qualunque altra operazione sul file il VFS andrà ad
utilizzare la funzione corrispondente attraverso il puntatore
\var{f\_op}. Dato che è cura del VFS quando crea la struttura all'apertura del
file assegnare a \var{f\_op} il puntatore alla versione di
\kstruct{file\_operation} corretta per quel file, sarà possibile scrivere allo
stesso modo sulla porta seriale come su un normale file di dati, e lavorare
sui file allo stesso modo indipendentemente dal filesystem.

Il VFS realizza la quasi totalità delle operazioni relative ai file grazie
alle funzioni presenti nelle due strutture \kstruct{inode\_operation} e
\kstruct{file\_operation}.  Ovviamente non è detto che tutte le operazioni
possibili siano poi disponibili in tutti i casi, ad esempio \code{llseek} non
sarà presente per un dispositivo come la porta seriale o per una
\textit{fifo}, mentre sui file del filesystem \texttt{vfat} non saranno
disponibili i permessi, ma resta il fatto che grazie al VFS le \textit{system
  call} per le operazioni sui file possono restare sempre le stesse nonostante
le enormi differenze che possono esserci negli oggetti a cui si applicano.
 

\itindend{Virtual~File~System~(VFS)}

% NOTE: documentazione interessante:
%       * sorgenti del kernel: Documentation/filesystems/vfs.txt
%       * http://thecoffeedesk.com/geocities/rkfs.html
%       * http://www.linux.it/~rubini/docs/vfs/vfs.html



\subsection{Il funzionamento di un filesystem Unix}
\label{sec:file_filesystem}

Come già accennato in sez.~\ref{sec:file_arch_overview} Linux (ed ogni sistema
unix-like) organizza i dati che tiene su disco attraverso l'uso di un
filesystem. Una delle caratteristiche di Linux rispetto agli altri Unix è
quella di poter supportare, grazie al VFS, una enorme quantità di filesystem
diversi, ognuno dei quali avrà una sua particolare struttura e funzionalità
proprie.  Per questo non entreremo nei dettagli di un filesystem specifico, ma
daremo una descrizione a grandi linee che si adatta alle caratteristiche
comuni di qualunque filesystem di un sistema unix-like.

\itindbeg{superblock}

Una possibile strutturazione dell'informazione su un disco è riportata in
fig.~\ref{fig:file_disk_filesys}, dove si hanno tre filesystem su tre
partizioni. In essa per semplicità si è fatto riferimento alla struttura del
filesystem \acr{ext2}, che prevede una suddivisione dei dati in \textit{block
  group}.  All'interno di ciascun \textit{block group} viene anzitutto
replicato il cosiddetto \textit{superblock}, (la struttura che contiene
l'indice iniziale del filesystem e che consente di accedere a tutti i dati
sottostanti) e creata una opportuna suddivisione dei dati e delle informazioni
per accedere agli stessi.  Sulle caratteristiche di \acr{ext2} e derivati
torneremo in sez.~\ref{sec:file_ext2}.

\itindend{superblock}
\itindbeg{inode}

È comunque caratteristica comune di tutti i filesystem per Unix,
indipendentemente da come poi viene strutturata nei dettagli questa
informazione, prevedere la presenza di due tipi di risorse: gli
\textit{inode}, cui abbiamo già accennato in sez.~\ref{sec:file_vfs_work}, che
sono le strutture che identificano i singoli oggetti sul filesystem, e i
blocchi, che invece attengono allo spazio disco che viene messo a disposizione
per i dati in essi contenuti.

\begin{figure}[!htb]
  \centering
  \includegraphics[width=11cm]{img/disk_struct}
  \caption{Organizzazione dello spazio su un disco in partizioni e
  filesystem.}
  \label{fig:file_disk_filesys}
\end{figure}

Se si va ad esaminare con maggiore dettaglio la strutturazione
dell'informazione all'interno del filesystem \acr{ext2}, tralasciando i
dettagli relativi al funzionamento del filesystem stesso come la
strutturazione in gruppi dei blocchi, il \textit{superblock} e tutti i dati di
gestione possiamo esemplificare la situazione con uno schema come quello
esposto in fig.~\ref{fig:file_filesys_detail}.

\begin{figure}[!htb]
  \centering
  \includegraphics[width=11cm]{img/filesys_struct}
  \caption{Strutturazione dei dati all'interno di un filesystem.}
  \label{fig:file_filesys_detail}
\end{figure}

Da fig.~\ref{fig:file_filesys_detail} si evidenziano alcune delle
caratteristiche di base di un filesystem, che restano le stesse anche su
filesystem la cui organizzazione dei dati è totalmente diversa da quella
illustrata, e sulle quali è bene porre attenzione visto che sono fondamentali
per capire il funzionamento delle funzioni che manipolano i file e le
directory che tratteremo nel prosieguo del capitolo. In particolare è
opportuno tenere sempre presente che:


\begin{enumerate*}
  
\item L'\textit{inode} contiene i cosiddetti \textsl{metadati}, vale dire le
  informazioni riguardanti le proprietà del file come oggetto del filesystem:
  il tipo di file, i permessi di accesso, le dimensioni, i puntatori ai
  blocchi fisici che contengono i dati e così via. Le informazioni che la
  funzione \func{stat} (vedi sez.~\ref{sec:file_stat}) fornisce provengono
  dall'\textit{inode}.  Dentro una directory si troverà solo il nome del file
  e il numero dell'\textit{inode} ad esso associato; il nome non è una
  proprietà del file e non viene mantenuto nell'\textit{inode}. Da da qui in
  poi chiameremo il nome del file contenuto in una directory
  ``\textsl{voce}'', come traduzione della nomenclatura inglese
  \textit{directory entry} che non useremo per evitare confusione con le
  \textit{dentry} del kernel viste in sez.~\ref{sec:file_vfs_work}.
  
\item Come mostrato in fig.~\ref{fig:file_filesys_detail} per i file
  \texttt{macro.tex} e \texttt{gapil\_macro.tex}, ci possono avere più voci
  che fanno riferimento allo stesso \textit{inode}. Fra le proprietà di un
  file mantenute nell'\textit{inode} c'è anche il contatore con il numero di
  riferimenti che sono stati fatti ad esso, il cosiddetto \textit{link
    count}.\footnote{mantenuto anche nel campo \var{i\_nlink} della struttura
    \kstruct{inode} di fig.~\ref{fig:kstruct_inode}.}  Solo quando questo
  contatore si annulla i dati del file possono essere effettivamente rimossi
  dal disco. Per questo la funzione per cancellare un file si chiama
  \func{unlink} (vedi sez.~\ref{sec:link_symlink_rename}), ed in realtà non
  cancella affatto i dati del file, ma si limita ad eliminare la relativa voce
  da una directory e decrementare il numero di riferimenti
  nell'\textit{inode}.
  
\item All'interno di ogni filesystem ogni \textit{inode} è identificato da un
  numero univoco. Il numero di \textit{inode} associato ad una voce in una
  directory si riferisce ad questo numero e non ci può essere una directory
  che contiene riferimenti ad \textit{inode} relativi ad altri filesystem.
  Questa è la ragione che limita l'uso del comando \cmd{ln}, che crea una
  nuova voce per un file esistente con la funzione \func{link} (vedi
  sez.~\ref{sec:link_symlink_rename}), a operare su file nel filesystem
  corrente.
  
\item Quando si cambia nome ad un file senza cambiare filesystem il contenuto
  del file non viene spostato fisicamente, viene semplicemente creata una
  nuova voce per l'\textit{inode} in questione e rimossa la precedente, questa
  è la modalità in cui opera normalmente il comando \cmd{mv} attraverso la
  funzione \func{rename} (vedi sez.~\ref{sec:link_symlink_rename}). Questa
  operazione non modifica minimamente neanche l'\textit{inode} del file, dato
  che non si opera sul file ma sulla directory che lo contiene.

\item Gli \textit{inode} dei file, che contengono i \textsl{metadati}, ed i
  blocchi di spazio disco, che contengono i dati, sono risorse indipendenti ed
  in genere vengono gestite come tali anche dai diversi filesystem; è pertanto
  possibile esaurire sia lo spazio disco (il caso più comune) che lo spazio
  per gli \textit{inode}. Nel primo caso non sarà possibile allocare ulteriore
  spazio, ma si potranno creare file (vuoti), nel secondo non si potranno
  creare nuovi file, ma si potranno estendere quelli che ci
  sono.\footnote{questo comportamento non è generale, alcuni filesystem più
    sofisticati possono evitare il problema dell'esaurimento degli
    \textit{inode} riallocando lo spazio disco libero per i blocchi.}

\end{enumerate*}

\begin{figure}[!htb]
  \centering 
  \includegraphics[width=12cm]{img/dir_links}
  \caption{Organizzazione dei \textit{link} per le directory.}
  \label{fig:file_dirs_link}
\end{figure}

Infine tenga presente che, essendo file pure loro, il numero di riferimenti
esiste anche per le directory. Per questo se a partire dalla situazione
mostrata in fig.~\ref{fig:file_filesys_detail} creiamo una nuova directory
\file{img} nella directory \file{gapil}, avremo una situazione come quella
illustrata in fig.~\ref{fig:file_dirs_link}.

La nuova directory avrà un numero di riferimenti pari a due, in quanto è
referenziata dalla directory da cui si era partiti (in cui è inserita la nuova
voce che fa riferimento a \texttt{img}) e dalla voce interna ``\texttt{.}''
che è presente in ogni directory.  Questo è il valore che si troverà sempre
per ogni directory che non contenga a sua volta altre directory. Al contempo,
la directory da cui si era partiti avrà un numero di riferimenti di almeno
tre, in quanto adesso sarà referenziata anche dalla voce ``\texttt{..}'' di
\texttt{img}. L'aggiunta di una sottodirectory fa cioè crescere di uno il
\textit{link count} della directory genitrice.

\itindend{inode}


\subsection{Alcuni dettagli sul filesystem \acr{ext2} e successori}
\label{sec:file_ext2}

Benché non esista ``il'' filesystem di Linux, dato che esiste un supporto
nativo di diversi filesystem che sono in uso da anni, quello che gli avvicina
di più è la famiglia di filesystem evolutasi a partire dal \textit{second
  extended filesystem}, o \acr{ext2}. Il filesystem \acr{ext2} ha subito un
grande sviluppo e diverse evoluzioni, fra cui l'aggiunta del
\textit{journaling} con il passaggio ad \acr{ext3}, che probabilmente è ancora
il filesystem più diffuso, ed una serie di ulteriori miglioramenti con il
successivo \acr{ext4}. In futuro è previsto che questo debba essere sostituito
da un filesystem completamente diverso, \acr{Btrfs}, che dovrebbe diventare il
filesystem standard di Linux, ma questo al momento è ancora in fase di
sviluppo.\footnote{si fa riferimento al momento dell'ultima revisione di
  questo paragrafo, l'inizio del 2012.}

Il filesystem \acr{ext2} nasce come filesystem nativo per Linux a partire
dalle prime versioni del kernel e supporta tutte le caratteristiche di un
filesystem standard Unix: è in grado di gestire nomi di file lunghi (256
caratteri, estensibili a 1012) e supporta una dimensione massima dei file fino
a 4~Tb. I successivi filesystem \acr{ext3} ed \acr{ext4} sono evoluzioni di
questo filesystem, e sia pure con molti miglioramenti ed estensioni
significative ne mantengono le caratteristiche fondamentali.

Oltre alle caratteristiche standard, \acr{ext2} fornisce alcune estensioni che
non sono presenti su un classico filesystem di tipo Unix; le principali sono
le seguenti:
\begin{itemize*}
\item gli attributi estesi (vedi sez.~\ref{sec:file_xattr}) che consentono di
  estendere le informazioni salvabili come metadati e le ACL (vedi
  sez.~\ref{sec:file_ACL}) che consentono di estendere il modello tradizionale
  dei permessi sui file.
\item sono supportate entrambe le semantiche di BSD e SVr4 come opzioni di
  montaggio. La semantica BSD comporta che i file in una directory sono creati
  con lo stesso identificatore di gruppo della directory che li contiene. La
  semantica SVr4 comporta che i file vengono creati con l'identificatore del
  gruppo primario del processo, eccetto il caso in cui la directory ha il bit
  di \acr{sgid} impostato (per una descrizione dettagliata del significato di
  questi termini si veda sez.~\ref{sec:file_access_control}), nel qual caso
  file e subdirectory ereditano sia il \ids{GID} che lo \acr{sgid}.
\item l'amministratore può scegliere la dimensione dei blocchi del filesystem
  in fase di creazione, a seconda delle sue esigenze: blocchi più grandi
  permettono un accesso più veloce, ma sprecano più spazio disco.
\item il filesystem implementa collegamenti simbolici veloci, in cui il nome
  del file non è salvato su un blocco, ma tenuto all'interno
  dell'\textit{inode} (evitando letture multiple e spreco di spazio), non
  tutti i nomi però possono essere gestiti così per limiti di spazio (il
  limite è 60 caratteri).
\item vengono supportati i cosiddetti \textit{file attributes} (vedi
  sez.~\ref{sec:file_perm_overview}) che attivano comportamenti specifici per
  i file su cui vengono attivati come marcarli come immutabili (che possono
  cioè essere soltanto letti) per la protezione di file di configurazione
  sensibili, o come \textit{append-only} (che possono essere aperti in
  scrittura solo per aggiungere dati) per la protezione dei file di log.
\end{itemize*}

La struttura di \acr{ext2} è stata ispirata a quella del filesystem di BSD: un
filesystem è composto da un insieme di blocchi, la struttura generale è quella
riportata in fig.~\ref{fig:file_filesys_detail}, in cui la partizione è divisa
in gruppi di blocchi.

Ciascun gruppo di blocchi contiene una copia delle informazioni essenziali del
filesystem (i \textit{superblock} sono quindi ridondati) per una maggiore
affidabilità e possibilità di recupero in caso di corruzione del
\textit{superblock} principale. L'utilizzo di raggruppamenti di blocchi ha
inoltre degli effetti positivi nelle prestazioni dato che viene ridotta la
distanza fra i dati e la tabella degli \textit{inode}.

\begin{figure}[!htb]
  \centering
  \includegraphics[width=9cm]{img/dir_struct}  
  \caption{Struttura delle directory nel \textit{second extended filesystem}.}
  \label{fig:file_ext2_dirs}
\end{figure}


Le directory sono implementate come una \textit{linked list} con voci di
dimensione variabile. Ciascuna voce della lista contiene il numero di
\textit{inode}, la sua lunghezza, il nome del file e la sua lunghezza, secondo
lo schema in fig.~\ref{fig:file_ext2_dirs}; in questo modo è possibile
implementare nomi per i file anche molto lunghi (fino a 1024 caratteri) senza
sprecare spazio disco.

Con l'introduzione del filesystem \acr{ext3} sono state introdotte diverse
modifiche strutturali, la principale di queste è quella che \acr{ext3} è un
filesystem \textit{journaled}, è cioè in grado di eseguire una registrazione
delle operazioni di scrittura su un giornale (uno speciale file interno) in
modo da poter garantire il ripristino della coerenza dei dati del
filesystem\footnote{si noti bene che si è parlato di dati \textsl{del}
  filesystem, non di dati \textsl{nel} filesystem, quello di cui viene
  garantito un veloce ripristino è relativo ai dati della struttura interna
  del filesystem, non di eventuali dati contenuti nei file che potrebbero
  essere stati persi.} in brevissimo tempo in caso di interruzione improvvisa
della corrente o di crollo del sistema che abbia causato una interruzione
della scrittura dei dati sul disco.

Oltre a questo \acr{ext3} introduce ulteriori modifiche volte a migliorare
sia le prestazioni che la semplicità di gestione del filesystem, in
particolare per le directory si è passato all'uso di alberi binari con
indicizzazione tramite \textit{hash} al posto delle \textit{linked list} che
abbiamo illustrato, ottenendo un forte guadagno di prestazioni in caso di
directory contenenti un gran numero di file.

% TODO (bassa priorità) portare a ext3, ext4 e btrfs ed illustrare le
% problematiche che si possono incontrare (in particolare quelle relative alla
% perdita di contenuti in caso di crash del sistema)
% TODO (media priorità) trattare btrfs quando sarà usato come stabile


\subsection{La gestione dell'uso dei filesystem}
\label{sec:filesystem_mounting}

Come accennato in sez.~\ref{sec:file_arch_overview} per poter accedere ai file
occorre rendere disponibile al sistema il filesystem su cui essi sono
memorizzati. L'operazione di attivazione del filesystem è chiamata
\textsl{montaggio} e per far questo in Linux si usa la funzione di sistema
\funcd{mount}, il cui prototipo è:\footnote{la funzione è una versione
  specifica di Linux che usa la omonima \textit{system call} e non è
  portabile.}

\begin{funcproto}{
\fhead{sys/mount.h} 
\fdecl{mount(const char *source, const char *target, const char
  *filesystemtype, \\ 
\phantom{mount(}unsigned long mountflags, const void *data)}
\fdesc{Monta un filesystem.} 
}
{La funzione ritorna $0$ in caso di successo e $-1$ per un errore, nel qual
  caso \var{errno} assumerà uno dei valori:
  \begin{errlist}
  \item[\errcode{EACCES}] non si ha il permesso di accesso su uno dei
    componenti del \textit{pathname}, o si è cercato di montare un filesystem
    disponibile in sola lettura senza aver specificato \const{MS\_RDONLY} o il
    device \param{source} è su un filesystem montato con l'opzione
    \const{MS\_NODEV}.
  \item[\errcode{EBUSY}] \param{source} è già montato, o non può essere
    rimontato in sola lettura perché ci sono ancora file aperti in scrittura,
    o non può essere montato su \param{target} perché la directory è ancora in
    uso.
  \item[\errcode{EINVAL}] il dispositivo \param{source} presenta un
    \textit{superblock} non valido, o si è cercato di rimontare un filesystem
    non ancora montato, o di montarlo senza che \param{target} sia un
    \textit{mount point} o di spostarlo quando \param{target} non è un
    \textit{mount point} o è la radice o si è usato un valore di
    \param{mountflags} non valido.
  \item[\errcode{ELOOP}] si è cercato di spostare un \textit{mount point} su
    una sottodirectory di \param{source} o si sono incontrati troppi
    collegamenti simbolici nella risoluzione di un nome.
  \item[\errcode{EMFILE}] in caso di filesystem virtuale, la tabella dei
    dispositivi fittizi (chiamati \textit{dummy} nella documentazione inglese)
    è piena.
  \item[\errcode{ENODEV}] il tipo \param{filesystemtype} non esiste o non è
    configurato nel kernel.
  \item[\errcode{ENOTBLK}] non si è usato un \textit{block device} per
    \param{source} quando era richiesto.
  \item[\errcode{ENXIO}] il \textit{major number} del
    dispositivo \param{source} è sbagliato.
  \item[\errcode{EPERM}] il processo non ha i privilegi di amministratore.
  \end{errlist} 
  ed inoltre \errval{EFAULT}, \errval{ENOMEM}, \errval{ENAMETOOLONG},
  \errval{ENOENT}, \errval{ENOTDIR} nel loro significato generico.}
\end{funcproto}

\itindbeg{mount~point}

L'uso più comune della funzione è quello di montare sulla directory indicata
da \param{target}, detta \textit{mount point}, il filesystem contenuto nel
file di dispositivo indicato da \param{source}. In entrambi i casi, come
daremo per assunto da qui in avanti tutte le volte che si parla di directory o
file nel passaggio di un argomento di una funzione, si intende che questi
devono essere indicati con la stringa contenente il loro \textit{pathname}.

Normalmente un filesystem è contenuto su un disco o una partizione, ma come
illustrato in sez.~\ref{sec:file_vfs_work} la struttura del \textit{Virtual
  File System} è estremamente flessibile e può essere usata anche per oggetti
diversi da un disco. Ad esempio usando il \textit{loop device} si può montare
un file qualunque (come l'immagine di un CD-ROM o di un floppy) che contiene
l'immagine di un filesystem, inoltre alcuni tipi di filesystem, come
\texttt{proc} o \texttt{sysfs} sono virtuali e non hanno un supporto che ne
contenga i dati che sono generati al volo dal kernel ad ogni lettura, e
inviati al kernel ad ogni scrittura (costituiscono quindi un meccanismo di
comunicazione, attraverso l'ordinaria interfaccia dei file, con il kernel).

Il tipo di filesystem che si vuole montare è specificato
dall'argomento \param{filesystemtype}, che deve essere una delle stringhe
riportate nel file \procfilem{/proc/filesystems} che, come accennato in
sez.~\ref{sec:file_vfs_work}, contiene l'elenco dei filesystem supportati dal
kernel. Nel caso si sia indicato un filesystem virtuale, che non è associato a
nessun file di dispositivo, il contenuto di \param{source} viene ignorato.

L'argomento \param{data} viene usato per passare le impostazioni relative alle
caratteristiche specifiche di ciascun filesystem. Si tratta di una stringa di
parole chiave (separate da virgole e senza spazi) che indicano le cosiddette
``\textsl{opzioni}'' del filesystem che devono essere impostate; in genere
viene usato direttamente il contenuto del parametro dell'opzione \texttt{-o}
del comando \texttt{mount}. I valori utilizzabili dipendono dal tipo di
filesystem e ciascuno ha i suoi, pertanto si rimanda alla documentazione della
pagina di manuale di questo comando e dei singoli filesystem.

Dopo l'esecuzione della funzione il contenuto del filesystem viene reso
disponibile nella directory specificata come \textit{mount point} ed il
precedente contenuto di detta directory viene mascherato dal contenuto della
directory radice del filesystem montato. Fino ai kernel della serie 2.2.x non
era possibile montare un filesystem se un \textit{mount point} era già in uso,
coi kernel successivi è possibile montare più filesystem sullo stesso
\textit{mount point} impilandoli l'uno sull'altro, anche in questo caso vale
quanto appena detto, e solo il contenuto dell'ultimo filesystem montato sarà
visibile, mascherando quelli sottostanti.

In realtà quella di montare un filesystem è solo una delle operazioni che si
possono effettuare con \func{mount}, la funzione infatti è dedicata a tutte le
operazioni relative alla gestione del montaggio dei filesystem e dei
\textit{mount point}. Ad esempio fin dalle sue origini poteva essere
utilizzata per effettuare il rimontaggio di un filesystem con opzioni diverse,
ed a partire dal kernel 2.4.x è divenuto possibile usarla per spostare
atomicamente un \textit{mount point} da una directory ad un'altra, per montare
lo stesso filesystem in diversi \textit{mount point}, per montare una
directory su un'altra (il cosiddetto \textit{bind mount}).

\itindend{mount~point}

Il tipo di operazione compiuto da \func{mount} viene stabilito in base al
valore dell'argomento \param{mountflags}, che oltre alla selezione del tipo di
operazione da compiere, consente anche di indicare alcune opzioni generiche
valide per qualunque filesystem.\footnote{benché queste siano espresse nel
  comando \cmd{mount} con l'opzione \texttt{-o} esse non vengono impostate nei
  valori di \param{data}, che serve solo per le opzioni specifiche di ogni
  filesystem.}  Il valore dell'argomento deve essere espresso come maschera
binaria e i vari bit che lo compongono, detti anche \textit{mount flags},
devono essere impostati con un OR aritmetico dei valori dalle opportune
costanti che illustreremo a breve.

In Linux \param{mountflags} deve essere un intero a 32 bit;
fino ai kernel della serie 2.2.x i 16 più significativi avevano un valore
riservato che doveva essere specificato obbligatoriamente,\footnote{il valore
  era il \textit{magic number} \code{0xC0ED}, si può usare la costante
  \constd{MS\_MGC\_MSK} per ottenere la parte di \param{mountflags} riservata
  al \textit{magic number}, mentre per specificarlo si può dare un OR
  aritmetico con la costante \constd{MS\_MGC\_VAL}.} e si potevano usare solo
i 16 meno significativi. Oggi invece, con un numero di opzioni superiore, sono
utilizzati tutti e 32 i bit, ma qualora nei 16 più significativi sia presente
detto valore, che non esprime una combinazione valida, esso viene ignorato.

Come accennato il tipo di operazione eseguito da \func{mount} viene stabilito
in base al contenuto di \param{mountflags}, la scelta viene effettuata
controllando nell'ordine:
\begin{enumerate*}
\item se è presente il flag \const{MS\_REMOUNT}, nel qual caso verrà eseguito
  il rimontaggio del filesystem, con le nuove opzioni indicate da \param{data}
  e dagli altri flag di \param{mountflags};
\item se è presente il flag \const{MS\_BIND}, nel qual caso verrà eseguito un
  \textit{bind mount} (argomento che tratteremo più avanti);
\item se è presente uno fra \const{MS\_SHARED}, \const{MS\_PRIVATE},
  \const{MS\_SLAVE}, \const{MS\_UNBINDABLE}, nel qual caso verrà cambiata la
  modalità di propagazione del montaggio (detti valori sono mutualmente
  esclusivi).
\item se è presente \const{MS\_MOVE}, nel qual caso verrà effettuato uno
  spostamento del \textit{mount point};
\item se nessuno dei precedenti è presente si tratta di una ordinaria
  operazione di montaggio di un filesystem.
\end{enumerate*}

Il fatto che questi valori vengano controllati in quest'ordine significa che
l'effetto di alcuni di questi flag possono cambiare se usati in combinazione
con gli altri che vengono prima nella sequenza (è quanto avviene ad esempio
per \const{MS\_BIND} usato con \const{MS\_REMOUNT}). Tratteremo questi
\textit{mount flags} speciali per primi, nell'ordine appena illustrato,
tornando sugli altri più avanti.

Usando il flag \constd{MS\_REMOUNT} si richiede a \func{mount} di rimontare un
filesystem già montato cambiandone le opzioni di montaggio in maniera atomica
(non è cioè necessario smontare e rimontare il filesystem per effettuare il
cambiamento). Questa operazione consente di modificare le opzioni del
filesystem anche se questo è in uso. Gli argomenti \param{source} e
\param{target} devono essere gli stessi usati per il montaggio originale,
mentre sia \param{data} che \param{mountflags} conterranno le nuove opzioni,
\param{filesystemtype} viene ignorato.  Perché l'operazione abbia successo
occorre comunque che il cambiamento sia possibile (ad esempio non sarà
possibile rimontare in sola lettura un filesystem su cui sono aperti file per
la lettura/scrittura).

Qualunque opzione specifica del filesystem indicata con \param{data} può
essere modificata (ma si dovranno rielencare tutte quelle volute), mentre con
\param{mountflags} possono essere modificate solo alcune opzioni generiche:
\const{MS\_LAZYTIME}, \const{MS\_MANDLOCK}, \const{MS\_NOATIME},
\const{MS\_NODEV}, \const{MS\_NODIRATIME}, \const{MS\_NOEXEC},
\const{MS\_NOSUID}, \const{MS\_RELATIME}, \const{MS\_RDONLY},
\const{MS\_STRICTATIME} e \const{MS\_SYNCHRONOUS}. Inoltre dal kernel 3.17 il
comportamento relativo alle opzioni che operano sui tempi di ultimo accesso
dei file (vedi sez.~\ref{sec:file_file_times}) è cambiato e se non si è
indicato nessuno dei vari \texttt{MS\_*ATIME} vengono mantenute le
impostazioni esistenti anziché forzare l'uso di \const{MS\_RELATIME}.

\itindbeg{bind~mount}

Usando il flag \constd{MS\_BIND} si richiede a \func{mount} di effettuare un
cosiddetto \textit{bind mount}, l'operazione che consente di montare una
directory di un filesystem in un'altra directory. L'opzione è disponibile a
partire dai kernel della serie 2.4. In questo caso verranno presi in
considerazione solo gli argomenti \param{source}, che stavolta indicherà la
directory che si vuole montare e non un file di dispositivo, e \param{target}
che indicherà la directory su cui verrà effettuato il \textit{bind mount}. Gli
argomenti \param{filesystemtype} e \param{data} vengono ignorati.

Quello che avviene con questa operazione è che in corrispondenza del
\textit{pathname} indicato da \param{target} viene montato l'\textit{inode} di
\param{source}, così che la porzione di albero dei file presente sotto
\param{source} diventi visibile allo stesso modo sotto
\param{target}. Trattandosi esattamente dei dati dello stesso filesystem, ogni
modifica fatta in uno qualunque dei due rami di albero sarà visibile
nell'altro, visto che entrambi faranno riferimento agli stessi \textit{inode}.

Dal punto di vista del VFS l'operazione è analoga al montaggio di un
filesystem proprio nel fatto che anche in questo caso si inserisce in
corrispondenza della \textit{dentry} di \texttt{target} un diverso
\textit{inode}, che stavolta, invece di essere quello della radice del
filesystem indicato da un file di dispositivo, è quello di una directory già
montata.

Si tenga presente che proprio per questo sotto \param{target} comparirà il
contenuto che è presente sotto \param{source} all'interno del filesystem in
cui quest'ultima è contenuta. Questo potrebbe non corrispondere alla porzione
di albero che sta sotto \param{source} qualora in una sottodirectory di
quest'ultima si fosse effettuato un altro montaggio. In tal caso infatti nella
porzione di albero sotto \param{source} si troverebbe il contenuto del nuovo
filesystem (o di un altro \textit{bind mount}) mentre sotto \param{target} ci
sarebbe il contenuto presente nel filesystem originale.

L'unico altro \textit{mount flag} usabile direttamente con \const{MS\_BIND} è
\const{MS\_REC} che consente di eseguire una operazione di \textit{bind mount}
ricorsiva, in cui sotto \param{target} vengono montati ricorsivamente anche
tutti gli eventuali ulteriori \textit{bind mount} già presenti sotto
\param{source}.

E' però possibile, a partire dal kernel 2.6.26, usare questo flag insieme a
\const{MS\_REMOUNT}, nel qual caso consente di effettuare una modifica delle
opzioni di montaggio del \textit{bind mount} ed in particolare effettuare il
cosiddetto \textit{read-only bind mount} in cui viene onorata anche la
presenza aggiuntiva del flag \const{MS\_RDONLY}. In questo modo si ottiene che
l'accesso ai file sotto \param{target} sia effettuabile esclusivamente in sola
lettura, mantenendo il normale accesso in lettura/scrittura sotto
\param{source}.

Il supporto per il \textit{bind mount} consente di superare i limiti presenti
per gli \textit{hard link} (di cui parleremo in
sez.~\ref{sec:link_symlink_rename}) con la possibilità di fare riferimento
alla porzione dell'albero dei file di un filesystem presente a partire da una
certa directory utilizzando una qualunque altra directory, anche se questa sta
su un filesystem diverso.\footnote{e non c'è neanche il problema di non esser
  più in grado di cancellare un \textit{hard link} ad una directory sullo
  stesso filesystem (vedi sez.~\ref{sec:link_symlink_rename}), per cui su
  Linux questi non sono possibili, dato che in questo caso per la rimozione
  del collegamento basta smontare \param{target}.} Si può così fornire una
alternativa all'uso dei collegamenti simbolici (di cui parleremo in
sez.~\ref{sec:link_symlink_rename}) che funziona correttamente anche
all'intero di un \textit{chroot} (argomento su cui torneremo in
sez.~\ref{sec:file_chroot}).

\itindend{bind~mount}
\itindbeg{shared~subtree}

I quattro flag \const{MS\_PRIVATE}, \const{MS\_SHARED}, \const{MS\_SLAVE} e
\const{MS\_UNBINDABLE} sono stati introdotti a partire dal kernel 2.6.15 per
realizzare l'infrastruttura dei cosiddetti \textit{shared subtree}, che
estendono le funzionalità dei \textit{bind mount}.  La funzionalità nasce
dalle esigenze di poter utilizzare a pieno le funzionalità di isolamento
fornite dal kernel per i processi (i \textit{namespace}, che tratteremo in
sez.~\ref{sec:process_namespaces}) in particolare per quanto riguarda la
possibilità di far avere ad un processo una visione ristretta dei filesystem
montati (il \textit{mount namespace}), ma l'applicazione è comunque rilevante
anche con un classico \textit{chroot} (vedi sez.~\ref{sec:file_chroot}).

\itindbeg{submount}

Abbiamo visto come nella modalità ordinaria in cui si esegue un
\textit{bind mount} sotto \param{target} compaia lo stesso ramo di albero dei
file presente sotto \param{source}, ma limitato a quanto presente nel
filesystem di \param{source}; i risultati di un eventuale
``\textit{submount}'' effettuato all'interno di \param{source} non saranno
visibili. Ed anche se quelli presenti al momento dell'uso di \const{MS\_BIND}
possono essere riottenuti usando \const{MS\_REC}, ogni eventuale
``\textit{submount}'' successivo (che avvenga sotto \param{source} o sotto
\param{target}) resterà ``\textsl{privato}'' al ramo di albero su cui è
avvenuto.

\itindend{submount}
\itindbeg{mount peer group}

Ci sono casi però in cui può risultare utile che eventuali
``\textit{submount}'' siano visibili sui rami di albero presenti al di sotto
di tutte le directory coinvolte in un \textit{bind mount}, anche se effettuati
in un secondo tempo. Per poter ottenere questa funzionalità i
\textit{bind mount} sono stati estesi introducendo i \textit{mount peer
  group}, che consentono di raggrupparli in modo da poter inviare a ciascuno
di essi tutti gli eventi relativi a montaggi o smontaggi effettuati al loro
interno ed avere sempre una propagazione degli stessi che li renda coerenti.

Quando si effettua un montaggio ordinario, o si esegue un \textit{bind mount},
di default non viene utilizzato nessun \textit{mount peer group} ed il
\textit{mount point} viene classificato come ``\textsl{privato}'', nel senso
che abbiamo appena visto.  Si può però marcare un \textit{mount point} come
``\textsl{condiviso}'', ed in questo modo esso verrà associato ad un
\textit{mount peer group} insieme a tutti gli altri ulteriori \textit{mount
  point} per i quali sia stato eseguito un \textit{bind mount}. Questo fa sì
che tutte le volte che si effettua un montaggio o uno smontaggio all'interno
di uno qualunque dei \textit{mount point} del gruppo, questo venga propagato
anche su tutti gli altri e sotto tutti sia visibile sempre lo stesso ramo di
albero dei file.

A completare l'infrastruttura degli \textit{shared subtree} sono state
previste due ulteriori funzionalità: la prima è quella di marcare un
\textit{mount point} come ``\textit{slave}'', in tal caso le operazioni di
montaggio e smontaggio effettuate al suo interno non verranno più propagate
agli altri membri del \textit{mount peer group} di cui fa parte, ma continuerà
a ricevere quelle eseguite negli altri membri.

La seconda funzionalità è quella di marcare un \textit{mount point} come
``\textit{unbindable}''; questo anzitutto impedirà che possa essere usato come
sorgente di un \textit{bind mount} ed inoltre lo renderà privato, con la
conseguenza che quando è presente all'interno di altri \textit{bind mount},
all'interno di questi si vedrà solo il contenuto originale e non quello
risultante da eventuali ulteriori montaggi effettuati al suo interno.

\itindend{mount peer group}

I \textit{mount flag} che controllano le operazioni relative agli
\textit{shared subtree} sono descritti nella lista seguente. Si ricordi che
sono mutuamente esclusivi, e compatibili solo con l'uso degli ulteriori flag
\const{MS\_REC} (che applica ricorsivamente l'operazione a tutti gli eventuali
\textit{mount point} sottostanti) e \const{MS\_SILENT}; in tutti gli altri
casi \func{mount} fallirà con un errore di \errval{EINVAL}. L'unico altro
argomento che deve essere specificato quando li si usano è \param{target};
\param{source}, \param{data} e \param{filesystem} sono ignorati.

\begin{basedescript}{\desclabelwidth{1.9cm}\desclabelstyle{\nextlinelabel}}

\item[\constd{MS\_PRIVATE}] Marca un \textit{mount point} come \textit{private
    mount}.  Di default, finché non lo si marca altrimenti con una delle altre
  opzioni dell'interfaccia, ogni \textit{mount point} è privato. Ogni
  \textit{bind mount} ottenuto da un \textit{mount point} privato si comporta
  come descritto nella trattazione di \const{MS\_BIND}. Si usa questo flag
  principalmente per revocare gli effetti delle altre opzioni e riportare il
  comportamento a quello ordinario.

\item[\constd{MS\_SHARED}] Marca un \textit{mount point} come \textit{shared
    mount}.  Lo scopo dell'opzione è ottenere che tutti i successivi
  \textit{bind mount} ottenuti da un \textit{mount point} così marcato siano
  di tipo \textit{shared} e vengano inseriti nello stesso \textit{mount peer
    group} in modo da ``\textsl{condividere}'' ogni ulteriore operazione di
  montaggio o smontaggio. Con questa opzione le operazioni di montaggio e
  smontaggio effettuate al di sotto di uno \textit{shared mount} vengono
  automaticamente ``\textsl{propagate}'' a tutti gli altri membri del
  \textit{mount peer group} di cui fa parte, in modo che la sezione di albero
  dei file visibile al di sotto di ciascuno di essi sia sempre la stessa.

\item[\constd{MS\_SLAVE}] Marca un \textit{mount point} come \textit{slave
    mount}. Se il \textit{mount point} è parte di un \textit{mount peer group}
  esso diventerà di tipo \textit{slave}: le operazioni di montaggio e
  smontaggio al suo interno non verranno più propagate agli altri membri del
  gruppo, ma continuerà a ricevere quelle eseguite negli altri membri. Se non
  esistono altri membri nel gruppo il \textit{mount point} diventerà privato,
  negli altri casi non subirà nessun cambiamento.

\item[\constd{MS\_UNBINDABLE}] Marca un \textit{mount point} come
  \textit{unbindable mount}.  Un \textit{mount point} marcato in questo modo
  non può essere usato per un \textit{bind mount} del suo contenuto. Si
  comporta cioè come allo stesso modo di un \textit{mount point} ordinario di
  tipo \textit{private} con in più la restrizione che nessuna sua
  sottodirectory (anche se relativa ad un ulteriore montaggio) possa essere
  utilizzata come sorgente di un \textit{bind mount}.
  
\end{basedescript}
\itindend{shared~subtree}

L'ultimo \textit{mount flag} che controlla una modalità operativa di
\func{mount} è \constd{MS\_MOVE}, che consente di effettuare lo spostamento
del \textit{mount point} di un filesystem. La directory del \textit{mount
  point} originale deve essere indicata nell'argomento \param{source}, e la
sua nuova posizione nell'argomento \param{target}. Tutti gli altri argomenti
della funzione vengono ignorati.

Lo spostamento avviene atomicamente, ed il ramo di albero presente sotto
\param{source} sarà immediatamente visibile sotto \param{target}. Non esiste
cioè nessun momento in cui il filesystem non risulti montato in una o
nell'altra directory e pertanto è garantito che la risoluzione di
\textit{pathname} relativi all'interno del filesystem non possa fallire.

Elenchiamo infine i restanti \textit{mount flag}, il cui utilizzo non attiene
alle operazioni di \func{mount}, ma soltanto l'impostazione di opzioni
generiche relative al funzionamento di un filesystem e che vengono per lo più
utilizzati solo in fase di montaggio:

\begin{basedescript}{\desclabelwidth{2.0cm}\desclabelstyle{\nextlinelabel}}
\item[\constd{MS\_DIRSYNC}] Richiede che ogni modifica al contenuto di una
  directory venga immediatamente registrata su disco in maniera sincrona
  (introdotta a partire dai kernel della serie 2.6). L'opzione si applica a
  tutte le directory del filesystem, ma su alcuni filesystem è possibile
  impostarla a livello di singole directory o per i sotto-rami di una directory
  con il comando \cmd{chattr}.\footnote{questo avviene tramite delle opportune
    \texttt{ioctl} (vedi sez.~\ref{sec:file_fcntl_ioctl}).}

  Questo consente di ridurre al minimo il rischio di perdita dei dati delle
  directory in caso di crollo improvviso del sistema, al costo di una certa
  perdita di prestazioni dato che le funzioni di scrittura relative ad
  operazioni sulle directory non saranno più bufferizzate e si bloccheranno
  fino all'arrivo dei dati sul disco prima che un programma possa proseguire.

\item[\constd{MS\_LAZYTIME}] Modifica la modalità di registrazione di tempi
  dei file (vedi sez.~\ref{sec:file_file_times}) per ridurre al massimo gli
  accessi a disco (particolarmente utile per i portatili). Attivandolo i tempi
  dei file vengono mantenuti in memoria e vengono salvati su disco solo in
  quattro casi: quando c'è da eseguire un aggiornamento dei dati
  dell'\textit{inode} per altri motivi, se viene usata una delle funzioni di
  sincronizzazione dei dati su disco (vedi sez.~\ref{sec:file_sync}), se
  l'\textit{inode} viene rimosso dalla memoria, o se è passato un giorno
  dall'ultima registrazione. Introdotto a partire dal kernel 4.0.

  In questo modo si possono ridurre significativamente le scritture su disco
  mantenendo tutte le informazioni riguardo ai tempi dei file, riducendo anche
  l'impatto dell'uso di \const{MS\_STRICTATIME}. Il costo da pagare è il
  rischio, in caso di crash del sistema, di avere dati vecchi fino a 24 ore
  per quel che riguarda i tempi dei file.
  
\item[\constd{MS\_MANDLOCK}] Consente l'uso del \textit{mandatory locking}
  (vedi sez.~\ref{sec:file_mand_locking}) sui file del filesystem. Per poterlo
  utilizzare effettivamente però esso dovrà essere comunque attivato
  esplicitamente per i singoli file impostando i permessi come illustrato in
  sez.~\ref{sec:file_mand_locking}.

\item[\constd{MS\_NOATIME}] Viene disabilitato sul filesystem l'aggiornamento
  dell'\textit{access time} (vedi sez.~\ref{sec:file_file_times}) per
  qualunque tipo di file. Dato che l'aggiornamento dell'\textit{access time} è
  una funzionalità la cui utilità è spesso irrilevante ma comporta un costo
  elevato visto che una qualunque lettura comporta comunque una scrittura su
  disco, questa opzione consente di disabilitarla completamente. La soluzione
  può risultare troppo drastica dato che l'informazione viene comunque
  utilizzata da alcuni programmi, per cui nello sviluppo del kernel sono state
  introdotte altre opzioni che forniscono soluzioni più appropriate e meno
  radicali.

\item[\constd{MS\_NODEV}] Viene disabilitato sul filesystem l'accesso ai file
  di dispositivo eventualmente presenti su di esso. L'opzione viene usata come
  misura di precauzione per rendere inutile la presenza di eventuali file di
  dispositivo su filesystem che non dovrebbero contenerne.\footnote{si ricordi
    che le convenzioni del \textit{Linux Filesystem Hierarchy Standard}
    richiedono che questi siano mantenuti esclusivamente sotto \texttt{/dev}.}

  Viene utilizzata, assieme a \const{MS\_NOEXEC} e \const{MS\_NOSUID}, per
  fornire un accesso più controllato a quei filesystem di cui gli utenti hanno
  il controllo dei contenuti, in particolar modo quelli posti su dispositivi
  rimuovibili. In questo modo si evitano alla radice possibili situazioni in
  cui un utente malizioso inserisce su uno di questi filesystem dei file di
  dispositivo con permessi ``opportunamente'' ampliati che gli consentirebbero
  di accedere anche a risorse cui non dovrebbe.

\item[\constd{MS\_NODIRATIME}] Viene disabilitato sul filesystem
  l'aggiornamento dell'\textit{access time} (vedi
  sez.~\ref{sec:file_file_times}) ma soltanto per le directory. Costituisce
  una alternativa per \const{MS\_NOATIME}, che elimina l'informazione per le
  directory, che in pratica che non viene mai utilizzata, mantenendola per i
  file in cui invece ha un impiego, sia pur limitato.

\item[\constd{MS\_NOEXEC}] Viene disabilitata sul filesystem l'esecuzione di un
  qualunque file eseguibile eventualmente presente su di esso. L'opzione viene
  usata come misura di precauzione per rendere impossibile l'uso di programmi
  posti su filesystem che non dovrebbero contenerne.

  Anche in questo caso viene utilizzata per fornire un accesso più controllato
  a quei filesystem di cui gli utenti hanno il controllo dei contenuti. Da
  questo punto di vista l'opzione è meno importante delle analoghe
  \const{MS\_NODEV} e \const{MS\_NOSUID} in quanto l'esecuzione di un
  programma creato dall'utente pone un livello di rischio nettamente
  inferiore, ed è in genere consentita per i file contenuti nella sua home
  directory.\footnote{cosa che renderebbe superfluo l'attivazione di questa
    opzione, il cui uso ha senso solo per ambienti molto controllati in cui si
    vuole che gli utenti eseguano solo i programmi forniti
    dall'amministratore.}

\item[\constd{MS\_NOSUID}] Viene disabilitato sul filesystem l'effetto dei bit
  dei permessi \acr{suid} e \acr{sgid} (vedi sez.~\ref{sec:file_special_perm})
  eventualmente presenti sui file in esso contenuti. L'opzione viene usata
  come misura di precauzione per rendere inefficace l'effetto di questi bit
  per filesystem in cui non ci dovrebbero essere file dotati di questi
  permessi.

  Di nuovo viene utilizzata, analogamente a \const{MS\_NOEXEC} e
  \const{MS\_NODEV}, per fornire un accesso più controllato a quei filesystem
  di cui gli utenti hanno il controllo dei contenuti. In questo caso si evita
  che un utente malizioso possa inserire su uno di questi filesystem un
  eseguibile con il bit \acr{suid} attivo e di proprietà dell'amministratore o
  di un altro utente, che gli consentirebbe di eseguirlo per conto di
  quest'ultimo.

\item[\constd{MS\_RDONLY}] Esegue il montaggio del filesystem in sola lettura,
  non sarà possibile nessuna modifica ai suoi contenuti. Viene usato tutte le
  volte che si deve accedere ai contenuti di un filesystem con la certezza che
  questo non venga modificato (ad esempio per ispezionare un filesystem
  corrotto). All'avvio di default il kernel monta la radice in questa
  modalità. Si tenga presente che se non viene indicato il filesystem verrà
  montato, o rimontato nel caso lo si usi con \const{MS\_REMOUNT}, in
  lettura/scrittura; questo significa in sostanza che non esiste una opzione
  separata per indicare il montaggio in lettura/scrittura.

\item[\constd{MS\_REC}] Applica ricorsivamente a tutti i \textit{mount point}
  presenti al di sotto del \textit{mount point} indicato gli effetti della
  opzione degli \textit{shared subtree} associata. In questo caso l'argomento
  \param{target} deve fare riferimento ad un \textit{mount point} e tutti gli
  altri argomenti sono ignorati, ed il flag deve essere indicato con uno fra
  \const{MS\_PRIVATE}, \const{MS\_SHARED}, \const{MS\_SLAVE} e
  \const{MS\_UNBINDABLE}. Può anche essere usato con \const{MS\_BIND} per
  richiedere il montaggio ricorsivo anche degli eventuali ulteriori
  \textit{bind mount} presenti sotto \param{target}.

\item[\constd{MS\_RELATIME}] Indica di effettuare l'aggiornamento
  dell'\textit{access time} sul filesystem soltanto quando questo risulti
  antecedente il valore corrente del \textit{modification time} o del
  \textit{change time} (per i tempi dei file si veda
  sez.~\ref{sec:file_file_times}). L'opzione è disponibile a partire dal
  kernel 2.6.20, mentre dal 2.6.30 questo è diventato il comportamento di
  default del sistema, che può essere riportato a quello tradizionale con
  l'uso di \const{MS\_STRICTATIME}. Sempre dal 2.6.30 il comportamento è stato
  anche modificato e l'\textit{access time} viene comunque aggiornato se è più
  vecchio di un giorno.

  L'opzione consente di evitare i problemi di prestazioni relativi
  all'aggiornamento dell'\textit{access time} senza avere impatti negativi
  riguardo le funzionalità, il comportamento adottato infatti consente di
  rendere evidente che vi è stato un accesso dopo la scrittura, ed evitando al
  contempo ulteriori operazioni su disco negli accessi successivi. In questo
  modo l'informazione relativa al fatto che un file sia stato letto resta
  disponibile, ed i programmi che ne fanno uso continuano a funzionare. Con
  l'introduzione di questo comportamento l'uso delle alternative
  \const{MS\_NOATIME} e \const{MS\_NODIRATIME} è sostanzialmente inutile.

\item[\constd{MS\_SILENT}] Richiede la soppressione di alcuni messaggi di
  avvertimento nei log del kernel (vedi sez.~\ref{sec:sess_daemon}). L'opzione
  è presente a partire dal kernel 2.6.17 e sostituisce, utilizzando un nome
  non fuorviante, la precedente \const{MS\_VERBOSE}, introdotta nel kernel
  2.6.12, che aveva lo stesso effetto.

\item[\constd{MS\_STRICTATIME}] Ripristina il comportamento tradizionale per
  cui l'\textit{access time} viene aggiornato ad ogni accesso al
  file. L'opzione è disponibile solo a partire dal kernel 2.6.30 quando il
  comportamento di default del kernel è diventato quello fornito da
  \const{MS\_RELATIME}.

\item[\constd{MS\_SYNCHRONOUS}] Abilita la scrittura sincrona richiedendo che
  ogni modifica al contenuto del filesystem venga immediatamente registrata su
  disco. Lo stesso comportamento può essere ottenuto con il flag
  \const{O\_SYNC} di \func{open} (vedi sez.~\ref{sec:file_open_close}).

  Questa opzione consente di ridurre al minimo il rischio di perdita dei dati
  in caso di crollo improvviso del sistema, al costo di una pesante perdita di
  prestazioni dato che tutte le funzioni di scrittura non saranno più
  bufferizzate e si bloccheranno fino all'arrivo dei dati sul disco. Per un
  compromesso in cui questo comportamento avviene solo per le directory, ed ha
  quindi una incidenza nettamente minore, si può usare \const{MS\_DIRSYNC}.

\end{basedescript}

% NOTE: per l'opzione \texttt{lazytime} introdotta con il kernel 4.0,
% vedi http://lwn.net/Articles/621046/

% NOTE per \const{MS\_SLAVE},\const{MS\_SHARE}, \const{MS\_PRIVATE} e
% \const{MS\_UNBINDABLE} dal 2.6.15 vedi shared subtrees, in particolare
%  * http://lwn.net/Articles/159077/ e
%  * Documentation/filesystems/sharedsubtree.txt

% TODO: (bassa priorità) non documentati ma presenti in sys/mount.h:
%       * MS_POSIXACL
%       * MS_KERNMOUNT
%       * MS_I_VERSION
%       * MS_ACTIVE
%       * MS_NOUSER


Una volta che non si voglia più utilizzare un certo filesystem è possibile
``\textsl{smontarlo}'' usando la funzione di sistema \funcd{umount}, il cui
prototipo è:

\begin{funcproto}{
\fhead{sys/mount.h}
\fdecl{umount(const char *target)}
\fdesc{Smonta un filesystem.} 
}
{La funzione ritorna  $0$ in caso di successo e $-1$ per un errore,
  nel qual caso \var{errno} assumerà uno dei valori: 
  \begin{errlist}
  \item[\errcode{EBUSY}] il filesystem è occupato.
  \item[\errcode{EINVAL}] \param{target} non è un \textit{mount point}.
  \item[\errcode{EPERM}] il processo non ha i privilegi di
    amministratore.\footnotemark 
  \end{errlist}
  ed inoltre \errval{EFAULT}, \errval{ELOOP}, \errval{ENAMETOOLONG},
  \errval{ENOENT}, \errval{ENOMEM} nel loro significato generico.  }
\end{funcproto}

\footnotetext{più precisamente la capacità \const{CAP\_SYS\_ADMIN}, vedi
  sez.~\ref{sec:proc_capabilities}.}

La funzione prende il nome della directory su cui il filesystem è montato e
non il file o il dispositivo che è stato montato,\footnote{questo è vero a
  partire dal kernel 2.3.99-pre7, prima esistevano due chiamate separate e la
  funzione poteva essere usata anche specificando il file di dispositivo.} in
quanto a partire dai kernel della serie 2.4.x è possibile montare lo stesso
dispositivo in più punti. Nel caso più di un filesystem sia stato montato
sullo stesso \textit{mount point} viene smontato quello che è stato montato
per ultimo. Si tenga presente che la funzione fallisce se il filesystem è
``\textsl{occupato}'', cioè quando ci sono ancora dei file aperti sul
filesystem, se questo contiene la directory di lavoro (vedi
sez.~\ref{sec:file_work_dir}) di un qualunque processo o il \textit{mount
  point} di un altro filesystem.

Linux provvede inoltre una seconda funzione di sistema, \funcd{umount2}, che
consente un maggior controllo delle operazioni, come forzare lo smontaggio di
un filesystem anche quando questo risulti occupato; il suo prototipo è:

\begin{funcproto}{
\fhead{sys/mount.h}
\fdecl{umount2(const char *target, int flags)}
\fdesc{Smonta un filesystem.} 
}
{La funzione ritorna  $0$ in caso di successo e $-1$ per un errore,
  nel qual caso \var{errno} assumerà uno dei valori: 
  \begin{errlist}
     \item[\errcode{EAGAIN}] si è chiamata la funzione con \const{MNT\_EXPIRE}
       ed il filesystem non era occupato.
     \item[\errcode{EBUSY}] \param{target} è la directory di lavoro di qualche
       processo, o contiene dei file aperti, o un altro \textit{mount point}.
     \item[\errcode{EINVAL}] \param{target} non è un \textit{mount point} o si
       è usato \const{MNT\_EXPIRE} con \const{MNT\_FORCE} o
       \const{MNT\_DETACH} o si è specificato un flag non esistente.
  \end{errlist}
  e tutti gli altri valori visti per \func{umount} con lo stesso significato.}
\end{funcproto}

Il valore di \param{flags} è una maschera binaria dei flag che controllano le
modalità di smontaggio, che deve essere specificato con un OR aritmetico delle
costanti illustrate in tab.~\ref{tab:umount2_flags}.  Specificando
\constd{MNT\_FORCE} la funzione cercherà di liberare il filesystem anche se è
occupato per via di una delle condizioni descritte in precedenza. A seconda
del tipo di filesystem alcune (o tutte) possono essere superate, evitando
l'errore di \errcode{EBUSY}. In tutti i casi prima dello smontaggio viene
eseguita una sincronizzazione dei dati.

\begin{table}[!htb]
  \centering
  \footnotesize
  \begin{tabular}[c]{|l|p{8cm}|}
    \hline
    \textbf{Costante} & \textbf{Descrizione}\\
    \hline
    \hline
    \const{MNT\_FORCE}  & Forza lo smontaggio del filesystem anche se questo è
                           occupato (presente dai kernel della serie 2.2).\\
    \const{MNT\_DETACH} & Esegue uno smontaggio ``\textsl{pigro}'', in cui si
                           blocca l'accesso ma si aspetta che il filesystem si
                           liberi (presente dal kernel 2.4.11 e dalla
                           \acr{glibc} 2.11).\\ 
    \const{MNT\_EXPIRE} & Se non occupato marca un \textit{mount point} come
                           ``\textsl{in scadenza}'' in modo che ad una
                           successiva chiamata senza utilizzo del filesystem
                           questo venga smontato (presente dal 
                           kernel 2.6.8 e dalla \acr{glibc} 2.11).\\ 
    \const{UMOUNT\_NOFOLLOW}& Non dereferenzia \param{target} se questo è un
                               collegamento simbolico (vedi
                               sez.~\ref{sec:link_symlink_rename}) evitando
                               problemi di sicurezza (presente dal kernel
                               2.6.34).\\  
    \hline
  \end{tabular}
  \caption{Costanti che identificano i bit dell'argomento \param{flags}
    della funzione \func{umount2}.} 
  \label{tab:umount2_flags}
\end{table}

Con l'opzione \constd{MNT\_DETACH} si richiede invece uno smontaggio
``\textsl{pigro}'' (o \textit{lazy umount}) in cui il filesystem diventa
inaccessibile per i nuovi processi subito dopo la chiamata della funzione, ma
resta accessibile per quelli che lo hanno ancora in uso e non viene smontato
fintanto che resta occupato.

Con \constd{MNT\_EXPIRE}, che non può essere specificato insieme agli altri
due, si marca il \textit{mount point} di un filesystem non occupato come
``\textsl{in scadenza}'', in tal caso \func{umount2} ritorna con un errore di
\errcode{EAGAIN}, mentre in caso di filesystem occupato si sarebbe ricevuto
\errcode{EBUSY}.  Una volta marcato, se nel frattempo non viene fatto nessun
uso del filesystem, ad una successiva chiamata con \const{MNT\_EXPIRE} questo
verrà smontato. Questo flag consente di realizzare un meccanismo che smonti
automaticamente i filesystem che restano inutilizzati per un certo periodo di
tempo.

Infine il flag \constd{UMOUNT\_NOFOLLOW} non dereferenzia \param{target} se
questo è un collegamento simbolico (vedi
sez.~\ref{sec:link_symlink_rename}). Questa è una misura di sicurezza
introdotta per evitare, per quei filesystem per il quale è prevista una
gestione diretta da parte degli utenti, come quelli basati su \itindex{FUSE}
FUSE,\footnote{il \textit{Filesystem in USEr space} (FUSE) è una delle più
  interessanti applicazioni del VFS che consente, tramite un opportuno modulo,
  di implementarne le funzioni in \textit{user space}, così da rendere
  possibile l'implementazione di un qualunque filesystem (con applicazioni di
  grande interesse come il filesystem cifrato \textit{encfs} o il filesystem
  di rete \textit{sshfs}) che possa essere usato direttamente per conto degli
  utenti.}  che si possano passare ai programmi che effettuano lo smontaggio
dei filesystem, che in genere sono privilegiati ma consentono di agire solo
sui propri \textit{mount point}, dei collegamenti simbolici che puntano ad
altri \textit{mount point}, ottenendo così la possibilità di smontare
qualunque filesystem.


Altre due funzioni di sistema specifiche di Linux,\footnote{esse si trovano
  anche su BSD, ma con una struttura diversa.} utili per ottenere in maniera
diretta informazioni riguardo al filesystem su cui si trova un certo file,
sono \funcd{statfs} e \funcd{fstatfs}, i cui prototipi sono:

\begin{funcproto}{
\fhead{sys/vfs.h}
\fdecl{int statfs(const char *path, struct statfs *buf)}
\fdecl{int fstatfs(int fd, struct statfs *buf)}
\fdesc{Restituiscono informazioni relative ad un filesystem.} 
}
{Le funzioni ritornano $0$ in caso di successo e $-1$ per un errore,
  nel qual caso \var{errno} assumerà uno dei valori: 
  \begin{errlist}
  \item[\errcode{ENOSYS}] il filesystem su cui si trova il file specificato
    non supporta la funzione.
  \end{errlist} ed inoltre \errval{EFAULT} ed \errval{EIO} per entrambe,
  \errval{EBADF} per \func{fstatfs}, \errval{ENOTDIR}, \errval{ENAMETOOLONG},
  \errval{ENOENT}, \errval{EACCES}, \errval{ELOOP} per \func{statfs} nel loro
  significato generico.}
\end{funcproto}

Queste funzioni permettono di ottenere una serie di informazioni generali
riguardo al filesystem su cui si trova il file specificato con un
\textit{pathname} per \func{statfs} e con un file descriptor (vedi
sez.~\ref{sec:file_fd}) per \func{statfs}.  Le informazioni vengono restituite
all'indirizzo \param{buf} di una struttura \struct{statfs} definita come in
fig.~\ref{fig:sys_statfs}, ed i campi che sono indefiniti per il filesystem in
esame sono impostati a zero.  I valori del campo \var{f\_type} sono definiti
per i vari filesystem nei relativi file di header dei sorgenti del kernel da
costanti del tipo \var{XXX\_SUPER\_MAGIC}, dove \var{XXX} in genere è il nome
del filesystem stesso.

\begin{figure}[!htb]
  \footnotesize \centering
  \begin{minipage}[c]{0.8\textwidth}
    \includestruct{listati/statfs.h}
  \end{minipage}
  \normalsize 
  \caption{La struttura \structd{statfs}.} 
  \label{fig:sys_statfs}
\end{figure}

\conffilebeg{/etc/mtab} 

La \acr{glibc} provvede infine una serie di funzioni per la gestione dei due
file \conffiled{/etc/fstab}\footnote{più precisamente \funcm{setfsent},
  \funcm{getfsent}, \funcm{getfsfile}, \funcm{getfsspec}, \funcm{endfsent}.}
ed \conffile{/etc/mtab}\footnote{più precisamente \funcm{setmntent},
  \funcm{getmntent},\funcm{getmntent\_r}, \funcm{addmntent},\funcm{endmntent},
  \funcm{hasmntopt}.} che convenzionalmente sono usati in quasi tutti i
sistemi unix-like per mantenere rispettivamente le informazioni riguardo ai
filesystem da montare e a quelli correntemente montati. Le funzioni servono a
leggere il contenuto di questi file in opportune strutture \structd{fstab} e
\structd{mntent}, e, nel caso di \conffile{/etc/mtab}, per inserire e
rimuovere le voci presenti nel file.

In generale si dovrebbero usare queste funzioni, in particolare quelle
relative a \conffile{/etc/mtab}, quando si debba scrivere un programma che
effettua il montaggio di un filesystem. In realtà in questi casi è molto più
semplice invocare direttamente il programma \cmd{mount}. Inoltre l'uso stesso
di \conffile{/etc/mtab} è considerato una pratica obsoleta, in quanto se non
aggiornato correttamente (cosa che è impossibile se la radice è montata in
sola lettura) il suo contenuto diventa fuorviante.

Per questo motivo il suo utilizzo viene deprecato ed in molti casi viene già
oggi sostituito da un collegamento simbolico a \procfile{/proc/mounts}, che
contiene una versione degli stessi contenuti (vale a dire l'elenco dei
filesystem montati) generata direttamente dal kernel, e quindi sempre
disponibile e sempre aggiornata. Per questo motivo tralasceremo la
trattazione, di queste funzioni, rimandando al manuale della \acr{glibc}
\cite{GlibcMan} per la documentazione completa.

\conffileend{/etc/mtab}

% TODO (bassa priorità) scrivere delle funzioni (getfsent e getmntent &C)
% TODO (bassa priorità) documentare ? swapon e swapoff (man 2 ...) 

% TODO con il 5.2 è stata introdotta una serie di nuove syscall per montare un
% filesystem, vedi https://lwn.net/Articles/759499/ e
% https://git.kernel.org/linus/f1b5618e013a 


\section{La gestione di file e directory}
\label{sec:file_dir}

In questa sezione esamineremo le funzioni usate per la manipolazione dei nomi
file e directory, per la creazione di collegamenti simbolici e diretti, per la
gestione e la lettura delle directory.  In particolare ci soffermeremo sulle
conseguenze che derivano dalla architettura di un filesystem unix-like
illustrata in sez.~\ref{sec:file_filesystem} per quanto attiene il
comportamento e gli effetti delle varie funzioni. Tratteremo infine la
directory di lavoro e le funzioni per la gestione di file speciali e
temporanei.


\subsection{La gestione dei nomi dei file}
\label{sec:link_symlink_rename}

% \subsection{Le funzioni \func{link} e \func{unlink}}
% \label{sec:file_link}

Una caratteristica comune a diversi sistemi operativi è quella di poter creare
dei nomi alternativi, come gli alias del vecchio MacOS o i collegamenti di
Windows o i nomi logici del VMS, che permettono di fare riferimento allo
stesso file chiamandolo con nomi diversi o accedendovi da directory diverse.
Questo è possibile anche in ambiente Unix, dove un nome alternativo viene
usualmente chiamato ``\textsl{collegamento}'' (o \textit{link}).  Data
l'architettura del sistema riguardo la gestione dei file vedremo però che ci
sono due metodi sostanzialmente diversi per fare questa operazione.

\itindbeg{hard~link}
\index{collegamento!diretto|(}

In sez.~\ref{sec:file_filesystem} abbiamo spiegato come la capacità di
chiamare un file con nomi diversi sia connaturata con l'architettura di un
filesystem per un sistema Unix, in quanto il nome del file che si trova in una
directory è solo un'etichetta associata ad un puntatore che permette di
ottenere il riferimento ad un \textit{inode}, e che è quest'ultimo che viene
usato dal kernel per identificare univocamente gli oggetti sul filesystem.

Questo significa che fintanto che si resta sullo stesso filesystem la
realizzazione di un \textit{link} è immediata: uno stesso file può avere tanti
nomi diversi, dati da altrettante associazioni diverse allo stesso
\textit{inode} effettuate tramite ``etichette'' diverse in directory
diverse. Si noti anche come nessuno di questi nomi possa assumere una
particolare preferenza o originalità rispetto agli altri, in quanto tutti
fanno comunque riferimento allo stesso \textit{inode} e quindi tutti
otterranno lo stesso file.

Quando si vuole aggiungere ad una directory una voce che faccia riferimento ad
un file già esistente come appena descritto, per ottenere quello che viene
denominato ``\textsl{collegamento diretto}'' (o \textit{hard link}), si deve
usare la funzione di sistema \funcd{link}, il cui prototipo è:

\begin{funcproto}{
\fhead{unistd.h}
\fdecl{int link(const char *oldpath, const char *newpath)}
\fdesc{Crea un nuovo collegamento diretto (\textit{hard link}).} 
}
{La funzione ritorna $0$ in caso di successo e $-1$ per un errore,
  nel qual caso \var{errno} assumerà uno dei valori: 
  \begin{errlist}
  \item[\errcode{EEXIST}] un file (o una directory) di nome \param{newpath}
    esiste già.
  \item[\errcode{EMLINK}] ci sono troppi collegamenti al file \param{oldpath}
    (il numero massimo è specificato dalla variabile \const{LINK\_MAX}, vedi
    sez.~\ref{sec:sys_limits}).
  \item[\errcode{EPERM}] il filesystem che contiene \param{oldpath} e
    \param{newpath} non supporta i collegamenti diretti, è una directory o per
    \param{oldpath} non si rispettano i criteri per i \textit{protected
      hardlink}.\footnotemark 
  \item[\errcode{EXDEV}] i file \param{oldpath} e \param{newpath} non fanno
    riferimento ad un filesystem montato sullo stesso 
    \textit{mount point}.
  \end{errlist} ed inoltre \errval{EACCES}, \errval{EDQUOT}, \errval{EFAULT},
  \errval{EIO}, \errval{ELOOP}, \errval{ENAMETOOLONG}, \errval{ENOENT},
  \errval{ENOMEM}, \errval{ENOSPC}, \errval{ENOTDIR}, \errval{EROFS} nel loro
  significato generico.}
\end{funcproto}

\footnotetext{i \textit{protected hardlink} sono una funzionalità di
  protezione introdotta con il kernel 3.16 (si veda
  sez.~\ref{sec:procadv_security_misc} per i dettagli) che limita la capacità
  di creare un \textit{hard link} ad un file qualunque.}

La funzione crea in \param{newpath} un collegamento diretto al file indicato
da \param{oldpath}. Per quanto detto la creazione di un nuovo collegamento
diretto non copia il contenuto del file, ma si limita a creare la voce
specificata da \param{newpath} nella directory corrispondente e l'unica
proprietà del file che verrà modificata sarà il numero di riferimenti al file
(il campo \var{i\_nlink} della struttura \kstruct{inode}, vedi
fig.~\ref{fig:kstruct_inode}) che verrà aumentato di uno. In questo modo lo
stesso file potrà essere acceduto sia con \param{newpath} che
con \param{oldpath}.

Per quanto dicevamo in sez.~\ref{sec:file_filesystem} la creazione di un
collegamento diretto è possibile solo se entrambi i \textit{pathname} sono
nello stesso filesystem ed inoltre esso deve supportare gli \textit{hard link}
(il meccanismo non è disponibile ad esempio con il filesystem \acr{vfat} di
Windows). In realtà la funzione ha un ulteriore requisito, e cioè che non solo
che i due file siano sullo stesso filesystem, ma anche che si faccia
riferimento ad essi all'interno dello stesso \textit{mount point}.\footnote{si
  tenga presente infatti, come detto in sez.~\ref{sec:filesystem_mounting},
  che a partire dal kernel 2.4 uno stesso filesystem può essere montato più
  volte su directory diverse.}
La funzione inoltre opera sia sui file ordinari che sugli altri oggetti del
filesystem, con l'eccezione delle directory. In alcune versioni di Unix solo
l'amministratore è in grado di creare un collegamento diretto ad un'altra
directory: questo viene fatto perché con una tale operazione è possibile
creare dei \textit{loop} nel filesystem (vedi fig.~\ref{fig:file_link_loop})
la cui rimozione diventerebbe piuttosto complicata.\footnote{occorrerebbe
  infatti eseguire il programma \cmd{fsck} per riparare il filesystem, perché
  in caso di \textit{loop} la directory non potrebbe essere più svuotata,
  contenendo comunque se stessa, e quindi non potrebbe essere rimossa.}

Data la pericolosità di questa operazione, e visto che i collegamenti
simbolici (che tratteremo a breve) ed i \textit{bind mount} (già visti in
sez.~\ref{sec:filesystem_mounting}) possono fornire la stessa funzionalità
senza questi problemi, nel caso di Linux questa capacità è stata completamente
disabilitata, e al tentativo di creare un collegamento diretto ad una
directory la funzione \func{link} restituisce sempre l'errore \errcode{EPERM}.

Un ulteriore comportamento peculiare di Linux è quello in cui si crea un
\textit{hard link} ad un collegamento simbolico. In questo caso lo standard
POSIX.1-2001 prevederebbe che quest'ultimo venga risolto e che il collegamento
sia effettuato rispetto al file cui esso punta, e che venga riportato un
errore qualora questo non esista o non sia un file. Questo era anche il
comportamento iniziale di Linux ma a partire dai kernel della serie
2.0.x\footnote{per la precisione il comportamento era quello previsto dallo
  standard POSIX fino al kernel di sviluppo 1.3.56, ed è stato temporaneamente
  ripristinato anche durante lo sviluppo della serie 2.1.x, per poi tornare al
  comportamento attuale fino ad oggi (per riferimento si veda
  \url{http://lwn.net/Articles/293902}).} è stato adottato un comportamento
che non segue più lo standard per cui l'\textit{hard link} viene creato nei
confronti del collegamento simbolico, e non del file cui questo punta. La
revisione POSIX.1-2008 lascia invece il comportamento dipendente
dall'implementazione, cosa che rende Linux conforme a questa versione
successiva dello standard.

\itindbeg{symbolic~link}
\index{collegamento!simbolico|(}

La ragione di questa differenza rispetto al vecchio standard, presente anche
in altri sistemi unix-like, è dovuta al fatto che un collegamento simbolico
può fare riferimento anche ad un file non esistente o a una directory, per i
quali l'\textit{hard link} non può essere creato, per cui la scelta di seguire
il collegamento simbolico è stata ritenuta una scelta scorretta nella
progettazione dell'interfaccia.  Infatti se non ci fosse il comportamento
adottato da Linux sarebbe impossibile creare un \textit{hard link} ad un
collegamento simbolico, perché la funzione lo risolverebbe e l'\textit{hard
  link} verrebbe creato verso la destinazione. Invece evitando di seguire lo
standard l'operazione diventa possibile, ed anche il comportamento della
funzione risulta molto più comprensibile. Tanto più che se proprio se si vuole
creare un \textit{hard link} rispetto alla destinazione di un collegamento
simbolico è sempre possibile farlo direttamente.\footnote{ciò non toglie che
  questo comportamento possa causare problemi, come nell'esempio descritto
  nell'articolo citato nella nota precedente, a programmi che non si aspettano
  questa differenza rispetto allo standard POSIX.}

Dato che \func{link} crea semplicemente dei nomi che fanno riferimenti agli
\textit{inode}, essa può funzionare soltanto per file che risiedono sullo
stesso filesystem e solo per un filesystem di tipo Unix.  Inoltre abbiamo
visto che in Linux non è consentito eseguire un collegamento diretto ad una
directory.

Per ovviare a queste limitazioni, come accennato all'inizio, i sistemi
unix-like supportano un'altra forma di collegamento, detta
``\textsl{collegamento simbolico}'' (o anche \textit{soft link} o
\textit{symbolic link}). In questo caso si tratta, come avviene in altri
sistemi operativi, di file speciali che contengono semplicemente il
riferimento ad un altro file (o directory). In questo modo è possibile
effettuare \textit{link} anche attraverso filesystem diversi, a file posti in
filesystem che non supportano i collegamenti diretti, a delle directory, ed
anche a file che non esistono ancora.

\itindend{hard~link}
\index{collegamento!diretto|)}

Il meccanismo funziona in quanto i \textit{symbolic link} sono riconosciuti
come tali dal kernel\footnote{è uno dei diversi tipi di file visti in
  tab.~\ref{tab:file_file_types}, contrassegnato come tale nell'\textit{inode}
  e riconoscibile dal valore del campo \var{st\_mode} della struttura
  \struct{stat} (vedi sez.~\ref{sec:file_stat}).} e tutta una serie di
funzioni di sistema (come \func{open} o \func{stat}) quando ricevono come
argomento il \textit{pathname} di un collegamento simbolico vanno
automaticamente ad operare sul file da esso specificato. La funzione di
sistema che permette di creare un nuovo collegamento simbolico è
\funcd{symlink}, ed il suo prototipo è:

\begin{funcproto}{
\fhead{unistd.h}
\fdecl{int symlink(const char *oldpath, const char *newpath)}
\fdesc{Crea un nuovo collegamento simbolico (\textit{symbolic link}).} 
}
{La funzione ritorna $0$ in caso di successo e $-1$ per un errore,
  nel qual caso \var{errno} assumerà uno dei valori: 
  \begin{errlist}
  \item[\errcode{EACCES}]  o non si hanno i permessi sulla directory in cui
    creare il \textit{link}.
  \item[\errcode{EEXIST}] esiste già un file \param{newpath}.
  \item[\errcode{ENOENT}] una componente di \param{newpath} non esiste o
    \param{oldpath} è una stringa vuota.
  \item[\errcode{EPERM}] il filesystem che contiene \param{newpath} non
    supporta i collegamenti simbolici.
  \item[\errcode{EROFS}] \param{newpath} è su un filesystem montato in sola
    lettura.
  \end{errlist} ed inoltre \errval{EDQUOT}, \errval{EFAULT}, \errval{EIO},
  \errval{ELOOP}, \errval{ENAMETOOLONG}, \errval{ENOMEM}, \errval{ENOSPC} e
  \errval{ENOTDIR} nel loro significato generico.}
\end{funcproto}


La funzione crea un nuovo collegamento simbolico \param{newpath} che fa
riferimento ad \param{oldpath}.  Si tenga presente che la funzione non
effettua nessun controllo sull'esistenza di un file di nome \param{oldpath},
ma si limita ad inserire il \textit{pathname} nel collegamento
simbolico. Pertanto un collegamento simbolico può anche riferirsi ad un file
che non esiste ed in questo caso si ha quello che viene chiamato un
\itindex{dangling~link} \textit{dangling link}, letteralmente un
\index{collegamento!ciondolante} ``\textsl{collegamento ciondolante}''. Ad
esempio possiamo usare il comando \cmd{ln} per creare un collegamento
simbolico nella nostra directory con:
\begin{Console}
piccardi@hain:~/gapil$ \textbf{ln -s /tmp/tmp_file symlink}
\end{Console}
%$
e questo avrà successo anche se \file{/tmp/tmp\_file} non esiste:
\begin{Console}
piccardi@hain:~/gapil$ \textbf{ls symlink}
symlink
\end{Console}
%$
ma questo può generare confusione, perché accedendo in lettura a
\file{symlink}, ad esempio con \cmd{cat}, otterremmo:
\begin{Console}
piccardi@hain:~/gapil$ \textbf{cat symlink}
cat: symlink: No such file or directory
\end{Console}
%$
con un errore che può sembrare sbagliato, dato che \cmd{ls} ci ha mostrato in
precedenza l'esistenza di \file{symlink}. Se invece andassimo a scrivere su
\file{symlink}, l'effetto sarebbe quello di ottenere la creazione di
\file{/tmp/tmp\_file} (che a quel punto verrebbe creato) senza errori.

Come accennato i collegamenti simbolici sono risolti automaticamente dal
kernel all'invocazione delle varie \textit{system call}. In
tab.~\ref{tab:file_symb_effect} si è riportato un elenco dei comportamenti
delle varie funzioni di sistema che operano sui file nei confronti della
risoluzione dei collegamenti simbolici, specificando quali li seguono e quali
invece possono operare direttamente sui loro contenuti.
\begin{table}[htb]
  \centering
  \footnotesize
  \begin{tabular}[c]{|l|c|c|}
    \hline
    \textbf{Funzione} & \textbf{Segue il link} & \textbf{Non segue il link} \\
    \hline 
    \hline 
    \func{access}   & $\bullet$ & --        \\
    \func{chdir}    & $\bullet$ & --        \\
    \func{chmod}    & $\bullet$ & --        \\
    \func{chown}    & --        & $\bullet$ \\
    \func{creat}    & $\bullet$ & --        \\
    \func{exec}     & $\bullet$ & --        \\
    \func{lchown}   & $\bullet$ & --        \\
    \func{link}\footnotemark & --        & $\bullet$ \\
    \func{lstat}    & --        & $\bullet$ \\
    \func{mkdir}    & $\bullet$ & --        \\
    \func{mkfifo}   & $\bullet$ & --        \\
    \func{mknod}    & $\bullet$ & --        \\
    \func{open}     & $\bullet$ & --        \\
    \func{opendir}  & $\bullet$ & --        \\
    \func{pathconf} & $\bullet$ & --        \\
    \func{readlink} & --        & $\bullet$ \\
    \func{remove}   & --        & $\bullet$ \\
    \func{rename}   & --        & $\bullet$ \\
    \func{stat}     & $\bullet$ & --        \\
    \func{truncate} & $\bullet$ & --        \\
    \func{unlink}   & --        & $\bullet$ \\
    \hline 
  \end{tabular}
  \caption{Uso dei collegamenti simbolici da parte di alcune funzioni.}
  \label{tab:file_symb_effect}
\end{table}

\footnotetext{a partire dalla serie 2.0, e contrariamente a quanto indicato
  dallo standard POSIX.1-2001.}

Si noti che non si è specificato il comportamento delle funzioni che operano
con i file descriptor (che tratteremo nel prossimo capitolo), in quanto la
risoluzione del collegamento simbolico viene in genere effettuata dalla
funzione che restituisce il file descriptor (normalmente la \func{open}, vedi
sez.~\ref{sec:file_open_close}) e tutte le operazioni seguenti fanno
riferimento solo a quest'ultimo.

Si tenga anche presente che a partire dal kernel 3.16, se si abilita la
funzionalità dei \textit{protected symlinks} (attiva di default in tutte le
distribuzioni più recenti) la risoluzione dei nomi attraverso un collegamento
simbolico può fallire per una serie di restrizione di sicurezza aggiuntive
imposte dal meccanismo (si consulti sez.~\ref{sec:procadv_security_misc} per i
dettagli).

Dato che, come indicato in tab.~\ref{tab:file_symb_effect}, funzioni come la
\func{open} seguono i collegamenti simbolici, occorrono funzioni apposite per
accedere alle informazioni del collegamento invece che a quelle del file a cui
esso fa riferimento. Quando si vuole leggere il contenuto di un collegamento
simbolico si usa la funzione di sistema \funcd{readlink}, il cui prototipo è:

\begin{funcproto}{
\fhead{unistd.h}
\fdecl{int readlink(const char *pathname, char *buff, size\_t size)}
\fdesc{Legge il contenuto di un collegamento simbolico.} 
}
{La funzione ritorna il numero di caratteri letti dentro \param{buff} in caso
  di successo e $-1$ per un errore,  nel qual caso \var{errno} assumerà uno
  dei valori:
  \begin{errlist}
  \item[\errcode{EACCES}] non si hanno i permessi di attraversamento di una
    delle directory del pathname
  \item[\errcode{EINVAL}] \param{pathname} non è un collegamento simbolico
    o \param{size} non è positiva.
  \end{errlist} ed inoltre \errval{EFAULT}, \errval{EIO}, \errval{ELOOP},
  \errval{ENAMETOOLONG}, \errval{ENOENT}, \errval{ENOMEM} e \errval{ENOTDIR}
  nel loro significato generico.}
\end{funcproto}

La funzione legge il \textit{pathname} a cui fa riferimento il collegamento
simbolico indicato dall'argomento \param{pathname} scrivendolo sul buffer
\param{buff} di dimensione \param{size}. Si tenga presente che la funzione non
termina la stringa con un carattere nullo e che se questa è troppo lunga la
tronca alla dimensione specificata da \param{size} per evitare di scrivere
dati oltre le dimensioni del buffer.

\begin{figure}[htb]
  \centering
  \includegraphics[width=8cm]{img/link_loop}
  \caption{Esempio di loop nel filesystem creato con un collegamento
    simbolico.}
  \label{fig:file_link_loop}
\end{figure}

Come accennato uno dei motivi per cui non sono consentiti \textit{hard link}
alle directory è che questi possono creare dei \textit{loop} nella risoluzione
dei nomi che non possono essere eliminati facilmente. Invece è sempre
possibile, ed in genere anche molto utile, creare un collegamento simbolico ad
una directory, anche se in questo caso si potranno ottenere anche dei
\textit{loop}.

La situazione è illustrata in fig.~\ref{fig:file_link_loop}, che riporta la
struttura della directory \file{/boot}. Come si vede si è creato al suo
interno un collegamento simbolico che punta di nuovo a
\file{/boot}.\footnote{il \textit{loop} mostrato in
  fig.~\ref{fig:file_link_loop} è stato usato per poter permettere a al
  \textit{bootloader} \cmd{grub} di vedere i file contenuti nella directory
  \file{/boot} con lo stesso \textit{pathname} con cui verrebbero visti dal
  sistema operativo, anche quando si trovano, come accade spesso, su una
  partizione separata (che \cmd{grub} all'avvio vedrebbe come \file{/}).} Un
\textit{loop} di di questo tipo però può causare problemi per tutti i
programmi che effettuano la scansione di una directory, e ad esempio se
lanciassimo un comando come \code{grep -r linux *}, il \textit{loop} nella
directory porterebbe ad esaminare \file{/boot}, \file{/boot/boot},
\file{/boot/boot/boot} e così via.

Per questo motivo il kernel e le librerie prevedono che nella risoluzione di
un \textit{pathname} possano essere seguiti fino ad un certo numero massimo di
collegamenti simbolici, il cui valore limite è specificato dalla costante
\constd{MAXSYMLINKS}. Se il limite viene superato si ha un errore ed
\var{errno} viene impostata al valore \errcode{ELOOP}, che nella quasi
totalità dei casi indica appunto che si è creato un collegamento simbolico che
fa riferimento ad una directory del suo stesso \textit{pathname}.


\itindend{symbolic~link}
\index{collegamento!simbolico|)}

Un'altra funzione relativa alla gestione dei nomi dei file, anche se a prima
vista parrebbe riguardare un argomento completamente diverso, è quella per la
cancellazione di un file. In realtà una \textit{system call} che serva proprio
a cancellare un file non esiste neanche perché, come accennato in
sez.~\ref{sec:file_filesystem}, quando in un sistema unix-like si richiede la
rimozione di un file, quello che si va a cancellare è soltanto la voce che
referenzia il suo \textit{inode} all'interno di una directory.

La funzione di sistema che consente di effettuare questa operazione, il cui
nome come si può notare ha poco a che fare con il concetto di rimozione, è
\funcd{unlink}, ed il suo prototipo è:

\begin{funcproto}{
\fhead{unistd.h}
\fdecl{int unlink(const char *pathname)}
\fdesc{Cancella un file.} 
}
{La funzione ritorna $0$ in caso di successo e $-1$ per un errore,
  nel qual caso \var{errno} assumerà uno dei valori:\footnotemark  
  \begin{errlist}
  \item[\errcode{EACCES}] non si ha il permesso di scrittura sulla directory
    che contiene \param{pathname} o quello di attraversamento per una delle
    directory superiori.
  \item[\errcode{EBUSY}] \param{pathname} non può essere rimosso perché è in
    uso da parte del sistema (in particolare per i cosiddetti \textit{silly
      renames} di NFS).
  \item[\errcode{EISDIR}] \param{pathname} si riferisce ad una
    directory.
  \item[\errcode{EPERM}] il filesystem non consente l'operazione, o la
    directory che contiene \param{pathname} ha lo \textit{sticky bit} e non si
    è il proprietario del file o non si hanno privilegi amministrativi. 
  \end{errlist} ed inoltre \errval{EFAULT}, \errval{EIO}, \errval{ELOOP},
  \errval{ENOENT}, \errval{ENOMEM}, \errval{ENOTDIR}, \errval{EROFS} nel loro
  significato generico.}
\end{funcproto}

\footnotetext{questa funzione su Linux ha alcune peculiarità nei codici di
  errore, in particolare riguardo la rimozione delle directory che non è mai
  permessa e che causa l'errore \errcode{EISDIR}; questo è un valore specifico
  di Linux non conforme allo standard POSIX che prescrive invece l'uso di
  \errcode{EPERM} in caso l'operazione non sia consentita o il processo non
  abbia privilegi sufficienti, valore che invece Linux usa anche se il
  filesystem non supporta la funzione, inoltre il codice \errcode{EBUSY} nel
  caso la directory sia occupata su Linux non esiste.}

La funzione elimina il nome specificato dall'argomento \param{pathname} nella
directory che lo contiene e decrementa il numero di riferimenti nel relativo
\textit{inode}; come per \func{link} queste due operazioni sono effettuate
all'interno della \textit{system call} in maniera atomica rispetto ai
processi.

Si ricordi che, anche se se ne è rimosso il nome, un file viene realmente
cancellato soltanto quando il numero di collegamenti mantenuto
nell'\textit{inode} diventa nullo; solo allora l'\textit{inode} viene
disallocato e lo spazio che il file occupava sul disco viene liberato.

Si tenga presente comunque che a questo si aggiunge sempre un'ulteriore
condizione e cioè che non ci siano processi che stiano ancora lavorando sul il
file. Come vedremo in sez.~\ref{sec:file_unix_interface} il kernel una tabella
di tutti file aperti da ciascun processo, che a sua volta contiene i
riferimenti agli \textit{inode} ad essi relativi. Prima di procedere alla
cancellazione dello spazio occupato su disco dal contenuto di un file il
kernel controlla anche questa tabella, per verificare che anche in essa non ci
sia più nessun riferimento all'\textit{inode} in questione, assicurandosi con
questo che nessun processo stia ancora usando il file.

Nel caso di socket, \textit{fifo} o file di dispositivo la funzione rimuove il
nome, e come per i file normali i processi che hanno aperto uno di questi
oggetti possono continuare ad utilizzarli.  Nel caso di cancellazione di un
\textit{link} simbolico, che consiste solo nel rimando ad un altro file,
questo viene immediatamente eliminato e non sarà più utilizzabile. 

Per cancellare una voce in una directory è necessario avere il permesso di
scrittura su di essa, dato che si va a rimuovere una voce dal suo contenuto, e
il diritto di esecuzione/attraversamento sulla directory che la contiene
(affronteremo in dettaglio l'argomento dei permessi di file e directory in
sez.~\ref{sec:file_access_control}). Se inoltre per la directory è impostato
lo \textit{sticky bit} (vedi sez.~\ref{sec:file_special_perm}), occorrerà
anche essere proprietari del file o proprietari della directory o avere i
privilegi di amministratore.

Questa caratteristica del sistema, che consente di usare un file anche se lo
si è ``cancellato'', può essere usata per essere sicuri di non lasciare file
temporanei su disco in caso di uscita imprevista di un programma. La tecnica è
quella di aprire un nuovo file e chiamare \func{unlink} su di esso subito
dopo, in questo modo il contenuto del file sarà sempre disponibile all'interno
del processo attraverso il suo file descriptor (vedi sez.~\ref{sec:file_fd}),
ma non ne resterà traccia in nessuna directory, inoltre lo spazio occupato su
disco verrà immediatamente rilasciato alla conclusione del processo, quando
tutti i file vengono chiusi.

Al contrario di quanto avviene con altri Unix, in Linux non è possibile usare
la funzione \func{unlink} sulle directory, che in tal caso fallisce con un
errore di \errcode{EISDIR}. Per cancellare una directory si deve usare la
apposita funzione di sistema \func{rmdir} (che vedremo in
sez.~\ref{sec:file_dir_creat_rem}), oppure la funzione \func{remove}.

Quest'ultima è la funzione prevista dallo standard ANSI C per effettuare una
cancellazione generica di un file o di una directory e viene usata in generale
anche per i sistemi operativi che non supportano gli \textit{hard link}. Nei
sistemi unix-like \funcd{remove} è equivalente ad usare in maniera trasparente
\func{unlink} per i file ed \func{rmdir} per le directory; il suo prototipo è:

\begin{funcproto}{
\fhead{stdio.h}
\fdecl{int remove(const char *pathname)}
\fdesc{Cancella un file o una directory.} 
}
{La funzione ritorna $0$ in caso di successo e $-1$ per un errore, nel qual
  caso \var{errno} assumerà uno dei valori relativi alla chiamata utilizzata,
  pertanto si può fare riferimento a quanto illustrato nelle descrizioni di
  \func{unlink} e \func{rmdir}.}
\end{funcproto}

La funzione utilizza la funzione \func{unlink} per cancellare i file (e si
applica anche a link simbolici, socket, \textit{fifo} e file di dispositivo) e
la funzione \func{rmdir} (vedi sez.~\ref{sec:file_dir_creat_rem}) per
cancellare le directory.\footnote{questo vale usando la \acr{glibc}; nella
  \acr{libc4} e nella \acr{libc5} la funzione \func{remove} era un semplice
  alias alla funzione \func{unlink} e quindi non poteva essere usata per le
  directory.}  Si tenga presente che, per alcune limitazioni del protocollo
NFS, utilizzare questa funzione su file che usano questo filesystem di rete
può comportare la scomparsa di file ancora in uso.

Infine per cambiare nome ad un file o a una directory si usa la funzione di
sistema \funcd{rename},\footnote{la funzione è definita dallo standard ANSI C,
  ma si applica solo per i file, lo standard POSIX estende la funzione anche
  alle directory.} il cui prototipo è:

\begin{funcproto}{
\fhead{stdio.h}
\fdecl{int rename(const char *oldpath, const char *newpath)}
\fdesc{Rinomina un file o una directory.} 
}
{La funzione ritorna $0$ in caso di successo e $-1$ per un errore,
  nel qual caso \var{errno} assumerà uno dei valori: 
  \begin{errlist}
  \item[\errcode{EACCESS}] manca il permesso di scrittura sulle directory
    contenenti \param{oldpath} e \param{newpath} o di attraversare 
    il loro \textit{pathname} o di scrivere su \param{newpath}
    se questa è una directory.
  \item[\errcode{EBUSY}] o \param{oldpath} o \param{newpath} sono in uso da
    parte di qualche processo (come directory di lavoro o come radice) o del
    sistema (come \textit{mount point}) ed il sistema non riesce a risolvere
    la situazione.
  \item[\errcode{EEXIST}] \param{newpath} è una directory già esistente e
    non è vuota (anche \errcode{ENOTEMPTY}).
  \item[\errcode{EINVAL}] \param{newpath} contiene un prefisso di
    \param{oldpath} o più in generale si è cercato di creare una directory come
    sotto-directory di sé stessa.
  \item[\errcode{EISDIR}] \param{newpath} è una directory mentre
    \param{oldpath} non è una directory.
  \item[\errcode{ENOTDIR}] uno dei componenti dei \textit{pathname} non è una
    directory o \param{oldpath} è una directory e 
    \param{newpath} esiste e non è una directory.
  \item[\errval{EPERM}] la directory contenente \param{oldpath} o quella
    contenente un \param{newpath} esistente hanno lo \textit{sticky bit} e non
    si è i proprietari dei rispettivi file (o non si hanno privilegi
    amministrativi) oppure il filesystem non supporta l'operazione. 
  \item[\errcode{EXDEV}] \param{oldpath} e \param{newpath} non sono sullo
    stesso filesystem e sotto lo stesso \textit{mount point}. 
  \end{errlist} ed inoltre \errval{EFAULT}, \errval{ELOOP}, \errval{EMLINK},
  \errval{ENAMETOOLONG}, \errval{ENOENT}, \errval{ENOMEM}, \errval{ENOSPC} e
  \errval{EROFS} nel loro significato generico.}
\end{funcproto}

La funzione rinomina in \param{newpath} il file o la directory indicati
dall'argomento \param{oldpath}. Il nome viene eliminato nella directory
originale e ricreato nella directory di destinazione mantenendo il riferimento
allo stesso \textit{inode}. Non viene spostato nessun dato e l'\textit{inode}
del file non subisce nessuna modifica in quanto le modifiche sono eseguite
sulle directory che contengono \param{newpath} ed \param{oldpath}.

Il vantaggio nell'uso di questa funzione al posto della chiamata successiva di
\func{link} e \func{unlink} è che l'operazione è eseguita atomicamente, non
c'è modifica, per quanto temporanea, al \textit{link count} del file e non può
esistere un istante in cui un altro processo possa trovare attivi entrambi i
nomi per lo stesso file se la destinazione non esiste o in cui questa sparisca
temporaneamente se già esiste.

Dato che opera in maniera analoga la funzione è soggetta alle stesse
restrizioni di \func{link}, quindi è necessario che \param{oldpath}
e \param{newpath} siano nello stesso filesystem e facciano riferimento allo
stesso \textit{mount point}, e che il filesystem supporti questo tipo di
operazione. Qualora questo non avvenga si dovrà effettuare l'operazione in
maniera non atomica copiando il file a destinazione e poi cancellando
l'originale.

Il comportamento della funzione è diverso a seconda che si voglia rinominare
un file o una directory. Se ci riferisce ad un file allora \param{newpath}, se
esiste, non deve essere una directory, altrimenti si avrà un errore di
\errcode{EISDIR}. Se \param{newpath} indica un file già esistente questo verrà
rimpiazzato atomicamente, ma nel caso in cui \func{rename} fallisca il kernel
assicura che esso non sarà toccato. I caso di sovrascrittura però potrà
esistere una breve finestra di tempo in cui sia \param{oldpath}
che \param{newpath} potranno fare entrambi riferimento al file che viene
rinominato.

Se \param{oldpath} è una directory allora \param{newpath}, se esistente, deve
essere una directory vuota, altrimenti si avranno gli errori \errcode{ENOTDIR}
(se non è una directory) o \errcode{ENOTEMPTY} o \errcode{EEXIST} (se non è
vuota). Chiaramente \param{newpath} non potrà contenere \param{oldpath} nel
suo \textit{pathname}, non essendo possibile rendere una directory una
sottodirectory di sé stessa, se questo fosse il caso si otterrebbe un errore
di \errcode{EINVAL}.

Se \param{oldpath} si riferisce ad un collegamento simbolico questo sarà
rinominato restando tale senza nessun effetto sul file a cui fa riferimento.
Se invece \param{newpath} esiste ed è un collegamento simbolico verrà
cancellato come qualunque altro file.  Infine qualora \param{oldpath}
e \param{newpath} siano due nomi che già fanno riferimento allo stesso file lo
standard POSIX prevede che la funzione non ritorni un errore, e semplicemente
non faccia nulla, lasciando entrambi i nomi.  Linux segue questo standard,
anche se, come fatto notare dal manuale della \acr{glibc}, il comportamento
più ragionevole sarebbe quello di cancellare \param{oldpath}.

In tutti i casi si dovranno avere i permessi di scrittura nelle directory
contenenti \param{oldpath} e \param{newpath}, e nel caso \param{newpath} sia
una directory vuota già esistente anche su di essa (perché dovrà essere
aggiornata la voce ``\texttt{..}''). Se poi le directory
contenenti \param{oldpath} o \param{newpath} hanno lo \textit{sticky bit}
attivo (vedi sez.~\ref{sec:file_special_perm}) si dovrà essere i proprietari
dei file (o delle directory) che si vogliono rinominare, o avere i permessi di
amministratore.


\subsection{La creazione e la cancellazione delle directory} 
\label{sec:file_dir_creat_rem}

Benché in sostanza le directory non siano altro che dei file contenenti
elenchi di nomi con riferimenti ai rispettivi \textit{inode}, non è possibile
trattarle come file ordinari e devono essere create direttamente dal kernel
attraverso una opportuna \textit{system call}.\footnote{questo è quello che
  permette anche, attraverso l'uso del VFS, l'utilizzo di diversi formati per
  la gestione dei suddetti elenchi, dalle semplici liste a strutture complesse
  come alberi binari, hash, ecc. che consentono una ricerca veloce quando il
  numero di file è molto grande.}  La funzione di sistema usata per creare una
directory è \funcd{mkdir}, ed il suo prototipo è:

\begin{funcproto}{
\fhead{sys/stat.h}
\fhead{sys/types.h}
\fdecl{int mkdir(const char *dirname, mode\_t mode)}
\fdesc{Crea una nuova directory.} 
}
{La funzione ritorna $0$ in caso di successo e $-1$ per un errore, nel qual
  caso \var{errno} assumerà uno dei valori: 
  \begin{errlist}
  \item[\errcode{EACCES}] non c'è il permesso di scrittura per la directory in
    cui si vuole inserire la nuova directory o di attraversamento per le
    directory al di sopra di essa.
  \item[\errcode{EEXIST}] un file o una directory o un collegamento simbolico
    con quel nome esiste già.
  \item[\errcode{EMLINK}] la directory in cui si vuole creare la nuova
    directory contiene troppi file; sotto Linux questo normalmente non avviene
    perché il filesystem standard consente la creazione di un numero di file
    maggiore di quelli che possono essere contenuti nel disco, ma potendo
    avere a che fare anche con filesystem di altri sistemi questo errore può
    presentarsi.
  \item[\errcode{ENOSPC}] non c'è abbastanza spazio sul file system per creare
    la nuova directory o si è esaurita la quota disco dell'utente.
  \end{errlist}
  ed inoltre \errval{EFAULT}, \errval{ELOOP}, \errval{ENAMETOOLONG},
  \errval{ENOENT}, \errval{ENOMEM}, \errval{ENOTDIR}, \errval{EPERM},
  \errval{EROFS} nel loro significato generico.}
\end{funcproto}

La funzione crea una nuova directory vuota, che contiene cioè solo le due voci
standard presenti in ogni directory (``\file{.}'' e ``\file{..}''), con il
nome indicato dall'argomento \param{dirname}. 

I permessi di accesso (vedi sez.~\ref{sec:file_access_control}) con cui la
directory viene creata sono specificati dall'argomento \param{mode}, i cui
possibili valori sono riportati in tab.~\ref{tab:file_permission_const}; si
tenga presente che questi sono modificati dalla maschera di creazione dei file
(si veda sez.~\ref{sec:file_perm_management}).  La titolarità della nuova
directory è impostata secondo quanto illustrato in
sez.~\ref{sec:file_ownership_management}.

Come accennato in precedenza per eseguire la cancellazione di una directory è
necessaria una specifica funzione di sistema, \funcd{rmdir}, il suo prototipo
è:

\begin{funcproto}{
\fhead{sys/stat.h}
\fdecl{int rmdir(const char *dirname)}
\fdesc{Cancella una directory.} 
}
{La funzione ritorna $0$ in caso di successo e $-1$ per un errore, nel qual
  caso \var{errno} assumerà uno dei valori: 
  \begin{errlist}
  \item[\errcode{EACCES}] non c'è il permesso di scrittura per la directory
    che contiene la directory che si vuole cancellare, o non c'è il permesso
    di attraversare (esecuzione) una delle directory specificate in
    \param{dirname}.
  \item[\errcode{EBUSY}] la directory specificata è la directory di lavoro o
    la radice di qualche processo o un \textit{mount point}.
  \item[\errcode{EINVAL}] si è usato ``\texttt{.}'' come ultimo componente
    di \param{dirname}.
  \item[\errcode{EPERM}] il filesystem non supporta la cancellazione di
    directory, oppure la directory che contiene \param{dirname} ha lo
    \textit{sticky bit} impostato e non si è i proprietari della directory o
    non si hanno privilegi amministrativi.
  \end{errlist}
  ed inoltre \errval{EFAULT}, \errval{ELOOP}, \errval{ENAMETOOLONG},
  \errval{ENOENT}, \errval{ENOMEM}, \errval{ENOTDIR}, \errcode{ENOTEMPTY} e
  \errval{EROFS} nel loro significato generico.}
\end{funcproto}


La funzione cancella la directory \param{dirname}, che deve essere vuota, la
directory deve cioè contenere le due voci standard ``\file{.}'' e
``\file{..}'' e niente altro.  Il nome può essere indicato con un
\textit{pathname} assoluto o relativo, ma si deve fare riferimento al nome
nella directory genitrice, questo significa che \textit{pathname} terminanti
in ``\file{.}'' e ``\file{..}'' anche se validi in altri contesti, causeranno
il fallimento della funzione.

Inoltre per eseguire la cancellazione, oltre ad essere vuota, occorre anche
che la directory non sia utilizzata, questo vuol dire anche che non deve
essere la directory di lavoro (vedi sez.~\ref{sec:file_work_dir}) o la
directory radice (vedi sez.~\ref{sec:file_chroot}) di nessun processo, od
essere utilizzata come \textit{mount point}.

Se la directory cancellata risultasse aperta in qualche processo per una
lettura dei suoi contenuti (vedi sez.~\ref{sec:file_dir_read}), pur
scomparendo dal filesystem e non essendo più possibile accedervi o crearvi
altri file, le risorse ad essa associate verrebbero disallocate solo alla
chiusura di tutti questi ulteriori riferimenti.


\subsection{Lettura e scansione delle directory}
\label{sec:file_dir_read}

Benché le directory alla fine non siano altro che dei file che contengono
delle liste di nomi associati ai relativi \textit{inode}, per il ruolo che
rivestono nella struttura del sistema non possono essere trattate come dei
normali file di dati. Ad esempio, onde evitare inconsistenze all'interno del
filesystem, solo il kernel può scrivere il contenuto di una directory, e non
può essere un processo a inserirvi direttamente delle voci con le usuali
funzioni di scrittura.

Ma se la scrittura e l'aggiornamento dei dati delle directory è compito del
kernel, sono molte le situazioni in cui i processi necessitano di poterne
leggere il contenuto. Benché questo possa essere fatto direttamente (vedremo
in sez.~\ref{sec:file_open_close} che è possibile aprire una directory come se
fosse un file, anche se solo in sola lettura) in generale il formato con cui
esse sono scritte può dipendere dal tipo di filesystem, tanto che, come
riportato in tab.~\ref{tab:file_file_operations}, il VFS prevede una apposita
funzione per la lettura delle directory.

\itindbeg{directory~stream}

Tutto questo si riflette nello standard POSIX\footnote{le funzioni erano
  presenti in SVr4 e 4.3BSD, la loro specifica è riportata in POSIX.1-2001.}
che ha introdotto una apposita interfaccia per la lettura delle directory,
basata sui cosiddetti \textit{directory stream}, chiamati così per l'analogia
con i \textit{file stream} dell'interfaccia standard ANSI C che vedremo in
sez.~\ref{sec:files_std_interface}. La prima funzione di questa interfaccia è
\funcd{opendir}, il cui prototipo è:

\begin{funcproto}{
\fhead{sys/types.h}
\fhead{dirent.h}
\fdecl{DIR *opendir(const char *dirname)}
\fdesc{Apre un \textit{directory stream}.} 
}
{La funzione ritorna un puntatore al \textit{directory stream} in caso di
  successo e \val{NULL} per un errore, nel qual caso \var{errno} assumerà uno
  dei valori \errval{EACCES}, \errval{EMFILE}, \errval{ENFILE},
  \errval{ENOENT}, \errval{ENOMEM} e \errval{ENOTDIR} nel loro significato
  generico.}
\end{funcproto}

La funzione apre un \textit{directory stream} per la directory
\param{dirname}, ritornando il puntatore ad un oggetto di tipo \typed{DIR} (che
è il tipo opaco usato dalle librerie per gestire i \textit{directory stream})
da usare per tutte le operazioni successive, la funzione inoltre posiziona lo
\textit{stream} sulla prima voce contenuta nella directory.

Si tenga presente che comunque la funzione opera associando il
\textit{directory stream} ad un opportuno file descriptor (vedi
sez.~\ref{sec:file_fd}) sottostante, sul quale vengono compiute le
operazioni. Questo viene sempre aperto impostando il flag di
\textit{close-on-exec} (si ricordi quanto detto in sez.~\ref{sec:proc_exec}),
così da evitare che resti aperto in caso di esecuzione di un altro programma.

Nel caso in cui sia necessario conoscere il file descriptor associato ad un
\textit{directory stream} si può usare la funzione
\funcd{dirfd},\footnote{questa funzione è una estensione introdotta con BSD
  4.3-Reno ed è presente in Linux con le libc5 (a partire dalla versione
  5.1.2) e con la \acr{glibc} ma non presente in POSIX fino alla revisione
  POSIX.1-2008, per questo per poterla utilizzare fino alla versione 2.10
  della \acr{glibc} era necessario definire le macro \macro{\_BSD\_SOURCE} o
  \macro{\_SVID\_SOURCE}, dalla versione 2.10 si possono usare anche
  \texttt{\macro{\_POSIX\_C\_SOURCE} >= 200809L} o
  \texttt{\macro{\_XOPEN\_SOURCE} >= 700}.}  il cui prototipo è:

\begin{funcproto}{
\fhead{sys/types.h}
\fhead{dirent.h}
\fdecl{int dirfd(DIR *dir)}
\fdesc{Legge il file descriptor associato ad un \textit{directory stream}.} 
}
{La funzione ritorna un valore positivo corrispondente al file descriptor in
  caso di successo e $-1$ per un errore, nel qual caso \var{errno} assumerà
  uno dei valori:
  \begin{errlist}
  \item[\errcode{EINVAL}] \param{dir} non è un puntatore ad un
    \textit{directory stream}. 
  \item[\errcode{ENOTSUP}] l'implementazione non supporta l'uso di un file
    descriptor per la directory.
  \end{errlist}
}
\end{funcproto}

La funzione restituisce il file descriptor associato al \textit{directory
  stream} \param{dir}. Di solito si utilizza questa funzione in abbinamento a
funzioni che operano sui file descriptor, ad esempio si potrà usare
\func{fstat} per ottenere le proprietà della directory, o \func{fchdir} per
spostare su di essa la directory di lavoro (vedi sez.~\ref{sec:file_work_dir}).

Viceversa se si è aperto un file descriptor corrispondente ad una directory è
possibile associarvi un \textit{directory stream} con la funzione
\funcd{fdopendir},\footnote{questa funzione è però disponibile solo a partire
  dalla versione 2.4 della \acr{glibc}, ed è stata introdotta nello standard
  POSIX solo a partire dalla revisione POSIX.1-2008, prima della versione 2.10
  della \acr{glibc} per poterla utilizzare era necessario definire la macro
  \macro{\_GNU\_SOURCE}, dalla versione 2.10 si possono usare anche
  \texttt{\macro{\_POSIX\_C\_SOURCE} >= 200809L} o \texttt{\_XOPEN\_SOURCE >=
    700} .}  il cui prototipo è:

\begin{funcproto}{
\fhead{sys/types.h}
\fhead{dirent.h}
\fdecl{DIR *fdopendir(int fd)}
\fdesc{Associa un \textit{directory stream} ad un file descriptor.} 
}
{La funzione ritorna un puntatore al \textit{directory stream} in caso di
  successo e \val{NULL} per un errore, nel qual caso \var{errno} assumerà uno
  dei valori \errval{EBADF} o \errval{ENOMEM} nel loro significato generico.}
\end{funcproto}

La funzione è identica a \func{opendir}, ma ritorna un \textit{directory
  stream} facendo riferimento ad un file descriptor \param{fd} che deve essere
stato aperto in precedenza; la funzione darà un errore qualora questo non
corrisponda ad una directory. L'uso di questa funzione permette di rispondere
agli stessi requisiti delle funzioni ``\textit{at}'' che vedremo in
sez.~\ref{sec:file_openat}.

Una volta utilizzata il file descriptor verrà usato internamente dalle
funzioni che operano sul \textit{directory stream} e non dovrà essere più
utilizzato all'interno del proprio programma. In particolare dovrà essere
chiuso attraverso il \textit{directory stream} con \func{closedir} e non
direttamente. Si tenga presente inoltre che \func{fdopendir} non modifica lo
stato di un eventuale flag di \textit{close-on-exec}, che pertanto dovrà
essere impostato esplicitamente in fase di apertura del file descriptor.

Una volta che si sia aperto un \textit{directory stream} la lettura del
contenuto della directory viene effettuata attraverso la funzione
\funcd{readdir}, il cui prototipo è:

\begin{funcproto}{
\fhead{sys/types.h}
\fhead{dirent.h}
\fdecl{struct dirent *readdir(DIR *dir)}
\fdesc{Legge una voce dal \textit{directory stream}.} 
}
{La funzione ritorna il puntatore alla struttura contenente i dati in caso di
  successo e \val{NULL} per un errore o se si è raggiunta la fine dello
  \textit{stream}. Il solo codice di errore restituito in \var{errno} è
  \errval{EBADF} qualora \param{dir} non indichi un \textit{directory stream}
  valido.}
\end{funcproto}

La funzione legge la voce corrente nella directory, posizionandosi sulla voce
successiva. Pertanto se si vuole leggere l'intero contenuto di una directory
occorrerà ripetere l'esecuzione della funzione fintanto che non si siano
esaurite tutte le voci in essa presenti, che viene segnalata dalla
restituzione di \val{NULL} come valore di ritorno. Si può distinguere questa
condizione da un errore in quanto in questo caso \var{errno} non verrebbe
modificata.

I dati letti da \func{readdir} vengono memorizzati in una struttura
\struct{dirent}, la cui definizione è riportata in
fig.~\ref{fig:file_dirent_struct}.\footnote{la definizione è quella usata da
  Linux, che si trova nel file \file{/usr/include/bits/dirent.h}, essa non
  contempla la presenza del campo \var{d\_namlen} che indica la lunghezza del
  nome del file.} La funzione non è rientrante e restituisce il puntatore ad
una struttura allocata staticamente, che viene sovrascritta tutte le volte che
si ripete la lettura di una voce sullo stesso \textit{directory stream}.

Di questa funzione esiste anche una versione rientrante,
\funcd{readdir\_r},\footnote{per usarla è necessario definire una qualunque
  delle macro \texttt{\macro{\_POSIX\_C\_SOURCE} >= 1},
  \macro{\_XOPEN\_SOURCE}, \macro{\_BSD\_SOURCE}, \macro{\_SVID\_SOURCE},
  \macro{\_POSIX\_SOURCE}.} che non usa una struttura allocata staticamente, e
può essere utilizzata anche con i \textit{thread}, il suo prototipo è:

\begin{funcproto}{
\fhead{sys/types.h}
\fhead{dirent.h}
\fdecl{int readdir\_r(DIR *dir, struct dirent *entry, struct dirent **result)}
\fdesc{Legge una voce dal \textit{directory stream}.} 
}
{La funzione ritorna $0$ in caso di successo ed un numero positivo per un
  errore, nel qual caso \var{errno} assumerà gli stessi valori di
  \func{readdir}.} 
\end{funcproto}

La funzione restituisce in \param{result} come \textit{value result argument}
l'indirizzo della struttura \struct{dirent} dove sono stati salvati i dati,
che deve essere allocata dal chiamante, ed il cui indirizzo deve essere
indicato con l'argomento \param{entry}.  Se si è raggiunta la fine del
\textit{directory stream} invece in \param{result} viene restituito il valore
\val{NULL}.

\begin{figure}[!htb]
  \footnotesize \centering
  \begin{minipage}[c]{0.8\textwidth}
    \includestruct{listati/dirent.c}
  \end{minipage} 
  \normalsize 
  \caption{La struttura \structd{dirent} per la lettura delle informazioni dei 
    file.}
  \label{fig:file_dirent_struct}
\end{figure}

% Lo spazio per la \struct{dirent} dove vengono restituiti i dati della
% directory deve essere allocato a cura del chiamante, dato che la dimensione


I vari campi di \struct{dirent} contengono le informazioni relative alle voci
presenti nella directory. Sia BSD che SVr4 che POSIX.1-2001\footnote{il
  vecchio standard POSIX prevedeva invece solo la presenza del campo
  \var{d\_fileno}, identico \var{d\_ino}, che in Linux era definito come alias
  di quest'ultimo, mentre il campo \var{d\_name} era considerato dipendente
  dall'implementazione.}  prevedono che siano sempre presenti il campo
\var{d\_name}, che contiene il nome del file nella forma di una stringa
terminata da uno zero, ed il campo \var{d\_ino}, che contiene il numero di
\textit{inode} cui il file è associato e corrisponde al campo \var{st\_ino} di
\struct{stat}.  La presenza di ulteriori campi opzionali oltre i due citati è
segnalata dalla definizione di altrettante macro nella forma
\code{\_DIRENT\_HAVE\_D\_XXX} dove \code{XXX} è il nome del relativo
campo. Come si può evincere da fig.~\ref{fig:file_dirent_struct} nel caso di
Linux sono pertanto definite le macro \macrod{\_DIRENT\_HAVE\_D\_TYPE},
\macrod{\_DIRENT\_HAVE\_D\_OFF} e \macrod{\_DIRENT\_HAVE\_D\_RECLEN}, mentre non
è definita la macro \macrod{\_DIRENT\_HAVE\_D\_NAMLEN}.

Dato che possono essere presenti campi opzionali e che lo standard
POSIX.1-2001 non specifica una dimensione definita per il nome dei file (che
può variare a seconda del filesystem), ma solo un limite superiore pari a
\const{NAME\_MAX} (vedi tab.~\ref{tab:sys_file_macro}), in generale per
allocare una struttura \struct{dirent} in maniera portabile occorre eseguire
un calcolo per ottenere le dimensioni appropriate per il proprio
sistema.\footnote{in SVr4 la lunghezza del campo è definita come
  \code{NAME\_MAX+1} che di norma porta al valore di 256 byte usato anche in
  fig.~\ref{fig:file_dirent_struct}.} Lo standard però richiede che il campo
\var{d\_name} sia sempre l'ultimo della struttura, questo ci consente di
ottenere la dimensione della prima parte con la macro di utilità generica
\macro{offsetof}, che si può usare con il seguente prototipo:

{\centering
\vspace{3pt}
\begin{funcbox}{
\fhead{stddef.h}
\fdecl{size\_t \macrod{offsetof}(type, member)}
\fdesc{Restituisce la posizione del campo \param{member} nella
  struttura \param{type}.}
} 
\end{funcbox}
}

Ottenuta allora con \code{offsetof(struct dirent, d\_name)} la dimensione
della parte iniziale della struttura, basterà sommarci la dimensione massima
dei nomi dei file nel filesystem che si sta usando, che si può ottenere
attraverso la funzione \func{pathconf} (per la quale si rimanda a
sez.~\ref{sec:sys_file_limits}) più un ulteriore carattere per la terminazione
della stringa.

Per quanto riguarda il significato dei campi opzionali, il campo \var{d\_type}
indica il tipo di file (se \textit{fifo}, directory, collegamento simbolico,
ecc.), e consente di evitare una successiva chiamata a \func{lstat} (vedi
sez.~\ref{sec:file_stat}) per determinarlo. I suoi possibili valori sono
riportati in tab.~\ref{tab:file_dtype_macro}. Si tenga presente che questo
valore è disponibile solo per i filesystem che ne supportano la restituzione
(fra questi i più noti sono \acr{Btrfs}, \acr{ext2}, \acr{ext3}, e
\acr{ext4}), per gli altri si otterrà sempre il valore
\const{DT\_UNKNOWN}.\footnote{inoltre fino alla versione 2.1 della
  \acr{glibc}, pur essendo il campo \var{d\_type} presente, il suo uso non era
  implementato, e veniva restituito comunque il valore \const{DT\_UNKNOWN}.}

\begin{table}[htb]
  \centering
  \footnotesize
  \begin{tabular}[c]{|l|l|}
    \hline
    \textbf{Valore} & \textbf{Tipo di file} \\
    \hline
    \hline
    \constd{DT\_UNKNOWN} & Tipo sconosciuto.\\
    \constd{DT\_REG}     & File normale.\\
    \constd{DT\_DIR}     & Directory.\\
    \constd{DT\_LNK}     & Collegamento simbolico.\\
    \constd{DT\_FIFO}    & \textit{Fifo}.\\
    \constd{DT\_SOCK}    & Socket.\\
    \constd{DT\_CHR}     & Dispositivo a caratteri.\\
    \constd{DT\_BLK}     & Dispositivo a blocchi.\\
    \hline    
  \end{tabular}
  \caption{Costanti che indicano i vari tipi di file nel campo \var{d\_type}
    della struttura \struct{dirent}.}
  \label{tab:file_dtype_macro}
\end{table}

Per la conversione da e verso l'analogo valore mantenuto dentro il campo
\var{st\_mode} di \struct{stat} (vedi fig.~\ref{fig:file_stat_struct}) sono
definite anche due macro di conversione, \macro{IFTODT} e \macro{DTTOIF}:

{\centering
\vspace{3pt}
\begin{funcbox}{
\fhead{dirent.h}
\fdecl{int \macrod{IFTODT}(mode\_t MODE)}
\fdesc{Converte il tipo di file dal formato di \var{st\_mode} a quello di
  \var{d\_type}.}
\fdecl{mode\_t \macrod{DTTOIF}(int DTYPE)}
\fdesc{Converte il tipo di file dal formato di \var{d\_type} a quello di
  \var{st\_mode}.}  
} 
\end{funcbox}
}

Il campo \var{d\_off} contiene invece la posizione della voce successiva della
directory, mentre il campo \var{d\_reclen} la lunghezza totale della voce
letta. Con questi due campi diventa possibile, determinando la posizione delle
varie voci, spostarsi all'interno dello \textit{stream} usando la funzione
\funcd{seekdir},\footnote{sia questa funzione che \func{telldir}, sono
  estensioni prese da BSD, ed introdotte nello standard POSIX solo a partire
  dalla revisione POSIX.1-2001, per poterle utilizzare deve essere definita
  una delle macro \macro{\_XOPEN\_SOURCE}, \macro{\_BSD\_SOURCE} o
  \macro{\_SVID\_SOURCE}.} il cui prototipo è:

\begin{funcproto}{
\fhead{dirent.h}
\fdecl{void seekdir(DIR *dir, off\_t offset)}
\fdesc{Cambia la posizione all'interno di un \textit{directory stream}.} 
}
{La funzione non ritorna niente e non imposta errori.}
\end{funcproto}

La funzione non ritorna nulla e non segnala errori, è però necessario che il
valore dell'argomento \param{offset} sia valido per lo
\textit{stream} \param{dir}; esso pertanto deve essere stato ottenuto o dal
valore di \var{d\_off} di \struct{dirent} o dal valore restituito dalla
funzione \funcd{telldir}, che legge la posizione corrente; il cui prototipo
è:\footnote{prima della \acr{glibc} 2.1.1 la funzione restituiva un valore di
  tipo \type{off\_t}, sostituito a partire dalla versione 2.1.2 da \ctyp{long}
  per conformità a POSIX.1-2001.}

\begin{funcproto}{
\fhead{dirent.h}
\fdecl{long telldir(DIR *dir)}
\fdesc{Ritorna la posizione corrente in un \textit{directory stream}.} 
}
{La funzione ritorna la posizione corrente nello \textit{stream} (un numero
  positivo) in caso di successo e $-1$ per un errore, nel qual caso
  \var{errno} assume solo il valore di \errval{EBADF}, corrispondente ad un
  valore errato per \param{dir}.  }
\end{funcproto}

La sola funzione di posizionamento per un \textit{directory stream} prevista
originariamente dallo standard POSIX è \funcd{rewinddir}, che riporta la
posizione a quella iniziale; il suo prototipo è:

\begin{funcproto}{
\fhead{sys/types.h}
\fhead{dirent.h}
\fdecl{void rewinddir(DIR *dir)}
\fdesc{Si posiziona all'inizio di un \textit{directory stream}.} 
}
{La funzione non ritorna niente e non imposta errori.}
\end{funcproto}

Una volta completate le operazioni si può chiudere il \textit{directory
  stream}, ed il file descriptor ad esso associato, con la funzione
\funcd{closedir}, il cui prototipo è:

\begin{funcproto}{
\fhead{sys/types.h}
\fhead{dirent.h}
\fdecl{int closedir(DIR *dir)}
\fdesc{Chiude un \textit{directory stream}.} 
}
{La funzione ritorna $0$ in caso di successo e $-1$ per un errore, nel qual
  caso \var{errno} assume solo il valore \errval{EBADF}.}
\end{funcproto}

A parte queste funzioni di base in BSD 4.3 venne introdotta un'altra funzione
che permette di eseguire una scansione completa, con tanto di ricerca ed
ordinamento, del contenuto di una directory; la funzione è
\funcd{scandir}\footnote{in Linux questa funzione è stata introdotta fin dalle
  \acr{libc4} e richiede siano definite le macro \macro{\_BSD\_SOURCE} o
  \macro{\_SVID\_SOURCE}.} ed il suo prototipo è:

\begin{funcproto}{
\fhead{dirent.h}
\fdecl{int scandir(const char *dir, struct dirent ***namelist, \\
\phantom{int scandir(}int(*filter)(const struct dirent *), \\
\phantom{int scandir(}int(*compar)(const struct dirent **, const struct dirent **))}
\fdesc{Esegue una scansione di un \textit{directory stream}.} 
}
{La funzione ritorna il numero di voci trovate in caso di successo e $-1$ per
  un errore, nel qual caso \var{errno} può assumere solo il valore
  \errval{ENOMEM}.}
\end{funcproto}

Al solito, per la presenza fra gli argomenti di due puntatori a funzione, il
prototipo non è molto comprensibile; queste funzioni però sono quelle che
controllano rispettivamente la selezione di una voce, passata con
l'argomento \param{filter}, e l'ordinamento di tutte le voci selezionate,
specificata dell'argomento \param{compar}.

La funzione legge tutte le voci della directory indicata dall'argomento
\param{dir}, passando un puntatore a ciascuna di esse (una struttura
\struct{dirent}) come argomento della funzione di selezione specificata da
\param{filter}; se questa ritorna un valore diverso da zero il puntatore viene
inserito in un vettore che viene allocato dinamicamente con \func{malloc}.
Qualora si specifichi un valore \val{NULL} per l'argomento \param{filter} non
viene fatta nessuna selezione e si ottengono tutte le voci presenti.

Le voci selezionate possono essere riordinate tramite \funcm{qsort}, le
modalità del riordinamento possono essere personalizzate usando la funzione
\param{compar} come criterio di ordinamento di \funcm{qsort}, la funzione
prende come argomenti le due strutture \struct{dirent} da confrontare
restituendo un valore positivo, nullo o negativo per indicarne l'ordinamento;
alla fine l'indirizzo della lista ordinata dei puntatori alle strutture
\struct{dirent} viene restituito nell'argomento
\param{namelist}.\footnote{la funzione alloca automaticamente la lista, e
  restituisce, come \textit{value result argument}, l'indirizzo della stessa;
  questo significa che \param{namelist} deve essere dichiarato come
  \code{struct dirent **namelist} ed alla funzione si deve passare il suo
  indirizzo.}

\itindend{directory~stream}

Per l'ordinamento, vale a dire come valori possibili per l'argomento
\param{compar}, sono disponibili due funzioni predefinite, \funcd{alphasort} e
\funcd{versionsort}, i cui prototipi sono:

\begin{funcproto}{
\fhead{dirent.h}
\fdecl{int alphasort(const void *a, const void *b)}
\fdecl{int versionsort(const void *a, const void *b)}
\fdesc{Riordinano le voci di \textit{directory stream}.} 
}
{Le funzioni restituiscono un valore minore, uguale o maggiore di zero qualora
  il primo argomento sia rispettivamente minore, uguale o maggiore del secondo
  e non forniscono errori.}
\end{funcproto}

La funzione \func{alphasort} deriva da BSD ed è presente in Linux fin dalle
\acr{libc4}\footnote{la versione delle \acr{libc4} e \acr{libc5} usa però come
  argomenti dei puntatori a delle strutture \struct{dirent}; la glibc usa il
  prototipo originario di BSD, mostrato anche nella definizione, che prevede
  puntatori a \ctyp{void}.} e deve essere specificata come argomento
\param{compar} per ottenere un ordinamento alfabetico secondo il valore del
campo \var{d\_name} delle varie voci. La \acr{glibc} prevede come
estensione\footnote{la \acr{glibc}, a partire dalla versione 2.1, effettua
  anche l'ordinamento alfabetico tenendo conto delle varie localizzazioni,
  usando \funcm{strcoll} al posto di \funcm{strcmp}.} anche
\func{versionsort}, che ordina i nomi tenendo conto del numero di versione,
cioè qualcosa per cui \texttt{file10} viene comunque dopo \texttt{file4}.

\begin{figure}[!htb]
  \footnotesize \centering
  \begin{minipage}[c]{\codesamplewidth}
    \includecodesample{listati/my_ls.c}
  \end{minipage}
  \caption{Esempio di codice per eseguire la lista dei file contenuti in una
    directory.} 
  \label{fig:file_my_ls}
\end{figure}

Un semplice esempio dell'uso di queste funzioni è riportato in
fig.~\ref{fig:file_my_ls}, dove si è riportata la sezione principale di un
programma che, usando la funzione di scansione illustrata in
fig.~\ref{fig:file_dirscan}, stampa i nomi dei file contenuti in una directory
e la relativa dimensione, in sostanza una versione semplificata del comando
\cmd{ls}.

Il programma è estremamente semplice; in fig.~\ref{fig:file_my_ls} si è omessa
la parte di gestione delle opzioni, che prevede solo l'uso di una funzione per
la stampa della sintassi, anch'essa omessa, ma il codice completo può essere
trovato coi sorgenti allegati alla guida nel file \file{myls.c}.

In sostanza tutto quello che fa il programma, dopo aver controllato
(\texttt{\small 12-15}) di avere almeno un argomento, che indicherà la
directory da esaminare, è chiamare (\texttt{\small 16}) la funzione
\myfunc{dir\_scan} per eseguire la scansione, usando la funzione \code{do\_ls}
(\texttt{\small 22-29}) per fare tutto il lavoro.

Quest'ultima si limita (\texttt{\small 26}) a chiamare \func{stat} sul file
indicato dalla directory entry passata come argomento (il cui nome è appunto
\var{direntry->d\_name}), memorizzando in una opportuna struttura \var{data} i
dati ad esso relativi, per poi provvedere (\texttt{\small 27}) a stampare il
nome del file e la dimensione riportata in \var{data}.

Dato che la funzione verrà chiamata all'interno di \myfunc{dir\_scan} per ogni
voce presente, questo è sufficiente a stampare la lista completa dei file e
delle relative dimensioni. Si noti infine come si restituisca sempre 0 come
valore di ritorno per indicare una esecuzione senza errori.

\begin{figure}[!htb]
  \footnotesize \centering
  \begin{minipage}[c]{\codesamplewidth}
    \includecodesample{listati/dir_scan.c}
  \end{minipage}
  \caption{Codice della funzione di scansione di una directory contenuta nel
    file \file{dir\_scan.c}.} 
  \label{fig:file_dirscan}
\end{figure}

Tutto il grosso del lavoro è svolto dalla funzione \myfunc{dir\_scan},
riportata in fig.~\ref{fig:file_dirscan}. La funzione è volutamente generica e
permette di eseguire una funzione, passata come secondo argomento, su tutte le
voci di una directory.  La funzione inizia con l'aprire (\texttt{\small
  18-22}) uno \textit{stream} sulla directory passata come primo argomento,
stampando un messaggio in caso di errore.

Il passo successivo (\texttt{\small 23-24}) è cambiare directory di lavoro
(vedi sez.~\ref{sec:file_work_dir}), usando in sequenza le funzioni
\func{dirfd} e \func{fchdir} (in realtà si sarebbe potuto usare direttamente
\func{chdir} su \var{dirname}), in modo che durante il successivo ciclo
(\texttt{\small 26-30}) sulle singole voci dello \textit{stream} ci si trovi
all'interno della directory.\footnote{questo è essenziale al funzionamento
  della funzione \code{do\_ls}, e ad ogni funzione che debba usare il campo
  \var{d\_name}, in quanto i nomi dei file memorizzati all'interno di una
  struttura \struct{dirent} sono sempre relativi alla directory in questione,
  e senza questo posizionamento non si sarebbe potuto usare \func{stat} per
  ottenere le dimensioni.}

Avendo usato lo stratagemma di fare eseguire tutte le manipolazioni necessarie
alla funzione passata come secondo argomento, il ciclo di scansione della
directory è molto semplice; si legge una voce alla volta (\texttt{\small 26})
all'interno di una istruzione di \code{while} e fintanto che si riceve una
voce valida, cioè un puntatore diverso da \val{NULL}, si esegue
(\texttt{\small 27}) la funzione di elaborazione \var{compare} (che nel nostro
caso sarà \code{do\_ls}), ritornando con un codice di errore (\texttt{\small
  28}) qualora questa presenti una anomalia, identificata da un codice di
ritorno negativo. Una volta terminato il ciclo la funzione si conclude con la
chiusura (\texttt{\small 31}) dello \textit{stream}\footnote{nel nostro caso,
  uscendo subito dopo la chiamata, questo non servirebbe, in generale però
  l'operazione è necessaria, dato che la funzione può essere invocata molte
  volte all'interno dello stesso processo, per cui non chiudere i
  \textit{directory stream} comporterebbe un consumo progressivo di risorse,
  con conseguente rischio di esaurimento delle stesse.} e la restituzione
(\texttt{\small 32}) del codice di operazioni concluse con successo.



\subsection{La directory di lavoro}
\label{sec:file_work_dir}

\index{directory~di~lavoro|(} 

Come accennato in sez.~\ref{sec:proc_fork} a ciascun processo è associata una
directory nell'albero dei file,\footnote{questa viene mantenuta all'interno
  dei dati della sua \kstruct{task\_struct} (vedi
  fig.~\ref{fig:proc_task_struct}), più precisamente nel campo \texttt{pwd}
  della sotto-struttura \kstruct{fs\_struct}.} che è chiamata
\textsl{directory corrente} o \textsl{directory di lavoro} (in inglese
\textit{current working directory}).  La directory di lavoro è quella da cui
si parte quando un \textit{pathname} è espresso in forma relativa, dove il
``\textsl{relativa}'' fa riferimento appunto a questa directory.

Quando un utente effettua il login, questa directory viene impostata alla
\textit{home directory} del suo account. Il comando \cmd{cd} della shell
consente di cambiarla a piacere, spostandosi da una directory ad un'altra, il
comando \cmd{pwd} la stampa sul terminale.  Siccome la directory di lavoro
resta la stessa quando viene creato un processo figlio (vedi
sez.~\ref{sec:proc_fork}), la directory di lavoro della shell diventa anche la
directory di lavoro di qualunque comando da essa lanciato.

Dato che è il kernel che tiene traccia dell'\textit{inode} della directory di
lavoro di ciascun processo, per ottenerne il \textit{pathname} occorre usare
una apposita funzione, \funcd{getcwd},\footnote{con Linux \func{getcwd} è una
  \textit{system call} dalla versione 2.1.9, in precedenza il valore doveva
  essere ottenuto tramite il filesystem \texttt{/proc} da
  \procfile{/proc/self/cwd}.} il cui prototipo è:

\begin{funcproto}{
\fhead{unistd.h}
\fdecl{char *getcwd(char *buffer, size\_t size)}
\fdesc{Legge il \textit{pathname} della directory di lavoro corrente.} 
}
{La funzione ritorna il puntatore a \param{buffer} in caso di successo e
  \val{NULL} per un errore, nel qual caso \var{errno} assumerà uno dei valori: 
  \begin{errlist}
  \item[\errcode{EACCES}] manca il permesso di lettura o di attraversamento  su
    uno dei componenti del \textit{pathname} (cioè su una delle directory
    superiori alla corrente).
  \item[\errcode{EINVAL}] l'argomento \param{size} è zero e \param{buffer} non
    è nullo.
  \item[\errcode{ENOENT}] la directory di lavoro è stata eliminata. 
  \item[\errcode{ERANGE}] l'argomento \param{size} è più piccolo della
    lunghezza del \textit{pathname}. 
  \end{errlist}
  ed inoltre \errcode{EFAULT} ed \errcode{ENOMEM} nel loro significato
  generico.}
\end{funcproto}

La funzione restituisce il \textit{pathname} completo della directory di
lavoro corrente nella stringa puntata da \param{buffer}, che deve essere
precedentemente allocata, per una dimensione massima di \param{size}.  Il
buffer deve essere sufficientemente largo da poter contenere il
\textit{pathname} completo più lo zero di terminazione della stringa. Qualora
esso ecceda le dimensioni specificate con \param{size} la funzione restituisce
un errore.

A partire dal kernel Linux 2.6.36 il nome può avere come prefisso la stringa
\texttt{(unreachable)} se la directory di lavoro resta fuori dalla directory
radice del processo dopo un \func{chroot} (torneremo su questi argomenti in
sez.~\ref{sec:file_chroot}); pertanto è sempre opportuno controllare il primo
carattere della stringa restituita dalla funzione per evitare di interpretare
mare un \textit{pathname} irraggiungibile.

Come estensione allo standard POSIX.1, supportata da Linux e dalla
\acr{glibc}, si può anche specificare un puntatore nullo come \param{buffer}
nel qual caso la stringa sarà allocata automaticamente per una dimensione pari
a \param{size} qualora questa sia diversa da zero, o della lunghezza esatta
del \textit{pathname} altrimenti. In questo caso ci si deve ricordare di
disallocare la stringa con \func{free} una volta cessato il suo utilizzo.

Un uso comune di \func{getcwd} è quello di salvarsi la directory di lavoro
all'avvio del programma per poi potervi tornare in un tempo successivo, un
metodo alternativo più veloce, se non si è a corto di file descriptor, è
invece quello di aprire all'inizio la directory corrente (vale a dire
``\texttt{.}'') e tornarvi in seguito con \func{fchdir}.

Di questa funzione esiste una versione alternativa per compatibilità
all'indietro con BSD, \funcm{getwd}, che non prevede l'argomento \param{size}
e quindi non consente di specificare la dimensione di \param{buffer} che
dovrebbe essere allocato in precedenza ed avere una dimensione sufficiente
(per BSD maggiore \const{PATH\_MAX}, che di solito 256 byte, vedi
sez.~\ref{sec:sys_limits}). Il problema è che su Linux non esiste una
dimensione superiore per la lunghezza di un \textit{pathname}, per cui non è
detto che il buffer sia sufficiente a contenere il nome del file, e questa è
la ragione principale per cui questa funzione è deprecata, e non la tratteremo.

Una seconda funzione usata per ottenere la directory di lavoro è
\funcm{get\_current\_dir\_name} (la funzione è una estensione GNU e presente
solo nella \acr{glibc}) che non prende nessun argomento ed è sostanzialmente
equivalente ad una \code{getcwd(NULL, 0)}, con la differenza che se
disponibile essa ritorna il valore della variabile di ambiente \envvar{PWD},
che essendo costruita dalla shell può contenere un \textit{pathname}
comprendente anche dei collegamenti simbolici. Usando \func{getcwd} infatti,
essendo il \textit{pathname} ricavato risalendo all'indietro l'albero della
directory, si perderebbe traccia di ogni passaggio attraverso eventuali
collegamenti simbolici.

Per cambiare la directory di lavoro si può usare la funzione di sistema
\funcd{chdir}, equivalente del comando di shell \cmd{cd}, il cui nome sta
appunto per \textit{change directory}, il suo prototipo è:

\begin{funcproto}{
\fhead{unistd.h}
\fdecl{int chdir(const char *pathname)}
\fdesc{Cambia la directory di lavoro per \textit{pathname}.} 
}
{La funzione ritorna $0$ in caso di successo e $-1$ per un errore, nel qual
  caso \var{errno} assumerà uno dei valori: 
  \begin{errlist}
  \item[\errcode{EACCES}] manca il permesso di ricerca su uno dei componenti
    di \param{pathname}.
  \item[\errcode{ENAMETOOLONG}] il nome indicato in \param{path} è troppo lungo.
  \item[\errcode{ENOTDIR}] non si è specificata una directory.
  \end{errlist}
  ed inoltre \errval{EFAULT}, \errval{EIO}, \errval{ELOOP}, \errval{ENOENT} e
  \errval{ENOMEM} nel loro significato generico.}
\end{funcproto}

La funzione cambia la directory di lavoro in \param{pathname} ed
ovviamente \param{pathname} deve indicare una directory per la quale si hanno
i permessi di accesso.

Dato che ci si può riferire ad una directory anche tramite un file descriptor,
per cambiare directory di lavoro è disponibile anche la funzione di sistema
\funcd{fchdir}, il cui prototipo è:

\begin{funcproto}{
\fhead{unistd.h}
\fdecl{int fchdir(int fd)}
\fdesc{Cambia la directory di lavoro per file descriptor.} 
}
{La funzione ritorna $0$ in caso di successo e $-1$ per un errore, nel qual
  caso \var{errno} assumerà i valori \errval{EBADF} o \errval{EACCES} nel loro
  significato generico.}
\end{funcproto}

La funzione è identica a \func{chdir}, ma prende come argomento un file
descriptor \param{fd} invece di un \textit{pathname}. Anche in questo
caso \param{fd} deve essere un file descriptor valido che fa riferimento ad
una directory. Inoltre l'unico errore di accesso possibile (tutti gli altri
sarebbero occorsi all'apertura di \param{fd}), è quello in cui il processo non
ha il permesso di attraversamento alla directory specificata da \param{fd}.

\index{directory~di~lavoro|)} 


\subsection{La creazione dei \textsl{file speciali}}
\label{sec:file_mknod}

\index{file!di~dispositivo|(} 
\index{file!speciali|(} 

Finora abbiamo parlato esclusivamente di file, directory e collegamenti
simbolici, ma in sez.~\ref{sec:file_file_types} abbiamo visto che il sistema
prevede anche degli altri tipi di file, che in genere vanno sotto il nome
generico di \textsl{file speciali}, come i file di dispositivo, le
\textit{fifo} ed i socket.

La manipolazione delle caratteristiche di questi file speciali, il cambiamento
di nome o la loro cancellazione può essere effettuata con le stesse funzioni
che operano sugli altri file, ma quando li si devono creare sono necessarie,
come per le directory, delle funzioni apposite. La prima di queste è la
funzione di sistema \funcd{mknod}, il cui prototipo è:

\begin{funcproto}{
\fhead{sys/types.h}
\fhead{sys/stat.h}
\fhead{fcntl.h}
\fhead{unistd.h}
\fdecl{int mknod(const char *pathname, mode\_t mode, dev\_t dev)}
\fdesc{Crea un file speciale sul filesystem.} 
}
{La funzione ritorna $0$ in caso di successo e $-1$ per un errore, nel qual
  caso \var{errno} assumerà uno dei valori: 
  \begin{errlist}
  \item[\errcode{EEXIST}] \param{pathname} esiste già o è un collegamento
    simbolico. 
  \item[\errcode{EINVAL}] il valore di \param{mode} non indica un file, una
    \textit{fifo}, un socket o un dispositivo.
  \item[\errcode{EPERM}] non si hanno privilegi sufficienti a creare
    l'\texttt{inode}, o il filesystem su cui si è cercato di
    creare \param{pathname} non supporta l'operazione.
  \end{errlist}
  ed inoltre \errval{EACCES}, \errval{EFAULT}, \errval{ELOOP},
  \errval{ENAMETOOLONG}, \errval{ENOENT}, \errval{ENOMEM}, \errval{ENOSPC},
  \errval{ENOTDIR} e \errval{EROFS} nel loro significato generico.}
\end{funcproto}

La funzione permette di creare un \textit{inode} di tipo generico sul
filesystem, e viene in genere utilizzata per creare i file di dispositivo, ma
si può usare anche per creare qualunque tipo di file speciale ed anche file
regolari. L'argomento \param{mode} specifica sia il tipo di file che si vuole
creare che i relativi permessi, secondo i valori riportati in
tab.~\ref{tab:file_mode_flags}, che vanno combinati con un OR aritmetico. I
permessi sono comunque modificati nella maniera usuale dal valore di
\textit{umask} (si veda sez.~\ref{sec:file_perm_management}).

Per il tipo di file può essere specificato solo uno fra i seguenti valori:
\const{S\_IFREG} per un file regolare (che sarà creato vuoto),
\const{S\_IFBLK} per un dispositivo a blocchi, \const{S\_IFCHR} per un
dispositivo a caratteri, \const{S\_IFSOCK} per un socket e \const{S\_IFIFO}
per una \textit{fifo};\footnote{con Linux la funzione non può essere usata per
  creare directory o collegamenti simbolici, si dovranno usare le funzioni
  \func{mkdir} e \func{symlink} a questo dedicate.} un valore diverso
comporterà l'errore \errcode{EINVAL}. Inoltre \param{pathname} non deve
esistere, neanche come collegamento simbolico.

Qualora si sia specificato in \param{mode} un file di dispositivo (vale a dire
o \const{S\_IFBLK} o \const{S\_IFCHR}), il valore di \param{dev} dovrà essere
usato per indicare a quale dispositivo si fa riferimento, altrimenti il suo
valore verrà ignorato.  Solo l'amministratore può creare un file di
dispositivo usando questa funzione (il processo deve avere la capacità
\const{CAP\_MKNOD}), ma in Linux\footnote{questo è un comportamento specifico
  di Linux, la funzione non è prevista dallo standard POSIX.1 originale,
  mentre è presente in SVr4 e 4.4BSD, ma esistono differenze nei comportamenti
  e nei codici di errore, tanto che questa è stata introdotta in POSIX.1-2001
  con una nota che la definisce portabile solo quando viene usata per creare
  delle \textit{fifo}, ma comunque deprecata essendo utilizzabile a tale scopo
  la specifica \func{mkfifo}.} l'uso per la creazione di un file ordinario, di
una \textit{fifo} o di un socket è consentito anche agli utenti normali.

Gli \textit{inode} creati con \func{mknod} apparterranno al proprietario e al
gruppo del processo (usando \ids{UID} e \ids{GID} del gruppo effettivo) che li
ha creati a meno non sia presente il bit \acr{sgid} per la directory o sia
stata attivata la semantica BSD per il filesystem (si veda
sez.~\ref{sec:file_ownership_management}) in cui si va a creare
l'\textit{inode}, nel qual caso per il gruppo verrà usato il \ids{GID} del
proprietario della directory.

\itindbeg{major~number}
\itindbeg{minor~number}

Nella creazione di un file di dispositivo occorre poi specificare
correttamente il valore di \param{dev}; questo infatti è di tipo
\type{dev\_t}, che è un tipo primitivo (vedi
tab.~\ref{tab:intro_primitive_types}) riservato per indicare un
\textsl{numero} di dispositivo. Il kernel infatti identifica ciascun
dispositivo con un valore numerico, originariamente questo era un intero a 16
bit diviso in due parti di 8 bit chiamate rispettivamente \textit{major
  number} e \textit{minor number}, che sono poi i due numeri mostrati dal
comando \texttt{ls -l} al posto della dimensione quando lo si esegue su un
file di dispositivo.

Il \textit{major number} identifica una classe di dispositivi (ad esempio la
seriale, o i dischi IDE) e serve in sostanza per indicare al kernel quale è il
modulo che gestisce quella classe di dispositivi. Per identificare uno
specifico dispositivo di quella classe (ad esempio una singola porta seriale,
o uno dei dischi presenti) si usa invece il \textit{minor number}. L'elenco
aggiornato di questi numeri con le relative corrispondenze ai vari dispositivi
può essere trovato nel file \texttt{Documentation/devices.txt} allegato alla
documentazione dei sorgenti del kernel.

Data la crescita nel numero di dispositivi supportati, ben presto il limite
massimo di 256 si è rivelato troppo basso, e nel passaggio dai kernel della
serie 2.4 alla serie 2.6 è stata aumentata a 32 bit la dimensione del tipo
\type{dev\_t}, con delle dimensioni passate a 12 bit per il \textit{major
  number} e 20 bit per il \textit{minor number}. La transizione però ha
comportato il fatto che \type{dev\_t} è diventato un tipo opaco, e la
necessità di specificare il numero tramite delle opportune macro, così da non
avere problemi di compatibilità con eventuali ulteriori estensioni.

Le macro sono definite nel file \headfiled{sys/sysmacros.h},\footnote{se si
  usa la \acr{glibc} dalla versione 2.3.3 queste macro sono degli alias alle
  versioni specifiche di questa libreria, \macrod{gnu\_dev\_major},
  \macrod{gnu\_dev\_minor} e \macrod{gnu\_dev\_makedev} che si possono usare
  direttamente, al costo di una minore portabilità.} che viene automaticamente
incluso quando si include \headfile{sys/types.h}. Si possono pertanto ottenere
i valori del \textit{major number} e \textit{minor number} di un dispositivo
rispettivamente con le macro \macro{major} e \macro{minor}:

{\centering
\vspace{3pt}
\begin{funcbox}{
\fhead{sys/types.h}
\fdecl{int \macrod{major}(dev\_t dev)}
\fdesc{Restituisce il \textit{major number} del dispositivo \param{dev}.}
\fdecl{int \macrod{minor}(dev\_t dev)}
\fdesc{Restituisce il \textit{minor number} del dispositivo \param{dev}.}  
} 
\end{funcbox}
}

\noindent mentre una volta che siano noti \textit{major number} e
\textit{minor number} si potrà costruire il relativo identificativo con la
macro \macro{makedev}:

{\centering
\vspace{3pt}
\begin{funcbox}{
\fhead{sys/types.h}
\fdecl{dev\_t \macrod{makedev}(int major, int minor)}
\fdesc{Dati \textit{major number} e \textit{minor number} restituisce
  l'identificativo di un dispositivo.} 
} 
\end{funcbox}
}


\itindend{major~number}
\itindend{minor~number}
\index{file!di~dispositivo|)}

Dato che la funzione di sistema \func{mknod} presenta diverse varianti nei
vari sistemi unix-like, lo standard POSIX.1-2001 la dichiara portabile solo in
caso di creazione delle \textit{fifo}, ma anche in questo caso alcune
combinazioni degli argomenti restano non specificate, per cui nello stesso
standard è stata introdotta una funzione specifica per creare una
\textit{fifo} deprecando l'uso di \func{mknod} a tale riguardo.  La funzione è
\funcd{mkfifo} ed il suo prototipo è:

\begin{funcproto}{
\fhead{sys/types.h}
\fhead{sys/stat.h}
\fdecl{int mkfifo(const char *pathname, mode\_t mode)}
\fdesc{Crea una \textit{fifo}.} 
}
{La funzione ritorna $0$ in caso di successo e $-1$ per un errore, nel qual
  caso \var{errno} assumerà \errval{EACCES}, \errval{EEXIST},
  \errval{ENAMETOOLONG}, \errval{ENOENT}, \errval{ENOSPC}, \errval{ENOTDIR} e
  \errval{EROFS} nel loro significato generico.}
\end{funcproto}

La funzione crea la \textit{fifo} \param{pathname} con i
permessi \param{mode}. Come per \func{mknod} il file \param{pathname} non deve
esistere (neanche come collegamento simbolico); al solito i permessi
specificati da \param{mode} vengono modificati dal valore di \textit{umask}
(vedi sez.~\ref{sec:file_perm_management}).

\index{file!speciali|)} 


\subsection{I file temporanei}
\label{sec:file_temp_file}

In molte occasioni è utile poter creare dei file temporanei; benché la cosa
sembri semplice, in realtà il problema è più sottile di quanto non appaia a
prima vista. Infatti anche se sembrerebbe banale generare un nome a caso e
creare il file dopo aver controllato che questo non esista, nel momento fra il
controllo e la creazione si ha giusto lo spazio per una possibile \textit{race
  condition} (si ricordi quanto visto in sez.~\ref{sec:proc_race_cond}).

\itindbeg{symlink~attack}

Molti problemi di sicurezza derivano proprio da una creazione non accorta di
file temporanei che lascia aperta questa \textit{race condition}. Un
attaccante allora potrà sfruttarla con quello che viene chiamato
``\textit{symlink attack}'' dove nell'intervallo fra la generazione di un nome
e l'accesso allo stesso, viene creato un collegamento simbolico con quel nome
verso un file diverso, ottenendo, se il programma sotto attacco ne ha la
capacità, un accesso privilegiato.\footnote{dal kernel 3.6 sono state
  introdotte delle contromisure, illustrate in
  sez.~\ref{sec:procadv_security_misc}, che rendono impraticabili questo tipo
  di attacchi, ma questa non è una buona scusa per ignorare il problema.}

\itindend{symlink~attack}

La \acr{glibc} provvede varie funzioni per generare nomi di file temporanei,
di cui si abbia certezza di unicità al momento della generazione; storicamente
la prima di queste funzioni create a questo scopo era
\funcd{tmpnam},\footnote{la funzione è stata deprecata nella revisione
  POSIX.1-2008 dello standard POSIX.} il cui prototipo è:

\begin{funcproto}{
\fhead{stdio.h}
\fdecl{char *tmpnam(char *string)}
\fdesc{Genera un nome univoco per un file temporaneo.} 
}
{La funzione ritorna il puntatore alla stringa con il nome in caso di successo
  e \val{NULL} in caso di fallimento, non sono definiti errori.}
\end{funcproto}

La funzione restituisce il puntatore ad una stringa contente un nome di file
valido e non esistente al momento dell'invocazione. Se si è passato come
argomento \param{string} un puntatore non nullo ad un buffer di caratteri
questo deve essere di dimensione \constd{L\_tmpnam} ed il nome generato vi
verrà copiato automaticamente, altrimenti il nome sarà generato in un buffer
statico interno che verrà sovrascritto ad una chiamata successiva.  Successive
invocazioni della funzione continueranno a restituire nomi unici fino ad un
massimo di \constd{TMP\_MAX} volte, limite oltre il quale il comportamento è
indefinito. Al nome viene automaticamente aggiunto come prefisso la directory
specificata dalla costante \constd{P\_tmpdir}.\footnote{le costanti
  \const{L\_tmpnam}, \const{P\_tmpdir} e \const{TMP\_MAX} sono definite in
  \headfile{stdio.h}.}

Di questa funzione esiste una versione rientrante, \funcm{tmpnam\_r}, che non
fa nulla quando si passa \val{NULL} come argomento.  Una funzione simile,
\funcd{tempnam}, permette di specificare un prefisso per il file
esplicitamente, il suo prototipo è:

\begin{funcproto}{
\fhead{stdio.h}
\fdecl{char *tempnam(const char *dir, const char *pfx)}
\fdesc{Genera un nome univoco per un file temporaneo.} 
}
{La funzione ritorna il puntatore alla stringa con il nome in caso di successo
  e \val{NULL} per un errore, nel qual caso \var{errno} potrà assumere solo il
  valore \errval{ENOMEM} qualora fallisca l'allocazione della stringa.}
\end{funcproto}

La funzione alloca con \code{malloc} la stringa in cui restituisce il nome,
per cui è sempre rientrante, occorre però ricordarsi di disallocare con
\code{free} il puntatore che restituisce.  L'argomento \param{pfx} specifica
un prefisso di massimo 5 caratteri per il nome provvisorio. La funzione
assegna come directory per il file temporaneo, verificando che esista e sia
accessibile, la prima valida fra le seguenti:
\begin{itemize*}
\item la variabile di ambiente \envvar{TMPDIR} (non ha effetto se non è
  definita o se il programma chiamante è \acr{suid} o \acr{sgid}, vedi
  sez.~\ref{sec:file_special_perm}),
\item il valore dell'argomento \param{dir} (se diverso da \val{NULL}),
\item il valore della costante \const{P\_tmpdir},
\item la directory \file{/tmp}.
\end{itemize*}

In ogni caso, anche se con queste funzioni la generazione del nome è casuale,
ed è molto difficile ottenere un nome duplicato, nulla assicura che un altro
processo non possa avere creato, fra l'ottenimento del nome e l'apertura del
file, un altro file o un collegamento simbolico con lo stesso nome. Per questo
motivo quando si usa il nome ottenuto da una di queste funzioni occorre sempre
assicurarsi che non si stia usando un collegamento simbolico e aprire il nuovo
file in modalità di esclusione (cioè con l'opzione \const{O\_EXCL} per i file
descriptor o con il flag ``\texttt{x}'' per gli \textit{stream}) che fa
fallire l'apertura in caso il file sia già esistente. Essendo disponibili
alternative migliori l'uso di queste funzioni è deprecato.

Per evitare di dovere effettuare a mano tutti questi controlli, lo standard
POSIX definisce la funzione \funcd{tmpfile}, che permette di ottenere in
maniera sicura l'accesso ad un file temporaneo, il suo prototipo è:

\begin{funcproto}{
\fhead{stdio.h}
\fdecl{FILE *tmpfile(void)}
\fdesc{Apre un file temporaneo in lettura/scrittura.} 
}
{La funzione ritorna il puntatore allo \textit{stream} associato al file
  temporaneo in caso di successo e \val{NULL} per un errore, nel qual caso
  \var{errno} assumerà uno dei valori:
  \begin{errlist}
    \item[\errcode{EEXIST}] non è stato possibile generare un nome univoco.
    \item[\errcode{EINTR}] la funzione è stata interrotta da un segnale.
  \end{errlist}
  ed inoltre \errval{EFAULT}, \errval{EMFILE}, \errval{ENFILE},
  \errval{ENOSPC}, \errval{EROFS} e \errval{EACCES} nel loro significato
  generico.}
\end{funcproto}


La funzione restituisce direttamente uno \textit{stream} già aperto (in
modalità \code{w+b}, si veda sez.~\ref{sec:file_fopen}) e pronto per l'uso,
che viene automaticamente cancellato alla sua chiusura o all'uscita dal
programma. Lo standard non specifica in quale directory verrà aperto il file,
ma la \acr{glibc} prima tenta con \const{P\_tmpdir} e poi con
\file{/tmp}. Questa funzione è rientrante e non soffre di problemi di
\textit{race condition}.

Alcune versioni meno recenti di Unix non supportano queste funzioni; in questo
caso si possono usare le vecchie funzioni \funcd{mktemp} e \func{mkstemp} che
modificano una stringa di input che serve da modello e che deve essere
conclusa da 6 caratteri ``\texttt{X}'' che verranno sostituiti da un codice
unico. La prima delle due è analoga a \func{tmpnam} e genera soltanto un nome
casuale, il suo prototipo è:

\begin{funcproto}{
\fhead{stlib.h}
\fdecl{char *mktemp(char *template)}
\fdesc{Genera un nome univoco per un file temporaneo.} 
}
{La funzione ritorna  il puntatore a \param{template} in caso di successo e
  \val{NULL} per un errore, nel qual caso \var{errno} assumerà uno dei valori: 
  \begin{errlist}
    \item[\errcode{EINVAL}] \param{template} non termina con \code{XXXXXX}.
  \end{errlist}}
\end{funcproto}

La funzione genera un nome univoco sostituendo le \code{XXXXXX} finali di
\param{template}; dato che \param{template} deve poter essere modificata dalla
funzione non si può usare una stringa costante.  Tutte le avvertenze riguardo
alle possibili \textit{race condition} date per \func{tmpnam} continuano a
valere; inoltre in alcune vecchie implementazioni il valore usato per
sostituire le \code{XXXXXX} viene formato con il \ids{PID} del processo più
una lettera, il che mette a disposizione solo 26 possibilità diverse per il
nome del file, e rende il nome temporaneo facile da indovinare.  Per tutti
questi motivi la funzione è deprecata e non dovrebbe mai essere usata.

La seconda funzione, \funcd{mkstemp} è sostanzialmente equivalente a
\func{tmpfile}, ma restituisce un file descriptor invece di un nome; il suo
prototipo è:

\begin{funcproto}{
\fhead{stlib.h}
\fdecl{int mkstemp(char *template)}
\fdesc{Apre un file temporaneo.} 
}

{La funzione ritorna il file descriptor in caso di successo e $-1$ per un
  errore, nel qual 
  caso \var{errno} assumerà uno dei valori: 
  \begin{errlist}
    \item[\errcode{EEXIST}] non è riuscita a creare un file temporaneo, il
      contenuto di \param{template} è indefinito.
    \item[\errcode{EINVAL}] \param{template} non termina con \code{XXXXXX}.
  \end{errlist}}
\end{funcproto}

Come per \func{mktemp} anche in questo caso \param{template} non può essere
una stringa costante. La funzione apre un file in lettura/scrittura con la
funzione \func{open}, usando l'opzione \const{O\_EXCL} (si veda
sez.~\ref{sec:file_open_close}), in questo modo al ritorno della funzione si
ha la certezza di essere stati i creatori del file, i cui permessi (si veda
sez.~\ref{sec:file_perm_overview}) sono impostati al valore \code{0600}
(lettura e scrittura solo per il proprietario).\footnote{questo è vero a
  partire dalla \acr{glibc} 2.0.7, le versioni precedenti della \acr{glibc} e
  le vecchie \acr{libc5} e \acr{libc4} usavano il valore \code{0666} che
  permetteva a chiunque di leggere e scrivere i contenuti del file.}  Di
questa funzione esiste una variante \funcd{mkostemp}, introdotta
specificamente dalla \acr{glibc},\footnote{la funzione è stata introdotta
  nella versione 2.7 delle librerie e richiede che sia definita la macro
  \macro{\_GNU\_SOURCE}.} il cui prototipo è:

\begin{funcproto}{
\fhead{stlib.h}
\fdecl{int mkostemp(char *template, int flags)}
\fdesc{Apre un file temporaneo.} 
}
{La funzione ritorna un file descriptor in caso di successo e $-1$ per un
  errore, nel qual caso \var{errno} assumerà  gli stessi valori di
  \func{mkstemp}.} 
\end{funcproto}
\noindent la cui sola differenza è la presenza dell'ulteriore argomento
\var{flags} che consente di specificare alcuni ulteriori flag (come
\const{O\_APPEND}, \const{O\_CLOEXEC}, \const{O\_SYNC}, il cui significato
vedremo in sez.~\ref{sec:file_open_close}) da passare ad \func{open}
nell'apertura del file.\footnote{si tenga presente che \func{mkostemp}
  utilizza già \const{O\_CREAT}, \const{O\_EXCL} e \const{O\_RDWR}, che non è
  il caso di riindicare, dato che ciò potrebbe portare ad errori in altri
  sistemi operativi.}

Di queste due funzioni sono state poi introdotte, a partire dalla \acr{glibc}
2.11 due varianti, \funcd{mkstemps} e \funcd{mkostemps}, che consentono di
indicare anche un suffisso, i loro prototipi sono:

\begin{funcproto}{
\fhead{stlib.h}
\fdecl{int mkstemps(char *template, int suffixlen)}
\fdesc{Apre un file temporaneo.} 
\fdecl{int mkostemps(char *template, int suffixlen, int flags)}
\fdesc{Apre un file temporaneo.} 
}

{Le funzioni hanno gli stessi valori di ritorno e gli stessi errori di
  \func{mkstemp} con lo stesso significato, tranne \errval{EINVAL} che viene
  restituito se \param{template} non è di lunghezza pari ad almeno
  $6+$\param{suffixlen} ed i 6 caratteri prima del suffisso non sono
  \code{XXXXXX}.}
\end{funcproto}

Le due funzioni, un'estensione non standard fornita dalla \acr{glibc}, sono
identiche a \funcd{mkstemp} e \funcd{mkostemp}, ma consentono di avere un nome
del file nella forma \texttt{prefissoXXXXXXsuffisso} dove la lunghezza del
suffisso deve essere indicata con \param{suffixlen}.

Infine con OpenBSD è stata introdotta un'altra funzione simile alle
precedenti, \funcd{mkdtemp}, che crea invece una directory
temporanea;\footnote{la funzione è stata introdotta nella \acr{glibc} a
  partire dalla versione 2.1.91 ed inserita nello standard POSIX.1-2008.}  il
suo prototipo è:

\begin{funcproto}{
\fhead{stlib.h}
\fdecl{char *mkdtemp(char *template)}
\fdesc{Crea una directory temporanea.} 
}
{La funzione ritorna il puntatore al nome della directory in caso di successo
  e \val{NULL} per un errore, nel qual caso \var{errno} assumerà uno dei
  valori:
  \begin{errlist}
    \item[\errcode{EINVAL}] \param{template} non termina con \code{XXXXXX}.
  \end{errlist}
  più gli altri eventuali codici di errore di \func{mkdir}.}
\end{funcproto}

La funzione crea una directory temporanea il cui nome è ottenuto sostituendo
le \code{XXXXXX} finali di \param{template} con permessi \code{0700} (si veda
sez.~\ref{sec:file_perm_overview} per i dettagli). Dato che la creazione della
directory è sempre atomica i precedenti problemi di \textit{race condition}
non si pongono.


\section{La manipolazione delle caratteristiche dei file}
\label{sec:file_infos}

Come spiegato in sez.~\ref{sec:file_filesystem} tutte le informazioni generali
relative alle caratteristiche di ciascun file, a partire dalle informazioni
relative al controllo di accesso, sono mantenute nell'\textit{inode}. Vedremo
in questa sezione come sia possibile leggere tutte queste informazioni usando
la funzione \func{stat}, che permette l'accesso a tutti i dati memorizzati
nell'\textit{inode}; esamineremo poi le varie funzioni usate per manipolare
tutte queste informazioni, eccetto quelle che riguardano la gestione del
controllo di accesso, trattate in sez.~\ref{sec:file_access_control}.


\subsection{La lettura delle caratteristiche dei file}
\label{sec:file_stat}

La lettura delle informazioni relative ai file è fatta attraverso la famiglia
delle funzioni \func{stat} che sono quelle che usa il comando \cmd{ls} per
poter ottenere e mostrare tutti i dati relativi ad un file; ne fanno parte le
funzioni di sistema \funcd{stat}, \funcd{fstat} e \funcd{lstat}, i cui
prototipi sono:

\begin{funcproto}{
\fhead{sys/types.h}
\fhead{sys/stat.h}
\fhead{unistd.h}
\fdecl{int stat(const char *file\_name, struct stat *buf)}
\fdecl{int lstat(const char *file\_name, struct stat *buf)}
\fdecl{int fstat(int filedes, struct stat *buf)}
\fdesc{Leggono le informazioni di un file.} 
}
{Le funzioni ritornano $0$ in caso di successo e $-1$ per un errore, nel qual
  caso \var{errno} assumerà uno dei valori:
  \begin{errlist}
    \item[\errcode{EOVERFLOW}] il file ha una dimensione che non può essere
      rappresentata nel tipo \type{off\_t} (può avvenire solo in caso di
      programmi compilati su piattaforme a 32 bit senza le estensioni
      (\texttt{-D \_FILE\_OFFSET\_BITS=64}) per file a 64 bit).
  \end{errlist}
  ed inoltre \errval{EFAULT} ed \errval{ENOMEM}, per \func{stat} e
  \func{lstat} anche \errval{EACCES}, \errval{ELOOP}, \errval{ENAMETOOLONG},
  \errval{ENOENT}, \errval{ENOTDIR}, per \func{fstat} anche \errval{EBADF}, 
  nel loro significato generico.}
\end{funcproto}

La funzione \func{stat} legge le informazioni del file indicato
da \param{file\_name} e le inserisce nel buffer puntato
dall'argomento \param{buf}; la funzione \func{lstat} è identica a \func{stat}
eccetto che se \param{file\_name} è un collegamento simbolico vengono lette le
informazioni relative ad esso e non al file a cui fa riferimento. Infine
\func{fstat} esegue la stessa operazione su un file già aperto, specificato
tramite il suo file descriptor \param{filedes}.

La struttura \struct{stat} usata da queste funzioni è definita nell'header
\headfiled{sys/stat.h} e in generale dipende dall'implementazione; la versione
usata da Linux è mostrata in fig.~\ref{fig:file_stat_struct}, così come
riportata dalla pagina di manuale di \func{stat}. In realtà la definizione
effettivamente usata nel kernel dipende dall'architettura e ha altri campi
riservati per estensioni come tempi dei file più precisi (vedi
sez.~\ref{sec:file_file_times}).

\begin{figure}[!htb]
  \footnotesize
  \centering
  \begin{minipage}[c]{0.8\textwidth}
    \includestruct{listati/stat.h}
  \end{minipage} 
  \normalsize 
  \caption{La struttura \structd{stat} per la lettura delle informazioni dei 
    file.}
  \label{fig:file_stat_struct}
\end{figure}

Si noti come i vari membri della struttura siano specificati come tipi
primitivi del sistema, di quelli definiti in
tab.~\ref{tab:intro_primitive_types}, e dichiarati in \headfile{sys/types.h},
con l'eccezione di \typed{blksize\_t} e \typed{blkcnt\_t} che sono nuovi tipi
introdotti per rendersi indipendenti dalla piattaforma. 

Benché la descrizione dei commenti di fig.~\ref{fig:file_stat_struct} sia
abbastanza chiara, vale la pena illustrare maggiormente il significato dei
campi di \struct{stat} su cui non torneremo in maggior dettaglio nel resto di
questa sezione:
\begin{itemize*}
\item Il campo \var{st\_nlink} contiene il numero di \textit{hard link} che
  fanno riferimento al file (il cosiddetto \textit{link count}) di cui abbiamo
  già parlato in numerose occasioni.
\item Il campo \var{st\_ino} contiene il numero di \textit{inode} del file,
  quello viene usato all'interno del filesystem per identificarlo e che può
  essere usato da un programma per determinare se due \textit{pathname} fanno
  riferimento allo stesso file.
\item Il campo \var{st\_dev} contiene il numero del dispositivo su cui risiede
  il file (o meglio il suo filesystem). Si tratta dello stesso numero che si
  usa con \func{mknod} e che può essere decomposto in \textit{major number} e
  \textit{minor number} con le macro \macro{major} e \macro{minor} viste in
  sez.~\ref{sec:file_mknod}.
\item Il campo \var{st\_rdev} contiene il numero di dispositivo associato al
  file stesso ed ovviamente ha un valore significativo soltanto quando il file
  è un dispositivo a caratteri o a blocchi.
\item Il campo \var{st\_blksize} contiene la dimensione dei blocchi di dati
  usati nell'I/O su disco, che è anche la dimensione usata per la
  bufferizzazione dei dati dalle librerie del C per l'interfaccia degli
  \textit{stream}.  Leggere o scrivere blocchi di dati in dimensioni inferiori
  a questo valore è inefficiente in quanto le operazioni su disco usano
  comunque trasferimenti di questa dimensione.
\end{itemize*}

Nell'evoluzione del kernel la \textit{system call} che fornisce \func{stat} è
stata modificata più volte per tener conto dei cambiamenti fatti alla
struttura \struct{stat},\footnote{questo ha significato l'utilizzo a basso
  livello di diverse \textit{system call} e diverse versioni della struttura.}
in particolare a riguardo ai tempi dei file, di cui è stata aumentata la
precisione (torneremo su questo in sez.~\ref{sec:file_file_times}) ma anche
per gli aggiornamenti fatti ai campi \var{st\_ino}, \var{st\_uid} e
\var{st\_gid}.

Sulle piattaforme a 32 bit questi cambiamenti, che han visto un aumento delle
dimensioni dei campi della struttura per adattarli alle nuove esigenze, sono
mascherati dalla \acr{glibc} che attraverso \func{stat} invoca la versione più
recente della \textit{system call} e rimpacchetta i dati se questo è
necessario per eseguire dei vecchi programmi. Nelle piattaforme a 64 bit
invece è presente un'unica versione della \textit{system call} e la struttura
\struct{stat} ha campi di dimensione sufficiente.

Infine a partire dal kernel 2.6.16 è stata introdotta una ulteriore funzione
della famiglia, \func{fstatat} che consente di trattare con sicurezza i
\textit{pathname} relativi, la tratteremo in sez.~\ref{sec:file_openat},
insieme alla nuova \textit{system call} \func{statx}, introdotta dal kernel
4.11 per estendere l'interfaccia di \func{stat} e le informazioni che essa può
restituire.


\subsection{I tipi di file}
\label{sec:file_types}

Abbiamo sottolineato fin dall'introduzione che Linux, come ogni sistema
unix-like, supporta oltre ai file ordinari e alle directory una serie di altri
``\textsl{tipi}'' di file che possono stare su un filesystem (elencati in
tab.~\ref{tab:file_file_types}).  Il tipo di file viene ritornato dalle
funzioni della famiglia \func{stat} all'interno del campo \var{st\_mode} di
una struttura \struct{stat}. 

\begin{table}[htb]
  \centering
  \footnotesize
  \begin{tabular}[c]{|l|l|}
    \hline
    \textbf{Macro} & \textbf{Tipo del file} \\
    \hline
    \hline
    \macrod{S\_ISREG}\texttt{(m)}  & File normale.\\
    \macrod{S\_ISDIR}\texttt{(m)}  & Directory.\\
    \macrod{S\_ISCHR}\texttt{(m)}  & Dispositivo a caratteri.\\
    \macrod{S\_ISBLK}\texttt{(m)}  & Dispositivo a blocchi.\\
    \macrod{S\_ISFIFO}\texttt{(m)} & \textit{Fifo}.\\
    \macrod{S\_ISLNK}\texttt{(m)}  & Collegamento simbolico.\\
    \macrod{S\_ISSOCK}\texttt{(m)} & Socket.\\
    \hline    
  \end{tabular}
  \caption{Macro per i tipi di file (definite in \headfile{sys/stat.h}).}
  \label{tab:file_type_macro}
\end{table}

Il campo \var{st\_mode} è una maschera binaria in cui l'informazione viene
suddivisa nei vari bit che compongono, ed oltre a quelle sul tipo di file,
contiene anche le informazioni relative ai permessi su cui torneremo in
sez.~\ref{sec:file_perm_overview}. Dato che i valori numerici usati per
definire il tipo di file possono variare a seconda delle implementazioni, lo
standard POSIX definisce un insieme di macro che consentono di verificare il
tipo di file in maniera standardizzata.

Queste macro vengono usate anche da Linux che supporta pure le estensioni allo
standard per i collegamenti simbolici e i socket definite da BSD.\footnote{le
  ultime due macro di tab.~\ref{tab:file_type_macro}, che non sono presenti
  nello standard POSIX fino alla revisione POSIX.1-1996.}  L'elenco completo
delle macro con cui è possibile estrarre da \var{st\_mode} l'informazione
relativa al tipo di file è riportato in tab.~\ref{tab:file_type_macro}, tutte
le macro restituiscono un valore intero da usare come valore logico e prendono
come argomento il valore di \var{st\_mode}.

\begin{table}[htb]
  \centering
  \footnotesize
  \begin{tabular}[c]{|l|c|l|}
    \hline
    \textbf{Flag} & \textbf{Valore} & \textbf{Significato} \\
    \hline
    \hline
    \constd{S\_IFMT}   &  0170000 & Maschera per i bit del tipo di file.\\
    \constd{S\_IFSOCK} &  0140000 & Socket.\\
    \constd{S\_IFLNK}  &  0120000 & Collegamento simbolico.\\
    \constd{S\_IFREG}  &  0100000 & File regolare.\\ 
    \constd{S\_IFBLK}  &  0060000 & Dispositivo a blocchi.\\
    \constd{S\_IFDIR}  &  0040000 & Directory.\\
    \constd{S\_IFCHR}  &  0020000 & Dispositivo a caratteri.\\
    \constd{S\_IFIFO}  &  0010000 & \textit{Fifo}.\\
    \hline
    \constd{S\_ISUID}  &  0004000 & Set user ID (\acr{suid}) bit, vedi
                                   sez.~\ref{sec:file_special_perm}).\\
    \constd{S\_ISGID}  &  0002000 & Set group ID (\acr{sgid}) bit, vedi
                                   sez.~\ref{sec:file_special_perm}).\\
    \constd{S\_ISVTX}  &  0001000 & \acr{Sticky} bit, vedi
                                   sez.~\ref{sec:file_special_perm}).\\
    \hline
    \constd{S\_IRWXU}  &  00700   & Maschera per i permessi del proprietario.\\
    \constd{S\_IRUSR}  &  00400   & Il proprietario ha permesso di lettura.\\
    \constd{S\_IWUSR}  &  00200   & Il proprietario ha permesso di scrittura.\\
    \constd{S\_IXUSR}  &  00100   & Il proprietario ha permesso di esecuzione.\\
    \hline
    \constd{S\_IRWXG}  &  00070   & Maschera per i permessi del gruppo.\\
    \constd{S\_IRGRP}  &  00040   & Il gruppo ha permesso di lettura.\\
    \constd{S\_IWGRP}  &  00020   & Il gruppo ha permesso di scrittura.\\
    \constd{S\_IXGRP}  &  00010   & Il gruppo ha permesso di esecuzione.\\
    \hline
    \constd{S\_IRWXO}  &  00007   & Maschera per i permessi di tutti gli altri\\
    \constd{S\_IROTH}  &  00004   & Gli altri hanno permesso di lettura.\\
    \constd{S\_IWOTH}  &  00002   & Gli altri hanno permesso di esecuzione.\\
    \constd{S\_IXOTH}  &  00001   & Gli altri hanno permesso di esecuzione.\\
    \hline    
  \end{tabular}
  \caption{Costanti per l'identificazione dei vari bit che compongono il campo
    \var{st\_mode} (definite in \headfile{sys/stat.h}).}
  \label{tab:file_mode_flags}
\end{table}

Oltre alle macro di tab.~\ref{tab:file_type_macro}, che semplificano
l'operazione, è possibile usare direttamente il valore di \var{st\_mode} per
ricavare il tipo di file controllando direttamente i vari bit in esso
memorizzati. Per questo sempre in \headfile{sys/stat.h} sono definite le varie
costanti numeriche riportate in tab.~\ref{tab:file_mode_flags}, che
definiscono le maschere che consentono di selezionare non solo i dati relativi
al tipo di file, ma anche le informazioni relative ai permessi su cui
torneremo in sez.~\ref{sec:file_access_control}, ed identificare i rispettivi
valori.

Le costanti che servono per la identificazione del tipo di file sono riportate
nella prima sezione di tab.~\ref{tab:file_mode_flags}, mentre le sezioni
successive attengono alle costanti usate per i permessi.  Il primo valore
dell'elenco è la maschera binaria \const{S\_IFMT} che permette di estrarre da
\var{st\_mode} (con un AND aritmetico) il blocco di bit nei quali viene
memorizzato il tipo di file. I valori successivi sono le costanti
corrispondenti ai vari tipi di file, e possono essere usate per verificare la
presenza del tipo di file voluto ed anche, con opportune combinazioni,
alternative fra più tipi di file. 

Si tenga presente però che a differenza dei permessi, l'informazione relativa
al tipo di file non è una maschera binaria, per questo motivo se si volesse
impostare una condizione che permetta di controllare se un file è una
directory o un file ordinario non si possono controllare dei singoli bit, ma
si dovrebbe usare una macro di preprocessore come:
\includecodesnip{listati/is_regdir.h}
in cui si estraggono ogni volta da \var{st\_mode} i bit relativi al tipo di
file e poi si effettua il confronto con i due possibili tipi di file.


\subsection{Le dimensioni dei file}
\label{sec:file_file_size}

Abbiamo visto in fig.~\ref{fig:file_stat_struct} che campo \var{st\_size} di
una struttura \struct{stat} contiene la dimensione del file in byte. In realtà
questo è vero solo se si tratta di un file regolare contenente dei dati; nel
caso di un collegamento simbolico invece la dimensione è quella del
\textit{pathname} che il collegamento stesso contiene, e per una directory
quella dello spazio occupato per le voci della stessa (che dipende da come
queste vengono mantenute dal filesystem), infine per le \textit{fifo}, i socket
ed i file di dispositivo questo campo è sempre nullo.

Il campo \var{st\_blocks} invece definisce la lunghezza del file espressa in
numero di blocchi di 512 byte. La differenza con \var{st\_size} è che in
questo caso si fa riferimento alla quantità di spazio disco allocata per il
file, e non alla dimensione dello stesso che si otterrebbe leggendolo
sequenzialmente.

Si deve tener presente infatti che la lunghezza del file riportata in
\var{st\_size} non è detto che corrisponda all'occupazione dello spazio su
disco, e non solo perché la parte finale del file potrebbe riempire
parzialmente un blocco. In un sistema unix-like infatti è possibile
l'esistenza dei cosiddetti \textit{sparse file}, cioè file in cui sono
presenti dei ``\textsl{buchi}'' (\textit{holes} nella nomenclatura inglese)
che si formano tutte le volte che si va a scrivere su un file dopo aver
eseguito uno spostamento oltre la sua fine (tratteremo in dettaglio
l'argomento in sez.~\ref{sec:file_lseek}).

In questo caso si avranno risultati differenti a seconda del modo in cui si
calcola la lunghezza del file, ad esempio il comando \cmd{du}, (che riporta il
numero di blocchi occupati) potrà dare una dimensione inferiore, mentre se si
legge il contenuto del file (ad esempio usando il comando \cmd{wc -c}), dato
che in tal caso per i ``\textsl{buchi}'' vengono restituiti degli zeri, si
avrà lo stesso risultato di \cmd{ls}.

Se è sempre possibile allargare un file, scrivendoci sopra o usando la
funzione \func{lseek} (vedi sez.~\ref{sec:file_lseek}) per spostarsi oltre la
sua fine, esistono anche casi in cui si può avere bisogno di effettuare un
troncamento, scartando i dati presenti al di là della dimensione scelta come
nuova fine del file.

Un file può sempre essere troncato a zero aprendolo con il flag
\const{O\_TRUNC} (vedi sez.~\ref{sec:file_open_close}), ma questo è un caso
particolare; per qualunque altra dimensione si possono usare le due funzioni
di sistema \funcd{truncate} e \funcd{ftruncate}, i cui prototipi sono:

\begin{funcproto}{
\fhead{unistd.h}
\fdecl{int ftruncate(int fd, off\_t length))}
\fdecl{int truncate(const char *file\_name, off\_t length)}
\fdesc{Troncano un file.} 
}
{Le funzioni ritornano $0$ in caso di successo e $-1$ per un errore, nel qual
  caso \var{errno} assumerà uno dei valori: 
  \begin{errlist}
  \item[\errcode{EINTR}] si è stati interrotti da un segnale.
  \item[\errcode{EINVAL}] \param{length} è negativa o maggiore delle
    dimensioni massime di un file.
  \item[\errcode{EPERM}] il filesystem non supporta il troncamento di un file.
  \item[\errcode{ETXTBSY}] il file è un programma in esecuzione.
  \end{errlist} 
  per entrambe, mentre per \func{ftruncate} si avranno anche: 
  \begin{errlist}
  \item[\errcode{EBADF}] \param{fd} non è un file descriptor.
  \item[\errcode{EINVAL}] \param{fd} non è un riferimento a un file o non è
    aperto in scrittura. 
  \end{errlist}
  e per \func{truncate} si avranno anche: 
  \begin{errlist}
  \item[\errcode{EACCES}] non si ha il permesso di scrittura sul file o il
    permesso di attraversamento di una delle componenti del \textit{pathname}.
  \item[\errcode{EISDIR}] \param{file\_name} fa riferimento ad una directory.
  \end{errlist}
  ed inoltre \errval{EFAULT}, \errval{EIO}, \errval{ELOOP},
  \errval{ENAMETOOLONG}, \errval{ENOENT}, \errval{ENOTDIR} e \errval{EROFS}
  nel loro significato generico.}
\end{funcproto}

Entrambe le funzioni fan sì che la dimensione del file sia troncata ad un
valore massimo specificato da \param{length}, e si distinguono solo per il
fatto che il file viene indicato con un \textit{pathname} per \func{truncate}
e con un file descriptor per \funcd{ftruncate}. Si tenga presente che se il
file è più lungo della lunghezza specificata i dati in eccesso saranno
perduti.

Il comportamento in caso di lunghezza del file inferiore a \param{length} non
è specificato e dipende dall'implementazione: il file può essere lasciato
invariato o esteso fino alla lunghezza scelta. Nel caso di Linux viene esteso
con la creazione di un \textsl{buco} nel file e ad una lettura si otterranno
degli zeri, si tenga presente però che questo comportamento è supportato solo
per filesystem nativi, ad esempio su un filesystem non nativo come il VFAT di
Windows questo non è possibile.


\subsection{I tempi dei file}
\label{sec:file_file_times}

Il sistema mantiene per ciascun file tre tempi, che sono registrati
nell'\textit{inode} insieme agli altri attributi del file. Questi possono
essere letti tramite la funzione \func{stat}, che li restituisce attraverso
tre campi della struttura \struct{stat} di fig.~\ref{fig:file_stat_struct}. Il
significato di questi tempi e dei relativi campi della struttura \struct{stat}
è illustrato nello schema di tab.~\ref{tab:file_file_times}, dove è anche
riportato un esempio delle funzioni che effettuano cambiamenti su di essi. Il
valore del tempo è espresso nel cosiddetto \textit{calendar time}, su cui
torneremo in dettaglio in sez.~\ref{sec:sys_time}.

\begin{table}[htb]
  \centering
  \footnotesize
  \begin{tabular}[c]{|c|l|l|c|}
    \hline
    \textbf{Membro} & \textbf{Significato} & \textbf{Funzione} 
    & \textbf{Opzione di \cmd{ls}} \\
    \hline
    \hline
    \var{st\_atime}& ultimo accesso ai dati del file    &
                     \func{read}, \func{utime}          & \cmd{-u}\\
    \var{st\_mtime}& ultima modifica ai dati del file   &
                     \func{write}, \func{utime}         & default\\
    \var{st\_ctime}& ultima modifica ai dati dell'\textit{inode} &
                     \func{chmod}, \func{utime}         & \cmd{-c}\\
    \hline
  \end{tabular}
  \caption{I tre tempi associati a ciascun file.}
  \label{tab:file_file_times}
\end{table}

Il primo punto da tenere presente è la differenza fra il cosiddetto tempo di
ultima modifica (il \textit{modification time}) riportato in \var{st\_mtime},
ed il tempo di ultimo cambiamento di stato (il \textit{change status time})
riportato in \var{st\_ctime}. Il primo infatti fa riferimento ad una modifica
del contenuto di un file, mentre il secondo ad una modifica dei metadati
dell'\textit{inode}. Dato che esistono molte operazioni, come la funzione
\func{link} e altre che vedremo in seguito, che modificano solo le
informazioni contenute nell'\textit{inode} senza toccare il contenuto del
file, diventa necessario l'utilizzo di questo secondo tempo.

Il tempo di ultima modifica viene usato ad esempio da programmi come
\cmd{make} per decidere quali file necessitano di essere ricompilati perché
più recenti dei loro sorgenti oppure dai programmi di backup, talvolta insieme
anche al tempo di cambiamento di stato, per decidere quali file devono essere
aggiornati nell'archiviazione.  Il tempo di ultimo accesso viene di solito
usato per identificare i file che non vengono più utilizzati per un certo
lasso di tempo. Ad esempio un programma come \texttt{leafnode} lo usa per
cancellare gli articoli letti più vecchi, mentre \texttt{mutt} lo usa per
marcare i messaggi di posta che risultano letti.  

Il sistema non tiene mai conto dell'ultimo accesso all'\textit{inode},
pertanto funzioni come \func{access} o \func{stat} non hanno alcuna influenza
sui tre tempi. Il comando \cmd{ls} (quando usato con le opzioni \cmd{-l} o
\cmd{-t}) mostra i tempi dei file secondo lo schema riportato nell'ultima
colonna di tab.~\ref{tab:file_file_times}. Si noti anche come in
tab.~\ref{tab:file_file_times} non venga riportato il \textsl{tempo di
  creazione} di un file; in un sistema unix-like infatti questo non esiste, e
non è previsto dall'interfaccia classica, ma essendo usato da altri sistemi
operativi (in particolare Windows) in tutti i filesystem più recenti ne viene
supportata la registrazione, ed a partire dal kernel 4.11 è diventato
possibile anche ottenerne la lettura con la nuova \textit{system call}
\func{statx} (che tratteremo in sez.~\ref{sec:file_openat}).

L'aggiornamento del tempo di ultimo accesso è stato a lungo considerato un
difetto progettuale di Unix, questo infatti comporta la necessità di
effettuare un accesso in scrittura sul disco anche in tutti i casi in cui
questa informazione non interessa e sarebbe possibile avere un semplice
accesso in lettura sui dati bufferizzati. Questo comporta un ovvio costo sia
in termini di prestazioni, che di consumo di risorse come la batteria per i
portatili, o i cicli di riscrittura per i dischi su memorie riscrivibili.

Per questo motivo abbiamo visto in sez.~\ref{sec:filesystem_mounting} come nel
corso dello sviluppo del kernel siano stati introdotti degli opportuni
\textit{mount flag} che consentono di evitare di aggiornare continuamente una
informazione che nella maggior parte dei casi non ha un interesse
rilevante. Per questo motivo i valori dell'\textit{access time} possono
dipendere dalle opzioni di montaggio, ed anche, essendo stato cambiato il
comportamento di default a partire dalla versione 2.6.30, dal kernel che si
sta usando.

In generale quello che avviene con i kernel più recenti è che il tempo di
ultimo accesso viene aggiornato solo se è precedente al tempo di ultima
modifica o cambiamento, o se è cambiato ed passato più di un giorno
dall'ultimo aggiornamento. Così si può rendere evidente che vi è stato un
accesso dopo una modifica, e che il file viene comunque osservato a cadenza
regolare, conservando le informazioni veramente utili senza consumare
inutilmente risorse in continue scritture per mantenere costantemente
aggiornata una informazione che a questo punto non ha più nessuna rilevanza
pratica.\footnote{qualora ce ne fosse la necessità è comunque possibile,
  tramite l'opzione di montaggio \texttt{strictatime}, richiedere in ogni caso
  il comportamento tradizionale.}

\begin{table}[htb]
  \centering
  \footnotesize
  \begin{tabular}[c]{|l|c|c|c|c|c|c|l|}
    \hline
    \multicolumn{1}{|p{2.3cm}|}{\centering{\vspace{6pt}\textbf{Funzione}}} &
    \multicolumn{3}{|p{3.2cm}|}{\centering{
        \textbf{File o directory del riferimento}}}&
    \multicolumn{3}{|p{3.2cm}|}{\centering{
        \textbf{Directory contenente il riferimento}}} 
    &\multicolumn{1}{|p{3.cm}|}{\centering{\vspace{6pt}\textbf{Note}}} \\
    \cline{2-7}
    \cline{2-7}
    \multicolumn{1}{|p{2.3cm}|}{} 
    &\multicolumn{1}{|p{.8cm}|}{\centering{\textsl{(a)}}}
    &\multicolumn{1}{|p{.8cm}|}{\centering{\textsl{(m)}}}
    &\multicolumn{1}{|p{.8cm}|}{\centering{\textsl{(c)}}}
    &\multicolumn{1}{|p{.8cm}|}{\centering{\textsl{(a)}}}
    &\multicolumn{1}{|p{.8cm}|}{\centering{\textsl{(m)}}}
    &\multicolumn{1}{|p{.8cm}|}{\centering{\textsl{(c)}}}
    &\multicolumn{1}{|p{3cm}|}{} \\
    \hline
    \hline
    \func{chmod}, \func{fchmod} 
             & --      & --      &$\bullet$& --      & --      & --      &\\
    \func{chown}, \func{fchown} 
             & --      & --      &$\bullet$& --      & --      & --      &\\
    \func{creat}  
             &$\bullet$&$\bullet$&$\bullet$& --      &$\bullet$&$\bullet$&  
             con \const{O\_CREATE} \\
    \func{creat}  
             & --      &$\bullet$&$\bullet$& --      &$\bullet$&$\bullet$&   
             con \const{O\_TRUNC} \\
    \func{exec}  
             &$\bullet$& --      & --      & --      & --      & --      &\\
    \func{lchown}  
             & --      & --      &$\bullet$& --      & --      & --      &\\
    \func{link}
             & --      & --      &$\bullet$& --      &$\bullet$&$\bullet$&\\
    \func{mkdir}
             &$\bullet$&$\bullet$&$\bullet$& --      &$\bullet$&$\bullet$&\\
    \func{mknod}
             &$\bullet$&$\bullet$&$\bullet$& --      &$\bullet$&$\bullet$&\\
    \func{mkfifo}
             &$\bullet$&$\bullet$&$\bullet$& --      &$\bullet$&$\bullet$&\\
    \func{open}
             &$\bullet$&$\bullet$&$\bullet$& --      &$\bullet$&$\bullet$& 
             con \const{O\_CREATE} \\
    \func{open}
             & --      &$\bullet$&$\bullet$& --      & --      & --      & 
             con \const{O\_TRUNC}  \\
    \func{pipe}
             &$\bullet$&$\bullet$&$\bullet$& --      & --      & --      &\\
    \func{read}
             &$\bullet$& --      & --      & --      & --      & --      &\\
    \func{remove}
             & --      & --      &$\bullet$& --      &$\bullet$&$\bullet$& 
             se esegue \func{unlink}\\
    \func{remove}
              & --      & --      & --      & --      &$\bullet$&$\bullet$& 
              se esegue \func{rmdir}\\
    \func{rename}
              & --      & --      &$\bullet$& --      &$\bullet$&$\bullet$& 
              per ambo gli argomenti\\
    \func{rmdir}
              & --      & --      & --      & --      &$\bullet$&$\bullet$&\\ 
    \func{truncate}
              & --      &$\bullet$&$\bullet$& --      & --      & --      &\\ 
    \func{ftruncate}
              & --      &$\bullet$&$\bullet$& --      & --      & --      &\\ 
    \func{unlink}
              & --      & --      &$\bullet$& --      &$\bullet$&$\bullet$&\\ 
    \func{utime}
              &$\bullet$&$\bullet$&$\bullet$& --      & --      & --      &\\ 
    \func{utimes}
              &$\bullet$&$\bullet$&$\bullet$& --      & --      & --      &\\ 
    \func{write}
              & --      &$\bullet$&$\bullet$& --      & --      & --      &\\ 
    \hline
  \end{tabular}
  \caption{Prospetto dei cambiamenti effettuati sui tempi di ultimo 
    accesso \textsl{(a)}, ultima modifica \textsl{(m)} e ultimo cambiamento di
    stato \textsl{(c)} dalle varie funzioni operanti su file e directory.}
  \label{tab:file_times_effects}  
\end{table}


L'effetto delle varie funzioni di manipolazione dei file sui relativi tempi è
illustrato in tab.~\ref{tab:file_times_effects}, facendo riferimento al
comportamento classico per quanto riguarda \var{st\_atime}. Si sono riportati
gli effetti sia per il file a cui si fa riferimento, sia per la directory che
lo contiene. Questi ultimi possono essere capiti immediatamente se si tiene
conto di quanto già detto e ripetuto a partire da
sez.~\ref{sec:file_filesystem}, e cioè che anche le directory sono anch'esse
file che contengono una lista di nomi, che il sistema tratta in maniera del
tutto analoga a tutti gli altri.

Per questo motivo tutte le volte che compiremo un'operazione su un file che
comporta una modifica del nome contenuto nella directory, andremo anche a
scrivere sulla directory che lo contiene cambiandone il tempo di ultima
modifica. Un esempio di questo tipo di operazione può essere la cancellazione
di un file, invece leggere o scrivere o cambiare i permessi di un file ha
effetti solo sui tempi di quest'ultimo.

Si ricordi infine come \var{st\_ctime} non è il tempo di creazione del file,
che in Unix non esiste, anche se può corrispondervi per file che non sono mai
stati modificati. Per questo motivo, a differenza di quanto avviene con altri
sistemi operativi, quando si copia un file, a meno di preservare
esplicitamente i tempi (ad esempio con l'opzione \cmd{-p} di \cmd{cp}) esso
avrà sempre il tempo corrente in cui si è effettuata la copia come data di
ultima modifica.

I tempi di ultimo accesso ed ultima modifica possono essere modificati
esplicitamente usando la funzione di sistema \funcd{utime}, il cui prototipo
è:

\begin{funcproto}{
\fhead{utime.h}
\fdecl{int utime(const char *filename, struct utimbuf *times)}
\fdesc{Modifica i tempi di ultimo accesso ed ultima modifica di un file.} 
}

{La funzione ritorna $0$ in caso di successo e $-1$ per un errore, nel qual
  caso \var{errno} assumerà uno dei valori: 
  \begin{errlist}
  \item[\errcode{EACCES}] non c'è il permesso di attraversamento per uno dei
    componenti di \param{filename} o \param{times} è \val{NULL} e non si ha il
    permesso di scrittura sul file, o non si è proprietari del file o non si
    hanno i privilegi di amministratore.
  \item[\errcode{EPERM}] \param{times} non è \val{NULL}, e non si è
    proprietari del file o non si hanno i privilegi di amministratore.
  \end{errlist}
  ed inoltre \errval{ENOENT} e \errval{EROFS} nel loro significato generico.}
\end{funcproto}

La funzione cambia i tempi di ultimo accesso e di ultima modifica del file
specificato dall'argomento \param{filename}, e richiede come secondo argomento
il puntatore ad una struttura \struct{utimbuf}, la cui definizione è riportata
in fig.~\ref{fig:struct_utimebuf}, con i nuovi valori di detti tempi
(rispettivamente nei campi \var{actime} e \var{modtime}). Se si passa un
puntatore nullo verrà impostato il tempo corrente.

\begin{figure}[!htb]
  \footnotesize \centering
  \begin{minipage}[c]{0.8\textwidth}
    \includestruct{listati/utimbuf.h}
  \end{minipage} 
  \normalsize 
  \caption{La struttura \structd{utimbuf}, usata da \func{utime} per modificare
    i tempi dei file.}
  \label{fig:struct_utimebuf}
\end{figure}

L'effetto della funzione ed i privilegi necessari per eseguirla dipendono dal
valore dell'argomento \param{times}. Se è \val{NULL} la funzione imposta il
tempo corrente ed è sufficiente avere accesso in scrittura al file o essere
proprietari del file o avere i privilegi di amministratore. Se invece si è
specificato un valore diverso la funzione avrà successo solo se si è
proprietari del file o se si hanno i privilegi di amministratore.\footnote{per
  essere precisi la capacità \const{CAP\_FOWNER}, vedi
  sez.~\ref{sec:proc_capabilities}.} In entrambi i casi per verificare la
proprietà del file viene utilizzato l'\ids{UID} effettivo del processo.

Si tenga presente che non è possibile modificare manualmente il tempo di
cambiamento di stato del file, che viene aggiornato direttamente dal kernel
tutte le volte che si modifica l'\textit{inode}, e quindi anche alla chiamata
di \func{utime}.  Questo serve anche come misura di sicurezza per evitare che
si possa modificare un file nascondendo completamente le proprie tracce. In
realtà la cosa resta possibile se si è in grado di accedere al file di
dispositivo, scrivendo direttamente sul disco senza passare attraverso il
filesystem, ma ovviamente in questo modo la cosa è più complicata da
realizzare.\footnote{esistono comunque molti programmi che consentono di farlo
  con relativa semplicità per cui non si dia per scontato che il valore sia
  credibile in caso di macchina compromessa.}

A partire dal kernel 2.6 la risoluzione dei tempi dei file, che nei campi di
tab.~\ref{tab:file_file_times} è espressa in secondi, è stata portata ai
nanosecondi per la gran parte dei filesystem. L'ulteriore informazione può
essere ottenuta attraverso altri campi appositamente aggiunti alla struttura
\struct{stat}. Se si sono definite le macro \macro{\_BSD\_SOURCE} o
\macro{\_SVID\_SOURCE} questi sono \var{st\_atim.tv\_nsec},
\var{st\_mtim.tv\_nsec} e \var{st\_ctim.tv\_nsec} se queste non sono definite,
\var{st\_atimensec}, \var{st\_mtimensec} e \var{st\_mtimensec}. Qualora il
supporto per questa maggior precisione sia assente questi campi aggiuntivi
saranno nulli.

Per la gestione di questi nuovi valori è stata definita una seconda funzione
di sistema, \funcd{utimes}, che consente di specificare tempi con maggior
precisione; il suo prototipo è:

\begin{funcproto}{
\fhead{sys/time.h}
\fdecl{int utimes(const char *filename, struct timeval times[2])}
\fdesc{Modifica i tempi di ultimo accesso e ultima modifica di un file.} 
}
{La funzione ritorna $0$ in caso di successo e $-1$ per un errore, nel qual
  caso \var{errno} assumerà gli stessi valori di \func{utime}.}  
\end{funcproto}
 
La funzione è del tutto analoga alla precedente \func{utime} ma usa come
secondo argomento un vettore di due strutture \struct{timeval}, la cui
definizione è riportata in fig.~\ref{fig:sys_timeval_struct}, che consentono
di indicare i tempi con una precisione del microsecondo. Il primo elemento
di \param{times} indica il valore per il tempo di ultimo accesso, il secondo
quello per il tempo di ultima modifica. Se si indica come secondo argomento un
puntatore nullo di nuovo verrà utilizzato il tempo corrente.

\begin{figure}[!htb]
  \footnotesize \centering
  \begin{minipage}[c]{0.8\textwidth}
    \includestruct{listati/timeval.h}
  \end{minipage} 
  \normalsize 
  \caption{La struttura \structd{timeval} usata per indicare valori di tempo
    con la precisione del microsecondo.}
  \label{fig:sys_timeval_struct}
\end{figure}


% Oltre ad \func{utimes} su Linux sono presenti altre due funzioni per la
% manipolazione dei tempi dei file,\footnote{le due funzioni non sono definite
%   in nessuno standard, ma sono presenti, oltre che su Linux, anche su BSD;
%   sono accessibili definendo \macro{\_DEFAULT\_SOURCE} dalla \acr{glibc} 2.19
%   o \macro{\_GNU\_SOURCE} prima.} la prima è \funcd{futimes} e consente di
% effettuare la modifica utilizzando un file già aperto, il suo prototipo è:

% \begin{funcproto}{
% \fhead{sys/time.h}
% \fdecl{int futimes(int fd, const struct timeval tv[2])}
% \fdesc{Cambia i tempi di un file già aperto.} 
% }

% {La funzione ritornano $0$ in caso di successo e $-1$ per un errore, nel qual
%   caso \var{errno} assumerà uno gli stessi valori di \func{utimes} ed inoltre:
%   \begin{errlist}
%   \item[\errcode{EBADF}] \param{fd} non è un file descriptor.
%   \item[\errcode{ENOSYS}] il filesystem \texttt{/proc} non è accessibile.
%   \end{errlist}}  
% \end{funcproto}

% La seconda funzione, introdotta a partire dal kernel 2.6.22, è
% \funcd{lutimes}, e consente rispettivamente di modificare i tempi di un
% collegamento simbolico; il suo prototipo è:

% \begin{funcproto}{
% \fhead{sys/time.h}
% \fdecl{int lutimes(const char *filename, const struct timeval tv[2])}
% \fdesc{Cambia i tempi di un collegamento simbolico.} 
% }

% {La funzione ritornano $0$ in caso di successo e $-1$ per un errore, nel qual
%   caso \var{errno} assumerà uno gli stessi valori di \func{utimes}, con in più:
%   \begin{errlist}
%   \item[\errcode{ENOSYS}] la funzione non è supportata.
%   \end{errlist}}
% \end{funcproto}

% Le due funzioni hanno lo stesso comportamento di \texttt{utimes} e richiedono
% gli stessi privilegi per poter operare, la differenza è che con \func{futimes}
% si può indicare il file su cui operare se questo è già aperto facendo
% riferimento al suo file descriptor, mentre con \func{lutimes}, nel caso in cui
% \param{filename} sia un collegamento simbolico, saranno modificati i suoi
% tempi invece di quelli del file a cui esso punta.

Oltre ad \func{utimes} su Linux sono presenti altre due funzioni,\footnote{le
  due funzioni non sono definite in nessuno standard, ma sono presenti, oltre
  che su Linux, anche su BSD; sono accessibili definendo
  \macro{\_DEFAULT\_SOURCE} dalla \acr{glibc} 2.19 o \macro{\_GNU\_SOURCE}
  prima.} \funcd{futimes} e \funcd{lutimes}, che consentono rispettivamente
di effettuare la modifica utilizzando un file già aperto o di eseguirla
direttamente su un collegamento simbolico. I relativi prototipi sono:

\begin{funcproto}{
\fhead{sys/time.h}
\fdecl{int futimes(int fd, const struct timeval tv[2])}
\fdesc{Cambia i tempi di un file già aperto.} 
\fdecl{int lutimes(const char *filename, const struct timeval tv[2])}
\fdesc{Cambia i tempi di un collegamento simbolico.} 
}

{Le funzioni ritornano $0$ in caso di successo e $-1$ per un errore, nel qual
  caso \var{errno} assumerà uno gli stessi valori di \func{utimes}, con in più
  per \func{futimes}:
  \begin{errlist}
  \item[\errcode{EBADF}] \param{fd} non è un file descriptor.
  \item[\errcode{ENOSYS}] il filesystem \texttt{/proc} non è accessibile per
    \func{futimes} o la funzione non è supportata per \func{lutimes}.
  \end{errlist}}  
\end{funcproto}

Le due funzioni hanno lo stesso comportamento di \texttt{utimes} e richiedono
gli stessi privilegi per poter operare, la differenza è che con \func{futimes}
si può indicare il file su cui operare se questo è già aperto facendo
riferimento al suo file descriptor, mentre con \func{lutimes} nel caso in
cui \param{filename} sia un collegamento simbolico saranno modificati i suoi
tempi invece di quelli del file a cui esso punta.

Nonostante il kernel nelle versioni più recenti supporti, come accennato,
risoluzioni dei tempi dei file fino al nanosecondo, le funzioni fin qui
esaminate non consentono di impostare valori con questa precisione. Per questo
sono state introdotte due nuove funzioni di sistema, \func{utimensat} (che
vedremo in sez.~\ref{sec:file_openat} insieme alle altre
\textit{at-functions}), e \funcd{futimens}, il cui prototipo è:

\begin{funcproto}{
\fhead{sys/time.h}
\fdecl{int futimens(int fd, const struct timespec times[2])}
\fdesc{Cambia i tempi di un file già aperto.} 
}

{La funzione ritorna $0$ in caso di successo e $-1$ per un errore, nel qual
  caso \var{errno} assumerà uno dei valori:
  \begin{errlist}
  \item[\errcode{EACCES}] si è richiesta l'impostazione del tempo corrente ma
    non si ha il permesso di scrittura sul file, o non si è proprietari del
    file o non si hanno i privilegi di amministratore; oppure il file è
    immutabile (vedi sez.~\ref{sec:file_perm_overview}).
  \item[\errcode{EBADF}] \param{fd} non è un file descriptor valido.
  \item[\errcode{EFAULT}] \param{times} non è un puntatore valido.
  \item[\errcode{EINVAL}] si sono usati dei valori non corretti per i tempi di
    \param{times}.
  \item[\errcode{EPERM}] si è richiesto un cambiamento nei tempi non al tempo
    corrente, ma non si è proprietari del file o non si hanno i privilegi di
    amministratore; oppure il file è immutabile o \textit{append-only} (vedi
    sez.~\ref{sec:file_perm_overview}).
  \end{errlist}
  ed inoltre per entrambe \errval{EROFS} nel suo significato generico.}
\end{funcproto}

La funzione è sostanzialmente una estensione di \func{futimes} che consente di
specificare i tempi con precisione maggiore. Per questo per indicare i valori
dei tempi da impostare utilizza un vettore \param{times} di due strutture
\struct{timespec} che permettono di indicare il valore del tempo con una
precisione fino al nanosecondo (se ne è riportata la definizione in
fig.~\ref{fig:sys_timespec_struct}).

\begin{figure}[!htb]
  \footnotesize \centering
  \begin{minipage}[c]{0.8\textwidth}
    \includestruct{listati/timespec.h}
  \end{minipage} 
  \normalsize 
  \caption{La struttura \structd{timespec} usata per indicare valori di tempo
    con la precisione del nanosecondo.}
  \label{fig:sys_timespec_struct}
\end{figure}

Come per le precedenti funzioni il primo elemento di \param{times} indica il
tempo di ultimo accesso ed il secondo quello di ultima modifica, e se si usa
il valore \val{NULL} verrà impostato il tempo corrente sia per l'ultimo
accesso che per l'ultima modifica.

Nei singoli elementi di \param{times} si possono inoltre utilizzare due valori
speciali per il campo \var{tv\_nsec}: con \constd{UTIME\_NOW} si richiede
l'uso del tempo corrente, mentre con \constd{UTIME\_OMIT} si richiede di non
impostare il tempo. Si può così aggiornare in maniera specifica soltanto uno
fra il tempo di ultimo accesso e quello di ultima modifica. Quando si usa uno
di questi valori speciali per \var{tv\_nsec} il corrispondente valore di
\var{tv\_sec} viene ignorato.

Questa funzione e \func{utimensat} sono una estensione definita nella
revisione POSIX.1-2008 dello standard POSIX; in Linux sono state introdotte a
partire dal kernel 2.6.22, e sono supportate dalla \acr{glibc} a partire dalla
versione 2.6, si tenga presente però che per kernel precedenti il 2.6.26 le
due funzioni sono difettose nel rispetto di alcuni requisiti minori dello
standard e nel controllo della correttezza dei tempi, per i dettagli dei quali
si rimanda alla pagina di manuale.


\section{Il controllo di accesso ai file}
\label{sec:file_access_control}

Una delle caratteristiche fondamentali di tutti i sistemi unix-like è quella
del controllo di accesso ai file, che viene implementato per qualunque
filesystem standard.\footnote{per filesystem standard si intende un filesystem
  che implementi le caratteristiche previste dallo standard POSIX; in Linux
  sono utilizzabili anche filesystem di altri sistemi operativi, che non
  supportano queste caratteristiche.} In questa sezione ne esamineremo i
concetti essenziali e le funzioni usate per gestirne i vari aspetti.


\subsection{I permessi per l'accesso ai file}
\label{sec:file_perm_overview}

Ad ogni file Linux associa sempre, oltre ad un insieme di permessi, l'utente
che ne è proprietario (il cosiddetto \textit{owner}) ed un gruppo di
appartenenza, indicati dagli identificatori di utente e gruppo (\ids{UID} e
\ids{GID}) di cui abbiamo già parlato in
sez.~\ref{sec:proc_access_id}.\footnote{questo è vero solo per filesystem di
  tipo Unix, ad esempio non è vero per il filesystem VFAT di Windows, che non
  fornisce nessun supporto per l'accesso multiutente, e per il quale queste
  proprietà vengono assegnate in maniera fissa con opportune opzioni di
  montaggio.}  Anche questi sono mantenuti sull'\textit{inode} insieme alle
altre proprietà e sono accessibili da programma tramite la funzione
\func{stat} (trattata in sez.~\ref{sec:file_stat}), che restituisce l'utente
proprietario nel campo \var{st\_uid} ed il gruppo proprietario nel campo
\var{st\_gid} della omonima struttura \struct{stat}.

Il controllo di accesso ai file segue un modello abbastanza semplice che
prevede tre permessi fondamentali strutturati su tre livelli di accesso.
Esistono varie estensioni a questo modello,\footnote{come le \textit{Access
    Control List} che sono state aggiunte ai filesystem standard con opportune
  estensioni (vedi sez.~\ref{sec:file_ACL}) per arrivare a meccanismi di
  controllo ancora più sofisticati come il \textit{Mandatory Access Control}
  di \textit{SELinux} e delle altre estensioni come \textit{Smack} o
  \textit{AppArmor}.} ma nella maggior parte dei casi il meccanismo standard è
più che sufficiente a soddisfare tutte le necessità più comuni.  I tre
permessi di base associati ad ogni file sono:
\begin{itemize*}
\item il permesso di lettura (indicato con la lettera \texttt{r}, dall'inglese
  \textit{read}).
\item il permesso di scrittura (indicato con la lettera \texttt{w},
  dall'inglese \textit{write}).
\item il permesso di esecuzione (indicato con la lettera \texttt{x},
  dall'inglese \textit{execute}).
\end{itemize*}
mentre i tre livelli su cui sono divisi i privilegi sono:
\begin{itemize*}
\item i privilegi per l'utente proprietario del file.
\item i privilegi per un qualunque utente faccia parte del gruppo cui
  appartiene il file.
\item i privilegi per tutti gli altri utenti.
\end{itemize*}

L'insieme dei permessi viene espresso con un numero a 12 bit; di questi i nove
meno significativi sono usati a gruppi di tre per indicare i permessi base di
lettura, scrittura ed esecuzione e sono applicati rispettivamente
rispettivamente al proprietario, al gruppo, a tutti gli altri.

\begin{figure}[htb]
  \centering
  \includegraphics[width=8cm]{img/fileperm}
  \caption{Lo schema dei bit utilizzati per specificare i permessi di un file
    contenuti nel campo \var{st\_mode} di \struct{stat}.}
  \label{fig:file_perm_bit}
\end{figure}

I restanti tre bit (noti come \textit{suid bit}, \textit{sgid bit}, e
\textit{sticky bit}) sono usati per indicare alcune caratteristiche più
complesse del meccanismo del controllo di accesso su cui torneremo in seguito
(in sez.~\ref{sec:file_special_perm}), lo schema di allocazione dei bit è
riportato in fig.~\ref{fig:file_perm_bit}.  Come tutte le altre proprietà di
un file anche i permessi sono memorizzati nell'\textit{inode}, e come
accennato in sez.~\ref{sec:file_types} essi vengono restituiti in una parte
del campo \var{st\_mode} della struttura \struct{stat} (si veda di nuovo
fig.~\ref{fig:file_stat_struct}).

\begin{table}[htb]
  \centering
    \footnotesize
  \begin{tabular}[c]{|c|l|}
    \hline
    \textbf{\var{st\_mode}} bit & \textbf{Significato} \\
    \hline 
    \hline 
    \const{S\_IRUSR} & \textit{user-read}, l'utente può leggere.\\
    \const{S\_IWUSR} & \textit{user-write}, l'utente può scrivere.\\
    \const{S\_IXUSR} & \textit{user-execute}, l'utente può eseguire.\\ 
    \hline            
    \const{S\_IRGRP} & \textit{group-read}, il gruppo può leggere.\\
    \const{S\_IWGRP} & \textit{group-write}, il gruppo può scrivere.\\
    \const{S\_IXGRP} & \textit{group-execute}, il gruppo può eseguire.\\
    \hline            
    \const{S\_IROTH} & \textit{other-read}, tutti possono leggere.\\
    \const{S\_IWOTH} & \textit{other-write}, tutti possono scrivere.\\
    \const{S\_IXOTH} & \textit{other-execute}, tutti possono eseguire.\\
    \hline              
  \end{tabular}
  \caption{I bit dei permessi di accesso ai file, come definiti in 
    \texttt{<sys/stat.h>}}
  \label{tab:file_bit_perm}
\end{table}


In genere ci si riferisce ai tre livelli dei privilegi usando le lettere
\texttt{u} (per \textit{user}), \texttt{g} (per \textit{group}) e \texttt{o}
(per \textit{other}), inoltre se si vuole indicare tutti i raggruppamenti
insieme si usa la lettera \texttt{a} (per \textit{all}). Si tenga ben presente
questa distinzione dato che in certi casi, mutuando la terminologia in uso a
suo tempo nel VMS, si parla dei permessi base come di permessi per
\textit{owner}, \textit{group} ed \textit{all}, le cui iniziali possono dar
luogo a confusione.  Le costanti che permettono di accedere al valore numerico
di questi bit nel campo \var{st\_mode}, già viste in
tab.~\ref{tab:file_mode_flags}, sono riportate per chiarezza una seconda volta
in tab.~\ref{tab:file_bit_perm}.

I permessi vengono usati in maniera diversa dalle varie funzioni, e a seconda
che si riferiscano a dei file, dei collegamenti simbolici o delle directory;
qui ci limiteremo ad un riassunto delle regole generali, entrando nei dettagli
più avanti.

La prima regola è che per poter accedere ad un file attraverso il suo
\textit{pathname} occorre il permesso di esecuzione in ciascuna delle
directory che compongono il \textit{pathname}; lo stesso vale per aprire un
file nella directory corrente (per la quale appunto serve il diritto di
esecuzione). Per una directory infatti il permesso di esecuzione significa che
essa può essere attraversata nella risoluzione del \textit{pathname}, e per
questo viene anche chiamato permesso di attraversamento. Esso è sempre
distinto dal permesso di lettura che invece implica che si può leggere il
contenuto della directory.

Questo significa che se si ha il permesso di esecuzione senza permesso di
lettura si potrà lo stesso aprire un file all'interno di una directory (se si
hanno i permessi adeguati per il medesimo) ma non si potrà vederlo con
\cmd{ls} mancando il permesso di leggere il contenuto della directory. Per
crearlo o rinominarlo o cancellarlo invece occorrerà avere anche il permesso
di scrittura per la directory.

Avere il permesso di lettura per un file consente di aprirlo con le opzioni
(si veda quanto riportato in sez.~\ref{sec:file_open_close}) di sola lettura o
di lettura/scrittura e leggerne il contenuto. Avere il permesso di scrittura
consente di aprire un file in sola scrittura o lettura/scrittura e modificarne
il contenuto, lo stesso permesso è necessario per poter troncare il file o per
aggiornare il suo tempo di ultima modifica al tempo corrente, ma non per
modificare arbitrariamente quest'ultimo, operazione per la quale, come per
buona parte delle modifiche effettuate sui metadati del file, occorre esserne
i proprietari.

Non si può creare un file fintanto che non si disponga del permesso di
esecuzione e di quello di scrittura per la directory di destinazione. Gli
stessi permessi occorrono per cancellare un file da una directory (si ricordi
che questo non implica necessariamente la rimozione del contenuto del file dal
disco).

Per la cancellazione non è necessario nessun tipo di permesso per il file
stesso dato che, come illustrato in sez.~\ref{sec:link_symlink_rename} esso
non viene toccato, nella cancellazione infatti viene solo modificato il
contenuto della directory, rimuovendo la voce che ad esso fa riferimento. Lo
stesso vale per poter rinominare o spostare il file in altra directory, in
entrambi i casi occorrerà il permesso di scrittura sulle directory che si
vanno a modificare.

Per poter eseguire un file, che sia un programma compilato od uno script di
shell, od un altro tipo di file eseguibile riconosciuto dal kernel, occorre
oltre al permesso di lettura per accedere al contenuto avere anche il permesso
di esecuzione. Inoltre solo i file regolari possono essere eseguiti. Per i
file di dispositivo i permessi validi sono solo quelli di lettura e scrittura,
che corrispondono al poter eseguire dette operazioni sulla periferica
sottostante. 

I permessi per un collegamento simbolico sono ignorati, contano quelli del
file a cui fa riferimento; per questo in genere il comando \cmd{ls} riporta
per un collegamento simbolico tutti i permessi come concessi. Utente e gruppo
a cui esso appartiene vengono pure ignorati quando il collegamento viene
risolto, vengono controllati solo quando viene richiesta la rimozione del
collegamento e quest'ultimo è in una directory con lo \textit{sticky bit}
impostato (si veda sez.~\ref{sec:file_special_perm}).

La procedura con cui il kernel stabilisce se un processo possiede un certo
permesso (di lettura, scrittura o esecuzione) si basa sul confronto fra
l'utente e il gruppo a cui il file appartiene (i valori di \var{st\_uid} e
\var{st\_gid} accennati in precedenza) e l'\ids{UID} effettivo, il \ids{GID}
effettivo e gli eventuali \ids{GID} supplementari del processo.\footnote{in
  realtà Linux, per quanto riguarda l'accesso ai file, utilizza gli
  identificatori del gruppo \textit{filesystem} (si ricordi quanto esposto in
  sez.~\ref{sec:proc_perms}), ma essendo questi del tutto equivalenti ai primi,
  eccetto il caso in cui si voglia scrivere un server NFS, ignoreremo questa
  differenza.}

Per una spiegazione dettagliata degli identificatori associati ai processi si
veda sez.~\ref{sec:proc_perms}; normalmente, a parte quanto vedremo in
sez.~\ref{sec:file_special_perm}, l'\ids{UID} effettivo e il \ids{GID} effettivo
corrispondono ai valori dell'\ids{UID} e del \ids{GID} dell'utente che ha
lanciato il processo, mentre i \ids{GID} supplementari sono quelli dei gruppi
cui l'utente appartiene.

I passi attraverso i quali viene stabilito se il processo possiede il diritto
di accesso sono i seguenti:
\begin{enumerate*}
\item Se l'\ids{UID} effettivo del processo è zero (corrispondente
  all'amministratore) l'accesso è sempre garantito senza nessun controllo. Per
  questo motivo l'amministratore ha piena libertà di accesso a tutti i file.
\item Se l'\ids{UID} effettivo del processo è uguale all'\ids{UID} del
  proprietario del file (nel qual caso si dice che il processo è proprietario
  del file) allora:
  \begin{itemize*}
  \item se il relativo\footnote{per relativo si intende il bit di
      \textit{user-read} se il processo vuole accedere in lettura, quello di
      \textit{user-write} per l'accesso in scrittura, ecc.} bit dei permessi
    d'accesso dell'utente è impostato, l'accesso è consentito;
  \item altrimenti l'accesso è negato.
  \end{itemize*}
\item Se il \ids{GID} effettivo del processo o uno dei \ids{GID} supplementari
  del processo corrispondono al \ids{GID} del file allora:
  \begin{itemize*}
  \item se il bit dei permessi d'accesso del gruppo è impostato, l'accesso è
    consentito;
  \item altrimenti l'accesso è negato.
  \end{itemize*}
\item Se il bit dei permessi d'accesso per tutti gli altri è impostato,
  l'accesso è consentito, altrimenti l'accesso è negato.
\end{enumerate*}

Si tenga presente che questi passi vengono eseguiti esattamente in
quest'ordine. Questo vuol dire che se un processo è il proprietario di un
file, l'accesso è consentito o negato solo sulla base dei permessi per
l'utente; i permessi per il gruppo non vengono neanche controllati. Lo stesso
vale se il processo appartiene ad un gruppo appropriato, in questo caso i
permessi per tutti gli altri non vengono controllati. 

Questo significa che se si è proprietari di un file ma non si ha il permesso
di scrittura, non vi si potrà scrivere anche se questo fosse scrivibile per
tutti gli altri. Permessi di questo tipo sono ovviamente poco ortodossi, e
comunque, come vedremo in sez.~\ref{sec:file_perm_management}, il proprietario
di un file può sempre modificarne i permessi, e riassegnarsi un eventuale
permesso di scrittura mancante.

\itindbeg{file~attributes} 

A questi che sono i permessi ordinari si aggiungono, per i filesystem che
supportano questa estensione, due permessi speciali mantenuti nei cosiddetti
\textit{file attributes}, che si possono leggere ed impostare con i comandi
\cmd{lsattr} e \cmd{chattr}.\footnote{per l'utilizzo di questi comandi e per
  le spiegazioni riguardo tutti gli altri \textit{file attributes} si rimanda
  alla sezione 1.4.4 di \cite{AGL}.}

Il primo è il cosiddetto attributo di immutabilità (\textit{immutable},
identificato dalla lettera \texttt{i}) che impedisce ogni modifica al file,
\textit{inode} compreso. Questo significa non solo che non se ne può cambiare
il contenuto, ma neanche nessuna delle sue proprietà, ed in particolare non si
può modificare nei permessi o nei tempi o nel proprietario ed inoltre, visto
che non se può modificare il \textit{link count}, non si può neanche
cancellare, rinominare, o creare \textit{hard link} verso di esso.

Il secondo è il cosiddetto attributo di \textit{append-only}, (identificato
dalla lettera \texttt{a}) che consente soltanto la scrittura in coda al
file. Il file cioè può essere soltanto esteso nel contenuto, ma i suoi
metadati, a parte i tempi che però possono essere impostati al valore
corrente, non possono essere modificati in alcun modo, quindi di nuovo non si
potrà cancellare, rinominare, o modificare nei permessi o nelle altre
proprietà.

Entrambi questi attributi attivano queste restrizioni a livello di filesystem,
per cui a differenza dei permessi ordinari esse varranno per qualunque utente
compreso l'amministratore. L'amministratore è l'unico che può attivare o
disattivare questi attributi,\footnote{più precisamente un processo con la
  capacità \const{CAP\_LINUX\_IMMUTABLE}, vedi
  sez.~\ref{sec:proc_capabilities}.} e potendo rimuoverli è comunque capace di
tornare in grado di eseguire qualunque operazione su un file immutabile o
\textit{append-only}.

\itindend{file~attributes}



\subsection{I bit dei permessi speciali}
\label{sec:file_special_perm}

\itindbeg{suid~bit}
\itindbeg{sgid~bit}

Come si è accennato (in sez.~\ref{sec:file_perm_overview}) nei dodici bit del
campo \var{st\_mode} di \struct{stat} che vengono usati per il controllo di
accesso oltre ai bit dei permessi veri e propri, ci sono altri tre bit che
vengono usati per indicare alcune proprietà speciali dei file.  Due di questi
sono i bit detti \acr{suid} (da \textit{set-user-ID bit}) e \acr{sgid} (da
\textit{set-group-ID bit}) che sono identificati dalle costanti
\const{S\_ISUID} e \const{S\_ISGID}.

Come spiegato in dettaglio in sez.~\ref{sec:proc_exec}, quando si lancia un
programma il comportamento normale del kernel è quello di impostare gli
identificatori del gruppo \textit{effective} del nuovo processo al valore dei
corrispondenti del gruppo \textit{real} del processo corrente, che normalmente
corrispondono a quelli dell'utente con cui si è entrati nel sistema.

Se però il file del programma, che ovviamente deve essere
eseguibile,\footnote{anzi più precisamente un binario eseguibile: per motivi
  di sicurezza il kernel ignora i bit \acr{suid} e \acr{sgid} per gli script
  eseguibili.} ha il bit \acr{suid} impostato, il kernel assegnerà come
\ids{UID} effettivo al nuovo processo l'\ids{UID} del proprietario del file al
posto dell'\ids{UID} del processo originario.  Avere il bit \acr{sgid}
impostato ha lo stesso effetto sul \ids{GID} effettivo del processo. É
comunque possibile riconoscere questa situazione perché il cambiamento viene
effettuato solo sugli identificativi del gruppo \textit{effective}, mentre
quelli dei gruppi \textit{real} e \textit{saved} restano quelli dell'utente
che ha eseguito il programma.

I bit \acr{suid} e \acr{sgid} vengono usati per permettere agli utenti normali
di usare programmi che richiedono privilegi speciali. L'esempio classico è il
comando \cmd{passwd} che ha la necessità di modificare il file delle password,
quest'ultimo ovviamente può essere scritto solo dall'amministratore, ma non è
necessario chiamare l'amministratore per cambiare la propria password. Infatti
il comando \cmd{passwd} appartiene in genere all'utente \textit{root} ma ha il
bit \acr{suid} impostato, per cui quando viene lanciato da un utente normale
ottiene comunque  i privilegi di amministratore.

Chiaramente avere un processo che ha privilegi superiori a quelli che avrebbe
normalmente l'utente che lo ha lanciato comporta vari rischi, e questo tipo di
programmi devono essere scritti accuratamente per evitare che possano essere
usati per guadagnare privilegi non consentiti (l'argomento è affrontato in
dettaglio in sez.~\ref{sec:proc_perms}).

La presenza dei bit \acr{suid} e \acr{sgid} su un file può essere rilevata con
il comando \cmd{ls -l}, che visualizza una lettera ``\cmd{s}'' al posto della
``\cmd{x}'' in corrispondenza dei permessi di utente o gruppo. La stessa
lettera ``\cmd{s}'' può essere usata nel comando \cmd{chmod} per impostare
questi bit.  Infine questi bit possono essere controllati all'interno di
\var{st\_mode} con l'uso delle due costanti \const{S\_ISUID} e
\const{S\_IGID}, i cui valori sono riportati in
tab.~\ref{tab:file_mode_flags}.

Gli stessi bit vengono ad assumere un significato completamente diverso per le
directory, normalmente infatti Linux usa la convenzione di SVr4 per indicare
con questi bit l'uso della semantica BSD nella creazione di nuovi file (si
veda sez.~\ref{sec:file_ownership_management} per una spiegazione dettagliata
al proposito).

Infine Linux utilizza il bit \acr{sgid} per un'ulteriore estensione mutuata da
SVr4. Il caso in cui un file ha il bit \acr{sgid} impostato senza che lo sia
anche il corrispondente bit di esecuzione viene utilizzato per attivare per
quel file il \textit{mandatory locking} (affronteremo questo argomento in
dettaglio più avanti, in sez.~\ref{sec:file_mand_locking}).

\itindend{suid~bit}
\itindend{sgid~bit}

\itindbeg{sticky~bit}

L'ultimo dei bit rimanenti, identificato dalla costante \const{S\_ISVTX}, è in
parte un rimasuglio delle origini dei sistemi Unix. A quell'epoca infatti la
memoria virtuale e l'accesso ai file erano molto meno sofisticati e per
ottenere la massima velocità possibile per i programmi usati più comunemente
si poteva impostare questo bit.

L'effetto di questo bit era che il segmento di testo del programma (si veda
sez.~\ref{sec:proc_mem_layout} per i dettagli) veniva scritto nella swap la
prima volta che questo veniva lanciato, e vi permaneva fino al riavvio della
macchina (da questo il nome di \textsl{sticky bit}); essendo la swap un file
continuo o una partizione indicizzata direttamente si poteva risparmiare in
tempo di caricamento rispetto alla ricerca attraverso la struttura del
filesystem. Lo \textsl{sticky bit} è indicato usando la lettera ``\texttt{t}''
al posto della ``\texttt{x}'' nei permessi per gli altri.

Ovviamente per evitare che gli utenti potessero intasare la swap solo
l'amministratore era in grado di impostare questo bit, che venne chiamato
anche con il nome di \textit{saved text bit}, da cui deriva quello della
costante.  Le attuali implementazioni di memoria virtuale e filesystem rendono
sostanzialmente inutile questo procedimento.

Benché ormai non venga più utilizzato per i file, lo \textit{sticky bit} ha
invece assunto un uso importante per le directory;\footnote{lo \textit{sticky
    bit} per le directory è un'estensione non definita nello standard POSIX,
  Linux però la supporta, così come BSD e SVr4.} in questo caso se tale bit è
impostato un file potrà essere rimosso dalla directory soltanto se l'utente ha
il permesso di scrittura su di essa ed inoltre è vera una delle seguenti
condizioni:
\begin{itemize*}
\item l'utente è proprietario del file,
\item l'utente è proprietario della directory,
\item l'utente è l'amministratore.
\end{itemize*}

Un classico esempio di directory che ha questo bit impostato è \file{/tmp}, i
cui permessi infatti di solito sono i seguenti:
\begin{Console}
piccardi@hain:~/gapil$ \textbf{ls -ld /tmp}
drwxrwxrwt    6 root     root         1024 Aug 10 01:03 /tmp
\end{Console}
%$
quindi con lo \textit{sticky bit} bit impostato. In questo modo qualunque
utente nel sistema può creare dei file in questa directory, che come
suggerisce il nome è normalmente utilizzata per la creazione di file
temporanei, ma solo l'utente che ha creato un certo file potrà cancellarlo o
rinominarlo. In questo modo si evita che un utente possa, più o meno
consapevolmente, cancellare i file temporanei creati degli altri utenti.

\itindend{sticky~bit}



\subsection{Le funzioni per la gestione dei permessi dei file}
\label{sec:file_perm_management}

Come visto in sez.~\ref{sec:file_perm_overview} il controllo di accesso ad un
file viene fatto utilizzando l'\ids{UID} ed il \ids{GID} effettivo del
processo; ci sono casi però in cui si può voler effettuare il controllo con
l'\ids{UID} reale ed il \ids{GID} reale, vale a dire usando i valori di
\ids{UID} e \ids{GID} relativi all'utente che ha lanciato il programma, e che,
come accennato in sez.~\ref{sec:file_special_perm} e spiegato in dettaglio in
sez.~\ref{sec:proc_perms}, non è detto siano uguali a quelli effettivi.

Per far questo si può usare la funzione di sistema \funcd{access}, il cui
prototipo è:

\begin{funcproto}{
\fhead{unistd.h}
\fdecl{int access(const char *pathname, int mode)}
\fdesc{Verifica i permessi di accesso.} 
}

{La funzione ritorna $0$ in caso di successo e $-1$ per un errore, nel qual
  caso \var{errno} assumerà uno dei valori: 
  \begin{errlist} 
  \item[\errcode{EACCES}] l'accesso al file non è consentito, o non si ha il
    permesso di attraversare una delle directory di \param{pathname}.
  \item[\errcode{EINVAL}] il valore di \param{mode} non è valido.
  \item[\errcode{EROFS}] si è richiesto l'accesso in scrittura per un file su
    un filesystem montato in sola lettura.
  \item[\errcode{ETXTBSY}] si è richiesto l'accesso in scrittura per un
    eseguibile binario correntemente in esecuzione.
  \end{errlist}
  ed inoltre \errval{EFAULT}, \errval{EIO}, \errval{ELOOP},
  \errval{ENAMETOOLONG}, \errval{ENOENT}, \errval{ENOTDIR} nel loro
  significato generico.}
\end{funcproto}

La funzione verifica i permessi di accesso, indicati da \param{mode}, per il
file indicato da \param{pathname}. I valori possibili per l'argomento
\param{mode} sono esprimibili come combinazione delle costanti numeriche
riportate in tab.~\ref{tab:file_access_mode_val} (attraverso un OR binario
delle stesse). I primi tre valori implicano anche la verifica dell'esistenza
del file, se si vuole verificare solo quest'ultima si può usare \const{F\_OK},
o anche direttamente \func{stat}. Nel caso in cui \param{pathname} si
riferisca ad un collegamento simbolico, questo viene seguito ed il controllo è
fatto sul file a cui esso fa riferimento.

La funzione controlla solo i bit dei permessi di accesso per \param{pathname},
ma occorre poter risolvere quest'ultimo, e se non c'è il permesso di
esecuzione per una qualunque delle sue componenti la funzione fallirà
indipendentemente dai permessi del file.  Si tenga presente poi che il fatto
che una directory abbia permesso di scrittura non significa che vi si possa
scrivere come fosse un file, e che se un file ha il permesso di esecuzione non
è detto che sia eseguibile. La funzione ha successo solo se tutti i permessi
controllati sono disponibili.

\begin{table}[htb]
  \centering
  \footnotesize
  \begin{tabular}{|c|l|}
    \hline
    \textbf{\param{mode}} & \textbf{Significato} \\
    \hline
    \hline
    \constd{R\_OK} & Verifica il permesso di lettura. \\
    \constd{W\_OK} & Verifica il permesso di scrittura. \\
    \constd{X\_OK} & Verifica il permesso di esecuzione. \\
    \constd{F\_OK} & Verifica l'esistenza del file. \\
    \hline
  \end{tabular}
  \caption{Valori possibili per l'argomento \param{mode} della funzione 
    \func{access}.}
  \label{tab:file_access_mode_val}
\end{table}

Un esempio tipico per l'uso di questa funzione è quello di un processo che sta
eseguendo un programma coi privilegi di un altro utente (ad esempio attraverso
l'uso del \textit{suid bit}) che vuole controllare se l'utente originale ha i
permessi per accedere ad un certo file, ma eseguire questo controllo prima di
aprire il file espone al rischio di una \textit{race condition} che apre ad un
possibile \textit{symlink attack} fra il controllo e l'apertura del file. In
questo caso è sempre opportuno usare invece la funzione \func{faccessat} che
tratteremo insieme alle altre \textit{at-functions} in
sez.~\ref{sec:file_openat}.


Del tutto analoghe a \func{access} sono le due funzioni \funcm{euidaccess} e
\funcm{eaccess} che ripetono lo stesso controllo usando però gli
identificatori del gruppo effettivo, verificando quindi le effettive capacità
di accesso ad un file. Le funzioni hanno entrambe lo stesso
prototipo\footnote{in realtà \funcm{eaccess} è solo un sinonimo di
  \funcm{euidaccess} fornita per compatibilità con l'uso di questo nome in
  altri sistemi.} che è del tutto identico a quello di \func{access}. Prendono
anche gli stessi valori e restituiscono gli stessi risultati e gli stessi
codici di errore.

Le due funzioni non sono previste da nessuno standard, ed utilizzabili solo
avendo definito \macro{\_GNU\_SOURCE}; inoltre qualora le si vogliano
utilizzare per verificare che un processo abbia i permessi per accedere ad un
file prima di farlo effettivamente, ci si esporrebbe ad una \textit{race
  condition}, dato che i permessi potrebbero cambiare nell'intervallo fra il
controllo e l'accesso effettivo. Per questo motivo in questo caso è molto più
semplice e sicuro tentare direttamente l'accesso, e trattare opportunamente
l'eventuale fallimento per mancanza di permessi.

Per cambiare i permessi di un file sono invece disponibili due funzioni di
sistema \funcd{chmod} e \funcd{fchmod}, che operano rispettivamente usando il
\textit{pathname} di un file o se questo è già aperto sul relativo file
descriptor; i loro prototipi sono:

\begin{funcproto}{
\fhead{sys/types.h}
\fhead{sys/stat.h}
\fdecl{int chmod(const char *path, mode\_t mode)}
\fdesc{Cambia i permessi del file indicato da \param{path} al valore indicato
  da \param{mode}.} 
\fdecl{int fchmod(int fd, mode\_t mode)}
\fdesc{Analoga alla precedente, ma usa il file descriptor \param{fd} per
  indicare il file.}  

}

{La funzione ritorna $0$ in caso di successo e $-1$ per un errore, nel qual
  caso \var{errno} assumerà uno dei valori: 
  \begin{errlist}
  \item[\errcode{EPERM}] l'\ids{UID} effettivo non corrisponde a quello del
    proprietario del file o non si hanno i privilegi di amministratore.
  \item[\errcode{EROFS}] il file è su un filesystem in sola lettura.
  \end{errlist}
  ed inoltre per entrambe \errval{EIO}, per \func{chmod} \errval{EACCES},
  \errval{EFAULT}, \errval{ELOOP}, \errval{ENAMETOOLONG}, \errval{ENOENT},
  \errval{ENOMEM}, \errval{ENOTDIR}, per \func{fchmod} \errval{EBADF} nel loro
  significato generico.}
\end{funcproto}


Entrambe le funzioni utilizzano come secondo argomento \param{mode}, una
variabile dell'apposito tipo primitivo \type{mode\_t} (vedi
tab.~\ref{tab:intro_primitive_types}), che è una maschera binaria da
utilizzare per specificare i permessi sui file.

\begin{table}[!htb]
  \centering
  \footnotesize
  \begin{tabular}[c]{|c|c|l|}
    \hline
    \textbf{\param{mode}} & \textbf{Valore} & \textbf{Significato} \\
    \hline
    \hline
    \const{S\_ISUID} & 04000 & Set user ID  bit.\\
    \const{S\_ISGID} & 02000 & Set group ID bit.\\
    \const{S\_ISVTX} & 01000 & Sticky bit.\\
    \hline
    \const{S\_IRWXU} & 00700 & L'utente ha tutti i permessi.\\
    \const{S\_IRUSR} & 00400 & L'utente ha il permesso di lettura.\\
    \const{S\_IWUSR} & 00200 & L'utente ha il permesso di scrittura.\\
    \const{S\_IXUSR} & 00100 & L'utente ha il permesso di esecuzione.\\
    \hline
    \const{S\_IRWXG} & 00070 & Il gruppo ha tutti i permessi.\\
    \const{S\_IRGRP} & 00040 & Il gruppo ha il permesso di lettura.\\
    \const{S\_IWGRP} & 00020 & Il gruppo ha il permesso di scrittura.\\
    \const{S\_IXGRP} & 00010 & Il gruppo ha il permesso di esecuzione.\\
    \hline
    \const{S\_IRWXO} & 00007 & Gli altri hanno tutti i permessi.\\
    \const{S\_IROTH} & 00004 & Gli altri hanno il permesso di lettura.\\
    \const{S\_IWOTH} & 00002 & Gli altri hanno il permesso di scrittura.\\
    \const{S\_IXOTH} & 00001 & Gli altri hanno il permesso di esecuzione.\\
    \hline
  \end{tabular}
  \caption{Valori delle costanti usate per indicare i vari bit di
    \param{mode} utilizzato per impostare i permessi dei file.}
  \label{tab:file_permission_const}
\end{table}

Si sono riportate in tab.~\ref{tab:file_permission_const} le costanti con cui
indicare i singoli bit di \param{mode}; si noti come corrispondano agli stessi
valori usati per \var{st\_mode} già visti in tab.~\ref{tab:file_mode_flags}.
Il valore di \param{mode} può essere ottenuto combinando fra loro con un OR
binario le costanti simboliche relative ai vari bit, o specificato
direttamente, come per l'omonimo comando di shell \texttt{chmod}, con un
valore numerico (con la shell in genere lo si scrive in ottale, dato che i bit
dei permessi sono divisibili in gruppi di tre), che si può calcolare
direttamente usando lo schema di utilizzo dei bit illustrato in
fig.~\ref{fig:file_perm_bit}.

Ad esempio i permessi standard assegnati ai nuovi file (lettura e scrittura
per il proprietario, sola lettura per il gruppo e gli altri) sono
corrispondenti al valore ottale \texttt{0644}, un programma invece avrebbe
anche il bit di esecuzione attivo, con un valore ottale di \texttt{0755}, se
infine si volesse attivare anche il bit \acr{suid} il valore ottale da fornire
sarebbe \texttt{4755}.

Il cambiamento dei permessi di un file eseguito attraverso queste funzioni ha
comunque alcune limitazioni, previste per motivi di sicurezza. L'uso delle
funzioni infatti è possibile solo se l'\ids{UID} effettivo del processo
chiamante corrisponde a quello del proprietario del file o se il processo ha i
privilegi di amministratore,\footnote{per la precisione la capacità
  \const{CAP\_FOWNER}, vedi sez.~\ref{sec:proc_capabilities}.} in caso
contrario esse falliranno con un errore di \errcode{EPERM}.

Ma oltre a questa regola generale, di immediata comprensione, esistono delle
limitazioni ulteriori. Per questo motivo, anche se si è proprietari del file,
non tutti i valori possibili di \param{mode} sono permessi o hanno effetto;
in particolare accade che:
\begin{enumerate*}
\item siccome solo l'amministratore può impostare lo \textit{sticky bit}, se
  l'\ids{UID} effettivo del processo non è zero esso viene automaticamente
  cancellato, senza notifica di errore, qualora sia stato indicato
  in \param{mode}.
\item per quanto detto in sez.~\ref{sec:file_ownership_management} riguardo la
  creazione dei nuovi file, si può avere il caso in cui il file creato da un
  processo è assegnato ad un gruppo per il quale il processo non ha privilegi.
  Per evitare che si possa assegnare il bit \acr{sgid} ad un file appartenente
  ad un gruppo per cui non si hanno diritti, questo viene automaticamente
  cancellato da \param{mode}, senza notifica di errore, qualora il gruppo del
  file non corrisponda a quelli associati al processo; la cosa non avviene
  quando l'\ids{UID} effettivo del processo è zero.
\end{enumerate*}

Per alcuni filesystem\footnote{i filesystem più comuni (\acr{ext2},
  \acr{ext3}, \acr{ext4}, \acr{XFS}, \acr{Btrfs}) supportano
  questa caratteristica, che è mutuata da BSD.} è inoltre prevista
un'ulteriore misura di sicurezza, volta a scongiurare l'abuso dei bit
\acr{suid} e \acr{sgid}; essa consiste nel cancellare automaticamente questi
bit dai permessi di un file qualora un processo che non appartenga
all'amministratore\footnote{per la precisione un processo che non dispone
  della capacità \const{CAP\_FSETID}, vedi sez.~\ref{sec:proc_capabilities}.}
effettui una scrittura. In questo modo anche se un utente malizioso scopre un
file \acr{suid} su cui può scrivere, un'eventuale modifica comporterà la
perdita di questo privilegio.

Le funzioni \func{chmod} e \func{fchmod} ci permettono di modificare i
permessi di un file, resta però il problema di quali sono i permessi assegnati
quando il file viene creato. Le funzioni dell'interfaccia nativa di Unix, come
vedremo in sez.~\ref{sec:file_open_close}, permettono di indicare
esplicitamente i permessi di creazione di un file, ma questo non è possibile
per le funzioni dell'interfaccia standard ANSI C che non prevede l'esistenza
di utenti e gruppi, ed inoltre il problema si pone anche per l'interfaccia
nativa quando i permessi non vengono indicati esplicitamente.

\itindbeg{umask} 

Per le funzioni dell'interfaccia standard ANSI C l'unico riferimento possibile
è quello della modalità di apertura del nuovo file (lettura/scrittura o sola
lettura), che però può fornire un valore che è lo stesso per tutti e tre i
permessi di sez.~\ref{sec:file_perm_overview} (cioè $666$ nel primo caso e
$222$ nel secondo). Per questo motivo il sistema associa ad ogni
processo\footnote{è infatti contenuta nel campo \var{umask} della struttura
  \kstruct{fs\_struct}, vedi fig.~\ref{fig:proc_task_struct}.}  una maschera
di bit, la cosiddetta \textit{umask}, che viene utilizzata per impedire che
alcuni permessi possano essere assegnati ai nuovi file in sede di creazione. I
bit indicati nella maschera vengono infatti cancellati dai permessi quando un
nuovo file viene creato.\footnote{l'operazione viene fatta sempre: anche
  qualora si indichi esplicitamente un valore dei permessi nelle funzioni di
  creazione che lo consentono, i permessi contenuti nella \textit{umask}
  verranno tolti.}

La funzione di sistema che permette di impostare il valore di questa maschera
di controllo è \funcd{umask}, ed il suo prototipo è:

\begin{funcproto}{
\fhead{stat.h}
\fdecl{mode\_t umask(mode\_t mask)}
\fdesc{Imposta la maschera dei permessi.} 
}

{La funzione ritorna il precedente valore della maschera, non sono previste
  condizioni di errore.}
\end{funcproto}

La funzione imposta la maschera dei permessi dei bit al valore specificato
da \param{mask}, di cui vengono presi solo i 9 bit meno significativi.  In
genere si usa questa maschera per impostare un valore predefinito che escluda
preventivamente alcuni permessi, il caso più comune è eliminare il permesso di
scrittura per il gruppo e gli altri, corrispondente ad un valore
per \param{mask} pari a $022$.  In questo modo è possibile cancellare
automaticamente i permessi non voluti.  Di norma questo valore viene impostato
una volta per tutte al login (a $022$ se non indicato altrimenti), e gli
utenti non hanno motivi per modificarlo.

\itindend{umask} 


\subsection{La gestione della titolarità dei file}
\label{sec:file_ownership_management}

Vedremo in sez.~\ref{sec:file_open_close} con quali funzioni si possono creare
nuovi file, in tale occasione vedremo che è possibile specificare in sede di
creazione quali permessi applicare ad un file, però non si può indicare a
quale utente e gruppo esso deve appartenere.  Lo stesso problema si presenta
per la creazione di nuove directory (procedimento descritto in
sez.~\ref{sec:file_dir_creat_rem}).

Lo standard POSIX prescrive che l'\ids{UID} del nuovo file corrisponda
all'\ids{UID} effettivo del processo che lo crea; per il \ids{GID} invece
prevede due diverse possibilità:
\begin{itemize*}
\item che il \ids{GID} del file corrisponda al \ids{GID} effettivo del
  processo.
\item che il \ids{GID} del file corrisponda al \ids{GID} della directory in
  cui esso è creato.
\end{itemize*}

In genere BSD usa sempre la seconda possibilità, che viene per questo chiamata
semantica BSD. Linux invece segue normalmente quella che viene chiamata
semantica SVr4: di norma un nuovo file viene creato, seguendo la prima
opzione, con il \ids{GID} del processo, se però la directory in cui viene
creato ha il bit \acr{sgid} dei permessi impostato allora viene usata la
seconda opzione. L'adozione di questa semantica però può essere controllata,
all'interno di alcuni filesystem,\footnote{con il kernel 2.6.25 questi erano
  \acr{ext2}, \acr{ext3}, \acr{ext4}, e \acr{XFS}.}  con l'uso dell'opzione di
montaggio \texttt{grpid}, che se attivata fa passare all'uso della semantica
BSD.

Usare la semantica BSD ha il vantaggio che il \ids{GID} viene sempre
automaticamente propagato, restando coerente a quello della directory di
partenza, in tutte le sotto-directory. La semantica SVr4 offre la possibilità
di scegliere, ma per ottenere lo stesso risultato di coerenza che si ha con
quella di BSD necessita che quando si creano nuove directory venga propagato
il bit \acr{sgid}. Questo è il comportamento predefinito del comando
\cmd{mkdir}, ed è in questo modo ad esempio che le varie distribuzioni
assicurano che le sotto-directory create nella home di un utente restino
sempre con il \ids{GID} del gruppo primario dello stesso.

La presenza del bit \acr{sgid} è inoltre molto comoda quando si hanno
directory contenenti file condivisi da un gruppo di utenti in cui possono
scrivere tutti i membri dello stesso, dato che assicura che i file che gli
utenti vi creano appartengano allo gruppo stesso. Questo non risolve però
completamente i problemi di accesso da parte di altri utenti dello stesso
gruppo, in quanto di default i permessi assegnati al gruppo non sono
sufficienti per un accesso in scrittura; in questo caso si deve aver cura di
usare prima della creazione dei file un valore per \textit{umask} lasci il
permesso di scrittura per il gruppo.\footnote{in tal caso si può assegnare
  agli utenti del gruppo una \textit{umask} di $002$, anche se la soluzione
  migliore in questo caso è usare una ACL di default (vedi
  sez.~\ref{sec:file_ACL}).}

Come avviene nel caso dei permessi esistono anche delle appropriate funzioni
di sistema, \funcd{chown}, \funcd{fchown} e \funcd{lchown}, che permettono di
cambiare sia l'utente che il gruppo a cui un file appartiene; i rispettivi
prototipi sono:

\begin{funcproto}{
\fhead{sys/types.h}
\fhead{sys/stat.h}
\fdecl{int chown(const char *path, uid\_t owner, gid\_t group)}
\fdecl{int fchown(int fd, uid\_t owner, gid\_t group)}
\fdecl{int lchown(const char *path, uid\_t owner, gid\_t group)}
\fdesc{Cambiano proprietario e gruppo proprietario di un file.} 
}

{Le funzioni ritornano $0$ in caso di successo e $-1$ per un errore, nel qual
  caso \var{errno} assumerà uno dei valori: 
  \begin{errlist}
  \item[\errcode{EPERM}] l'\ids{UID} effettivo non corrisponde a quello del
    proprietario del file o non è zero, o utente e gruppo non sono validi.
  \end{errlist}
  ed inoltre per tutte \errval{EROFS} e \errval{EIO}, per \func{chown}
  \errval{EACCES}, \errval{EFAULT}, \errval{ELOOP}, \errval{ENAMETOOLONG},
  \errval{ENOENT}, \errval{ENOMEM}, \errval{ENOTDIR}, per \func{fchown}
  \errval{EBADF} nel loro significato generico.}
\end{funcproto}

Le funzioni cambiano utente e gruppo di appartenenza di un file ai valori
specificati dalle variabili \param{owner} e \param{group}. Con Linux solo
l'amministratore\footnote{o in generale un processo con la capacità
  \const{CAP\_CHOWN}, vedi sez.~\ref{sec:proc_capabilities}.} può cambiare il
proprietario di un file; in questo viene seguita la semantica usata da BSD che
non consente agli utenti di assegnare i loro file ad altri utenti evitando
eventuali aggiramenti delle quote.  L'amministratore può cambiare sempre il
gruppo di un file, il proprietario può cambiare il gruppo solo dei file che
gli appartengono e solo se il nuovo gruppo è il suo gruppo primario o uno dei
gruppi di cui fa parte.

La funzione \func{chown} segue i collegamenti simbolici, per operare
direttamente su un collegamento simbolico si deve usare la funzione
\func{lchown}.\footnote{fino alla versione 2.1.81 in Linux \func{chown} non
  seguiva i collegamenti simbolici, da allora questo comportamento è stato
  assegnato alla funzione \func{lchown}, introdotta per l'occasione, ed è
  stata creata una nuova \textit{system call} per \func{chown} che seguisse i
  collegamenti simbolici.} La funzione \func{fchown} opera su un file aperto,
essa è mutuata da BSD, ma non è nello standard POSIX.  Un'altra estensione
rispetto allo standard POSIX è che specificando $-1$ come valore per
\param{owner} e \param{group} i valori restano immutati.

Quando queste funzioni sono chiamate con successo da un processo senza i
privilegi di amministratore entrambi i bit \acr{suid} e \acr{sgid} vengono
cancellati. Questo non avviene per il bit \acr{sgid} nel caso in cui esso sia
usato (in assenza del corrispondente permesso di esecuzione) per indicare che
per il file è attivo il \textit{mandatory locking} (vedi
sez.~\ref{sec:file_mand_locking}).


\subsection{Un quadro d'insieme sui permessi}
\label{sec:file_riepilogo}

Avendo affrontato in maniera separata il comportamento delle varie funzioni
che operano sui permessi dei file ed avendo trattato in sezioni diverse il
significato dei singoli bit dei permessi, vale la pena di fare un riepilogo in
cui si riassumano le caratteristiche di ciascuno di essi, in modo da poter
fornire un quadro d'insieme.

\begin{table}[!htb]
  \centering
  \footnotesize
  \begin{tabular}[c]{|c|c|c|c|c|c|c|c|c|c|c|c|l|}
    \hline
    \multicolumn{3}{|c|}{special}&
    \multicolumn{3}{|c|}{user}&
    \multicolumn{3}{|c|}{group}&
    \multicolumn{3}{|c|}{other}&
    \multirow{2}{*}{\textbf{Significato per i file}} \\
    \cline{1-12}
    \texttt{s}&\texttt{s}&\texttt{t}&
    \texttt{r}&\texttt{w}&\texttt{x}&
    \texttt{r}&\texttt{w}&\texttt{x}&
    \texttt{r}&\texttt{w}&\texttt{x}& \\
    \hline
    \hline
   1&-&-&-&-&1&-&-&-&-&-&-&Eseguito conferisce l'\ids{UID} effettivo dell'utente.\\
   -&1&-&-&-&1&-&-&-&-&-&-&Eseguito conferisce il \ids{GID} effettivo del gruppo.\\
   -&1&-&-&-&0&-&-&-&-&-&-&Il \textit{mandatory locking} è abilitato.\\
   -&-&1&-&-&-&-&-&-&-&-&-&Non utilizzato.\\
   -&-&-&1&-&-&-&-&-&-&-&-&Permesso di lettura per l'utente.\\
   -&-&-&-&1&-&-&-&-&-&-&-&Permesso di scrittura per l'utente.\\
   -&-&-&-&-&1&-&-&-&-&-&-&Permesso di esecuzione per l'utente.\\
   -&-&-&-&-&-&1&-&-&-&-&-&Permesso di lettura per il gruppo.\\
   -&-&-&-&-&-&-&1&-&-&-&-&Permesso di scrittura per il gruppo.\\
   -&-&-&-&-&-&-&-&1&-&-&-&Permesso di esecuzione per il gruppo.\\
   -&-&-&-&-&-&-&-&-&1&-&-&Permesso di lettura per tutti gli altri.\\
   -&-&-&-&-&-&-&-&-&-&1&-&Permesso di scrittura per tutti gli altri.\\
   -&-&-&-&-&-&-&-&-&-&-&1&Permesso di esecuzione per tutti gli altri.\\
    \hline
    \hline
    \multicolumn{3}{|c|}{special}&
    \multicolumn{3}{|c|}{user}&
    \multicolumn{3}{|c|}{group}&
    \multicolumn{3}{|c|}{other}&
    \multirow{2}{*}{\textbf{Significato per le directory}} \\
    \cline{1-12}
    \texttt{s}&\texttt{s}&\texttt{t}&
    \texttt{r}&\texttt{w}&\texttt{x}&
    \texttt{r}&\texttt{w}&\texttt{x}&
    \texttt{r}&\texttt{w}&\texttt{x}& \\
    \hline
    \hline
    1&-&-&-&-&-&-&-&-&-&-&-&Non utilizzato.\\
    -&1&-&-&-&-&-&-&-&-&-&-&Propaga il gruppo ai nuovi file creati.\\
    -&-&1&-&-&-&-&-&-&-&-&-&Solo il proprietario di un file può rimuoverlo.\\
    -&-&-&1&-&-&-&-&-&-&-&-&Permesso di visualizzazione per l'utente.\\
    -&-&-&-&1&-&-&-&-&-&-&-&Permesso di aggiornamento per l'utente.\\
    -&-&-&-&-&1&-&-&-&-&-&-&Permesso di attraversamento per l'utente.\\
    -&-&-&-&-&-&1&-&-&-&-&-&Permesso di visualizzazione per il gruppo.\\
    -&-&-&-&-&-&-&1&-&-&-&-&Permesso di aggiornamento per il gruppo.\\
    -&-&-&-&-&-&-&-&1&-&-&-&Permesso di attraversamento per il gruppo.\\
    -&-&-&-&-&-&-&-&-&1&-&-&Permesso di visualizzazione per tutti gli altri.\\
    -&-&-&-&-&-&-&-&-&-&1&-&Permesso di aggiornamento per tutti gli altri.\\
    -&-&-&-&-&-&-&-&-&-&-&1&Permesso di attraversamento per tutti gli altri.\\
    \hline
  \end{tabular}
  \caption{Tabella riassuntiva del significato dei bit dei permessi per un
    file e directory.} 
  \label{tab:file_fileperm_bits}
\end{table}

Nella parte superiore di tab.~\ref{tab:file_fileperm_bits} si è riassunto il
significato dei vari bit dei permessi per un file ordinario; per quanto
riguarda l'applicazione dei permessi per proprietario, gruppo ed altri si
ricordi quanto illustrato in sez.~\ref{sec:file_perm_overview}.  Per
compattezza, nella tabella si sono specificati i bit di \textit{suid},
\textit{sgid} e \textit{sticky} con la notazione illustrata anche in
fig.~\ref{fig:file_perm_bit}.  Nella parte inferiore si sono invece riassunti
i significati dei vari bit dei permessi per una directory; anche in questo
caso si è riapplicato ai bit di \textit{suid}, \textit{sgid} e \textit{sticky}
la notazione illustrata in fig.~\ref{fig:file_perm_bit}.

Si ricordi infine che i permessi non hanno alcun significato per i
collegamenti simbolici, mentre per i file di dispositivo hanno senso soltanto
i permessi di lettura e scrittura, che si riflettono sulla possibilità di
compiere dette operazioni sul dispositivo stesso.

Nella tabella si è indicato con il carattere ``-'' il fatto che il valore del
bit in questione non è influente rispetto a quanto indicato nella riga della
tabella; la descrizione del significato fa riferimento soltanto alla
combinazione di bit per i quali è stato riportato esplicitamente un valore.
Si rammenti infine che il valore dei bit dei permessi non ha alcun effetto
qualora il processo possieda i privilegi di amministratore.


\section{Caratteristiche e funzionalità avanzate}
\label{sec:file_dir_advances}

Tratteremo qui alcune caratteristiche e funzionalità avanzate della gestione
di file e directory, affrontando anche una serie di estensioni
dell'interfaccia classica dei sistemi unix-like, principalmente utilizzate a
scopi di sicurezza, che sono state introdotte nelle versioni più recenti di
Linux.

\subsection{Gli attributi estesi}
\label{sec:file_xattr}

\itindbeg{Extended~Attributes}

Nelle sezioni precedenti abbiamo trattato in dettaglio le varie informazioni
che il sistema mantiene negli \textit{inode}, e le varie funzioni che
permettono di modificarle.  Si sarà notato come in realtà queste informazioni
siano estremamente ridotte.  Questo è dovuto al fatto che Unix origina negli
anni '70, quando le risorse di calcolo e di spazio disco erano minime. Con il
venir meno di queste restrizioni è incominciata ad emergere l'esigenza di
poter associare ai file delle ulteriori informazioni astratte (quelli che
abbiamo chiamato genericamente \textsl{metadati}) che però non potevano
trovare spazio nei dati classici mantenuti negli \textit{inode}.

Per risolvere questo problema alcuni sistemi unix-like (e fra questi anche
Linux) hanno introdotto un meccanismo generico, detto \textit{Extended
  Attributes} che consente di associare delle informazioni ulteriori ai
singoli file; ad esempio vengono usati per la gestione delle ACL, che
tratteremo in sez.~\ref{sec:file_ACL} e per le \textit{file capabilities}, che
vedremo in sez.~\ref{sec:proc_capabilities}. Gli \textsl{attributi estesi}
(abbreviati in \textsl{xattr}) non sono altro che delle coppie nome/valore che
sono associate permanentemente ad un oggetto sul filesystem, analoghi di
quello che sono le variabili di ambiente (vedi sez.~\ref{sec:proc_environ})
per un processo.

Altri sistemi (come Solaris, MacOS e Windows) hanno adottato un meccanismo
diverso in cui ad un file sono associati diversi flussi di dati, su cui
possono essere mantenute ulteriori informazioni, che possono essere accedute
con le normali operazioni di lettura e scrittura. Questi non vanno confusi con
gli \textit{Extended Attributes} (anche se su Solaris hanno lo stesso nome),
che sono un meccanismo molto più semplice, che pur essendo limitato (potendo
contenere solo una quantità limitata di informazione) hanno il grande
vantaggio di essere molto più semplici da realizzare, più
efficienti,\footnote{cosa molto importante, specie per le applicazioni che
  richiedono una gran numero di accessi, come le ACL.} e di garantire
l'atomicità di tutte le operazioni.

In Linux gli attributi estesi sono sempre associati al singolo \textit{inode}
e l'accesso viene sempre eseguito in forma atomica, in lettura il valore
corrente viene scritto su un buffer in memoria, mentre la scrittura prevede
che ogni valore precedente sia sovrascritto.

Si tenga presente che non tutti i filesystem supportano gli \textit{Extended
  Attributes}; al momento della scrittura di queste dispense essi sono
presenti solo sui vari \acr{extN}, \acr{ReiserFS}, \acr{JFS},
\acr{XFS}, \acr{Btrfs}, \acr{Lustre} e \acr{OCFS2}.  Inoltre a
seconda della implementazione ci possono essere dei limiti sulla quantità di
attributi che si possono utilizzare.\footnote{ad esempio nel caso di
  \acr{ext2} ed \acr{ext3} è richiesto che essi siano contenuti
  all'interno di un singolo blocco, pertanto con dimensioni massime pari a
  1024, 2048 o 4096 byte a seconda delle dimensioni di quest'ultimo impostate
  in fase di creazione del filesystem, mentre con \textsl{XFS} e
  \textsl{Btrfs} non ci sono limiti ed i dati vengono memorizzati in maniera
  diversa (nell'\textit{inode} stesso, in un blocco a parte, o in una
  struttura ad albero dedicata) per mantenerne la scalabilità; lasciamo i
  dettagli dei vari filesystem alla documentazione accessibile con \texttt{man
    xattr}.} Infine lo spazio utilizzato per mantenere gli attributi estesi
viene tenuto in conto per il calcolo delle quote di utente e gruppo
proprietari del file.

Come meccanismo per mantenere informazioni aggiuntive associate al singolo
file, gli \textit{Extended Attributes} possono avere usi anche molto diversi
fra loro.  Per poterli distinguere allora sono stati suddivisi in
\textsl{classi}, a cui poter applicare requisiti diversi per l'accesso e la
gestione. Per questo motivo il nome di un attributo deve essere sempre
specificato nella forma \texttt{namespace.attribute}, dove \texttt{namespace}
fa riferimento alla classe a cui l'attributo appartiene, mentre
\texttt{attribute} è il nome ad esso assegnato. In tale forma il nome di un
attributo esteso deve essere univoco. Al momento sono state definite le
quattro classi di attributi riportate in
tab.~\ref{tab:extended_attribute_class}.

\begin{table}[htb]
  \centering
  \footnotesize
  \begin{tabular}{|l|p{8cm}|}
    \hline
    \textbf{Nome} & \textbf{Descrizione} \\
    \hline
    \hline
    \texttt{security}&Gli \textit{extended security attributes}: vengono
                      utilizzati dalle estensioni di sicurezza del kernel (i
                      \textit{Linux Security Modules}), per le realizzazione
                      di meccanismi evoluti di controllo di accesso come
                      \textit{SELinux} o le \textit{capabilities} dei
                      file di sez.~\ref{sec:proc_capabilities}.\\ 
    \texttt{system} & Gli \textit{extended system attributes}: sono usati
                      dal kernel per memorizzare dati di sistema associati ai
                      file come le ACL (vedi sez.~\ref{sec:file_ACL}) o le
                      \textit{capabilities} (vedi
                      sez.~\ref{sec:proc_capabilities}).\\
    \texttt{trusted}& I \textit{trusted extended attributes}: vengono
                      utilizzati per poter realizzare in \textit{user space} 
                      meccanismi che consentano di mantenere delle
                      informazioni sui file che non devono essere accessibili
                      ai processi ordinari.\\
    \texttt{user}   & Gli \textit{extended user attributes}: utilizzati per
                      mantenere informazioni aggiuntive sui file (come il
                      \textit{mime-type}, la codifica dei caratteri o del
                      file) accessibili dagli utenti.\\
    \hline
  \end{tabular}
  \caption{I nomi utilizzati valore di \texttt{namespace} per distinguere le
    varie classi di \textit{Extended Attributes}.}
  \label{tab:extended_attribute_class}
\end{table}


Dato che uno degli usi degli \textit{Extended Attributes} è di impiegarli per
realizzare delle estensioni al tradizionale meccanismo dei controlli di
accesso di Unix (come le ACL, \textit{SELinux}, ecc.), l'accesso ai loro
valori viene regolato in maniera diversa a seconda sia della loro classe che
di quali, fra le estensioni che li utilizzano, sono poste in uso. In
particolare, per ciascuna delle classi riportate in
tab.~\ref{tab:extended_attribute_class}, si hanno i seguenti casi:
\begin{basedescript}{\desclabelwidth{1.7cm}\desclabelstyle{\nextlinelabel}}
\item[\texttt{security}] L'accesso agli \textit{extended security attributes}
  dipende dalle politiche di sicurezza stabilite da loro stessi tramite
  l'utilizzo di un sistema di controllo basato sui \textit{Linux Security
    Modules} (ad esempio \textit{SELinux}). Pertanto l'accesso in lettura o
  scrittura dipende dalle politiche di sicurezza implementate all'interno dal
  modulo di sicurezza che si sta utilizzando al momento (ciascuno avrà le
  sue). Se non è stato caricato nessun modulo di sicurezza l'accesso in
  lettura sarà consentito a tutti i processi, mentre quello in scrittura solo
  ai processi con privilegi amministrativi (per la precisione dotati della
  capacità \const{CAP\_SYS\_ADMIN}, vedi sez.~\ref{sec:proc_capabilities}).

\item[\texttt{system}] Anche l'accesso agli \textit{extended system
    attributes} dipende dalle politiche di accesso che il kernel realizza
  anche utilizzando gli stessi valori in essi contenuti. Ad esempio nel caso
  delle ACL (vedi sez.~\ref{sec:file_ACL}) l'accesso è consentito in lettura
  ai processi che hanno la capacità di eseguire una ricerca sul file (cioè
  hanno il permesso di lettura sulla directory che contiene il file) ed in
  scrittura al proprietario del file o ai processi dotati della capacità
  \const{CAP\_FOWNER}, vale a dire una politica di accesso analoga a quella
  impiegata per gli ordinari permessi dei file.

\item[\texttt{trusted}] L'accesso ai \textit{trusted extended attributes}, sia
  per la lettura che per la scrittura, è consentito soltanto ai processi con
  privilegi amministrativi dotati della capacità \const{CAP\_SYS\_ADMIN}. In
  questo modo si possono utilizzare questi attributi per realizzare in
  \textit{user space} dei meccanismi di controllo che accedono ad informazioni
  non disponibili ai processi ordinari.

\item[\texttt{user}] L'accesso agli \textit{extended user attributes} è
  regolato dai normali permessi dei file: occorre avere il permesso di lettura
  per leggerli e quello di scrittura per scriverli o modificarli. Dato l'uso
  di questi attributi si è scelto di applicare al loro accesso gli stessi
  criteri che si usano per l'accesso al contenuto dei file (o delle directory)
  cui essi fanno riferimento. Questa scelta vale però soltanto per i file e le
  directory ordinarie, se valesse in generale infatti si avrebbe un serio
  problema di sicurezza dato che esistono diversi oggetti sul filesystem per i
  quali è normale avere il permesso di scrittura consentito a tutti gli
  utenti, come i collegamenti simbolici, o alcuni file di dispositivo come
  \texttt{/dev/null}. Se fosse possibile usare su di essi gli \textit{extended
    user attributes} un utente qualunque potrebbe inserirvi dati a
  piacere.\footnote{la cosa è stata notata su XFS, dove questo comportamento
    permetteva, non essendovi limiti sullo spazio occupabile dagli
    \textit{Extended Attributes}, di bloccare il sistema riempiendo il disco.}

  La semantica del controllo di accesso indicata inoltre non avrebbe alcun
  senso al di fuori di file e directory: i permessi di lettura e scrittura per
  un file di dispositivo attengono alle capacità di accesso al dispositivo
  sottostante,\footnote{motivo per cui si può formattare un disco anche se
    \texttt{/dev} è su un filesystem in sola lettura.} mentre per i
  collegamenti simbolici questi vengono semplicemente ignorati: in nessuno dei
  due casi hanno a che fare con il contenuto del file, e nella discussione
  relativa all'uso degli \textit{extended user attributes} nessuno è mai stato
  capace di indicare una qualche forma sensata di utilizzo degli stessi per
  collegamenti simbolici o file di dispositivo, e neanche per le \textit{fifo}
  o i socket.

  Per questo motivo essi sono stati completamente disabilitati per tutto ciò
  che non sia un file regolare o una directory.\footnote{si può verificare la
    semantica adottata consultando il file \texttt{fs/xattr.c} dei sorgenti
    del kernel.} Inoltre per le directory è stata introdotta una ulteriore
  restrizione, dovuta di nuovo alla presenza ordinaria di permessi di
  scrittura completi su directory come \texttt{/tmp}. Per questo motivo, per
  evitare eventuali abusi, se una directory ha lo \textit{sticky bit} attivo
  sarà consentito scrivere i suoi \textit{extended user attributes} soltanto
  se si è proprietari della stessa, o si hanno i privilegi amministrativi
  della capacità \const{CAP\_FOWNER}.
\end{basedescript}

Le funzioni per la gestione degli attributi estesi, come altre funzioni di
gestione avanzate specifiche di Linux, non fanno parte della \acr{glibc}, e
sono fornite da una apposita libreria, \texttt{libattr}, che deve essere
installata a parte;\footnote{la versione corrente della libreria è
  \texttt{libattr1}.}  pertanto se un programma le utilizza si dovrà indicare
esplicitamente l'uso della suddetta libreria invocando il compilatore con
l'opzione \texttt{-lattr}.  

Per leggere gli attributi estesi sono disponibili tre diverse funzioni di
sistema, \funcd{getxattr}, \funcd{lgetxattr} e \funcd{fgetxattr}, che
consentono rispettivamente di richiedere gli attributi relativi a un file, a
un collegamento simbolico e ad un file descriptor; i rispettivi prototipi
sono:

\begin{funcproto}{
\fhead{sys/types.h}
\fhead{attr/xattr.h}
\fdecl{ssize\_t getxattr(const char *path, const char *name, void *value,
  size\_t size)}
\fdecl{ssize\_t lgetxattr(const char *path, const char *name, void *value,
  size\_t size)}
\fdecl{ssize\_t fgetxattr(int filedes, const char *name, void *value,
  size\_t size)}
\fdesc{Leggono il valore di un attributo esteso.} 
}

{Le funzioni ritornano un intero positivo che indica la dimensione
  dell'attributo richiesto in caso di successo e $-1$ per un errore, nel qual
  caso \var{errno} assumerà uno dei valori:
  \begin{errlist}
  \item[\errcode{ENOATTR}] l'attributo richiesto non esiste.
  \item[\errcode{ENOTSUP}] gli attributi estesi non sono supportati dal
    filesystem o sono disabilitati.
  \item[\errcode{ERANGE}] la dimensione \param{size} del buffer \param{value}
    non è sufficiente per contenere il risultato.
  \end{errlist}
  ed inoltre tutti gli errori delle analoghe della famiglia \func{stat} con lo
  stesso significato, ed in particolare \errcode{EPERM} se non si hanno i
  permessi di accesso all'attributo.}
\end{funcproto}

Le funzioni \func{getxattr} e \func{lgetxattr} prendono come primo argomento
un \textit{pathname} che indica il file di cui si vuole richiedere un
attributo, la sola differenza è che la seconda, se il \textit{pathname} indica
un collegamento simbolico, restituisce gli attributi di quest'ultimo e non
quelli del file a cui esso fa riferimento. La funzione \func{fgetxattr} prende
invece come primo argomento un numero di file descriptor, e richiede gli
attributi del file ad esso associato.

Tutte e tre le funzioni richiedono di specificare nell'argomento \param{name}
il nome dell'attributo di cui si vuole ottenere il valore. Il nome deve essere
indicato comprensivo di prefisso del \textit{namespace} cui appartiene (uno
dei valori di tab.~\ref{tab:extended_attribute_class}) nella forma
\texttt{namespace.attributename}, come stringa terminata da un carattere NUL.
Il valore dell'attributo richiesto verrà restituito nel buffer puntato
dall'argomento \param{value} per una dimensione massima di \param{size} byte
(gli attributi estesi possono essere costituiti arbitrariamente da dati
testuali o binari); se quest'ultima non è sufficiente si avrà un errore di
\errcode{ERANGE}.

Per evitare di dover indovinare la dimensione di un attributo per tentativi si
può eseguire una interrogazione utilizzando un valore nullo per \param{size};
in questo caso non verrà letto nessun dato, ma verrà restituito come valore di
ritorno della funzione chiamata la dimensione totale dell'attributo esteso
richiesto, che si potrà usare come stima per allocare un buffer di dimensioni
sufficienti.

Si tenga conto che questa è comunque una stima perché anche se le funzioni
restituiscono la dimensione esatta dell'attributo al momento in cui sono
eseguite, questa potrebbe essere modificata in qualunque momento da un
successivo accesso eseguito da un altro processo, pertanto si verifichi sempre
il valore di ritorno ed il codice di errore della funzione usata, senza dare
per scontato che essa abbia sempre successo.

Un secondo gruppo di funzioni sono quelle che consentono di impostare il
valore di un attributo esteso, le funzioni sono \funcd{setxattr},
\funcd{lsetxattr} e \funcd{fsetxattr} e consentono di operare rispettivamente
su un file, su un collegamento simbolico o utilizzando un file descriptor; i
rispettivi prototipi sono:

\begin{funcproto}{
\fhead{sys/types.h}
\fhead{attr/xattr.h}
\fdecl{int setxattr(const char *path, const char *name, const void *value,\\
  \phantom{int setxattr(}size\_t size, int flags)} 
\fdecl{int lsetxattr(const char *path, const char *name, const void *value,\\
  \phantom{int lsetxattr(}size\_t size, int flags)} 
\fdecl{int fsetxattr(int filedes, const char *name, const void *value,\\
  \phantom{int fsetxattr(}size\_t size, int flags)} 
\fdesc{Impostano il valore di un attributo esteso.} 
}

{Le funzioni ritornano un $0$ in caso di successo e $-1$ per un errore, nel qual
  caso \var{errno} assumerà uno dei valori:
  \begin{errlist}
  \item[\errcode{EEXIST}] si è usato il flag \const{XATTR\_CREATE} ma
    l'attributo esiste già.
  \item[\errcode{ENOATTR}] si è usato il flag \const{XATTR\_REPLACE} e
    l'attributo richiesto non esiste.
  \item[\errcode{ENOTSUP}] gli attributi estesi non sono supportati dal
    filesystem o sono disabilitati.
  \end{errlist}
  ed inoltre tutti gli errori delle analoghe della famiglia \func{stat} con lo
  stesso significato ed in particolare \errcode{EPERM} se non si hanno i
  permessi di accesso all'attributo.}
\end{funcproto}

Le tre funzioni prendono come primo argomento un valore adeguato al loro scopo
(un \textit{pathname} le prime due, un file descriptor la terza), usato in
maniera del tutto identica a quanto visto in precedenza per le analoghe che
leggono gli attributi estesi.

Il secondo argomento, \param{name}, deve indicare, anche in questo caso con
gli stessi criteri appena visti per le analoghe \func{getxattr},
\func{lgetxattr} e \func{fgetxattr}, il nome (completo di suffisso)
dell'attributo su cui si vuole operare.  Il valore che verrà assegnato
all'attributo dovrà essere preparato nel buffer puntato da \param{value}, e la
sua dimensione totale (in byte) dovrà essere indicata dall'argomento
\param{size}.

Infine l'argomento \param{flag} consente di controllare le modalità di
sovrascrittura dell'attributo esteso, esso può prendere due valori: con
\constd{XATTR\_REPLACE} si richiede che l'attributo esista, nel qual caso
verrà sovrascritto ed altrimenti si avrà errore; con \constd{XATTR\_CREATE} si
richiede che l'attributo non esista, nel qual caso verrà creato, altrimenti si
avrà errore ed il valore attuale non sarà modificato.  Utilizzando per
\param{flag} un valore nullo l'attributo verrà modificato se è già presente, o
creato se non c'è.

Le funzioni finora illustrate permettono di leggere o scrivere gli attributi
estesi presenti su un file, ma sarebbe altrettanto utile poter sapere quali
sono questi attributi; per questo sono disponibili le ulteriori tre funzioni
di sistema \funcd{listxattr}, \funcd{llistxattr} e \funcd{flistxattr} i cui
prototipi sono:

\begin{funcproto}{
\fhead{sys/types.h}
\fhead{attr/xattr.h}
\fdecl{ssize\_t listxattr(const char *path, char *list, size\_t size)} 
\fdecl{ssize\_t llistxattr(const char *path, char *list, size\_t size)} 
\fdecl{ssize\_t flistxattr(int filedes, char *list, size\_t size)} 
\fdesc{Leggono la lista degli attributi estesi di un file.} 
}

{Le funzioni ritornano un intero positivo che indica la dimensione della lista
  in caso di successo e $-1$ per un errore, nel qual caso \var{errno} assumerà
  uno dei valori:
  \begin{errlist}
  \item[\errcode{ENOTSUP}] gli attributi estesi non sono supportati dal
    filesystem o sono disabilitati.
  \item[\errcode{ERANGE}] la dimensione \param{size} del buffer \param{value}
    non è sufficiente per contenere il risultato.
  \end{errlist}
  ed inoltre tutti gli errori delle analoghe della famiglia \func{stat} con lo
  stesso significato ed in particolare \errcode{EPERM} se non si hanno i
  permessi di accesso all'attributo.}
\end{funcproto}

Come per le precedenti le tre funzioni leggono gli attributi rispettivamente
di un file, un collegamento simbolico o specificando un file descriptor, da
specificare con il loro primo argomento. Gli altri due argomenti, identici per
tutte e tre, indicano rispettivamente il puntatore \param{list} al buffer dove
deve essere letta la lista e la dimensione \param{size} di quest'ultimo.

La lista viene fornita come sequenza non ordinata dei nomi dei singoli
attributi estesi (sempre comprensivi del prefisso della loro classe) ciascuno
dei quali è terminato da un carattere nullo. I nomi sono inseriti nel buffer
uno di seguito all'altro. Il valore di ritorno della funzione indica la
dimensione totale della lista in byte.

Come per le funzioni di lettura dei singoli attributi se le dimensioni del
buffer non sono sufficienti si avrà un errore, ma è possibile ottenere dal
valore di ritorno della funzione una stima della dimensione totale della lista
usando per \param{size} un valore nullo. 

Infine per rimuovere semplicemente un attributo esteso, si ha a disposizione
un ultimo gruppo di funzioni di sistema: \funcd{removexattr},
\funcd{lremovexattr} e \funcd{fremovexattr}; i rispettivi prototipi sono:

\begin{funcproto}{
\fhead{sys/types.h}
\fhead{attr/xattr.h}
\fdecl{int removexattr(const char *path, const char *name)} 
\fdecl{int lremovexattr(const char *path, const char *name)} 
\fdecl{int fremovexattr(int filedes, const char *name)} 
\fdesc{Rimuovono un attributo esteso di un file.} 
}

{Le funzioni ritornano $0$ in caso di successo e $-1$ per un errore, nel qual
  caso \var{errno} assumerà uno dei valori:
  \begin{errlist}
  \item[\errcode{ENOATTR}] l'attributo richiesto non esiste.
  \item[\errcode{ENOTSUP}] gli attributi estesi non sono supportati dal
    filesystem o sono disabilitati.
  \end{errlist}
  ed inoltre tutti gli errori delle analoghe della famiglia \func{stat} con lo
  stesso significato ed in particolare \errcode{EPERM} se non si hanno i
  permessi di accesso all'attributo.}
\end{funcproto}

Le tre funzioni rimuovono un attributo esteso operando rispettivamente su di
un file, su un collegamento simbolico o un file descriptor, che vengono
specificati dal valore passato con il loro primo argomento.  L'attributo da
rimuovere deve essere anche in questo caso indicato con
l'argomento \param{name} secondo le stesse modalità già illustrate in
precedenza per le altre funzioni relative alla gestione degli attributi
estesi.

\itindend{Extended~Attributes}


\subsection{Le \textit{Access  Control List}}
\label{sec:file_ACL}

% la documentazione di sistema è nei pacchetti libacl1-dev e acl 
% vedi anche http://www.suse.de/~agruen/acl/linux-acls/online/

\itindbeg{Access~Control~List~(ACL)}

Il modello classico dei permessi di Unix, per quanto funzionale ed efficiente,
è comunque piuttosto limitato e per quanto possa aver coperto per lunghi anni
le esigenze più comuni con un meccanismo semplice e potente, non è in grado di
rispondere in maniera adeguata a situazioni che richiedono una gestione più
complessa dei permessi di accesso.\footnote{già un requisito come quello di
  dare accesso in scrittura ad alcune persone ed in sola lettura ad altre non
  si può soddisfare in maniera semplice.}

Per questo motivo erano state progressivamente introdotte nelle varie versioni
di Unix dei meccanismi di gestione dei permessi dei file più flessibili, nella
forma delle cosiddette \textit{Access Control List} (indicate usualmente con
la sigla ACL).  Nello sforzo di standardizzare queste funzionalità era stato
creato un gruppo di lavoro il cui scopo era estendere lo standard POSIX 1003
attraverso due nuovi insiemi di specifiche, la POSIX 1003.1e per l'interfaccia
di programmazione e la POSIX 1003.2c per i comandi di shell.

Gli obiettivi del gruppo di lavoro erano però forse troppo ambizioni, e nel
gennaio del 1998 i finanziamenti vennero ritirati senza che si fosse arrivati
alla definizione dello standard richiesto. Dato però che una parte della
documentazione prodotta era di alta qualità venne deciso di rilasciare al
pubblico la diciassettesima bozza del documento, quella che va sotto il nome
di \textit{POSIX 1003.1e Draft 17}, che è divenuta la base sulla quale si
definiscono le cosiddette \textit{Posix ACL}.

A differenza di altri sistemi, come ad esempio FreeBSD, nel caso di Linux si è
scelto di realizzare le ACL attraverso l'uso degli \textit{Extended
  Attributes} (appena trattati in sez.~\ref{sec:file_xattr}), e fornire tutte
le relative funzioni di gestione tramite una libreria, \texttt{libacl} che
nasconde i dettagli implementativi delle ACL e presenta ai programmi una
interfaccia che fa riferimento allo standard POSIX 1003.1e.

Anche in questo caso le funzioni di questa libreria non fanno parte della
\acr{glibc} e devono essere installate a parte;\footnote{la versione corrente
  della libreria è \texttt{libacl1}, e nel caso si usi Debian la si può
  installare con il pacchetto omonimo e con il collegato \texttt{libacl1-dev}
  per i file di sviluppo.}  pertanto se un programma le utilizza si dovrà
indicare esplicitamente l'uso della libreria \texttt{libacl} invocando il
compilatore con l'opzione \texttt{-lacl}. Si tenga presente inoltre che le ACL
devono essere attivate esplicitamente montando il filesystem\footnote{che deve
  supportarle, ma questo è ormai vero per praticamente tutti i filesystem più
  comuni, con l'eccezione di NFS per il quale esiste però il supporto per le
  versioni NFSv2 e NFSv3 del protocollo, con NFSv4 esistono invece delle ACL
  native che hanno una semantica diversa, su di esse possono mappare le ACL
  POSIX, ma l'inverso è possibile solo in forma incompleta.} su cui le si
vogliono utilizzare con l'opzione \texttt{acl} attiva. Dato che si tratta di
una estensione è infatti opportuno utilizzarle soltanto laddove siano
necessarie.

Una ACL è composta da un insieme di voci, e ciascuna voce è a sua volta
costituita da un \textsl{tipo}, da un eventuale \textsl{qualificatore} (deve
essere presente soltanto per le voci di tipo \const{ACL\_USER} e
\const{ACL\_GROUP}) e da un insieme di permessi.  Ad ogni oggetto sul
filesystem si può associare una ACL che ne governa i permessi di accesso,
detta \textit{access ACL}.  Inoltre per le directory si può impostare una ACL
aggiuntiva, detta ``\textit{Default ACL}'', che serve ad indicare quale dovrà
essere la ACL assegnata di default nella creazione di un file all'interno
della directory stessa. Come avviene per i permessi le ACL possono essere
impostate solo del proprietario del file, o da un processo con la capacità
\const{CAP\_FOWNER}.

\begin{table}[htb]
  \centering
  \footnotesize
  \begin{tabular}{|l|p{8cm}|}
    \hline
    \textbf{Tipo} & \textbf{Descrizione} \\
    \hline
    \hline
    \constd{ACL\_USER\_OBJ}& Voce che contiene i diritti di accesso del
                             proprietario del file.\\
    \constd{ACL\_USER}     & Voce che contiene i diritti di accesso per
                             l'utente indicato dal rispettivo
                             qualificatore.\\  
    \constd{ACL\_GROUP\_OBJ}&Voce che contiene i diritti di accesso del
                             gruppo proprietario del file.\\
    \constd{ACL\_GROUP}    & Voce che contiene i diritti di accesso per
                             il gruppo indicato dal rispettivo
                             qualificatore.\\
    \constd{ACL\_MASK}     & Voce che contiene la maschera dei massimi
                             permessi di accesso che possono essere garantiti
                             da voci del tipo \const{ACL\_USER},
                             \const{ACL\_GROUP} e \const{ACL\_GROUP\_OBJ}.\\
    \constd{ACL\_OTHER}    & Voce che contiene i diritti di accesso di chi
                             non corrisponde a nessuna altra voce dell'ACL.\\
    \hline
  \end{tabular}
  \caption{Le costanti che identificano i tipi delle voci di una ACL.}
  \label{tab:acl_tag_types}
\end{table}

L'elenco dei vari tipi di voci presenti in una ACL, con una breve descrizione
del relativo significato, è riportato in tab.~\ref{tab:acl_tag_types}. Tre di
questi tipi, \const{ACL\_USER\_OBJ}, \const{ACL\_GROUP\_OBJ} e
\const{ACL\_OTHER}, corrispondono direttamente ai tre permessi ordinari dei
file (proprietario, gruppo proprietario e tutti gli altri) e per questo una
ACL valida deve sempre contenere una ed una sola voce per ciascuno di questi
tipi.

Una ACL può poi contenere un numero arbitrario di voci di tipo
\const{ACL\_USER} e \const{ACL\_GROUP}, ciascuna delle quali indicherà i
permessi assegnati all'utente e al gruppo indicato dal relativo qualificatore.
Ovviamente ciascuna di queste voci dovrà fare riferimento a un utente o a un
gruppo diverso, e non corrispondenti a quelli proprietari del file. Inoltre se
in una ACL esiste una voce di uno di questi due tipi, è obbligatoria anche la
presenza di una ed una sola voce di tipo \const{ACL\_MASK}, che negli altri
casi è opzionale.

Una voce di tipo \const{ACL\_MASK} serve a mantenere la maschera dei permessi
che possono essere assegnati tramite voci di tipo \const{ACL\_USER},
\const{ACL\_GROUP} e \const{ACL\_GROUP\_OBJ}. Se in una di queste voci si
fosse specificato un permesso non presente in \const{ACL\_MASK} questo
verrebbe ignorato. L'uso di una ACL di tipo \const{ACL\_MASK} è di particolare
utilità quando essa associata ad una \textit{Default ACL} su una directory, in
quanto i permessi così specificati verranno ereditati da tutti i file creati
nella stessa directory. Si ottiene così una sorta di \textit{umask} associata
ad un oggetto sul filesystem piuttosto che a un processo.

Dato che le ACL vengono a costituire una estensione dei permessi ordinari, uno
dei problemi che si erano posti nella loro standardizzazione era appunto
quello della corrispondenza fra questi e le ACL. Come accennato i permessi
ordinari vengono mappati nelle tre voci di tipo \const{ACL\_USER\_OBJ},
\const{ACL\_GROUP\_OBJ} e \const{ACL\_OTHER} che devono essere presenti in
qualunque ACL; un cambiamento ad una di queste voci viene automaticamente
riflesso sui permessi ordinari dei file e viceversa.\footnote{per permessi
  ordinari si intende quelli mantenuti nell'\textit{inode}, che devono restare
  dato che un filesystem può essere montato senza abilitare le ACL.}

In realtà la mappatura è diretta solo per le voci \const{ACL\_USER\_OBJ} e
\const{ACL\_OTHER}, nel caso di \const{ACL\_GROUP\_OBJ} questo vale soltanto
se non è presente una voce di tipo \const{ACL\_MASK}, che è quanto avviene
normalmente se non sono presenti ACL aggiuntive rispetto ai permessi
ordinari. Se invece questa è presente verranno tolti dai permessi di
\const{ACL\_GROUP\_OBJ} (cioè dai permessi per il gruppo proprietario del
file) tutti quelli non presenti in \const{ACL\_MASK}.\footnote{questo diverso
  comportamento a seconda delle condizioni è stato introdotto dalla
  standardizzazione \textit{POSIX 1003.1e Draft 17} per mantenere il
  comportamento invariato sui sistemi dotati di ACL per tutte quelle
  applicazioni che sono conformi soltanto all'ordinario standard \textit{POSIX
    1003.1}.}

Un secondo aspetto dell'incidenza delle ACL sul comportamento del sistema è
quello relativo alla creazione di nuovi file,\footnote{o oggetti sul
  filesystem, il comportamento discusso vale per le funzioni \func{open} e
  \func{creat} (vedi sez.~\ref{sec:file_open_close}), \func{mkdir} (vedi
  sez.~\ref{sec:file_dir_creat_rem}), \func{mknod} e \func{mkfifo} (vedi
  sez.~\ref{sec:file_mknod}).} che come accennato può essere modificato dalla
presenza di una \textit{Default ACL} sulla directory che andrà a contenerli.
Se questa non c'è valgono le regole usuali illustrate in
sez.~\ref{sec:file_perm_management}, per cui essi sono determinati dalla
\textit{umask} del processo, e la sola differenza è che i permessi ordinari da
esse risultanti vengono automaticamente rimappati anche su una ACL di accesso
assegnata automaticamente al nuovo file, che contiene soltanto le tre
corrispondenti voci di tipo \const{ACL\_USER\_OBJ}, \const{ACL\_GROUP\_OBJ} e
\const{ACL\_OTHER}.

Se invece è presente una ACL di default sulla directory che contiene il nuovo
file, essa diventerà automaticamente anche la ACL di accesso di quest'ultimo,
a meno di non aver indicato, nelle funzioni di creazione che lo consentono,
uno specifico valore per i permessi ordinari.\footnote{tutte le funzioni
  citate in precedenza supportano un argomento \var{mode} che indichi un
  insieme di permessi iniziale.} In tal caso saranno eliminati dalle voci
corrispondenti che deriverebbero dalla ACL di default, tutti i permessi non
presenti in tale indicazione.

Dato che questa è la ragione che ha portato alla loro creazione, la principale
modifica introdotta nel sistema con la presenza della ACL è quella alle regole
del controllo di accesso ai file che si sono illustrate in
sez.~\ref{sec:file_perm_overview}.  Come nel caso ordinario per il controllo
vengono sempre utilizzati gli identificatori del gruppo \textit{effective} del
processo, ma in caso di presenza di una ACL sul file, i passi attraverso i
quali viene stabilito se il processo ha il diritto di accesso sono i seguenti:
\begin{enumerate*}
\item Se l'\ids{UID} del processo è nullo (se cioè si è l'amministratore)
  l'accesso è sempre garantito senza nessun controllo.\footnote{più
    precisamente se si devono avere le capacità \const{CAP\_DAC\_OVERRIDE} per
    i file e \const{CAP\_DAC\_READ\_SEARCH} per le directory, vedi
    sez.~\ref{sec:proc_capabilities}.}
\item Se l'\ids{UID} del processo corrisponde al proprietario del file allora:
  \begin{itemize*}
  \item se la voce \const{ACL\_USER\_OBJ} contiene il permesso richiesto,
    l'accesso è consentito;
  \item altrimenti l'accesso è negato.
  \end{itemize*}
\item Se l'\ids{UID} del processo corrisponde ad un qualunque qualificatore
  presente in una voce \const{ACL\_USER} allora:
  \begin{itemize*}
  \item se la voce \const{ACL\_USER} corrispondente e la voce
    \const{ACL\_MASK} contengono entrambe il permesso richiesto, l'accesso è
    consentito;
  \item altrimenti l'accesso è negato.
  \end{itemize*}
\item Se è il \ids{GID} del processo o uno dei \ids{GID} supplementari
  corrisponde al gruppo proprietario del file allora: 
  \begin{itemize*}
  \item se la voce \const{ACL\_GROUP\_OBJ} e una eventuale voce
    \const{ACL\_MASK} (se non vi sono voci di tipo \const{ACL\_GROUP} questa
    può non essere presente) contengono entrambe il permesso richiesto,
    l'accesso è consentito;
  \item altrimenti l'accesso è negato.
  \end{itemize*}
\item Se è il \ids{GID} del processo o uno dei \ids{GID} supplementari
  corrisponde ad un qualunque qualificatore presente in una voce
  \const{ACL\_GROUP} allora:
  \begin{itemize*}
  \item se la voce \const{ACL\_GROUP} corrispondente e la voce
    \const{ACL\_MASK} contengono entrambe il permesso richiesto, l'accesso è
    consentito;
  \item altrimenti l'accesso è negato.
  \end{itemize*}
\item Se la voce \const{ACL\_USER\_OBJ} contiene il permesso richiesto,
  l'accesso è consentito, altrimenti l'accesso è negato.
\end{enumerate*}

I passi di controllo vengono eseguiti esattamente in questa sequenza, e la
decisione viene presa non appena viene trovata una corrispondenza con gli
identificatori del processo. Questo significa che i permessi presenti in una
voce di tipo \const{ACL\_USER} hanno la precedenza sui permessi ordinari
associati al gruppo proprietario del file (vale a dire su
\const{ACL\_GROUP\_OBJ}).

Per la gestione delle ACL lo standard \textit{POSIX 1003.1e Draft 17} ha
previsto delle apposite funzioni ed tutta una serie di tipi di dati dedicati,
arrivando fino a definire un tipo di dato e delle costanti apposite per
identificare i permessi standard di lettura, scrittura ed esecuzione.  Tutte
le operazioni devono essere effettuate attraverso tramite questi tipi di dati,
che incapsulano tutte le informazioni contenute nelle ACL.

La prima di queste funzioni che prendiamo in esame (si ricordi che come per
tutte le altre per poterla usare occorre invocare il compilatore con l'opzione
\texttt{-l acl}) è \funcd{acl\_init}, il cui prototipo è:

\begin{funcproto}{
\fhead{sys/types.h}
\fhead{sys/acl.h}
\fdecl{acl\_t acl\_init(int count)}
\fdesc{Inizializza un'area di lavoro per una ACL.} 
}

{La funzione ritorna un oggetto di tipo \type{acl\_t} in caso di successo e
  \val{NULL} per un errore, nel qual caso \var{errno} assumerà uno dei valori:
  \begin{errlist}
  \item[\errcode{EINVAL}] il valore di \param{count} è negativo.
  \item[\errcode{ENOMEM}] non c'è sufficiente memoria disponibile.
  \end{errlist}
}
\end{funcproto}

La funzione alloca ed inizializza un'area di memoria che verrà usata per
mantenere i dati di una ACL contenente fino ad un massimo di \param{count}
voci. La funzione ritorna un valore di tipo \typed{acl\_t} da usare in tutte le
altre funzioni che operano sulla ACL. La funzione si limita alla allocazione
iniziale e non inserisce nessun valore nella ACL che resta vuota. 

Si tenga presente che pur essendo \typed{acl\_t} un tipo opaco che identifica
``\textsl{l'oggetto}'' ACL, il valore restituito dalla funzione non è altro
che un puntatore all'area di memoria allocata per i dati richiesti. Pertanto
in caso di fallimento verrà restituito un puntatore nullo di tipo
``\code{(acl\_t) NULL}'' e si dovrà, in questa come in tutte le funzioni
seguenti che restituiscono un oggetto di tipo \type{acl\_t}, confrontare il
valore di ritorno della funzione con \val{NULL} (anche se, a voler essere
estremamente pignoli, si dovrebbe usare ``\code{(acl\_t) NULL}'', ma è
sufficiente fare un confronto direttamente con \val{NULL} essendo cura del
compilatore fare le conversioni necessarie).

Una volta che si siano completate le operazioni sui dati di una ACL la memoria
allocata per un oggetto \type{acl\_t} dovrà essere liberata esplicitamente
attraverso una chiamata alla funzione \funcd{acl\_free}, il cui prototipo è:

\begin{funcproto}{
\fhead{sys/types.h}
\fhead{sys/acl.h}
\fdecl{int acl\_free(void *obj\_p)}
\fdesc{Disalloca la memoria riservata per una ACL.} 
}

{La funzione ritorna $0$ in caso di successo e $-1$ per un errore, nel qual
  caso \var{errno} può assumere solo il valore:
  \begin{errlist}
  \item[\errcode{EINVAL}] \param{obj\_p} non è valido.
  \end{errlist}
}
\end{funcproto}

Si noti come la funzione usi come argomento un puntatore di tipo ``\ctyp{void
  *}'', essa infatti può essere usata non solo per liberare la memoria
allocata per i dati di una ACL, ma anche per quella usata per creare le
stringhe di descrizione testuale delle ACL o per ottenere i valori dei
qualificatori della una voce di una ACL. L'uso del tipo generico ``\ctyp{void
  *}'' consente di evitare di eseguire un \textit{cast} al tipo di dato di cui
si vuole effettuare la disallocazione.

Si tenga presente poi che oltre a \func{acl\_init} ci sono molte altre
funzioni che possono allocare memoria per i dati delle ACL, è pertanto
opportuno tenere traccia di tutte le chiamate a queste funzioni perché alla
fine delle operazioni tutta la memoria allocata dovrà essere liberata con
\func{acl\_free}.

Una volta che si abbiano a disposizione i dati di una ACL tramite il
riferimento ad oggetto di tipo \type{acl\_t} questi potranno essere copiati
con la funzione \funcd{acl\_dup}, il cui prototipo è:

\begin{funcproto}{
\fhead{sys/types.h}
\fhead{sys/acl.h}
\fdecl{acl\_t acl\_dup(acl\_t acl)}
\fdesc{Crea una copia di una ACL.} 
}

{La funzione ritorna un oggetto di tipo \type{acl\_t} in caso di successo in
  caso di successo e \val{NULL} per un errore, nel qual caso \var{errno}
  assumerà uno dei valori:
  \begin{errlist}
  \item[\errcode{EINVAL}] l'argomento \param{acl} non è un puntatore valido
    per una ACL.
  \item[\errcode{ENOMEM}] non c'è sufficiente memoria disponibile per eseguire
    la copia.
  \end{errlist}
}
\end{funcproto}

La funzione crea una copia dei dati della ACL indicata tramite l'argomento
\param{acl}, allocando autonomamente tutto spazio necessario alla copia e
restituendo un secondo oggetto di tipo \type{acl\_t} come riferimento a
quest'ultima.  Valgono per questo le stesse considerazioni fatte per il valore
di ritorno di \func{acl\_init}, ed in particolare il fatto che occorrerà
prevedere una ulteriore chiamata esplicita a \func{acl\_free} per liberare la
memoria occupata dalla copia.

Se si deve creare una ACL manualmente l'uso di \func{acl\_init} è scomodo,
dato che la funzione restituisce una ACL vuota, una alternativa allora è usare
\funcd{acl\_from\_mode} che consente di creare una ACL a partire da un valore
di permessi ordinari, il prototipo della funzione è:

\begin{funcproto}{
\fhead{sys/types.h}
\fhead{sys/acl.h}
\fdecl{acl\_t acl\_from\_mode(mode\_t mode)}
\fdesc{Crea una ACL inizializzata con i permessi ordinari.} 
}

{La funzione ritorna un oggetto di tipo \type{acl\_t} in caso di successo e
  \val{NULL} per un errore, nel qual caso \var{errno} può assumere solo
  il valore \errval{ENOMEM}.}
\end{funcproto}


La funzione restituisce una ACL inizializzata con le tre voci obbligatorie
\const{ACL\_USER\_OBJ}, \const{ACL\_GROUP\_OBJ} e \const{ACL\_OTHER} già
impostate secondo la corrispondenza ai valori dei permessi ordinari indicati
dalla maschera passata nell'argomento \param{mode}. Questa funzione è una
estensione usata dalle ACL di Linux e non è portabile, ma consente di
semplificare l'inizializzazione in maniera molto comoda. 

Altre due funzioni che consentono di creare una ACL già inizializzata sono
\funcd{acl\_get\_fd} e \funcd{acl\_get\_file}, che consentono di leggere la
ACL di un file; i rispettivi prototipi sono:

\begin{funcproto}{
\fhead{sys/types.h}
\fhead{sys/acl.h}
\fdecl{acl\_t acl\_get\_file(const char *path\_p, acl\_type\_t type)}
\fdecl{acl\_t acl\_get\_fd(int fd)}
\fdesc{Leggono i dati delle ACL di un file.} 
}

{Le funzioni ritornano un oggetto di tipo \type{acl\_t} in caso di successo e
  \val{NULL} per un errore, nel qual caso \var{errno} assumerà uno dei valori:
  \begin{errlist}
  \item[\errcode{EACCESS}] non c'è accesso per una componente di
    \param{path\_p} o si è richiesta una ACL di default per un file (solo per
    \func{acl\_get\_file}).
  \item[\errcode{EINVAL}] \param{type} non ha un valore valido (solo per
    \func{acl\_get\_file}).
  \item[\errcode{ENOTSUP}] il filesystem cui fa riferimento il file non
    supporta le ACL.
  \end{errlist}
  ed inoltre \errval{ENOMEM} per entrambe, \errval{EBADF} per
  \func{acl\_get\_fd}, e \errval{ENAMETOOLONG}, \errval{ENOENT},
  \errval{ENOTDIR}, per \func{acl\_get\_file} nel loro significato generico. }
\end{funcproto}

Le due funzioni ritornano, con un oggetto di tipo \type{acl\_t}, il valore
della ACL correntemente associata ad un file, che può essere identificato
tramite un file descriptor usando \func{acl\_get\_fd} o con un
\textit{pathname} usando \func{acl\_get\_file}.

Nel caso di quest'ultima funzione, che può richiedere anche la ACL relativa ad
una directory, il secondo argomento \param{type} consente di specificare se si
vuole ottenere la ACL di default o quella di accesso. Questo argomento deve
essere di tipo \typed{acl\_type\_t} e può assumere solo i due valori riportati
in tab.~\ref{tab:acl_type}.

\begin{table}[htb]
  \centering
  \footnotesize
  \begin{tabular}{|l|l|}
    \hline
    \textbf{Tipo} & \textbf{Descrizione} \\
    \hline
    \hline
    \constd{ACL\_TYPE\_ACCESS} & Indica una ACL di accesso.\\
    \constd{ACL\_TYPE\_DEFAULT}& Indica una ACL di default.\\  
    \hline
  \end{tabular}
  \caption{Le costanti che identificano il tipo di ACL.}
  \label{tab:acl_type}
\end{table}

Si tenga presente che nel caso di \func{acl\_get\_file} occorrerà che il
processo chiamante abbia privilegi di accesso sufficienti a poter leggere gli
attributi estesi dei file (come illustrati in sez.~\ref{sec:file_xattr});
inoltre una ACL di tipo \const{ACL\_TYPE\_DEFAULT} potrà essere richiesta
soltanto per una directory, e verrà restituita solo se presente, altrimenti
verrà restituita una ACL vuota.

Infine si potrà creare una ACL direttamente dalla sua rappresentazione
testuale con la funzione \funcd{acl\_from\_text}, il cui prototipo è:

\begin{funcproto}{
\fhead{sys/types.h}
\fhead{sys/acl.h}
\fdecl{acl\_t acl\_from\_text(const char *buf\_p)}
\fdesc{Crea una ACL a partire dalla sua rappresentazione testuale.} 
}

{La funzione ritorna un oggetto di tipo \type{acl\_t} in caso di successo e
  \val{NULL} per un errore, nel qual caso \var{errno} assumerà uno
  dei valori:
  \begin{errlist}
  \item[\errcode{EINVAL}] la rappresentazione testuale all'indirizzo
    \param{buf\_p} non è valida.
  \item[\errcode{ENOMEM}] non c'è memoria sufficiente per allocare i dati.
   \end{errlist}
}
\end{funcproto}

La funzione prende come argomento il puntatore ad un buffer dove si è inserita
la rappresentazione testuale della ACL che si vuole creare, la memoria
necessaria viene automaticamente allocata ed in caso di successo viene
restituito come valore di ritorno un oggetto di tipo \type{acl\_t} con il
contenuto della stessa, che, come per le precedenti funzioni, dovrà essere
disallocato esplicitamente al termine del suo utilizzo.

La rappresentazione testuale di una ACL è quella usata anche dai comandi
ordinari per la gestione delle ACL (\texttt{getfacl} e \texttt{setfacl}), che
prevede due diverse forme, estesa e breve, entrambe supportate da
\func{acl\_from\_text}. La forma estesa prevede che sia specificata una voce
per riga, nella forma:
\begin{Example}
tipo:qualificatore:permessi
\end{Example}
dove il tipo può essere uno fra \texttt{user}, \texttt{group}, \texttt{other}
e \texttt{mask}. Il qualificatore è presente solo per \texttt{user} e
\texttt{group} ed indica l'utente o il gruppo a cui la voce si riferisce,
mentre i permessi sono espressi con una tripletta di lettere analoga a quella
usata per i permessi dei file, vale a dire ``\texttt{r}'' per il permesso di
lettura, ``\texttt{w}'' per il permesso di scrittura, ``\texttt{x}'' per il
permesso di esecuzione (scritti in quest'ordine) e ``\texttt{-}'' per
l'assenza del permesso.

Un possibile esempio di rappresentazione della ACL di un file ordinario a cui,
oltre ai permessi ordinari, si è aggiunto un altro utente con un accesso in
lettura, è il seguente:
\begin{Example}
user::rw-
group::r--
other::r--
user:piccardi:r--
group:gapil:r--
\end{Example}

Va precisato che i due tipi \texttt{user} e \texttt{group} sono usati
rispettivamente per indicare delle voci relative ad utenti e gruppi (cioè per
voci di tipo \const{ACL\_USER\_OBJ} e \const{ACL\_USER} per \texttt{user} e
\const{ACL\_GROUP\_OBJ} e \const{ACL\_GROUP} per \texttt{group}) applicate sia
a quelli proprietari del file che a quelli generici. Quelle dei proprietari si
riconoscono per l'assenza di un qualificatore, ed in genere si scrivono per
prima delle altre.  Il significato delle voci di tipo \texttt{mask} e
\texttt{mark} è evidente. Usando questa forma estesa si possono anche inserire
dei commenti nel testo precedendoli con il carattere ``\texttt{\#}''.

La forma breve prevede invece la scrittura delle singole voci su una riga,
separate da virgole; come specificatori del tipo di voce si possono usare le
iniziali dei valori usati nella forma estesa (cioè ``\texttt{u}'',
``\texttt{g}'', ``\texttt{o}'' e ``\texttt{m}''), mentre le altri parte della
voce sono le stesse. In questo caso non sono consentiti permessi.

Per la conversione inversa, che consente di ottenere la rappresentazione
testuale di una ACL, sono invece disponibili due funzioni. La prima delle due,
di uso più immediato, è \funcd{acl\_to\_text}, ed il suo prototipo è:

\begin{funcproto}{
\fhead{sys/types.h}
\fhead{sys/acl.h}
\fdecl{char *acl\_to\_text(acl\_t acl, ssize\_t *len\_p)}
\fdesc{Produce la rappresentazione testuale di una ACL.} 
}

{La funzione ritorna il puntatore ad una stringa con la rappresentazione
  testuale della ACL in caso di successo e \var{NULL} per un errore, nel qual
  caso \var{errno} assumerà uno dei valori:
  \begin{errlist}
  \item[\errcode{EINVAL}] o \param{acl} non è un puntatore ad una ACL o la ACL
    che esso indica non è valida o non può esser tradotta in forma testuale.
  \item[\errcode{ENOMEM}] non c'è memoria sufficiente per allocare i dati.
  \end{errlist}
}  
\end{funcproto}

La funzione restituisce il puntatore ad una stringa, terminata da un NUL,
contenente la rappresentazione testuale in forma estesa della ACL passata come
argomento, ed alloca automaticamente la memoria necessaria. Questa dovrà poi
essere liberata, quando non più necessaria, con \func{acl\_free}.  Se
nell'argomento \param{len\_p} si passa come valore il puntatore ad una
variabile intera, in questa verrà restituita (come \textit{value result
  argument}) la dimensione della stringa con la rappresentazione testuale, non
comprendente il carattere nullo finale.

La seconda funzione, che permette di controllare con una gran dovizia di
particolari la generazione della stringa contenente la rappresentazione
testuale della ACL, è \funcd{acl\_to\_any\_text}, ed il suo prototipo è:

\begin{funcproto}{
\fhead{sys/types.h}
\fhead{sys/acl.h}
\fdecl{char *acl\_to\_any\_text(acl\_t acl, const char *prefix, char
    separator, int options)}
\fdesc{Produce la rappresentazione testuale di una ACL.} 
}

{La funzione ritorna il puntatore ad una stringa con la rappresentazione
  testuale della ACL in caso di successo e \val{NULL} per un errore, nel qual
  caso \var{errno} assumerà uno dei valori:
  \begin{errlist}
  \item[\errcode{EINVAL}] la ACL indicata da \param{acl} non è valida.
  \item[\errcode{ENOMEM}] non c'è memoria sufficiente per allocare i dati.
  \end{errlist}
}  
\end{funcproto}

La funzione converte in formato testo la ACL indicata dall'argomento
\param{acl}, usando il carattere \param{separator} come separatore delle
singole voci; se l'argomento \param{prefix} non è nullo la stringa da esso
indicata viene utilizzata come prefisso per le singole voci. 

L'ultimo argomento, \param{options}, consente di controllare la modalità con
cui viene generata la rappresentazione testuale. Un valore nullo fa sì che
vengano usati gli identificatori standard \texttt{user}, \texttt{group},
\texttt{other} e \texttt{mask} con i nomi di utenti e gruppi risolti rispetto
ai loro valori numerici. Altrimenti si può specificare un valore in forma di
maschera binaria, da ottenere con un OR aritmetico dei valori riportati in
tab.~\ref{tab:acl_to_text_options}.

\begin{table}[htb]
  \centering
  \footnotesize
  \begin{tabular}{|l|p{8cm}|}
    \hline
    \textbf{Tipo} & \textbf{Descrizione} \\
    \hline
    \hline
    \constd{TEXT\_ABBREVIATE}    & Stampa le voci in forma abbreviata.\\
    \constd{TEXT\_NUMERIC\_IDS}  & Non effettua la risoluzione di
                                   \ids{UID} e \ids{GID} lasciando i valori
                                   numerici.\\
    \constd{TEXT\_SOME\_EFFECTIVE}&Per ciascuna voce che contiene permessi che
                                   vengono eliminati dalla \const{ACL\_MASK}
                                   viene generato un commento con i permessi 
                                   effettivamente risultanti; il commento è
                                   separato con un tabulatore.\\
    \constd{TEXT\_ALL\_EFFECTIVE}& Viene generato un commento con i permessi
                                   effettivi per ciascuna voce che contiene
                                   permessi citati nella \const{ACL\_MASK},
                                   anche quando questi non vengono modificati
                                   da essa; il commento è separato con un
                                   tabulatore.\\
    \constd{TEXT\_SMART\_INDENT} & Da usare in combinazione con le precedenti
                                   opzioni \const{TEXT\_SOME\_EFFECTIVE} e
                                   \const{TEXT\_ALL\_EFFECTIVE}, aumenta
                                   automaticamente il numero di spaziatori
                                   prima degli eventuali commenti in modo da
                                   mantenerli allineati.\\
    \hline
  \end{tabular}
  \caption{Possibili valori per l'argomento \param{options} di
    \func{acl\_to\_any\_text}.} 
  \label{tab:acl_to_text_options}
\end{table}

Come per \func{acl\_to\_text} anche in questo caso il buffer contenente la
rappresentazione testuale dell'ACL, di cui la funzione restituisce
l'indirizzo, viene allocato automaticamente, e dovrà essere esplicitamente
disallocato con una chiamata ad \func{acl\_free}. Si tenga presente infine che
questa funzione è una estensione specifica di Linux, e non è presente nella
bozza dello standard POSIX.1e.

Per quanto utile per la visualizzazione o l'impostazione da riga di comando
delle ACL, la forma testuale non è la più efficiente per poter memorizzare i
dati relativi ad una ACL, ad esempio quando si vuole eseguirne una copia a
scopo di archiviazione. Per questo è stata prevista la possibilità di
utilizzare una rappresentazione delle ACL in una apposita forma binaria
contigua e persistente. È così possibile copiare il valore di una ACL in un
buffer e da questa rappresentazione tornare indietro e generare una ACL.

Lo standard POSIX.1e prevede a tale scopo tre funzioni, la prima e più
semplice è \funcd{acl\_size}, che consente di ottenere la dimensione che avrà
la citata rappresentazione binaria, in modo da poter allocare per essa un
buffer di dimensione sufficiente, il suo prototipo è:

\begin{funcproto}{
\fhead{sys/types.h}
\fhead{sys/acl.h}
\fdecl{ssize\_t acl\_size(acl\_t acl)}
\fdesc{Determina la dimensione della rappresentazione binaria di una ACL.} 
}

{La funzione ritorna la dimensione in byte della rappresentazione binaria
  della ACL in caso di successo e $-1$ per un errore, nel qual caso
  \var{errno} può assumere solo il valore:
  \begin{errlist}
  \item[\errcode{EINVAL}] la ACL indicata da \param{acl} non è valida.
  \end{errlist}
}  
\end{funcproto}

Ottenuta con \func{acl\_size} la dimensione del buffer necessaria per potervi
memorizzare una ACL questo dovrà potrà essere allocato direttamente con
\func{malloc}, ed a questo punto vi si potrà salvare la rappresentazione
binaria della precedente ACL utilizzando la funzione \funcd{acl\_copy\_ext},
il cui prototipo è:

\begin{funcproto}{
\fhead{sys/types.h}
\fhead{sys/acl.h}
\fdecl{ssize\_t acl\_copy\_ext(void *buf\_p, acl\_t acl, ssize\_t size)}
\fdesc{Ottiene la rappresentazione binaria di una ACL.} 
}

{La funzione ritorna la dimensione in byte della rappresentazione binaria
  della ACL in caso di successo e $-1$ per un errore, nel qual caso
  \var{errno} assumerà uno dei valori:
  \begin{errlist}
  \item[\errcode{EINVAL}] la ACL indicata da \param{acl} non è valida, o
    \param{acl} non è un puntatore ad una ACL o \param{size} è negativo o
    nullo.
  \item[\errcode{ERANGE}] il valore di \param{size} è più piccolo della
    dimensione della rappresentazione binaria della ACL.
  \end{errlist}
}  
\end{funcproto}

La funzione scriverà la rappresentazione binaria della ACL indicata da
\param{acl} sul buffer di dimensione \param{size} all'indirizzo
\param{buf\_p}, restituendo la dimensione della stessa come valore di
ritorno. Qualora la dimensione della rappresentazione ecceda il valore di
\param{size} la funzione fallirà con un errore di \errcode{ERANGE}. La
funzione non ha nessun effetto sulla ACL indicata da \param{acl}.

Viceversa se si vuole ripristinare una ACL a partire da una rappresentazione
binaria si potrà usare la funzione \funcd{acl\_copy\_int}, il cui prototipo è:

\begin{funcproto}{
\fhead{sys/types.h} 
\fhead{sys/acl.h}
\fdecl{acl\_t acl\_copy\_int(const void *buf\_p)}
\fdesc{Ripristina la rappresentazione binaria di una ACL.} 
}

{La funzione ritorna un oggetto di tipo \type{acl\_t} in caso di successo e
  \val{NULL} per un errore, nel qual caso \var{errno} assumerà uno dei valori:
  \begin{errlist}
  \item[\errcode{EINVAL}] il buffer all'indirizzo \param{buf\_p} non contiene
    una rappresentazione corretta di una ACL.
  \item[\errcode{ENOMEM}] non c'è memoria sufficiente per allocare un oggetto
    \type{acl\_t} per la ACL richiesta.
  \end{errlist}
}
\end{funcproto}

La funzione alloca autonomamente un oggetto di tipo \type{acl\_t}, restituito
come valore di ritorno, con il contenuto della ACL rappresentata dai dati del
buffer puntato da \param{buf\_p}. Al solito l'oggetto \type{acl\_t} dovrà
essere disallocato esplicitamente al termine del suo utilizzo.

Una volta che si disponga della ACL desiderata, questa potrà essere impostata
su un file o una directory. Per impostare una ACL sono disponibili due
funzioni: \funcd{acl\_set\_file}, che opera sia su file che su directory
usando un \textit{pathname}, e \funcd{acl\_set\_file} che opera solo su file
usando un file descriptor; i rispettivi prototipi sono:

\begin{funcproto}{
\fhead{sys/types.h}
\fhead{sys/acl.h}
\fdecl{int acl\_set\_file(const char *path, acl\_type\_t type, acl\_t acl)}
\fdesc{Imposta una ACL su un file o una directory.} 
\fdecl{int acl\_set\_fd(int fd, acl\_t acl)}
\fdesc{Imposta una ACL su un file descriptor.} 
}

{Le funzioni ritornano $0$ in caso di successo e $-1$ per un errore, nel qual
  caso \var{errno} assumerà uno dei valori: 
  \begin{errlist}
  \item[\errcode{EACCES}] o un generico errore di accesso a \param{path} o il
    valore di \param{type} specifica una ACL il cui tipo non può essere
    assegnato a \param{path}.
  \item[\errcode{EINVAL}] o \param{acl} non è una ACL valida, o \param{type}
    ha un valore non corretto per \func{acl\_set\_file} o o ha più voci di
    quante se ne possono assegnare al file per \func{acl\_set\_fd}.
  \item[\errcode{ENOSPC}] non c'è spazio disco sufficiente per contenere i
    dati aggiuntivi della ACL.
  \item[\errcode{ENOTSUP}] si è cercato di impostare una ACL su un file
    contenuto in un filesystem che non supporta le ACL.
  \end{errlist}
  ed inoltre nel loro significato generico \errval{EPERM}, \errval{EROFS} per
  entrambe, \errval{EBADF} per \func{acl\_set\_fd}, \errval{ENAMETOOLONG},
  \errval{ENOENT}, \errval{ENOTDIR} per \func{acl\_set\_file}.}
\end{funcproto}

Con \func{acl\_set\_file} si assegna la ACL contenuta in \param{acl} al file o
alla directory indicate da \param{path}, con \param{type} che indica il tipo
di ACL con le costanti di tab.~\ref{tab:acl_type}; si tenga presente però che
le ACL di default possono essere solo impostate qualora \param{path} indichi
una directory. Inoltre perché la funzione abbia successo la ACL dovrà essere
valida, e contenere tutti le voci necessarie, con l'eccezione si specifica una
ACL vuota per cancellare la ACL di default associata a \param{path}, valida
solo in caso di directory.\footnote{questo però è una estensione della
  implementazione delle ACL di Linux, la bozza di standard POSIX.1e prevedeva
  l'uso della apposita funzione \funcd{acl\_delete\_def\_file}, che prende
  come unico argomento il \textit{pathname} della directory di cui si vuole
  cancellare l'ACL di default, per i dettagli si ricorra alla pagina di
  manuale.}

La seconda funzione, \func{acl\_set\_fd}, è del tutto è analoga alla prima, ma
non dovendo avere a che fare con directory (e la conseguente possibilità di
avere una ACL di default) non necessita che si specifichi il tipo di ACL, che
sarà sempre di accesso, e prende come unico argomento, a parte il file
descriptor, la ACL da impostare.

Le funzioni viste finora operano a livello di una intera ACL, eseguendo in una
sola volta tutte le operazioni relative a tutte le voci in essa contenuta. In
generale è possibile modificare un singolo valore all'interno di una singola
voce direttamente con le funzioni previste dallo standard POSIX.1e.  Queste
funzioni però sono alquanto macchinose da utilizzare per cui è molto più
semplice operare direttamente sulla rappresentazione testuale. Questo è il
motivo per non tratteremo nei dettagli dette funzioni, fornendone solo una
descrizione sommaria; chi fosse interessato può ricorrere alle pagine di
manuale.

Se si vuole operare direttamente sui contenuti di un oggetto di tipo
\type{acl\_t} infatti occorre fare riferimento alle singole voci tramite gli
opportuni puntatori di tipo \typed{acl\_entry\_t}, che possono essere ottenuti
dalla funzione \funcm{acl\_get\_entry} (per una voce esistente) o dalla
funzione \funcm{acl\_create\_entry} per una voce da aggiungere. Nel caso della
prima funzione si potrà poi ripetere la lettura per ottenere i puntatori alle
singole voci successive alla prima.

\begin{figure}[!htb]
  \footnotesize \centering
  \begin{minipage}[c]{\codesamplewidth}
    \includecodesample{listati/mygetfacl.c}
  \end{minipage} 
  \normalsize
  \caption{Corpo principale del programma \texttt{mygetfacl.c}.}
  \label{fig:proc_mygetfacl}
\end{figure}

Una volta ottenuti detti puntatori si potrà operare sui contenuti delle
singole voci: con le funzioni \funcm{acl\_get\_tag\_type},
\funcm{acl\_get\_qualifier}, \funcm{acl\_get\_permset} si potranno leggere
rispettivamente tipo, qualificatore e permessi, mentre con le corrispondenti
\funcm{acl\_set\_tag\_type}, \funcm{acl\_set\_qualifier},
\funcm{acl\_set\_permset} si potranno impostare i valori; in entrambi i casi
vengono utilizzati tipi di dato ad hoc, descritti nelle pagine di manuale. Si
possono poi copiare i valori di una voce da una ACL ad un altra con
\funcm{acl\_copy\_entry} o eliminare una voce da una ACL con
\funcm{acl\_delete\_entry} e verificarne la validità prima di usarla con
\funcm{acl\_valid} o \funcm{acl\_check}.

\itindend{Access~Control~List~(ACL)}

Come esempio di utilizzo di queste funzioni nei sorgenti allegati alla guida
si è distribuito il programma \texttt{mygetfacl.c}, che consente di leggere le
ACL di un file, passato come argomento.

La sezione principale del programma, da cui si è rimossa la sezione sulla
gestione delle opzioni, è riportata in fig.~\ref{fig:proc_mygetfacl}. Il
programma richiede un unico argomento (\texttt{\small 16-20}) che indica il
file di cui si vuole leggere la ACL. Se questo è presente si usa
(\texttt{\small 22}) la funzione \func{get\_acl\_file} per leggerne la ACL, e
si controlla (\texttt{\small 23-26}) se l'operazione ha successo, uscendo con
un messaggio di errore in caso contrario. 

Ottenuta la ACL la si converte in formato testuale (\texttt{\small 27}) con la
funzione \func{acl\_to\_text}, controllando di nuovo se l'operazione ha
successo (\texttt{\small 28-31}) ed uscendo in caso contrario.  Si provvede
infine a stampare la rappresentazione testuale (\texttt{\small 32}) e dopo
aver liberato (\texttt{\small 33-34}) le risorse allocate automaticamente, si
conclude l'esecuzione.


\subsection{La gestione delle quote disco}
\label{sec:disk_quota}

Quella delle quote disco è una funzionalità introdotta inizialmente da BSD e
presente in Linux fino dai kernel dalla serie 2.0, che consente di porre dei
tetti massimi al consumo delle risorse di un filesystem (spazio disco e
\textit{inode}) da parte di utenti e gruppi.

Dato che la funzionalità ha senso solo per i filesystem su cui si mantengono i
dati degli utenti essa deve essere attivata esplicitamente.\footnote{in genere
  la si attiva sul filesystem che contiene le \textit{home} degli utenti, dato
  che non avrebbe senso per i file di sistema che in genere appartengono
  all'amministratore.} Questo si fa, per tutti i filesystem che le supportano,
tramite due distinte opzioni di montaggio, \texttt{usrquota} e
\texttt{grpquota}, che abilitano le quote rispettivamente per gli utenti e per
i gruppi. Così è possibile usare le limitazioni delle quote o sugli utenti o
sui gruppi o su entrambi.

Dal kernel 4.1, ed inizialmente solo per il filesystem XFS, sono diventate
disponibili un terzo tipo di quote, dette \textit{project quota}, che
consentono di applicare delle quote ad un ``\textsl{progetto}'', identificato
come ramo di albero sotto una directory, per il quale possono essere imposti
dei limiti (di nuovo in termini di spazio disco o \textit{inode}) per i file
che ne fanno parte. Si può così porre dei limiti sul contenuto di un ramo di
albero.

Il meccanismo prevede che per ciascun filesystem che supporta le quote disco
(i vari \acr{extN}, \acr{Btrfs}, \acr{XFS}, \acr{JFS}, \acr{ReiserFS}) il
kernel provveda sia a mantenere aggiornati i dati relativi al consumo delle
risorse da parte degli utenti e dei gruppi (e del progetto), che a far
rispettare i limiti imposti dal sistema, con la generazione di un errore di
\errcode{EDQUOT} per tutte le operazioni sui file che porterebbero ad un
superamento degli stessi. Si tenga presente che questi due compiti sono
separati, il primo si attiva al montaggio del filesystem con il supporto per
le quote, il secondo deve essere abilitato esplicitamente.

Per il mantenimento dei dati di consumo delle risorse vengono usati due file
riservati nella directory radice del filesystem su cui si sono attivate le
quote, uno per le quote utente e l'altro per le quote gruppo.\footnote{la cosa
  vale per tutti i filesystem tranne \textit{XFS} che mantiene i dati
  internamente, compresi quelli per le \textit{project quota}, che pertanto,
  essendo questo l'unico filesyste che le supporta, non hanno un file ad esse
  riservato.} Con la versione 2 del supporto delle quote, che da anni è
l'unica rimasta in uso, questi file sono \texttt{aquota.user} e
\texttt{aquota.group}, in precedenza erano \texttt{quota.user} e
\texttt{quota.group}.

Dato che questi file vengono aggiornati soltanto se il filesystem è stato
montato attivando il supporto delle quote, se si abilita il supporto in un
secondo tempo e nel frattempo sono state eseguite delle operazioni sul
filesystem quando il supporto era disabilitato, i dati contenuti possono non
corrispondere esattamente allo stato corrente del consumo delle risorse. Per
questo motivo prima di montare in scrittura un filesystem su cui sono
abilitate le quote viene richiesto di utilizzare il comando \cmd{quotacheck}
per verificare e aggiornare i dati.

Le restrizioni sul consumo delle risorse previste dal sistema delle quote
prevedono sempre la presenza di due diversi limiti, il primo viene detto
\textit{soft limit} e può essere superato per brevi periodi di tempo senza che
causare errori per lo sforamento delle quote, il secondo viene detto
\textit{hard limit} e non può mai essere superato.

Il periodo di tempo per cui è possibile eccedere rispetto alle restrizioni
indicate dal \textit{soft limit} è detto ``\textsl{periodo di grazia}''
(\textit{grace period}), che si attiva non appena si supera la quota da esso
indicata. Se si continua a restare al di sopra del \textit{soft limit} una
volta scaduto il \textit{grace period} questo verrà trattato allo stesso modo
dell'\textit{hard limit} e si avrà l'emissione immediata di un errore.

Si tenga presente infine che entrambi i tipi di limiti (\textit{soft limit} e
\textit{hard limit}) possono essere disposti separatamente su entrambe le
risorse di un filesystem, essi cioè possono essere presenti in maniera
indipendente sia sullo spazio disco, con un massimo per il numero di blocchi,
che sui file, con un massimo per il numero di \textit{inode}.

La funzione di sistema che consente di controllare tutti i vari aspetti della
gestione delle quote è \funcd{quotactl}, ed il suo prototipo è:

\begin{funcproto}{
\fhead{sys/types.h}
\fhead{sys/quota.h}
\fdecl{int quotactl(int cmd, const char *dev, int id, caddr\_t addr)}
\fdesc{Esegue una operazione di controllo sulle quote disco.} 
}

{La funzione ritorna $0$ in caso di successo e $-1$ per un errore, nel qual
  caso \var{errno} assumerà uno dei valori: 
  \begin{errlist}
  \item[\errcode{EACCES}] si è richiesto \const{Q\_QUOTAON}, ma il file delle
    quote indicato da \param{addr} esiste ma non è un file ordinario o no sta
    su \param{dev}.
  \item[\errcode{EBUSY}] si è richiesto \const{Q\_QUOTAON}, ma le quote sono
    già attive.
  \item[\errcode{EFAULT}] \param{addr} o \param{dev} non sono un puntatori
    validi. 
  \item[\errcode{EINVAL}] o \param{cmd} non è un comando valido,
    o il dispositivo \param{dev} non esiste.
  \item[\errcode{EIO}] errore di lettura/scrittura sul file delle quote.
  \item[\errcode{EMFILE}] non si può aprire il file delle quote avendo
    superato il limite sul numero di file aperti nel sistema. 
  \item[\errcode{ENODEV}] \param{dev} non corrisponde ad un \textit{mount
      point} attivo.
  \item[\errcode{ENOPKG}] il kernel è stato compilato senza supporto per le
    quote. 
  \item[\errcode{ENOTBLK}] \param{dev} non è un dispositivo a blocchi.
  \item[\errcode{EPERM}] non si hanno i permessi per l'operazione richiesta.
  \item[\errcode{ESRCH}] è stato richiesto uno fra \const{Q\_GETQUOTA},
    \const{Q\_SETQUOTA}, \const{Q\_SETUSE}, \const{Q\_SETQLIM} per un
    filesystem senza quote attivate.
  \end{errlist}
}
\end{funcproto}

% TODO rivedere gli errori

La funzione richiede che il filesystem sul quale si vuole operare, che deve
essere specificato con il nome del relativo file di dispositivo nell'argomento
\param{dev}, sia montato con il supporto delle quote abilitato. Per le
operazioni che lo richiedono inoltre si dovrà indicare con l'argomento
\param{id} l'utente o il gruppo o il progetto (specificati rispettivamente per
\ids{UID}, \ids{GID} o identificativo) su cui si vuole operare,\footnote{nel
  caso di \textit{project quota} gli identificativi vengono associati alla
  directory base del progetto nel file \file{/etc/projects}, ed impostati con
  \cmd{xfs\_quota}, l'argomento è di natura sistemistica e va al di là dello
  scopo di questo testo.} o altri dati relativi all'operazione. Alcune
operazioni più complesse usano infine l'argomento \param{addr} per indicare un
indirizzo ad un area di memoria il cui utilizzo dipende dall'operazione
stessa.

La funzione prevede la possibilità di eseguire una serie operazioni sulle
quote molto diverse fra loro, la scelta viene effettuata tramite il primo
argomento, \param{cmd}, che però oltre all'operazione indica anche a quale
tipo di quota (utente o gruppo) l'operazione deve applicarsi. Per questo il
valore di questo argomento viene costruito con l'ausilio della di una apposita
macro \macro{QCMD}:

{\centering
\vspace{3pt}
\begin{funcbox}{
\fhead{sys/quota.h}
\fdecl{int \macrod{QCMD}(subcmd,type)}
\fdesc{Imposta il comando \param{subcmd} per il tipo di quote (utente o
  gruppo) \param{type}.}
} 
\end{funcbox}
}

La macro consente di specificare, oltre al tipo di operazione, da indicare con
l'argomento \param{subcmd}, se questa deve applicarsi alle quote utente o alle
quote gruppo o alle quote progetto. Questo viene indicato dall'argomento
\param{type} che deve essere sempre definito ed assegnato ad uno fra i valori
\const{USRQUOTA}, \const{GRPQUOTA} e \const{PRJQUOTA}.

\begin{table}[htb]
  \centering
  \footnotesize
  \begin{tabular}{|l|p{10cm}|}
    \hline
    \textbf{Comando} & \textbf{Descrizione} \\
    \hline
    \hline
    \constd{Q\_QUOTAON} & Attiva l'applicazione delle quote disco per il
                          filesystem indicato da \param{dev}, si deve passare
                          in \param{addr} il \textit{pathname} al file che
                          mantiene le quote, che deve esistere, e \param{id}
                          deve indicare la versione del formato con uno dei
                          valori di tab.~\ref{tab:quotactl_id_format};
                          l'operazione richiede i privilegi di
                          amministratore.\\
    \constd{Q\_QUOTAOFF}& Disattiva l'applicazione delle quote disco per il
                          filesystem indicato da \param{dev}, \param{id}
                          e \param{addr} vengono ignorati; l'operazione
                          richiede i privilegi di amministratore.\\  
    \constd{Q\_GETQUOTA}& Legge i limiti ed i valori correnti delle quote nel
                          filesystem indicato da \param{dev} per l'utente o
                          il gruppo specificato da \param{id}; si devono avere
                          i privilegi di amministratore per leggere i dati
                          relativi ad altri utenti o a gruppi di cui non si fa
                          parte, il risultato viene restituito in una struttura
                          \struct{dqblk} all'indirizzo indicato
                          da \param{addr}.\\
    \constd{Q\_SETQUOTA}& Imposta i limiti per le quote nel filesystem
                          indicato da \param{dev} per l'utente o il gruppo
                          specificato da \param{id} secondo i valori ottenuti
                          dalla struttura \struct{dqblk} puntata
                          da \param{addr}; l'operazione richiede i privilegi
                          di amministratore.\\ 
    \constd{Q\_GETINFO} & Legge le informazioni (in sostanza i \textit{grace
                            time}) delle quote del filesystem indicato
                          da \param{dev} sulla struttura \struct{dqinfo} 
                          puntata da \param{addr}, \param{id} viene ignorato.\\
    \constd{Q\_SETINFO} & Imposta le informazioni delle quote del filesystem
                          indicato da \param{dev} come ottenuti dalla
                          struttura \struct{dqinfo} puntata
                          da \param{addr}, \param{id} viene ignorato;  
                          l'operazione richiede i privilegi di amministratore.\\
    \constd{Q\_GETFMT}  & Richiede il valore identificativo (quello di
                          tab.~\ref{tab:quotactl_id_format}) per il formato
                          delle quote attualmente in uso sul filesystem
                          indicato da \param{dev}, che sarà memorizzato
                          sul buffer di 4 byte puntato da \param{addr}.\\
    \constd{Q\_SYNC}    & Aggiorna la copia su disco dei dati delle quote del
                          filesystem indicato da \param{dev}; in questo
                          caso \param{dev} può anche essere \val{NULL} nel
                          qual caso verranno aggiornati i dati per tutti i
                          filesystem con quote attive, \param{id}
                          e \param{addr} vengono comunque ignorati.\\ 
    \constd{Q\_GETSTATS}& Ottiene statistiche ed altre informazioni generali 
                          relative al sistema delle quote per il filesystem
                          indicato da \param{dev}, richiede che si
                          passi come argomento \param{addr} l'indirizzo di una
                          struttura \struct{dqstats}, mentre i valori
                          di \param{id} e \param{dev} vengono ignorati;
                          l'operazione è obsoleta e non supportata nei kernel
                          più recenti, che espongono la stessa informazione
                          nei file sotto \procfile{/proc/self/fs/quota/}.\\
%    \const{} & .\\
    \hline
  \end{tabular}
  \caption{Possibili valori per l'argomento \param{subcmd} di
    \macro{QCMD}.} 
  \label{tab:quotactl_commands}
\end{table}

I possibili valori per l'argomento \param{subcmd} di \macro{QCMD} sono
riportati in tab.~\ref{tab:quotactl_commands}, che illustra brevemente il
significato delle operazioni associate a ciascuno di essi. In generale le
operazioni di attivazione, disattivazione e di modifica dei limiti delle quote
sono riservate e richiedono i privilegi di amministratore.\footnote{per essere
  precisi tutte le operazioni indicate come privilegiate in
  tab.~\ref{tab:quotactl_commands} richiedono la capacità
  \const{CAP\_SYS\_ADMIN}.} Inoltre gli utenti possono soltanto richiedere i
dati relativi alle proprie quote, solo l'amministratore può ottenere i dati di
tutti.


Alcune delle operazioni di tab.~\ref{tab:quotactl_commands} sono alquanto
complesse e richiedono un approfondimento maggiore. Le due più rilevanti sono
probabilmente \const{Q\_GETQUOTA} e \const{Q\_SETQUOTA}, che consentono la
gestione dei limiti delle quote. Entrambe fanno riferimento ad una specifica
struttura \struct{dqblk}, la cui definizione è riportata in
fig.~\ref{fig:dqblk_struct},\footnote{la definizione mostrata è quella usata
  fino dal kernel 2.4.22, non prenderemo in considerazione le versioni
  obsolete.} nella quale vengono inseriti i dati relativi alle quote di un
singolo utente o gruppo.

\begin{figure}[!htb]
  \footnotesize \centering
  \begin{minipage}[c]{0.95\textwidth}
    \includestruct{listati/dqblk.h}
  \end{minipage} 
  \normalsize 
  \caption{La struttura \structd{dqblk} per i dati delle quote disco.}
  \label{fig:dqblk_struct}
\end{figure}

La struttura \struct{dqblk} viene usata sia con \const{Q\_GETQUOTA} per
ottenere i valori correnti dei limiti e dell'occupazione delle risorse, che
con \const{Q\_SETQUOTA} per effettuare modifiche ai limiti. Come si può notare
ci sono alcuni campi (in sostanza \val{dqb\_curspace}, \val{dqb\_curinodes},
\val{dqb\_btime}, \val{dqb\_itime}) che hanno senso solo in lettura, in quanto
riportano uno stato non modificabile da \func{quotactl} come l'uso corrente di
spazio disco ed \textit{inode}, o il tempo che resta nel caso si sia superato
un \textit{soft limit}.

Inoltre in caso di modifica di un limite si può voler operare solo su una
delle risorse (blocchi o \textit{inode}),\footnote{non è possibile modificare
  soltanto uno dei limiti (\textit{hard} o \textit{soft}) occorre sempre
  rispecificarli entrambi.} per questo la struttura prevede un campo apposito,
\val{dqb\_valid}, il cui scopo è quello di indicare quali sono gli altri campi
che devono essere considerati validi. Questo campo è una maschera binaria che
deve essere espressa nei termini di OR aritmetico delle apposite costanti di
tab.~\ref{tab:quotactl_qif_const}, dove si è riportato il significato di
ciascuna di esse ed i campi a cui fanno riferimento.

\begin{table}[!htb]
  \centering
  \footnotesize
  \begin{tabular}{|l|p{10cm}|}
    \hline
    \textbf{Costante} & \textbf{Descrizione} \\
    \hline
    \hline
    \constd{QIF\_BLIMITS}& Limiti sui blocchi di spazio disco
                           (\val{dqb\_bhardlimit} e \val{dqb\_bsoftlimit}).\\
    \constd{QIF\_SPACE}  & Uso corrente dello spazio disco
                           (\val{dqb\_curspace}).\\
    \constd{QIF\_ILIMITS}& Limiti sugli \textit{inode}
                           (\val{dqb\_ihardlimit} e \val{dqb\_isoftlimit}).\\
    \constd{QIF\_INODES} & Uso corrente degli \textit{inode}
                           (\val{dqb\_curinodes}).\\
    \constd{QIF\_BTIME}  & Tempo di sforamento del \textit{soft limit} sul
                           numero di blocchi (\val{dqb\_btime}).\\
    \constd{QIF\_ITIME}  & Tempo di sforamento del \textit{soft limit} sul
                           numero di \textit{inode} (\val{dqb\_itime}).\\ 
    \constd{QIF\_LIMITS} & L'insieme di \const{QIF\_BLIMITS} e
                           \const{QIF\_ILIMITS}.\\
    \constd{QIF\_USAGE}  & L'insieme di \const{QIF\_SPACE} e
                           \const{QIF\_INODES}.\\
    \constd{QIF\_TIMES}  & L'insieme di \const{QIF\_BTIME} e
                           \const{QIF\_ITIME}.\\ 
    \constd{QIF\_ALL}    & Tutti i precedenti.\\
    \hline
  \end{tabular}
  \caption{Costanti per il campo \val{dqb\_valid} di \struct{dqblk}.} 
  \label{tab:quotactl_qif_const}
\end{table}

In lettura con \const{Q\_SETQUOTA} eventuali valori presenti in \struct{dqblk}
vengono comunque ignorati, al momento la funzione sovrascrive tutti i campi
che restituisce e li marca come validi in \val{dqb\_valid}. Si possono invece
usare \const{QIF\_BLIMITS} o \const{QIF\_ILIMITS} per richiedere di impostare
solo la rispettiva tipologia di limiti con \const{Q\_SETQUOTA}. Si tenga
presente che il sistema delle quote richiede che l'occupazione di spazio disco
sia indicata in termini di blocchi e non di byte, dato che la dimensione dei
blocchi dipende da come si è creato il filesystem potrà essere necessario
effettuare qualche conversione per avere un valore in byte.\footnote{in genere
  viene usato un default di 1024 byte per blocco, ma quando si hanno file di
  dimensioni medie maggiori può convenire usare valori più alti per ottenere
  prestazioni migliori in conseguenza di un minore frazionamento dei dati e di
  indici più corti.}

Come accennato la realizzazione delle quote disco ha visto diverse revisioni,
con modifiche sia del formato delle stesse che dei nomi dei file
utilizzati. Per questo alcune operazioni di gestione (in particolare
\const{Q\_QUOTAON} e \const{Q\_GETFMT}) e possono fare riferimento a queste
versioni, che vengono identificate tramite le costanti di
tab.~\ref{tab:quotactl_id_format}.

\begin{table}[htb]
  \centering
  \footnotesize
  \begin{tabular}{|l|p{10cm}|}
    \hline
    \textbf{Identificatore} & \textbf{Descrizione} \\
    \hline
    \hline
    \constd{QFMT\_VFS\_OLD}& Il vecchio (ed obsoleto) formato delle quote.\\
    \constd{QFMT\_VFS\_V0} & La versione 0 usata dal VFS di Linux, supporta
                             \ids{UID} e \ids{GID} a 32 bit e limiti fino a
                             $2^{42}$ byte e $2^{32}$ file.\\
    \constd{QFMT\_VFS\_V1} & La versione 1 usata dal VFS di Linux, supporta
                             \ids{UID} e \ids{GID} a 32 bit e limiti fino a
                             $2^{64}$ byte e $2^{64}$ file.\\
    \hline
  \end{tabular}
  \caption{Valori di identificazione del formato delle quote.} 
  \label{tab:quotactl_id_format}
\end{table}

Altre due operazioni che necessitano di ulteriori spiegazioni sono
\const{Q\_GETINFO} e \const{Q\_SETINFO}, che consentono di ottenere i dati
relativi alle impostazioni delle altre proprietà delle quote, che al momento
sono solo la durata del \textit{grace time} per i due tipi di limiti. Queste
sono due proprietà generali identiche per tutti gli utenti (e i gruppi), per
cui viene usata una operazione distinta dalle precedenti. Anche in questo caso
le due operazioni richiedono l'uso di una apposita struttura \struct{dqinfo},
la cui definizione è riportata in fig.~\ref{fig:dqinfo_struct}.

\begin{figure}[!htb]
  \footnotesize \centering
  \begin{minipage}[c]{0.8\textwidth}
    \includestruct{listati/dqinfo.h}
  \end{minipage} 
  \normalsize 
  \caption{La struttura \structd{dqinfo} per i dati delle quote disco.}
  \label{fig:dqinfo_struct}
\end{figure}

Come per \struct{dqblk} anche in questo caso viene usato un campo della
struttura, \val{dqi\_valid} come maschera binaria per dichiarare quale degli
altri campi sono validi; le costanti usate per comporre questo valore sono
riportate in tab.~\ref{tab:quotactl_iif_const} dove si è riportato il
significato di ciascuna di esse ed i campi a cui fanno riferimento.

\begin{table}[htb]
  \centering
  \footnotesize
  \begin{tabular}{|l|l|}
    \hline
    \textbf{Costante} & \textbf{Descrizione} \\
    \hline
    \hline
    \constd{IIF\_BGRACE}& Il \textit{grace period} per i blocchi
                         (\val{dqi\_bgrace}).\\
    \constd{IIF\_IGRACE}& Il \textit{grace period} per gli \textit{inode} 
                         (\val{dqi\_igrace}).\\ 
    \constd{IIF\_FLAGS} & I flag delle quote (\val{dqi\_flags}) (inusato ?).\\
    \constd{IIF\_ALL}   & Tutti i precedenti.\\
    \hline
  \end{tabular}
  \caption{Costanti per il campo \val{dqi\_valid} di \struct{dqinfo}.} 
  \label{tab:quotactl_iif_const}
\end{table}

Come in precedenza con \const{Q\_GETINFO} tutti i valori vengono letti
sovrascrivendo il contenuto di \struct{dqinfo} e marcati come validi in
\val{dqi\_valid}. In scrittura con \const{Q\_SETINFO} si può scegliere quali
impostare, si tenga presente che i tempi dei campi \val{dqi\_bgrace} e
\val{dqi\_igrace} devono essere specificati in secondi.

Come esempi dell'uso di \func{quotactl} utilizzeremo estratti del codice di un
modulo Python usato per fornire una interfaccia diretta a \func{quotactl}
senza dover passare dalla scansione dei risultati di un comando. Il modulo si
trova fra i pacchetti Debian messi a disposizione da Truelite Srl,
all'indirizzo \url{http://labs.truelite.it/projects/packages}.\footnote{in
  particolare il codice C del modulo è nel file \texttt{quotamodule.c}
  visionabile a partire dall'indirizzo indicato nella sezione
  \textit{Repository}.}

\begin{figure}[!htbp]
  \footnotesize \centering
  \begin{minipage}[c]{\codesamplewidth}
    \includecodesample{listati/get_quota.c}
  \end{minipage}
  \caption{Esempio di codice per ottenere i dati delle quote.} 
  \label{fig:get_quota}
\end{figure}

Il primo esempio, riportato in fig.~\ref{fig:get_quota}, riporta il codice
della funzione che consente di leggere le quote. La funzione fa uso
dell'interfaccia dal C verso Python, che definisce i vari simboli \texttt{Py*}
(tipi di dato e funzioni). Non staremo ad approfondire i dettagli di questa
interfaccia, per la quale esistono numerose trattazioni dettagliate, ci
interessa solo esaminare l'uso di \func{quotactl}. 

In questo caso la funzione prende come argomenti (\texttt{\small 1}) l'intero
\texttt{who} che indica se si vuole operare sulle quote utente o gruppo,
l'identificatore \texttt{id} dell'utente o del gruppo scelto, ed il nome del
file di dispositivo del filesystem su cui si sono attivate le
quote.\footnote{questi vengono passati come argomenti dalle funzioni mappate
  come interfaccia pubblica del modulo (una per gruppi ed una per gli utenti)
  che si incaricano di decodificare i dati passati da una chiamata nel codice
  Python.} Questi argomenti vengono passati direttamente alla chiamata a
\func{quotactl} (\texttt{\small 5}), a parte \texttt{who} che viene abbinato
con \macro{QCMD} al comando \const{Q\_GETQUOTA} per ottenere i dati.

La funzione viene eseguita all'interno di un condizionale (\texttt{\small
  5-16}) che in caso di successo provvede a costruire (\texttt{\small 6-12})
opportunamente una risposta restituendo tramite la opportuna funzione di
interfaccia un oggetto Python contenente i dati della struttura \struct{dqblk}
relativi a uso corrente e limiti sia per i blocchi che per gli
\textit{inode}. In caso di errore (\texttt{\small 13-15}) si usa un'altra
funzione dell'interfaccia per passare il valore di \var{errno} come eccezione.

\begin{figure}[!htbp]
  \footnotesize \centering
  \begin{minipage}[c]{\codesamplewidth}
    \includecodesample{listati/set_block_quota.c}
  \end{minipage}
  \caption{Esempio di codice per impostare i limiti sullo spazio disco.}
  \label{fig:set_block_quota}
\end{figure}

Per impostare i limiti sullo spazio disco si potrà usare una seconda funzione,
riportata in fig.~\ref{fig:set_block_quota}, che prende gli stessi argomenti
della precedente, con lo stesso significato, a cui si aggiungono i valori per
il \textit{soft limit} e l'\textit{hard limit}. In questo caso occorrerà,
prima di chiamare \func{quotactl}, inizializzare opportunamente
(\texttt{\small 5-7}) i campi della struttura \struct{dqblk} che si vogliono
utilizzare (quelli relativi ai limiti sui blocchi) e specificare gli stessi
con \const{QIF\_BLIMITS} in \var{dq.dqb\_valid}. 

Fatto questo la chiamata a \func{quotactl}, stavolta con il comando
\const{Q\_SETQUOTA}, viene eseguita come in precedenza all'interno di un
condizionale (\texttt{\small 9-14}). In questo caso non essendovi da
restituire nessun dato in caso di successo si usa (\texttt{\small 10}) una
apposita funzione di uscita, mentre si restituisce come prima una eccezione
con il valore di \var{errno} in caso di errore (\texttt{\small 12-13}).

\subsection{La gestione dei {chroot}}
\label{sec:file_chroot}

% TODO: valutare se introdurre una nuova sezione sulle funzionalità di
% sicurezza avanzate, con dentro chroot SELinux e AppArmor, Tomoyo, Smack,
% cgroup o altro

% TODO: spostare chroot e le funzioni affini relative ai container da qualche
% parte diversa se è il caso. 

% TODO Inheriting capabilities vedi http://lwn.net/Articles/632520/ eambient
% capabilities introdotte con il kernel 4.3, vedi 
% http://git.kernel.org/cgit/linux/kernel/git/torvalds/linux.git/commit/?id=58319057b7847667f0c9585b9de0e8932b0fdb08

Benché non abbia niente a che fare con permessi, utenti e gruppi, la funzione
\func{chroot} viene usata spesso per restringere le capacità di accesso di un
programma ad una sezione limitata del filesystem, per cui ne parleremo in
questa sezione.

Come accennato in sez.~\ref{sec:proc_fork} ogni processo oltre ad una
directory di lavoro, ha anche una directory \textsl{radice}\footnote{entrambe
  sono contenute in due campi (rispettivamente \var{pwd} e \var{root}) di
  \kstruct{fs\_struct}; vedi fig.~\ref{fig:proc_task_struct}.} che, pur
essendo di norma corrispondente alla radice dell'albero dei file dell'intero
sistema, ha per il processo il significato specifico di directory rispetto
alla quale vengono risolti i \textit{pathname} assoluti.\footnote{cioè quando
  un processo chiede la risoluzione di un \textit{pathname}, il kernel usa
  sempre questa directory come punto di partenza.} Il fatto che questo valore
sia specificato per ogni processo apre allora la possibilità di modificare le
modalità di risoluzione dei \textit{pathname} assoluti da parte di un processo
cambiando questa directory, così come si fa coi \textit{pathname} relativi
cambiando la directory di lavoro.

Normalmente la directory radice di un processo coincide con la radice generica
dell'albero dei file, che è la directory che viene montata direttamente dal
kernel all'avvio secondo quanto illustrato in sez.~\ref{sec:file_pathname}.
Questo avviene perché, come visto in sez.~\ref{cha:process_handling} la
directory radice di un processo viene ereditata dal padre attraverso una
\func{fork} e mantenuta attraverso una \func{exec}, e siccome tutti i processi
derivano da \cmd{init}, che ha come radice quella montata dal kernel, questa
verrà mantenuta.

In certe situazioni però è utile poter impedire che un processo possa accedere
a tutto l'albero dei file iniziale; per far questo si può cambiare la sua
directory radice con la funzione di sistema \funcd{chroot}, il cui prototipo
è:

\begin{funcproto}{
\fhead{unistd.h}
\fdecl{int chroot(const char *path)}
\fdesc{Cambia la directory radice del processo.} 
}

{La funzione ritorna $0$ in caso di successo e $-1$ per un errore, nel qual
  caso \var{errno} assumerà uno dei valori: 
  \begin{errlist}
  \item[\errcode{EPERM}] non si hanno i privilegi di amministratore.
  \end{errlist}
  ed inoltre \errval{EFAULT}, \errval{ENAMETOOLONG}, \errval{ENOENT},
  \errval{ENOMEM}, \errval{ENOTDIR}, \errval{EACCES}, \errval{ELOOP};
  \errval{EROFS} e \errval{EIO} nel loro significato generico.}
\end{funcproto}

La funzione imposta la directory radice del processo a quella specificata da
\param{path} (che ovviamente deve esistere) ed ogni \textit{pathname} assoluto
usato dalle funzioni chiamate nel processo sarà risolto a partire da essa,
rendendo impossibile accedere alla parte di albero sovrastante. Si ha così
quella che viene chiamata una \textit{chroot jail}, in quanto il processo non
può più accedere a file al di fuori della sezione di albero in cui è stato
\textsl{imprigionato}.

Solo un processo con i privilegi di amministratore può usare questa
funzione,\footnote{più precisamente se possiede la capacità
  \const{CAP\_SYS\_CHROOT}.} e la nuova radice, per quanto detto in
sez.~\ref{sec:proc_fork}, sarà ereditata da tutti i suoi processi figli. Si
tenga presente però che la funzione non cambia la directory di lavoro del
processo, che potrebbe restare fuori dalla \textit{chroot jail}.

Questo è il motivo per cui la funzione è efficace nel restringere un processo
ad un ramo di albero solo se dopo averla eseguita si cedono i privilegi di
amministratore. Infatti se per un qualunque motivo il processo resta con la
sua directory di lavoro al di fuori dalla \textit{chroot jail}, potrà accedere
a tutto il resto del filesystem usando dei \textit{pathname} relativi, dato
che in tal caso è possibile, grazie all'uso di ``\texttt{..}'', risalire
all'indietro fino alla radice effettiva dell'albero dei file.

Potrebbe sembrare che per risolvere il problema sia sufficiente ricordarsi di
eseguire preventivamente anche una \func{chdir} sulla directory su cui si
andrà ad eseguire \func{chroot}, così da assicurarsi che le directory di
lavoro sia all'interno della \textit{chroot jail}.  Ma se ad un processo
restano i privilegi di amministratore esso potrà comunque portare la sua
directory di lavoro fuori dalla \textit{chroot jail} in cui si trova. Basterà
infatti eseguire di nuovo \func{chroot} su una qualunque directory contenuta
nell'attuale directory di lavoro perché quest'ultima risulti al di fuori della
nuova \textit{chroot jail}.  Per questo motivo l'uso di questa funzione non ha
molto senso quando un processo di cui si vuole limitare l'accesso necessita
comunque dei privilegi di amministratore per le sue normali operazioni.

Nonostante queste limitazioni la funzione risulta utile qualora la si possa
applicare correttamente cedendo completamente i privilegi di amministratore
una volta eseguita.  Ed esempio caso tipico di uso di \func{chroot} è quello
di un server FTP anonimo in si vuole che il server veda solo i file che deve
trasferire. In tal caso si esegue una \func{chroot} sulla directory che
contiene i file, che il server dovrà in grado di leggere come utente
ordinario, e poi si cedono tutti i privilegi di amministratore.  Si tenga
presente però che in casi come questo occorrerà fornire all'interno della
\textit{chroot jail} un accesso anche a tutti i file (in genere programmi e
librerie) di cui il server potrebbe avere bisogno.


% LocalWords:  sez like filesystem unlink MacOS Windows VMS inode kernel unistd
% LocalWords:  int const char oldpath newpath errno EXDEV EPERM st Smack SysV
% LocalWords:  EEXIST EMLINK EACCES ENAMETOOLONG ENOTDIR EFAULT ENOMEM EROFS ls
% LocalWords:  ELOOP ENOSPC EIO pathname nlink stat vfat fsck EISDIR ENOENT cap
% LocalWords:  POSIX socket fifo sticky root system call count crash init linux
% LocalWords:  descriptor remove rename rmdir stdio glibc libc NFS DT obj dup
% LocalWords:  ENOTEMPTY EBUSY mount point EINVAL soft symbolic tab symlink fig
% LocalWords:  dangling access chdir chmod chown creat exec lchown lstat mkdir
% LocalWords:  mkfifo mknod opendir pathconf readlink truncate path buff size
% LocalWords:  grub bootloader grep MAXSYMLINKS cat VFS sys dirname fcntl tv Py
% LocalWords:  dev umask IFREG IFBLK IFCHR IFIFO SVr sgid BSD SVID NULL from to
% LocalWords:  stream dirent EMFILE ENFILE dirfd SOURCE fchdir readdir struct
% LocalWords:  EBADF namlen HAVE thread entry result value argument fileno ext
% LocalWords:  name TYPE OFF RECLEN UNKNOWN REG SOCK CHR BLK type IFTODT DTTOIF
% LocalWords:  DTYPE off reclen seekdir telldir void rewinddir closedir select
% LocalWords:  namelist compar malloc qsort alphasort versionsort strcoll myls
% LocalWords:  strcmp direntry while current working home shell pwd get stddef
% LocalWords:  getcwd ERANGE getwd change fd race condition tmpnam getfacl mark
% LocalWords:  string tmpdir TMP tempnam pfx TMPNAME suid tmp EXCL tmpfile pid
% LocalWords:  EINTR mktemp mkstemp stlib template filename XXXXXX OpenBSD buf
% LocalWords:  mkdtemp fstat filedes nell'header padding ISREG ISDIR ISCHR IFMT
% LocalWords:  ISBLK ISFIFO ISLNK ISSOCK IFSOCK IFLNK IFDIR ISUID UID ISGID GID
% LocalWords:  ISVTX IRUSR IWUSR IXUSR IRGRP IWGRP IXGRP IROTH IWOTH IXOTH  OLD
% LocalWords:  blocks blksize holes lseek TRUNC ftruncate ETXTBSY length QCMD
% LocalWords:  hole atime read utime mtime write ctime modification leafnode cp
% LocalWords:  make fchmod fchown utimbuf times actime modtime Mac owner uid fs
% LocalWords:  gid Control List patch mandatory control execute group other all
% LocalWords:  effective passwd IGID locking swap saved text IRWXU IRWXG subcmd
% LocalWords:  IRWXO capability FSETID mask capabilities chroot jail QUOTAOFF
% LocalWords:  FTP filter Attributes Solaris FreeBSD libacl hash at dqblk SYNC
% LocalWords:  XFS SELinux namespace attribute security trusted Draft Modules
% LocalWords:  attributes mime ADMIN FOWNER libattr lattr getxattr lgetxattr of
% LocalWords:  fgetxattr attr ssize ENOATTR ENOTSUP NUL setxattr lsetxattr list
% LocalWords:  fsetxattr flags XATTR REPLACE listxattr llistxattr flistxattr by
% LocalWords:  removexattr lremovexattr fremovexattr attributename acl GETINFO
% LocalWords:  OBJ setfacl len any prefix separator options NUMERIC IDS SMART
% LocalWords:  INDENT major number IDE Documentation makedev proc copy LNK long
% LocalWords:  euidaccess eaccess delete def tag qualifier permset calendar NOW
% LocalWords:  mutt noatime relatime strictatime atim nsec mtim ctim atimensec
% LocalWords:  mtimensec utimes timeval futimes lutimes ENOSYS futimens OMIT PR
% LocalWords:  utimensat timespec sec futimesat LIDS DAC OVERRIDE SEARCH chattr
% LocalWords:  Discrectionary KILL SETGID domain SETUID setuid setreuid SETPCAP
% LocalWords:  setresuid setfsuid IMMUTABLE immutable append only BIND SERVICE
% LocalWords:  BROADCAST broadcast multicast multicasting RAW PACKET IPC LOCK
% LocalWords:  memory mlock mlockall shmctl mmap MODULE RAWIO ioperm iopl PACCT
% LocalWords:  ptrace accounting NICE RESOURCE TTY CONFIG hangup vhangup dell'
% LocalWords:  LEASE lease SETFCAP AUDIT permitted inherited inheritable AND nn
% LocalWords:  bounding execve fork capget capset header hdrp datap ESRCH undef
% LocalWords:  version libcap lcap clear ncap caps pag capgetp CapInh CapPrm RT
% LocalWords:  fffffeff CapEff getcap scheduling lookup  dqinfo SETINFO GETFMT
% LocalWords:  NEWNS unshare nice NUMA ioctl journaling close XOPEN fdopendir
% LocalWords:  Btrfs mkostemp extN ReiserFS JFS Posix usrquota grpquota EDQUOT
% LocalWords:  aquota quotacheck limit grace period quotactl cmd caddr addr dqb
% LocalWords:  QUOTAON ENODEV ENOPKG ENOTBLK GETQUOTA SETQUOTA SETUSE SETQLIM
% LocalWords:  forced allowed sendmail SYSLOG WAKE ALARM CLOCK BOOTTIME dqstats
% LocalWords:  REALTIME securebits GETSTATS QFMT curspace curinodes btime itime
% LocalWords:  QIF BLIMITS bhardlimit bsoftlimit ILIMITS ihardlimit isoftlimit
% LocalWords:  INODES LIMITS USAGE valid dqi IIF BGRACE bgrace IGRACE igrace is
% LocalWords:  Python Truelite Srl quotamodule Repository who nell' dall' KEEP
% LocalWords:  SECURE KEEPCAPS prctl FIXUP NOROOT LOCKED dell'IPC dell'I IOPRIO
% LocalWords:  CAPBSET CLASS IDLE dcookie overflow DIFFERS Virtual everything
% LocalWords:  dentry register resolution cache dcache operation llseek poll ln
% LocalWords:  multiplexing fsync fasync seek block superblock gapil tex img du
% LocalWords:  second linked journaled source filesystemtype unsigned device
% LocalWords:  mountflags NODEV ENXIO dummy devfs magic MGC RDONLY NOSUID MOVE
% LocalWords:  NOEXEC SYNCHRONOUS REMOUNT MANDLOCK NODIRATIME umount MNT statfs
% LocalWords:  fstatfs fstab mntent ino bound orig new setpcap metadati sysfs
% LocalWords:  bind DIRSYNC lsattr Hierarchy FHS SHARED UNBINDABLE shared REC
% LocalWords:  subtree SILENT log unbindable BUSY EAGAIN EXPIRE DETACH NOFOLLOW
% LocalWords:  lazy encfs sshfs setfsent getfsent getfsfile getfsspec endfsent
% LocalWords:  setmntent getmntent addmntent endmntent hasmntopt such offsetof
% LocalWords:  member scan attack EOVERFLOW BITS blkcnt rdev FDCWD functions
% LocalWords:  faccessat grpid lacl AppArmor capsetp mygetfacl table Tb MSK
%  LocalWords:  LAZYTIME submount peer protected hardlink symlinks silly RDWR
%  LocalWords:  renames unreachable CLOEXEC mkstemps mkostemps suffixlen Aug
%  LocalWords:  prefissoXXXXXXsuffisso nell'I fstatat statx sull' drwxrwxrwt
%  LocalWords:  Disalloca

%%% Local Variables: 
%%% mode: latex
%%% TeX-master: "gapil"
%%% End: 

%% fileio.tex (merge fileunix.tex - filestd.tex)
%%
%% Copyright (C) 2000-2018 Simone Piccardi.  Permission is granted to
%% copy, distribute and/or modify this document under the terms of the GNU Free
%% Documentation License, Version 1.1 or any later version published by the
%% Free Software Foundation; with the Invariant Sections being "Un preambolo",
%% with no Front-Cover Texts, and with no Back-Cover Texts.  A copy of the
%% license is included in the section entitled "GNU Free Documentation
%% License".
%%

\chapter{La gestione dell'I/O su file}
\label{cha:file_IO_interface}

Esamineremo in questo capitolo le due interfacce di programmazione che
consentono di gestire i dati mantenuti nei file. Cominceremo con quella nativa
del sistema, detta dei \textit{file descriptor}, che viene fornita
direttamente dalle \textit{system call} e che non prevede funzionalità evolute
come la bufferizzazione o funzioni di lettura o scrittura
formattata. Esamineremo poi anche l'interfaccia definita dallo standard ANSI
C, che viene chiamata dei \textit{file stream} o anche più brevemente degli
\textit{stream}. Per entrambe dopo una introduzione alle caratteristiche
generali tratteremo le funzioni base per la gestione dell'I/O, lasciando per
ultime le caratteristiche più avanzate.


\section{L'interfaccia dei \textit{file descriptor}}
\label{sec:file_unix_interface}


Come visto in sez.~\ref{sec:file_vfs_work} il kernel mette a disposizione
tramite il \textit{Virtual File System} una serie di \textit{system call} che
consentono di operare sui file in maniera generale. Abbiamo trattato quelle
relative alla gestione delle proprietà dei file nel precedente capitolo,
vedremo quelle che si applicano al contenuto dei file in questa sezione,
iniziando con una breve introduzione sull'architettura dei \textit{file
  descriptor} per poi trattare le funzioni di base e le modalità con cui
consentono di gestire i dati memorizzati sui file.


\subsection{I \textit{file descriptor}}
\label{sec:file_fd}

\itindbeg{file~descriptor} 

L'accesso al contenuto dei file viene effettuato, sia pure con differenze
nella realizzazione pratica, in maniera sostanzialmente identica in tutte le
implementazioni di un sistema unix-like, ricorrendo a quella che viene
chiamata l'interfaccia dei \textit{file descriptor}.

Per poter accedere al contenuto di un file occorre creare un canale di
comunicazione con il kernel che renda possibile operare su di esso. Questo si
fa aprendo il file con la funzione \func{open} (vedi
sez.~\ref{sec:file_open_close}) che provvederà a localizzare l'\textit{inode}
del file e inizializzare i puntatori che rendono disponibili le funzioni che
il VFS mette a disposizione (quelle di
tab.~\ref{tab:file_file_operations}). Una volta terminate le operazioni, il 
file dovrà essere chiuso, e questo chiuderà il canale di comunicazione
impedendo ogni ulteriore operazione.

All'interno di ogni processo i file aperti sono identificati da un numero
intero non negativo, che viene chiamato appunto \textit{file descriptor}.
Quando un file viene aperto la funzione \func{open} restituisce questo numero,
tutte le ulteriori operazioni dovranno essere compiute specificando questo
stesso numero come argomento alle varie funzioni dell'interfaccia.

\itindbeg{process~table}
\itindbeg{file~table}

Per capire come funziona il meccanismo occorre spiegare a grandi linee come il
kernel gestisce l'interazione fra processi e file.  Abbiamo già accennato in
sez.~\ref{sec:proc_hierarchy} come il kernel mantenga un elenco di tutti
processi nella cosiddetta \textit{process table}. Lo stesso, come accennato in
sez.~\ref{sec:file_vfs_work}, vale anche per tutti i file aperti, il cui
elenco viene mantenuto nella cosiddetta \textit{file table}.

La \textit{process table} è una tabella che contiene una voce per ciascun
processo attivo nel sistema. Ciascuna voce è costituita dal puntatore a una
struttura di tipo \kstruct{task\_struct} nella quale sono raccolte tutte le
informazioni relative al processo, fra queste informazioni c'è anche il
puntatore ad una ulteriore struttura di tipo
\kstruct{files\_struct},\footnote{la definizione corrente di questa struttura
  si trova nel file \texttt{include/linux/fdtable.h} dei sorgenti del kernel,
  quella mostrata in fig.~\ref{fig:file_proc_file} è una versione pesantemente
  semplificata.} che contiene le informazioni relative ai file che il processo
ha aperto.

La \textit{file table} è una tabella che contiene una voce per ciascun file
che è stato aperto nel sistema. Come accennato in sez.~\ref{sec:file_vfs_work}
per ogni file aperto viene allocata una struttura \kstruct{file} e la
\textit{file table} è costituita da un elenco di puntatori a ciascuna di
queste strutture, che, come illustrato in fig.~\ref{fig:kstruct_file},
contengono le informazioni necessarie per la gestione dei file, ed in
particolare:
\begin{itemize*}
\item i flag di stato del file nel campo \var{f\_flags}.
\item la posizione corrente nel file, il cosiddetto \textit{offset}, nel campo
  \var{f\_pos}.
\item un puntatore alla struttura \kstruct{inode} che identifica
  l'\textit{inode} del file.\footnote{nel kernel 2.4.x si è in realtà passati
    ad un puntatore ad una struttura \kstruct{dentry} che punta a sua volta
    all'\textit{inode} passando per la nuova struttura del VFS.}
\item un puntatore \var{f\_op} alla tabella delle funzioni che si possono
  usare sul file.\footnote{quelle della struttura \kstruct{file\_operation},
    descritte sommariamente in tab.~\ref{tab:file_file_operations}.}
\end{itemize*}

\begin{figure}[!htb]
  \centering
  \includegraphics[width=12cm]{img/procfile}
  \caption{Schema della architettura dell'accesso ai file attraverso
  l'interfaccia dei file descriptor.}
  \label{fig:file_proc_file}
\end{figure}

In fig.~\ref{fig:file_proc_file} si è riportato uno schema semplificato in cui
è illustrata questa architettura, ed in cui si sono evidenziate le
interrelazioni fra la \textit{file table}, la \textit{process table} e le
varie strutture di dati che il kernel mantiene per ciascun file e ciascun
processo.

\itindend{process~table}

Come si può notare alla fine il collegamento che consente di porre in
relazione i file ed i processi è effettuato attraverso i dati mantenuti nella
struttura \kstruct{files\_struct}, essa infatti contiene alcune informazioni
essenziali come:
\begin{itemize*}
\item i flag relativi ai file aperti dal processo.
\item il numero di file aperti dal processo.
\item la \itindex{file~descriptor~table} \textit{file descriptor table}, una
  tabella con i puntatori, per ciascun file aperto, alla relativa voce nella
  \textit{file table}.
\end{itemize*}

In questa infrastruttura un file descriptor non è altro che l'intero positivo
che indicizza quest'ultima tabella, e che consente di recuperare il puntatore
alla struttura \kstruct{file} corrispondente al file aperto dal processo a cui
era stato assegnato questo indice. Una volta ottenuta grazie al file
descriptor la struttura \kstruct{file} corrispondente al file voluto nella
\textit{file table}, il kernel potrà usare le funzioni messe disposizione dal
VFS per eseguire sul file tutte le operazioni necessarie.

Il meccanismo dell'apertura dei file prevede che venga sempre fornito il primo
file descriptor libero nella tabella, e per questo motivo essi vengono
assegnati in successione tutte le volte che si apre un nuovo file, posto che
non ne sia stato chiuso nessuno in precedenza.

\itindbeg{standard~input} 
\itindbeg{standard~output}
\itindbeg{standard~error}

In tutti i sistemi unix-like esiste una convenzione generale per cui ogni
processo si aspetta di avere sempre tre file aperti che, per quanto appena
detto, avranno come \textit{file descriptor} i valori 0, 1 e 2.  Il primo file
è sempre associato al cosiddetto \textit{standard input}, è cioè il file da
cui un processo si aspetta di dover leggere i dati in ingresso. Il secondo
file è il cosiddetto \textit{standard output}, cioè quello su cui ci si
aspetta di dover scrivere i dati in uscita. Il terzo è lo \textit{standard
  error}, su cui vengono scritti i dati relativi agli errori.

\itindend{file~descriptor} 


Benché questa sia alla fine soltanto una convenzione, essa è seguita dalla
totalità delle applicazioni, e non aderirvi potrebbe portare a problemi di
interoperabilità.  Nel caso della shell tutti questi file sono associati al
terminale di controllo, e corrispondono quindi alla lettura della tastiera per
l'ingresso e alla scrittura sul terminale per l'uscita.  Lo standard POSIX.1
provvede, al posto dei valori numerici, tre costanti simboliche, definite in
tab.~\ref{tab:file_std_files}.

\begin{table}[htb]
  \centering
  \footnotesize
  \begin{tabular}[c]{|l|l|}
    \hline
    \textbf{File} & \textbf{Significato} \\
    \hline
    \hline
    \constd{STDIN\_FILENO}  & file descriptor dello \textit{standard input}.\\ 
    \constd{STDOUT\_FILENO} & file descriptor dello \textit{standard output}.\\
    \constd{STDERR\_FILENO} & file descriptor dello \textit{standard error}.\\
    \hline
  \end{tabular}
  \caption{Costanti definite in \headfile{unistd.h} per i file standard.}
  \label{tab:file_std_files}
\end{table}

\itindend{standard~input} 
\itindend{standard~output}
\itindend{standard~error}

In fig.~\ref{fig:file_proc_file} si è rappresentata una situazione diversa
rispetto a quella usuale della shell, in cui tutti e tre questi file fanno
riferimento al terminale su cui si opera. Nell'esempio invece viene illustrata
la situazione di un programma in cui lo \textit{standard input} è associato ad
un file mentre lo \textit{standard output} e lo \textit{standard error} sono
associati ad un altro file.  Si noti poi come per questi ultimi le strutture
\kstruct{file} nella \textit{file table}, pur essendo distinte, fanno
riferimento allo stesso \textit{inode}, dato che il file che è stato aperto lo
stesso. Questo è quello che avviene normalmente quando si apre più volte lo
stesso file.

Si ritrova quindi anche con le voci della \textit{file table} una situazione
analoga di quella delle voci di una directory, con la possibilità di avere più
voci che fanno riferimento allo stesso \textit{inode}. L'analogia è in realtà
molto stretta perché quando si cancella un file, il kernel verifica anche che
non resti nessun riferimento in una qualunque voce della \textit{file table}
prima di liberare le risorse ad esso associate e disallocare il relativo
\textit{inode}.

Nelle vecchie versioni di Unix (ed anche in Linux fino al kernel 2.0.x) il
numero di file aperti era anche soggetto ad un limite massimo dato dalle
dimensioni del vettore di puntatori con cui era realizzata la tabella dei file
descriptor dentro \kstruct{files\_struct}. Questo limite intrinseco nei kernel
più recenti non sussiste più, dato che si è passati da un vettore ad una
lista, ma restano i limiti imposti dall'amministratore (vedi
sez.~\ref{sec:sys_limits}).

\itindend{file~table}


\subsection{Apertura, creazione e chiusura di un file}
\label{sec:file_open_close}

La funzione di sistema \funcd{open} è la principale funzione dell'interfaccia
di gestione dei file, quella che dato un \textit{pathname} consente di
ottenere un file descriptor ``\textsl{aprendo}'' il file
corrispondente,\footnote{è \func{open} che alloca \kstruct{file}, la inserisce
  nella \textit{file table} e crea il riferimento nella
  \kstruct{files\_struct} del processo.} il suo prototipo è:

\begin{funcproto}{
\fhead{sys/types.h}
\fhead{sys/stat.h}
\fhead{fcntl.h}
\fdecl{int open(const char *pathname, int flags)}
\fdecl{int open(const char *pathname, int flags, mode\_t mode)}

\fdesc{Apre un file.} 
}

{La funzione ritorna il file descriptor in caso di successo e $-1$ per un
  errore, nel qual caso \var{errno} assumerà uno dei valori:
  \begin{errlist}
  \item[\errcode{EEXIST}] \param{pathname} esiste e si è specificato
    \const{O\_CREAT} e \const{O\_EXCL}.
  \item[\errcode{EINTR}] la funzione era bloccata ed è stata interrotta da un
    segnale (vedi sez.~\ref{sec:sig_gen_beha}).
  \item[\errcode{EINVAL}] si è usato \const{O\_CREAT} indicando un pathname
    con caratteri non supportati dal filesystem sottostante o si è richiesto
    \const{O\_TMPFILE} senza indicare \const{O\_WRONLY} o \const{O\_RDWR} o si
    è usato \const{O\_DIRECT} per un filesystem che non lo supporta.
  \item[\errcode{EISDIR}] \param{pathname} indica una directory e o si è
    tentato un accesso che prevede la scrittura o si è usato
    \const{O\_TMPFILE} con accesso che prevede la scrittura ma il kernel non
    supporta la funzionalità.
  \item[\errcode{EFBIG}] il file è troppo grande per essere aperto, in genere
    dovuto al fatto che si è compilata una applicazione a 32 bit senza
    abilitare il supporto per le dimensioni a 64 bit; questo è il valore
    restituito fino al kernel 2.6.23, coi successivi viene restituito
    \errcode{EOVERFLOW} come richiesto da POSIX.1.
  \item[\errcode{ELOOP}] si sono incontrati troppi collegamenti simbolici nel
    risolvere \param{pathname} o si è indicato \const{O\_NOFOLLOW} e
    \param{pathname} è un collegamento simbolico (e non si è usato
    \const{O\_PATH}).
  \item[\errcode{ENODEV}] \param{pathname} si riferisce a un file di
    dispositivo che non esiste.
  \item[\errcode{ENOENT}] \param{pathname} non esiste e non si è richiesto
    \const{O\_CREAT}, o non esiste un suo componente, o si riferisce ad una
    directory inesistente, si è usato \const{O\_TMPFILE} con accesso che
    prevede la scrittura ma il kernel non supporta la funzionalità.
  \item[\errcode{ENOTDIR}] si è specificato \const{O\_DIRECTORY} e
    \param{pathname} non è una directory.
  \item[\errcode{ENXIO}] si sono impostati \const{O\_NONBLOCK} o
    \const{O\_WRONLY} ed il file è una \textit{fifo} che non viene letta da
    nessun processo o \param{pathname} è un file di dispositivo ma il
    dispositivo è assente.
  \item[\errcode{EPERM}] si è specificato \const{O\_NOATIME} e non si è né
    amministratori né proprietari del file.
  \item[\errcode{ETXTBSY}] si è cercato di accedere in scrittura all'immagine
    di un programma in esecuzione.
  \item[\errcode{EWOULDBLOCK}] la funzione si sarebbe bloccata ma si è
    richiesto \const{O\_NONBLOCK}.
  \end{errlist}
  ed inoltre \errval{EACCES}, \errval{EDQUOT}, \errval{EFAULT}, \errval{EMFILE},
  \errval{ENAMETOOLONG}, \errval{ENFILE}, \errval{ENOMEM}, \errval{ENOSPC},
  \errval{EROFS}, nel loro significato generico.}
\end{funcproto}

La funzione apre il file indicato da \param{pathname} nella modalità indicata
da \param{flags}. Essa può essere invocata in due modi diversi, specificando
opzionalmente un terzo argomento \param{mode}. Qualora il file non esista e
venga creato, questo argomento consente di indicare quali permessi dovranno
essergli assegnati.\footnote{questo è possibile solo se si è usato in
  \param{flags} uno fra \const{O\_CREATE} e \const{O\_TMPFILE}, in tutti gli
  altri casi sarà ignorato.} I valori possibili sono gli stessi già visti in
sez.~\ref{sec:file_perm_overview} e possono essere specificati come OR binario
delle costanti descritte in tab.~\ref{tab:file_bit_perm}. Questi permessi sono
comunque filtrati dal valore della \textit{umask} (vedi
sez.~\ref{sec:file_perm_management}) del processo.

La funzione restituisce sempre il primo file descriptor libero, una
caratteristica che permette di prevedere qual è il valore del file descriptor
che si otterrà al ritorno di \func{open}, e che viene spesso usata dalle
applicazioni per sostituire i file corrispondenti ai file standard visti in
tab.~\ref{tab:file_std_files}. Se ad esempio si chiude lo \textit{standard
  input} e si apre subito dopo un nuovo file questo diventerà il nuovo
\textit{standard input} dato che avrà il file descriptor 0.

Al momento dell'apertura il nuovo file descriptor non è condiviso con nessun
altro processo (torneremo sul significato della condivisione dei file
descriptor, che in genere si ottiene dopo una \func{fork}, in
sez.~\ref{sec:file_shared_access}) ed è impostato, come accennato in
sez.~\ref{sec:proc_exec}, per restare aperto attraverso una
\func{exec}. Inoltre la posizione sul file, il cosiddetto \textit{offset}, è
impostata all'inizio del file. Una volta aperto un file si potrà operare su di
esso direttamente tramite il file descriptor, e quanto avviene al
\textit{pathname} con cui lo si è aperto sarà del tutto ininfluente.

\itindbeg{file~status~flags}

Il comportamento della funzione, e le diverse modalità con cui può essere
aperto il file, vengono controllati dall'argomento \param{flags} il cui valore
deve essere indicato come maschera binaria in cui ciascun bit ha un
significato specifico.  Alcuni di questi bit vanno anche a costituire i
cosiddetti \textit{file status flags} (i \textsl{flag di stato} del file), che
vengono mantenuti nel campo \var{f\_flags} della struttura \kstruct{file} che
abbiamo riportato anche in fig.~\ref{fig:file_proc_file}).

Ciascun flag viene identificato da una apposita costante, ed il valore
di \param{flags} deve essere specificato come OR aritmetico di queste
costanti. Inoltre per evitare problemi di compatibilità con funzionalità che
non sono previste o non ancora supportate in versioni meno recenti del kernel,
la \func{open} di Linux ignora i flag che non riconosce, pertanto
l'indicazione di un flag inesistente non provoca una condizione di errore.

I vari bit che si possono usare come componenti di \param{flags} sono divisi
in tre gruppi principali. Il primo gruppo è quello dei cosiddetti flag delle
\textsl{modalità di accesso} (o \textit{access mode flags}), che specificano
che tipo di accesso si effettuerà sul file, fra lettura, scrittura e
lettura/scrittura. Questa modalità deve essere indicata usando una delle
costanti di tab.~\ref{tab:open_access_mode_flag}.

\begin{table}[htb]
  \centering
  \footnotesize
    \begin{tabular}[c]{|l|l|}
      \hline
      \textbf{Flag} & \textbf{Significato} \\
      \hline
      \hline
      \constd{O\_RDONLY} & Apre il file in sola lettura.\\
      \constd{O\_WRONLY} & Apre il file in sola scrittura.\\
      \constd{O\_RDWR}   & Apre il file sia in lettura che in scrittura.\\
      \hline
    \end{tabular}
    \caption{Le tre costanti che identificano le modalità di accesso
      nell'apertura di un file.}
  \label{tab:open_access_mode_flag}
\end{table}

A differenza di tutti gli altri flag che vedremo in seguito, in questo caso
non si ha a che fare con singoli bit separati dell'argomento \param{flags}, ma
con un numero composto da due bit. Questo significa ad esempio che la
combinazione \code{\const{O\_RDONLY}|\const{O\_WRONLY}} non è affatto
equivalente a \const{O\_RDWR}, e non deve essere usata.\footnote{in realtà
  su Linux, dove i valori per le tre costanti di
  tab.~\ref{tab:open_access_mode_flag} sono rispettivamente $0$, $1$ e $2$, il
  valore $3$ viene usato con un significato speciale, ed assolutamente fuori
  standard, disponibile solo per i file di dispositivo e solo per alcuni
  driver, in cui si richiede la verifica della capacità di accesso in lettura
  e scrittura ma viene restituito un file descriptor che non può essere letto
  o scritto, ma solo usato con una \func{ioctl} (vedi
  sez.~\ref{sec:file_fcntl_ioctl}).}

La modalità di accesso deve sempre essere specificata quando si apre un file,
il valore indicato in \param{flags} viene salvato nei \textit{file status
  flags}, e può essere riletto con \func{fcntl} (vedi
sez.~\ref{sec:file_fcntl_ioctl}), il relativo valore può essere poi ottenuto
un AND aritmetico della maschera binaria \constd{O\_ACCMODE}, ma non può essere
modificato. Nella \acr{glibc} sono definite inoltre \constd{O\_READ} come
sinonimo di \const{O\_RDONLY} e \constd{O\_WRITE} come sinonimo di
\const{O\_WRONLY}.\footnote{si tratta di definizioni completamente fuori
  standard, attinenti, insieme a \constd{O\_EXEC} che permetterebbe l'apertura
  di un file per l'esecuzione, ad un non meglio precisato ``\textit{GNU
    system}''; pur essendo equivalenti alle definizioni classiche non è
  comunque il caso di utilizzarle.}

\itindend{file~status~flags}

Il secondo gruppo di flag è quello delle \textsl{modalità di
  apertura},\footnote{la pagina di manuale di \func{open} parla di
  \textit{file creation flags}, ma alcuni di questi flag non hanno nulla a che
  fare con la creazione dei file, mentre il manuale dalla \acr{glibc} parla di
  più correttamente di \textit{open-time flags}, dato che si tratta di flag il
  cui significato ha senso solo al momento dell'apertura del file.} che
permettono di specificare alcune delle caratteristiche del comportamento di
\func{open} nel momento in viene eseguita per aprire un file. Questi flag
hanno effetto solo nella chiamata della funzione, non sono memorizzati fra i
\textit{file status flags} e non possono essere riletti da \func{fcntl} (vedi
sez.~\ref{sec:file_fcntl_ioctl}).

\begin{table}[htb]
  \centering
  \footnotesize
    \begin{tabular}[c]{|l|p{10 cm}|}
      \hline
      \textbf{Flag} & \textbf{Significato} \\
      \hline
      \hline
      \constd{O\_CREAT}  & Se il file non esiste verrà creato, con le regole
                           di titolarità del file viste in
                           sez.~\ref{sec:file_ownership_management}. Se si
                           imposta questo flag l'argomento \param{mode} deve
                           essere sempre specificato.\\  
      \constd{O\_DIRECTORY}& Se \param{pathname} non è una directory la
                             chiamata fallisce. Questo flag, introdotto con il
                             kernel 2.1.126, è specifico di Linux e
                             serve ad evitare dei possibili
                             \itindex{Denial~of~Service~(DoS)}
                             \textit{DoS}\footnotemark quando \func{opendir} 
                             viene chiamata su una \textit{fifo} o su un
                             dispositivo associato ad una unità a nastri. Non
                             viene usato al di fuori dell'implementazione di
                             \func{opendir}, ed è utilizzabile soltanto se si è
                             definita la macro \macro{\_GNU\_SOURCE}.\\
      \constd{O\_EXCL}   & Deve essere usato in congiunzione con
                           \const{O\_CREAT} ed in tal caso impone che il file
                           indicato da \param{pathname} non sia già esistente
                           (altrimenti causa il fallimento della chiamata con
                           un errore di \errcode{EEXIST}).\\
      \constd{O\_LARGEFILE}& Viene usato sui sistemi a 32 bit per richiedere
                             l'apertura di file molto grandi, la cui
                             dimensione non è rappresentabile con la versione a
                             32 bit del tipo \type{off\_t}, utilizzando
                             l'interfaccia alternativa abilitata con la
                             macro \macro{\_LARGEFILE64\_SOURCE}. Come
                             illustrato in sez.~\ref{sec:intro_gcc_glibc_std} è
                             sempre preferibile usare la conversione automatica
                             delle funzioni che si attiva assegnando a $64$ la
                             macro \macro{\_FILE\_OFFSET\_BITS}, e non usare mai
                             questo flag.\\
      \constd{O\_NOCTTY} & Se \param{pathname} si riferisce ad un dispositivo
                           di terminale, questo non diventerà il terminale di
                           controllo, anche se il processo non ne ha ancora
                           uno (si veda sez.~\ref{sec:sess_ctrl_term}).\\ 
      \constd{O\_NOFOLLOW}& Se \param{pathname} è un collegamento simbolico
                            la chiamata fallisce. Questa è un'estensione BSD
                            aggiunta in Linux a partire dal kernel
                            2.1.126, ed utilizzabile soltanto se si è definita
                            la macro \macro{\_GNU\_SOURCE}.\\
      \const{O\_TMPFILE} & Consente di creare un file temporaneo anonimo, non
                           visibile con un pathname sul filesystem, ma
                           leggibile e scrivibile all'interno del processo.
                           Introdotto con il kernel 3.11, è specifico di
                           Linux.\\ 
      \constd{O\_TRUNC}  & Se usato su un file di dati aperto in scrittura,
                           ne tronca la lunghezza a zero; con un terminale o
                           una \textit{fifo} viene ignorato, negli altri casi
                           il comportamento non è specificato.\\ 
      \hline
    \end{tabular}
    \caption{Le costanti che identificano le \textit{modalità di apertura} di
      un file.} 
  \label{tab:open_time_flag}
\end{table}

\footnotetext{acronimo di \itindex{Denial~of~Service~(DoS)} \textit{Denial of
    Service}, si chiamano così attacchi miranti ad impedire un servizio
  causando una qualche forma di carico eccessivo per il sistema, che resta
  bloccato nelle risposte all'attacco.}

Si è riportato in tab.~\ref{tab:open_time_flag} l'elenco dei flag delle
\textsl{modalità di apertura}.\footnote{la \acr{glibc} definisce anche i due
  flag \constd{O\_SHLOCK}, che aprirebbe il file con uno \textit{shared lock}
  e \constd{O\_EXLOCK} che lo aprirebbe con un \textit{exclusive lock} (vedi
  sez.~\ref{sec:file_locking}), si tratta di opzioni specifiche di BSD, che non
  esistono con Linux.}  Uno di questi, \const{O\_EXCL}, ha senso solo se usato
in combinazione a \const{O\_CREAT} quando si vuole creare un nuovo file per
assicurarsi che questo non esista di già, e lo si usa spesso per creare i
cosiddetti ``\textsl{file di lock}'' (vedi sez.~\ref{sec:ipc_file_lock}).

Si tenga presente che questa opzione è supportata su NFS solo a partire da
NFSv3 e con il kernel 2.6, nelle versioni precedenti la funzionalità viene
emulata controllando prima l'esistenza del file per cui usarla per creare un
file di lock potrebbe dar luogo a una \textit{race condition}, in tal caso
infatti un file potrebbe venir creato fra il controllo la successiva apertura
con \const{O\_CREAT}; la cosa si può risolvere comunque creando un file con un
nome univoco ed usando la funzione \func{link} per creare il file di lock,
(vedi sez.~\ref{sec:ipc_file_lock}).

Se si usa \const{O\_EXCL} senza \const{O\_CREAT} il comportamento è
indefinito, escluso il caso in cui viene usato con \const{O\_TMPFILE} su cui
torneremo più avanti; un'altra eccezione è quello in cui lo si usa da solo su
un file di dispositivo, in quel caso se questo è in uso (ad esempio se è
relativo ad un filesystem che si è montato) si avrà un errore di
\errval{EBUSY}, e pertanto può essere usato in questa modalità per rilevare lo
stato del dispositivo.

Nella creazione di un file con \const{O\_CREAT} occorre sempre specificare
l'argomento di \param{mode}, che altrimenti è ignorato. Si tenga presente che
indipendentemente dai permessi che si possono assegnare, che in seguito
potrebbero non consentire lettura o scrittura, quando il file viene aperto
l'accesso viene garantito secondo quanto richiesto con i flag di
tab.~\ref{tab:open_access_mode_flag}.  Quando viene creato un nuovo file
\const{O\_CREAT} con tutti e tre i tempi del file di
tab.~\ref{tab:file_file_times} vengono impostati al tempo corrente. Se invece
si tronca il file con \const{O\_TRUNC} verranno impostati soltanto il
\textit{modification time} e lo \textit{status change time}.

Il flag \constd{O\_TMPFILE}, introdotto con il kernel
3.11,\footnote{inizialmente solo su alcuni filesystem (i vari \acr{extN},
  \acr{Minix}, \acr{UDF}, \acr{shmem}) poi progressivamente esteso ad altri
  (\acr{XFS} con il 3.15, \acr{Btrfs} e \acr{F2FS} con il 3.16, \acr{ubifs}
  con il 4.9).}  consente di aprire un file temporaneo senza che questo venga
associato ad un nome e compaia nel filesystem. In questo caso la funzione
restituirà un file descriptor da poter utilizzare per leggere e scrivere dati,
ma il contenuto dell'argomento \param{path} verrà usato solamente per
determinare, in base alla directory su cui si verrebbe a trovare il
\textit{pathname} indicato, il filesystem all'interno del quale deve essere
allocato l'\textit{inode} e lo spazio disco usato dal file
descriptor. L'\textit{inode} resterà anonimo e l'unico riferimento esistente
sarà quello contenuto nella \textit{file table} del processo che ha chiamato
\func{open}.

Lo scopo principale del flag è quello fornire una modalità atomica, semplice e
sicura per applicare la tecnica della creazione di un file temporaneo seguita
dalla sua immediata cancellazione con \func{unlink} per non lasciare rimasugli
sul filesystem, di cui è parlato in sez.~\ref{sec:link_symlink_rename}.
Inoltre, dato che il file non compare nel filesystem, si evitano alla radice
tutti gli eventuali problemi di \textit{race condition} o \textit{symlink
  attack} e non ci si deve neanche preoccupare di ottenere un opportuno nome
univoco con l'uso delle funzioni di sez.~\ref{sec:file_temp_file}.

Una volta aperto il file vi si potrà leggere o scrivere a seconda che siano
utilizzati \const{O\_RDWR} o \const{O\_WRONLY}, mentre l'uso di
\func{O\_RDONLY} non è consentito, non avendo molto senso ottenere un file
descriptor su un file che nasce vuoto per cui non si potrà comunque leggere
nulla. L'unico altro flag che può essere utilizzato insieme a
\const{O\_TMPFILE} è \const{O\_EXCL}, che in questo caso assume però un
significato diverso da quello ordinario, dato che in questo caso il file
associato al file descriptor non esiste comunque.

L'uso di \const{O\_EXCL} attiene infatti all'altro possibile impiego di
\const{O\_TMPFILE} oltre a quello citato della creazione sicura di un file
temporaneo come sostituto sicuro di \func{tmpfile}: la possibilità di creare
un contenuto iniziale per un file ed impostarne permessi, proprietario e
attributi estesi con \func{fchmod}, \func{fchown} e \func{fsetxattr}, senza
possibilità di \textit{race condition} ed interferenze esterne, per poi far
apparire il tutto sul filesystem in un secondo tempo utilizzando \func{linkat}
sul file descriptor (torneremo su questo in sez.~\ref{sec:file_openat}) per
dargli un nome. Questa operazione però non sarà possibile se si è usato
\const{O\_EXCL}, che in questo caso viene ad assumere il significato di
escludere la possibilità di far esistere il file anche in un secondo tempo.

% NOTE: per O_TMPFILE vedi: http://kernelnewbies.org/Linux_3.11
% https://lwn.net/Articles/558598/ http://lwn.net/Articles/619146/


\begin{table}[!htb]
  \centering
  \footnotesize
    \begin{tabular}[c]{|l|p{10 cm}|}
      \hline
      \textbf{Flag} & \textbf{Significato} \\
      \hline
      \hline
      \constd{O\_APPEND} & Il file viene aperto in \textit{append mode}. La
                           posizione sul file (vedi sez.~\ref{sec:file_lseek})
                           viene sempre mantenuta sulla sua coda, per cui
                           quanto si scrive viene sempre aggiunto al contenuto
                           precedente. Con NFS questa funzionalità non è
                           supportata e viene emulata, per questo possono
                           verificarsi \textit{race condition} con una
                           sovrapposizione dei dati se più di un processo
                           scrive allo stesso tempo.\\ 
      \constd{O\_ASYNC}  & Apre il file per l'I/O in modalità asincrona (vedi
                           sez.~\ref{sec:signal_driven_io}). Quando è
                           impostato viene generato il segnale \signal{SIGIO}
                           tutte le volte che il file è pronto per le
                           operazioni di lettura o scrittura. Questo flag si
                           può usare solo terminali, pseudo-terminali e socket
                           e, a partire dal kernel 2.6, anche sulle
                           \textit{fifo}. Per un bug dell'implementazione non
                           è opportuno usarlo in fase di apertura del file,
                           deve invece essere attivato successivamente con
                           \func{fcntl}.\\
      \constd{O\_CLOEXEC}& Attiva la modalità di \textit{close-on-exec} (vedi
                           sez.~\ref{sec:proc_exec}) sul file. Il flag è 
                           previsto dallo standard POSIX.1-2008, ed è stato
                           introdotto con il kernel 2.6.23 per evitare una
                           \textit{race condition} che si potrebbe verificare
                           con i \textit{thread} fra l'apertura del file e
                           l'impostazione della suddetta modalità con
                           \func{fcntl} (vedi
                           sez.~\ref{sec:file_fcntl_ioctl}).\\ 
      \const{O\_DIRECT}  & Esegue l'I/O direttamente dalla memoria in
                           \textit{user space} in maniera sincrona, in modo da
                           scavalcare i meccanismi di bufferizzazione del
                           kernel. Introdotto con il kernel 2.4.10 ed
                           utilizzabile soltanto se si è definita la 
                           macro \macro{\_GNU\_SOURCE}.\\ 
      \constd{O\_NOATIME}& Blocca l'aggiornamento dei tempi di accesso dei
                           file (vedi sez.~\ref{sec:file_file_times}). Per
                           molti filesystem questa funzionalità non è
                           disponibile per il singolo file ma come opzione
                           generale da specificare in fase di
                           montaggio. Introdotto con il kernel 2.6.8 ed 
                           utilizzabile soltanto se si è definita la 
                           macro \macro{\_GNU\_SOURCE}.\\ 
      \constd{O\_NONBLOCK}& Apre il file in \textsl{modalità non bloccante} per
                            le operazioni di I/O (vedi
                            sez.~\ref{sec:file_noblocking}). Questo significa
                            il fallimento delle successive operazioni di
                            lettura o scrittura qualora il file non sia pronto
                            per la loro esecuzione immediata, invece del 
                            blocco delle stesse in attesa di una successiva
                            possibilità di esecuzione come avviene
                            normalmente. Questa modalità ha senso solo per le
                            \textit{fifo}, vedi sez.~\ref{sec:ipc_named_pipe}),
                            o quando si vuole aprire un file di dispositivo
                            per eseguire una \func{ioctl} (vedi
                            sez.~\ref{sec:file_fcntl_ioctl}).\\ 
      \constd{O\_NDELAY} & In Linux è un sinonimo di \const{O\_NONBLOCK}, ma
                           origina da SVr4, dove però causava il ritorno da
                           una \func{read} con un valore nullo e non con un
                           errore, questo introduce un'ambiguità, dato che
                           come vedremo in sez.~\ref{sec:file_read} il ritorno
                           di un valore nullo da parte di \func{read} ha 
                           il significato di una \textit{end-of-file}.\\
      \const{O\_PATH}    & Ottiene un file descriptor io cui uso è limitato
                           all'indicare una posizione sul filesystem o
                           eseguire operazioni che operano solo a livello del
                           file descriptor (e non di accesso al contenuto del
                           file). Introdotto con il kernel 2.6.39, è specifico
                           di Linux.\\
      \constd{O\_SYNC}   & Apre il file per l'input/output sincrono. Ogni
                           scrittura si bloccherà fino alla conferma
                           dell'arrivo di tutti i dati e di tutti i metadati
                           sull'hardware sottostante (in questo significato
                           solo dal kernel 2.6.33).\\
      \constd{O\_DSYNC}  & Apre il file per l'input/output sincrono. Ogni
                           scrittura di dati si bloccherà fino alla conferma
                           dell'arrivo degli stessi e della parte di metadati
                           ad essi relativa sull'hardware sottostante (in
                           questo significato solo dal kernel 2.6.33).\\
      \hline
    \end{tabular}
    \caption{Le costanti che identificano le \textit{modalità di operazione} di
      un file.} 
  \label{tab:open_operation_flag}
\end{table}

Il terzo gruppo è quello dei flag delle \textsl{modalità di operazione},
riportati in tab.~\ref{tab:open_operation_flag}, che permettono di specificare
varie caratteristiche del comportamento delle operazioni di I/O che verranno
eseguite sul file o le modalità in cui lo si potrà utilizzare. Tutti questi,
tranne \const{O\_CLOEXEC}, che viene mantenuto per ogni singolo file
descriptor, vengono salvati nel campo \var{f\_flags} della struttura
\kstruct{file} insieme al valore della \textsl{modalità di accesso}, andando
far parte dei \textit{file status flags}. Il loro valore viene impostato alla
chiamata di \func{open}, ma possono venire riletti in un secondo tempo con
\func{fcntl}, inoltre alcuni di essi possono anche essere modificati tramite
questa funzione, con conseguente effetto sulle caratteristiche operative che
controllano (torneremo sull'argomento in sez.~\ref{sec:file_fcntl_ioctl}).

Il flag \const{O\_ASYNC} (che, per compatibilità con BSD, si può indicare
anche con la costante \constd{FASYNC}) è definito come possibile valore per
\func{open}, ma per un bug dell'implementazione,\footnote{segnalato come
  ancora presente nella pagina di manuale, almeno fino al novembre 2018.} non
solo non attiva il comportamento citato, ma se usato richiede di essere
esplicitamente disattivato prima di essere attivato in maniera effettiva con
l'uso di \func{fcntl}. Per questo motivo, non essendovi nessuna necessità
specifica di definirlo in fase di apertura del file, è sempre opportuno
attivarlo in un secondo tempo con \func{fcntl} (vedi
sez.~\ref{sec:file_fcntl_ioctl}).

Il flag \constd{O\_DIRECT} non è previsto da nessuno standard, anche se è
presente in alcuni kernel unix-like.\footnote{il flag è stato introdotto dalla
  SGI in IRIX, ma è presente senza limiti di allineamento dei buffer anche in
  FreeBSD.} Per i kernel della serie 2.4 si deve garantire che i buffer in
\textit{user space} da cui si effettua il trasferimento diretto dei dati siano
allineati alle dimensioni dei blocchi del filesystem. Con il kernel 2.6 in
genere basta che siano allineati a multipli di 512 byte, ma le restrizioni
possono variare a seconda del filesystem, ed inoltre su alcuni filesystem
questo flag può non essere supportato, nel qual caso si avrà un errore di
\errval{EINVAL}.

Lo scopo di \const{O\_DIRECT} è consentire un completo controllo sulla
bufferizzazione dei propri dati per quelle applicazioni (in genere database)
che hanno esigenze specifiche che non vengono soddisfatte nella maniera più
efficiente dalla politica generica utilizzata dal kernel. In genere l'uso di
questo flag peggiora le prestazioni tranne quando le applicazioni sono in
grado di ottimizzare la propria bufferizzazione in maniera adeguata. Se lo si
usa si deve avere cura di non mescolare questo tipo di accesso con quello
ordinario, in quante le esigenze di mantenere coerenti i dati porterebbero ad
un peggioramento delle prestazioni. Lo stesso dicasi per l'interazione con
eventuale mappatura in memoria del file (vedi sez.~\ref{sec:file_memory_map}).

Si tenga presente infine che anche se l'uso di \const{O\_DIRECT} comporta una
scrittura sincrona dei dati dei buffer in \textit{user space}, questo non è
completamente equivalente all'uso di \const{O\_SYNC} che garantisce anche
sulla scrittura sincrona dei metadati associati alla scrittura dei dati del
file.\footnote{la situazione si complica ulteriormente per NFS, in cui l'uso
  del flag disabilita la bufferizzazione solo dal lato del client, e può
  causare problemi di prestazioni.} Per questo in genere se si usa
\const{O\_DIRECT} è opportuno richiedere anche \const{O\_SYNC}.

Si tenga presente infine che la implementazione di \const{O\_SYNC} di Linux
differisce da quanto previsto dallo standard POSIX.1 che prevede, oltre a
questo flag che dovrebbe indicare la sincronizzazione completa di tutti i dati
e di tutti i metadati, altri due flag \const{O\_DSYNC} e \const{O\_RSYNC}. 

Il primo dei due richiede la scrittura sincrona di tutti i dati del file e dei
metadati che ne consentono l'immediata rilettura, ma non di tutti i metadati,
per evitare la perdita di prestazioni relativa alla sincronizzazione di
informazioni ausiliarie come i tempi dei file.  Il secondo, da usare in
combinazione con \const{O\_SYNC} o \const{O\_DSYNC} ne sospende l'effetto,
consentendo al kernel di bufferizzare le scritture, ma soltanto finché non
avviene una lettura, in quel caso i dati ed i metadati dovranno essere
sincronizzati immediatamente (secondo le modalità indicate da \const{O\_SYNC}
e \const{O\_DSYNC}) e la lettura verrà bloccata fintanto che detta
sincronizzazione non sia completata.

Nel caso di Linux, fino al kernel 2.6.33, esisteva solo \const{O\_SYNC}, ma
con il comportamento previsto dallo standard per \const{O\_DSYNC}, e sia
questo che \const{O\_RSYNC} erano definiti (fin dal kernel 2.1.130) come
sinonimi di \const{O\_SYNC}.  Con il kernel 2.6.33 il significato di
\const{O\_SYNC} è diventato quello dello standard, ma gli è stato assegnato un
valore diverso, mantenendo quello originario, con il comportamento
corrispondete, per \const{O\_DSYNC} in modo che applicazioni compilate con
versioni precedenti delle librerie e del kernel non trovassero un
comportamento diverso.  Inoltre il nuovo \const{O\_SYNC} è stato definito in
maniera opportuna in modo che su versioni del kernel precedenti la 2.6.33
torni a corrispondere al valore di \const{O\_DSYNC}.

% NOTE: per le differenze fra O_DSYNC, O_SYNC e O_RSYNC introdotte nella  
% nello sviluppo del kernel 2.6.33, vedi http://lwn.net/Articles/350219/ 

Il flag \constd{O\_PATH},\label{open_o_path_flag} introdotto con il kernel
2.6.39, viene usato per limitare l'uso del file descriptor restituito da
\func{open} o all'identificazione di una posizione sul filesystem (ad uso
delle \textit{at-functions} che tratteremo in sez.~\ref{sec:file_openat}) o
alle operazioni che riguardano il file descriptor in quanto tale, senza
consentire operazioni sul file; in sostanza se si apre un file con
\const{O\_PATH} si potrà soltanto:
\begin{itemize*}
\item usare il file descriptor come indicatore della directory di partenza con
  una delle \textit{at-functions} (vedi sez.~\ref{sec:file_openat});
\item cambiare directory di lavoro con \func{fchdir} se il file descriptor fa
  riferimento a una directory (dal kernel 3.5);
\item usare le funzioni che duplicano il file descriptor (vedi
  sez.~\ref{sec:file_dup});
\item passare il file descriptor ad un altro processo usando un socket
  \const{PF\_UNIX} (vedi sez.~\ref{sec:unix_socket})
\item ottenere le informazioni relative al file con \func{fstat} (dal kernel
  3.6) o al filesystem con \func{fstatfs} (dal kernel 3.12);
\item ottenere il valore dei \textit{file descriptor flags} (fra cui comparirà
  anche lo stesso \const{O\_PATH}) o impostare o leggere i \textit{file status
    flags} con \func{fcntl} (rispettivamente con le operazioni
  \const{F\_GETFL}, \const{F\_SETFD} e \const{F\_GETFD}, vedi
  sez.~\ref{sec:file_fcntl_ioctl}).
\item chiudere il file con \func{close}.
\end{itemize*}

In realtà usando \const{O\_PATH} il file non viene effettivamente aperto, per
cui ogni tentativo di usare il file descriptor così ottenuto con funzioni che
operano effettivamente sul file (come ad esempio \func{read}, \func{write},
\func{fchown}, \func{fchmod}, \func{ioctl}, ecc.) fallirà con un errore di
\errval{EBADF}, come se questo non fosse un file descriptor valido. Per questo
motivo usando questo flag non è necessario avere nessun permesso per aprire un
file, neanche quello di lettura (ma occorre ovviamente avere il permesso di
esecuzione per le directory sovrastanti).

Questo consente di usare il file descriptor con funzioni che non richiedono
permessi sul file, come \func{fstat}, laddove un'apertura con
\const{O\_RDONLY} sarebbe fallita. I permessi verranno eventualmente
controllati, se necessario, nelle operazioni seguenti, ad esempio per usare
\func{fchdir} con il file descriptor (se questo fa riferimento ad una
directory) occorrerà avere il permesso di esecuzione.

Se si usa \const{O\_PATH} tutti gli altri flag eccettuati \const{O\_CLOEXEC},
\const{O\_DIRECTORY} e \const{O\_NOFOLLOW} verranno ignorati. I primi due
mantengono il loro significato usuale, mentre \const{O\_NOFOLLOW} fa si che se
il file indicato è un un link simbolico venga aperto quest'ultimo (cambiando
quindi il comportamento ordinario che prova il fallimento della chiamata),
così da poter usare il file descriptor ottenuto per le funzioni
\func{fchownat}, \func{fstatat}, \func{linkat} e \func{readlinkat} che ne
supportano l'uso come come primo argomento (torneremo su questo in
sez.~\ref{sec:file_openat}).

Nelle prime versioni di Unix i valori di \param{flag} specificabili per
\func{open} erano solo quelli relativi alle modalità di accesso del file.  Per
questo motivo per creare un nuovo file c'era una \textit{system call}
apposita, \funcd{creat}, nel caso di Linux questo non è più necessario ma la
funzione è definita ugualmente; il suo prototipo è:

\begin{funcproto}{
\fhead{fcntl.h}
\fdecl{int creat(const char *pathname, mode\_t mode)}
\fdesc{Crea un nuovo file vuoto.} 
}

{La funzione ritorna $0$ in caso di successo e $-1$ per un errore, nel qual
  caso \var{errno} assumerà gli stessi valori che si otterrebbero con
  \func{open}.}
\end{funcproto}

La funzione crea un nuovo file vuoto, con i permessi specificati
dall'argomento \param{mode}. È del tutto equivalente a \code{open(filedes,
  O\_CREAT|O\_WRONLY|O\_TRUNC, mode)} e resta solo per compatibilità con i
vecchi programmi.

Una volta che l'accesso ad un file non sia più necessario la funzione di
sistema \funcd{close} permette di ``\textsl{chiuderlo}'', in questo modo il
file non sarà più accessibile ed il relativo file descriptor ritornerà
disponibile; il suo prototipo è:

\begin{funcproto}{
\fhead{unistd.h}
\fdecl{int close(int fd)}
\fdesc{Chiude un file.} 
}

{La funzione ritorna $0$ in caso di successo e $-1$ per un errore, nel qual
  caso \var{errno} assumerà uno dei valori: 
  \begin{errlist}
    \item[\errcode{EBADF}]  \param{fd} non è un descrittore valido.
    \item[\errcode{EINTR}] la funzione è stata interrotta da un segnale.
  \end{errlist}
  ed inoltre \errval{EIO} nel suo significato generico.}
\end{funcproto}

La funzione chiude il file descriptor \param{fd}. La chiusura rilascia ogni
eventuale blocco (il \textit{file locking} è trattato in
sez.~\ref{sec:file_locking}) che il processo poteva avere acquisito su di
esso. Se \param{fd} è l'ultimo riferimento (di eventuali copie, vedi
sez.~\ref{sec:file_shared_access} e \ref{sec:file_dup}) ad un file aperto,
tutte le risorse nella \textit{file table} vengono rilasciate. Infine se il
file descriptor era l'ultimo riferimento ad un file su disco quest'ultimo
viene cancellato.

Si ricordi che quando un processo termina tutti i suoi file descriptor vengono
automaticamente chiusi, molti programmi sfruttano questa caratteristica e non
usano esplicitamente \func{close}. In genere comunque chiudere un file senza
controllare lo stato di uscita di \func{close} un è errore; molti filesystem
infatti implementano la tecnica del cosiddetto \itindex{write-behind}
\textit{write-behind}, per cui una \func{write} può avere successo anche se i
dati non sono stati effettivamente scritti su disco. In questo caso un
eventuale errore di I/O avvenuto in un secondo tempo potrebbe sfuggire, mentre
verrebbe riportato alla chiusura esplicita del file. Per questo motivo non
effettuare il controllo può portare ad una perdita di dati
inavvertita.\footnote{in Linux questo comportamento è stato osservato con NFS
  e le quote su disco.}

In ogni caso una \func{close} andata a buon fine non garantisce che i dati
siano stati effettivamente scritti su disco, perché il kernel può decidere di
ottimizzare l'accesso a disco ritardandone la scrittura. L'uso della funzione
\func{sync} (vedi sez.~\ref{sec:file_sync}) effettua esplicitamente lo scarico
dei dati, ma anche in questo caso resta l'incertezza dovuta al comportamento
dell'hardware, che a sua volta può introdurre ottimizzazioni dell'accesso al
disco che ritardano la scrittura dei dati. Da questo deriva l'abitudine di
alcuni sistemisti di ripetere tre volte il comando omonimo prima di eseguire
lo shutdown di una macchina.

Si tenga comunque presente che ripetere la chiusura in caso di fallimento non
è opportuno, una volta chiamata \func{close} il file descriptor viene comunque
rilasciato, indipendentemente dalla presenza di errori, e se la riesecuzione
non comporta teoricamente problemi (se non la sua inutilità) se fatta
all'interno di un processo singolo, nel caso si usino i \textit{thread} si
potrebbe chiudere un file descriptor aperto nel contempo da un altro
\textit{thread}.

\subsection{La gestione della posizione nel file}
\label{sec:file_lseek}

Come già accennato in sez.~\ref{sec:file_fd} a ciascun file aperto è associata
una \textsl{posizione corrente nel file} (il cosiddetto \textit{file offset},
mantenuto nel campo \var{f\_pos} di \kstruct{file}) espressa da un numero
intero positivo che esprime il numero di byte dall'inizio del file. Tutte le
operazioni di lettura e scrittura avvengono a partire da questa posizione che
viene automaticamente spostata in avanti del numero di byte letti o scritti.

In genere, a meno di non avere richiesto la modalità di scrittura in
\textit{append} (vedi sez.~\ref{sec:file_write}) con \const{O\_APPEND}, questa
posizione viene impostata a zero all'apertura del file. È possibile impostarla
ad un valore qualsiasi con la funzione di sistema \funcd{lseek}, il cui
prototipo è:

\begin{funcproto}{
\fhead{sys/types.h}
\fhead{unistd.h}
\fdecl{off\_t lseek(int fd, off\_t offset, int whence)}
\fdesc{Imposta la posizione sul file.} 
}

{La funzione ritorna il valore della posizione sul file in caso di successo e
  $-1$ per un errore, nel qual caso \var{errno} assumerà uno dei valori:
  \begin{errlist}
    \item[\errcode{EINVAL}] \param{whence} non è un valore valido.
    \item[\errcode{EOVERFLOW}] \param{offset} non può essere rappresentato nel
      tipo \type{off\_t}.
    \item[\errcode{ESPIPE}] \param{fd} è una \textit{pipe}, un socket o una
      \textit{fifo}.
  \end{errlist}
  ed inoltre \errval{EBADF} nel suo significato generico.}
\end{funcproto}

La funzione imposta la nuova posizione sul file usando il valore indicato
da \param{offset}, che viene sommato al riferimento dato
dall'argomento \param{whence}, che deve essere indicato con una delle costanti
riportate in tab.~\ref{tab:lseek_whence_values}.\footnote{per compatibilità
  con alcune vecchie notazioni questi valori possono essere rimpiazzati
  rispettivamente con 0, 1 e 2 o con \constd{L\_SET}, \constd{L\_INCR} e
  \constd{L\_XTND}.} Si tenga presente che la chiamata a \func{lseek} non causa
nessun accesso al file, si limita a modificare la posizione corrente (cioè il
campo \var{f\_pos} della struttura \kstruct{file}, vedi
fig.~\ref{fig:file_proc_file}).  Dato che la funzione ritorna la nuova
posizione, usando il valore zero per \param{offset} si può riottenere la
posizione corrente nel file con \code{lseek(fd, 0, SEEK\_CUR)}.

\begin{table}[htb]
  \centering
  \footnotesize
  \begin{tabular}[c]{|l|p{10cm}|}
    \hline
    \textbf{Costante} & \textbf{Significato} \\
    \hline
    \hline
    \constd{SEEK\_SET}& Si fa riferimento all'inizio del file: il valore, che 
                        deve essere positivo, di \param{offset} indica
                        direttamente la nuova posizione corrente.\\
    \constd{SEEK\_CUR}& Si fa riferimento alla posizione corrente del file:
                        ad essa viene sommato \param{offset}, che può essere
                        negativo e positivo, per ottenere la nuova posizione
                        corrente.\\
    \constd{SEEK\_END}& Si fa riferimento alla fine del file: alle dimensioni
                        del file viene sommato \param{offset}, che può essere
                        negativo e positivo, per ottenere la nuova posizione
                        corrente.\\
    \hline
    \constd{SEEK\_DATA}&Sposta la posizione nel file sull'inizio del primo
                        blocco di dati dopo un \textit{hole} che segue (o
                        coincide) con la posizione indicata da \param{offset}
                        (dal kernel 3.1).\\
    \constd{SEEK\_HOLE}&Sposta la posizione sul file all'inizio del primo
                        \textit{hole} nel file che segue o inizia
                        con \param{offset}, oppure si porta su \param{offset} 
                        se questo è all'interno di un \textit{hole}, oppure si
                        porta alla fine del file se non ci sono \textit{hole}
                        dopo \param{offset} (dal kernel 3.1).\\ 
    \hline
  \end{tabular}  
  \caption{Possibili valori per l'argomento \param{whence} di \func{lseek}.} 
  \label{tab:lseek_whence_values}
\end{table}


% NOTE: per SEEK_HOLE e SEEK_DATA, inclusi nel kernel 3.1, vedi
% http://lwn.net/Articles/439623/ 

Si tenga presente inoltre che usare \const{SEEK\_END} non assicura affatto che
la successiva scrittura avvenga alla fine del file, infatti se questo è stato
aperto anche da un altro processo che vi ha scritto, la fine del file può
essersi spostata, ma noi scriveremo alla posizione impostata in precedenza
(questa è una potenziale sorgente di \textit{race condition}, vedi
sez.~\ref{sec:file_shared_access}).

Non tutti i file supportano la capacità di eseguire una \func{lseek}, in
questo caso la funzione ritorna l'errore \errcode{ESPIPE}. Questo, oltre che
per i tre casi citati nel prototipo, vale anche per tutti quei dispositivi che
non supportano questa funzione, come ad esempio per i file di
terminale.\footnote{altri sistemi, usando \const{SEEK\_SET}, in questo caso
  ritornano il numero di caratteri che vi sono stati scritti.} Lo standard
POSIX però non specifica niente in proposito. Inoltre alcuni file speciali, ad
esempio \file{/dev/null}, non causano un errore ma restituiscono un valore
indefinito.

\itindbeg{sparse~file} 
\index{file!\textit{hole}|(} 

Infine si tenga presente che, come accennato in sez.~\ref{sec:file_file_size},
con \func{lseek} è possibile impostare una posizione anche oltre la corrente
fine del file. In tal caso alla successiva scrittura il file sarà esteso a
partire da detta posizione, con la creazione di quello che viene chiamato un
``\textsl{buco}'' (in gergo \textit{hole}) nel file.  Il nome deriva dal fatto
che nonostante la dimensione del file sia cresciuta in seguito alla scrittura
effettuata, lo spazio vuoto fra la precedente fine del file e la nuova parte,
scritta dopo lo spostamento, non corrisponde ad una allocazione effettiva di
spazio su disco, che sarebbe inutile dato che quella zona è effettivamente
vuota.

Questa è una delle caratteristiche specifiche della gestione dei file di un
sistema unix-like e quando si ha questa situazione si dice che il file in
questione è uno \textit{sparse file}. In sostanza, se si ricorda la struttura
di un filesystem illustrata in fig.~\ref{fig:file_filesys_detail}, quello che
accade è che nell'\textit{inode} del file viene segnata l'allocazione di un
blocco di dati a partire dalla nuova posizione, ma non viene allocato nulla
per le posizioni intermedie. In caso di lettura sequenziale del contenuto del
file il kernel si accorgerà della presenza del buco, e restituirà degli zeri
come contenuto di quella parte del file.

Questa funzionalità comporta una delle caratteristiche della gestione dei file
su Unix che spesso genera più confusione in chi non la conosce, per cui
sommando le dimensioni dei file si può ottenere, se si hanno molti
\textit{sparse file}, un totale anche maggiore della capacità del proprio
disco e comunque maggiore della dimensione che riporta un comando come
\cmd{du}, che calcola lo spazio disco occupato in base al numero dei blocchi
effettivamente allocati per il file.

Tutto ciò avviene proprio perché in un sistema unix-like la dimensione di un
file è una caratteristica del tutto indipendente dalla quantità di spazio
disco effettivamente allocato, e viene registrata sull'\textit{inode} come le
altre proprietà del file. La dimensione viene aggiornata automaticamente
quando si estende un file scrivendoci, e viene riportata dal campo
\var{st\_size} di una struttura \struct{stat} quando si effettua la chiamata
ad una delle funzioni \texttt{*stat} viste in sez.~\ref{sec:file_stat}.

Questo comporta che la dimensione di un file, fintanto che lo si è scritto
sequenzialmente, sarà corrispondente alla quantità di spazio disco da esso
occupato, ma possono esistere dei casi, come questo in cui ci si sposta in una
posizione oltre la fine corrente del file, o come quello accennato in
sez.~\ref{sec:file_file_size} in cui si estende la dimensione di un file con
una \func{truncate}, in cui si modifica soltanto il valore della dimensione di
\var{st\_size} senza allocare spazio su disco. Così è possibile creare
inizialmente file di dimensioni anche molto grandi, senza dover occupare da
subito dello spazio disco che in realtà sarebbe inutilizzato.

\itindend{sparse~file}

A partire dal kernel 3.1, riprendendo una interfaccia adottata su Solaris,
sono state aggiunti due nuovi valori per l'argomento \param{whence}, riportati
nella seconda sezione di tab.~\ref{tab:lseek_whence_values}, che consentono di
riconoscere la presenza di \textit{hole} all'interno dei file ad uso di quelle
applicazioni (come i programmi di backup) che possono salvare spazio disco
nella copia degli \textit{sparse file}. Una applicazione può così determinare
la presenza di un \textit{hole} usando \const{SEEK\_HOLE} all'inizio del file
e determinare poi l'inizio della successiva sezione di dati usando
\const{SEEK\_DATA}. Per compatibilità con i filesystem che non supportano
questa funzionalità è previsto comunque che in tal caso \const{SEEK\_HOLE}
riporti sempre la fine del file e \const{SEEK\_DATA} il valore
di \param{offset}.

Inoltre la decisione di come riportare (o di non riportare) la presenza di un
buco in un file è lasciata all'implementazione del filesystem, dato che oltre
a quelle classiche appena esposte esistono vari motivi per cui una sezione di
un file può non contenere dati ed essere riportata come tale (ad esempio può
essere stata preallocata con \func{fallocate}, vedi
sez.~\ref{sec:file_fadvise}). Questo significa che l'uso di questi nuovi
valori non garantisce la mappatura della effettiva allocazione dello spazio
disco di un file, per il quale esiste una specifica operazione di controllo
(vedi sez.~\ref{sec:file_fcntl_ioctl}).

\index{file!\textit{hole}|)} 


\subsection{Le funzioni per la lettura di un file}
\label{sec:file_read}

Una volta che un file è stato aperto (con il permesso in lettura) si possono
leggere i dati che contiene utilizzando la funzione di sistema \funcd{read},
il cui prototipo è:

\begin{funcproto}{
\fhead{unistd.h}
\fdecl{ssize\_t read(int fd, void * buf, size\_t count)}
\fdesc{Legge i dati da un file.} 
}

{La funzione ritorna $0$ in caso di successo e $-1$ per un errore, nel qual
  caso \var{errno} assumerà uno dei valori: 
  \begin{errlist}
  \item[\errcode{EAGAIN}] la funzione non ha nessun dato da restituire e si è
    aperto il file con \const{O\_NONBLOCK}.
  \item[\errcode{EINTR}] la funzione è stata interrotta da un segnale.
  \item[\errcode{EINVAL}] \param{fd} è associato ad un oggetto non leggibile,
    o lo si è ottenuto da \func{timerfd\_create} (vedi
    sez.~\ref{sec:sig_signalfd_eventfd}) e si è usato un valore sbagliato
    per \param{size} o si è usato \const{O\_DIRECT} ed il buffer non è
    allineato.
  \item[\errval{EIO}] si è tentata la lettura dal terminale di controllo
    essendo in background ignorando o bloccando \const{SIGTTIN} (vedi
    sez.~\ref{sec:term_io_design}) o per errori di basso livello sul supporto.
  \end{errlist}
  ed inoltre \errval{EBADF}, \errval{EFAULT} e \errval{EISDIR}, nel loro
  significato generico.}
\end{funcproto}

La funzione tenta di leggere \param{count} byte dal file \param{fd} a partire
dalla posizione corrente, scrivendoli nel buffer \param{buf}.\footnote{fino ad
  un massimo di \const{0x7ffff000} byte, indipendentemente che l'architettura
  sia a 32 o 64 bit.} Dopo la lettura la posizione sul file è spostata
automaticamente in avanti del numero di byte letti. Se \param{count} è zero la
funzione restituisce zero senza nessun altro risultato. Inoltre che non è
detto che la funzione \func{read} restituisca il numero di byte richiesto, ci
sono infatti varie ragioni per cui la funzione può restituire un numero di
byte inferiore: questo è un comportamento normale, e non un errore, che
bisogna sempre tenere presente.

La prima e più ovvia di queste ragioni è che si è chiesto di leggere più byte
di quanto il file ne contenga. In questo caso il file viene letto fino alla
sua fine, e la funzione ritorna regolarmente il numero di byte letti
effettivamente. Raggiunta la fine del file, alla ripetizione di un'operazione
di lettura, otterremmo il ritorno immediato di \func{read} con uno zero.  La
condizione di raggiungimento della fine del file non è un errore, e viene
segnalata appunto da un valore di ritorno di \func{read} nullo. Ripetere
ulteriormente la lettura non avrebbe nessun effetto se non quello di
continuare a ricevere zero come valore di ritorno.

Con i \textsl{file regolari} questa è l'unica situazione in cui si può avere
un numero di byte letti inferiore a quello richiesto, ma questo non è vero
quando si legge da un terminale, da una \textit{fifo} o da una
\textit{pipe}. In tal caso infatti, se non ci sono dati in ingresso, la
\func{read} si blocca (a meno di non aver selezionato la modalità non
bloccante, vedi sez.~\ref{sec:file_noblocking}) e ritorna solo quando ne
arrivano; se il numero di byte richiesti eccede quelli disponibili la funzione
ritorna comunque, ma con un numero di byte inferiore a quelli richiesti.

Lo stesso comportamento avviene caso di lettura dalla rete (cioè su un socket,
come vedremo in sez.~\ref{sec:sock_io_behav}), o per la lettura da certi file
di dispositivo, come le unità a nastro, che restituiscono sempre i dati ad un
singolo blocco alla volta, o come le linee seriali, che restituiscono solo i
dati ricevuti fino al momento della lettura, o i terminali, per i quali si
applicano anche ulteriori condizioni che approfondiremo in
sez.~\ref{sec:sess_terminal_io}.

Infine anche le due condizioni segnalate dagli errori \errcode{EINTR} ed
\errcode{EAGAIN} non sono propriamente degli errori. La prima si verifica
quando la \func{read} è bloccata in attesa di dati in ingresso e viene
interrotta da un segnale. In tal caso l'azione da intraprendere è quella di
rieseguire la funzione, torneremo in dettaglio sull'argomento in
sez.~\ref{sec:sig_gen_beha}.  La seconda si verifica quando il file è aperto
in modalità non bloccante (con \const{O\_NONBLOCK}) e non ci sono dati in
ingresso: la funzione allora ritorna immediatamente con un errore
\errcode{EAGAIN}\footnote{in BSD si usa per questo errore la costante
  \errcode{EWOULDBLOCK}, in Linux, con la \acr{glibc}, questa è sinonima di
  \errcode{EAGAIN}, ma se si vuole essere completamente portabili occorre
  verificare entrambi i valori, dato che POSIX.1-2001 non richiede che siano
  coincidenti.} che indica soltanto che non essendoci al momento dati
disponibili occorre provare a ripetere la lettura in un secondo tempo,
torneremo sull'argomento in sez.~\ref{sec:file_noblocking}.

La funzione \func{read} è una delle \textit{system call} fondamentali,
esistenti fin dagli albori di Unix, ma nella seconda versione delle
\textit{Single Unix Specification}\footnote{questa funzione, e l'analoga
  \func{pwrite} sono state aggiunte nel kernel 2.1.60, il supporto nella
  \acr{glibc}, compresa l'emulazione per i vecchi kernel che non hanno la
  \textit{system call}, è stato aggiunto con la versione 2.1, in versioni
  precedenti sia del kernel che delle librerie la funzione non è disponibile.}
(quello che viene chiamato normalmente Unix98, vedi
sez.~\ref{sec:intro_xopen}) è stata introdotta la definizione di un'altra
funzione di sistema, \funcd{pread}, il cui prototipo è:

\begin{funcproto}{
\fhead{unistd.h}
\fdecl{ssize\_t pread(int fd, void * buf, size\_t count, off\_t offset)}
\fdesc{Legge a partire da una posizione sul file.} 
}

{La funzione ritorna il numero di byte letti in caso di successo e $-1$ per un
  errore, nel qual caso \var{errno} assumerà i valori già visti per
  \func{read} e \func{lseek}.}
\end{funcproto}

La funzione prende esattamente gli stessi argomenti di \func{read} con lo
stesso significato, a cui si aggiunge l'argomento \param{offset} che indica
una posizione sul file a partire dalla quale verranno letti i \param{count}
byte. Identico è il comportamento ed il valore di ritorno, ma la posizione
corrente sul file resterà invariata.  Il valore di \param{offset} fa sempre
riferimento all'inizio del file.

L'uso di \func{pread} è equivalente all'esecuzione di una \func{lseek} alla
posizione indicata da \param{offset} seguita da una \func{read}, seguita da
un'altra \func{lseek} che riporti al valore iniziale della posizione corrente
sul file, ma permette di eseguire l'operazione atomicamente. Questo può essere
importante quando la posizione sul file viene condivisa da processi diversi
(vedi sez.~\ref{sec:file_shared_access}) ed è particolarmente utile in caso di
programmazione \textit{multi-thread} (vedi sez.~\ref{cha:threads}) quando
all'interno di un processo si vuole che le operazioni di un \textit{thread}
non possano essere influenzata da eventuali variazioni della posizione sul
file effettuate da altri \textit{thread}.

La funzione \func{pread} è disponibile anche in Linux, però diventa
accessibile solo attivando il supporto delle estensioni previste dalle
\textit{Single Unix Specification} con un valore della macro
\macro{\_XOPEN\_SOURCE} maggiore o uguale a 500 o a partire dalla \acr{glibc}
2.12 con un valore dalla macro \macro{\_POSIX\_C\_SOURCE} maggiore o uguale al
valore \val{200809L}.  Si ricordi di definire queste macro prima
dell'inclusione del file di dichiarazione \headfile{unistd.h}.


\subsection{Le funzioni per la scrittura di un file}
\label{sec:file_write}

Una volta che un file è stato aperto (con il permesso in scrittura) si può
scrivere su di esso utilizzando la funzione di sistema \funcd{write}, il cui
prototipo è:

\begin{funcproto}{
\fhead{unistd.h}
\fdecl{ssize\_t write(int fd, void * buf, size\_t count)}
\fdesc{Scrive i dati su un file.} 
}

{La funzione ritorna il numero di byte scritti in caso di successo e $-1$ per
  un errore, nel qual caso \var{errno} assumerà uno dei valori:
  \begin{errlist}
  \item[\errcode{EAGAIN}] ci si sarebbe bloccati, ma il file era aperto in
    modalità \const{O\_NONBLOCK}.
  \item[\errcode{EDESTADDRREQ}] si è eseguita una scrittura su un socket di
    tipo \textit{datagram} (vedi sez.~\ref{sec:sock_type}) senza aver prima
    connesso il corrispondente con \func{connect} (vedi
    sez.~\ref{sec:UDP_sendto_recvfrom}).
  \item[\errcode{EFBIG}] si è cercato di scrivere oltre la dimensione massima
    consentita dal filesystem o il limite per le dimensioni dei file del
    processo o su una posizione oltre il massimo consentito.
  \item[\errcode{EINTR}] si è stati interrotti da un segnale prima di aver
    potuto scrivere qualsiasi dato.
  \item[\errcode{EINVAL}] \param{fd} è connesso ad un oggetto che non consente
    la scrittura o si è usato \const{O\_DIRECT} ed il buffer non è allineato.
  \item[\errcode{EPERM}] la scrittura è proibita da un \textit{file seal}
    (vedi sez.~\ref{sec:file_fcntl_ioctl}).
  \item[\errcode{EPIPE}] \param{fd} è connesso ad una \textit{pipe} il cui
    altro capo è chiuso in lettura; in questo caso viene anche generato il
    segnale \signal{SIGPIPE}, se questo viene gestito (o bloccato o ignorato)
    la funzione ritorna questo errore.
  \end{errlist}
  ed inoltre \errval{EBADF}, \errval{EDQUOT}, \errval{EFAULT}, \errval{EIO},
  \errval{EISDIR}, \errval{ENOSPC} nel loro significato generico.}
\end{funcproto}


\itindbeg{append~mode}

Come nel caso di \func{read} la funzione tenta di scrivere \param{count} byte
a partire dalla posizione corrente nel file e sposta automaticamente la
posizione in avanti del numero di byte scritti. Se il file è aperto in
modalità \textit{append} con \const{O\_APPEND} i dati vengono sempre scritti
alla fine del file.  Lo standard POSIX richiede che i dati scritti siano
immediatamente disponibili ad una \func{read} chiamata dopo che la
\func{write} che li ha scritti è ritornata; ma dati i meccanismi di caching
non è detto che tutti i filesystem supportino questa capacità.

\itindend{append~mode}

Se \param{count} è zero la funzione restituisce zero senza fare nient'altro.
Per i file ordinari il numero di byte scritti è sempre uguale a quello
indicato da \param{count}, a meno di un errore. Negli altri casi si ha lo
stesso comportamento di \func{read}.

Anche per \func{write} lo standard Unix98 (ed i successivi POSIX.1-2001 e
POSIX.1-2008) definiscono un'analoga \funcd{pwrite} per scrivere alla
posizione indicata senza modificare la posizione corrente nel file, il suo
prototipo è:

\begin{funcproto}{
\fhead{unistd.h}
\fdecl{ssize\_t pwrite(int fd, void * buf, size\_t count, off\_t offset)}
\fdesc{Scrive a partire da una posizione sul file.} 
}

{La funzione ritorna il numero di byte letti in caso di successo e $-1$ per un
  errore, nel qual caso \var{errno} assumerà i valori già visti per
  \func{write} e \func{lseek}.}
\end{funcproto}

\noindent per questa funzione valgono le stesse considerazioni fatte per
\func{pread}, a cui si aggiunge il fatto che su Linux, a differenza di quanto
previsto dallo standard POSIX che richiederebbe di ignorarlo, se si è aperto
il file con \const{O\_APPEND} i dati saranno comunque scritti in coda al file,
ignorando il valore di \param{offset}.


\section{Caratteristiche avanzate}
\label{sec:file_adv_func}

In questa sezione approfondiremo alcune delle caratteristiche più sottili
della gestione file in un sistema unix-like, esaminando in dettaglio il
comportamento delle funzioni base, inoltre tratteremo le funzioni che
permettono di eseguire alcune operazioni avanzate con i file (il grosso
dell'argomento sarà comunque affrontato nel cap.~\ref{cha:file_advanced}).


\subsection{La gestione dell'accesso concorrente ai files}
\label{sec:file_shared_access}

In sez.~\ref{sec:file_fd} abbiamo descritto brevemente l'architettura
dell'interfaccia con i file da parte di un processo, mostrando in
fig.~\ref{fig:file_proc_file} le principali strutture usate dal kernel;
esamineremo ora in dettaglio le conseguenze che questa architettura ha nei
confronti dell'accesso concorrente allo stesso file da parte di processi
diversi.

\begin{figure}[!htb]
  \centering
  \includegraphics[width=11cm]{img/filemultacc}
  \caption{Schema dell'accesso allo stesso file da parte di due processi 
    diversi}
  \label{fig:file_mult_acc}
\end{figure}

Il primo caso è quello in cui due processi indipendenti aprono lo stesso file
su disco; sulla base di quanto visto in sez.~\ref{sec:file_fd} avremo una
situazione come quella illustrata in fig.~\ref{fig:file_mult_acc}: ciascun
processo avrà una sua voce nella \textit{file table} referenziata da un
diverso file descriptor nella sua \kstruct{file\_struct}. Entrambe le voci
nella \textit{file table} faranno però riferimento allo stesso \textit{inode}
su disco.

Questo significa che ciascun processo avrà la sua posizione corrente sul file,
la sua modalità di accesso e versioni proprie di tutte le proprietà che
vengono mantenute nella sua voce della \textit{file table}. Questo ha
conseguenze specifiche sugli effetti della possibile azione simultanea sullo
stesso file, in particolare occorre tenere presente che:
\begin{itemize*}
\item ciascun processo può scrivere indipendentemente, dopo ciascuna
  \func{write} la posizione corrente sarà cambiata solo nel processo
  scrivente. Se la scrittura eccede la dimensione corrente del file questo
  verrà esteso automaticamente con l'aggiornamento del campo \var{i\_size}
  della struttura \kstruct{inode}.
\item se un file è in modalità \const{O\_APPEND} tutte le volte che viene
  effettuata una scrittura la posizione corrente viene prima impostata alla
  dimensione corrente del file letta dalla struttura \kstruct{inode}. Dopo la
  scrittura il file viene automaticamente esteso. Questa operazione avviene
  atomicamente, ogni altro processo che usi \const{O\_APPEND} vedrà la
  dimensione estesa e continuerà a scrivere in coda al file.
\item l'effetto di \func{lseek} è solo quello di cambiare il campo
  \var{f\_pos} nella struttura \kstruct{file} della \textit{file table}, non
  c'è nessuna operazione sul file su disco. Quando la si usa per porsi alla
  fine del file la posizione viene impostata leggendo la attuale dimensione
  corrente dalla struttura \kstruct{inode}.
\end{itemize*}

\begin{figure}[!htb]
  \centering
  \includegraphics[width=11cm]{img/fileshar}
  \caption{Schema dell'accesso ai file da parte di un processo figlio}
  \label{fig:file_acc_child}
\end{figure}

Il secondo caso è quello in cui due file descriptor di due processi diversi
puntano alla stessa voce nella \textit{file table}.  Questo è ad esempio il
caso dei file aperti che vengono ereditati dal processo figlio all'esecuzione
di una \func{fork} (si ricordi quanto detto in sez.~\ref{sec:proc_fork}). La
situazione è illustrata in fig.~\ref{fig:file_acc_child}; dato che il processo
figlio riceve una copia dello spazio di indirizzi del padre, riceverà anche
una copia di \kstruct{file\_struct} e della relativa tabella dei file aperti.

Questo significa che il figlio avrà gli stessi file aperti del padre in
quanto la sua \kstruct{file\_struct}, pur essendo allocata in maniera
indipendente, contiene gli stessi valori di quella del padre e quindi i suoi
file descriptor faranno riferimento alla stessa voce nella \textit{file
  table}, condividendo così la posizione corrente sul file. Questo ha le
conseguenze descritte a suo tempo in sez.~\ref{sec:proc_fork}: in caso di
scrittura o lettura da parte di uno dei due processi, la posizione corrente
nel file varierà per entrambi, in quanto verrà modificato il campo
\var{f\_pos} della struttura \kstruct{file}, che è la stessa per
entrambi. Questo consente una sorta di ``\textsl{sincronizzazione}''
automatica della posizione sul file fra padre e figlio che occorre tenere
presente.

Si noti inoltre che in questo caso anche i flag di stato del file, essendo
mantenuti nella struttura \kstruct{file} della \textit{file table}, vengono
condivisi, per cui una modifica degli stessi con \func{fcntl} (vedi
sez.~\ref{sec:file_fcntl_ioctl}) si applicherebbe a tutti processi che
condividono la voce nella \textit{file table}. Ai file però sono associati
anche altri flag, dei quali l'unico usato al momento è \constd{FD\_CLOEXEC},
detti \itindex{file~descriptor~flags} \textit{file descriptor flags}; questi
invece sono mantenuti in \kstruct{file\_struct}, e perciò sono locali per
ciascun processo e non vengono modificati dalle azioni degli altri anche in
caso di condivisione della stessa voce della \textit{file table}.

Si tenga presente dunque che in un sistema unix-like è sempre possibile per
più processi accedere in contemporanea allo stesso file e che non esistono, a
differenza di altri sistemi operativi, dei meccanismi di blocco o di
restrizione dell'accesso impliciti quando più processi vogliono accedere allo
stesso file. Questo significa che le operazioni di lettura e scrittura vengono
sempre fatte da ogni processo in maniera indipendente, utilizzando una
posizione corrente nel file che normalmente, a meno di non trovarsi nella
situazione di fig.~\ref{fig:file_acc_child}, è locale a ciascuno di essi.

Dal punto di vista della lettura dei dati questo comporta la possibilità di
poter leggere dati non coerenti in caso di scrittura contemporanea da parte di
un altro processo. Dal punto di vista della scrittura invece si potranno avere
sovrapposizioni imprevedibili quando due processi scrivono nella stessa
sezione di file, dato che ciascuno lo farà in maniera indipendente.  Il
sistema però fornisce in alcuni casi la possibilità di eseguire alcune
operazioni di scrittura in maniera coordinata anche senza utilizzare dei
meccanismi di sincronizzazione espliciti come il \textit{file locking}, che
esamineremo in sez.~\ref{sec:file_locking}.

Un caso tipico di necessità di accesso condiviso in scrittura è quello in cui
vari processi devono scrivere alla fine di un file (ad esempio un file di
log). Come accennato in sez.~\ref{sec:file_lseek} impostare la posizione alla
fine del file con \func{lseek} e poi scrivere con \func{write} può condurre ad
una \textit{race condition}; infatti può succedere che un secondo processo
scriva alla fine del file fra la \func{lseek} e la \func{write}. In questo
caso, come abbiamo appena visto, il file sarà esteso, ma il primo processo,
avrà una posizione corrente che aveva impostato con \func{lseek} che non
corrisponde più alla fine del file, e la sua successiva \func{write}
sovrascriverà i dati del secondo processo.

Il problema deriva dal fatto che usare due \textit{system call} in successione
non è mai un'operazione atomica dato che il kernel può interrompere
l'esecuzione del processo fra le due. Nel caso specifico il problema è stato
risolto introducendo la modalità di scrittura in \textit{append}, attivabile
con il flag \const{O\_APPEND}. In questo caso infatti, come abbiamo illustrato
in sez.~\ref{sec:file_open_close}, è il kernel che aggiorna automaticamente la
posizione alla fine del file prima di effettuare la scrittura, e poi estende
il file.  Tutto questo avviene all'interno di una singola \textit{system
  call}, la \func{write}, che non essendo interrompibile da un altro processo
realizza un'operazione atomica.


\subsection{La duplicazione dei file descriptor}
\label{sec:file_dup}

Abbiamo già visto in sez.~\ref{sec:file_shared_access} come un processo figlio
condivida gli stessi file descriptor del padre; è possibile però ottenere un
comportamento analogo all'interno di uno stesso processo con la cosiddetta
\textit{duplicazione} di un file descriptor. Per far questo si usa la funzione
di sistema \funcd{dup}, il cui prototipo è:

\begin{funcproto}{
\fhead{unistd.h}
\fdecl{int dup(int oldfd)}
\fdesc{Crea un file descriptor duplicato.} 
}

{La funzione ritorna il nuovo file descriptor in caso di successo e $-1$ per
  un errore, nel qual caso \var{errno} assumerà uno dei valori:
  \begin{errlist}
  \item[\errcode{EBADF}] \param{oldfd} non è un file aperto.
  \item[\errcode{EMFILE}] si è raggiunto il numero massimo consentito di file
    descriptor aperti (vedi sez.~\ref{sec:sys_resource_limit}).
  \end{errlist}
}  
\end{funcproto}

La funzione ritorna, come \func{open}, il primo file descriptor libero. Il
file descriptor è una copia esatta del precedente ed entrambi possono essere
interscambiati nell'uso. Per capire meglio il funzionamento della funzione si
può fare riferimento a fig.~\ref{fig:file_dup}. L'effetto della funzione è
semplicemente quello di copiare il valore di un certo file descriptor in
un altro all'interno della struttura \kstruct{file\_struct}, cosicché anche
questo faccia riferimento alla stessa voce nella \textit{file table}. Per
questo motivo si dice che il nuovo file descriptor è ``\textsl{duplicato}'',
da cui il nome della funzione.

\begin{figure}[!htb]
  \centering \includegraphics[width=11cm]{img/filedup}
  \caption{Schema dell'accesso ai file duplicati}
  \label{fig:file_dup}
\end{figure}

Si noti che per quanto illustrato in fig.~\ref{fig:file_dup} i file descriptor
duplicati condivideranno eventuali lock (vedi sez.~\ref{sec:file_locking}), i
flag di stato, e la posizione corrente sul file. Se ad esempio si esegue una
\func{lseek} per modificare la posizione su uno dei due file descriptor, essa
risulterà modificata anche sull'altro, dato che quello che viene modificato è
lo stesso campo nella voce della \textit{file table} a cui entrambi fanno
riferimento.

L'unica differenza fra due file descriptor duplicati è che ciascuno avrà un
suo \textit{file descriptor flag} indipendente. A questo proposito deve essere
tenuto presente che nel caso in cui si usi \func{dup} per duplicare un file
descriptor, se questo ha il flag di \textit{close-on-exec} attivo (vedi
sez.~\ref{sec:proc_exec} e sez.~\ref{sec:file_fcntl_ioctl}), questo verrà
cancellato nel file descriptor restituito come copia.

L'uso principale di questa funzione è nella shell per la redirezione dei file
standard di tab.~\ref{tab:file_std_files} fra l'esecuzione di una \func{fork}
e la successiva \func{exec}. Diventa così possibile associare un file (o una
\textit{pipe}) allo \textit{standard input} o allo \textit{standard output}
(vedremo un esempio in sez.~\ref{sec:ipc_pipe_use}, quando tratteremo le
\textit{pipe}).

Ci si può chiedere perché non sia in questo caso sufficiente chiudere il file
standard che si vuole redirigere e poi aprire direttamente con \func{open} il
file vi si vuole far corrispondere, invece di duplicare un file descriptor che
si è già aperto. La risposta sta nel fatto che il file che si vuole redirigere
non è detto sia un file regolare, ma potrebbe essere, come accennato, anche
una \textit{fifo} o un socket, oppure potrebbe essere un file associato ad un
file descriptor che si è ereditato già aperto (ad esempio attraverso una
\func{exec}) da un processo antenato del padre, del quale non si conosce il
nome. Operando direttamente con i file descriptor \func{dup} consente di
ignorare le origini del file descriptor che si duplica e funziona in maniera
generica indipendentemente dall'oggetto a cui questo fa riferimento.

Per ottenere la redirezione occorre pertanto disporre del file descriptor
associato al file che si vuole usare e chiudere il file descriptor che si
vuole sostituire, cosicché esso possa esser restituito alla successiva
chiamata di \func{dup} come primo file descriptor disponibile.  Dato che
questa è l'operazione più comune, è prevista un'altra funzione di sistema,
\funcd{dup2}, che permette di specificare esplicitamente qual è il numero di
file descriptor che si vuole ottenere come duplicato; il suo prototipo è:

\begin{funcproto}{
\fhead{unistd.h}
\fdecl{int dup2(int oldfd, int newfd)}
\fdesc{Duplica un file descriptor su un altro.} 
}

{La funzione ritorna il nuovo file descriptor in caso di successo e $-1$ per
  un errore, nel qual caso \var{errno} assumerà uno dei valori:
  \begin{errlist}
  \item[\errcode{EBADF}] \param{oldfd} non è un file aperto o \param{newfd} ha
    un valore fuori dall'intervallo consentito per i file descriptor.
  \item[\errcode{EBUSY}] si è rilevata la possibilità di una \textit{race
      condition}.
  \item[\errcode{EINTR}] la funzione è stata interrotta da un segnale.
  \item[\errcode{EMFILE}] si è raggiunto il numero massimo consentito di file
    descriptor aperti.
  \end{errlist}
}  
\end{funcproto}

La funzione duplica il file descriptor \param{oldfd} su un altro file
descriptor di valore \param{newfd}. Qualora il file descriptor \param{newfd}
sia già aperto, come avviene ad esempio nel caso della duplicazione di uno dei
file standard di tab.~\ref{tab:file_std_files}, esso sarà prima chiuso e poi
duplicato. Se \param{newfd} è uguale a \param{oldfd} la funzione non fa nulla
e si limita a restituire \param{newfd}.

L'uso di \func{dup2} ha vari vantaggi rispetto alla combinazione di
\func{close} e \func{dup}; anzitutto se \param{oldfd} è uguale \param{newfd}
questo verrebbe chiuso e \func{dup} fallirebbe, ma soprattutto l'operazione è
atomica e consente di evitare una \textit{race condition} in cui dopo la
chiusura del file si potrebbe avere la ricezione di un segnale il cui gestore
(vedi sez.~\ref{sec:sig_signal_handler}) potrebbe a sua volta aprire un file,
per cui alla fine \func{dup} restituirebbe un file descriptor diverso da
quello voluto.

Con Linux inoltre la funzione prevede la possibilità di restituire l'errore
\errcode{EBUSY}, che non è previsto dallo standard, quando viene rilevata la
possibilità di una \textit{race condition} interna in cui si cerca di
duplicare un file descriptor che è stato allocato ma per il quale non sono
state completate le operazioni di apertura.\footnote{la condizione è
  abbastanza peculiare e non attinente al tipo di utilizzo indicato, quanto
  piuttosto ad un eventuale tentativo di duplicare file descriptor non ancora
  aperti, la condizione di errore non è prevista dallo standard, ma in
  condizioni simili FreeBSD risponde con un errore di \errval{EBADF}, mentre
  OpenBSD elimina la possibilità di una \textit{race condition} al costo di
  una perdita di prestazioni.} In tal caso occorre ritentare l'operazione.

La duplicazione dei file descriptor può essere effettuata anche usando la
funzione di controllo dei file \func{fcntl} (che esamineremo in
sez.~\ref{sec:file_fcntl_ioctl}) con il parametro \const{F\_DUPFD}.
L'operazione ha la sintassi \code{fcntl(oldfd, F\_DUPFD, newfd)} e se si usa 0
come valore per \param{newfd} diventa equivalente a \func{dup}.  La sola
differenza fra le due funzioni (a parte la sintassi ed i diversi codici di
errore) è che \func{dup2} chiude il file descriptor \param{newfd} se questo è
già aperto, garantendo che la duplicazione sia effettuata esattamente su di
esso, invece \func{fcntl} restituisce il primo file descriptor libero di
valore uguale o maggiore di \param{newfd}, per cui se \param{newfd} è aperto
la duplicazione avverrà su un altro file descriptor.

Su Linux inoltre è presente una terza funzione di sistema non
standard,\footnote{la funzione è stata introdotta con il kernel 2.6.27 e resa
  disponibile con la \acr{glibc} 2.9.} \funcd{dup3}, che consente di duplicare
un file descriptor reimpostandone i flag, per usarla occorre definire la macro
\macro{\_GNU\_SOURCE} ed il suo prototipo è:

\begin{funcproto}{
\fhead{unistd.h}
\fdecl{int dup3(int oldfd, int newfd, int flags)}
\fdesc{Duplica un file descriptor su un altro.} 
}

{La funzione ritorna il nuovo file descriptor in caso di successo e $-1$ per
  un errore, nel qual caso \var{errno} assumerà gli stessi valori di
  \func{dup2} più \errcode{EINVAL} qualora \param{flags} contenga un valore
  non valido o \param{newfd} sia uguale a \param{oldfd}.
}  
\end{funcproto}

La funzione è identica a \func{dup2} ma prevede la possibilità di mantenere il
flag di \textit{close-on-exec} sul nuovo file descriptor specificando
\const{O\_CLOEXEC} in \param{flags} (che è l'unico flag usabile in questo
caso). Inoltre rileva esplicitamente la possibile coincidenza
fra \param{newfd} e \param{oldfd}, fallendo con un errore di \errval{EINVAL}.


\subsection{Le funzioni di sincronizzazione dei dati}
\label{sec:file_sync}

Come accennato in sez.~\ref{sec:file_open_close} tutte le operazioni di
scrittura sono in genere bufferizzate dal kernel, che provvede ad effettuarle
in maniera asincrona per ottimizzarle, ad esempio accorpando gli accessi alla
stessa zona del disco in un secondo tempo rispetto al momento della esecuzione
della \func{write}.

Per questo motivo quando è necessaria una sincronizzazione immediata dei dati
il sistema mette a disposizione delle funzioni che provvedono a forzare lo
scarico dei dati dai buffer del kernel.  La prima di queste funzioni di
sistema è \funcd{sync}, il cui prototipo è:\footnote{questo è il prototipo
  usato a partire dalla \acr{glibc} 2.2.2 seguendo gli standard, in precedenza
  la funzione era definita come \code{int sync(void)} e ritornava sempre $0$.}

\begin{funcproto}{
\fhead{unistd.h}
\fdecl{void sync(void)}
\fdesc{Sincronizza il buffer della cache dei file col disco.} 
}

{La funzione non ritorna nulla e non prevede condizioni di errore.}  
\end{funcproto}

I vari standard prevedono che la funzione si limiti a far partire le
operazioni ritornando immediatamente, con Linux invece, fin dal kernel 1.3.20,
la funzione aspetta la conclusione delle operazioni di sincronizzazione. Si
tenga presente comunque che questo non dà la garanzia assoluta che i dati
siano integri dopo la chiamata, l'hardware dei dischi è in genere dotato di un
suo meccanismo interno di bufferizzazione che a sua volta può ritardare
ulteriormente la scrittura effettiva.

La funzione viene usata dal comando \cmd{sync} quando si vuole forzare
esplicitamente lo scarico dei dati su disco, un tempo era invocata da un
apposito demone di sistema (in genere chiamato \cmd{update}) che eseguiva lo
scarico dei dati ad intervalli di tempo fissi.  Con le nuove versioni del
kernel queste operazioni vengono gestite direttamente dal sistema della
memoria virtuale, attraverso opportuni \textit{task} interni al kernel. Nei
kernel recenti questo comportamento può essere controllato con l'uso dei vari
file \texttt{dirty\_*} in \sysctlfiled{vm/}.\footnote{si consulti la
  documentazione allegata ai sorgenti del kernel nel file
  \file{Documentation/sysctl/vm.txt}, trattandosi di argomenti di natura
  sistemistica non li prenderemo in esame.}

Si tenga presente che la funzione di sistema \funcm{bdflush}, che un tempo
veniva usata per controllare lo scaricamento dei dati, è deprecata a partire
dal kernel 2.6 e causa semplicemente la stampa di un messaggio nei log del
kernel, e non è più presente dalle \acr{glibc} 2.23, pertanto non la
prenderemo in esame.

Quando si vogliano scaricare i dati di un singolo file, ad esempio essere
sicuri che i dati di un database siano stati registrati su disco, si possono
usare le due funzioni di sistema \funcd{fsync} e \funcd{fdatasync}, i cui
prototipi sono:

\begin{funcproto}{
\fhead{unistd.h}
\fdecl{int fsync(int fd)}
\fdesc{Sincronizza dati e metadati di un file.} 
\fdecl{int fdatasync(int fd)}
\fdesc{Sincronizza i dati di un file.} 
}

{Le funzioni ritornano $0$ in caso di successo e $-1$ per un errore, nel qual
  caso \var{errno} assumerà uno dei valori: 
  \begin{errlist}
  \item[\errcode{EDQUOT}] si è superata un quota disco durante la
    sincronizzazione.
  \item[\errcode{EINVAL}] \param{fd} è un file speciale che non supporta la
    sincronizzazione (talvolta anche \errval{EROFS}).
  \item[\errcode{EIO}] c'è stato un errore di I/O durante la sincronizzazione,
    che in questo caso può derivare anche da scritture sullo stesso file
    eseguite su altri file descriptor.
  \item[\errcode{ENOSPC}] si è esaurito lo spazio disco durante la
    sincronizzazione.
  \end{errlist}
  ed inoltre \errval{EBADF} nel suo significato generico.}
\end{funcproto}

Entrambe le funzioni forzano la sincronizzazione col disco di tutti i dati del
file specificato, ed attendono fino alla conclusione delle operazioni. La
prima, \func{fsync} forza anche la sincronizzazione dei meta-dati del file,
che riguardano sia le modifiche alle tabelle di allocazione dei settori, che
gli altri dati contenuti nell'\textit{inode} che si leggono con \func{fstat},
come i tempi del file. Se lo scopo dell'operazione, come avviene spesso per i
database, è assicurarsi che i dati raggiungano il disco e siano rileggibili
immediatamente in maniera corretta, è sufficiente l'uso di \func{fdatasync}
che evita le scritture non necessarie per avere l'integrità dei dati, come
l'aggiornamento dei tempi di ultima modifica ed ultimo accesso.

Si tenga presente che l'uso di queste funzioni non comporta la
sincronizzazione della directory che contiene il file e la scrittura della
relativa voce su disco, che se necessaria deve essere effettuata
esplicitamente con \func{fsync} sul file descriptor della
directory.\footnote{in realtà per il filesystem \acr{ext2}, quando lo si monta
  con l'opzione \cmd{sync}, il kernel provvede anche alla sincronizzazione
  automatica delle voci delle directory.}

La funzione può restituire anche \errval{ENOSPC} e \errval{EDQUOT} per quei
casi in cui l'allocazione dello spazio disco non viene effettuata
all'esecuzione di una \func{write} (come NFS o altri filesystem di rete) per
cui l'errore viene rilevato quando la scrittura viene effettivamente
eseguita.

L'uso di \func{sync} può causare, quando ci sono più filesystem montati,
problemi di prestazioni dovuti al fatto che effettua la sincronizzazione dei
dati su tutti i filesystem, anche quando sarebbe sufficiente eseguirla
soltanto su quello dei file su cui si sta lavorando; quando i dati in attesa
sono molti questo può causare una alta attività di I/O ed i relativi problemi
di prestazioni.

Per questo motivo è stata introdotta una nuova funzione di sistema,
\funcd{syncfs},\footnote{la funzione è stata introdotta a partire dal kernel
  2.6.39 ed è accessibile solo se è definita la macro \macro{\_GNU\_SOURCE}, è
  specifica di Linux e non prevista da nessuno standard.} che effettua lo
scarico dei dati soltanto per il filesystem su cui si sta operando, il suo
prototipo è:

\begin{funcproto}{
\fhead{unistd.h}
\fdecl{int syncfs(int fd)}
\fdesc{Sincronizza il buffer della cache dei file del singolo filesystem col
  disco.}
}

{La funzione ritorna $0$ in caso di successo e $-1$ per un errore, nel qual
  caso \var{errno} assumerà uno dei valori: 
  \begin{errlist}
    \item[\errcode{EBADF}] \param{fd} non è un descrittore valido.
  \end{errlist}
}  
\end{funcproto}

La funzione richiede che si specifichi nell'argomento \param{fd} un file
descriptor su cui si sta operando, e la registrazione immediata dei dati sarà
limitata al filesystem su cui il file ad esso corrispondente si trova.


\subsection{Le \textit{at-functions}: \func{openat} e le altre}
\label{sec:file_openat}

\itindbeg{at-functions}

Un problema generico che si pone con l'uso della funzione \func{open}, così
come con le altre funzioni che prendono come argomenti dei \textit{pathname},
è la possibilità, quando si usa un \textit{pathname} che non fa riferimento
diretto ad un file posto nella directory di lavoro corrente, che alcuni dei
componenti dello stesso vengano modificati in parallelo alla chiamata a
\func{open}, cosa che lascia aperta la possibilità di una \textit{race
  condition} in cui c'è spazio per un \textit{symlink attack} (si ricordi
quanto visto per \func{access} in sez.~\ref{sec:file_perm_management})
cambiando una delle directory sovrastanti il file fra un controllo e la
successiva apertura. 

Inoltre, come già accennato, la directory di lavoro corrente è una proprietà
associata al singolo processo; questo significa che quando si lavora con i
\textit{thread} questa è la stessa per tutti, per cui se la si cambia
all'interno di un \textit{thread} il cambiamento varrà anche per tutti gli
altri. Non esiste quindi con le funzioni classiche un modo semplice per far sì
che i singoli \textit{thread} possano aprire file usando una propria directory
per risolvere i \textit{pathname} relativi.

Per risolvere questi problemi, riprendendo una interfaccia già presente in
Solaris, a fianco delle normali funzioni che operano sui file (come
\func{open}, \func{mkdir}, ecc.) sono state introdotte delle ulteriori
funzioni di sistema, chiamate genericamente ``\textit{at-functions}'' in
quanto usualmente contraddistinte dal suffisso \texttt{at}, che permettono
l'apertura di un file (o le rispettive altre operazioni) usando un
\textit{pathname} relativo ad una directory
specificata.\footnote{l'introduzione è avvenuta su proposta dello sviluppatore
  principale della \acr{glibc} Urlich Drepper e le corrispondenti
  \textit{system call} sono state inserite nel kernel a partire dalla versione
  2.6.16, in precedenza era disponibile una emulazione che, sia pure con
  prestazioni inferiori, funzionava facendo ricorso all'uso del filesystem
  \textit{proc} con l'apertura del file attraverso il riferimento a
  \textit{pathname} del tipo di \texttt{/proc/self/fd/dirfd/relative\_path}.}
Essendo accomunate dalla stessa interfaccia le tratteremo insieme in questa
sezione pur non essendo strettamente attinenti l'I/O su file.


Benché queste funzioni non siano presenti negli standard tradizionali esse
sono state adottate da altri sistemi unix-like come Solaris, i vari BSD, fino
ad essere incluse in una recente revisione dello standard POSIX.1 (la
POSIX.1-2008). Con la \acr{glibc} per l'accesso a queste funzioni è necessario
definire la macro \macro{\_ATFILE\_SOURCE} (comunque attiva di default).

L'uso di queste funzioni richiede una apertura preliminare della directory che
si intende usare come base per la risoluzione dei \textit{pathname} relativi
(ad esempio usando \func{open} con il flag \const{O\_PATH} visto in
sez.~\ref{sec:file_open_close}) per ottenere un file descriptor che dovrà
essere passato alle stesse.  Tutte queste funzioni infatti prevedono la
presenza un apposito argomento, in genere il primo, che negli esempi seguenti
chiameremo sempre \param{dirfd}.

In questo modo, una volta aperta la directory di partenza, si potranno
effettuare controlli ed aperture solo con \textit{pathname} relativi alla
stessa, e tutte le \textit{race condition} dovute al possibile cambiamento di
uno dei componenti posti al di sopra della stessa cesseranno di esistere.
Inoltre, pur restando la directory di lavoro una proprietà comune del
processo, si potranno usare queste funzioni quando si lavora con i
\textit{thread} per eseguire la risoluzione dei \textit{pathname} relativi ed
avere una directory di partenza diversa in ciascuno di essi.

Questo metodo consente inoltre di ottenere aumenti di prestazioni
significativi quando si devono eseguire molte operazioni su sezioni
dell'albero dei file che prevedono delle gerarchie di sottodirectory molto
profonde. Infatti in questo caso basta eseguire la risoluzione del
\textit{pathname} di una qualunque directory di partenza una sola volta
(nell'apertura iniziale) e non tutte le volte che si deve accedere a ciascun
file che essa contiene. Infine poter identificare una directory di partenza
tramite il suo file descriptor consente di avere un riferimento stabile alla
stessa anche qualora venisse rinominata, e tiene occupato il filesystem dove
si trova, come per la directory di lavoro di un processo.

La sintassi generica di queste nuove funzioni prevede l'utilizzo come primo
argomento del file descriptor della directory da usare come base per la
risoluzione dei nomi, mentre gli argomenti successivi restano identici a
quelli della corrispondente funzione ordinaria. Come esempio prendiamo in
esame la nuova funzione di sistema \funcd{openat}, il cui prototipo è:

\begin{funcproto}{
\fhead{fcntl.h}
\fdecl{int openat(int dirfd, const char *pathname, int flags)}
\fdecl{int openat(int dirfd, const char *pathname, int flags, mode\_t mode)}
\fdesc{Apre un file a partire da una directory di lavoro.} 
}

{La funzione ritorna gli stessi valori e gli stessi codici di errore di
  \func{open}, ed in più:
  \begin{errlist}
  \item[\errcode{EBADF}] \param{dirfd} non è un file descriptor valido.
  \item[\errcode{ENOTDIR}] \param{pathname} è un \textit{pathname} relativo,
    ma \param{dirfd} fa riferimento ad un file.
   \end{errlist}
}  
\end{funcproto}

Il comportamento di \func{openat} è del tutto analogo a quello di \func{open},
con la sola eccezione del fatto che se per l'argomento \param{pathname} si
utilizza un \textit{pathname} relativo questo sarà risolto rispetto alla
directory indicata da \param{dirfd}; qualora invece si usi un
\textit{pathname} assoluto \param{dirfd} verrà semplicemente ignorato. Infine
se per \param{dirfd} si usa il valore speciale \constd{AT\_FDCWD} la
risoluzione sarà effettuata rispetto alla directory di lavoro corrente del
processo. Questa, come le altre costanti \texttt{AT\_*}, è definita in
\headfile{fcntl.h}, per cui per usarla occorrerà includere comunque questo
file, anche per le funzioni che non sono definite in esso.

Così come il comportamento, anche i valori di ritorno e le condizioni di
errore delle nuove funzioni sono gli stessi delle funzioni classiche, agli
errori si aggiungono però quelli dovuti a valori errati per \param{dirfd}; in
particolare si avrà un errore di \errcode{EBADF} se esso non è un file
descriptor valido, ed un errore di \errcode{ENOTDIR} se esso non fa
riferimento ad una directory, tranne il caso in cui si sia specificato un
\textit{pathname} assoluto, nel qual caso, come detto, il valore di
\param{dirfd} sarà completamente ignorato.

\begin{table}[htb]
  \centering
  \footnotesize
  \begin{tabular}[c]{|l|c|l|}
    \hline
    \textbf{Funzione} &\textbf{Flags} &\textbf{Corrispondente} \\
    \hline
    \hline
     \func{execveat}  &$\bullet$&\func{execve}  \\
     \func{faccessat} &$\bullet$&\func{access}  \\
     \func{fchmodat}  &$\bullet$&\func{chmod}   \\
     \func{fchownat}  &$\bullet$&\func{chown},\func{lchown}\\
     \func{fstatat}   &$\bullet$&\func{stat},\func{lstat}  \\
     \funcm{futimesat}& --      & obsoleta  \\
     \func{linkat}    &$\bullet$&\func{link}    \\
     \funcm{mkdirat}  & --      &\func{mkdir}   \\
     \funcm{mkfifoat} & --      &\func{mkfifo}  \\
     \funcm{mknodat}  & --      &\func{mknod}   \\
     \func{openat}    & --      &\func{open}    \\
     \funcm{readlinkat}& --     &\func{readlink}\\
     \func{renameat}  & --      &\func{rename}  \\
     \func{renameat2}\footnotemark& -- &\func{rename}  \\
     \funcm{scandirat}& --      &\func{scandir}  \\
     \func{statx}     &$\bullet$&\func{stat}  \\
     \funcm{symlinkat}& --      &\func{symlink} \\
     \func{unlinkat}  &$\bullet$&\func{unlink},\func{rmdir}  \\
     \func{utimensat} &$\bullet$&\func{utimes},\func{lutimes}\\
    \hline
  \end{tabular}
  \caption{Corrispondenze fra le nuove funzioni ``\textit{at}'' e le
    corrispettive funzioni classiche.}
  \label{tab:file_atfunc_corr}
\end{table}

\footnotetext{anche se la funzione ha un argomento \param{flags} questo
  attiene a funzionalità specifiche della stessa e non all'uso generico fatto
  nelle altre \textit{at-functions}, pertanto lo si è indicato come assente.}

In tab.~\ref{tab:file_atfunc_corr} si sono riportate le funzioni introdotte
con questa nuova interfaccia, con a fianco la corrispondente funzione
classica. Tutte seguono la convenzione appena vista per \func{openat}, in cui
agli argomenti della funzione classica viene anteposto l'argomento
\param{dirfd}. Per alcune, indicate dal contenuto della omonima colonna di
tab.~\ref{tab:file_atfunc_corr}, oltre al nuovo argomento iniziale, è prevista
anche l'aggiunta di un argomento finale, \param{flags}, che è stato introdotto
per fornire un meccanismo con cui modificarne il comportamento.

Per tutte quelle che non hanno un argomento aggiuntivo il comportamento è
identico alla corrispondente funzione ordinaria, pertanto non le tratteremo
esplicitamente, vale per loro quanto detto con \func{openat} per l'uso del
nuovo argomento \param{dirfd}. Tratteremo invece esplicitamente tutte quelle
per cui l'argomento è presente, in quanto il loro comportamento viene
modificato a seconda del valore assegnato a \param{flags}; questo deve essere
passato come maschera binaria con una opportuna combinazione delle costanti
elencate in tab.~\ref{tab:at-functions_constant_values}, in quanto sono
possibili diversi valori a seconda della funzione usata.

\begin{table}[htb]
  \centering
  \footnotesize
  \begin{tabular}[c]{|l|p{8cm}|}
    \hline
    \textbf{Costante} & \textbf{Significato} \\
    \hline
    \hline
    \constd{AT\_EMPTY\_PATH}    & Usato per operare direttamente (specificando
                                  una stringa vuota  per il \texttt{pathname})
                                  sul file descriptor \param{dirfd} che in
                                  questo caso può essere un file qualunque.\\
    \constd{AT\_SYMLINK\_NOFOLLOW}& Se impostato la funzione non esegue la
                                    dereferenziazione dei collegamenti
                                    simbolici.\\
    \hline
    \constd{AT\_EACCES}         & Usato solo da \func{faccessat}, richiede che
                                  il controllo dei permessi sia fatto usando
                                  l'\ids{UID} effettivo invece di quello
                                  reale.\\
    \constd{AT\_NO\_AUTOMOUNT}  & Usato solo da \func{fstatat} e \func{statx},
                                  evita il montaggio automatico qualora 
                                  \param{pathname} faccia riferimento ad una
                                  directory marcata per
                                  l'\textit{automount}\footnotemark
                                  (dal kernel 2.6.38).\\
    \constd{AT\_REMOVEDIR}      & Usato solo da \func{unlinkat}, richiede che
                                  la funzione si comporti come \func{rmdir}
                                  invece che come \func{unlink}.\\
    \constd{AT\_SYMLINK\_FOLLOW}& Usato solo da \func{linkat}, se impostato la
                                  funzione esegue la dereferenziazione dei
                                  collegamenti simbolici.\\
    \hline
  \end{tabular}  
  \caption{Le costanti utilizzate per i bit dell'argomento aggiuntivo
    \param{flags} delle \textit{at-functions}, definite in
    \headfile{fcntl.h}.}
  \label{tab:at-functions_constant_values}
\end{table}

\footnotetext{si tratta di una funzionalità fornita dal kernel che consente di
  montare automaticamente una directory quando si accede ad un
  \textit{pathname} al di sotto di essa, per i dettagli, più di natura
  sistemistica, si può consultare sez.~5.1.6 di \cite{AGL}.}

Si tenga presente che non tutte le funzioni che prevedono l'argomento
aggiuntivo sono \textit{system call}, ad esempio \func{faccessat} e
\func{fchmodat} sono realizzate con dei \textit{wrapper} nella \acr{glibc} per
aderenza allo standard POSIX.1-2008, dato che la \textit{system call}
sottostante non prevede l'argomento \param{flags}. 

In tab.~\ref{tab:at-functions_constant_values} si sono elencati i valori
utilizzabili per i flag (tranne quelli specifici di \func{statx} su cui
torneremo più avanti), mantenendo nella prima parte quelli comuni usati da più
funzioni. Il primo di questi è \const{AT\_SYMLINK\_NOFOLLOW}, che viene usato
da tutte le funzioni tranne \func{linkat} e \func{unlinkat}, e che consente di
scegliere, quando si sta operando su un collegamento simbolico, se far agire
la funzione direttamente sullo stesso o sul file da esso referenziato. Si
tenga presente però che per \funcm{fchmodat} questo, che è l'unico flag
consentito e previsto dallo standard, non è attualmente implementato (anche
perché non avrebbe molto senso cambiare i permessi di un link simbolico) e
pertanto l'uso della funzione è analogo a quello delle altre funzioni che non
hanno l'argomento \param{flags} (e non la tratteremo esplicitamente).

L'altro flag comune è \const{AT\_EMPTY\_PATH}, utilizzabile a partire dal
kernel 2.6.39, che consente di usare per \param{dirfd} un file descriptor
associato ad un file qualunque e non necessariamente ad una directory; in
particolare si può usare un file descriptor ottenuto aprendo un file con il
flag \param{O\_PATH} (vedi quanto illustrato a
pag.~\pageref{open_o_path_flag}). Quando si usa questo flag \param{pathname}
deve essere vuoto, da cui il nome della costante, ed in tal caso la funzione
agirà direttamente sul file associato al file descriptor \param{dirfd}.

Una prima funzione di sistema che utilizza l'argomento \param{flag} è
\funcd{fchownat}, che può essere usata per sostituire sia \func{chown} che
\func{lchown}; il suo prototipo è:

\begin{funcproto}{
\fhead{fcntl.h} 
\fhead{unistd.h}
\fdecl{int fchownat(int dirfd, const char *pathname, uid\_t owner, gid\_t
    group, int flags)}
\fdesc{Modifica il proprietario di un file.} 
}

{La funzione ritorna gli stessi valori e gli stessi codici di errore di
  \func{chown}, ed in più:
  \begin{errlist}
  \item[\errcode{EBADF}] \param{dirfd} non è un file descriptor valido.
  \item[\errcode{EINVAL}] \param{flags} non ha un valore valido.
  \item[\errcode{ENOTDIR}] \param{pathname} è un \textit{pathname} relativo,
    ma \param{dirfd} fa riferimento ad un file.
  \end{errlist}
}  
\end{funcproto}

In questo caso, oltre a quanto già detto per \func{openat} riguardo all'uso di
\param{dirfd}, se si è impostato \const{AT\_SYMLINK\_NOFOLLOW} in
\param{flags}, si indica alla funzione di non eseguire la dereferenziazione di
un eventuale collegamento simbolico, facendo comportare \func{fchownat} come
\func{lchown} invece che come \func{chown}. La funzione supporta anche l'uso
di \const{AT\_EMPTY\_PATH}, con il significato illustrato in precedenza e non
ha flag specifici.

Una seconda funzione di sistema che utilizza l'argomento \param{flags}, in
questo caso anche per modificare il suo comportamento, è \funcd{faccessat}, ed
il suo prototipo è:

\begin{funcproto}{
\fhead{fcntl.h} 
\fhead{unistd.h}
\fdecl{int faccessat(int dirfd, const char *path, int mode, int flags)}
\fdesc{Controlla i permessi di accesso.} 
}

{La funzione ritorna gli stessi valori e gli stessi codici di errore di
  \func{access}, ed in più:
  \begin{errlist}
  \item[\errcode{EBADF}] \param{dirfd} non è un file descriptor valido.
  \item[\errcode{EINVAL}] \param{flags} non ha un valore valido.
  \item[\errcode{ENOTDIR}] \param{pathname} è un \textit{pathname} relativo,
    ma \param{dirfd} fa riferimento ad un file.
  \end{errlist}
}  
\end{funcproto}

La funzione esegue il controllo di accesso ad un file, e \param{flags}
consente di modificarne il comportamento rispetto a quello ordinario di
\func{access} (cui è analoga, e con cui condivide i problemi di sicurezza
visti in sez.~\ref{sec:file_stat}) usando il valore \const{AT\_EACCES} per
indicare alla funzione di eseguire il controllo dei permessi con l'\ids{UID}
\textsl{effettivo} invece di quello \textsl{reale}. L'unico altro valore
consentito è \const{AT\_SYMLINK\_NOFOLLOW}, con il significato già spiegato.

Un utilizzo specifico dell'argomento \param{flags} viene fatto anche dalla
funzione di sistema \funcd{unlinkat}, in questo caso l'argomento viene
utilizzato perché tramite esso si può indicare alla funzione di comportarsi
sia come analogo di \func{unlink} che di \func{rmdir}; il suo prototipo è:

\begin{funcproto}{
\fhead{fcntl.h}
\fhead{unistd.h}
\fdecl{int unlinkat(int dirfd, const char *pathname, int flags)}
\fdesc{Rimuove una voce da una directory.} 
}

{La funzione ritorna gli stessi valori e gli stessi codici di errore di
  \func{unlink} o di \func{rmdir} a seconda del valore di \param{flags}, ed in
  più:
  \begin{errlist}
  \item[\errcode{EBADF}] \param{dirfd} non è un file descriptor valido.
  \item[\errcode{EINVAL}] \param{flags} non ha un valore valido.
  \item[\errcode{ENOTDIR}] \param{pathname} è un \textit{pathname} relativo,
    ma \param{dirfd} fa riferimento ad un file.
  \end{errlist}
}  
\end{funcproto}

Di default il comportamento di \func{unlinkat} è equivalente a quello che
avrebbe \func{unlink} applicata a \param{pathname}, fallendo in tutti i casi
in cui questo è una directory, se però si imposta \param{flags} al valore di
\const{AT\_REMOVEDIR}, essa si comporterà come \func{rmdir}, in tal caso
\param{pathname} deve essere una directory, che sarà rimossa qualora risulti
vuota.  Non essendo in questo caso prevista la possibilità di usare altri
valori (la funzione non segue comunque i collegamenti simbolici e
\const{AT\_EMPTY\_PATH} non è supportato) anche se \param{flags} è una
maschera binaria, essendo \const{AT\_REMOVEDIR} l'unico flag disponibile per
questa funzione, lo si può assegnare direttamente.

Un'altra funzione di sistema che usa l'argomento \param{flags} è
\func{utimensat}, che però non è una corrispondente esatta delle funzioni
classiche \func{utimes} e \func{lutimes}, in quanto ha una maggiore precisione
nella indicazione dei tempi dei file, per i quali, come per \func{futimens},
si devono usare strutture \struct{timespec} che consentono una precisione fino
al nanosecondo; la funzione è stata introdotta con il kernel
2.6.22,\footnote{in precedenza, a partire dal kernel 2.6.16, era stata
  introdotta una \textit{system call} \funcm{futimesat} seguendo una bozza
  della revisione dello standard poi modificata; questa funzione, sostituita
  da \func{utimensat}, è stata dichiarata obsoleta, non è supportata da
  nessuno standard e non deve essere più utilizzata: pertanto non ne
  parleremo.} ed il suo prototipo è:

\begin{funcproto}{
\fhead{fcntl.h}
\fhead{sys/stat.h}
\fdecl{int utimensat(int dirfd, const char *pathname, const struct
    timespec times[2],\\
\phantom{int utimensat(}int flags)}
\fdesc{Cambia i tempi di un file.} 
}

{La funzione ritorna $0$ in caso di successo e $-1$ per un errore, nel qual
  caso \var{errno} assumerà i valori di \func{utimes}, \func{lutimes} e
  \func{futimens} con lo stesso significato ed inoltre:
  \begin{errlist}
  \item[\errcode{EBADF}] \param{dirfd} non è \const{AT\_FDCWD} o un file
    descriptor valido.
  \item[\errcode{EFAULT}] \param{dirfd} è \const{AT\_FDCWD} ma
    \param{pathname} è \var{NULL} o non è un puntatore valido.
  \item[\errcode{EINVAL}] si usato un valore non valido per \param{flags},
    oppure \param{pathname} è \var{NULL}, \param{dirfd} non è
    \const{AT\_FDCWD} e \param{flags} contiene \const{AT\_SYMLINK\_NOFOLLOW}.
  \item[\errcode{ESRCH}] non c'è il permesso di attraversamento per una delle
    componenti di \param{pathname}.
  \end{errlist}
}
\end{funcproto}

La funzione imposta i tempi dei file utilizzando i valori passati nel vettore
di strutture \struct{timespec} ed ha in questo lo stesso comportamento di
\func{futimens}, vista in sez.~\ref{sec:file_file_times}, ma al contrario di
questa può essere applicata anche direttamente ad un file come \func{utimes};
l'unico valore consentito per \param{flags} è \const{AT\_SYMLINK\_NOFOLLOW}
che indica alla funzione di non dereferenziare i collegamenti simbolici, cosa
che le permette di riprodurre anche le funzionalità di \func{lutimes} (con una
precisione dei tempi maggiore).

Su Linux solo \func{utimensat} è una \textit{system call} mentre
\func{futimens} è una funzione di libreria, infatti \func{utimensat} ha un
comportamento speciale se \param{pathname} è \var{NULL}, in tal caso
\param{dirfd} viene considerato un file descriptor ordinario e il cambiamento
del tempo viene applicato al file sottostante, qualunque esso sia. Viene cioè
sempre usato il comportamento che per altre funzioni deve essere attivato con
\const{AT\_EMPTY\_PATH} (che non è previsto per questa funzione) per cui
\code{futimens(fd, times}) è del tutto equivalente a \code{utimensat(fd, NULL,
  times, 0)}. Si tenga presente che nella \acr{glibc} questo comportamento è
disabilitato, e la funzione, seguendo lo standard POSIX, ritorna un errore di
\errval{EINVAL} se invocata in questo modo.

Come corrispondente di \func{stat}, \func{fstat} e \func{lstat} si può
utilizzare invece la funzione di sistema \funcd{fstatat}, il cui prototipo è:

\begin{funcproto}{
\fhead{fcntl.h}
\fhead{sys/stat.h}
\fdecl{int fstatat(int dirfd, const char *pathname, struct stat *statbuf, int
  flags)} 
\fdesc{Rimuove una voce da una directory.} 
}

{La funzione ritorna gli stessi valori e gli stessi codici di errore di
  \func{stat}, \func{fstat}, o \func{lstat} a seconda del valore di
  \param{flags}, ed in più:
  \begin{errlist}
  \item[\errcode{EBADF}] \param{dirfd} non è un file descriptor valido.
  \item[\errcode{EINVAL}] \param{flags} non ha un valore valido.
  \item[\errcode{ENOTDIR}] \param{pathname} è un \textit{pathname} relativo,
    ma \param{dirfd} fa riferimento ad un file.
  \end{errlist}
}  
\end{funcproto}

La funzione ha lo stesso comportamento delle sue equivalenti classiche, l'uso
di \param{flags} consente di farla comportare come \func{lstat} se si usa
\const{AT\_SYMLINK\_NOFOLLOW}, o come \func{fstat} se si usa con
\const{AT\_EMPTY\_PATH} e si passa il file descriptor in \param{dirfd}. Viene
però supportato l'ulteriore valore \const{AT\_NO\_AUTOMOUNT} che qualora
\param{pathname} faccia riferimento ad una directory marcata per
l'\textit{automount} ne evita il montaggio automatico.
            
Ancora diverso è il caso di \funcd{linkat} anche se in questo caso l'utilizzo
continua ad essere attinente al comportamento con i collegamenti simbolici, il
suo prototipo è:

\begin{funcproto}{
\fhead{fcntl.h}
\fdecl{int linkat(int olddirfd, const char *oldpath, int newdirfd, \\
\phantom{int linkat(}const char *newpath, int flags)}
\fdesc{Crea un nuovo collegamento diretto (\textit{hard link}).} 
}

{La funzione ritorna gli stessi valori e gli stessi codici di errore di
  \func{link}, ed in più:
  \begin{errlist}
  \item[\errcode{EBADF}] \param{olddirfd} o \param{newdirfd} non sono un file
    descriptor valido.
  \item[\errcode{EINVAL}] \param{flags} non ha un valore valido.
  \item[\errcode{ENOTDIR}] \param{oldpath} e \param{newpath} sono
    \textit{pathname} relativi, ma \param{oldirfd} o \param{newdirfd} fa
    riferimento ad un file.
  \end{errlist}
}  
\end{funcproto}

Anche in questo caso la funzione è sostanzialmente identica alla classica
\func{link}, ma dovendo specificare due \textit{pathname} (sorgente e
destinazione) aggiunge a ciascuno di essi un argomento (rispettivamente
\param{olddirfd} e \param{newdirfd}) per poter indicare entrambi come relativi
a due directory aperte in precedenza.

In questo caso, dato che su Linux il comportamento di \func{link} è quello di
non seguire mai i collegamenti simbolici, \const{AT\_SYMLINK\_NOFOLLOW} non
viene utilizzato. A partire dal kernel 2.6.18 è stato aggiunto a questa
funzione la possibilità di usare il valore \const{AT\_SYMLINK\_FOLLOW} per
l'argomento \param{flags},\footnote{nei kernel precendenti, dall'introduzione
  nel 2.6.16, l'argomento \param{flags} era presente, ma senza alcun valore
  valido, e doveva essere passato sempre con valore nullo.}  che richiede di
dereferenziare i collegamenti simbolici. Inoltre a partire dal kernel 3.11 si
può usare \const{AT\_EMPTY\_PATH} per creare un nuovo \textit{hard link} al
file associato al file descriptor \param{olddirfd}.


% NOTE per la discussione sui problemi di sicurezza relativi a questa
% funzionalità vedi http://lwn.net/Articles/562488/

La funzione prevede inoltre un comportamento specifico nel caso che
\param{olddirfd} faccia riferimento ad un file anonimo ottenuto usando
\func{open} con \const{O\_TMPFILE}. In generale quando il file associato ad
\param{olddirfd} ha un numero di link nullo (come in questo caso), la funzione
fallisce, c'è però una 


l'uso di \const{AT\_EMPTY\_PATH} assume un significato
ulteriore, e richiede i privilegi di amministratore (la \textit{capability}
\const{CAP\_DAC\_READ\_SEARCH}) quando viene usato con un file descriptor
anomino ottenuto usando \const{O\_TMPFILE} con \func{open}. In generale


Altre due funzioni che utilizzano due \textit{pathname} (e due file
descriptor) sono \funcd{renameat} e \funcd{renameat2}, corrispondenti alla
classica \func{rename}; i rispettivi prototipi sono:

\begin{funcproto}{
\fhead{fcntl.h}
\fdecl{int renameat(int olddirfd, const char *oldpath, int newdirfd, const
  char *newpath)} 
\fdecl{int renameat2(int olddirfd, const char *oldpath, int newdirfd, \\
\phantom{int renameat2(}const char *newpath, int flags)}
\fdesc{Rinomina o sposta un file o una directory.} 
}

{La funzioni ritornano gli stessi valori e gli stessi codici di errore di
  \func{rename}, ed in più per entrambe:
  \begin{errlist}
  \item[\errcode{EBADF}] \param{olddirfd} o \param{newdirfd} non sono un file
    descriptor valido.
  \item[\errcode{ENOTDIR}] \param{oldpath} e \param{newpath} sono
    \textit{pathname} relativi, ma i corrispondenti \param{oldirfd} o
    \param{newdirfd} fan riferimento ad un file e non a una directory.
  \end{errlist}
  e per \func{renameat2} anche:
  \begin{errlist}
  \item[\errcode{EEXIST}] si è richiesto \macro{RENAME\_NOREPLACE} ma
    \param{newpath} esiste già.
  \item[\errcode{EINVAL}] Si è usato un flag non valido in \param{flags}, o si
    sono usati insieme a \macro{RENAME\_EXCHANGE} o \macro{RENAME\_NOREPLACE}
    o \macro{RENAME\_WHITEOUT}, o non c'è il supporto nel filesystem per una
    delle operazioni richieste in \param{flags}.
  \item[\errcode{ENOENT}] si è richiesto \macro{RENAME\_EXCHANGE} e
    \param{newpath} non esiste.
  \item[\errcode{EPERM}] si è richiesto \macro{RENAME\_WHITEOUT} ma il
    chiamante non ha i privilegi di amministratore.
  \end{errlist}
}  
\end{funcproto}

In realtà la corrispondente di \func{rename}, prevista dallo standard
POSIX.1-2008 e disponibile dal kernel 2.6.16 come le altre
\textit{at-functions}, sarebbe soltanti \func{renameat}, su Linux però, a
partire dal kernel dal 3.15, questa è stata realizzata in termini della nuova
funzione di sistema \func{renameat2} che prevede l'uso dell'argomento
aggiuntivo \param{flags}; in questo caso \func{renameat} è totalmente
equivalente all'utilizzo di \func{renamat2} con un valore nullo per
\param{flags}.

L'uso di \func{renameat} è identico a quello di \func{rename}, con le solite
note relativie alle estensioni delle \textit{at-functions}, applicate ad
entrambi i \textit{pathname} passati come argomenti alla funzione. Con
\func{renameat2} l'introduzione dell'argomento \func{flags} (i cui valori
possibili sono riportati in tab.~\ref{tab:renameat2_flag_values}) ha permesso
di aggiungere alcune funzionalità non previste al momento da nessuno standard,
e specifiche di Linux. Si tenga conto che questa funzione di sistema non viene
definita nella \acr{glibc} per cui deve essere chiamata utilizzando
\func{syscall} come illustrato in sez.~\ref{sec:intro_syscall}.

\begin{table}[htb]
  \centering
  \footnotesize
  \begin{tabular}[c]{|l|p{8cm}|}
    \hline
    \textbf{Costante} & \textbf{Significato} \\
    \hline
    \hline
    \const{RENAME\_EXCHANGE} & richiede uno scambio di nomi fra
                               \param{oldpath} e \param{newpath}, non è
                               usabile con \const{RENAME\_NOREPLACE}.\\
    \const{RENAME\_NOREPLACE}& non sovrascrive  \param{newpath} se questo
                               esiste dando un errore.\\
    \const{RENAME\_WHITEOUT} & crea un oggetto di \textit{whiteout}
                               contestualmente al cambio di nome 
                               (disponibile a partire dal kernel 3.18).\\ 
    \hline
  \end{tabular}  
  \caption{I valori specifici dei bit dell'argomento \param{flags} per l'uso
    con \func{renameat2}.}
  \label{tab:renameat2_flag_values}
\end{table}

L'uso dell'argomento \param{flags} in questo caso non attiene alle
funzionalità relative alla \textit{at-functions}, ma consente di estendere le
funzionalità di \func{rename}. In particolare \func{renameat2} consente di
eseguire uno scambio di nomi in maniera atomica usando il flag
\constd{RENAME\_EXCHANGE}; quando viene specificato la funzione non solo
rinomina \param{oldpath} in \param{newpath}, ma rinomina anche, senza dover
effettuare un passaggio intermedio, \param{newpath} in \param{oldpath}. Quando
si usa questo flag, entrambi i \textit{pathname} passati come argomenti alla
funzione devono esistere, e non è possibile usare \const{RENAME\_NOREPLACE},
non ci sono infine restrizioni sul tipo dei file (regolari, directory, link
simbolici, ecc.) di cui si scambia il nome.

Il flag \constd{RENAME\_NOREPLACE} consente di richiedere la generazione di un
errore nei casi in cui \func{rename} avrebbe causato una sovrascrittura della
destinazione, rendendo possibile evitare la stessa in maniera atomica; un
controllo preventivo dell'esistenza del file infatti avrebbe aperto alla
possibilità di una \textit{race condition} fra il momento del controllo e
quella del cambio di nome.

\itindbeg{overlay~filesytem}
\itindbeg{union~filesytem}

Infine il flag \constd{RENAME\_WHITEOUT}, introdotto con il kernel 3.18,
richiede un approfomdimento specifico, in quanto attiene all'uso della
funzione con dei filesystem di tipo \textit{overlay}/\textit{union}, dato che
il flag ha senso solo quando applicato a file che stanno su questo tipo di
filesystem. 

Un \textit{overlay} o \texttt{union filesystem} è un filesystem speciale
strutturato in livelli, in cui si rende scrivibile un filesystem accessibile
in sola lettura, \textsl{sovrapponendogli} un filesystem scrivibile su cui
vanno tutte le modifiche. Un tale tipo di filesystem serve ad esempio a
rendere scrivibili i dati processati quando si fa partire una distribuzione
\textit{Live} basata su CD o DVD, ad esempio usando una chiavetta o uno spazio
disco aggiuntivo.

In questo caso quando si rinomina un file che sta nello strato in sola lettura
quello che succede è questo viene copiato a destinazione sulla parte
accessibile in scrittura, ma l'originale non può essere cancellato, per far si
che esso non appaia più è possibile creare 

\itindend{overlay~filesytem}
\itindend{union~filesytem}



% TODO trattare anche statx, aggiunta con il kernel 4.11 (vedi
% https://lwn.net/Articles/707602/ e
% https://git.kernel.org/pub/scm/linux/kernel/git/torvalds/linux.git/commit/?id=a528d35e8bfcc521d7cb70aaf03e1bd296c8493f) 


% TODO: Trattare esempio di inzializzazione di file e successivo collegamento
% con l'uso di O_TMPFILE e linkat, vedi man open

% TODO: manca prototipo e motivazione di fexecve, da trattare qui in quanto
% inserita nello stesso standard e da usare con openat, vedi 
% http://pubs.opengroup.org/onlinepubs/9699939699/toc.pdf
% TODO manca prototipo di execveat, introdotta nel 3.19, vedi
% https://lwn.net/Articles/626150/ cerca anche fexecve
% TODO: manca prototipo e motivazione di execveat, vedi
% http://man7.org/linux/man-pages/man2/execveat.2.html 

% TODO manca prototipo di renameat2, introdotta nel 3.15, vedi
% http://lwn.net/Articles/569134/ 


% TODO: trattare i nuovi AT_flags quando e se arriveranno, vedi
% https://lwn.net/Articles/767547/ 



\itindend{at-functions}


\subsection{Le operazioni di controllo}
\label{sec:file_fcntl_ioctl}

Oltre alle operazioni base esaminate in sez.~\ref{sec:file_unix_interface}
esistono tutta una serie di operazioni ausiliarie che è possibile eseguire su
un file descriptor, che non riguardano la normale lettura e scrittura di dati,
ma la gestione sia delle loro proprietà, che di tutta una serie di ulteriori
funzionalità che il kernel può mettere a disposizione.

% TODO: trattare qui i file seal 

Per le operazioni di manipolazione e di controllo delle varie proprietà e
caratteristiche di un file descriptor, viene usata la funzione di sistema
\funcd{fcntl},\footnote{ad esempio si gestiscono con questa funzione varie
  modalità di I/O asincrono (vedi sez.~\ref{sec:file_asyncronous_operation}) e
  il \textit{file locking} (vedi sez.~\ref{sec:file_locking}).} il cui
prototipo è:

\begin{funcproto}{
\fhead{unistd.h}
\fhead{fcntl.h}
\fdecl{int fcntl(int fd, int cmd)}
\fdecl{int fcntl(int fd, int cmd, long arg)}
\fdecl{int fcntl(int fd, int cmd, struct flock * lock)}
\fdecl{int fcntl(int fd, int cmd, struct f\_owner\_ex * owner)}
\fdesc{Esegue una operazione di controllo sul file.} 
}

{La funzione ha valori di ritorno diversi a seconda dell'operazione richiesta
  in caso di successo mentre ritorna sempre $-1$ per un errore, nel qual caso
  \var{errno} assumerà valori diversi che dipendono dal tipo di operazione,
  l'unico valido in generale è:
  \begin{errlist}
  \item[\errcode{EBADF}] \param{fd} non è un file aperto.
  \end{errlist}
}  
\end{funcproto}

Il primo argomento della funzione è sempre il numero di file descriptor
\var{fd} su cui si vuole operare. Il comportamento di questa funzione, il
numero e il tipo degli argomenti, il valore di ritorno e gli eventuali errori
aggiuntivi, sono determinati dal valore dell'argomento \param{cmd} che in
sostanza corrisponde all'esecuzione di un determinato \textsl{comando}. A
seconda del comando specificato il terzo argomento può essere assente (ma se
specificato verrà ignorato), può assumere un valore intero di tipo
\ctyp{long}, o essere un puntatore ad una struttura \struct{flock}.

In sez.~\ref{sec:file_dup} abbiamo incontrato un esempio dell'uso di
\func{fcntl} per la duplicazione dei file descriptor, una lista di tutti i
possibili valori per \var{cmd}, e del relativo significato, dei codici di
errore restituiti e del tipo del terzo argomento (cui faremo riferimento con
il nome indicato nel precedente prototipo), è riportata di seguito:
\begin{basedescript}{\desclabelwidth{1.8cm}}
\item[\constd{F\_DUPFD}] trova il primo file descriptor disponibile di valore
  maggiore o uguale ad \param{arg}, e ne fa un duplicato
  di \param{fd}, ritorna il nuovo file descriptor in caso di successo e $-1$
  in caso di errore. Oltre a \errval{EBADF} gli errori possibili sono
  \errcode{EINVAL} se \param{arg} è negativo o maggiore del massimo consentito
  o \errcode{EMFILE} se il processo ha già raggiunto il massimo numero di
  descrittori consentito.

\itindbeg{close-on-exec}

\item[\constd{F\_DUPFD\_CLOEXEC}] ha lo stesso effetto di \const{F\_DUPFD}, ma
  in più attiva il flag di \textit{close-on-exec} sul file descriptor
  duplicato, in modo da evitare una successiva chiamata con
  \const{F\_SETFD}. La funzionalità è stata introdotta con il kernel 2.6.24 ed
  è prevista nello standard POSIX.1-2008 (si deve perciò definire
  \macro{\_POSIX\_C\_SOURCE} ad un valore adeguato secondo quanto visto in
  sez.~\ref{sec:intro_gcc_glibc_std}).

\item[\constd{F\_GETFD}] restituisce il valore dei \textit{file descriptor
    flags} di \param{fd} in caso di successo o $-1$ in caso di errore, il
  terzo argomento viene ignorato. Non sono previsti errori diversi da
  \errval{EBADF}. Al momento l'unico flag usato è quello di
  \textit{close-on-exec}, identificato dalla costante \const{FD\_CLOEXEC}, che
  serve a richiedere che il file venga chiuso nella esecuzione di una
  \func{exec} (vedi sez.~\ref{sec:proc_exec}). Un valore nullo significa
  pertanto che il flag non è impostato.

\item[\constd{F\_SETFD}] imposta il valore dei \textit{file descriptor flags}
  al valore specificato con \param{arg}, ritorna un valore nullo in caso di
  successo e $-1$ in caso di errore. Non sono previsti errori diversi da
  \errval{EBADF}. Dato che l'unico flag attualmente usato è quello di
  \textit{close-on-exec}, identificato dalla costante \const{FD\_CLOEXEC},
  tutti gli altri bit di \param{arg}, anche se impostati, vengono
  ignorati.\footnote{questo almeno è quanto avviene fino al kernel 3.2, come
    si può evincere dal codice della funzione \texttt{do\_fcntl} nel file
    \texttt{fs/fcntl.c} dei sorgenti del kernel.}
\itindend{close-on-exec}

\item[\constd{F\_GETFL}] ritorna il valore dei \textit{file status flags} di
  \param{fd} in caso di successo o $-1$ in caso di errore, il terzo argomento
  viene ignorato. Non sono previsti errori diversi da \errval{EBADF}. Il
  comando permette di rileggere il valore di quei bit
  dell'argomento \param{flags} di \func{open} che vengono memorizzati nella
  relativa voce della \textit{file table} all'apertura del file, vale a dire
  quelli riportati in tab.~\ref{tab:open_access_mode_flag} e
  tab.~\ref{tab:open_operation_flag}). Si ricordi che quando si usa la
  funzione per determinare le modalità di accesso con cui è stato aperto il
  file è necessario estrarre i bit corrispondenti nel \textit{file status
    flag} con la maschera \const{O\_ACCMODE} come già accennato in
  sez.~\ref{sec:file_open_close}. 

\item[\constd{F\_SETFL}] imposta il valore dei \textit{file status flags} al
  valore specificato da \param{arg}, ritorna un valore nullo in caso di
  successo o $-1$ in caso di errore. In generale possono essere impostati solo
  i flag riportati in tab.~\ref{tab:open_operation_flag}, su Linux si possono
  modificare soltanto \const{O\_APPEND}, \const{O\_ASYNC}, \const{O\_DIRECT},
  \const{O\_NOATIME} e \const{O\_NONBLOCK}. Oltre a \errval{EBADF} si otterrà
  \errcode{EPERM} se si cerca di rimuovere \const{O\_APPEND} da un file
  marcato come \textit{append-only} o se di cerca di impostare
  \const{O\_NOATIME} su un file di cui non si è proprietari (e non si hanno i
  permessi di amministratore) ed \errcode{EINVAL} se si cerca di impostare
  \const{O\_DIRECT} su un file che non supporta questo tipo di operazioni.

\item[\constd{F\_GETLK}] richiede un controllo sul file lock specificato da
  \param{lock}, sovrascrivendo la struttura da esso puntata con il risultato,
  ritorna un valore nullo in caso di successo o $-1$ in caso di errore. Come
  per i due successivi comandi oltre a \errval{EBADF} se \param{lock} non è un
  puntatore valido restituisce l'errore generico \errcode{EFAULT}. Questa
  funzionalità è trattata in dettaglio in sez.~\ref{sec:file_posix_lock}.

\item[\constd{F\_SETLK}] richiede o rilascia un file lock a seconda di quanto
  specificato nella struttura puntata da \param{lock}, ritorna un valore nullo
  in caso di successo e $-1$ se il file lock è tenuto da qualcun altro, nel
  qual caso si ha un errore di \errcode{EACCES} o \errcode{EAGAIN}.  Questa
  funzionalità è trattata in dettaglio in sez.~\ref{sec:file_posix_lock}.

\item[\constd{F\_SETLKW}] identica a \const{F\_SETLK} eccetto per il fatto che
  la funzione non ritorna subito ma attende che il blocco sia rilasciato, se
  l'attesa viene interrotta da un segnale la funzione restituisce $-1$ e
  imposta \var{errno} a \errcode{EINTR}.  Questa funzionalità è trattata in
  dettaglio in sez.~\ref{sec:file_posix_lock}.

\item[\constd{F\_GETOWN}] restituisce in caso di successo l'identificatore del
  processo o del \textit{process group} (vedi sez.~\ref{sec:sess_proc_group})
  che è preposto alla ricezione del segnale \signal{SIGIO} (o l'eventuale
  segnale alternativo impostato con \const{F\_SETSIG}) per gli eventi
  asincroni associati al file descriptor \param{fd} e del segnale
  \signal{SIGURG} per la notifica dei dati urgenti di un socket (vedi
  sez.~\ref{sec:TCP_urgent_data}). Restituisce $-1$ in caso di errore ed il
  terzo argomento viene ignorato. Non sono previsti errori diversi da
  \errval{EBADF}.

  Per distinguerlo dal caso in cui il segnale viene inviato a un singolo
  processo, nel caso di un \textit{process group} viene restituito un valore
  negativo il cui valore assoluto corrisponde all'identificatore del
  \textit{process group}. Con Linux questo comporta un problema perché se il
  valore restituito dalla \textit{system call} è compreso nell'intervallo fra
  $-1$ e $-4095$ in alcune architetture questo viene trattato dalla
  \acr{glibc} come un errore,\footnote{il problema deriva dalle limitazioni
    presenti in architetture come quella dei normali PC (i386) per via delle
    modalità in cui viene effettuata l'invocazione delle \textit{system call}
    che non consentono di restituire un adeguato codice di ritorno.} per cui
  in tal caso \func{fcntl} ritornerà comunque $-1$ mentre il valore restituito
  dalla \textit{system call} verrà assegnato ad \var{errno}, cambiato di
  segno.

  Per questo motivo con il kernel 2.6.32 è stato introdotto il comando
  alternativo \const{F\_GETOWN\_EX}, che vedremo a breve, che consente di
  evitare il problema. A partire dalla versione 2.11 la \acr{glibc}, se
  disponibile, usa questa versione alternativa per mascherare il problema
  precedente e restituire un valore corretto in tutti i casi.\footnote{in cui
    cioè viene restituito un valore negativo corretto qualunque sia
    l'identificatore del \textit{process group}, che non potendo avere valore
    unitario (non esiste infatti un \textit{process group} per \cmd{init}) non
    può generare ambiguità con il codice di errore.} Questo però comporta che
  il comportamento del comando può risultare diverso a seconda delle versioni
  della \acr{glibc} e del kernel.

\item[\constd{F\_SETOWN}] imposta, con il valore dell'argomento \param{arg},
  l'identificatore del processo o del \textit{process group} che riceverà i
  segnali \signal{SIGIO} e \signal{SIGURG} per gli eventi associati al file
  descriptor \param{fd}. Ritorna un valore nullo in caso di successo o $-1$ in
  caso di errore. Oltre a \errval{EBADF} gli errori possibili sono
  \errcode{ESRCH} se \param{arg} indica un processo o un \textit{process
    group} inesistente.

  L'impostazione è soggetta alle stesse restrizioni presenti sulla funzione
  \func{kill} (vedi sez.~\ref{sec:sig_kill_raise}), per cui un utente non
  privilegiato può inviare i segnali solo ad un processo che gli appartiene,
  in genere comunque si usa il processo corrente.  Come per \const{F\_GETOWN},
  per indicare un \textit{process group} si deve usare per \param{arg} un
  valore negativo, il cui valore assoluto corrisponda all'identificatore del
  \textit{process group}.

  A partire dal kernel 2.6.12 se si sta operando con i \textit{thread} della
  implementazione nativa di Linux (quella della NTPL, vedi
  sez.~\ref{sec:linux_ntpl}) e se si è impostato un segnale specifico con
  \const{F\_SETSIG}, un valore positivo di \param{arg} viene interpretato come
  indicante un \textit{Thread ID} e non un \textit{Process ID}.  Questo
  consente di inviare il segnale impostato con \const{F\_SETSIG} ad uno
  specifico \textit{thread}. In genere questo non comporta differenze
  significative per il processi ordinari, in cui non esistono altri
  \textit{thread}, dato che su Linux il \textit{thread} principale, che in tal
  caso è anche l'unico, mantiene un valore del \textit{Thread ID} uguale al
  \ids{PID} del processo. Il problema è però che questo comportamento non si
  applica a \signal{SIGURG}, per il quale \param{arg} viene sempre
  interpretato come l'identificatore di un processo o di un \textit{process
    group}.

\item[\constd{F\_GETOWN\_EX}] legge nella struttura puntata
  dall'argomento \param{owner} l'identificatore del processo, \textit{thread}
  o \textit{process group} (vedi sez.~\ref{sec:sess_proc_group}) che è
  preposto alla ricezione dei segnali \signal{SIGIO} e \signal{SIGURG} per gli
  eventi associati al file descriptor \param{fd}.  Ritorna un valore nullo in
  caso di successo o $-1$ in caso di errore. Oltre a \errval{EBADF} e da
  \errval{EFAULT} se \param{owner} non è un puntatore valido.

  Il comando, che è disponibile solo a partire dal kernel 2.6.32, effettua lo
  stesso compito di \const{F\_GETOWN} di cui costituisce una evoluzione che
  consente di superare i limiti e le ambiguità relative ai valori restituiti
  come identificativo. A partire dalla versione 2.11 della \acr{glibc} esso
  viene usato dalla libreria per realizzare una versione di \func{fcntl} che
  non presenti i problemi illustrati in precedenza per la versione precedente
  di \const{F\_GETOWN}.  Il comando è specifico di Linux ed utilizzabile solo
  se si è definita la macro \macro{\_GNU\_SOURCE}.

\item[\constd{F\_SETOWN\_EX}] imposta con il valore della struttura
  \struct{f\_owner\_ex} puntata \param{owner}, l'identificatore del processo o
  del \textit{process group} che riceverà i segnali \signal{SIGIO} e
  \signal{SIGURG} per gli eventi associati al file
  descriptor \param{fd}. Ritorna un valore nullo in caso di successo o $-1$ in
  caso di errore, con gli stessi errori di \const{F\_SETOWN} più
  \errcode{EINVAL} se il campo \var{type} di \struct{f\_owner\_ex} non indica
  un tipo di identificatore valido.

  \begin{figure}[!htb]
    \footnotesize \centering
    \begin{varwidth}[c]{0.5\textwidth}
      \includestruct{listati/f_owner_ex.h}
    \end{varwidth}
    \normalsize 
    \caption{La struttura \structd{f\_owner\_ex}.} 
    \label{fig:f_owner_ex}
  \end{figure}

  Come \const{F\_GETOWN\_EX} il comando richiede come terzo argomento il
  puntatore ad una struttura \struct{f\_owner\_ex} la cui definizione è
  riportata in fig.~\ref{fig:f_owner_ex}, in cui il primo campo indica il tipo
  di identificatore il cui valore è specificato nel secondo campo, che assume
  lo stesso significato di \param{arg} per \const{F\_SETOWN}. Per il campo
  \var{type} i soli valori validi sono \constd{F\_OWNER\_TID},
  \constd{F\_OWNER\_PID} e \constd{F\_OWNER\_PGRP}, che indicano
  rispettivamente che si intende specificare con \var{pid} un \textit{Tread
    ID}, un \textit{Process ID} o un \textit{Process Group ID}. A differenza
  di \const{F\_SETOWN} se si specifica un \textit{Tread ID} questo riceverà
  sia \signal{SIGIO} (o il segnale impostato con \const{F\_SETSIG}) che
  \signal{SIGURG}. Il comando è specifico di Linux, è disponibile solo a
  partire dal kernel 2.6.32, ed è utilizzabile solo se si è definita la macro
  \macro{\_GNU\_SOURCE}.

\item[\constd{F\_GETSIG}] restituisce il valore del segnale inviato dai vari
  meccanismi di I/O asincrono associati al file descriptor \param{fd} (quelli
  trattati in sez.~\ref{sec:file_asyncronous_operation}) in caso di successo o
  $-1$ in caso di errore, il terzo argomento viene ignorato. Non sono previsti
  errori diversi da \errval{EBADF}.  Un valore nullo indica che si sta usando
  il segnale predefinito, che è \signal{SIGIO}. Un valore diverso da zero
  indica il segnale che è stato impostato con \const{F\_SETSIG}, che può
  essere anche lo stesso \signal{SIGIO}. Il comando è specifico di Linux ed
  utilizzabile solo se si è definita la macro \macro{\_GNU\_SOURCE}.

\item[\constd{F\_SETSIG}] imposta il segnale inviato dai vari meccanismi di
  I/O asincrono associati al file descriptor \param{fd} (quelli trattati in
  sez.~\ref{sec:file_asyncronous_operation}) al valore indicato
  da \param{arg}, ritorna un valore nullo in caso di successo o $-1$ in caso
  di errore.  Oltre a \errval{EBADF} gli errori possibili sono
  \errcode{EINVAL} se \param{arg} indica un numero di segnale non valido.  Un
  valore nullo di \param{arg} indica di usare il segnale predefinito, cioè
  \signal{SIGIO}. Un valore diverso da zero, compreso lo stesso
  \signal{SIGIO}, specifica il segnale voluto.  Il comando è specifico di
  Linux ed utilizzabile solo se si è definita la macro \macro{\_GNU\_SOURCE}.

  L'impostazione di un valore diverso da zero permette inoltre, se si è
  installato il gestore del segnale come \var{sa\_sigaction} usando
  \const{SA\_SIGINFO}, (vedi sez.~\ref{sec:sig_sigaction}), di rendere
  disponibili al gestore informazioni ulteriori riguardo il file che ha
  generato il segnale attraverso i valori restituiti in
  \struct{siginfo\_t}. Se inoltre si imposta un segnale \textit{real-time} si
  potranno sfruttare le caratteristiche di avanzate di questi ultimi (vedi
  sez.~\ref{sec:sig_real_time}), ed in particolare la capacità di essere
  accumulati in una coda prima della notifica.

\item[\constd{F\_GETLEASE}] restituisce il tipo di \textit{file lease} che il
  processo detiene nei confronti del file descriptor \var{fd} o $-1$ in caso
  di errore, il terzo argomento viene ignorato. Non sono previsti errori
  diversi da \errval{EBADF}.  Il comando è specifico di Linux ed utilizzabile
  solo se si è definita la macro \macro{\_GNU\_SOURCE}.  Questa funzionalità è
  trattata in dettaglio in sez.~\ref{sec:file_asyncronous_lease}.

\item[\constd{F\_SETLEASE}] imposta o rimuove a seconda del valore
  di \param{arg} un \textit{file lease} sul file descriptor \var{fd} a seconda
  del valore indicato da \param{arg}. Ritorna un valore nullo in caso di
  successo o $-1$ in caso di errore. Oltre a \errval{EBADF} si otterrà
  \errcode{EINVAL} se si è specificato un valore non valido per \param{arg}
  (deve essere usato uno dei valori di tab.~\ref{tab:file_lease_fctnl}),
  \errcode{ENOMEM} se non c'è memoria sufficiente per creare il \textit{file
    lease}, \errcode{EACCES} se non si è il proprietario del file e non si
  hanno i privilegi di amministratore.\footnote{per la precisione occorre la
    capacità \const{CAP\_LEASE}.}

  Il supporto il supporto per i \textit{file lease}, che consente ad un
  processo che detiene un \textit{lease} su un file di riceve una notifica
  qualora un altro processo cerchi di eseguire una \func{open} o una
  \func{truncate} su di esso è stato introdotto a partire dai kernel della
  serie 2.4 Il comando è specifico di Linux ed utilizzabile solo se si è
  definita la macro \macro{\_GNU\_SOURCE}. Questa funzionalità è trattata in
  dettaglio in sez.~\ref{sec:file_asyncronous_lease}.

\item[\constd{F\_NOTIFY}] attiva il meccanismo di notifica asincrona per cui
  viene riportato al processo chiamante, tramite il segnale \signal{SIGIO} (o
  altro segnale specificato con \const{F\_SETSIG}) ogni modifica eseguita o
  direttamente sulla directory cui \var{fd} fa riferimento, o su uno dei file
  in essa contenuti; ritorna un valore nullo in caso di successo o $-1$ in
  caso di errore. Il comando è specifico di Linux ed utilizzabile solo se si è
  definita la macro \macro{\_GNU\_SOURCE}.  Questa funzionalità, disponibile
  dai kernel della serie 2.4.x, è trattata in dettaglio in
  sez.~\ref{sec:file_asyncronous_lease}.

\item[\constd{F\_GETPIPE\_SZ}] restituisce in caso di successo la dimensione
  del buffer associato alla \textit{pipe} \param{fd} (vedi
  sez.~\ref{sec:ipc_pipes}) o $-1$ in caso di errore, il terzo argomento viene
  ignorato. Non sono previsti errori diversi da \errval{EBADF}, che viene
  restituito anche se il file descriptor non è una \textit{pipe}. Il comando è
  specifico di Linux, è disponibile solo a partire dal kernel 2.6.35, ed è
  utilizzabile solo se si è definita la macro \macro{\_GNU\_SOURCE}.

\item[\constd{F\_SETPIPE\_SZ}] imposta la dimensione del buffer associato alla
  \textit{pipe} \param{fd} (vedi sez.~\ref{sec:ipc_unix}) ad un valore uguale
  o superiore a quello indicato dall'argomento \param{arg}. Ritorna un valore
  nullo in caso di successo o $-1$ in caso di errore. Oltre a \errval{EBADF}
  gli errori possibili sono \errcode{EBUSY} se si cerca di ridurre la
  dimensione del buffer al di sotto della quantità di dati effettivamente
  presenti su di esso ed \errcode{EPERM} se un processo non privilegiato cerca
  di impostare un valore troppo alto.  La dimensione minima del buffer è pari
  ad una pagina di memoria, a cui verrà comunque arrotondata ogni dimensione
  inferiore, il valore specificato viene in genere arrotondato per eccesso al
  valore ritenuto più opportuno dal sistema, pertanto una volta eseguita la
  modifica è opportuno rileggere la nuova dimensione con
  \const{F\_GETPIPE\_SZ}. I processi non privilegiati\footnote{per la
    precisione occorre la capacità \const{CAP\_SYS\_RESOURCE}.} non possono
  impostare un valore superiore a quello indicato da
  \sysctlfiled{fs/pipe-size-max}.  Il comando è specifico di Linux, è
  disponibile solo a partire dal kernel 2.6.35, ed è utilizzabile solo se si è
  definita la macro \macro{\_GNU\_SOURCE}.

\end{basedescript}

% TODO: trattare RWH_WRITE_LIFE_EXTREME e RWH_WRITE_LIFE_SHORT aggiunte con
% il kernel 4.13 (vedi https://lwn.net/Articles/727385/)

La maggior parte delle funzionalità controllate dai comandi di \func{fcntl}
sono avanzate e richiedono degli approfondimenti ulteriori, saranno pertanto
riprese più avanti quando affronteremo le problematiche ad esse relative. In
particolare le tematiche relative all'I/O asincrono e ai vari meccanismi di
notifica saranno trattate in maniera esaustiva in
sez.~\ref{sec:file_asyncronous_operation} mentre quelle relative al
\textit{file locking} saranno esaminate in sez.~\ref{sec:file_locking}). L'uso
di questa funzione con i socket verrà trattato in
sez.~\ref{sec:sock_ctrl_func}.

La gran parte dei comandi di \func{fcntl} (come \const{F\_DUPFD},
\const{F\_GETFD}, \const{F\_SETFD}, \const{F\_GETFL}, \const{F\_SETFL},
\const{F\_GETLK}, \const{F\_SETLK} e \const{F\_SETLKW}) sono previsti da SVr4
e 4.3BSD e standardizzati in POSIX.1-2001 che inoltre prevede gli ulteriori
\const{F\_GETOWN} e \const{F\_SETOWN}. Pertanto nell'elenco si sono indicate
esplicitamente soltanto le ulteriori richieste in termini delle macro di
funzionalità di sez.~\ref{sec:intro_gcc_glibc_std} soltanto per le
funzionalità inserite in standard successivi o specifiche di Linux.


% \subsection{La funzione \func{ioctl}}
% \label{sec:file_ioctl}

Benché l'interfaccia di gestione dell'I/O sui file di cui abbiamo parlato
finora si sia dimostrata valida anche per l'interazione diretta con le
periferiche attraverso i loro file di dispositivo, consentendo di usare le
stesse funzioni utilizzate per i normali file di dati, esistono però
caratteristiche peculiari, specifiche dell'hardware e delle funzionalità che
ciascun dispositivo può provvedere, che non possono venire comprese in questa
interfaccia astratta come ad esempio l'impostazione della velocità di una
porta seriale, o le dimensioni di un framebuffer.

Per questo motivo nell'architettura del sistema è stata prevista l'esistenza
di una apposita funzione di sistema, \funcd{ioctl}, come meccanismo generico
per compiere operazioni specializzate; il suo prototipo è:

\begin{funcproto}{
\fhead{sys/ioctl.h}
\fdecl{int ioctl(int fd, int request, ...)}
\fdesc{Esegue una operazione speciale.} 
}

{La funzione ritorna $0$ in caso di successo nella maggior parte dei casi, ma
  alcune operazioni possono restituire un valore positivo, mentre ritorna
  sempre $-1$ per un errore, nel qual caso \var{errno} assumerà uno dei
  valori:
  \begin{errlist}
  \item[\errcode{EINVAL}] gli argomenti \param{request} o \param{argp} non sono
    validi.
  \item[\errcode{ENOTTY}] il file \param{fd} non è associato con un
    dispositivo, o la richiesta non è applicabile all'oggetto a cui fa
    riferimento \param{fd}.
  \end{errlist}
  ed inoltre \errval{EBADF} e \errval{EFAULT} nel loro significato generico.}
\end{funcproto}


La funzione richiede che si passi come primo argomento un file
descriptor \param{fd} regolarmente aperto, mentre l'operazione da compiere
deve essere indicata dal valore dell'argomento \param{request}. Il terzo
argomento dipende dall'operazione prescelta; tradizionalmente è specificato
come \code{char * argp}, da intendersi come puntatore ad un area di memoria
generica (all'epoca della creazione di questa funzione infatti ancora non era
stato introdotto il tipo \ctyp{void}) ma per certe operazioni può essere
omesso, e per altre è un semplice intero.

Normalmente la funzione ritorna zero in caso di successo e $-1$ in caso di
errore, ma per alcune operazioni il valore di ritorno, che nel caso viene
impostato ad un valore positivo, può essere utilizzato come indicazione del
risultato della stessa. È più comune comunque restituire i risultati
all'indirizzo puntato dal terzo argomento.

Data la genericità dell'interfaccia non è possibile classificare in maniera
sistematica le operazioni che si possono gestire con \func{ioctl}, un breve
elenco di alcuni esempi di esse è il seguente:
\begin{itemize*}
\item il cambiamento dei font di un terminale.
\item l'esecuzione di una traccia audio di un CDROM.
\item i comandi di avanti veloce e di riavvolgimento di un nastro.
\item il comando di espulsione di un dispositivo rimovibile.
\item l'impostazione della velocità trasmissione di una linea seriale.
\item l'impostazione della frequenza e della durata dei suoni emessi dallo
  speaker.
\item l'impostazione degli attributi dei file su un filesystem
  ext2.\footnote{i comandi \texttt{lsattr} e \texttt{chattr} fanno questo con
    delle \func{ioctl} dedicate, usabili solo su questo filesystem e derivati
    successivi (come ext3).}
\end{itemize*}

In generale ogni dispositivo ha un suo insieme di operazioni specifiche
effettuabili attraverso \func{ioctl}, tutte queste sono definite nell'header
file \headfiled{sys/ioctl.h}, e devono essere usate solo sui dispositivi cui
fanno riferimento. Infatti anche se in genere i valori di \param{request} sono
opportunamente differenziati a seconda del dispositivo\footnote{il kernel usa
  un apposito \textit{magic number} per distinguere ciascun dispositivo nella
  definizione delle macro da usare per \param{request}, in modo da essere
  sicuri che essi siano sempre diversi, ed il loro uso per dispositivi diversi
  causi al più un errore.  Si veda il capitolo quinto di \cite{LinDevDri} per
  una trattazione dettagliata dell'argomento.} così che la richiesta di
operazioni relative ad altri dispositivi usualmente provoca il ritorno della
funzione con una condizione di errore, in alcuni casi, relativi a valori
assegnati prima che questa differenziazione diventasse pratica corrente, si
potrebbero usare valori validi anche per il dispositivo corrente, con effetti
imprevedibili o indesiderati.

Data la assoluta specificità della funzione, il cui comportamento varia da
dispositivo a dispositivo, non è possibile fare altro che dare una descrizione
sommaria delle sue caratteristiche; torneremo ad esaminare in seguito quelle
relative ad alcuni casi specifici, ad esempio la gestione dei terminali è
effettuata attraverso \func{ioctl} in quasi tutte le implementazioni di Unix,
mentre per l'uso di \func{ioctl} con i socket si veda
sez.~\ref{sec:sock_ctrl_func}. 

Riportiamo qui solo l'elenco delle operazioni che sono predefinite per
qualunque file, caratterizzate dal prefisso \texttt{FIO}. Queste operazioni
sono definite nel kernel a livello generale, e vengono sempre interpretate per
prime, per cui, come illustrato in \cite{LinDevDri}, eventuali operazioni
specifiche che usino lo stesso valore verrebbero ignorate:
\begin{basedescript}{\desclabelwidth{2.0cm}}
\item[\constd{FIOCLEX}] imposta il flag di \textit{close-on-exec} sul file, in
  questo caso, essendo usata come operazione logica, \func{ioctl} non richiede
  un terzo argomento, il cui eventuale valore viene ignorato.
\item[\constd{FIONCLEX}] cancella il flag di \textit{close-on-exec} sul file,
  in questo caso, essendo usata come operazione logica, \func{ioctl} non
  richiede un terzo argomento, il cui eventuale valore viene ignorato.
\item[\constd{FIOASYNC}] abilita o disabilita la modalità di I/O asincrono sul
  file (vedi sez.~\ref{sec:signal_driven_io}); il terzo argomento
  deve essere un puntatore ad un intero (cioè di tipo \texttt{const int *})
  che contiene un valore logico (un valore nullo disabilita, un valore non
  nullo abilita).
\item[\constd{FIONBIO}] abilita o disabilita sul file l'I/O in modalità non
  bloccante; il terzo argomento deve essere un puntatore ad un intero (cioè di
  tipo \texttt{const int *}) che contiene un valore logico (un valore nullo
  disabilita, un valore non nullo abilita).
\item[\constd{FIOSETOWN}] imposta il processo che riceverà i segnali
  \signal{SIGURG} e \signal{SIGIO} generati sul file; il terzo argomento deve
  essere un puntatore ad un intero (cioè di tipo \texttt{const int *}) il cui
  valore specifica il PID del processo.
\item[\constd{FIOGETOWN}] legge il processo che riceverà i segnali
  \signal{SIGURG} e \signal{SIGIO} generati sul file; il terzo argomento deve
  essere un puntatore ad un intero (cioè di tipo \texttt{int *}) su cui sarà
  scritto il PID del processo.
\item[\constd{FIONREAD}] legge il numero di byte disponibili in lettura sul
  file descriptor; questa operazione è disponibile solo su alcuni file
  descriptor, in particolare sui socket (vedi sez.~\ref{sec:sock_ioctl_IP}) o
  sui file descriptor di \textit{epoll} (vedi sez.~\ref{sec:file_epoll}), il
  terzo argomento deve essere un puntatore ad un intero (cioè di tipo
  \texttt{int *}) su cui sarà restituito il valore.
\item[\constd{FIOQSIZE}] restituisce la dimensione corrente di un file o di una
  directory, mentre se applicata ad un dispositivo fallisce con un errore di
  \errcode{ENOTTY}; il terzo argomento deve essere un puntatore ad un intero
  (cioè di tipo \texttt{int *}) su cui sarà restituito il valore.
\end{basedescript}

% TODO aggiungere FIBMAP e FIEMAP, vedi http://lwn.net/Articles/260795/,
% http://lwn.net/Articles/429345/ 

Si noti però come la gran parte di queste operazioni specifiche dei file (per
essere precisi le prime sei dell'elenco) siano effettuabili in maniera
generica anche tramite l'uso di \func{fcntl}. Le due funzioni infatti sono
molto simili e la presenza di questa sovrapposizione è principalmente dovuta
al fatto che alle origini di Unix i progettisti considerarono che era
necessario trattare diversamente rispetto alle operazione di controllo delle
modalità di I/O file e dispositivi usando \func{fcntl} per i primi e
\func{ioctl} per i secondi, all'epoca tra l'altro i dispositivi che usavano
\func{ioctl} erano sostanzialmente solo i terminali, il che spiega l'uso
comune di \errcode{ENOTTY} come codice di errore. Oggi non è più così ma le
due funzioni sono rimaste.

% TODO trovare qualche posto per la eventuale documentazione delle seguenti
% (bassa/bassissima priorità)
% EXT4_IOC_MOVE_EXT (dal 2.6.31)
%  EXT4_IOC_SHUTDOWN (dal 4.10), XFS_IOC_GOINGDOWN e futura FS_IOC_SHUTDOWN
% ioctl di btrfs, vedi http://lwn.net/Articles/580732/

% \chapter{}

\section{L'interfaccia standard ANSI C}
\label{sec:files_std_interface}


Come visto in sez.~\ref{sec:file_unix_interface} le operazioni di I/O sui file
sono gestibili a basso livello con l'interfaccia standard unix, che ricorre
direttamente alle \textit{system call} messe a disposizione dal kernel.

Questa interfaccia però non provvede le funzionalità previste dallo standard
ANSI C, che invece sono realizzate attraverso opportune funzioni di libreria.
Queste funzioni di libreria, insieme alle altre funzioni definite dallo
standard (che sono state implementate la prima volta da Ritchie nel 1976 e da
allora sono rimaste sostanzialmente immutate), vengono a costituire il nucleo
della \acr{glibc} per la gestione dei file.

Esamineremo in questa sezione le funzioni base dell'interfaccia degli
\textit{stream}, analoghe a quelle di sez.~\ref{sec:file_unix_interface} per i
file descriptor. In particolare vedremo come aprire, leggere, scrivere e
cambiare la posizione corrente in uno \textit{stream}.


\subsection{I \textit{file stream}}
\label{sec:file_stream}

\itindbeg{file~stream}

Come più volte ribadito, l'interfaccia dei file descriptor è un'interfaccia di
basso livello, che non provvede nessuna forma di formattazione dei dati e
nessuna forma di bufferizzazione per ottimizzare le operazioni di I/O.

In \cite{APUE} Stevens descrive una serie di test sull'influenza delle
dimensioni del blocco di dati (l'argomento \param{buf} di \func{read} e
\func{write}) nell'efficienza nelle operazioni di I/O con i file descriptor,
evidenziando come le prestazioni ottimali si ottengano a partire da dimensioni
del buffer dei dati pari a quelle dei blocchi del filesystem (il valore dato
dal campo \var{st\_blksize} di \struct{stat}), che di norma corrispondono alle
dimensioni dei settori fisici in cui è suddiviso il disco.

Se il programmatore non si cura di effettuare le operazioni in blocchi di
dimensioni adeguate, le prestazioni sono inferiori.  La caratteristica
principale dell'interfaccia degli \textit{stream} è che essa provvede da sola
alla gestione dei dettagli della bufferizzazione e all'esecuzione delle
operazioni di lettura e scrittura in blocchi di dimensioni appropriate
all'ottenimento della massima efficienza.

Per questo motivo l'interfaccia viene chiamata anche interfaccia dei
\textit{file stream}, dato che non è più necessario doversi preoccupare dei
dettagli con cui viene gestita la comunicazione con l'hardware sottostante
(come nel caso della dimensione dei blocchi del filesystem), ed un file può
essere sempre considerato come composto da un flusso continuo di dati, da cui
deriva appunto il nome \textit{stream}.

A parte i dettagli legati alla gestione delle operazioni di lettura e
scrittura, sia per quel che riguarda la bufferizzazione che le formattazioni,
per tutto il resto i \textit{file stream} restano del tutto equivalenti ai
file descriptor (sui quali sono basati), ed in particolare continua a valere
quanto visto in sez.~\ref{sec:file_shared_access} a proposito dell'accesso
concorrente ed in sez.~\ref{sec:file_access_control} per il controllo di
accesso.

Per ragioni storiche la struttura di dati che rappresenta uno \textit{stream}
è stata chiamata \typed{FILE}, questi oggetti sono creati dalle funzioni di
libreria e contengono tutte le informazioni necessarie a gestire le operazioni
sugli \textit{stream}, come la posizione corrente, lo stato del buffer e degli
indicatori di stato e di fine del file.

Per questo motivo gli utenti non devono mai utilizzare direttamente o allocare
queste strutture (che sono dei \textsl{tipi opachi}) ma usare sempre puntatori
del tipo \texttt{FILE *} ottenuti dalla libreria stessa, tanto che in certi
casi il termine di puntatore a file è diventato sinonimo di \textit{stream}.
Tutte le funzioni della libreria che operano sui file accettano come argomenti
solo variabili di questo tipo, che diventa accessibile includendo l'header
file \headfile{stdio.h}.

\itindend{file~stream}

Ai tre file descriptor standard (vedi tab.~\ref{tab:file_std_files}) aperti
per ogni processo, corrispondono altrettanti \textit{stream}, che
rappresentano i canali standard di input/output prestabiliti; anche questi tre
\textit{stream} sono identificabili attraverso dei nomi simbolici definiti
nell'header \headfile{stdio.h} che sono:

\begin{basedescript}{\desclabelwidth{3.0cm}}
\item[\var{FILE *stdin}] Lo \textit{standard input} cioè il \textit{file
    stream} da cui il processo riceve ordinariamente i dati in
  ingresso. Normalmente è associato dalla shell all'input del terminale e
  prende i caratteri dalla tastiera.
\item[\var{FILE *stdout}] Lo \textit{standard output} cioè il \textit{file
    stream} su cui il processo invia ordinariamente i dati in
  uscita. Normalmente è associato dalla shell all'output del terminale e
  scrive sullo schermo.
\item[\var{FILE *stderr}] Lo \textit{standard error} cioè il \textit{file
    stream} su cui il processo è supposto inviare i messaggi di
  errore. Normalmente anch'esso è associato dalla shell all'output del
  terminale e scrive sullo schermo.
\end{basedescript}

Nella \acr{glibc} \var{stdin}, \var{stdout} e \var{stderr} sono effettivamente
tre variabili di tipo \type{FILE}\texttt{ *} che possono essere usate come
tutte le altre, ad esempio si può effettuare una redirezione dell'output di un
programma con il semplice codice: \includecodesnip{listati/redir_stdout.c} ma
in altri sistemi queste variabili possono essere definite da macro, e se si
hanno problemi di portabilità e si vuole essere sicuri, diventa opportuno
usare la funzione \func{freopen}.


\subsection{Le modalità di bufferizzazione}
\label{sec:file_buffering}

La bufferizzazione è una delle caratteristiche principali dell'interfaccia
degli \textit{stream}; lo scopo è quello di ridurre al minimo il numero di
\textit{system call} (\func{read} o \func{write}) eseguite nelle operazioni di
input/output. Questa funzionalità è assicurata automaticamente dalla libreria,
ma costituisce anche uno degli aspetti più comunemente fraintesi, in
particolare per quello che riguarda l'aspetto della scrittura dei dati sul
file.

I dati che vengono scritti su di uno \textit{stream} normalmente vengono
accumulati in un buffer e poi trasmessi in blocco, con l'operazione che viene
usualmente chiamata \textsl{scaricamento} del buffer (dal termine inglese
\textit{flush}) tutte le volte che questo viene riempito. Questa operazione
avviene perciò in maniera asincrona rispetto alla scrittura. Un comportamento
analogo avviene anche in lettura (cioè dal file viene letto un blocco di dati,
anche se ne sono richiesti una quantità inferiore), ma la cosa ovviamente ha
rilevanza inferiore, dato che i dati letti sono sempre gli stessi. In caso di
scrittura invece, quando si ha un accesso contemporaneo allo stesso file (ad
esempio da parte di un altro processo) si potranno vedere solo le parti
effettivamente scritte, e non quelle ancora presenti nel buffer.

Per lo stesso motivo, in tutte le situazioni in cui si sta facendo
dell'input/output interattivo, bisognerà tenere presente le caratteristiche
delle operazioni di scaricamento dei dati, poiché non è detto che ad una
scrittura sullo \textit{stream} corrisponda una immediata scrittura sul
dispositivo, e la cosa è particolarmente evidente con le operazioni di
input/output sul terminale.

Per rispondere ad esigenze diverse lo standard definisce tre distinte modalità
in cui può essere eseguita la bufferizzazione, delle quali occorre essere ben
consapevoli, specie in caso di lettura e scrittura da dispositivi interattivi:
\begin{itemize}
\item \textit{unbuffered}: in questo caso non c'è bufferizzazione ed i
  caratteri vengono trasmessi direttamente al file non appena possibile
  (effettuando immediatamente una \func{write});
\item \textit{line buffered}: in questo caso i caratteri vengono normalmente
  trasmessi al file in blocco ogni volta che viene incontrato un carattere di
  \textit{newline} (il carattere ASCII \verb|\n|) cioè un a capo (in sostanza
  quando si preme invio);
\item \textit{fully buffered}: in questo caso i caratteri vengono
  trasmessi da e verso il file in blocchi di dimensione opportuna.
\end{itemize}

Lo standard ANSI C specifica inoltre che lo \textit{standard output} e lo
\textit{standard input} siano aperti in modalità \textit{fully buffered}
quando non fanno riferimento ad un dispositivo interattivo, e che lo standard
error non sia mai aperto in modalità \textit{fully buffered}.

Linux, come BSD e SVr4, specifica il comportamento predefinito in maniera
ancora più precisa, e cioè impone che lo standard error sia sempre
\textit{unbuffered}, in modo che i messaggi di errore siano mostrati il più
rapidamente possibile, e che \textit{standard input} \textit{standard output}
siano aperti in modalità \textit{line buffered} quando sono associati ad un
terminale (od altro dispositivo interattivo) ed in modalità \textit{fully
  buffered} altrimenti.

Il comportamento specificato per \textit{standard input} e \textit{standard
  output} vale anche per tutti i nuovi \textit{stream} aperti da un processo;
la selezione comunque avviene automaticamente, e la libreria apre lo
\textit{stream} nella modalità più opportuna a seconda del file o del
dispositivo scelto.

La modalità \textit{line buffered} è quella che necessita di maggiori
chiarimenti e attenzioni per quel che concerne il suo funzionamento. Come già
accennato nella descrizione, \emph{di norma} i dati vengono inviati al kernel
alla ricezione di un carattere di \textsl{a capo} (il \textit{newline});
questo non è vero in tutti i casi, infatti, dato che le dimensioni del buffer
usato dalle librerie sono fisse, se le si eccedono si può avere uno scarico
dei dati anche prima che sia stato inviato un carattere di \textit{newline}.

Un secondo punto da tenere presente, particolarmente quando si ha a che fare
con I/O interattivo, è che quando si effettua una lettura da uno
\textit{stream} che comporta l'accesso alle \textit{system call} del kernel,
ad esempio se lo \textit{stream} da cui si legge è in modalità
\textit{unbuffered}, viene anche eseguito lo scarico di tutti i buffer degli
\textit{stream} in scrittura. In sez.~\ref{sec:file_buffering_ctrl} vedremo
come la libreria definisca delle opportune funzioni per controllare le
modalità di bufferizzazione e lo scarico dei dati.



\subsection{Apertura e chiusura di uno \textit{stream}}
\label{sec:file_fopen}

Le funzioni che si possono usare per aprire uno \textit{stream} sono solo tre:
\funcd{fopen}, \funcd{fdopen} e \funcd{freopen},\footnote{\func{fopen} e
  \func{freopen} fanno parte dello standard ANSI C, \func{fdopen} è parte
  dello standard POSIX.1.} ed i rispettivi prototipi sono:

\begin{funcproto}{
\fhead{stdio.h}
\fdecl{FILE *fopen(const char *path, const char *mode)}
\fdesc{Apre uno \textit{stream} da un \texttt{pathname}.} 
\fdecl{FILE *fdopen(int fildes, const char *mode)}
\fdesc{Associa uno \textit{stream} a un file descriptor.} 
\fdecl{FILE *freopen(const char *path, const char *mode, FILE *stream)}
\fdesc{Chiude uno \textit{stream} e lo riapre su un file diverso.} 
}

{Le funzioni ritornano un puntatore ad un oggetto \type{FILE} in caso di
  successo e \val{NULL} per un errore, nel qual caso \var{errno} assumerà il
  valore ricevuto dalla funzione sottostante di cui è fallita l'esecuzione,
  gli errori pertanto possono essere quelli di \func{malloc} per tutte e tre
  le funzioni, quelli \func{open} per \func{fopen}, quelli di \func{fcntl} per
  \func{fdopen} e quelli di \func{fopen}, \func{fclose} e \func{fflush} per
  \func{freopen}.}
\end{funcproto}

Normalmente la funzione che si usa per aprire uno \textit{stream} è
\func{fopen}, essa apre il file specificato dal \textit{pathname} \param{path}
nella modalità specificata da \param{mode}, che è una stringa che deve
iniziare con almeno uno dei valori indicati in tab.~\ref{tab:file_fopen_mode},
anche se sono possibili varie estensioni che vedremo in seguito.

L'uso più comune di \func{freopen} è per redirigere uno dei tre file standard
(vedi sez.~\ref{sec:file_stream}): il file \param{path} viene aperto nella
modalità indicata da \param{mode} ed associato allo \textit{stream} indicato
dall'argomento \param{stream}, e se questo era uno \textit{stream} già aperto
esso viene preventivamente chiuso e tutti i dati pendenti vengono scaricati.

Infine \func{fdopen} viene usata per associare uno \textit{stream} ad un file
descriptor esistente ottenuto tramite una altra funzione (ad esempio con una
\func{open}, una \func{dup}, o una \func{pipe}) e serve quando si vogliono
usare gli \textit{stream} con file come le \textit{fifo} o i socket, che non possono
essere aperti con le funzioni delle librerie standard del C.

\begin{table}[htb]
  \centering
  \footnotesize
  \begin{tabular}[c]{|l|p{8cm}|}
    \hline
    \textbf{Valore} & \textbf{Significato}\\
    \hline
    \hline
    \texttt{r} & Il file viene aperto, l'accesso viene posto in sola
                 lettura, lo \textit{stream} è posizionato all'inizio del
                 file.\\ 
    \texttt{r+}& Il file viene aperto, l'accesso viene posto in lettura e
                 scrittura, lo \textit{stream} è posizionato all'inizio del
                 file.\\ 
%    \hline
    \texttt{w} & Il file viene aperto e troncato a lunghezza nulla (o
                 creato se non esiste), l'accesso viene posto in sola
                 scrittura, lo \textit{stream} è posizionato all'inizio del
                 file.\\ 
    \texttt{w+}& Il file viene aperto e troncato a lunghezza nulla (o
                 creato se non esiste), l'accesso viene posto in scrittura e
                 lettura, lo \textit{stream} è posizionato all'inizio del
                 file.\\ 
%    \hline
    \texttt{a} & Il file viene aperto (o creato se non esiste) in
                 \textit{append mode}, l'accesso viene posto in sola
                 scrittura.\\
    \texttt{a+}& Il file viene aperto (o creato se non esiste) in
                 \textit{append mode}, l'accesso viene posto in lettura e
                 scrittura.\\
    \hline
    \texttt{b} & Specifica che il file è binario, non ha alcun effetto. \\
    \texttt{x} & L'apertura fallisce se il file esiste già. \\
    \hline
  \end{tabular}
  \caption{Modalità di apertura di uno \textit{stream} dello standard ANSI C
    che sono sempre presenti in qualunque sistema POSIX.}
  \label{tab:file_fopen_mode}
\end{table}

In realtà lo standard ANSI C prevede un totale di 15 possibili valori
diversi per \param{mode}, ma in tab.~\ref{tab:file_fopen_mode} si sono
riportati solo i sei valori effettivi, ad essi può essere aggiunto pure
il carattere \texttt{b} (come ultimo carattere o nel mezzo agli altri per
le stringhe di due caratteri) che in altri sistemi operativi serve a
distinguere i file binari dai file di testo; in un sistema POSIX questa
distinzione non esiste e il valore viene accettato solo per
compatibilità, ma non ha alcun effetto.

La \acr{glibc} supporta alcune estensioni, queste devono essere sempre
indicate dopo aver specificato il \param{mode} con uno dei valori di
tab.~\ref{tab:file_fopen_mode}. L'uso del carattere \texttt{x} serve per
evitare di sovrascrivere un file già esistente (è analoga all'uso dell'opzione
\const{O\_EXCL} in \func{open}): se il file specificato già esiste e si
aggiunge questo carattere a \param{mode} la \func{fopen} fallisce.

Un'altra estensione serve a supportare la localizzazione, quando si
aggiunge a \param{mode} una stringa della forma \verb|",ccs=STRING"| il
valore \verb|STRING| è considerato il nome di una codifica dei caratteri
e \func{fopen} marca il file per l'uso dei caratteri estesi e abilita le
opportune funzioni di conversione in lettura e scrittura.

Nel caso si usi \func{fdopen} i valori specificati da \param{mode} devono
essere compatibili con quelli con cui il file descriptor è stato aperto.
Inoltre i modi \cmd{w} e \cmd{w+} non troncano il file. La posizione nello
\textit{stream} viene impostata a quella corrente nel file descriptor, e le
variabili di errore e di fine del file (vedi sez.~\ref{sec:file_io}) sono
cancellate. Il file non viene duplicato e verrà chiuso automaticamente alla
chiusura dello \textit{stream}.

I nuovi file saranno creati secondo quanto visto in
sez.~\ref{sec:file_ownership_management} ed avranno i permessi di accesso
impostati al valore
\code{S\_IRUSR|S\_IWUSR|S\_IRGRP|S\_IWGRP|S\_IROTH|S\_IWOTH} (pari a
\val{0666}) modificato secondo il valore della \textit{umask} per il processo
(si veda sez.~\ref{sec:file_perm_management}). Una volta aperto lo
\textit{stream}, si può cambiare la modalità di bufferizzazione (si veda
sez.~\ref{sec:file_buffering_ctrl}) fintanto che non si è effettuato alcuna
operazione di I/O sul file.

In caso di file aperti in lettura e scrittura occorre ricordarsi che c'è
di mezzo una bufferizzazione; per questo motivo lo standard ANSI C
richiede che ci sia un'operazione di posizionamento fra un'operazione
di output ed una di input o viceversa (eccetto il caso in cui l'input ha
incontrato la fine del file), altrimenti una lettura può ritornare anche
il risultato di scritture precedenti l'ultima effettuata. 

Per questo motivo è una buona pratica (e talvolta necessario) far seguire ad
una scrittura una delle funzioni \func{fflush}, \func{fseek}, \func{fsetpos} o
\func{rewind} prima di eseguire una rilettura; viceversa nel caso in cui si
voglia fare una scrittura subito dopo aver eseguito una lettura occorre prima
usare una delle funzioni \func{fseek}, \func{fsetpos} o \func{rewind}. Anche
un'operazione nominalmente nulla come \code{fseek(file, 0, SEEK\_CUR)} è
sufficiente a garantire la sincronizzazione.

Una volta completate le operazioni su di esso uno \textit{stream} può essere
chiuso con la funzione \funcd{fclose}, il cui prototipo è:

\begin{funcproto}{
\fhead{stdio.h}
\fdecl{int fclose(FILE *stream)}
\fdesc{Chiude uno \textit{stream}.} 
}

{La funzione ritorna $0$ in caso di successo e \val{EOF} per un errore, nel
  qual caso \var{errno} assumerà il valore \errval{EBADF} se il file
  descriptor indicato da \param{stream} non è valido, o uno dei valori
  specificati dalla sottostante funzione che è fallita (\func{close},
  \func{write} o \func{fflush}).
}
\end{funcproto}

La funzione chiude lo \textit{stream} \param{stream} ed effettua lo scarico di
tutti i dati presenti nei buffer di uscita e scarta tutti i dati in ingresso;
se era stato allocato un buffer per lo \textit{stream} questo verrà
rilasciato. La funzione effettua lo scarico solo per i dati presenti nei
buffer in \textit{user space} usati dalla \acr{glibc}; se si vuole essere
sicuri che il kernel forzi la scrittura su disco occorrerà effettuare una
\func{sync} (vedi sez.~\ref{sec:file_sync}).

Linux supporta anche un'altra funzione, \funcd{fcloseall}, come estensione
GNU implementata dalla \acr{glibc}, accessibile avendo definito
\macro{\_GNU\_SOURCE}, il suo prototipo è:

\begin{funcproto}{
\fhead{stdio.h}
\fdecl{int fcloseall(void)}
\fdesc{Chiude tutti gli \textit{stream}.} 
}

{La funzione ritorna $0$ in caso di successo e \val{EOF} per un errore, nel
  qual caso \var{errno} assumerà gli stessi valori di \func{fclose}.}  
\end{funcproto}

La funzione esegue lo scarico dei dati bufferizzati in uscita e scarta quelli
in ingresso, chiudendo tutti i file. Questa funzione è provvista solo per i
casi di emergenza, quando si è verificato un errore ed il programma deve
essere abortito, ma si vuole compiere qualche altra operazione dopo aver
chiuso i file e prima di uscire (si ricordi quanto visto in
sez.~\ref{sec:proc_conclusion}).


\subsection{Gestione dell'I/O e posizionamento su uno \textit{stream}}
\label{sec:file_io}

Una delle caratteristiche più utili dell'interfaccia degli \textit{stream} è
la ricchezza delle funzioni disponibili per le operazioni di lettura e
scrittura sui file. Sono infatti previste ben tre diverse modalità di
input/output non formattato:
\begin{itemize}
\item\textsl{binario} in cui si leggono e scrivono blocchi di dati di
   dimensione arbitraria, (analogo della modalità ordinaria dell'I/O sui file
   descriptor), trattato in sez.~\ref{sec:file_binary_io}.
\item\textsl{a caratteri} in cui si legge e scrive un carattere alla volta,
   con la bufferizzazione che viene gestita automaticamente dalla libreria,
   trattato in sez.~\ref{sec:file_char_io}.
\item\textsl{di linea} in cui si legge e scrive una linea alla volta,
   (terminata dal carattere di newline \verb|'\n'|), trattato in
   sez.~\ref{sec:file_line_io}.
\end{itemize}
a cui si aggiunge la modalità di input/output formattato, trattato in
sez.~\ref{sec:file_formatted_io}.

Ognuna di queste modalità utilizza per l'I/O delle funzioni specifiche che
vedremo nelle sezioni citate, affronteremo qui tutte gli argomenti e le
funzioni che si applicano in generale a tutte le modalità di I/O.

A differenza di quanto avviene con l'interfaccia dei file descriptor, con gli
\textit{stream} il raggiungimento della fine del file viene considerato un
errore, e viene notificato come tale dai valori di uscita delle varie
funzioni. Nella maggior parte dei casi questo avviene con la restituzione del
valore intero (di tipo \ctyp{int}) \val{EOF} definito anch'esso nell'header
\headfile{stdlib.h}. La costante deve essere negativa perché in molte funzioni
un valore positivo indica la quantità di dati scritti, la \acr{glibc} usa il
valore $-1$, ma altre implementazioni possono avere valori diversi.

Dato che le funzioni dell'interfaccia degli \textit{stream} sono funzioni di
libreria che si appoggiano a delle \textit{system call}, esse non impostano
direttamente la variabile \var{errno}, che mantiene sempre il valore impostato
dalla \textit{system call} invocata internamente che ha riportato l'errore.

Siccome la condizione di \textit{end-of-file} è anch'essa segnalata come
errore, nasce il problema di come distinguerla da un errore effettivo; basarsi
solo sul valore di ritorno della funzione e controllare il valore di
\var{errno} infatti non basta, dato che quest'ultimo potrebbe essere stato
impostato in una altra occasione, (si veda sez.~\ref{sec:sys_errno} per i
dettagli del funzionamento di \var{errno}).

Per questo motivo tutte le implementazioni delle librerie standard mantengono
per ogni \textit{stream} almeno due flag all'interno dell'oggetto \type{FILE},
il flag di \textit{end-of-file}, che segnala che si è raggiunta la fine del
file in lettura, e quello di errore, che segnala la presenza di un qualche
errore nelle operazioni di input/output; questi due flag possono essere
riletti dalle funzioni \funcd{feof} e \funcd{ferror}, i cui prototipi sono:

\begin{funcproto}{
\fhead{stdio.h}
\fdecl{int feof(FILE *stream)}
\fdesc{Controlla il flag di \textit{end-of-file} di uno \textit{stream}.} 
\fdecl{int ferror(FILE *stream)}
\fdesc{Controlla il flag di errore di uno \textit{stream}.} 
}

{Le funzioni ritornano un valore diverso da zero se i relativi flag sono
  impostati, e non prevedono condizioni di errore.}
\end{funcproto}

Si tenga presente comunque che la lettura di questi flag segnala soltanto che
c'è stato un errore o che si è raggiunta la fine del file in una qualunque
operazione sullo \textit{stream}, il controllo su quanto avvenuto deve quindi
essere effettuato ogni volta che si chiama una funzione di libreria.

Entrambi i flag (di errore e di \textit{end-of-file}) possono essere
cancellati usando la funzione \funcd{clearerr}, il cui prototipo è:

\begin{funcproto}{
\fhead{stdio.h}
\fdecl{void clearerr(FILE *stream)}
\fdesc{Cancella i flag di errore ed \textit{end-of-file} di uno
  \textit{stream}.}
}

{La funzione non ritorna nulla e prevede condizioni di errore.}  
\end{funcproto}

In genere si usa questa funzione una volta che si sia identificata e corretta
la causa di un errore per evitare di mantenere i flag attivi, così da poter
rilevare una successiva ulteriore condizione di errore. Di questa funzione
esiste una analoga \funcm{clearerr\_unlocked} (con lo stesso argomento e
stessi valori di ritorno) che non esegue il blocco dello \textit{stream}
(tratteremo il significato di blocco di uno \textit{stream} in
sez.~\ref{sec:file_stream_thread}).

Come per i file descriptor anche per gli \textit{stream} è possibile spostarsi
all'interno di un file per effettuare operazioni di lettura o scrittura in un
punto prestabilito, sempre che l'operazione di riposizionamento sia supportata
dal file sottostante lo \textit{stream}, nel caso cioè in cui si ha a che fare
con quello che viene detto un file ad \textsl{accesso casuale}. Dato che in un
sistema Unix esistono vari tipi di file, come le \textit{fifo} ed i file di
dispositivo (ad esempio i terminali), non è scontato che questo sia vero in
generale, pur essendolo sempre nel caso di file di dati.

Con Linux ed in generale in ogni sistema unix-like la posizione nel file, come
abbiamo già visto in sez.~\ref{sec:file_lseek}, è espressa da un intero
positivo, rappresentato dal tipo \type{off\_t}. Il problema è che alcune delle
funzioni usate per il riposizionamento sugli \textit{stream} originano dalle
prime versioni di Unix, in cui questo tipo non era ancora stato definito, e
che in altri sistemi non è detto che la posizione su un file venga sempre
rappresentata con il numero di caratteri dall'inizio: ad esempio nel VMS dove
esistono i file a record (in cui cioè l'I/O avviene per blocchi, i record, di
dimensione fissa), essa può essere rappresentata come un numero di record, più
l'offset rispetto al record corrente.

Tutto questo comporta la presenza di diverse funzioni che eseguono
sostanzialmente le stesse operazioni, ma usano argomenti di tipo diverso. Le
funzioni tradizionali usate per eseguire una modifica della posizione corrente
sul file con uno \textit{stream} sono \funcd{fseek} e \funcd{rewind}, i
rispettivi prototipi sono:

\begin{funcproto}{
\fhead{stdio.h}
\fdecl{int fseek(FILE *stream, long offset, int whence)}
\fdesc{Sposta la posizione nello \textit{stream}.} 
\fdecl{void rewind(FILE *stream)}
\fdesc{Riporta la posizione nello \textit{stream} all'inizio del file.} 
}

{La funzione \func{fseek} ritorna $0$ in caso di successo e $-1$ per un
  errore, nel qual caso \var{errno} assumerà i valori di \func{lseek},
  \func{rewind} non ritorna nulla e non ha condizioni di errore.}
\end{funcproto}

L'uso di \func{fseek} è del tutto analogo a quello di \func{lseek} per i file
descriptor (vedi sez.~\ref{sec:file_lseek}). Anche gli argomenti, a parte il
tipo, hanno esattamente lo stesso significato. In particolare \param{whence}
deve assumere gli stessi valori già visti nella prima parte di
tab.~\ref{tab:lseek_whence_values}.  La funzione restituisce 0 in caso di
successo e -1 in caso di errore.

La funzione \func{rewind} riporta semplicemente la posizione corrente sul file
all'inizio dello \textit{stream}, ma non è esattamente equivalente ad aver
eseguito una \code{fseek(stream, 0L, SEEK\_SET)}, in quanto con l'uso della
funzione vengono cancellati anche i flag di errore e di fine del file.

Per ottenere la posizione corrente sul file di uno \textit{stream} lo standard
ANSI C prescrive l'uso della funzione \funcd{ftell}, il cui prototipo è:

\begin{funcproto}{
\fhead{stdio.h}
\fdecl{long ftell(FILE *stream)} 
\fdesc{Legge la posizione attuale nello \textit{stream}.} 
}

{La funzione ritorna la posizione corrente in caso di successo e $-1$ per un
  errore, nel qual caso \var{errno} assumerà  i valori di \func{lseek}.}  
\end{funcproto}

\noindent che restituisce la posizione come numero di byte dall'inizio dello
\textit{stream}.

Sia \func{fseek} che \func{ftell} esprimono la posizione nel file con un
intero di tipo \ctyp{long}. Dato che in certi casi, ad esempio quando si usa
un filesystem indicizzato a 64 bit su una macchina con architettura a 32 bit,
questo può non essere possibile lo standard POSIX ha introdotto le nuove
funzioni \funcd{fgetpos} e \funcd{fsetpos}, che invece usano il nuovo tipo
\typed{fpos\_t}, ed i cui prototipi sono:

\begin{funcproto}{
\fhead{stdio.h}
\fdecl{int fsetpos(FILE *stream, fpos\_t *pos)}
\fdesc{Imposta la posizione corrente sul file.} 
\fdecl{int fgetpos(FILE *stream, fpos\_t *pos)}
\fdesc{Legge la posizione corrente sul file.} 
}

{La funzione ritorna $0$ in caso di successo e $-1$ per un errore, nel qual
  caso \var{errno} assumerà i valori di \func{lseek}.}
\end{funcproto}

In Linux, a partire dalle glibc 2.1, sono presenti anche le due funzioni
\func{fseeko} e \func{ftello}, che sono assolutamente identiche alle
precedenti \func{fseek} e \func{ftell} ma hanno argomenti di tipo
\type{off\_t} anziché di tipo \ctyp{long int}. Dato che \ctyp{long} è nella
gran parte dei casi un intero a 32 bit, questo diventa un problema quando la
posizione sul file viene espressa con un valore a 64 bit come accade nei
sistemi più moderni.

% TODO: mettere prototipi espliciti fseeko e ftello o menzione?


\subsection{Input/output binario}
\label{sec:file_binary_io}

La prima modalità di input/output non formattato ricalca quella della
interfaccia dei file descriptor, e provvede semplicemente la scrittura e la
lettura dei dati da un buffer verso un file e viceversa. In generale questa è
la modalità che si usa quando si ha a che fare con dati non formattati. Le due
funzioni che si usano per l'I/O binario sono \funcd{fread} ed \funcd{fwrite};
i rispettivi prototipi sono:

\begin{funcproto}{
\fhead{stdio.h} 
\fdecl{size\_t fread(void *ptr, size\_t size, size\_t nmemb, FILE *stream)}
\fdesc{Legge i dati da uno \textit{stream}.} 
\fdecl{size\_t fwrite(const void *ptr, size\_t size, size\_t nmemb, 
  FILE *stream)}
\fdesc{Scrive i dati su uno \textit{stream}.} 
}

{Le funzioni ritornano il numero di elementi letti o scritti, in caso di
  errore o fine del file viene restituito un numero di elementi inferiore al
  richiesto.}
\end{funcproto}

Le funzioni rispettivamente leggono e scrivono \param{nmemb} elementi di
dimensione \param{size} dal buffer \param{ptr} al file \param{stream}.  In
genere si usano queste funzioni quando si devono trasferire su file blocchi di
dati binari in maniera compatta e veloce; un primo caso di uso tipico è quello
in cui si salva un vettore (o un certo numero dei suoi elementi) con una
chiamata del tipo:
\includecodesnip{listati/WriteVect.c}
in questo caso devono essere specificate le dimensioni di ciascun
elemento ed il numero di quelli che si vogliono scrivere. Un secondo
caso è invece quello in cui si vuole trasferire su file una struttura;
si avrà allora una chiamata tipo:
\includecodesnip{listati/WriteStruct.c}
in cui si specifica la dimensione dell'intera struttura ed un solo
elemento. 

In realtà quello che conta nel trasferimento dei dati sono le dimensioni
totali, che sono sempre pari al prodotto \code{size * nelem}, la differenza
sta nel fatto che le funzioni non ritornano il numero di byte scritti, ma il
numero di elementi (e con questo possono facilitare i conti).

La funzione \func{fread} legge sempre un numero intero di elementi, se
incontra la fine del file l'oggetto letto parzialmente viene scartato (lo
stesso avviene in caso di errore). In questo caso la posizione dello
\textit{stream} viene impostata alla fine del file (e non a quella
corrispondente alla quantità di dati letti).

In caso di errore (o fine del file per \func{fread}) entrambe le
funzioni restituiscono il numero di oggetti effettivamente letti o
scritti, che sarà inferiore a quello richiesto. Contrariamente a quanto
avviene per i file descriptor, questo segnala una condizione di errore e
occorrerà usare \func{feof} e \func{ferror} per stabilire la natura del
problema.

Benché queste funzioni assicurino la massima efficienza per il
salvataggio dei dati, i dati memorizzati attraverso di esse presentano
lo svantaggio di dipendere strettamente dalla piattaforma di sviluppo
usata ed in genere possono essere riletti senza problemi solo dallo
stesso programma che li ha prodotti.

Infatti diversi compilatori possono eseguire ottimizzazioni diverse delle
strutture dati e alcuni compilatori (come il \cmd{gcc}) possono anche
scegliere se ottimizzare l'occupazione di spazio, impacchettando più
strettamente i dati, o la velocità inserendo opportuni \textit{padding} per
l'allineamento dei medesimi generando quindi output binari diversi. Inoltre
altre incompatibilità si possono presentare quando entrano in gioco differenze
di architettura hardware, come la dimensione del bus o la modalità di
ordinamento dei bit o il formato delle variabili in floating point.

Per questo motivo quando si usa l'input/output binario occorre sempre prendere
le opportune precauzioni come usare un formato di più alto livello che
permetta di recuperare l'informazione completa, per assicurarsi che versioni
diverse del programma siano in grado di rileggere i dati, tenendo conto delle
eventuali differenze.

La \acr{glibc} definisce infine due ulteriori funzioni per l'I/O binario,
\funcd{fread\_unlocked} e \funcd{fwrite\_unlocked}, che evitano il lock
implicito dello \textit{stream} usato per dalla librerie per la gestione delle
applicazioni \textit{multi-thread} (si veda sez.~\ref{sec:file_stream_thread}
per i dettagli), i loro prototipi sono:

\begin{funcproto}{
\fhead{stdio.h}
\fdecl{size\_t fread\_unlocked(void *ptr, size\_t size, size\_t
    nmemb, FILE *stream)}
\fdecl{size\_t fwrite\_unlocked(const void *ptr, size\_t size,
    size\_t nmemb, FILE *stream)}
\fdesc{Leggono o scrivono dati su uno \textit{stream} senza acquisire il lock
  implicito sullo stesso.} 
}

{Le funzioni ritornano gli stessi valori delle precedenti \func{fread} e
  \func{fwrite}.}
\end{funcproto}

% TODO: trattare in generale le varie *_unlocked


\subsection{Input/output a caratteri}
\label{sec:file_char_io}

La seconda modalità di input/output è quella a caratteri, in cui si
trasferisce un carattere alla volta.  Le funzioni per la lettura a
caratteri sono tre, \funcd{fgetc}, \funcd{getc} e \funcd{getchar}, ed i
rispettivi prototipi sono:

\begin{funcproto}{
\fhead{stdio.h}
\fdecl{int getc(FILE *stream)}
\fdecl{int fgetc(FILE *stream)}
\fdesc{Leggono un singolo byte da uno \textit{stream}.} 
\fdecl{int getchar(void)}
\fdesc{Legge un byte dallo \textit{standard input}.} 
}

{Le funzioni ritornano il byte letto in caso di successo e \val{EOF} per un
  errore o se si arriva alla fine del file.}  
\end{funcproto}

La funzione \func{getc} legge un byte da \param{stream} e lo restituisce come
intero, ed in genere è implementata come una macro per cui può avere
\textit{side effects}, mentre \func{fgetc} è assicurato essere sempre una
funzione. Infine \func{getchar} è equivalente a \code{getc(stdin)}.

A parte \func{getchar}, che si usa in genere per leggere un carattere da
tastiera, le altre due funzioni sono sostanzialmente equivalenti. La
differenza è che \func{getc} è ottimizzata al massimo e normalmente
viene implementata con una macro, per cui occorre stare attenti a cosa
le si passa come argomento, infatti \param{stream} può essere valutato
più volte nell'esecuzione, e non viene passato in copia con il
meccanismo visto in sez.~\ref{sec:proc_var_passing}; per questo motivo se
si passa un'espressione si possono avere effetti indesiderati.

Invece \func{fgetc} è assicurata essere sempre una funzione, per questo motivo
la sua esecuzione normalmente è più lenta per via dell'overhead della
chiamata, ma è altresì possibile ricavarne l'indirizzo, che può essere passato
come argomento ad un altra funzione (e non si hanno i problemi accennati in
precedenza nel tipo di argomento).

Le tre funzioni restituiscono tutte un \ctyp{unsigned char} convertito
ad \ctyp{int} (si usa \ctyp{unsigned char} in modo da evitare
l'espansione del segno). In questo modo il valore di ritorno è sempre
positivo, tranne in caso di errore o fine del file.

Nelle estensioni GNU che provvedono la localizzazione sono definite tre
funzioni equivalenti alle precedenti, \funcd{getwc}, \funcd{fgetwc} e
\funcd{getwchar}, che invece di un carattere di un byte restituiscono un
carattere in formato esteso (cioè di tipo \ctyp{wint\_t}), il loro prototipo
è:

\begin{funcproto}{
\fhead{stdio.h} 
\fhead{wchar.h}
\fdecl{wint\_t getwc(FILE *stream)}
\fdecl{wint\_t fgetwc(FILE *stream)}
\fdesc{Leggono un carattere da uno \textit{stream}.} 
\fdecl{wint\_t getwchar(void)}
\fdesc{Legge un carattere dallo \textit{standard input}.} 
}

{Le funzioni ritornano il carattere letto in caso di successo e \val{WEOF} per
  un errore o se si arriva alla fine del file.}  
\end{funcproto}

La funzione \func{getwc} legge un carattere esteso da \param{stream} e lo
restituisce come intero, ed in genere è implementata come una macro, mentre
\func{fgetwc} è assicurata essere sempre una funzione. Infine \func{getwchar}
è equivalente a \code{getwc(stdin)}.

Per scrivere un carattere si possono usare tre funzioni, analoghe alle
precedenti usate per leggere: \funcd{putc}, \funcd{fputc} e \funcd{putchar}; i
loro prototipi sono:

\begin{funcproto}{
\fhead{stdio.h} 
\fdecl{int putc(int c, FILE *stream)}
\fdecl{int fputc(int c, FILE *stream)}
\fdesc{Scrive un byte su uno \textit{stream}.}
\fdecl{int putchar(int c)}
\fdesc{Scrive un byte sullo \textit{standard output}.}
}

{Le funzioni ritornano il valore del byte scritto in caso di successo e
  \val{EOF} per un errore.}  
\end{funcproto}

La funzione \func{putc} scrive un byte su \param{stream} e lo restituisce come
intero, ed in genere è implementata come una macro, mentre \func{fputc} è
assicurata essere sempre una funzione. Infine \func{putchar} è equivalente a
\code{putc(stdout)}.  Tutte queste funzioni scrivono sempre un byte alla
volta, anche se prendono come argomento un \ctyp{int} (che pertanto deve
essere ottenuto con un cast da un \ctyp{unsigned char}). Anche il valore di
ritorno è sempre un intero; in caso di errore o fine del file il valore di
ritorno è \val{EOF}.

Come nel caso dell'I/O binario con \func{fread} e \func{fwrite} la \acr{glibc}
provvede come estensione, per ciascuna delle funzioni precedenti,
un'ulteriore funzione, il cui nome è ottenuto aggiungendo un
\code{\_unlocked}, che esegue esattamente le stesse operazioni, evitando però
il lock implicito dello \textit{stream}.

Per compatibilità con SVID sono inoltre provviste anche due funzioni,
\funcd{getw} e \funcd{putw}, da usare per leggere e scrivere una \textit{word}
(cioè due byte in una volta); i loro prototipi sono:

\begin{funcproto}{
\fhead{stdio.h} 
\fdecl{getw(FILE *stream)}
\fdesc{Legge una parola da uno \textit{stream}.} 
\fdecl{int putw(int w, FILE *stream)}
\fdesc{Scrive una parola su uno \textit{stream}.} 
}

{Le funzioni ritornano la parola letta o scritta in caso di successo e
  \val{EOF} per un errore.}
\end{funcproto}

Le funzioni leggono e scrivono una \textit{word} di due byte, usando comunque
una variabile di tipo \ctyp{int}; il loro uso è deprecato in favore dell'uso
di \func{fread} e \func{fwrite}, in quanto non è possibile distinguere il
valore -1 da una condizione di errore che restituisce \val{EOF}.

Uno degli usi più frequenti dell'input/output a caratteri è nei programmi di
\textit{parsing} in cui si analizza il testo; in questo contesto diventa utile
poter analizzare il carattere successivo da uno \textit{stream} senza estrarlo
effettivamente (la tecnica è detta \textit{peeking ahead}) in modo che il
programma possa regolarsi avendo dato una \textsl{sbirciatina} a quello che
viene dopo.

Nel nostro caso questo tipo di comportamento può essere realizzato prima
leggendo il carattere, e poi rimandandolo indietro, cosicché ridiventi
disponibile per una lettura successiva; la funzione che inverte la
lettura si chiama \funcd{ungetc} ed il suo prototipo è:

\begin{funcproto}{
\fhead{stdio.h}
\fdecl{int ungetc(int c, FILE *stream)}
\fdesc{Manda indietro un byte su uno \textit{stream}.} 
}

{La funzione ritorna il byte inviato in caso di successo e \val{EOF} per un
  errore.}  
\end{funcproto}
 
La funzione rimanda indietro il carattere \param{c}, con un cast a
\ctyp{unsigned char}, sullo \textit{stream} \param{stream}. Benché lo standard
ANSI C preveda che l'operazione possa essere ripetuta per un numero arbitrario
di caratteri, alle implementazioni è richiesto di garantire solo un livello;
questo è quello che fa la \acr{glibc}, che richiede che avvenga un'altra
operazione fra due \func{ungetc} successive.

Non è necessario che il carattere che si manda indietro sia l'ultimo che si è
letto, e non è necessario neanche avere letto nessun carattere prima di usare
\func{ungetc}, ma di norma la funzione è intesa per essere usata per rimandare
indietro l'ultimo carattere letto.  Nel caso \param{c} sia un \val{EOF} la
funzione non fa nulla, e restituisce sempre \val{EOF}; così si può usare
\func{ungetc} anche con il risultato di una lettura alla fine del file.

Se si è alla fine del file si può comunque rimandare indietro un carattere, il
flag di \textit{end-of-file} verrà automaticamente cancellato perché c'è un
nuovo carattere disponibile che potrà essere riletto successivamente.

Infine si tenga presente che \func{ungetc} non altera il contenuto del file,
ma opera esclusivamente sul buffer interno. Se si esegue una qualunque delle
operazioni di riposizionamento (vedi sez.~\ref{sec:file_io}) i caratteri
rimandati indietro vengono scartati.


\subsection{Input/output di linea}
\label{sec:file_line_io}

La terza ed ultima modalità di input/output non formattato è quella di linea,
in cui si legge o si scrive una riga alla volta. Questa è la modalità usata
normalmente per l'I/O da terminale, ed è anche quella che presenta le
caratteristiche più controverse.

Le funzioni previste dallo standard ANSI C per leggere una linea sono
sostanzialmente due, \funcd{gets} e \funcd{fgets}, i cui rispettivi
prototipi sono:

\begin{funcproto}{
\fhead{stdio.h}
\fdecl{char *gets(char *string)}
\fdesc{Legge una linea di testo dallo \textit{standard input}.}
\fdecl{char *fgets(char *string, int size, FILE *stream)}
\fdesc{Legge una linea di testo da uno \textit{stream}.} 
}

{Le funzioni ritornano l'indirizzo della stringa con la linea di testo letta o
  scritta in caso di successo e \val{NULL} per un errore.}
\end{funcproto}
 
Entrambe le funzioni effettuano la lettura, dal file specificato \func{fgets},
dallo \textit{standard input} \func{gets}, di una linea di caratteri terminata
dal carattere ASCII di \textit{newline}, che come detto corrisponde a quello
generato dalla pressione del tasto di invio sulla tastiera. Si tratta del
carattere che indica la terminazione di una riga (in sostanza del carattere di
``\textsl{a capo}'') che viene rappresentato nelle stringhe di formattazione
che vedremo in sez.~\ref{sec:file_formatted_io} come
``\verb|\n|''. Nell'esecuzione delle funzioni \func{gets} sostituisce
``\verb|\n|'' con uno zero, mentre \func{fgets} aggiunge uno zero dopo il
\textit{newline}, che resta dentro la stringa.

\itindbeg{buffer~overflow}

Se la lettura incontra la fine del file (o c'è un errore) viene restituito un
puntatore \val{NULL}, ed il buffer \param{buf} non viene toccato.  L'uso di
\func{gets} è deprecato e deve essere assolutamente evitato, la funzione
infatti non controlla il numero di byte letti, per cui nel caso la stringa
letta superi le dimensioni del buffer, si avrà un \textit{buffer overflow},
con sovrascrittura della memoria del processo adiacente al
buffer.\footnote{questa tecnica è spiegata in dettaglio e con molta efficacia
  nell'ormai famoso articolo di Aleph1 \cite{StS}.}

Questa è una delle vulnerabilità più sfruttate per guadagnare accessi non
autorizzati al sistema (i cosiddetti \textit{exploit}), basta infatti inviare
una stringa sufficientemente lunga ed opportunamente forgiata per
sovrascrivere gli indirizzi di ritorno nello \textit{stack} (supposto che la
\func{gets} sia stata chiamata da una subroutine), in modo da far ripartire
l'esecuzione nel codice inviato nella stringa stessa, che in genere contiene
uno \textit{shell code}, cioè una sezione di programma che lancia una shell da
cui si potranno poi eseguire altri programmi.

\itindend{buffer~overflow}

La funzione \func{fgets} non ha i precedenti problemi di \func{gets} in quanto
prende in ingresso la dimensione del buffer \param{size}, che non verrà mai
ecceduta in lettura. La funzione legge fino ad un massimo di \param{size}
caratteri (\textit{newline} compreso), ed aggiunge uno zero di terminazione;
questo comporta che la stringa possa essere al massimo di \code{size-1}
caratteri.  Se la linea eccede la dimensione del buffer verranno letti solo
\code{size-1} caratteri, ma la stringa sarà sempre terminata correttamente con
uno zero finale; sarà possibile leggere i rimanenti caratteri in una chiamata
successiva.

Per la scrittura di una linea lo standard ANSI C prevede altre due
funzioni, \funcd{fputs} e \funcd{puts}, analoghe a quelle di lettura, i
rispettivi prototipi sono:

\begin{funcproto}{
\fhead{stdio.h}
\fdecl{int puts(char *string)}
\fdesc{Scrive una linea di testo sullo \textit{standard output}.}
\fdecl{int fputs(char *string, int size, FILE *stream)}
\fdesc{Scrive una linea di testo su uno \textit{stream}.} 
}

{Le funzioni ritornano un valore non negativo in caso di successo e \val{EOF}
  per un errore.}
\end{funcproto}

La funzione \func{puts} scrive una linea di testo mantenuta
all'indirizzo \param{string} sullo \textit{standard output} mentre \func{puts}
la scrive sul file indicato da \param{stream}.  Dato che in questo caso si
scrivono i dati in uscita \func{puts} non ha i problemi di \func{gets} ed è in
genere la forma più immediata per scrivere messaggi sullo \textit{standard
  output}; la funzione prende una stringa terminata da uno zero ed aggiunge
automaticamente il ritorno a capo. La differenza con \func{fputs} (a parte la
possibilità di specificare un file diverso da \var{stdout}) è che quest'ultima
non aggiunge il \textit{newline}, che deve essere previsto esplicitamente.

Come per le analoghe funzioni di input/output a caratteri, anche per l'I/O di
linea esistono delle estensioni per leggere e scrivere linee di caratteri
estesi, le funzioni in questione sono \funcd{fgetws} e \funcd{fputws} ed i
loro prototipi sono:

\begin{funcproto}{
\fhead{wchar.h}
\fdecl{wchar\_t *fgetws(wchar\_t *ws, int n, FILE *stream)}
\fdesc{Legge una stringa di carattere estesi da uno \textit{stream}.} 
\fdecl{int fputws(const wchar\_t *ws, FILE *stream)}
\fdesc{Scrive una stringa di carattere estesi da uno \textit{stream}.} 
}

{Le funzioni ritornano rispettivamente l'indirizzo della stringa o un non
  negativo in caso di successo e \val{NULL} o \val{EOF} per un errore o per la
  fine del file.}
\end{funcproto}


La funzione \func{fgetws} legge un massimo di \param{n} caratteri estesi dal
file \param{stream} al buffer \param{ws}, mentre la funzione \func{fputws}
scrive la linea \param{ws} di caratteri estesi sul file indicato
da \param{stream}.  Il comportamento di queste due funzioni è identico a
quello di \func{fgets} e \func{fputs}, a parte il fatto che tutto (numero di
caratteri massimo, terminatore della stringa, \textit{newline}) è espresso in
termini di caratteri estesi anziché di normali caratteri ASCII.

Come per l'I/O binario e quello a caratteri, anche per l'I/O di linea la
\acr{glibc} supporta una serie di altre funzioni, estensioni di tutte quelle
illustrate finora (eccetto \func{gets} e \func{puts}), che eseguono
esattamente le stesse operazioni delle loro equivalenti, evitando però il lock
implicito dello \textit{stream} (vedi sez.~\ref{sec:file_stream_thread}). Come
per le altre forma di I/O, dette funzioni hanno lo stesso nome della loro
analoga normale, con l'aggiunta dell'estensione \code{\_unlocked}.

Come abbiamo visto, le funzioni di lettura per l'input/output di linea
previste dallo standard ANSI C presentano svariati inconvenienti. Benché
\func{fgets} non abbia i gravissimi problemi di \func{gets}, può comunque dare
risultati ambigui se l'input contiene degli zeri; questi infatti saranno
scritti sul buffer di uscita e la stringa in output apparirà come più corta
dei byte effettivamente letti. Questa è una condizione che è sempre possibile
controllare (deve essere presente un \textit{newline} prima della effettiva
conclusione della stringa presente nel buffer), ma a costo di una
complicazione ulteriore della logica del programma. Lo stesso dicasi quando si
deve gestire il caso di stringa che eccede le dimensioni del buffer.

Per questo motivo la \acr{glibc} prevede, come estensione GNU, due nuove
funzioni per la gestione dell'input/output di linea, il cui uso permette di
risolvere questi problemi. L'uso di queste funzioni deve essere attivato
definendo la macro \macro{\_GNU\_SOURCE} prima di includere
\headfile{stdio.h}. La prima delle due, \funcd{getline}, serve per leggere una
linea terminata da un \textit{newline}, esattamente allo stesso modo di
\func{fgets}, il suo prototipo è:

\begin{funcproto}{
\fhead{stdio.h}
\fdecl{ssize\_t getline(char **buffer, size\_t *n, FILE *stream)}
\fdesc{Legge una riga da uno \textit{stream}.} 
}

{La funzione ritorna il numero di caratteri letti in caso di successo e $-1$
  per un errore o per il raggiungimento della fine del file.}
\end{funcproto}

La funzione legge una linea dal file \param{stream} copiandola sul buffer
indicato da \param{buffer} riallocandolo se necessario (l'indirizzo del buffer
e la sua dimensione vengono sempre riscritte). Permette così di eseguire una
lettura senza doversi preoccupare della eventuale lunghezza eccessiva della
stringa da leggere. 

Essa prende come primo argomento l'indirizzo del puntatore al buffer su cui si
vuole copiare la linea. Quest'ultimo \emph{deve} essere stato allocato in
precedenza con una \func{malloc}, non si può cioè passare come argomento primo
argomento l'indirizzo di un puntatore ad una variabile locale. Come secondo
argomento la funzione vuole l'indirizzo della variabile contenente le
dimensioni del buffer suddetto.

Se il buffer di destinazione è sufficientemente ampio la stringa viene scritta
subito, altrimenti il buffer viene allargato usando \func{realloc} e la nuova
dimensione ed il nuovo puntatore vengono restituiti indietro, si noti infatti
come entrambi gli argomenti siano dei \textit{value result argument}, per i
quali vengono passati dei puntatori anziché i valori delle variabili, secondo
quanto abbiamo descritto in sez.~\ref{sec:proc_var_passing}).

Se si passa alla funzione l'indirizzo di un puntatore impostato a \val{NULL} e
\var{*n} è zero, la funzione provvede da sola all'allocazione della memoria
necessaria a contenere la linea. In tutti i casi si ottiene dalla funzione un
puntatore all'inizio del testo della linea letta. Un esempio di codice può
essere il seguente: 
\includecodesnip{listati/getline.c} 
e per evitare \textit{memory leak} occorre ricordarsi di liberare la memoria
allocata dalla funzione eseguendo una \func{free} su \var{ptr}.

Il valore di ritorno di \func{getline} indica il numero di caratteri letti
dallo \textit{stream}, quindi compreso il \textit{newline}, ma non lo zero di
terminazione. Questo permette anche di distinguere anche gli eventuali zeri
letti come dati dallo \textit{stream} da quello inserito dalla funzione dopo
il \textit{newline} per terminare la stringa.  Se si è alla fine del file e
non si è potuto leggere nulla o se c'è stato un errore la funzione restituisce
$-1$.

La seconda estensione GNU per la lettura con l'I/O di linea è una
generalizzazione di \func{getline} per poter usare come separatore delle linee
un carattere qualsiasi al posto del \textit{newline}. La funzione si chiama
\funcd{getdelim} ed il suo prototipo è:

\begin{funcproto}{
\fhead{stdio.h}
\fdecl{size\_t getdelim(char **buffer, size\_t *n, int delim, FILE *stream)} 
\fdesc{Legge da uno \textit{stream} una riga delimitata da un carattere
  scelto.} 
}

{La funzione ha gli stessi valori di ritorno e gli stessi errori di
  \func{getline}.}
\end{funcproto}

La funzione è identica a \func{getline} solo che usa \param{delim} al posto
del carattere di \textit{newline} come separatore di linea. Il comportamento
di \func{getdelim} è identico a quello di \func{getline}, che può essere
implementata da \func{getdelim} passando ``\verb|\n|'' come valore
dell'argomento \param{delim}.


\subsection{Input/output formattato}
\label{sec:file_formatted_io}

L'ultima modalità di input/output è quella formattata, che è una delle
caratteristiche più utilizzate delle librerie standard del C; in genere questa
è la modalità in cui si esegue normalmente l'output su terminale poiché
permette di stampare in maniera facile e veloce dati, tabelle e messaggi.

L'output formattato viene eseguito con una delle 13 funzioni della famiglia
\func{printf}; le tre più usate sono \funcd{printf}, \funcd{fprintf} e
\funcd{sprintf}, i cui prototipi sono:

\begin{funcproto}{
\fhead{stdio.h} 
\fdecl{int printf(const char *format, ...)}
\fdesc{Scrive una stringa formattata sullo \textit{standard output}.}
\fdecl{int fprintf(FILE *stream, const char *format, ...)}
\fdesc{Scrive una stringa formattata su uno \textit{stream}.} 
\fdecl{int sprintf(char *str, const char *format, ...)} 
\fdesc{Scrive una stringa formattata su un buffer.} 
}

{Le funzioni ritornano il numero di caratteri scritti in caso di successo e un
  valore negativo per un errore.}  
\end{funcproto}


Le funzioni usano la stringa \param{format} come indicatore del formato con
cui dovrà essere scritto il contenuto degli argomenti, il cui numero è
variabile e dipende dal formato stesso.

Le prime due servono per scrivere su file (lo \textit{standard output} o
quello specificato) la terza permette di scrivere su una stringa, in genere
l'uso di \func{sprintf} è sconsigliato in quanto è possibile, se non si ha la
sicurezza assoluta sulle dimensioni del risultato della stampa, eccedere le
dimensioni di \param{str}, con conseguente sovrascrittura di altre variabili e
possibili \textit{buffer overflow}. Per questo motivo si consiglia l'uso
dell'alternativa \funcd{snprintf}, il cui prototipo è:

\begin{funcproto}{
\fhead{stdio.h}
\fdecl{snprintf(char *str, size\_t size, const char *format, ...)} 
\fdesc{Scrive una stringa formattata su un buffer.} 
}

{La funzione ha lo stesso valore di ritorno e gli stessi errori di
  \func{sprintf}.}
\end{funcproto}

\noindent la funzione è identica a \func{sprintf}, ma non scrive
su \param{str} più di \param{size} caratteri, garantendo così che il buffer
non possa essere sovrascritto.

\begin{table}[!htb]
  \centering
  \footnotesize
  \begin{tabular}[c]{|l|l|p{10cm}|}
    \hline
    \textbf{Valore} & \textbf{Tipo} & \textbf{Significato} \\
    \hline
    \hline
   \cmd{\%d} &\ctyp{int}         & Stampa un numero intero in formato decimale
                                   con segno.\\
   \cmd{\%i} &\ctyp{int}         & Identico a \cmd{\%d} in output.\\
   \cmd{\%o} &\ctyp{unsigned int}& Stampa un numero intero come ottale.\\
   \cmd{\%u} &\ctyp{unsigned int}& Stampa un numero intero in formato
                                   decimale senza segno.\\
   \cmd{\%x}, 
   \cmd{\%X} &\ctyp{unsigned int}& Stampano un intero in formato esadecimale,
                                   rispettivamente con lettere minuscole e
                                   maiuscole.\\
   \cmd{\%f} &\ctyp{double}      & Stampa un numero in virgola mobile con la
                                   notazione a virgola fissa.\\
   \cmd{\%e}, 
   \cmd{\%E} &\ctyp{double} & Stampano un numero in virgola mobile con la
                              notazione esponenziale, rispettivamente con
                              lettere minuscole e maiuscole.\\
   \cmd{\%g}, 
   \cmd{\%G} &\ctyp{double} & Stampano un numero in virgola mobile con la
                              notazione più appropriate delle due precedenti,
                              rispettivamente con lettere minuscole e
                              maiuscole.\\
   \cmd{\%a}, 
   \cmd{\%A} &\ctyp{double} & Stampano un numero in virgola mobile in
                              notazione esadecimale frazionaria.\\
   \cmd{\%c} &\ctyp{int}    & Stampa un carattere singolo.\\
   \cmd{\%s} &\ctyp{char *} & Stampa una stringa.\\
   \cmd{\%p} &\ctyp{void *} & Stampa il valore di un puntatore.\\
   \cmd{\%n} &\ctyp{\&int}  & Prende il numero di caratteri stampati finora.\\
   \cmd{\%\%}&              & Stampa un ``\texttt{\%}''.\\
    \hline
  \end{tabular}
  \caption{Valori possibili per gli specificatori di conversione in una
    stringa di formato di \func{printf}.} 
  \label{tab:file_format_spec}
\end{table}

La parte più complessa delle funzioni di scrittura formattata è il formato
della stringa \param{format} che indica le conversioni da fare, e da cui
deriva anche il numero degli argomenti che dovranno essere passati a seguire:
si noti come tutte queste funzioni siano ``\textit{variadic}'', prendendo un
numero di argomenti variabile che dipende appunto da quello che si è
specificato in \param{format}.

La stringa di formato è costituita da caratteri normali (tutti eccetto
``\texttt{\%}''), che vengono passati invariati in uscita, e da direttive di
conversione, in cui devono essere sempre presenti il carattere
``\texttt{\%}'', che introduce la direttiva, ed uno degli specificatori di
conversione (riportati in tab.~\ref{tab:file_format_spec}) che la conclude.

Il formato di una direttiva di conversione prevede una serie di possibili
elementi opzionali oltre al carattere ``\cmd{\%}'' e allo specificatore di
conversione. In generale essa è sempre del tipo:
\begin{Example}
% [n. parametro $] [flag] [[larghezza] [. precisione]] [tipo] conversione
\end{Example}
in cui tutti i valori tranne il ``\texttt{\%}'' e lo specificatore di
conversione sono opzionali (e per questo sono indicati fra parentesi quadre);
si possono usare più elementi opzionali, nel qual caso devono essere
specificati in questo ordine:
\begin{itemize*}
\item uno specificatore del parametro da usare (terminato da un carattere
  ``\val{\$}''),
\item uno o più flag (i cui valori possibili sono riassunti in
  tab.~\ref{tab:file_format_flag}) che controllano il formato di stampa della
  conversione,
\item uno specificatore di larghezza (un numero decimale), eventualmente
  seguito (per i numeri in virgola mobile) da un specificatore di precisione
  (un altro numero decimale),
\item uno specificatore del tipo di dato, che ne indica la dimensione (i cui
  valori possibili sono riassunti in tab.~\ref{tab:file_format_type}).
\end{itemize*}

\begin{table}[htb]
  \centering
  \footnotesize
  \begin{tabular}[c]{|l|p{10cm}|}
    \hline
    \textbf{Valore} & \textbf{Significato}\\
    \hline
    \hline
    \val{\#} & Chiede la conversione in forma alternativa.\\
    \val{0}  & La conversione è riempita con zeri alla sinistra del valore.\\
    \val{-}  & La conversione viene allineata a sinistra sul bordo del campo.\\
    \val{' '}& Mette uno spazio prima di un numero con segno di valore 
               positivo.\\
    \val{+}  & Mette sempre il segno ($+$ o $-$) prima di un numero.\\
    \hline
  \end{tabular}
  \caption{I valori dei flag per il formato di \func{printf}}
  \label{tab:file_format_flag}
\end{table}

Dettagli ulteriori sulle varie opzioni di stampa e su tutte le casistiche
dettagliate dei vari formati possono essere trovati nella pagina di manuale di
\func{printf} e nella documentazione della \acr{glibc}.

\begin{table}[htb]
  \centering
  \footnotesize
  \begin{tabular}[c]{|l|p{10cm}|}
    \hline
    \textbf{Valore} & \textbf{Significato} \\
    \hline
    \hline
    \cmd{hh} & Una conversione intera corrisponde a un \ctyp{char} con o senza
               segno, o il puntatore per il numero dei parametri \cmd{n} è di 
               tipo \ctyp{char}.\\
    \cmd{h}  & Una conversione intera corrisponde a uno \ctyp{short} con o 
               senza segno, o il puntatore per il numero dei parametri \cmd{n}
               è di tipo \ctyp{short}.\\
    \cmd{l}  & Una conversione intera corrisponde a un \ctyp{long} con o 
               senza segno, o il puntatore per il numero dei parametri \cmd{n}
               è di tipo \ctyp{long}, o il carattere o la stringa seguenti
               sono in formato esteso.\\ 
    \cmd{ll} & Una conversione intera corrisponde a un \ctyp{long long} con o 
               senza segno, o il puntatore per il numero dei parametri \cmd{n}
               è di tipo \ctyp{long long}.\\
    \cmd{L}  & Una conversione in virgola mobile corrisponde a un
               \ctyp{double}.\\
    \cmd{q}  & Sinonimo di \cmd{ll}.\\
    \cmd{j}  & Una conversione intera corrisponde a un \ctyp{intmax\_t} o 
               \ctyp{uintmax\_t}.\\
    \cmd{z}  & Una conversione intera corrisponde a un \ctyp{size\_t} o 
               \ctyp{ssize\_t}.\\
    \cmd{t}  & Una conversione intera corrisponde a un \ctyp{ptrdiff\_t}.\\
    \hline
  \end{tabular}
  \caption{Il modificatore di tipo di dato per il formato di \func{printf}}
  \label{tab:file_format_type}
\end{table}

Una versione alternativa delle funzioni di output formattato, che permettono
di usare il puntatore ad una lista variabile di argomenti (vedi
sez.~\ref{sec:proc_variadic}), sono \funcd{vprintf}, \funcd{vfprintf} e
\funcd{vsprintf}, i cui prototipi sono:

\begin{funcproto}{
\fhead{stdio.h}
\fdecl{int vprintf(const char *format, va\_list ap)}
\fdesc{Scrive una stringa formattata sullo \textit{standard output}.} 
\fdecl{int vfprintf(FILE *stream, const char *format, va\_list ap)}
\fdesc{Scrive una stringa formattata su uno \textit{stream}.}
\fdecl{int vsprintf(char *str, const char *format, va\_list ap)}
\fdesc{Scrive una stringa formattata su un buffer.}
}

{Le funzioni ritornano il numero di caratteri scritti in caso di successo e un
  valore negativo per un errore.}  
\end{funcproto}

Con queste funzioni diventa possibile selezionare gli argomenti che si
vogliono passare ad una funzione di stampa, passando direttamente la lista
tramite l'argomento \param{ap}. Per poter far questo ovviamente la lista
variabile degli argomenti dovrà essere opportunamente trattata (l'argomento è
esaminato in sez.~\ref{sec:proc_variadic}), e dopo l'esecuzione della funzione
l'argomento \param{ap} non sarà più utilizzabile (in generale dovrebbe essere
eseguito un \code{va\_end(ap)} ma in Linux questo non è necessario).

Come per \func{sprintf} anche per \func{vsprintf} esiste una analoga
\funcd{vsnprintf} che pone un limite sul numero di caratteri che vengono
scritti sulla stringa di destinazione:

\begin{funcproto}{
\fhead{stdio.h}
\fdecl{vsnprintf(char *str, size\_t size, const char *format, va\_list ap)}
\fdesc{Scrive una stringa formattata su un buffer.} 
}

{La funzione ha lo stesso valore di ritorno e gli stessi errori di
  \func{vsprintf}.}
\end{funcproto}

\noindent in modo da evitare possibili \textit{buffer overflow}.


Per eliminare alla radice questi problemi, la \acr{glibc} supporta una
specifica estensione GNU che alloca dinamicamente tutto lo spazio necessario;
l'estensione si attiva al solito definendo \macro{\_GNU\_SOURCE}, le due
funzioni sono \funcd{asprintf} e \funcd{vasprintf}, ed i rispettivi prototipi
sono:

\begin{funcproto}{
\fhead{stdio.h}
\fdecl{int asprintf(char **strptr, const char *format, ...)}
\fdecl{int vasprintf(char **strptr, const char *format, va\_list ap)}
\fdesc{Scrive una stringa formattata su un buffer.} 
}

{Le funzioni hanno lo stesso valore di ritorno e gli stessi errori di
  \func{vsprintf}.}
\end{funcproto}


Entrambe le funzioni prendono come argomento \param{strptr} che deve essere
l'indirizzo di un puntatore ad una stringa di caratteri, in cui verrà
restituito (si ricordi quanto detto in sez.~\ref{sec:proc_var_passing} a
proposito dei \textit{value result argument}) l'indirizzo della stringa
allocata automaticamente dalle funzioni. Occorre inoltre ricordarsi di
invocare \func{free} per liberare detto puntatore quando la stringa non serve
più, onde evitare \textit{memory leak}.

% TODO verificare se mettere prototipi di \func{dprintf} e \func{vdprintf}

Infine una ulteriore estensione GNU definisce le due funzioni \funcm{dprintf} e
\funcm{vdprintf}, che prendono un file descriptor al posto dello
\textit{stream}. Altre estensioni permettono di scrivere con caratteri
estesi. Anche queste funzioni, il cui nome è generato dalle precedenti
funzioni aggiungendo una \texttt{w} davanti a \texttt{print}, sono trattate in
dettaglio nella documentazione della \acr{glibc}.

In corrispondenza alla famiglia di funzioni \func{printf} che si usano per
l'output formattato, l'input formattato viene eseguito con le funzioni della
famiglia \func{scanf}; fra queste le tre più importanti sono \funcd{scanf},
\funcd{fscanf} e \funcd{sscanf}, i cui prototipi sono:

\begin{funcproto}{
\fhead{stdio.h}
\fdecl{int scanf(const char *format, ...)}
\fdesc{Esegue la scansione di dati dallo \textit{standard input}.}
\fdecl{int fscanf(FILE *stream, const char *format, ...)}
\fdesc{Esegue la scansione di dati da uno \textit{stream}. } 
\fdecl{int sscanf(char *str, const char *format, ...)}
\fdesc{Esegue la scansione di dati da un buffer.} 
}

{La funzione ritorna il numero di elementi assegnati in caso di successo e
  \val{EOF} per un errore o se si raggiunta la fine del file.}
\end{funcproto}

Le funzioni eseguono una scansione della rispettiva fonte di input cercando
una corrispondenza di quanto letto con il formato dei dati specificato
da \param{format}, ed effettua le relative conversioni memorizzando il
risultato negli argomenti seguenti, il cui numero è variabile e dipende dal
valore di \param{format}. Come per le analoghe funzioni di scrittura esistono
le relative \funcm{vscanf}, \funcm{vfscanf} e \funcm{vsscanf} che usano un
puntatore ad una lista di argomenti. Le funzioni ritornano il numero di
elementi assegnati. Questi possono essere in numero inferiore a quelli
specificati, ed anche zero. Quest'ultimo valore significa che non si è trovata
corrispondenza.

Tutte le funzioni della famiglia delle \func{scanf} vogliono come argomenti i
puntatori alle variabili che dovranno contenere le conversioni; questo è un
primo elemento di disagio in quanto è molto facile dimenticarsi di questa
caratteristica.

Le funzioni leggono i caratteri dallo \textit{stream} (o dalla stringa) di
input ed eseguono un confronto con quanto indicato in \param{format}, la
sintassi di questo argomento è simile a quella usata per l'analogo di
\func{printf}, ma ci sono varie differenze.  Le funzioni di input infatti sono
più orientate verso la lettura di testo libero che verso un input formattato
in campi fissi. Uno spazio in \param{format} corrisponde con un numero
qualunque di caratteri di separazione (che possono essere spazi, tabulatori,
virgole ecc.), mentre caratteri diversi richiedono una corrispondenza
esatta. Le direttive di conversione sono analoghe a quelle di \func{printf} e
si trovano descritte in dettaglio nelle pagine di manuale e nel manuale della
\acr{glibc}.

Le funzioni eseguono la lettura dall'input, scartano i separatori (e gli
eventuali caratteri diversi indicati dalla stringa di formato) effettuando le
conversioni richieste; in caso la corrispondenza fallisca (o la funzione non
sia in grado di effettuare una delle conversioni richieste) la scansione viene
interrotta immediatamente e la funzione ritorna lasciando posizionato lo
\textit{stream} al primo carattere che non corrisponde.

Data la notevole complessità di uso di queste funzioni, che richiedono molta
cura nella definizione delle corrette stringhe di formato e sono facilmente
soggette ad errori, e considerato anche il fatto che è estremamente macchinoso
recuperare in caso di fallimento nelle corrispondenze, l'input formattato non
è molto usato. In genere infatti quando si ha a che fare con un input
relativamente semplice si preferisce usare l'input di linea ed effettuare
scansione e conversione di quanto serve direttamente con una delle funzioni di
conversione delle stringhe; se invece il formato è più complesso diventa più
facile utilizzare uno strumento come \cmd{flex}\footnote{il programma
  \cmd{flex}, è una implementazione libera di \cmd{lex} un generatore di
  analizzatori lessicali. Per i dettagli si può fare riferimento al manuale
  \cite{flex}.} per generare un analizzatore lessicale o 
\cmd{bison}\footnote{il programma \cmd{bison} è un clone del generatore di
  parser \cmd{yacc}, maggiori dettagli possono essere trovati nel relativo
  manuale \cite{bison}.} per generare un parser.



\section{Funzioni avanzate}
\label{sec:file_stream_adv_func}

In questa sezione esamineremo alcune funzioni avanzate che permettono di
eseguire operazioni di basso livello nella gestione degli \textit{stream},
come leggerne gli attributi, controllarne le modalità di bufferizzazione,
gestire in maniera esplicita i lock impliciti presenti ad uso della
programmazione \textit{multi-thread}.


\subsection{Le funzioni di controllo}
\label{sec:file_stream_cntrl}

Al contrario di quanto avviene con i file descriptor, le librerie standard del
C non prevedono nessuna funzione come la \func{fcntl} per il controllo degli
attributi dei file. Però, dato che ogni \textit{stream} si appoggia ad un file
descriptor, si può usare la funzione \funcd{fileno} per ottenere il valore di
quest'ultimo; il suo prototipo è:

\begin{funcproto}{
\fhead{stdio.h}
\fdecl{int fileno(FILE *stream)}
\fdesc{Legge il file descriptor sottostante lo \textit{stream}.} 
}

{La funzione ritorna il numero del file descriptor in caso di successo e $-1$
  per un errore, nel qual caso \var{errno} assumerà il valore \errval{EBADF}
  se \param{stream} non è valido.}
\end{funcproto}

In questo modo diventa possibile usare direttamente \func{fcntl} sul file
descriptor sottostante, ma anche se questo permette di accedere agli attributi
del file descriptor sottostante lo \textit{stream}, non ci dà nessuna
informazione riguardo alle proprietà dello \textit{stream} medesimo.  La
\acr{glibc} però supporta alcune estensioni derivate da Solaris, che
permettono di ottenere informazioni utili relative allo \textit{stream}.

Ad esempio in certi casi può essere necessario sapere se un certo
\textit{stream} è accessibile in lettura o scrittura. In genere questa
informazione non è disponibile, e ci si deve ricordare come è stato aperto il
file. La cosa può essere complessa se le operazioni vengono effettuate in una
subroutine, che a questo punto necessiterà di informazioni aggiuntive rispetto
al semplice puntatore allo \textit{stream}. Questo problema può essere risolto
con le due funzioni \funcd{\_\_freadable} e \funcd{\_\_fwritable} i cui
prototipi sono:

\begin{funcproto}{
\fhead{stdio\_ext.h}
\fdecl{int \_\_freadable(FILE *stream)}
\fdesc{Controlla se uno \textit{stream} consente la lettura.} 
\fdecl{int \_\_fwritable(FILE *stream)}
\fdesc{Controlla se uno \textit{stream} consente la scrittura.} 
}

{Le funzioni ritornano un valore diverso da $0$ se l'operazione richiesta è
  consentita, non sono previste condizioni di errore.}  
\end{funcproto}

\noindent che permettono di ottenere questa informazione.

La conoscenza dell'ultima operazione effettuata su uno \textit{stream} aperto
è utile in quanto permette di trarre conclusioni sullo stato del buffer e del
suo contenuto. Altre due funzioni, \funcd{\_\_freading} e \funcd{\_\_fwriting}
servono a tale scopo, il loro prototipo è:

\begin{funcproto}{
\fhead{stdio\_ext.h}
\fdecl{int \_\_freading(FILE *stream)}
\fdesc{Controlla l'ultima operazione di lettura.}
\fdecl{int \_\_fwriting(FILE *stream)}
\fdesc{Controlla l'ultima operazione di scrittura.}
}

{Le funzioni ritornano un valore diverso da $0$ se l'operazione richiesta è
  consentita, non sono previste condizioni di errore.}
\end{funcproto}

La funzione \func{\_\_freading} restituisce un valore diverso da zero
se \param{stream} è aperto in sola lettura o se l'ultima operazione è stata di
lettura mentre \func{\_\_fwriting} restituisce un valore diverso da zero
se \param{stream} è aperto in sola scrittura o se l'ultima operazione è stata
di scrittura.

Le due funzioni permettono di determinare di che tipo è stata l'ultima
operazione eseguita su uno \textit{stream} aperto in lettura/scrittura;
ovviamente se uno \textit{stream} è aperto in sola lettura (o sola scrittura)
la modalità dell'ultima operazione è sempre determinata; l'unica ambiguità è
quando non sono state ancora eseguite operazioni, in questo caso le funzioni
rispondono come se una operazione ci fosse comunque stata.


\subsection{Il controllo della bufferizzazione}
\label{sec:file_buffering_ctrl}

Come accennato in sez.~\ref{sec:file_buffering} le librerie definiscono una
serie di funzioni che permettono di controllare il comportamento degli
\textit{stream}; se non si è specificato nulla, la modalità di buffering viene
decisa autonomamente sulla base del tipo di file sottostante, ed i buffer
vengono allocati automaticamente.

Però una volta che si sia aperto lo \textit{stream} (ma prima di aver compiuto
operazioni su di esso) è possibile intervenire sulle modalità di buffering; la
funzione che permette di controllare la bufferizzazione è \funcd{setvbuf}, il
cui prototipo è:

\begin{funcproto}{
\fhead{stdio.h}
\fdecl{int setvbuf(FILE *stream, char *buf, int mode, size\_t size)}
\fdesc{Imposta la bufferizzazione dello \textit{stream}.} 
}

{La funzione ritorna $0$ in caso di successo e un altro valore qualunque per
  un errore, nel qual caso \var{errno} assumerà un valore appropriato.}  
\end{funcproto}

La funzione imposta la bufferizzazione dello \textit{stream} \param{stream}
nella modalità indicata da \param{mode} con uno dei valori di
tab.~\ref{tab:file_stream_buf_mode}, usando \param{buf} come buffer di
lunghezza \param{size} e permette di controllare tutti gli aspetti della
bufferizzazione. L'utente può specificare un buffer da usare al posto di
quello allocato dal sistema passandone alla funzione l'indirizzo
in \param{buf} e la dimensione in \param{size}.

\begin{table}[htb]
  \centering
  \footnotesize
    \begin{tabular}[c]{|l|l|}
      \hline
      \textbf{Valore} & \textbf{Modalità} \\
      \hline
      \hline
      \constd{\_IONBF} & \textit{unbuffered}\\
      \constd{\_IOLBF} & \textit{line buffered}\\
      \constd{\_IOFBF} & \textit{fully buffered}\\
      \hline
    \end{tabular}
    \caption{Valori dell'argomento \param{mode} di \func{setvbuf} 
      per l'impostazione delle modalità di bufferizzazione.}
  \label{tab:file_stream_buf_mode}
\end{table}

Ovviamente se si usa un buffer specificato dall'utente questo deve essere
stato allocato e rimanere disponibile per tutto il tempo in cui si opera sullo
\textit{stream}. In genere conviene allocarlo con \func{malloc} e disallocarlo
dopo la chiusura del file; ma fintanto che il file è usato all'interno di una
funzione, può anche essere usata una variabile automatica. In
\headfile{stdio.h} è definita la costante \constd{BUFSIZ}, che indica le
dimensioni generiche del buffer di uno \textit{stream}, queste vengono usate
dalla funzione \func{setbuf}.  Non è detto però che tale dimensione
corrisponda sempre al valore ottimale (che può variare a seconda del
dispositivo).

Dato che la procedura di allocazione manuale è macchinosa, comporta dei
rischi, come delle scritture accidentali sul buffer, e non assicura la scelta
delle dimensioni ottimali, è sempre meglio lasciare allocare il buffer alle
funzioni di libreria, che sono in grado di farlo in maniera ottimale e
trasparente all'utente (in quanto la deallocazione avviene
automaticamente). Inoltre siccome alcune implementazioni usano parte del
buffer per mantenere delle informazioni di controllo, non è detto che le
dimensioni dello stesso coincidano con quelle su cui viene effettuato l'I/O.

Per evitare che \func{setvbuf} imposti il buffer basta passare un valore
\val{NULL} per \param{buf} e la funzione ignorerà l'argomento \param{size}
usando il buffer allocato automaticamente dal sistema.  Si potrà comunque
modificare la modalità di bufferizzazione, passando in \param{mode} uno degli
opportuni valori elencati in tab.~\ref{tab:file_stream_buf_mode}. Qualora si
specifichi la modalità non bufferizzata i valori di \param{buf} e \param{size}
vengono sempre ignorati.

Oltre a \func{setvbuf} la \acr{glibc} definisce altre tre funzioni per la
gestione della bufferizzazione di uno \textit{stream}: \funcd{setbuf},
\funcd{setbuffer} e \funcd{setlinebuf}, i rispettivi prototipi sono:

\begin{funcproto}{
\fhead{stdio.h}
\fdecl{void setbuf(FILE *stream, char *buf)}
\fdecl{void setbuffer(FILE *stream, char *buf, size\_t size)}
\fdesc{Impostano il buffer per uno \textit{stream}.} 
\fdecl{void setlinebuf(FILE *stream)}
\fdesc{Porta uno \textit{stream} in modalità \textit{line buffered}.}
}

{Le funzioni non ritornano niente e non hanno condizioni di errore.}  
\end{funcproto}


La funzione \func{setbuf} disabilita la bufferizzazione se \param{buf} è
\val{NULL}, altrimenti usa \param{buf} come buffer di dimensione
\const{BUFSIZ} in modalità \textit{fully buffered}, mentre \func{setbuffer}
disabilita la bufferizzazione se \param{buf} è \val{NULL}, altrimenti
usa \param{buf} come buffer di dimensione \param{size} in modalità
\textit{fully buffered}.  Tutte queste funzioni sono realizzate con opportune
chiamate a \func{setvbuf} e sono definite solo per compatibilità con le
vecchie librerie BSD, pertanto non è il caso di usarle se non per la
portabilità su vecchi sistemi.

Infine la \acr{glibc} provvede le funzioni non standard, anche queste
originarie di Solaris, \funcd{\_\_flbf} e \funcd{\_\_fbufsize} che permettono
di leggere le proprietà di bufferizzazione di uno \textit{stream}; i cui
prototipi sono:

\begin{funcproto}{
\fhead{stdio\_ext.h}
\fdecl{size\_t \_\_fbufsize(FILE *stream)}
\fdesc{Restituisce le dimensioni del buffer di uno \textit{stream}.}
\fdecl{int \_\_flbf(FILE *stream)}
\fdesc{Controlla la modalità di bufferizzazione di uno \textit{stream}.}
}

{Le funzioni ritornano rispettivamente la dimensione del buffer o un valore
  non nullo se lo \textit{stream} è in modalità \textit{line-buffered}, non
  sono previste condizioni di errore.}
\end{funcproto}

Come già accennato, indipendentemente dalla modalità di bufferizzazione
scelta, si può forzare lo scarico dei dati sul file con la funzione
\funcd{fflush}, il cui prototipo è:

\begin{funcproto}{
\fhead{stdio.h}
\fdecl{int fflush(FILE *stream)}
\fdesc{Forza la scrittura dei dati bufferizzati di uno \textit{stream}.} 
}

{La funzione ritorna $0$ in caso di successo e \val{EOF} per un errore, nel
  qual caso \var{errno} assumerà il valore \errval{EBADF} se \param{stream}
  non è aperto o non è aperto in scrittura, o ad uno degli errori di
  \func{write}.}
\end{funcproto}

\noindent anche di questa funzione esiste una analoga \func{fflush\_unlocked}
(accessibile definendo una fra \macro{\_BSD\_SOURCE}, \macro{\_SVID\_SOURCE} o
\macro{\_GNU\_SOURCE}) che non effettua il blocco dello \textit{stream}.

% TODO aggiungere prototipo \func{fflush\_unlocked}?

Se \param{stream} è \val{NULL} lo scarico dei dati è forzato per tutti gli
\textit{stream} aperti. Esistono però circostanze, ad esempio quando si vuole
essere sicuri che sia stato eseguito tutto l'output su terminale, in cui serve
poter effettuare lo scarico dei dati solo per gli \textit{stream} in modalità
\textit{line buffered}. Per fare questo la \acr{glibc} supporta una
estensione di Solaris, la funzione \funcd{\_flushlbf}, il cui prototipo è:

\begin{funcproto}{
\fhead{stdio-ext.h}
\fdecl{void \_flushlbf(void)}
\fdesc{Forza la scrittura dei dati bufferizzati degli \textit{stream} in
  modalità \textit{line buffered}.} 
}

{La funzione non ritorna nulla e non presenta condizioni di errore.}  
\end{funcproto}

Si ricordi comunque che lo scarico dei dati dai buffer effettuato da queste
funzioni non comporta la scrittura di questi su disco; se si vuole che il
kernel dia effettivamente avvio alle operazioni di scrittura su disco occorre
usare \func{sync} o \func{fsync} (si veda~sez.~\ref{sec:file_sync}).

Infine esistono anche circostanze in cui si vuole scartare tutto l'output
pendente; per questo si può usare \funcd{fpurge}, il cui prototipo è:

\begin{funcproto}{
\fhead{stdio.h}
\fdecl{int fpurge(FILE *stream)}
\fdesc{Cancella i buffer di uno \textit{stream}.} 
}

{La funzione ritorna $0$ in caso di successo e \val{EOF} per un errore.}  
\end{funcproto}

La funzione scarta tutti i dati non ancora scritti (se il file è aperto in
scrittura), e tutto l'input non ancora letto (se è aperto in lettura),
compresi gli eventuali caratteri rimandati indietro con \func{ungetc}.


\subsection{Gli \textit{stream} e i \textit{thread}}
\label{sec:file_stream_thread}


Gli \textit{stream} possono essere usati in applicazioni \textit{multi-thread}
allo stesso modo in cui sono usati nelle applicazioni normali, ma si deve
essere consapevoli delle possibili complicazioni anche quando non si usano i
\textit{thread}, dato che l'implementazione delle librerie è influenzata
pesantemente dalle richieste necessarie per garantirne l'uso con i
\textit{thread}.

Lo standard POSIX richiede che le operazioni sui file siano atomiche rispetto
ai \textit{thread}, per questo le operazioni sui buffer effettuate dalle
funzioni di libreria durante la lettura e la scrittura di uno \textit{stream}
devono essere opportunamente protette, in quanto il sistema assicura
l'atomicità solo per le \textit{system call}. Questo viene fatto associando ad
ogni \textit{stream} un opportuno blocco che deve essere implicitamente
acquisito prima dell'esecuzione di qualunque operazione.

Ci sono comunque situazioni in cui questo non basta, come quando un
\textit{thread} necessita di compiere più di una operazione sullo
\textit{stream} atomicamente. Per questo motivo le librerie provvedono anche
le funzioni \funcd{flockfile} e \funcd{funlockfile} che permettono la gestione
esplicita dei blocchi sugli \textit{stream}. Esse sono disponibili definendo
\macrod{\_POSIX\_THREAD\_SAFE\_FUNCTIONS} ed i loro prototipi sono:

\begin{funcproto}{
\fhead{stdio.h}
\fdecl{void flockfile(FILE *stream)}
\fdesc{Acquisisce il lock su uno \textit{stream}.} 
\fdecl{void funlockfile(FILE *stream)}
\fdesc{Rilascia  il lock su uno \textit{stream}.} 
}
{Le funzioni non ritornano nulla e non sono previste condizioni di errore.}  
\end{funcproto}

La funzione \func{flockfile} esegue l'acquisizione del lock dello
\textit{stream} \param{stream}, bloccandosi se questo risulta non è
disponibile, mentre \func{funlockfile} rilascia un lock che si è
precedentemente acquisito.

Una terza funzione, che serve a provare ad acquisire un lock senza bloccarsi
qualora non sia possibile, è \funcd{ftrylockfile}, il cui prototipo è:

\begin{funcproto}{
\fhead{stdio.h}
\fdecl{int ftrylockfile(FILE *stream)}
\fdesc{Tenta l'acquisizione del lock di uno \textit{stream}.} 
}

{La funzione ritorna $0$ in caso di acquisizione del lock ed un altro valore
  qualunque altrimenti, non sono previste condizioni di errore.}
\end{funcproto}

Con queste funzioni diventa possibile acquisire un blocco ed eseguire tutte le
operazioni volute, per poi rilasciarlo. Ma, vista la complessità delle
strutture di dati coinvolte, le operazioni di blocco non sono del tutto
indolori, e quando il locking dello \textit{stream} non è necessario (come in
tutti i programmi che non usano i \textit{thread}), tutta la procedura può
comportare dei costi pesanti in termini di prestazioni. 

Per questo motivo abbiamo visto come alle usuali funzioni di I/O non
formattato siano associate delle versioni \code{\_unlocked} (alcune previste
dallo stesso standard POSIX, altre aggiunte come estensioni dalla \acr{glibc})
che possono essere usate quando il locking non serve\footnote{in certi casi
  dette funzioni possono essere usate, visto che sono molto più efficienti,
  anche in caso di necessità di locking, una volta che questo sia stato
  acquisito manualmente.}  con prestazioni molto più elevate, dato che spesso
queste versioni (come accade per \func{getc} e \func{putc}) sono realizzate
come macro.

La sostituzione di tutte le funzioni di I/O con le relative versioni
\code{\_unlocked} in un programma che non usa i \textit{thread} è però un
lavoro abbastanza noioso. Per questo motivo la \acr{glibc} fornisce al
programmatore pigro un'altra via, anche questa mutuata da estensioni
introdotte in Solaris, da poter utilizzare per disabilitare in blocco il
locking degli \textit{stream}: l'uso della funzione \funcd{\_\_fsetlocking},
il cui prototipo è:

\begin{funcproto}{
\fhead{stdio\_ext.h}
\fdecl{int \_\_fsetlocking(FILE *stream, int type)}
\fdesc{Specifica se abilitare il locking su uno \textit{stream}.}
}

{La funzione ritorna stato di locking interno dello \textit{stream}, non sono
  previste condizioni di errore.}  
\end{funcproto}

La funzione imposta o legge lo stato della modalità in cui le operazioni di
I/O su \param{stream} vengono effettuate rispetto all'acquisizione implicita
del locking a seconda del valore specificato con \param{type}, che può
assumere uno dei valori indicati in tab.~\ref{tab:file_fsetlocking_type}.

\begin{table}[htb]
  \centering
  \footnotesize
    \begin{tabular}[c]{|l|p{8cm}|}
      \hline
      \textbf{Valore} & \textbf{Significato} \\
      \hline
      \hline
      \constd{FSETLOCKING\_INTERNAL}& Lo \textit{stream} userà da ora in poi il
                                      blocco implicito predefinito.\\
      \constd{FSETLOCKING\_BYCALLER}& Al ritorno della funzione sarà l'utente a
                                      dover gestire da solo il locking dello
                                      \textit{stream}.\\
      \constd{FSETLOCKING\_QUERY}   & Restituisce lo stato corrente della
                                      modalità di blocco dello
                                      \textit{stream}.\\
      \hline
    \end{tabular}
    \caption{Valori dell'argomento \param{type} di \func{\_\_fsetlocking} 
      per l'impostazione delle modalità di bufferizzazione.}
  \label{tab:file_fsetlocking_type}
\end{table}

La funzione, se usata con \const{FSETLOCKING\_QUERY}, non modifica la modalità
di operazione ma restituisce lo stato di locking interno dello \textit{stream}
con uno dei valori \const{FSETLOCKING\_INTERNAL} o
\const{FSETLOCKING\_BYCALLER}.

% TODO trattare \func{clearerr\_unlocked} 



%%% Local Variables: 
%%% mode: latex
%%% TeX-master: "gapil"
%%% End: 

% LocalWords:  stream cap system call kernel Ritchie glibc descriptor Stevens
% LocalWords:  buf read write filesystem st blksize stat sez l'header stdio BSD
% LocalWords:  nell'header stdin shell stdout stderr error freopen flush line
% LocalWords:  unbuffered buffered newline fully SVr fopen fdopen POSIX const
% LocalWords:  char path int fildes NULL errno malloc fcntl fclose fflush tab
% LocalWords:  dup fifo socket append EXCL ccs IRUSR IWUSR IRGRP IWGRP inode fd
% LocalWords:  IROTH IWOTH umask fseek fsetpos rewind SEEK CUR EOF EBADF close
% LocalWords:  sync fcloseall void stdlib of feof ferror clearerr ws VFS table
% LocalWords:  unlocked fread fwrite size ptr nmemb nelem gcc padding point str
% LocalWords:  lock thread fgetc getc getchar dell'overhead unsigned ap process
% LocalWords:  getwc fgetwc getwchar wint wchar WEOF putc fputc putchar  struct
% LocalWords:  SVID getw putw parsing peeking ahead ungetc gets fgets string Di
% LocalWords:  overflow Aleph stack fputs puts fgetws fputws getline ssize leak
% LocalWords:  realloc value result argument memory getdelim delim printf short
% LocalWords:  fprintf sprintf format snprintf variadic long double intmax list
% LocalWords:  uintmax ptrdiff vprintf vfprintf vsprintf vsnprintf asprintf lex
% LocalWords:  vasprintf strptr dprintf vdprintf print scanf fscanf sscanf flex
% LocalWords:  vscanf vfscanf vsscanf bison parser yacc like off VMS whence pos
% LocalWords:  lseek ftell fgetpos fpos fseeko ftello fileno Solaris freadable
% LocalWords:  fwritable ext freading fwriting buffering setvbuf BUFSIZ setbuf
% LocalWords:  IONBF IOLBF IOFBF setbuffer setlinebuf flbf fbufsize flushlbf hh
% LocalWords:  fsync fpurge flockfile ftrylockfile funlockfile  files fig flags
% LocalWords:  locking fsetlocking type virtual operation dentry unistd sys AT
% LocalWords:  modification hole functions pathname EEXIST CREAT EINTR attack
% LocalWords:  EISDIR EFBIG EOVERFLOW ELOOP NOFOLLOW ENODEV ENOENT ENOTDIR fork
% LocalWords:  EMFILE ENAMETOOLONG ENFILE ENOMEM ENOSPC EROFS exec access RDWR
% LocalWords:  RDONLY ioctl AND ACCMODE creation Denial Service DoS opendir NFS
% LocalWords:  SOURCE LARGEFILE BITS NOCTTY TRUNC SHLOCK shared EXLOCK race SGI
% LocalWords:  exclusive condition change ASYNC SIGIO CLOEXEC DIRECT NDELAY EIO
% LocalWords:  DSYNC FASYNC IRIX FreeBSD EINVAL client RSYNC creat filedes INCR
% LocalWords:  behind shutdown ESPIPE XTND truncate fallocate count EAGAIN log
% LocalWords:  timerfd Specification pwrite pread define XOPEN EPIPE SIGPIPE at
% LocalWords:  caching cache update bdflush fdatasync fstat oldfd newfd DUPFD
% LocalWords:  openat mkdirat mkdir proc ATFILE dirfd FDCWD utimes lutimes uid
% LocalWords:  utimensat faccessat fchmodat chmod fchownat chown lchown fstatat
% LocalWords:  lstat linkat mknodat mknod readlinkat readlink renameat rename
% LocalWords:  symlinkat symlink unlinkat unlink rmdir mkfifoat mkfifo owner is
% LocalWords:  gid group FOLLOW REMOVEDIR cmd arg flock SETFD GETFD GETFL SETFL
% LocalWords:  GETLK SETLK SETLKW GETOWN PID Signal SIGURG SETOWN GETSIG SETSIG
% LocalWords:  sigaction SIGINFO siginfo SETLEASE lease GETLEASE NOTIFY request
% LocalWords:  everything framebuffer ENOTTY argp CDROM lsattr chattr magic TID
% LocalWords:  number FIOCLEX FIONCLEX FIOASYNC FIONBIO FIOSETOWN FIOGETOWN pid
% LocalWords:  FIONREAD epoll FIOQSIZE side effects SAFE BYCALLER QUERY EACCES
% LocalWords:  EBUSY OpenBSD syncfs futimes timespec only init ESRCH kill NTPL
% LocalWords:  ENXIO  NONBLOCK WRONLY EPERM NOATIME ETXTBSY EWOULDBLOCK PGRP SZ
% LocalWords:  EFAULT capabilities GETPIPE SETPIPE RESOURCE dell'I all' NFSv

%%% Local Variables: 
%%% mode: latex
%%% TeX-master: "gapil"
%%% End: 

% LocalWords:  nell' du vm Documentation Urlich Drepper futimesat times l'I
%  LocalWords:  futimens fs Tread all'I ll TMPFILE EDQUOT extN Minix UDF XFS
%  LocalWords:  shmem Btrfs ubifs tmpfile fchmod fchown fsetxattr fchdir PF
%  LocalWords:  fstatfs sull' SIGTTIN EDESTADDRREQ datagram connect seal
%  LocalWords:  dirty execveat execve scandirat statx

%% system.tex
%%
%% Copyright (C) 2000-2018 Simone Piccardi.  Permission is granted to
%% copy, distribute and/or modify this document under the terms of the GNU Free
%% Documentation License, Version 1.1 or any later version published by the
%% Free Software Foundation; with the Invariant Sections being "Un preambolo",
%% with no Front-Cover Texts, and with no Back-Cover Texts.  A copy of the
%% license is included in the section entitled "GNU Free Documentation
%% License".
%%

\chapter{La gestione del sistema, del tempo e degli errori}
\label{cha:system}

In questo capitolo tratteremo varie interfacce che attengono agli aspetti più
generali del sistema, come quelle per la gestione dei parametri e della
configurazione dello stesso, quelle per la lettura dei limiti e delle
caratteristiche, quelle per il controllo dell'uso delle risorse dei processi,
quelle per la gestione ed il controllo dei filesystem, degli utenti, dei tempi
e degli errori.


\section{La gestione di caratteristiche e parametri del sistema}
\label{sec:sys_characteristics}

In questa sezione tratteremo le varie modalità con cui un programma può
ottenere informazioni riguardo alle capacità del sistema, e, per quelle per
cui è possibile, sul come modificarle. Ogni sistema unix-like infatti è
contraddistinto da un gran numero di limiti e costanti che lo caratterizzano,
e che possono dipendere da fattori molteplici, come l'architettura hardware,
l'implementazione del kernel e delle librerie, le opzioni di
configurazione. Il kernel inoltre mette a disposizione l'accesso ad alcuni
parametri che possono modificarne il comportamento.

La definizione di queste caratteristiche ed il tentativo di provvedere dei
meccanismi generali che i programmi possono usare per ricavarle è uno degli
aspetti più complessi e controversi con cui le diverse standardizzazioni si
sono dovute confrontare, spesso con risultati spesso tutt'altro che chiari.
Daremo comunque una descrizione dei principali metodi previsti dai vari
standard per ricavare sia le caratteristiche specifiche del sistema, che
quelle della gestione dei file e prenderemo in esame le modalità con cui è
possibile intervenire sui parametri del kernel.

\subsection{Limiti e caratteristiche del sistema}
\label{sec:sys_limits}

Quando si devono determinare le caratteristiche generali del sistema ci si
trova di fronte a diverse possibilità; alcune di queste infatti possono
dipendere dall'architettura dell'hardware (come le dimensioni dei tipi
interi), o dal sistema operativo (come la presenza o meno del gruppo degli
identificatori \textit{saved}), altre invece possono dipendere dalle opzioni
con cui si è costruito il sistema (ad esempio da come si è compilato il
kernel), o dalla configurazione del medesimo; per questo motivo in generale
sono necessari due tipi diversi di funzionalità:
\begin{itemize*}
\item la possibilità di determinare limiti ed opzioni al momento della
  compilazione.
\item la possibilità di determinare limiti ed opzioni durante l'esecuzione.
\end{itemize*}

La prima funzionalità si può ottenere includendo gli opportuni header file che
contengono le costanti necessarie definite come macro di preprocessore, per la
seconda invece sono ovviamente necessarie delle funzioni. La situazione è
complicata dal fatto che ci sono molti casi in cui alcuni di questi limiti
sono fissi in un'implementazione mentre possono variare in un altra: tutto
questo crea una ambiguità che non è sempre possibile risolvere in maniera
chiara. In generale quello che succede è che quando i limiti del sistema sono
fissi essi vengono definiti come macro di preprocessore nel file
\headfile{limits.h}, se invece possono variare, il loro valore sarà ottenibile
tramite la funzione \func{sysconf} (che esamineremo a breve).

\begin{table}[htb]
  \centering
  \footnotesize
  \begin{tabular}[c]{|l|r|l|}
    \hline
    \textbf{Costante}&\textbf{Valore}&\textbf{Significato}\\
    \hline
    \hline
    \constd{MB\_LEN\_MAX}&       16  & Massima dimensione di un 
                                       carattere esteso.\\
    \constd{CHAR\_BIT} &          8  & Numero di bit di \ctyp{char}.\\
    \constd{UCHAR\_MAX}&        255  & Massimo di \ctyp{unsigned char}.\\
    \constd{SCHAR\_MIN}&       -128  & Minimo di \ctyp{signed char}.\\
    \constd{SCHAR\_MAX}&        127  & Massimo di \ctyp{signed char}.\\
    \constd{CHAR\_MIN} &   0 o -128  & Minimo di \ctyp{char}.\footnotemark\\
    \constd{CHAR\_MAX} &  127 o 255  & Massimo di \ctyp{char}.\footnotemark\\
    \constd{SHRT\_MIN} &     -32768  & Minimo di \ctyp{short}.\\
    \constd{SHRT\_MAX} &      32767  & Massimo di \ctyp{short}.\\
    \constd{USHRT\_MAX}&      65535  & Massimo di \ctyp{unsigned short}.\\
    \constd{INT\_MAX}  & 2147483647  & Minimo di \ctyp{int}.\\
    \constd{INT\_MIN}  &-2147483648  & Minimo di \ctyp{int}.\\
    \constd{UINT\_MAX} & 4294967295  & Massimo di \ctyp{unsigned int}.\\
    \constd{LONG\_MAX} & 2147483647  & Massimo di \ctyp{long}.\\
    \constd{LONG\_MIN} &-2147483648  & Minimo di \ctyp{long}.\\
    \constd{ULONG\_MAX}& 4294967295  & Massimo di \ctyp{unsigned long}.\\
    \hline                
  \end{tabular}
  \caption{Costanti definite in \headfile{limits.h} in conformità allo standard
    ANSI C.}
  \label{tab:sys_ansic_macro}
\end{table}

\footnotetext[1]{il valore può essere 0 o \const{SCHAR\_MIN} a seconda che il
  sistema usi caratteri con segno o meno.} 

\footnotetext[2]{il valore può essere \const{UCHAR\_MAX} o \const{SCHAR\_MAX}
  a seconda che il sistema usi caratteri con segno o meno.}

Lo standard ANSI C definisce dei limiti che sono tutti fissi, pertanto questo
saranno sempre disponibili al momento della compilazione. Un elenco, ripreso
da \headfile{limits.h}, è riportato in tab.~\ref{tab:sys_ansic_macro}. Come si
può vedere per la maggior parte questi limiti attengono alle dimensioni dei
dati interi, che sono in genere fissati dall'architettura hardware, le
analoghe informazioni per i dati in virgola mobile sono definite a parte, ed
accessibili includendo \headfiled{float.h}. 

\begin{table}[htb]
  \centering
  \footnotesize
  \begin{tabular}[c]{|l|r|l|}
    \hline
    \textbf{Costante}&\textbf{Valore}&\textbf{Significato}\\
    \hline
    \hline
    \constd{LLONG\_MAX} & 9223372036854775807& Massimo di \ctyp{long long}.\\
    \constd{LLONG\_MIN} &-9223372036854775808& Minimo di \ctyp{long long}.\\
    \constd{ULLONG\_MAX}&18446744073709551615& Massimo di \ctyp{unsigned long
                                               long}.\\ 
    \hline                
  \end{tabular}
  \caption{Macro definite in \headfile{limits.h} in conformità allo standard
    ISO C90.}
  \label{tab:sys_isoc90_macro}
\end{table}

Lo standard prevede anche un'altra costante, \constd{FOPEN\_MAX}, che può non
essere fissa e che pertanto non è definita in \headfile{limits.h}, essa deve
essere definita in \headfile{stdio.h} ed avere un valore minimo di 8. A questi
valori lo standard ISO C90 ne aggiunge altri tre, relativi al tipo \ctyp{long
  long} introdotto con il nuovo standard, i relativi valori sono in
tab.~\ref{tab:sys_isoc90_macro}.

Ovviamente le dimensioni dei vari tipi di dati sono solo una piccola parte
delle caratteristiche del sistema; mancano completamente tutte quelle che
dipendono dalla implementazione dello stesso. Queste, per i sistemi unix-like,
sono state definite in gran parte dallo standard POSIX.1, che tratta anche i
limiti relativi alle caratteristiche dei file che vedremo in
sez.~\ref{sec:sys_file_limits}.

\begin{table}[htb]
  \centering
  \footnotesize
  \begin{tabular}[c]{|l|r|p{8cm}|}
    \hline
    \textbf{Costante}&\textbf{Valore}&\textbf{Significato}\\
    \hline
    \hline
    \constd{ARG\_MAX} &131072& Dimensione massima degli argomenti
                               passati ad una funzione della famiglia
                               \func{exec}.\\ 
    \constd{CHILD\_MAX} & 999& Numero massimo di processi contemporanei
                               che un utente può eseguire.\\
    \constd{OPEN\_MAX}  & 256& Numero massimo di file che un processo
                               può mantenere aperti in contemporanea.\\
    \constd{STREAM\_MAX}&   8& Massimo numero di stream aperti per
                               processo in contemporanea.\\
    \constd{TZNAME\_MAX}&   6& Dimensione massima del nome di una
                               \textit{timezone} (vedi
                               sez.~\ref{sec:sys_time_base})).\\  
    \constd{NGROUPS\_MAX}& 32& Numero di gruppi supplementari per
                               processo (vedi sez.~\ref{sec:proc_access_id}).\\
    \constd{SSIZE\_MAX}&32767& Valore massimo del tipo \type{ssize\_t}.\\
    \hline
  \end{tabular}
  \caption{Costanti per i limiti del sistema.}
  \label{tab:sys_generic_macro}
\end{table}

Purtroppo la sezione dello standard che tratta questi argomenti è una delle
meno chiare, tanto che Stevens, in \cite{APUE}, la porta come esempio di
``\textsl{standardese}''. Lo standard prevede che ci siano 13 macro che
descrivono le caratteristiche del sistema: 7 per le caratteristiche generiche,
riportate in tab.~\ref{tab:sys_generic_macro}, e 6 per le caratteristiche dei
file, riportate in tab.~\ref{tab:sys_file_macro}.

\begin{table}[htb]
  \centering
  \footnotesize
  \begin{tabular}[c]{|l|r|p{8cm}|}
    \hline
    \textbf{Costante}&\textbf{Valore}&\textbf{Significato}\\
    \hline
    \hline
    \macrod{\_POSIX\_ARG\_MAX}   & 4096& Dimensione massima degli argomenti
                                         passati ad una funzione della famiglia
                                         \func{exec}.\\ 
    \macrod{\_POSIX\_CHILD\_MAX} &    6& Numero massimo di processi
                                         contemporanei che un utente può 
                                         eseguire.\\
    \macrod{\_POSIX\_OPEN\_MAX}  &   16& Numero massimo di file che un processo
                                         può mantenere aperti in 
                                         contemporanea.\\
    \macrod{\_POSIX\_STREAM\_MAX}&    8& Massimo numero di stream aperti per
                                         processo in contemporanea.\\
    \macrod{\_POSIX\_TZNAME\_MAX}&    6& Dimensione massima del nome di una
                                         \textit{timezone}
                                         (vedi sez.~\ref{sec:sys_date}). \\ 
    \macrod{\_POSIX\_RTSIG\_MAX} &    8& Numero massimo di segnali
                                         \textit{real-time} (vedi
                                         sez.~\ref{sec:sig_real_time}).\\
    \macrod{\_POSIX\_NGROUPS\_MAX}&   0& Numero di gruppi supplementari per
                                         processo (vedi 
                                         sez.~\ref{sec:proc_access_id}).\\
    \macrod{\_POSIX\_SSIZE\_MAX} &32767& Valore massimo del tipo 
                                         \type{ssize\_t}.\\
    % \macrod{\_POSIX\_AIO\_LISTIO\_MAX}&2& \\
    % \macrod{\_POSIX\_AIO\_MAX}    &    1& \\
    \hline                
  \end{tabular}
  \caption{Macro dei valori minimi di alcune caratteristiche generali del
    sistema per la conformità allo standard POSIX.1.}
  \label{tab:sys_posix1_general}
\end{table}

Lo standard dice che queste macro devono essere definite in
\headfile{limits.h} quando i valori a cui fanno riferimento sono fissi, e
altrimenti devono essere lasciate indefinite, ed i loro valori dei limiti
devono essere accessibili solo attraverso \func{sysconf}.  In realtà queste
vengono sempre definite ad un valore generico. Si tenga presente poi che
alcuni di questi limiti possono assumere valori molto elevati (come
\const{CHILD\_MAX}), e non è pertanto il caso di utilizzarli per allocare
staticamente della memoria.

A complicare la faccenda si aggiunge il fatto che POSIX.1 prevede una serie di
altre costanti (il cui nome inizia sempre con \code{\_POSIX\_}) che
definiscono i valori minimi le stesse caratteristiche devono avere, perché una
implementazione possa dichiararsi conforme allo standard, alcuni dei questi
valori sono riportati in tab.~\ref{tab:sys_posix1_general}.

In genere questi valori non servono a molto, la loro unica utilità è quella di
indicare un limite superiore che assicura la portabilità senza necessità di
ulteriori controlli. Tuttavia molti di essi sono ampiamente superati in tutti
i sistemi POSIX in uso oggigiorno. Per questo è sempre meglio utilizzare i
valori ottenuti da \func{sysconf}.

\begin{table}[htb]
  \centering
  \footnotesize
  \begin{tabular}[c]{|l|p{8cm}|}
    \hline
    \textbf{Macro}&\textbf{Significato}\\
    \hline
    \hline
    \macrod{\_POSIX\_JOB\_CONTROL}& Il sistema supporta il 
                                    \textit{job control} (vedi 
                                    sez.~\ref{sec:sess_job_control}).\\
    \macrod{\_POSIX\_SAVED\_IDS}  & Il sistema supporta gli identificatori del 
                                    gruppo \textit{saved} (vedi 
                                    sez.~\ref{sec:proc_access_id})
                                    per il controllo di accesso dei processi.\\
    \macrod{\_POSIX\_VERSION}     & Fornisce la versione dello standard POSIX.1
                                    supportata nel formato YYYYMML (ad esempio 
                                    199009L).\\
    \hline
  \end{tabular}
  \caption{Alcune macro definite in \headfile{limits.h} in conformità allo
    standard POSIX.1.}
  \label{tab:sys_posix1_other}
\end{table}

Oltre ai precedenti valori e a quelli relativi ai file elencati in
tab.~\ref{tab:sys_posix1_file}, che devono essere obbligatoriamente definiti,
lo standard POSIX.1 ne prevede molti altri. La lista completa si trova
dall'header file \file{bits/posix1\_lim.h}, da non usare mai direttamente (è
incluso automaticamente all'interno di \headfile{limits.h}). Di questi vale la
pena menzionarne alcune macro di uso comune, riportate in
tab.~\ref{tab:sys_posix1_other}, che non indicano un valore specifico, ma
denotano la presenza di alcune funzionalità nel sistema, come il supporto del
\textit{job control} o degli identificatori del gruppo \textit{saved}.

Oltre allo standard POSIX.1, anche lo standard POSIX.2 definisce una serie di
altre costanti. Siccome queste sono principalmente attinenti a limiti relativi
alle applicazioni di sistema presenti, come quelli su alcuni parametri delle
espressioni regolari o del comando \cmd{bc}, non li tratteremo esplicitamente,
se ne trova una menzione completa nell'header file \file{bits/posix2\_lim.h},
e alcuni di loro sono descritti nella pagina di manuale di \func{sysconf} e
nel manuale della \acr{glibc}.

Quando uno dei limiti o delle caratteristiche del sistema può variare, per non
dover essere costretti a ricompilare un programma tutte le volte che si
cambiano le opzioni con cui è compilato il kernel, o alcuni dei parametri
modificabili al momento dell'esecuzione, è necessario ottenerne il valore
attraverso la funzione \funcd{sysconf}, cui prototipo è:

\begin{funcproto}{
\fhead{unistd.h}
\fdecl{long sysconf(int name)}
\fdesc{Restituisce il valore di un parametro di sistema.} 
}

{La funzione ritorna in caso di successo il valore del parametro richiesto, o
  1 se si tratta di un'opzione disponibile, 0 se l'opzione non è disponibile e
  $-1$ per un errore, nel qual caso però \var{errno} non viene impostata.}  
\end{funcproto}

La funzione prende come argomento un intero che specifica quale dei limiti si
vuole conoscere. Uno specchietto contenente i principali valori disponibili in
Linux è riportato in tab.~\ref{tab:sys_sysconf_par}, l'elenco completo è
contenuto in \file{bits/confname.h}, ed una lista più esaustiva, con le
relative spiegazioni, si può trovare nel manuale della \acr{glibc}.

\begin{table}[htb]
  \centering
  \footnotesize
    \begin{tabular}[c]{|l|l|p{8cm}|}
      \hline
      \textbf{Parametro}&\textbf{Macro sostituita} &\textbf{Significato}\\
      \hline
      \hline
      \texttt{\_SC\_ARG\_MAX}   & \const{ARG\_MAX}&
                                  La dimensione massima degli argomenti passati
                                  ad una funzione della famiglia \func{exec}.\\
      \texttt{\_SC\_CHILD\_MAX} & \const{CHILD\_MAX}&
                                  Il numero massimo di processi contemporanei
                                  che un utente può eseguire.\\
      \texttt{\_SC\_OPEN\_MAX}  & \const{OPEN\_MAX}&
                                  Il numero massimo di file che un processo può
                                  mantenere aperti in contemporanea.\\
      \texttt{\_SC\_STREAM\_MAX}& \const{STREAM\_MAX}&
                                  Il massimo numero di stream che un processo
                                  può mantenere aperti in contemporanea. Questo
                                  limite è previsto anche dallo standard ANSI C,
                                  che specifica la macro \const{FOPEN\_MAX}.\\
      \texttt{\_SC\_TZNAME\_MAX}& \const{TZNAME\_MAX}&
                                  La dimensione massima di un nome di una
                                  \texttt{timezone} (vedi
                                  sez.~\ref{sec:sys_date}).\\
      \texttt{\_SC\_NGROUPS\_MAX}&\const{NGROUP\_MAX}&
                                  Massimo numero di gruppi supplementari che
                                  può avere un processo (vedi
                                  sez.~\ref{sec:proc_access_id}).\\
      \texttt{\_SC\_SSIZE\_MAX} & \const{SSIZE\_MAX}& 
                                  Valore massimo del tipo di dato
                                  \type{ssize\_t}.\\ 
      \texttt{\_SC\_CLK\_TCK}   & \const{CLK\_TCK} &
                                  Il numero di \textit{clock tick} al secondo, 
                                  cioè l'unità di misura del
                                  \textit{process time} (vedi
                                  sez.~\ref{sec:sys_unix_time}).\\  
      \texttt{\_SC\_JOB\_CONTROL}&\macro{\_POSIX\_JOB\_CONTROL}&
                                  Indica se è supportato il \textit{job
                                    control} (vedi
                                  sez.~\ref{sec:sess_job_control}) in stile
                                  POSIX.\\ 
      \texttt{\_SC\_SAVED\_IDS} & \macro{\_POSIX\_SAVED\_IDS}&
                                  Indica se il sistema supporta i
                                  \textit{saved id} (vedi
                                  sez.~\ref{sec:proc_access_id}).\\  
      \texttt{\_SC\_VERSION}    & \macro{\_POSIX\_VERSION} &
                                  Indica il mese e l'anno di approvazione
                                  della revisione dello standard POSIX.1 a cui
                                  il sistema fa riferimento, nel formato
                                  YYYYMML, la revisione più recente è 199009L,
                                  che indica il Settembre 1990.\\ 
     \hline
    \end{tabular}
  \caption{Parametri del sistema leggibili dalla funzione \func{sysconf}.}
  \label{tab:sys_sysconf_par}
\end{table}

In generale ogni limite o caratteristica del sistema per cui è definita una
macro, sia dagli standard ANSI C e ISO C90, che da POSIX.1 e POSIX.2, può
essere ottenuto attraverso una chiamata a \func{sysconf}. Il nome della
costante da utilizzare come valore dell'argomento \param{name} si otterrà
aggiungendo \code{\_SC\_} ai nomi delle costanti definite dai primi due
standard (quelle di tab.~\ref{tab:sys_generic_macro}), o sostituendolo a
\code{\_POSIX\_} per le costanti definite dagli altri due standard (quelle di
tab.~\ref{tab:sys_posix1_general}).

In linea teorica si dovrebbe fare uso di \func{sysconf} solo quando la
relativa costante di sistema non è definita, quindi con un codice analogo al
seguente:
\includecodesnip{listati/get_child_max.c}
ma in realtà con Linux queste costanti sono comunque definite, indicando però
un limite generico che non è detto sia corretto; per questo motivo è sempre
meglio usare i valori restituiti da \func{sysconf}.


\subsection{Limiti e caratteristiche dei file}
\label{sec:sys_file_limits}

Come per le caratteristiche generali del sistema anche per i file esistono una
serie di limiti (come la lunghezza del nome del file o il numero massimo di
link) che dipendono sia dall'implementazione che dal filesystem in uso. Anche
in questo caso lo standard prevede alcune macro che ne specificano il valore,
riportate in tab.~\ref{tab:sys_file_macro}.

\begin{table}[htb]
  \centering
  \footnotesize
  \begin{tabular}[c]{|l|r|l|}
    \hline
    \textbf{Costante}&\textbf{Valore}&\textbf{Significato}\\
    \hline
    \hline                
    \constd{LINK\_MAX}   &8 & Numero massimo di link a un file.\\
    \constd{NAME\_MAX}&  14 & Lunghezza in byte di un nome di file. \\
    \constd{PATH\_MAX}& 256 & Lunghezza in byte di un \textit{pathname}.\\
    \constd{PIPE\_BUF}&4096 & Byte scrivibili atomicamente in una \textit{pipe}
                              (vedi sez.~\ref{sec:ipc_pipes}).\\
    \constd{MAX\_CANON}&255 & Dimensione di una riga di terminale in modo 
                              canonico (vedi sez.~\ref{sec:term_io_design}).\\
    \constd{MAX\_INPUT}&255 & Spazio disponibile nella coda di input 
                              del terminale (vedi 
                              sez.~\ref{sec:term_io_design}).\\
    \hline                
  \end{tabular}
  \caption{Costanti per i limiti sulle caratteristiche dei file.}
  \label{tab:sys_file_macro}
\end{table}

Come per i limiti di sistema, lo standard POSIX.1 detta una serie di valori
minimi anche per queste caratteristiche, che ogni sistema che vuole essere
conforme deve rispettare. Le relative macro sono riportate in
tab.~\ref{tab:sys_posix1_file} e per esse vale lo stesso discorso fatto per le
analoghe di tab.~\ref{tab:sys_posix1_general}.

\begin{table}[htb]
  \centering
  \footnotesize
  \begin{tabular}[c]{|l|r|l|}
    \hline
    \textbf{Macro}&\textbf{Valore}&\textbf{Significato}\\
    \hline
    \hline
    \macrod{\_POSIX\_LINK\_MAX}   &8 & Numero massimo di link a un file.\\
    \macrod{\_POSIX\_NAME\_MAX}&  14 & Lunghezza in byte di un nome di file.\\
    \macrod{\_POSIX\_PATH\_MAX}& 256 & Lunghezza in byte di un 
                                       \textit{pathname}.\\
    \macrod{\_POSIX\_PIPE\_BUF}& 512 & Byte scrivibili atomicamente in una
                                       \textit{pipe}.\\
    \macrod{\_POSIX\_MAX\_CANON}&255 & Dimensione di una riga di
                                       terminale in modo canonico.\\
    \macrod{\_POSIX\_MAX\_INPUT}&255 & Spazio disponibile nella coda di input 
                                       del terminale.\\
%    \macrod{\_POSIX\_MQ\_OPEN\_MAX}&  8& \\
%    \macrod{\_POSIX\_MQ\_PRIO\_MAX}& 32& \\
%    \macrod{\_POSIX\_FD\_SETSIZE}& 16 & \\
%    \macrod{\_POSIX\_DELAYTIMER\_MAX}& 32 & \\
    \hline
  \end{tabular}
  \caption{Costanti dei valori minimi delle caratteristiche dei file per la
    conformità allo standard POSIX.1.}
  \label{tab:sys_posix1_file}
\end{table}

Tutti questi limiti sono definiti in \headfile{limits.h}; come nel caso
precedente il loro uso è di scarsa utilità in quanto ampiamente superati in
tutte le implementazioni moderne. In generale i limiti per i file sono molto
più soggetti ad essere variabili rispetto ai limiti generali del sistema; ad
esempio parametri come la lunghezza del nome del file o il numero di link
possono variare da filesystem a filesystem.

Per questo motivo quando si ha a che fare con limiti relativi ai file questi
devono essere sempre controllati con la funzione \funcd{pathconf}, il cui
prototipo è:

\begin{funcproto}{
\fhead{unistd.h}
\fdecl{long pathconf(char *path, int name)}
\fdesc{Restituisce il valore di un parametro dei file.} 
}

{La funzione ritorna il valore del parametro richiesto in caso di successo e
  $-1$ per un errore, nel qual caso \var{errno} viene impostata ad uno degli
  errori possibili relativi all'accesso a \param{path}.}
\end{funcproto}

La funzione richiede che si specifichi il limite che si vuole controllare con
l'argomento \param{name}, per il quale si deve usare la relativa costante
identificativa, il cui nome si ottiene da quelle descritte in
tab.~\ref{tab:sys_file_macro} e tab.~\ref{tab:sys_posix1_file} con la stessa
convenzione già vista con \func{sysconf}, ma un questo caso con l'uso del
suffisso ``\texttt{\_PC\_}''.

In questo caso la funzione richiede anche un secondo argomento \param{path}
che specifichi a quale file si fa riferimento, dato che il valore del limite
cercato può variare a seconda del filesystem su cui si trova il file. Una
seconda versione della funzione, \funcd{fpathconf}, opera su un file
descriptor invece che su un \textit{pathname}, il suo prototipo è:

\begin{funcproto}{
\fhead{unistd.h}
\fdecl{long fpathconf(int fd, int name)}
\fdesc{Restituisce il valore di un parametro dei file.} 
}

{È identica a \func{pathconf} solo che utilizza un file descriptor invece di
  un \textit{pathname}; pertanto gli errori restituiti in \var{errno} cambiano
  di conseguenza.}
\end{funcproto}
\noindent ed il suo comportamento è identico a quello di \func{pathconf} a
parte quello di richiedere l'indicazione di un file descriptor
nell'argomento \param{fd}.



\subsection{I parametri del kernel ed il filesystem \texttt{/proc}}
\label{sec:sys_sysctl}

Tradizionalmente la funzione che permette la lettura ed l'impostazione dei
parametri del sistema è \funcm{sysctl}. Si tratta di una funzione derivata da
BSD4.4 ed introdotta su Linux a partire dal kernel 1.3.57, ma oggi il suo uso
è totalmente deprecato.  Una \textit{system call} \funcm{\_sysctl} continua ad
esistere, ma non dispone più di una interfaccia nella \acr{glibc} ed il suo
utilizzo può essere effettuato solo tramite \func{syscall}, ma di nuovo questo
viene sconsigliato in quanto la funzionalità non è più mantenuta e molto
probabilmente sarà rimossa nel prossimo futuro. Per questo motivo eviteremo di
trattarne i particolari.

Lo scopo di \funcm{sysctl} era quello di fornire ai programmi una modalità per
modificare i parametri di sistema. Questi erano organizzati in maniera
gerarchica all'interno di un albero e per accedere a ciascuno di essi
occorreva specificare un percorso attraverso i vari nodi dell'albero, in
maniera analoga a come avviene per la risoluzione di un \textit{pathname}.

I parametri accessibili e modificabili attraverso questa funzione sono
moltissimi, dipendendo anche dallo stato corrente del kernel, ad esempio dai
moduli che sono stati caricati nel sistema. Inoltre non essendo standardizzati
i loro nomi possono variare da una versione di kernel all'altra, alcuni esempi
di questi parametri sono:
\begin{itemize*}
\item il nome di dominio
\item i parametri del meccanismo di \textit{paging}.
\item il filesystem montato come radice
\item la data di compilazione del kernel
\item i parametri dello stack TCP
\item il numero massimo di file aperti
\end{itemize*}



\index{file!filesystem~\texttt  {/proc}!definizione|(}

Dato che fin dall'inizio i parametri erano organizzati in una struttura
albero, è parso naturale rimappare questa organizzazione utilizzando il
filesystem \file{/proc}. Questo è un filesystem completamente virtuale, il cui
contenuto è generato direttamente dal kernel, che non fa riferimento a nessun
dispositivo fisico, ma presenta in forma di file e directory i dati di alcune
delle strutture interne del kernel stesso. Il suo utilizzo principale, come
denuncia il nome stesso, è quello di fornire una interfaccia per ottenere i
dati relativi ai processi (venne introdotto a questo scopo su BSD), ma nel
corso del tempo il suo uso è stato ampliato.

All'interno di questo filesystem sono pertanto presenti una serie di file che
riflettono il contenuto dei parametri del kernel (molti dei quali accessibili
in sola lettura) e in altrettante directory, nominate secondo il relativo
\ids{PID}, vengono mantenute le informazioni relative a ciascun processo
attivo nel sistema.

In particolare l'albero dei valori dei parametri di sistema impostabili con
\func{sysctl} viene presentato in forma di una gerarchia di file e directory a
partire dalla directory \file{/proc/sys}, cosicché è possibile accedere al
valore di un parametro del kernel tramite il \textit{pathname} ad un file
sotto \file{/proc/sys} semplicemente leggendone il contenuto, così come si può
modificare un parametro scrivendo sul file ad esso corrispondente.

Il kernel si occupa di generare al volo il contenuto ed i nomi dei file
corrispondenti ai vari parametri che sono presenti, e questo ha il grande
vantaggio di rendere accessibili gli stessi ad un qualunque comando di shell e
di permettere la navigazione dell'albero in modo da riconoscere quali
parametri sono presenti senza dover cercare un valore all'interno di una
pagina di manuale.

Inizialmente l'uso del filesystem \file{/proc} serviva soltanto a replicare
l'accesso, con altrettante corrispondenze ai file presenti in
\file{/proc/sys}, ai parametri impostabili tradizionalmente con \func{sysctl},
ma vista la assoluta naturalità dell'interfaccia, e la sua maggiore
efficienza, nelle versioni più recenti del kernel questa è diventata la
modalità canonica per modificare i parametri del kernel, evitando di dover
ricorrere all'uso di una \textit{system call} specifica che pur essendo ancora
presente, prima o poi verrà eliminata.

Nonostante la semplificazione nella gestione ottenuta con l'uso di
\file{/proc/sys} resta il problema generale di conoscere il significato di
ciascuno degli innumerevoli parametri che vi si trovano. Purtroppo la
documentazione degli stessi spesso risulta incompleta e non aggiornata, ma
buona parte di quelli può importanti sono descritti dalla documentazione
inclusa nei sorgenti del kernel, nella directory \file{Documentation/sysctl}.

Ma oltre alle informazioni che sostituiscono quelle ottenibili dalla ormai
deprecata \func{sysctl} dentro \file{/proc} sono disponibili moltissime altre
informazioni, fra cui ad esempio anche quelle fornite dalla funzione di
sistema \funcd{uname},\footnote{con Linux ci sono in realtà 3 \textit{system
    call} diverse per le dimensioni delle stringhe restituite, le prime due
  usano rispettivamente delle lunghezze di 9 e 65 byte, la terza usa anch'essa
  65 byte, ma restituisce anche l'ultimo campo, \var{domainname}, con una
  lunghezza di 257 byte, la \acr{glibc} provvede a mascherare questi dettagli
  usando la versione più recente disponibile.} il cui prototipo è:

\begin{funcproto}{
\fhead{sys/utsname.h}
\fdecl{int uname(struct utsname *info)}
\fdesc{Restituisce informazioni generali sul sistema.} 
}

{La funzione ritorna $0$ in caso di successo e $-1$ per un errore, nel qual
  caso \var{errno} può assumere solo il valore \errval{EFAULT}.}  
\end{funcproto}

La funzione, che viene usata dal comando \cmd{uname}, restituisce una serie di
informazioni relative al sistema nelle stringhe che costituiscono i campi
della struttura \struct{utsname} (la cui definizione è riportata in
fig.~\ref{fig:sys_utsname}) che viene scritta nel buffer puntato
dall'argomento \param{info}.

\begin{figure}[!ht!b]
  \footnotesize \centering
  \begin{minipage}[c]{0.85\textwidth}
    \includestruct{listati/ustname.h}
  \end{minipage}
  \normalsize 
  \caption{La struttura \structd{utsname}.} 
  \label{fig:sys_utsname}
\end{figure}

Si noti come in fig.~\ref{fig:sys_utsname} le dimensioni delle stringhe di
\struct{utsname} non sono specificate.  Il manuale della \acr{glibc} indica
due costanti per queste dimensioni, \constd{\_UTSNAME\_LENGTH} per i campi
standard e \constd{\_UTSNAME\_DOMAIN\_LENGTH} per quello relativo al nome di
dominio, altri sistemi usano nomi diversi come \constd{SYS\_NMLN} o
\constd{\_SYS\_NMLN} o \constd{UTSLEN} che possono avere valori diversi. Dato
che il buffer per \struct{utsname} deve essere preallocato l'unico modo per
farlo in maniera sicura è allora usare come dimensione il valore ottenuto con
\code{sizeof(utsname)}.

Le informazioni vengono restituite in ciascuno dei singoli campi di
\struct{utsname} in forma di stringhe terminate dal carattere NUL. In
particolare dette informazioni sono:
\begin{itemize*}
\item il nome del sistema operativo;
\item il nome della macchine (l'\textit{hostname});
\item il nome della release del kernel;
\item il nome della versione del kernel;
\item il tipo di hardware della macchina;
\item il nome del domino (il \textit{domainname});
\end{itemize*}
ma l'ultima di queste informazioni è stata aggiunta di recente e non è
prevista dallo standard POSIX, per questo essa è accessibile, come mostrato in
fig.~\ref{fig:sys_utsname}, solo se si è definita la macro
\macro{\_GNU\_SOURCE}.

Come accennato queste stesse informazioni, anche se a differenza di
\func{sysctl} la funzione continua ad essere mantenuta, si possono ottenere
direttamente tramite il filesystem \file{/proc}, esse infatti sono mantenute
rispettivamente nei file \sysctlrelfiled{kernel}{ostype},
\sysctlrelfiled{kernel}{hostname}, \sysctlrelfiled{kernel}{osrelease},
\sysctlrelfiled{kernel}{version} e \sysctlrelfiled{kernel}{domainname} che si
trovano sotto la directory \file{/proc/sys/kernel/}.

\index{file!filesystem~\texttt  {/proc}!definizione|)}



\section{La gestione del sistema}
\label{sec:sys_management}

In questa sezione prenderemo in esame le interfacce di programmazione messe a
disposizione per affrontare una serie di tematiche attinenti la gestione
generale del sistema come quelle relative alla gestione di utenti e gruppi, al
trattamento delle informazioni relative ai collegamenti al sistema, alle
modalità per effettuare lo spegnimento o il riavvio di una macchina.


\subsection{La gestione delle informazioni su utenti e gruppi}
\label{sec:sys_user_group}

Tradizionalmente le informazioni utilizzate nella gestione di utenti e gruppi
(password, corrispondenze fra nomi simbolici e \ids{UID} numerici, home
directory, ecc.)  venivano registrate all'interno dei due file di testo
\conffiled{/etc/passwd} ed \conffiled{/etc/group}, il cui formato è descritto
dalle relative pagine del manuale\footnote{nella quinta sezione, quella dei
  file di configurazione (esistono comandi corrispondenti), una trattazione
  sistemistica dell'intero argomento coperto in questa sezione si consulti
  sez.~4.3 di \cite{AGL}.} e tutte le funzioni che richiedevano l'accesso a
queste informazione andavano a leggere direttamente il contenuto di questi
file.

In realtà oltre a questi due file da molto tempo gran parte dei sistemi
unix-like usano il cosiddetto sistema delle \textit{shadow password} che
prevede anche i due file \conffiled{/etc/shadow} e \conffiled{/etc/gshadow}, in
cui sono state spostate le informazioni di autenticazione (ed inserite alcune
estensioni di gestione avanzata) per toglierle dagli altri file che devono
poter essere letti da qualunque processo per poter effettuare l'associazione
fra username e \ids{UID}.

Col tempo però questa impostazione ha incominciato a mostrare dei limiti. Da
una parte il meccanismo classico di autenticazione è stato ampliato, ed oggi
la maggior parte delle distribuzioni di GNU/Linux usa la libreria PAM (sigla
che sta per \textit{Pluggable Authentication Method}) che fornisce una
interfaccia comune per i processi di autenticazione, svincolando completamente
le singole applicazioni dai dettagli del come questa viene eseguita e di dove
vengono mantenuti i dati relativi. Si tratta di un sistema modulare, in cui è
possibile utilizzare anche più meccanismi insieme, diventa così possibile
avere vari sistemi di riconoscimento (biometria, chiavi hardware, ecc.),
diversi formati per le password e diversi supporti per le informazioni. Il
tutto avviene in maniera trasparente per le applicazioni purché per ciascun
meccanismo si disponga della opportuna libreria che implementa l'interfaccia
di PAM.

Dall'altra parte, il diffondersi delle reti e la necessità di centralizzare le
informazioni degli utenti e dei gruppi per insiemi di macchine e servizi
all'interno di una stessa organizzazione, in modo da mantenere coerenti i
dati, ha portato anche alla necessità di poter recuperare e memorizzare dette
informazioni su supporti diversi dai file citati, introducendo il sistema del
\textit{Name Service Switch} (che tratteremo brevemente in
sez.~\ref{sec:sock_resolver}) dato che la sua applicazione è cruciale nella
procedura di risoluzione di nomi di rete.

In questo paragrafo ci limiteremo comunque a trattare le funzioni classiche
per la lettura delle informazioni relative a utenti e gruppi tralasciando
completamente quelle relative all'autenticazione. 
%  Per questo non tratteremo
% affatto l'interfaccia di PAM, ma approfondiremo invece il sistema del
% \textit{Name Service Switch}, un meccanismo messo a disposizione dalla
% \acr{glibc} per modularizzare l'accesso a tutti i servizi in cui sia
% necessario trovare una corrispondenza fra un nome ed un numero (od altra
% informazione) ad esso associato, come appunto, quella fra uno username ed un
% \ids{UID} o fra un \ids{GID} ed il nome del gruppo corrispondente.
Le prime funzioni che vedremo sono quelle previste dallo standard POSIX.1;
queste sono del tutto generiche e si appoggiano direttamente al \textit{Name
  Service Switch}, per cui sono in grado di ricevere informazioni qualunque
sia il supporto su cui esse vengono mantenute.  Per leggere le informazioni
relative ad un utente si possono usare due funzioni, \funcd{getpwuid} e
\funcd{getpwnam}, i cui prototipi sono:

\begin{funcproto}{
\fhead{pwd.h} 
\fhead{sys/types.h} 
\fdecl{struct passwd *getpwuid(uid\_t uid)}
\fdecl{struct passwd *getpwnam(const char *name)} 
\fdesc{Restituiscono le informazioni relative all'utente specificato.} 
}

{Le funzioni ritornano il puntatore alla struttura contenente le informazioni
  in caso di successo e \val{NULL} nel caso non sia stato trovato nessun
  utente corrispondente a quanto specificato, nel qual caso \var{errno}
  assumerà il valore riportato dalle funzioni di sistema sottostanti.}
\end{funcproto}

Le due funzioni forniscono le informazioni memorizzate nel registro degli
utenti (che nelle versioni più recenti per la parte di credenziali di
autenticazione vengono ottenute attraverso PAM) relative all'utente
specificato attraverso il suo \ids{UID} o il nome di login. Entrambe le
funzioni restituiscono un puntatore ad una struttura di tipo \struct{passwd}
la cui definizione (anch'essa eseguita in \headfiled{pwd.h}) è riportata in
fig.~\ref{fig:sys_passwd_struct}, dove è pure brevemente illustrato il
significato dei vari campi.

\begin{figure}[!htb]
  \footnotesize
  \centering
  \begin{minipage}[c]{0.8\textwidth}
    \includestruct{listati/passwd.h}
  \end{minipage} 
  \normalsize 
  \caption{La struttura \structd{passwd} contenente le informazioni relative
    ad un utente del sistema.}
  \label{fig:sys_passwd_struct}
\end{figure}

La struttura usata da entrambe le funzioni è allocata staticamente, per questo
motivo viene sovrascritta ad ogni nuova invocazione, lo stesso dicasi per la
memoria dove sono scritte le stringhe a cui i puntatori in essa contenuti
fanno riferimento. Ovviamente questo implica che dette funzioni non possono
essere rientranti; per questo motivo ne esistono anche due versioni
alternative (denotate dalla solita estensione \code{\_r}), i cui prototipi
sono:

\begin{funcproto}{
\fhead{pwd.h} 
\fhead{sys/types.h} 
\fdecl{struct passwd *getpwuid\_r(uid\_t uid, struct passwd *password,
    char *buffer,\\
\phantom{struct passwd *getpwuid\_r(}size\_t buflen, struct passwd **result)}
\fdecl{struct passwd *getpwnam\_r(const char *name, struct passwd
    *password, char *buffer,\\
\phantom{struct passwd *getpwnam\_r(}size\_t buflen, struct passwd **result)}
\fdesc{Restituiscono le informazioni relative all'utente specificato.} 
}

{Le funzioni ritornano $0$ in caso di successo e $-1$ per un errore, nel qual
  caso \var{errno} assumerà il valore riportato dalle di sistema funzioni
  sottostanti.}
\end{funcproto}

In questo caso l'uso è molto più complesso, in quanto bisogna prima allocare
la memoria necessaria a contenere le informazioni. In particolare i valori
della struttura \struct{passwd} saranno restituiti all'indirizzo
\param{password} mentre la memoria allocata all'indirizzo \param{buffer}, per
un massimo di \param{buflen} byte, sarà utilizzata per contenere le stringhe
puntate dai campi di \param{password}. Infine all'indirizzo puntato da
\param{result} viene restituito il puntatore ai dati ottenuti, cioè
\param{buffer} nel caso l'utente esista, o \val{NULL} altrimenti.  Qualora i
dati non possano essere contenuti nei byte specificati da \param{buflen}, la
funzione fallirà restituendo \errcode{ERANGE} (e \param{result} sarà comunque
impostato a \val{NULL}).

Sia queste versioni rientranti che precedenti gli errori eventualmente
riportati in \var{errno} in caso di fallimento dipendono dalla sottostanti
funzioni di sistema usate per ricavare le informazioni (si veda quanto
illustrato in sez.~\ref{sec:sys_errno}) per cui se lo si vuole utilizzare è
opportuno inizializzarlo a zero prima di invocare le funzioni per essere
sicuri di non avere un residuo di errore da una chiamata precedente. Il non
aver trovato l'utente richiesto infatti può essere dovuto a diversi motivi (a
partire dal fatto che non esista) per cui si possono ottenere i valori di
errore più vari a seconda dei casi.

Del tutto analoghe alle precedenti sono le funzioni \funcd{getgrnam} e
\funcd{getgrgid} che permettono di leggere le informazioni relative ai gruppi,
i loro prototipi sono:

\begin{funcproto}{
\fhead{grp.h}
\fhead{sys/types.h}
\fdecl{struct group *getgrgid(gid\_t gid)} 
\fdecl{struct group *getgrnam(const char *name)} 
\fdesc{Restituiscono le informazioni relative al gruppo specificato.} 
}

{Le funzioni ritornano il puntatore alla struttura contenente le informazioni
  in caso di successo e \val{NULL} nel caso non sia stato trovato nessun
  utente corrispondente a quanto specificato, nel qual caso \var{errno}
  assumerà il valore riportato dalle funzioni di sistema sottostanti.}
\end{funcproto}

Come per le precedenti per gli utenti esistono anche le analoghe versioni
rientranti che di nuovo utilizzano la stessa estensione \code{\_r}; i loro
prototipi sono:

\begin{funcproto}{
\fhead{grp.h}
\fhead{sys/types.h}
\fdecl{int getgrgid\_r(gid\_t gid, struct group *grp, char *buf, 
  size\_t buflen,\\
\phantom{int getgrgid\_r(}struct group **result)}
\fdecl{int getgrnam\_r(const char *name, struct group *grp, char *buf, 
  size\_t buflen,\\
\phantom{int getgrnam\_r(}struct group **result)}
\fdesc{Restituiscono le informazioni relative al gruppo specificato.} 
}

{Le funzioni ritornano $0$ in caso di successo e $-1$ per un errore, nel qual
  caso \var{errno} assumerà il valore riportato dalle funzioni di sistema
  sottostanti.}
\end{funcproto}

Il comportamento di tutte queste funzioni è assolutamente identico alle
precedenti che leggono le informazioni sugli utenti, l'unica differenza è che
in questo caso le informazioni vengono restituite in una struttura di tipo
\struct{group}, la cui definizione è riportata in
fig.~\ref{fig:sys_group_struct}.

\begin{figure}[!htb]
  \footnotesize
  \centering
  \begin{minipage}[c]{0.8\textwidth}
    \includestruct{listati/group.h}
  \end{minipage} 
  \normalsize 
  \caption{La struttura \structd{group} contenente le informazioni relative ad
    un gruppo del sistema.}
  \label{fig:sys_group_struct}
\end{figure}

Le funzioni viste finora sono in grado di leggere le informazioni sia
direttamente dal file delle password in \conffile{/etc/passwd} che tramite il
sistema del \textit{Name Service Switch} e sono completamente generiche. Si
noti però che non c'è una funzione che permetta di impostare direttamente una
password.\footnote{in realtà questo può essere fatto ricorrendo alle funzioni
  della libreria PAM, ma questo non è un argomento che tratteremo qui.} Dato
che POSIX non prevede questa possibilità esiste un'altra interfaccia che lo
fa, derivata da SVID le cui funzioni sono riportate in
tab.~\ref{tab:sys_passwd_func}. Questa interfaccia però funziona soltanto
quando le informazioni sono mantenute su un apposito file di \textsl{registro}
di utenti e gruppi, con il formato classico di \conffile{/etc/passwd} e
\conffile{/etc/group}.

\begin{table}[htb]
  \footnotesize
  \centering
  \begin{tabular}[c]{|l|p{8cm}|}
    \hline
    \textbf{Funzione} & \textbf{Significato}\\
    \hline
    \hline
    \funcm{fgetpwent}   & Legge una voce dal file di registro degli utenti
                          specificato.\\
    \funcm{fgetpwent\_r}& Come la precedente, ma rientrante.\\ 
    \funcm{putpwent}    & Immette una voce in un file di registro degli
                          utenti.\\ 
    \funcm{getpwent}    & Legge una voce da \conffile{/etc/passwd}.\\
    \funcm{getpwent\_r} & Come la precedente, ma rientrante.\\ 
    \funcm{setpwent}    & Ritorna all'inizio di \conffile{/etc/passwd}.\\
    \funcm{endpwent}    & Chiude \conffile{/etc/passwd}.\\
    \funcm{fgetgrent}   & Legge una voce dal file di registro dei gruppi 
                         specificato.\\
    \funcm{fgetgrent\_r}& Come la precedente, ma rientrante.\\
    \funcm{putgrent}    & Immette una voce in un file di registro dei gruppi.\\
    \funcm{getgrent}    & Legge una voce da \conffile{/etc/group}.\\ 
    \funcm{getgrent\_r} & Come la precedente, ma rientrante.\\
    \funcm{setgrent}    & Ritorna all'inizio di \conffile{/etc/group}.\\
    \funcm{endgrent}    & Chiude \conffile{/etc/group}.\\
    \hline
  \end{tabular}
  \caption{Funzioni per la manipolazione dei campi di un file usato come
    registro per utenti o gruppi nel formato di \conffile{/etc/passwd} e
    \conffile{/etc/group}.} 
  \label{tab:sys_passwd_func}
\end{table}

% TODO mancano i prototipi di alcune delle funzioni

Dato che oramai tutte le distribuzioni di GNU/Linux utilizzano le
\textit{shadow password} (quindi con delle modifiche rispetto al formato
classico del file \conffile{/etc/passwd}), si tenga presente che le funzioni
di questa interfaccia che permettono di scrivere delle voci in un
\textsl{registro} degli utenti (cioè \func{putpwent} e \func{putgrent}) non
hanno la capacità di farlo specificando tutti i contenuti necessari rispetto a
questa estensione.

Per questo motivo l'uso di queste funzioni è deprecato, in quanto comunque non
funzionale rispetto ad un sistema attuale, pertanto ci limiteremo a fornire
soltanto l'elenco di tab.~\ref{tab:sys_passwd_func}, senza nessuna spiegazione
ulteriore.  Chi volesse insistere ad usare questa interfaccia può fare
riferimento alle pagine di manuale delle rispettive funzioni ed al manuale
della \acr{glibc} per i dettagli del funzionamento.



\subsection{Il registro della \textsl{contabilità} degli utenti}
\label{sec:sys_accounting}

Un altro insieme di funzioni relative alla gestione del sistema che
esamineremo è quello che permette di accedere ai dati del registro della
cosiddetta \textsl{contabilità} (o \textit{accounting}) degli utenti.  In esso
vengono mantenute una serie di informazioni storiche relative sia agli utenti
che si sono collegati al sistema, tanto per quelli correntemente collegati,
che per la registrazione degli accessi precedenti, sia relative all'intero
sistema, come il momento di lancio di processi da parte di \cmd{init}, il
cambiamento dell'orologio di sistema, il cambiamento di runlevel o il riavvio
della macchina.

I dati vengono usualmente memorizzati nei due file \file{/var/run/utmp} e
\file{/var/log/wtmp}. che sono quelli previsti dal \textit{Linux Filesystem
  Hierarchy Standard}, adottato dalla gran parte delle distribuzioni.  Quando
un utente si collega viene aggiunta una voce a \file{/var/run/utmp} in cui
viene memorizzato il nome di login, il terminale da cui ci si collega,
l'\ids{UID} della shell di login, l'orario della connessione ed altre
informazioni.  La voce resta nel file fino al logout, quando viene cancellata
e spostata in \file{/var/log/wtmp}.

In questo modo il primo file viene utilizzato per registrare chi sta
utilizzando il sistema al momento corrente, mentre il secondo mantiene la
registrazione delle attività degli utenti. A quest'ultimo vengono anche
aggiunte delle voci speciali per tenere conto dei cambiamenti del sistema,
come la modifica del runlevel, il riavvio della macchina, ecc. Tutte queste
informazioni sono descritte in dettaglio nel manuale della \acr{glibc}.

Questi file non devono mai essere letti direttamente, ma le informazioni che
contengono possono essere ricavate attraverso le opportune funzioni di
libreria. Queste sono analoghe alle precedenti funzioni (vedi
tab.~\ref{tab:sys_passwd_func}) usate per accedere al registro degli utenti,
solo che in questo caso la struttura del registro della \textsl{contabilità} è
molto più complessa, dato che contiene diversi tipi di informazione.

Le prime tre funzioni, \funcd{setutent}, \funcd{endutent} e \funcd{utmpname}
servono rispettivamente a aprire e a chiudere il file che contiene il registro
della \textsl{contabilità} degli, e a specificare su quale file esso viene
mantenuto. I loro prototipi sono:

\begin{funcproto}{
\fhead{utmp.h} 
\fdecl{void utmpname(const char *file)}
\fdesc{Specifica il file da usare come registro.} 
\fdecl{void setutent(void)}
\fdesc{Apre il file del registro.} 
\fdecl{void endutent(void)}
\fdesc{Chiude il file del registro.} 
}

{Le funzioni non ritornano nulla.}  
\end{funcproto}

Si tenga presente che le funzioni non restituiscono nessun valore, pertanto
non è possibile accorgersi di eventuali errori, ad esempio se si è impostato
un nome di file sbagliato con \func{utmpname}.

Nel caso non si sia utilizzata \func{utmpname} per specificare un file di
registro alternativo, sia \func{setutent} che \func{endutent} operano usando
il default che è \sysfile{/var/run/utmp} il cui nome, così come una serie di
altri valori di default per i \textit{pathname} di uso più comune, viene
mantenuto nei valori di una serie di costanti definite includendo
\headfiled{paths.h}, in particolare quelle che ci interessano sono:
\begin{basedescript}{\desclabelwidth{2.0cm}}
\item[\constd{\_PATH\_UTMP}] specifica il file che contiene il registro per
  gli utenti correntemente collegati, questo è il valore che viene usato se
  non si è utilizzato \func{utmpname} per modificarlo;
\item[\constd{\_PATH\_WTMP}] specifica il file che contiene il registro per
  l'archivio storico degli utenti collegati;
\end{basedescript}
che nel caso di Linux hanno un valore corrispondente ai file
\sysfile{/var/run/utmp} e \sysfile{/var/log/wtmp} citati in precedenza.

Una volta aperto il file del registro degli utenti si può eseguire una
scansione leggendo o scrivendo una voce con le funzioni \funcd{getutent},
\funcd{getutid}, \funcd{getutline} e \funcd{pututline}, i cui prototipi sono:


\begin{funcproto}{
\fhead{utmp.h}
\fdecl{struct utmp *getutent(void)}
\fdesc{Legge una voce dalla posizione corrente nel registro.} 
\fdecl{struct utmp *getutid(struct utmp *ut)}
\fdesc{Ricerca una voce sul registro.} 
\fdecl{struct utmp *getutline(struct utmp *ut)}
\fdesc{Ricerca una voce sul registro attinente a un terminale.} 
\fdecl{struct utmp *pututline(struct utmp *ut)}
\fdesc{Scrive una voce nel registro.} 
}

{Le funzioni ritornano il puntatore ad una struttura \struct{utmp} in caso di
  successo e \val{NULL} in caso di errore, nel qual caso \var{errno} assumerà
  il valore riportato dalle funzioni di sistema sottostanti.}
\end{funcproto}

Tutte queste funzioni fanno riferimento ad una struttura di tipo
\struct{utmp}, la cui definizione in Linux è riportata in
fig.~\ref{fig:sys_utmp_struct}. Le prime tre funzioni servono per leggere una
voce dal registro: \func{getutent} legge semplicemente la prima voce
disponibile, le altre due permettono di eseguire una ricerca. Aprendo il
registro con \func{setutent} ci si posiziona al suo inizio, ogni chiamata di
queste funzioni eseguirà la lettura sulle voci seguenti, pertanto la posizione
sulla voce appena letta, in modo da consentire una scansione del file. Questo
vale anche per \func{getutid} e \func{getutline}, il che comporta che queste
funzioni effettuano comunque una ricerca ``\textsl{in avanti}''.

\begin{figure}[!htb]
  \footnotesize
  \centering
  \begin{minipage}[c]{0.9\textwidth}
    \includestruct{listati/utmp.h}
  \end{minipage} 
  \normalsize 
  \caption{La struttura \structd{utmp} contenente le informazioni di una voce
    del registro di \textsl{contabilità}.}
  \label{fig:sys_utmp_struct}
\end{figure}

Con \func{getutid} si può cercare una voce specifica, a seconda del valore del
campo \var{ut\_type} dell'argomento \param{ut}.  Questo può assumere i valori
riportati in tab.~\ref{tab:sys_ut_type}, quando assume i valori
\const{RUN\_LVL}, \const{BOOT\_TIME}, \const{OLD\_TIME}, \const{NEW\_TIME},
verrà restituito la prima voce che corrisponde al tipo determinato; quando
invece assume i valori \const{INIT\_PROCESS}, \const{LOGIN\_PROCESS},
\const{USER\_PROCESS} o \const{DEAD\_PROCESS} verrà restituita la prima voce
corrispondente al valore del campo \var{ut\_id} specificato in \param{ut}.

\begin{table}[htb]
  \footnotesize
  \centering
  \begin{tabular}[c]{|l|p{8cm}|}
    \hline
    \textbf{Valore} & \textbf{Significato}\\
    \hline
    \hline
    \constd{EMPTY}         & Non contiene informazioni valide.\\
    \constd{RUN\_LVL}      & Identica il runlevel del sistema.\\
    \constd{BOOT\_TIME}    & Identifica il tempo di avvio del sistema.\\
    \constd{OLD\_TIME}     & Identifica quando è stato modificato l'orologio di
                             sistema.\\
    \constd{NEW\_TIME}     & Identifica da quanto è stato modificato il 
                             sistema.\\
    \constd{INIT\_PROCESS} & Identifica un processo lanciato da \cmd{init}.\\
    \constd{LOGIN\_PROCESS}& Identifica un processo di login.\\
    \constd{USER\_PROCESS} & Identifica un processo utente.\\
    \constd{DEAD\_PROCESS} & Identifica un processo terminato.\\
%    \constd{ACCOUNTING}    & ??? \\
    \hline
  \end{tabular}
  \caption{Classificazione delle voci del registro a seconda dei
    possibili valori del campo \var{ut\_type}.} 
  \label{tab:sys_ut_type}
\end{table}

La funzione \func{getutline} esegue la ricerca sulle voci che hanno un
\var{ut\_type} con valore uguale a \const{LOGIN\_PROCESS} o
\const{USER\_PROCESS}, restituendo la prima che corrisponde al valore di
\var{ut\_line}, che specifica il dispositivo di terminale che interessa, da
indicare senza il \file{/dev/} iniziale. Lo stesso criterio di ricerca è usato
da \func{pututline} per trovare uno spazio dove inserire la voce specificata;
qualora questo spazio non venga trovato la voce viene aggiunta in coda al
registro.

In generale occorre però tenere conto che queste funzioni non sono
completamente standardizzate, e che in sistemi diversi possono esserci
differenze; ad esempio \func{pututline} restituisce \code{void} in vari
sistemi (compreso Linux, fino alle \acr{libc5}). Qui seguiremo la sintassi
fornita dalla \acr{glibc}, ma gli standard POSIX 1003.1-2001 e XPG4.2 hanno
introdotto delle nuove strutture (e relativi file) di tipo \struct{utmpx}, che
sono un sovrainsieme della \struct{utmp} usata tradizionalmente ed altrettante
funzioni che le usano al posto di quelle citate.

La \acr{glibc} utilizzava già una versione estesa di \struct{utmp}, che
rende inutili queste nuove strutture, per questo su Linux \struct{utmpx} viene
definita esattamente come \struct{utmp}, con gli stessi campi di
fig.~\ref{fig:sys_utmp_struct}. Altrettanto dicasi per le nuove funzioni di
gestione previste dallo standard: \funcm{getutxent}, \funcm{getutxid},
\funcm{getutxline}, \funcm{pututxline}, \funcm{setutxent} e \funcm{endutxent}.

Tutte queste funzioni, definite con \struct{utmpx} dal file di dichiarazione
\headfile{utmpx.h}, su Linux sono ridefinite come sinonimi delle funzioni
appena viste, con argomento di tipo \struct{utmpx} anziché \struct{utmp} ed
hanno lo stesso identico comportamento. Per completezza viene definita anche
\funcm{utmpxname} che non è prevista da POSIX.1-2001.

Come già visto in sez.~\ref{sec:sys_user_group}, l'uso di strutture allocate
staticamente rende le funzioni di lettura dei dati appena illustrate non
rientranti. Per questo motivo la \acr{glibc} fornisce anche delle versioni
rientranti: \func{getutent\_r}, \func{getutid\_r}, \func{getutline\_r}, che
invece di restituire un puntatore restituiscono un intero e prendono due
argomenti aggiuntivi, i rispettivi prototipi sono:

\begin{funcproto}{
\fhead{utmp.h}
\fdecl{int *getutent\_r(struct utmp *buffer, struct utmp **result)}
\fdesc{Legge una voce dalla posizione corrente nel registro.} 
\fdecl{int *getutid\_r(struct utmp *buffer, struct utmp **result, struct utmp
  *ut)} 
\fdesc{Ricerca una voce sul registro.} 
\fdecl{int *getutline\_r(struct utmp *buffer, struct utmp **result, struct utmp
  *ut)} 
\fdesc{Ricerca una voce sul registro attinente a un terminale.}
}

{Le funzioni ritornano $0$ in caso di successo e $-1$ per un errore, nel qual
  caso \var{errno} assumerà il valore riportato dalle funzioni di sistema
  sottostanti.}
\end{funcproto}

Le funzioni si comportano esattamente come le precedenti analoghe non
rientranti, solo che restituiscono il risultato all'indirizzo specificato dal
primo argomento aggiuntivo \param{buffer} mentre il secondo, \param{result)}
viene usato per restituire il puntatore al buffer stesso.

Infine la \acr{glibc} fornisce altre due funzioni, \funcd{updwtmp} e
\funcd{logwtmp}, come estensione per scrivere direttamente delle voci nel file
sul registro storico \sysfile{/var/log/wtmp}; i rispettivi prototipi sono:

\begin{funcproto}{
\fhead{utmp.h}
\fdecl{void updwtmp(const char *wtmp\_file, const struct utmp *ut)}
\fdesc{Aggiunge una voce in coda al registro.} 
\fdecl{void logwtmp(const char *line, const char *name, const char *host)}
\fdesc{Aggiunge nel registro una voce con i valori specificati.} 
}

{Le funzioni non restituiscono nulla.}
\end{funcproto}

La prima funzione permette l'aggiunta di una voce in coda al file del registro
storico, indicato dal primo argomento, specificando direttamente una struttura
\struct{utmp}.  La seconda invece utilizza gli argomenti \param{line},
\param{name} e \param{host} per costruire la voce che poi aggiunge chiamando
\func{updwtmp}.

Queste funzioni non sono previste da POSIX.1-2001, anche se sono presenti in
altri sistemi (ad esempio Solaris e NetBSD), per mantenere una coerenza con le
altre funzioni definite nello standard che usano la struttura \struct{utmpx}
la \acr{glibc} definisce anche una funzione \funcm{updwtmpx}, che come in
precedenza è identica a \func{updwtmp} con la sola differenza di richiedere
l'uso di \headfiled{utmpx.h} e di una struttura \struct{utmpx} come secondo
argomento. 


\subsection{La gestione dello spegnimento e del riavvio}
\label{sec:sys_reboot}

Una delle operazioni di gestione generale del sistema è quella che attiene
alle modalità con cui se ne può gestire lo spegnimento ed il riavvio.  Perché
questo avvenga in maniera corretta, in particolare per le parti che comportano
lo spegnimento effettivo della macchina, occorre che il kernel effettui le
opportune operazioni interagendo con il BIOS ed i dispositivi che controllano
l'erogazione della potenza.

La funzione di sistema che controlla lo spegnimento ed il riavvio (ed altri
aspetti della relativa procedura) è \funcd{reboot},\footnote{la funzione
  illustrata è quella fornita dalla \acr{glibc} che maschera i dettagli di
  basso livello della \textit{system call} la quale richiede attualmente tre
  argomenti; fino al kernel 2.1.30 la \textit{system call} richiedeva un
  ulteriore quarto argomento, i primi due indicano dei \textit{magic number}
  interi che possono assumere solo alcuni valori predefiniti, il terzo un
  comando, corrispondente all'unico argomento della funzione della \acr{glibc}
  ed il quarto argomento aggiuntivo, ora ignorato, un puntatore generico ad
  ulteriori dati.}  il cui prototipo è:

\begin{funcproto}{
\fhead{unistd.h}
\fhead{sys/reboot.h}
\fdecl{int reboot(int cmd)}
\fdesc{Controlla il riavvio o l'arresto della macchina.}
}

{La funzione non ritorna o ritorna $0$ in caso di successo e $-1$ per un
  errore, nel qual caso \var{errno} assumerà uno dei valori:
  \begin{errlist}
  \item[\errcode{EFAULT}] c'è un indirizzo non valido nel passaggio degli
    argomenti con il comando \const{LINUX\_REBOOT\_CMD\_RESTART2} (obsoleto).
  \item[\errcode{EINVAL}] si sono specificati valori non validi per gli
    argomenti.
  \item[\errcode{EPERM}] il chiamante non ha i privilegi di amministratore (la
    \textit{capability} \const{CAP\_SYS\_BOOT}).
  \end{errlist}
}  
\end{funcproto}

La funzione, oltre al riavvio ed allo spegnimento, consente anche di
controllare l'uso della combinazione di tasti tradizionalmente usata come
scorciatoia da tastiera per richiedere il riavvio (\texttt{Ctrl-Alt-Del},
denominata in breve nella documentazione CAD) ed i suoi effetti specifici
dipendono dalla architettura hardware. Se si è richiesto un riavvio o uno
spegnimento in caso di successo la funzione, non esistendo più il programma,
ovviamente non ritorna, pertanto bisogna avere cura di aver effettuato tutte
le operazioni preliminari allo spegnimento prima di eseguirla.

Il comportamento della funzione viene controllato dall'argomento \param{cmd}
e deve assumere indicato con una delle costanti seguente elenco, che
illustra i comandi attualmente disponibili:

\begin{basedescript}{\desclabelwidth{2.cm}\desclabelstyle{\nextlinelabel}}
\item[\constd{LINUX\_REBOOT\_CMD\_CAD\_OFF}] Disabilita l'uso diretto della
  combinazione \texttt{Ctrl-Alt-Del}, la cui pressione si traduce nell'invio
  del segnale \signal{SIGINT} a \texttt{init} (o più in generale al processo
  con \ids{PID} 1) il cui effetto dipende dalla configurazione di
  quest'ultimo.
\item[\constd{LINUX\_REBOOT\_CMD\_CAD\_ON}] Attiva l'uso diretto della
  combinazione \texttt{Ctrl-Alt-Del}, la cui pressione si traduce
  nell'esecuzione dell'azione che si avrebbe avuto chiamando \func{reboot} con
  il comando \const{LINUX\_REBOOT\_CMD\_RESTART}.
\item[\constd{LINUX\_REBOOT\_CMD\_HALT}] Viene inviato sulla console il
  messaggio ``\textit{System halted.}'' l'esecuzione viene bloccata
  immediatamente ed il controllo passato al monitor nella ROM (se esiste e
  l'architettura lo consente). Se non si è eseguita una sincronizzazione dei
  dati su disco con \func{sync} questi saranno perduti.
\item[\constd{LINUX\_REBOOT\_CMD\_KEXEC}] viene eseguito direttamente il nuovo
  kernel che è stato opportunamente caricato in memoria da una
  \func{kexec\_load} (che tratteremo a breve) eseguita in precedenza. La
  funzionalità è disponibile solo a partire dal kernel 2.6.13 e se il kernel
  corrente è stato compilato includendo il relativo supporto.\footnote{deve
    essere stata abilitata l'opzione di compilazione \texttt{CONFIG\_KEXEC}.}
  Questo meccanismo consente di eseguire una sorta di riavvio rapido che evita
  di dover ripassare dalla inizializzazione da parte del BIOS ed il lancio del
  kernel attraverso un bootloader. Se non si è eseguita una sincronizzazione
  dei dati su disco con \func{sync} questi saranno perduti.
\item[\constd{LINUX\_REBOOT\_CMD\_POWER\_OFF}] Viene inviato sulla console il
  messaggio ``\textit{Power down.}'' l'esecuzione viene bloccata
  immediatamente e la macchina, se possibile, viene spenta.  Se non si è
  eseguita una sincronizzazione dei dati su disco con \func{sync} questi
  saranno perduti.
\item[\constd{LINUX\_REBOOT\_CMD\_RESTART}] Viene inviato sulla console il
  messaggio ``\textit{Restarting system.}'' ed avviata immediatamente la
  procedura di riavvio ordinaria. Se non si è eseguita una sincronizzazione
  dei dati su disco con \func{sync} questi saranno perduti.
\item[\constd{LINUX\_REBOOT\_CMD\_RESTART2}] Viene inviato sulla console il
  messaggio ``\textit{Restarting system with command '\%s'.}'' ed avviata
  immediatamente la procedura di riavvio usando il comando fornito
  nell'argomento \param{arg} che viene stampato al posto di \textit{'\%s'}
  (veniva usato per lanciare un altro programma al posto di \cmd{init}). Nelle
  versioni recenti questo argomento viene ignorato ed il riavvio può essere
  controllato dall'argomento di avvio del kernel \texttt{reboot=...}  Se non
  si è eseguita una sincronizzazione dei dati su disco con \func{sync} questi
  saranno perduti.
\end{basedescript}


Come appena illustrato usando il comando \const{LINUX\_REBOOT\_CMD\_KEXEC} si
può eseguire un riavvio immediato pre-caricando una immagine del kernel, che
verrà eseguita direttamente. Questo meccanismo consente di evitare la
reinizializzazione della macchina da parte del BIOS, ed oltre a velocizzare un
eventuale riavvio, ha il vantaggio poter accedere allo stato corrente della
macchina e della memoria, per cui viene usato spesso per installare un kernel
di emergenza da eseguire in caso di crollo del sistema per recuperare il
maggior numero di informazioni possibili.

La funzione di sistema che consente di caricare questa immagine del kernel è
\funcd{kexec\_load}, la funzione non viene definita nella \acr{glibc} e deve
pertanto essere invocata con \func{syscall}, il suo prototipo è:

\begin{funcproto}{
\fhead{linux/kexec.h}
\fdecl{long kexec\_load(unsigned long entry, unsigned long nr\_segments,
struct kexec\_segment\\
\phantom{long kexec\_load(}*segments, unsigned long flags)} 

\fdesc{Carica un kernel per un riavvio immediato.}
}

{La funzione ritorna $0$ in caso di successo e $-1$ per un errore, nel qual
  caso \var{errno} assumerà uno dei valori:
  \begin{errlist}
  \item[\errcode{EBUSY}] c'è già un caricamento in corso, o un altro kernel è
    già in uso.
  \item[\errcode{EINVAL}] il valore di \param{flags} non è valido o si è
    indicato un valore eccessivo per \param{nr\_segments}.
  \item[\errcode{EPERM}] il chiamante non ha i privilegi di amministratore (la
    \textit{capability} \const{CAP\_SYS\_BOOT}).
  \end{errlist}
}  
\end{funcproto}

Il primo argomento indica l'indirizzo fisico di esecuzione del nuovo kernel
questo viene caricato usando un vettore di strutture \struct{kexec\_segment}
(la cui definizione è riportata in fig.~\ref{fig:kexec_segment}) che
contengono i singoli segmenti dell'immagine. I primi due campi indicano
indirizzo e dimensione del segmento di memoria in \textit{user space}, i
secondi indirizzo e dimensione in \textit{kernel space}. 


\begin{figure}[!htb]
  \footnotesize
  \centering
  \begin{minipage}[c]{0.8\textwidth}
    \includestruct{listati/kexec_segment.h}
  \end{minipage} 
  \normalsize 
  \caption{La struttura \structd{kexec\_segment} per il caricamento di un
    segmento di immagine del kernel.}
  \label{fig:kexec_segment}
\end{figure}

L'argomento \param{flags} è una maschera binaria contenente i flag che
consentono di indicare le modalità con cui dovrà essere eseguito il nuovo
kernel. La parte meno significativa viene usata per impostare l'architettura
di esecuzione. Il valore \const{KEXEC\_ARCH\_DEFAULT} indica l'architettura
corrente, ma se ne può specificare anche una diversa, con i valori della
seconda parte di tab.~\ref{tab:kexec_load_flags}, e questa verrà usato posto
che sia effettivamente eseguibile sul proprio processore.

\begin{table}[htb]
  \footnotesize
  \centering
  \begin{tabular}[c]{|l|p{8cm}|}
    \hline
    \textbf{Valore} & \textbf{Significato}\\
    \hline
    \hline
    \constd{KEXEC\_ON\_CRASH}        & Il kernel caricato sarà eseguito
                                      automaticamente in caso di crollo del
                                      sistema.\\
    \constd{KEXEC\_PRESERVE\_CONTEXT}& Viene preservato lo stato dei programmi 
                                      e dei dispositivi prima dell'esecuzione
                                      del nuovo kernel. Viene usato
                                      principalmente per l'ibernazione del
                                      sistema ed ha senso solo se si è
                                      indicato un numero di segmento maggiore
                                      di zero.\\
    \hline
    \constd{KEXEC\_ARCH\_DEFAULT}    & Il kernel caricato verrà eseguito nella
                                      architettura corrente. \\
    \texttt{KEXEC\_ARCH\_XXX}       & Il kernel caricato verrà eseguito nella
                                      architettura indicata (con \texttt{XXX}
                                      che può essere: \texttt{386},
                                      \texttt{X86\_64}, \texttt{PPC}, 
                                      \texttt{PPC64}, \texttt{IA\_64},
                                      \texttt{ARM}, \texttt{S390},
                                      \texttt{SH}\texttt{MIPS}
                                      e \texttt{MIPS\_LE}).\\ 
%    \const{}    &  \\
    \hline
  \end{tabular}
  \caption{Valori per l'argomento \param{flags} di \func{kexec\_load}.} 
  \label{tab:kexec_load_flags}
\end{table}

I due valori più importanti sono però quelli della parte più significativa
(riportati nella prima sezione di tab.~\ref{tab:kexec_load_flags}). Il primo,
\const{KEXEC\_ON\_CRASH}, consente di impostare l'esecuzione automatica del
nuovo kernel caricato in caso di crollo del sistema, e viene usato quando si
carica un kernel di emergenza da utilizzare per poter raccogliere informazioni
diagnostiche che altrimenti verrebbero perdute non essendo il kernel ordinario
più in grado di essere eseguito in maniera coerente.  Il secondo valore,
\const{KEXEC\_PRESERVE\_CONTEXT}, indica invece di preservare lo stato dei
programmi e dei dispositivi, e viene in genere usato per realizzare la
cosiddetta ibernazione in RAM.

% TODO: introdotta con il kernel 3.17 è stata introdotta
% kexec_file_load, per caricare immagine firmate per il secure boot,
% vedi anche http://lwn.net/Articles/603116/

% TODO documentare keyctl ????
% (fare sezione dedicata ????)

% TODO documentare la Crypto API del kernel

% TODO documentare la syscall getrandom, introdotta con il kernel 3.17, vedi
% http://lwn.net/Articles/606141/, ed introdotta con le glibc solo con la
% versione 2.25, vedi https://lwn.net/Articles/711013/

%\subsection{La gestione delle chiavi crittografiche}
%\label{sec:keyctl_management}

%TODO non è chiaro se farlo qui, ma documentare la syscall bpf aggiunta con il
% kernel 3.18, vedi http://lwn.net/Articles/612878/; al riguardo vedi anche
% https://lwn.net/Articles/660331/ 

\section{Il controllo dell'uso delle risorse}
\label{sec:sys_res_limits}


Dopo aver esaminato in sez.~\ref{sec:sys_management} le funzioni che
permettono di controllare le varie caratteristiche, capacità e limiti del
sistema a livello globale, in questa sezione tratteremo le varie funzioni che
vengono usate per quantificare le risorse (CPU, memoria, ecc.) utilizzate da
ogni singolo processo e quelle che permettono di imporre a ciascuno di essi
vincoli e limiti di utilizzo.


\subsection{L'uso delle risorse}
\label{sec:sys_resource_use}

Come abbiamo accennato in sez.~\ref{sec:proc_wait} le informazioni riguardo
l'utilizzo delle risorse da parte di un processo è mantenuto in una struttura
di tipo \struct{rusage}, la cui definizione (che si trova in
\headfiled{sys/resource.h}) è riportata in
fig.~\ref{fig:sys_rusage_struct}. Si ricordi che questa è una delle
informazioni preservate attraverso una \func{exec}.

\begin{figure}[!htb]
  \footnotesize
  \centering
  \begin{minipage}[c]{0.8\textwidth}
    \includestruct{listati/rusage.h}
  \end{minipage} 
  \normalsize 
  \caption{La struttura \structd{rusage} per la lettura delle informazioni dei 
    delle risorse usate da un processo.}
  \label{fig:sys_rusage_struct}
\end{figure}

La definizione della struttura in fig.~\ref{fig:sys_rusage_struct} è ripresa
da BSD 4.3,\footnote{questo non ha a nulla a che fare con il cosiddetto
  \textit{BSD accounting} (vedi sez. \ref{sec:sys_bsd_accounting}) che si
  trova nelle opzioni di compilazione del kernel (e di norma è disabilitato)
  che serve per mantenere una contabilità delle risorse usate da ciascun
  processo in maniera molto più dettagliata.} ma attualmente solo alcuni dei
campi definiti sono effettivamente mantenuti. Con i kernel della serie 2.4 i
soli campi che sono mantenuti sono: \var{ru\_utime}, \var{ru\_stime},
\var{ru\_minflt} e \var{ru\_majflt}. Con i kernel della serie 2.6 si
aggiungono anche \var{ru\_nvcsw} e \var{ru\_nivcsw}, a partire dal 2.6.22
anche \var{ru\_inblock} e \var{ru\_oublock} e dal 2.6.32 anche
\var{ru\_maxrss}.

In genere includere esplicitamente \file{<sys/time.h>} non è più strettamente
necessario, ma aumenta la portabilità, e serve comunque quando, come nella
maggior parte dei casi, si debba accedere ai campi di \struct{rusage} relativi
ai tempi di utilizzo del processore, che sono definiti come strutture di tipo
\struct{timeval} (vedi fig.~\ref{fig:sys_timeval_struct}).

La struttura \struct{rusage} è la struttura utilizzata da \func{wait4} (si
ricordi quando visto in sez.~\ref{sec:proc_wait}) per ricavare la quantità di
risorse impiegate dal processo di cui si è letto lo stato di terminazione, ma
essa può anche essere letta direttamente utilizzando la funzione di sistema
\funcd{getrusage}, il cui prototipo è:

\begin{funcproto}{
\fhead{sys/time.h}
\fhead{sys/resource.h}
\fhead{unistd.h}
\fdecl{int getrusage(int who, struct rusage *usage)}

\fdesc{Legge la quantità di risorse usate da un processo.}
}

{La funzione ritorna $0$ in caso di successo e $-1$ per un errore, nel qual
  caso \var{errno} assumerà uno dei valori:
  \begin{errlist}
  \item[\errcode{EINVAL}] l'argomento \param{who} non è valido
  \end{errlist}
  ed inoltre \errval{EFAULT} nel suo significato generico.
}  
\end{funcproto}

La funzione ritorna i valori per l'uso delle risorse nella struttura
\struct{rusage} puntata dall'argomento \param{usage}.  L'argomento \param{who}
permette di specificare il soggetto di cui si vuole leggere l'uso delle
risorse; esso può assumere solo i valori illustrati in
tab.~\ref{tab:getrusage_who}, di questi \const{RUSAGE\_THREAD} è specifico di
Linux ed è disponibile solo a partire dal kernel 2.6.26. La funzione è stata
recepita nello standard POSIX.1-2001, che però indica come campi di
\struct{rusage} soltanto \var{ru\_utime} e \var{ru\_stime}.

\begin{table}[htb]
  \footnotesize
  \centering
  \begin{tabular}[c]{|l|p{8cm}|}
    \hline
    \textbf{Valore} & \textbf{Significato}\\
    \hline
    \hline
    \constd{RUSAGE\_SELF}     & Ritorna l'uso delle risorse del processo
                               corrente, che in caso di uso dei
                               \textit{thread} ammonta alla somma delle 
                               risorse utilizzate da tutti i \textit{thread}
                               del processo.\\ 
    \constd{RUSAGE\_CHILDREN} & Ritorna l'uso delle risorse dell'insieme dei
                               processi figli di cui è ricevuto lo stato di
                               terminazione, che a loro volta comprendono
                               quelle dei loro figli e così via.\\ 
    \constd{RUSAGE\_THREAD}   & Ritorna l'uso delle risorse del \textit{thread}
                               chiamante.\\ 
    \hline
  \end{tabular}
  \caption{Valori per l'argomento \param{who} di \func{getrusage}.} 
  \label{tab:getrusage_who}
\end{table}

I campi più utilizzati sono comunque \var{ru\_utime} e \var{ru\_stime} che
indicano rispettivamente il tempo impiegato dal processo nell'eseguire le
istruzioni in \textit{user space}, e quello impiegato dal kernel nelle
\textit{system call} eseguite per conto del processo (vedi
sez.~\ref{sec:sys_unix_time}). I campi \var{ru\_minflt} e \var{ru\_majflt}
servono a quantificare l'uso della memoria virtuale e corrispondono
rispettivamente al numero di \textit{page fault} (vedi
sez.~\ref{sec:proc_mem_gen}) avvenuti senza richiedere I/O su disco (i
cosiddetti \textit{minor page fault}), a quelli che invece han richiesto I/O
su disco (detti invece \textit{major page
  fault}).% mentre \var{ru\_nswap} ed al numero di volte che
% il processo è stato completamente tolto dalla memoria per essere inserito
% nello swap.
% TODO verificare \var{ru\_nswap} non citato nelle pagine di manuali recenti e
% dato per non utilizzato.

I campi \var{ru\_nvcsw} e \var{ru\_nivcsw} indicano il numero di volte che un
processo ha subito un \textit{context switch} da parte dello
\textit{scheduler} rispettivamente nel caso un cui questo avviene prima
dell'esaurimento della propria \textit{time-slice} (in genere a causa di una
\textit{system call} bloccante), o per averla esaurita o essere stato
interrotto da un processo a priorità maggiore. I campi \var{ru\_inblock} e
\var{ru\_oublock} indicano invece il numero di volte che è stata eseguita una
attività di I/O su un filesystem (rispettivamente in lettura e scrittura) ed
infine \var{ru\_maxrss} indica il valore più alto della \textit{Resident Set
  Size} raggiunto dal processo stesso o, nel caso sia stato usato
\const{RUSAGE\_CHILDREN}, da uno dei suoi figli.
 
Si tenga conto che per un errore di implementazione nei i kernel precedenti il
2.6.9, nonostante questo fosse esplicitamente proibito dallo standard POSIX.1,
l'uso di \const{RUSAGE\_CHILDREN} comportava l'inserimento dell'ammontare
delle risorse usate dai processi figli anche quando si era impostata una
azione di \const{SIG\_IGN} per il segnale \signal{SIGCHLD} (per i segnali si
veda cap.~\ref{cha:signals}). Il comportamento è stato corretto per aderire
allo standard a partire dal kernel 2.6.9.


\subsection{Limiti sulle risorse}
\label{sec:sys_resource_limit}

Come accennato nell'introduzione il kernel mette a disposizione delle
funzionalità che permettono non solo di mantenere dati statistici relativi
all'uso delle risorse, ma anche di imporre dei limiti precisi sul loro
utilizzo da parte sia dei singoli processi che degli utenti.

Per far questo sono definite una serie di risorse e ad ogni processo vengono
associati due diversi limiti per ciascuna di esse; questi sono il
\textsl{limite corrente} (o \textit{current limit}) che esprime un valore
massimo che il processo non può superare ad un certo momento, ed il
\textsl{limite massimo} (o \textit{maximum limit}) che invece esprime il
valore massimo che può assumere il \textsl{limite corrente}. In generale il
primo viene chiamato anche \textit{soft limit} dato che il suo valore può
essere aumentato dal processo stesso durante l'esecuzione, ciò può però essere
fatto solo fino al valore del secondo, che per questo viene detto \textit{hard
  limit}.

In generale il superamento di un limite corrente comporta o l'emissione di uno
specifico segnale o il fallimento della \textit{system call} che lo ha
provocato. A questo comportamento generico fanno eccezione \const{RLIMIT\_CPU}
in cui si ha in comportamento diverso per il superamento dei due limiti e
\const{RLIMIT\_CORE} che influenza soltanto la dimensione o l'eventuale
creazione dei file di \textit{core dump} (vedi sez.~\ref{sec:sig_standard}).

Per permettere di leggere e di impostare i limiti di utilizzo delle risorse da
parte di un processo sono previste due funzioni di sistema, \funcd{getrlimit}
e \funcd{setrlimit}, i cui prototipi sono:

\begin{funcproto}{
\fhead{sys/time.h}
\fhead{sys/resource.h}
\fhead{unistd.h}
\fdecl{int getrlimit(int resource, struct rlimit *rlim)}
\fdesc{Legge i limiti di una risorsa.}
\fdecl{int setrlimit(int resource, const struct rlimit *rlim)}
\fdesc{Imposta i limiti di una risorsa.}
}

{Le funzioni ritornano $0$ in caso di successo e $-1$ per un errore, nel qual
  caso \var{errno} assumerà uno dei valori:
  \begin{errlist}
  \item[\errcode{EINVAL}] i valori per \param{resource} non sono validi o
    nell'impostazione si è specificato \var{rlim->rlim\_cur} maggiore di
    \var{rlim->rlim\_max}.
    \item[\errcode{EPERM}] un processo senza i privilegi di amministratore ha
    cercato di innalzare i propri limiti.
  \end{errlist}
  ed inoltre \errval{EFAULT} nel suo significato generico.  
}  
\end{funcproto}

Entrambe le funzioni permettono di specificare attraverso l'argomento
\param{resource} su quale risorsa si vuole operare. L'accesso (rispettivamente
in lettura e scrittura) ai valori effettivi dei limiti viene poi effettuato
attraverso la struttura \struct{rlimit} puntata da
\param{rlim}, la cui definizione è riportata in
fig.~\ref{fig:sys_rlimit_struct}, ed i cui campi corrispondono appunto a
limite corrente e limite massimo.

\begin{figure}[!htb]
  \footnotesize
  \centering
  \begin{minipage}[c]{0.8\textwidth}
    \includestruct{listati/rlimit.h}
  \end{minipage} 
  \normalsize 
  \caption{La struttura \structd{rlimit} per impostare i limiti di utilizzo 
    delle risorse usate da un processo.}
  \label{fig:sys_rlimit_struct}
\end{figure}

Come accennato processo ordinario può alzare il proprio limite corrente fino
al valore del limite massimo, può anche ridurre, irreversibilmente, il valore
di quest'ultimo.  Nello specificare un limite, oltre a fornire dei valori
specifici, si può anche usare la costante \const{RLIM\_INFINITY} che permette
di sbloccare completamente l'uso di una risorsa. Si ricordi però che solo un
processo con i privilegi di amministratore\footnote{per essere precisi in
  questo caso quello che serve è la \textit{capability}
  \const{CAP\_SYS\_RESOURCE} (vedi sez.~\ref{sec:proc_capabilities}).} può
innalzare un limite al di sopra del valore corrente del limite massimo ed
usare un valore qualsiasi per entrambi i limiti.

Ciascuna risorsa su cui si possono applicare dei limiti è identificata da uno
specifico valore dell'argomento \param{resource}, i valori possibili per
questo argomento, ed il significato della risorsa corrispondente, dei
rispettivi limiti e gli effetti causati dal superamento degli stessi sono
riportati nel seguente elenco:

\begin{basedescript}{\desclabelwidth{2.2cm}}%\desclabelstyle{\nextlinelabel}}
\item[\constd{RLIMIT\_AS}] Questa risorsa indica, in byte, la dimensione
  massima consentita per la memoria virtuale di un processo, il cosiddetto
  \textit{Address Space}, (vedi sez.~\ref{sec:proc_mem_gen}). Se il limite
  viene superato dall'uso di funzioni come \func{brk}, \func{mremap} o
  \func{mmap} esse falliranno con un errore di \errcode{ENOMEM}, mentre se il
  superamento viene causato dalla crescita dello \textit{stack} il processo
  riceverà un segnale di \signal{SIGSEGV}. Dato che il valore usato è un
  intero di tipo \ctyp{long} nelle macchine a 32 bit questo può assumere un
  valore massimo di 2Gb (anche se la memoria disponibile può essere maggiore),
  in tal caso il limite massimo indicabile resta 2Gb, altrimenti la risorsa si
  dà per non limitata.

\item[\constd{RLIMIT\_CORE}] Questa risorsa indica, in byte, la massima
  dimensione per un file di \textit{core dump} (vedi
  sez.~\ref{sec:sig_standard}) creato nella terminazione di un processo. File
  di dimensioni maggiori verranno troncati a questo valore, mentre con un
  valore nullo si bloccherà la creazione dei \textit{core dump}.

\item[\constd{RLIMIT\_CPU}] Questa risorsa indica, in secondi, il massimo tempo
  di CPU (vedi sez.~\ref{sec:sys_cpu_times}) che il processo può usare. Il
  superamento del limite corrente comporta l'emissione di un segnale di
  \signal{SIGXCPU}, la cui azione predefinita (vedi
  sez.~\ref{sec:sig_classification}) è terminare il processo. Il segnale però
  può essere intercettato e ignorato, in tal caso esso verrà riemesso una
  volta al secondo fino al raggiungimento del limite massimo. Il superamento
  del limite massimo comporta comunque l'emissione di un segnale di
  \signal{SIGKILL}. Si tenga presente che questo è il comportamento presente
  su Linux dai kernel della serie 2.2 ad oggi, altri kernel possono avere
  comportamenti diversi per quanto avviene quando viene superato il
  \textit{soft limit}, pertanto per avere operazioni portabili è suggerito di
  intercettare sempre \signal{SIGXCPU} e terminare in maniera ordinata il
  processo con la prima ricezione.

\item[\constd{RLIMIT\_DATA}] Questa risorsa indica, in byte, la massima
  dimensione del segmento dati di un processo (vedi
  sez.~\ref{sec:proc_mem_layout}).  Il tentativo di allocare più memoria di
  quanto indicato dal limite corrente causa il fallimento della funzione di
  allocazione eseguita (\func{brk} o \func{sbrk}) con un errore di
  \errcode{ENOMEM}.

\item[\constd{RLIMIT\_FSIZE}] Questa risorsa indica, in byte, la massima
  dimensione di un file che un processo può usare. Se il processo cerca di
  scrivere o di estendere il file oltre questa dimensione riceverà un segnale
  di \signal{SIGXFSZ}, che di norma termina il processo. Se questo segnale
  viene intercettato la \textit{system call} che ha causato l'errore fallirà
  con un errore di \errcode{EFBIG}.

\item[\constd{RLIMIT\_LOCKS}] Questa risorsa indica il numero massimo di
  \textit{file lock} (vedi sez.~\ref{sec:file_locking}) e di \textit{file
    lease} (vedi sez.~\ref{sec:file_asyncronous_lease}) che un processo poteva
  effettuare.  È un limite presente solo nelle prime versioni del kernel 2.4,
  pertanto non deve essere più utilizzato.

\item[\constd{RLIMIT\_MEMLOCK}] Questa risorsa indica, in byte, l'ammontare
  massimo di memoria che può essere bloccata in RAM da un processo (vedi
  sez.~\ref{sec:proc_mem_lock}). Dato che il \textit{memory locking} viene
  effettuato sulle pagine di memoria, il valore indicato viene automaticamente
  arrotondato al primo multiplo successivo della dimensione di una pagina di
  memoria. Il limite comporta il fallimento delle \textit{system call} che
  eseguono il \textit{memory locking} (\func{mlock}, \func{mlockall} ed anche,
  vedi sez.~\ref{sec:file_memory_map}, \func{mmap} con l'operazione
  \const{MAP\_LOCKED}).

  Dal kernel 2.6.9 questo limite comprende anche la memoria che può essere
  bloccata da ciascun utente nell'uso della memoria condivisa (vedi
  sez.~\ref{sec:ipc_sysv_shm}) con \func{shmctl}, che viene contabilizzata
  separatamente ma sulla quale viene applicato questo stesso limite. In
  precedenza invece questo limite veniva applicato sulla memoria condivisa per
  processi con privilegi amministrativi, il limite su questi è stato rimosso e
  la semantica della risorsa cambiata.


\item[\constd{RLIMIT\_MSGQUEUE}] Questa risorsa indica il numero massimo di
  byte che possono essere utilizzati da un utente, identificato con
  l'\ids{UID} reale del processo chiamante, per le code di messaggi POSIX
  (vedi sez.~\ref{sec:ipc_posix_mq}). Per ciascuna coda che viene creata viene
  calcolata un'occupazione pari a:
\includecodesnip{listati/mq_occupation.c}
dove \var{attr} è la struttura \struct{mq\_attr} (vedi
fig.~\ref{fig:ipc_mq_attr}) usata nella creazione della coda. Il primo addendo
consente di evitare la creazione di una coda con un numero illimitato di
messaggi vuoti che comunque richiede delle risorse di gestione. Questa risorsa
è stata introdotta con il kernel 2.6.8.

\item[\constd{RLIMIT\_NICE}] Questa risorsa indica il numero massimo a cui può
  essere il portato il valore di \textit{nice} (vedi
  sez.~\ref{sec:proc_sched_stand}). Dato che non possono essere usati numeri
  negativi per specificare un limite, il valore di \textit{nice} viene
  calcolato come \code{20-rlim\_cur}. Questa risorsa è stata introdotta con il
  kernel 2.6.12.

\item[\constd{RLIMIT\_NOFILE}] Questa risorsa indica il numero massimo di file
  che un processo può aprire. Il tentativo di creazione di un ulteriore file
  descriptor farà fallire la funzione (\func{open}, \func{dup}, \func{pipe},
  ecc.) con un errore \errcode{EMFILE}.

\item[\constd{RLIMIT\_NPROC}] Questa risorsa indica il numero massimo di
  processi che possono essere creati dallo stesso utente, che viene
  identificato con l'\ids{UID} reale (vedi sez.~\ref{sec:proc_access_id}) del
  processo chiamante. Se il limite viene raggiunto \func{fork} fallirà con un
  \errcode{EAGAIN}.

\item[\constd{RLIMIT\_RSS}] Questa risorsa indica, in pagine di memoria, la
  dimensione massima della memoria residente (il cosiddetto RSS
  \itindex{Resident~Set~Size~(RSS)} \textit{Resident Set Size}) cioè
  l'ammontare della memoria associata al processo che risiede effettivamente
  in RAM e non a quella eventualmente portata sulla \textit{swap} o non ancora
  caricata dal filesystem per il segmento testo del programma.  Ha effetto
  solo sulle chiamate a \func{madvise} con \const{MADV\_WILLNEED} (vedi
  sez.~\ref{sec:file_memory_map}). Presente solo sui i kernel precedenti il
  2.4.30.

\item[\constd{RLIMIT\_RTPRIO}] Questa risorsa indica il valore massimo della
  priorità statica che un processo può assegnarsi o assegnare con
  \func{sched\_setscheduler} e \func{sched\_setparam} (vedi
  sez.~\ref{sec:proc_real_time}). Il limite è stato introdotto a partire dal
  kernel 2.6.12 (ma per un bug è effettivo solo a partire dal 2.6.13). In
  precedenza solo i processi con privilegi amministrativi potevano avere una
  priorità statica ed utilizzare una politica di \textit{scheduling} di tipo
  \textit{real-time}.

\item[\constd{RLIMIT\_RTTIME}] Questa risorsa indica, in microsecondi, il tempo
  massimo di CPU che un processo eseguito con una priorità statica può
  consumare. Il superamento del limite corrente comporta l'emissione di un
  segnale di \signal{SIGXCPU}, e quello del limite massimo di \signal{SIGKILL}
  con le stesse regole viste \const{RLIMIT\_CPU}: se \signal{SIGXCPU} viene
  intercettato ed ignorato il segnale verrà riemesso ogni secondo fino al
  superamento del limite massimo. Questo limite è stato introdotto con il
  kernel 2.6.25 per impedire che un processo \textit{real-time} possa bloccare
  il sistema.

\item[\constd{RLIMIT\_SIGPENDING}] Questa risorsa indica il numero massimo di
  segnali che possono essere mantenuti in coda per ciascun utente,
  identificato per \ids{UID} reale. Il limite comprende sia i segnali normali
  che quelli \textit{real-time} (vedi sez.~\ref{sec:sig_real_time}) ed è
  attivo solo per \func{sigqueue}, con \func{kill} si potrà sempre inviare un
  segnale che non sia già presente su una coda. Questo limite è stato
  introdotto con il kernel 2.6.8.

\item[\constd{RLIMIT\_STACK}] Questa risorsa indica, in byte, la massima
  dimensione dello \textit{stack} del processo. Se il processo esegue
  operazioni che estendano lo \textit{stack} oltre questa dimensione riceverà
  un segnale di \signal{SIGSEGV}.

  A partire dal kernel 2.6.23 questo stesso limite viene applicato per la gran
  parte delle architetture anche ai dati che possono essere passati come
  argomenti e variabili di ambiente ad un programma posto in esecuzione con
  \func{execve}, nella misura di un quarto del valore indicato per lo
  \textit{stack}.  Questo valore in precedenza era fisso e pari a 32 pagine di
  memoria, corrispondenti per la gran parte delle architetture a 128kb di
  dati, dal 2.6.25, per evitare problemi di compatibilità quando
  \const{RLIMIT\_STACK} è molto basso, viene comunque garantito uno spazio
  base di 32 pagine qualunque sia l'architettura.

\end{basedescript}

Si tenga conto infine che tutti i limiti eventualmente presenti su un processo
vengono ereditati dai figli da esso creati attraverso una \func{fork} (vedi
sez.~\ref{sec:proc_fork}) e mantenuti invariati per i programmi messi in
esecuzione attraverso una \func{exec} (vedi sez.~\ref{sec:proc_exec}).

Si noti come le due funzioni \func{getrlimit} e \func{setrlimit} consentano di
operare solo sul processo corrente. Per questo motivo a partire dal kernel
2.6.36 (e dalla \acr{glibc} 2.13) è stata introdotta un'altra funzione di
sistema \funcd{prlimit} il cui scopo è quello di estendere e sostituire le
precedenti.  Il suo prototipo è:

\begin{funcproto}{
\fhead{sys/resource.h}
\fdecl{int prlimit(pid\_t pid, int resource, const struct rlimit *new\_limit,\\
\phantom{int prlimit(}struct rlimit *old\_limit}
\fdesc{Legge e imposta i limiti di una risorsa.} 
}

{La funzione ritorna $0$ in caso di successo e $-1$ per un errore, nel qual
  caso \var{errno} assumerà uno dei valori:
  \begin{errlist}
  \item[\errcode{EINVAL}] i valori per \param{resource} non sono validi o
    nell'impostazione si è specificato \var{rlim->rlim\_cur} maggiore di
    \var{rlim->rlim\_max}.
  \item[\errcode{EPERM}] un processo senza i privilegi di amministratore ha
    cercato di innalzare i propri limiti o si è cercato di modificare i limiti
    di un processo di un altro utente.
  \item [\errcode{ESRCH}] il process \param{pid} non esiste.
  \end{errlist}
  ed inoltre \errval{EFAULT} nel suo significato generico.
}
\end{funcproto}

La funzione è specifica di Linux e non portabile; per essere usata richiede
che sia stata definita la macro \macro{\_GNU\_SOURCE}. Il primo argomento
indica il \ids{PID} del processo di cui si vogliono cambiare i limiti e si può
usare un valore nullo per indicare il processo chiamante.  Per modificare i
limiti di un altro processo, a meno di non avere privilegi
amministrativi,\footnote{anche in questo caso la \textit{capability}
  necessaria è \const{CAP\_SYS\_RESOURCE} (vedi
  sez.~\ref{sec:proc_capabilities}).}  l'\ids{UID} ed il \ids{GID} reale del
chiamante devono coincidere con \ids{UID} e \ids{GID} del processo indicato
per i tre gruppi reale, effettivo e salvato.

Se \param{new\_limit} non è \val{NULL} verrà usato come puntatore alla
struttura \struct{rlimit} contenente i valori dei nuovi limiti da impostare,
mentre se \param{old\_limit} non è \val{NULL} verranno letti i valori correnti
del limiti nella struttura \struct{rlimit} da esso puntata. In questo modo è
possibile sia leggere che scrivere, anche in contemporanea, i valori dei
limiti. Il significato dell'argomento \param{resource} resta identico rispetto
a \func{getrlimit} e \func{setrlimit}, così come i restanti requisiti. 


\subsection{Le informazioni sulle risorse di memoria e processore}
\label{sec:sys_memory_res}

La gestione della memoria è già stata affrontata in dettaglio in
sez.~\ref{sec:proc_memory}; abbiamo visto allora che il kernel provvede il
meccanismo della memoria virtuale attraverso la divisione della memoria fisica
in pagine.  In genere tutto ciò è del tutto trasparente al singolo processo,
ma in certi casi, come per l'I/O mappato in memoria (vedi
sez.~\ref{sec:file_memory_map}) che usa lo stesso meccanismo per accedere ai
file, è necessario conoscere le dimensioni delle pagine usate dal kernel. Lo
stesso vale quando si vuole gestire in maniera ottimale l'interazione della
memoria che si sta allocando con il meccanismo della paginazione.

Un tempo la dimensione delle pagine di memoria era fissata una volta per tutte
dall'architettura hardware, per cui il relativo valore veniva mantenuto in una
costante che bastava utilizzare in fase di compilazione. Oggi invece molte
architetture permettono di variare questa dimensione (ad esempio sui PC
recenti si possono usare pagine di 4kb e di 4 Mb) per cui per non dover
ricompilare i programmi per ogni possibile caso e relativa scelta di
dimensioni, è necessario poter utilizzare una funzione che restituisca questi
valori quando il programma viene eseguito.

Dato che si tratta di una caratteristica generale del sistema come abbiamo
visto in sez.~\ref{sec:sys_characteristics} questa dimensione può essere
ottenuta come tutte le altre attraverso una chiamata a \func{sysconf}, nel
caso specifico si dovrebbe utilizzare il parametro \const{\_SC\_PAGESIZE}. Ma
in BSD 4.2 è stata introdotta una apposita funzione di sistema
\funcd{getpagesize} che restituisce la dimensione delle pagine di memoria. La
funzione è disponibile anche su Linux (ma richiede che sia definita la macro
\macro{\_BSD\_SOURCE}) ed il suo prototipo è:

\begin{funcproto}{
\fhead{unistd.h}
\fdecl{int getpagesize(void)}
\fdesc{Legge la dimensione delle pagine di memoria.} 
}

{La funzione ritorna la dimensione di una pagina in byte, e non sono previsti
  errori.}
\end{funcproto}

La funzione è prevista in SVr4, BSD 4.4 e SUSv2, anche se questo ultimo
standard la etichetta come obsoleta, mentre lo standard POSIX 1003.1-2001 la
ha eliminata, ed i programmi che intendono essere portabili devono ricorrere
alla chiamata a \func{sysconf}. 

In Linux è implementata come una \textit{system call} nelle architetture in
cui essa è necessaria, ed in genere restituisce il valore del simbolo
\const{PAGE\_SIZE} del kernel, che dipende dalla architettura hardware, anche
se le versioni delle librerie del C precedenti la \acr{glibc} 2.1
implementavano questa funzione restituendo sempre un valore statico.

% TODO verificare meglio la faccenda di const{PAGE\_SIZE} 

La \textsl{glibc} fornisce, come specifica estensione GNU, altre due
funzioni, \funcd{get\_phys\_pages} e \funcd{get\_avphys\_pages} che permettono
di ottenere informazioni riguardo le pagine di memoria; i loro prototipi sono:

\begin{funcproto}{
\fhead{sys/sysinfo.h} 
\fdecl{long int get\_phys\_pages(void)}
\fdesc{Legge il numero totale di pagine di memoria.} 
\fdecl{long int get\_avphys\_pages(void)} 
\fdesc{Legge il numero di pagine di memoria disponibili nel sistema.} 
}

{La funzioni ritornano il numero di pagine, e non sono previsti
  errori.}  
\end{funcproto}

Queste funzioni sono equivalenti all'uso della funzione \func{sysconf}
rispettivamente con i parametri \const{\_SC\_PHYS\_PAGES} e
\const{\_SC\_AVPHYS\_PAGES}. La prima restituisce il numero totale di pagine
corrispondenti alla RAM della macchina; la seconda invece la memoria
effettivamente disponibile per i processi.

La \acr{glibc} supporta inoltre, come estensioni GNU, due funzioni che
restituiscono il numero di processori della macchina (e quello dei processori
attivi); anche queste sono informazioni comunque ottenibili attraverso
\func{sysconf} utilizzando rispettivamente i parametri
\const{\_SC\_NPROCESSORS\_CONF} e \const{\_SC\_NPROCESSORS\_ONLN}.

Infine la \acr{glibc} riprende da BSD la funzione \funcd{getloadavg} che
permette di ottenere il carico di processore della macchina, in questo modo è
possibile prendere decisioni su quando far partire eventuali nuovi processi.
Il suo prototipo è:

\begin{funcproto}{
\fhead{stdlib.h}
\fdecl{int getloadavg(double loadavg[], int nelem)}
\fdesc{Legge il carico medio della macchina.} 
}

{La funzione ritorna il numero di campionamenti restituiti e $-1$ se non
  riesce ad ottenere il carico medio, \var{errno} non viene modificata.}
\end{funcproto}

La funzione restituisce in ciascun elemento di \param{loadavg} il numero medio
di processi attivi sulla coda dello \textit{scheduler}, calcolato su diversi
intervalli di tempo.  Il numero di intervalli che si vogliono leggere è
specificato da \param{nelem}, dato che nel caso di Linux il carico viene
valutato solo su tre intervalli (corrispondenti a 1, 5 e 15 minuti), questo è
anche il massimo valore che può essere assegnato a questo argomento.


\subsection{La \textsl{contabilità} in stile BSD}
\label{sec:sys_bsd_accounting}

Una ultima modalità per monitorare l'uso delle risorse è, se si è compilato il
kernel con il relativo supporto,\footnote{se cioè si è abilitata l'opzione di
  compilazione \texttt{CONFIG\_BSD\_PROCESS\_ACCT}.} quella di attivare il
cosiddetto \textit{BSD accounting}, che consente di registrare su file una
serie di informazioni\footnote{contenute nella struttura \texttt{acct}
  definita nel file \texttt{include/linux/acct.h} dei sorgenti del kernel.}
riguardo alla \textsl{contabilità} delle risorse utilizzate da ogni processo
che viene terminato.

Linux consente di salvare la contabilità delle informazioni relative alle
risorse utilizzate dai processi grazie alla funzione \funcd{acct}, il cui
prototipo è:

\begin{funcproto}{
\fhead{unistd.h}
\fdecl{int acct(const char *filename)}
\fdesc{Abilita il \textit{BSD accounting}.} 
}

{La funzione ritorna $0$ in caso di successo e $-1$ per un errore, nel qual
  caso \var{errno} assumerà uno dei valori: 
  \begin{errlist}
    \item[\errcode{EACCES}] non si hanno i permessi per accedere a
      \param{pathname}.
    \item[\errcode{ENOSYS}] il kernel non supporta il \textit{BSD accounting}.
    \item[\errcode{EPERM}] il processo non ha privilegi sufficienti ad
      abilitare il \textit{BSD accounting}.
    \item[\errcode{EUSERS}] non sono disponibili nel kernel strutture per il
      file o si è finita la memoria.
    \end{errlist}
    ed inoltre \errval{EFAULT}, \errval{EIO}, \errval{ELOOP},
    \errval{ENAMETOOLONG}, \errval{ENFILE}, \errval{ENOENT}, \errval{ENOMEM},
    \errval{ENOTDIR}, \errval{EROFS} nel loro significato generico.}
\end{funcproto}

La funzione attiva il salvataggio dei dati sul file indicato dal
\textit{pathname} contenuti nella stringa puntata da \param{filename}; la
funzione richiede che il processo abbia i privilegi di amministratore (è
necessaria la \textit{capability} \const{CAP\_SYS\_PACCT}, vedi
sez.~\ref{sec:proc_capabilities}). Se si specifica il valore \val{NULL} per
\param{filename} il \textit{BSD accounting} viene invece disabilitato. Un
semplice esempio per l'uso di questa funzione è riportato nel programma
\texttt{AcctCtrl.c} dei sorgenti allegati alla guida.

Quando si attiva la contabilità, il file che si indica deve esistere; esso
verrà aperto in sola scrittura e le informazioni verranno registrate in
\textit{append} in coda al file tutte le volte che un processo termina. Le
informazioni vengono salvate in formato binario, e corrispondono al contenuto
della apposita struttura dati definita all'interno del kernel.

Il funzionamento di \func{acct} viene inoltre modificato da uno specifico
parametro di sistema, modificabile attraverso \sysctlfiled{kernel/acct} (o
tramite la corrispondente \func{sysctl}). Esso contiene tre valori interi, il
primo indica la percentuale di spazio disco libero sopra il quale viene
ripresa una registrazione che era stata sospesa per essere scesi sotto il
minimo indicato dal secondo valore (sempre in percentuale di spazio disco
libero). Infine l'ultimo valore indica la frequenza in secondi con cui deve
essere controllata detta percentuale.

% TODO: bassa priorità, trattare la lettura del file di accounting, da
% programma, vedi man 5 acct


\section{La gestione dei tempi del sistema}
\label{sec:sys_time}

In questa sezione, una volta introdotti i concetti base della gestione dei
tempi da parte del sistema, tratteremo le varie funzioni attinenti alla
gestione del tempo in un sistema unix-like, a partire da quelle per misurare i
veri tempi di sistema associati ai processi, a quelle per convertire i vari
tempi nelle differenti rappresentazioni che vengono utilizzate, a quelle della
gestione di data e ora.


\subsection{La misura del tempo in Unix}
\label{sec:sys_unix_time}

\itindbeg{calendar~time}
\itindbeg{process~time}

Tradizionalmente nei sistemi unix-like sono sempre stati previsti due tipi
distinti di tempi, caratterizzati da altrettante modalità di misura ed
espressi con diversi tipi di dati, chiamati rispettivamente \textit{calendar
  time} e \textit{process time}, secondo le seguenti definizioni:
\begin{basedescript}{\desclabelwidth{1.5cm}\desclabelstyle{\nextlinelabel}}

\item[\textit{calendar time}] detto anche \textsl{tempo di calendario}, 
  \textsl{tempo d'orologio} o \textit{tempo reale}. Si tratta di un
  tempo assoluto o di un intervallo di tempo come lo intende
  normalmente per le misure fatte con un orologio. Per esprimere
  questo tempo è stato riservato il tipo \type{time\_t}, e viene
  tradizionalmente misurato in secondi a partire dalla mezzanotte del
  primo gennaio 1970, data che viene chiamata \textit{the Epoch}.

\item[\textit{process time}] detto anche \textsl{tempo di processore} o
  \textsl{tempo di CPU}. Si tratta del tempo impiegato da un processore
  nell'esecuzione del codice di un programma all'interno di un processo. Per
  esprimere questo tempo è stato riservato il tipo \type{clock\_t}, e viene
  misurato nei cosiddetti \textit{clock tick}, tradizionalmente corrispondenti
  al numero di interruzioni del processore da parte del timer di sistema. A
  differenza del precedente indica soltanto un intervallo di durata.
\end{basedescript}

Il \textit{calendar time} viene sempre mantenuto facendo riferimento
al cosiddetto \textit{tempo universale coordinato} UTC, anche se
talvolta viene usato il cosiddetto GMT (\textit{Greenwich Mean Time})
dato che l'UTC corrisponde all'ora locale di Greenwich. Si tratta del
tempo su cui viene mantenuto il cosiddetto \textsl{orologio di
  sistema}, e viene usato per indicare i tempi dei file (quelli di
sez.~\ref{sec:file_file_times}) o le date di avvio dei processi, ed è
il tempo che viene usato dai demoni che compiono lavori amministrativi
ad orari definito, come \cmd{cron}.

Si tenga presente che questo tempo è mantenuto dal kernel e non è detto che
corrisponda al tempo misurato dall'orologio hardware presente su praticamente
tutte le piastre madri dei computer moderni (il cosiddetto \textit{hardware
  clock}), il cui valore viene gestito direttamente dall'hardware in maniera
indipendente e viene usato dal kernel soltanto all'avvio per impostare un
valore iniziale dell'orologio di sistema. La risoluzione tradizionale data dal
tipo di dato \type{time\_t} è di un secondo, ma nei sistemi più recenti sono
disponibili altri tipi di dati con precisioni maggiori.

Si tenga presente inoltre che a differenza di quanto avviene con altri sistemi
operativi,\footnote{è possibile, ancorché assolutamente sconsigliabile,
  forzare l'orologio di sistema all'ora locale per compatibilità con quei
  sistemi operativi che han fatto questa deprecabile scelta.}  l'orologio di
sistema viene mantenuto sempre in UTC e che la conversione all'ora locale del
proprio fuso orario viene effettuata dalle funzioni di libreria utilizzando le
opportune informazioni di localizzazione (specificate in
\conffiled{/etc/timezone}). In questo modo si ha l'assicurazione che l'orologio
di sistema misuri sempre un tempo monotono crescente come nella realtà, anche
in presenza di cambi di fusi orari.

\itindend{calendar~time}

Il \textit{process time} invece indica sempre una misura di un lasso di tempo
e viene usato per tenere conto dei tempi di esecuzione dei processi. Esso
viene sempre diviso in \textit{user time} e \textit{system time}, per misurare
la durata di ciascun processo il kernel infatti calcola tre tempi:
\begin{basedescript}{\desclabelwidth{2.2cm}\desclabelstyle{\nextlinelabel}}
\item[\textit{clock time}] il tempo \textsl{reale}, viene chiamato anche
  \textit{wall clock time} o \textit{elapsed time}, passato dall'avvio del
  processo. Questo tempo fa riferimento al 
  \textit{calendar time} e dice la durata effettiva dell'esecuzione del
  processo, ma chiaramente dipende dal carico del sistema e da quanti altri
  processi stanno girando nello stesso momento.
  
\item[\textit{user time}] il tempo effettivo che il processore ha impiegato
  nell'esecuzione delle istruzioni del programma in \textit{user space}. È
  anche quello riportato nella risorsa \var{ru\_utime} di \struct{rusage}
  vista in sez.~\ref{sec:sys_resource_use}.
  
\item[\textit{system time}] il tempo effettivo che il processore ha impiegato
  per eseguire codice delle \textit{system call} nel kernel per conto del
  processo.  È anche quello riportato nella risorsa \var{ru\_stime} di
  \struct{rusage} vista in sez.~\ref{sec:sys_resource_use}.
\end{basedescript}

La somma di \textit{user time} e \textit{system time} indica il
\textit{process time}, vale a dire il tempo di processore totale che il
sistema ha effettivamente utilizzato per eseguire il programma di un certo
processo. Si può ottenere un riassunto dei valori di questi tempi quando si
esegue un qualsiasi programma lanciando quest'ultimo come argomento del
comando \cmd{time}.

\itindend{process~time}
\itindbeg{clock~tick}

Come accennato il \textit{process time} viene misurato nei cosiddetti
\textit{clock tick}. Un tempo questo corrispondeva al numero di interruzioni
effettuate dal timer di sistema, oggi lo standard POSIX richiede che esso sia
espresso come multiplo della costante \constd{CLOCKS\_PER\_SEC} che deve
essere definita come 1000000, qualunque sia la risoluzione reale dell'orologio
di sistema e la frequenza delle interruzioni del timer che, come accennato in
sez.~\ref{sec:proc_hierarchy} e come vedremo a breve, è invece data dalla
costante \const{HZ}.

Il tipo di dato usato per questo tempo, \type{clock\_t}, con questa
convenzione ha una risoluzione del microsecondo. Ma non tutte le funzioni di
sistema come vedremo seguono questa convenzione, in tal caso il numero di
\textit{clock tick} al secondo può essere ricavato anche attraverso
\func{sysconf} richiedendo il valore della costante \const{\_SC\_CLK\_TCK}
(vedi sez.~\ref{sec:sys_limits}).  Il vecchio simbolo \const{CLK\_TCK}
definito in \headfile{time.h} è ormai considerato obsoleto e non deve essere
usato.

\constbeg{HZ}

In realtà tutti calcoli dei tempi vengono effettuati dal kernel per il
cosiddetto \textit{software clock}, utilizzando il \textit{timer di sistema} e
facendo i conti in base al numero delle interruzioni generate dello stesso, i
cosiddetti \itindex{jiffies} ``\textit{jiffies}''. La durata di un
``\textit{jiffy}'' è determinata dalla frequenza di interruzione del timer,
indicata in Hertz, come accennato in sez.~\ref{sec:proc_hierarchy}, dal valore
della costante \const{HZ} del kernel, definita in \file{asm/param.h}.

Fino al kernel 2.4 il valore di \const{HZ} era 100 su tutte le architetture
tranne l'alpha, per cui era 1000. Con il 2.6.0 è stato portato a 1000 su tutte
le architetture, ma dal 2.6.13 il valore è diventato una opzione di
compilazione del kernel, con un default di 250 e valori possibili di 100, 250,
1000. Dal 2.6.20 è stato aggiunto anche il valore 300 che è divisibile per le
frequenze di refresh della televisione (50 o 60 Hz). Si può pensare che questi
valori determinino anche la corrispondente durata dei \textit{clock tick}, ma
in realtà questa granularità viene calcolata in maniera indipendente usando la
costante del kernel \const{USER\_HZ}.

\constend{HZ}

Fino al kernel 2.6.21 la durata di un \textit{jiffy} costituiva la risoluzione
massima ottenibile nella misura dei tempi impiegabile in una \textit{system
  call} (ad esempio per i timeout). Con il 2.6.21 e l'introduzione degli
\textit{high-resolution timers} (HRT) è divenuto possibile ottenere, per le
funzioni di attesa ed i timer, la massima risoluzione possibile fornita
dall'hardware. Torneremo su questo in sez.~\ref{sec:sig_timer_adv}.

\itindend{clock~tick}


\subsection{La gestione del \textit{process time}}
\label{sec:sys_cpu_times}

Di norma tutte le operazioni del sistema fanno sempre riferimento al
\textit{calendar time}, l'uso del \textit{process time} è riservato a
quei casi in cui serve conoscere i tempi di esecuzione di un processo
(ad esempio per valutarne l'efficienza). In tal caso infatti fare
ricorso al \textit{calendar time} è inutile in quanto il tempo può
essere trascorso mentre un altro processo era in esecuzione o in
attesa del risultato di una operazione di I/O.

La funzione più semplice per leggere il \textit{process time} di un processo è
\funcd{clock}, che da una valutazione approssimativa del tempo di CPU
utilizzato dallo stesso; il suo prototipo è:

\begin{funcproto}{
\fhead{time.h}
\fdecl{clock\_t clock(void)}
\fdesc{Legge il valore corrente del tempo di CPU.} 
}

{La funzione ritorna il tempo di CPU in caso di successo e $-1$ se questo non
  è ottenibile o rappresentabile in un valore di tipo \type{clock\_t},
  \var{errno} non viene usata.}  
\end{funcproto}

La funzione restituisce il tempo in \textit{clock tick} ma la \acr{glibc}
segue lo standard POSIX e quindi se si vuole il tempo in secondi occorre
dividere il risultato per la costante \const{CLOCKS\_PER\_SEC}. In genere
\type{clock\_t} viene rappresentato come intero a 32 bit, il che comporta un
valore massimo corrispondente a circa 72 minuti, dopo i quali il contatore
riprenderà lo stesso valore iniziale.

La funzione è presente anche nello standard ANSI C, ma in tal caso non è
previsto che il valore ritornato indichi un intervallo di tempo ma solo un
valore assoluto, per questo se si vuole la massima portabilità anche al di
fuori di kernel unix-like, può essere opportuno chiamare la funzione
all'inizio del programma ed ottenere il valore del tempo con una differenza.

Si tenga presente inoltre che con altri kernel unix-like il valore riportato
dalla funzione può includere anche il tempo di processore usato dai processi
figli di cui si è ricevuto lo stato di terminazione con \func{wait} e
affini. Questo non vale per Linux, in cui questa informazione deve essere
ottenuta separatamente.

Come accennato in sez.~\ref{sec:sys_unix_time} il tempo di processore è la
somma di altri due tempi, l'\textit{user time} ed il \textit{system time}, che
sono quelli effettivamente mantenuti dal kernel per ciascun processo. Questi
possono essere letti separatamente attraverso la funzione \funcd{times}, il
cui prototipo è:

\begin{funcproto}{
\fhead{sys/times.h}
\fdecl{clock\_t times(struct tms *buf)}
\fdesc{Legge il valore corrente dei tempi di processore.} 
}

{La funzione ritorna un numero di \textit{clock tick} in caso di successo e
  $-1$ per un errore, nel qual caso \var{errno} potrà assumere solo il valore
  \errval{EFAULT} nel suo significato generico.}
\end{funcproto} 

La funzione restituisce i valori di \textit{process time} del processo
corrente in una struttura di tipo \struct{tms}, la cui definizione è riportata
in fig.~\ref{fig:sys_tms_struct}. La struttura prevede quattro campi; i primi
due, \var{tms\_utime} e \var{tms\_stime}, sono l'\textit{user time} ed il
\textit{system time} del processo, così come definiti in
sez.~\ref{sec:sys_unix_time}.  Gli altri due campi, \var{tms\_cutime} e
\var{tms\_cstime}, riportano la somma dell'\textit{user time} e del
\textit{system time} di tutti processi figli di cui si è ricevuto lo stato di
terminazione. 

\begin{figure}[!htb]
  \footnotesize
  \centering
  \begin{minipage}[c]{0.8\textwidth}
    \includestruct{listati/tms.h}
  \end{minipage} 
  \normalsize 
  \caption{La struttura \structd{tms} dei tempi di processore associati a un
    processo.} 
  \label{fig:sys_tms_struct}
\end{figure}


Si tenga presente che i tempi di processore dei processi figli di un processo
vengono sempre sommati al valore corrente ogni volta che se ne riceve lo stato
di terminazione, e detto valore è quello che viene a sua volta ottenuto dal
processo padre. Pertanto nei campi \var{tms\_cutime} e \var{tms\_cstime} si
sommano anche i tempi di ulteriori discendenti di cui i rispettivi genitori
abbiano ricevuto lo stato di terminazione.

Si tenga conto che l'aggiornamento di \var{tms\_cutime} e \var{tms\_cstime}
viene eseguito solo quando una chiamata a \func{wait} o \func{waitpid} è
ritornata. Per questo motivo se un processo figlio termina prima di ricevere
lo stato di terminazione di tutti i suoi figli, questi processi
``\textsl{nipoti}'' non verranno considerati nel calcolo di questi tempi e
così via per i relativi ``\textsl{discendenti}''. 

Come accennato in sez.~\ref{sec:sys_resource_use} per i kernel precedenti la
versione 2.6.9 il tempo di processore dei processi figli veniva sommato
comunque chiedendo di ignorare \signal{SIGCHLD} anche se lo standard POSIX
richiede esplicitamente che questo avvenga solo quando si riceve lo stato di
uscita con una funzione della famiglia delle \func{wait}, anche in questo caso
il comportamento è stato adeguato allo standard a partire dalla versione
2.6.9.

A differenza di quanto avviene per \func{clock} i valori restituiti nei campi
di una struttura \struct{tms} sono misurati in numero di \textit{clock tick}
effettivi e non in multipli di \const{CLOCKS\_PER\_SEC}, pertanto per ottenere
il valore effettivo del tempo in secondi occorrerà dividere per il risultato
di \code{sysconf(\_SC\_CLK\_TCK)}.

Lo stesso vale per il valore di ritorno della funzione, il cui significato fa
riferimento ad un tempo relativo ad un certo punto nel passato la cui
definizione dipende dalle diverse implementazioni, e varia anche fra diverse
versioni del kernel. Fino al kernel 2.4 si faceva infatti riferimento al
momento dell'avvio del kernel. Con il kernel 2.6 si fa riferimento a
$2^{32}/\mathtt{HZ}-300$ secondi prima dell'avvio. 

Considerato che il numero dei \textit{clock tick} per un kernel che è attivo
da molto tempo può eccedere le dimensioni per il tipo \type{clock\_t} il
comportamento più opportuno per i programmi è di ignorare comunque il valore
di ritorno della funzione e ricorrere alle funzioni per il tempo di calendario
del prossimo paragrafo qualora si voglia calcolare il tempo effettivamente
trascorso dall'inizio del programma.

Infine si tenga presente che per dei limiti nelle convenzioni per il ritorno
dei valori delle \textit{system call} su alcune architetture hardware (ed in
particolare la \texttt{i386} dei PC a 32 bit) nel kernel della serie 2.6 il
valore di ritorno della funzione può risultare erroneamente uguale a $-1$,
indicando un errore, nei primi secondi dopo il boot (per la precisione nei
primi 41 secondi) e se il valore del contatore eccede le dimensione del tipo
\type{clock\_t}.


\subsection{Le funzioni per il \textit{calendar time}}
\label{sec:sys_time_base}

\itindbeg{calendar~time}

Come anticipato in sez.~\ref{sec:sys_unix_time} il \textit{calendar time}
viene espresso normalmente con una variabile di tipo \type{time\_t}, che
usualmente corrisponde ad un tipo elementare; in Linux è definito come
\ctyp{long int}, che di norma corrisponde a 32 bit. Il valore corrente del
\textit{calendar time}, che indicheremo come \textsl{tempo di sistema}, può
essere ottenuto con la funzione \funcd{time} che lo restituisce nel suddetto
formato, il suo prototipo è:

\begin{funcproto}{
\fhead{time.h}
\fdecl{time\_t time(time\_t *t)}
\fdesc{Legge il valore corrente del \textit{calendar time}.} 
}

{La funzione ritorna il valore del \textit{calendar time} in caso di successo
  e $-1$ per un errore, nel qual caso \var{errno} potrà assumere solo il
  valore \errval{EFAULT} nel suo significato generico.}
\end{funcproto}

L'argomento \param{t}, se non nullo, deve essere l'indirizzo di una variabile
su cui duplicare il valore di ritorno.

Analoga a \func{time} è la funzione \funcd{stime} che serve per effettuare
l'operazione inversa, e cioè per impostare il tempo di sistema qualora questo
sia necessario; il suo prototipo è:

\begin{funcproto}{
\fhead{time.h}
\fdecl{int stime(time\_t *t)}
\fdesc{Imposta il valore corrente del \textit{calendar time}.} 
}

{La funzione ritorna $0$ in caso di successo e $-1$ per un errore, nel qual
  caso \var{errno} assumerà uno dei valori: 
  \begin{errlist}
  \item[\errcode{EPERM}] non si hanno i permessi di amministrazione.
  \end{errlist}
  ed inoltre \errval{EFAULT} nel suo significato generico.}
\end{funcproto}


Dato che modificare l'ora ha un impatto su tutto il sistema il cambiamento
dell'orologio è una operazione privilegiata e questa funzione può essere usata
solo da un processo con i privilegi di amministratore (per la precisione la
\textit{capability} \const{CAP\_SYS\_TIME}), altrimenti la chiamata fallirà
con un errore di \errcode{EPERM}.

Data la scarsa precisione nell'uso di \type{time\_t}, che ha una risoluzione
massima di un secondo, quando si devono effettuare operazioni sui tempi di
norma l'uso delle due funzioni precedenti è sconsigliato, ed esse sono di
solito sostituite da \funcd{gettimeofday} e \funcd{settimeofday},\footnote{le
  due funzioni \func{time} e \func{stime} sono più antiche e derivano da SVr4,
  \func{gettimeofday} e \func{settimeofday} sono state introdotte da BSD, ed
  in BSD4.3 sono indicate come sostitute delle precedenti, \func{gettimeofday}
  viene descritta anche in POSIX.1-2001.} i cui prototipi sono:

\begin{funcproto}{
\fhead{sys/time.h}
\fhead{time.h}
\fdecl{int gettimeofday(struct timeval *tv, struct timezone *tz)} 
\fdesc{Legge il tempo corrente del sistema.}
\fdecl{int settimeofday(const struct timeval *tv, const struct timezone *tz)} 
\fdesc{Imposta il tempo di sistema.} 
}

{La funzioni ritornano $0$ in caso di successo e $-1$ per un errore, nel qual
  caso \var{errno} assumerà i valori \errval{EINVAL}, \errval{EFAULT} e per
  \func{settimeofday} anche \errval{EPERM}, nel loro significato generico.}
\end{funcproto}


Si noti come queste funzioni utilizzino per indicare il tempo una struttura di
tipo \struct{timeval}, la cui definizione si è già vista in
fig.~\ref{fig:sys_timeval_struct}, questa infatti permette una espressione
alternativa dei valori del \textit{calendar time}, con una precisione,
rispetto a \type{time\_t}, fino al microsecondo, ma la precisione è solo
teorica, e la precisione reale della misura del tempo dell'orologio di sistema
non dipende dall'uso di queste strutture.

Come nel caso di \func{stime} anche \func{settimeofday} può essere utilizzata
solo da un processo coi privilegi di amministratore e più precisamente con la
capacità \const{CAP\_SYS\_TIME}. Si tratta comunque di una condizione generale
che continua a valere per qualunque funzione che vada a modificare l'orologio
di sistema, comprese tutte quelle che tratteremo in seguito.

\itindbeg{timezone}

Il secondo argomento di entrambe le funzioni è una struttura
\struct{timezone}, che storicamente veniva utilizzata per specificare appunto
la cosiddetta \textit{timezone}, cioè l'insieme del fuso orario e delle
convenzioni per l'ora legale che permettevano il passaggio dal tempo
universale all'ora locale. Questo argomento oggi è obsoleto ed in Linux non è
mai stato utilizzato; esso non è supportato né dalla vecchia \textsl{libc5},
né dalla \textsl{glibc}: pertanto quando si chiama questa funzione deve essere
sempre impostato a \val{NULL}.

\itindbeg{timezone}

Modificare l'orologio di sistema con queste funzioni è comunque problematico,
in quanto esse effettuano un cambiamento immediato. Questo può creare dei
buchi o delle ripetizioni nello scorrere dell'orologio di sistema, con
conseguenze indesiderate.  Ad esempio se si porta avanti l'orologio si possono
perdere delle esecuzioni di \cmd{cron} programmate nell'intervallo che si è
saltato. Oppure se si porta indietro l'orologio si possono eseguire due volte
delle operazioni previste nell'intervallo di tempo che viene ripetuto. 

Per questo motivo la modalità più corretta per impostare l'ora è quella di
usare la funzione \funcd{adjtime}, il cui prototipo è:

\begin{funcproto}{
\fhead{sys/time.h}
\fdecl{int adjtime(const struct timeval *delta, struct timeval *olddelta)}
\fdesc{Aggiusta l'orologio di sistema.} 
}

{La funzione ritorna $0$ in caso di successo e $-1$ per un errore, nel qual
  caso \var{errno} assumerà uno dei valori: 
  \begin{errlist}
  \item[\errcode{EINVAL}] il valore di \param{delta} eccede il massimo
    consentito.
  \item[\errcode{EPERM}] il processo non i privilegi di amministratore.
  \end{errlist}
}  
\end{funcproto}


Questa funzione permette di avere un aggiustamento graduale del tempo di
sistema in modo che esso sia sempre crescente in maniera monotona. Il valore
indicato nella struttura \struct{timeval} puntata da \param{delta} esprime il
valore di cui si vuole spostare l'orologio. Se è positivo l'orologio sarà
accelerato per un certo tempo in modo da guadagnare il tempo richiesto,
altrimenti sarà rallentato. 

La funzione è intesa per piccoli spostamenti del tempo di sistema, ed esistono
pertanto dei limiti massimi per i valori che si possono specificare
per \param{delta}. La \acr{glibc} impone un intervallo compreso fra
\code{INT\_MIN/1000000 + 2} e \code{INT\_MAX/1000000 - 2}, corrispondente, su
una architettura PC ordinaria a 32 bit, ad un valore compreso fra $-2145$ e
$2145$ secondi.

Inoltre se si invoca la funzione prima che una precedente richiesta di
aggiustamento sia stata completata, specificando un altro valore, il
precedente aggiustamento viene interrotto, ma la parte dello stesso che è già
stata completata non viene rimossa. Però è possibile in questo caso farsi
restituire nella struttura puntata da \param{olddelta} il tempo restante della
precedente richiesta. Fino al kernel 2.6.26 ed alla \acr{glibc} 2.8 questo
però era possibile soltanto specificando un diverso aggiustamento
per \param{delta}, il bug è stato corretto a partire dalle versioni citate e
si può ottenere l'informazione relativa alla frazione di aggiustamento
mancante usando il valore \val{NULL} per \param{delta}.

Linux poi prevede una specifica funzione di sistema che consente un
aggiustamento molto più dettagliato del tempo, permettendo ad esempio anche di
regolare anche la velocità e le derive dell'orologio di sistema.  La funzione
è \funcd{adjtimex} ed il suo prototipo è:

\begin{funcproto}{
\fhead{sys/timex.h}
\fdecl{int adjtimex(struct timex *buf)} 
\fdesc{Regola l'orologio di sistema.} 
}

{La funzione ritorna lo stato dell'orologio (un valore $\ge 0$) in caso di
  successo e $-1$ per un errore, nel qual caso \var{errno} assumerà uno dei
  valori:
  \begin{errlist}
  \item[\errcode{EINVAL}] si sono indicati valori fuori dall'intervallo
    consentito per qualcuno dei campi di \param{buf}.
  \item[\errcode{EPERM}] si è richiesta una modifica dei parametri ed il
    processo non ha i privilegi di amministratore.
  \end{errlist}
  ed inoltre \errval{EFAULT} nel suo significato generico.}
\end{funcproto}

In caso di successo la funzione restituisce un valore numerico non negativo
che indica lo stato dell'orologio, che può essere controllato con i valori
delle costanti elencate in tab.~\ref{tab:adjtimex_return}.

\begin{table}[!htb]
  \footnotesize
  \centering
  \begin{tabular}[c]{|l|c|l|}
    \hline
    \textbf{Nome} & \textbf{Valore} & \textbf{Significato}\\
    \hline
    \hline
    \constd{TIME\_OK}   & 0 & Orologio sincronizzato.\\ 
    \constd{TIME\_INS}  & 1 & Inserimento di un \textit{leap second}.\\ 
    \constd{TIME\_DEL}  & 2 & Cancellazione di un \textit{leap second}.\\ 
    \constd{TIME\_OOP}  & 3 & \textit{leap second} in corso.\\ 
    \constd{TIME\_WAIT} & 4 & \textit{leap second} avvenuto.\\ 
    \constd{TIME\_BAD}  & 5 & Orologio non sincronizzato.\\ 
    \hline
  \end{tabular}
  \caption{Possibili valori ritornati da \func{adjtimex} in caso di successo.} 
  \label{tab:adjtimex_return}
\end{table}

La funzione richiede come argomento il puntatore ad una struttura di tipo
\struct{timex}, la cui definizione, effettuata in \headfiled{sys/timex.h}, è
riportata in fig.~\ref{fig:sys_timex_struct} per i campi che interessano la
possibilità di essere modificati documentati anche nella pagina di manuale. In
realtà la struttura è stata estesa con ulteriori campi, i cui valori sono
utilizzabili solo in lettura, la cui definizione si può trovare direttamente

\begin{figure}[!htb]
  \footnotesize \centering
  \begin{minipage}[c]{\textwidth}
    \includestruct{listati/timex.h}
  \end{minipage} 
  \normalsize 
  \caption{La struttura \structd{timex} per il controllo dell'orologio di
    sistema.} 
  \label{fig:sys_timex_struct}
\end{figure}

L'azione della funzione dipende dal valore del campo \var{mode}
di \param{buf}, che specifica quale parametro dell'orologio di sistema,
specificato nel corrispondente campo di \struct{timex}, deve essere
impostato. Un valore nullo serve per leggere i parametri correnti, i valori
diversi da zero devono essere specificati come OR binario delle costanti
riportate in tab.~\ref{tab:sys_timex_mode}.

\begin{table}[!htb]
  \footnotesize
  \centering
  \begin{tabular}[c]{|l|c|p{8cm}|}
    \hline
    \textbf{Nome} & \textbf{Valore} & \textbf{Significato}\\
    \hline
    \hline
    \constd{ADJ\_OFFSET}        & 0x0001 & Imposta la differenza fra il tempo
                                           reale e l'orologio di sistema: 
                                           deve essere indicata in microsecondi
                                           nel campo \var{offset} di
                                           \struct{timex}.\\ 
    \constd{ADJ\_FREQUENCY}     & 0x0002 & Imposta la differenza in frequenza
                                           fra il tempo reale e l'orologio di
                                           sistema: deve essere indicata
                                           in parti per milione nel campo
                                           \var{frequency} di \struct{timex}.\\
    \constd{ADJ\_MAXERROR}      & 0x0004 & Imposta il valore massimo 
                                           dell'errore sul tempo, espresso in
                                           microsecondi nel campo
                                           \var{maxerror} di \struct{timex}.\\ 
    \constd{ADJ\_ESTERROR}      & 0x0008 & Imposta la stima dell'errore
                                           sul tempo, espresso in microsecondi 
                                           nel campo \var{esterror} di
                                           \struct{timex}.\\
    \constd{ADJ\_STATUS}        & 0x0010 & Imposta alcuni valori di stato
                                           interni usati dal 
                                           sistema nella gestione
                                           dell'orologio specificati nel campo
                                           \var{status} di \struct{timex}.\\ 
    \constd{ADJ\_TIMECONST}     & 0x0020 & Imposta la larghezza di banda del 
                                           PLL implementato dal kernel,
                                           specificato nel campo
                                           \var{constant} di \struct{timex}.\\ 
    \constd{ADJ\_TICK}          & 0x4000 & Imposta il valore dei \textit{tick}
                                           del timer in
                                           microsecondi, espresso nel campo
                                           \var{tick} di \struct{timex}.\\  
    \constd{ADJ\_OFFSET\_SINGLESHOT}&0x8001&Chiede uno spostamento una tantum 
                                           dell'orologio secondo il valore del
                                           campo \var{offset} simulando il
                                           comportamento di \func{adjtime}.\\ 
    \hline
  \end{tabular}
  \caption{Costanti per l'assegnazione del valore del campo \var{mode} della
    struttura \struct{timex}.} 
  \label{tab:sys_timex_mode}
\end{table}

La funzione utilizza il meccanismo di David L. Mills, descritto
nell'\href{http://www.ietf.org/rfc/rfc1305.txt}{RFC~1305}, che è alla base del
protocollo NTP. La funzione è specifica di Linux e non deve essere usata se la
portabilità è un requisito, la \acr{glibc} provvede anche un suo omonimo
\func{ntp\_adjtime}.  La trattazione completa di questa funzione necessita di
una lettura approfondita del meccanismo descritto nell'RFC~1305, ci limitiamo
a descrivere in tab.~\ref{tab:sys_timex_mode} i principali valori utilizzabili
per il campo \var{mode}, un elenco più dettagliato del significato dei vari
campi della struttura \struct{timex} può essere ritrovato in \cite{GlibcMan}.

Il valore delle costanti per \var{mode} può essere anche espresso, secondo la
sintassi specificata per la forma equivalente di questa funzione definita come
\func{ntp\_adjtime}, utilizzando il prefisso \code{MOD} al posto di
\code{ADJ}.

Si tenga presente infine che con l'introduzione a partire dal kernel 2.6.21
degli \textit{high-resolution timer} ed il supporto per i cosiddetti POSIX
\textit{real-time clock}, si può ottenere il \textit{calendar time}
direttamente da questi, come vedremo in sez.~\ref{sec:sig_timer_adv}, con la
massima risoluzione possibile per l'hardware della macchina.



\subsection{La gestione delle date.}
\label{sec:sys_date}

\itindbeg{broken-down~time}

Le funzioni viste al paragrafo precedente sono molto utili per trattare le
operazioni elementari sui tempi, però le rappresentazioni del tempo ivi
illustrate, se han senso per specificare un intervallo, non sono molto
intuitive quando si deve esprimere un'ora o una data.  Per questo motivo è
stata introdotta una ulteriore rappresentazione, detta \textit{broken-down
  time}, che permette appunto di \textsl{suddividere} il \textit{calendar
  time} usuale in ore, minuti, secondi, ecc. e  viene usata tenendo conto
anche dell'eventuale utilizzo di un fuso orario.

\begin{figure}[!htb]
  \footnotesize \centering
  \begin{minipage}[c]{.8\textwidth}
    \includestruct{listati/tm.h}
  \end{minipage} 
  \normalsize 
  \caption{La struttura \structd{tm} per una rappresentazione del tempo in
    termini di ora, minuti, secondi, ecc.}
  \label{fig:sys_tm_struct}
\end{figure}

Questo viene effettuato attraverso una opportuna struttura \struct{tm}, la cui
definizione è riportata in fig.~\ref{fig:sys_tm_struct}, ed è in genere questa
struttura che si utilizza quando si deve specificare un tempo a partire dai
dati naturali (ora e data), dato che essa consente anche di tenere conto della
gestione del fuso orario e dell'ora legale. In particolare gli ultimi due
campi, \var{tm\_gmtoff} e \var{tm\_zone}, sono estensioni previste da BSD e
supportate dalla \acr{glibc} quando è definita la macro \macro{\_BSD\_SOURCE}.

Ciascuno dei campi di \struct{tm} ha dei precisi intervalli di valori
possibili, con convenzioni purtroppo non troppo coerenti. Ad esempio
\var{tm\_sec} che indica i secondi deve essere nell'intervallo da 0 a 59, ma è
possibile avere anche il valore 60 per un cosiddetto \textit{leap second} (o
\textsl{secondo intercalare}), cioè uno di quei secondi aggiunti al calcolo
dell'orologio per effettuare gli aggiustamenti del calendario per tenere conto
del disallineamento con il tempo solare.\footnote{per dettagli si consulti
  \url{http://it.wikipedia.org/wiki/Leap_second}.}

I campi \var{tm\_min} e\var{tm\_hour} che indicano rispettivamente minuti ed
ore hanno valori compresi rispettivamente fra 0 e 59 e fra 0 e 23. Il campo
\var{tm\_mday} che indica il giorno del mese prevede invece un valore compreso
fra 1 e 31, ma la \acr{glibc} supporta pure il valore 0 come indicazione
dell'ultimo giorno del mese precedente. Il campo \var{tm\_mon} indica il mese
dell'anno a partire da gennaio con valori compresi fra 0 e 11.

I campi \var{tm\_wday} e \var{tm\_yday} indicano invece rispettivamente il
giorno della settimana, a partire dalla Domenica, ed il giorno dell'anno, a
partire del primo gennaio, ed hanno rispettivamente valori compresi fra 0 e 6
e fra 0 e 365. L'anno espresso da \var{tm\_year} viene contato come numero di
anni a partire dal 1900. Infine \var{tm\_isdst} è un valore che indica se per
gli altri campi si intende come attiva l'ora legale ed influenza il
comportamento di \func{mktime}.


Le funzioni per la gestione del \textit{broken-down time} sono varie e vanno
da quelle usate per convertire gli altri formati in questo, usando o meno
l'ora locale o il tempo universale, a quelle per trasformare il valore di un
tempo in una stringa contenente data ed ora. Le prime due funzioni,
\funcd{asctime} e \funcd{ctime} servono per poter stampare in forma leggibile
un tempo, i loro prototipi sono:

\begin{funcproto}{
\fhead{time.h}
\fdecl{char * asctime(const struct tm *tm)}
\fdesc{Converte un \textit{broken-down time} in una stringa.} 
\fdecl{char * ctime(const time\_t *timep)}
\fdesc{Converte un \textit{calendar time} in una stringa.} 
}

{Le funzioni ritornano un puntatore al risultato in caso di successo e
  \val{NULL} per un errore, \var{errno} non viene modificata.}
\end{funcproto}

Le funzioni prendono rispettivamente come argomenti i puntatori ad una
struttura \struct{tm} contenente un \textit{broken-down time} o ad una
variabile di tipo \type{time\_t} che esprime il \textit{calendar time},
restituendo il puntatore ad una stringa che esprime la data, usando le
abbreviazioni standard di giorni e mesi in inglese, nella forma:
\begin{Example}
Sun Apr 29 19:47:44 2012\n"
\end{Example}

Nel caso di \func{ctime} la funzione tiene conto della eventuale impostazione
di una \textit{timezone} e effettua una chiamata preventiva a \func{tzset}
(che vedremo a breve), in modo che la data espressa tenga conto del fuso
orario. In realtà \func{ctime} è banalmente definita in termini di
\func{asctime} come \code{asctime(localtime(t)}.

Dato che l'uso di una stringa statica rende le funzioni non rientranti
POSIX.1c e SUSv2 prevedono due sostitute rientranti, il cui nome è al solito
ottenuto aggiungendo un \code{\_r}, che prendono un secondo argomento
\code{char *buf}, in cui l'utente deve specificare il buffer su cui la stringa
deve essere copiata (deve essere di almeno 26 caratteri).

Per la conversione fra \textit{broken-down time} e \textit{calendar time} sono
invece disponibili altre tre funzioni, \funcd{gmtime}, \funcd{localtime} e
\funcd{mktime} i cui prototipi sono:

\begin{funcproto}{
\fdecl{struct tm * gmtime(const time\_t *timep)}
\fdesc{Converte un \textit{calendar time} in un \textit{broken-down time} in
  UTC.}  
\fdecl{struct tm * localtime(const time\_t *timep)}
\fdesc{Converte un \textit{calendar time} in un \textit{broken-down time}
  nell'ora locale.} 
\fdecl{time\_t mktime(struct tm *tm)} 
\fdesc{Converte un \textit{broken-down time} in un \textit{calendar time}.} 

}

{Le funzioni ritornano un puntatore al risultato in caso di successo e
  \val{NULL} per un errore, tranne che \func{mktime} che restituisce
  direttamente il valore o $-1$ in caso di errore, \var{errno} non viene
  modificata.}
\end{funcproto}

Le le prime funzioni, \func{gmtime}, \func{localtime} servono per convertire
il tempo in \textit{calendar time} specificato da un argomento di tipo
\type{time\_t} restituendo un \textit{broken-down time} con il puntatore ad
una struttura \struct{tm}. La prima effettua la conversione senza tenere conto
del fuso orario, esprimendo la data in tempo coordinato universale (UTC), cioè
l'ora di Greenwich, mentre \func{localtime} usa l'ora locale e per questo
effettua una chiamata preventiva a \func{tzset}.

Anche in questo caso le due funzioni restituiscono l'indirizzo di una
struttura allocata staticamente, per questo sono state definite anche altre
due versioni rientranti (con la solita estensione \code{\_r}), che prevedono
un secondo argomento \code{struct tm *result}, fornito dal chiamante, che deve
preallocare la struttura su cui sarà restituita la conversione. La versione
rientrante di \func{localtime} però non effettua la chiamata preventiva a
\func{tzset} che deve essere eseguita a cura dell'utente.

Infine \func{mktime} esegue la conversione di un \textit{broken-down time} a
partire da una struttura \struct{tm} restituendo direttamente un valore di
tipo \type{time\_t} con il \textit{calendar time}. La funzione ignora i campi
\var{tm\_wday} e \var{tm\_yday} e per gli altri campi normalizza eventuali
valori fuori degli intervalli specificati in precedenza: se cioè si indica un
12 per \var{tm\_mon} si prenderà il gennaio dell'anno successivo. Inoltre la
funzione tiene conto del valore di \var{tm\_isdst} per effettuare le
correzioni relative al fuso orario: un valore positivo indica che deve essere
tenuta in conto l'ora legale, un valore nullo che non deve essere applicata
nessuna correzione, un valore negativo che si deve far ricorso alle
informazioni relative al proprio fuso orario per determinare lo stato dell'ora
legale.  

La funzione inoltre modifica i valori della struttura \struct{tm} in forma di
\textit{value result argument}, normalizzando i valori dei vari campi,
impostando i valori risultanti per \var{tm\_wday} e \var{tm\_yday} e
assegnando a \var{tm\_isdst} il valore (positivo o nullo) corrispondente allo
stato dell'ora legale. La funzione inoltre provvede ad impostare il valore
della variabile globale \var{tzname}.

\itindend{calendar~time}

\begin{figure}[!htb]
  \footnotesize
  \centering
  \begin{minipage}[c]{.75\textwidth}
    \includestruct{listati/time_zone_var.c}
  \end{minipage} 
  \normalsize 
  \caption{Le variabili globali usate per la gestione delle
    \textit{timezone}.}
  \label{fig:sys_tzname}
\end{figure}

Come accennato l'uso del \textit{broken-down time} permette di tenere conto
anche della differenza fra tempo universale e ora locale, compresa l'eventuale
ora legale.  Questo viene fatto dalle funzioni di conversione grazie alle
informazioni riguardo la propria \textit{timezone} mantenute nelle tre
variabili globali mostrate in fig.~\ref{fig:sys_tzname}, cui si può accedere
direttamente includendo \headfile{time.h}. Come illustrato queste variabili
vengono impostate internamente da alcune delle precedenti funzioni di
conversione, ma lo si può fare esplicitamente chiamando direttamente la
funzione \funcd{tzset}, il cui prototipo è:

\begin{funcproto}{
\fhead{time.h}
\fdecl{void tzset(void)} 
\fdesc{Imposta le variabili globali della \textit{timezone}.} 
}

{La funzione non ritorna niente e non dà errori.}  
\end{funcproto}

La funzione inizializza le variabili di fig.~\ref{fig:sys_tzname} a partire
dal valore della variabile di ambiente \envvar{TZ}, se quest'ultima non è
definita verrà usato il file \conffiled{/etc/localtime}. La variabile
\var{tzname} contiene due stringhe, che indicano i due nomi standard della
\textit{timezone} corrente. La prima è il nome per l'ora solare, la seconda
per l'ora legale. Anche se in fig.~\ref{fig:sys_tzname} sono indicate come
\code{char *} non è il caso di modificare queste stringhe. La variabile
\var{timezone} indica la differenza di fuso orario in secondi, mentre
\var{daylight} indica se è attiva o meno l'ora legale.

Benché la funzione \func{asctime} fornisca la modalità più immediata per
stampare un tempo o una data, la flessibilità non fa parte delle sue
caratteristiche; quando si vuole poter stampare solo una parte (l'ora, o il
giorno) di un tempo si può ricorrere alla più sofisticata \funcd{strftime},
il cui prototipo è:

\begin{funcproto}{
\fhead{time.h}
\fdecl{size\_t strftime(char *s, size\_t max, const char *format, 
  const struct tm *tm)}
\fdesc{Crea una stringa con una data secondo il formato indicato.} 
}

{La funzione ritorna il numero di caratteri inseriti nella stringa \param{s}
  oppure $0$, \var{errno} non viene modificata.}  
\end{funcproto}


La funzione converte il \textit{broken-down time} indicato nella struttura
puntata dall'argomento \param{tm} in una stringa di testo da salvare
all'indirizzo puntato dall'argomento \param{s}, purché essa sia di dimensione
inferiore al massimo indicato dall'argomento \param{max}. Il numero di
caratteri generati dalla funzione viene restituito come valore di ritorno,
senza tener però conto del terminatore finale, che invece viene considerato
nel computo della dimensione. Se quest'ultima è eccessiva viene restituito 0 e
lo stato di \param{s} è indefinito.

\begin{table}[htb]
  \footnotesize
  \centering
  \begin{tabular}[c]{|c|l|p{6cm}|}
    \hline
    \textbf{Modificatore} & \textbf{Esempio} & \textbf{Significato}\\
    \hline
    \hline
    \var{\%a}&\texttt{Wed}        & Nome del giorno, abbreviato.\\ 
    \var{\%A}&\texttt{Wednesday}  & Nome del giorno, completo.\\ 
    \var{\%b}&\texttt{Apr}        & Nome del mese, abbreviato.\\ 
    \var{\%B}&\texttt{April}      & Nome del mese, completo.\\ 
    \var{\%c}&\texttt{Wed Apr 24 18:40:50 2002}& Data e ora.\\ 
    \var{\%d}&\texttt{24}         & Giorno del mese.\\ 
    \var{\%H}&\texttt{18}         & Ora del giorno, da 0 a 24.\\ 
    \var{\%I}&\texttt{06}         & Ora del giorno, da 0 a 12.\\ 
    \var{\%j}&\texttt{114}        & Giorno dell'anno.\\ 
    \var{\%m}&\texttt{04}         & Mese dell'anno.\\ 
    \var{\%M}&\texttt{40}         & Minuto.\\ 
    \var{\%p}&\texttt{PM}         & AM/PM.\\ 
    \var{\%S}&\texttt{50}         & Secondo.\\ 
    \var{\%U}&\texttt{16}         & Settimana dell'anno (partendo dalla
                                    domenica).\\ 
    \var{\%w}&\texttt{3}          & Giorno della settimana.\\ 
    \var{\%W}&\texttt{16}         & Settimana dell'anno (partendo dal
                                    lunedì).\\ 
    \var{\%x}&\texttt{04/24/02}   & La data.\\ 
    \var{\%X}&\texttt{18:40:50}   & L'ora.\\ 
    \var{\%y}&\texttt{02}         & Anno nel secolo.\\ 
    \var{\%Y}&\texttt{2002}       & Anno.\\ 
    \var{\%Z}&\texttt{CEST}       & Nome della \textit{timezone}.\\  
    \var{\%\%}&\texttt{\%}        & Il carattere \%.\\ 
    \hline
  \end{tabular}
  \caption{Valori previsti dallo standard ANSI C per modificatore della
    stringa di formato di \func{strftime}.}  
  \label{tab:sys_strftime_format}
\end{table}

Il risultato della funzione è controllato dalla stringa di formato
\param{format}, tutti i caratteri restano invariati eccetto \texttt{\%} che
viene utilizzato come modificatore. Alcuni dei possibili valori che esso può
assumere sono riportati in tab.~\ref{tab:sys_strftime_format}.\footnote{per la
  precisione si sono riportati definiti dallo standard ANSI C, che sono anche
  quelli ripresi in POSIX.1; la \acr{glibc} fornisce anche le estensioni
  introdotte da POSIX.2 per il comando \cmd{date}, i valori introdotti da
  SVID3 e ulteriori estensioni GNU; l'elenco completo dei possibili valori è
  riportato nella pagina di manuale della funzione.} La funzione tiene conto
anche delle eventuali impostazioni di localizzazione per stampare i vari nomi
in maniera adeguata alla lingua scelta, e con le convenzioni nazionali per i
formati di data ed ora.

Infine per effettuare l'operazione di conversione inversa, da una stringa ad
un \textit{broken-down time}, si può utilizzare la funzione \funcd{strptime},
il cui prototipo è:

\begin{funcproto}{
\fhead{time.h}
\fdecl{char *strptime(const char *s, const char *format, struct tm *tm)}
\fdesc{Converte una stringa con in un \textit{broken-down time} secondo un
  formato.} 
}

{La funzione ritorna il puntatore al primo carattere non processato della
  stringa o al terminatore finale qualora questa sia processata interamente,
  \var{errno} non viene modificata.}
\end{funcproto}

La funzione processa la stringa puntata dall'argomento \param{s} da sinistra a
destra, utilizzando il formato contenuto nella stringa puntata
dall'argomento \param{format}, avvalorando volta volta i corrispondenti campi
della struttura puntata dall'argomento \param{tm}. La scansione si interrompe
immediatamente in caso di mancata corrispondenza a quanto indicato nella
stringa di formato, che usa una sintassi analoga a quella già vista per
\func{strftime}. La funzione supporta i modificatori di
tab.~\ref{tab:sys_strftime_format} più altre estensioni, ma per i dettagli a
questo riguardo si rimanda alla lettura della pagina di manuale.

Si tenga presente comunque che anche in caso di scansione completamente
riuscita la funzione sovrascrive soltanto i campi di \param{tm} indicati dal
formato, la struttura originaria infatti non viene inizializzati e gli altri
campi restano ai valori che avevano in precedenza.


\itindend{broken-down~time}

\section{La gestione degli errori}
\label{sec:sys_errors}

In questa sezione esamineremo le caratteristiche principali della gestione
degli errori in un sistema unix-like. Infatti a parte il caso particolare di
alcuni segnali (che tratteremo in cap.~\ref{cha:signals}) in un sistema
unix-like il kernel non avvisa mai direttamente un processo dell'occorrenza di
un errore nell'esecuzione di una funzione, ma di norma questo viene riportato
semplicemente usando un opportuno valore di ritorno della funzione invocata.
Inoltre il sistema di classificazione degli errori è stato progettato
sull'architettura a processi, e presenta una serie di problemi nel caso lo si
debba usare con i \textit{thread}.


\subsection{La variabile \var{errno}}
\label{sec:sys_errno}

Quasi tutte le funzioni delle librerie del C sono in grado di individuare e
riportare condizioni di errore, ed è una norma fondamentale di buona
programmazione controllare \textsl{sempre} che le funzioni chiamate si siano
concluse correttamente.

In genere le funzioni di libreria usano un valore speciale per indicare che
c'è stato un errore. Di solito questo valore, a seconda della funzione, è $-1$
o un puntatore nullo o la costante \val{EOF}; ma questo valore segnala solo
che c'è stato un errore, e non il tipo di errore.

Per riportare il tipo di errore il sistema usa la variabile globale
\var{errno}, definita nell'header \headfile{errno.h}.  Come accennato l'uso di
una variabile globale può comportare problemi nel caso dei \textit{thread}, ma
lo standard ISO C consente anche di definire \var{errno} come un cosiddetto
``\textit{modifiable lvalue}'', cosa che consente di usare anche una macro, e
questo è infatti il metodo usato da Linux per renderla locale ai singoli
\textit{thread}.

La variabile è in genere definita come \dirct{volatile} dato che può essere
cambiata in modo asincrono da un segnale, per un esempio si veda
sez.~\ref{sec:sig_sigchld} ricordando quanto trattato in
sez.~\ref{sec:proc_race_cond}). Dato che un gestore di segnale scritto bene si
cura di salvare e ripristinare il valore della variabile all'uscita, nella
programmazione normale, quando si può fare l'assunzione che i gestori di
segnali siano ben scritti, di questo non è necessario preoccuparsi.

I valori che può assumere \var{errno} sono riportati in app.~\ref{cha:errors},
nell'header \headfile{errno.h} sono anche definiti i nomi simbolici per le
costanti numeriche che identificano i vari errori che abbiamo citato fin
dall'inizio nelle descrizioni delle funzioni.  Essi iniziano tutti per \val{E}
e si possono considerare come nomi riservati, per questo abbiamo sempre fatto
riferimento a questi nomi, e lo faremo più avanti quando descriveremo i
possibili errori restituiti dalle funzioni. Il programma di esempio
\cmd{errcode} stampa il codice relativo ad un valore numerico con l'opzione
\cmd{-l}.

Il valore di \var{errno} viene sempre impostato a zero all'avvio di un
programma, e la gran parte delle funzioni di libreria impostano \var{errno} ad
un valore diverso da zero in caso di errore. Il valore è invece indefinito in
caso di successo, perché anche se una funzione di libreria ha successo,
potrebbe averne chiamate altre al suo interno che potrebbero essere fallite
anche senza compromettere il risultato finale, modificando però \var{errno}.

Pertanto un valore non nullo di \var{errno} non è sintomo di errore (potrebbe
essere il risultato di un errore precedente) e non lo si può usare per
determinare quando o se una chiamata a funzione è fallita.  La procedura
corretta da seguire per identificare un errore è sempre quella di controllare
\var{errno} immediatamente dopo aver verificato il fallimento della funzione
attraverso il suo codice di ritorno.


\subsection{Le funzioni \func{strerror} e \func{perror}}
\label{sec:sys_strerror}

Benché gli errori siano identificati univocamente dal valore numerico di
\var{errno} le librerie provvedono alcune funzioni e variabili utili per
riportare in opportuni messaggi le condizioni di errore verificatesi.  La
prima funzione che si può usare per ricavare i messaggi di errore è
\funcd{strerror}, il cui prototipo è:

\begin{funcproto}{
\fhead{string.h}
\fdecl{char *strerror(int errnum)} 
\fdesc{Restituisce una stringa con un messaggio di errore.} 
}

{La funzione ritorna il puntatore alla stringa con il messaggio di errore,
  \var{errno} non viene modificato.}
\end{funcproto}

La funzione ritorna il puntatore alla stringa contenente il messaggio di
errore corrispondente al valore di \param{errnum}, se questo non è un valore
valido verrà comunque restituita una stringa valida contenente un messaggio
che dice che l'errore è sconosciuto nella forma. La versione della \acr{glibc}
non modifica il valore di \var{errno} in caso di errore, ma questo non è detto
valga per altri sistemi in quanto lo standard POSIX.1-2001 permette che ciò
avvenga. Non si faccia affidamento su questa caratteristica se si vogliono
scrivere programmi portabili.

In generale \func{strerror} viene usata passando direttamente \var{errno} come
argomento, ed il valore di quest'ultima non verrà modificato. La funzione
inoltre tiene conto del valore della variabile di ambiente
\envvar{LC\_MESSAGES} per usare le appropriate traduzioni dei messaggi
d'errore nella localizzazione presente.

La funzione \func{strerror} utilizza una stringa statica che non deve essere
modificata dal programma; essa è utilizzabile solo fino ad una chiamata
successiva a \func{strerror} o \func{perror} e nessun'altra funzione di
libreria tocca questa stringa. In ogni caso l'uso di una stringa statica rende
la funzione non rientrante, per cui nel caso si usino i \textit{thread} la
\acr{glibc} fornisce una apposita versione rientrante \funcd{strerror\_r}, il
cui prototipo è:

\begin{funcproto}{
\fhead{string.h}
\fdecl{char * strerror\_r(int errnum, char *buf, size\_t size)} 
\fdesc{Restituisce una stringa con un messaggio di errore.} 
}

{La funzione ritorna l'indirizzo del messaggio in caso di successo e
  \val{NULL} per un errore, nel qual caso \var{errno} assumerà uno dei valori:
  \begin{errlist}
  \item[\errcode{EINVAL}] si è specificato un valore di \param{errnum} non
    valido.
  \item[\errcode{ERANGE}] la lunghezza di \param{buf} è insufficiente a
    contenere la stringa di errore.
  \end{errlist}
}  
\end{funcproto}

Si tenga presente che questa è la versione prevista normalmente nella
\acr{glibc}, ed effettivamente definita in \headfile{string.h}, ne esiste una
analoga nello standard SUSv3 (riportata anche nella pagina di manuale), che
restituisce \code{int} al posto di \code{char *}, e che tronca la stringa
restituita a \param{size}, a cui si accede definendo le opportune macro (per
le quali si rimanda alla lettura della pagina di manuale). 

La funzione è analoga a \func{strerror} ma restituisce la stringa di errore
nel buffer \param{buf} che il singolo \textit{thread} deve allocare
autonomamente per evitare i problemi connessi alla condivisione del buffer
statico. Il messaggio è copiato fino alla dimensione massima del buffer,
specificata dall'argomento \param{size}, che deve comprendere pure il
carattere di terminazione; altrimenti la stringa risulterà troncata.

Una seconda funzione usata per riportare i codici di errore in maniera
automatizzata sullo standard error è \funcd{perror}, il cui prototipo è:

\begin{funcproto}{
\fhead{stdio.h}
\fdecl{void perror(const char *message)}
\fdesc{Stampa un messaggio di errore personalizzato.} 
}

{La funzione non ritorna nulla e non modifica \var{errno}.}
\end{funcproto}


I messaggi di errore stampati sono gli stessi di \func{strerror}, (riportati
in app.~\ref{cha:errors}), e, usando il valore corrente di \var{errno}, si
riferiscono all'ultimo errore avvenuto. La stringa specificata con
\param{message} viene stampata prima del messaggio d'errore, consentono una
personalizzazione (ad esempio l'indicazione del contesto in cui si è
verificato), seguita dai due punti e da uno spazio, il messaggio è terminato
con un a capo.  Il messaggio può essere riportato anche usando le due
variabili globali:
\includecodesnip{listati/errlist.c} 
dichiarate in \headfile{errno.h}. La prima contiene i puntatori alle stringhe
di errore indicizzati da \var{errno}; la seconda esprime il valore più alto
per un codice di errore, l'utilizzo di una di queste stringhe è
sostanzialmente equivalente a quello di \func{strerror}.

\begin{figure}[!htbp]
  \footnotesize \centering
  \begin{minipage}[c]{\codesamplewidth}
    \includecodesample{listati/errcode_mess.c}
  \end{minipage}
  \normalsize
  \caption{Codice per la stampa del messaggio di errore standard.}
  \label{fig:sys_err_mess}
\end{figure}

In fig.~\ref{fig:sys_err_mess} è riportata la sezione attinente del codice del
programma \cmd{errcode}, che può essere usato per stampare i messaggi di
errore e le costanti usate per identificare i singoli errori. Il sorgente
completo del programma è allegato nel file \file{ErrCode.c} e contiene pure la
gestione delle opzioni e tutte le definizioni necessarie ad associare il
valore numerico alla costante simbolica. In particolare si è riportata la
sezione che converte la stringa passata come argomento in un intero
(\texttt{\small 1-2}), controllando con i valori di ritorno di \funcm{strtol}
che la conversione sia avvenuta correttamente (\texttt{\small 4-10}), e poi
stampa, a seconda dell'opzione scelta il messaggio di errore (\texttt{\small
  11-14}) o la macro (\texttt{\small 15-17}) associate a quel codice.



\subsection{Alcune estensioni GNU}
\label{sec:sys_err_GNU}

Le precedenti funzioni sono quelle definite ed usate nei vari standard; la
\acr{glibc} ha però introdotto una serie di estensioni ``GNU'' che
forniscono alcune funzionalità aggiuntive per una gestione degli errori
semplificata e più efficiente. 

La prima estensione consiste in due variabili, \code{char *
  program\_invocation\_name} e \code{char * program\_invocation\_short\_name}
che consentono di ricavare il nome del proprio programma.  Queste sono utili
quando si deve aggiungere il nome del programma al messaggio d'errore, cosa
comune quando si ha un programma che non viene lanciato da linea di comando e
salva gli errori in un file di log. La prima contiene il nome usato per
lanciare il programma dalla shell ed in sostanza è equivalente ad
\code{argv[0]}; la seconda mantiene solo il nome del programma eliminando
eventuali directory qualora questo sia stato lanciato con un
\textit{pathname}.

Una seconda estensione cerca di risolvere uno dei problemi che si hanno con
l'uso di \func{perror}, dovuto al fatto che non c'è flessibilità su quello che
si può aggiungere al messaggio di errore, che può essere solo una stringa. In
molte occasioni invece serve poter scrivere dei messaggi con maggiori
informazioni. Ad esempio negli standard di programmazione GNU si richiede che
ogni messaggio di errore sia preceduto dal nome del programma, ed in generale
si può voler stampare il contenuto di qualche variabile per facilitare la
comprensione di un eventuale problema. Per questo la \acr{glibc} definisce
la funzione \funcd{error}, il cui prototipo è:

\begin{funcproto}{
\fhead{stdio.h}
\fdecl{void error(int status, int errnum, const char *format, ...)}
\fdesc{Stampa un messaggio di errore formattato.} 
}

{La funzione non ritorna nulla e non riporta errori.}  
\end{funcproto}

La funzione fa parte delle estensioni GNU per la gestione degli errori,
l'argomento \param{format} segue la stessa sintassi di \func{printf} (vedi
sez.~\ref{sec:file_formatted_io}), ed i relativi argomenti devono essere
forniti allo stesso modo, mentre \param{errnum} indica l'errore che si vuole
segnalare (non viene quindi usato il valore corrente di \var{errno}).

La funzione stampa sullo \textit{standard error} il nome del programma, come
indicato dalla variabile globale \var{program\_name}, seguito da due punti ed
uno spazio, poi dalla stringa generata da \param{format} e dagli argomenti
seguenti, seguita da due punti ed uno spazio infine il messaggio di errore
relativo ad \param{errnum}, il tutto è terminato da un a capo.

Il comportamento della funzione può essere ulteriormente controllato se si
definisce una variabile \var{error\_print\_progname} come puntatore ad una
funzione \ctyp{void} che restituisce \ctyp{void} che si incarichi di stampare
il nome del programma. 

L'argomento \param{status} può essere usato per terminare direttamente il
programma in caso di errore, nel qual caso \func{error} dopo la stampa del
messaggio di errore chiama \func{exit} con questo stato di uscita. Se invece
il valore è nullo \func{error} ritorna normalmente ma viene incrementata
un'altra variabile globale, \var{error\_message\_count}, che tiene conto di
quanti errori ci sono stati.

Un'altra funzione per la stampa degli errori, ancora più sofisticata, che
prende due argomenti aggiuntivi per indicare linea e file su cui è avvenuto
l'errore è \funcd{error\_at\_line}; il suo prototipo è:

\begin{funcproto}{
\fhead{stdio.h}
\fdecl{void error\_at\_line(int status, int errnum, const char *fname, 
  unsigned int lineno, \\
\phantom{void error\_at\_line(}const char *format, ...)}
\fdesc{Stampa un messaggio di errore formattato.} 
}

{La funzione non ritorna nulla e non riporta errori.}  
\end{funcproto}

\noindent ed il suo comportamento è identico a quello di \func{error} se non
per il fatto che, separati con il solito due punti-spazio, vengono inseriti un
nome di file indicato da \param{fname} ed un numero di linea subito dopo la
stampa del nome del programma. Inoltre essa usa un'altra variabile globale,
\var{error\_one\_per\_line}, che impostata ad un valore diverso da zero fa sì
che errori relativi alla stessa linea non vengano ripetuti.


% LocalWords:  filesystem like kernel saved header limits sysconf sez tab float
% LocalWords:  FOPEN stdio MB LEN CHAR char UCHAR unsigned SCHAR MIN signed INT
% LocalWords:  SHRT short USHRT int UINT LONG long ULONG LLONG ULLONG POSIX ARG
% LocalWords:  Stevens exec CHILD STREAM stream TZNAME timezone NGROUPS SSIZE
% LocalWords:  ssize LISTIO JOB CONTROL job control IDS VERSION YYYYMML bits bc
% LocalWords:  dall'header posix lim nell'header glibc run unistd name errno SC
% LocalWords:  NGROUP CLK TCK clock tick process PATH pathname BUF CANON path
% LocalWords:  pathconf fpathconf descriptor fd uname sys struct utsname info
% LocalWords:  EFAULT fig SOURCE NUL LENGTH DOMAIN NMLN UTSLEN system call proc
% LocalWords:  domainname sysctl BSD nlen void oldval size oldlenp newval EPERM
% LocalWords:  newlen ENOTDIR EINVAL ENOMEM linux array oldvalue paging stack
% LocalWords:  TCP shell Documentation ostype hostname osrelease version mount
% LocalWords:  const source filesystemtype mountflags ENODEV ENOTBLK block read
% LocalWords:  device EBUSY only EACCES NODEV ENXIO major RTSIG syscall PID 
% LocalWords:  number EMFILE dummy ENAMETOOLONG ENOENT ELOOP virtual devfs MGC
% LocalWords:  magic MSK RDONLY NOSUID suid sgid NOEXEC SYNCHRONOUS REMOUNT MNT
% LocalWords:  MANDLOCK mandatory locking WRITE APPEND append IMMUTABLE NOATIME
% LocalWords:  access NODIRATIME BIND MOVE umount flags FORCE statfs fstatfs ut
% LocalWords:  buf ENOSYS EIO EBADF type fstab mntent home shadow username uid
% LocalWords:  passwd PAM Pluggable Authentication Method Service Switch pwd ru
% LocalWords:  getpwuid getpwnam NULL buflen result ERANGE getgrnam getgrgid AS
% LocalWords:  grp group gid SVID fgetpwent putpwent getpwent setpwent endpwent
% LocalWords:  fgetgrent putgrent getgrent setgrent endgrent accounting init HZ
% LocalWords:  runlevel Hierarchy logout setutent endutent utmpname utmp paths
% LocalWords:  WTMP getutent getutid getutline pututline LVL OLD DEAD EMPTY dev
% LocalWords:  line libc XPG utmpx getutxent getutxid getutxline pututxline who
% LocalWords:  setutxent endutxent wmtp updwtmp logwtmp wtmp host rusage utime
% LocalWords:  minflt majflt nswap fault swap timeval wait getrusage usage SELF
% LocalWords:  CHILDREN current limit soft RLIMIT address brk mremap mmap dump
% LocalWords:  SIGSEGV SIGXCPU SIGKILL sbrk FSIZE SIGXFSZ EFBIG LOCKS lock dup
% LocalWords:  MEMLOCK NOFILE NPROC fork EAGAIN SIGPENDING sigqueue kill RSS tv
% LocalWords:  resource getrlimit setrlimit rlimit rlim INFINITY capabilities
% LocalWords:  capability CAP Sun Sparc PAGESIZE getpagesize SVr SUSv get IGN
% LocalWords:  phys pages avphys NPROCESSORS CONF ONLN getloadavg stdlib double
% LocalWords:  loadavg nelem scheduler CONFIG ACCT acct filename EUSER sizeof
% LocalWords:  ENFILE EROFS PACCT AcctCtrl cap calendar UTC Jan the Epoch GMT
% LocalWords:  Greenwich Mean l'UTC timer CLOCKS SEC cron wall elapsed times tz
% LocalWords:  tms cutime cstime waitpid gettimeofday settimeofday timex NetBSD
% LocalWords:  timespec adjtime olddelta adjtimex David Mills RFC NTP ntp cmd
% LocalWords:  nell'RFC ADJ FREQUENCY frequency MAXERROR maxerror ESTERROR PLL
% LocalWords:  esterror TIMECONST constant SINGLESHOT MOD INS insert leap OOP
% LocalWords:  second delete progress has occurred BAD broken tm gmtoff asctime
% LocalWords:  ctime timep gmtime localtime mktime tzname tzset daylight format
% LocalWords:  strftime thread EOF modifiable lvalue app errcode strerror LC at
% LocalWords:  perror string errnum MESSAGES error message strtol log jiffy asm
% LocalWords:  program invocation argv printf print progname exit count fname
% LocalWords:  lineno one standardese Di page Wed Wednesday Apr April PM AM CAD
% LocalWords:  CEST utmpxname Solaris updwtmpx reboot RESTART Ctrl OFF SIGINT
% LocalWords:  HALT halted sync KEXEC kexec load bootloader POWER Power with nr
% LocalWords:  Restarting command arg entry segments segment ARCH CRASH CONTEXT
% LocalWords:  PRESERVE PPC IA ARM SH MIPS nvcsw nivcsw inblock oublock maxrss
% LocalWords:  context switch slice Resident SIG SIGCHLD cur Gb lease mlock Hz
% LocalWords:  memory mlockall MAP LOCKED shmctl MSGQUEUE attr NICE nice MADV
% LocalWords:  madvise WILLNEED RTPRIO sched setscheduler setparam scheduling
% LocalWords:  RTTIME execve kb prlimit pid new old ESRCH EUSERS refresh high
% LocalWords:  resolution HRT jiffies strptime

%%% Local Variables: 
%%% mode: latex
%%% TeX-master: "gapil"
%%% End: 

%% signal.tex
%%
%% Copyright (C) 2000-2018 Simone Piccardi.  Permission is granted to
%% copy, distribute and/or modify this document under the terms of the GNU Free
%% Documentation License, Version 1.1 or any later version published by the
%% Free Software Foundation; with the Invariant Sections being "Un preambolo",
%% with no Front-Cover Texts, and with no Back-Cover Texts.  A copy of the
%% license is included in the section entitled "GNU Free Documentation
%% License".
%%

\chapter{I segnali}
\label{cha:signals}

I segnali sono il primo e più semplice meccanismo di comunicazione nei
confronti dei processi. Nella loro versione originale essi portano con sé
nessuna informazione che non sia il loro tipo; si tratta in sostanza di
un'interruzione software portata ad un processo.

In genere essi vengono usati dal kernel per riportare ai processi situazioni
eccezionali (come errori di accesso, eccezioni aritmetiche, ecc.) ma possono
anche essere usati come forma elementare di comunicazione fra processi (ad
esempio vengono usati per il controllo di sessione), per notificare eventi
(come la terminazione di un processo figlio), ecc.

In questo capitolo esamineremo i vari aspetti della gestione dei segnali,
partendo da una introduzione relativa ai concetti base con cui essi vengono
realizzati, per poi affrontarne la classificazione a secondo di uso e modalità
di generazione fino ad esaminare in dettaglio le funzioni e le metodologie di
gestione avanzate e le estensioni fatte all'interfaccia classica nelle nuovi
versioni dello standard POSIX.


\section{Introduzione}
\label{sec:sig_intro}

In questa sezione esamineremo i concetti generali relativi ai segnali, vedremo
le loro caratteristiche di base, introdurremo le nozioni di fondo relative
all'architettura del funzionamento dei segnali e alle modalità con cui il
sistema gestisce l'interazione fra di essi ed i processi.


\subsection{I concetti base}
\label{sec:sig_base}

Come il nome stesso indica i segnali sono usati per notificare ad un processo
l'occorrenza di un qualche evento. Gli eventi che possono generare un segnale
sono vari; un breve elenco di possibili cause per l'emissione di un segnale è
il seguente:

\begin{itemize*}
\item un errore del programma, come una divisione per zero o un tentativo di
  accesso alla memoria fuori dai limiti validi;
\item la terminazione di un processo figlio;
\item la scadenza di un timer o di un allarme;
\item il tentativo di effettuare un'operazione di input/output che non può
  essere eseguita;
\item una richiesta dell'utente dal terminale di terminare o fermare il
  programma.
\item l'invio esplicito da parte del processo stesso o di un altro.
\end{itemize*}

Ciascuno di questi eventi, compresi gli ultimi due che pure sono controllati
dall'utente o da un altro processo, comporta l'intervento diretto da parte del
kernel che causa la generazione di un particolare tipo di segnale.

Quando un processo riceve un segnale, invece del normale corso del programma,
viene eseguita una azione predefinita o una apposita funzione di gestione che
può essere stata specificata dall'utente, nel qual caso si dice che si
\textsl{intercetta} il segnale. Riprendendo la terminologia originale da qui
in avanti faremo riferimento a questa funzione come al \textsl{gestore} del
segnale, traduzione approssimata dell'inglese \textit{signal handler}.


\subsection{Le \textsl{semantiche} del funzionamento dei segnali}
\label{sec:sig_semantics}

Negli anni il comportamento del sistema in risposta ai segnali è stato
modificato in vari modi nelle differenti implementazioni di Unix.  Si possono
individuare due tipologie fondamentali di comportamento dei segnali (dette
\textsl{semantiche}) che vengono chiamate rispettivamente \textsl{semantica
  affidabile} (o \textit{reliable}) e \textsl{semantica inaffidabile} (o
\textit{unreliable}).

Nella \textsl{semantica inaffidabile}, che veniva implementata dalle prime
versioni di Unix, la funzione di gestione del segnale specificata dall'utente
non restava attiva una volta che era stata eseguita; era perciò compito
dell'utente ripetere l'installazione dello stesso all'interno del
\textsl{gestore} del segnale in tutti quei casi in cui si voleva che esso
restasse attivo.

\begin{figure}[!htbp]
  \footnotesize \centering
  \begin{minipage}[c]{\codesamplewidth}
    \includecodesample{listati/unreliable_sig.c}
  \end{minipage} 
  \normalsize 
  \caption{Esempio di codice di un gestore di segnale per la semantica
    inaffidabile.} 
  \label{fig:sig_old_handler}
\end{figure}

In questo caso però è possibile una situazione in cui i segnali possono essere
perduti. Si consideri il segmento di codice riportato in
fig.~\ref{fig:sig_old_handler}: nel programma principale viene installato un
gestore (\texttt{\small 5}), la cui prima operazione (\texttt{\small 11}) è
quella di reinstallare se stesso. Se nell'esecuzione del gestore fosse
arrivato un secondo segnale prima che esso abbia potuto eseguire la
reinstallazione di se stesso per questo secondo segnale verrebbe eseguito il
comportamento predefinito, il che può comportare, a seconda dei casi, la
perdita del segnale (se l'impostazione predefinita è quella di ignorarlo) o la
terminazione immediata del processo; in entrambi i casi l'azione prevista dal
gestore non verrebbe eseguita.

Questa è la ragione per cui l'implementazione dei segnali secondo questa
semantica viene chiamata \textsl{inaffidabile}: infatti la ricezione del
segnale e la reinstallazione del suo gestore non sono operazioni atomiche, e
sono sempre possibili delle \textit{race condition} (si ricordi
sez.~\ref{sec:proc_multi_prog}).  Un altro problema è che in questa semantica
non esiste un modo per bloccare i segnali quando non si vuole che arrivino; i
processi possono ignorare il segnale, ma non è possibile istruire il sistema a
non fare nulla in occasione di un segnale, pur mantenendo memoria del fatto
che è avvenuto.

Nella semantica \textsl{affidabile} (quella utilizzata da Linux e da ogni Unix
moderno) il gestore una volta installato resta attivo e non si hanno tutti i
problemi precedenti. In questa semantica i segnali vengono \textsl{generati}
dal kernel per un processo all'occorrenza dell'evento che causa il segnale. In
genere questo viene fatto dal kernel impostando un apposito campo della
\struct{task\_struct} del processo nella \textit{process table} (si veda
fig.~\ref{fig:proc_task_struct}).

Si dice che il segnale viene \textsl{consegnato} al processo (dall'inglese
\textit{delivered}) quando viene eseguita l'azione per esso prevista, mentre
per tutto il tempo che passa fra la generazione del segnale e la sua consegna
esso è detto \textsl{pendente} (o \textit{pending}). In genere questa
procedura viene effettuata dallo \textit{scheduler} quando, riprendendo
l'esecuzione del processo in questione, verifica la presenza del segnale nella
\struct{task\_struct} e mette in esecuzione il gestore.

In questa semantica un processo ha la possibilità di bloccare la consegna dei
segnali, in questo caso, se l'azione per il suddetto segnale non è quella di
ignorarlo, il segnale resta \textsl{pendente} fintanto che il processo non lo
sblocca (nel qual caso viene consegnato) o imposta l'azione corrispondente per
ignorarlo.

Si tenga presente che il kernel stabilisce cosa fare con un segnale che è
stato bloccato al momento della consegna, non quando viene generato; questo
consente di cambiare l'azione per il segnale prima che esso venga consegnato,
e si può usare la funzione \func{sigpending} (vedi sez.~\ref{sec:sig_sigmask})
per determinare quali segnali sono bloccati e quali sono pendenti.

Infine occorre precisare che i segnali predatano il supporto per i
\textit{thread} e vengono sempre inviati al processo come insieme, cosa che
può creare incertezza nel caso questo sia multi-\textit{thread}. In tal caso
quando è possibile determinare quale è il \textit{thread} specifico che deve
ricevere il segnale, come avviene per i segnali di errore, questo sarà inviato
solo a lui, altrimenti sarà inviato a discrezione del kernel ad uno qualunque
dei \textit{thread} del processo che possa riceverlo (che cioè non blocchi il
segnale), torneremo sull'argomento in sez.~\ref{sec:thread_signal}.

\subsection{Tipi di segnali}
\label{sec:sig_types}

In generale si tende a classificare gli eventi che possono generare dei
segnali in tre categorie principali: errori, eventi esterni e richieste
esplicite.

Un errore significa che un programma ha fatto qualcosa di sbagliato e non può
continuare ad essere eseguito. Non tutti gli errori causano dei segnali, in
genere le condizioni di errore più comuni comportano la restituzione di un
codice di errore da parte di una funzione di libreria. Sono gli errori che
possono avvenire nell'esecuzione delle istruzioni di un programma, come le
divisioni per zero o l'uso di indirizzi di memoria non validi, che causano
l'emissione di un segnale.

Un evento esterno ha in genere a che fare con le operazioni di lettura e
scrittura su file, o con l'interazione con dispositivi o con altri processi;
esempi di segnali di questo tipo sono quelli legati all'arrivo di dati in
ingresso, scadenze di un timer, terminazione di processi figli, la pressione
dei tasti di stop o di suspend su un terminale.

Una richiesta esplicita significa l'uso da parte di un programma delle
apposite funzioni di sistema, come \func{kill} ed affini (vedi
sez.~\ref{sec:sig_kill_raise}) per la generazione ``\textsl{manuale}'' di un
segnale.

Si dice poi che i segnali possono essere \textsl{asincroni} o
\textsl{sincroni}. Un segnale \textsl{sincrono} è legato ad una azione
specifica di un programma ed è inviato (a meno che non sia bloccato) durante
tale azione. Molti errori generano segnali \textsl{sincroni}, così come la
richiesta esplicita da parte del processo tramite le chiamate al sistema.
Alcuni errori come la divisione per zero non sono completamente sincroni e
possono arrivare dopo qualche istruzione.

I segnali \textsl{asincroni} sono generati da eventi fuori dal controllo del
processo che li riceve, e arrivano in tempi impredicibili nel corso
dell'esecuzione del programma. Eventi esterni come la terminazione di un
processo figlio generano segnali \textsl{asincroni}, così come le richieste di
generazione di un segnale effettuate da altri processi.

In generale un tipo di segnale o è sincrono o è asincrono, salvo il caso in
cui esso sia generato attraverso una richiesta esplicita tramite chiamata al
sistema, nel qual caso qualunque tipo di segnale (quello scelto nella
chiamata) può diventare sincrono o asincrono a seconda che sia generato
internamente o esternamente al processo.


\subsection{La notifica dei segnali}
\label{sec:sig_notification}

Come accennato quando un segnale viene generato, se la sua azione predefinita
non è quella di essere ignorato, il kernel prende nota del fatto nella
\struct{task\_struct} del processo; si dice così che il segnale diventa
\textsl{pendente} (o \textit{pending}), e rimane tale fino al momento in cui
verrà notificato al processo o verrà specificata come azione quella di
ignorarlo.

Normalmente l'invio al processo che deve ricevere il segnale è immediato ed
avviene non appena questo viene rimesso in esecuzione dallo \textit{scheduler}
che esegue l'azione specificata. Questo a meno che il segnale in questione non
sia stato bloccato prima della notifica, nel qual caso l'invio non avviene ed
il segnale resta \textsl{pendente} indefinitamente.

Quando lo si sblocca un segnale \textsl{pendente} sarà subito notificato. Si
tenga presente però che tradizionalmente i segnali \textsl{pendenti} non si
accodano, alla generazione infatti il kernel marca un flag nella
\struct{task\_struct} del processo, per cui se prima della notifica ne vengono
generati altri il flag è comunque marcato, ed il gestore viene eseguito sempre
una sola volta. In realtà questo non vale nel caso dei cosiddetti segnali
\textit{real-time}, che vedremo in sez.~\ref{sec:sig_real_time}, ma questa è
una funzionalità avanzata che per ora tralasceremo.

Si ricordi inoltre che se l'azione specificata per un segnale è quella di
essere ignorato questo sarà scartato immediatamente al momento della sua
generazione, e questo anche se in quel momento il segnale è bloccato, perché
bloccare su un segnale significa bloccarne la notifica. Per questo motivo un
segnale, fintanto che viene ignorato, non sarà mai notificato, anche se prima
è stato bloccato ed in seguito si è specificata una azione diversa, nel qual
caso solo i segnali successivi alla nuova specificazione saranno notificati.

Una volta che un segnale viene notificato, che questo avvenga subito o dopo
una attesa più o meno lunga, viene eseguita l'azione specificata per il
segnale. Per alcuni segnali (per la precisione \signal{SIGKILL} e
\signal{SIGSTOP}) questa azione è predeterminata dal kernel e non può essere
mai modificata, ma per tutti gli altri si può selezionare una delle tre
possibilità seguenti:

\begin{itemize*}
\item ignorare il segnale;
\item intercettare il segnale, ed utilizzare il gestore specificato;
\item accettare l'azione predefinita per quel segnale.
\end{itemize*}

Un programma può specificare queste scelte usando le due funzioni
\func{signal} e \func{sigaction}, che tratteremo rispettivamente in
sez.~\ref{sec:sig_signal} e sez.~\ref{sec:sig_sigaction}. Se si è installato
un gestore sarà quest'ultimo ad essere eseguito alla notifica del segnale.
Inoltre il sistema farà si che mentre viene eseguito il gestore di un segnale,
quest'ultimo venga automaticamente bloccato, così si possono evitare alla
radice possibili \textit{race condition}.

Nel caso non sia stata specificata un'azione, viene utilizzata la cosiddetta
azione predefinita che, come vedremo in sez.~\ref{sec:sig_standard}, è propria
di ciascun segnale. Nella maggior parte dei casi questa azione comporta la
terminazione immediata del processo, ma per alcuni segnali che rappresentano
eventi innocui l'azione predefinita è di essere ignorati. Inoltre esistono
alcuni segnali la cui azione è semplicemente quella di fermare l'esecuzione
del programma, vale a dire portarlo nello stato di \textit{stopped} (lo stato
\texttt{T}, si ricordi tab.~\ref{tab:proc_proc_states} e quanto illustrato in
sez.~\ref{sec:proc_sched}).

Quando un segnale termina un processo il padre può determinare la causa della
terminazione esaminandone lo stato di uscita così come viene riportato dalle
funzioni \func{wait} e \func{waitpid} (vedi sez.~\ref{sec:proc_wait}). Questo
ad esempio è il modo in cui la shell determina i motivi della terminazione di
un programma e scrive un eventuale messaggio di errore.

\itindbeg{core~dump}

I segnali che rappresentano errori del programma (divisione per zero o
violazioni di accesso) hanno come ulteriore caratteristica della loro azione
predefinita, oltre a terminare il processo, quella di scrivere nella directory
di lavoro corrente del processo di un file \file{core} su cui viene salvata
l'immagine della memoria del processo.

Questo file costituisce il cosiddetto \textit{core dump}, e contenendo
l'immagine della memoria del processo, consente di risalire allo stato dello
\textit{stack} (vedi sez.~\ref{sec:proc_mem_layout}) prima della
terminazione. Questo permette di esaminare il contenuto del file un secondo
tempo con un apposito programma (un \textit{debugger} come \cmd{gdb}) per
investigare sulla causa dell'errore, ed in particolare, grazie appunto ai dati
dello \textit{stack}, consente di identificare quale funzione ha causato
l'errore.

Si tenga presente che il \textit{core dump} viene creato non solo in caso di
errore effettivo, ma anche se il segnale per cui la sua creazione è prevista
nell'azione dell'azione predefinita viene inviato al programma con una delle
funzioni \func{kill}, \func{raise}, ecc.

\itindend{core~dump}


\section{La classificazione dei segnali}
\label{sec:sig_classification}

Esamineremo in questa sezione quali sono i vari segnali definiti nel sistema,
quali sono le loro caratteristiche e la loro tipologia, tratteremo le varie
macro e costanti che permettono di identificarli, e illustreremo le funzioni
che ne stampano la descrizione.


\subsection{I segnali standard}
\label{sec:sig_standard}

Ciascun segnale è identificato dal kernel con un numero, ma benché per alcuni
segnali questi numeri siano sempre gli stessi, tanto da essere usati come
sinonimi, l'uso diretto degli identificativi numerici da parte dei programmi è
comunque da evitare, in quanto essi non sono mai stati standardizzati e
possono variare a seconda dell'implementazione del sistema, e nel caso di
Linux anche a seconda della architettura hardware e della versione del kernel.

Quelli che invece sono stati, almeno a grandi linee, standardizzati, sono i
nomi dei segnali e le costanti di preprocessore che li identificano, che sono
tutte nella forma \texttt{SIG\textsl{nome}}, e sono queste che devono essere
usate nei programmi. Come tutti gli altri nomi e le funzioni che concernono i
segnali, esse sono definite nell'header di sistema \headfile{signal.h}.

\begin{table}[!htb]
  \footnotesize
  \centering
  \begin{tabular}[c]{|l|c|c|l|}
    \hline
    \textbf{Segnale} &\textbf{Standard}&\textbf{Azione}&\textbf{Descrizione} \\
    \hline
    \hline
    \signal{SIGHUP}  &P & T & Hangup o terminazione del processo di 
                              controllo.\\
    \signal{SIGINT}  &PA& T & Interrupt da tastiera (\cmd{C-c}).\\
    \signal{SIGQUIT} &P & C & Quit da tastiera (\cmd{C-y}).\\
    \signal{SIGILL}  &PA& C & Istruzione illecita.\\
    \signal{SIGTRAP} &S & C & Trappole per un Trace/breakpoint.\\
    \signal{SIGABRT} &PA& C & Segnale di abort da \func{abort}.\\
    \signald{SIGIOT}  &B & C & Trappola di I/O. Sinonimo di \signal{SIGABRT}.\\
    \signal{SIGBUS}  &BS& C & Errore sul bus (bad memory access).\\
    \signal{SIGFPE}  &AP& C & Errore aritmetico.\\
    \signal{SIGKILL} &P & T& Segnale di terminazione forzata.\\
    \signal{SIGUSR1} &P & T & Segnale utente numero 1.\\
    \signal{SIGSEGV} &AP& C & Errore di accesso in memoria.\\
    \signal{SIGUSR2} &P & T & Segnale utente numero 2.\\
    \signal{SIGPIPE} &P & T & \textit{Pipe} spezzata.\\
    \signal{SIGALRM} &P & T & Segnale del timer da \func{alarm}.\\
    \signal{SIGTERM} &AP& T & Segnale di terminazione (\texttt{C-\bslash}).\\
    \signal{SIGCHLD} &P & I & Figlio terminato o fermato.\\
    \signal{SIGCONT} &P &-- & Continua se fermato.\\
    \signal{SIGSTOP} &P & S & Ferma il processo.\\
    \signal{SIGTSTP} &P & S & Pressione del tasto di stop sul terminale.\\
    \signal{SIGTTIN} &P & S & Input sul terminale per un processo 
                              in background.\\
    \signal{SIGTTOU} &P & S & Output sul terminale per un processo          
                              in background.\\
    \signal{SIGURG}  &BS& I & Ricezione di una \textit{urgent condition} su 
                              un socket.\\
    \signal{SIGXCPU} &BS& C & Ecceduto il limite sul tempo di CPU.\\
    \signal{SIGXFSZ} &BS& C & Ecceduto il limite sulla dimensione dei file.\\
    \signal{SIGVTALRM}&BS& T& Timer di esecuzione scaduto.\\
    \signal{SIGPROF} &BS& T & Timer del \textit{profiling} scaduto.\\
    \signal{SIGWINCH}&B & I & Finestra ridimensionata (4.3BSD, Sun).\\
    \signal{SIGIO}   &B & T & L'I/O è possibile.\\
    \signal{SIGPOLL} &VS& T & \textit{Pollable event}, sinonimo di
                              \signal{SIGIO}.\\
    \signal{SIGPWR}  &V & T & Fallimento dell'alimentazione.\\
    \signal{SIGSYS}  &VS& C & \textit{system call} sbagliata.\\
    \hline
    \signal{SIGSTKFLT}&?& T & Errore sullo stack del coprocessore (inusato).\\
    \signald{SIGUNUSED}&?& C & Segnale inutilizzato (sinonimo di
                               \signal{SIGSYS}).\\
    \hline
    \signal{SIGCLD}  &V & I & Sinonimo di \signal{SIGCHLD}.\\
    \signal{SIGEMT}  &V & C & Trappola di emulatore.\\
    \signal{SIGINFO} &B & T & Sinonimo di \signal{SIGPWR}.\\
    \signal{SIGLOST} &? & T & Perso un lock sul file, sinonimo
                              di \signal{SIGIO} (inusato).\\
    \hline
  \end{tabular}
  \caption{Lista dei segnali ordinari in Linux.}
  \label{tab:sig_signal_list}
\end{table}

In tab.~\ref{tab:sig_signal_list} si è riportato l'elenco completo dei segnali
ordinari definiti su Linux per tutte le possibili architetture (tratteremo
quelli \textit{real-time} in sez.~\ref{sec:sig_real_time}). Ma si tenga
presente che solo quelli elencati nella prima sezione della tabella sono
presenti su tutte le architetture. Nelle sezioni successive si sono riportati
rispettivamente quelli che esistono solo sull'architettura PC e quelli che non
esistono sull'architettura PC, ma sono definiti sulle architetture
\textit{alpha} o \textit{mips}.

Alcuni segnali erano previsti fin dallo standard ANSI C, ed i segnali sono
presenti in tutti i sistemi unix-like, ma l'elenco di quelli disponibili non è
uniforme, ed alcuni di essi sono presenti solo su alcune implementazioni o
architetture hardware, ed anche il loro significato può variare. Per questo si
sono riportati nella seconda colonna della tabella riporta gli standard in cui
ciascun segnale è stato definito, indicati con altrettante lettere da
interpretare secondo la legenda di tab.~\ref{tab:sig_standard_leg}. Si tenga
presente che il significato dei segnali è abbastanza indipendente dalle
implementazioni solo per quelli definiti negli standard POSIX.1-1990 e
POSIX.1-2001. 

\begin{table}[htb]
  \footnotesize
  \centering
  \begin{tabular}[c]{|c|l|}
    \hline
    \textbf{Sigla} & \textbf{Standard} \\
    \hline
    \hline
    P & POSIX.1-1990.\\
    B & BSD (4.2 BSD e Sun).\\
    A & ANSI C.\\
    S & SUSv2 (e POSIX.1-2001).\\
    V & System V.\\
    ? & Ignoto.\\
    \hline
  \end{tabular}
  \caption{Legenda dei valori degli standard riportati nella seconda colonna
    di tab.~\ref{tab:sig_signal_list}.} 
  \label{tab:sig_standard_leg}
\end{table}

Come accennato in sez.~\ref{sec:sig_notification} a ciascun segnale è
associata una specifica azione predefinita che viene eseguita quando nessun
gestore è installato. Le azioni predefinite possibili, che abbiamo già
descritto in sez.~\ref{sec:sig_notification}, sono state riportate in
tab.~\ref{tab:sig_signal_list} nella terza colonna, e di nuovo sono state
indicate con delle lettere la cui legenda completa è illustrata in
tab.~\ref{tab:sig_action_leg}).

\begin{table}[htb]
  \footnotesize
  \centering
  \begin{tabular}[c]{|c|l|}
    \hline
    \textbf{Sigla} & \textbf{Significato} \\
    \hline
    \hline
    T & L'azione predefinita è terminare il processo.\\
    C & L'azione predefinita è terminare il processo e scrivere un 
        \textit{core dump}.\\
    I & L'azione predefinita è ignorare il segnale.\\
    S & L'azione predefinita è fermare il processo.\\
    \hline
  \end{tabular}
  \caption{Legenda delle azioni predefinite dei segnali riportate nella terza
    colonna di tab.~\ref{tab:sig_signal_list}.}
  \label{tab:sig_action_leg}
\end{table}


Si inoltre noti come \signal{SIGCONT} sia l'unico segnale a non avere
l'indicazione di una azione predefinita nella terza colonna di
tab.~\ref{tab:sig_signal_list}, questo perché il suo effetto è sempre quello
di far ripartire un programma in stato \texttt{T} fermato da un segnale di
stop. Inoltre i segnali \signal{SIGSTOP} e \signal{SIGKILL} si distinguono da
tutti gli altri per la specifica caratteristica di non potere essere né
intercettati, né bloccati, né ignorati.

Il numero totale di segnali presenti è dato dalla macro \macrod{NSIG} (e tiene
conto anche di quelli \textit{real-time}) e dato che i numeri dei segnali sono
allocati progressivamente, essa corrisponde anche al successivo del valore
numerico assegnato all'ultimo segnale definito.  La descrizione dettagliata
del significato dei precedenti segnali, raggruppati per tipologia, verrà
affrontata nei paragrafi successivi.


\subsection{I segnali di errore}
\label{sec:sig_prog_error}

Questi segnali sono generati quando il sistema, o in certi casi direttamente
l'hardware (come per i \textit{page fault} non validi o le eccezioni del
processore) rileva un qualche errore insanabile nel programma in
esecuzione. In generale la generazione di questi segnali significa che il
programma ha dei gravi problemi (ad esempio ha dereferenziato un puntatore non
valido o ha eseguito una operazione aritmetica proibita) e l'esecuzione non
può essere proseguita.

In genere si intercettano questi segnali per permettere al programma di
terminare in maniera pulita, ad esempio per ripristinare le impostazioni della
console o eliminare i file di lock prima dell'uscita.  In questo caso il
gestore deve concludersi ripristinando l'azione predefinita e rialzando il
segnale, in questo modo il programma si concluderà senza effetti spiacevoli,
ma riportando lo stesso stato di uscita che avrebbe avuto se il gestore non ci
fosse stato.

L'azione predefinita per tutti questi segnali è causare la terminazione del
processo che li ha causati. In genere oltre a questo il segnale provoca pure
la registrazione su disco di un file di \textit{core dump}, che un debugger
può usare per ricostruire lo stato del programma al momento della
terminazione.  Questi segnali sono:
\begin{basedescript}{\desclabelwidth{2.0cm}}
\item[\signald{SIGFPE}] Riporta un errore aritmetico fatale. Benché il nome
  derivi da \textit{floating point exception} si applica a tutti gli errori
  aritmetici compresa la divisione per zero e l'overflow.  Se il gestore
  ritorna il comportamento del processo è indefinito, ed ignorare questo
  segnale può condurre ad un ciclo infinito.

%   Per questo segnale le cose sono complicate dal fatto che possono esserci
%   molte diverse eccezioni che \signal{SIGFPE} non distingue, mentre lo
%   standard IEEE per le operazioni in virgola mobile definisce varie eccezioni
%   aritmetiche e richiede che esse siano notificate.
% TODO trovare altre info su SIGFPE e trattare la notifica delle eccezioni 
  
\item[\signald{SIGILL}] Il nome deriva da \textit{illegal instruction},
  significa che il programma sta cercando di eseguire una istruzione
  privilegiata o inesistente, in generale del codice illecito. Poiché il
  compilatore del C genera del codice valido si ottiene questo segnale se il
  file eseguibile è corrotto o si stanno cercando di eseguire dei dati.
  Quest'ultimo caso può accadere quando si passa un puntatore sbagliato al
  posto di un puntatore a funzione, o si eccede la scrittura di un vettore di
  una variabile locale, andando a corrompere lo \textit{stack}. Lo stesso
  segnale viene generato in caso di overflow dello \textit{stack} o di
  problemi nell'esecuzione di un gestore.  Se il gestore ritorna il
  comportamento del processo è indefinito.

\item[\signald{SIGSEGV}] Il nome deriva da \textit{segment violation}, e
  significa che il programma sta cercando di leggere o scrivere in una zona di
  memoria protetta al di fuori di quella che gli è stata riservata dal
  sistema. In genere è il meccanismo della protezione della memoria che si
  accorge dell'errore ed il kernel genera il segnale.  È tipico ottenere
  questo segnale dereferenziando un puntatore nullo o non inizializzato
  leggendo al di là della fine di un vettore.  Se il gestore ritorna il
  comportamento del processo è indefinito.

\item[\signald{SIGBUS}] Il nome deriva da \textit{bus error}. Come
  \signal{SIGSEGV} questo è un segnale che viene generato di solito quando si
  dereferenzia un puntatore non inizializzato, la differenza è che
  \signal{SIGSEGV} indica un accesso non permesso su un indirizzo esistente
  (al di fuori dallo \textit{heap} o dallo \textit{stack}), mentre
  \signal{SIGBUS} indica l'accesso ad un indirizzo non valido, come nel caso
  di un puntatore non allineato.

\item[\signald{SIGABRT}] Il nome deriva da \textit{abort}. Il segnale indica
  che il programma stesso ha rilevato un errore che viene riportato chiamando
  la funzione \func{abort}, che genera questo segnale.

\item[\signald{SIGTRAP}] È il segnale generato da un'istruzione di breakpoint o
  dall'attivazione del tracciamento per il processo. È usato dai programmi per
  il debugging e un programma normale non dovrebbe ricevere questo segnale.

\item[\signald{SIGSYS}] Sta ad indicare che si è eseguita una istruzione che
  richiede l'esecuzione di una \textit{system call}, ma si è fornito un codice
  sbagliato per quest'ultima. 

\item[\signald{SIGEMT}] Il nome sta per \textit{emulation trap}. Il segnale non
  è previsto da nessuno standard ed è definito solo su alcune architetture che
  come il vecchio PDP11 prevedono questo tipo di interruzione, non è presente
  sui normali PC.
\end{basedescript}


\subsection{I segnali di terminazione}
\label{sec:sig_termination}

Questo tipo di segnali sono usati per terminare un processo; hanno vari nomi a
causa del differente uso che se ne può fare, ed i programmi possono
trattarli in maniera differente. 

La ragione per cui può essere necessario intercettare questi segnali è che il
programma può dover eseguire una serie di azioni di pulizia prima di
terminare, come salvare informazioni sullo stato in cui si trova, cancellare
file temporanei, o ripristinare delle condizioni alterate durante il
funzionamento (come il modo del terminale o le impostazioni di una qualche
periferica). L'azione predefinita di questi segnali è di terminare il
processo, questi segnali sono:
\begin{basedescript}{\desclabelwidth{2.0cm}}
\item[\signald{SIGTERM}] Il nome sta per \textit{terminate}. È un segnale
  generico usato per causare la conclusione di un programma. È quello che
  viene generato di default dal comando \cmd{kill}.  Al contrario di
  \signal{SIGKILL} può essere intercettato, ignorato, bloccato. In genere lo
  si usa per chiedere in maniera ``\textsl{educata}'' ad un processo di
  concludersi.

\item[\signald{SIGINT}] Il nome sta per \textit{interrupt}. È il segnale di
  interruzione per il programma. È quello che viene generato di default dal
  dall'invio sul terminale del carattere di controllo ``\textit{INTR}'',
  \textit{interrupt} appunto, che viene generato normalmente dalla sequenza
  \cmd{C-c} sulla tastiera.

\item[\signald{SIGQUIT}] È analogo a \signal{SIGINT} con la differenza che è
  controllato da un altro carattere di controllo, ``\textit{QUIT}'',
  corrispondente alla sequenza \texttt{C-\bslash} sulla tastiera. A differenza
  del precedente l'azione predefinita, oltre alla terminazione del processo,
  comporta anche la creazione di un \textit{core dump}.  In genere lo si può
  pensare come corrispondente ad una condizione di errore del programma
  rilevata dall'utente. Per questo motivo non è opportuno fare eseguire al
  gestore di questo segnale le operazioni di pulizia normalmente previste
  (tipo la cancellazione di file temporanei), dato che in certi casi esse
  possono eliminare informazioni utili nell'esame dei \textit{core dump}.

\item[\signald{SIGKILL}] Il nome è utilizzato per terminare in maniera immediata
  qualunque programma. Questo segnale non può essere né intercettato, né
  ignorato, né bloccato, per cui causa comunque la terminazione del processo.
  In genere esso viene generato solo per richiesta esplicita dell'utente dal
  comando (o tramite la funzione) \cmd{kill}. Dato che non lo si può
  intercettare è sempre meglio usarlo come ultima risorsa quando metodi meno
  brutali, come \signal{SIGTERM} o \cmd{C-c} non funzionano. 

  Se un processo non risponde a nessun altro segnale \signal{SIGKILL} ne causa
  sempre la terminazione (in effetti il fallimento della terminazione di un
  processo da parte di \signal{SIGKILL} costituirebbe un malfunzionamento del
  kernel). Talvolta è il sistema stesso che può generare questo segnale quando
  per condizioni particolari il processo non può più essere eseguito neanche
  per eseguire un gestore.

\item[\signald{SIGHUP}] Il nome sta per \textit{hang-up}. Segnala che il
  terminale dell'utente si è disconnesso, ad esempio perché si è interrotta la
  rete. Viene usato anche per riportare la terminazione del processo di
  controllo di un terminale a tutti i processi della sessione (vedi
  sez.~\ref{sec:sess_job_control}), in modo che essi possano disconnettersi
  dal relativo terminale.  Viene inoltre usato in genere per segnalare ai
  programmi di servizio (i cosiddetti \textsl{demoni}, vedi
  sez.~\ref{sec:sess_daemon}), che non hanno un terminale di controllo, la
  necessità di reinizializzarsi e rileggere il file (o i file) di
  configurazione.
\end{basedescript}


\subsection{I segnali di allarme}
\label{sec:sig_alarm}

Questi segnali sono generati dalla scadenza di un timer (vedi
sez.~\ref{sec:sig_alarm_abort}). Il loro comportamento predefinito è quello di
causare la terminazione del programma, ma con questi segnali la scelta
predefinita è irrilevante, in quanto il loro uso presuppone sempre la
necessità di un gestore.  Questi segnali sono:
\begin{basedescript}{\desclabelwidth{2.0cm}}
\item[\signald{SIGALRM}] Il nome sta per \textit{alarm}. Segnale la scadenza di
  un timer misurato sul tempo reale o sull'orologio di sistema. È normalmente
  usato dalla funzione \func{alarm}.

\item[\signald{SIVGTALRM}] Il nome sta per \textit{virtual alarm}. È analogo al
  precedente ma segnala la scadenza di un timer sul tempo di CPU usato dal
  processo. 

\item[\signald{SIGPROF}] Il nome sta per \textit{profiling}. Indica la scadenza
  di un timer che misura sia il tempo di CPU speso direttamente dal processo
  che quello che il sistema ha speso per conto di quest'ultimo. In genere
  viene usato dagli strumenti che servono a fare la profilazione dell'utilizzo
  del tempo di CPU da parte del processo.
\end{basedescript}


\subsection{I segnali di I/O asincrono}
\label{sec:sig_asyncio}

Questi segnali operano in congiunzione con le funzioni di I/O asincrono. Per
questo occorre comunque usare \func{fcntl} per abilitare un file descriptor a
generare questi segnali.  L'azione predefinita è di essere ignorati. Questi
segnali sono:
\begin{basedescript}{\desclabelwidth{2.0cm}}
\item[\signald{SIGIO}] Questo segnale viene inviato quando un file descriptor è
  pronto per eseguire dell'input/output. In molti sistemi solo i socket e i
  terminali possono generare questo segnale, in Linux questo può essere usato
  anche per i file, posto che la chiamata a \func{fcntl} che lo attiva abbia
  avuto successo.

\item[\signald{SIGURG}] Questo segnale è inviato quando arrivano dei dati
  urgenti o \textit{out-of-band} su di un socket; per maggiori dettagli al
  proposito si veda sez.~\ref{sec:TCP_urgent_data}.

\item[\signald{SIGPOLL}] Questo segnale è definito nella standard POSIX.1-2001,
  ed è equivalente a \signal{SIGIO} che invece deriva da BSD. Su Linux è
  definito per compatibilità con i sistemi System V.
\end{basedescript}


\subsection{I segnali per il controllo di sessione}
\label{sec:sig_job_control}

Questi sono i segnali usati dal controllo delle sessioni e dei processi, il
loro uso è specializzato e viene trattato in maniera specifica nelle sezioni
in cui si trattano gli argomenti relativi.  Questi segnali sono:
\begin{basedescript}{\desclabelwidth{2.0cm}}
\item[\signald{SIGCHLD}] Questo è il segnale mandato al processo padre quando un
  figlio termina o viene fermato. L'azione predefinita è di ignorare il
  segnale, la sua gestione è trattata in sez.~\ref{sec:proc_wait}.

\item[\signald{SIGCLD}] Per Linux questo è solo un segnale identico al
  precedente e definito come sinonimo. Il nome è obsoleto, deriva dalla
  definizione del segnale su System V, ed oggi deve essere evitato.

\item[\signald{SIGCONT}] Il nome sta per \textit{continue}. Il segnale viene
  usato per fare ripartire un programma precedentemente fermato da
  \signal{SIGSTOP}. Questo segnale ha un comportamento speciale, e fa sempre
  ripartire il processo prima della sua consegna. Il comportamento predefinito
  è di fare solo questo; il segnale non può essere bloccato. Si può anche
  installare un gestore, ma il segnale provoca comunque il riavvio del
  processo.
  
  La maggior pare dei programmi non hanno necessità di intercettare il
  segnale, in quanto esso è completamente trasparente rispetto all'esecuzione
  che riparte senza che il programma noti niente. Si possono installare dei
  gestori per far sì che un programma produca una qualche azione speciale
  se viene fermato e riavviato, come per esempio riscrivere un prompt, o
  inviare un avviso. 

\item[\signald{SIGSTOP}] Il segnale ferma l'esecuzione di un processo, lo porta
  cioè nello stato \textit{stopped} (vedi sez.~\ref{sec:proc_sched}). Il
  segnale non può essere né intercettato, né ignorato, né bloccato.

\item[\signald{SIGTSTP}] Il nome sta per \textit{interactive stop}. Il segnale
  ferma il processo interattivamente, ed è generato dal carattere
  ``\textit{SUSP}'', prodotto dalla combinazione di tasti \cmd{C-z}, ed al
  contrario di \signal{SIGSTOP} può essere intercettato e ignorato. In genere
  un programma installa un gestore per questo segnale quando vuole lasciare il
  sistema o il terminale in uno stato definito prima di fermarsi; se per
  esempio un programma ha disabilitato l'eco sul terminale può installare un
  gestore per riabilitarlo prima di fermarsi.

\item[\signald{SIGTTIN}] Un processo non può leggere dal terminale se esegue
  una sessione di lavoro in \textit{background}. Quando un processo in
  \textit{background} tenta di leggere da un terminale viene inviato questo
  segnale a tutti i processi della sessione di lavoro. L'azione predefinita è
  di fermare il processo.  L'argomento è trattato in
  sez.~\ref{sec:sess_job_control_overview}.

\item[\signald{SIGTTOU}] Segnale analogo al precedente \signal{SIGTTIN}, ma
  generato quando si tenta di scrivere sul terminale o modificarne uno dei
  modi con un processo in \textit{background}. L'azione predefinita è di
  fermare il processo, l'argomento è trattato in
  sez.~\ref{sec:sess_job_control_overview}.
\end{basedescript}


\subsection{I segnali di operazioni errate}
\label{sec:sig_oper_error}

Questi segnali sono usati per riportare al programma errori generati da
operazioni da lui eseguite; non indicano errori del programma quanto errori
che impediscono il completamento dell'esecuzione dovute all'interazione con il
resto del sistema.  L'azione predefinita di questi segnali è normalmente
quella di terminare il processo, questi segnali sono:
\begin{basedescript}{\desclabelwidth{2.0cm}}
\item[\signald{SIGPIPE}] Sta per \textit{Broken pipe}. Se si usano delle
  \textit{pipe}, (o delle FIFO o dei socket) è necessario, prima che un
  processo inizi a scrivere su una di esse, che un altro l'abbia aperta in
  lettura (si veda sez.~\ref{sec:ipc_pipes}). Se il processo in lettura non è
  partito o è terminato inavvertitamente alla scrittura sulla \textit{pipe} il
  kernel genera questo segnale. Se il segnale è bloccato, intercettato o
  ignorato la chiamata che lo ha causato fallisce, restituendo l'errore
  \errcode{EPIPE}.

\item[\signald{SIGXCPU}] Sta per \textit{CPU time limit exceeded}. Questo
  segnale è generato quando un processo eccede il limite impostato per il
  tempo di CPU disponibile, vedi sez.~\ref{sec:sys_resource_limit}. Fino al
  kernel 2.2 terminava semplicemente il processo, a partire dal kernel 2.4,
  seguendo le indicazioni dello standard POSIX.1-2001 viene anche generato un
  \textit{core dump}.

\item[\signald{SIGXFSZ}] Sta per \textit{File size limit exceeded}. Questo
  segnale è generato quando un processo tenta di estendere un file oltre le
  dimensioni specificate dal limite impostato per le dimensioni massime di un
  file, vedi sez.~\ref{sec:sys_resource_limit}.  Fino al kernel 2.2 terminava
  semplicemente il processo, a partire dal kernel 2.4, seguendo le indicazioni
  dello standard POSIX.1-2001 viene anche generato un \textit{core dump}.

\item[\signald{SIGLOST}] Sta per \textit{Resource lost}. Tradizionalmente è il
  segnale che viene generato quando si perde un advisory lock su un file su
  NFS perché il server NFS è stato riavviato. Il progetto GNU lo utilizza per
  indicare ad un client il crollo inaspettato di un server. In Linux è
  definito come sinonimo di \signal{SIGIO} e non viene più usato.
\end{basedescript}


\subsection{Ulteriori segnali}
\label{sec:sig_misc_sig}

Raccogliamo qui infine una serie di segnali che hanno scopi differenti non
classificabili in maniera omogenea. Questi segnali sono:
\begin{basedescript}{\desclabelwidth{2.0cm}}
\item[\signald{SIGUSR1}] Insieme a \signal{SIGUSR2} è un segnale a disposizione
  dell'utente che lo può usare per quello che vuole. Viene generato solo
  attraverso l'invocazione della funzione \func{kill}. Entrambi i segnali
  possono essere utili per implementare una comunicazione elementare fra
  processi diversi, o per eseguire a richiesta una operazione utilizzando un
  gestore. L'azione predefinita è di terminare il processo.
\item[\signald{SIGUSR2}] È il secondo segnale a disposizione degli utenti. Per
  il suo utilizzo vale esattamente quanto appena detto per \signal{SIGUSR1}.
\item[\signald{SIGWINCH}] Il nome sta per \textit{window (size) change} e viene
  generato in molti sistemi (GNU/Linux compreso) quando le dimensioni (in
  righe e colonne) di un terminale vengono cambiate. Viene usato da alcuni
  programmi testuali per riformattare l'uscita su schermo quando si cambia
  dimensione a quest'ultimo. L'azione predefinita è di essere ignorato.
\item[\signald{SIGINFO}] Il segnale indica una richiesta di informazioni. È
  usato con il controllo di sessione, causa la stampa di informazioni da parte
  del processo leader del gruppo associato al terminale di controllo, gli
  altri processi lo ignorano. Su Linux però viene utilizzato come sinonimo di
  \signal{SIGPWR} e l'azione predefinita è di terminare il processo.
\item[\signald{SIGPWR}] Il segnale indica un cambio nello stato di
  alimentazione di un eventuale gruppo di continuità e viene usato
  principalmente per segnalare l'assenza ed il ritorno della corrente. Viene
  usato principalmente con \cmd{init} per attivare o fermare le procedure di
  spegnimento automatico all'esaurimento delle batterie. L'azione predefinita
  è di terminare il processo.
\item[\signald{SIGSTKFLT}] Indica un errore nello stack del coprocessore
  matematico, è definito solo per le architetture PC, ma è completamente
  inusato. L'azione predefinita è di terminare il processo.
\end{basedescript}


\subsection{Le funzioni \func{strsignal} e \func{psignal}}
\label{sec:sig_strsignal}

Per la descrizione dei segnali il sistema mette a disposizione due funzioni
che stampano un messaggio di descrizione specificando il numero del segnale
con una delle costanti di tab.~\ref{tab:sig_signal_list}.  In genere si usano
quando si vuole notificare all'utente il segnale ricevuto, ad esempio nel caso
di terminazione di un processo figlio o di un gestore che gestisce più
segnali.

La prima funzione, \funcd{strsignal}, è una estensione GNU fornita dalla
\acr{glibc}, ed è accessibile solo avendo definito la macro
\macro{\_GNU\_SOURCE}, il suo comportamento è analogo a quello della funzione
\func{strerror} (si veda sez.~\ref{sec:sys_strerror}) usata per notificare gli
errori:

\begin{funcproto}{
\fhead{string.h}
\fdecl{char *strsignal(int signum)} 
\fdesc{Ottiene la descrizione di un segnale.} 
}

{La funzione ritorna puntatore ad una stringa che descrive il segnale, non
  sono previste condizioni di errore ed \var{errno} non viene modificata.}
\end{funcproto}


La funzione ritorna sempre il puntatore ad una stringa che contiene la
descrizione del segnale indicato dall'argomento \param{signum}, se questo non
indica un segnale valido viene restituito il puntatore ad una stringa che
segnale che il valore indicato non è valido.  Dato che la stringa è allocata
staticamente non se ne deve modificare il contenuto, che resta valido solo
fino alla successiva chiamata di \func{strsignal}. Nel caso si debba mantenere
traccia del messaggio sarà necessario copiarlo.

La seconda funzione, \funcd{psignal}, deriva da BSD ed è analoga alla funzione
\func{perror} descritta in sez.~\ref{sec:sys_strerror}, il suo prototipo è:

\begin{funcproto}{
\fhead{signal.h}
\fdecl{void psignal(int sig, const char *s)}
\fdesc{Stampa un messaggio di descrizione di un segnale.} 
}
{La funzione non ritorna nulla e non prevede errori.}  
\end{funcproto}

La funzione stampa sullo \textit{standard error} un messaggio costituito dalla
stringa passata nell'argomento \param{s}, seguita dal carattere di due punti
ed una descrizione del segnale indicato dall'argomento \param{sig}. 

Una modalità alternativa per utilizzare le descrizioni restituite da
\func{strsignal} e \func{psignal} è quello di usare la variabile globale
\var{sys\_siglist}, che è definita in \headfile{signal.h} e può essere
acceduta con la dichiarazione:
\includecodesnip{listati/siglist.c}

L'array \var{sys\_siglist} contiene i puntatori alle stringhe di descrizione,
indicizzate per numero di segnale, per cui una chiamata del tipo di \code{char
  *decr = strsignal(SIGINT)} può essere sostituita dall'equivalente \code{char
  *decr = sys\_siglist[SIGINT]}.



\section{La gestione di base dei segnali}
\label{sec:sig_management}

I segnali sono il primo e più classico esempio di eventi asincroni, cioè di
eventi che possono accadere in un qualunque momento durante l'esecuzione di un
programma. Per questa loro caratteristica la loro gestione non può essere
effettuata all'interno del normale flusso di esecuzione dello stesso, ma è
delegata appunto agli eventuali gestori che si sono installati.

In questa sezione vedremo come si effettua la gestione dei segnali, a partire
dalla loro interazione con le \textit{system call}, passando per le varie
funzioni che permettono di installare i gestori e controllare le reazioni di
un processo alla loro occorrenza.


\subsection{Il comportamento generale del sistema}
\label{sec:sig_gen_beha}

Abbiamo già trattato in sez.~\ref{sec:sig_intro} le modalità con cui il
sistema gestisce l'interazione fra segnali e processi, ci resta da esaminare
però il comportamento delle \textit{system call}; in particolare due di esse,
\func{fork} ed \func{exec}, dovranno essere prese esplicitamente in
considerazione, data la loro stretta relazione con la creazione di nuovi
processi.

Come accennato in sez.~\ref{sec:proc_fork} quando viene creato un nuovo
processo esso eredita dal padre sia le azioni che sono state impostate per i
singoli segnali, che la maschera dei segnali bloccati (vedi
sez.~\ref{sec:sig_sigmask}).  Invece tutti i segnali pendenti e gli allarmi
vengono cancellati; essi infatti devono essere recapitati solo al padre, al
figlio dovranno arrivare solo i segnali dovuti alle sue azioni.

Quando si mette in esecuzione un nuovo programma con \func{exec} (si ricordi
quanto detto in sez.~\ref{sec:proc_exec}) tutti i segnali per i quali è stato
installato un gestore vengono reimpostati a \constd{SIG\_DFL}. Non ha più
senso infatti fare riferimento a funzioni definite nel programma originario,
che non sono presenti nello spazio di indirizzi del nuovo programma.

Si noti che questo vale solo per le azioni per le quali è stato installato un
gestore, viene mantenuto invece ogni eventuale impostazione dell'azione a
\constd{SIG\_IGN}. Questo permette ad esempio alla shell di impostare ad
\const{SIG\_IGN} le risposte per \signal{SIGINT} e \signal{SIGQUIT} per i
programmi eseguiti in background, che altrimenti sarebbero interrotti da una
successiva pressione di \texttt{C-c} o \texttt{C-y}.

Per quanto riguarda il comportamento di tutte le altre \textit{system call} si
danno sostanzialmente due casi, a seconda che esse siano \textsl{lente}
(\textit{slow}) o \textsl{veloci} (\textit{fast}). La gran parte di esse
appartiene a quest'ultima categoria, che non è influenzata dall'arrivo di un
segnale. Esse sono dette \textsl{veloci} in quanto la loro esecuzione è
sostanzialmente immediata. La risposta al segnale viene sempre data dopo che
la \textit{system call} è stata completata, in quanto attendere per eseguire
un gestore non comporta nessun inconveniente.

\index{system~call~lente|(}

In alcuni casi però alcune \textit{system call} possono bloccarsi
indefinitamente e per questo motivo vengono chiamate \textsl{lente} o
\textsl{bloccanti}. In questo caso non si può attendere la conclusione della
\textit{system call}, perché questo renderebbe impossibile una risposta pronta
al segnale, per cui il gestore viene eseguito prima che la \textit{system
  call} sia ritornata.  Un elenco dei casi in cui si presenta questa
situazione è il seguente:
\begin{itemize*}
\item la lettura da file che possono bloccarsi in attesa di dati non ancora
  presenti (come per certi file di dispositivo, i socket o le \textit{pipe});
\item la scrittura sugli stessi file, nel caso in cui dati non possano essere
  accettati immediatamente (di nuovo comune per i socket);
\item l'apertura di un file di dispositivo che richiede operazioni non
  immediate per una risposta (ad esempio l'apertura di un nastro che deve
  essere riavvolto);
\item le operazioni eseguite con \func{ioctl} che non è detto possano essere
  eseguite immediatamente;
\item l'uso di funzioni di intercomunicazione fra processi (vedi
  cap.~\ref{cha:IPC}) che si bloccano in attesa di risposte da altri processi;
\item l'uso della funzione \func{pause} (vedi sez.~\ref{sec:sig_pause_sleep})
  e le analoghe \func{sigsuspend}, \func{sigtimedwait}, e \func{sigwaitinfo}
  (vedi sez.~\ref{sec:sig_real_time}), usate appunto per attendere l'arrivo di
  un segnale;
\item l'uso delle funzioni associate al \textit{file locking} (vedi
  sez.~\ref{sec:file_locking})
\item l'uso della funzione \func{wait} e le analoghe funzioni di attesa se
  nessun processo figlio è ancora terminato.
\end{itemize*}

In questo caso si pone il problema di cosa fare una volta che il gestore sia
ritornato. La scelta originaria dei primi Unix era quella di far ritornare
anche la \textit{system call} restituendo l'errore di \errcode{EINTR}. Questa
è a tutt'oggi una scelta corrente, ma comporta che i programmi che usano dei
gestori controllino lo stato di uscita delle funzioni che eseguono una system
call lenta per ripeterne la chiamata qualora l'errore fosse questo.

Dimenticarsi di richiamare una \textit{system call} interrotta da un segnale è
un errore comune, tanto che la \acr{glibc} provvede una macro
\code{TEMP\_FAILURE\_RETRY(expr)} che esegue l'operazione automaticamente,
ripetendo l'esecuzione dell'espressione \var{expr} fintanto che il risultato
non è diverso dall'uscita con un errore \errcode{EINTR}.

La soluzione è comunque poco elegante e BSD ha scelto un approccio molto
diverso, che è quello di fare ripartire automaticamente una \textit{system
  call} interrotta invece di farla fallire. In questo caso ovviamente non c'è
bisogno di preoccuparsi di controllare il codice di errore; si perde però la
possibilità di eseguire azioni specifiche all'occorrenza di questa particolare
condizione.

Linux e la \acr{glibc} consentono di utilizzare entrambi gli approcci,
attraverso una opportuna opzione di \func{sigaction} (vedi
sez.~\ref{sec:sig_sigaction}). È da chiarire comunque che nel caso di
interruzione nel mezzo di un trasferimento parziale di dati, le \textit{system
  call} ritornano sempre indicando i byte trasferiti.

Si tenga presente però che alcune \textit{system call} vengono comunque
interrotte con un errore di \errcode{EINTR} indipendentemente dal fatto che ne
possa essere stato richiesto il riavvio automatico, queste funzioni sono:

\begin{itemize*}
\item le funzioni di attesa di un segnale: \func{pause} (vedi
  sez.~\ref{sec:sig_pause_sleep}) o \func{sigsuspend}, \func{sigtimedwait}, e
  \func{sigwaitinfo} (vedi sez.~\ref{sec:sig_real_time}).
\item le funzioni di attesa dell'\textit{I/O multiplexing} (vedi
  sez.~\ref{sec:file_multiplexing}) come \func{select}, \func{pselect},
  \func{poll}, \func{ppoll}, \func{epoll\_wait} e \func{epoll\_pwait}.
\item le funzioni del System V IPC che prevedono attese: \func{msgrcv},
  \func{msgsnd} (vedi sez.~\ref{sec:ipc_sysv_mq}), \func{semop} e
  \func{semtimedop} (vedi sez.~\ref{sec:ipc_sysv_sem}).
\item le funzioni per la messa in attesa di un processo come \func{usleep},
  \func{nanosleep} (vedi sez.~\ref{sec:sig_pause_sleep}) e
  \func{clock\_nanosleep} (vedi sez.~\ref{sec:sig_timer_adv}).
\item le funzioni che operano sui socket quando è stato impostato un
  \textit{timeout} sugli stessi con \func{setsockopt} (vedi
  sez.~\ref{sec:sock_generic_options}) ed in particolare \func{accept},
  \func{recv}, \func{recvfrom}, \func{recvmsg} per un \textit{timeout} in
  ricezione e \func{connect}, \func{send}, \func{sendto} e \func{sendmsg} per
  un \textit{timeout} in trasmissione.
%\item la funzione \func{io\_getevents} per l'I/O asincrono (vedi sez.??)
\end{itemize*}


\index{system~call~lente|)}


\subsection{L'installazione di un gestore}
\label{sec:sig_signal}

L'interfaccia più semplice per la gestione dei segnali è costituita dalla
funzione di sistema \funcd{signal} che è definita fin dallo standard ANSI C.
Quest'ultimo però non considera sistemi multitasking, per cui la definizione è
tanto vaga da essere del tutto inutile in un sistema Unix. Per questo motivo
ogni implementazione successiva ne ha modificato e ridefinito il
comportamento, pur mantenendone immutato il prototipo\footnote{in realtà in
  alcune vecchie implementazioni (SVr4 e 4.3+BSD in particolare) vengono usati
  alcuni argomenti aggiuntivi per definire il comportamento della funzione,
  vedremo in sez.~\ref{sec:sig_sigaction} che questo è possibile usando la
  funzione \func{sigaction}.}  che è:

\begin{funcproto}{
\fhead{signal.h}
\fdecl{sighandler\_t signal(int signum, sighandler\_t handler)}
\fdesc{Installa un gestore di segnale (\textit{signal handler}).} 
}

{La funzione ritorna il precedente gestore in caso di successo in caso di
  successo e \constd{SIG\_ERR} per un errore, nel qual caso \var{errno}
  assumerà il valore:
  \begin{errlist}
  \item[\errcode{EINVAL}] il numero di segnale \param{signum} non è valido.
  \end{errlist}
}  
\end{funcproto}

In questa definizione per l'argomento \param{handler} che indica il gestore da
installare si è usato un tipo di dato, \type{sighandler\_t}, che è una
estensione GNU, definita dalla \acr{glibc}, che permette di riscrivere il
prototipo di \func{signal} nella forma appena vista, molto più leggibile di
quanto non sia la versione originaria, che di norma è definita come:
\includecodesnip{listati/signal.c}
questa infatti, per la poca chiarezza della sintassi del C quando si vanno a
trattare puntatori a funzioni, è molto meno comprensibile.  Da un confronto
con il precedente prototipo si può dedurre la definizione di
\typed{sighandler\_t} che è:
\includecodesnip{listati/sighandler_t.c}
e cioè un puntatore ad una funzione \ctyp{void} (cioè senza valore di ritorno)
e che prende un argomento di tipo \ctyp{int}. Si noti come si devono usare le
parentesi intorno al nome della funzione per via delle precedenze degli
operatori del C, senza di esse si sarebbe definita una funzione che ritorna un
puntatore a \ctyp{void} e non un puntatore ad una funzione \ctyp{void}.

La funzione \func{signal} quindi restituisce e prende come secondo argomento
un puntatore a una funzione di questo tipo, che è appunto la funzione che
verrà usata come gestore del segnale.  Il numero di segnale passato
nell'argomento \param{signum} può essere indicato direttamente con una delle
costanti definite in sez.~\ref{sec:sig_standard}.  

L'argomento \param{handler} che indica il gestore invece, oltre all'indirizzo
della funzione da chiamare all'occorrenza del segnale, può assumere anche i
due valori costanti \const{SIG\_IGN} e \const{SIG\_DFL}. Il primo indica che
il segnale deve essere ignorato. Il secondo ripristina l'azione predefinita, e
serve a tornare al comportamento di default quando non si intende più gestire
direttamente un segnale. Si ricordi però che i due segnali \signal{SIGKILL} e
\signal{SIGSTOP} non possono essere né ignorati né intercettati e per loro
l'uso di \func{signal} non ha alcun effetto, qualunque cosa si specifichi
per \param{handler}.

La funzione restituisce l'indirizzo dell'azione precedente, che può essere
salvato per poterlo ripristinare (con un'altra chiamata a \func{signal}) in un
secondo tempo. Si ricordi che se si imposta come azione \const{SIG\_IGN} o si
imposta \const{SIG\_DFL} per un segnale la cui azione predefinita è di essere
ignorato, tutti i segnali pendenti saranno scartati, e non verranno mai
notificati.

L'uso di \func{signal} è soggetto a problemi di compatibilità, dato che essa
si comporta in maniera diversa per sistemi derivati da BSD o da System V. In
questi ultimi infatti la funzione è conforme al comportamento originale dei
primi Unix in cui il gestore viene disinstallato alla sua chiamata, secondo la
semantica inaffidabile; anche Linux seguiva questa convenzione con le vecchie
librerie del C come la \acr{libc4} e la \acr{libc5}.\footnote{nelle
  \acr{libc5} esiste però la possibilità di includere \file{bsd/signal.h} al
  posto di \headfile{signal.h}, nel qual caso la funzione \func{signal} viene
  ridefinita per seguire la semantica affidabile usata da BSD.}

Al contrario BSD segue la semantica affidabile, non disinstallando il gestore
e bloccando il segnale durante l'esecuzione dello stesso. Con l'utilizzo della
\acr{glibc} dalla versione 2 anche Linux è passato a questo comportamento.  Il
comportamento della versione originale della funzione, il cui uso è deprecato
per i motivi visti in sez.~\ref{sec:sig_semantics}, può essere ottenuto
chiamando \funcm{sysv\_signal}, una volta che si sia definita la macro
\macro{\_XOPEN\_SOURCE}.  In generale, per evitare questi problemi, l'uso di
\func{signal}, che tra l'altro ha un comportamento indefinito in caso di
processo multi-\textit{thread}, è da evitare: tutti i nuovi programmi devono
usare \func{sigaction}.

È da tenere presente che, seguendo lo standard POSIX, il comportamento di un
processo che ignora i segnali \signal{SIGFPE}, \signal{SIGILL}, o
\signal{SIGSEGV}, qualora questi non originino da una chiamata ad una
\func{kill} o altra funzione affine, è indefinito. Un gestore che ritorna da
questi segnali può dare luogo ad un ciclo infinito.


\subsection{Le funzioni per l'invio di segnali}
\label{sec:sig_kill_raise}

Come accennato in sez.~\ref{sec:sig_types} un segnale può anche essere
generato direttamente nell'esecuzione di un programma, attraverso la chiamata
ad una opportuna \textit{system call}. Le funzioni che si utilizzano di solito
per inviare un segnale generico ad un processo sono \func{raise} e
\func{kill}.

La funzione \funcd{raise}, definita dallo standard ANSI C, serve per inviare
un segnale al processo corrente,\footnote{non prevedendo la presenza di un
  sistema multiutente lo standard ANSI C non poteva che definire una funzione
  che invia il segnale al programma in esecuzione, nel caso di Linux questa
  viene implementata come funzione di compatibilità.}  il suo prototipo è:

\begin{funcproto}{
\fhead{signal.h}
\fdecl{int raise(int sig)}
\fdesc{Invia un segnale al processo corrente.} 
}

{La funzione ritorna $0$ in caso di successo e $-1$ per un errore, nel qual
  caso \var{errno} assumerà il valore: 
  \begin{errlist}
  \item[\errcode{EINVAL}] il segnale \param{sig} non è valido.
  \end{errlist}
}
\end{funcproto}

Il valore di \param{sig} specifica il segnale che si vuole inviare e può
essere specificato con una delle costanti illustrate in
tab.~\ref{tab:sig_signal_list}.  In genere questa funzione viene usata per
riprodurre il comportamento predefinito di un segnale che sia stato
intercettato. In questo caso, una volta eseguite le operazioni volute, il
gestore dovrà prima reinstallare l'azione predefinita, per poi attivarla
chiamando \func{raise}.

In realtà \func{raise} è una funzione di libreria, che per i processi ordinari
viene implementata attraverso la funzione di sistema \funcd{kill} che è quella
che consente effettivamente di inviare un segnale generico ad un processo, il
 suo prototipo è:

\begin{funcproto}{
\fhead{sys/types.h}
\fhead{signal.h}
\fdecl{int kill(pid\_t pid, int sig)}
\fdesc{Invia un segnale ad uno o più processi.} 
}

{La funzione ritorna $0$ in caso di successo e $-1$ per un errore, nel qual
  caso \var{errno} assumerà uno dei valori: 
  \begin{errlist}
    \item[\errcode{EINVAL}] il segnale specificato non esiste.
    \item[\errcode{EPERM}] non si hanno privilegi sufficienti ad inviare il
      segnale.
    \item[\errcode{ESRCH}] il processo o il gruppo di processi indicato non
      esiste.
  \end{errlist}
}
\end{funcproto}

La funzione invia il segnale specificato dall'argomento \param{sig} al
processo o ai processi specificati con l'argomento \param{pid}.  Lo standard
POSIX prevede che il valore 0 per \param{sig} sia usato per specificare il
segnale nullo.  Se la funzione viene chiamata con questo valore non viene
inviato nessun segnale, ma viene eseguito il controllo degli errori, in tal
caso si otterrà un errore \errcode{EPERM} se non si hanno i permessi necessari
ed un errore \errcode{ESRCH} se il processo o i processi specificati
con \param{pid} non esistono.

\begin{table}[htb]
  \footnotesize
  \centering
  \begin{tabular}[c]{|r|p{8cm}|}
    \hline
    \textbf{Valore} & \textbf{Significato} \\
    \hline
    \hline
    $>0$ & Il segnale è mandato al processo con \ids{PID} uguale
           a \param{pid}.\\
    0    & Il segnale è mandato ad ogni processo del \textit{process group}
           del chiamante.\\  
    $-1$ & Il segnale è mandato ad ogni processo (eccetto \cmd{init}).\\
    $<-1$& Il segnale è mandato ad ogni processo del \textit{process group} 
           con \ids{PGID} uguale a $|\param{pid}|$.\\
    \hline
  \end{tabular}
  \caption{Valori dell'argomento \param{pid} per la funzione
    \func{kill}.}
  \label{tab:sig_kill_values}
\end{table}

A seconda del valore dell'argomento \param{pid} si può inviare il segnale ad
uno specifico processo, ad un \textit{process group} (vedi
sez.~\ref{sec:sess_proc_group}) o a tutti i processi, secondo quanto
illustrato in tab.~\ref{tab:sig_kill_values} che riporta i valori possibili
per questo argomento. Si tenga conto però che il sistema ricicla i \ids{PID}
(come accennato in sez.~\ref{sec:proc_pid}) per cui l'esistenza di un processo
non significa che esso sia realmente quello a cui si intendeva mandare il
segnale.

Indipendentemente dalla funzione specifica che viene usata solo
l'amministratore può inviare un segnale ad un processo qualunque, in tutti gli
altri casi l'\ids{UID} reale o l'\ids{UID} effettivo del processo chiamante
devono corrispondere all'\ids{UID} reale o all'\ids{UID} salvato della
destinazione. Fa eccezione il caso in cui il segnale inviato sia
\signal{SIGCONT}, nel quale occorre anche che entrambi i processi appartengano
alla stessa sessione.

Si tenga presente che, per il ruolo fondamentale che riveste nel sistema, non
è possibile inviare al processo 1 (cioè a \cmd{init}) segnali per i quali esso
non abbia un gestore installato.  Infine, seguendo le specifiche POSIX
1003.1-2001, l'uso della chiamata \code{kill(-1, sig)} comporta che il segnale
sia inviato (con la solita eccezione di \cmd{init}) a tutti i processi per i
quali i permessi lo consentano. Lo standard permette comunque alle varie
implementazioni di escludere alcuni processi specifici: nel caso in questione
Linux non invia il segnale al processo che ha effettuato la chiamata.

Si noti pertanto che la funzione \code{raise(sig)} può essere definita in
termini di \func{kill}, ed è sostanzialmente equivalente ad una
\code{kill(getpid(), sig)}. Siccome \func{raise}, che è definita nello
standard ISO C, non esiste in alcune vecchie versioni di Unix, in generale
l'uso di \func{kill} finisce per essere più portabile.  Una seconda funzione
che può essere definita in termini di \func{kill} è \funcd{killpg}, il suo
prototipo è:

\begin{funcproto}{
\fhead{signal.h}
\fdecl{int killpg(pid\_t pidgrp, int signal)}
\fdesc{Invia un segnale ad un \textit{process group}.} 
}

{ La funzione ritorna $0$ in caso di successo e $-1$ per un errore, e gli
  errori sono gli stessi di \func{kill}.
}
\end{funcproto}


La funzione invia il segnale \param{signal} al \textit{process group} il cui
\acr{PGID} (vedi sez.~\ref{sec:sess_proc_group}) è indicato
dall'argomento \param{pidgrp}, che deve essere un intero positivo. Il suo
utilizzo è sostanzialmente equivalente all'esecuzione di \code{kill(-pidgrp,
  signal)}.

Oltre alle precedenti funzioni di base, vedremo più avanti che esistono altre
funzioni per inviare segnali generici, come \func{sigqueue} per i segnali
\textit{real-time} (vedi sez.~\ref{sec:sig_real_time}) e le specifiche
funzioni per i \textit{thread} che tratteremo in sez.~\ref{sec:thread_signal}.

Esiste però un'ultima funzione che permette l'invio diretto di un segnale che
vale la pena di trattare a parte per le sue peculiarità. La funzione in
questione è \funcd{abort} che, come accennato in
sez.~\ref{sec:proc_termination}, permette di abortire l'esecuzione di un
programma tramite l'invio del segnale \signal{SIGABRT}. Il suo prototipo è:

\begin{funcproto}{
\fhead{stdlib.h}
\fdecl{void abort(void)}
\fdesc{Abortisce il processo corrente.} 
}

{La funzione non ritorna, il processo viene terminato.}
\end{funcproto}

La differenza fra questa funzione e l'uso di \func{raise} o di un'altra
funzione per l'invio di \signal{SIGABRT} è che anche se il segnale è bloccato
o ignorato, la funzione ha effetto lo stesso. Il segnale può però essere
intercettato per effettuare eventuali operazioni di chiusura prima della
terminazione del processo.

Lo standard ANSI C richiede inoltre che anche se il gestore ritorna, la
funzione non ritorni comunque. Lo standard POSIX.1 va oltre e richiede che se
il processo non viene terminato direttamente dal gestore sia la stessa
\func{abort} a farlo al ritorno dello stesso. Inoltre, sempre seguendo lo
standard POSIX, prima della terminazione tutti i file aperti e gli stream
saranno chiusi ed i buffer scaricati su disco. Non verranno invece eseguite le
eventuali funzioni registrate con \func{atexit} e \func{on\_exit}.

% TODO trattare pidfd_send_signal, aggiunta con il kernel 5.1 (vedi
% https://lwn.net/Articles/783052/) per mandare segnali a processi senza dover
% usare un PID, vedi anche https://lwn.net/Articles/773459/,
% https://git.kernel.org/linus/3eb39f47934f 


\subsection{Le funzioni di allarme ed i \textit{timer}}
\label{sec:sig_alarm_abort}

Un caso particolare di segnali generati a richiesta è quello che riguarda i
vari segnali usati per la temporizzazione, per ciascuno di essi infatti sono
previste delle funzioni specifiche che ne effettuino l'invio. La più comune, e
la più semplice, delle funzioni usate per la temporizzazione è la funzione di
sistema \funcd{alarm}, il cui prototipo è:

\begin{funcproto}{
\fhead{unistd.h}
\fdecl{unsigned int alarm(unsigned int seconds)}
\fdesc{Predispone l'invio di un allarme.} 
}

{La funzione ritorna il numero di secondi rimanenti ad un precedente allarme,
  o $0$ se non c'erano allarmi pendenti, non sono previste condizioni di
  errore.}
\end{funcproto}

La funzione fornisce un meccanismo che consente ad un processo di predisporre
un'interruzione nel futuro, ad esempio per effettuare una qualche operazione
dopo un certo periodo di tempo, programmando l'emissione di un segnale (nel
caso in questione \signal{SIGALRM}) dopo il numero di secondi specificato
dall'argomento \param{seconds}. Se si specifica per \param{seconds} un valore
nullo non verrà inviato nessun segnale. Siccome alla chiamata viene cancellato
ogni precedente allarme, questo valore può essere usato per cancellare una
programmazione precedente.

La funzione inoltre ritorna il numero di secondi rimanenti all'invio
dell'allarme programmato in precedenza. In questo modo è possibile controllare
se non si è cancellato un precedente allarme e predisporre eventuali misure
che permettano di gestire il caso in cui servono più interruzioni.

In sez.~\ref{sec:sys_unix_time} abbiamo visto che ad ogni processo sono
associati tre tempi diversi: il \textit{clock time}, l'\textit{user time} ed
il \textit{system time}.  Per poterli calcolare il kernel mantiene per ciascun
processo tre diversi timer:
\begin{itemize*}
\item un \textit{real-time timer} che calcola il tempo reale trascorso (che
  corrisponde al \textit{clock time}). La scadenza di questo timer provoca
  l'emissione di \signal{SIGALRM};
\item un \textit{virtual timer} che calcola il tempo di processore usato dal
  processo in \textit{user space} (che corrisponde all'\textit{user time}). La
  scadenza di questo timer provoca l'emissione di \signal{SIGVTALRM};
\item un \textit{profiling timer} che calcola la somma dei tempi di processore
  utilizzati direttamente dal processo in \textit{user space}, e dal kernel
  nelle \textit{system call} ad esso relative (che corrisponde a quello che in
  sez.~\ref{sec:sys_unix_time} abbiamo chiamato \textit{processor time}). La
  scadenza di questo timer provoca l'emissione di \signal{SIGPROF}.
\end{itemize*}

Il timer usato da \func{alarm} è il \textit{clock time}, e corrisponde cioè al
tempo reale. La funzione come abbiamo visto è molto semplice, ma proprio per
questo presenta numerosi limiti: non consente di usare gli altri timer, non
può specificare intervalli di tempo con precisione maggiore del secondo e
genera il segnale una sola volta.

Per ovviare a questi limiti Linux deriva da BSD la funzione \funcd{setitimer}
che permette di usare un timer qualunque e l'invio di segnali periodici, al
costo però di una maggiore complessità d'uso e di una minore portabilità. Il
suo prototipo è:

\begin{funcproto}{
\fhead{sys/time.h}
\fdecl{int setitimer(int which, const struct itimerval *value, struct
  itimerval *ovalue)}
  
\fdesc{Predispone l'invio di un segnale di allarme.} 
}

{La funzione ritorna $0$ in caso di successo e $-1$ per un errore, nel qual
  caso \var{errno} assumerà uno dei valori \errval{EINVAL} o \errval{EFAULT}
  nel loro significato generico.}
\end{funcproto}


La funzione predispone l'invio di un segnale di allarme alla scadenza
dell'intervallo indicato dall'argomento \param{value}.  Il valore
dell'argomento \param{which} permette di specificare quale dei tre timer
illustrati in precedenza usare; i possibili valori sono riportati in
tab.~\ref{tab:sig_setitimer_values}.
\begin{table}[htb]
  \footnotesize
  \centering
  \begin{tabular}[c]{|l|l|}
    \hline
    \textbf{Valore} & \textbf{Timer} \\
    \hline
    \hline
    \constd{ITIMER\_REAL}    & \textit{real-time timer}\\
    \constd{ITIMER\_VIRTUAL} & \textit{virtual timer}\\
    \constd{ITIMER\_PROF}    & \textit{profiling timer}\\
    \hline
  \end{tabular}
  \caption{Valori dell'argomento \param{which} per la funzione
    \func{setitimer}.}
  \label{tab:sig_setitimer_values}
\end{table}

Il valore della struttura specificata \param{value} viene usato per impostare
il timer, se il puntatore \param{ovalue} non è nullo il precedente valore
viene salvato qui. I valori dei timer devono essere indicati attraverso una
struttura \struct{itimerval}, definita in fig.~\ref{fig:file_stat_struct}.

La struttura è composta da due membri, il primo, \var{it\_interval} definisce
il periodo del timer; il secondo, \var{it\_value} il tempo mancante alla
scadenza. Entrambi esprimono i tempi tramite una struttura \struct{timeval} che
permette una precisione fino al microsecondo.

Ciascun timer decrementa il valore di \var{it\_value} fino a zero, poi invia
il segnale e reimposta \var{it\_value} al valore di \var{it\_interval}, in
questo modo il ciclo verrà ripetuto; se invece il valore di \var{it\_interval}
è nullo il timer si ferma.

\begin{figure}[!htb]
  \footnotesize \centering
  \begin{minipage}[c]{0.8\textwidth}
    \includestruct{listati/itimerval.h}
  \end{minipage} 
  \normalsize 
  \caption{La struttura \structd{itimerval}, che definisce i valori dei timer
    di sistema.}
  \label{fig:sig_itimerval}
\end{figure}

L'uso di \func{setitimer} consente dunque un controllo completo di tutte le
caratteristiche dei timer, ed in effetti la stessa \func{alarm}, benché
definita direttamente nello standard POSIX.1, può a sua volta essere espressa
in termini di \func{setitimer}, come evidenziato dal manuale della \acr{glibc}
\cite{GlibcMan} che ne riporta la definizione mostrata in
fig.~\ref{fig:sig_alarm_def}.\footnote{questo comporta anche che non è il caso
  di mescolare chiamate ad \func{abort} e a \func{setitimer}.}

\begin{figure}[!htb]
  \footnotesize \centering
  \begin{minipage}[c]{0.8\textwidth}
    \includestruct{listati/alarm_def.c}
  \end{minipage} 
  \normalsize 
  \caption{Definizione di \func{alarm} in termini di \func{setitimer}.} 
  \label{fig:sig_alarm_def}
\end{figure}

Si deve comunque tenere presente che fino al kernel 2.6.16 la precisione di
queste funzioni era limitata dalla frequenza del timer di sistema, determinato
dal valore della costante \texttt{HZ} di cui abbiamo già parlato in
sez.~\ref{sec:proc_hierarchy}, in quanto le temporizzazioni erano calcolate in
numero di interruzioni del timer (i cosiddetti ``\textit{jiffies}''), ed era
assicurato soltanto che il segnale non sarebbe stato mai generato prima della
scadenza programmata (l'arrotondamento cioè era effettuato per
eccesso).\footnote{questo in realtà non è del tutto vero a causa di un bug,
  presente fino al kernel 2.6.12, che in certe circostanze causava l'emissione
  del segnale con un arrotondamento per difetto.}

L'uso del contatore dei \textit{jiffies}, un intero a 32 bit nella maggior
parte dei casi, comportava inoltre l'impossibilità di specificare tempi molto
lunghi. superiori al valore della costante \constd{MAX\_SEC\_IN\_JIFFIES},
pari, nel caso di default di un valore di \const{HZ} di 250, a circa 99 giorni
e mezzo. Con il cambiamento della rappresentazione effettuato nel kernel
2.6.16 questo problema è scomparso e con l'introduzione dei timer ad alta
risoluzione (vedi sez.~\ref{sec:sig_timer_adv}) nel kernel 2.6.21 la
precisione è diventata quella fornita dall'hardware disponibile.

Una seconda causa di potenziali ritardi è che il segnale viene generato alla
scadenza del timer, ma poi deve essere consegnato al processo; se quest'ultimo
è attivo (questo è sempre vero per \const{ITIMER\_VIRTUAL}) la consegna è
immediata, altrimenti può esserci un ulteriore ritardo che può variare a
seconda del carico del sistema.

Questo ha una conseguenza che può indurre ad errori molto subdoli, si tenga
conto poi che in caso di sistema molto carico, si può avere il caso patologico
in cui un timer scade prima che il segnale di una precedente scadenza sia
stato consegnato. In questo caso, per il comportamento dei segnali descritto
in sez.~\ref{sec:sig_sigchld}, un solo segnale sarà consegnato. Per questo
oggi l'uso di questa funzione è deprecato a favore degli
\textit{high-resolution timer} e della cosiddetta \textit{POSIX Timer API},
che tratteremo in sez.~\ref{sec:sig_timer_adv}.

Dato che sia \func{alarm} che \func{setitimer} non consentono di leggere il
valore corrente di un timer senza modificarlo, è possibile usare la funzione
\funcd{getitimer}, il cui prototipo è:

\begin{funcproto}{
\fhead{sys/time.h}
\fdecl{int getitimer(int which, struct itimerval *value)}
\fdesc{Legge il valore di un timer.} 
}

{ La funzione ritorna $0$ in caso di successo e $-1$ per un errore, nel qual
  caso \var{errno} assumerà gli stessi valori di \func{getitimer}.  }
\end{funcproto}

La funzione legge nella struttura \struct{itimerval} puntata da \param{value}
il valore del timer specificato da \param{which} ed i suoi argomenti hanno lo
stesso significato e formato di quelli di \func{setitimer}.


\subsection{Le funzioni di pausa e attesa}
\label{sec:sig_pause_sleep}

Sono parecchie le occasioni in cui si può avere necessità di sospendere
temporaneamente l'esecuzione di un processo. Nei sistemi più elementari in
genere questo veniva fatto con un ciclo di attesa in cui il programma ripete
una operazione un numero sufficiente di volte per far passare il tempo
richiesto.

Ma in un sistema multitasking un ciclo di attesa è solo un inutile spreco di
tempo di processore, dato che altri programmi possono essere eseguiti nel
frattempo, per questo ci sono delle apposite funzioni che permettono di
mantenere un processo in attesa per il tempo voluto, senza impegnare il
processore. In pratica si tratta di funzioni che permettono di portare
esplicitamente il processo nello stato di \textit{sleep} (si ricordi quanto
illustrato in tab.~\ref{tab:proc_proc_states}) per un certo periodo di tempo.

La prima di queste è la funzione di sistema \funcd{pause}, che viene usata per
mettere un processo in attesa per un periodo di tempo indefinito, fino
all'arrivo di un segnale, il suo prototipo è:

\begin{funcproto}{
\fhead{unistd.h}
\fdecl{int pause(void)}
\fdesc{Pone il processo in pausa fino al ricevimento di un segnale.} 
}

{La funzione ritorna solo dopo che un segnale è stato ricevuto ed il relativo
  gestore è ritornato, nel qual caso restituisce $-1$ e \var{errno} assume il
  valore \errval{EINTR}.}
\end{funcproto}

La funzione ritorna sempre con una condizione di errore, dato che il successo
sarebbe quello di continuare ad aspettare indefinitamente. In genere si usa
questa funzione quando si vuole mettere un processo in attesa di un qualche
evento specifico che non è sotto il suo diretto controllo, ad esempio la si
può usare per interrompere l'esecuzione del processo fino all'arrivo di un
segnale inviato da un altro processo.

Quando invece si vuole fare attendere un processo per un intervallo di tempo
già noto in partenza, lo standard POSIX.1 prevede una funzione di attesa
specifica, \funcd{sleep}, il cui prototipo è:

\begin{funcproto}{

\fhead{unistd.h}
\fdecl{unsigned int sleep(unsigned int seconds)}
\fdesc{Pone il processo in pausa per un tempo in secondi.} 
}

{La funzione ritorna $0$ se l'attesa viene completata  o il
  numero di secondi restanti se viene interrotta da un segnale, non sono
  previsti codici di errore.}
\end{funcproto}

La funzione pone il processo in stato di \textit{sleep} per il numero di
secondi specificato dall'argomento \param{seconds}, a meno di non essere
interrotta da un segnale. Alla terminazione del periodo di tempo indicato la
funzione ritorna riportando il processo in stato \textit{runnable} così che
questo possa riprendere l'esecuzione.

In caso di interruzione della funzione non è una buona idea ripetere la
chiamata per il tempo rimanente restituito dalla stessa, in quanto la
riattivazione del processo può avvenire in un qualunque momento, ma il valore
restituito sarà sempre arrotondato al secondo. Questo può avere la conseguenza
che se la successione dei segnali è particolarmente sfortunata e le differenze
si accumulano, si possono avere ritardi anche di parecchi secondi rispetto a
quanto programmato inizialmente. In genere la scelta più sicura in questo caso
è quella di stabilire un termine per l'attesa, e ricalcolare tutte le volte il
numero di secondi che restano da aspettare.

Si tenga presente che alcune implementazioni l'uso di \func{sleep} può avere
conflitti con quello di \signal{SIGALRM}, dato che la funzione può essere
realizzata con l'uso di \func{pause} e \func{alarm}, in una maniera analoga a
quella dell'esempio che vedremo in sez.~\ref{sec:sig_example}. In tal caso
mescolare chiamate di \func{alarm} e \func{sleep} o modificare l'azione
associata \signal{SIGALRM}, può portare a dei risultati indefiniti. Nel caso
della \acr{glibc} è stata usata una implementazione completamente indipendente
e questi problemi non ci sono, ma un programma portabile non può fare questa
assunzione.

La granularità di \func{sleep} permette di specificare attese soltanto in
secondi, per questo sia sotto BSD4.3 che in SUSv2 è stata definita un'altra
funzione con una precisione teorica del microsecondo. I due standard hanno
delle definizioni diverse, ma la \acr{glibc} segue (secondo la pagina di
manuale almeno dalla versione 2.2.2)  quella di SUSv2 per cui la
funzione \funcd{usleep} (dove la \texttt{u} è intesa come sostituzione di
$\mu$), ha il seguente prototipo:

\begin{funcproto}{
\fhead{unistd.h}
\fdecl{int usleep(unsigned long usec)}
\fdesc{Pone il processo in pausa per un tempo in microsecondi.} 
}

{La funzione ritorna $0$ se l'attesa viene completata e $-1$ per un errore,
  nel qual caso \var{errno} assumerà uno dei valori:
  \begin{errlist}
  \item[\errcode{EINTR}] la funzione è stata interrotta da un segnale.
  \item[\errcode{EINVAL}] si è indicato un valore di \param{usec} maggiore di
    1000000.
  \end{errlist}
}
\end{funcproto}

Anche questa funzione, a seconda delle implementazioni, può presentare
problemi nell'interazione con \func{alarm} e \signal{SIGALRM}, per questo
motivo, pur essendovi citata, nello standard POSIX.1-2001 viene deprecata in
favore della nuova funzione di sistema \funcd{nanosleep}, il cui prototipo è:

\begin{funcproto}{
\fhead{unistd.h}
\fdecl{int nanosleep(const struct timespec *req, struct timespec *rem)}
\fdesc{Pone il processo in pausa per un intervallo di tempo.} 
}

{La funzione ritorna $0$ se l'attesa viene completata e $-1$ per un errore,
  nel qual caso \var{errno} assumerà uno dei valori:
  \begin{errlist}
    \item[\errcode{EINTR}] la funzione è stata interrotta da un segnale.
    \item[\errcode{EINVAL}] si è specificato un numero di secondi negativo o un
      numero di nanosecondi maggiore di 999.999.999.
  \end{errlist}
}
\end{funcproto}

La funzione pone il processo in pausa portandolo nello stato di \textit{sleep}
per il tempo specificato dall'argomento \param{req}, ed in caso di
interruzione restituisce il tempo restante nell'argomento \param{rem}.  Lo
standard richiede che la funzione sia implementata in maniera del tutto
indipendente da \func{alarm}, e nel caso di Linux questo è fatto utilizzando
direttamente il timer del kernel. Lo standard richiede inoltre che la funzione
sia utilizzabile senza interferenze con l'uso di \signal{SIGALRM}. La funzione
prende come argomenti delle strutture di tipo \struct{timespec}, la cui
definizione è riportata in fig.~\ref{fig:sys_timespec_struct}, il che permette
di specificare un tempo con una precisione teorica fino al nanosecondo.

La funzione risolve anche il problema di proseguire l'attesa dopo
l'interruzione dovuta ad un segnale; infatti in tal caso in \param{rem} viene
restituito il tempo rimanente rispetto a quanto richiesto
inizialmente,\footnote{con l'eccezione, valida solo nei kernel della serie
  2.4, in cui, per i processi riavviati dopo essere stati fermati da un
  segnale, il tempo passato in stato \texttt{T} non viene considerato nel
  calcolo della rimanenza.} e basta richiamare la funzione per completare
l'attesa.

Anche qui però occorre tenere presente che i tempi sono arrotondati, per cui
la precisione, per quanto migliore di quella ottenibile con \func{sleep}, è
relativa e in caso di molte interruzioni si può avere una deriva, per questo
esiste la funzione \func{clock\_nanosleep} (vedi sez.~\ref{sec:sig_timer_adv})
che permette di specificare un tempo assoluto anziché un tempo relativo.

Chiaramente, anche se il tempo può essere specificato con risoluzioni fino al
nanosecondo, la precisione di \func{nanosleep} è determinata dalla risoluzione
temporale del timer di sistema. Perciò la funzione attenderà comunque il tempo
specificato, ma prima che il processo possa tornare ad essere eseguito
occorrerà almeno attendere la successiva interruzione del timer di sistema,
cioè un tempo che a seconda dei casi può arrivare fino a 1/\const{HZ}, (sempre
che il sistema sia scarico ed il processa venga immediatamente rimesso in
esecuzione). Per questo motivo il valore restituito in \param{rem} è sempre
arrotondato al multiplo successivo di 1/\const{HZ}. 

Con i kernel della serie 2.4 in realtà era possibile ottenere anche pause più
precise del centesimo di secondo usando politiche di \textit{scheduling}
\textit{real-time} come \const{SCHED\_FIFO} o \const{SCHED\_RR} (vedi
sez.~\ref{sec:proc_real_time}); in tal caso infatti il calcolo sul numero di
interruzioni del timer veniva evitato utilizzando direttamente un ciclo di
attesa con cui si raggiungevano pause fino ai 2~ms con precisioni del
$\mu$s. Questa estensione è stata rimossa con i kernel della serie 2.6, che
consentono una risoluzione più alta del timer di sistema; inoltre a partire
dal kernel 2.6.21, \func{nanosleep} può avvalersi del supporto dei timer ad
alta risoluzione, ottenendo la massima precisione disponibile sull'hardware
della propria macchina.


\subsection{Un esempio elementare}
\label{sec:sig_sigchld}

Un semplice esempio per illustrare il funzionamento di un gestore di segnale è
quello della gestione di \signal{SIGCHLD}. Abbiamo visto in
sez.~\ref{sec:proc_termination} che una delle azioni eseguite dal kernel alla
conclusione di un processo è quella di inviare questo segnale al padre. In
generale dunque, quando non interessa elaborare lo stato di uscita di un
processo, si può completare la gestione della terminazione installando un
gestore per \signal{SIGCHLD} il cui unico compito sia quello di chiamare
\func{waitpid} per completare la procedura di terminazione in modo da evitare
la formazione di \textit{zombie}.\footnote{si ricordi comunque che dal kernel
  2.6 seguendo lo standard POSIX.1-2001 per evitare di dover ricevere gli
  stati di uscita che non interessano basta impostare come azione predefinita
  quella di ignorare \signal{SIGCHLD}, nel qual caso viene assunta la
  semantica di System V, in cui il segnale non viene inviato, il sistema non
  genera \textit{zombie} e lo stato di terminazione viene scartato senza dover
  chiamare una \func{wait}.}

In fig.~\ref{fig:sig_sigchld_handl} è mostrato il codice contenente una
implementazione generica di una funzione di gestione per \signal{SIGCHLD},
(che si trova nei sorgenti allegati nel file \file{SigHand.c}); se ripetiamo i
test di sez.~\ref{sec:proc_termination}, invocando \cmd{forktest} con
l'opzione \cmd{-s} (che si limita ad effettuare l'installazione di questa
funzione come gestore di \signal{SIGCHLD}) potremo verificare che non si ha
più la creazione di \textit{zombie}.

\begin{figure}[!htbp]
  \footnotesize  \centering
  \begin{minipage}[c]{\codesamplewidth}
    \includecodesample{listati/hand_sigchild.c}
  \end{minipage}
  \normalsize 
  \caption{Codice di una funzione generica di gestione per il segnale
    \signal{SIGCHLD}.}
  \label{fig:sig_sigchld_handl}
\end{figure}

Il codice del gestore è di lettura immediata, come buona norma di
programmazione (si ricordi quanto accennato sez.~\ref{sec:sys_errno}) si
comincia (\texttt{\small 6--7}) con il salvare lo stato corrente di
\var{errno}, in modo da poterlo ripristinare prima del ritorno del gestore
(\texttt{\small 16--17}). In questo modo si preserva il valore della variabile
visto dal corso di esecuzione principale del processo, che altrimenti sarebbe
sovrascritto dal valore restituito nella successiva chiamata di
\func{waitpid}.

Il compito principale del gestore è quello di ricevere lo stato di
terminazione del processo, cosa che viene eseguita nel ciclo in
(\texttt{\small 9--15}).  Il ciclo è necessario a causa di una caratteristica
fondamentale della gestione dei segnali: abbiamo già accennato come fra la
generazione di un segnale e l'esecuzione del gestore possa passare un certo
lasso di tempo e niente ci assicura che il gestore venga eseguito prima della
generazione di ulteriori segnali dello stesso tipo. In questo caso normalmente
i segnali successivi vengono ``\textsl{fusi}'' col primo ed al processo ne
viene recapitato soltanto uno.

Questo può essere un caso comune proprio con \signal{SIGCHLD}, qualora capiti
che molti processi figli terminino in rapida successione. Esso inoltre si
presenta tutte le volte che un segnale viene bloccato: per quanti siano i
segnali emessi durante il periodo di blocco, una volta che quest'ultimo sarà
rimosso verrà recapitato un solo segnale.

Allora, nel caso della terminazione dei processi figli, se si chiamasse
\func{waitpid} una sola volta, essa leggerebbe lo stato di terminazione per un
solo processo, anche se i processi terminati sono più di uno, e gli altri
resterebbero in stato di \textit{zombie} per un tempo indefinito.

Per questo occorre ripetere la chiamata di \func{waitpid} fino a che essa non
ritorni un valore nullo, segno che non resta nessun processo di cui si debba
ancora ricevere lo stato di terminazione (si veda sez.~\ref{sec:proc_wait} per
la sintassi della funzione). Si noti anche come la funzione venga invocata con
il parametro \const{WNOHANG} che permette di evitare il suo blocco quando
tutti gli stati di terminazione sono stati ricevuti.



\section{La gestione avanzata dei segnali}
\label{sec:sig_adv_control}

Le funzioni esaminate finora fanno riferimento alle modalità più elementari
della gestione dei segnali; non si sono pertanto ancora prese in
considerazione le tematiche più complesse, collegate alle varie \textit{race
  condition} che i segnali possono generare e alla natura asincrona degli
stessi.

Affronteremo queste problematiche in questa sezione, partendo da un esempio
che le evidenzi, per poi prendere in esame le varie funzioni che permettono di
risolvere i problemi più complessi connessi alla programmazione con i segnali,
fino a trattare le caratteristiche generali della gestione dei medesimi nella
casistica ordinaria.


\subsection{Alcune problematiche aperte}
\label{sec:sig_example}

Come accennato in sez.~\ref{sec:sig_pause_sleep} è possibile implementare
\func{sleep} a partire dall'uso di \func{pause} e \func{alarm}. A prima vista
questo può sembrare di implementazione immediata; ad esempio una semplice
versione di \func{sleep} potrebbe essere quella illustrata in
fig.~\ref{fig:sig_sleep_wrong}.

\begin{figure}[!htb]
  \footnotesize \centering
  \begin{minipage}[c]{\codesamplewidth}
    \includecodesample{listati/sleep_danger.c}
  \end{minipage}
  \normalsize 
  \caption{Una implementazione pericolosa di \func{sleep}.} 
  \label{fig:sig_sleep_wrong}
\end{figure}

Dato che è nostra intenzione utilizzare \signal{SIGALRM} il primo passo della
nostra implementazione sarà quello di installare il relativo gestore salvando
il precedente (\texttt{\small 14--17}).  Si effettuerà poi una chiamata ad
\func{alarm} per specificare il tempo d'attesa per l'invio del segnale a cui
segue la chiamata a \func{pause} per fermare il programma (\texttt{\small
  18--20}) fino alla sua ricezione.  Al ritorno di \func{pause}, causato dal
ritorno del gestore (\texttt{\small 1--9}), si ripristina il gestore originario
(\texttt{\small 21--22}) restituendo l'eventuale tempo rimanente
(\texttt{\small 23--24}) che potrà essere diverso da zero qualora
l'interruzione di \func{pause} venisse causata da un altro segnale.

Questo codice però, a parte il non gestire il caso in cui si è avuta una
precedente chiamata a \func{alarm} (che si è tralasciato per brevità),
presenta una pericolosa \textit{race condition}.  Infatti, se il processo
viene interrotto fra la chiamata di \func{alarm} e \func{pause}, può capitare
(ad esempio se il sistema è molto carico) che il tempo di attesa scada prima
dell'esecuzione di quest'ultima, cosicché essa sarebbe eseguita dopo l'arrivo
di \signal{SIGALRM}. In questo caso ci si troverebbe di fronte ad un
\textit{deadlock}, in quanto \func{pause} non verrebbe mai più interrotta (se
non in caso di un altro segnale).

Questo problema può essere risolto (ed è la modalità con cui veniva fatto in
SVr2) usando la funzione \func{longjmp} (vedi sez.~\ref{sec:proc_longjmp}) per
uscire dal gestore. In questo modo, con una condizione sullo stato di
uscita di quest'ultima, si può evitare la chiamata a \func{pause}, usando un
codice del tipo di quello riportato in fig.~\ref{fig:sig_sleep_incomplete}.

\begin{figure}[!htb]
  \footnotesize \centering
  \begin{minipage}[c]{\codesamplewidth}
    \includecodesample{listati/sleep_defect.c}
  \end{minipage}
  \normalsize 
  \caption{Una implementazione ancora malfunzionante di \func{sleep}.} 
  \label{fig:sig_sleep_incomplete}
\end{figure}

In questo caso il gestore (\texttt{\small 18--27}) non ritorna come in
fig.~\ref{fig:sig_sleep_wrong}, ma usa la funzione \func{longjmp}
(\texttt{\small 25}) per rientrare direttamente nel corpo principale del
programma. Dato che in questo caso il valore di uscita che verrà restituito da
\func{setjmp} è 1, grazie alla condizione impostata in (\texttt{\small 9--12})
si potrà evitare comunque che \func{pause} sia chiamata a vuoto.

Ma anche questa implementazione comporta dei problemi, in questo caso infatti
non viene gestita correttamente l'interazione con gli altri segnali. Se
infatti il segnale di allarme interrompe un altro gestore, l'esecuzione non
riprenderà nel gestore in questione, ma nel ciclo principale, interrompendone
inopportunamente l'esecuzione.  Lo stesso tipo di problemi si presenterebbero
se si volesse usare questa implementazione di \func{alarm} per stabilire un
timeout su una qualunque \textit{system call} bloccante.

Un secondo esempio dei problemi a cui si può andare incontro è quello in cui
si usa un segnale per notificare una qualche forma di evento. In genere quello
che si fa in questo caso è impostare all'interno del gestore un opportuno flag
da controllare nel corpo principale del programma, con un codice del tipo di
quello riportato in fig.~\ref{fig:sig_event_wrong}.

La logica del programma è quella di far impostare al gestore (\texttt{\small
  14--19}) una variabile globale, preventivamente inizializzata nel programma
principale, ad un diverso valore. In questo modo dal corpo principale del
programma si potrà determinare, osservandone il contenuto di detta variabile,
l'occorrenza o meno del segnale, ed eseguire le azioni conseguenti
(\texttt{\small 6--11}) relative.

\begin{figure}[!htbp]
  \footnotesize\centering
  \begin{minipage}[c]{\codesamplewidth}
    \includecodesample{listati/sig_alarm.c}
  \end{minipage}
  \normalsize 
  \caption{Un esempio non funzionante del codice per il controllo di un
    evento generato da un segnale.}
  \label{fig:sig_event_wrong}
\end{figure}

Questo è il tipico esempio di caso, già citato in
sez.~\ref{sec:proc_race_cond}, in cui si genera una \textit{race
  condition}. Infatti, in una situazione in cui un segnale è già arrivato (e
quindi \var{flag} è già stata impostata ad 1 nel gestore) se un altro segnale
arriva immediatamente dopo l'esecuzione del controllo (\texttt{\small 6}) ma
prima della cancellazione di \var{flag} fatta subito dopo (\texttt{\small 7}),
la sua occorrenza sarà perduta.

Questi esempi ci mostrano come per poter eseguire una gestione effettiva dei
segnali occorrono delle funzioni più sofisticate di quelle finora
illustrate. La funzione \func{signal} infatti ha la sua origine nella
interfaccia alquanto primitiva che venne adottata nei primi sistemi Unix, ma
con questa funzione è sostanzialmente impossibile gestire in maniera adeguata
di tutti i possibili aspetti con cui un processo deve reagire alla ricezione
di un segnale.



\subsection{Gli \textsl{insiemi di segnali} o \textit{signal set}}
\label{sec:sig_sigset}

\itindbeg{signal~set} 

Come evidenziato nel paragrafo precedente, le funzioni di gestione dei segnali
originarie, nate con la semantica inaffidabile, hanno dei limiti non
superabili; in particolare non è prevista nessuna funzione che permetta di
gestire il blocco dei segnali o di verificare lo stato dei segnali pendenti.

Per questo motivo lo standard POSIX.1, insieme alla nuova semantica dei
segnali ha introdotto una interfaccia di gestione completamente nuova, che
permette di ottenere un controllo molto più dettagliato. In particolare lo
standard ha introdotto un nuovo tipo di dato \type{sigset\_t}, che permette di
rappresentare un \textsl{insieme di segnali} (un \textit{signal set}, come
viene usualmente chiamato), tale tipo di dato viene usato per gestire il
blocco dei segnali.

Inizialmente un \textsl{insieme di segnali} veniva rappresentato da un intero
di dimensione opportuna, di solito pari al numero di bit dell'architettura
della macchina, ciascun bit del quale era associato ad uno specifico
segnale. Nel caso di architetture a 32 bit questo comporta un massimo di 32
segnali distinti e dato che a lungo questi sono stati sufficienti non c'era
necessità di nessuna struttura più complicata, in questo modo era possibile
implementare le operazioni direttamente con istruzioni elementari del
processore. 

Oggi questo non è più vero, in particolare con l'introduzione dei segnali
\textit{real-rime} (che vedremo in sez.~\ref{sec:sig_real_time}).  Dato che in
generale non si può fare conto sulle caratteristiche di una implementazione,
perché non è detto che si disponga di un numero di bit sufficienti per mettere
tutti i segnali in un intero, o perché in \type{sigset\_t} possono essere
immagazzinate ulteriori informazioni, tutte le operazioni devono essere
effettuate tramite le opportune funzioni di libreria che si curano di
mascherare i dettagli di basso livello.

Lo standard POSIX.1 definisce cinque funzioni per la manipolazione degli
insiemi di segnali. Le prime quattro, che consentono di manipolare i contenuti
di un \textit{signal set}, sono \funcd{sigemptyset}, \funcd{sigfillset},
\funcd{sigaddset} e \funcd{sigdelset}; i rispettivi prototipi sono:

\begin{funcproto}{
\fhead{signal.h}
\fdecl{int sigemptyset(sigset\_t *set)}
\fdesc{Inizializza un insieme di segnali vuoto.}
\fdecl{int sigfillset(sigset\_t *set)}
\fdesc{Inizializza un insieme di segnali pieno.}
\fdecl{int sigaddset(sigset\_t *set, int signum)}
\fdesc{Aggiunge un segnale ad un insieme di segnali.}
\fdecl{int sigdelset(sigset\_t *set, int signum)}
\fdesc{Rimuove un segnale da un insieme di segnali.}
}

{Le funzioni ritornano $0$ in caso di successo, e $-1$ per un errore, nel qual
  caso \var{errno} assumerà il valore:
  \begin{errlist}
  \item[\errcode{EINVAL}] \param{signum} non è un segnale valido.
  \end{errlist}
}
\end{funcproto}

Le prime due funzioni inizializzano l'insieme di segnali indicato
dall'argomento \param{set} rispettivamente ad un contenuto vuoto (in cui cioè
non c'è nessun segnale) e pieno (in cui cioè ci sono tutti i segnali). Le
altre due funzioni consentono di inserire o rimuovere uno specifico segnale
indicato con l'argomento \param{signum} in un insieme. 

A queste funzioni si aggiunge l'ulteriore \funcd{sigismember}, che consente di
verificare la presenza di un segnale in un insieme, il suo prototipo è:

\begin{funcproto}{
\fhead{signal.h}
\fdecl{int sigismember(const sigset\_t *set, int signum)}
\fdesc{Controlla se un segnale è in un insieme di segnali.}
}

{La funzione ritorna $1$ il segnale è nell'insieme e $0$ altrimenti, e $-1$
  per un errore, nel qual caso \var{errno} assumerà il valore \errval{EINVAL}
  se si è specificato un puntatore \var{NULL}.}
\end{funcproto}

La \acr{glibc} prevede inoltre altre funzioni non standardizzate, accessibili
definendo la macro \macro{\_GNU\_SOURCE}. La prima di queste è
\funcd{sigisemptyset}, che consente di verificare un insieme è vuoto, il suo
prototipo è:

\begin{funcproto}{
\fhead{signal.h}
\fdecl{int sigisemptyset(sigset\_t *set)}
\fdesc{Controlla se un insieme di segnali è vuoto.}
}

{La funzione ritorna $1$ l'insieme è vuoto e $0$ altrimenti, non sono previste
  condizioni di errore.}
\end{funcproto}

Alla precedente si aggiungono altre due funzioni consentono di effettuare
delle operazioni logiche con gli insiemi di segnali, esse sono
\funcd{sigorset} e \funcd{sigandset}, ed i rispettivi prototipi sono:

\begin{funcproto}{
\fhead{signal.h}
\fdecl{sigorset(sigset\_t *dest, sigset\_t *left, sigset\_t *right)}
\fdesc{Crea l'unione di due insieme di segnali.}
\fdecl{sigandset(sigset\_t *dest, sigset\_t *left, sigset\_t *right)}
\fdesc{Crea l'intersezione di due insieme di segnali.} 
}

{Le funzioni ritornano $0$ in caso di successo e $-1$ per un errore, nel qual
  caso \var{errno} assumerà il valore \errcode{EINVAL}.}
\end{funcproto}


In genere si usa un insieme di segnali per specificare quali segnali si vuole
bloccare, o per riottenere dalle varie funzioni di gestione la maschera dei
segnali attivi (vedi sez.~\ref{sec:sig_sigmask}). La modalità più comune, che
è anche quella più portabile, prevede che possano essere definiti aggiungendo
i segnali voluti ad un insieme vuoto ottenuto con \func{sigemptyset} o
togliendo quelli che non servono da un insieme completo ottenuto con
\func{sigfillset}.

\itindend{signal~set} 


\subsection{La funzione \func{sigaction}}
\label{sec:sig_sigaction}

Abbiamo già accennato in sez.~\ref{sec:sig_signal} i problemi di compatibilità
relativi all'uso di \func{signal}. Per ovviare a tutto questo lo standard
POSIX.1 ha ridefinito completamente l'interfaccia per la gestione dei segnali,
rendendola molto più flessibile e robusta, anche se leggermente più complessa.

La principale funzione di sistema prevista dall'interfaccia POSIX.1 per la
gestione dei segnali è \funcd{sigaction}. Essa ha sostanzialmente lo stesso
uso di \func{signal}, permette cioè di specificare le modalità con cui un
segnale può essere gestito da un processo. Il suo prototipo è:

\begin{funcproto}{
\fhead{signal.h}
\fdecl{int sigaction(int signum, const struct sigaction *act, struct sigaction
  *oldact)}  
\fdesc{Installa una nuova azione per un segnale.} 
}

{La funzione ritorna $0$ in caso di successo e $-1$ per un errore, nel qual
  caso \var{errno} assumerà uno dei valori: 
  \begin{errlist}
  \item[\errcode{EFAULT}] si sono specificati indirizzi non validi.
  \item[\errcode{EINVAL}] si è specificato un numero di segnale invalido o si è
    cercato di installare il gestore per \signal{SIGKILL} o
    \signal{SIGSTOP}.
  \end{errlist}
}
\end{funcproto}

La funzione serve ad installare una nuova \textsl{azione} per il segnale
indicato dall'argomento \param{signum}. Si parla di \textsl{azione} e non di
\textsl{gestore} come nel caso di \func{signal}, in quanto la funzione
consente di specificare le varie caratteristiche della risposta al segnale,
non solo la funzione che verrà eseguita alla sua occorrenza.  

Per questo motivo lo standard POSIX.1 raccomanda di usare sempre questa
funzione al posto della precedente \func{signal}, che in genere viene
ridefinita in termini di \func{sigaction}, in quanto la nuova interfaccia
permette un controllo completo su tutti gli aspetti della gestione di un
segnale, sia pure al prezzo di una maggiore complessità d'uso.

Se il puntatore \param{act} non è nullo, la funzione installa la nuova azione
da esso specificata, se \param{oldact} non è nullo il valore dell'azione
corrente viene restituito indietro.  Questo permette (specificando \param{act}
nullo e \param{oldact} non nullo) di superare uno dei limiti di \func{signal},
che non consente di ottenere l'azione corrente senza installarne una nuova. Se
sia \param{act} che \param{oldact} la funzione può essere utilizzata per
verificare, se da luogo ad un errore, se il segnale indicato è valido per la
piattaforma che si sta usando.

Entrambi i puntatori fanno riferimento alla struttura \struct{sigaction},
tramite la quale si specificano tutte le caratteristiche dell'azione associata
ad un segnale.  Anch'essa è descritta dallo standard POSIX.1 ed in Linux è
definita secondo quanto riportato in fig.~\ref{fig:sig_sigaction}. Il campo
\var{sa\_restorer}, non previsto dallo standard, è obsoleto e non deve essere
più usato.

\begin{figure}[!htb]
  \footnotesize \centering
  \begin{minipage}[c]{0.8\textwidth}
    \includestruct{listati/sigaction.h}
  \end{minipage} 
  \normalsize 
  \caption{La struttura \structd{sigaction}.} 
  \label{fig:sig_sigaction}
\end{figure}

Il campo \var{sa\_mask} serve ad indicare l'insieme dei segnali che devono
essere bloccati durante l'esecuzione del gestore, ad essi viene comunque
sempre aggiunto il segnale che ne ha causato la chiamata, a meno che non si
sia specificato con \var{sa\_flag} un comportamento diverso. Quando il
gestore ritorna comunque la maschera dei segnali bloccati (vedi
sez.~\ref{sec:sig_sigmask}) viene ripristinata al valore precedente
l'invocazione.

L'uso di questo campo permette ad esempio di risolvere il problema residuo
dell'implementazione di \code{sleep} mostrata in
fig.~\ref{fig:sig_sleep_incomplete}. In quel caso infatti se il segnale di
allarme avesse interrotto un altro gestore questo non sarebbe stato eseguito
correttamente, la cosa poteva essere prevenuta installando gli altri gestori
usando \var{sa\_mask} per bloccare \signal{SIGALRM} durante la loro
esecuzione.  Il valore di \var{sa\_flag} permette di specificare vari aspetti
del comportamento di \func{sigaction}, e della reazione del processo ai vari
segnali; i valori possibili ed il relativo significato sono riportati in
tab.~\ref{tab:sig_sa_flag}.

\begin{table}[!htb]
  \footnotesize
  \centering
  \begin{tabular}[c]{|l|p{8cm}|}
    \hline
    \textbf{Valore} & \textbf{Significato} \\
    \hline
    \hline
    \constd{SA\_NOCLDSTOP}& Se il segnale è \signal{SIGCHLD} allora non deve
                            essere notificato quando il processo figlio viene
                            fermato da uno dei segnali \signal{SIGSTOP},
                            \signal{SIGTSTP}, \signal{SIGTTIN} o 
                            \signal{SIGTTOU}, questo flag ha significato solo
                            quando si imposta un gestore per \signal{SIGCHLD}.\\
    \constd{SA\_NOCLDWAIT}& Se il segnale è \signal{SIGCHLD} e si richiede di
                            ignorare il segnale con \const{SIG\_IGN} allora i
                            processi figli non diventano \textit{zombie} quando
                            terminano; questa funzionalità è stata introdotta
                            nel kernel 2.6 e va a modificare il comportamento
                            di \func{waitpid} come illustrato in
                            sez.~\ref{sec:proc_wait}, se si installa un gestore
                            con questo flag attivo il segnale \signal{SIGCHLD}
                            viene comunque generato.\\
    \constd{SA\_NODEFER}  & Evita che il segnale corrente sia bloccato durante
                            l'esecuzione del gestore.\\
    \constd{SA\_NOMASK}   & Nome obsoleto e sinonimo non standard di
                            \const{SA\_NODEFER}, non deve essere più
                            utilizzato.\\ 
    \constd{SA\_ONESHOT}  & Nome obsoleto e sinonimo non standard di
                            \const{SA\_RESETHAND}, non deve essere più
                            utilizzato.\\ 
    \constd{SA\_ONSTACK}  & Stabilisce l'uso di uno \textit{stack} alternativo
                            per l'esecuzione del gestore (vedi
                            sez.~\ref{sec:sig_specific_features}).\\  
    \constd{SA\_RESETHAND}& Ristabilisce l'azione per il segnale al valore 
                            predefinito una volta che il gestore è stato
                            lanciato, riproduce cioè il comportamento della
                            semantica inaffidabile.\\  
    \constd{SA\_RESTART}  & Riavvia automaticamente le \textit{slow system
                            call} quando vengono interrotte dal suddetto
                            segnale, riproduce cioè il comportamento standard
                            di BSD.\\ 
    \constd{SA\_SIGINFO}  & Deve essere specificato quando si vuole usare un
                            gestore in forma estesa usando
                            \var{sa\_sigaction} al posto di
                            \var{sa\_handler}.\\
    \hline
  \end{tabular}
  \caption{Valori del campo \var{sa\_flag} della struttura \struct{sigaction}.}
  \label{tab:sig_sa_flag}
\end{table}

Come si può notare in fig.~\ref{fig:sig_sigaction} \func{sigaction} permette
di utilizzare due forme diverse di gestore,\footnote{la possibilità è prevista
  dallo standard POSIX.1b, ed è stata aggiunta nei kernel della serie 2.1.x
  con l'introduzione dei segnali \textit{real-time} (vedi
  sez.~\ref{sec:sig_real_time}); in precedenza era possibile ottenere alcune
  informazioni addizionali usando \var{sa\_handler} con un secondo parametro
  addizionale di tipo \var{sigcontext}, che adesso è deprecato.}  da
specificare, a seconda dell'uso o meno del flag \const{SA\_SIGINFO},
rispettivamente attraverso i campi \var{sa\_sigaction} o \var{sa\_handler}.
Quest'ultima è quella classica usata anche con \func{signal}, mentre la prima
permette di usare un gestore più complesso, in grado di ricevere informazioni
più dettagliate dal sistema, attraverso la struttura \struct{siginfo\_t},
riportata in fig.~\ref{fig:sig_siginfo_t}.  I due campi devono essere usati in
maniera alternativa, in certe implementazioni questi campi vengono addirittura
definiti come una \direct{union}.\footnote{la direttiva \direct{union} del
  linguaggio C definisce una variabile complessa, analoga a una stuttura, i
  cui campi indicano i diversi tipi di valori che possono essere salvati, in
  maniera alternativa, all'interno della stessa.}

Installando un gestore di tipo \var{sa\_sigaction} diventa allora possibile
accedere alle informazioni restituite attraverso il puntatore a questa
struttura. Tutti i segnali impostano i campi \var{si\_signo}, che riporta il
numero del segnale ricevuto, \var{si\_errno}, che riporta, quando diverso da
zero, il codice dell'errore associato al segnale, e \var{si\_code}, che viene
usato dal kernel per specificare maggiori dettagli riguardo l'evento che ha
causato l'emissione del segnale.

\begin{figure}[!htb]
  \footnotesize \centering
  \begin{minipage}[c]{0.9\textwidth}
    \includestruct{listati/siginfo_t.h}
  \end{minipage} 
  \normalsize 
  \caption{La struttura \structd{siginfo\_t}.} 
  \label{fig:sig_siginfo_t}
\end{figure}
 
In generale \var{si\_code} contiene, per i segnali generici, per quelli
\textit{real-time} e per tutti quelli inviati tramite da un processo con
\func{kill} o affini, le informazioni circa l'origine del segnale stesso, ad
esempio se generato dal kernel, da un timer, da \func{kill}, ecc. Il valore
viene sempre espresso come una costante,\footnote{le definizioni di tutti i
  valori possibili si trovano in \file{bits/siginfo.h}.} ed i valori possibili
in questo caso sono riportati in tab.~\ref{tab:sig_si_code_generic}.

Nel caso di alcuni segnali però il valore di \var{si\_code} viene usato per
fornire una informazione specifica relativa alle motivazioni della ricezione
dello stesso; ad esempio i vari segnali di errore (\signal{SIGILL},
\signal{SIGFPE}, \signal{SIGSEGV} e \signal{SIGBUS}) lo usano per fornire
maggiori dettagli riguardo l'errore, come il tipo di errore aritmetico, di
istruzione illecita o di violazione di memoria; mentre alcuni segnali di
controllo (\signal{SIGCHLD}, \signal{SIGTRAP} e \signal{SIGPOLL}) forniscono
altre informazioni specifiche.

\begin{table}[!htb]
  \footnotesize
  \centering
  \begin{tabular}[c]{|l|p{10cm}|}
    \hline
    \textbf{Valore} & \textbf{Significato} \\
    \hline
    \hline
    \constd{SI\_USER}   & Generato da \func{kill} o \func{raise} o affini.\\
    \constd{SI\_KERNEL} & Inviato direttamente dal kernel.\\
    \constd{SI\_QUEUE}  & Inviato con \func{sigqueue} (vedi
                          sez.~\ref{sec:sig_real_time}).\\ 
    \constd{SI\_TIMER}  & Scadenza di un \textit{POSIX timer} (vedi
                          sez.~\ref{sec:sig_timer_adv}).\\ 
    \constd{SI\_MESGQ}  & Inviato al cambiamento di stato di una coda di
                          messaggi POSIX (vedi sez.~\ref{sec:ipc_posix_mq}),
                          introdotto con il kernel 2.6.6.\\ 
    \constd{SI\_ASYNCIO}& Una operazione di I/O asincrono (vedi
                          sez.~\ref{sec:file_asyncronous_io}) è stata
                          completata.\\
    \constd{SI\_SIGIO}  & Segnale di \signal{SIGIO} da una coda (vedi
                          sez.~\ref{sec:file_asyncronous_operation}).\\ 
    \constd{SI\_TKILL}  & Inviato da \func{tkill} o \func{tgkill} (vedi
                          sez.~\ref{cha:thread_xxx}), introdotto con il kernel
                          2.4.19.\\ 
    \hline
  \end{tabular}
  \caption{Valori del campo \var{si\_code} della struttura \struct{sigaction}
    per i segnali generici.}
  \label{tab:sig_si_code_generic}
\end{table}


In questo caso il valore del campo \var{si\_code} deve essere verificato nei
confronti delle diverse costanti previste per ciascuno di detti segnali; dato
che si tratta di costanti, e non di una maschera binaria, i valori numerici
vengono riutilizzati e ciascuno di essi avrà un significato diverso a seconda
del segnale a cui è associato. 

L'elenco dettagliato dei nomi di queste costanti è riportato nelle diverse
sezioni di tab.~\ref{tab:sig_si_code_special} che sono state ordinate nella
sequenza in cui si sono appena citati i rispettivi segnali, il prefisso del
nome indica comunque in maniera diretta il segnale a cui le costanti fanno
riferimento.

\begin{table}[!htb]
  \footnotesize
  \centering
  \begin{tabular}[c]{|l|p{8cm}|}
    \hline
    \textbf{Valore} & \textbf{Significato} \\
    \hline
    \hline
    \constd{ILL\_ILLOPC}  & Codice di operazione illegale.\\
    \constd{ILL\_ILLOPN}  & Operando illegale.\\
    \constd{ILL\_ILLADR}  & Modo di indirizzamento illegale.\\
    \constd{ILL\_ILLTRP}  & Trappola di processore illegale.\\
    \constd{ILL\_PRVOPC}  & Codice di operazione privilegiato.\\
    \constd{ILL\_PRVREG}  & Registro privilegiato.\\
    \constd{ILL\_COPROC}  & Errore del coprocessore.\\
    \constd{ILL\_BADSTK}  & Errore nello stack interno.\\
    \hline
    \constd{FPE\_INTDIV}  & Divisione per zero intera.\\
    \constd{FPE\_INTOVF}  & Overflow intero.\\
    \constd{FPE\_FLTDIV}  & Divisione per zero in virgola mobile.\\
    \constd{FPE\_FLTOVF}  & Overflow in virgola mobile.\\
    \constd{FPE\_FLTUND}  & Underflow in virgola mobile.\\
    \constd{FPE\_FLTRES}  & Risultato in virgola mobile non esatto.\\
    \constd{FPE\_FLTINV}  & Operazione in virgola mobile non valida.\\
    \constd{FPE\_FLTSUB}  & Mantissa? fuori intervallo.\\
    \hline
    \constd{SEGV\_MAPERR} & Indirizzo non mappato.\\
    \constd{SEGV\_ACCERR} & Permessi non validi per l'indirizzo.\\
    \hline
    \constd{BUS\_ADRALN}  & Allineamento dell'indirizzo non valido.\\
    \constd{BUS\_ADRERR}  & Indirizzo fisico inesistente.\\
    \constd{BUS\_OBJERR}  & Errore hardware sull'indirizzo.\\
    \hline
    \constd{TRAP\_BRKPT}  & Breakpoint sul processo.\\
    \constd{TRAP\_TRACE}  & Trappola di tracciamento del processo.\\
    \hline
    \constd{CLD\_EXITED}  & Il figlio è uscito.\\
    \constd{CLD\_KILLED}  & Il figlio è stato terminato.\\
    \constd{CLD\_DUMPED}  & Il figlio è terminato in modo anormale.\\
    \constd{CLD\_TRAPPED} & Un figlio tracciato ha raggiunto una trappola.\\
    \constd{CLD\_STOPPED} & Il figlio è stato fermato.\\
    \constd{CLD\_CONTINUED}& Il figlio è ripartito.\\
    \hline
    \constd{POLL\_IN}   & Disponibili dati in ingresso.\\
    \constd{POLL\_OUT}  & Spazio disponibile sul buffer di uscita.\\
    \constd{POLL\_MSG}  & Disponibili messaggi in ingresso.\\
    \constd{POLL\_ERR}  & Errore di I/O.\\
    \constd{POLL\_PRI}  & Disponibili dati di alta priorità in ingresso.\\
    \constd{POLL\_HUP}  & Il dispositivo è stato disconnesso.\\
    \hline
  \end{tabular}
  \caption{Valori del campo \var{si\_code} della struttura \struct{sigaction}
    impostati rispettivamente dai segnali \signal{SIGILL}, \signal{SIGFPE},
    \signal{SIGSEGV}, \signal{SIGBUS}, \signal{SIGCHLD}, \signal{SIGTRAP} e
    \signal{SIGPOLL}/\signal{SIGIO}.}
  \label{tab:sig_si_code_special}
\end{table}

Il resto della struttura \struct{siginfo\_t} è definito come una \dirct{union}
ed i valori eventualmente presenti dipendono dal segnale ricevuto, così
\signal{SIGCHLD} ed i segnali \textit{real-time} (vedi
sez.~\ref{sec:sig_real_time}) inviati tramite \func{kill} avvalorano
\var{si\_pid} e \var{si\_uid} coi valori corrispondenti al processo che ha
emesso il segnale, \signal{SIGCHLD} avvalora anche i campi \var{si\_status},
\var{si\_utime} e \var{si\_stime} che indicano rispettivamente lo stato di
uscita, l'\textit{user time} e il \textit{system time} (vedi
sez.~\ref{sec:sys_cpu_times}) usati dal processo; \signal{SIGILL},
\signal{SIGFPE}, \signal{SIGSEGV} e \signal{SIGBUS} avvalorano \var{si\_addr}
con l'indirizzo in cui è avvenuto l'errore, \signal{SIGIO} (vedi
sez.~\ref{sec:file_asyncronous_io}) avvalora \var{si\_fd} con il numero del
file descriptor e \var{si\_band} per i dati urgenti (vedi
sez.~\ref{sec:TCP_urgent_data}) su un socket, il segnale inviato alla scadenza
di un POSIX timer (vedi sez.~\ref{sec:sig_timer_adv}) avvalora i campi
\var{si\_timerid} e \var{si\_overrun}.

Benché sia possibile usare nello stesso programma sia \func{sigaction} che
\func{signal} occorre molta attenzione, in quanto le due funzioni possono
interagire in maniera anomala. Infatti l'azione specificata con
\struct{sigaction} contiene un maggior numero di informazioni rispetto al
semplice indirizzo del gestore restituito da \func{signal}.  Per questo motivo
se si usa quest'ultima per installare un gestore sostituendone uno
precedentemente installato con \func{sigaction}, non sarà possibile effettuare
un ripristino corretto dello stesso.

Per questo è sempre opportuno usare \func{sigaction}, che è in grado di
ripristinare correttamente un gestore precedente, anche se questo è stato
installato con \func{signal}. In generale poi non è il caso di usare il valore
di ritorno di \func{signal} come campo \var{sa\_handler}, o viceversa, dato
che in certi sistemi questi possono essere diversi. In definitiva dunque, a
meno che non si sia vincolati all'aderenza stretta allo standard ISO C, è
sempre il caso di evitare l'uso di \func{signal} a favore di \func{sigaction}.

\begin{figure}[!htbp]
  \footnotesize  \centering
  \begin{minipage}[c]{\codesamplewidth}
    \includecodesample{listati/Signal.c}
  \end{minipage} 
  \normalsize 
  \caption{La funzione \func{Signal}, equivalente a \func{signal}, definita
    attraverso \func{sigaction}.}
  \label{fig:sig_Signal_code}
\end{figure}

Per questo motivo si è provveduto, per mantenere un'interfaccia semplificata
che abbia le stesse caratteristiche di \func{signal}, a definire attraverso
\func{sigaction} una funzione equivalente \func{Signal}, il cui codice è
riportato in fig.~\ref{fig:sig_Signal_code} (il codice completo si trova nel
file \file{SigHand.c} nei sorgenti allegati). Anche in questo caso, per
semplificare la definizione si è poi definito un apposito tipo
\texttt{SigFunc} per esprimere in modo più comprensibile la forma di un
gestore di segnale.

Si noti come, essendo la funzione estremamente semplice, essa è definita come
\dirct{inline}. Questa direttiva viene usata per dire al compilatore di
trattare la funzione cui essa fa riferimento in maniera speciale inserendo il
codice direttamente nel testo del programma.  Anche se i compilatori più
moderni sono in grado di effettuare da soli queste manipolazioni (impostando
le opportune ottimizzazioni) questa è una tecnica usata per migliorare le
prestazioni per le funzioni piccole ed usate di frequente, in particolare nel
kernel, dove in certi casi le ottimizzazioni dal compilatore, tarate per l'uso
in \textit{user space}, non sono sempre adatte.

In tal caso infatti le istruzioni per creare un nuovo frame nello
\textit{stack} per chiamare la funzione costituirebbero una parte rilevante
del codice, appesantendo inutilmente il programma.  Originariamente questo
comportamento veniva ottenuto con delle macro, ma queste hanno tutta una serie
di problemi di sintassi nel passaggio degli argomenti (si veda ad esempio
\cite{PratC}) che in questo modo possono essere evitati.



\subsection{La gestione della \textsl{maschera dei segnali} o 
  \textit{signal mask}}
\label{sec:sig_sigmask}

\index{maschera dei segnali|(}

Come spiegato in sez.~\ref{sec:sig_semantics} tutti i moderni sistemi
unix-like permettono di bloccare temporaneamente (o di eliminare
completamente, impostando come azione \const{SIG\_IGN}) la consegna dei
segnali ad un processo. Questo è fatto specificando la cosiddetta
\textsl{maschera dei segnali} (o \textit{signal mask}) del
processo\footnote{nel caso di Linux essa è mantenuta dal campo \var{blocked}
  della \struct{task\_struct} del processo.} cioè l'insieme dei segnali la cui
consegna è bloccata.

Abbiamo accennato in sez.~\ref{sec:proc_fork} che la maschera dei segnali
viene ereditata dal padre alla creazione di un processo figlio, e abbiamo
visto al paragrafo precedente che essa può essere modificata durante
l'esecuzione di un gestore ed automaticamente ripristinata quando questo
ritorna, attraverso l'uso dal campo \var{sa\_mask} di \struct{sigaction}.

Uno dei problemi evidenziatisi con l'esempio di fig.~\ref{fig:sig_event_wrong}
è che in molti casi è necessario proteggere delle sezioni di codice, in modo
da essere sicuri che essi siano eseguite senza interruzioni da parte di un
segnale.  Nel caso in questione si trattava della sezione di codice fra il
controllo e la eventuale cancellazione del flag impostato dal gestore di un
segnale che testimoniava l'avvenuta occorrenza dello stesso.

Come illustrato in sez.~\ref{sec:proc_atom_oper} le operazioni più semplici,
come l'assegnazione o il controllo di una variabile, di norma sono atomiche, e
qualora si voglia essere sicuri si può usare il tipo \type{sig\_atomic\_t}. Ma
quando si devono eseguire più operazioni su delle variabili (nell'esempio
citato un controllo ed una assegnazione) o comunque eseguire una serie di
istruzioni, l'atomicità non è più possibile.

In questo caso, se si vuole essere sicuri di non poter essere interrotti da un
segnale durante l'esecuzione di una sezione di codice, lo si può bloccare
esplicitamente modificando la maschera dei segnali del processo con la
funzione di sistema \funcd{sigprocmask}, il cui prototipo è:

\begin{funcproto}{
\fhead{signal.h}
\fdecl{int sigprocmask(int how, const sigset\_t *set, sigset\_t *oldset)}
\fdesc{Imposta la maschera dei segnali del processo corrente.} 
}

{La funzione ritorna $0$ in caso di successo e $-1$ per un errore, nel qual
  caso \var{errno} assumerà uno dei valori: 
  \begin{errlist}
  \item[\errcode{EFAULT}] si sono specificati indirizzi non validi.
  \item[\errcode{EINVAL}] si è specificato un numero di segnale invalido.
  \end{errlist}
}
\end{funcproto}

La funzione usa l'insieme di segnali posto all'indirizzo passato
nell'argomento \param{set} per modificare la maschera dei segnali del processo
corrente. La modifica viene effettuata a seconda del valore
dell'argomento \param{how}, secondo le modalità specificate in
tab.~\ref{tab:sig_procmask_how}. Qualora si specifichi un valore non nullo
per \param{oldset} la maschera dei segnali corrente viene salvata a
quell'indirizzo.

\begin{table}[htb]
  \footnotesize
  \centering
  \begin{tabular}[c]{|l|p{8cm}|}
    \hline
    \textbf{Valore} & \textbf{Significato} \\
    \hline
    \hline
    \constd{SIG\_BLOCK}   & L'insieme dei segnali bloccati è l'unione fra
                            quello specificato e quello corrente.\\
    \constd{SIG\_UNBLOCK} & I segnali specificati in \param{set} sono rimossi
                            dalla maschera dei segnali, specificare la
                            cancellazione di un segnale non bloccato è legale.\\
    \constd{SIG\_SETMASK} & La maschera dei segnali è impostata al valore
                            specificato da \param{set}.\\
    \hline
  \end{tabular}
  \caption{Valori e significato dell'argomento \param{how} della funzione
    \func{sigprocmask}.}
  \label{tab:sig_procmask_how}
\end{table}

In questo modo diventa possibile proteggere delle sezioni di codice bloccando
l'insieme di segnali voluto per poi riabilitarli alla fine della sezione
critica. La funzione permette di risolvere problemi come quelli mostrati in
fig.~\ref{fig:sig_event_wrong}, proteggendo la sezione fra il controllo del
flag e la sua cancellazione.  La funzione può essere usata anche all'interno
di un gestore, ad esempio per riabilitare la consegna del segnale che l'ha
invocato, in questo caso però occorre ricordare che qualunque modifica alla
maschera dei segnali viene perduta al ritorno dallo stesso.

Benché con l'uso di \func{sigprocmask} si possano risolvere la maggior parte
dei casi di \textit{race condition} restano aperte alcune possibilità legate
all'uso di \func{pause}.  Il caso è simile a quello del problema illustrato
nell'esempio di fig.~\ref{fig:sig_sleep_incomplete}, e cioè la possibilità che
il processo riceva il segnale che si intende usare per uscire dallo stato di
attesa invocato con \func{pause} immediatamente prima dell'esecuzione di
quest'ultima. Per poter effettuare atomicamente la modifica della maschera dei
segnali (di solito attivandone uno specifico) insieme alla sospensione del
processo lo standard POSIX ha previsto la funzione di sistema
\funcd{sigsuspend}, il cui prototipo è:

\begin{funcproto}{
\fhead{signal.h}
\fdecl{int sigsuspend(const sigset\_t *mask)} 
\fdesc{Imposta la maschera dei segnali mettendo in attesa il processo.} 
}

{La funzione ritorna $0$ in caso di successo e $-1$ per un errore, nel qual
  caso \var{errno} assumerà uno dei valori: 
  \begin{errlist}
  \item[\errcode{EFAULT}] si sono specificati indirizzi non validi.
  \item[\errcode{EINVAL}] si è specificato un numero di segnale invalido.
  \end{errlist}
}
\end{funcproto}

Come esempio dell'uso di queste funzioni proviamo a riscrivere un'altra volta
l'esempio di implementazione di \code{sleep}. Abbiamo accennato in
sez.~\ref{sec:sig_sigaction} come con \func{sigaction} sia possibile bloccare
\signal{SIGALRM} nell'installazione dei gestori degli altri segnali, per poter
usare l'implementazione vista in fig.~\ref{fig:sig_sleep_incomplete} senza
interferenze.  Questo però comporta una precauzione ulteriore al semplice uso
della funzione, vediamo allora come usando la nuova interfaccia è possibile
ottenere un'implementazione, riportata in fig.~\ref{fig:sig_sleep_ok} che non
presenta neanche questa necessità.

\begin{figure}[!htbp]
  \footnotesize \centering
  \begin{minipage}[c]{\codesamplewidth}
    \includecodesample{listati/sleep.c}
  \end{minipage} 
  \normalsize 
  \caption{Una implementazione completa di \func{sleep}.} 
  \label{fig:sig_sleep_ok}
\end{figure}
 
Per evitare i problemi di interferenza con gli altri segnali in questo caso
non si è usato l'approccio di fig.~\ref{fig:sig_sleep_incomplete} evitando
l'uso di \func{longjmp}. Come in precedenza il gestore (\texttt{\small
  27--30}) non esegue nessuna operazione, limitandosi a ritornare per
interrompere il programma messo in attesa.

La prima parte della funzione (\texttt{\small 6--10}) provvede ad installare
l'opportuno gestore per \signal{SIGALRM}, salvando quello originario, che
sarà ripristinato alla conclusione della stessa (\texttt{\small 23}); il passo
successivo è quello di bloccare \signal{SIGALRM} (\texttt{\small 11--14}) per
evitare che esso possa essere ricevuto dal processo fra l'esecuzione di
\func{alarm} (\texttt{\small 16}) e la sospensione dello stesso. Nel fare
questo si salva la maschera corrente dei segnali, che sarà ripristinata alla
fine (\texttt{\small 22}), e al contempo si prepara la maschera dei segnali
\var{sleep\_mask} per riattivare \signal{SIGALRM} all'esecuzione di
\func{sigsuspend}.  

In questo modo non sono più possibili \textit{race condition} dato che
\signal{SIGALRM} viene disabilitato con \func{sigprocmask} fino alla chiamata
di \func{sigsuspend}.  Questo metodo è assolutamente generale e può essere
applicato a qualunque altra situazione in cui si deve attendere per un
segnale, i passi sono sempre i seguenti:
\begin{enumerate*}
\item leggere la maschera dei segnali corrente e bloccare il segnale voluto
  con \func{sigprocmask};
\item mandare il processo in attesa con \func{sigsuspend} abilitando la
  ricezione del segnale voluto;
\item ripristinare la maschera dei segnali originaria.
\end{enumerate*}
Per quanto possa sembrare strano bloccare la ricezione di un segnale per poi
riabilitarla immediatamente dopo, in questo modo si evita il \textit{deadlock}
dovuto all'arrivo del segnale prima dell'esecuzione di \func{sigsuspend}.

\index{maschera dei segnali|)}


\subsection{Criteri di programmazione per i gestori dei segnali}
\label{sec:sig_signal_handler}

Abbiamo finora parlato dei gestori dei segnali come funzioni chiamate in
corrispondenza della consegna di un segnale. In realtà un gestore non può
essere una funzione qualunque, in quanto esso può essere eseguito in
corrispondenza all'interruzione in un punto qualunque del programma
principale, cosa che ad esempio può rendere problematico chiamare all'interno
di un gestore di segnali la stessa funzione che dal segnale è stata
interrotta.

\index{funzioni!\textit{signal safe}|(}

Il concetto è comunque più generale e porta ad una distinzione fra quelle che
POSIX chiama \textsl{funzioni insicure} (\textit{signal unsafe function}) e
\textsl{funzioni sicure} (o più precisamente \textit{signal safe function}).
Quando un segnale interrompe una funzione insicura ed il gestore chiama al suo
interno una funzione insicura il sistema può dare luogo ad un comportamento
indefinito, la cosa non avviene invece per le funzioni sicure.

Tutto questo significa che la funzione che si usa come gestore di segnale deve
essere programmata con molta cura per evirare questa evenienza e che non è
possibile utilizzare al suo interno una qualunque funzione di sistema, se si
vogliono evitare questi problemi si può ricorrere soltanto all'uso delle
funzioni considerate sicure.

L'elenco delle funzioni considerate sicure varia a seconda della
implementazione utilizzata e dello standard a cui si fa riferimento. Non è
riportata una lista specifica delle funzioni sicure per Linux, e si suppone
pertanto che siano quelle richieste dallo standard. Secondo quanto richiesto
dallo standard POSIX 1003.1 nella revisione del 2003, le ``\textit{signal safe
  function}'' che possono essere chiamate anche all'interno di un gestore di
segnali sono tutte quelle della lista riportata in
fig.~\ref{fig:sig_safe_functions}.

\begin{figure}[!htb]
  \footnotesize \centering
  \begin{minipage}[c]{14cm}
    \func{\_exit}, \func{abort}, \func{accept}, \func{access},
    \func{aio\_error} \func{aio\_return}, \func{aio\_suspend}, \func{alarm},
    \func{bind}, \func{cfgetispeed}, \func{cfgetospeed}, \func{cfsetispeed},
    \func{cfsetospeed}, \func{chdir}, \func{chmod}, \func{chown},
    \func{clock\_gettime}, \func{close}, \func{connect}, \func{creat},
    \func{dup}, \func{dup2}, \func{execle}, \func{execve}, \func{fchmod},
    \func{fchown}, \func{fcntl}, \func{fdatasync}, \func{fork},
    \func{fpathconf}, \func{fstat}, \func{fsync}, \func{ftruncate},
    \func{getegid}, \func{geteuid}, \func{getgid}, \func{getgroups},
    \func{getpeername}, \func{getpgrp}, \func{getpid}, \func{getppid},
    \func{getsockname}, \func{getsockopt}, \func{getuid}, \func{kill},
    \func{link}, \func{listen}, \func{lseek}, \func{lstat}, \func{mkdir},
    \func{mkfifo}, \func{open}, \func{pathconf}, \func{pause}, \func{pipe},
    \func{poll}, \funcm{posix\_trace\_event}, \func{pselect}, \func{raise},
    \func{read}, \func{readlink}, \func{recv}, \func{recvfrom},
    \func{recvmsg}, \func{rename}, \func{rmdir}, \func{select},
    \func{sem\_post}, \func{send}, \func{sendmsg}, \func{sendto},
    \func{setgid}, \func{setpgid}, \func{setsid}, \func{setsockopt},
    \func{setuid}, \func{shutdown}, \func{sigaction}, \func{sigaddset},
    \func{sigdelset}, \func{sigemptyset}, \func{sigfillset},
    \func{sigismember}, \func{signal}, \func{sigpause}, \func{sigpending},
    \func{sigprocmask}, \func{sigqueue}, \funcm{sigset}, \func{sigsuspend},
    \func{sleep}, \func{socket}, \func{socketpair}, \func{stat},
    \func{symlink}, \func{sysconf}, \func{tcdrain}, \func{tcflow},
    \func{tcflush}, \func{tcgetattr}, \func{tcgetgrp}, \func{tcsendbreak},
    \func{tcsetattr}, \func{tcsetpgrp}, \func{time}, \func{timer\_getoverrun},
    \func{timer\_gettime}, \func{timer\_settime}, \func{times}, \func{umask},
    \func{uname}, \func{unlink}, \func{utime}, \func{wait}, \func{waitpid},
    \func{write}.
  \end{minipage} 
  \normalsize 
  \caption{Elenco delle funzioni sicure secondo lo standard POSIX
    1003.1-2003.}
  \label{fig:sig_safe_functions}
\end{figure}

\index{funzioni!\textit{signal safe}|)}

Lo standard POSIX.1-2004 modifica la lista di
fig.~\ref{fig:sig_safe_functions} aggiungendo le funzioni \func{\_Exit} e
\func{sockatmark}, mentre lo standard POSIX.1-2008 rimuove della lista le tre
funzioni \func{fpathconf}, \func{pathconf}, \func{sysconf} e vi aggiunge le
ulteriori funzioni in fig.~\ref{fig:sig_safe_functions_posix_2008}.

\begin{figure}[!htb]
  \footnotesize \centering
  \begin{minipage}[c]{14cm}
     \func{execl}, \func{execv}, \func{faccessat}, \func{fchmodat},
     \func{fchownat}, \func{fexecve}, \func{fstatat}, \func{futimens},
     \func{linkat}, \func{mkdirat}, \func{mkfifoat}, \func{mknod},
     \func{mknodat}, \func{openat}, \func{readlinkat}, \func{renameat},
     \func{symlinkat}, \func{unlinkat}, \func{utimensat}, \func{utimes}.
  \end{minipage} 
  \normalsize 
  \caption{Ulteriori funzioni sicure secondo lo standard POSIX.1-2008.}
  \label{fig:sig_safe_functions_posix_2008}
\end{figure}


Per questo motivo è opportuno mantenere al minimo indispensabile le operazioni
effettuate all'interno di un gestore di segnali, qualora si debbano compiere
operazioni complesse è sempre preferibile utilizzare la tecnica in cui si usa
il gestore per impostare il valore di una qualche variabile globale, e poi si
eseguono le operazioni complesse nel programma verificando (con tutti gli
accorgimenti visti in precedenza) il valore di questa variabile tutte le volte
che si è rilevata una interruzione dovuta ad un segnale.


\section{Funzionalità avanzate}
\label{sec:sig_advanced_signal}

Tratteremo in questa ultima sezione alcune funzionalità avanzate relativa ai
segnali ed in generale ai meccanismi di notifica, a partire dalla funzioni
introdotte per la gestione dei cosiddetti ``\textsl{segnali real-time}'', alla
gestione avanzata delle temporizzazioni e le nuove interfacce per la gestione
di segnali ed eventi attraverso l'uso di file descriptor.

\subsection{I segnali \textit{real-time}}
\label{sec:sig_real_time}

Lo standard POSIX.1b, nel definire una serie di nuove interfacce per i servizi
\textit{real-time}, ha introdotto una estensione del modello classico dei
segnali che presenta dei significativi miglioramenti,\footnote{questa
  estensione è stata introdotta in Linux a partire dal kernel 2.1.43, e dalla
  versione 2.1 della \acr{glibc}.} in particolare sono stati superati tre
limiti fondamentali dei segnali classici:
\begin{basedescript}{\desclabelwidth{1cm}\desclabelstyle{\nextlinelabel}}
\item[\textbf{I segnali non sono accumulati}] 
  se più segnali vengono generati prima dell'esecuzione di un gestore
  questo sarà eseguito una sola volta, ed il processo non sarà in grado di
  accorgersi di quante volte l'evento che ha generato il segnale è accaduto.
\item[\textbf{I segnali non trasportano informazione}]   
  i segnali classici non prevedono altra informazione sull'evento
  che li ha generati se non il fatto che sono stati emessi (tutta
  l'informazione che il kernel associa ad un segnale è il suo numero).
\item[\textbf{I segnali non hanno un ordine di consegna}] 
  l'ordine in cui diversi segnali vengono consegnati è casuale e non
  prevedibile. Non è possibile stabilire una priorità per cui la reazione a
  certi segnali ha la precedenza rispetto ad altri.
\end{basedescript}

Per poter superare queste limitazioni lo standard POSIX.1b ha introdotto delle
nuove caratteristiche, che sono state associate ad una nuova classe di
segnali, che vengono chiamati \textsl{segnali real-time}, in particolare le
funzionalità aggiunte sono:

\begin{enumerate}
\item i segnali sono inseriti in una coda che permette di consegnare istanze
  multiple dello stesso segnale qualora esso venga inviato più volte prima
  dell'esecuzione del gestore; si assicura così che il processo riceva un
  segnale per ogni occorrenza dell'evento che lo genera;
\item è stata introdotta una priorità nella consegna dei segnali: i segnali
  vengono consegnati in ordine a seconda del loro valore, partendo da quelli
  con un numero minore, che pertanto hanno una priorità maggiore;
\item è stata introdotta la possibilità di restituire dei dati al gestore,
  attraverso l'uso di un apposito campo \var{si\_value} nella struttura
  \struct{siginfo\_t}, accessibile tramite gestori di tipo
  \var{sa\_sigaction}.
\end{enumerate}

Tutte queste nuove funzionalità eccetto l'ultima, che, come illustrato in
sez.~\ref{sec:sig_sigaction}, è disponibile anche con i segnali ordinari, si
applicano solo ai nuovi segnali \textit{real-time}; questi ultimi sono
accessibili in un intervallo di valori specificati dalle due costanti
\constd{SIGRTMIN} e \constd{SIGRTMAX}, che specificano il numero minimo e
massimo associato ad un segnale \textit{real-time}.

Su Linux di solito il primo valore è 33, mentre il secondo è \code{\_NSIG-1},
che di norma (vale a dire sulla piattaforma i386) è 64. Questo dà un totale di
32 segnali disponibili, contro gli almeno 8 richiesti da POSIX.1b. Si tenga
presente però che i primi segnali \textit{real-time} disponibili vengono usati
dalla \acr{glibc} per l'implementazione dei \textit{thread} POSIX (vedi
sez.~\ref{sec:thread_posix_intro}), ed il valore di \const{SIGRTMIN} viene
modificato di conseguenza.\footnote{per la precisione vengono usati i primi
  tre per la vecchia implementazione dei \textit{LinuxThread} ed i primi due
  per la nuova NTPL (\textit{New Thread Posix Library}), il che comporta che
  \const{SIGRTMIN} a seconda dei casi può assumere i valori 34 o 35.}

Per questo motivo nei programmi che usano i segnali \textit{real-time} non si
deve mai usare un valore assoluto dato che si correrebbe il rischio di
utilizzare un segnale in uso alle librerie, ed il numero del segnale deve
invece essere sempre specificato in forma relativa a \const{SIGRTMIN} (come
\code{SIGRTMIN + n}) avendo inoltre cura di controllare di non aver mai
superato \const{SIGRTMAX}.

I segnali con un numero più basso hanno una priorità maggiore e vengono
consegnati per primi, inoltre i segnali \textit{real-time} non possono
interrompere l'esecuzione di un gestore di un segnale a priorità più alta; la
loro azione predefinita è quella di terminare il programma.  I segnali
ordinari hanno tutti la stessa priorità, che è più alta di quella di qualunque
segnale \textit{real-time}. Lo standard non definisce niente al riguardo ma
Linux, come molte altre implementazioni, adotta questa politica.

Si tenga presente che questi nuovi segnali non sono associati a nessun evento
specifico, a meno di non richiedere specificamente il loro utilizzo in
meccanismi di notifica come quelli per l'I/O asincrono (vedi
sez.~\ref{sec:file_asyncronous_io}) o per le code di messaggi POSIX (vedi
sez.~\ref{sec:ipc_posix_mq}), pertanto devono essere inviati esplicitamente.

Inoltre, per poter usufruire della capacità di restituire dei dati, i relativi
gestori devono essere installati con \func{sigaction}, specificando per
\var{sa\_flags} la modalità \const{SA\_SIGINFO} che permette di utilizzare la
forma estesa \var{sa\_sigaction} del gestore (vedi
sez.~\ref{sec:sig_sigaction}).  In questo modo tutti i segnali
\textit{real-time} possono restituire al gestore una serie di informazioni
aggiuntive attraverso l'argomento \struct{siginfo\_t}, la cui definizione è
stata già vista in fig.~\ref{fig:sig_siginfo_t}, nella trattazione dei gestori
in forma estesa.

In particolare i campi utilizzati dai segnali \textit{real-time} sono
\var{si\_pid} e \var{si\_uid} in cui vengono memorizzati rispettivamente il
\ids{PID} e l'\ids{UID} effettivo del processo che ha inviato il segnale,
mentre per la restituzione dei dati viene usato il campo \var{si\_value}.

\begin{figure}[!htb]
  \footnotesize \centering
  \begin{minipage}[c]{0.8\textwidth}
    \includestruct{listati/sigval_t.h}
  \end{minipage} 
  \normalsize 
  \caption{La definizione dell'unione \structd{sigval}, definita anche come
    tipo \typed{sigval\_t}.}
  \label{fig:sig_sigval}
\end{figure}

Detto campo, identificato con il tipo di dato \type{sigval\_t}, è una
\dirct{union} di tipo \struct{sigval} (la sua definizione è in
fig.~\ref{fig:sig_sigval}) in cui può essere memorizzato o un valore numerico,
se usata nella forma \var{sival\_int}, o un puntatore, se usata nella forma
\var{sival\_ptr}. L'unione viene usata dai segnali \textit{real-time} e da
vari meccanismi di notifica per restituire dati al gestore del segnale in
\var{si\_value}. Un campo di tipo \type{sigval\_t} è presente anche nella
struttura \struct{sigevent} (definita in fig.~\ref{fig:struct_sigevent}) che
viene usata dai meccanismi di notifica come quelli per i timer POSIX (vedi
sez.~\ref{sec:sig_timer_adv}), l'I/O asincrono (vedi
sez.~\ref{sec:file_asyncronous_io}) o le code di messaggi POSIX (vedi
sez.~\ref{sec:ipc_posix_mq}).

A causa delle loro caratteristiche, la funzione \func{kill} non è adatta ad
inviare segnali \textit{real-time}, poiché non è in grado di fornire alcun
valore per il campo \var{si\_value} restituito nella struttura
\struct{siginfo\_t} prevista da un gestore in forma estesa. Per questo motivo
lo standard ha previsto una nuova funzione, \funcd{sigqueue}, il cui prototipo
è:

\begin{funcproto}{
\fhead{signal.h}
\fdecl{int sigqueue(pid\_t pid, int signo, const union sigval value)}
\fdesc{Invia un segnale con un valore di informazione.} 
}

{La funzione ritorna $0$ in caso di successo e $-1$ per un errore, nel qual
  caso \var{errno} assumerà uno dei valori: 
  \begin{errlist}
  \item[\errcode{EAGAIN}] la coda è esaurita, ci sono già
    \const{SIGQUEUE\_MAX} segnali in attesa si consegna.
  \item[\errcode{EINVAL}] si è specificato un valore non valido per
    \param{signo}.
  \item[\errcode{EPERM}] non si hanno privilegi appropriati per inviare il
    segnale al processo specificato.
  \item[\errcode{ESRCH}] il processo \param{pid} non esiste.
  \end{errlist}
}
\end{funcproto}


La funzione invia il segnale indicato dall'argomento \param{signo} al processo
indicato dall'argomento \param{pid}. Per il resto il comportamento della
funzione è analogo a quello di \func{kill}, ed i privilegi occorrenti ad
inviare il segnale ad un determinato processo sono gli stessi; un valore nullo
di \param{signo} permette di verificare le condizioni di errore senza inviare
nessun segnale.

Se il segnale è bloccato la funzione ritorna immediatamente, se si è
installato un gestore con \const{SA\_SIGINFO} e ci sono risorse disponibili,
(vale a dire che c'è posto nella coda dei segnali \textit{real-time}) esso
viene inserito e diventa pendente. Una volta consegnato il segnale il gestore
otterrà nel campo \var{si\_code} di \struct{siginfo\_t} il valore
\const{SI\_QUEUE} e nel campo \var{si\_value} il valore indicato
nell'argomento \param{value}. Se invece si è installato un gestore nella forma
classica il segnale sarà generato, ma tutte le caratteristiche tipiche dei
segnali \textit{real-time} (priorità e coda) saranno perse.

Per lo standard POSIX la profondità della coda è indicata dalla costante
\constd{SIGQUEUE\_MAX}, una della tante costanti di sistema definite dallo
standard POSIX che non abbiamo riportato esplicitamente in
sez.~\ref{sec:sys_limits}. Il suo valore minimo secondo lo standard,
\macrod{\_POSIX\_SIGQUEUE\_MAX}, è pari a 32. Nel caso di Linux la coda ha una
dimensione variabile; fino alla versione 2.6.7 c'era un limite massimo globale
che poteva essere impostato come parametro del kernel in
\sysctlfiled{kernel/rtsig-max} ed il valore predefinito era pari a 1024. A
partire dal kernel 2.6.8 il valore globale è stato rimosso e sostituito dalla
risorsa \const{RLIMIT\_SIGPENDING} associata al singolo utente, che può essere
modificata con \func{setrlimit} come illustrato in
sez.~\ref{sec:sys_resource_limit}.

Lo standard POSIX.1b definisce inoltre delle nuove funzioni di sistema che
permettono di gestire l'attesa di segnali specifici su una coda, esse servono
in particolar modo nel caso dei \textit{thread}, in cui si possono usare i
segnali \textit{real-time} come meccanismi di comunicazione elementare; la
prima di queste è \funcd{sigwait}, il cui prototipo è:

\begin{funcproto}{
\fhead{signal.h}
\fdecl{int sigwait(const sigset\_t *set, int *sig)}
\fdesc{Attende la ricezione di un segnale.} 
}
{La funzione ritorna $0$ in caso di successo e $-1$ per un errore, nel qual
  caso \var{errno} assumerà uno dei valori: 
  \begin{errlist}
  \item[\errcode{EINTR}] la funzione è stata interrotta.
  \item[\errcode{EINVAL}] si è specificato un valore non valido per
  \end{errlist}
  ed inoltre \errval{EFAULT} nel suo significato generico.}
\end{funcproto}

La funzione estrae dall'insieme dei segnali pendenti uno qualunque fra quelli
indicati nel \textit{signal set} specificato in \param{set}, il cui valore
viene restituito nella variabile puntata da \param{sig}.  Se sono pendenti più
segnali, viene estratto quello a priorità più alta, cioè quello con il numero
più basso. Se, nel caso di segnali \textit{real-time}, c'è più di un segnale
pendente, ne verrà estratto solo uno. Una volta estratto il segnale non verrà
più consegnato, e se era in una coda il suo posto sarà liberato. Se non c'è
nessun segnale pendente il processo viene bloccato fintanto che non ne arriva
uno.

Per un funzionamento corretto la funzione richiede che alla sua chiamata i
segnali di \param{set} siano bloccati. In caso contrario si avrebbe un
conflitto con gli eventuali gestori: pertanto non si deve utilizzare per
lo stesso segnale questa funzione e \func{sigaction}. Se questo non avviene il
comportamento del sistema è indeterminato: il segnale può sia essere
consegnato che essere ricevuto da \func{sigwait}, il tutto in maniera non
prevedibile.

Lo standard POSIX.1b definisce altre due funzioni di sistema, anch'esse usate
prevalentemente con i \textit{thread}; \funcd{sigwaitinfo} e
\funcd{sigtimedwait}, i relativi prototipi sono:

\begin{funcproto}{
\fhead{signal.h}
\fdecl{int sigwaitinfo(const sigset\_t *set, siginfo\_t *info)}
\fdesc{Attende un segnale con le relative informazioni.}
\fdecl{int sigtimedwait(const sigset\_t *set, siginfo\_t *info, const
  struct timespec *timeout)}
\fdesc{Attende un segnale con le relative informazioni per un tempo massimo.}
}

{Le funzioni ritornano $0$ in caso di successo e $-1$ per un errore, nel qual
  caso \var{errno} assumerà uno gli stessi valori di \func{sigwait} ai quali
  si aggiunge per \func{sigtimedwait}:
  \begin{errlist}
  \item[\errcode{EAGAIN}] si è superato il timeout senza che un segnale atteso
    sia stato ricevuto.
  \end{errlist}
}
\end{funcproto}


Entrambe le funzioni sono estensioni di \func{sigwait}. La prima permette di
ricevere, oltre al numero del segnale, anche le informazioni ad esso associate
tramite l'argomento \param{info}; in particolare viene restituito il numero
del segnale nel campo \var{si\_signo}, la sua causa in \var{si\_code}, e se il
segnale è stato immesso sulla coda con \func{sigqueue}, il valore di ritorno
ad esso associato viene riportato in \var{si\_value}, che altrimenti è
indefinito.

La seconda è identica alla prima ma in più permette di specificare un timeout
con l'argomento omonimo, scaduto il quale ritornerà con un errore. Se si
specifica per \param{timeout} un puntatore nullo il comportamento sarà
identico a \func{sigwaitinfo}. Se si specifica un tempo di timeout nullo e non
ci sono segnali pendenti la funzione ritornerà immediatamente, in questo modo
si può eliminare un segnale dalla coda senza dover essere bloccati qualora
esso non sia presente.

L'uso di queste funzioni è principalmente associato alla gestione dei segnali
con i \textit{thread}. In genere esse vengono chiamate dal \textit{thread}
incaricato della gestione, che al ritorno della funzione esegue il codice che
usualmente sarebbe messo nel gestore, per poi ripetere la chiamata per
mettersi in attesa del segnale successivo. Questo ovviamente comporta che non
devono essere installati gestori, che solo il \textit{thread} di gestione deve
usare \func{sigwait} e che i segnali gestiti in questa maniera, per evitare
che venga eseguita l'azione predefinita, devono essere mascherati per tutti i
\textit{thread}, compreso quello dedicato alla gestione, che potrebbe
riceverlo fra due chiamate successive.


\subsection{La gestione avanzata delle temporizzazioni}
\label{sec:sig_timer_adv}

Sia le funzioni per la gestione dei tempi viste in
sez.~\ref{sec:sys_cpu_times} che quelle per la gestione dei timer di
sez.~\ref{sec:sig_alarm_abort} sono state a lungo limitate dalla risoluzione
massima dei tempi dell'orologio interno del kernel, che era quella ottenibile
dal timer di sistema che governa lo \textit{scheduler}, e quindi limitate
dalla frequenza dello stesso che si ricordi, come già illustrato in
sez.~\ref{sec:proc_hierarchy}, è data dal valore della costante \texttt{HZ}. 

I contatori usati per il calcolo dei tempi infatti erano basati sul numero di
\textit{jiffies} che vengono incrementati ad ogni \textit{clock tick} del
timer di sistema, il che comportava anche, come accennato in
sez.~\ref{sec:sig_alarm_abort} per \func{setitimer}, problemi per il massimo
periodo di tempo copribile da alcuni di questi orologi, come quelli associati
al \textit{process time} almeno fino a quando, con il kernel 2.6.16, non è
stato rimosso il limite di un valore a 32 bit per i \textit{jiffies}.

\itindbeg{POSIX~Timer~API}

Nelle architetture moderne però tutti i computer sono dotati di temporizzatori
hardware che possono supportare risoluzioni molto elevate, ed in maniera del
tutto indipendente dalla frequenza scelta per il timer di sistema che governa
lo \textit{scheduler}, normalmente si possono ottenere precisioni fino al
microsecondo, andando molto oltre in caso di hardware dedicato. 

Per questo lo standard POSIX.1-2001 ha previsto una serie di nuove funzioni
relative a quelli che vengono chiamati ``\textsl{orologi}
\textit{real-time}'', in grado di supportare risoluzioni fino al
nanosecondo. Inoltre le CPU più moderne sono dotate a loro volta di contatori
ad alta definizione che consentono una grande accuratezza nella misura del
tempo da esse dedicato all'esecuzione di un processo.

Per usare queste funzionalità ed ottenere risoluzioni temporali più accurate,
occorre però un opportuno supporto da parte del kernel, ed i cosiddetti
\itindex{High~Resolution~Timer~(HRT)} \textit{high resolution timer} che
consentono di fare ciò sono stati introdotti nel kernel ufficiale solo a
partire dalla versione 2.6.21.\footnote{per il supporto deve essere stata
  abilitata l'opzione di compilazione \texttt{CONFIG\_HIGH\_RES\_TIMERS}, il
  supporto era però disponibile anche in precedenza nei patch facenti parte
  dello sviluppo delle estensioni \textit{real-time} del kernel, per cui
  alcune distribuzioni possono averlo anche con versioni precedenti del
  kernel.} Le funzioni definite dallo standard POSIX per gestire orologi ad
alta definizione però erano già presenti, essendo stata introdotte insieme ad
altre funzioni per il supporto delle estensioni \textit{real-time} con il
rilascio del kernel 2.6, ma la risoluzione effettiva era nominale.

A tutte le implementazioni che si rifanno a queste estensioni è richiesto di
disporre di una versione \textit{real-time} almeno per l'orologio generale di
sistema, quello che mantiene il \textit{calendar time} (vedi
sez.~\ref{sec:sys_time_base}), che in questa forma deve indicare il numero di
secondi e nanosecondi passati a partire dal primo gennaio 1970 (\textit{The
  Epoch}). Si ricordi infatti che l'orologio ordinario usato dal
\textit{calendar time} riporta solo un numero di secondi, e che la risoluzione
effettiva normalmente non raggiunge il nanosecondo (a meno di hardware
specializzato).  Oltre all'orologio generale di sistema possono essere
presenti altri tipi di orologi \textit{real-time}, ciascuno dei quali viene
identificato da un opportuno valore di una variabile di tipo
\type{clockid\_t}; un elenco di quelli disponibili su Linux è riportato in
tab.~\ref{tab:sig_timer_clockid_types}.

\begin{table}[htb]
  \footnotesize
  \centering
  \begin{tabular}[c]{|l|p{8cm}|}
    \hline
    \textbf{Valore} & \textbf{Significato} \\
    \hline
    \hline
    \constd{CLOCK\_REALTIME}    & Orologio \textit{real-time} di sistema, può
                                  essere impostato solo con privilegi
                                  amministrativi.\\ 
    \constd{CLOCK\_MONOTONIC}   & Orologio che indica un tempo monotono
                                  crescente (a partire da un tempo iniziale non
                                  specificato) che non può essere modificato e
                                  non cambia neanche in caso di reimpostazione
                                  dell'orologio di sistema.\\
    \constd{CLOCK\_PROCESS\_CPUTIME\_ID}& Contatore del tempo di CPU usato 
                                  da un processo (il \textit{process time} di
                                  sez.~\ref{sec:sys_cpu_times}, nel totale di
                                  \textit{system time} e \textit{user time})
                                  comprensivo di tutto il tempo di CPU usato
                                  da eventuali \textit{thread}.\\
    \constd{CLOCK\_THREAD\_CPUTIME\_ID}& Contatore del tempo di CPU
                                  (\textit{user time} e \textit{system time})
                                  usato da un singolo \textit{thread}.\\
    \hline
    \constd{CLOCK\_MONOTONIC\_RAW}&Simile al precedente, ma non subisce gli
                                  aggiustamenti dovuti all'uso di NTP (viene
                                  usato per fare riferimento ad una fonte
                                  hardware). Questo orologio è specifico di
                                  Linux, ed è disponibile a partire dal kernel
                                  2.6.28.\\
    \constd{CLOCK\_BOOTTIME}    & Identico a \const{CLOCK\_MONOTONIC} ma tiene
                                  conto anche del tempo durante il quale il
                                  sistema è stato sospeso (nel caso di
                                  sospensione in RAM o \textsl{ibernazione} su
                                  disco. Questo orologio è specifico di Linux,
                                  ed è disponibile a partire dal kernel
                                  2.6.39.\\
    \constd{CLOCK\_REALTIME\_ALARM}&Identico a \const{CLOCK\_REALTIME}, ma se
                                  usato per un timer il sistema sarà riattivato 
                                  anche se è in sospensione. Questo orologio è
                                  specifico di Linux, ed è disponibile a
                                  partire dal kernel 3.0.\\
    \constd{CLOCK\_BOOTTIME\_ALARM}&Identico a \const{CLOCK\_BOOTTIME}, ma se
                                  usato per un timer il sistema sarà riattivato 
                                  anche se è in sospensione. Questo orologio è
                                  specifico di Linux, ed è disponibile a
                                  partire dal kernel 3.0.\\
%    \const{}   & .\\
    \hline
  \end{tabular}
  \caption{Valori possibili per una variabile di tipo \typed{clockid\_t} 
    usata per indicare a quale tipo di orologio si vuole fare riferimento.}
  \label{tab:sig_timer_clockid_types}
\end{table}


% TODO: dal 4.17 CLOCK_MONOTONIC e CLOCK_BOOTTIME sono identici vedi
% https://lwn.net/Articles/751651/ e
% https://git.kernel.org/linus/d6ed449afdb38f89a7b38ec50e367559e1b8f71f
% change reverted, vedi: https://lwn.net/Articles/752757/

% NOTE: dal 3.0 anche i cosiddetti Posix Alarm Timers, con
% CLOCK_REALTIME_ALARM vedi http://lwn.net/Articles/429925/
% TODO: dal 3.10 anche CLOCK_TAI 

% TODO seguire l'evoluzione delle nuove syscall per il problema del 2038,
% iniziate ad entrare nel kernel dal 5.1, vedi
% https://lwn.net/Articles/776435/, https://lwn.net/Articles/782511/,
% https://git.kernel.org/linus/b1b988a6a035 

Per poter utilizzare queste funzionalità la \acr{glibc} richiede che la
macro \macro{\_POSIX\_C\_SOURCE} sia definita ad un valore maggiore o uguale
di \texttt{199309L} (vedi sez.~\ref{sec:intro_gcc_glibc_std}), inoltre i
programmi che le usano devono essere collegati con la libreria delle
estensioni \textit{real-time} usando esplicitamente l'opzione \texttt{-lrt}.

Si tenga presente inoltre che la disponibilità di queste funzionalità avanzate
può essere controllato dalla definizione della macro \macrod{\_POSIX\_TIMERS}
ad un valore maggiore di 0, e che le ulteriori macro
\macrod{\_POSIX\_MONOTONIC\_CLOCK}, \macrod{\_POSIX\_CPUTIME} e
\macrod{\_POSIX\_THREAD\_CPUTIME} indicano la presenza dei rispettivi orologi
di tipo \const{CLOCK\_MONOTONIC}, \const{CLOCK\_PROCESS\_CPUTIME\_ID} e
\const{CLOCK\_THREAD\_CPUTIME\_ID}; tutte queste macro sono definite in
\headfile{unistd.h}, che pertanto deve essere incluso per poterle
controllarle. Infine se il kernel ha il supporto per gli \textit{high
  resolution timer} un elenco degli orologi e dei timer può essere ottenuto
tramite il file \procfile{/proc/timer\_list}.

Le due funzioni che ci consentono rispettivamente di modificare o leggere il
valore per uno degli orologi \textit{real-time} sono \funcd{clock\_settime} e
\funcd{clock\_gettime}; i rispettivi prototipi sono:

\begin{funcproto}{
\fhead{time.h}
\fdecl{int clock\_settime(clockid\_t clockid, const struct timespec *tp)}
\fdesc{Imposta un orologio \textit{real-time}.} 
\fdecl{int clock\_gettime(clockid\_t clockid, struct timespec *tp)}
\fdesc{Legge un orologio \textit{real-time}.} 
}

{Le funzioni ritornano $0$ in caso di successo e $-1$ per un errore, nel qual
  caso \var{errno} assumerà uno dei valori: 
  \begin{errlist}
  \item[\errcode{EFAULT}] l'indirizzo \param{tp} non è valido.
  \item[\errcode{EINVAL}] il valore specificato per \param{clockid} non è
    valido o il relativo orologio \textit{real-time} non è supportato dal
    sistema.
  \item[\errcode{EPERM}] non si ha il permesso di impostare l'orologio
    indicato (solo per \func{clock\_settime}).
  \end{errlist}
}
\end{funcproto}

Entrambe le funzioni richiedono che si specifichi come primo argomento il tipo
di orologio su cui si vuole operare con uno dei valori di
tab.~\ref{tab:sig_timer_clockid_types} o con il risultato di una chiamata a
\func{clock\_getcpuclockid} (che tratteremo a breve), il secondo argomento
invece è sempre il puntatore \param{tp} ad una struttura \struct{timespec}
(vedi fig.~\ref{fig:sys_timespec_struct}) che deve essere stata
precedentemente allocata.  Per \func{clock\_settime} questa dovrà anche essere
stata inizializzata con il valore che si vuole impostare sull'orologio, mentre
per \func{clock\_gettime} verrà restituito al suo interno il valore corrente
dello stesso.

Si tenga presente inoltre che per eseguire un cambiamento sull'orologio
generale di sistema \const{CLOCK\_REALTIME} occorrono i privilegi
amministrativi;\footnote{ed in particolare la \textit{capability}
  \const{CAP\_SYS\_TIME}.} inoltre ogni cambiamento ad esso apportato non avrà
nessun effetto sulle temporizzazioni effettuate in forma relativa, come quelle
impostate sulle quantità di \textit{process time} o per un intervallo di tempo
da trascorrere, ma solo su quelle che hanno richiesto una temporizzazione ad
un istante preciso (in termini di \textit{calendar time}). Si tenga inoltre
presente che nel caso di Linux \const{CLOCK\_REALTIME} è l'unico orologio per
cui si può effettuare una modifica, infatti nonostante lo standard preveda la
possibilità di modifiche anche per \const{CLOCK\_PROCESS\_CPUTIME\_ID} e
\const{CLOCK\_THREAD\_CPUTIME\_ID}, il kernel non le consente.

Oltre alle due funzioni precedenti, lo standard POSIX prevede una terza
funzione di sistema che consenta di ottenere la risoluzione effettiva fornita
da un certo orologio, la funzione è \funcd{clock\_getres} ed il suo prototipo
è:

\begin{funcproto}{
\fhead{time.h}
\fdecl{int clock\_getres(clockid\_t clockid, struct timespec *res)}
\fdesc{Legge la risoluzione di un orologio \textit{real-time}.} 
}

{La funzione ritorna $0$ in caso di successo e $-1$ per un errore, nel qual
  caso \var{errno} assumerà uno dei valori: 
  \begin{errlist}
  \item[\errcode{EFAULT}] l'indirizzo di \param{res} non è valido.
  \item[\errcode{EINVAL}] il valore specificato per \param{clockid} non è
    valido.
  \end{errlist}
}
\end{funcproto}

La funzione richiede come primo argomento l'indicazione dell'orologio di cui
si vuole conoscere la risoluzione (effettuata allo stesso modo delle due
precedenti) e questa verrà restituita in una struttura \struct{timespec}
all'indirizzo puntato dall'argomento \param{res}.

Come accennato il valore di questa risoluzione dipende sia dall'hardware
disponibile che dalla implementazione delle funzioni, e costituisce il limite
minimo di un intervallo di tempo che si può indicare. Qualunque valore si
voglia utilizzare nelle funzioni di impostazione che non corrisponda ad un
multiplo intero di questa risoluzione, sarà troncato in maniera automatica. 

Gli orologi elencati nella seconda sezione di
tab.~\ref{tab:sig_timer_clockid_types} sono delle estensioni specifiche di
Linux, create per rispondere ad alcune esigenze specifiche, come quella di
tener conto di eventuali periodi di sospensione del sistema, e presenti solo
nelle versioni più recenti del kernel. In particolare gli ultimi due,
contraddistinti dal suffisso \texttt{\_ALARM}, hanno un impiego particolare,
derivato dalle esigenze emerse con Android per l'uso di Linux sui cellulari,
che consente di creare timer che possono scattare, riattivando il sistema,
anche quando questo è in sospensione. Per il loro utilizzo è prevista la
necessità di una capacità specifica, \const{CAP\_WAKE\_ALARM} (vedi
sez.~\ref{sec:proc_capabilities}).

Si tenga presente inoltre che con l'introduzione degli \textit{high resolution
  timer} i due orologi \const{CLOCK\_PROCESS\_CPUTIME\_ID} e
\const{CLOCK\_THREAD\_CPUTIME\_ID} fanno riferimento ai contatori presenti in
opportuni registri interni del processore; questo sui sistemi multiprocessore
può avere delle ripercussioni sulla precisione delle misure di tempo che vanno
al di là della risoluzione teorica ottenibile con \func{clock\_getres}, che
può essere ottenuta soltanto quando si è sicuri che un processo (o un
\textit{thread}) sia sempre stato eseguito sullo stesso processore.

Con i sistemi multiprocessore infatti ogni singola CPU ha i suoi registri
interni, e se ciascuna di esse utilizza una base di tempo diversa (se cioè il
segnale di temporizzazione inviato ai processori non ha una sola provenienza)
in genere ciascuna di queste potrà avere delle frequenze leggermente diverse,
e si otterranno pertanto dei valori dei contatori scorrelati fra loro, senza
nessuna possibilità di sincronizzazione.

Il problema si presenta, in forma più lieve, anche se la base di tempo è la
stessa, dato che un sistema multiprocessore non avvia mai tutte le CPU allo
stesso istante, si potrà così avere di nuovo una differenza fra i contatori,
soggetta però soltanto ad uno sfasamento costante. Per questo caso il kernel
per alcune architetture ha del codice che consente di ridurre al minimo la
differenza, ma non può essere comunque garantito che questa si annulli (anche
se in genere risulta molto piccola e trascurabile nella gran parte dei casi).

Per poter gestire questo tipo di problematiche lo standard ha previsto una
apposita funzione che sia in grado di ottenere l'identificativo dell'orologio
associato al \textit{process time} di un processo, la funzione è
\funcd{clock\_getcpuclockid} ed il suo prototipo è:

\begin{funcproto}{
\fhead{time.h}
\fdecl{int clock\_getcpuclockid(pid\_t pid, clockid\_t *clockid)}
\fdesc{Ottiene l'identificatore dell'orologio di CPU usato da un processo.} 
}

{La funzione ritorna $0$ in caso di successo ed un numero positivo per un
  errore, nel qual caso \var{errno} assumerà uno dei valori:
  \begin{errlist}
  \item[\errcode{ENOSYS}] non c'è il supporto per ottenere l'orologio relativo
    al \textit{process time} di un altro processo, e \param{pid} non
    corrisponde al processo corrente.
  \item[\errcode{EPERM}] il chiamante non ha il permesso di accedere alle
    informazioni relative al processo \param{pid}, avviene solo se è
    disponibile il supporto per leggere l'orologio relativo ad un altro
    processo.
  \item[\errcode{ESRCH}] non esiste il processo \param{pid}.
  \end{errlist}
}
\end{funcproto}

La funzione ritorna l'identificativo di un orologio di sistema associato ad un
processo indicato tramite l'argomento \param{pid}. Un utente normale, posto
che il kernel sia sufficientemente recente da supportare questa funzionalità,
può accedere soltanto ai dati relativi ai propri processi.

Del tutto analoga a \func{clock\_getcpuclockid}, ma da utilizzare per ottenere
l'orologio associato ad un \textit{thread} invece che a un processo, è
\funcd{pthread\_getcpuclockid},\footnote{per poterla utilizzare, come per
  qualunque funzione che faccia riferimento ai \textit{thread}, occorre
  effettuare il collegamento alla relativa libreria di gestione compilando il
  programma con \texttt{-lpthread}.} il cui prototipo è:

\begin{funcproto}{
\fhead{pthread.h}
\fhead{time.h}
\fdecl{int pthread\_getcpuclockid(pthread\_t thread, clockid\_t *clockid)}
\fdesc{Ottiene l'identificatore dell'orologio di CPU associato ad un
  \textit{thread}.} 
}

{La funzione ritorna $0$ in caso di successo ed un numero positivo per un
  errore, nel qual caso \var{errno} assumerà uno dei valori:
  \begin{errlist}
  \item[\errcode{ENOENT}] la funzione non è supportata dal sistema.
  \item[\errcode{ESRCH}] non esiste il \textit{thread} identificato
  \end{errlist}
 }
\end{funcproto}


% TODO, dal 2.6.39 aggiunta clock_adjtime 

Con l'introduzione degli orologi ad alta risoluzione è divenuto possibile
ottenere anche una gestione più avanzata degli allarmi; abbiamo già visto in
sez.~\ref{sec:sig_alarm_abort} come l'interfaccia di \func{setitimer} derivata
da BSD presenti delle serie limitazioni, come la possibilità di perdere un
segnale sotto carico, tanto che nello standard POSIX.1-2008 questa viene
marcata come obsoleta, e ne viene fortemente consigliata la sostituzione con
nuova interfaccia definita dallo standard POSIX.1-2001 che va sotto il nome di
\textit{POSIX Timer API}. Questa interfaccia è stata introdotta a partire dal
kernel 2.6, anche se il supporto di varie funzionalità da essa previste è
stato aggiunto solo in un secondo tempo.

Una delle principali differenze della nuova interfaccia è che un processo può
utilizzare un numero arbitrario di timer; questi vengono creati (ma non
avviati) tramite la funzione di sistema \funcd{timer\_create}, il cui
prototipo è:

\begin{funcproto}{
\fhead{signal.h}
\fhead{time.h}
\fdecl{int timer\_create(clockid\_t clockid, struct sigevent *evp,
    timer\_t *timerid)}
\fdesc{Crea un nuovo timer POSIX.} 
}

{La funzione ritorna $0$ in caso di successo e $-1$ per un errore, nel qual
  caso \var{errno} assumerà uno dei valori: 
  \begin{errlist}
  \item[\errcode{EAGAIN}] fallimento nel tentativo di allocare le strutture
    dei timer.
  \item[\errcode{EINVAL}] uno dei valori specificati per \param{clockid} o per
    i campi \var{sigev\_notify}, \var{sigev\_signo} o
    \var{sigev\_notify\_thread\_id} di \param{evp} non è valido.
  \item[\errcode{ENOMEM}] errore di allocazione della memoria.
  \end{errlist}
}
\end{funcproto}

La funzione richiede tre argomenti: il primo argomento serve ad indicare quale
tipo di orologio si vuole utilizzare e prende uno dei valori di
tab.~\ref{tab:sig_timer_clockid_types}; di detti valori però non è previsto
l'uso di \const{CLOCK\_MONOTONIC\_RAW} mentre
\const{CLOCK\_PROCESS\_CPUTIME\_ID} e \const{CLOCK\_THREAD\_CPUTIME\_ID} sono
disponibili solo a partire dal kernel 2.6.12. Si può così fare riferimento sia
ad un tempo assoluto che al tempo utilizzato dal processo (o \textit{thread})
stesso. Si possono inoltre utilizzare, posto di avere un kernel che li
supporti, gli orologi aggiuntivi della seconda parte di
tab.~\ref{tab:sig_timer_clockid_types}. 

Il secondo argomento richiede una trattazione più dettagliata, in quanto
introduce una struttura di uso generale, \struct{sigevent}, che viene
utilizzata anche da altre funzioni, come quelle per l'I/O asincrono (vedi
sez.~\ref{sec:file_asyncronous_io}) o le code di messaggi POSIX (vedi
sez.~\ref{sec:ipc_posix_mq}) e che serve ad indicare in maniera generica un
meccanismo di notifica.

\begin{figure}[!htb]
  \footnotesize \centering
  \begin{minipage}[c]{0.8\textwidth}
    \includestruct{listati/sigevent.h}
  \end{minipage} 
  \normalsize 
  \caption{La struttura \structd{sigevent}, usata per specificare in maniera
    generica diverse modalità di notifica degli eventi.}
  \label{fig:struct_sigevent}
\end{figure}

La struttura \struct{sigevent} (accessibile includendo \headfile{time.h}) è
riportata in fig.~\ref{fig:struct_sigevent}, la definizione effettiva dipende
dall'implementazione, quella mostrata è la versione descritta nella pagina di
manuale di \func{timer\_create}. Il campo \var{sigev\_notify} è il più
importante essendo quello che indica le modalità della notifica, gli altri
dipendono dal valore che si è specificato per \var{sigev\_notify}, si sono
riportati in tab.~\ref{tab:sigevent_sigev_notify}. La scelta del meccanismo di
notifica viene fatta impostando uno dei valori di
tab.~\ref{tab:sigevent_sigev_notify} per \var{sigev\_notify}, e fornendo gli
eventuali ulteriori argomenti necessari a secondo della scelta
effettuata. Diventa così possibile indicare l'uso di un segnale o l'esecuzione
(nel caso di uso dei \textit{thread}) di una funzione di modifica in un
\textit{thread} dedicato.

\begin{table}[htb]
  \footnotesize
  \centering
  \begin{tabular}[c]{|l|p{10cm}|}
    \hline
    \textbf{Valore} & \textbf{Significato} \\
    \hline
    \hline
    \constd{SIGEV\_NONE}   & Non viene inviata nessuna notifica.\\
    \constd{SIGEV\_SIGNAL} & La notifica viene effettuata inviando al processo
                             chiamante il segnale specificato dal campo
                             \var{sigev\_signo}; se il gestore di questo
                             segnale è stato installato con
                             \const{SA\_SIGINFO} gli verrà restituito il
                             valore specificato con \var{sigev\_value} (una
                             \dirct{union} \texttt{sigval}, la cui definizione
                             è in fig.~\ref{fig:sig_sigval}) come valore del
                             campo \var{si\_value} di \struct{siginfo\_t}.\\
    \constd{SIGEV\_THREAD} & La notifica viene effettuata creando un nuovo
                             \textit{thread} che esegue la funzione di
                             notifica specificata da 
                             \var{sigev\_notify\_function} con argomento
                             \var{sigev\_value}. Se questo è diverso da
                             \val{NULL}, il \textit{thread} viene creato con
                             gli attributi specificati da
                             \var{sigev\_notify\_attribute}.\footnotemark\\
    \constd{SIGEV\_THREAD\_ID}& Invia la notifica come segnale (con le stesse
                             modalità di \const{SIGEV\_SIGNAL}) che però viene
                             recapitato al \textit{thread} indicato dal campo
                             \var{sigev\_notify\_thread\_id}. Questa modalità
                             è una estensione specifica di Linux, creata come
                             supporto per le librerie di gestione dei
                             \textit{thread}, pertanto non deve essere usata
                             da codice normale.\\
    \hline
  \end{tabular}
  \caption{Valori possibili per il campo \var{sigev\_notify} in una struttura
    \struct{sigevent}.} 
  \label{tab:sigevent_sigev_notify}
\end{table}

\footnotetext{nel caso dei \textit{timer} questa funzionalità è considerata un
  esempio di pessima implementazione di una interfaccia, richiesta dallo
  standard POSIX, ma da evitare totalmente nell'uso ordinario, a causa della
  possibilità di creare disservizi generando una gran quantità di processi,
  tanto che ne è stata richiesta addirittura la rimozione.}

Nel caso di \func{timer\_create} occorrerà passare alla funzione come secondo
argomento l'indirizzo di una di queste strutture per indicare le modalità con
cui si vuole essere notificati della scadenza del timer, se non si specifica
nulla (passando un valore \val{NULL}) verrà inviato il segnale
\signal{SIGALRM} al processo corrente, o per essere più precisi verrà
utilizzato un valore equivalente all'aver specificato \const{SIGEV\_SIGNAL}
per \var{sigev\_notify}, \signal{SIGALRM} per \var{sigev\_signo} e
l'identificatore del timer come valore per \var{sigev\_value.sival\_int}.

Il terzo argomento deve essere l'indirizzo di una variabile di tipo
\typed{timer\_t} dove sarà scritto l'identificativo associato al timer appena
creato, da usare in tutte le successive funzioni di gestione. Una volta creato
questo identificativo resterà univoco all'interno del processo stesso fintanto
che il timer non viene cancellato.

Si tenga presente che eventuali POSIX timer creati da un processo non vengono
ereditati dai processi figli creati con \func{fork} e che vengono cancellati
nella esecuzione di un programma diverso attraverso una delle funzioni
\func{exec}. Si tenga presente inoltre che il kernel prealloca l'uso di un
segnale \textit{real-time} per ciascun timer che viene creato con
\func{timer\_create}; dato che ciascuno di essi richiede un posto nella coda
dei segnali \textit{real-time}, il numero massimo di timer utilizzabili da un
processo è limitato dalle dimensioni di detta coda, ed anche, qualora questo
sia stato impostato, dal limite \const{RLIMIT\_SIGPENDING}.

Una volta creato il timer \func{timer\_create} ed ottenuto il relativo
identificatore, si può attivare o disattivare un allarme (in gergo
\textsl{armare} o \textsl{disarmare} il timer) con la funzione di sistema
\funcd{timer\_settime}, il cui prototipo è:

\begin{funcproto}{
\fhead{signal.h}
\fhead{time.h}
\fdecl{int timer\_settime(timer\_t timerid, int flags, const struct
    itimerspec *new\_value, struct itimerspec *old\_value)}
\fdesc{Arma o disarma un timer POSIX.} 
}

{La funzione ritorna $0$ in caso di successo e $-1$ per un errore, nel qual
  caso \var{errno} assumerà uno dei valori: 
  \begin{errlist}
  \item[\errcode{EFAULT}] si è specificato un indirizzo non valido
    per \param{new\_value} o \param{old\_value}.
  \item[\errcode{EINVAL}] all'interno di \param{new\_value.value} si è
    specificato un tempo negativo o un numero di nanosecondi maggiore di
    999999999.
  \end{errlist}
}
\end{funcproto}

La funzione richiede che si indichi la scadenza del timer con
l'argomento \param{new\_value}, che deve essere specificato come puntatore ad
una struttura di tipo \struct{itimerspec}, la cui definizione è riportata in
fig.~\ref{fig:struct_itimerspec}; se il puntatore \param{old\_value} è diverso
da \val{NULL} il valore corrente della scadenza verrà restituito in una
analoga struttura, ovviamente in entrambi i casi le strutture devono essere
state allocate.

\begin{figure}[!htb]
  \footnotesize \centering
  \begin{minipage}[c]{0.8\textwidth}
    \includestruct{listati/itimerspec.h}
  \end{minipage} 
  \normalsize 
  \caption{La struttura \structd{itimerspec}, usata per specificare la
    scadenza di un allarme.}
  \label{fig:struct_itimerspec}
\end{figure}

Ciascuno dei due campi di \struct{itimerspec} indica un tempo, da specificare
con una precisione fino al nanosecondo tramite una struttura \struct{timespec}
(la cui definizione è riportata fig.~\ref{fig:sys_timespec_struct}). Il campo
\var{it\_value} indica la prima scadenza dell'allarme. Di default, quando il
valore di \param{flags} è nullo, questo valore viene considerato come un
intervallo relativo al tempo corrente, il primo allarme scatterà cioè dopo il
numero di secondi e nanosecondi indicati da questo campo. Se invece si usa
per \param{flags} il valore \constd{TIMER\_ABSTIME}, che al momento è l'unico
valore valido per \param{flags}, allora \var{it\_value} viene considerato come
un valore assoluto rispetto al valore usato dall'orologio a cui è associato il
timer. 

Quindi a seconda dei casi si potrà impostare un timer o con un tempo assoluto,
quando si opera rispetto all'orologio di sistema (nel qual caso il valore deve
essere in secondi e nanosecondi dalla \textit{epoch}) o con un numero di
secondi o nanosecondi rispetto alla partenza di un orologio di CPU, quando si
opera su uno di questi.  Infine un valore nullo di \var{it\_value}, dove per
nullo si intende con valori nulli per entrambi i campi \var{tv\_sec} e
\var{tv\_nsec}, può essere utilizzato, indipendentemente dal tipo di orologio
utilizzato, per disarmare l'allarme.

Il campo \var{it\_interval} di \struct{itimerspec} viene invece utilizzato per
impostare un allarme periodico.  Se il suo valore è nullo, se cioè sono nulli
tutti e due i due campi \var{tv\_sec} e \var{tv\_nsec} di detta struttura
\struct{timespec}, l'allarme scatterà una sola volta secondo quando indicato
con \var{it\_value}, altrimenti il valore specificato nella struttura verrà
preso come l'estensione del periodo di ripetizione della generazione
dell'allarme, che proseguirà indefinitamente fintanto che non si disarmi il
timer.

Se il timer era già stato armato la funzione sovrascrive la precedente
impostazione, se invece si indica come prima scadenza un tempo già passato,
l'allarme verrà notificato immediatamente e al contempo verrà incrementato il
contatore dei superamenti. Questo contatore serve a fornire una indicazione al
programma che riceve l'allarme su un eventuale numero di scadenze che sono
passate prima della ricezione della notifica dell'allarme. 

É infatti possibile, qualunque sia il meccanismo di notifica scelto, che
quest'ultima venga ricevuta dopo che il timer è scaduto più di una volta,
specialmente se si imposta un timer con una ripetizione a frequenza
elevata. Nel caso dell'uso di un segnale infatti il sistema mette in coda un
solo segnale per timer,\footnote{questo indipendentemente che si tratti di un
  segnale ordinario o \textit{real-time}, per questi ultimi sarebbe anche
  possibile inviare un segnale per ogni scadenza, questo però non viene fatto
  per evitare il rischio, tutt'altro che remoto, di riempire la coda.} e se il
sistema è sotto carico o se il segnale è bloccato, prima della sua ricezione
può passare un intervallo di tempo sufficientemente lungo ad avere scadenze
multiple, e lo stesso può accadere anche se si usa un \textit{thread} di
notifica.

Per questo motivo il gestore del segnale o il \textit{thread} di notifica può
ottenere una indicazione di quante volte il timer è scaduto dall'invio della
notifica utilizzando la funzione di sistema \funcd{timer\_getoverrun}, il cui
prototipo è:

\begin{funcproto}{
\fhead{time.h}
\fdecl{int timer\_getoverrun(timer\_t timerid)}
\fdesc{Ottiene il numero di scadenze di un timer POSIX.} 
}

{La funzione ritorna il numero di scadenze di un timer in caso di successo e
  $-1$ per un errore, nel qual caso \var{errno} assumerà uno dei valori:
  \begin{errlist}
  \item[\errcode{EINVAL}] \param{timerid} non indica un timer valido.
  \end{errlist}
}
\end{funcproto}

La funzione ritorna il numero delle scadenze avvenute, che può anche essere
nullo se non ve ne sono state. Come estensione specifica di Linux,\footnote{in
  realtà lo standard POSIX.1-2001 prevede gli \textit{overrun} solo per i
  segnali e non ne parla affatto in riferimento ai \textit{thread}.}  quando
si usa un segnale come meccanismo di notifica, si può ottenere direttamente
questo valore nel campo \var{si\_overrun} della struttura \struct{siginfo\_t}
(illustrata in fig.~\ref{fig:sig_siginfo_t}) restituita al gestore del segnale
installato con \func{sigaction}; in questo modo non è più necessario eseguire
successivamente una chiamata a questa funzione per ottenere il numero delle
scadenze. Al gestore del segnale viene anche restituito, come ulteriore
informazione, l'identificativo del timer, in questo caso nel campo
\var{si\_timerid}.

Qualora si voglia rileggere lo stato corrente di un timer, ed ottenere il
tempo mancante ad una sua eventuale scadenza, si deve utilizzare la funzione
di sistema \funcd{timer\_gettime}, il cui prototipo è:

\begin{funcproto}{
\fhead{time.h}
\fdecl{int timer\_gettime(timer\_t timerid, int flags, struct
    itimerspec *curr\_value)}
\fdesc{Legge lo stato di un timer POSIX.} 
}

{La funzione ritorna $0$ in caso di successo e $-1$ per un errore, nel qual
  caso \var{errno} assumerà uno dei valori: 
  \begin{errlist}
  \item[\errcode{EFAULT}] si è specificato un indirizzo non valido
    per \param{curr\_value}.
  \item[\errcode{EINVAL}] \param{timerid} non indica un timer valido.
  \end{errlist}
}
\end{funcproto}

La funzione restituisce nella struttura \struct{itimerspec} puntata
da \param{curr\_value} il tempo restante alla prossima scadenza nel campo
\var{it\_value}. Questo tempo viene sempre indicato in forma relativa, anche
nei casi in cui il timer era stato precedentemente impostato con
\const{TIMER\_ABSTIME} indicando un tempo assoluto.  Il ritorno di un valore
nullo nel campo \var{it\_value} significa che il timer è disarmato o è
definitivamente scaduto. 

Nel campo \var{it\_interval} di \param{curr\_value} viene invece restituito,
se questo era stato impostato, il periodo di ripetizione del timer.  Anche in
questo caso il ritorno di un valore nullo significa che il timer non era stato
impostato per una ripetizione e doveva operare, come suol dirsi, a colpo
singolo (in gergo \textit{one shot}).

Infine, quando un timer non viene più utilizzato, lo si può cancellare,
rimuovendolo dal sistema e recuperando le relative risorse, effettuando in
sostanza l'operazione inversa rispetto a \func{timer\_create}. Per questo
compito lo standard prevede una apposita funzione di sistema,
\funcd{timer\_delete}, il cui prototipo è:

\begin{funcproto}{
\fhead{time.h}
\fdecl{int timer\_delete(timer\_t timerid)}
\fdesc{Cancella un timer POSIX.} 
}

{La funzione ritorna $0$ in caso di successo e $-1$ per un errore, nel qual
  caso \var{errno} assumerà uno dei valori: 
  \begin{errlist}
    \item[\errcode{EINVAL}] \param{timerid} non indica un timer valido.
  \end{errlist}
}
\end{funcproto}

La funzione elimina il timer identificato da \param{timerid}, disarmandolo se
questo era stato attivato. Nel caso, poco probabile ma comunque possibile, che
un timer venga cancellato prima della ricezione del segnale pendente per la
notifica di una scadenza, il comportamento del sistema è indefinito.

Infine a partire dal kernel 2.6 e per le versioni della \acr{libc} superiori
alla 2.1, si può utilizzare la nuova interfaccia dei timer POSIX anche per le
funzioni di attesa, per questo è disponibile la funzione di sistema
\funcd{clock\_nanosleep}, il cui prototipo è:

\begin{funcproto}{
\fhead{time.h}
\fdecl{int clock\_nanosleep(clockid\_t clock\_id, int flags, const struct
    timespec *request,\\
\phantom{int clock\_nanosleep(}struct timespec *remain)}
\fdesc{Pone il processo in pausa per un tempo specificato.}
}

{La funzione ritorna $0$ in caso di successo ed un valore positivo per un
  errore, espresso dai valori:
  \begin{errlist}
    \item[\errcode{EINTR}] la funzione è stata interrotta da un segnale.
    \item[\errcode{EINVAL}] si è specificato un numero di secondi negativo o
      un numero di nanosecondi maggiore di 999.999.999 o indicato un orologio
      non valido.
  \end{errlist}
  ed inoltre \errval{EFAULT} nel suo significato generico.}
\end{funcproto}

I due argomenti \param{request} e \param{remain} sono identici agli analoghi di
\func{nanosleep} che abbiamo visto in sez.~\ref{sec:sig_pause_sleep}, ed hanno
lo stesso significato.  L'argomento \param{clock\_id} consente di indicare
quale orologio si intende utilizzare per l'attesa con uno dei valori della
prima parte di tab.~\ref{tab:sig_timer_clockid_types} (eccetto
\const{CLOCK\_THREAD\_CPUTIME\_ID}).  L'argomento \param{flags} consente di
modificare il comportamento della funzione, il suo unico valore valido al
momento è \const{TIMER\_ABSTIME} che, come per \func{timer\_settime} indica di
considerare il tempo indicato in \param{request} come assoluto anziché
relativo.

Il comportamento della funzione è analogo a \func{nanosleep}, se la chiamata
viene interrotta il tempo rimanente viene restituito in \param{remain}.
Utilizzata normalmente con attese relative può soffrire degli stessi problemi
di deriva di cui si è parlato in sez.~\ref{sec:sig_pause_sleep} dovuti ad
interruzioni ripetute per via degli arrotondamenti fatti a questo tempo.  Ma
grazie alla possibilità di specificare tempi assoluti con \param{flags} si può
ovviare a questo problema ricavando il tempo corrente con
\func{clock\_gettime}, aggiungendovi l'intervallo di attesa, ed impostando
questa come tempo assoluto.

Si tenga presente che se si è usato il valore \const{TIMER\_ABSTIME}
per \param{flags} e si è indicato un tempo assoluto che è già passato la
funzione ritorna immediatamente senza nessuna sospensione. In caso di
interruzione da parte di un segnale il tempo rimanente viene restituito
in \param{remain} soltanto se questo non è un puntatore \val{NULL} e non si è
specificato \const{TIMER\_ABSTIME} per \param{flags}.


\itindend{POSIX~Timer~API}



\subsection{Ulteriori funzioni di gestione}
\label{sec:sig_specific_features}

In questo ultimo paragrafo esamineremo le rimanenti funzioni di gestione dei
segnali non descritte finora, relative agli aspetti meno utilizzati e più
``\textsl{esoterici}'' della interfaccia.

% TODO: trattare (qui?) pidfd_send_signal() introdotta con il kernel 5.1 vedi
% https://lwn.net/Articles/784831/ e https://lwn.net/Articles/773459/

La prima di queste funzioni è la funzione di sistema \funcd{sigpending},
anch'essa introdotta dallo standard POSIX.1, il suo prototipo è:

\begin{funcproto}{
\fhead{signal.h}
\fdecl{int sigpending(sigset\_t *set)}
\fdesc{Legge l'insieme dei segnali pendenti.} 
}

{La funzione ritorna $0$ in caso di successo e $-1$ per un errore, nel qual
  caso \var{errno} assumerà solo il valore \errcode{EFAULT} nel suo
  significato generico.}
\end{funcproto}

La funzione permette di ricavare quali sono i segnali pendenti per il processo
in corso, cioè i segnali che sono stati inviati dal kernel ma non sono stati
ancora ricevuti dal processo in quanto bloccati. Non esiste una funzione
equivalente nella vecchia interfaccia, ma essa è tutto sommato poco utile,
dato che essa può solo assicurare che un segnale è stato inviato, dato che
escluderne l'avvenuto invio al momento della chiamata non significa nulla
rispetto a quanto potrebbe essere in un qualunque momento successivo.

Una delle caratteristiche di BSD, disponibile anche in Linux, è la possibilità
di usare uno \textit{stack} alternativo per i segnali; è cioè possibile fare
usare al sistema un altro \textit{stack} (invece di quello relativo al
processo, vedi sez.~\ref{sec:proc_mem_layout}) solo durante l'esecuzione di un
gestore.  L'uso di uno \textit{stack} alternativo è del tutto trasparente ai
gestori, occorre però seguire una certa procedura:
\begin{enumerate*}
\item allocare un'area di memoria di dimensione sufficiente da usare come
  \textit{stack} alternativo;
\item usare la funzione \func{sigaltstack} per rendere noto al sistema
  l'esistenza e la locazione dello \textit{stack} alternativo;
\item quando si installa un gestore occorre usare \func{sigaction}
  specificando il flag \const{SA\_ONSTACK} (vedi tab.~\ref{tab:sig_sa_flag})
  per dire al sistema di usare lo \textit{stack} alternativo durante
  l'esecuzione del gestore.
\end{enumerate*}

In genere il primo passo viene effettuato allocando un'opportuna area di
memoria con \code{malloc}; in \headfile{signal.h} sono definite due costanti,
\constd{SIGSTKSZ} e \constd{MINSIGSTKSZ}, che possono essere utilizzate per
allocare una quantità di spazio opportuna, in modo da evitare overflow. La
prima delle due è la dimensione canonica per uno \textit{stack} di segnali e
di norma è sufficiente per tutti gli usi normali.

La seconda è lo spazio che occorre al sistema per essere in grado di lanciare
il gestore e la dimensione di uno \textit{stack} alternativo deve essere
sempre maggiore di questo valore. Quando si conosce esattamente quanto è lo
spazio necessario al gestore gli si può aggiungere questo valore per allocare
uno \textit{stack} di dimensione sufficiente.

Come accennato, per poter essere usato, lo \textit{stack} per i segnali deve
essere indicato al sistema attraverso la funzione \funcd{sigaltstack}; il suo
prototipo è:

\begin{funcproto}{
\fhead{signal.h}
\fdecl{int sigaltstack(const stack\_t *ss, stack\_t *oss)}
\fdesc{Installa uno \textit{stack} alternativo per i gestori di segnali.} 
}

{La funzione ritorna $0$ in caso di successo e $-1$ per un errore, nel qual
  caso \var{errno} assumerà uno dei valori: 
  \begin{errlist}
  \item[\errcode{EFAULT}] uno degli indirizzi degli argomenti non è valido.
  \item[\errcode{EINVAL}] \param{ss} non è nullo e \var{ss\_flags} contiene un
  valore diverso da zero che non è \const{SS\_DISABLE}.
  \item[\errcode{ENOMEM}] la dimensione specificata per il nuovo
    \textit{stack} è minore di \const{MINSIGSTKSZ}.
  \item[\errcode{EPERM}] si è cercato di cambiare lo \textit{stack}
    alternativo mentre questo è attivo (cioè il processo è in esecuzione su di
    esso).
  \end{errlist}
}
\end{funcproto}

La funzione prende come argomenti puntatori ad una struttura di tipo
\var{stack\_t}, definita in fig.~\ref{fig:sig_stack_t}. I due valori
\param{ss} e \param{oss}, se non nulli, indicano rispettivamente il nuovo
\textit{stack} da installare e quello corrente (che viene restituito dalla
funzione per un successivo ripristino).

\begin{figure}[!htb]
  \footnotesize \centering
  \begin{minipage}[c]{0.8\textwidth}
    \includestruct{listati/stack_t.h}
  \end{minipage} 
  \normalsize 
  \caption{La struttura \structd{stack\_t}.} 
  \label{fig:sig_stack_t}
\end{figure}

Il campo \var{ss\_sp} di \struct{stack\_t} indica l'indirizzo base dello
\textit{stack}, mentre \var{ss\_size} ne indica la dimensione; il campo
\var{ss\_flags} invece indica lo stato dello \textit{stack}.  Nell'indicare un
nuovo \textit{stack} occorre inizializzare \var{ss\_sp} e \var{ss\_size}
rispettivamente al puntatore e alla dimensione della memoria allocata, mentre
\var{ss\_flags} deve essere nullo.  Se invece si vuole disabilitare uno
\textit{stack} occorre indicare \constd{SS\_DISABLE} come valore di
\var{ss\_flags} e gli altri valori saranno ignorati.

Se \param{oss} non è nullo verrà restituito dalla funzione indirizzo e
dimensione dello \textit{stack} corrente nei relativi campi, mentre
\var{ss\_flags} potrà assumere il valore \constd{SS\_ONSTACK} se il processo è
in esecuzione sullo \textit{stack} alternativo (nel qual caso non è possibile
cambiarlo) e \const{SS\_DISABLE} se questo non è abilitato.

In genere si installa uno \textit{stack} alternativo per i segnali quando si
teme di avere problemi di esaurimento dello \textit{stack} standard o di
superamento di un limite (vedi sez.~\ref{sec:sys_resource_limit}) imposto con
chiamate del tipo \code{setrlimit(RLIMIT\_STACK, \&rlim)}.  In tal caso
infatti si avrebbe un segnale di \signal{SIGSEGV}, che potrebbe essere gestito
soltanto avendo abilitato uno \textit{stack} alternativo.

Si tenga presente che le funzioni chiamate durante l'esecuzione sullo
\textit{stack} alternativo continueranno ad usare quest'ultimo, che, al
contrario di quanto avviene per lo \textit{stack} ordinario dei processi, non
si accresce automaticamente (ed infatti eccederne le dimensioni può portare a
conseguenze imprevedibili).  Si ricordi infine che una chiamata ad una
funzione della famiglia \func{exec} cancella ogni \textit{stack} alternativo.

Abbiamo visto in fig.~\ref{fig:sig_sleep_incomplete} come si possa usare
\func{longjmp} per uscire da un gestore rientrando direttamente nel corpo
del programma, sappiamo però che nell'esecuzione di un gestore il segnale
che l'ha invocato viene bloccato, e abbiamo detto che possiamo ulteriormente
modificarlo con \func{sigprocmask}. 

Resta quindi il problema di cosa succede alla maschera dei segnali quando si
esce da un gestore usando questa funzione. Il comportamento dipende
dall'implementazione. In particolare la semantica usata da BSD prevede che sia
ripristinata la maschera dei segnali precedente l'invocazione, come per un
normale ritorno, mentre quella usata da System V no.

Lo standard POSIX.1 non specifica questo comportamento per \func{setjmp} e
\func{longjmp}, ed il comportamento della \acr{glibc} dipende da quale delle
caratteristiche si sono abilitate con le macro viste in
sez.~\ref{sec:intro_gcc_glibc_std}.

Lo standard POSIX però prevede anche la presenza di altre due funzioni
\funcd{sigsetjmp} e \funcd{siglongjmp}, che permettono di decidere in maniera
esplicita quale dei due comportamenti il programma deve assumere; i loro
prototipi sono:

\begin{funcproto}{
\fhead{setjmp.h}
\fdecl{int sigsetjmp(sigjmp\_buf env, int savesigs)}
\fdesc{Salva il contesto dello \textit{stack} e la maschera dei segnali.}  
\fdecl{void siglongjmp(sigjmp\_buf env, int val)}
\fdesc{Ripristina il contesto dello \textit{stack} e la maschera dei segnali.} 
}

{La funzioni sono identiche alle analoghe \func{setjmp} e \func{longjmp} di
  sez.~\ref{sec:proc_longjmp} ed hanno gli stessi errori e valori di uscita.}
\end{funcproto}

Le due funzioni prendono come primo argomento la variabile su cui viene
salvato il contesto dello \textit{stack} per permettere il salto non-locale;
nel caso specifico essa è di tipo \typed{sigjmp\_buf}, e non \type{jmp\_buf}
come per le analoghe di sez.~\ref{sec:proc_longjmp} in quanto in questo caso
viene salvata anche la maschera dei segnali.

Nel caso di \func{sigsetjmp}, se si specifica un valore di \param{savesigs}
diverso da zero la maschera dei valori verrà salvata in \param{env} insieme al
contesto dello \textit{stack} usato per il salto non locale. Se così si è
fatto la maschera dei segnali verrà ripristinata in una successiva chiamata a
\func{siglongjmp}. Quest'ultima, a parte l'uso di un valore di \param{env} di
tipo \type{sigjmp\_buf}, è assolutamente identica a \func{longjmp}.


% TODO: se e quando si troverà un argomento adeguato inserire altre funzioni
% sparse attinenti ai segnali, al momento sono note solo:
% * sigreturn (funzione interna, scarsamente interessante)
% argomento a priorità IDLE (fare quando non resta niente altro da trattare)


% LocalWords:  kernel POSIX timer shell control ctrl kill raise signal handler
% LocalWords:  reliable unreliable fig race condition sez struct process table
% LocalWords:  delivered pending scheduler sigpending l'I suspend SIGKILL wait
% LocalWords:  SIGSTOP sigaction waitpid dump stack debugger nell'header NSIG
% LocalWords:  tab BSD SUSv SIGHUP PL Hangup SIGINT Interrupt SIGQUIT Quit AEF
% LocalWords:  SIGILL SIGABRT abort SIGFPE SIGSEGV SIGPIPE SIGALRM alarm SIGUSR
% LocalWords:  SIGTERM SIGCHLD SIGCONT SIGTSTP SIGTTIN SIGTTOU SIGBUS bad SL of
% LocalWords:  memory access SIGPOLL Pollable event Sys SIGIO SIGPROF profiling
% LocalWords:  SIGSYS SVID SIGTRAP breakpoint SIGURG urgent socket Virtual IOT
% LocalWords:  clock SIGXCPU SIGXFSZ SIGIOT trap SIGEMT SIGSTKFLT SIGCLD SIGPWR
% LocalWords:  SIGINFO SIGLOST lock NFS SIGWINCH Sun SIGUNUSED fault point heap
% LocalWords:  exception l'overflow illegal instruction overflow segment error
% LocalWords:  violation system call interrupt INTR hang SIGVTALRM virtual SUSP
% LocalWords:  profilazione fcntl descriptor sleep interactive Broken FIFO lost
% LocalWords:  EPIPE Resource advisory client limit exceeded size window change
% LocalWords:  strsignal psignal SOURCE strerror string char int signum perror
% LocalWords:  void sig const sys siglist L'array decr fork exec DFL IGN ioctl
% LocalWords:  EINTR glibc TEMP FAILURE RETRY expr multitasking SVr sighandler
% LocalWords:  ERR libc bsd sysv XOPEN EINVAL pid errno ESRCH EPERM getpid init
% LocalWords:  killpg pidgrp group unistd unsigned seconds all' setitimer which
% LocalWords:  itimerval value ovalue EFAULT ITIMER it interval timeval ms VIRT
% LocalWords:  getitimer stdlib stream atexit exit usleep long usec nanosleep
% LocalWords:  timespec req rem HZ scheduling SCHED RR SigHand forktest WNOHANG
% LocalWords:  deadlock longjmp setjmp sigset sigemptyset sigfillset sigaddset
% LocalWords:  sigdelset sigismember act oldact restorer mask NOCLDSTOP ONESHOT
% LocalWords:  RESETHAND RESTART NOMASK NODEFER ONSTACK sigcontext union signo
% LocalWords:  siginfo bits uid addr fd inline like blocked atomic sigprocmask
% LocalWords:  how oldset BLOCK UNBLOCK SETMASK sigsuspend sigaltstack malloc
% LocalWords:  SIGSTKSZ MINSIGSTKSZ ss oss ENOMEM flags DISABLE sp setrlimit LB
% LocalWords:  RLIMIT rlim sigsetjmp siglongjmp sigjmp buf env savesigs jmp ptr
% LocalWords:  SIGRTMIN SIGRTMAX sigval sigevent sigqueue EAGAIN sysctl safe tp
% LocalWords:  QUEUE thread sigwait sigwaitinfo sigtimedwait info DEF SLB bind
% LocalWords:  function accept return cfgetispeed cfgetospeed cfsetispeed chdir
% LocalWords:  cfsetospeed chmod chown gettime close connect creat dup execle
% LocalWords:  execve fchmod fchown fdatasync fpathconf fstat fsync ftruncate
% LocalWords:  getegid geteuid getgid getgroups getpeername getpgrp getppid sem
% LocalWords:  getsockname getsockopt getuid listen lseek lstat mkdir mkfifo tv
% LocalWords:  pathconf poll posix pselect read readlink recv recvfrom recvmsg
% LocalWords:  rename rmdir select send sendmsg sendto setgid setpgid setsid
% LocalWords:  setsockopt setuid shutdown sigpause socketpair stat symlink page
% LocalWords:  sysconf tcdrain tcflow tcflush tcgetattr tcgetgrp tcsendbreak
% LocalWords:  tcsetattr tcsetpgrp getoverrun times umask uname unlink utime
% LocalWords:  write sival SIVGTALRM NOCLDWAIT MESGQ ASYNCIO TKILL tkill tgkill
% LocalWords:  ILL ILLOPC ILLOPN ILLADR ILLTRP PRVOPC PRVREG COPROC BADSTK FPE
% LocalWords:  INTDIV INTOVF FLTDIV FLTOVF FLTUND underflow FLTRES FLTINV SEGV
% LocalWords:  FLTSUB MAPERR ACCERR ADRALN ADRERR OBJERR BRKPT CLD EXITED MSG
% LocalWords:  KILLED DUMPED TRAPPED STOPPED CONTINUED PRI HUP SigFunc jiffies
% LocalWords:  SEC unsafe sockatmark execl execv faccessat fchmodat fchownat
% LocalWords:  fexecve fstatat futimens linkat mkdirat mkfifoat mknod mknodat
% LocalWords:  openat readlinkat renameat symlinkat unlinkat utimensat utimes
% LocalWords:  LinuxThread NTPL Library clockid evp timerid sigev notify high
% LocalWords:  resolution CONFIG RES patch REALTIME MONOTONIC RAW NTP CPUTIME
% LocalWords:  tick calendar The Epoch list getcpuclockid capability CAP getres
% LocalWords:  ENOSYS pthread ENOENT NULL attribute itimerspec new old ABSTIME
% LocalWords:  epoch multiplexing overrun res lpthread sec nsec curr one shot
% LocalWords:  delete stopped gdb alpha mips emulation locking ppoll epoll PGID
% LocalWords:  pwait msgrcv msgsnd semop semtimedop runnable sigisemptyset HRT
% LocalWords:  sigorset sigandset BOOTTIME Android request remain


%%% Local Variables: 
%%% mode: latex
%%% TeX-master: "gapil"
%%% End: 

%% session.tex
%%
%% Copyright (C) 2000-2018 Simone Piccardi.  Permission is granted to
%% copy, distribute and/or modify this document under the terms of the GNU Free
%% Documentation License, Version 1.1 or any later version published by the
%% Free Software Foundation; with the Invariant Sections being "Un preambolo",
%% with no Front-Cover Texts, and with no Back-Cover Texts.  A copy of the
%% license is included in the section entitled "GNU Free Documentation
%% License".
%%

\chapter{Terminali e sessioni di lavoro}
\label{cha:session}


A lungo l'unico modo per interagire con sistema di tipo Unix è stato tramite
l'interfaccia dei terminali, ma anche oggi, nonostante la presenza di diverse
interfacce grafiche, essi continuano ad essere estensivamente usati per il
loro stretto legame la linea di comando.

Nella prima parte esamineremo i concetti base in cui si articola l'interfaccia
dei terminali, a partire dal sistema del \textit{job control} e delle sessioni
di lavoro, toccando infine anche le problematiche dell'interazione con
programmi non interattivi. Nella seconda parte tratteremo il funzionamento
dell'I/O su terminale, e delle varie peculiarità che esso viene ad assumere
nell'uso come interfaccia di accesso al sistema da parte degli utenti. La
terza parte coprirà le tematiche relative alla creazione e gestione dei
terminali virtuali, che consentono di replicare via software l'interfaccia dei
terminali.



\section{L'interazione con i terminali}
\label{sec:sess_job_control}

I terminali sono l'interfaccia con cui fin dalla loro nascita i sistemi
unix-like hanno gestito l'interazione con gli utenti, tramite quella riga di
comando che li caratterizza da sempre. Ma essi hanno anche una rilevanza
particolare perché quella dei terminali è l'unica interfaccia hardware usata
dal kernel per comunicare direttamente con gli utenti, con la cosiddetta
\textit{console} di sistema, senza dover passare per un programma.

Originariamente si trattava di dispositivi specifici (i terminali seriali, se
non addirittura le telescriventi). Oggi questa interfaccia viene in genere
emulata o tramite programmi o con le cosiddette console virtuali associate a
monitor e tastiera, ma esiste sempre la possibilità di associarla direttamente
ad alcuni dispositivi, come eventuali linee seriali, ed in certi casi, come
buona parte dei dispositivi embedded su cui gira Linux (come router, access
point, ecc.) questa resta anche l'unica opzione per una \textit{console} di
sistema.


\subsection{Il \textit{job control}}
\label{sec:sess_job_control_overview}

Viene comunemente chiamato \textit{job control} quell'insieme di funzionalità
il cui scopo è quello di permettere ad un utente di poter sfruttare le
capacità multitasking di un sistema Unix per eseguire in contemporanea più
processi, pur potendo accedere, di solito, ad un solo terminale, avendo cioè
un solo punto in cui si può avere accesso all'input ed all'output degli
stessi. Con le interfacce grafiche di \textit{X Window} e con i terminali
virtuali via rete oggi tutto questo non è più vero, dato che si può accedere a
molti terminali in contemporanea da una singola postazione di lavoro, ma il
sistema è nato prima dell'esistenza di tutto ciò.

Il \textit{job control} è una caratteristica opzionale, introdotta in BSD
negli anni '80, e successivamente standardizzata da POSIX.1. La sua
disponibilità nel sistema è verificabile attraverso il controllo della macro
\macro{\_POSIX\_JOB\_CONTROL}. In generale il \textit{job control} richiede il
supporto sia da parte della shell (quasi tutte ormai lo hanno), che da parte
del kernel. In particolare il kernel deve assicurare sia la presenza di un
driver per i terminali abilitato al \textit{job control} che quella dei
relativi segnali illustrati in sez.~\ref{sec:sig_job_control}. 

In un sistema che supporta il \textit{job control}, una volta completato il
login, l'utente avrà a disposizione una shell dalla quale eseguire i comandi e
potrà iniziare quella che viene chiamata una \textsl{sessione di lavoro}, che
riunisce (vedi sez.~\ref{sec:sess_proc_group}) tutti i processi eseguiti
all'interno dello stesso login (esamineremo tutto il processo in dettaglio in
sez.~\ref{sec:sess_login}).

Siccome la shell è collegata ad un solo terminale, che viene usualmente
chiamato \textsl{terminale di controllo}, (vedi sez.~\ref{sec:sess_ctrl_term})
un solo comando alla volta, quello che viene detto in \textit{foreground} o in
\textsl{primo piano}, potrà scrivere e leggere dal terminale. La shell però
può eseguire, aggiungendo una ``\cmd{\&}'' alla fine del comando, più
programmi in contemporanea, mandandoli come si dice, ``in
\textit{background}'' (letteralmente ``\textsl{sullo sfondo}''), nel qual caso
essi saranno eseguiti senza essere collegati al terminale.

Si noti come si sia parlato di comandi e non di programmi o processi. Fra le
funzionalità della shell infatti c'è anche quella di consentire di concatenare
più comandi in una sola riga con il \textit{pipelining}, ed in tal caso
verranno eseguiti più programmi. Inoltre, anche quando si invoca un singolo
programma, questo potrà sempre lanciare eventuali sotto-processi per eseguire
dei compiti specifici.

Per questo l'esecuzione di una riga di comando può originare più di un
processo, quindi nella gestione del \textit{job control} non si può far
riferimento ai singoli processi.  Per questo il kernel prevede la possibilità
di raggruppare più processi in un cosiddetto \textit{process group} (detto
anche \textsl{raggruppamento di processi}, vedi
sez.~\ref{sec:sess_proc_group}). Deve essere cura della shell far sì che tutti
i processi che originano da una stessa riga di comando appartengano allo
stesso raggruppamento di processi, in modo che le varie funzioni di controllo,
ed i segnali inviati dal terminale, possano fare riferimento ad esso.

In generale all'interno di una sessione avremo un eventuale (può non esserci)
\textit{process group} in \textit{foreground}, che riunisce i processi che
possono accedere al terminale, e più \textit{process group} in
\textit{background}, che non possono accedervi. Il \textit{job control}
prevede che quando un processo appartenente ad un raggruppamento in
\textit{background} cerca di accedere al terminale, venga inviato un segnale a
tutti i processi del raggruppamento, in modo da bloccarli (vedi
sez.~\ref{sec:sess_ctrl_term}).

Un comportamento analogo si ha anche per i segnali generati dai comandi di
tastiera inviati dal terminale, che vengono inviati a tutti i processi del
raggruppamento in \textit{foreground}. In particolare \cmd{C-z} interrompe
l'esecuzione del comando, che può poi essere mandato in \textit{background}
con il comando \cmd{bg}. Il comando \cmd{fg} consente invece di mettere in
\textit{foreground} un comando precedentemente lanciato in
\textit{background}. Si tenga presente che \cmd{bg} e \cmd{fg} sono comandi
interni alla shell, che non comportano l'esecuzione di un programma esterno,
ma operazioni di gestione compiute direttamente dalla shell stessa.

Di norma la shell si cura anche di notificare all'utente, di solito prima
della stampa a video del prompt, lo stato dei vari processi. Essa infatti sarà
in grado, grazie all'uso della funzione di sistema \func{waitpid} (vedi
sez.~\ref{sec:proc_wait}), di rilevare sia i processi che sono terminati, sia
i raggruppamenti che sono bloccati, in quest'ultimo caso si dovrà usare la
specifica opzione \const{WUNTRACED}, secondo quanto già illustrato in
sez.~\ref{sec:proc_wait}.


\subsection{I \textit{process group} e le \textsl{sessioni}}
\label{sec:sess_proc_group}

\itindbeg{process~group}

Come accennato in sez.~\ref{sec:sess_job_control_overview} nel job control i
processi vengono raggruppati in \textit{process group} e \textsl{sessioni};
per far questo vengono utilizzati due ulteriori identificatori (oltre quelli
visti in sez.~\ref{sec:proc_pid}) che il kernel associa a ciascun
processo:\footnote{in Linux questi identificatori sono mantenuti nei campi
  \var{pgrp} e \var{session} della struttura \kstruct{task\_struct} definita
  in \file{include/linux/sched.h}.}  l'identificatore del \textit{process
  group} e l'identificatore della \textsl{sessione}, che vengono indicati
rispettivamente con le sigle \ids{PGID} e \ids{SID}, e sono mantenuti in
variabili di tipo \type{pid\_t}. I valori di questi identificatori possono
essere visualizzati dal comando \cmd{ps} usando l'opzione \cmd{-j}.

Un \textit{process group} è pertanto definito da tutti i processi che hanno lo
stesso \ids{PGID}; è possibile leggere il valore di questo identificatore con
le funzioni di sistema \funcd{getpgid} e \funcd{getpgrp}, i cui prototipi
sono:

\begin{funcproto}{
\fhead{unistd.h}
\fdecl{pid\_t getpgid(pid\_t pid)} 
\fdesc{Legge il \ids{PGID} di un processo.} 
\fdecl{pid\_t getpgrp(void)}
\fdesc{Legge il \ids{PGID} del processo corrente.} 
}

{Le funzioni ritornano il \ids{PGID} richiesto in caso di successo,
  \func{getpgrp} ha sempre successo mentre \func{getpgid} restituisce $-1$ per
  un errore, nel qual caso \var{errno} potrà assumere solo il valore
  \errval{ESRCH} se il processo indicato non esiste.
}
\end{funcproto}

Le due funzioni sono definite nello standard POSIX.1-2001, ma la prima deriva
da SVr4 e la seconda da BSD4.2 dove però è previsto possa prendere un
argomento per indicare il \ids{PID} di un altro processo. Si può riottenere
questo comportamento se di definisce la macro \macro{\_BSD\_SOURCE} e non sono
definite le altre macro che richiedono la conformità a POSIX, X/Open o SysV
(vedi sez.~\ref{sec:intro_standard}).

La funzione \func{getpgid} permette di specificare il \ids{PID} del processo
di cui si vuole sapere il \ids{PGID}. Un valore nullo per \param{pid}
restituisce il \ids{PGID} del processo corrente, che è il comportamento
ordinario di \func{getpgrp}, che di norma equivalente a \code{getpgid(0)}.

In maniera analoga l'identificatore della sessione di un processo (il
\ids{SID}) può essere letto dalla funzione di sistema \funcd{getsid}, il cui
prototipo è:

\begin{funcproto}{
\fhead{unistd.h}
\fdecl{pid\_t getsid(pid\_t pid)}
\fdesc{Legge il \ids{SID} di un processo.} 
}

{La funzione ritorna l'identificatore (un numero positivo) in caso di successo
  e $-1$ per un errore, nel qual caso \var{errno} assumerà uno dei valori:
  \begin{errlist}
    \item[\errcode{EPERM}] il processo selezionato non fa parte della stessa
      sessione del processo corrente (solo in alcune implementazioni).
    \item[\errcode{ESRCH}] il processo selezionato non esiste.
  \end{errlist}
}
\end{funcproto}

La funzione è stata introdotta in Linux a partire dal kernel 1.3.44, il
supporto nelle librerie del C è iniziato dalla versione 5.2.19. La funzione
non era prevista originariamente da POSIX.1, che parla solo di processi leader
di sessione, e non di identificatori di sessione, ma è prevista da SVr4 e fa
parte di POSIX.1-2001.  Per poterla utilizzare occorre definire la macro
\macro{\_XOPEN\_SOURCE} ad un valore maggiore o uguale di 500. Su Linux
l'errore \errval{EPERM} non viene mai restituito.

Entrambi gli identificatori, \ids{SID} e \ids{PGID}, vengono inizializzati
nella creazione di ciascun processo con lo stesso valore che hanno nel
processo padre, per cui un processo appena creato appartiene sempre allo
stesso raggruppamento e alla stessa sessione del padre. Vedremo a breve come
sia possibile creare più \textit{process group} all'interno della stessa
sessione, e spostare i processi dall'uno all'altro, ma sempre all'interno di
una stessa sessione.

Ciascun raggruppamento di processi ha sempre un processo principale, il
cosiddetto \itindex{process~group~leader} \textit{process group leader} o più
brevemente \textit{group leader}, che è identificato dall'avere un \ids{PGID}
uguale al suo \ids{PID}. In genere questo è il primo processo del
raggruppamento, che si incarica di lanciare tutti gli altri. Un nuovo
raggruppamento si crea con la funzione di sistema \funcd{setpgrp}, il cui
prototipo è:

\begin{funcproto}{
\fhead{unistd.h}
\fdecl{int setpgrp(void)}
\fdesc{Rende un processo \textit{group leader} di un nuovo gruppo.} 
}

{La funzione ritorna il valore del nuovo \textit{process group} e non sono
  previsti errori.}
\end{funcproto}

La funzione assegna al \ids{PGID} il valore del \ids{PID} del processo
corrente, rendendolo in tal modo \textit{group leader} di un nuovo
raggruppamento. Tutti i successivi processi da esso creati apparterranno (a
meno di non cambiare di nuovo il \ids{PGID}) al nuovo raggruppamento.

La versione illustrata è quella usata nella definizione di POSIX.1, in BSD
viene usata una funzione con questo nome, che però è identica a
\func{setpgid}, che vedremo a breve, negli argomenti e negli effetti. Nella
\acr{glibc} viene sempre usata sempre questa definizione, a meno di non
richiedere esplicitamente la compatibilità all'indietro con BSD, definendo la
macro \macro{\_BSD\_SOURCE} ed evitando di definire le macro che richiedono
gli altri standard, come per \func{getpgrp}.

È inoltre possibile spostare un processo da un raggruppamento di processi ad
un altro cambiandone il \ids{PGID} con la funzione di sistema \funcd{setpgid},
il cui prototipo è:

\begin{funcproto}{
\fhead{unistd.h}
\fdecl{int setpgid(pid\_t pid, pid\_t pgid)}
\fdesc{Modifica il \ids{PGID} di un processo.} 
}

{La funzione ritorna il valore del nuovo \textit{process group}  in caso di
  successo e $-1$ per un errore, nel qual caso \var{errno} assumerà uno dei
  valori:
  \begin{errlist}
    \item[\errcode{EACCES}] il processo di cui si vuole cambiare il \ids{PGID}
      ha già eseguito una \func{exec}.
    \item[\errcode{EINVAL}] il valore di \param{pgid} è negativo.
    \item[\errcode{EPERM}] il cambiamento non è consentito.
    \item[\errcode{ESRCH}] il processo selezionato non esiste.
  \end{errlist}
}
\end{funcproto}


La funzione permette di cambiare il \ids{PGID} del processo indicato
dall'argomento \param{pid}, ma il cambiamento può essere effettuato solo se
l'argomento \param{pgid} indica un \textit{process group} che è nella stessa
sessione del processo chiamante.  Inoltre la funzione può essere usata
soltanto sul processo corrente o su uno dei suoi figli, ed in quest'ultimo
caso ha successo soltanto se questo non ha ancora eseguito una
\func{exec}.\footnote{questa caratteristica è implementata dal kernel che
  mantiene allo scopo un altro campo, \var{did\_exec}, nella struttura
  \kstruct{task\_struct}.}  Specificando un valore nullo per \param{pid} si
indica il processo corrente, mentre specificando un valore nullo
per \param{pgid} si imposta il \textit{process group} al valore del \ids{PID}
del processo selezionato, questo significa che \func{setpgrp} è equivalente a
\code{setpgid(0, 0)}.

Di norma questa funzione viene usata dalla shell quando si usano delle
pipeline, per mettere nello stesso \textit{process group} tutti i programmi
lanciati su ogni linea di comando; essa viene chiamata dopo una \func{fork}
sia dal processo padre, per impostare il valore nel figlio, che da
quest'ultimo, per sé stesso, in modo che il cambiamento di \textit{process
  group} sia immediato per entrambi; una delle due chiamate sarà ridondante,
ma non potendo determinare quale dei due processi viene eseguito per primo,
occorre eseguirle comunque entrambe per evitare di esporsi ad una \textit{race
  condition}.

Si noti come nessuna delle funzioni esaminate finora permetta di spostare un
processo da una sessione ad un altra; infatti l'unico modo di far cambiare
sessione ad un processo è quello di crearne una nuova con l'uso della funzione
di sistema \funcd{setsid}, il cui prototipo è:

\begin{funcproto}{
\fhead{unistd.h}
\fdecl{pid\_t setsid(void)}
\fdesc{Crea una nuova sessione sul processo corrente.} 
}

{La funzione ritorna il valore del nuovo \ids{SID} in caso di successo e $-1$
  per un errore, nel qual caso \var{errno} assumerà uno dei valori:
  \begin{errlist}
  \item[\errcode{EPERM}] il \ids{PGID} e \ids{PID} del processo coincidono.
  \end{errlist}
}
\end{funcproto}


La funzione imposta il \ids{PGID} ed il \ids{SID} del processo corrente al
valore del suo \ids{PID}, creando così una nuova sessione ed un nuovo
\textit{process group} di cui esso diventa leader (come per i \textit{process
  group} un processo si dice \textsl{leader di sessione} se il suo \ids{SID} è
uguale al suo \ids{PID}) ed unico componente.\footnote{in Linux la proprietà è
  mantenuta in maniera indipendente con un apposito campo \var{leader} in
  \kstruct{task\_struct}.} Inoltre la funzione distacca il processo da ogni
terminale di controllo (torneremo sull'argomento in
sez.~\ref{sec:sess_ctrl_term}) cui fosse in precedenza associato.

La funzione ha successo soltanto se il processo non è già leader di un
\textit{process group}, per cui per usarla di norma si esegue una \func{fork}
e si esce, per poi chiamare \func{setsid} nel processo figlio, in modo che,
avendo questo lo stesso \ids{PGID} del padre ma un \ids{PID} diverso, non ci
siano possibilità di errore.

Potrebbe sorgere il dubbio che, per il riutilizzo dei valori dei \ids{PID}
fatto nella creazione dei nuovi processi (vedi sez.~\ref{sec:proc_pid}), il
figlio venga ad assumere un valore corrispondente ad un \textit{process group}
esistente; questo viene evitato dal kernel che considera come disponibili per
un nuovo \ids{PID} solo valori che non corrispondono ad altri \ids{PID},
\ids{PGID} o \ids{SID} in uso nel sistema. Questa funzione viene usata di
solito nel processo di login (per i dettagli vedi sez.~\ref{sec:sess_login})
per raggruppare in una sessione tutti i comandi eseguiti da un utente dalla
sua shell.

\itindend{process~group}

\subsection{Il terminale di controllo e il controllo di sessione}
\label{sec:sess_ctrl_term}

Come accennato in sez.~\ref{sec:sess_job_control_overview}, nel sistema del
\textit{job control} i processi all'interno di una sessione fanno riferimento
ad un terminale di controllo (ad esempio quello su cui si è effettuato il
login), sul quale effettuano le operazioni di lettura e scrittura, e dal quale
ricevono gli eventuali segnali da tastiera. Nel caso di login grafico la cosa
può essere più complessa, e di norma l'I/O è effettuato tramite il server X,
ma ad esempio per i programmi, anche grafici, lanciati da un qualunque
emulatore di terminale, sarà quest'ultimo a fare da terminale (virtuale) di
controllo.

Per realizzare questa funzionalità lo standard POSIX.1 prevede che ad ogni
sessione possa essere associato un terminale di controllo; in Linux questo
viene realizzato mantenendo fra gli attributi di ciascun processo anche qual'è
il suo terminale di controllo.\footnote{lo standard POSIX.1 non specifica
  nulla riguardo l'implementazione; in Linux anch'esso viene mantenuto nella
  solita struttura \kstruct{task\_struct}, nel campo \var{tty}.}  In generale
ogni processo eredita dal padre, insieme al \ids{PGID} e al \ids{SID} anche il
terminale di controllo (vedi sez.~\ref{sec:proc_fork}). In questo modo tutti
processi originati dallo stesso leader di sessione mantengono lo stesso
terminale di controllo.

Alla creazione di una nuova sessione con \func{setsid} ogni associazione con
il precedente terminale di controllo viene cancellata, ed il processo che è
divenuto un nuovo leader di sessione dovrà essere associato ad un nuovo
terminale di controllo.\footnote{questo però solo quando necessario, cosa che,
  come vedremo in sez.~\ref{sec:sess_daemon}, non è sempre vera.} Questo viene
fatto dal kernel la prima volta che il processo apre un terminale (cioè uno
dei vari file di dispositivo \file{/dev/tty*}). In tal caso questo diventa
automaticamente il terminale di controllo, ed il processo diventa il
\textsl{processo di controllo} di quella sessione. Questo avviene
automaticamente a meno di non avere richiesto esplicitamente il contrario
aprendo il terminale di controllo con il flag \const{O\_NOCTTY} (vedi
sez.~\ref{sec:file_open_close}). In questo Linux segue la semantica di SVr4;
BSD invece richiede che il terminale venga allocato esplicitamente con una
\func{ioctl} con il comando \constd{TIOCSCTTY}.

In genere, a meno di redirezioni, nelle sessioni di lavoro il terminale di
controllo è associato ai file standard (input, output ed error) dei processi
nella sessione, ma solo quelli che fanno parte del cosiddetto raggruppamento
di \textit{foreground}, possono leggere e scrivere in certo istante. Per
impostare il raggruppamento di \textit{foreground} di un terminale si usa la
funzione \funcd{tcsetpgrp}, il cui prototipo è:

\begin{funcproto}{
\fhead{unistd.h}
\fhead{termios.h}
\fdecl{int tcsetpgrp(int fd, pid\_t pgrpid)}
\fdesc{Imposta il \textit{process group} di \textit{foreground}.}
}

{La funzione ritorna $0$ in caso di successo e $-1$ per un errore, nel qual
  caso \var{errno} assumerà uno dei valori: 
  \begin{errlist}
    \item[\errcode{ENOSYS}] il sistema non supporta il \textit{job control}.
    \item[\errcode{ENOTTY}] il file \param{fd} non corrisponde al terminale di
      controllo del processo chiamante.
    \item[\errcode{EPERM}] il \textit{process group} specificato non è nella
    stessa sessione del processo chiamante.
  \end{errlist}
  ed inoltre \errval{EBADF} ed \errval{EINVAL} nel loro significato generico.}
\end{funcproto}

La funzione imposta a \param{pgrpid} il \textit{process group} di
\textit{foreground} del terminale associato al file descriptor \param{fd}. La
funzione può essere eseguita con successo solo da un processo che
ha \param{fd} come terminale di controllo della propria sessione;
inoltre \param{pgrpid} deve essere un \textit{process group} (non vuoto)
appartenente alla stessa sessione del processo chiamante.

Come accennato in sez.~\ref{sec:sess_job_control_overview}, tutti i processi
(e relativi raggruppamenti) che non fanno parte del gruppo di
\textit{foreground} sono detti in \textit{background}; se uno si essi cerca di
accedere al terminale di controllo provocherà l'invio da parte del kernel di
uno dei due segnali \signal{SIGTTIN} o \signal{SIGTTOU} (a seconda che
l'accesso sia stato in lettura o scrittura) a tutto il suo \textit{process
  group}; dato che il comportamento di default di questi segnali (si riveda
quanto esposto in sez.~\ref{sec:sig_job_control}) è di fermare il processo, di
norma questo comporta che tutti i membri del gruppo verranno fermati, ma non
si avranno condizioni di errore.\footnote{la shell in genere notifica comunque
  un avvertimento, avvertendo la presenza di processi bloccati grazie all'uso
  di \func{waitpid}.} Se però si bloccano o ignorano i due segnali citati, le
funzioni di lettura e scrittura falliranno con un errore di \errcode{EIO}.

In genere la funzione viene chiamata da un processo che è gruppo di
\textit{foreground} per passare l'accesso al terminale ad altri processi,
quando viene chiamata da un processo che non è nel gruppo di
\textit{foreground}, a meno che questi non stia bloccando o ignorando il
segnale \signal{SIGTTOU}, detto segnale verrà inviato a tutti i processi del
gruppo di cui il chiamante fa parte.

Un processo può controllare qual è il gruppo di \textit{foreground} associato
ad un terminale con la funzione \funcd{tcgetpgrp}, il cui prototipo è:

\begin{funcproto}{
\fhead{unistd.h}
\fhead{termios.h}
\fdecl{pid\_t tcgetpgrp(int fd)}
\fdesc{Legge il \textit{process group} di \textit{foreground}.} 
}

{La funzione ritorna il \ids{PGID} del gruppo di \textit{foreground} in caso
  di successo e $-1$ per un errore, nel qual caso \var{errno} assumerà uno dei
  valori:
  \begin{errlist}
    \item[\errcode{ENOTTY}] non c'è un terminale di controllo o \param{fd} non
      corrisponde al terminale di controllo del processo chiamante.
  \end{errlist}
  ed inoltre \errval{EBADF} nel suo significato generico.}
\end{funcproto}

La funzione legge il \textit{process group} di \textit{foreground} associato
al file descriptor \param{fd}, che deve essere un terminale, restituendolo
come risultato. Sia questa funzione che la precedente sono state introdotte
con lo standard POSIX.1-2001, ma su Linux sono realizzate utilizzando le
operazioni di \func{ioctl} con i comandi \constd{TIOCGPGRP} e
\constd{TIOCSPGRP}.

Si noti come entrambe le funzioni usino come argomento il valore di un file
descriptor, il risultato comunque non dipende dal file descriptor che si usa
ma solo dal terminale cui fa riferimento. Il kernel inoltre permette a ciascun
processo di accedere direttamente al suo terminale di controllo attraverso il
file speciale \file{/dev/tty}, che per ogni processo è un sinonimo per il
proprio terminale di controllo. Questo consente anche a processi che possono
aver rediretto l'output di accedere al terminale di controllo, pur non
disponendo più del file descriptor originario; un caso tipico è il programma
\cmd{crypt} che accetta la redirezione sullo \textit{standard input} di un
file da decifrare, ma deve poi leggere la password dal terminale.

Un'altra caratteristica del terminale di controllo usata nel \textit{job
  control} è che utilizzando su di esso le combinazioni di tasti speciali
(\texttt{C-z}, \texttt{C-c}, \texttt{C-y} e \texttt{C-|}) si farà sì che il
kernel invii i corrispondenti segnali (rispettivamente \signal{SIGTSTP},
\signal{SIGINT}, \signal{SIGQUIT} e \signal{SIGTERM}, trattati in
sez.~\ref{sec:sig_job_control}) a tutti i processi del raggruppamento di
\textit{foreground}; in questo modo la shell può gestire il blocco e
l'interruzione dei vari comandi.

Per completare la trattazione delle caratteristiche del \textit{job control}
legate al terminale di controllo, occorre prendere in considerazione i vari
casi legati alla terminazione anomala dei processi, che sono di norma gestite
attraverso il segnale \signal{SIGHUP}. Il nome del segnale deriva da
\textit{hungup}, termine che viene usato per indicare la condizione in cui il
terminale diventa inutilizzabile, (letteralmente sarebbe
\textsl{impiccagione}).

Quando si verifica questa condizione, ad esempio se si interrompe la linea, o
va giù la rete o più semplicemente si chiude forzatamente la finestra di
terminale su cui si stava lavorando, il kernel provvederà ad inviare il
segnale di \signal{SIGHUP} al processo di controllo. L'azione preimpostata in
questo caso è la terminazione del processo, il problema che si pone è cosa
accade agli altri processi nella sessione, che non han più un processo di
controllo che possa gestire l'accesso al terminale, che potrebbe essere
riutilizzato per qualche altra sessione.

Lo standard POSIX.1 prevede che quando il processo di controllo termina, che
ciò avvenga o meno per un \textit{hungup} del terminale (ad esempio si
potrebbe terminare direttamente la shell con \cmd{kill}) venga inviato un
segnale di \signal{SIGHUP} ai processi del raggruppamento di
\textit{foreground}. In questo modo essi potranno essere avvisati che non
esiste più un processo in grado di gestire il terminale (di norma tutto ciò
comporta la terminazione anche di questi ultimi).

Restano però gli eventuali processi in \textit{background}, che non ricevono
il segnale; in effetti se il terminale non dovesse più servire essi potrebbero
proseguire fino al completamento della loro esecuzione; ma si pone il problema
di come gestire quelli che sono bloccati, o che si bloccano nell'accesso al
terminale, in assenza di un processo che sia in grado di effettuare il
controllo dello stesso.

Questa è la situazione in cui si ha quello che viene chiamato un
\itindex{process~group~orphaned} \textit{orphaned process group}. Lo standard
POSIX.1 lo definisce come un \textit{process group} i cui processi hanno come
padri esclusivamente o altri processi nel raggruppamento, o processi fuori
della sessione.  Lo standard prevede inoltre che se la terminazione di un
processo fa sì che un raggruppamento di processi diventi orfano e se i suoi
membri sono bloccati, ad essi vengano inviati in sequenza i segnali di
\signal{SIGHUP} e \signal{SIGCONT}.

La definizione può sembrare complicata, e a prima vista non è chiaro cosa
tutto ciò abbia a che fare con il problema della terminazione del processo di
controllo.  Consideriamo allora cosa avviene di norma nel \textit{job
  control}: una sessione viene creata con \func{setsid} che crea anche un
nuovo \textit{process group}. Per definizione quest'ultimo è sempre
\textsl{orfano}, dato che il padre del leader di sessione è fuori dalla stessa
e il nuovo \textit{process group} contiene solo il leader di sessione. Questo
è un caso limite, e non viene emesso nessun segnale perché quanto previsto
dallo standard riguarda solo i raggruppamenti che diventano orfani in seguito
alla terminazione di un processo.\footnote{l'emissione dei segnali infatti
  avviene solo nella fase di uscita del processo, come una delle operazioni
  legate all'esecuzione di \func{\_exit}, secondo quanto illustrato in
  sez.~\ref{sec:proc_termination}.}

Il leader di sessione provvederà a creare nuovi raggruppamenti che a questo
punto non sono orfani in quanto esso resta padre per almeno uno dei processi
del gruppo (gli altri possono derivare dal primo). Alla terminazione del
leader di sessione però avremo che, come visto in
sez.~\ref{sec:proc_termination}, tutti i suoi figli vengono adottati da
\cmd{init}, che è fuori dalla sessione. Questo renderà orfani tutti i
\textit{process group} creati direttamente dal leader di sessione a meno di
non aver spostato con \func{setpgid} un processo da un gruppo ad un altro,
(cosa che di norma non viene fatta) i quali riceveranno, nel caso siano
bloccati, i due segnali; \signal{SIGCONT} ne farà proseguire l'esecuzione, ed
essendo stato nel frattempo inviato anche \signal{SIGHUP}, se non c'è un
gestore per quest'ultimo, i processi bloccati verranno automaticamente
terminati.



\subsection{Dal login alla shell}
\label{sec:sess_login}

L'organizzazione del sistema del job control è strettamente connessa alle
modalità con cui un utente accede al sistema per dare comandi, collegandosi ad
esso con un terminale, che sia questo realmente tale, come un VT100 collegato
ad una seriale o virtuale, come quelli associati a schermo e tastiera o ad una
connessione di rete. Dato che i concetti base sono gli stessi, e dato che alla
fine le differenze sono nel dispositivo cui il kernel associa i file standard
(vedi tab.~\ref{tab:file_std_files}) per l'I/O, tratteremo solo il caso
classico del terminale, in generale nel caso di login via rete o di terminali
lanciati dall'interfaccia grafica cambia anche il processo da cui ha origine
l'esecuzione della shell.

Abbiamo già brevemente illustrato in sez.~\ref{sec:intro_kern_and_sys} le
modalità con cui il sistema si avvia, e di come, a partire da \cmd{init},
vengano lanciati tutti gli altri processi. Adesso vedremo in maniera più
dettagliata le modalità con cui il sistema arriva a fornire ad un utente la
shell che gli permette di lanciare i suoi comandi su un terminale.

Nella maggior parte delle distribuzioni di GNU/Linux viene usata la procedura
di avvio di System V;\footnote{in realtà negli ultimi tempi questa situazione
  sta cambiando, e sono state proposte diversi possibili rimpiazzi per il
  tradizionale \texttt{init} di System V, come \texttt{upstart} o
  \texttt{systemd}, ma per quanto trattato in questa sezione il risultato
  finale non cambia, si avrà comunque il lancio di un programma che consenta
  l'accesso al terminale.}  questa prevede che \cmd{init} legga dal file di
configurazione \conffiled{/etc/inittab} quali programmi devono essere lanciati,
ed in quali modalità, a seconda del cosiddetto \textit{run level}, anch'esso
definito nello stesso file.

Tralasciando la descrizione del sistema dei \textit{run level}, (per il quale
si rimanda alla lettura delle pagine di manuale di \cmd{init} e di
\file{inittab} o alla trattazione in sez.~5.3 di \cite{AGL}) quello che
comunque viene sempre fatto è di eseguire almeno una istanza di un programma
che permetta l'accesso ad un terminale. Uno schema di massima della procedura
è riportato in fig.~\ref{fig:sess_term_login}.

\begin{figure}[!htb]
  \centering \includegraphics[width=13cm]{img/tty_login}
  \caption{Schema della procedura di login su un terminale.}
  \label{fig:sess_term_login}
\end{figure}

Un terminale, che esso sia un terminale effettivo, attaccato ad una seriale o
ad un altro tipo di porta di comunicazione, o una delle console virtuali
associate allo schermo, o un terminale virtuale ad uso degli emulatori o delle
sessioni di rete, viene sempre visto attraverso un apposito file di
dispositivo che presenta una serie di caratteristiche comuni che vanno a
costituire l'interfaccia generica di accesso ai terminali.

Per controllare un terminale fisico come la seriale o le console virtuali
dello schermo si usa di solito il programma \cmd{getty} (o una delle sue
varianti). Alla radice della catena che porta ad una shell per i comandi
perciò c'è sempre \cmd{init} che esegue prima una \func{fork} e poi una
\func{exec} per lanciare una istanza di questo programma, il tutto ripetuto
per ciascuno dei terminali che si vogliono attivare, secondo quanto indicato
dall'amministratore nel file di configurazione del programma,
\conffile{/etc/inittab}.

Quando viene lanciato da \cmd{init} il programma parte con i privilegi di
amministratore e con un ambiente vuoto; \cmd{getty} si cura di chiamare
\func{setsid} per creare una nuova sessione ed un nuovo \textit{process
  group}, e di aprire il terminale (che così diventa il terminale di controllo
della sessione) in lettura sullo \textit{standard input} ed in scrittura sullo
\textit{standard output} e sullo \textit{standard error}; inoltre effettuerà,
qualora servano, ulteriori impostazioni.\footnote{ad esempio, come qualcuno si
  sarà accorto scrivendo un nome di login in maiuscolo, può effettuare la
  conversione automatica dell'input in minuscolo, ponendosi in una modalità
  speciale che non distingue fra i due tipi di caratteri (a beneficio di
  alcuni vecchi terminali che non supportavano le minuscole).}  Alla fine il
programma stamperà un messaggio di benvenuto per poi porsi in attesa
dell'immissione del nome di un utente.

Una volta che si sia immesso il nome di login \cmd{getty} esegue direttamente
il programma \cmd{login} con una \func{execle}, passando come argomento la
stringa con il nome, ed un ambiente opportunamente costruito che contenga
quanto necessario; ad esempio di solito viene opportunamente inizializzata la
variabile di ambiente \envvar{TERM} per identificare il terminale su cui si
sta operando, a beneficio dei programmi che verranno lanciati in seguito.

A sua volta \cmd{login}, che mantiene i privilegi di amministratore, usa il
nome dell'utente per effettuare una ricerca nel database degli utenti (in
genere viene chiamata \func{getpwnam}, che abbiamo visto in
sez.~\ref{sec:sys_user_group}, per leggere la password e gli altri dati dal
database degli utenti) e richiede una password. Se l'utente non esiste o se la
password non corrisponde\footnote{il confronto non viene effettuato con un
  valore in chiaro; quanto immesso da terminale viene invece a sua volta
  criptato, ed è il risultato che viene confrontato con il valore che viene
  mantenuto nel database degli utenti.} la richiesta viene ripetuta un certo
numero di volte dopo di che \cmd{login} esce ed \cmd{init} provvede a
rilanciare un'altra istanza di \cmd{getty}.

Se invece la password corrisponde \cmd{login} esegue \func{chdir} per
impostare come directory di lavoro la \textit{home directory} dell'utente,
cambia i diritti di accesso al terminale (con \func{chown} e \func{chmod}) per
assegnarne la titolarità all'utente ed al suo gruppo principale, assegnandogli
al contempo i diritti di lettura e scrittura.\footnote{oggi queste operazioni,
  insieme ad altre relative alla contabilità ed alla tracciatura degli
  accessi, vengono gestite dalle distribuzioni più recenti in una maniera
  generica appoggiandosi a servizi di sistema come \textit{ConsoleKit}, ma il
  concetto generale resta sostanzialmente lo stesso.}  Inoltre il programma
provvede a costruire gli opportuni valori per le variabili di ambiente, come
\envvar{HOME}, \envvar{SHELL}, ecc.  Infine attraverso l'uso di \func{setuid},
\func{setgid} e \func{initgroups} verrà cambiata l'identità del proprietario
del processo, infatti, come spiegato in sez.~\ref{sec:proc_setuid}, avendo
invocato tali funzioni con i privilegi di amministratore, tutti gli \ids{UID}
ed i \ids{GID} (reali, effettivi e salvati) saranno impostati a quelli
dell'utente.

A questo punto \cmd{login} provvederà (fatte salve eventuali altre azioni
iniziali, come la stampa di messaggi di benvenuto o il controllo della posta)
ad eseguire con un'altra \func{exec} la shell, che si troverà con un ambiente
già pronto con i file standard di tab.~\ref{tab:file_std_files} impostati sul
terminale, e pronta, nel ruolo di leader di sessione e di processo di
controllo per il terminale, a gestire l'esecuzione dei comandi come illustrato
in sez.~\ref{sec:sess_job_control_overview}.

Dato che il processo padre resta sempre \cmd{init} quest'ultimo potrà
provvedere, ricevendo un \signal{SIGCHLD} all'uscita della shell quando la
sessione di lavoro è terminata, a rilanciare \cmd{getty} sul terminale per
ripetere da capo tutto il procedimento.



\subsection{Interazione senza terminale: i \textsl{demoni} ed il
  \textit{syslog}}
\label{sec:sess_daemon}

Come sottolineato fin da sez.~\ref{sec:intro_base_concept}, in un sistema
unix-like tutte le operazioni sono eseguite tramite processi, comprese quelle
operazioni di sistema (come l'esecuzione dei comandi periodici, o la consegna
della posta, ed in generale tutti i programmi di servizio) che non hanno
niente a che fare con la gestione diretta dei comandi dell'utente.

Questi programmi, che devono essere eseguiti in modalità non interattiva e
senza nessun intervento dell'utente, sono normalmente chiamati
\textsl{demoni}, (o \textit{daemons}), nome ispirato dagli omonimi spiritelli
della mitologia greca che svolgevano compiti che gli dei trovavano noiosi, di
cui parla anche Socrate (che sosteneva di averne uno al suo servizio).

%TODO ricontrollare, i miei ricordi di filosofia sono piuttosto datati.

Se però si lancia un programma demone dalla riga di comando in un sistema che
supporta, come Linux, il \textit{job control} esso verrà comunque associato ad
un terminale di controllo e mantenuto all'interno di una sessione, e anche se
può essere mandato in background e non eseguire più nessun I/O su terminale,
si avranno comunque tutte le conseguenze che abbiamo trattato in
sez.~\ref{sec:sess_ctrl_term}, in particolare l'invio dei segnali in
corrispondenza dell'uscita del leader di sessione.

Per questo motivo un programma che deve funzionare come demone deve sempre
prendere autonomamente i provvedimenti opportuni (come distaccarsi dal
terminale e dalla sessione) ad impedire eventuali interferenze da parte del
sistema del \textit{job control}; questi sono riassunti in una lista di
prescrizioni\footnote{ad esempio sia Stevens in \cite{APUE}, che la
  \textit{Unix Programming FAQ} \cite{UnixFAQ} ne riportano di sostanzialmente
  identiche.} da seguire quando si scrive un demone.

Pertanto, quando si lancia un programma che deve essere eseguito come demone
occorrerà predisporlo in modo che esso compia le seguenti azioni:
\begin{enumerate*}
\item Eseguire una \func{fork} e terminare immediatamente il processo padre
  proseguendo l'esecuzione nel figlio.  In questo modo si ha la certezza che
  il figlio non è un \textit{process group leader}, (avrà il \ids{PGID} del
  padre, ma un \ids{PID} diverso) e si può chiamare \func{setsid} con
  successo. Inoltre la shell considererà terminato il comando all'uscita del
  padre.
\item Eseguire \func{setsid} per creare una nuova sessione ed un nuovo
  raggruppamento di cui il processo diventa automaticamente il leader, che
  però non ha associato nessun terminale di controllo.
\item Assicurarsi che al processo non venga associato in seguito nessun nuovo
  terminale di controllo; questo può essere fatto sia avendo cura di usare
  sempre l'opzione \const{O\_NOCTTY} nell'aprire i file di terminale, che
  eseguendo una ulteriore \func{fork} uscendo nel padre e proseguendo nel
  figlio. In questo caso, non essendo più quest'ultimo un leader di sessione
  non potrà ottenere automaticamente un terminale di controllo.
\item Eseguire una \func{chdir} per impostare la directory di lavoro del
  processo (su \file{/} o su una directory che contenga dei file necessari per
  il programma), per evitare che la directory da cui si è lanciato il processo
  resti in uso e non sia possibile rimuoverla o smontare il filesystem che la
  contiene.
\item Impostare la maschera dei permessi (di solito con \code{umask(0)}) in
  modo da non essere dipendenti dal valore ereditato da chi ha lanciato
  originariamente il processo.
\item Chiudere tutti i file aperti che non servono più (in generale tutti); in
  particolare vanno chiusi i file standard che di norma sono ancora associati
  al terminale (un'altra opzione è quella di redirigerli verso
  \file{/dev/null}).
\end{enumerate*}


In Linux buona parte di queste azioni possono venire eseguite invocando la
funzione \funcd{daemon} (fornita dalla \acr{glibc}), introdotta per la prima
volta in BSD4.4; il suo prototipo è:

\begin{funcproto}{
\fhead{unistd.h}
\fdecl{int daemon(int nochdir, int noclose)}
\fdesc{Rende il processo un demone.} 
}

{La funzione ritorna (nel nuovo processo) $0$ in caso di successo e $-1$ per
  un errore, nel qual caso \var{errno} assumerà uno dei valori i
    valori impostati dalle sottostanti \func{fork} e \func{setsid}.}
\end{funcproto}

La funzione esegue una \func{fork}, per uscire subito, con \func{\_exit}, nel
padre, mentre l'esecuzione prosegue nel figlio che esegue subito una
\func{setsid}. In questo modo si compiono automaticamente i passi 1 e 2 della
precedente lista. Se \param{nochdir} è nullo la funzione imposta anche la
directory di lavoro su \file{/}, se \param{noclose} è nullo i file standard
vengono rediretti su \file{/dev/null} (corrispondenti ai passi 4 e 6); in caso
di valori non nulli non viene eseguita nessuna altra azione.

Dato che un programma demone non può più accedere al terminale, si pone il
problema di come fare per la notifica di eventuali errori, non potendosi più
utilizzare lo \textit{standard error}. Per il normale I/O infatti ciascun
demone avrà le sue modalità di interazione col sistema e gli utenti a seconda
dei compiti e delle funzionalità che sono previste; ma gli errori devono
normalmente essere notificati all'amministratore del sistema.

\itindbeg{syslog}

Una soluzione può essere quella di scrivere gli eventuali messaggi su uno
specifico file (cosa che a volte viene fatta comunque) ma questo comporta il
grande svantaggio che l'amministratore dovrà tenere sotto controllo un file
diverso per ciascun demone, e che possono anche generarsi conflitti di nomi.
Per questo in BSD4.2 venne introdotto un servizio di sistema, il
\textit{syslog}, che oggi si trova su tutti i sistemi Unix, e che permette ai
demoni di inviare messaggi all'amministratore in una maniera standardizzata.

Il servizio prevede vari meccanismi di notifica, e, come ogni altro servizio
in un sistema unix-like, viene gestito attraverso un apposito programma, che è
anch'esso un \textsl{demone}. In generale i messaggi di errore vengono
raccolti dal file speciale \file{/dev/log}, un socket locale (vedi
sez.~\ref{sec:sock_sa_local}) dedicato a questo scopo, o via rete, con un
socket UDP e trattati dal demone che gestisce il servizio. Il più comune di
questi è \texttt{syslogd}, che consente un semplice smistamento dei messaggi
sui file in base alle informazioni in esse presenti; oggi però
\texttt{syslogd} è in sostanziale disuso, sostituito da programmi più
sofisticati come \texttt{rsyslog} o \texttt{syslog-ng}.

Il servizio del \textit{syslog} permette infatti di trattare i vari messaggi
classificandoli attraverso due indici: il primo, chiamato \textit{facility},
suddivide in diverse categorie i messaggi in modo di raggruppare quelli
provenienti da operazioni che hanno attinenza fra loro, ed è organizzato in
sottosistemi (kernel, posta elettronica, demoni di stampa, ecc.). Il secondo,
chiamato \textit{priority}, identifica l'importanza dei vari messaggi, e
permette di classificarli e differenziare le modalità di notifica degli
stessi.

Il sistema del \textit{syslog} attraverso il proprio demone di gestione
provvede poi a riportare i messaggi all'amministratore attraverso una serie
differenti meccanismi come:
\begin{itemize*}
\item scriverli su un file (comunemente detto \textit{log file}, o giornale),
\item scriverli sulla console,
\item scriverli sui terminali degli utenti connessi,
\item inviarli via mail ad uno specifico utente,
\item inviarli ad un altro programma,
\item inviarli via rete ad una macchina di raccolta,
\item ignorarli completamente;
\end{itemize*}
le modalità con cui queste azioni vengono realizzate dipendono ovviamente dal
demone che si usa, per la gestione del quale si rimanda ad un testo di
amministrazione di sistema.\footnote{l'argomento è ad esempio coperto dal
  capitolo 3.2.3 si \cite{AGL}.}

La \acr{glibc} definisce una serie di funzioni standard con cui un processo
può accedere in maniera generica al servizio di \textit{syslog}, che però
funzionano solo localmente; se si vogliono inviare i messaggi ad un altro
sistema occorre farlo esplicitamente con un socket UDP, o utilizzare le
capacità di reinvio del servizio.

La prima funzione definita dall'interfaccia è \funcd{openlog}, che inizializza
una connessione al servizio di \textit{syslog}. Essa in generale non è
necessaria per l'uso del servizio, ma permette di impostare alcuni valori che
controllano gli effetti delle chiamate successive; il suo prototipo è:

\begin{funcproto}{
\fhead{syslog.h}
\fdecl{void openlog(const char *ident, int option, int facility)}
\fdesc{Inizializza una connessione al sistema del \textit{syslog}.} 
}

{La funzione non restituisce nulla.}
\end{funcproto}

La funzione permette di specificare, tramite \param{ident}, l'identità di chi
ha inviato il messaggio (di norma si passa il nome del programma, come
specificato da \code{argv[0]}), e la stringa verrà preposta all'inizio di ogni
messaggio. Si tenga presente che il valore di \param{ident} che si passa alla
funzione è un puntatore, se la stringa cui punta viene cambiata lo sarà pure
nei successivi messaggi, e se viene cancellata i risultati potranno essere
impredicibili, per questo è sempre opportuno usare una stringa costante.

L'argomento \param{facility} permette invece di preimpostare per le successive
chiamate l'omonimo indice che classifica la categoria del messaggio.
L'argomento è interpretato come una maschera binaria, e pertanto è possibile
inviare i messaggi su più categorie alla volta. I valori delle costanti che
identificano ciascuna categoria sono riportati in
tab.~\ref{tab:sess_syslog_facility}, il valore di \param{facility} deve essere
specificato con un OR aritmetico.

\begin{table}[htb]
  \footnotesize
  \centering
  \begin{tabular}[c]{|l|p{8cm}|}
    \hline
    \textbf{Valore}& \textbf{Significato}\\
    \hline
    \hline
    \constd{LOG\_AUTH}     & Messaggi relativi ad autenticazione e sicurezza,
                             obsoleto, è sostituito da \const{LOG\_AUTHPRIV}.\\
    \constd{LOG\_AUTHPRIV} & Sostituisce \const{LOG\_AUTH}.\\
    \constd{LOG\_CRON}     & Messaggi dei demoni di gestione dei comandi
                             programmati (\cmd{cron} e \cmd{at}).\\
    \constd{LOG\_DAEMON}   & Demoni di sistema.\\
    \constd{LOG\_FTP}      & Servizio FTP.\\
    \constd{LOG\_KERN}     & Messaggi del kernel.\\
    \constd{LOG\_LOCAL0}   & Riservato all'amministratore per uso locale.\\
    \hspace{.5cm}$\vdots$ &   \hspace{3cm}$\vdots$\\
    \constd{LOG\_LOCAL7}   & Riservato all'amministratore per uso locale.\\
    \constd{LOG\_LPR}      & Messaggi del sistema di gestione delle stampanti.\\
    \constd{LOG\_MAIL}     & Messaggi del sistema di posta elettronica.\\
    \constd{LOG\_NEWS}     & Messaggi del sistema di gestione delle news 
                            (USENET).\\
    \constd{LOG\_SYSLOG}   & Messaggi generati dal demone di gestione del
                            \textit{syslog}.\\
    \constd{LOG\_USER}     & Messaggi generici a livello utente.\\
    \constd{LOG\_UUCP}     & Messaggi del sistema UUCP (\textit{Unix to Unix
                             CoPy}), ormai in disuso.\\
\hline
\end{tabular}
\caption{Valori possibili per l'argomento \param{facility} di \func{openlog}.}
\label{tab:sess_syslog_facility}
\end{table}

L'argomento \param{option} serve invece per controllare il comportamento della
funzione \func{openlog} e delle modalità con cui le successive chiamate
scriveranno i messaggi, esso viene specificato come maschera binaria composta
con un OR aritmetico di una qualunque delle costanti riportate in
tab.~\ref{tab:sess_openlog_option}.

\begin{table}[htb]
  \footnotesize
\centering
\begin{tabular}[c]{|l|p{8cm}|}
\hline
\textbf{Valore}& \textbf{Significato}\\
\hline
\hline
  \constd{LOG\_CONS}   & Scrive sulla console in caso di errore nell'invio del
                         messaggio al sistema del \textit{syslog}. \\
  \constd{LOG\_NDELAY} & Apre la connessione al sistema del \textit{syslog}
                         subito invece di attendere l'invio del primo
                         messaggio.\\ 
  \constd{LOG\_NOWAIT} & Non usato su Linux, su altre piattaforme non attende i
                         processi figli creati per inviare il messaggio.\\
  \constd{LOG\_ODELAY} & Attende il primo messaggio per aprire la connessione al
                         sistema del \textit{syslog}.\\ 
  \constd{LOG\_PERROR} & Stampa anche su \file{stderr} (non previsto in
                         POSIX.1-2001).\\ 
  \constd{LOG\_PID}    & Inserisce nei messaggi il \ids{PID} del processo
                         chiamante.\\
\hline
\end{tabular}
\caption{Valori possibili per l'argomento \param{option} di \func{openlog}.}
\label{tab:sess_openlog_option}
\end{table}

La funzione che si usa per generare un messaggio è \funcd{syslog}, dato che
l'uso di \func{openlog} è opzionale, sarà quest'ultima a provvede a chiamare la
prima qualora ciò non sia stato fatto (nel qual caso il valore di
\param{ident} è \val{NULL}). Il suo prototipo è:

\begin{funcproto}{
\fhead{syslog.h}
\fdecl{void syslog(int priority, const char *format, ...)}
\fdesc{Genera un messaggio per il \textit{syslog}.} 
}

{La funzione non restituisce nulla.}
\end{funcproto}

La funzione genera un messaggio le cui caratteristiche sono indicate
da \param{priority}. Per i restanti argomenti il suo comportamento è analogo a
quello di \func{printf}, e il valore dell'argomento \param{format} è identico
a quello descritto nella pagina di manuale di quest'ultima (per i valori
principali si può vedere la trattazione sommaria che se ne è fatto in
sez.~\ref{sec:file_formatted_io}).  L'unica differenza è che la sequenza
\val{\%m} viene rimpiazzata dalla stringa restituita da
\code{strerror(errno)}. Gli argomenti seguenti i primi due devono essere
forniti secondo quanto richiesto da \param{format}.

\begin{table}[htb]
  \footnotesize
  \centering
  \begin{tabular}[c]{|l|p{8cm}|}
    \hline
    \textbf{Valore}& \textbf{Significato}\\
    \hline
    \hline
    \constd{LOG\_EMERG}   & Il sistema è inutilizzabile.\\
    \constd{LOG\_ALERT}   & C'è una emergenza che richiede intervento
                            immediato.\\
    \constd{LOG\_CRIT}    & Si è in una condizione critica.\\
    \constd{LOG\_ERR}     & Si è in una condizione di errore.\\
    \constd{LOG\_WARNING} & Messaggio di avvertimento.\\
    \constd{LOG\_NOTICE}  & Notizia significativa relativa al comportamento.\\
    \constd{LOG\_INFO}    & Messaggio informativo.\\
    \constd{LOG\_DEBUG}   & Messaggio di debug.\\
    \hline
  \end{tabular}
  \caption{Valori possibili per l'indice di importanza del messaggio da
    specificare nell'argomento \param{priority} di \func{syslog}.}
  \label{tab:sess_syslog_priority}
\end{table}

L'argomento \param{priority} permette di impostare sia la \textit{facility}
che la \textit{priority} del messaggio. In realtà viene prevalentemente usato
per specificare solo quest'ultima in quanto la prima viene di norma
preimpostata con \func{openlog}. La priorità è indicata con un valore numerico
specificabile attraverso le costanti riportate in
tab.~\ref{tab:sess_syslog_priority}.  

La \acr{glibc}, seguendo POSIX.1-2001, prevede otto diverse priorità
ordinate da 0 a 7, in ordine di importanza decrescente; questo comporta che i
tre bit meno significativi dell'argomento \param{priority} sono occupati da
questo valore, mentre i restanti bit più significativi vengono usati per
specificare la \textit{facility}.  Nel caso si voglia specificare anche la
\textit{facility} basta eseguire un OR aritmetico del valore della priorità
con la maschera binaria delle costanti di tab.~\ref{tab:sess_syslog_facility}.

Una funzione sostanzialmente identica a \func{syslog} è \funcd{vsyslog}. La
funzione è originaria di BSD e per utilizzarla deve essere definito
\macro{\_BSD\_SOURCE}; il suo prototipo è:

\begin{funcproto}{
\fhead{syslog.h}
\fdecl{void vsyslog(int priority, const char *format, va\_list src)}
\fdesc{Genera un messaggio per il \textit{syslog}.} 
}

{La funzione non restituisce nulla.}
\end{funcproto}

La sola differenza con \func{syslog} è quella di prendere invece di una lista
di argomenti esplicita un unico argomento finale passato nella forma di una
\macro{va\_list}; la funzione risulta utile qualora si ottengano gli argomenti
dalla invocazione di un'altra funzione \textit{variadic} (si ricordi quanto
visto in sez.~\ref{sec:proc_variadic}).

Per semplificare la gestione della scelta del livello di priorità a partire
dal quale si vogliono registrare i messaggi, le funzioni di gestione
mantengono per ogni processo una maschera che determina quale delle chiamate
effettuate a \func{syslog} verrà effettivamente registrata.  In questo modo
sarà possibile escludere i livelli di priorità che non interessa registrare,
impostando opportunamente la maschera una volta per tutte.

Questo significa che in genere nei programmi vengono comunque previste le
chiamate a \func{syslog} per tutti i livelli di priorità, ma poi si imposta
questa maschera per registrare solo quello che effettivamente interessa. La
funzione che consente di fare questo è \funcd{setlogmask}, ed il suo prototipo
è:

\begin{funcproto}{
\fhead{syslog.h}
\fdecl{int setlogmask(int mask)}
\fdesc{Imposta la maschera dei messaggi del \textit{syslog}.} 
}

{La funzione ritorna il precedente valore della maschera dei messaggi e non
  prevede errori.}
\end{funcproto}

La funzione restituisce il valore della maschera corrente, e se si passa un
valore nullo per \param{mask} la maschera corrente non viene modificata; in
questo modo si può leggere il valore della maschera corrente. Indicando un
valore non nullo per \param{mask} la registrazione dei messaggi viene
disabilitata per tutte quelle priorità che non rientrano nella maschera.

In genere il valore viene impostato usando la macro
\macrod{LOG\_MASK}\texttt{(p)} dove \code{p} è una delle costanti di
tab.~\ref{tab:sess_syslog_priority}. É inoltre disponibile anche la macro
\macrod{LOG\_UPTO}\texttt{(p)} che permette di specificare automaticamente
tutte le priorità fino a quella indicata da \code{p}.

Una volta che si sia certi che non si intende registrare più nessun messaggio
si può chiudere esplicitamente la connessione al \textit{syslog} (l'uso di
questa funzione è comunque completamente opzionale) con la funzione
\funcd{closelog}, il cui prototipo è:

\begin{funcproto}{
\fhead{syslog.h}
\fdecl{void closelog(void)}
\fdesc{Chiude la connessione al \textit{syslog}.} 
}

{La funzione non ritorna nulla.}
\end{funcproto}


Come si evince anche dalla presenza della facility \const{LOG\_KERN} in
tab.~\ref{tab:sess_syslog_facility}, uno dei possibili utenti del servizio del
\textit{syslog} è anche il kernel, che a sua volta può avere necessità di
inviare messaggi verso l'\textit{user space}. I messaggi del kernel sono
mantenuti in un apposito buffer circolare e generati all'interno del kernel
tramite la funzione \texttt{printk}, analoga alla \func{printf} usata in
\textit{user space}.\footnote{una trattazione eccellente dell'argomento si
  trova nel quarto capitolo di \cite{LinDevDri}.}

Come per i messaggi ordinari anche i messaggi del kernel hanno una priorità ma
in questo caso non si può contare sulla coincidenza con le costanti di
tab.~\ref{tab:sess_syslog_priority} dato che il codice del kernel viene
mantenuto in maniera indipendente dalle librerie del C. Per questo motivo le
varie priorità usate dal kernel sono associate ad un valore numerico che viene
tradotto in una stringa preposta ad ogni messaggio, secondo i valori che si
sono riportati in fig.~\ref{fig:printk_priority}

\begin{figure}[!htb]
  \footnotesize \centering
  \begin{minipage}[c]{0.80\textwidth}
    \includestruct{listati/printk_prio.c}
  \end{minipage} 
  \normalsize 
  \caption{Definizione delle stringhe coi relativi valori numerici che
    indicano le priorità dei messaggi del kernel (ripresa da
    \file{include/linux/kernel.h}).}
  \label{fig:printk_priority}
\end{figure}

Dato che i messaggi generati da \texttt{printk} hanno un loro specifico
formato tradizionalmente si usava un demone ausiliario, \cmd{klogd}, per
leggerli, rimappare le loro priorità sui valori di
tab.~\ref{tab:sess_syslog_priority} ed inviarli al sistema del \textit{syslog}
nella facility \const{LOG\_KERN}.  Oggi i nuovi demoni più avanzati che
realizzano il servizio (come \texttt{rsyslog} o \texttt{syslog-ng}) sono in
grado di fare tutto questo da soli leggendoli direttamente senza necessità di
un intermediario.

Ma i messaggi del kernel non sono necessariamente connessi al sistema del
\textit{syslog}; ad esempio possono anche essere letti direttamente dal buffer
circolare con il comando \texttt{dmesg}. Inoltre è previsto che se superano
una certa priorità essi vengano stampati direttamente sul terminale indicato
come \textit{console} di sistema,\footnote{quello che viene indicato con il
  parametro di avvio \texttt{console} del kernel, si consulti al riguardo
  sez.~5.3.1 di \cite{AGL}.}  in modo che sia possibile vederli anche in caso
di blocco totale del sistema (nell'assunzione che la console sia collegata).

In particolare la stampa dei messaggi sulla console è controllata dal
contenuto del file \sysctlfiled{kernel/printk} (o con l'equivalente parametro
di \func{sysctl}) che prevede quattro valori numerici interi: il primo,
\textit{console\_loglevel}, indica la priorità corrente oltre la quale vengono
stampati i messaggi sulla console, il secondo,
\textit{default\_message\_loglevel}, la priorità di default assegnata ai
messaggi che non ne hanno impostata una, il terzo,
\textit{minimum\_console\_level}, il valore minimo che si può assegnare al
primo valore,\footnote{che può essere usato con una delle operazioni di
  gestione che vedremo a breve per ``\textsl{silenziare}'' il kernel.} ed il
quarto, \textit{default\_console\_loglevel}, il valore di
default.\footnote{anch'esso viene usato nelle operazioni di controllo per
  tornare ad un valore predefinito.}

Per la lettura dei messaggi del kernel e la gestione del relativo buffer
circolare esiste una apposita \textit{system call} chiamata anch'essa
 \texttt{syslog}, ma dato il conflitto di nomi questa viene rimappata su
 un'altra funzione di libreria, in particolare nella \acr{glibc} essa viene
 invocata tramite la funzione \funcd{klogctl},\footnote{nella \acr{libc4} e
   nella \acr{libc5} la funzione invece era \code{SYS\_klog}.} il cui prototipo
 è:

\begin{funcproto}{
\fhead{sys/klog.h}
\fdecl{int klogctl(int op, char *buffer, int len)}
\fdesc{Gestisce i messaggi di log del kernel.} 
}

{La funzione ritorna un intero positivo o nullo dipendente dall'operazione
  scelta in caso di successo e $-1$ per un errore, nel qual caso \var{errno}
  assumerà uno dei valori:
  \begin{errlist}
  \item[\errcode{EINVAL}] l'argomento \param{op} non ha un valore valido, o si
    sono specificati valori non validi per gli altri argomenti quando questi
    sono richiesti.
  \item[\errcode{ENOSYS}] il supporto per \texttt{printk} non è stato compilato
    nel kernel.
  \item[\errcode{EPERM}] non si hanno i privilegi richiesti per l'operazione
    richiesta.
  \item[\errcode{ERESTARTSYS}] l'operazione è stata interrotta da un segnale.
   \end{errlist}
}
\end{funcproto}

La funzione prevede che si passi come primo argomento \param{op} un codice
numerico che indica l'operazione richiesta, il secondo argomento deve essere,
per le operazioni che compiono una lettura di dati, l'indirizzo del buffer su
cui copiarli, ed il terzo quanti byte leggere. L'effettivo uso di questi due
argomenti dipende comunque dall'operazione richiesta, ma essi devono essere
comunque specificati, anche quando non servono, nel qual caso verranno
semplicemente ignorati.

\begin{table}[htb]
  \footnotesize
  \centering
  \begin{tabular}[c]{|l|p{10cm}|}
    \hline
    \textbf{Valore}& \textbf{Significato}\\
    \hline
    \hline
    \texttt{0} & apre il log (attualmente non fa niente), \param{buffer}
                 e \param{len} sono ignorati.\\
    \texttt{1} & chiude il log (attualmente non fa niente), \param{buffer}
                 e \param{len} sono ignorati.\\
    \texttt{2} & legge \param{len} byte nel buffer \param{buffer} dal log dei
                 messaggi.\\   
    \texttt{3} & legge \param{len} byte nel buffer \param{buffer} dal buffer
                 circolare dei messaggi.\\
    \texttt{4} & legge \param{len} byte nel buffer \param{buffer} dal buffer
                 circolare dei messaggi e lo svuota.\\
    \texttt{5} & svuota il buffer circolare dei messaggi, \param{buffer}
                 e \param{len} sono ignorati.\\
    \texttt{6} & disabilita la stampa dei messaggi sulla console, \param{buffer}
                 e \param{len} sono ignorati.\\
    \texttt{7} & abilita la stampa dei messaggi sulla console, \param{buffer}
                 e \param{len} sono ignorati.\\
    \texttt{8} & imposta a \param{len} il livello dei messaggi stampati sulla
                 console, \param{buffer} è ignorato.\\ 
    \texttt{9} & ritorna il numero di byte da leggere presenti sul buffer di
                 log, \param{buffer} e \param{len} sono ignorati (dal kernel
                 2.4.10).\\ 
    \texttt{10}& ritorna la dimensione del buffer di log, \param{buffer}
                 e \param{len} sono ignorati (dal kernel 2.6.6).\\
\hline
\end{tabular}
\caption{Valori possibili per l'argomento \param{op} di \func{klogctl}.}
\label{tab:klogctl_operation}
\end{table}

Si sono riportati in tab.~\ref{tab:klogctl_operation} i valori utilizzabili
per \param{op}, con una breve spiegazione della relativa operazione e di come
vengono usati gli altri due argomenti. Come si può notare la funzione è una
sorta di interfaccia comune usata per eseguire operazioni completamente
diverse fra loro.

L'operazione corrispondente al valore 2 di \param{op} consente di leggere un
messaggio dal cosiddetto \textit{log} del kernel. Eseguire questa operazione è
equivalente ad eseguire una lettura dal file
\procfile{/proc/kmsg},\footnote{in realtà è vero l'opposto, è questa funzione
  che viene eseguita quando si legge da questo file.} se non vi sono messaggi
la funzione si blocca in attesa di dati e ritorna soltanto quando questi
diventano disponibili. In tal caso verranno letti ed
estratti\footnote{estratti in quanti i dati del \textit{log} del kernel si
  possono leggere una volta sola, se più processi eseguono l'operazione di
  lettura soltanto uno riceverà i dati, a meno che completata la propria
  operazione di lettura non restino altri messaggi pendenti che a questo punto
  potrebbero essere letti da un altro processo in attesa.} dal log \param{len}
byte che verranno scritti su \param{buffer}; in questo caso il valore di
ritorno di \func{klogctl} corrisponderà al numero di byte ottenuti.

Se invece si usa l'operazione 3 i dati vengono letti dal buffer circolare
usato da \texttt{printk}, che mantiene tutti i messaggi stampati dal kernel
fino al limite delle sue dimensioni, in questo caso i messaggi possono essere
letti più volte. Usando invece l'operazione 4 si richiede di cancellare il
buffer dopo la lettura, che così risulterà vuoto ad una lettura
successiva. Anche con queste operazioni \param{len} indica il numero di byte
da leggere e \param{buffer} l'indirizzo dove leggerli, e la funzione ritorna
il numero di byte effettivamente letti. L'operazione 5 esegue soltanto la
cancellazione del buffer circolare, \param{len} e \param{buffer} sono ignorati
e la funzione ritorna un valore nullo.

Le operazioni corrispondenti ai valori 6, 7 ed 8 consentono di modificare la
priorità oltre la quale i messaggi vengono stampati direttamente sulla
\textit{console} e fanno riferimento ai parametri del kernel gestiti con le
variabili contenute in \sysctlfile{kernel/printk} di cui abbiamo
parlato prima, ed in particolare con 6 si imposta come corrente il valore
minimo della terza variabile (\textit{minimum\_console\_level}), ottenendo
l'effetto di ridurre al minimo i messaggi che arrivano in console, mentre con
7 si ripristina il valore di default.\footnote{secondo la documentazione
  questo sarebbe quello indicato della quarta variabile,
  \textit{default\_console\_loglevel} in genere pari a 7, ma alcune prove con
  il programma \texttt{mydmesg} che si trova nei sorgenti allegati alla guida
  rivelano che l'unico effetto di questa operazione è riportare il valore a
  quello precedente se lo si è ridotto al minimo con l'operazione 6.}  Per
impostare direttamente un valore specifico infine si può usare 8, nel qual
caso il valore numerico del livello da impostare deve essere specificato
con \param{len}, che può assumere solo un valore fra 1 e 8.

Infine le due operazioni 9 e 10 consentono di ottenere rispettivamente il
numero di byte ancora non letti dal log del kernel, e la dimensione totale di
questo. Per entrambe i dati sono restituiti come valore di ritorno, e gli
argomento \param{buffer} e \param{len} sono ignorati.

Si tenga presente che la modifica del livello minimo per cui i messaggi
vengono stampati sulla console (operazioni 6, 7 e 8) e la cancellazione del
buffer circolare di \texttt{printk} (operazioni 4 e 5) sono privilegiate; fino
al kernel 2.6.30 era richiesta la capacità \const{CAP\_SYS\_ADMIN}, a partire
dal 2.6.38 detto privilegio è stato assegnato ad una capacità aggiuntiva,
\const{CAP\_SYSLOG}. Tutto questo è stato fatto per evitare che processi
eseguiti all'interno di un sistema di virtualizzazione ``\textsl{leggera}''
(come i \textit{Linux Container} di LXC) che necessitano di
\const{CAP\_SYS\_ADMIN} per operare all'interno del proprio ambiente
ristretto, potessero anche avere la capacità di influire sui log del kernel
al di fuori di questo.

\itindend{syslog}



\section{L'I/O su terminale}
\label{sec:sess_terminal_io}

Benché come ogni altro dispositivo i terminali siano accessibili come file,
essi hanno assunto storicamente, essendo stati a lungo l'unico modo di
accedere al sistema, una loro rilevanza specifica, che abbiamo già avuto modo
di incontrare nella precedente sezione.

Esamineremo qui le peculiarità dell'I/O eseguito sui terminali, che per la
loro particolare natura presenta delle differenze rispetto ai normali file su
disco e agli altri dispositivi.



\subsection{L'architettura dell'I/O su terminale}
\label{sec:term_io_design}

I terminali sono una classe speciale di dispositivi a caratteri (si ricordi la
classificazione di sez.~\ref{sec:file_file_types}). Un terminale ha infatti
una caratteristica che lo contraddistingue da un qualunque altro dispositivo,
è infatti destinato a gestire l'interazione con un utente e deve perciò essere
in grado di fare da terminale di controllo per una sessione; tutto questo
comporta la presenza di una serie di capacità specifiche.

L'interfaccia per i terminali è una delle più oscure e complesse, essendosi
stratificata dagli inizi dei sistemi Unix fino ad oggi. Questo comporta una
grande quantità di opzioni e controlli relativi ad un insieme di
caratteristiche (come ad esempio la velocità della linea) necessarie per
dispositivi, come i terminali seriali, che al giorno d'oggi sono praticamente
in disuso.

Storicamente i primi terminali erano appunto terminali di telescriventi
(\textit{teletype}), da cui deriva sia il nome dell'interfaccia, \textit{TTY},
che quello dei relativi file di dispositivo, che sono sempre della forma
\texttt{/dev/tty*}.\footnote{ciò vale solo in parte per i terminali virtuali,
  essi infatti hanno due lati, un \textit{master}, che può assumere i nomi
  \file{/dev/pty[p-za-e][0-9a-f]} ed un corrispondente \textit{slave} con nome
  \file{/dev/tty[p-za-e][0-9a-f]}.}  Oggi essi includono le porte seriali, le
console virtuali dello schermo, i terminali virtuali che vengono creati come
canali di comunicazione dal kernel e che di solito vengono associati alle
connessioni di rete (ad esempio per trattare i dati inviati con \cmd{telnet} o
\cmd{ssh}).

L'I/O sui terminali si effettua con le stesse modalità dei file normali: si
apre il relativo file di dispositivo, e si leggono e scrivono i dati con le
usuali funzioni di lettura e scrittura. Così se apriamo una console virtuale
avremo che \func{read} leggerà quanto immesso dalla tastiera, mentre
\func{write} scriverà sullo schermo.  In realtà questo è vero solo a grandi
linee, perché non tiene conto delle caratteristiche specifiche dei terminali;
una delle principali infatti è che essi prevedono due modalità di operazione,
dette rispettivamente ``\textsl{modo canonico}'' e ``\textsl{modo non
  canonico}'', che hanno dei comportamenti nettamente
diversi. \index{modo~canonico}\index{modo~non~canonico}

% TODO: inserire qui il comportamento di read relativo all'errore EIO sulla
% lettura in background???

La modalità preimpostata all'apertura del terminale è quella canonica, in cui
le operazioni di lettura vengono sempre effettuate assemblando i dati in una
linea. Questo significa che eseguendo una \func{read} su un terminale in modo
canonico la funzione si bloccherà, anche se si sono scritti dei caratteri,
fintanto che non si preme il tasto di ritorno a capo: a questo punto la linea
sarà completata e la funzione ritornerà leggendola per intero.

Inoltre in modalità canonica alcuni dei caratteri che si scrivono sul
terminale vengono interpretati direttamente dal kernel per compiere operazioni
(come la generazione dei segnali associati al \textit{job control} illustrata
in sez.~\ref{sec:sig_job_control}), e non saranno mai letti dal
dispositivo. Quella canonica è di norma la modalità in cui opera la shell.

Un terminale in modo non canonico invece non effettua nessun accorpamento dei
dati in linee né li interpreta; esso viene di solito usato da quei programmi,
come ad esempio gli editor, che necessitano di poter leggere un carattere alla
volta e che gestiscono al loro interno l'interpretazione dei caratteri
ricevuti impiegandoli opportunamente come comandi o come dati.

\begin{figure}[!htb]
  \centering \includegraphics[width=12cm]{img/term_struct}
  \caption{Struttura interna generica del kernel per l'accesso ai dati di un
    terminale.}
  \label{fig:term_struct}
\end{figure}

Per capire le caratteristiche dell'I/O sui terminali occorre esaminare le
modalità con cui esso viene effettuato. L'accesso, come per tutti i
dispositivi, viene gestito dal kernel, ma per tutti i terminali viene
utilizzata una architettura generica che si è schematizzata in
fig.~\ref{fig:term_struct}.  

Ad ogni terminale sono sempre associate due code
per gestire l'input e l'output, che ne implementano una bufferizzazione
all'interno del kernel che è completamente indipendente dalla eventuale
ulteriore bufferizzazione fornita dall'interfaccia standard dei file.

La coda di ingresso mantiene i caratteri che sono stati letti dal terminale ma
non ancora letti da un processo, la sua dimensione è definita dal parametro di
sistema \const{MAX\_INPUT} (si veda sez.~\ref{sec:sys_file_limits}), che ne
specifica il limite minimo, in realtà la coda può essere più grande e cambiare
dimensione dinamicamente. 

Se è stato abilitato il controllo di flusso in ingresso il kernel emette i
caratteri di STOP e START per bloccare e sbloccare l'ingresso dei dati;
altrimenti i caratteri immessi oltre le dimensioni massime vengono persi; in
alcuni casi il kernel provvede ad inviare automaticamente un avviso (un
carattere di BELL, che provoca un beep) sull'output quando si eccedono le
dimensioni della coda.  

Se è abilitato la modalità canonica i caratteri in ingresso restano nella coda
fintanto che non viene ricevuto un a capo; un altro parametro del sistema,
\const{MAX\_CANON}, specifica la dimensione massima di una riga in modalità
canonica.

La coda di uscita è analoga a quella di ingresso e contiene i caratteri
scritti dai processi ma non ancora inviati al terminale. Se è abilitato il
controllo di flusso in uscita il kernel risponde ai caratteri di START e STOP
inviati dal terminale. Le dimensioni della coda non sono specificate, ma non
hanno molta importanza, in quanto qualora esse vengano eccedute il driver
provvede automaticamente a bloccare la funzione chiamante.



\subsection{La gestione delle caratteristiche di un terminale}
\label{sec:term_attr}

Data le loro peculiarità, fin dalla realizzazione dei primi sistemi unix-like
si è posto il problema di come gestire le caratteristiche specifiche dei
terminali. Storicamente i vari dialetti di Unix hanno utilizzato delle
funzioni specifiche diverse fra loro, ma con la realizzazione dello standard
POSIX.1-2001 è stata effettuata opportuna unificazione delle funzioni
attinenti i terminali, sintetizzando le differenze fra BSD e System V in una
unica interfaccia, che è quella adottata da Linux.

Molte delle funzioni previste dallo standard POSIX.1-2001 prevedono come
argomento un file descriptor, dato che in origine le relative operazioni
venivano effettuate con delle opportune chiamate a \func{ioctl}.  Ovviamente
dette funzioni potranno essere usate correttamente soltanto con dei file
descriptor che corrispondono ad un terminale, in caso contrario lo standard
richiede che venga restituito un errore di \errcode{ENOTTY}.

Per evitare l'errore, ed anche semplicemente per verificare se un file
descriptor corrisponde ad un terminale, è disponibile la funzione
\funcd{isatty}, il cui prototipo è:

\begin{funcproto}{
\fhead{unistd.h}
\fdecl{int isatty(int fd)}
\fdesc{Controlla se un file è un terminale.} 
}

{La funzione ritorna $1$ se \param{fd} è connesso ad un terminale e $0$
  altrimenti, nel qual caso \var{errno} assumerà uno dei valori:
  \begin{errlist}
  \item[\errcode{EBADF}] \param{fd} non è un file descriptor valido.
  \item[\errcode{EINVAL}] \param{fd} non è associato a un terminale (non
    ottempera a POSIX.1-2001 che richiederebbe \errcode{ENOTTY}).
  \end{errlist}
}
\end{funcproto}

Un'altra funzione per avere informazioni su un terminale è \funcd{ttyname},
che permette di ottenere il nome del file di dispositivo del terminale
associato ad un file descriptor; il suo prototipo è:

\begin{funcproto}{
\fhead{unistd.h}
\fdecl{char *ttyname(int fd)}
\fdesc{Restituisce il nome del terminale associato ad un file descriptor.} 
}

{La funzione ritorna il puntatore alla stringa contenente il nome del
  terminale in caso di successo e \val{NULL} per un errore, nel qual caso
  \var{errno} assumerà uno dei valori:
  \begin{errlist}
  \item[\errcode{EBADF}] \param{fd} non è un file descriptor valido.
  \item[\errcode{ENOTTY}] \param{fd} non è associato a un terminale.
  \end{errlist}
}
\end{funcproto}

La funzione restituisce il puntatore alla stringa contenente il nome del file
di dispositivo del terminale associato a \param{fd}, che però è allocata
staticamente e può essere sovrascritta da successive chiamate. Per questo
della funzione esiste anche una versione rientrante, \funcd{ttyname\_r}, che
non presenta il problema dell'uso di una zona di memoria statica; il suo
prototipo è:

\begin{funcproto}{
\fhead{unistd.h}
\fdecl{int ttyname\_r(int fd, char *buff, size\_t len)}
\fdesc{Restituisce il nome del terminale associato ad un file descriptor.} 
}

{La funzione ritorna $0$ in caso di successo e $-1$ per un errore, nel qual
  caso \var{errno} assumerà uno dei valori: 
  \begin{errlist}
    \item[\errcode{ERANGE}] la lunghezza del buffer \param{len} non è
      sufficiente per contenere la stringa restituita.
  \end{errlist}
  ed inoltre \errval{EBADF} ed \errval{ENOTTY} con lo stesso significato di
  \func{ttyname}.}
\end{funcproto}

La funzione prende due argomenti in più, il puntatore \param{buff} alla zona
di memoria in cui l'utente vuole che il risultato venga scritto, che dovrà
essere allocata in precedenza, e la relativa dimensione, \param{len}. Se la
stringa che deve essere restituita, compreso lo zero di terminazione finale,
eccede questa dimensione si avrà un errore.

Una funzione analoga alle precedenti prevista da POSIX.1, che restituisce il
nome di un file di dispositivo, è \funcd{ctermid}, il cui prototipo è:

\begin{funcproto}{
\fhead{stdio.h}
\fdecl{char *ctermid(char *s)}
\fdesc{Restituisce il nome del terminale di controllo del processo.} 
}

{La funzione ritorna il puntatore alla stringa contenente il \textit{pathname}
  del terminale o \val{NULL} se non riesce ad eseguire l'operazione, non sono
  previsti errori.}
\end{funcproto}

La funzione restituisce un puntatore al \textit{pathname} del file di
dispositivo del terminale di controllo del processo chiamante.  Se si passa
come argomento \val{NULL} la funzione restituisce il puntatore ad una stringa
statica che può essere sovrascritta da chiamate successive, e non è
rientrante. Indicando invece un puntatore ad una zona di memoria già allocata
la stringa sarà scritta su di essa, ma in questo caso il buffer preallocato
deve avere una dimensione di almeno \constd{L\_ctermid}
caratteri.\footnote{\const{L\_ctermid} è una delle varie costanti del sistema,
  non trattata esplicitamente in sez.~\ref{sec:sys_characteristics}, che
  indica la dimensione che deve avere una stringa per poter contenere il nome
  di un terminale.}

Si tenga presente che il \textit{pathname} restituito dalla funzione potrebbe
non identificare univocamente il terminale (ad esempio potrebbe essere
\file{/dev/tty}), inoltre non è detto che il processo possa effettivamente
essere in grado di aprire il terminale.

I vari attributi associati ad un terminale vengono mantenuti per ciascuno di
essi in una struttura \struct{termios} che viene usata dalle varie funzioni
dell'interfaccia. In fig.~\ref{fig:term_termios} si sono riportati tutti i
campi della definizione di questa struttura usata in Linux; di questi solo i
primi cinque sono previsti dallo standard POSIX.1, ma le varie implementazioni
ne aggiungono degli altri per mantenere ulteriori informazioni.\footnote{la
  definizione della struttura si trova in \file{bits/termios.h}, da non
  includere mai direttamente; Linux, seguendo l'esempio di BSD, aggiunge i due
  campi \var{c\_ispeed} e \var{c\_ospeed} per mantenere le velocità delle
  linee seriali, ed un campo ulteriore, \var{c\_line} per indicare la
  disciplina di linea.}

\begin{figure}[!htb] 
  \footnotesize \centering
  \begin{minipage}[c]{0.80\textwidth}
    \includestruct{listati/termios.h}
  \end{minipage} 
  \normalsize 
  \caption{La struttura \structd{termios}, che identifica le proprietà di un
    terminale.}
  \label{fig:term_termios}
\end{figure}

% per le definizioni dei dati vedi anche:
% http://www.lafn.org/~dave/linux/termios.txt 

I primi quattro campi sono quattro flag che controllano il comportamento del
terminale; essi sono realizzati come maschera binaria, pertanto il tipo
\typed{tcflag\_t} è di norma realizzato con un intero senza segno di lunghezza
opportuna. I valori devono essere specificati bit per bit, avendo cura di non
modificare i bit su cui non si interviene.

\begin{table}[b!ht]
  \footnotesize
  \centering
  \begin{tabular}[c]{|l|p{10cm}|}
    \hline
    \textbf{Valore}& \textbf{Significato}\\
    \hline
    \hline
    \constd{IGNBRK} & Ignora le condizioni di BREAK sull'input. Una
                      \textit{condizione di BREAK} è definita nel contesto di
                      una trasmissione seriale asincrona come una sequenza di
                      bit nulli più lunga di un byte.\\
    \constd{BRKINT} & Controlla la reazione ad un BREAK quando
                      \const{IGNBRK} non è impostato. Se \const{BRKINT} è
                      impostato il BREAK causa lo scarico delle code, 
                      e se il terminale è il terminale di controllo per un 
                      gruppo in foreground anche l'invio di \signal{SIGINT} ai
                      processi di quest'ultimo. Se invece \const{BRKINT} non è
                      impostato un BREAK viene letto come un carattere
                      NUL, a meno che non sia impostato \const{PARMRK}
                      nel qual caso viene letto come la sequenza di caratteri
                      \texttt{0xFF 0x00 0x00}.\\
    \constd{IGNPAR} & Ignora gli errori di parità, il carattere viene passato
                      come ricevuto. Ha senso solo se si è impostato 
                      \const{INPCK}.\\
    \constd{PARMRK} & Controlla come vengono riportati gli errori di parità. Ha 
                      senso solo se \const{INPCK} è impostato e \const{IGNPAR}
                      no. Se impostato inserisce una sequenza \texttt{0xFF
                      0x00} prima di ogni carattere che presenta errori di
                      parità, se non impostato un carattere con errori di
                      parità viene letto come uno \texttt{0x00}. Se un
                      carattere ha il valore \texttt{0xFF} e \const{ISTRIP} 
                      non è impostato, per evitare ambiguità esso viene sempre
                      riportato come \texttt{0xFF 0xFF}.\\
    \constd{INPCK}  & Abilita il controllo di parità in ingresso. Se non viene
                      impostato non viene fatto nessun controllo ed i caratteri
                      vengono passati in input direttamente.\\
    \constd{ISTRIP} & Se impostato i caratteri in input sono tagliati a sette
                      bit mettendo a zero il bit più significativo, altrimenti 
                      vengono passati tutti gli otto bit.\\
    \constd{INLCR}  & Se impostato in ingresso il carattere di a capo
                      (\verb|'\n'|) viene automaticamente trasformato in un
                      ritorno carrello (\verb|'\r'|).\\
    \constd{IGNCR}  & Se impostato il carattere di ritorno carrello 
                      (\textit{carriage return}, \verb|'\r'|) viene scartato 
                      dall'input. Può essere utile per i terminali che inviano 
                      entrambi i caratteri di ritorno carrello e a capo 
                      (\textit{newline}, \verb|'\n'|).\\
    \constd{ICRNL}  & Se impostato un carattere di ritorno carrello  
                      (\verb|'\r'|) sul terminale viene automaticamente 
                      trasformato in un a capo (\verb|'\n'|) sulla coda di
                      input.\\
    \constd{IUCLC}  & Se impostato trasforma i caratteri maiuscoli dal
                      terminale in minuscoli sull'ingresso (opzione non 
                      POSIX).\\
    \constd{IXON}   & Se impostato attiva il controllo di flusso in uscita con i
                      caratteri di START e STOP. se si riceve
                      uno STOP l'output viene bloccato, e viene fatto
                      ripartire solo da uno START, e questi due
                      caratteri non vengono passati alla coda di input. Se non
                      impostato i due caratteri sono passati alla coda di input
                      insieme agli altri.\\
    \constd{IXANY}  & Se impostato con il controllo di flusso permette a
                      qualunque carattere di far ripartire l'output bloccato da
                      un carattere di STOP.\\
    \constd{IXOFF}  & Se impostato abilita il controllo di flusso in
                      ingresso. Il computer emette un carattere di STOP per
                      bloccare l'input dal terminale e lo sblocca con il
                      carattere START.\\
    \constd{IMAXBEL}& Se impostato fa suonare il cicalino se si riempie la cosa
                      di ingresso; in Linux non è implementato e il kernel si
                      comporta cose se fosse sempre impostato (è una estensione
                      BSD).\\
    \constd{IUTF8}  & Indica che l'input è in UTF-8, cosa che consente di
                      utilizzare la cancellazione dei caratteri in maniera
                      corretta (dal kernel 2.6.4 e non previsto in POSIX).\\
    \hline
  \end{tabular}
  \caption{Costanti identificative dei vari bit del flag di controllo
    \var{c\_iflag} delle modalità di input di un terminale.}
  \label{tab:sess_termios_iflag}
\end{table}

Il primo flag, mantenuto nel campo \var{c\_iflag}, è detto \textsl{flag di
  input} e controlla le modalità di funzionamento dell'input dei caratteri sul
terminale, come il controllo di parità, il controllo di flusso, la gestione
dei caratteri speciali; un elenco dei vari bit, del loro significato e delle
costanti utilizzate per identificarli è riportato in
tab.~\ref{tab:sess_termios_iflag}.

Si noti come alcuni di questi flag (come quelli per la gestione del flusso)
fanno riferimento a delle caratteristiche che ormai sono completamente
obsolete; la maggior parte inoltre è tipica di terminali seriali, e non ha
alcun effetto su dispositivi diversi come le console virtuali o gli
pseudo-terminali usati nelle connessioni di rete.

Il secondo flag, mantenuto nel campo \var{c\_oflag}, è detto \textsl{flag di
  output} e controlla le modalità di funzionamento dell'output dei caratteri,
come l'impacchettamento dei caratteri sullo schermo, la traslazione degli a
capo, la conversione dei caratteri speciali; un elenco dei vari bit, del loro
significato e delle costanti utilizzate per identificarli è riportato in
tab.~\ref{tab:sess_termios_oflag}, di questi solo \const{OPOST} era previsto
da POSIX.1, buona parte degli altri sono stati aggiunti in POSIX.1-2001,
quelli ancora assenti sono stati indicati esplicitamente.

\begin{table}[!htb]
  \footnotesize
  \centering
  \begin{tabular}[c]{|l|p{10cm}|}
    \hline
    \textbf{Valore}& \textbf{Significato}\\
    \hline
    \hline
    \constd{OPOST} & Se impostato i caratteri vengono convertiti opportunamente
                    (in maniera dipendente dall'implementazione) per la 
                    visualizzazione sul terminale, ad esempio al
                    carattere di a capo (NL) può venire aggiunto un ritorno
                    carrello (CR).\\
    \constd{OLCUC} & Se impostato trasforma i caratteri minuscoli in ingresso 
                    in caratteri maiuscoli sull'uscita (non previsto da
                    POSIX).\\
    \constd{ONLCR} & Se impostato converte automaticamente il carattere di a 
                    capo (NL) in un carattere di ritorno carrello (CR).\\
    \constd{OCRNL} & Se impostato converte automaticamente il carattere di a
                    capo (NL) nella coppia di caratteri ritorno carrello, a 
                    capo (CR-NL).\\
    \constd{ONOCR} & Se impostato converte il carattere di ritorno carrello
                    (CR) nella coppia di caratteri CR-NL.\\
    \constd{ONLRET}& Se impostato rimuove dall'output il carattere di ritorno
                    carrello (CR).\\
    \constd{OFILL} & Se impostato in caso di ritardo sulla linea invia dei
                    caratteri di riempimento invece di attendere.\\
    \constd{OFDEL} & Se impostato il carattere di riempimento è DEL
                    (\texttt{0x3F}), invece che NUL (\texttt{0x00}), (non
                    previsto da POSIX e non implementato su Linux).\\
    \constd{NLDLY} & Maschera per i bit che indicano il ritardo per il
                    carattere di a capo (NL), i valori possibili sono 
                    \val{NL0} o \val{NL1}.\\
    \constd{CRDLY} & Maschera per i bit che indicano il ritardo per il
                    carattere ritorno carrello (CR), i valori possibili sono
                    \val{CR0}, \val{CR1}, \val{CR2} o \val{CR3}.\\
    \constd{TABDLY}& Maschera per i bit che indicano il ritardo per il
                    carattere di tabulazione, i valori possibili sono
                    \val{TAB0}, \val{TAB1}, \val{TAB2} o \val{TAB3}.\\
    \constd{BSDLY} & Maschera per i bit che indicano il ritardo per il
                    carattere di ritorno indietro (\textit{backspace}), i
                    valori possibili sono \val{BS0} o \val{BS1}.\\
    \constd{VTDLY} & Maschera per i bit che indicano il ritardo per il
                    carattere di tabulazione verticale, i valori possibili sono
                    \val{VT0} o \val{VT1}.\\
    \constd{FFDLY} & Maschera per i bit che indicano il ritardo per il
                    carattere di pagina nuova (\textit{form feed}), i valori
                    possibili sono \val{FF0} o \val{FF1}.\\
    \hline
  \end{tabular}
  \caption{Costanti identificative dei vari bit del flag di controllo
    \var{c\_oflag} delle modalità di output di un terminale.}
  \label{tab:sess_termios_oflag}
\end{table}

Si noti come alcuni dei valori riportati in tab.~\ref{tab:sess_termios_oflag}
fanno riferimento a delle maschere di bit; essi infatti vengono utilizzati per
impostare alcuni valori numerici relativi ai ritardi nell'output di alcuni
caratteri: una caratteristica originaria dei primi terminali su telescrivente,
che avevano bisogno di tempistiche diverse per spostare il carrello in
risposta ai caratteri speciali, e che oggi sono completamente in disuso.

Si tenga presente inoltre che nel caso delle maschere il valore da inserire in
\var{c\_oflag} deve essere fornito avendo cura di cancellare prima tutti i bit
della maschera, i valori da immettere infatti (quelli riportati nella
spiegazione corrispondente) sono numerici e non per bit, per cui possono
sovrapporsi fra di loro. Occorrerà perciò utilizzare un codice del tipo:

\includecodesnip{listati/oflag.c}

\noindent che prima cancella i bit della maschera in questione e poi imposta
il valore.


\begin{table}[!htb]
  \footnotesize
  \centering
  \begin{tabular}[c]{|l|p{10cm}|}
    \hline
    \textbf{Valore}& \textbf{Significato}\\
    \hline
    \hline
    \constd{CBAUD}  & Maschera dei bit (4+1) usati per impostare della velocità
                      della linea (il \textit{baud rate}) in ingresso; non è
                      presente in POSIX ed in Linux non è implementato in
                      quanto viene usato un apposito campo di
                      \struct{termios}.\\ 
    \constd{CBAUDEX}& Bit aggiuntivo per l'impostazione della velocità della
                      linea, non è presente in POSIX e per le stesse
                      motivazioni del precedente non è implementato in Linux.\\
    \constd{CSIZE}  & Maschera per i bit usati per specificare la dimensione 
                      del carattere inviato lungo la linea di trasmissione, i
                      valore ne indica la lunghezza (in bit), ed i valori   
                      possibili sono \val{CS5}, \val{CS6}, \val{CS7} e \val{CS8}
                      corrispondenti ad un analogo numero di bit.\\
    \constd{CSTOPB} & Se impostato vengono usati due bit di stop sulla linea
                      seriale, se non impostato ne viene usato soltanto uno.\\
    \constd{CREAD}  & Se è impostato si può leggere l'input del terminale,
                      altrimenti i caratteri in ingresso vengono scartati
                      quando arrivano.\\
    \constd{PARENB} & Se impostato abilita la generazione il controllo di
                      parità. La reazione in caso di errori dipende dai
                      relativi valori per \var{c\_iflag}, riportati in 
                      tab.~\ref{tab:sess_termios_iflag}. Se non è impostato i
                      bit di parità non vengono generati e i caratteri non
                      vengono controllati.\\
    \constd{PARODD} & Ha senso solo se è attivo anche \const{PARENB}. Se 
                      impostato viene usata una parità è dispari, altrimenti 
                      viene usata una parità pari.\\
    \constd{HUPCL}  & Se è impostato viene distaccata la connessione del
                      modem quando l'ultimo dei processi che ha ancora un file
                      aperto sul terminale lo chiude o esce.\\
    \constd{LOBLK}  & Se impostato blocca l'output su un strato di shell non
                      corrente, non è presente in POSIX e non è implementato
                      da Linux.\\
    \constd{CLOCAL} & Se impostato indica che il terminale è connesso in locale
                      e che le linee di controllo del modem devono essere
                      ignorate. Se non impostato effettuando una chiamata ad
                      \func{open} senza aver specificato il flag di
                      \const{O\_NONBLOCK} si bloccherà il processo finché 
                      non si è stabilita una connessione con il modem; inoltre 
                      se viene rilevata una disconnessione viene inviato un
                      segnale di \signal{SIGHUP} al processo di controllo del
                      terminale. La lettura su un terminale sconnesso comporta
                      una condizione di \textit{end of file} e la scrittura un
                      errore di \errcode{EIO}.\\
    \constd{CIBAUD} & Maschera dei bit della velocità della linea in
                      ingresso; analogo a \const{CBAUD}, non è previsto da
                      POSIX e non è implementato in Linux dato che è
                      mantenuto in un apposito campo di \struct{termios}.\\
    \constd{CMSPAR} & imposta un bit di parità costante: se \const{PARODD} è
                      impostato la parità è sempre 1 (\textit{MARK}) se non è
                      impostato la parità è sempre 0 (\textit{SPACE}), non è
                      previsto da POSIX.\\ 
% vedi: http://www.lothosoft.ch/thomas/libmip/markspaceparity.php
    \constd{CRTSCTS}& Abilita il controllo di flusso hardware sulla seriale,
                      attraverso l'utilizzo delle dei due fili di RTS e CTS.\\
    \hline
  \end{tabular}
  \caption{Costanti identificative dei vari bit del flag di controllo
    \var{c\_cflag} delle modalità di controllo di un terminale.}
  \label{tab:sess_termios_cflag}
\end{table}

Il terzo flag, mantenuto nel campo \var{c\_cflag}, è detto \textsl{flag di
  controllo} ed è legato al funzionamento delle linee seriali, permettendo di
impostarne varie caratteristiche, come il numero di bit di stop, le
impostazioni della parità, il funzionamento del controllo di flusso; esso ha
senso solo per i terminali connessi a linee seriali. Un elenco dei vari bit,
del loro significato e delle costanti utilizzate per identificarli è riportato
in tab.~\ref{tab:sess_termios_cflag}.

I valori previsti per questo flag sono molto specifici, e completamente
attinenti al controllo delle modalità operative di un terminale che opera
attraverso una linea seriale; essi pertanto non hanno nessuna rilevanza per i
terminali che usano un'altra interfaccia fisica, come le console virtuali e
gli pseudo-terminali usati dalle connessioni di rete.

Inoltre alcuni valori di questi flag sono previsti solo per quelle
implementazioni (lo standard POSIX non specifica nulla riguardo
l'implementazione, ma solo delle funzioni di lettura e scrittura) che
mantengono le velocità delle linee seriali all'interno dei flag; come
accennato in Linux questo viene fatto (seguendo l'esempio di BSD) attraverso
due campi aggiuntivi, \var{c\_ispeed} e \var{c\_ospeed}, nella struttura
\struct{termios} (mostrati in fig.~\ref{fig:term_termios}).

\begin{table}[b!ht]
  \footnotesize
  \centering
  \begin{tabular}[c]{|l|p{10cm}|}
    \hline
    \textbf{Valore}& \textbf{Significato}\\
    \hline
    \hline
    \constd{ISIG}   & Se impostato abilita il riconoscimento dei caratteri
                      INTR, QUIT, e SUSP generando il relativo segnale.\\
    \constd{ICANON} & Se impostato il terminale opera in modalità canonica,
                      altrimenti opera in modalità non canonica.\\
    \constd{XCASE}  & Se impostato il terminale funziona solo con le
                      maiuscole. L'input è convertito in minuscole tranne per i
                      caratteri preceduti da una ``\texttt{\bslash}''. In output
                      le maiuscole sono precedute da una ``\texttt{\bslash}'' e 
                      le minuscole convertite in maiuscole. Non è presente in
                      POSIX.\\
    \constd{ECHO}   & Se è impostato viene attivato l'eco dei caratteri in
                      input sull'output del terminale.\\
    \constd{ECHOE}  & Se è impostato l'eco mostra la cancellazione di un
                      carattere in input (in reazione al carattere ERASE)
                      cancellando l'ultimo carattere della riga corrente dallo
                      schermo; altrimenti il carattere è rimandato in eco per
                      mostrare quanto accaduto (usato per i terminali con
                      l'uscita su una stampante).\\
    \constd{ECHOK}  & Se impostato abilita il trattamento della visualizzazione
                      del carattere KILL, andando a capo dopo aver visualizzato
                      lo stesso, altrimenti viene solo mostrato il carattere e
                      sta all'utente ricordare che l'input precedente è stato
                      cancellato.\\
    \constd{ECHONL} & Se impostato viene effettuato l'eco di un a
                      capo (\verb|\n|) anche se non è stato impostato
                      \const{ECHO}.\\
    \constd{ECHOCTL}& Se impostato insieme ad \const{ECHO} i caratteri di
                      controllo ASCII (tranne TAB, NL, START, e STOP) sono
                      mostrati nella forma che prepone un ``\texttt{\circonf}'' 
                      alla lettera ottenuta sommando \texttt{0x40} al valore del
                      carattere (di solito questi si possono ottenere anche
                      direttamente premendo il tasto \texttt{ctrl} più la
                      relativa lettera). Non è presente in POSIX.\\
    \constd{ECHOPRT}& Se impostato abilita la visualizzazione del carattere di
                      cancellazione in una modalità adatta ai terminali con
                      l'uscita su stampante; l'invio del carattere di ERASE
                      comporta la stampa di un ``\texttt{|}'' seguito dal
                      carattere cancellato, e così via in caso di successive
                      cancellazioni, quando si riprende ad immettere carattere 
                      normali prima verrà stampata una ``\texttt{/}''. Non è
                      presente in POSIX.\\
    \constd{ECHOKE} & Se impostato abilita il trattamento della visualizzazione
                      del carattere KILL cancellando i caratteri precedenti
                      nella linea secondo le modalità specificate dai valori di
                      \const{ECHOE} e \const{ECHOPRT}. Non è presente in
                      POSIX.\\ 
    \constd{DEFECHO}& Se impostato effettua l'eco solo se c'è un processo in
                      lettura. Non è presente in POSIX e non è supportato da
                      Linux.\\
    \constd{FLUSHO} & Effettua la cancellazione della coda di uscita. Viene
                      attivato dal carattere DISCARD. Non è presente in POSIX e
                      non è supportato da Linux.\\
    \constd{NOFLSH} & Se impostato disabilita lo scarico delle code di ingresso
                      e uscita quando vengono emessi i segnali \signal{SIGINT}, 
                      \signal{SIGQUIT} e \signal{SIGSUSP}.\\
    \constd{TOSTOP} & Se abilitato, con il supporto per il job control presente,
                      genera il segnale \signal{SIGTTOU} per un processo in
                      background che cerca di scrivere sul terminale.\\
    \constd{PENDIN} & Indica che la linea deve essere ristampata, viene
                      attivato dal carattere REPRINT e resta attivo fino alla
                      fine della ristampa. Non è presente in POSIX e
                      non è supportato in Linux.\\
    \constd{IEXTEN} & Abilita alcune estensioni previste dalla
                      implementazione. Deve essere impostato perché caratteri
                      speciali come EOL2, LNEXT, REPRINT e WERASE possano
                      essere interpretati.\\
    \hline
  \end{tabular}
  \caption{Costanti identificative dei vari bit del flag di controllo
    \var{c\_lflag} delle modalità locali di un terminale.}
  \label{tab:sess_termios_lflag}
\end{table}

Il quarto flag, mantenuto nel campo \var{c\_lflag}, è detto \textsl{flag
  locale}, e serve per controllare il funzionamento dell'interfaccia fra il
driver e l'utente, come abilitare l'eco, gestire i caratteri di controllo e
l'emissione dei segnali, impostare modo canonico o non canonico. Un elenco dei
vari bit, del loro significato e delle costanti utilizzate per identificarli è
riportato in tab.~\ref{tab:sess_termios_lflag}. Con i terminali odierni
l'unico flag con cui probabilmente si può avere a che fare è questo, in quanto
è con questo che si impostano le caratteristiche generiche comuni a tutti i
terminali.

Si tenga presente che i flag che riguardano le modalità di eco dei caratteri
(\const{ECHOE}, \const{ECHOPRT}, \const{ECHOK}, \const{ECHOKE},
\const{ECHONL}) controllano solo il comportamento della visualizzazione, il
riconoscimento dei vari caratteri dipende dalla modalità di operazione, ed
avviene solo in modo canonico, pertanto questi flag non hanno significato se
non è impostato \const{ICANON}.

Oltre ai vari flag per gestire le varie caratteristiche dei terminali,
\struct{termios} contiene pure il campo \var{c\_cc} che viene usato per
impostare i caratteri speciali associati alle varie funzioni di controllo. Il
numero di questi caratteri speciali è indicato dalla costante \constd{NCCS},
POSIX ne specifica almeno 11, ma molte implementazioni ne definiscono molti
altri.\footnote{in Linux il valore della costante è 32, anche se i caratteri
  effettivamente definiti sono solo 17.}

\begin{table}[!htb]
  \footnotesize
  \centering
  \begin{tabular}[c]{|l|c|c|p{9cm}|}
    \hline
    \textbf{Indice} & \textbf{Valore}&\textbf{Codice} & \textbf{Funzione}\\
    \hline
    \hline
    \constd{VINTR} &\texttt{0x03}&(\texttt{C-c})& Carattere di interrupt, 
                                                  provoca l'emissione di 
                                                  \signal{SIGINT}.\\
    \constd{VQUIT} &\texttt{0x1C}&(\texttt{C-\bslash})& Carattere di uscita,  
                                                        provoca l'emissione di 
                                                        \signal{SIGQUIT}.\\
    \constd{VERASE}&\texttt{0x7f}&DEL,\texttt{C-?}& Carattere di ERASE, cancella
                                                    l'ultimo carattere
                                                    precedente nella linea.\\
    \constd{VKILL} &\texttt{0x15}&(\texttt{C-u})& Carattere di KILL, cancella
                                                  l'intera riga.\\
    \constd{VEOF}  &\texttt{0x04}&(\texttt{C-d})& Carattere di
                                                  \textit{end-of-file}. Causa
                                                  l'invio del contenuto del
                                                  buffer di ingresso al
                                                  processo in lettura anche se
                                                  non è ancora stato ricevuto
                                                  un a capo. Se è il primo
                                                  carattere immesso comporta il
                                                  ritorno di \func{read} con
                                                  zero caratteri, cioè la
                                                  condizione di
                                                  \textit{end-of-file}.\\
    \constd{VMIN}  &     ---     &  ---   & Numero minimo di caratteri per una 
                                            lettura in modo non canonico.\\
    \constd{VEOL}  &\texttt{0x00}& NUL &    Carattere di fine riga. Agisce come
                                         un a capo, ma non viene scartato ed
                                         è letto come l'ultimo carattere
                                         nella riga.\\
    \constd{VTIME} &     ---     &  ---   & Timeout, in decimi di secondo, per
                                            una lettura in modo non canonico.\\
    \constd{VEOL2} &\texttt{0x00}&   NUL  & Ulteriore carattere di fine
                                            riga. Ha lo stesso effetto di
                                            \const{VEOL} ma può essere un
                                            carattere diverso. \\
    \constd{VSWTC} &\texttt{0x00}&   NUL  & Carattere di switch. Non supportato
                                            in Linux.\\
    \constd{VSTART}&\texttt{0x17}&(\texttt{C-q})& Carattere di START. Riavvia un
                                                 output bloccato da uno STOP.\\
    \constd{VSTOP} &\texttt{0x19}&(\texttt{C-s})& Carattere di STOP. Blocca
                                                 l'output fintanto che non
                                                 viene premuto un carattere di
                                                 START.\\
    \constd{VSUSP} &\texttt{0x1A}&(\texttt{C-z})& Carattere di
                                                 sospensione. Invia il segnale
                                                 \signal{SIGTSTP}.\\
    \constd{VDSUSP}&\texttt{0x19}&(\texttt{C-y})& Carattere di sospensione
                                                 ritardata. Invia il segnale 
                                                 \signal{SIGTSTP} quando il
                                                 carattere viene letto dal
                                                 programma, (non presente in
                                                 POSIX e non riconosciuto in
                                                 Linux).\\ 
    \constd{VLNEXT}&\texttt{0x16}&(\texttt{C-v})& Carattere di escape, serve a
                                                 quotare il carattere
                                                 successivo che non viene
                                                 interpretato ma passato
                                                 direttamente all'output.\\
    \constd{VWERASE}&\texttt{0x17}&(\texttt{C-w})&Cancellazione di una
                                                 parola.\\
    \constd{VREPRINT}&\texttt{0x12}&(\texttt{C-r})& Ristampa i caratteri non
                                                 ancora letti (non presente in
                                                 POSIX).\\
    \constd{VDISCARD}&\texttt{0x0F}&(\texttt{C-o})& Non riconosciuto in Linux.\\
    \constd{VSTATUS} &\texttt{0x13}&(\texttt{C-t})& Non riconosciuto in Linux.\\
    \hline
  \end{tabular}
  \caption{Valori dei caratteri di controllo mantenuti nel campo \var{c\_cc}
    della struttura \struct{termios}.} 
  \label{tab:sess_termios_cc}
\end{table}


A ciascuna di queste funzioni di controllo corrisponde un elemento del vettore
\var{c\_cc} che specifica quale è il carattere speciale associato; per
portabilità invece di essere indicati con la loro posizione numerica nel
vettore, i vari elementi vengono indicizzati attraverso delle opportune
costanti, il cui nome corrisponde all'azione ad essi associata. Un elenco
completo dei caratteri di controllo, con le costanti e delle funzionalità
associate è riportato in tab.~\ref{tab:sess_termios_cc}, usando quelle
definizioni diventa possibile assegnare un nuovo carattere di controllo con un
codice del tipo:
\includecodesnip{listati/value_c_cc.c}

La maggior parte di questi caratteri (tutti tranne \const{VTIME} e
\const{VMIN}) hanno effetto solo quando il terminale viene utilizzato in modo
canonico; per alcuni devono essere soddisfatte ulteriori richieste, ad esempio
\const{VINTR}, \const{VSUSP}, e \const{VQUIT} richiedono sia impostato
\const{ISIG}; \const{VSTART} e \const{VSTOP} richiedono sia impostato
\const{IXON}; \const{VLNEXT}, \const{VWERASE}, \const{VREPRINT} richiedono sia
impostato \const{IEXTEN}.  In ogni caso quando vengono attivati i caratteri
vengono interpretati e non sono passati sulla coda di ingresso.

Per leggere ed scrivere tutte le varie impostazioni dei terminali viste finora
lo standard POSIX prevede due funzioni che utilizzano come argomento un
puntatore ad una struttura \struct{termios} che sarà quella in cui andranno
immagazzinate le impostazioni.  Le funzioni sono \funcd{tcgetattr} e
\funcd{tcsetattr} ed il loro prototipo è:

\begin{funcproto}{
\fhead{unistd.h}
\fhead{termios.h}

\fdecl{int tcgetattr(int fd, struct termios *termios\_p)} 
\fdesc{Legge il valore delle impostazioni di un terminale.} 
\fdecl{int tcsetattr(int fd, int optional\_actions, struct termios *termios\_p)}
\fdesc{Scrive le impostazioni di un terminale.}
}

{Le funzioni ritornano $0$ in caso di successo e $-1$ per un errore, nel qual
  caso \var{errno} assumerà uno dei valori: 
  \begin{errlist}
    \item[\errcode{EINTR}] la funzione è stata interrotta. 
  \end{errlist}
  ed inoltre \errval{EBADF}, \errval{ENOTTY} ed \errval{EINVAL} nel loro
  significato generico.}
\end{funcproto}


Le funzioni operano sul terminale cui fa riferimento il file descriptor
\param{fd} utilizzando la struttura indicata dal puntatore \param{termios\_p}
per lo scambio dei dati. Si tenga presente che le impostazioni sono associate
al terminale e non al file descriptor; questo significa che se si è cambiata
una impostazione un qualunque altro processo che apra lo stesso terminale, od
un qualunque altro file descriptor che vi faccia riferimento, vedrà le nuove
impostazioni pur non avendo nulla a che fare con il file descriptor che si è
usato per effettuare i cambiamenti.

Questo significa che non è possibile usare file descriptor diversi per
utilizzare automaticamente il terminale in modalità diverse, se esiste una
necessità di accesso differenziato di questo tipo occorrerà cambiare
esplicitamente la modalità tutte le volte che si passa da un file descriptor
ad un altro.

La funzione \func{tcgetattr} legge i valori correnti delle impostazioni di un
terminale qualunque nella struttura puntata da \param{termios\_p};
\func{tcsetattr} invece effettua la scrittura delle impostazioni e quando
viene invocata sul proprio terminale di controllo può essere eseguita con
successo solo da un processo in foreground.  Se invocata da un processo in
background infatti tutto il gruppo riceverà un segnale di \signal{SIGTTOU} come
se si fosse tentata una scrittura, a meno che il processo chiamante non abbia
\signal{SIGTTOU} ignorato o bloccato, nel qual caso l'operazione sarà eseguita.

La funzione \func{tcsetattr} prevede tre diverse modalità di funzionamento,
specificabili attraverso l'argomento \param{optional\_actions}, che permette
di stabilire come viene eseguito il cambiamento delle impostazioni del
terminale, i valori possibili sono riportati in
tab.~\ref{tab:sess_tcsetattr_option}; di norma (come fatto per le due funzioni
di esempio) si usa sempre \const{TCSANOW}, le altre opzioni possono essere
utili qualora si cambino i parametri di output.

\begin{table}[htb]
  \footnotesize
  \centering
  \begin{tabular}[c]{|l|p{8cm}|}
    \hline
    \textbf{Valore}& \textbf{Significato}\\
    \hline
    \hline
    \constd{TCSANOW}  & Esegue i cambiamenti in maniera immediata.\\
    \constd{TCSADRAIN}& I cambiamenti vengono eseguiti dopo aver atteso che
                        tutto l'output presente sulle code è stato scritto.\\
    \constd{TCSAFLUSH}& È identico a \const{TCSADRAIN}, ma in più scarta
                        tutti i dati presenti sulla coda di input.\\
    \hline
  \end{tabular}
  \caption{Possibili valori per l'argomento \param{optional\_actions} della
    funzione \func{tcsetattr}.} 
  \label{tab:sess_tcsetattr_option}
\end{table}

Occorre infine tenere presente che \func{tcsetattr} ritorna con successo anche
se soltanto uno dei cambiamenti richiesti è stato eseguito. Pertanto se si
effettuano più cambiamenti è buona norma controllare con una ulteriore
chiamata a \func{tcgetattr} che essi siano stati eseguiti tutti quanti.

\begin{figure}[!htbp]
  \footnotesize \centering
  \begin{minipage}[c]{\codesamplewidth}
    \includecodesample{listati/SetTermAttr.c}
  \end{minipage} 
  \normalsize 
  \caption{Codice della funzione \func{SetTermAttr} che permette di
    impostare uno dei flag di controllo locale del terminale.}
  \label{fig:term_set_attr}
\end{figure}

Come già accennato per i cambiamenti effettuati ai vari flag di controllo
occorre che i valori di ciascun bit siano specificati avendo cura di mantenere
intatti gli altri; per questo motivo in generale si deve prima leggere il
valore corrente delle impostazioni con \func{tcgetattr} per poi modificare i
valori impostati.

In fig.~\ref{fig:term_set_attr} e fig.~\ref{fig:term_unset_attr} si è
riportato rispettivamente il codice delle due funzioni \func{SetTermAttr} e
\func{UnSetTermAttr}, che possono essere usate per impostare o rimuovere, con
le dovute precauzioni, un qualunque bit di \var{c\_lflag}. Il codice completo
di entrambe le funzioni può essere trovato nel file \file{SetTermAttr.c} dei
sorgenti allegati alla guida.

La funzione \func{SetTermAttr} provvede ad impostare il bit specificato
dall'argomento \param{flag}; prima si leggono i valori correnti
(\texttt{\small 8}) con \func{tcgetattr}, uscendo con un messaggio in caso di
errore (\texttt{\small 9--10}), poi si provvede a impostare solo i bit
richiesti (possono essere più di uno) con un OR binario (\texttt{\small 12});
infine si scrive il nuovo valore modificato con \func{tcsetattr}
(\texttt{\small 13}), notificando un eventuale errore (\texttt{\small 14--15})
o uscendo normalmente.

\begin{figure}[!htbp]
  \footnotesize \centering
  \begin{minipage}[c]{\codesamplewidth}
    \includecodesample{listati/UnSetTermAttr.c}
  \end{minipage} 
  \normalsize 
  \caption{Codice della funzione \func{UnSetTermAttr} che permette di
    rimuovere uno dei flag di controllo locale del terminale.} 
  \label{fig:term_unset_attr}
\end{figure}

La seconda funzione, \func{UnSetTermAttr}, è assolutamente identica alla
prima, solo che in questo caso (\texttt{\small 9}) si rimuovono i bit
specificati dall'argomento \param{flag} usando un AND binario del valore
negato.

Al contrario di tutte le altre caratteristiche dei terminali, che possono
essere impostate esplicitamente utilizzando gli opportuni campi di
\struct{termios}, per le velocità della linea (il cosiddetto \textit{baud
  rate}) non è prevista una implementazione standardizzata, per cui anche se
in Linux sono mantenute in due campi dedicati nella struttura, questi non
devono essere acceduti direttamente ma solo attraverso le apposite funzioni di
interfaccia provviste da POSIX.1.

Lo standard prevede due funzioni per scrivere la velocità delle linee seriali,
\funcd{cfsetispeed} per la velocità della linea di ingresso e
\funcd{cfsetospeed} per la velocità della linea di uscita; i loro prototipi
sono:

\begin{funcproto}{
\fhead{unistd.h}
\fhead{termios.h}
\fdecl{int cfsetispeed(struct termios *termios\_p, speed\_t speed)}
\fdesc{Imposta la velocità delle linee seriali in ingresso.}
\fdecl{int cfsetospeed(struct termios *termios\_p, speed\_t speed)} 
\fdesc{Imposta la velocità delle linee seriali in uscita.} 
}

{Le funzioni ritornano $0$ in caso di successo e $-1$ per un errore, che
  avviene solo quando il valore specificato non è valido.}
\end{funcproto}

Si noti che le funzioni si limitano a scrivere opportunamente il valore della
velocità prescelta \param{speed} all'interno della struttura puntata da
\param{termios\_p}; per effettuare l'impostazione effettiva occorrerà poi
chiamare \func{tcsetattr}.

Si tenga presente che per le linee seriali solo alcuni valori di velocità sono
validi; questi possono essere specificati direttamente (la \acr{glibc}
prevede che i valori siano indicati in bit per secondo), ma in generale
altre versioni di libreria possono utilizzare dei valori diversi. Per questo
POSIX.1 prevede una serie di costanti che però servono solo per specificare le
velocità tipiche delle linee seriali:
\begin{verbatim}
     B0       B50      B75      B110     B134     B150     B200     
     B300     B600     B1200    B1800    B2400    B4800    B9600    
     B19200   B38400   B57600   B115200  B230400  B460800
\end{verbatim}

Un terminale può utilizzare solo alcune delle velocità possibili, le funzioni
però non controllano se il valore specificato è valido, dato che non possono
sapere a quale terminale le velocità saranno applicate; sarà l'esecuzione di
\func{tcsetattr} a fallire quando si cercherà di eseguire l'impostazione.

Di norma il valore ha senso solo per i terminali seriali dove indica appunto
la velocità della linea di trasmissione; se questa non corrisponde a quella
del terminale quest'ultimo non potrà funzionare: quando il terminale non è
seriale il valore non influisce sulla velocità di trasmissione dei dati. 

In generale impostare un valore nullo (\val{B0}) sulla linea di output fa sì
che il modem non asserisca più le linee di controllo, interrompendo di fatto
la connessione, qualora invece si utilizzi questo valore per la linea di input
l'effetto sarà quello di rendere la sua velocità identica a quella della linea
di output.

Dato che in genere si imposta sempre la stessa velocità sulle linee di uscita
e di ingresso è supportata anche la funzione \funcd{cfsetspeed}, una
estensione di BSD (la funzione origina da 4.4BSD e richiede sia definita la
macro \macro{\_BSD\_SOURCE}) il cui prototipo è:

\begin{funcproto}{
\fhead{unistd.h} 
\fhead{termios.h}
\fdecl{int cfsetspeed(struct termios *termios\_p, speed\_t speed)}
\fdesc{Imposta la velocità delle linee seriali.} 
}

{La funzione ritorna $0$ in caso di successo e $-1$ per un errore, che avviene
  solo quando il valore specificato non è valido.}
\end{funcproto}

\noindent la funzione è identica alle due precedenti ma imposta la stessa
velocità sia per la linea di ingresso che per quella di uscita.

Analogamente a quanto avviene per l'impostazione, le velocità possono essere
lette da una struttura \struct{termios} utilizzando altre due funzioni,
\funcd{cfgetispeed} e \funcd{cfgetospeed}, i cui prototipi sono:

\begin{funcproto}{
\fhead{unistd.h} 
\fhead{termios.h}
\fdecl{speed\_t cfgetispeed(struct termios *termios\_p)} 
\fdesc{Legge la velocità delle linee seriali in ingresso.}
\fdecl{speed\_t cfgetospeed(struct termios *termios\_p)} 
\fdesc{Legge la velocità delle linee seriali in uscita.}
}

{Le funzioni ritornano la velocità della linea, non sono previste condizioni
  di errore.}
\end{funcproto}

Anche in questo caso le due funzioni estraggono i valori della velocità della
linea da una struttura, il cui indirizzo è specificato dall'argomento
\param{termios\_p} che deve essere stata letta in precedenza con
\func{tcgetattr}.

Infine sempre da BSD è stata ripresa una funzione che consente di impostare il
terminale in una modalità analoga alla cosiddetta modalità ``\textit{raw}'' di
System V, in cui i dati in input vengono resi disponibili un carattere alla
volta, e l'eco e tutte le interpretazioni dei caratteri in entrata e uscita
sono disabilitate. La funzione è \funcd{cfmakeraw} ed il suo prototipo è:

\begin{funcproto}{
\fhead{unistd.h} 
\fhead{termios.h}
\fdecl{void cfmakeraw(struct termios *termios\_p)} 
\fdesc{Imposta il terminale in  modalità ``\textit{raw}''.}
}

{La funzione imposta solo i valori in \param{termios\_p}, e non
    sono previste condizioni di errore.}
\end{funcproto}

Anche in questo caso la funzione si limita a preparare i valori che poi
saranno impostato con una successiva chiamata a \func{tcsetattr}, in sostanza
la funzione è equivalente a:
\includecodesnip{listati/cfmakeraw.c}


\subsection{La gestione della disciplina di linea.}
\label{sec:term_line_discipline}

Come illustrato dalla struttura riportata in fig.~\ref{fig:term_struct} tutti
i terminali hanno un insieme di funzionalità comuni, che prevedono la presenza
di code di ingresso ed uscita; in generale si fa riferimento a queste
funzionalità con il nome di \textsl{disciplina di
  linea}.\index{disciplina~di~linea} Lo standard POSIX prevede alcune funzioni
che permettono di intervenire direttamente sulla gestione della disciplina di
linea e sull'interazione fra i dati in ingresso ed uscita e le relative code.

In generale tutte queste funzioni vengono considerate, dal punto di vista
dell'accesso al terminale, come delle funzioni di scrittura, pertanto se usate
da processi in background sul loro terminale di controllo provocano
l'emissione di \signal{SIGTTOU}, come illustrato in
sez.~\ref{sec:sess_ctrl_term}, con la stessa eccezione, già vista per
\func{tcsetattr}, che quest'ultimo sia bloccato o ignorato dal processo
chiamante.

Una prima funzione, che è efficace solo in caso di terminali seriali
asincroni, e non fa niente per tutti gli altri terminali, è
\funcd{tcsendbreak}; il suo prototipo è:

\begin{funcproto}{
\fhead{unistd.h} 
\fhead{termios.h}
\fdecl{int tcsendbreak(int fd, int duration)}
\fdesc{Genera una condizione di \textit{break}.}  

}

{La funzione ritorna $0$ in caso di successo e $-1$ per un errore, nel qual
  caso \var{errno} assumerà uno dei valori \errval{EBADF} o \errval{ENOTTY}
  nel loro significato generico.}
\end{funcproto}

La funzione invia un flusso di bit nulli, che genera una condizione di
\textit{break}, sul terminale associato a \param{fd}. Un valore nullo
di \param{duration} implica una durata del flusso fra 0.25 e 0.5 secondi, un
valore diverso da zero implica una durata pari a \code{duration*T} dove
\code{T} è un valore compreso fra 0.25 e 0.5 secondi. Lo standard POSIX
specifica il comportamento solo nel caso si sia impostato un valore nullo
per \param{duration}, il comportamento negli altri casi può dipendere
dall'implementazione.

Le altre funzioni previste dallo standard POSIX servono a controllare il
comportamento dell'interazione fra le code associate al terminale e l'utente;
la prima di queste è \funcd{tcdrain}, il cui prototipo è:


\begin{funcproto}{
\fhead{unistd.h} 
\fhead{termios.h}
\fdecl{int tcdrain(int fd)}
\fdesc{Attende lo svuotamento della coda di uscita.} 
}

{La funzione ritorna $0$ in caso di successo e $-1$ per un errore, nel qual
  caso \var{errno} assumerà i valori \errval{EBADF} o \errval{ENOTTY}.}
\end{funcproto}

La funzione blocca il processo fino a che tutto l'output presente sulla coda
di uscita non è stato trasmesso al terminale associato ad \param{fd}. % La
                                % funzione è  un punto di cancellazione per i
                                % programmi multi-thread, in tal caso le
                                % chiamate devono essere protette con dei
                                % gestori di cancellazione. 

Una seconda funzione, \funcd{tcflush}, permette svuotare immediatamente le code
di cancellando tutti i dati presenti al loro interno; il suo prototipo è:

\begin{funcproto}{
\fhead{unistd.h}
\fhead{termios.h}
\fdecl{int tcflush(int fd, int queue)}
\fdesc{Cancella i dati presenti nelle code di ingresso o di uscita.}  }

{La funzione ritorna $0$ in caso di successo e $-1$ per un errore, nel qual
  caso \var{errno} assumerà i valori \errval{EBADF} o \errval{ENOTTY}.}
\end{funcproto}

La funzione agisce sul terminale associato a \param{fd}, l'argomento
\param{queue} permette di specificare su quale coda (ingresso, uscita o
entrambe), operare. Esso può prendere i valori riportati in
tab.~\ref{tab:sess_tcflush_queue}, nel caso si specifichi la coda di ingresso
cancellerà i dati ricevuti ma non ancora letti, nel caso si specifichi la coda
di uscita cancellerà i dati scritti ma non ancora trasmessi.

\begin{table}[htb]
  \footnotesize
  \centering
  \begin{tabular}[c]{|l|l|}
    \hline
    \textbf{Valore}& \textbf{Significato}\\
    \hline
    \hline
    \constd{TCIFLUSH} & Cancella i dati sulla coda di ingresso.\\
    \constd{TCOFLUSH} & Cancella i dati sulla coda di uscita. \\
    \constd{TCIOFLUSH}& Cancella i dati su entrambe le code.\\
    \hline
  \end{tabular}
  \caption{Possibili valori per l'argomento \param{queue} della
    funzione \func{tcflush}.} 
  \label{tab:sess_tcflush_queue}
\end{table}


L'ultima funzione dell'interfaccia che interviene sulla disciplina di linea è
\funcd{tcflow}, che viene usata per sospendere la trasmissione e la ricezione
dei dati sul terminale; il suo prototipo è:

\begin{funcproto}{
\fhead{unistd.h}
\fhead{termios.h}
\fdecl{int tcflow(int fd, int action)}
\fdesc{Sospende e riavvia il flusso dei dati sul terminale.} 
}

{La funzione ritorna $0$ in caso di successo e $-1$ per un errore, nel qual
  caso \var{errno} assumerà i valori \errval{EBADF} o \errval{ENOTTY}.}
\end{funcproto}

La funzione permette di controllare (interrompendo e facendo riprendere) il
flusso dei dati fra il terminale ed il sistema sia in ingresso che in uscita.
Il comportamento della funzione è regolato dall'argomento \param{action}, i
cui possibili valori, e relativa azione eseguita dalla funzione, sono
riportati in tab.~\ref{tab:sess_tcflow_action}.

\begin{table}[htb]
   \footnotesize
  \centering
  \begin{tabular}[c]{|l|p{8cm}|}
    \hline
    \textbf{Valore}& \textbf{Azione}\\
    \hline
    \hline
    \constd{TCOOFF}& Sospende l'output.\\
    \constd{TCOON} & Riprende un output precedentemente sospeso.\\
    \constd{TCIOFF}& Il sistema trasmette un carattere di STOP, che 
                     fa interrompere la trasmissione dei dati dal terminale.\\
    \constd{TCION} & Il sistema trasmette un carattere di START, che 
                     fa riprendere la trasmissione dei dati dal terminale.\\
    \hline
  \end{tabular}
  \caption{Possibili valori per l'argomento \param{action} della
    funzione \func{tcflow}.} 
  \label{tab:sess_tcflow_action}
\end{table}



\subsection{Operare in \textsl{modo non canonico}}
\label{sec:term_non_canonical}

Operare con un terminale in modo canonico è relativamente semplice; basta
eseguire una lettura e la funzione ritornerà quando il terminale avrà
completato una linea di input. Non è detto che la linea sia letta interamente
(si può aver richiesto un numero inferiore di byte) ma in ogni caso nessun
dato verrà perso, e il resto della linea sarà letto alla chiamata successiva.

Inoltre in modo canonico la gestione dei dati in ingresso è di norma eseguita
direttamente dal kernel, che si incarica (a seconda di quanto impostato con le
funzioni viste nei paragrafi precedenti) di cancellare i caratteri, bloccare e
riavviare il flusso dei dati, terminare la linea quando viene ricevuti uno dei
vari caratteri di terminazione (NL, EOL, EOL2, EOF).

In modo non canonico è invece compito del programma gestire tutto quanto, i
caratteri NL, EOL, EOL2, EOF, ERASE, KILL, CR, REPRINT non vengono
interpretati automaticamente ed inoltre, non dividendo più l'input in linee,
il sistema non ha più un limite definito su quando ritornare i dati ad un
processo. Per questo motivo abbiamo visto che in \var{c\_cc} sono previsti due
caratteri speciali, MIN e TIME (specificati dagli indici \const{VMIN} e
\const{VTIME} in \var{c\_cc}) che dicono al sistema di ritornare da una
\func{read} quando è stata letta una determinata quantità di dati o è passato
un certo tempo.

Come accennato nella relativa spiegazione in tab.~\ref{tab:sess_termios_cc},
TIME e MIN non sono in realtà caratteri ma valori numerici. Il comportamento
del sistema per un terminale in modalità non canonica prevede quattro casi
distinti:
\begin{description}
\item[MIN$>0$, TIME$>0$] In questo caso MIN stabilisce il numero minimo di
  caratteri desiderati e TIME un tempo di attesa, in decimi di secondo, fra un
  carattere e l'altro. Una \func{read} ritorna se vengono ricevuti almeno MIN
  caratteri prima della scadenza di TIME (MIN è solo un limite inferiore, se
  la funzione ha richiesto un numero maggiore di caratteri ne possono essere
  restituiti di più); se invece TIME scade vengono restituiti i byte ricevuti
  fino ad allora (un carattere viene sempre letto, dato che il timer inizia a
  scorrere solo dopo la ricezione del primo carattere).
\item[MIN$>0$, TIME$=0$] Una \func{read} ritorna solo dopo che sono stati
  ricevuti almeno MIN caratteri. Questo significa che una \func{read} può
  bloccarsi indefinitamente. 
\item[MIN$=0$, TIME$>0$] In questo caso TIME indica un tempo di attesa dalla
  chiamata di \func{read}, la funzione ritorna non appena viene ricevuto un
  carattere o scade il tempo. Si noti che è possibile che \func{read} ritorni
  con un valore nullo.
\item[MIN$=0$, TIME$=0$] In questo caso una \func{read} ritorna immediatamente
  restituendo tutti i caratteri ricevuti. Anche in questo caso può ritornare
  con un valore nullo.
\end{description}



\section{La gestione dei terminali virtuali}
\label{sec:sess_virtual_terminal}

% 
% TODO terminali virtuali 
% Qui c'è da mettere tutta la parte sui terminali virtuali, e la gestione
% degli stessi
%

Da fare.

\subsection{I terminali virtuali}
\label{sec:sess_pty}

Qui vanno spiegati i terminali virtuali, \file{/dev/pty} e compagnia.
% vedi man pts
% vedi 


\subsection{Allocazione dei terminali virtuali}
\label{sec:sess_openpty}

Qui vanno le cose su \func{openpty} e compagnia.

% TODO le ioctl dei terminali (man tty_ioctl)
% e http://www.net-security.org/article.php?id=83
% TODO trattare \func{posix\_openpt}
% vedi http://lwn.net/Articles/688809/,
% http://man7.org/linux/man-pages/man3/ptsname.3.html


% TODO materiale sulle seriali
% vedi http://www.easysw.com/~mike/serial/serial.html
% TODO materiale generico sul layer TTY
% vedi http://www.linusakesson.net/programming/tty/index.php

% LocalWords:  kernel multitasking job control BSD POSIX shell sez group
% LocalWords:  foreground process bg fg waitpid WUNTRACED pgrp session sched
% LocalWords:  struct pgid sid pid ps getpgid getpgrp unistd void errno int
% LocalWords:  ESRCH getsid glibc system call XOPEN SOURCE EPERM setpgrp EACCES
% LocalWords:  setpgid exec EINVAL did fork race condition setsid tty ioctl
% LocalWords:  NOCTTY TIOCSCTTY error tcsetpgrp termios fd pgrpid descriptor VT
% LocalWords:  ENOTTY ENOSYS EBADF SIGTTIN SIGTTOU EIO tcgetpgrp crypt SIGTSTP
% LocalWords:  SIGINT SIGQUIT SIGTERM SIGHUP hungup kill orphaned SIGCONT exit
% LocalWords:  init Slackware run level inittab fig device getty TERM at execle
% LocalWords:  getpwnam chdir home chown chmod setuid setgid initgroups SIGCHLD
% LocalWords:  daemon like daemons NdT Stevens Programming FAQ filesystem umask
% LocalWords:  noclose syslog syslogd socket UDP klogd printk printf facility
% LocalWords:  priority log openlog const char ident option argv tab AUTH CRON
% LocalWords:  AUTHPRIV cron FTP KERN LOCAL LPR NEWS news USENET UUCP CONS CRIT
% LocalWords:  NDELAY NOWAIT ODELAY PERROR stderr format strerror EMERG ALERT
% LocalWords:  ERR WARNING NOTICE INFO DEBUG debug setlogmask mask UPTO za ssh
% LocalWords:  teletype telnet read write BELL beep CANON isatty desc ttyname
% LocalWords:  NULL ctermid stdio pathname buff size len ERANGE bits ispeed xFF
% LocalWords:  ospeed line tcflag INPCK IGNPAR PARMRK ISTRIP IGNBRK BREAK NUL
% LocalWords:  BRKINT IGNCR carriage return newline ICRNL INLCR IUCLC IXON NL
% LocalWords:  IXANY IXOFF IMAXBEL iflag OPOST CR OCRNL OLCUC ONLCR ONOCR OFILL
% LocalWords:  ONLRET OFDEL NLDLY CRDLY TABDLY BSDLY backspace BS VTDLY FFDLY
% LocalWords:  form feed FF oflag CLOCAL of HUPCL CREAD CSTOPB PARENB TIOCGPGRP
% LocalWords:  PARODD CSIZE CS CBAUD CBAUDEX CIBAUD CRTSCTS RTS CTS cflag ECHO
% LocalWords:  ICANON ECHOE ERASE ECHOPRT ECHOK ECHOKE ECHONL ECHOCTL ctrl ISIG
% LocalWords:  INTR QUIT SUSP IEXTEN EOL LNEXT REPRINT WERASE NOFLSH and TOSTOP
% LocalWords:  SIGSUSP XCASE DEFECHO FLUSHO DISCARD PENDIN lflag NCCS VINTR EOF
% LocalWords:  interrupt VQUIT VERASE VKILL VEOF VTIME VMIN VSWTC switch VSTART
% LocalWords:  VSTOP VSUSP VEOL VREPRINT VDISCARD VWERASE VLNEXT escape actions
% LocalWords:  tcgetattr tcsetattr EINTR TCSANOW TCSADRAIN TCSAFLUSH speed MIN
% LocalWords:  SetTermAttr UnSetTermAttr cfsetispeed cfsetospeed cfgetispeed ng
% LocalWords:  cfgetospeed quest'ultime tcsendbreak duration break tcdrain list
% LocalWords:  tcflush queue TCIFLUSH TCOFLUSH TCIOFLUSH tcflow action TCOOFF
% LocalWords:  TCOON TCIOFF TCION timer openpty Window nochdir embedded router
% LocalWords:  access point upstart systemd rsyslog vsyslog variadic src linux
% LocalWords:  closelog dmesg sysctl klogctl sys ERESTARTSYS ConsoleKit to CoPy
% LocalWords:  loglevel message libc klog mydmesg CAP ADMIN LXC pipelining UID
% LocalWords:  TIOCSPGRP GID IUTF UTF LOBLK NONBLOCK CMSPAR MARK VDSUSP VSTATUS
% LocalWords:  cfsetspeed raw cfmakeraw

%%% Local Variables: 
%%% mode: latex
%%% TeX-master: "gapil"
%%% End: 

\chapter{La gestione avanzata dei processi}
\label{cha:proc_advanced}

In questo capitolo affronteremo gli argomenti relativi alla gestione avanzata
dei processi. Inizieremo con le funzioni che attengono alla gestione avanzata
della sicurezza, passando poi a quelle relative all'analisi ed al controllo
dell'esecuzione, e alle funzioni per le modalità avanzate di creazione dei
processi e l'uso dei cosiddetti \textit{namespace}. Infine affronteremo le
\textit{sytem call} attinenti ad una serie di funzionalità specialistiche come
la gestione della virgola mobile, le porte di I/O ecc.

\section{La gestione avanzata della sicurezza}
\label{sec:process_security}

Tratteremo in questa sezione le funzionalità più avanzate relative alla
gestione della sicurezza ed il controllo degli accessi all'interno dei
processi, a partire dalle \textit{capabilities} e dalle funzionalità del
cosiddetto \textit{Secure Computing}. Esamineremo inoltre le altre
funzionalità relative alla sicurezza come la gestione delle chiavi
crittografiche e varie estensioni e funzionalità disponibili su questo
argomento.


\subsection{La gestione delle \textit{capabilities}}
\label{sec:proc_capabilities}

\itindbeg{capabilities} 

Come accennato in sez.~\ref{sec:proc_access_id} l'architettura classica della
gestione dei privilegi in un sistema unix-like ha il sostanziale problema di
fornire all'amministratore dei poteri troppo ampi, il che comporta che anche
quando si siano predisposte delle misure di protezione per in essere in grado
di difendersi dagli effetti di una eventuale compromissione del sistema (come
montare un filesystem in sola lettura per impedirne modifiche, o marcare un
file come immutabile) una volta che questa sia stata effettuata e si siano
ottenuti i privilegi di amministratore, queste misure potranno essere comunque
rimosse (nei casi elencati nella precedente nota si potrà sempre rimontare il
sistema in lettura-scrittura, o togliere l'attributo di immutabilità).

Il problema consiste nel fatto che nell'architettura tradizionale di un
sistema unix-like i controlli di accesso sono basati su un solo livello di
separazione: per i processi normali essi sono posti in atto, mentre per i
processi con i privilegi di amministratore essi non vengono neppure eseguiti.
Per questo motivo non era previsto alcun modo per evitare che un processo con
diritti di amministratore non potesse eseguire certe operazioni, o per cedere
definitivamente alcuni privilegi da un certo momento in poi. 

Per risolvere questo problema sono possibili varie soluzioni, ad esempio dai
kernel della serie 2.5 è stata introdotta la struttura dei
\itindex{Linux~Security~Modules~(LSM)} \textit{Linux Security Modules} che han
permesso di aggiungere varie forme di \itindex{Mandatory~Access~Control~(DAC)}
\textit{Mandatory Access Control} (MAC), in cui si potessero parcellizzare e
controllare nei minimi dettagli tutti i privilegi e le modalità in cui questi
possono essere usati dai programmi e trasferiti agli utenti, con la creazione
di varie estensioni (come \textit{SELinux}, \textit{Smack}, \textit{Tomoyo},
\textit{AppArmor}) che consentono di superare l'architettura tradizionale dei
permessi basati sul modello classico del controllo di accesso chiamato
\itindex{Discrectionary~Access~Control~(DAC)} \textit{Discrectionary Access
  Control} (DAC).

Ma già in precedenza, a partire dai kernel della serie 2.2, era stato
introdotto un meccanismo, detto \textit{capabilities}, per consentire di
suddividere i vari privilegi tradizionalmente associati all'amministratore in
un insieme di \textsl{capacità} distinte.  L'idea era che queste capacità
potessero essere abilitate e disabilitate in maniera indipendente per ciascun
processo con privilegi di amministratore, permettendo così una granularità
molto più fine nella distribuzione degli stessi, che evitasse la situazione
originaria di ``\textsl{tutto o nulla}''.

\itindbeg{file~capabilities}

Il meccanismo completo delle \textit{capabilities} (l'implementazione si rifà
ad una bozza di quello che doveva diventare lo standard POSIX.1e, poi
abbandonato) prevede inoltre la possibilità di associare le stesse ai singoli
file eseguibili, in modo da poter stabilire quali capacità possono essere
utilizzate quando viene messo in esecuzione uno specifico programma; ma il
supporto per questa funzionalità, chiamata \textit{file capabilities}, è stato
introdotto soltanto a partire dal kernel 2.6.24. Fino ad allora doveva essere
il programma stesso ad eseguire una riduzione esplicita delle sue capacità,
cosa che ha reso l'uso di questa funzionalità poco diffuso, vista la presenza
di meccanismi alternativi per ottenere limitazioni delle capacità
dell'amministratore a livello di sistema operativo, come \textit{SELinux}.

Con questo supporto e con le ulteriori modifiche introdotte con il kernel
2.6.25 il meccanismo delle \textit{capabilities} è stato totalmente
rivoluzionato, rendendolo più aderente alle intenzioni originali dello
standard POSIX, rimuovendo il significato che fino ad allora aveva avuto la
capacità \const{CAP\_SETPCAP}, e cambiando le modalità di funzionamento del
cosiddetto \textit{capabilities bounding set}. Ulteriori modifiche sono state
apportate con il kernel 2.6.26 per consentire la rimozione non ripristinabile
dei privilegi di amministratore. Questo fa sì che il significato ed il
comportamento del kernel finisca per dipendere dalla versione dello stesso e
dal fatto che le nuove \textit{file capabilities} siano abilitate o meno. Per
capire meglio la situazione e cosa è cambiato conviene allora spiegare con
maggiori dettagli come funziona il meccanismo delle \textit{capabilities}.

Il primo passo per frazionare i privilegi garantiti all'amministratore,
supportato fin dalla introduzione iniziale del kernel 2.2, è stato quello in
cui a ciascun processo sono stati associati tre distinti insiemi di
\textit{capabilities}, denominati rispettivamente \textit{permitted},
\textit{inheritable} ed \textit{effective}. Questi insiemi vengono mantenuti
in forma di tre diverse maschere binarie,\footnote{il kernel li mantiene, come
  i vari identificatori di sez.~\ref{sec:proc_setuid}, all'interno della
  \texttt{task\_struct} di ciascun processo (vedi
  fig.~\ref{fig:proc_task_struct}), nei tre campi \texttt{cap\_effective},
  \texttt{cap\_inheritable}, \texttt{cap\_permitted} del tipo
  \texttt{kernel\_cap\_t}; questo era, fino al kernel 2.6.25 definito come
  intero a 32 bit per un massimo di 32 \textit{capabilities} distinte,
  attualmente è stato aggiornato ad un vettore in grado di mantenerne fino a
  64.} in cui ciascun bit corrisponde ad una capacità diversa.

L'utilizzo di tre distinti insiemi serve a fornire una interfaccia flessibile
per l'uso delle \textit{capabilities}, con scopi analoghi a quelli per cui
sono mantenuti i diversi insiemi di identificatori di
sez.~\ref{sec:proc_setuid}; il loro significato, che è rimasto sostanzialmente
lo stesso anche dopo le modifiche seguite alla introduzione delle
\textit{file capabilities} è il seguente:
\begin{basedescript}{\desclabelwidth{2.1cm}\desclabelstyle{\nextlinelabel}}
\item[\textit{permitted}] l'insieme delle \textit{capabilities}
  ``\textsl{permesse}'', cioè l'insieme di quelle capacità che un processo
  \textsl{può} impostare come \textsl{effettive} o come
  \textsl{ereditabili}. Se un processo cancella una capacità da questo insieme
  non potrà più riassumerla.\footnote{questo nei casi ordinari, sono
    previste però una serie di eccezioni, dipendenti anche dal tipo di
    supporto, che vedremo meglio in seguito dato il notevole intreccio nella
    casistica.}
\item[\textit{inheritable}] l'insieme delle \textit{capabilities}
  ``\textsl{ereditabili}'', cioè di quelle che verranno trasmesse come insieme
  delle \textsl{permesse} ad un nuovo programma eseguito attraverso una
  chiamata ad \func{exec}.
\item[\textit{effective}] l'insieme delle \textit{capabilities}
  ``\textsl{effettive}'', cioè di quelle che vengono effettivamente usate dal
  kernel quando deve eseguire il controllo di accesso per le varie operazioni
  compiute dal processo.
\label{sec:capabilities_set}
\end{basedescript}

Con l'introduzione delle \textit{file capabilities} sono stati introdotti
altri tre insiemi associabili a ciascun file.\footnote{la realizzazione viene
  eseguita con l'uso di uno specifico attributo esteso,
  \texttt{security.capability}, la cui modifica è riservata, (come illustrato
  in sez.~\ref{sec:file_xattr}) ai processi dotato della capacità
  \const{CAP\_SYS\_ADMIN}.} Le \textit{file capabilities} hanno effetto
soltanto quando il file che le porta viene eseguito come programma con una
\func{exec}, e forniscono un meccanismo che consente l'esecuzione dello stesso
con maggiori privilegi; in sostanza sono una sorta di estensione del
\acr{suid} bit limitato ai privilegi di amministratore. Anche questi tre
insiemi sono identificati con gli stessi nomi di quello dei processi, ma il
loro significato è diverso:
\begin{basedescript}{\desclabelwidth{2.1cm}\desclabelstyle{\nextlinelabel}}
\item[\textit{permitted}] (chiamato originariamente \textit{forced}) l'insieme
  delle capacità che con l'esecuzione del programma verranno aggiunte alle
  capacità \textsl{permesse} del processo.
\item[\textit{inheritable}] (chiamato originariamente \textit{allowed})
  l'insieme delle capacità che con l'esecuzione del programma possono essere
  ereditate dal processo originario (che cioè non vengono tolte
  dall'\textit{inheritable set} del processo originale all'esecuzione di
  \func{exec}).
\item[\textit{effective}] in questo caso non si tratta di un insieme ma di un
  unico valore logico; se attivo all'esecuzione del programma tutte le
  capacità che risulterebbero \textsl{permesse} verranno pure attivate,
  inserendole automaticamente nelle \textsl{effettive}, se disattivato nessuna
  capacità verrà attivata (cioè l'\textit{effective set} resterà vuoto).
\end{basedescript}

\itindbeg{capabilities~bounding~set}

Infine come accennato, esiste un ulteriore insieme, chiamato
\textit{capabilities bounding set}, il cui scopo è quello di costituire un
limite alle capacità che possono essere attivate per un programma. Il suo
funzionamento però è stato notevolmente modificato con l'introduzione delle
\textit{file capabilities} e si deve pertanto prendere in considerazione una
casistica assai complessa.

Per i kernel fino al 2.6.25, o se non si attiva il supporto per le
\textit{file capabilities}, il \textit{capabilities bounding set} è un
parametro generale di sistema, il cui valore viene riportato nel file
\sysctlfiled{kernel/cap-bound}. Il suo valore iniziale è definito in sede di
compilazione del kernel, e da sempre ha previsto come default la presenza di
tutte le \textit{capabilities} eccetto \const{CAP\_SETPCAP}. In questa
situazione solo il primo processo eseguito nel sistema (quello con
\textsl{pid} 1, di norma \texttt{/sbin/init}) ha la possibilità di
modificarlo; ogni processo eseguito successivamente, se dotato dei privilegi
di amministratore, è in grado soltanto di rimuovere una delle
\textit{capabilities} già presenti dell'insieme.\footnote{per essere precisi
  occorre la capacità \const{CAP\_SYS\_MODULE}.}

In questo caso l'effetto complessivo del \textit{capabilities bounding set} è
che solo le capacità in esso presenti possono essere trasmesse ad un altro
programma attraverso una \func{exec}. Questo in sostanza significa che se un
qualunque programma elimina da esso una capacità, considerato che
\texttt{init} (almeno nelle versioni ordinarie) non supporta la reimpostazione
del \textit{bounding set}, questa non sarà più disponibile per nessun processo
a meno di un riavvio, eliminando così in forma definitiva quella capacità per
tutti, compreso l'amministratore.\footnote{la qual cosa, visto il default
  usato per il \textit{capabilities bounding set}, significa anche che
  \const{CAP\_SETPCAP} non è stata praticamente mai usata nella sua forma
  originale.}

Con il kernel 2.6.25 e le \textit{file capabilities} il \textit{bounding set}
è diventato una proprietà di ciascun processo, che viene propagata invariata
sia attraverso una \func{fork} che una \func{exec}. In questo caso il file
\sysctlfile{kernel/cap-bound} non esiste e \texttt{init} non ha nessun
ruolo speciale, inoltre in questo caso all'avvio il valore iniziale prevede la
presenza di tutte le capacità (compresa \const{CAP\_SETPCAP}). 

Con questo nuovo meccanismo il \textit{bounding set} continua a ricoprire un
ruolo analogo al precedente nel passaggio attraverso una \func{exec}, come
limite alle capacità che possono essere aggiunte al processo in quanto
presenti nel \textit{permitted set} del programma messo in esecuzione, in
sostanza il nuovo programma eseguito potrà ricevere una capacità presente nel
suo \textit{permitted set} (quello del file) solo se questa è anche nel
\textit{bounding set} (del processo). In questo modo si possono rimuovere
definitivamente certe capacità da un processo, anche qualora questo dovesse
eseguire un programma privilegiato che prevede di riassegnarle.

Si tenga presente però che in questo caso il \textit{bounding set} blocca
esclusivamente le capacità indicate nel \textit{permitted set} del programma
che verrebbero attivate in caso di esecuzione, e non quelle eventualmente già
presenti nell'\textit{inheritable set} del processo (ad esempio perché
presenti prima di averle rimosse dal \textit{bounding set}). In questo caso
eseguendo un programma che abbia anche lui dette capacità nel suo
\textit{inheritable set} queste verrebbero assegnate.

In questa seconda versione inoltre il \textit{bounding set} costituisce anche
un limite per le capacità che possono essere aggiunte all'\textit{inheritable
  set} del processo stesso con \func{capset}, sempre nel senso che queste
devono essere presenti nel \textit{bounding set} oltre che nel
\textit{permitted set} del processo. Questo limite vale anche per processi con
i privilegi di amministratore,\footnote{si tratta sempre di avere la
  \textit{capability} \const{CAP\_SETPCAP}.} per i quali invece non vale la
condizione che le \textit{capabilities} da aggiungere nell'\textit{inheritable
  set} debbano essere presenti nel proprio \textit{permitted set}.\footnote{lo
  scopo anche in questo caso è ottenere una rimozione definitiva della
  possibilità di passare una capacità rimossa dal \textit{bounding set}.}

Come si può notare per fare ricorso alle \textit{capabilities} occorre
comunque farsi carico di una notevole complessità di gestione, aggravata dalla
presenza di una radicale modifica del loro funzionamento con l'introduzione
delle \textit{file capabilities}. Considerato che il meccanismo originale era
incompleto e decisamente problematico nel caso di programmi che non ne
sapessero tener conto,\footnote{il problema di sicurezza originante da questa
  caratteristica venne alla ribalta con \texttt{sendmail}, in cui, riuscendo a
  rimuovere \const{CAP\_SETGID} dall'\textit{inheritable set} di un processo,
  si ottenne di far fallire \func{setuid} in maniera inaspettata per il
  programma (che aspettandosi sempre il successo della funzione non ne
  controllava lo stato di uscita) con la conseguenza di fargli fare come
  amministratore operazioni che altrimenti sarebbero state eseguite, senza
  poter apportare danni, da utente normale.}  ci soffermeremo solo sulla
implementazione completa presente a partire dal kernel 2.6.25, tralasciando
ulteriori dettagli riguardo la versione precedente.

Riassumendo le regole finora illustrate tutte le \textit{capabilities} vengono
ereditate senza modifiche attraverso una \func{fork} mentre, indicati con
\texttt{orig\_*} i valori degli insiemi del processo chiamante, con
\texttt{file\_*} quelli del file eseguito e con \texttt{bound\_set} il
\textit{capabilities bounding set}, dopo l'invocazione di \func{exec} il
processo otterrà dei nuovi insiemi di capacità \texttt{new\_*} secondo la
formula espressa dal seguente pseudo-codice C:

\includecodesnip{listati/cap-results.c}

% \begin{figure}[!htbp]
%   \footnotesize \centering
%   \begin{minipage}[c]{12cm}
%     \includecodesnip{listati/cap-results.c}
%   \end{minipage}
%   \caption{Espressione della modifica delle \textit{capabilities} attraverso
%     una \func{exec}.}
%   \label{fig:cap_across_exec}
% \end{figure}

\noindent e si noti come in particolare il \textit{capabilities bounding set}
non venga comunque modificato e resti lo stesso sia attraverso una \func{fork}
che attraverso una \func{exec}.


\itindend{capabilities~bounding~set}

A queste regole se ne aggiungono delle altre che servono a riprodurre il
comportamento tradizionale di un sistema unix-like in tutta una serie di
circostanze. La prima di queste è relativa a quello che avviene quando si
esegue un file senza \textit{capabilities}; se infatti si considerasse questo
equivalente al non averne assegnata alcuna, non essendo presenti capacità né
nel \textit{permitted set} né nell'\textit{inheritable set} del file,
nell'esecuzione di un qualunque programma l'amministratore perderebbe tutti i
privilegi originali dal processo.

Per questo motivo se un programma senza \textit{capabilities} assegnate viene
eseguito da un processo con \ids{UID} reale 0, esso verrà trattato come
se tanto il \textit{permitted set} che l'\textit{inheritable set} fossero con
tutte le \textit{capabilities} abilitate, con l'\textit{effective set} attivo,
col risultato di fornire comunque al processo tutte le capacità presenti nel
proprio \textit{bounding set}. Lo stesso avviene quando l'eseguibile ha attivo
il \acr{suid} bit ed appartiene all'amministratore, in entrambi i casi si
riesce così a riottenere il comportamento classico di un sistema unix-like.

Una seconda circostanza è quella relativa a cosa succede alle
\textit{capabilities} di un processo nelle possibili transizioni da \ids{UID}
nullo a \ids{UID} non nullo o viceversa (corrispondenti rispettivamente a
cedere o riottenere i privilegi di amministratore) che si possono effettuare
con le varie funzioni viste in sez.~\ref{sec:proc_setuid}. In questo caso la
casistica è di nuovo alquanto complessa, considerata anche la presenza dei
diversi gruppi di identificatori illustrati in tab.~\ref{tab:proc_uid_gid}, si
avrà allora che:
\begin{enumerate*}
\item se si passa da \ids{UID} effettivo nullo a non nullo
  l'\textit{effective set} del processo viene totalmente azzerato, se
  viceversa si passa da \ids{UID} effettivo non nullo a nullo il
  \textit{permitted set} viene copiato nell'\textit{effective set};
\item se si passa da \textit{file system} \ids{UID} nullo a non nullo verranno
  cancellate dall'\textit{effective set} del processo tutte le capacità
  attinenti i file, e cioè \const{CAP\_LINUX\_IMMUTABLE}, \const{CAP\_MKNOD},
  \const{CAP\_DAC\_OVERRIDE}, \const{CAP\_DAC\_READ\_SEARCH},
  \const{CAP\_MAC\_OVERRIDE}, \const{CAP\_CHOWN}, \const{CAP\_FSETID} e
  \const{CAP\_FOWNER} (le prime due a partire dal kernel 2.2.30), nella
  transizione inversa verranno invece inserite nell'\textit{effective set}
  quelle capacità della precedente lista che sono presenti nel suo
  \textit{permitted set}.
\item se come risultato di una transizione riguardante gli identificativi dei
  gruppi \textit{real}, \textit{saved} ed \textit{effective} in cui si passa
  da una situazione in cui uno di questi era nullo ad una in cui sono tutti
  non nulli,\footnote{in sostanza questo è il caso di quando si chiama
    \func{setuid} per rimuovere definitivamente i privilegi di amministratore
    da un processo.} verranno azzerati completamente sia il \textit{permitted
    set} che l'\textit{effective set}.
\end{enumerate*}
\label{sec:capability-uid-transition}

La combinazione di tutte queste regole consente di riprodurre il comportamento
ordinario di un sistema di tipo Unix tradizionale, ma può risultare
problematica qualora si voglia passare ad una configurazione di sistema
totalmente basata sull'applicazione delle \textit{capabilities}; in tal caso
infatti basta ad esempio eseguire un programma con \acr{suid} bit di proprietà
dell'amministratore per far riottenere ad un processo tutte le capacità
presenti nel suo \textit{bounding set}, anche se si era avuta la cura di
cancellarle dal \textit{permitted set}.

\itindbeg{securebits}

Per questo motivo a partire dal kernel 2.6.26, se le \textit{file
  capabilities} sono abilitate, ad ogni processo viene stata associata una
ulteriore maschera binaria, chiamata \textit{securebits flags}, su cui sono
mantenuti una serie di flag (vedi tab.~\ref{tab:securebits_values}) il cui
valore consente di modificare queste regole speciali che si applicano ai
processi con \ids{UID} nullo. La maschera viene sempre mantenuta
attraverso una \func{fork}, mentre attraverso una \func{exec} viene sempre
cancellato il flag \const{SECURE\_KEEP\_CAPS}.

\begin{table}[htb]
  \centering
  \footnotesize
  \begin{tabular}{|l|p{10cm}|}
    \hline
    \textbf{Flag} & \textbf{Descrizione} \\
    \hline
    \hline
    \constd{SECURE\_KEEP\_CAPS}&Il processo non subisce la cancellazione delle
                                sue \textit{capabilities} quando tutti i suoi
                                \ids{UID} passano ad un valore non
                                nullo (regola di compatibilità per il cambio
                                di \ids{UID} n.~3 del precedente
                                elenco), sostituisce il precedente uso
                                dell'operazione \const{PR\_SET\_KEEPCAPS} di
                                \func{prctl}.\\
    \constd{SECURE\_NO\_SETUID\_FIXUP}&Il processo non subisce le modifiche
                                delle sue \textit{capabilities} nel passaggio
                                da nullo a non nullo degli \ids{UID}
                                dei gruppi \textit{effective} e
                                \textit{file system} (regole di compatibilità
                                per il cambio di \ids{UID} nn.~1 e 2 del
                                precedente elenco).\\
    \constd{SECURE\_NOROOT}   & Il processo non assume nessuna capacità
                                aggiuntiva quando esegue un programma, anche
                                se ha \ids{UID} nullo o il programma ha
                                il \acr{suid} bit attivo ed appartiene
                                all'amministratore (regola di compatibilità
                                per l'esecuzione di programmi senza
                                \textit{capabilities}).\\
    \hline
  \end{tabular}
  \caption{Costanti identificative dei flag che compongono la maschera dei
    \textit{securebits}.}  
  \label{tab:securebits_values}
\end{table}

A ciascuno dei flag di tab.~\ref{tab:securebits_values} è inoltre abbinato un
corrispondente flag di blocco, identificato da una costante omonima con
l'estensione \texttt{\_LOCKED}, la cui attivazione è irreversibile ed ha
l'effetto di rendere permanente l'impostazione corrente del corrispondente
flag ordinario; in sostanza con \constd{SECURE\_KEEP\_CAPS\_LOCKED} si rende
non più modificabile \const{SECURE\_KEEP\_CAPS}, ed analogamente avviene con
\constd{SECURE\_NO\_SETUID\_FIXUP\_LOCKED} per
\const{SECURE\_NO\_SETUID\_FIXUP} e con \constd{SECURE\_NOROOT\_LOCKED} per
\const{SECURE\_NOROOT}.

Per l'impostazione di questi flag sono state predisposte due specifiche
operazioni di \func{prctl} (vedi sez.~\ref{sec:process_prctl}),
\const{PR\_GET\_SECUREBITS}, che consente di ottenerne il valore, e
\const{PR\_SET\_SECUREBITS}, che consente di modificarne il valore; per
quest'ultima sono comunque necessari i privilegi di amministratore ed in
particolare la capacità \const{CAP\_SETPCAP}. Prima dell'introduzione dei
\textit{securebits} era comunque possibile ottenere lo stesso effetto di
\const{SECURE\_KEEP\_CAPS} attraverso l'uso di un'altra operazione di
\func{prctl}, \const{PR\_SET\_KEEPCAPS}.

\itindend{securebits}

Oltre alla gestione dei \textit{securebits} la nuova versione delle
\textit{file capabilities} prevede l'uso di \func{prctl} anche per la gestione
del \textit{capabilities bounding set}, attraverso altre due operazioni
dedicate, \const{PR\_CAPBSET\_READ} per controllarne il valore e
\const{PR\_CAPBSET\_DROP} per modificarlo; quest'ultima di nuovo è una
operazione privilegiata che richiede la capacità \const{CAP\_SETPCAP} e che,
come indica chiaramente il nome, permette solo la rimozione di una
\textit{capability} dall'insieme; per i dettagli sull'uso di tutte queste
operazioni si rimanda alla rilettura di sez.~\ref{sec:process_prctl}.

\itindend{file~capabilities}


% NOTE per dati relativi al process capability bounding set, vedi:
% http://git.kernel.org/git/?p=linux/kernel/git/torvalds/linux-2.6.git;a=commit;h=3b7391de67da515c91f48aa371de77cb6cc5c07e

% NOTE riferimenti ai vari cambiamenti vedi:
% http://lwn.net/Articles/280279/  
% http://lwn.net/Articles/256519/
% http://lwn.net/Articles/211883/


Un elenco delle \textit{capabilities} disponibili su Linux, con una breve
descrizione ed il nome delle costanti che le identificano, è riportato in
tab.~\ref{tab:proc_capabilities};\footnote{l'elenco presentato questa tabella,
  ripreso dalla pagina di manuale (accessibile con \texttt{man capabilities})
  e dalle definizioni in \texttt{include/linux/capabilities.h}, è aggiornato
  al kernel 3.2.} la tabella è divisa in due parti, la prima riporta le
\textit{capabilities} previste anche nella bozza dello standard POSIX1.e, la
seconda quelle specifiche di Linux.  Come si può notare dalla tabella alcune
\textit{capabilities} attengono a singole funzionalità e sono molto
specializzate, mentre altre hanno un campo di applicazione molto vasto, che è
opportuno dettagliare maggiormente.

\begin{table}[!hbtp]
  \centering
  \footnotesize
  \begin{tabular}{|l|p{10cm}|}
    \hline
    \textbf{Capacità}&\textbf{Descrizione}\\
    \hline
    \hline
%
% POSIX-draft defined capabilities.
%
    \constd{CAP\_AUDIT\_CONTROL}& Abilitare e disabilitare il
                              controllo dell'auditing (dal kernel 2.6.11).\\ 
    \constd{CAP\_AUDIT\_WRITE}&Scrivere dati nel giornale di
                              auditing del kernel (dal kernel 2.6.11).\\ 
    % TODO verificare questa roba dell'auditing
    \constd{CAP\_BLOCK\_SUSPEND}&Utilizzare funzionalità che possono bloccare 
                              la sospensione del sistema (dal kernel 3.5).\\ 
    \constd{CAP\_CHOWN}     & Cambiare proprietario e gruppo
                              proprietario di un file (vedi
                              sez.~\ref{sec:file_ownership_management}).\\
    \constd{CAP\_DAC\_OVERRIDE}& Evitare il controllo dei
                               permessi di lettura, scrittura ed esecuzione dei
                               file, (vedi sez.~\ref{sec:file_access_control}).\\ 
    \constd{CAP\_DAC\_READ\_SEARCH}& Evitare il controllo dei
                              permessi di lettura ed esecuzione per
                              le directory (vedi
                              sez.~\ref{sec:file_access_control}).\\
    \const{CAP\_FOWNER}     & Evitare il controllo della proprietà di un file
                              per tutte le operazioni privilegiate non coperte
                              dalle precedenti \const{CAP\_DAC\_OVERRIDE} e
                              \const{CAP\_DAC\_READ\_SEARCH}.\\
    \constd{CAP\_FSETID}    & Evitare la cancellazione automatica dei bit
                              \acr{suid} e \acr{sgid} quando un file
                              per i quali sono impostati viene modificato da
                              un processo senza questa capacità e la capacità
                              di impostare il bit \acr{sgid} su un file anche
                              quando questo è relativo ad un gruppo cui non si
                              appartiene (vedi
                              sez.~\ref{sec:file_perm_management}).\\ 
    \constd{CAP\_KILL}      & Mandare segnali a qualunque
                              processo (vedi sez.~\ref{sec:sig_kill_raise}).\\
    \constd{CAP\_SETFCAP}   & Impostare le \textit{capabilities} di un file
                              (dal kernel 2.6.24).\\ 
    \constd{CAP\_SETGID}    & Manipolare i group ID dei
                              processi, sia il principale che i supplementari,
                              (vedi sez.~\ref{sec:proc_setgroups}) che quelli
                              trasmessi tramite i socket \textit{unix domain}
                              (vedi sez.~\ref{sec:unix_socket}).\\
    \constd{CAP\_SETUID}    & Manipolare gli user ID del
                              processo (vedi sez.~\ref{sec:proc_setuid}) e di
                              trasmettere un user ID arbitrario nel passaggio
                              delle credenziali coi socket \textit{unix
                                domain} (vedi sez.~\ref{sec:unix_socket}).\\ 
%
% Linux specific capabilities
%
\hline
    \constd{CAP\_IPC\_LOCK} & Effettuare il \textit{memory locking} con le
                              funzioni \func{mlock}, \func{mlockall},
                              \func{shmctl}, \func{mmap} (vedi
                              sez.~\ref{sec:proc_mem_lock} e 
                              sez.~\ref{sec:file_memory_map}). \\ 
% TODO verificare l'interazione con SHM_HUGETLB
    \constd{CAP\_IPC\_OWNER}& Evitare il controllo dei permessi
                              per le operazioni sugli oggetti di
                              intercomunicazione fra processi (vedi
                              sez.~\ref{sec:ipc_sysv}).\\  
    \constd{CAP\_LEASE}     & Creare dei \textit{file lease} (vedi
                              sez.~\ref{sec:file_asyncronous_lease})
                              pur non essendo proprietari del file (dal kernel
                              2.4).\\ 
    \constd{CAP\_LINUX\_IMMUTABLE}& Impostare sui file gli attributi 
                             \textit{immutable} e \textit{append-only} (vedi
                             sez.~\ref{sec:file_perm_overview}) se
                             supportati.\\
    \constd{CAP\_MAC\_ADMIN}& Amministrare il \textit{Mandatory
                               Access Control} di \textit{Smack} (dal kernel
                              2.6.25).\\
    \constd{CAP\_MAC\_OVERRIDE}& Evitare il \textit{Mandatory
                               Access Control} di \textit{Smack} (dal kernel
                              2.6.25).\\   
    \constd{CAP\_MKNOD}     & Creare file di dispositivo con \func{mknod} (vedi
                              sez.~\ref{sec:file_mknod}) (dal kernel 2.4).\\ 
    \const{CAP\_NET\_ADMIN} & Eseguire alcune operazioni
                              privilegiate sulla rete.\\
    \constd{CAP\_NET\_BIND\_SERVICE}& Porsi in ascolto su porte riservate (vedi 
                              sez.~\ref{sec:TCP_func_bind}).\\ 
    \constd{CAP\_NET\_BROADCAST}& Consentire l'uso di socket in
                              \textit{broadcast} e \textit{multicast}.\\ 
    \constd{CAP\_NET\_RAW}  & Usare socket \texttt{RAW} e \texttt{PACKET}
                              (vedi sez.~\ref{sec:sock_type}).\\ 
    \const{CAP\_SETPCAP}    & Effettuare modifiche privilegiate alle
                              \textit{capabilities}.\\   
    \const{CAP\_SYS\_ADMIN} & Eseguire una serie di compiti amministrativi.\\
    \constd{CAP\_SYS\_BOOT} & Eseguire un riavvio del sistema (vedi
                              sez.~\ref{sec:sys_reboot}).\\ 
    \constd{CAP\_SYS\_CHROOT}& Eseguire la funzione \func{chroot} (vedi 
                              sez.~\ref{sec:file_chroot}).\\
    \constd{CAP\_SYS\_MODULE}& Caricare e rimuovere moduli del kernel.\\ 
    \const{CAP\_SYS\_NICE}  & Modificare le varie priorità dei processi (vedi 
                              sez.~\ref{sec:proc_priority}).\\
    \constd{CAP\_SYS\_PACCT}& Usare le funzioni di \textit{accounting} dei 
                              processi (vedi
                              sez.~\ref{sec:sys_bsd_accounting}).\\  
    \constd{CAP\_SYS\_PTRACE}& La capacità di tracciare qualunque processo con
                              \func{ptrace} (vedi 
                              sez.~\ref{sec:process_ptrace}).\\
    \constd{CAP\_SYS\_RAWIO}& Operare sulle porte di I/O con \func{ioperm} e
                               \func{iopl} (vedi
                              sez.~\ref{sec:process_io_port}).\\
    \const{CAP\_SYS\_RESOURCE}& Superare le varie limitazioni sulle risorse.\\ 
    \constd{CAP\_SYS\_TIME} & Modificare il tempo di sistema (vedi 
                              sez.~\ref{sec:sys_time}).\\ 
    \constd{CAP\_SYS\_TTY\_CONFIG}&Simulare un \textit{hangup} della console,
                              con la funzione \func{vhangup}.\\
    \constd{CAP\_SYSLOG}    & Gestire il buffer dei messaggi
                              del kernel, (vedi sez.~\ref{sec:sess_daemon}),
                              introdotta dal kernel 2.6.38 come capacità
                              separata da \const{CAP\_SYS\_ADMIN}.\\
    \constd{CAP\_WAKE\_ALARM}&Usare i timer di tipo
                              \const{CLOCK\_BOOTTIME\_ALARM} e
                              \const{CLOCK\_REALTIME\_ALARM}, vedi
                              sez.~\ref{sec:sig_timer_adv} (dal kernel 3.0).\\  
    \hline
  \end{tabular}
  \caption{Le costanti che identificano le \textit{capabilities} presenti nel
    kernel.}
\label{tab:proc_capabilities}
\end{table}

\constbeg{CAP\_SETPCAP}

Prima di dettagliare il significato della capacità più generiche, conviene
però dedicare un discorso a parte a \const{CAP\_SETPCAP}, il cui significato è
stato completamente cambiato con l'introduzione delle \textit{file
  capabilities} nel kernel 2.6.24. In precedenza questa capacità era quella
che permetteva al processo che la possedeva di impostare o rimuovere le
\textit{capabilities} presenti nel suo \textit{permitted set} su un qualunque
altro processo. In realtà questo non è mai stato l'uso inteso nelle bozze
dallo standard POSIX, ed inoltre, come si è già accennato, dato che questa
capacità è sempre stata assente (a meno di specifiche ricompilazioni del
kernel) nel \textit{capabilities bounding set} usato di default, essa non è
neanche mai stata realmente disponibile.

Con l'introduzione \textit{file capabilities} e il cambiamento del significato
del \textit{capabilities bounding set} la possibilità di modificare le
capacità di altri processi è stata completamente rimossa, e
\const{CAP\_SETPCAP} ha acquisito quello che avrebbe dovuto essere il suo
significato originario, e cioè la capacità del processo di poter inserire nel
suo \textit{inheritable set} qualunque capacità presente nel \textit{bounding
  set}. Oltre a questo la disponibilità di \const{CAP\_SETPCAP} consente ad un
processo di eliminare una capacità dal proprio \textit{bounding set} (con la
conseguente impossibilità successiva di eseguire programmi con quella
capacità), o di impostare i \textit{securebits} delle \textit{capabilities}.

\constend{CAP\_SETPCAP}
\constbeg{CAP\_FOWNER}

La prima fra le capacità ``\textsl{ampie}'' che occorre dettagliare
maggiormente è \const{CAP\_FOWNER}, che rimuove le restrizioni poste ad un
processo che non ha la proprietà di un file in un vasto campo di
operazioni;\footnote{vale a dire la richiesta che l'\ids{UID} effettivo del
  processo (o meglio l'\ids{UID} di filesystem, vedi
  sez.~\ref{sec:proc_setuid}) coincida con quello del proprietario.}  queste
comprendono i cambiamenti dei permessi e dei tempi del file (vedi
sez.~\ref{sec:file_perm_management} e sez.~\ref{sec:file_file_times}), le
impostazioni degli attributi dei file e delle ACL (vedi
sez.~\ref{sec:file_xattr} e \ref{sec:file_ACL}), poter ignorare lo
\textit{sticky bit} nella cancellazione dei file (vedi
sez.~\ref{sec:file_special_perm}), la possibilità di impostare il flag di
\const{O\_NOATIME} con \func{open} e \func{fcntl} (vedi
sez.~\ref{sec:file_open_close} e sez.~\ref{sec:file_fcntl_ioctl}) senza
restrizioni.

\constend{CAP\_FOWNER}
\constbeg{CAP\_NET\_ADMIN}

Una seconda capacità che copre diverse operazioni, in questo caso riguardanti
la rete, è \const{CAP\_NET\_ADMIN}, che consente di impostare le opzioni
privilegiate dei socket (vedi sez.~\ref{sec:sock_generic_options}), abilitare
il \textit{multicasting} (vedi sez.\ref{sec:sock_ipv4_options}), eseguire la
configurazione delle interfacce di rete (vedi
sez.~\ref{sec:sock_ioctl_netdevice}) ed impostare la tabella di instradamento.

\constend{CAP\_NET\_ADMIN}
\constbeg{CAP\_SYS\_ADMIN}

Una terza \textit{capability} con vasto campo di applicazione è
\const{CAP\_SYS\_ADMIN}, che copre una serie di operazioni amministrative,
come impostare le quote disco (vedi sez.\ref{sec:disk_quota}), attivare e
disattivare la \textit{swap}, montare, rimontare e smontare filesystem (vedi
sez.~\ref{sec:filesystem_mounting}), effettuare operazioni di controllo su
qualunque oggetto dell'IPC di SysV (vedi sez.~\ref{sec:ipc_sysv}), operare
sugli attributi estesi dei file di classe \texttt{security} o \texttt{trusted}
(vedi sez.~\ref{sec:file_xattr}), specificare un \ids{UID} arbitrario nella
trasmissione delle credenziali dei socket (vedi
sez.~\ref{sec:socket_credential_xxx}), assegnare classi privilegiate
(\const{IOPRIO\_CLASS\_RT} e prima del kernel 2.6.25 anche
\const{IOPRIO\_CLASS\_IDLE}) per lo scheduling dell'I/O (vedi
sez.~\ref{sec:io_priority}), superare il limite di sistema sul numero massimo
di file aperti,\footnote{quello indicato da \sysctlfiled{fs/file-max}.}
effettuare operazioni privilegiate sulle chiavi mantenute dal kernel (vedi
sez.~\ref{sec:keyctl_management}), usare la funzione \func{lookup\_dcookie},
usare \const{CLONE\_NEWNS} con \func{unshare} e \func{clone}, (vedi
sez.~\ref{sec:process_clone}).

\constend{CAP\_SYS\_ADMIN}
\constbeg{CAP\_SYS\_NICE}

Originariamente \const{CAP\_SYS\_NICE} riguardava soltanto la capacità di
aumentare le priorità di esecuzione dei processi, come la diminuzione del
valore di \textit{nice} (vedi sez.~\ref{sec:proc_sched_stand}), l'uso delle
priorità \textit{real-time} (vedi sez.~\ref{sec:proc_real_time}), o
l'impostazione delle affinità di processore (vedi
sez.~\ref{sec:proc_sched_multiprocess}); ma con l'introduzione di priorità
anche riguardo le operazioni di accesso al disco, e, nel caso di sistemi NUMA,
alla memoria, essa viene a coprire anche la possibilità di assegnare priorità
arbitrarie nell'accesso a disco (vedi sez.~\ref{sec:io_priority}) e nelle
politiche di allocazione delle pagine di memoria ai nodi di un sistema NUMA.

\constend{CAP\_SYS\_NICE}
\constbeg{CAP\_SYS\_RESOURCE}

Infine la \textit{capability} \const{CAP\_SYS\_RESOURCE} attiene alla
possibilità di superare i limiti imposti sulle risorse di sistema, come usare
lo spazio disco riservato all'amministratore sui filesystem che lo supportano,
usare la funzione \func{ioctl} per controllare il \textit{journaling} sul
filesystem \acr{ext3}, non subire le quote disco, aumentare i limiti sulle
risorse di un processo (vedi sez.~\ref{sec:sys_resource_limit}) e quelle sul
numero di processi, ed i limiti sulle dimensioni dei messaggi delle code del
SysV IPC (vedi sez.~\ref{sec:ipc_sysv_mq}).

\constend{CAP\_SYS\_RESOURCE}

Per la gestione delle \textit{capabilities} il kernel mette a disposizione due
funzioni che permettono rispettivamente di leggere ed impostare i valori dei
tre insiemi illustrati in precedenza. Queste due funzioni di sistema sono
\funcd{capget} e \funcd{capset} e costituiscono l'interfaccia di gestione
basso livello; i loro rispettivi prototipi sono:

\begin{funcproto}{
\fhead{sys/capability.h}
\fdecl{int capget(cap\_user\_header\_t hdrp, cap\_user\_data\_t datap)}
\fdesc{Legge le \textit{capabilities}.} 
\fdecl{int capset(cap\_user\_header\_t hdrp, const cap\_user\_data\_t datap)} 
\fdesc{Imposta le \textit{capabilities}.} 
}

{Le funzioni ritornano $0$ in caso di successo e $-1$ per un errore, nel qual
  caso \var{errno} assumerà uno dei valori: 
  \begin{errlist}
  \item[\errcode{EFAULT}] si è indicato un puntatore sbagliato o nullo
    per \param{hdrp} o \param{datap} (quest'ultimo può essere nullo solo se si
    usa \func{capget} per ottenere la versione delle \textit{capabilities}
    usata dal kernel).
  \item[\errcode{EINVAL}] si è specificato un valore non valido per uno dei
    campi di \param{hdrp}, in particolare una versione non valida della
    versione delle \textit{capabilities}.
  \item[\errcode{EPERM}] si è tentato di aggiungere una capacità nell'insieme
    delle \textit{capabilities} permesse, o di impostare una capacità non
    presente nell'insieme di quelle permesse negli insieme delle effettive o
    ereditate, o si è cercato di impostare una \textit{capability} di un altro
    processo senza avare \const{CAP\_SETPCAP}.
  \item[\errcode{ESRCH}] si è fatto riferimento ad un processo inesistente.
  \end{errlist}
}
\end{funcproto}

Queste due funzioni prendono come argomenti due tipi di dati dedicati,
definiti come puntatori a due strutture specifiche di Linux, illustrate in
fig.~\ref{fig:cap_kernel_struct}.  Per un certo periodo di tempo era anche
indicato che per poterle utilizzare fosse necessario che la macro
\macro{\_POSIX\_SOURCE} risultasse non definita (ed era richiesto di inserire
una istruzione \texttt{\#undef \_POSIX\_SOURCE} prima di includere
\headfiled{sys/capability.h}) requisito che non risulta più
presente.\footnote{e non è chiaro neanche quanto sia mai stato davvero
  necessario.}

\begin{figure}[!htb]
  \footnotesize
  \centering
  \begin{minipage}[c]{0.8\textwidth}
    \includestruct{listati/cap_user_header_t.h}
  \end{minipage}
  \normalsize 
  \caption{Definizione delle strutture a cui fanno riferimento i puntatori
    \structd{cap\_user\_header\_t} e \structd{cap\_user\_data\_t} usati per
    l'interfaccia di gestione di basso livello delle \textit{capabilities}.}
  \label{fig:cap_kernel_struct}
\end{figure}

Si tenga presente che le strutture di fig.~\ref{fig:cap_kernel_struct}, come i
prototipi delle due funzioni \func{capget} e \func{capset}, sono soggette ad
essere modificate con il cambiamento del kernel (in particolare i tipi di dati
delle strutture) ed anche se finora l'interfaccia è risultata stabile, non c'è
nessuna assicurazione che questa venga mantenuta,\footnote{viene però
  garantito che le vecchie funzioni continuino a funzionare.} Pertanto se si
vogliono scrivere programmi portabili che possano essere eseguiti senza
modifiche o adeguamenti su qualunque versione del kernel è opportuno
utilizzare le interfacce di alto livello che vedremo più avanti.

La struttura a cui deve puntare l'argomento \param{hdrp} serve ad indicare,
tramite il campo \var{pid}, il \ids{PID} del processo del quale si vogliono
leggere o modificare le \textit{capabilities}. Con \func{capset} questo, se si
usano le \textit{file capabilities}, può essere solo 0 o il \ids{PID} del
processo chiamante, che sono equivalenti. Non tratteremo, essendo comunque di
uso irrilevante, il caso in cui, in mancanza di tale supporto, la funzione può
essere usata per modificare le \textit{capabilities} di altri processi, per il
quale si rimanda, se interessati, alla lettura della pagina di manuale.

Il campo \var{version} deve essere impostato al valore della versione delle
stesse usata dal kernel (quello indicato da una delle costanti
\texttt{\_LINUX\_CAPABILITY\_VERSION\_n} di fig.~\ref{fig:cap_kernel_struct})
altrimenti le funzioni ritorneranno con un errore di \errcode{EINVAL},
restituendo nel campo stesso il valore corretto della versione in uso. La
versione due è comunque deprecata e non deve essere usata, ed il kernel
stamperà un avviso se lo si fa.

I valori delle \textit{capabilities} devono essere passati come maschere
binarie;\footnote{e si tenga presente che i valori di
  tab.~\ref{tab:proc_capabilities} non possono essere combinati direttamente,
  indicando il numero progressivo del bit associato alla relativa capacità.}
con l'introduzione delle \textit{capabilities} a 64 bit inoltre il
puntatore \param{datap} non può essere più considerato come relativo ad una
singola struttura, ma ad un vettore di due strutture.\footnote{è questo cambio
  di significato che ha portato a deprecare la versione 2, che con
  \func{capget} poteva portare ad un buffer overflow per vecchie applicazioni
  che continuavano a considerare \param{datap} come puntatore ad una singola
  struttura.}

Dato che le precedenti funzioni, oltre ad essere specifiche di Linux, non
garantiscono la stabilità nell'interfaccia, è sempre opportuno effettuare la
gestione delle \textit{capabilities} utilizzando le funzioni di libreria a
questo dedicate. Queste funzioni, che seguono quanto previsto nelle bozze
dello standard POSIX.1e, non fanno parte della \acr{glibc} e sono fornite in
una libreria a parte,\footnote{la libreria è \texttt{libcap2}, nel caso di
  Debian può essere installata con il pacchetto omonimo.} pertanto se un
programma le utilizza si dovrà indicare esplicitamente al compilatore l'uso
della suddetta libreria attraverso l'opzione \texttt{-lcap}.

\itindbeg{capability~state}

Le funzioni dell'interfaccia alle \textit{capabilities} definite nelle bozze
dello standard POSIX.1e prevedono l'uso di un tipo di dato opaco,
\typed{cap\_t}, come puntatore ai dati mantenuti nel cosiddetto
\textit{capability state},\footnote{si tratta in sostanza di un puntatore ad
  una struttura interna utilizzata dalle librerie, i cui campi non devono mai
  essere acceduti direttamente.} in sono memorizzati tutti i dati delle
\textit{capabilities}.

In questo modo è possibile mascherare i dettagli della gestione di basso
livello, che potranno essere modificati senza dover cambiare le funzioni
dell'interfaccia, che fanno riferimento soltanto ad oggetti di questo tipo.
L'interfaccia pertanto non soltanto fornisce le funzioni per modificare e
leggere le \textit{capabilities}, ma anche quelle per gestire i dati
attraverso i \textit{capability state}, che presentano notevoli affinità,
essendo parte di bozze dello stesso standard, con quelle già viste per le ACL.

La prima funzione dell'interfaccia è quella che permette di inizializzare un
\textit{capability state}, allocando al contempo la memoria necessaria per i
relativi dati. La funzione è \funcd{cap\_init} ed il suo prototipo è:

\begin{funcproto}{
\fhead{sys/capability.h}
\fdecl{cap\_t cap\_init(void)}
\fdesc{Crea ed inizializza un \textit{capability state}.} 
}

{La funzione ritorna un \textit{capability state} in caso di successo e
  \val{NULL} per un errore, nel qual caso \var{errno} potrà assumere solo il
  valore \errval{ENOMEM}.  }
\end{funcproto}

La funzione restituisce il puntatore \type{cap\_t} ad uno stato inizializzato
con tutte le \textit{capabilities} azzerate. In caso di errore (cioè quando
non c'è memoria sufficiente ad allocare i dati) viene restituito \val{NULL}
ed \var{errno} viene impostata a \errval{ENOMEM}.  

La memoria necessaria a mantenere i dati viene automaticamente allocata da
\func{cap\_init}, ma dovrà essere disallocata esplicitamente quando non è più
necessaria utilizzando, per questo l'interfaccia fornisce una apposita
funzione, \funcd{cap\_free}, il cui prototipo è:

\begin{funcproto}{
\fhead{sys/capability.h}
\fdecl{int cap\_free(void *obj\_d)}
\fdesc{Disalloca la memoria allocata per i dati delle \textit{capabilities}..} 
}

{La funzione ritorna $0$ in caso di successo e $-1$ per un errore, nel qual
  caso \var{errno} potrà assumere solo il valore \errval{EINVAL}.
}
\end{funcproto}


La funzione permette di liberare la memoria allocata dalle altre funzioni
della libreria sia per un \textit{capability state}, nel qual caso l'argomento
sarà un dato di tipo \type{cap\_t}, che per una descrizione testuale dello
stesso,\footnote{cioè quanto ottenuto tramite la funzione
  \func{cap\_to\_text}.} nel qual caso l'argomento sarà un dato di tipo
\texttt{char *}. Per questo motivo l'argomento \param{obj\_d} è dichiarato
come \texttt{void *}, per evitare la necessità di eseguire un \textit{cast},
ma dovrà comunque corrispondere ad un puntatore ottenuto tramite le altre
funzioni della libreria, altrimenti la funzione fallirà con un errore di
\errval{EINVAL}.

Infine si può creare una copia di un \textit{capability state} ottenuto in
precedenza tramite la funzione \funcd{cap\_dup}, il cui prototipo è:

\begin{funcproto}{
\fhead{sys/capability.h}
\fdecl{cap\_t cap\_dup(cap\_t cap\_p)}
\fdesc{Duplica un \textit{capability state} restituendone una copia.} 
}

{La funzione ritorna un \textit{capability state} in caso di successo e
  \val{NULL} per un errore, nel qual caso \var{errno} assumerà i valori
  \errval{ENOMEM} o \errval{EINVAL} nel loro significato generico.}
\end{funcproto}


La funzione crea una copia del \textit{capability state} posto all'indirizzo
\param{cap\_p} che si è passato come argomento, restituendo il puntatore alla
copia, che conterrà gli stessi valori delle \textit{capabilities} presenti
nell'originale. La memoria necessaria viene allocata automaticamente dalla
funzione. Una volta effettuata la copia i due \textit{capability state}
potranno essere modificati in maniera completamente indipendente, ed alla fine
delle operazioni si dovrà disallocare anche la copia, oltre all'originale.

Una seconda classe di funzioni di servizio previste dall'interfaccia sono
quelle per la gestione dei dati contenuti all'interno di un \textit{capability
  state}; la prima di queste è \funcd{cap\_clear}, il cui prototipo è:

\begin{funcproto}{
\fhead{sys/capability.h}
\fdecl{int cap\_clear(cap\_t cap\_p)}
\fdesc{Inizializza un \textit{capability state} cancellando tutte le
  \textit{capabilities}.}
}

{La funzione ritorna $0$ in caso di successo e $-1$ per un errore, nel qual
  caso \var{errno} potrà assumere solo il valore \errval{EINVAL}.
}
\end{funcproto}

La funzione si limita ad azzerare tutte le \textit{capabilities} presenti nel
\textit{capability state} all'indirizzo \param{cap\_p} passato come argomento,
restituendo uno stato \textsl{vuoto}, analogo a quello che si ottiene nella
creazione con \func{cap\_init}.

Una variante di \func{cap\_clear} è \funcd{cap\_clear\_flag} che cancella da
un \textit{capability state} tutte le \textit{capabilities} di un certo
insieme fra quelli elencati a pag.~\pageref{sec:capabilities_set}, il suo
prototipo è:

\begin{funcproto}{
\fhead{sys/capability.h}
\fdecl{int cap\_clear\_flag(cap\_t cap\_p, cap\_flag\_t flag)} 
\fdesc{Cancella delle \textit{capabilities} da un \textit{capability state}.} 
}

{La funzione ritorna $0$ in caso di successo e $-1$ per un errore, nel qual
  caso \var{errno}  potrà assumere solo il valore \errval{EINVAL}.
}
\end{funcproto}

La funzione richiede che si indichi quale degli insiemi si intente cancellare
da \param{cap\_p} con l'argomento \param{flag}. Questo deve essere specificato
con una variabile di tipo \type{cap\_flag\_t} che può assumere
esclusivamente\footnote{si tratta in effetti di un tipo enumerato, come si può
  verificare dalla sua definizione che si trova in
  \headfile{sys/capability.h}.} uno dei valori illustrati in
tab.~\ref{tab:cap_set_identifier}.

\begin{table}[htb]
  \centering
  \footnotesize
  \begin{tabular}[c]{|l|l|}
    \hline
    \textbf{Valore} & \textbf{Significato} \\
    \hline
    \hline
    \constd{CAP\_EFFECTIVE}  & Capacità dell'insieme \textsl{effettivo}.\\
    \constd{CAP\_PERMITTED}  & Capacità dell'insieme \textsl{permesso}.\\ 
    \constd{CAP\_INHERITABLE}& Capacità dell'insieme \textsl{ereditabile}.\\
    \hline
  \end{tabular}
  \caption{Valori possibili per il tipo di dato \typed{cap\_flag\_t} che
    identifica gli insiemi delle \textit{capabilities}.}
  \label{tab:cap_set_identifier}
\end{table}

Si possono inoltre confrontare in maniera diretta due diversi
\textit{capability state} con la funzione \funcd{cap\_compare}; il suo
prototipo è:

\begin{funcproto}{
\fhead{sys/capability.h}
\fdecl{int cap\_compare(cap\_t cap\_a, cap\_t cap\_b)}
\fdesc{Confronta due \textit{capability state}.} 
}

{La funzione ritorna $0$ se i \textit{capability state} sono identici
    ed un valore positivo se differiscono, non sono previsti errori.}
\end{funcproto}


La funzione esegue un confronto fra i due \textit{capability state} passati
come argomenti e ritorna in un valore intero il risultato, questo è nullo se
sono identici o positivo se vi sono delle differenze. Il valore di ritorno
della funzione consente inoltre di per ottenere ulteriori informazioni su
quali sono gli insiemi di \textit{capabilities} che risultano differenti.  Per
questo si può infatti usare la apposita macro \macro{CAP\_DIFFERS}:

{\centering
\vspace{3pt}
\begin{funcbox}{
\fhead{sys/capability.h}
\fdecl{int \macrod{CAP\_DIFFERS}(value, flag)}
\fdesc{Controlla lo stato di eventuali differenze delle \textit{capabilities}
  nell'insieme \texttt{flag}.}
}
\end{funcbox}
}

La macro richiede che si passi nell'argomento \texttt{value} il risultato
della funzione \func{cap\_compare} e in \texttt{flag} l'indicazione (coi
valori di tab.~\ref{tab:cap_set_identifier}) dell'insieme che si intende
controllare; restituirà un valore diverso da zero se le differenze rilevate da
\func{cap\_compare} sono presenti nell'insieme indicato.

Per la gestione dei singoli valori delle \textit{capabilities} presenti in un
\textit{capability state} l'interfaccia prevede due funzioni specifiche,
\funcd{cap\_get\_flag} e \funcd{cap\_set\_flag}, che permettono
rispettivamente di leggere o impostare il valore di una capacità all'interno
in uno dei tre insiemi già citati; i rispettivi prototipi sono:

\begin{funcproto}{
\fhead{sys/capability.h}
\fdecl{int cap\_get\_flag(cap\_t cap\_p, cap\_value\_t cap, cap\_flag\_t 
flag,\\
\phantom{int cap\_get\_flag(}cap\_flag\_value\_t *value\_p)}
\fdesc{Legge il valore di una \textit{capability}.}
\fdecl{int cap\_set\_flag(cap\_t cap\_p, cap\_flag\_t flag, int ncap,
  cap\_value\_t *caps, \\
\phantom{int cap\_set\_flag(}cap\_flag\_value\_t value)} 
\fdesc{Imposta il valore di una \textit{capability}.} 
}

{Le funzioni ritornano $0$ in caso di successo e $-1$ per un errore, nel qual
  caso \var{errno} potrà assumere solo il valore \errval{EINVAL}.  
}
\end{funcproto}

In entrambe le funzioni l'argomento \param{cap\_p} indica il puntatore al
\textit{capability state} su cui operare, mentre l'argomento \param{flag}
indica su quale dei tre insiemi si intende operare, sempre con i valori di
tab.~\ref{tab:cap_set_identifier}.  La capacità che si intende controllare o
impostare invece deve essere specificata attraverso una variabile di tipo
\typed{cap\_value\_t}, che può prendere come valore uno qualunque di quelli
riportati in tab.~\ref{tab:proc_capabilities}, in questo caso però non è
possibile combinare diversi valori in una maschera binaria, una variabile di
tipo \type{cap\_value\_t} può indicare una sola capacità.\footnote{in
  \headfile{sys/capability.h} il tipo \type{cap\_value\_t} è definito come
  \ctyp{int}, ma i valori validi sono soltanto quelli di
  tab.~\ref{tab:proc_capabilities}.}

Infine lo stato di una capacità è descritto ad una variabile di tipo
\type{cap\_flag\_value\_t}, che a sua volta può assumere soltanto
uno\footnote{anche questo è un tipo enumerato.} dei valori di
tab.~\ref{tab:cap_value_type}.

\begin{table}[htb]
  \centering
  \footnotesize
  \begin{tabular}[c]{|l|l|}
    \hline
    \textbf{Valore} & \textbf{Significato} \\
    \hline
    \hline
    \constd{CAP\_CLEAR}& La capacità non è impostata.\\ 
    \constd{CAP\_SET}  & La capacità è impostata.\\
    \hline
  \end{tabular}
  \caption{Valori possibili per il tipo di dato \typed{cap\_flag\_value\_t} che
    indica lo stato di una capacità.}
  \label{tab:cap_value_type}
\end{table}

La funzione \func{cap\_get\_flag} legge lo stato della capacità indicata
dall'argomento \param{cap} all'interno dell'insieme indicato dall'argomento
\param{flag} e lo restituisce come \textit{value result argument} nella
variabile puntata dall'argomento \param{value\_p}. Questa deve essere di tipo
\type{cap\_flag\_value\_t} ed assumerà uno dei valori di
tab.~\ref{tab:cap_value_type}. La funzione consente pertanto di leggere solo
lo stato di una capacità alla volta.

La funzione \func{cap\_set\_flag} può invece impostare in una sola chiamata
più \textit{capabilities}, anche se solo all'interno dello stesso insieme ed
allo stesso valore. Per questo motivo essa prende un vettore di valori di tipo
\type{cap\_value\_t} nell'argomento \param{caps}, la cui dimensione viene
specificata dall'argomento \param{ncap}. Il tipo di impostazione da eseguire
(cancellazione o attivazione) per le capacità elencate in \param{caps} viene
indicato dall'argomento \param{value} sempre con i valori di
tab.~\ref{tab:cap_value_type}.

Per semplificare la gestione delle \textit{capabilities} l'interfaccia prevede
che sia possibile utilizzare anche una rappresentazione testuale del contenuto
di un \textit{capability state} e fornisce le opportune funzioni di
gestione;\footnote{entrambe erano previste dalla bozza dello standard
  POSIX.1e.} la prima di queste, che consente di ottenere la rappresentazione
testuale, è \funcd{cap\_to\_text}, il cui prototipo è:

\begin{funcproto}{
\fhead{sys/capability.h}
\fdecl{char *cap\_to\_text(cap\_t caps, ssize\_t *length\_p)}
\fdesc{Genera una visualizzazione testuale delle \textit{capabilities}.} 
}

{La funzione ritorna un puntatore alla stringa con la descrizione delle
  \textit{capabilities} in caso di successo e \val{NULL} per un errore, nel
  qual caso \var{errno} assumerà i valori \errval{EINVAL} o \errval{ENOMEM}
  nel loro significato generico.}
\end{funcproto}

La funzione ritorna l'indirizzo di una stringa contente la descrizione
testuale del contenuto del \textit{capability state} \param{caps} passato come
argomento, e, qualora l'argomento \param{length\_p} sia diverso da \val{NULL},
restituisce come \textit{value result argument} nella variabile intera da
questo puntata la lunghezza della stringa. La stringa restituita viene
allocata automaticamente dalla funzione e pertanto dovrà essere liberata con
\func{cap\_free}.

La rappresentazione testuale, che viene usata anche dai programmi di gestione a
riga di comando, prevede che lo stato venga rappresentato con una stringa di
testo composta da una serie di proposizioni separate da spazi, ciascuna delle
quali specifica una operazione da eseguire per creare lo stato finale. Nella
rappresentazione si fa sempre conto di partire da uno stato in cui tutti gli
insiemi sono vuoti e si provvede a impostarne i contenuti.

Ciascuna proposizione è nella forma di un elenco di capacità, espresso con i
nomi di tab.~\ref{tab:proc_capabilities} separati da virgole, seguito da un
operatore, e dall'indicazione degli insiemi a cui l'operazione si applica. I
nomi delle capacità possono essere scritti sia maiuscoli che minuscoli, viene
inoltre riconosciuto il nome speciale \texttt{all} che è equivalente a
scrivere la lista completa. Gli insiemi sono identificati dalle tre lettere
iniziali: ``\texttt{p}'' per il \textit{permitted}, ``\texttt{i}'' per
l'\textit{inheritable} ed ``\texttt{e}'' per l'\textit{effective} che devono
essere sempre minuscole, e se ne può indicare più di uno.

Gli operatori possibili sono solo tre: ``\texttt{+}'' che aggiunge le capacità
elencate agli insiemi indicati, ``\texttt{-}'' che le toglie e ``\texttt{=}''
che le assegna esattamente. I primi due richiedono che sia sempre indicato sia
un elenco di capacità che gli insiemi a cui esse devono applicarsi, e
rispettivamente attiveranno o disattiveranno le capacità elencate nell'insieme
o negli insiemi specificati, ignorando tutto il resto. I due operatori possono
anche essere combinati nella stessa proposizione, per aggiungere e togliere le
capacità dell'elenco da insiemi diversi.

L'assegnazione si applica invece su tutti gli insiemi allo stesso tempo,
pertanto l'uso di ``\texttt{=}'' è equivalente alla cancellazione preventiva
di tutte le capacità ed alla impostazione di quelle elencate negli insiemi
specificati, questo significa che in genere lo si usa una sola volta
all'inizio della stringa. In tal caso l'elenco delle capacità può non essere
indicato e viene assunto che si stia facendo riferimento a tutte quante senza
doverlo scrivere esplicitamente.

Come esempi avremo allora che un processo non privilegiato di un utente, che
non ha nessuna capacità attiva, avrà una rappresentazione nella forma
``\texttt{=}'' che corrisponde al fatto che nessuna capacità viene assegnata a
nessun insieme (vale la cancellazione preventiva), mentre un processo con
privilegi di amministratore avrà una rappresentazione nella forma
``\texttt{=ep}'' in cui tutte le capacità vengono assegnate agli insiemi
\textit{permitted} ed \textit{effective} (e l'\textit{inheritable} è ignorato
in quanto per le regole viste a pag.~\ref{sec:capability-uid-transition} le
capacità verranno comunque attivate attraverso una \func{exec}). Infine, come
esempio meno banale dei precedenti, otterremo per \texttt{init} una
rappresentazione nella forma ``\texttt{=ep cap\_setpcap-e}'' dato che come
accennato tradizionalmente \const{CAP\_SETPCAP} è sempre stata rimossa da
detto processo.

Viceversa per ottenere un \textit{capability state} dalla sua rappresentazione
testuale si può usare la funzione \funcd{cap\_from\_text}, il cui prototipo è:

\begin{funcproto}{
\fhead{sys/capability.h}
\fdecl{cap\_t cap\_from\_text(const char *string)}
\fdesc{Crea un \textit{capability state} dalla sua rappresentazione testuale.} 
}

{La funzione ritorna un \textit{capability state} in caso di successo e
  \val{NULL} per un errore, nel qual caso \var{errno} assumerà i valori
  \errval{EINVAL} o \errval{ENOMEM} nel loro significato generico.}
\end{funcproto}


La funzione restituisce il puntatore ad un \textit{capability state}
inizializzato con i valori indicati nella stringa \param{string} che ne
contiene la rappresentazione testuale. La memoria per il \textit{capability
  state} viene allocata automaticamente dalla funzione e dovrà essere liberata
con \func{cap\_free}.

Alle due funzioni citate se ne aggiungono altre due che consentono di
convertire i valori delle costanti di tab.~\ref{tab:proc_capabilities} nelle
stringhe usate nelle rispettive rappresentazioni e viceversa. Le due funzioni,
\funcd{cap\_to\_name} e \funcd{cap\_from\_name}, sono estensioni specifiche di
Linux ed i rispettivi prototipi sono:

\begin{funcproto}{
\fhead{sys/capability.h}
\fdecl{char *cap\_to\_name(cap\_value\_t cap)}
\fdesc{Converte il valore numerico di una \textit{capabilities} alla sua
  rappresentazione testuale.} 
\fdecl{int cap\_from\_name(const char *name, cap\_value\_t *cap\_p)}

\fdesc{Converte la rappresentazione testuale di una \textit{capabilities} al
  suo valore numerico.} 
}

{La funzione \func{cap\_to\_name} ritorna un puntatore ad una stringa in caso
  di successo e \val{NULL} per un errore, mentre \func{cap\_to\_name} ritorna
  $0$ in caso di successo e $-1$ per un errore, per entrambe in caso di errore
  \var{errno} assumerà i valori \errval{EINVAL} o \errval{ENOMEM} nel loro
  significato generico.  
}
\end{funcproto}

La prima funzione restituisce la stringa (allocata automaticamente e che dovrà
essere liberata con \func{cap\_free}) che corrisponde al valore della
capacità \param{cap}, mentre la seconda restituisce nella variabile puntata
da \param{cap\_p}, come \textit{value result argument}, il valore della
capacità rappresentata dalla stringa \param{name}.

Fin quei abbiamo trattato solo le funzioni di servizio relative alla
manipolazione dei \textit{capability state} come strutture di dati;
l'interfaccia di gestione prevede però anche le funzioni per trattare le
\textit{capabilities} presenti nei processi. La prima di queste funzioni è
\funcd{cap\_get\_proc} che consente la lettura delle \textit{capabilities} del
processo corrente, il suo prototipo è:

\begin{funcproto}{
\fhead{sys/capability.h}
\fdecl{cap\_t cap\_get\_proc(void)}
\fdesc{Legge le \textit{capabilities} del processo corrente.} 
}

{La funzione ritorna un \textit{capability state} in caso di successo e
  \val{NULL} per un errore, nel qual caso \var{errno} assumerà i valori
  \errval{EINVAL}, \errval{EPERM} o \errval{ENOMEM} nel loro significato
  generico.}
\end{funcproto}

La funzione legge il valore delle \textit{capabilities} associate al processo
da cui viene invocata, restituendo il risultato tramite il puntatore ad un
\textit{capability state} contenente tutti i dati che provvede ad allocare
autonomamente e che di nuovo occorrerà liberare con \func{cap\_free} quando
non sarà più utilizzato.

Se invece si vogliono leggere le \textit{capabilities} di un processo
specifico occorre usare la funzione \funcd{cap\_get\_pid}, il cui
prototipo\footnote{su alcune pagine di manuale la funzione è descritta con un
  prototipo sbagliato, che prevede un valore di ritorno di tipo \type{cap\_t},
  ma il valore di ritorno è intero, come si può verificare anche dalla
  dichiarazione della stessa in \headfile{sys/capability.h}.} è:

\begin{funcproto}{
\fhead{sys/capability.h}
\fdecl{cap\_t cap\_get\_pid(pid\_t pid)}
\fdesc{Legge le \textit{capabilities} di un processo.} 
}

{La funzione ritorna un \textit{capability state} in caso di successo e
  \val{NULL} per un errore, nel qual caso \var{errno} assumerà i valori
  \errval{ESRCH} o \errval{ENOMEM} nel loro significato generico.  }
\end{funcproto}

La funzione legge il valore delle \textit{capabilities} del processo indicato
con l'argomento \param{pid}, e restituisce il risultato tramite il puntatore
ad un \textit{capability state} contenente tutti i dati che provvede ad
allocare autonomamente e che al solito deve essere disallocato con
\func{cap\_free}. Qualora il processo indicato non esista si avrà un errore di
\errval{ESRCH}. Gli stessi valori possono essere letti direttamente nel
filesystem \textit{proc}, nei file \texttt{/proc/<pid>/status}; ad esempio per
\texttt{init} si otterrà qualcosa del tipo:
\begin{Console}
piccardi@hain:~/gapil$ \textbf{cat /proc/1/status}
...
CapInh: 0000000000000000
CapPrm: 00000000fffffeff
CapEff: 00000000fffffeff  
...
\end{Console}
%$

\itindend{capability~state}

Infine per impostare le \textit{capabilities} del processo corrente (nella
bozza dello standard POSIX.1e non esiste una funzione che permetta di cambiare
le \textit{capabilities} di un altro processo) si deve usare la funzione
\funcd{cap\_set\_proc}, il cui prototipo è:

\begin{funcproto}{
\fhead{sys/capability.h}
\fdecl{int cap\_set\_proc(cap\_t cap\_p)}
\fdesc{Imposta le \textit{capabilities} del processo corrente.} 
}

{La funzione ritorna $0$ in caso di successo e $-1$ per un errore, nel qual
  caso \var{errno} assumerà i valori:
  \begin{errlist}
  \item[\errcode{EPERM}] si è cercato di attivare una capacità non permessa.
  \end{errlist} ed inoltre \errval{EINVAL} nel suo significato generico.}
\end{funcproto}

La funzione modifica le \textit{capabilities} del processo corrente secondo
quanto specificato con l'argomento \param{cap\_p}, posto che questo sia
possibile nei termini spiegati in precedenza (non sarà ad esempio possibile
impostare capacità non presenti nell'insieme di quelle permesse). 

In caso di successo i nuovi valori saranno effettivi al ritorno della
funzione, in caso di fallimento invece lo stato delle capacità resterà
invariato. Si tenga presente che \textsl{tutte} le capacità specificate
tramite \param{cap\_p} devono essere permesse; se anche una sola non lo è la
funzione fallirà, e per quanto appena detto, lo stato delle
\textit{capabilities} non verrà modificato (neanche per le parti eventualmente
permesse).

Oltre a queste funzioni su Linux sono presenti due ulteriori funzioni,
\funcm{capgetp} e \funcm{capsetp}, che svolgono un compito analogo. Queste
funzioni risalgono alla implementazione iniziale delle \textit{capabilities}
ed in particolare \funcm{capsetp} consentirebbe anche, come possibile in quel
caso, di cambiare le capacità di un altro processo. Le due funzioni oggi sono
deprecate e pertanto eviteremo di trattarle, per chi fosse interessato si
rimanda alla lettura della loro pagina di manuale.

Come esempio di utilizzo di queste funzioni nei sorgenti allegati alla guida
si è distribuito il programma \texttt{getcap.c}, che consente di leggere le
\textit{capabilities} del processo corrente\footnote{vale a dire di sé stesso,
  quando lo si lancia, il che può sembrare inutile, ma serve a mostrarci quali
  sono le \textit{capabilities} standard che ottiene un processo lanciato
  dalla riga di comando.} o tramite l'opzione \texttt{-p}, quelle di un
processo qualunque il cui \ids{PID} viene passato come parametro dell'opzione.

\begin{figure}[!htbp]
  \footnotesize \centering
  \begin{minipage}[c]{\codesamplewidth}
    \includecodesample{listati/getcap.c}
  \end{minipage} 
  \normalsize
  \caption{Corpo principale del programma \texttt{getcap.c}.}
  \label{fig:proc_getcap}
\end{figure}

La sezione principale del programma è riportata in fig.~\ref{fig:proc_getcap},
e si basa su una condizione sulla variabile \var{pid} che se si è usato
l'opzione \texttt{-p} è impostata (nella sezione di gestione delle opzioni,
che si è tralasciata) al valore del \ids{PID} del processo di cui si vuole
leggere le \textit{capabilities} e nulla altrimenti. Nel primo caso
(\texttt{\small 1-6}) si utilizza (\texttt{\small 2}) \func{cap\_get\_proc}
per ottenere lo stato delle capacità del processo, nel secondo (\texttt{\small
  7-13}) si usa invece \func{cap\_get\_pid} (\texttt{\small 8}) per leggere
il valore delle capacità del processo indicato.

Il passo successivo è utilizzare (\texttt{\small 15}) \func{cap\_to\_text} per
tradurre in una stringa lo stato, e poi (\texttt{\small 16}) stamparlo; infine
(\texttt{\small 18-19}) si libera la memoria allocata dalle precedenti
funzioni con \func{cap\_free} per poi ritornare dal ciclo principale della
funzione.

\itindend{capabilities}

% TODO vedi http://lwn.net/Articles/198557/ e 
% http://www.madore.org/~david/linux/newcaps/




\subsection{La gestione del \textit{Secure Computing}.}
\label{sec:procadv_seccomp}

\itindbeg{secure~computing~mode}

Il \textit{secure computing mode} è un meccanismo ideato per fornire un
supporto per l'esecuzione di codice esterno non fidato e non verificabile a
scopo di calcolo. L'idea era quella di disporre di una modalità di esecuzione
dei programmi che permettesse di vendere la capacità di calcolo della propria
macchina ad un qualche servizio di calcolo distribuito, senza comprometterne
la sicurezza eseguendo codice non sotto il proprio controllo.

La prima versione del meccanismo è stata introdotta con il kernel
2.6.23,\footnote{e disponibile solo avendo abilitato il supporto nel kernel
  con l'opzione di configurazione \texttt{CONFIG\_SECCOMP}.} è molto semplice,
il \textit{secure computing mode} viene attivato con \func{prctl} usando
l'opzione \const{PR\_SET\_SECCOMP}, ed indicando \const{SECCOMP\_MODE\_STRICT}
come valore per \param{arg2} (all'epoca unico valore possibile).  Una volta
abilitato in questa modalità (in seguito denominata \textit{strict mode}) il
processo o il \textit{thread} chiamante potrà utilizzare soltanto un insieme
estremamente limitato di \textit{system call}: \func{read}, \func{write},
\func{\_exit} e \funcm{sigreturn}; l'esecuzione di qualsiasi altra
\textit{system call} comporta l'emissione di un \signal{SIGKILL} e conseguente
terminazione immediata del processo.

Si tenga presente che in questo caso, con versioni recenti della \acr{glibc}
(il comportamento è stato introdotto con la 2.3), diventa impossibile usare
anche \func{\_exit} in \textit{strict mode}, in quanto questa funzione viene
intercettata ed al suo posto viene chiamata \func{exit\_group} (vedi
sez.~\ref{sec:pthread_management}) che non è consentita e comporta un
\signal{SIGKILL}.

Si tenga presente che, non essendo \func{execve} fra le funzioni permesse, per
poter eseguire un programma terzo essendo in \textit{strict mode} questo dovrà
essere fornito in una forma di codice interpretabile fornito attraverso un
socket o una \textit{pipe}, creati prima di lanciare il processo che eseguirà
il codice non fidato. 


% TODO a partire dal kernel 3.5 è stato introdotto la possibilità di usare un
% terzo argomento se il secondo è SECCOMP_MODE_FILTER, vedi
% Documentation/prctl/seccomp_filter.txt 
% vedi anche http://lwn.net/Articles/600250/

% TODO documentare PR_SET_SECCOMP introdotto a partire dal kernel 3.5. Vedi:
% * Documentation/prctl/seccomp_filter.txt
% * http://lwn.net/Articles/475043/

% TODO a partire dal kernel 3.17 è stata introdotta la nuova syscall seccomp,
% vedi http://lwn.net/Articles/600250/ e http://lwn.net/Articles/603321/


\itindend{secure~computing~mode}

\subsection{Altre funzionalità di sicurezza.}
\label{sec:procadv_security_misc}

Oltre alle funzionalità specifiche esaminate nelle sezioni precedenti, il
kernel supporta una varietà di ulteriori impostazioni di sicurezza,
accessibili nelle maniere più varie, che abbiamo raccolto in questa sezione.

Una serie di modalità di sicurezza sono attivabili a richiesta attraverso
alcune opzioni di controllo attivabili via \func{sysctl} o il filesystem
\texttt{/proc}, un elenco delle stesse e dei loro effetti è il seguente:

\begin{basedescript}{\desclabelwidth{1cm}\desclabelstyle{\nextlinelabel}}
\item[\sysctlrelfiled{fs}{protected\_hardlinks}] Un valore nullo, il default,
  mantiene il comportamento standard che non pone restrizioni alla creazione
  di \textit{hard link}. Se il valore viene posto ad 1 vengono invece attivate
  una serie di restrizioni protettive, denominate
  \itindex{protected~hardlinks} \textit{protected hardlinks}, che se non
  soddisfatte causano il fallimento di \func{link} con un errore di
  \errval{EPERM}. Perché questo non avvenga almeno una delle seguenti
  condizioni deve essere soddisfatta:
  \begin{itemize*}
  \item il chiamante deve avere privilegi amministrativi (la
    \textit{capability} \const{CAP\_FOWNER}). In caso di utilizzo
    dell'\textit{user namespace} oltre a possedere \const{CAP\_FOWNER} è
    necessario che l'\ids{UID} del proprietario del file sia mappato nel
    \textit{namespace}.
  \item il \textit{filesystem} \ids{UID} del chiamante (normalmente
    equivalente all'\ids{UID} effettivo) deve corrispondere a quello del
    proprietario del file a cui si vuole effettuare il collegamento.
  \item devono essere soddisfatte tutte le seguenti condizioni:
    \begin{itemize*}
    \item il file è un file ordinario
    \item il file non ha il \acr{suid} bit attivo
    \item il file non ha lo \acr{sgid} bit attivo ed il permesso di esecuzione
      per il gruppo
    \item il chiamante ha i permessi di lettura e scrittura sul file
    \end{itemize*}
  \end{itemize*}

  In sostanza in questo caso un utente potrà creare un collegamento diretto ad
  un altro file solo se ne è il proprietario o se questo è un file ordinario
  senza permessi speciali ed a cui ha accesso in lettura e scrittura.

  Questa funzionalità fornisce una protezione generica che non inficia l'uso
  ordinario di \func{link}, ma rende impraticabili una serie di possibili
  abusi della stessa; oltre ad impedire l'uso di un \textit{hard link} come
  variante in un attacco di \textit{symlink race} (eludendo i
  \textit{protected symlinks} di cui al punto successivo), evita anche che si
  possa lasciare un riferimento ad un eventuale programma \acr{suid}
  vulnerabile, creando un collegamento diretto allo stesso.


\item[\sysctlrelfiled{fs}{protected\_symlinks}] Un valore nullo, il default,
  mantiene il comportamento standard che non pone restrizioni nel seguire i
  link simbolici. Se il valore viene posto ad 1 vengono attivate delle
  restrizioni protettive, denominate \itindex{protected~symlinks}
  \textit{protected symlinks}. Quando vengono attivate una qualunque funzione
  che esegua la risoluzione di un \textit{pathname} contenente un link
  simbolico non conforme alle restrizioni fallirà con un errore di
  \errval{EACCESS}. Per evitare l'errore deve essere soddisfatta una delle
  seguenti condizioni:
  \begin{itemize*}
  \item il link non è in una directory con permessi analoghi a \file{/tmp}
    (scrivibile a tutti e con lo \textit{sticky bit} attivo);
  \item il link è in una directory con permessi analoghi a \file{/tmp} ma è
    soddisfatta una delle condizioni seguenti: 
    \begin{itemize*}
    \item il link simbolico appartiene al chiamante: il controllo viene fatto
      usando il \textit{filesystem} \ids{UID} (che normalmente corrisponde
      all'\ids{UID} effettivo).
    \item il link simbolico ha lo stesso proprietario della directory.
    \end{itemize*}
  \end{itemize*}

  Questa funzionalità consente di rendere impraticabili alcuni attacchi in cui
  si approfitta di una differenza di tempo fra il controllo e l'uso di un
  file, ed in particolare quella classe di attacchi viene usualmente chiamati
  \textit{symlink attack},\footnote{si tratta di un sottoinsieme di quella
    classe di attacchi chiamata genericamente \textit{TOCTTOU}, acronimo
    appunto di \textit{Time of check to time of use}.} di cui abbiamo parlato
  in sez.~\ref{sec:file_temp_file}.

  Un possibile esempio di questo tipo di attacco è quello contro un programma
  che viene eseguito per conto di un utente privilegiato (ad esempio un
  programma con il \acr{suid} o lo \acr{sgid} bit attivi) che prima controlla
  l'esistenza di un file e se non esiste lo crea. Se questa procedura, che è
  tipica della creazione di file temporanei sotto \file{/tmp}, non viene
  eseguita in maniera corretta,\footnote{ad esempio con le modalità che
    abbiamo trattato in sez.~\ref{sec:file_temp_file}, che per quanto note da
    tempo continuano ad essere ignorate.} un attaccante ha una finestra di
  tempo in cui può creare prima del programma un \textit{link simbolico} ad un
  file di sua scelta, compresi file di dispositivo o file a cui non avrebbe
  accesso, facendolo poi utilizzare al programma.

  Attivando la funzionalità si rende impossibile seguire un link simbolico in
  una directory temporanea come \texttt{/tmp}, a meno che questo non sia di
  proprietà del chiamante, o che questo non appartenga al proprietario della
  directory. Questo impedisce che i link simbolici creati da un attaccante
  possano essere seguiti da un programma privilegiato (perché apparterranno
  all'attaccante) mentre quelli creati dall'amministratore (che i genere è il
  proprietario di \texttt{/tmp}) saranno seguiti comunque.

\end{basedescript}


% TODO: trattare pure protected_regular e protected_fifos introdotti con il
% 4.19 (vedi https://lwn.net/Articles/763106/)



% TODO: trattare keyctl (man 2 keyctl)



% TODO trattare le funzioni di protezione della memoria pkey_alloc, pkey_free,
% pkey_mprotect, introdotte con il kernel 4.8, vedi
% http://lwn.net/Articles/689395/ e Documentation/x86/protection-keys.txt 

\section{Funzioni di gestione e controllo}
\label{sec:proc_manage_control}

In questa sezione prenderemo in esame alcune specifiche \textit{system call}
dedicate al controllo dei processi sia per quanto riguarda l'impostazione di
caratteristiche specialistiche, che per quanto riguarda l'analisi ed il
controllo della loro esecuzione.

\subsection{La funzione \func{prctl}}
\label{sec:process_prctl}

Benché la gestione ordinaria dei processi possa essere effettuata attraverso
le funzioni che abbiamo già esaminato nei capitoli \ref{cha:process_interface}
e \ref{cha:process_handling}, esistono una serie di proprietà e
caratteristiche specifiche dei processi per la cui gestione è stata
predisposta una apposita \textit{system call} che fornisce una interfaccia
generica per tutte le operazioni specialistiche. La funzione di sistema è
\funcd{prctl} ed il suo prototipo è:\footnote{la funzione non è standardizzata
  ed è specifica di Linux, anche se ne esiste una analoga in IRIX; è stata
  introdotta con il kernel 2.1.57.}

\begin{funcproto}{ 
\fhead{sys/prctl.h}
\fdecl{int prctl(int option, unsigned long arg2, unsigned long arg3, unsigned
  long arg4, \\
\phantom{int prctl(}unsigned long arg5)}
\fdesc{Esegue una operazione speciale sul processo corrente.} 
}

{La funzione ritorna $0$ o un valore positivo dipendente dall'operazione in
  caso di successo e $-1$ per un errore, nel qual caso \var{errno} assumerà
  valori diversi a seconda del tipo di operazione richiesta, sono possibili:
  \errval{EACCESS}, \errval{EBADF}, \errval{EBUSY}, \errval{EFAULT},
  \errval{EINVAL}, \errval{ENXIO}, \errval{EOPNOTSUPP} o \errval{EPERM}.}
\end{funcproto}

La funzione ritorna in caso di successo un valore nullo o positivo, e $-1$ in
caso di errore. Il significato degli argomenti della funzione successivi al
primo, il valore di ritorno in caso di successo, il tipo di errore restituito
in \var{errno} dipendono dall'operazione eseguita, indicata tramite il primo
argomento, \param{option}. Questo è un valore intero che identifica
l'operazione, e deve essere specificato con l'uso di una delle costanti
predefinite del seguente elenco.\footnote{l'elenco potrebbe non risultare
  aggiornato, in quanto nuove operazioni vengono aggiunte nello sviluppo del
  kernel.} Tratteremo esplicitamente per ciascuna di esse il significato del
il valore di ritorno in caso di successo, ma solo quando non corrisponde
all'ordinario valore nullo (dato per implicito).

%TODO: trattare PR_CAP_AMBIENT, dal 4.3
%TODO: trattare PR_CAP_FP_*, dal 4.0, solo per MIPS
%TODO: trattare PR_MPX_*_MANAGEMENT, dal 3.19
%TODO: trattare PR_*NO_NEW_PRIVS, dal 3.5

\begin{basedescript}{\desclabelwidth{1.5cm}\desclabelstyle{\nextlinelabel}}
\item[\constd{PR\_CAPBSET\_READ}] Controlla la disponibilità di una delle
  \textit{capability} (vedi sez.~\ref{sec:proc_capabilities}). La funzione
  ritorna 1 se la capacità specificata nell'argomento \param{arg2} (con una
  delle costanti di tab.~\ref{tab:proc_capabilities}) è presente nel
  \textit{capabilities bounding set} del processo e zero altrimenti,
  se \param{arg2} non è un valore valido si avrà un errore di \errval{EINVAL}.
  Introdotta a partire dal kernel 2.6.25.

\item[\constd{PR\_CAPBSET\_DROP}] Rimuove permanentemente una delle
  \textit{capabilities} (vedi sez.~\ref{sec:proc_capabilities}) dal processo e
  da tutti i suoi discendenti. La funzione cancella la capacità specificata
  nell'argomento \param{arg2} con una delle costanti di
  tab.~\ref{tab:proc_capabilities} dal \textit{capabilities bounding set} del
  processo. L'operazione richiede i privilegi di amministratore (la capacità
  \const{CAP\_SETPCAP}), altrimenti la chiamata fallirà con un errore di
  \errcode{EPERM}; se il valore di \param{arg2} non è valido o se il supporto
  per le \textit{file capabilities} non è stato compilato nel kernel la
  chiamata fallirà con un errore di \errval{EINVAL}. Introdotta a partire dal
  kernel 2.6.25.

\item[\constd{PR\_SET\_DUMPABLE}] Imposta il flag che determina se la
  terminazione di un processo a causa di un segnale per il quale è prevista la
  generazione di un file di \textit{core dump} (vedi
  sez.~\ref{sec:sig_standard}) lo genera effettivamente. In genere questo flag
  viene attivato automaticamente, ma per evitare problemi di sicurezza (la
  generazione di un file da parte di processi privilegiati può essere usata
  per sovrascriverne altri) viene cancellato quando si mette in esecuzione un
  programma con i bit \acr{suid} e \acr{sgid} attivi (vedi
  sez.~\ref{sec:file_special_perm}) o con l'uso delle funzioni per la modifica
  degli \ids{UID} dei processi (vedi sez.~\ref{sec:proc_setuid}).

  L'operazione è stata introdotta a partire dal kernel 2.3.20, fino al kernel
  2.6.12 e per i kernel successivi al 2.6.17 era possibile usare solo un
  valore 0 (espresso anche come \constd{SUID\_DUMP\_DISABLE}) di \param{arg2}
  per disattivare il flag ed un valore 1 (espresso anche come
  \constd{SUID\_DUMP\_USER}) per attivarlo. Nei kernel dal 2.6.13 al 2.6.17 è
  stato supportato anche il valore 2, che causava la generazione di un
  \textit{core dump} leggibile solo dall'amministratore, ma questa
  funzionalità è stata rimossa per motivi di sicurezza, in quanto consentiva
  ad un utente normale di creare un file di \textit{core dump} appartenente
  all'amministratore in directory dove l'utente avrebbe avuto permessi di
  accesso. Specificando un valore diverso da 0 o 1 si ottiene un errore di
  \errval{EINVAL}.

\item[\constd{PR\_GET\_DUMPABLE}] Ottiene come valore di ritorno della funzione
  lo stato corrente del flag che controlla la effettiva generazione dei
  \textit{core dump}. Introdotta a partire dal kernel 2.3.20.

\item[\constd{PR\_SET\_ENDIAN}] Imposta la \textit{endianness} del processo
  chiamante secondo il valore fornito in \param{arg2}. I valori possibili sono
  sono: \constd{PR\_ENDIAN\_BIG} (\textit{big endian}),
  \constd{PR\_ENDIAN\_LITTLE} (\textit{little endian}), e
  \constd{PR\_ENDIAN\_PPC\_LITTLE} (lo pseudo \textit{little endian} del
  PowerPC). Introdotta a partire dal kernel 2.6.18, solo per architettura
  PowerPC.

\item[\constd{PR\_GET\_ENDIAN}] Ottiene il valore della \textit{endianness} del
  processo chiamante, salvato sulla variabile puntata da \param{arg2} che deve
  essere passata come di tipo ``\ctyp{int *}''. Introdotta a partire dal
  kernel 2.6.18, solo su PowerPC.

\item[\constd{PR\_SET\_FPEMU}] Imposta i bit di controllo per l'emulazione
  della virgola mobile su architettura ia64, secondo il valore
  di \param{arg2}, si deve passare \constd{PR\_FPEMU\_NOPRINT} per emulare in
  maniera trasparente l'accesso alle operazioni in virgola mobile, o
  \constd{PR\_FPEMU\_SIGFPE} per non emularle ed inviare il segnale
  \signal{SIGFPE} (vedi sez.~\ref{sec:sig_prog_error}). Introdotta a partire
  dal kernel 2.4.18, solo su architettura ia64.

\item[\constd{PR\_GET\_FPEMU}] Ottiene il valore dei flag di controllo
  dell'emulazione della virgola mobile, salvato all'indirizzo puntato
  da \param{arg2}, che deve essere di tipo ``\ctyp{int *}''. Introdotta a
  partire dal kernel 2.4.18, solo su architettura ia64.

\item[\constd{PR\_SET\_FPEXC}] Imposta la modalità delle eccezioni in virgola
  mobile (\textit{floating-point exception mode}) al valore di \param{arg2}.
  I valori possibili sono: 
  \begin{itemize*}
  \item \constd{PR\_FP\_EXC\_SW\_ENABLE} per usare FPEXC per le eccezioni,
  \item \constd{PR\_FP\_EXC\_DIV} per la divisione per zero in virgola mobile,
  \item \constd{PR\_FP\_EXC\_OVF} per gli overflow,
  \item \constd{PR\_FP\_EXC\_UND} per gli underflow,
  \item \constd{PR\_FP\_EXC\_RES} per risultati non esatti,
  \item \constd{PR\_FP\_EXC\_INV} per operazioni invalide,
  \item \constd{PR\_FP\_EXC\_DISABLED} per disabilitare le eccezioni,
  \item \constd{PR\_FP\_EXC\_NONRECOV} per usare la modalità di eccezione
    asincrona non recuperabile,
  \item \constd{PR\_FP\_EXC\_ASYNC} per usare la modalità di eccezione
    asincrona recuperabile,
  \item \constd{PR\_FP\_EXC\_PRECISE} per la modalità precisa di
    eccezione.\footnote{trattasi di gestione specialistica della gestione
      delle eccezioni dei calcoli in virgola mobile che, i cui dettagli al
      momento vanno al di là dello scopo di questo testo.}
  \end{itemize*}
Introdotta a partire dal kernel 2.4.21, solo su PowerPC.

\item[\constd{PR\_GET\_FPEXC}] Ottiene il valore della modalità delle eccezioni
  delle operazioni in virgola mobile, salvata all'indirizzo
  puntato \param{arg2}, che deve essere di tipo ``\ctyp{int *}''.  Introdotta
  a partire dal kernel 2.4.21, solo su PowerPC.

\item[\constd{PR\_SET\_KEEPCAPS}] Consente di controllare quali
  \textit{capabilities} vengono cancellate quando si esegue un cambiamento di
  \ids{UID} del processo (per i dettagli si veda
  sez.~\ref{sec:proc_capabilities}, in particolare quanto illustrato a
  pag.~\pageref{sec:capability-uid-transition}). Un valore nullo (il default)
  per \param{arg2} comporta che vengano cancellate, il valore 1 che vengano
  mantenute, questo valore viene sempre cancellato attraverso una \func{exec}.
  L'uso di questo flag è stato sostituito, a partire dal kernel 2.6.26, dal
  flag \const{SECURE\_KEEP\_CAPS} dei \textit{securebits} (vedi
  sez.~\ref{sec:proc_capabilities} e l'uso di \const{PR\_SET\_SECUREBITS} più
  avanti) e si è impostato con essi \const{SECURE\_KEEP\_CAPS\_LOCKED} si
  otterrà un errore di \errval{EPERM}.  Introdotta a partire dal kernel
  2.2.18.

\item[\constd{PR\_GET\_KEEPCAPS}] Ottiene come valore di ritorno della funzione
  il valore del flag di controllo delle \textit{capabilities} impostato con
  \const{PR\_SET\_KEEPCAPS}. Introdotta a partire dal kernel 2.2.18.

\item[\constd{PR\_SET\_NAME}] Imposta il nome del processo chiamante alla
  stringa puntata da \param{arg2}, che deve essere di tipo ``\ctyp{char *}''. Il
  nome può essere lungo al massimo 16 caratteri, e la stringa deve essere
  terminata da NUL se più corta.  Introdotta a partire dal kernel 2.6.9.

\item[\constd{PR\_GET\_NAME}] Ottiene il nome del processo chiamante nella
  stringa puntata da \param{arg2}, che deve essere di tipo ``\ctyp{char *}'';
  si devono allocare per questo almeno 16 byte, e il nome sarà terminato da
  NUL se più corto. Introdotta a partire dal kernel 2.6.9.

\item[\constd{PR\_SET\_PDEATHSIG}] Consente di richiedere l'emissione di un
  segnale, che sarà ricevuto dal processo chiamante, in occorrenza della
  terminazione del proprio processo padre; in sostanza consente di invertire
  il ruolo di \signal{SIGCHLD}. Il valore di \param{arg2} deve indicare il
  numero del segnale, o 0 per disabilitare l'emissione. Il valore viene
  automaticamente cancellato per un processo figlio creato con \func{fork}.
  Introdotta a partire dal kernel 2.1.57.

\item[\constd{PR\_GET\_PDEATHSIG}] Ottiene il valore dell'eventuale segnale
  emesso alla terminazione del padre, salvato all'indirizzo
  puntato \param{arg2}, che deve essere di tipo ``\ctyp{int *}''. Introdotta a
  partire dal kernel 2.3.15.

\item[\constd{PR\_SET\_PTRACER}] Imposta un \ids{PID} per il ``\textit{tracer
    process}'' usando \param{arg2}. Una impostazione successiva sovrascrive la
  precedente, ed un valore nullo cancella la disponibilità di un
  ``\textit{tracer process}''. Questa è una funzionalità fornita da
  \textit{``Yama''}, uno specifico \textit{Linux Security Modules}, e serve a
  consentire al processo indicato, quando le restrizioni introdotte da questo
  modulo sono attive, di usare \func{ptrace} (vedi
  sez.~\ref{sec:process_ptrace}) sul processo chiamante, anche se quello
  indicato non ne è un progenitore. Il valore \constd{PR\_SET\_PTRACER\_ANY}
  consente a tutti i processi l'uso di \func{ptrace}. L'uso si \textit{Yama}
  attiene alla gestione della sicurezza dei processi, e consente di introdurre
  una restrizione all'uso di \func{ptrace}, che è spesso sorgente di
  compromissioni. Si tratta di un uso specialistico che va al di là dello
  scopo di queste dispense, per i dettagli si consulti la documentazione su
  \textit{Yama} nei sorgenti del kernel. Introdotta a partire dal kernel 3.4.

\item[\constd{PR\_SET\_SECCOMP}] Attiva il \textit{secure computing mode} per
  il processo corrente. Introdotta a partire dal kernel 2.6.23 la funzionalità
  è stata ulteriormente estesa con il kernel 3.5, salvo poi diventare un
  sottoinsieme della \textit{system call} \func{seccomp} a partire dal kernel
  3.17. Prevede che si indichi per \param{arg2} il valore
  \const{SECCOMP\_MODE\_STRICT} (unico possibile fino al kernel 2.6.23) per
  selezionare il cosiddetto \textit{strict mode} o, dal kernel 3.5,
  \const{SECCOMP\_MODE\_FILTER} per usare il \textit{filter mode}. Tratteremo
  questa opzione nei dettagli più avanti, in sez.~\ref{sec:procadv_seccomp},
  quando affronteremo l'argomento del \textit{Secure Computing}.

\item[\constd{PR\_GET\_SECCOMP}] Ottiene come valore di ritorno della funzione
  lo stato corrente del \textit{secure computing mode}. Fino al kernel 3.5,
  quando era possibile solo lo \textit{strict mode}, la funzione era
  totalmente inutile in quanto l'unico valore ottenibile era 0 in assenza di
  \textit{secure computing}, dato che la chiamata di questa funzione in
  \textit{strict mode} avrebbe comportato l'emissione di \signal{SIGKILL} per
  il chiamante. La funzione però, a partire dal kernel 2.6.23, era stata
  comunque definita per eventuali estensioni future, ed infatti con
  l'introduzione del \textit{filter mode} con il kernel 3.5, se essa viene
  inclusa nelle funzioni consentite restituisce il valore 2 quando il
  \textit{secure computing mode} è attivo (se non inclusa si avrà di nuovo un
  \signal{SIGKILL}).

\item[\constd{PR\_SET\_SECUREBITS}] Imposta i \textit{securebits} per il
  processo chiamante al valore indicato da \param{arg2}; per i dettagli sul
  significato dei \textit{securebits} si veda
  sez.~\ref{sec:proc_capabilities}, ed in particolare i valori di
  tab.~\ref{tab:securebits_values} e la relativa trattazione. L'operazione
  richiede i privilegi di amministratore (la capacità \const{CAP\_SETPCAP}),
  altrimenti la chiamata fallirà con un errore di \errval{EPERM}. Introdotta a
  partire dal kernel 2.6.26.

\item[\constd{PR\_GET\_SECUREBITS}] Ottiene come valore di ritorno della
  funzione l'impostazione corrente per i \textit{securebits}. Introdotta a
  partire dal kernel 2.6.26.

\item[\constd{PR\_SET\_TIMING}] Imposta il metodo di temporizzazione del
  processo da indicare con il valore di \param{arg2}, attualmente i valori
  possibili sono due, con \constd{PR\_TIMING\_STATISTICAL} si usa il metodo
  statistico tradizionale, con \constd{PR\_TIMING\_TIMESTAMP} il più accurato
  basato su dei \textit{timestamp}, quest'ultimo però non è ancora
  implementato ed il suo uso comporta la restituzione di un errore di
  \errval{EINVAL}. Introdotta a partire dal kernel 2.6.0-test4.

\item[\constd{PR\_GET\_TIMING}] Ottiene come valore di ritorno della funzione
  il metodo di temporizzazione del processo attualmente in uso (uno dei due
  valori citati per \const{PR\_SET\_TIMING}). Introdotta a partire dal kernel
  2.6.0-test4.

\item[\constd{PR\_SET\_TSC}] Imposta il flag che indica se il processo
  chiamante può leggere il registro di processore contenente il contatore dei
  \textit{timestamp} (TSC, o \textit{Time Stamp Counter}) da indicare con il
  valore di \param{arg2}. Si deve specificare \constd{PR\_TSC\_ENABLE} per
  abilitare la lettura o \constd{PR\_TSC\_SIGSEGV} per disabilitarla con la
  generazione di un segnale di \signal{SIGSEGV} (vedi
  sez.~\ref{sec:sig_prog_error}). La lettura viene automaticamente
  disabilitata se si attiva il \textit{secure computing mode} (vedi
  \const{PR\_SET\_SECCOMP} e sez.~\ref{sec:procadv_seccomp}).  Introdotta a
  partire dal kernel 2.6.26, solo su x86.

\item[\constd{PR\_GET\_TSC}] Ottiene il valore del flag che controlla la
  lettura del contattore dei \textit{timestamp}, salvato all'indirizzo
  puntato \param{arg2}, che deve essere di tipo ``\ctyp{int *}''. Introdotta a
  partire dal kernel 2.6.26, solo su x86.
% articoli sul TSC e relativi problemi: http://lwn.net/Articles/209101/,
% http://blog.cr0.org/2009/05/time-stamp-counter-disabling-oddities.html,
% http://en.wikipedia.org/wiki/Time_Stamp_Counter 

\item[\constd{PR\_SET\_UNALIGN}] Imposta la modalità di controllo per l'accesso
  a indirizzi di memoria non allineati, che in varie architetture risultano
  illegali, da indicare con il valore di \param{arg2}. Si deve specificare il
  valore \constd{PR\_UNALIGN\_NOPRINT} per ignorare gli accessi non allineati,
  ed il valore \constd{PR\_UNALIGN\_SIGBUS} per generare un segnale di
  \signal{SIGBUS} (vedi sez.~\ref{sec:sig_prog_error}) in caso di accesso non
  allineato.  Introdotta con diverse versioni su diverse architetture.

\item[\const{PR\_GET\_UNALIGN}] Ottiene il valore della modalità di controllo
  per l'accesso a indirizzi di memoria non allineati, salvato all'indirizzo
  puntato \param{arg2}, che deve essere di tipo \code{(int *)}. Introdotta con
  diverse versioni su diverse architetture.
\item[\const{PR\_MCE\_KILL}] Imposta la politica di gestione degli errori
  dovuti a corruzione della memoria per problemi hardware. Questo tipo di
  errori vengono riportati dall'hardware di controllo della RAM e vengono
  gestiti dal kernel,\footnote{la funzionalità è disponibile solo sulle
    piattaforme più avanzate che hanno il supporto hardware per questo tipo di
    controlli.} ma devono essere opportunamente riportati ai processi che
  usano quella parte di RAM che presenta errori; nel caso specifico questo
  avviene attraverso l'emissione di un segnale di \signal{SIGBUS} (vedi
  sez.~\ref{sec:sig_prog_error}).\footnote{in particolare viene anche
    impostato il valore di \var{si\_code} in \struct{siginfo\_t} a
    \const{BUS\_MCEERR\_AO}; per il significato di tutto questo si faccia
    riferimento alla trattazione di sez.~\ref{sec:sig_sigaction}.}

  Il comportamento di default prevede che per tutti i processi si applichi la
  politica generale di sistema definita nel file
  \sysctlfiled{vm/memory\_failure\_early\_kill}, ma specificando
  per \param{arg2} il valore \constd{PR\_MCE\_KILL\_SET} è possibile impostare
  con il contenuto di \param{arg3} una politica specifica del processo
  chiamante. Si può tornare alla politica di default del sistema utilizzando
  invece per \param{arg2} il valore \constd{PR\_MCE\_KILL\_CLEAR}. In tutti i
  casi, per compatibilità con eventuali estensioni future, tutti i valori
  degli argomenti non utilizzati devono essere esplicitamente posti a zero,
  pena il fallimento della chiamata con un errore di \errval{EINVAL}.
  
  In caso di impostazione di una politica specifica del processo con
  \const{PR\_MCE\_KILL\_SET} i valori di \param{arg3} possono essere soltanto
  due, che corrispondono anche al valore che si trova nell'impostazione
  generale di sistema di \texttt{memory\_failure\_early\_kill}, con
  \constd{PR\_MCE\_KILL\_EARLY} si richiede l'emissione immediata di
  \signal{SIGBUS} non appena viene rilevato un errore, mentre con
  \constd{PR\_MCE\_KILL\_LATE} il segnale verrà inviato solo quando il processo
  tenterà un accesso alla memoria corrotta. Questi due valori corrispondono
  rispettivamente ai valori 1 e 0 di
  \texttt{memory\_failure\_early\_kill}.\footnote{in sostanza nel primo caso
    viene immediatamente inviato il segnale a tutti i processi che hanno la
    memoria corrotta mappata all'interno del loro spazio degli indirizzi, nel
    secondo caso prima la pagina di memoria viene tolta dallo spazio degli
    indirizzi di ciascun processo, mentre il segnale viene inviato solo quei
    processi che tentano di accedervi.} Si può usare per \param{arg3} anche un
  terzo valore, \constd{PR\_MCE\_KILL\_DEFAULT}, che corrisponde a impostare
  per il processo la politica di default.\footnote{si presume la politica di
    default corrente, in modo da non essere influenzati da un eventuale
    successivo cambiamento della stessa.} Introdotta a partire dal kernel
  2.6.32.
\item[\constd{PR\_MCE\_KILL\_GET}] Ottiene come valore di ritorno della
  funzione la politica di gestione degli errori dovuti a corruzione della
  memoria. Tutti gli argomenti non utilizzati (al momento tutti) devono essere
  nulli pena la ricezione di un errore di \errval{EINVAL}. Introdotta a
  partire dal kernel 2.6.32.
\itindbeg{child~reaper}
\item[\constd{PR\_SET\_CHILD\_SUBREAPER}] Se \param{arg2} è diverso da zero
  imposta l'attributo di \textit{child reaper} per il processo, se nullo lo
  cancella. Lo stato di \textit{child reaper} è una funzionalità, introdotta
  con il kernel 3.4, che consente di far svolgere al processo che ha questo
  attributo il ruolo di ``\textsl{genitore adottivo}'' per tutti i processi
  suoi ``\textsl{discendenti}'' che diventano orfani, in questo modo il
  processo potrà ricevere gli stati di terminazione alla loro uscita,
  sostituendo in questo ruolo \cmd{init} (si ricordi quanto illustrato in
  sez.~\ref{sec:proc_termination}). Il meccanismo è stato introdotto ad uso
  dei programmi di gestione dei servizi, per consentire loro di ricevere gli
  stati di terminazione di tutti i processi che lanciano, anche se questi
  eseguono una doppia \func{fork}; nel comportamento ordinario infatti questi
  verrebbero adottati da \cmd{init} ed il programma che li ha lanciati non
  sarebbe più in grado di riceverne lo stato di terminazione. Se un processo
  con lo stato di \textit{child reaper} termina prima dei suoi discendenti,
  svolgerà questo ruolo il più prossimo antenato ad avere lo stato di
  \textit{child reaper}, 
\item[\constd{PR\_GET\_CHILD\_SUBREAPER}] Ottiene l'impostazione relativa allo
  lo stato di \textit{child reaper} del processo chiamante, salvata come
  \textit{value result} all'indirizzo puntato da \param{arg2} (da indicare
  come di tipo \code{int *}). Il valore viene letto come valore logico, se
  diverso da 0 lo stato di \textit{child reaper} è attivo altrimenti è
  disattivo. Introdotta a partire dal kernel 3.4.
\itindend{child~reaper}


% TODO documentare PR_MPX_INIT e PR_MPX_RELEASE, vedi
% http://lwn.net/Articles/582712/ 

% TODO documentare PR_SET_MM_MAP aggiunta con il kernel 3.18, per impostare i
% parametri di base del layout dello spazio di indirizzi di un processo (area
% codice e dati, stack, brack pointer ecc. vedi
% http://git.kernel.org/linus/f606b77f1a9e362451aca8f81d8f36a3a112139e 

% TODO documentare ARCH_SET_CPUID e ARCH_GET_CPUID, introdotte con il kernel
% 4.12, vedi https://lwn.net/Articles/721182/

% TODO documentare PR_SPEC_DISABLE_NOEXEC in 5.1, vedi
% https://lwn.net/Articles/782511/ 

\label{sec:prctl_operation}
\end{basedescript}


\subsection{La funzione \func{ptrace}}
\label{sec:process_ptrace}

%Da fare

% TODO: trattare PTRACE_SEIZE, aggiunta con il kernel 3.1
% TODO: trattare PTRACE_O_EXITKILL, aggiunta con il kernel 3.8 (vedi
% http://lwn.net/Articles/529060/) 
% TODO: trattare PTRACE_GETSIGMASK e PTRACE_SETSIGMASK introdotte con il
% kernel 3.11
% TODO: trattare PTRACE_O_SUSPEND_SECCOMP, aggiunta con il kernel 4.3, vedi
% http://lwn.net/Articles/656675/ 

\subsection{La funzione \func{kcmp}}
\label{sec:process_kcmp}

% TODO: trattare kcmp aggiunta con il kernel 3.5, vedi
% https://lwn.net/Articles/478111/
% vedi man kcmp e man 2 open



\section{La gestione avanzata della creazione dei processi}
\label{sec:process_adv_creation}

In questa sezione tratteremo le funzionalità avanzate relative alla creazione
dei processi e del loro ambiente, sia per quanto riguarda l'utilizzo delle
stesse per la creazione dei \textit{thread} che per la gestione dei
\textit{namespace} che sono alla base dei cosiddetti \textit{container}.


\subsection{La \textit{system call} \func{clone}}
\label{sec:process_clone}

La funzione tradizionale con cui creare un nuovo processo in un sistema
Unix-like, come illustrato in sez.~\ref{sec:proc_fork}, è \func{fork}, ma con
l'introduzione del supporto del kernel per i \textit{thread}\unavref{ (vedi
  cap.~\ref{cha:threads})}, si è avuta la necessità di una interfaccia che
consentisse un maggiore controllo sulla modalità con cui vengono creati nuovi
processi, che poi è stata utilizzata anche per fornire supporto per le
tecnologie di virtualizzazione dei processi (i cosiddetti \textit{container})
su cui torneremo in sez.~\ref{sec:process_namespaces}.

Per questo l'interfaccia per la creazione di un nuovo processo è stata
delegata ad una nuova \textit{system call}, \funcm{sys\_clone}, che consente
di reimplementare anche la tradizionale \func{fork}. In realtà in questo caso
più che di nuovi processi si può parlare della creazioni di nuovi
``\textit{task}'' del kernel che possono assumere la veste sia di un processo
classico isolato dagli altri come quelli trattati finora, che di un
\textit{thread} in cui la memoria viene condivisa fra il processo chiamante ed
il nuovo processo creato, come quelli che vedremo in
sez.~\ref{sec:linux_thread}. Per evitare confusione fra \textit{thread} e
processi ordinari, abbiamo deciso di usare la nomenclatura \textit{task} per
indicare la unità di esecuzione generica messa a disposizione del kernel che
\texttt{sys\_clone} permette di creare.

La \textit{system call} richiede soltanto due argomenti: il
primo, \param{flags}, consente di controllare le modalità di creazione del
nuovo \textit{task}, il secondo, \param{child\_stack}, imposta l'indirizzo
dello \textit{stack} per il nuovo \textit{task}, e deve essere indicato quando
si intende creare un \textit{thread}. L'esecuzione del programma creato da
\func{sys\_clone} riprende, come per \func{fork}, da dopo l'esecuzione della
stessa.

La necessità di avere uno \textit{stack} alternativo c'è solo quando si
intende creare un \textit{thread}, in tal caso infatti il nuovo \textit{task}
vede esattamente la stessa memoria del \textit{task}
``\textsl{padre}'',\footnote{in questo caso per padre si intende semplicemente
  il \textit{task} che ha eseguito \func{sys\_clone} rispetto al \textit{task}
  da essa creato, senza nessuna delle implicazioni che il concetto ha per i
  processi.} e nella sua esecuzione alla prima chiamata di una funzione
andrebbe a scrivere sullo \textit{stack} usato anche dal padre (si ricordi
quanto visto in sez.~\ref{sec:proc_mem_layout} riguardo all'uso dello
\textit{stack}).

Per evitare di doversi garantire contro la evidente possibilità di
\textit{race condition} che questa situazione comporta (vedi
sez.~\ref{sec:proc_race_cond} per una spiegazione della problematica) è
necessario che il chiamante allochi preventivamente un'area di memoria.  In
genere lo si fa con una \func{malloc} che allochi un buffer che la funzione
imposterà come \textit{stack} del nuovo processo, avendo ovviamente cura di
non utilizzarlo direttamente nel processo chiamante.

In questo modo i due \textit{task} avranno degli \textit{stack} indipendenti e
non si dovranno affrontare problematiche di \textit{race condition}.  Si tenga
presente inoltre che in molte architetture di processore lo \textit{stack}
cresce verso il basso, pertanto in tal caso non si dovrà specificare
per \param{child\_stack} il puntatore restituito da \func{malloc}, ma un
puntatore alla fine del buffer da essa allocato.

Dato che tutto ciò è necessario solo per i \textit{thread} che condividono la
memoria, la \textit{system call}, a differenza della funzione di libreria che
vedremo a breve, consente anche di passare per \param{child\_stack} il valore
\val{NULL}, che non imposta un nuovo \textit{stack}. Se infatti si crea un
processo, questo ottiene un suo nuovo spazio degli indirizzi (è sottinteso
cioè che non si stia usando il flag \const{CLONE\_VM} che vedremo a breve) ed
in questo caso si applica la semantica del \textit{copy on write} illustrata
in sez.~\ref{sec:proc_fork}, per cui le pagine dello \textit{stack} verranno
automaticamente copiate come le altre e il nuovo processo avrà un suo
\textit{stack} totalmente indipendente da quello del padre.

Dato che l'uso principale della nuova \textit{system call} è quello relativo
alla creazione dei \textit{thread}, la \acr{glibc} definisce una funzione di
libreria con una sintassi diversa, orientata a questo scopo, e la
\textit{system call} resta accessibile solo se invocata esplicitamente come
visto in sez.~\ref{sec:proc_syscall}.\footnote{ed inoltre per questa
  \textit{system call} non è disponibile la chiamata veloce con
  \texttt{vsyscall}.} La funzione di libreria si chiama semplicemente
\funcd{clone} ed il suo prototipo è:

\begin{funcproto}{ 
\fhead{sched.h}
\fdecl{int clone(int (*fn)(void *), void *child\_stack, int flags, void *arg,
  ...  \\
\phantom{int clone(}/* pid\_t *ptid, struct user\_desc *tls, pid\_t *ctid */ )}
\fdesc{Crea un nuovo processo o \textit{thread}.} 
}
{La funzione ritorna il \textit{Thread ID} assegnato al nuovo processo in caso
  di successo e $-1$ per un errore, nel qual caso \var{errno} assumerà uno dei
  valori: 
\begin{errlist}
    \item[\errcode{EAGAIN}] sono già in esecuzione troppi processi.
    \item[\errcode{EINVAL}] si è usata una combinazione non valida di flag o
      un valore nullo per \param{child\_stack}.
    \item[\errcode{ENOMEM}] non c'è memoria sufficiente per creare una nuova
      \texttt{task\_struct} o per copiare le parti del contesto del chiamante
      necessarie al nuovo \textit{task}.
    \item[\errcode{EPERM}] non si hanno i privilegi di amministratore
      richiesti dai flag indicati.
\end{errlist}}
\end{funcproto}

% NOTE: una pagina con la descrizione degli argomenti:
% * http://www.lindevdoc.org/wiki/Clone 

La funzione prende come primo argomento \param{fn} il puntatore alla funzione
che verrà messa in esecuzione nel nuovo processo, che può avere un unico
argomento di tipo puntatore a \ctyp{void}, il cui valore viene passato dal
terzo argomento \param{arg}. Per quanto il precedente prototipo possa
intimidire nella sua espressione, in realtà l'uso è molto semplice basterà
definire una qualunque funzione \param{fn} che restituisce un intero ed ha
come argomento un puntatore a \ctyp{void}, e \code{fn(arg)} sarà eseguita in
un nuovo processo.

Il nuovo processo resterà in esecuzione fintanto che la funzione \param{fn}
non ritorna, o esegue \func{exit} o viene terminata da un segnale. Il valore
di ritorno della funzione (o quello specificato con \func{exit}) verrà
utilizzato come stato di uscita della funzione. I tre
argomenti \param{ptid}, \param{tls} e \param{ctid} sono opzionali e sono
presenti solo a partire dal kernel 2.6 e sono stati aggiunti come supporto per
le funzioni di gestione dei \textit{thread} (la \textit{Native Thread Posix
  Library}, vedi sez.~\ref{sec:linux_ntpl}) nella \acr{glibc}, essi vengono
utilizzati soltanto se si sono specificati rispettivamente i flag
\const{CLONE\_PARENT\_SETTID}, \const{CLONE\_SETTLS} e
\const{CLONE\_CHILD\_SETTID}. 

La funzione ritorna un l'identificatore del nuovo \textit{task}, denominato
\texttt{Thread ID} (da qui in avanti \ids{TID}) il cui significato è analogo
al \ids{PID} dei normali processi e che a questo corrisponde qualora si crei
un processo ordinario e non un \textit{thread}.

Il comportamento di \func{clone}, che si riflette sulle caratteristiche del
nuovo processo da essa creato, è controllato principalmente
dall'argomento \param{flags}, che deve essere specificato come maschera
binaria, ottenuta con un OR aritmetico di una delle costanti del seguente
elenco, che illustra quelle attualmente disponibili:\footnote{si fa
  riferimento al momento della stesura di questa sezione, cioè con il kernel
  3.2.}

\begin{basedescript}{\desclabelwidth{1.5 cm}\desclabelstyle{\nextlinelabel}}

\item[\constd{CLONE\_CHILD\_CLEARTID}] cancella il valore del \textit{thread
    ID} posto all'indirizzo dato dall'argomento \param{ctid}, eseguendo un
  riattivazione del \textit{futex} (vedi sez.~\ref{sec:xxx_futex}) a
  quell'indirizzo. Questo flag viene utilizzato dalla librerie di gestione dei
  \textit{thread} ed è presente dal kernel 2.5.49.

\item[\constd{CLONE\_CHILD\_SETTID}] scrive il \ids{TID} del \textit{thread}
  figlio all'indirizzo dato dall'argomento \param{ctid}. Questo flag viene
  utilizzato dalla librerie di gestione dei \textit{thread} ed è presente dal
  kernel 2.5.49.

\item[\constd{CLONE\_FILES}] se impostato il nuovo processo condividerà con il
  padre la \textit{file descriptor table} (vedi sez.~\ref{sec:file_fd}),
  questo significa che ogni \textit{file descriptor} aperto da un processo
  verrà visto anche dall'altro e che ogni chiusura o cambiamento dei
  \textit{file descriptor flag} di un \textit{file descriptor} verrà per
  entrambi.

  Se non viene impostato il processo figlio eredita una copia della
  \textit{file descriptor table} del padre e vale la semantica classica della
  gestione dei \textit{file descriptor}, che costituisce il comportamento
  ordinario di un sistema unix-like e che illustreremo in dettaglio in
  sez.~\ref{sec:file_shared_access}.

\item[\constd{CLONE\_FS}] se questo flag viene impostato il nuovo processo
  condividerà con il padre le informazioni relative all'albero dei file, ed in
  particolare avrà la stessa radice (vedi sez.~\ref{sec:file_chroot}), la
  stessa directory di lavoro (vedi sez.~\ref{sec:file_work_dir}) e la stessa
  \textit{umask} (sez.~\ref{sec:file_perm_management}). Una modifica di una
  qualunque di queste caratteristiche in un processo, avrà effetto anche
  sull'altro. Se assente il nuovo processo riceverà una copia delle precedenti
  informazioni, che saranno così indipendenti per i due processi, come avviene
  nel comportamento ordinario di un sistema unix-like.

\item[\constd{CLONE\_IO}] se questo flag viene impostato il nuovo processo
  condividerà con il padre il contesto dell'I/O, altrimenti, come avviene nel
  comportamento ordinario con una \func{fork} otterrà un suo contesto
  dell'I/O.

  Il contesto dell'I/O viene usato dagli \textit{scheduler} di I/O (visti in
  sez.~\ref{sec:io_priority}) e se questo è lo stesso per diversi processi
  questi vengono trattati come se fossero lo stesso, condividendo il tempo per
  l'accesso al disco, e possono interscambiarsi nell'accesso a disco. L'uso di
  questo flag consente, quando più \textit{thread} eseguono dell'I/O per conto
  dello stesso processo (ad esempio con le funzioni di I/O asincrono di
  sez.~\ref{sec:file_asyncronous_io}), migliori prestazioni.

%TODO : tutti i CLONE_NEW* attengono ai namespace, ed è meglio metterli nella
%relativa sezione da creare a parte

% \item[\constd{CLONE\_NEWIPC}] è uno dei flag ad uso dei \textit{container},
%   introdotto con il kernel 2.6.19. L'uso di questo flag crea per il nuovo
%   processo un nuovo \textit{namespace} per il sistema di IPC, sia per quello
%   di SysV (vedi sez.~\ref{sec:ipc_sysv}) che, dal kernel 2.6.30, per le code
%   di messaggi POSIX (vedi sez.~\ref{sec:ipc_posix_mq}); si applica cioè a
%   tutti quegli oggetti che non vegono identificati con un \textit{pathname}
%   sull'albero dei file.

%   L'uso di questo flag richiede privilegi di amministratore (più precisamente
%   la capacità \const{CAP\_SYS\_ADMIN}) e non può essere usato in combinazione
%   con \const{CLONE\_SYSVSEM}. 

% \item[\constd{CLONE\_NEWNET}]
% \item[\constd{CLONE\_NEWNS}]
% \item[\constd{CLONE\_NEWPID}]
% \item[\constd{CLONE\_NEWUTS}]


% TODO trattare CLONE_NEWCGROUP introdotto con il kernel 4.6, vedi
% http://lwn.net/Articles/680566/ 

\item[\constd{CLONE\_PARENT}]
\item[\constd{CLONE\_PARENT\_SETTID}]
\item[\constd{CLONE\_PID}]

% TODO trattare CLONE_PIDFD introdotto con il kernel 5.2, vedi
% https://lwn.net/Articles/787963/ e anche https://lwn.net/Articles/789023/
  
\item[\constd{CLONE\_PTRACE}] se questo flag viene impostato ed il processo
  chiamante viene tracciato (vedi sez.~\ref{sec:process_ptrace}) anche il
  figlio viene tracciato. 

\item[\constd{CLONE\_SETTLS}]
\item[\constd{CLONE\_SIGHAND}]
\item[\constd{CLONE\_STOPPED}]
\item[\constd{CLONE\_SYSVSEM}]
\item[\constd{CLONE\_THREAD}]

\item[\constd{CLONE\_UNTRACED}] se questo flag viene impostato un processo non
  può più forzare \const{CLONE\_PTRACE} su questo processo.

\item[\constd{CLONE\_VFORK}] se questo flag viene impostato il chiamante viene
  fermato fintato che il figlio appena creato non rilascia la sua memoria
  virtuale con una chiamata a \func{exec} o \func{exit}, viene quindi
  replicato il comportamento di \func{vfork}.

\item[\constd{CLONE\_VM}] se questo flag viene impostato il nuovo processo
  condividerà con il padre la stessa memoria virtuale, e le scritture in
  memoria fatte da uno qualunque dei processi saranno visibili dall'altro,
  così come ogni mappatura in memoria (vedi sez.~\ref{sec:file_memory_map}). 

  Se non viene impostato il processo figlio otterrà una copia dello spazio
  degli indirizzi e si otterrà il comportamento ordinario di un processo di un
  sistema unix-like creato con la funzione \func{fork}.
\end{basedescript}



\subsection{La gestione dei \textit{namespace}}
\label{sec:process_namespaces}

\itindbeg{namespace}
Come accennato all'inizio di sez.~\ref{sec:process_clone} oltre al controllo
delle caratteristiche dei processi usate per la creazione dei \textit{thread},
l'uso di \func{clone} consente, ad uso delle nuove funzionalità di
virtualizzazione dei processi, di creare nuovi ``\textit{namespace}'' per una
serie di proprietà generali (come l'elenco dei \ids{PID}, l'albero dei file, i
\textit{mount point}, la rete, il sistema di IPC, ecc.).

L'uso dei ``\textit{namespace}'' consente creare gruppi di processi che vedono
le suddette proprietà in maniera indipendente fra loro. I processi di ciascun
gruppo vengono così eseguiti come in una sorta di spazio separato da quello
degli altri gruppi, che costituisce poi quello che viene chiamato un
\textit{container}.

\itindend{namespace}


\itindbeg{container}

\itindend{container}


%TODO sezione separata sui namespace 

%TODO trattare unshare, vedi anche http://lwn.net/Articles/532748/

%TODO: trattare la funzione setns e i namespace file descriptors (vedi
% http://lwn.net/Articles/407495/) introdotti con il kernel 3.0, altre
% informazioni su setns qui: http://lwn.net/Articles/532748/
% http://lwn.net/Articles/531498/



\section{Funzionalità avanzate e specialistiche}
\label{sec:process_special}


% TODO: trattare userfaultfd, introdotta con il 4.23, vedi
% http://man7.org/linux/man-pages/man2/userfaultfd.2.html 

% TODO: trattare process_vm_readv/process_vm_writev introdotte con il kernel
% 3.2, vedi http://man7.org/linux/man-pages/man2/process_vm_readv.2.html e i
% precedenti tentativi https://lwn.net/Articles/405346/


\subsection{La gestione delle operazioni in virgola mobile}
\label{sec:process_fenv}

Da fare.

% TODO eccezioni ed arrotondamenti per la matematica in virgola mobile 
% consultare la manpage di fenv, math_error, fpclassify, matherr, isgreater,
% isnan, nan, INFINITY


\subsection{L'accesso alle porte di I/O}
\label{sec:process_io_port}

%
% TODO l'I/O sulle porte di I/O 
% consultare le manpage di ioperm, iopl e outb
% non c'entra nulla qui, va trovato un altro posto (altri meccanismi di I/O in
% fileintro ?)

Da fare


%\subsection{La gestione di architetture a nodi multipli}
%\label{sec:process_NUMA}

% TODO trattare i cpuset, che attiene anche a NUMA, e che possono essere usati
% per associare l'uso di gruppi di processori a gruppi di processi (vedi
% manpage omonima)
% TODO trattare getcpu, che attiene anche a NUMA, mettere qui anche
% sched_getcpu, che potrebbe essere indipendente ma richiama getcpu

%TODO trattare le funzionalità per il NUMA
% vedi man numa e, mbind, get_mempolicy, set_mempolicy, 
% le pagine di manuale relative
% vedere anche dove metterle...

% \subsection{La gestione dei moduli}
% \label{sec:kernel_modules}

% da fare

%TODO trattare init_module e finit_module (quest'ultima introdotta con il
%kernel 3.8)

%%%% Altre cose di cui non è chiara la collocazione:

%TODO trattare membarrier, introdotta con il kernel 4.3
% vedi http://lwn.net/Articles/369567/ http://lwn.net/Articles/369640/
% http://git.kernel.org/cgit/linux/kernel/git/torvalds/linux.git/commit/?id=5b25b13ab08f616efd566347d809b4ece54570d1 
% vedi anche l'ulteriore opzione "expedited" introdotta con il kernel 4.14
% (https://lwn.net/Articles/728795/) 



%%% Local Variables:
%%% mode: latex
%%% TeX-master: "gapil"
%%% End:

%  LocalWords:  system call namespace prctl IRIX kernel sys int option long
%  LocalWords:  unsigned arg errno EACCESS EBADF EBUSY EFAULT EINVAL ENXIO PR
%  LocalWords:  EOPNOTSUPP EPERM CAPBSET READ capability sez tab capabilities
%  LocalWords:  bounding CAP SETPCAP DUMPABLE dump suid sgid UID DISABLE GET
%  LocalWords:  ENDIAN endianness BIG big endian LITTLE little PPC PowerPC ia
%  LocalWords:  FPEMU NOPRINT SIGFPE FPEXC point exception FP EXC SW ENABLE
%  LocalWords:  OVF overflow UND underflow RES INV DISABLED NONRECOV ASYNC AO
%  LocalWords:  KEEPCAPS pag exec SECURE KEEP CAPS securebits LOCKED NAME NUL
%  LocalWords:  char PDEATHSIG SIGCHLD fork PTRACER PID tracer process ptrace
%  LocalWords:  Security Modules ANY Yama SECCOMP secure computing seccomp vm
%  LocalWords:  STRICT strict FILTER filter SIGKILL TIMING STATISTICAL TSC fn
%  LocalWords:  TIMESTAMP timestamp Stamp Counter SIGSEGV UNALIGN SIGBUS MCE
%  LocalWords:  KILL siginfo MCEERR memory failure early kill CLEAR child cap
%  LocalWords:  reaper SUBREAPER init value result thread like flags stack FS
%  LocalWords:  race condition malloc NULL copy write glibc vsyscall sched RT
%  LocalWords:  void pid ptid struct desc tls ctid EAGAIN ENOMEM exit Posix
%  LocalWords:  Library PARENT SETTID SETTLS TID CLEARTID futex FILES table
%  LocalWords:  descriptor umask dell'I scheduler SIGHAND STOPPED SYSVSEM IPC
%  LocalWords:  UNTRACED VFORK vfork mount filesystem LSM Mandatory Access fs
%  LocalWords:  Control DAC MAC SELinux Smack Tomoyo AppArmor Discrectionary
%  LocalWords:  permitted inheritable effective fig security ADMIN forced new
%  LocalWords:  allowed dall' bound MODULE nell' all' capset sendmail SETGID
%  LocalWords:  setuid orig IMMUTABLE MKNOD OVERRIDE SEARCH CHOWN FSETID LOCK
%  LocalWords:  FOWNER saved FIXUP NOROOT AUDIT BLOCK SUSPEND SETFCAP group
%  LocalWords:  socket domain locking mlock mlockall shmctl mmap OWNER LEASE
%  LocalWords:  lease immutable append only mknod BIND SERVICE BROADCAST RAW
%  LocalWords:  broadcast multicast PACKET CHROOT chroot NICE PACCT RAWIO TTY
%  LocalWords:  accounting ioperm iopl RESOURCE CONFIG hangup vhangup SYSLOG
%  LocalWords:  WAKE ALARM CLOCK BOOTTIME REALTIME sticky NOATIME fcntl swap
%  LocalWords:  multicasting dell'IPC SysV trusted IOPRIO CLASS IDLE lookup
%  LocalWords:  scheduling dcookie NEWNS unshare nice NUMA ioctl journaling
%  LocalWords:  ext capget header hdrp datap const ESRCH SOURCE undef version
%  LocalWords:  libcap lcap obj to text dup clear DIFFERS get ncap caps ssize
%  LocalWords:  argument length all setpcap from string name proc cat capgetp
%  LocalWords:  capsetp getcap read sigreturn sysctl protected hardlinks tmp
%  LocalWords:  dell' symlink symlinks pathname TOCTTOU of

%% fileadv.tex
%%
%% Copyright (C) 2000-2018 Simone Piccardi.  Permission is granted to
%% copy, distribute and/or modify this document under the terms of the GNU Free
%% Documentation License, Version 1.1 or any later version published by the
%% Free Software Foundation; with the Invariant Sections being "Un preambolo",
%% with no Front-Cover Texts, and with no Back-Cover Texts.  A copy of the
%% license is included in the section entitled "GNU Free Documentation
%% License".
%%
\chapter{La gestione avanzata dei file}
\label{cha:file_advanced}

In questo capitolo affronteremo le tematiche relative alla gestione avanzata
dei file. Inizieremo con la trattazione delle problematiche del \textit{file
  locking} e poi prenderemo in esame le varie funzionalità avanzate che
permettono una gestione più sofisticata dell'I/O su file, a partire da quelle
che consentono di gestire l'accesso contemporaneo a più file esaminando le
varie modalità alternative di gestire l'I/O per concludere con la gestione dei
file mappati in memoria e le altre funzioni avanzate che consentono un
controllo più dettagliato delle modalità di I/O.


\section{Il \textit{file locking}}
\label{sec:file_locking}

\itindbeg{file~locking}

In sez.~\ref{sec:file_shared_access} abbiamo preso in esame le modalità in cui
un sistema unix-like gestisce l'accesso concorrente ai file da parte di
processi diversi. In quell'occasione si è visto come, con l'eccezione dei file
aperti in \textit{append mode}, quando più processi scrivono
contemporaneamente sullo stesso file non è possibile determinare la sequenza
in cui essi opereranno.

Questo causa la possibilità di una \textit{race condition}; in generale le
situazioni più comuni sono due: l'interazione fra un processo che scrive e
altri che leggono, in cui questi ultimi possono leggere informazioni scritte
solo in maniera parziale o incompleta; o quella in cui diversi processi
scrivono, mescolando in maniera imprevedibile il loro output sul file.

In tutti questi casi il \textit{file locking} è la tecnica che permette di
evitare le \textit{race condition}, attraverso una serie di funzioni che
permettono di bloccare l'accesso al file da parte di altri processi, così da
evitare le sovrapposizioni, e garantire la atomicità delle operazioni di
lettura o scrittura.


\subsection{L'\textit{advisory locking}}
\label{sec:file_record_locking}

La prima modalità di \textit{file locking} che è stata implementata nei
sistemi unix-like è quella che viene usualmente chiamata \textit{advisory
  locking},\footnote{Stevens in \cite{APUE} fa riferimento a questo argomento
  come al \textit{record locking}, dizione utilizzata anche dal manuale della
  \acr{glibc}; nelle pagine di manuale si parla di \textit{discrectionary file
    lock} per \func{fcntl} e di \textit{advisory locking} per \func{flock},
  mentre questo nome viene usato da Stevens per riferirsi al \textit{file
    locking} POSIX. Dato che la dizione \textit{record locking} è quantomeno
  ambigua, in quanto in un sistema Unix non esiste niente che possa fare
  riferimento al concetto di \textit{record}, alla fine si è scelto di
  mantenere il nome \textit{advisory locking}.} in quanto sono i singoli
processi, e non il sistema, che si incaricano di asserire e verificare se
esistono delle condizioni di blocco per l'accesso ai file. 

Questo significa che le funzioni \func{read} o \func{write} vengono eseguite
comunque e non risentono affatto della presenza di un eventuale \textit{lock};
pertanto è sempre compito dei vari processi che intendono usare il
\textit{file locking}, controllare esplicitamente lo stato dei file condivisi
prima di accedervi, utilizzando le relative funzioni.

In generale si distinguono due tipologie di \textit{file lock};\footnote{di
  seguito ci riferiremo sempre ai blocchi di accesso ai file con la
  nomenclatura inglese di \textit{file lock}, o più brevemente con
  \textit{lock}, per evitare confusioni linguistiche con il blocco di un
  processo (cioè la condizione in cui il processo viene posto in stato di
  \textit{sleep}).} la prima è il cosiddetto \textit{shared lock}, detto anche
\textit{read lock} in quanto serve a bloccare l'accesso in scrittura su un
file affinché il suo contenuto non venga modificato mentre lo si legge. Si
parla appunto di \textsl{blocco condiviso} in quanto più processi possono
richiedere contemporaneamente uno \textit{shared lock} su un file per
proteggere il loro accesso in lettura.

La seconda tipologia è il cosiddetto \textit{exclusive lock}, detto anche
\textit{write lock} in quanto serve a bloccare l'accesso su un file (sia in
lettura che in scrittura) da parte di altri processi mentre lo si sta
scrivendo. Si parla di \textsl{blocco esclusivo} appunto perché un solo
processo alla volta può richiedere un \textit{exclusive lock} su un file per
proteggere il suo accesso in scrittura.

In Linux sono disponibili due interfacce per utilizzare l'\textit{advisory
  locking}, la prima è quella derivata da BSD, che è basata sulla funzione
\func{flock}, la seconda è quella recepita dallo standard POSIX.1 (che è
derivata dall'interfaccia usata in System V), che è basata sulla funzione
\func{fcntl}.  I \textit{file lock} sono implementati in maniera completamente
indipendente nelle due interfacce (in realtà con Linux questo avviene solo
dalla serie 2.0 dei kernel) che pertanto possono coesistere senza
interferenze.

Entrambe le interfacce prevedono la stessa procedura di funzionamento: si
inizia sempre con il richiedere l'opportuno \textit{file lock} (un
\textit{exclusive lock} per una scrittura, uno \textit{shared lock} per una
lettura) prima di eseguire l'accesso ad un file.  Se il blocco viene acquisito
il processo prosegue l'esecuzione, altrimenti (a meno di non aver richiesto un
comportamento non bloccante) viene posto in stato di \textit{sleep}. Una volta
finite le operazioni sul file si deve provvedere a rimuovere il blocco.

La situazione delle varie possibilità che si possono verificare è riassunta in
tab.~\ref{tab:file_file_lock}, dove si sono riportati, a seconda delle varie
tipologie di blocco già presenti su un file, il risultato che si avrebbe in
corrispondenza di una ulteriore richiesta da parte di un processo di un blocco
nelle due tipologie di \textit{file lock} menzionate, con un successo o meno
della richiesta.

\begin{table}[htb]
  \centering
  \footnotesize
   \begin{tabular}[c]{|l|c|c|c|}
    \hline
    \textbf{Richiesta} & \multicolumn{3}{|c|}{\textbf{Stato del file}}\\
    \cline{2-4}
                &Nessun \textit{lock}&\textit{Read lock}&\textit{Write lock}\\
    \hline
    \hline
    \textit{Read lock} & esecuzione & esecuzione & blocco \\
    \textit{Write lock}& esecuzione & blocco & blocco \\
    \hline    
  \end{tabular}
  \caption{Tipologie di \textit{file locking}.}
  \label{tab:file_file_lock}
\end{table}

Si tenga presente infine che il controllo di accesso e la gestione dei
permessi viene effettuata quando si apre un file, l'unico controllo residuo
che si può avere riguardo il \textit{file locking} è che il tipo di blocco che
si vuole ottenere su un file deve essere compatibile con le modalità di
apertura dello stesso (in lettura per un \textit{read lock} e in scrittura per
un \textit{write lock}).

%%  Si ricordi che
%% la condizione per acquisire uno \textit{shared lock} è che il file non abbia
%% già un \textit{exclusive lock} attivo, mentre per acquisire un
%% \textit{exclusive lock} non deve essere presente nessun tipo di blocco.


\subsection{La funzione \func{flock}} 
\label{sec:file_flock}

La prima interfaccia per il \textit{file locking}, quella derivata da BSD,
permette di eseguire un blocco solo su un intero file; la funzione di sistema
usata per richiedere e rimuovere un \textit{file lock} è \funcd{flock}, ed il
suo prototipo è:

\begin{funcproto}{
\fhead{sys/file.h}
\fdecl{int flock(int fd, int operation)}
\fdesc{Applica o rimuove un \textit{file lock}.} 
}

{La funzione ritorna $0$ in caso di successo e $-1$ per un errore, nel qual
  caso \var{errno} assumerà uno dei valori: 
  \begin{errlist}
  \item[\errcode{EINTR}] la funzione è stata interrotta da un segnale
    nell'attesa dell'acquisizione di un \textit{file lock}.
  \item[\errcode{EINVAL}] si è specificato un valore non valido
    per \param{operation}.
  \item[\errcode{ENOLCK}] il kernel non ha memoria sufficiente per gestire il
    \textit{file lock}.
  \item[\errcode{EWOULDBLOCK}] il file ha già un blocco attivo, e si è
    specificato \const{LOCK\_NB}.
  \end{errlist}
  ed inoltre \errval{EBADF} nel suo significato generico.
}
\end{funcproto}

La funzione può essere usata per acquisire o rilasciare un \textit{file lock}
a seconda di quanto specificato tramite il valore dell'argomento
\param{operation}; questo viene interpretato come maschera binaria, e deve
essere passato costruendo il valore con un OR aritmetico delle costanti
riportate in tab.~\ref{tab:file_flock_operation}.

\begin{table}[htb]
  \centering
  \footnotesize
  \begin{tabular}[c]{|l|p{6cm}|}
    \hline
    \textbf{Valore} & \textbf{Significato} \\
    \hline
    \hline
    \constd{LOCK\_SH} & Richiede uno \textit{shared lock} sul file.\\ 
    \constd{LOCK\_EX} & Richiede un \textit{esclusive lock} sul file.\\
    \constd{LOCK\_UN} & Rilascia il \textit{file lock}.\\
    \constd{LOCK\_NB} & Impedisce che la funzione si blocchi nella
                        richiesta di un \textit{file lock}.\\
    \hline    
  \end{tabular}
  \caption{Valori dell'argomento \param{operation} di \func{flock}.}
  \label{tab:file_flock_operation}
\end{table}

I primi due valori, \const{LOCK\_SH} e \const{LOCK\_EX} permettono di
richiedere un \textit{file lock} rispettivamente condiviso o esclusivo, ed
ovviamente non possono essere usati insieme. Se con essi si specifica anche
\const{LOCK\_NB} la funzione non si bloccherà qualora il \textit{file lock}
non possa essere acquisito, ma ritornerà subito con un errore di
\errcode{EWOULDBLOCK}. Per rilasciare un \textit{file lock} si dovrà invece
usare direttamente \const{LOCK\_UN}.

Si tenga presente che non esiste una modalità per eseguire atomicamente un
cambiamento del tipo di blocco (da \textit{shared lock} a \textit{esclusive
  lock}), il blocco deve essere prima rilasciato e poi richiesto, ed è sempre
possibile che nel frattempo abbia successo un'altra richiesta pendente,
facendo fallire la riacquisizione.

Si tenga presente infine che \func{flock} non è supportata per i file
mantenuti su NFS, in questo caso, se si ha la necessità di utilizzare il
\textit{file locking}, occorre usare l'interfaccia del \textit{file locking}
POSIX basata su \func{fcntl} che è in grado di funzionare anche attraverso
NFS, a condizione ovviamente che sia il client che il server supportino questa
funzionalità.

La semantica del \textit{file locking} di BSD inoltre è diversa da quella del
\textit{file locking} POSIX, in particolare per quanto riguarda il
comportamento dei \textit{file lock} nei confronti delle due funzioni
\func{dup} e \func{fork}.  Per capire queste differenze occorre descrivere con
maggiore dettaglio come viene realizzato dal kernel il \textit{file locking}
per entrambe le interfacce.

In fig.~\ref{fig:file_flock_struct} si è riportato uno schema essenziale
dell'implementazione del \textit{file locking} in stile BSD su Linux. Il punto
fondamentale da capire è che un \textit{file lock}, qualunque sia
l'interfaccia che si usa, anche se richiesto attraverso un file descriptor,
agisce sempre su di un file; perciò le informazioni relative agli eventuali
\textit{file lock} sono mantenute dal kernel a livello di \textit{inode}, dato
che questo è l'unico riferimento in comune che possono avere due processi
diversi che aprono lo stesso file.

In particolare, come accennato in fig.~\ref{fig:file_flock_struct}, i
\textit{file lock} sono mantenuti in una \textit{linked list} di strutture
\kstructd{file\_lock}. La lista è referenziata dall'indirizzo di partenza
mantenuto dal campo \var{i\_flock} della struttura \kstruct{inode} (per le
definizioni esatte si faccia riferimento al file \file{include/linux/fs.h} nei
sorgenti del kernel).  Un bit del campo \var{fl\_flags} di specifica se si
tratta di un lock in semantica BSD (\constd{FL\_FLOCK}) o POSIX
(\constd{FL\_POSIX}) o un \textit{file lease} (\constd{FL\_LEASE}, vedi
sez.~\ref{sec:file_asyncronous_lease}).

\begin{figure}[!htb]
  \centering
  \includegraphics[width=12cm]{img/file_flock}
  \caption{Schema dell'architettura del \textit{file locking}, nel caso
    particolare del suo utilizzo da parte dalla funzione \func{flock}.}
  \label{fig:file_flock_struct}
\end{figure}

La richiesta di un \textit{file lock} prevede una scansione della lista per
determinare se l'acquisizione è possibile, ed in caso positivo l'aggiunta di
un nuovo elemento (cioè l'aggiunta di una nuova struttura
\kstruct{file\_lock}).  Nel caso dei blocchi creati con \func{flock} la
semantica della funzione prevede che sia \func{dup} che \func{fork} non creino
ulteriori istanze di un \textit{file lock} quanto piuttosto degli ulteriori
riferimenti allo stesso. Questo viene realizzato dal kernel secondo lo schema
di fig.~\ref{fig:file_flock_struct}, associando ad ogni nuovo \textit{file
  lock} un puntatore alla voce nella \textit{file table} da cui si è richiesto
il blocco, che così ne identifica il titolare. Il puntatore è mantenuto nel
campo \var{fl\_file} di \kstruct{file\_lock}, e viene utilizzato solo per i
\textit{file lock} creati con la semantica BSD.

Questa struttura prevede che, quando si richiede la rimozione di un
\textit{file lock}, il kernel acconsenta solo se la richiesta proviene da un
file descriptor che fa riferimento ad una voce nella \textit{file table}
corrispondente a quella registrata nel blocco.  Allora se ricordiamo quanto
visto in sez.~\ref{sec:file_dup} e sez.~\ref{sec:file_shared_access}, e cioè
che i file descriptor duplicati e quelli ereditati in un processo figlio
puntano sempre alla stessa voce nella \textit{file table}, si può capire
immediatamente quali sono le conseguenze nei confronti delle funzioni
\func{dup} e \func{fork}.

Sarà così possibile rimuovere un \textit{file lock} attraverso uno qualunque
dei file descriptor che fanno riferimento alla stessa voce nella \textit{file
  table}, anche se questo è diverso da quello con cui lo si è
creato,\footnote{attenzione, questo non vale se il file descriptor fa
  riferimento allo stesso file, ma attraverso una voce diversa della
  \textit{file table}, come accade tutte le volte che si apre più volte lo
  stesso file.} o se si esegue la rimozione in un processo figlio. Inoltre una
volta tolto un \textit{file lock} su un file, la rimozione avrà effetto su
tutti i file descriptor che condividono la stessa voce nella \textit{file
  table}, e quindi, nel caso di file descriptor ereditati attraverso una
\func{fork}, anche per processi diversi.

Infine, per evitare che la terminazione imprevista di un processo lasci attivi
dei \textit{file lock}, quando un file viene chiuso il kernel provvede anche a
rimuovere tutti i blocchi ad esso associati. Anche in questo caso occorre
tenere presente cosa succede quando si hanno file descriptor duplicati; in tal
caso infatti il file non verrà effettivamente chiuso (ed il blocco rimosso)
fintanto che non viene rilasciata la relativa voce nella \textit{file table};
e questo avverrà solo quando tutti i file descriptor che fanno riferimento
alla stessa voce sono stati chiusi.  Quindi, nel caso ci siano duplicati o
processi figli che mantengono ancora aperto un file descriptor, il
\textit{file lock} non viene rilasciato.
 

\subsection{Il \textit{file locking} POSIX}
\label{sec:file_posix_lock}

La seconda interfaccia per l'\textit{advisory locking} disponibile in Linux è
quella standardizzata da POSIX, basata sulla funzione di sistema
\func{fcntl}. Abbiamo già trattato questa funzione nelle sue molteplici
possibilità di utilizzo in sez.~\ref{sec:file_fcntl_ioctl}. Quando la si
impiega per il \textit{file locking} essa viene usata solo secondo il seguente
prototipo:

\begin{funcproto}{
\fhead{fcntl.h}
\fdecl{int fcntl(int fd, int cmd, struct flock *lock)}
\fdesc{Applica o rimuove un \textit{file lock}.} 
}

{La funzione ritorna $0$ in caso di successo e $-1$ per un errore, nel qual
  caso \var{errno} assumerà uno dei valori: 
  \begin{errlist}
    \item[\errcode{EACCES}] l'operazione è proibita per la presenza di
      \textit{file lock} da parte di altri processi.
    \item[\errcode{EDEADLK}] si è richiesto un \textit{lock} su una regione
      bloccata da un altro processo che è a sua volta in attesa dello sblocco
      di un \textit{lock} mantenuto dal processo corrente; si avrebbe pertanto
      un \textit{deadlock}. Non è garantito che il sistema riconosca sempre
      questa situazione.
    \item[\errcode{EINTR}] la funzione è stata interrotta da un segnale prima
      di poter acquisire un \textit{file lock}.
    \item[\errcode{ENOLCK}] il sistema non ha le risorse per il blocco: ci
      sono troppi segmenti di \textit{lock} aperti, si è esaurita la tabella
      dei \textit{file lock}, o il protocollo per il blocco remoto è fallito.
  \end{errlist}
  ed inoltre \errval{EBADF}, \errval{EFAULT} nel loro significato generico.}
\end{funcproto}

Al contrario di quanto avviene con l'interfaccia basata su \func{flock} con
\func{fcntl} è possibile bloccare anche delle singole sezioni di un file, fino
al singolo byte. Inoltre la funzione permette di ottenere alcune informazioni
relative agli eventuali blocchi preesistenti.  Per poter fare tutto questo la
funzione utilizza come terzo argomento una apposita struttura \struct{flock}
(la cui definizione è riportata in fig.~\ref{fig:struct_flock}) nella quale
inserire tutti i dati relativi ad un determinato blocco. Si tenga presente poi
che un \textit{file lock} fa sempre riferimento ad una regione, per cui si
potrà avere un conflitto anche se c'è soltanto una sovrapposizione parziale
con un'altra regione bloccata.

\begin{figure}[!htb]
  \footnotesize \centering
  \begin{minipage}[c]{0.90\textwidth}
    \includestruct{listati/flock.h}
  \end{minipage} 
  \normalsize 
  \caption{La struttura \structd{flock}, usata da \func{fcntl} per il
    \textit{file locking}.}
  \label{fig:struct_flock}
\end{figure}

I primi tre campi della struttura, \var{l\_whence}, \var{l\_start} e
\var{l\_len}, servono a specificare la sezione del file a cui fa riferimento
il blocco: \var{l\_start} specifica il byte di partenza, \var{l\_len} la
lunghezza della sezione e infine \var{l\_whence} imposta il riferimento da cui
contare \var{l\_start}. Il valore di \var{l\_whence} segue la stessa semantica
dell'omonimo argomento di \func{lseek}, coi tre possibili valori
\const{SEEK\_SET}, \const{SEEK\_CUR} e \const{SEEK\_END}, (si vedano le
relative descrizioni in tab.~\ref{tab:lseek_whence_values}).

Si tenga presente che un \textit{file lock} può essere richiesto anche per una
regione al di là della corrente fine del file, così che una eventuale
estensione dello stesso resti coperta dal blocco. Inoltre se si specifica un
valore nullo per \var{l\_len} il blocco si considera esteso fino alla
dimensione massima del file; in questo modo è possibile bloccare una qualunque
regione a partire da un certo punto fino alla fine del file, coprendo
automaticamente quanto eventualmente aggiunto in coda allo stesso.

Lo standard POSIX non richiede che \var{l\_len} sia positivo, ed a partire dal
kernel 2.4.21 è possibile anche indicare valori di \var{l\_len} negativi, in
tal caso l'intervallo coperto va da \var{l\_start}$+$\var{l\_len} a
\var{l\_start}$-1$, mentre per un valore positivo l'intervallo va da
\var{l\_start} a \var{l\_start}$+$\var{l\_len}$-1$. Si può però usare un
valore negativo soltanto se l'inizio della regione indicata non cade prima
dell'inizio del file, mentre come accennato con un valore positivo  si
può anche indicare una regione che eccede la dimensione corrente del file.

Il tipo di \textit{file lock} richiesto viene specificato dal campo
\var{l\_type}, esso può assumere i tre valori definiti dalle costanti
riportate in tab.~\ref{tab:file_flock_type}, che permettono di richiedere
rispettivamente uno \textit{shared lock}, un \textit{esclusive lock}, e la
rimozione di un blocco precedentemente acquisito. Infine il campo \var{l\_pid}
viene usato solo in caso di lettura, quando si chiama \func{fcntl} con
\const{F\_GETLK}, e riporta il \ids{PID} del processo che detiene il
\textit{file lock}.

\begin{table}[htb]
  \centering
  \footnotesize
  \begin{tabular}[c]{|l|l|}
    \hline
    \textbf{Valore} & \textbf{Significato} \\
    \hline
    \hline
    \constd{F\_RDLCK} & Richiede un blocco condiviso (\textit{read lock}).\\
    \constd{F\_WRLCK} & Richiede un blocco esclusivo (\textit{write lock}).\\
    \constd{F\_UNLCK} & Richiede l'eliminazione di un \textit{file lock}.\\
    \hline    
  \end{tabular}
  \caption{Valori possibili per il campo \var{l\_type} di \struct{flock}.}
  \label{tab:file_flock_type}
\end{table}

Oltre a quanto richiesto tramite i campi di \struct{flock}, l'operazione
effettivamente svolta dalla funzione è stabilita dal valore dall'argomento
\param{cmd} che, come già riportato in sez.~\ref{sec:file_fcntl_ioctl},
specifica l'azione da compiere; i valori utilizzabili relativi al \textit{file
  locking} sono tre:
\begin{basedescript}{\desclabelwidth{2.0cm}}
\item[\constd{F\_GETLK}] verifica se il \textit{file lock} specificato dalla
  struttura puntata da \param{lock} può essere acquisito: in caso negativo
  sovrascrive la struttura \param{flock} con i valori relativi al blocco già
  esistente che ne blocca l'acquisizione, altrimenti si limita a impostarne il
  campo \var{l\_type} con il valore \const{F\_UNLCK}.
\item[\constd{F\_SETLK}] se il campo \var{l\_type} della struttura puntata da
  \param{lock} è \const{F\_RDLCK} o \const{F\_WRLCK} richiede il
  corrispondente \textit{file lock}, se è \const{F\_UNLCK} lo rilascia; nel
  caso la richiesta non possa essere soddisfatta a causa di un blocco
  preesistente la funzione ritorna immediatamente con un errore di
  \errcode{EACCES} o di \errcode{EAGAIN}.
\item[\constd{F\_SETLKW}] è identica a \const{F\_SETLK}, ma se la richiesta di
  non può essere soddisfatta per la presenza di un altro blocco, mette il
  processo in stato di attesa fintanto che il blocco precedente non viene
  rilasciato; se l'attesa viene interrotta da un segnale la funzione ritorna
  con un errore di \errcode{EINTR}.
\end{basedescript}

Si noti che per quanto detto il comando \const{F\_GETLK} non serve a rilevare
una presenza generica di blocco su un file, perché se ne esistono altri
compatibili con quello richiesto, la funzione ritorna comunque impostando
\var{l\_type} a \const{F\_UNLCK}.  Inoltre a seconda del valore di
\var{l\_type} si potrà controllare o l'esistenza di un qualunque tipo di
blocco (se è \const{F\_WRLCK}) o di \textit{write lock} (se è
\const{F\_RDLCK}). Si consideri poi che può esserci più di un blocco che
impedisce l'acquisizione di quello richiesto (basta che le regioni si
sovrappongano), ma la funzione ne riporterà sempre soltanto uno, impostando
\var{l\_whence} a \const{SEEK\_SET} ed i valori \var{l\_start} e \var{l\_len}
per indicare quale è la regione bloccata.

Infine si tenga presente che effettuare un controllo con il comando
\const{F\_GETLK} e poi tentare l'acquisizione con \const{F\_SETLK} non è una
operazione atomica (un altro processo potrebbe acquisire un blocco fra le due
chiamate) per cui si deve sempre verificare il codice di ritorno di
\func{fcntl}\footnote{controllare il codice di ritorno delle funzioni invocate
  è comunque una buona norma di programmazione, che permette di evitare un
  sacco di errori difficili da tracciare proprio perché non vengono rilevati.}
quando la si invoca con \const{F\_SETLK}, per controllare che il blocco sia
stato effettivamente acquisito.

\begin{figure}[!htb]
  \centering \includegraphics[width=9cm]{img/file_lock_dead}
  \caption{Schema di una situazione di \textit{deadlock}.}
  \label{fig:file_flock_dead}
\end{figure}

Non operando a livello di interi file, il \textit{file locking} POSIX
introduce un'ulteriore complicazione; consideriamo la situazione illustrata in
fig.~\ref{fig:file_flock_dead}, in cui il processo A blocca la regione 1 e il
processo B la regione 2. Supponiamo che successivamente il processo A richieda
un lock sulla regione 2 che non può essere acquisito per il preesistente lock
del processo 2; il processo 1 si bloccherà fintanto che il processo 2 non
rilasci il blocco. Ma cosa accade se il processo 2 nel frattempo tenta a sua
volta di ottenere un lock sulla regione A? Questa è una tipica situazione che
porta ad un \textit{deadlock}, dato che a quel punto anche il processo 2 si
bloccherebbe, e niente potrebbe sbloccare l'altro processo.  Per questo motivo
il kernel si incarica di rilevare situazioni di questo tipo, ed impedirle
restituendo un errore di \errcode{EDEADLK} alla funzione che cerca di
acquisire un blocco che porterebbe ad un \textit{deadlock}.

Per capire meglio il funzionamento del \textit{file locking} in semantica
POSIX (che differisce alquanto rispetto da quello di BSD, visto
sez.~\ref{sec:file_flock}) esaminiamo più in dettaglio come viene gestito dal
kernel. Lo schema delle strutture utilizzate è riportato in
fig.~\ref{fig:file_posix_lock}; come si vede esso è molto simile all'analogo
di fig.~\ref{fig:file_flock_struct}. In questo caso nella figura si sono
evidenziati solo i campi di \kstructd{file\_lock} significativi per la
semantica POSIX, in particolare adesso ciascuna struttura contiene, oltre al
\ids{PID} del processo in \var{fl\_pid}, la sezione di file che viene bloccata
grazie ai campi \var{fl\_start} e \var{fl\_end}.  La struttura è comunque la
stessa, solo che in questo caso nel campo \var{fl\_flags} è impostato il bit
\const{FL\_POSIX} ed il campo \var{fl\_file} non viene usato. Il blocco è
sempre associato all'\textit{inode}, solo che in questo caso la titolarità non
viene identificata con il riferimento ad una voce nella \textit{file table},
ma con il valore del \ids{PID} del processo.

\begin{figure}[!htb]
  \centering \includegraphics[width=12cm]{img/file_posix_lock}
  \caption{Schema dell'architettura del \textit{file locking}, nel caso
    particolare del suo utilizzo secondo l'interfaccia standard POSIX.}
  \label{fig:file_posix_lock}
\end{figure}

Quando si richiede un \textit{file lock} il kernel effettua una scansione di
tutti i blocchi presenti sul file\footnote{scandisce cioè la \textit{linked
    list} delle strutture \kstruct{file\_lock}, scartando automaticamente
  quelle per cui \var{fl\_flags} non è \const{FL\_POSIX}, così che le due
  interfacce restano ben separate.}  per verificare se la regione richiesta
non si sovrappone ad una già bloccata, in caso affermativo decide in base al
tipo di blocco, in caso negativo il nuovo blocco viene comunque acquisito ed
aggiunto alla lista.

Nel caso di rimozione invece questa viene effettuata controllando che il
\ids{PID} del processo richiedente corrisponda a quello contenuto nel blocco.
Questa diversa modalità ha delle conseguenze precise riguardo il comportamento
dei \textit{file lock} POSIX. La prima conseguenza è che un \textit{file lock}
POSIX non viene mai ereditato attraverso una \func{fork}, dato che il processo
figlio avrà un \ids{PID} diverso, mentre passa indenne attraverso una
\func{exec} in quanto il \ids{PID} resta lo stesso.  Questo comporta che, al
contrario di quanto avveniva con la semantica BSD, quando un processo termina
tutti i \textit{file lock} da esso detenuti vengono immediatamente rilasciati.

La seconda conseguenza è che qualunque file descriptor che faccia riferimento
allo stesso file (che sia stato ottenuto con una \func{dup} o con una
\func{open} in questo caso non fa differenza) può essere usato per rimuovere
un blocco, dato che quello che conta è solo il \ids{PID} del processo. Da
questo deriva una ulteriore sottile differenza di comportamento: dato che alla
chiusura di un file i blocchi ad esso associati vengono rimossi, nella
semantica POSIX basterà chiudere un file descriptor qualunque per cancellare
tutti i blocchi relativi al file cui esso faceva riferimento, anche se questi
fossero stati creati usando altri file descriptor che restano aperti.

Dato che il controllo sull'accesso ai blocchi viene eseguito sulla base del
\ids{PID} del processo, possiamo anche prendere in considerazione un altro
degli aspetti meno chiari di questa interfaccia e cioè cosa succede quando si
richiedono dei blocchi su regioni che si sovrappongono fra loro all'interno
stesso processo. Siccome il controllo, come nel caso della rimozione, si basa
solo sul \ids{PID} del processo che chiama la funzione, queste richieste
avranno sempre successo.  Nel caso della semantica BSD, essendo i lock
relativi a tutto un file e non accumulandosi,\footnote{questa ultima
  caratteristica è vera in generale, se cioè si richiede più volte lo stesso
  \textit{file lock}, o più blocchi sulla stessa sezione di file, le richieste
  non si cumulano e basta una sola richiesta di rilascio per cancellare il
  blocco.}  la cosa non ha alcun effetto; la funzione ritorna con successo,
senza che il kernel debba modificare la lista dei \textit{file lock}.

Con i \textit{file lock} POSIX invece si possono avere una serie di situazioni
diverse: ad esempio è possibile rimuovere con una sola chiamata più
\textit{file lock} distinti (indicando in una regione che si sovrapponga
completamente a quelle di questi ultimi), o rimuovere solo una parte di un
blocco preesistente (indicando una regione contenuta in quella di un altro
blocco), creando un buco, o coprire con un nuovo blocco altri \textit{file
  lock} già ottenuti, e così via, a secondo di come si sovrappongono le
regioni richieste e del tipo di operazione richiesta.

Il comportamento seguito in questo caso è che la funzione ha successo ed
esegue l'operazione richiesta sulla regione indicata; è compito del kernel
preoccuparsi di accorpare o dividere le voci nella lista dei \textit{file
  lock} per far sì che le regioni bloccate da essa risultanti siano coerenti
con quanto necessario a soddisfare l'operazione richiesta.

\begin{figure}[!htbp]
  \footnotesize \centering
  \begin{minipage}[c]{\codesamplewidth}
    \includecodesample{listati/Flock.c}
  \end{minipage}
  \normalsize 
  \caption{Sezione principale del codice del programma \file{Flock.c}.}
  \label{fig:file_flock_code}
\end{figure}

Per fare qualche esempio sul \textit{file locking} si è scritto un programma che
permette di bloccare una sezione di un file usando la semantica POSIX, o un
intero file usando la semantica BSD; in fig.~\ref{fig:file_flock_code} è
riportata il corpo principale del codice del programma, (il testo completo è
allegato nella directory dei sorgenti, nel file \texttt{Flock.c}).

La sezione relativa alla gestione delle opzioni al solito si è omessa, come la
funzione che stampa le istruzioni per l'uso del programma, essa si cura di
impostare le variabili \var{type}, \var{start} e \var{len}; queste ultime due
vengono inizializzate al valore numerico fornito rispettivamente tramite gli
switch \code{-s} e \cmd{-l}, mentre il valore della prima viene impostato con
le opzioni \cmd{-w} e \cmd{-r} si richiede rispettivamente o un \textit{write
  lock} o \textit{read lock} (i due valori sono esclusivi, la variabile
assumerà quello che si è specificato per ultimo). Oltre a queste tre vengono
pure impostate la variabile \var{bsd}, che abilita la semantica omonima quando
si invoca l'opzione \cmd{-f} (il valore preimpostato è nullo, ad indicare la
semantica POSIX), e la variabile \var{cmd} che specifica la modalità di
richiesta del \textit{file lock} (bloccante o meno), a seconda dell'opzione
\cmd{-b}.

Il programma inizia col controllare (\texttt{\small 11-14}) che venga passato
un argomento (il file da bloccare), che sia stato scelto (\texttt{\small
  15-18}) il tipo di blocco, dopo di che apre (\texttt{\small 19}) il file,
uscendo (\texttt{\small 20-23}) in caso di errore. A questo punto il
comportamento dipende dalla semantica scelta; nel caso sia BSD occorre
reimpostare il valore di \var{cmd} per l'uso con \func{flock}; infatti il
valore preimpostato fa riferimento alla semantica POSIX e vale rispettivamente
\const{F\_SETLKW} o \const{F\_SETLK} a seconda che si sia impostato o meno la
modalità bloccante.

Nel caso si sia scelta la semantica BSD (\texttt{\small 25-34}) prima si
controlla (\texttt{\small 27-31}) il valore di \var{cmd} per determinare se
si vuole effettuare una chiamata bloccante o meno, reimpostandone il valore
opportunamente, dopo di che a seconda del tipo di blocco al valore viene
aggiunta la relativa opzione, con un OR aritmetico, dato che \func{flock}
vuole un argomento \param{operation} in forma di maschera binaria.  Nel caso
invece che si sia scelta la semantica POSIX le operazioni sono molto più
immediate si prepara (\texttt{\small 36-40}) la struttura per il lock, e lo
si esegue (\texttt{\small 41}).

In entrambi i casi dopo aver richiesto il blocco viene controllato il
risultato uscendo (\texttt{\small 44-46}) in caso di errore, o stampando un
messaggio (\texttt{\small 47-49}) in caso di successo. Infine il programma si
pone in attesa (\texttt{\small 50}) finché un segnale (ad esempio un \cmd{C-c}
dato da tastiera) non lo interrompa; in questo caso il programma termina, e
tutti i blocchi vengono rilasciati.

Con il programma possiamo fare varie verifiche sul funzionamento del
\textit{file locking}; cominciamo con l'eseguire un \textit{read lock} su un
file, ad esempio usando all'interno di un terminale il seguente comando:

\begin{Console}
[piccardi@gont sources]$ \textbf{./flock -r Flock.c}
Lock acquired
\end{Console}
%$
il programma segnalerà di aver acquisito un blocco e si bloccherà; in questo
caso si è usato il \textit{file locking} POSIX e non avendo specificato niente
riguardo alla sezione che si vuole bloccare sono stati usati i valori
preimpostati che bloccano tutto il file. A questo punto se proviamo ad
eseguire lo stesso comando in un altro terminale, e avremo lo stesso
risultato. Se invece proviamo ad eseguire un \textit{write lock} avremo:

\begin{Console}
[piccardi@gont sources]$ \textbf{./flock -w Flock.c}
Failed lock: Resource temporarily unavailable
\end{Console}
%$
come ci aspettiamo il programma terminerà segnalando l'indisponibilità del
blocco, dato che il file è bloccato dal precedente \textit{read lock}. Si noti
che il risultato è lo stesso anche se si richiede il blocco su una sola parte
del file con il comando:

\begin{Console}
[piccardi@gont sources]$ \textbf{./flock -w -s0 -l10 Flock.c}
Failed lock: Resource temporarily unavailable
\end{Console}
%$
se invece blocchiamo una regione con: 

\begin{Console}
[piccardi@gont sources]$ \textbf{./flock -r -s0 -l10 Flock.c}
Lock acquired
\end{Console}
%$
una volta che riproviamo ad acquisire il \textit{write lock} i risultati
dipenderanno dalla regione richiesta; ad esempio nel caso in cui le due
regioni si sovrappongono avremo che:

\begin{Console}
[piccardi@gont sources]$ \textbf{./flock -w -s5 -l15  Flock.c}
Failed lock: Resource temporarily unavailable
\end{Console}
%$
ed il blocco viene rifiutato, ma se invece si richiede una regione distinta
avremo che:

\begin{Console}
[piccardi@gont sources]$ \textbf{./flock -w -s11 -l15  Flock.c}
Lock acquired
\end{Console}
%$
ed il blocco viene acquisito. Se a questo punto si prova ad eseguire un
\textit{read lock} che comprende la nuova regione bloccata in scrittura:

\begin{Console}
[piccardi@gont sources]$ \textbf{./flock -r -s10 -l20 Flock.c}
Failed lock: Resource temporarily unavailable
\end{Console}
%$
come ci aspettiamo questo non sarà consentito.

Il programma di norma esegue il tentativo di acquisire il lock in modalità non
bloccante, se però usiamo l'opzione \cmd{-b} possiamo impostare la modalità
bloccante, riproviamo allora a ripetere le prove precedenti con questa
opzione:

\begin{Console}
[piccardi@gont sources]$ \textbf{./flock -r -b -s0 -l10 Flock.c} Lock acquired
\end{Console}
%$
il primo comando acquisisce subito un \textit{read lock}, e quindi non cambia
nulla, ma se proviamo adesso a richiedere un \textit{write lock} che non potrà
essere acquisito otterremo:

\begin{Console}
[piccardi@gont sources]$ \textbf{./flock -w -s0 -l10 Flock.c}
\end{Console}
%$
il programma cioè si bloccherà nella chiamata a \func{fcntl}; se a questo
punto rilasciamo il precedente blocco (terminando il primo comando un
\texttt{C-c} sul terminale) potremo verificare che sull'altro terminale il
blocco viene acquisito, con la comparsa di una nuova riga:

\begin{Console}
[piccardi@gont sources]$ \textbf{./flock -w -s0 -l10 Flock.c}
Lock acquired
\end{Console}
%$

Un'altra cosa che si può controllare con il nostro programma è l'interazione
fra i due tipi di blocco; se ripartiamo dal primo comando con cui si è
ottenuto un blocco in lettura sull'intero file, possiamo verificare cosa
succede quando si cerca di ottenere un blocco in scrittura con la semantica
BSD:

\begin{Console}
[root@gont sources]# \textbf{./flock -f -w Flock.c}
Lock acquired
\end{Console}
%$
che ci mostra come i due tipi di blocco siano assolutamente indipendenti; per
questo motivo occorre sempre tenere presente quale, fra le due semantiche
disponibili, stanno usando i programmi con cui si interagisce, dato che i
blocchi applicati con l'altra non avrebbero nessun effetto.

% \subsection{La funzione \func{lockf}}
% \label{sec:file_lockf}

Abbiamo visto come l'interfaccia POSIX per il \textit{file locking} sia molto
più potente e flessibile di quella di BSD, questo comporta anche una maggiore
complessità per via delle varie opzioni da passare a \func{fcntl}. Per questo
motivo è disponibile anche una interfaccia semplificata che utilizza la
funzione \funcd{lockf},\footnote{la funzione è ripresa da System V e per
  poterla utilizzare è richiesta che siano definite le opportune macro, una
  fra \macro{\_BSD\_SOURCE} o \macro{\_SVID\_SOURCE}, oppure
  \macro{\_XOPEN\_SOURCE} ad un valore di almeno 500, oppure
  \macro{\_XOPEN\_SOURCE} e \macro{\_XOPEN\_SOURCE\_EXTENDED}.} il cui
prototipo è:

\begin{funcproto}{
\fhead{unistd.h}
\fdecl{int lockf(int fd, int cmd, off\_t len)}
\fdesc{Applica, controlla o rimuove un \textit{file lock}.} 
}

{La funzione ritorna $0$ in caso di successo e $-1$ per un errore, nel qual
  caso \var{errno} assumerà uno dei valori: 
  \begin{errlist}
  \item[\errcode{EAGAIN}] il file è bloccato, e si sono richiesti
    \const{F\_TLOCK} o \const{F\_TEST} (in alcuni casi può dare anche
    \errcode{EACCESS}.
  \item[\errcode{EBADF}] \param{fd} non è un file descriptor aperto o si sono
    richiesti \const{F\_LOCK} o \const{F\_TLOCK} ma il file non è scrivibile.
  \item[\errcode{EINVAL}] si è usato un valore non valido per \param{cmd}.
  \end{errlist}
  ed inoltre \errcode{EDEADLK} e \errcode{ENOLCK} con lo stesso significato
  che hanno con \func{fcntl}.
}
\end{funcproto}
  
La funzione opera sul file indicato dal file descriptor \param{fd}, che deve
essere aperto in scrittura, perché utilizza soltanto \textit{lock}
esclusivi. La sezione di file bloccata viene controllata dal valore
di \param{len}, che indica la lunghezza della stessa, usando come riferimento
la posizione corrente sul file. La sezione effettiva varia a secondo del
segno, secondo lo schema illustrato in fig.~\ref{fig:file_lockf_boundary}, se
si specifica un valore nullo il file viene bloccato a partire dalla posizione
corrente fino alla sua fine presente o futura (nello schema corrisponderebbe
ad un valore infinito positivo).

\begin{figure}[!htb] 
  \centering
  \includegraphics[width=10cm]{img/lockf_boundary}
  \caption{Schema della sezione di file bloccata con \func{lockf}.}
  \label{fig:file_lockf_boundary}
\end{figure}

Il comportamento della funzione viene controllato dal valore
dell'argomento \param{cmd}, che specifica quale azione eseguire, i soli valori
consentiti sono i seguenti:

\begin{basedescript}{\desclabelwidth{2.0cm}}
\item[\constd{F\_LOCK}] Richiede un \textit{lock} esclusivo sul file, e blocca
  il processo chiamante se, anche parzialmente, la sezione indicata si
  sovrappone ad una che è già stata bloccata da un altro processo; in caso di
  sovrapposizione con un altro blocco già ottenuto le sezioni vengono unite.
\item[\constd{F\_TLOCK}] Richiede un \textit{exclusive lock}, in maniera
  identica a \const{F\_LOCK}, ma in caso di indisponibilità non blocca il
  processo restituendo un errore di \errval{EAGAIN}.
\item[\constd{F\_ULOCK}] Rilascia il blocco sulla sezione indicata, questo può
  anche causare la suddivisione di una sezione bloccata in precedenza nelle
  due parti eccedenti nel caso si sia indicato un intervallo più limitato.
\item[\constd{F\_TEST}] Controlla la presenza di un blocco sulla sezione di
  file indicata, \func{lockf} ritorna $0$ se la sezione è libera o bloccata
  dal processo stesso, o $-1$ se è bloccata da un altro processo, nel qual
  caso \var{errno} assume il valore \errval{EAGAIN} (ma su alcuni sistemi può
  essere restituito anche \errval{EACCESS}).
\end{basedescript}

La funzione è semplicemente una diversa interfaccia al \textit{file locking}
POSIX ed è realizzata utilizzando \func{fcntl}; pertanto la semantica delle
operazioni è la stessa di quest'ultima e quindi la funzione presenta lo stesso
comportamento riguardo gli effetti della chiusura dei file, ed il
comportamento sui file duplicati e nel passaggio attraverso \func{fork} ed
\func{exec}. Per questo stesso motivo la funzione non è equivalente a
\func{flock} e può essere usata senza interferenze insieme a quest'ultima.

% TODO trattare i POSIX file-private lock introdotti con il 3.15, 
% vedi http://lwn.net/Articles/586904/ correlato:
% http://www.samba.org/samba/news/articles/low_point/tale_two_stds_os2.html 

\subsection{Il \textit{mandatory locking}}
\label{sec:file_mand_locking}

\itindbeg{mandatory~locking}

Il \textit{mandatory locking} è una opzione introdotta inizialmente in SVr4,
per introdurre un \textit{file locking} che, come dice il nome, fosse
effettivo indipendentemente dai controlli eseguiti da un processo. Con il
\textit{mandatory locking} infatti è possibile far eseguire il blocco del file
direttamente al sistema, così che, anche qualora non si predisponessero le
opportune verifiche nei processi, questo verrebbe comunque rispettato.

Per poter utilizzare il \textit{mandatory locking} è stato introdotto un
utilizzo particolare del bit \acr{sgid} dei permessi dei file. Se si ricorda
quanto esposto in sez.~\ref{sec:file_special_perm}), esso viene di norma
utilizzato per cambiare il \ids{GID} effettivo con cui viene eseguito un
programma, ed è pertanto sempre associato alla presenza del permesso di
esecuzione per il gruppo. Impostando questo bit su un file senza permesso di
esecuzione in un sistema che supporta il \textit{mandatory locking}, fa sì che
quest'ultimo venga attivato per il file in questione. In questo modo una
combinazione dei permessi originariamente non contemplata, in quanto senza
significato, diventa l'indicazione della presenza o meno del \textit{mandatory
  locking}.\footnote{un lettore attento potrebbe ricordare quanto detto in
  sez.~\ref{sec:file_perm_management} e cioè che il bit \acr{sgid} viene
  cancellato (come misura di sicurezza) quando di scrive su un file, questo
  non vale quando esso viene utilizzato per attivare il \textit{mandatory
    locking}.}

L'uso del \textit{mandatory locking} presenta vari aspetti delicati, dato che
neanche l'amministratore può passare sopra ad un \textit{file lock}; pertanto
un processo che blocchi un file cruciale può renderlo completamente
inaccessibile, rendendo completamente inutilizzabile il sistema\footnote{il
  problema si potrebbe risolvere rimuovendo il bit \acr{sgid}, ma non è detto
  che sia così facile fare questa operazione con un sistema bloccato.}
inoltre con il \textit{mandatory locking} si può bloccare completamente un
server NFS richiedendo una lettura su un file su cui è attivo un blocco. Per
questo motivo l'abilitazione del \textit{mandatory locking} è di norma
disabilitata, e deve essere attivata filesystem per filesystem in fase di
montaggio, specificando l'apposita opzione di \func{mount} riportata in
sez.~\ref{sec:filesystem_mounting}, o con l'opzione \code{-o mand} per il
comando omonimo.

Si tenga presente inoltre che il \textit{mandatory locking} funziona solo
sull'interfaccia POSIX di \func{fcntl}. Questo ha due conseguenze: che non si
ha nessun effetto sui \textit{file lock} richiesti con l'interfaccia di
\func{flock}, e che la granularità del blocco è quella del singolo byte, come
per \func{fcntl}.

La sintassi di acquisizione dei blocchi è esattamente la stessa vista in
precedenza per \func{fcntl} e \func{lockf}, la differenza è che in caso di
\textit{mandatory lock} attivato non è più necessario controllare la
disponibilità di accesso al file, ma si potranno usare direttamente le
ordinarie funzioni di lettura e scrittura e sarà compito del kernel gestire
direttamente il \textit{file locking}.

Questo significa che in caso di \textit{read lock} la lettura dal file potrà
avvenire normalmente con \func{read}, mentre una \func{write} si bloccherà
fino al rilascio del blocco, a meno di non aver aperto il file con
\const{O\_NONBLOCK}, nel qual caso essa ritornerà immediatamente con un errore
di \errcode{EAGAIN}.

Se invece si è acquisito un \textit{write lock} tutti i tentativi di leggere o
scrivere sulla regione del file bloccata fermeranno il processo fino al
rilascio del blocco, a meno che il file non sia stato aperto con
\const{O\_NONBLOCK}, nel qual caso di nuovo si otterrà un ritorno immediato
con l'errore di \errcode{EAGAIN}.

Infine occorre ricordare che le funzioni di lettura e scrittura non sono le
sole ad operare sui contenuti di un file, e che sia \func{creat} che
\func{open} (quando chiamata con \const{O\_TRUNC}) effettuano dei cambiamenti,
così come \func{truncate}, riducendone le dimensioni (a zero nei primi due
casi, a quanto specificato nel secondo). Queste operazioni sono assimilate a
degli accessi in scrittura e pertanto non potranno essere eseguite (fallendo
con un errore di \errcode{EAGAIN}) su un file su cui sia presente un qualunque
blocco (le prime due sempre, la terza solo nel caso che la riduzione delle
dimensioni del file vada a sovrapporsi ad una regione bloccata).

L'ultimo aspetto della interazione del \textit{mandatory locking} con le
funzioni di accesso ai file è quello relativo ai file mappati in memoria (vedi
sez.~\ref{sec:file_memory_map}); anche in tal caso infatti, quando si esegue
la mappatura con l'opzione \const{MAP\_SHARED}, si ha un accesso al contenuto
del file. Lo standard SVID prevede che sia impossibile eseguire il
\textit{memory mapping} di un file su cui sono presenti dei
blocchi\footnote{alcuni sistemi, come HP-UX, sono ancora più restrittivi e lo
  impediscono anche in caso di \textit{advisory locking}, anche se questo
  comportamento non ha molto senso, dato che comunque qualunque accesso
  diretto al file è consentito.} in Linux è stata però fatta la scelta
implementativa\footnote{per i dettagli si possono leggere le note relative
  all'implementazione, mantenute insieme ai sorgenti del kernel nel file
  \file{Documentation/mandatory.txt}.}  di seguire questo comportamento
soltanto quando si chiama \func{mmap} con l'opzione \const{MAP\_SHARED} (nel
qual caso la funzione fallisce con il solito \errcode{EAGAIN}) che comporta la
possibilità di modificare il file.

Si tenga conto infine che su Linux l'implementazione corrente del
\textit{mandatory locking} è difettosa e soffre di una \textit{race
  condition}, per cui una scrittura con \func{write} che si sovrapponga alla
richiesta di un \textit{read lock} può modificare i dati anche dopo che questo
è stato ottenuto, ed una lettura con \func{read} può restituire dati scritti
dopo l'ottenimento di un \textit{write lock}. Lo stesso tipo di problema si
può presentare anche con l'uso di file mappati in memoria; pertanto allo stato
attuale delle cose è sconsigliabile fare affidamento sul \textit{mandatory
  locking}.

% TODO il supporto è stato reso opzionale nel 4.5, verrà eliminato nel futuro
% (vedi http://lwn.net/Articles/667210/)

\itindend{file~locking}

\itindend{mandatory~locking}


\section{L'\textit{I/O multiplexing}}
\label{sec:file_multiplexing}


Uno dei problemi che si presentano quando si deve operare contemporaneamente
su molti file usando le funzioni illustrate in
sez.~\ref{sec:file_unix_interface} e sez.~\ref{sec:files_std_interface} è che
si può essere bloccati nelle operazioni su un file mentre un altro potrebbe
essere disponibile. L'\textit{I/O multiplexing} nasce risposta a questo
problema. In questa sezione forniremo una introduzione a questa problematica
ed analizzeremo le varie funzioni usate per implementare questa modalità di
I/O.


\subsection{La problematica dell'\textit{I/O multiplexing}}
\label{sec:file_noblocking}

Abbiamo visto in sez.~\ref{sec:sig_gen_beha}, affrontando la suddivisione fra
\textit{fast} e \textit{slow} \textit{system call}, che in certi casi le
funzioni di I/O eseguite su un file descriptor possono bloccarsi
indefinitamente. Questo non avviene mai per i file normali, per i quali le
funzioni di lettura e scrittura ritornano sempre subito, ma può avvenire per
alcuni file di dispositivo, come ad esempio una seriale o un terminale, o con
l'uso di file descriptor collegati a meccanismi di intercomunicazione come le
\textit{pipe} (vedi sez.~\ref{sec:ipc_unix}) ed i socket (vedi
sez.~\ref{sec:sock_socket_def}). In casi come questi ad esempio una operazione
di lettura potrebbe bloccarsi se non ci sono dati disponibili sul descrittore
su cui la si sta effettuando.

Questo comportamento è alla radice di una delle problematiche più comuni che
ci si trova ad affrontare nella gestione delle operazioni di I/O: la necessità
di operare su più file descriptor eseguendo funzioni che possono bloccarsi
indefinitamente senza che sia possibile prevedere quando questo può
avvenire. Un caso classico è quello di un server di rete (tratteremo la
problematica in dettaglio nella seconda parte della guida) in attesa di dati
in ingresso prevenienti da vari client.

In un caso di questo tipo, se si andasse ad operare sui vari file descriptor
aperti uno dopo l'altro, potrebbe accadere di restare bloccati nell'eseguire
una lettura su uno di quelli che non è ``\textsl{pronto}'', quando ce ne
potrebbe essere un altro con dati disponibili. Questo comporta nel migliore
dei casi una operazione ritardata inutilmente nell'attesa del completamento di
quella bloccata, mentre nel peggiore dei casi, quando la conclusione
dell'operazione bloccata dipende da quanto si otterrebbe dal file descriptor
``\textsl{disponibile}'', si potrebbe addirittura arrivare ad un
\textit{deadlock}.

\itindbeg{polling}

Abbiamo già accennato in sez.~\ref{sec:file_open_close} che è possibile
prevenire questo tipo di comportamento delle funzioni di I/O aprendo un file
in \textsl{modalità non-bloccante}, attraverso l'uso del flag
\const{O\_NONBLOCK} nella chiamata di \func{open}. In questo caso le funzioni
di lettura o scrittura eseguite sul file che si sarebbero bloccate ritornano
immediatamente, restituendo l'errore \errcode{EAGAIN}.  L'utilizzo di questa
modalità di I/O permette di risolvere il problema controllando a turno i vari
file descriptor, in un ciclo in cui si ripete l'accesso fintanto che esso non
viene garantito. Ovviamente questa tecnica, detta \textit{polling}, è
estremamente inefficiente: si tiene costantemente impiegata la CPU solo per
eseguire in continuazione delle \textit{system call} che nella gran parte dei
casi falliranno.

\itindend{polling}

É appunto per superare questo problema è stato introdotto il concetto di
\textit{I/O multiplexing}, una nuova modalità per la gestione dell'I/O che
consente di tenere sotto controllo più file descriptor in contemporanea,
permettendo di bloccare un processo quando le operazioni di lettura o
scrittura non sono immediatamente effettuabili, e di riprenderne l'esecuzione
una volta che almeno una di quelle che erano state richieste diventi
possibile, in modo da poterla eseguire con la sicurezza di non restare
bloccati.

Dato che, come abbiamo già accennato, per i normali file su disco non si ha
mai un accesso bloccante, l'uso più comune delle funzioni che esamineremo nei
prossimi paragrafi è per i server di rete, in cui esse vengono utilizzate per
tenere sotto controllo dei socket; pertanto ritorneremo su di esse con
ulteriori dettagli e qualche esempio di utilizzo concreto in
sez.~\ref{sec:TCP_sock_multiplexing}.


\subsection{Le funzioni \func{select} e \func{pselect}}
\label{sec:file_select}

Il primo kernel unix-like ad introdurre una interfaccia per l'\textit{I/O
  multiplexing} è stato BSD, con la funzione \funcd{select} che è apparsa in
BSD4.2 ed è stata standardizzata in BSD4.4, in seguito è stata portata su
tutti i sistemi che supportano i socket, compreso le varianti di System V ed
inserita in POSIX.1-2001; il suo prototipo è:\footnote{l'header
  \texttt{sys/select.h} è stato introdotto con POSIX.1-2001, è ed presente con
  la \acr{glibc} a partire dalla versione 2.0, in precedenza, con le
  \acr{libc4} e \acr{libc5}, occorreva includere \texttt{sys/time.h},
  \texttt{sys/types.h} e \texttt{unistd.h}.}

\begin{funcproto}{
\fhead{sys/select.h}
\fdecl{int select(int ndfs, fd\_set *readfds, fd\_set *writefds, fd\_set
    *exceptfds, \\
\phantom{int select(}struct timeval *timeout)}
\fdesc{Attende che uno fra i file descriptor degli insiemi specificati diventi
  attivo.} 
}
{La funzione ritorna $0$ in caso di successo e $-1$ per un errore, nel qual
  caso \var{errno} assumerà uno dei valori: 
  \begin{errlist}
  \item[\errcode{EBADF}] si è specificato un file descriptor non valido
    (chiuso o con errori) in uno degli insiemi.
  \item[\errcode{EINTR}] la funzione è stata interrotta da un segnale.
  \item[\errcode{EINVAL}] si è specificato per \param{ndfs} un valore negativo
    o un valore non valido per \param{timeout}.
  \end{errlist}
  ed inoltre \errval{ENOMEM} nel suo significato generico.}
\end{funcproto}

La funzione mette il processo in stato di \textit{sleep} (vedi
tab.~\ref{tab:proc_proc_states}) fintanto che almeno uno dei file descriptor
degli insiemi specificati (\param{readfds}, \param{writefds} e
\param{exceptfds}), non diventa attivo, per un tempo massimo specificato da
\param{timeout}.

\itindbeg{file~descriptor~set} 

Per specificare quali file descriptor si intende selezionare la funzione usa
un particolare oggetto, il \textit{file descriptor set}, identificato dal tipo
\typed{fd\_set}, che serve ad identificare un insieme di file descriptor, in
maniera analoga a come un \textit{signal set} (vedi sez.~\ref{sec:sig_sigset})
identifica un insieme di segnali. Per la manipolazione di questi \textit{file
  descriptor set} si possono usare delle opportune macro di preprocessore:

{\centering
\vspace{3pt}
\begin{funcbox}{
\fhead{sys/select.h}
\fdecl{void \macrod{FD\_ZERO}(fd\_set *set)}
\fdesc{Inizializza l'insieme (vuoto).} 
\fdecl{void \macrod{FD\_SET}(int fd, fd\_set *set)}
\fdesc{Inserisce il file descriptor \param{fd} nell'insieme.} 
\fdecl{void \macrod{FD\_CLR}(int fd, fd\_set *set)}
\fdesc{Rimuove il file descriptor \param{fd} dall'insieme.} 
\fdecl{int \macrod{FD\_ISSET}(int fd, fd\_set *set)}
\fdesc{Controlla se il file descriptor \param{fd} è nell'insieme.} 
}
\end{funcbox}}


In genere un \textit{file descriptor set} può contenere fino ad un massimo di
\macrod{FD\_SETSIZE} file descriptor.  Questo valore in origine corrispondeva
al limite per il numero massimo di file aperti (ad esempio in Linux, fino alla
serie 2.0.x, c'era un limite di 256 file per processo), ma da quando, nelle
versioni più recenti del kernel, questo limite è stato rimosso, esso indica le
dimensioni massime dei numeri usati nei \textit{file descriptor set}, ed il
suo valore, secondo lo standard POSIX 1003.1-2001, è definito in
\headfile{sys/select.h}, ed è pari a 1024.

Si tenga presente che i \textit{file descriptor set} devono sempre essere
inizializzati con \macro{FD\_ZERO}; passare a \func{select} un valore non
inizializzato può dar luogo a comportamenti non prevedibili. Allo stesso modo
usare \macro{FD\_SET} o \macro{FD\_CLR} con un file descriptor il cui valore
eccede \macro{FD\_SETSIZE} può dare luogo ad un comportamento indefinito.

La funzione richiede di specificare tre insiemi distinti di file descriptor;
il primo, \param{readfds}, verrà osservato per rilevare la disponibilità di
effettuare una lettura,\footnote{per essere precisi la funzione ritornerà in
  tutti i casi in cui la successiva esecuzione di \func{read} risulti non
  bloccante, quindi anche in caso di \textit{end-of-file}.} il secondo,
\param{writefds}, per verificare la possibilità di effettuare una scrittura ed
il terzo, \param{exceptfds}, per verificare l'esistenza di eccezioni come i
dati urgenti su un socket, (vedi sez.~\ref{sec:TCP_urgent_data}).

Dato che in genere non si tengono mai sotto controllo fino a
\macro{FD\_SETSIZE} file contemporaneamente, la funzione richiede di
specificare qual è il valore più alto fra i file descriptor indicati nei tre
insiemi precedenti. Questo viene fatto per efficienza, per evitare di passare
e far controllare al kernel una quantità di memoria superiore a quella
necessaria. Questo limite viene indicato tramite l'argomento \param{ndfs}, che
deve corrispondere al valore massimo aumentato di uno. Si ricordi infatti che
i file descriptor sono numerati progressivamente a partire da zero, ed il
valore indica il numero più alto fra quelli da tenere sotto controllo,
dimenticarsi di aumentare di uno il valore di \param{ndfs} è un errore comune.

Infine l'argomento \param{timeout}, espresso con il puntatore ad una struttura
di tipo \struct{timeval} (vedi fig.~\ref{fig:sys_timeval_struct}) specifica un
tempo massimo di attesa prima che la funzione ritorni; se impostato a
\val{NULL} la funzione attende indefinitamente. Si può specificare anche un
tempo nullo (cioè una struttura \struct{timeval} con i campi impostati a
zero), qualora si voglia semplicemente controllare lo stato corrente dei file
descriptor, e così può essere utilizzata eseguire il \textit{polling} su un
gruppo di file descriptor. Usare questo argomento con tutti i \textit{file
  descriptor set} vuoti è un modo portabile, disponibile anche su sistemi in
cui non sono disponibili le funzioni avanzate di sez.~\ref{sec:sig_timer_adv},
per tenere un processo in stato di \textit{sleep} con precisioni inferiori al
secondo.

In caso di successo la funzione restituisce il numero di file descriptor
pronti, seguendo il comportamento previsto dallo standard
POSIX.1-2001,\footnote{si tenga però presente che esistono alcune versioni di
  Unix che non si comportano in questo modo, restituendo un valore positivo
  generico.}  e ciascun insieme viene sovrascritto per indicare quali sono i
file descriptor pronti per le operazioni ad esso relative, in modo da poterli
controllare con \macro{FD\_ISSET}.  Se invece scade il tempo indicato
da \param{timout} viene restituito un valore nullo e i \textit{file descriptor
  set} non vengono modificati. In caso di errore la funzione restituisce $-1$, i
valori dei tre insiemi e di \param{timeout} sono indefiniti e non si può fare
nessun affidamento sul loro contenuto; nelle versioni più recenti della
funzione invece i \textit{file descriptor set} non vengono modificati anche in
caso di errore.

Si tenga presente infine che su Linux, in caso di programmazione
\textit{multi-thread} se un file descriptor viene chiuso in un altro
\textit{thread} rispetto a quello in cui si sta usando \func{select}, questa
non subisce nessun effetto. In altre varianti di sistemi unix-like invece
\func{select} ritorna indicando che il file descriptor è pronto, con
conseguente possibile errore nel caso lo si usi senza che sia stato
riaperto. Lo standard non prevede niente al riguardo e non si deve dare per
assunto nessuno dei due comportamenti se si vogliono scrivere programmi
portabili.

\itindend{file~descriptor~set}

Una volta ritornata la funzione, si potrà controllare quali sono i file
descriptor pronti, ed operare su di essi. Si tenga presente però che
\func{select} fornisce solo di un suggerimento, esistono infatti condizioni in
cui \func{select} può riportare in maniera spuria che un file descriptor è
pronto, ma l'esecuzione di una operazione di I/O si bloccherebbe: ad esempio
con Linux questo avviene quando su un socket arrivano dei dati che poi vengono
scartati perché corrotti (ma sono possibili pure altri casi); in tal caso pur
risultando il relativo file descriptor pronto in lettura una successiva
esecuzione di una \func{read} si bloccherebbe. Per questo motivo quando si usa
l'\textit{I/O multiplexing} è sempre raccomandato l'uso delle funzioni di
lettura e scrittura in modalità non bloccante.

Su Linux quando la \textit{system call} \func{select} viene interrotta da un
segnale modifica il valore nella struttura puntata da \param{timeout},
impostandolo al tempo restante. In tal caso infatti si ha un errore di
\errcode{EINTR} ed occorre rilanciare la funzione per proseguire l'attesa, ed
in questo modo non è necessario ricalcolare tutte le volte il tempo
rimanente. Questo può causare problemi di portabilità sia quando si usa codice
scritto su Linux che legge questo valore, sia quando si usano programmi
scritti per altri sistemi che non dispongono di questa caratteristica e
ricalcolano \param{timeout} tutte le volte. In genere questa caratteristica è
disponibile nei sistemi che derivano da System V e non è disponibile per
quelli che derivano da BSD; lo standard POSIX.1-2001 non permette questo
comportamento e per questo motivo la \acr{glibc} nasconde il comportamento
passando alla \textit{system call} una copia dell'argomento \param{timeout}.

Uno dei problemi che si presentano con l'uso di \func{select} è che il suo
comportamento dipende dal valore del file descriptor che si vuole tenere sotto
controllo.  Infatti il kernel riceve con \param{ndfs} un limite massimo per
tale valore, e per capire quali sono i file descriptor da tenere sotto
controllo dovrà effettuare una scansione su tutto l'intervallo, che può anche
essere molto ampio anche se i file descriptor sono solo poche unità; tutto ciò
ha ovviamente delle conseguenze ampiamente negative per le prestazioni.

Inoltre c'è anche il problema che il numero massimo dei file che si possono
tenere sotto controllo, la funzione è nata quando il kernel consentiva un
numero massimo di 1024 file descriptor per processo, adesso che il numero può
essere arbitrario si viene a creare una dipendenza del tutto artificiale dalle
dimensioni della struttura \type{fd\_set}, che può necessitare di essere
estesa, con ulteriori perdite di prestazioni. 

Lo standard POSIX è rimasto a lungo senza primitive per l'\textit{I/O
  multiplexing}, introdotto solo con le ultime revisioni dello standard (POSIX
1003.1g-2000 e POSIX 1003.1-2001). La scelta è stata quella di seguire
l'interfaccia creata da BSD, ma prevede che tutte le funzioni ad esso relative
vengano dichiarate nell'header \headfiled{sys/select.h}, che sostituisce i
precedenti, ed inoltre aggiunge a \func{select} una nuova funzione
\funcd{pselect},\footnote{il supporto per lo standard POSIX 1003.1-2001, ed
  l'header \headfile{sys/select.h}, compaiono in Linux a partire dalla
  \acr{glibc} 2.1. Le \acr{libc4} e \acr{libc5} non contengono questo header,
  la \acr{glibc} 2.0 contiene una definizione sbagliata di \func{psignal},
  senza l'argomento \param{sigmask}, la definizione corretta è presente dalle
  \acr{glibc} 2.1-2.2.1 se si è definito \macro{\_GNU\_SOURCE} e nelle
  \acr{glibc} 2.2.2-2.2.4 se si è definito \macro{\_XOPEN\_SOURCE} con valore
  maggiore di 600.} il cui prototipo è:

\begin{funcproto}{
\fhead{sys/select.h}
\fdecl{int pselect(int n, fd\_set *readfds, fd\_set *writefds, 
  fd\_set *exceptfds, \\ 
\phantom{int pselect(}struct timespec *timeout, sigset\_t *sigmask)}
\fdesc{Attende che uno dei file descriptor degli insiemi specificati diventi
  attivo.} 
}
{La funzione ritorna il numero (anche nullo) di file descriptor che sono
  attivi in caso di successo e $-1$ per un errore, nel qual caso \var{errno}
  assumerà uno dei valori:
  \begin{errlist}
  \item[\errcode{EBADF}] si è specificato un file descriptor sbagliato in uno
    degli insiemi.
  \item[\errcode{EINTR}] la funzione è stata interrotta da un segnale.
  \item[\errcode{EINVAL}] si è specificato per \param{ndfs} un valore negativo
    o un valore non valido per \param{timeout}.
   \end{errlist}
   ed inoltre \errval{ENOMEM} nel suo significato generico.
}
\end{funcproto}

La funzione è sostanzialmente identica a \func{select}, solo che usa una
struttura \struct{timespec} (vedi fig.~\ref{fig:sys_timespec_struct}) per
indicare con maggiore precisione il timeout e non ne aggiorna il valore in
caso di interruzione. In realtà anche in questo caso la \textit{system call}
di Linux aggiorna il valore al tempo rimanente, ma la funzione fornita dalla
\acr{glibc} modifica questo comportamento passando alla \textit{system call}
una variabile locale, in modo da mantenere l'aderenza allo standard POSIX che
richiede che il valore di \param{timeout} non sia modificato. 

Rispetto a \func{select} la nuova funzione prende un argomento
aggiuntivo \param{sigmask}, un puntatore ad una maschera di segnali (si veda
sez.~\ref{sec:sig_sigmask}).  Nell'esecuzione la maschera dei segnali corrente
viene sostituita da quella così indicata immediatamente prima di eseguire
l'attesa, e viene poi ripristinata al ritorno della funzione. L'uso
di \param{sigmask} è stato introdotto allo scopo di prevenire possibili
\textit{race condition} quando oltre alla presenza di dati sui file descriptor
come nella \func{select} ordinaria, ci si deve porre in attesa anche
dell'arrivo di un segnale.

Come abbiamo visto in sez.~\ref{sec:sig_example} la tecnica classica per
rilevare l'arrivo di un segnale è quella di utilizzare il gestore per
impostare una variabile globale e controllare questa nel corpo principale del
programma; abbiamo visto in quell'occasione come questo lasci spazio a
possibili \textit{race condition}, per cui diventa essenziale utilizzare
\func{sigprocmask} per disabilitare la ricezione del segnale prima di eseguire
il controllo e riabilitarlo dopo l'esecuzione delle relative operazioni, onde
evitare l'arrivo di un segnale immediatamente dopo il controllo, che andrebbe
perso.

Nel nostro caso il problema si pone quando, oltre al segnale, si devono tenere
sotto controllo anche dei file descriptor con \func{select}, in questo caso si
può fare conto sul fatto che all'arrivo di un segnale essa verrebbe interrotta
e si potrebbero eseguire di conseguenza le operazioni relative al segnale e
alla gestione dati con un ciclo del tipo:
\includecodesnip{listati/select_race.c} 
qui però emerge una \textit{race condition}, perché se il segnale arriva prima
della chiamata a \func{select}, questa non verrà interrotta, e la ricezione
del segnale non sarà rilevata.

Per questo è stata introdotta \func{pselect} che attraverso l'argomento
\param{sigmask} permette di riabilitare la ricezione il segnale
contestualmente all'esecuzione della funzione,\footnote{in Linux però, fino al
  kernel 2.6.16, non era presente la relativa \textit{system call}, e la
  funzione era implementata nella \acr{glibc} attraverso \func{select} (vedi
  \texttt{man select\_tut}) per cui la possibilità di \textit{race condition}
  permaneva; in tale situazione si può ricorrere ad una soluzione alternativa,
  chiamata \itindex{self-pipe~trick} \textit{self-pipe trick}, che consiste
  nell'aprire una \textit{pipe} (vedi sez.~\ref{sec:ipc_pipes}) ed usare
  \func{select} sul capo in lettura della stessa; si può indicare l'arrivo di
  un segnale scrivendo sul capo in scrittura all'interno del gestore dello
  stesso; in questo modo anche se il segnale va perso prima della chiamata di
  \func{select} questa lo riconoscerà comunque dalla presenza di dati sulla
  \textit{pipe}.} ribloccandolo non appena essa ritorna, così che il
precedente codice potrebbe essere riscritto nel seguente modo:
\includecodesnip{listati/pselect_norace.c} 
in questo caso utilizzando \var{oldmask} durante l'esecuzione di
\func{pselect} la ricezione del segnale sarà abilitata, ed in caso di
interruzione si potranno eseguire le relative operazioni.


\subsection{Le funzioni \func{poll} e \func{ppoll}}
\label{sec:file_poll}

Nello sviluppo di System V, invece di utilizzare l'interfaccia di
\func{select}, che è una estensione tipica di BSD, è stata introdotta una
interfaccia completamente diversa, basata sulla funzione di sistema
\funcd{poll},\footnote{la funzione è prevista dallo standard XPG4, ed è stata
  introdotta in Linux come \textit{system call} a partire dal kernel 2.1.23 ed
  inserita nelle \acr{libc} 5.4.28, originariamente l'argomento \param{nfds}
  era di tipo \ctyp{unsigned int}, la funzione è stata inserita nello standard
  POSIX.1-2001 in cui è stato introdotto il tipo nativo \typed{nfds\_t}.} il
cui prototipo è:

\begin{funcproto}{
\fhead{sys/poll.h}
\fdecl{int poll(struct pollfd *ufds, nfds\_t nfds, int timeout)}
\fdesc{Attende un cambiamento di stato su un insieme di file
  descriptor.} 
}

{La funzione ritorna $0$ in caso di successo e $-1$ per un errore, nel qual
  caso \var{errno} assumerà uno dei valori: 
  \begin{errlist}
  \item[\errcode{EBADF}] si è specificato un file descriptor sbagliato in uno
    degli insiemi.
  \item[\errcode{EINTR}] la funzione è stata interrotta da un segnale.
  \item[\errcode{EINVAL}] il valore di \param{nfds} eccede il limite
    \const{RLIMIT\_NOFILE}.
  \end{errlist}
  ed inoltre \errval{EFAULT} e \errval{ENOMEM} nel loro significato generico.}
\end{funcproto}

La funzione permette di tenere sotto controllo contemporaneamente \param{ndfs}
file descriptor, specificati attraverso il puntatore \param{ufds} ad un
vettore di strutture \struct{pollfd}.  Come con \func{select} si può
interrompere l'attesa dopo un certo tempo, questo deve essere specificato con
l'argomento \param{timeout} in numero di millisecondi: un valore negativo
indica un'attesa indefinita, mentre un valore nullo comporta il ritorno
immediato, e può essere utilizzato per impiegare \func{poll} in modalità
\textsl{non-bloccante}.

\begin{figure}[!htb]
  \footnotesize \centering
  \begin{minipage}[c]{0.90\textwidth}
    \includestruct{listati/pollfd.h}
  \end{minipage} 
  \normalsize 
  \caption{La struttura \structd{pollfd}, utilizzata per specificare le
    modalità di controllo di un file descriptor alla funzione \func{poll}.}
  \label{fig:file_pollfd}
\end{figure}

Per ciascun file da controllare deve essere inizializzata una struttura
\struct{pollfd} nel vettore indicato dall'argomento \param{ufds}.  La
struttura, la cui definizione è riportata in fig.~\ref{fig:file_pollfd},
prevede tre campi: in \var{fd} deve essere indicato il numero del file
descriptor da controllare, in \var{events} deve essere specificata una
maschera binaria di flag che indichino il tipo di evento che si vuole
controllare, mentre in \var{revents} il kernel restituirà il relativo
risultato. 

Usando un valore negativo per \param{fd} la corrispondente struttura sarà
ignorata da \func{poll} ed il campo \var{revents} verrà azzerato, questo
consente di eliminare temporaneamente un file descriptor dalla lista senza
dover modificare il vettore \param{ufds}. Dato che i dati in ingresso sono del
tutto indipendenti da quelli in uscita (che vengono restituiti in
\var{revents}) non è necessario reinizializzare tutte le volte il valore delle
strutture \struct{pollfd} a meno di non voler cambiare qualche condizione.

Le costanti che definiscono i valori relativi ai bit usati nelle maschere
binarie dei campi \var{events} e \var{revents} sono riportate in
tab.~\ref{tab:file_pollfd_flags}, insieme al loro significato. Le si sono
suddivise in tre gruppi principali, nel primo gruppo si sono indicati i bit
utilizzati per controllare l'attività in ingresso, nel secondo quelli per
l'attività in uscita, infine il terzo gruppo contiene dei valori che vengono
utilizzati solo nel campo \var{revents} per notificare delle condizioni di
errore.

\begin{table}[htb]
  \centering
  \footnotesize
  \begin{tabular}[c]{|l|l|}
    \hline
    \textbf{Flag}  & \textbf{Significato} \\
    \hline
    \hline
    \constd{POLLIN}    & È possibile la lettura.\\
    \constd{POLLRDNORM}& Sono disponibili in lettura dati normali.\\ 
    \constd{POLLRDBAND}& Sono disponibili in lettura dati prioritari.\\
    \constd{POLLPRI}   & È possibile la lettura di dati urgenti.\\ 
    \hline
    \constd{POLLOUT}   & È possibile la scrittura immediata.\\
    \constd{POLLWRNORM}& È possibile la scrittura di dati normali.\\ 
    \constd{POLLWRBAND}& È possibile la scrittura di dati prioritari.\\
    \hline
    \constd{POLLERR}   & C'è una condizione di errore.\\
    \constd{POLLHUP}   & Si è verificato un hung-up.\\
    \constd{POLLRDHUP} & Si è avuta una \textsl{half-close} su un
                        socket.\footnotemark\\ 
    \constd{POLLNVAL}  & Il file descriptor non è aperto.\\
    \hline
    \constd{POLLMSG}   & Definito per compatibilità con SysV.\\
    \hline    
  \end{tabular}
  \caption{Costanti per l'identificazione dei vari bit dei campi
    \var{events} e \var{revents} di \struct{pollfd}.}
  \label{tab:file_pollfd_flags}
\end{table}

\footnotetext{si tratta di una estensione specifica di Linux, disponibile a
  partire dal kernel 2.6.17 definendo la marco \macro{\_GNU\_SOURCE}, che
  consente di riconoscere la chiusura in scrittura dell'altro capo di un
  socket, situazione che si viene chiamata appunto \textit{half-close}
  (\textsl{mezza chiusura}) su cui torneremo con maggiori dettagli in
  sez.~\ref{sec:TCP_shutdown}.}

Il valore \const{POLLMSG} non viene utilizzato ed è definito solo per
compatibilità con l'implementazione di System V che usa i cosiddetti
``\textit{stream}''. Si tratta di una interfaccia specifica di SysV non
presente in Linux, che non ha nulla a che fare con gli \textit{stream} delle
librerie standard del C visti in sez.~\ref{sec:file_stream}. Da essa derivano
i nomi di alcune costanti poiché per quegli \textit{stream} sono definite tre
classi di dati: \textsl{normali}, \textit{prioritari} ed \textit{urgenti}.  In
Linux la distinzione ha senso solo per i dati urgenti dei socket (vedi
sez.~\ref{sec:TCP_urgent_data}), ma su questo e su come \func{poll} reagisce
alle varie condizioni dei socket torneremo in sez.~\ref{sec:TCP_serv_poll},
dove vedremo anche un esempio del suo utilizzo.

Le costanti relative ai diversi tipi di dati normali e prioritari che fanno
riferimento alle implementazioni in stile System V sono \const{POLLRDNORM},
\const{POLLWRNORM}, \const{POLLRDBAND} e \const{POLLWRBAND}. Le prime due sono
equivalenti rispettivamente a \const{POLLIN} e \const{POLLOUT},
\const{POLLRDBAND} non viene praticamente mai usata su Linux mentre
\const{POLLWRBAND} ha senso solo sui socket. In ogni caso queste costanti sono
utilizzabili soltanto qualora si sia definita la macro
\macro{\_XOPEN\_SOURCE}.

In caso di successo \func{poll} ritorna restituendo il numero di file (un
valore positivo) per i quali si è verificata una delle condizioni di attesa
richieste o per i quali si è verificato un errore, avvalorando i relativi bit
di \var{revents}. In caso di errori sui file vengono utilizzati i valori della
terza sezione di tab.~\ref{tab:file_pollfd_flags} che hanno significato solo
per \var{revents} (se specificati in \var{events} vengono ignorati). Un valore
di ritorno nullo indica che si è raggiunto il timeout, mentre un valore
negativo indica un errore nella chiamata, il cui codice viene riportato al
solito tramite \var{errno}.

L'uso di \func{poll} consente di superare alcuni dei problemi illustrati in
precedenza per \func{select}; anzitutto, dato che in questo caso si usa un
vettore di strutture \struct{pollfd} di dimensione arbitraria, non esiste il
limite introdotto dalle dimensioni massime di un \textit{file descriptor set}
e la dimensione dei dati passati al kernel dipende solo dal numero dei file
descriptor che si vogliono controllare, non dal loro valore. Infatti, anche se
usando dei bit un \textit{file descriptor set} può essere più efficiente di un
vettore di strutture \struct{pollfd}, qualora si debba osservare un solo file
descriptor con un valore molto alto ci si troverà ad utilizzare inutilmente un
maggiore quantitativo di memoria.

Inoltre con \func{select} lo stesso \textit{file descriptor set} è usato sia
in ingresso che in uscita, e questo significa che tutte le volte che si vuole
ripetere l'operazione occorre reinizializzarlo da capo. Questa operazione, che
può essere molto onerosa se i file descriptor da tenere sotto osservazione
sono molti, non è invece necessaria con \func{poll}.

Abbiamo visto in sez.~\ref{sec:file_select} come lo standard POSIX preveda una
variante di \func{select} che consente di gestire correttamente la ricezione
dei segnali nell'attesa su un file descriptor.  Con l'introduzione di una
implementazione reale di \func{pselect} nel kernel 2.6.16, è stata aggiunta
anche una analoga funzione che svolga lo stesso ruolo per \func{poll}.

In questo caso si tratta di una estensione che è specifica di Linux e non è
prevista da nessuno standard; essa può essere utilizzata esclusivamente se si
definisce la macro \macro{\_GNU\_SOURCE} ed ovviamente non deve essere usata
se si ha a cuore la portabilità. La funzione è \funcd{ppoll}, ed il suo
prototipo è:

\begin{funcproto}{
\fhead{sys/poll.h}
\fdecl{int ppoll(struct pollfd *fds, nfds\_t nfds, 
  const struct timespec *timeout, \\
\phantom{int ppoll(}const sigset\_t *sigmask)} 

\fdesc{Attende un cambiamento di stato su un insieme di file descriptor.}
}

{La funzione ritorna il numero di file descriptor con attività in caso di
  successo, $0$ se c'è stato un timeout e $-1$ per un errore, nel qual caso
  \var{errno} assumerà uno dei valori:
  \begin{errlist}
  \item[\errcode{EBADF}] si è specificato un file descriptor sbagliato in uno
    degli insiemi.
  \item[\errcode{EINTR}] la funzione è stata interrotta da un segnale.
  \item[\errcode{EINVAL}] il valore di \param{nfds} eccede il limite
    \const{RLIMIT\_NOFILE}.
  \end{errlist}
ed inoltre \errval{EFAULT} e \errval{ENOMEM} nel loro significato generico.
}  
\end{funcproto}

La funzione ha lo stesso comportamento di \func{poll}, solo che si può
specificare, con l'argomento \param{sigmask}, il puntatore ad una maschera di
segnali; questa sarà la maschera utilizzata per tutto il tempo che la funzione
resterà in attesa, all'uscita viene ripristinata la maschera originale.  L'uso
di questa funzione è cioè equivalente, come illustrato nella pagina di
manuale, all'esecuzione atomica del seguente codice:
\includecodesnip{listati/ppoll_means.c} 

Eccetto per \param{timeout}, che come per \func{pselect} deve essere un
puntatore ad una struttura \struct{timespec}, gli altri argomenti comuni con
\func{poll} hanno lo stesso significato, e la funzione restituisce gli stessi
risultati illustrati in precedenza. Come nel caso di \func{pselect} la
\textit{system call} che implementa \func{ppoll} restituisce, se la funzione
viene interrotta da un segnale, il tempo mancante in \param{timeout}, e come
per \func{pselect} la funzione di libreria fornita dalla \acr{glibc} maschera
questo comportamento non modificando mai il valore di \param{timeout} anche se
in questo caso non esiste nessuno standard che richieda questo comportamento.

Infine anche per \func{poll} e \func{ppoll} valgono le considerazioni relative
alla possibilità di avere delle notificazione spurie della disponibilità di
accesso ai file descriptor illustrate per \func{select} in
sez.~\ref{sec:file_select}, che non staremo a ripetere qui.

\subsection{L'interfaccia di \textit{epoll}}
\label{sec:file_epoll}

\itindbeg{epoll}

Nonostante \func{poll} presenti alcuni vantaggi rispetto a \func{select},
anche questa funzione non è molto efficiente quando deve essere utilizzata con
un gran numero di file descriptor,\footnote{in casi del genere \func{select}
  viene scartata a priori, perché può avvenire che il numero di file
  descriptor ecceda le dimensioni massime di un \textit{file descriptor set}.}
in particolare nel caso in cui solo pochi di questi diventano attivi. Il
problema in questo caso è che il tempo impiegato da \func{poll} a trasferire i
dati da e verso il kernel è proporzionale al numero di file descriptor
osservati, non a quelli che presentano attività.

Quando ci sono decine di migliaia di file descriptor osservati e migliaia di
eventi al secondo (il caso classico è quello di un server web di un sito con
molti accessi) l'uso di \func{poll} comporta la necessità di trasferire avanti
ed indietro da \textit{user space} a \textit{kernel space} una lunga lista di
strutture \struct{pollfd} migliaia di volte al secondo. A questo poi si
aggiunge il fatto che la maggior parte del tempo di esecuzione sarà impegnato
ad eseguire una scansione su tutti i file descriptor tenuti sotto controllo
per determinare quali di essi (in genere una piccola percentuale) sono
diventati attivi. In una situazione come questa l'uso delle funzioni classiche
dell'interfaccia dell'\textit{I/O multiplexing} viene a costituire un collo di
bottiglia che degrada irrimediabilmente le prestazioni.

Per risolvere questo tipo di situazioni sono state ideate delle interfacce
specialistiche (come \texttt{/dev/poll} in Solaris, o \texttt{kqueue} in BSD)
il cui scopo fondamentale è quello di restituire solamente le informazioni
relative ai file descriptor osservati che presentano una attività, evitando
così le problematiche appena illustrate. In genere queste prevedono che si
registrino una sola volta i file descriptor da tenere sotto osservazione, e
forniscono un meccanismo che notifica quali di questi presentano attività.

Le modalità con cui avviene la notifica sono due, la prima è quella classica
(quella usata da \func{poll} e \func{select}) che viene chiamata \textit{level
  triggered}.\footnote{la nomenclatura è stata introdotta da Jonathan Lemon in
  un articolo su \texttt{kqueue} al BSDCON 2000, e deriva da quella usata
  nell'elettronica digitale.} In questa modalità vengono notificati i file
descriptor che sono \textsl{pronti} per l'operazione richiesta, e questo
avviene indipendentemente dalle operazioni che possono essere state fatte su
di essi a partire dalla precedente notifica.  Per chiarire meglio il concetto
ricorriamo ad un esempio: se su un file descriptor sono diventati disponibili
in lettura 2000 byte ma dopo la notifica ne sono letti solo 1000 (ed è quindi
possibile eseguire una ulteriore lettura dei restanti 1000), in modalità
\textit{level triggered} questo sarà nuovamente notificato come
\textsl{pronto}.

La seconda modalità, è detta \textit{edge triggered}, e prevede che invece
vengano notificati solo i file descriptor che hanno subito una transizione da
\textsl{non pronti} a \textsl{pronti}. Questo significa che in modalità
\textit{edge triggered} nel caso del precedente esempio il file descriptor
diventato pronto da cui si sono letti solo 1000 byte non verrà nuovamente
notificato come pronto, nonostante siano ancora disponibili in lettura 1000
byte. Solo una volta che si saranno esauriti tutti i dati disponibili, e che
il file descriptor sia tornato non essere pronto, si potrà ricevere una
ulteriore notifica qualora ritornasse pronto.

Nel caso di Linux al momento la sola interfaccia che fornisce questo tipo di
servizio è chiamata \textit{epoll},\footnote{l'interfaccia è stata creata da
  Davide Libenzi, ed è stata introdotta per la prima volta nel kernel 2.5.44,
  ma la sua forma definitiva è stata raggiunta nel kernel 2.5.66, il supporto
  è stato aggiunto nella \acr{glibc} a partire dalla versione 2.3.2.} anche se
sono state in discussione altre interfacce con le quali effettuare lo stesso
tipo di operazioni; \textit{epoll} è in grado di operare sia in modalità
\textit{level triggered} che \textit{edge triggered}.

La prima versione di \textit{epoll} prevedeva l'uso di uno speciale file di
dispositivo, \texttt{/dev/epoll}, per ottenere un file descriptor da
utilizzare con le funzioni dell'interfaccia ma poi si è passati all'uso di
apposite \textit{system call}.  Il primo passo per usare l'interfaccia di
\textit{epoll} è pertanto quello ottenere detto file descriptor chiamando una
delle due funzioni di sistema \funcd{epoll\_create} e \funcd{epoll\_create1},
i cui prototipi sono:

\begin{funcproto}{
\fhead{sys/epoll.h}
\fdecl{int epoll\_create(int size)}
\fdecl{int epoll\_create1(int flags)}

\fdesc{Apre un file descriptor per \textit{epoll}.}
}
{Le funzioni ritornano un file descriptor per \textit{epoll} in caso di
  successo e $-1$ per un errore, nel qual caso \var{errno} assumerà uno dei
  valori:
  \begin{errlist}
  \item[\errcode{EINVAL}] si è specificato un valore di \param{size} non
    positivo o non valido per \param{flags}.
  \item[\errcode{EMFILE}] si è raggiunto il limite sul numero massimo di
    istanze di \textit{epoll} per utente stabilito da
    \sysctlfiled{fs/epoll/max\_user\_instances}.
  \item[\errcode{ENFILE}] si è raggiunto il massimo di file descriptor aperti
    nel sistema.
  \item[\errcode{ENOMEM}] non c'è sufficiente memoria nel kernel per creare
    l'istanza.
  \end{errlist}
}  
\end{funcproto}

Entrambe le funzioni restituiscono un file descriptor, detto anche
\textit{epoll descriptor}; si tratta di un file descriptor speciale (per cui
\func{read} e \func{write} non sono supportate) che viene associato alla
infrastruttura utilizzata dal kernel per gestire la notifica degli eventi, e
che può a sua volta essere messo sotto osservazione con una chiamata a
\func{select}, \func{poll} o \func{epoll\_ctl}; in tal caso risulterà pronto
quando saranno disponibili eventi da notificare riguardo i file descriptor da
lui osservati.\footnote{è anche possibile inviarlo ad un altro processo
  attraverso un socket locale (vedi sez.~\ref{sec:sock_fd_passing}) ma
  l'operazione non ha alcun senso dato che il nuovo processo non avrà a
  disposizione le copie dei file descriptor messe sotto osservazione tramite
  esso.} Una volta che se ne sia terminato l'uso si potranno rilasciare tutte
le risorse allocate chiudendolo semplicemente con \func{close}.

Nel caso di \func{epoll\_create} l'argomento \param{size} serviva a dare
l'indicazione del numero di file descriptor che si vorranno tenere sotto
controllo, e costituiva solo un suggerimento per semplificare l'allocazione di
risorse sufficienti, non un valore massimo, ma a partire dal kernel 2.6.8 esso
viene totalmente ignorato e l'allocazione è sempre dinamica.

La seconda versione della funzione, \func{epoll\_create1} è stata introdotta
come estensione della precedente (è disponibile solo a partire dal kernel
2.6.27) per poter passare dei flag di controllo come maschera binaria in fase
di creazione del file descriptor. Al momento l'unico valore legale
per \param{flags} (a parte lo zero) è \constd{EPOLL\_CLOEXEC}, che consente di
impostare in maniera atomica sul file descriptor il flag di
\textit{close-on-exec} (si è trattato il significato di \const{O\_CLOEXEC} in
sez.~\ref{sec:file_open_close}), senza che sia necessaria una successiva
chiamata a \func{fcntl}.

Una volta ottenuto un file descriptor per \textit{epoll} il passo successivo è
indicare quali file descriptor mettere sotto osservazione e quali operazioni
controllare, per questo si deve usare la seconda funzione di sistema
dell'interfaccia, \funcd{epoll\_ctl}, il cui prototipo è:

\begin{funcproto}{
\fhead{sys/epoll.h}
\fdecl{int epoll\_ctl(int epfd, int op, int fd, struct epoll\_event *event)}

\fdesc{Esegue le operazioni di controllo di \textit{epoll}.}
}

{La funzione ritorna $0$ in caso di successo e $-1$ per un errore, nel qual
  caso \var{errno} assumerà uno dei valori:
  \begin{errlist}
  \item[\errcode{EBADF}] i file descriptor \param{epfd} o \param{fd} non sono
    validi.
  \item[\errcode{EEXIST}] l'operazione richiesta è \const{EPOLL\_CTL\_ADD} ma
    \param{fd} è già stato inserito in \param{epfd}.
  \item[\errcode{EINVAL}] il file descriptor \param{epfd} non è stato ottenuto
    con \func{epoll\_create}, o \param{fd} è lo stesso \param{epfd} o
    l'operazione richiesta con \param{op} non è supportata.
  \item[\errcode{ENOENT}] l'operazione richiesta è \const{EPOLL\_CTL\_MOD} o
    \const{EPOLL\_CTL\_DEL} ma \param{fd} non è inserito in \param{epfd}.
  \item[\errcode{ENOMEM}] non c'è sufficiente memoria nel kernel gestire
    l'operazione richiesta.
  \item[\errcode{ENOSPC}] si è raggiunto il limite massimo di registrazioni
    per utente di file descriptor da osservare imposto da
    \sysctlfiled{fs/epoll/max\_user\_watches}.
  \item[\errcode{EPERM}] il file associato a \param{fd} non supporta l'uso di
    \textit{epoll}.
  \end{errlist}
  }  
\end{funcproto}

La funzione prende sempre come primo argomento un file descriptor di
\textit{epoll}, \param{epfd}, che indica quale istanza di \textit{epoll} usare
e deve pertanto essere stato ottenuto in precedenza con una chiamata a
\func{epoll\_create} o \func{epoll\_create1}. L'argomento \param{fd} indica
invece il file descriptor che si vuole tenere sotto controllo, quest'ultimo
può essere un qualunque file descriptor utilizzabile con \func{poll}, ed anche
un altro file descriptor di \textit{epoll}, ma non lo stesso \param{epfd}.

Il comportamento della funzione viene controllato dal valore dall'argomento
\param{op} che consente di specificare quale operazione deve essere eseguita.
Le costanti che definiscono i valori utilizzabili per \param{op}
sono riportate in tab.~\ref{tab:epoll_ctl_operation}, assieme al significato
delle operazioni cui fanno riferimento.

\begin{table}[htb]
  \centering
  \footnotesize
  \begin{tabular}[c]{|l|p{8cm}|}
    \hline
    \textbf{Valore}  & \textbf{Significato} \\
    \hline
    \hline
    \constd{EPOLL\_CTL\_ADD}& Aggiunge un nuovo file descriptor da osservare
                              \param{fd} alla lista dei file descriptor
                              controllati tramite \param{epfd}, in
                              \param{event} devono essere specificate le
                              modalità di osservazione.\\
    \constd{EPOLL\_CTL\_MOD}& Modifica le modalità di osservazione del file
                              descriptor \param{fd} secondo il contenuto di
                              \param{event}.\\
    \constd{EPOLL\_CTL\_DEL}& Rimuove il file descriptor \param{fd} dalla lista
                              dei file controllati tramite \param{epfd}.\\
   \hline    
  \end{tabular}
  \caption{Valori dell'argomento \param{op} che consentono di scegliere quale
    operazione di controllo effettuare con la funzione \func{epoll\_ctl}.} 
  \label{tab:epoll_ctl_operation}
\end{table}

% era stata aggiunta EPOLL_CTL_DISABLE in previsione del kernel 3.7, vedi
% http://lwn.net/Articles/520012/ e http://lwn.net/Articles/520198/
% ma non è mai stata inserita.

Le modalità di utilizzo di \textit{epoll} prevedono che si definisca qual'è
l'insieme dei file descriptor da tenere sotto controllo utilizzando una serie
di chiamate a \const{EPOLL\_CTL\_ADD}.\footnote{un difetto dell'interfaccia è
  che queste chiamate devono essere ripetute per ciascun file descriptor,
  incorrendo in una perdita di prestazioni qualora il numero di file
  descriptor sia molto grande; per questo è stato proposto di introdurre come
  estensione una funzione \code{epoll\_ctlv} che consenta di effettuare con
  una sola chiamata le impostazioni per un blocco di file descriptor.} L'uso
di \const{EPOLL\_CTL\_MOD} consente in seguito di modificare le modalità di
osservazione di un file descriptor che sia già stato aggiunto alla lista di
osservazione. Qualora non si abbia più interesse nell'osservazione di un file
descriptor lo si può rimuovere dalla lista associata a \param{epfd} con
\const{EPOLL\_CTL\_DEL}.

Anche se è possibile tenere sotto controllo lo stesso file descriptor in due
istanze distinte di \textit{epoll} in genere questo è sconsigliato in quanto
entrambe riceveranno le notifiche, e gestire correttamente le notifiche
multiple richiede molta attenzione. Se invece si cerca di inserire due volte
lo stesso file descriptor nella stessa istanza di \textit{epoll} la funzione
fallirà con un errore di \errval{EEXIST}.  Tuttavia è possibile inserire nella
stessa istanza file descriptor duplicati (si ricordi quanto visto in
sez.~\ref{sec:file_dup}), una tecnica che può essere usata per registrarli con
un valore diverso per \param{events} e classificare così diversi tipi di
eventi.

Si tenga presente che quando si chiude un file descriptor questo, se era stato
posto sotto osservazione da una istanza di \textit{epoll}, viene rimosso
automaticamente solo nel caso esso sia l'unico riferimento al file aperto
sottostante (più precisamente alla struttura \kstruct{file}, si ricordi
fig.~\ref{fig:file_dup}) e non è necessario usare
\const{EPOLL\_CTL\_DEL}. Questo non avviene qualora esso sia stato duplicato
(perché la suddetta struttura non viene disallocata) e si potranno ricevere
eventi ad esso relativi anche dopo che lo si è chiuso; per evitare
l'inconveniente è necessario rimuoverlo esplicitamente con
\const{EPOLL\_CTL\_DEL}.

L'ultimo argomento, \param{event}, deve essere un puntatore ad una struttura
di tipo \struct{epoll\_event}, ed ha significato solo con le operazioni
\const{EPOLL\_CTL\_MOD} e \const{EPOLL\_CTL\_ADD}, per le quali serve ad
indicare quale tipo di evento relativo ad \param{fd} si vuole che sia tenuto
sotto controllo.  L'argomento viene ignorato con l'operazione
\const{EPOLL\_CTL\_DEL}.\footnote{fino al kernel 2.6.9 era comunque richiesto
  che questo fosse un puntatore valido, anche se poi veniva ignorato; a
  partire dal 2.6.9 si può specificare anche un valore \val{NULL} ma se si
  vuole mantenere la compatibilità con le versioni precedenti occorre usare un
  puntatore valido.}

\begin{figure}[!htb]
  \footnotesize \centering
  \begin{minipage}[c]{0.90\textwidth}
    \includestruct{listati/epoll_event.h}
  \end{minipage} 
  \normalsize 
  \caption{La struttura \structd{epoll\_event}, che consente di specificare
    gli eventi associati ad un file descriptor controllato con
    \textit{epoll}.}
  \label{fig:epoll_event}
\end{figure}

La struttura \struct{epoll\_event} è l'analoga di \struct{pollfd} e come
quest'ultima serve sia in ingresso (quando usata con \func{epoll\_ctl}) ad
impostare quali eventi osservare, che in uscita (nei risultati ottenuti con
\func{epoll\_wait}) per ricevere le notifiche degli eventi avvenuti.  La sua
definizione è riportata in fig.~\ref{fig:epoll_event}. 

Il primo campo, \var{events}, è una maschera binaria in cui ciascun bit
corrisponde o ad un tipo di evento, o una modalità di notifica; detto campo
deve essere specificato come OR aritmetico delle costanti riportate in
tab.~\ref{tab:epoll_events}. Nella prima parte della tabella si sono indicate
le costanti che permettono di indicare il tipo di evento, che sono le
equivalenti delle analoghe di tab.~\ref{tab:file_pollfd_flags} per
\func{poll}. Queste sono anche quelle riportate nella struttura
\struct{epoll\_event} restituita da \func{epoll\_wait} per indicare il tipo di
evento presentatosi, insieme a quelle della seconda parte della tabella, che
vengono comunque riportate anche se non le si sono impostate con
\func{epoll\_ctl}. La terza parte della tabella contiene le costanti che
modificano le modalità di notifica.

\begin{table}[htb]
  \centering
  \footnotesize
  \begin{tabular}[c]{|l|p{10cm}|}
    \hline
    \textbf{Valore}  & \textbf{Significato} \\
    \hline
    \hline
    \constd{EPOLLIN}     & Il file è pronto per le operazioni di lettura
                          (analogo di \const{POLLIN}).\\
    \constd{EPOLLOUT}    & Il file è pronto per le operazioni di scrittura
                          (analogo di \const{POLLOUT}).\\
    \constd{EPOLLRDHUP}  & L'altro capo di un socket di tipo
                          \const{SOCK\_STREAM} (vedi sez.~\ref{sec:sock_type})
                          ha chiuso la connessione o il capo in scrittura
                          della stessa (vedi
                          sez.~\ref{sec:TCP_shutdown}).\footnotemark\\
    \constd{EPOLLPRI}    & Ci sono dati urgenti disponibili in lettura (analogo
                          di \const{POLLPRI}); questa condizione viene comunque
                          riportata in uscita, e non è necessaria impostarla
                          in ingresso.\\ 
    \hline
    \constd{EPOLLERR}    & Si è verificata una condizione di errore 
                          (analogo di \const{POLLERR}); questa condizione
                          viene comunque riportata in uscita, e non è
                          necessaria impostarla in ingresso.\\
    \constd{EPOLLHUP}    & Si è verificata una condizione di hung-up; questa
                          condizione viene comunque riportata in uscita, e non
                          è necessaria impostarla in ingresso.\\
    \hline
    \constd{EPOLLET}     & Imposta la notifica in modalità \textit{edge
                            triggered} per il file descriptor associato.\\ 
    \constd{EPOLLONESHOT}& Imposta la modalità \textit{one-shot} per il file
                          descriptor associato (questa modalità è disponibile
                          solo a partire dal kernel 2.6.2).\\
    \constd{EPOLLWAKEUP} & Attiva la prevenzione della sospensione del sistema
                          se il file descriptor che si è marcato con esso
                          diventa pronto (aggiunto a partire dal kernel 3.5),
                          può essere impostato solo dall'amministratore (o da
                          un processo con la capacità
                          \const{CAP\_BLOCK\_SUSPEND}).\\ 
    \hline
  \end{tabular}
  \caption{Costanti che identificano i bit del campo \param{events} di
    \struct{epoll\_event}.}
  \label{tab:epoll_events}
\end{table}

\footnotetext{questa modalità è disponibile solo a partire dal kernel 2.6.17,
  ed è utile per riconoscere la chiusura di una connessione dall'altro capo di
  un socket quando si lavora in modalità \textit{edge triggered}.}

% TODO aggiunto con il kernel 4.5  EPOLLEXCLUSIVE, vedi
% http://lwn.net/Articles/633422/#excl 

Il secondo campo, \var{data}, è una \dirct{union} che serve a identificare il
file descriptor a cui si intende fare riferimento, ed in astratto può
contenere un valore qualsiasi (specificabile in diverse forme) che ne permetta
una indicazione univoca. Il modo più comune di usarlo però è quello in cui si
specifica il terzo argomento di \func{epoll\_ctl} nella forma
\var{event.data.fd}, assegnando come valore di questo campo lo stesso valore
dell'argomento \param{fd}, cosa che permette una immediata identificazione del
file descriptor.

% TODO verificare se prima o poi epoll_ctlv verrà introdotta

Le impostazioni di default prevedono che la notifica degli eventi richiesti
sia effettuata in modalità \textit{level triggered}, a meno che sul file
descriptor non si sia impostata la modalità \textit{edge triggered},
registrandolo con \const{EPOLLET} attivo nel campo \var{events}.  

Infine una particolare modalità di notifica è quella impostata con
\const{EPOLLONESHOT}: a causa dell'implementazione di \textit{epoll} infatti
quando si è in modalità \textit{edge triggered} l'arrivo in rapida successione
di dati in blocchi separati (questo è tipico con i socket di rete, in quanto i
dati arrivano a pacchetti) può causare una generazione di eventi (ad esempio
segnalazioni di dati in lettura disponibili) anche se la condizione è già
stata rilevata (si avrebbe cioè una rottura della logica \textit{edge
  triggered}).

Anche se la situazione è facile da gestire, la si può evitare utilizzando
\const{EPOLLONESHOT} per impostare la modalità \textit{one-shot}, in cui la
notifica di un evento viene effettuata una sola volta, dopo di che il file
descriptor osservato, pur restando nella lista di osservazione, viene
automaticamente disattivato (la cosa avviene contestualmente al ritorno di
\func{epoll\_wait} a causa dell'evento in questione) e per essere riutilizzato
dovrà essere riabilitato esplicitamente con una successiva chiamata con
\const{EPOLL\_CTL\_MOD}.

Una volta impostato l'insieme di file descriptor che si vogliono osservare con
i relativi eventi, la funzione di sistema che consente di attendere
l'occorrenza di uno di tali eventi è \funcd{epoll\_wait}, il cui prototipo è:

\begin{funcproto}{
\fhead{sys/epoll.h}
\fdecl{int epoll\_wait(int epfd, struct epoll\_event * events, int maxevents,
  int timeout)}

\fdesc{Attende che uno dei file descriptor osservati sia pronto.}
}

{La funzione ritorna il numero di file descriptor pronti in caso di successo e
  $-1$ per un errore, nel qual caso \var{errno} assumerà uno dei valori:
  \begin{errlist}
  \item[\errcode{EBADF}] il file descriptor \param{epfd} non è valido.
  \item[\errcode{EFAULT}] il puntatore \param{events} non è valido.
  \item[\errcode{EINTR}] la funzione è stata interrotta da un segnale prima
    della scadenza di \param{timeout}.
  \item[\errcode{EINVAL}] il file descriptor \param{epfd} non è stato ottenuto
    con \func{epoll\_create}, o \param{maxevents} non è maggiore di zero.
  \end{errlist}
}  
\end{funcproto}

La funzione si blocca in attesa di un evento per i file descriptor registrati
nella lista di osservazione di \param{epfd} fino ad un tempo massimo
specificato in millisecondi tramite l'argomento \param{timeout}. Gli eventi
registrati vengono riportati in un vettore di strutture \struct{epoll\_event}
(che deve essere stato allocato in precedenza) all'indirizzo indicato
dall'argomento \param{events}, fino ad un numero massimo di eventi impostato
con l'argomento \param{maxevents}.

La funzione ritorna il numero di eventi rilevati, o un valore nullo qualora
sia scaduto il tempo massimo impostato con \param{timeout}. Per quest'ultimo,
oltre ad un numero di millisecondi, si può utilizzare il valore nullo, che
indica di non attendere e ritornare immediatamente (anche in questo caso il
valore di ritorno sarà nullo) o il valore $-1$, che indica un'attesa
indefinita. L'argomento \param{maxevents} dovrà invece essere sempre un intero
positivo.

Come accennato la funzione restituisce i suoi risultati nel vettore di
strutture \struct{epoll\_event} puntato da \param{events}; in tal caso nel
campo \param{events} di ciascuna di esse saranno attivi i flag relativi agli
eventi accaduti, mentre nel campo \var{data} sarà restituito il valore che era
stato impostato per il file descriptor per cui si è verificato l'evento quando
questo era stato registrato con le operazioni \const{EPOLL\_CTL\_MOD} o
\const{EPOLL\_CTL\_ADD}, in questo modo il campo \var{data} consente di
identificare il file descriptor, ed è per questo che, come accennato, è
consuetudine usare per \var{data} il valore del file descriptor stesso.

Si ricordi che le occasioni per cui \func{epoll\_wait} ritorna dipendono da
come si è impostata la modalità di osservazione (se \textit{level triggered} o
\textit{edge triggered}) del singolo file descriptor. L'interfaccia assicura
che se arrivano più eventi fra due chiamate successive ad \func{epoll\_wait}
questi vengano combinati. Inoltre qualora su un file descriptor fossero
presenti eventi non ancora notificati, e si effettuasse una modifica
dell'osservazione con \const{EPOLL\_CTL\_MOD}, questi verrebbero riletti alla
luce delle modifiche.

Si tenga presente infine che con l'uso della modalità \textit{edge triggered}
il ritorno di \func{epoll\_wait} avviene solo quando il file descriptor ha
cambiato stato diventando pronto. Esso non sarà riportato nuovamente fino ad
un altro cambiamento di stato, per cui occorre assicurarsi di aver
completamente esaurito le operazioni su di esso.  Questa condizione viene
generalmente rilevata dall'occorrere di un errore di \errcode{EAGAIN} al
ritorno di una \func{read} o una \func{write}, (è opportuno ricordare ancora
una volta che l'uso dell'\textit{I/O multiplexing} richiede di operare sui
file in modalità non bloccante) ma questa non è la sola modalità possibile, ad
esempio la condizione può essere riconosciuta anche per il fatto che sono
stati restituiti meno dati di quelli richiesti.

Si tenga presente che in modalità \textit{edge triggered}, dovendo esaurire le
attività di I/O dei file descriptor risultati pronti per poter essere
rinotificati, la gestione elementare per cui li si trattano uno per uno in
sequenza può portare ad un effetto denominato \textit{starvation}
(``\textsl{carestia}'').  Si rischia cioè di concentrare le operazioni sul
primo file descriptor che dispone di molti dati, prolungandole per tempi molto
lunghi con un ritardo che può risultare eccessivo nei confronti di quelle da
eseguire sugli altri che verrebbero dopo.  Per evitare questo tipo di
problematiche viene consigliato di usare \func{epoll\_wait} per registrare un
elenco dei file descriptor da gestire, e di trattarli a turno in maniera più
equa.

Come già per \func{select} e \func{poll} anche per l'interfaccia di
\textit{epoll} si pone il problema di gestire l'attesa di segnali e di dati
contemporaneamente.  Valgono le osservazioni fatte in
sez.~\ref{sec:file_select}, e per poterlo fare di nuovo è necessaria una
variante della funzione di attesa che consenta di reimpostare all'uscita una
maschera di segnali, analoga alle estensioni \func{pselect} e \func{ppoll} che
abbiamo visto in precedenza per \func{select} e \func{poll}. In questo caso la
funzione di sistema si chiama \funcd{epoll\_pwait}\footnote{la funzione è
  stata introdotta a partire dal kernel 2.6.19, ed è, come tutta l'interfaccia
  di \textit{epoll}, specifica di Linux.} ed il suo prototipo è:

\begin{funcproto}{
\fhead{sys/epoll.h}
\fdecl{int epoll\_pwait(int epfd, struct epoll\_event * events, int maxevents, 
    int timeout, \\
\phantom{int epoll\_pwait(}const sigset\_t *sigmask)}

\fdesc{Attende che uno dei file descriptor osservati sia pronto, mascherando
    i segnali.}  }

{La funzione ritorna il numero di file descriptor pronti in caso di successo e
  $-1$ per un errore, nel qual caso \var{errno} assumerà uno dei valori già
  visti con \func{epoll\_wait}.

}  
\end{funcproto}

La funzione è del tutto analoga \func{epoll\_wait}, soltanto che alla sua
uscita viene ripristinata la maschera di segnali originale, sostituita durante
l'esecuzione da quella impostata con l'argomento \param{sigmask}; in sostanza
la chiamata a questa funzione è equivalente al seguente codice, eseguito però
in maniera atomica:
\includecodesnip{listati/epoll_pwait_means.c} 

Si tenga presente che come le precedenti funzioni di \textit{I/O multiplexing}
anche le funzioni dell'interfaccia di \textit{epoll} vengono utilizzate
prevalentemente con i server di rete, quando si devono tenere sotto
osservazione un gran numero di socket; per questo motivo rimandiamo anche in
questo caso la trattazione di un esempio concreto a quando avremo esaminato in
dettaglio le caratteristiche dei socket; in particolare si potrà trovare un
programma che utilizza questa interfaccia in sez.~\ref{sec:TCP_serv_epoll}.

\itindend{epoll}


\subsection{La notifica di eventi tramite file descriptor}
\label{sec:sig_signalfd_eventfd}

Abbiamo visto in sez.~\ref{sec:file_select} come il meccanismo classico delle
notifiche di eventi tramite i segnali, presente da sempre nei sistemi
unix-like, porti a notevoli problemi nell'interazione con le funzioni per
l'\textit{I/O multiplexing}, tanto che per evitare possibili \textit{race
  condition} sono state introdotte estensioni dello standard POSIX e funzioni
apposite come \func{pselect}, \func{ppoll} e \func{epoll\_pwait}.

Benché i segnali siano il meccanismo più usato per effettuare notifiche ai
processi, la loro interfaccia di programmazione, che comporta l'esecuzione di
una funzione di gestione in maniera asincrona e totalmente scorrelata
dall'ordinario flusso di esecuzione del processo, si è però dimostrata quasi
subito assai problematica. Oltre ai limiti relativi ai limiti al cosa si può
fare all'interno della funzione del gestore di segnali (quelli illustrati in
sez.~\ref{sec:sig_signal_handler}), c'è il problema più generale consistente
nel fatto che questa modalità di funzionamento cozza con altre interfacce di
programmazione previste dal sistema in cui si opera in maniera
\textsl{sincrona}, come quelle dell'\textit{I/O multiplexing} appena
illustrate.

In questo tipo di interfacce infatti ci si aspetta che il processo gestisca
gli eventi a cui deve reagire in maniera sincrona generando le opportune
risposte, mentre con l'arrivo di un segnale si possono avere interruzioni
asincrone in qualunque momento.  Questo comporta la necessità di dover
gestire, quando si deve tener conto di entrambi i tipi di eventi, le
interruzioni delle funzioni di attesa sincrone, ed evitare possibili
\textit{race conditions}. In sostanza se non ci fossero i segnali non ci
sarebbe da preoccuparsi, fintanto che si effettuano operazioni all'interno di
un processo, della non atomicità delle \textit{system call} lente che vengono
interrotte e devono essere riavviate.

Abbiamo visto però in sez.~\ref{sec:sig_real_time} che insieme ai segnali
\textit{real-time} sono state introdotte anche delle interfacce di gestione
sincrona dei segnali, con la funzione \func{sigwait} e le sue affini. Queste
funzioni consentono di gestire i segnali bloccando un processo fino alla
avvenuta ricezione e disabilitando l'esecuzione asincrona rispetto al resto
del programma del gestore del segnale. Questo consente di risolvere i problemi
di atomicità nella gestione degli eventi associati ai segnali, avendo tutto il
controllo nel flusso principale del programma, ottenendo così una gestione
simile a quella dell'\textit{I/O multiplexing}, ma non risolve i problemi
delle interazioni con quest'ultimo, perché o si aspetta la ricezione di un
segnale o si aspetta che un file descriptor sia accessibile e nessuna delle
rispettive funzioni consente di fare contemporaneamente entrambe le cose.

Per risolvere questo problema nello sviluppo del kernel si è pensato di
introdurre un meccanismo alternativo per la notifica dei segnali (esteso anche
ad altri eventi generici) che, ispirandosi di nuovo alla filosofia di Unix per
cui tutto è un file, consentisse di eseguire la notifica con l'uso di
opportuni file descriptor. Ovviamente si tratta di una funzionalità specifica
di Linux, non presente in altri sistemi unix-like, e non prevista da nessuno
standard, per cui va evitata se si ha a cuore la portabilità.

In sostanza, come per \func{sigwait}, si può disabilitare l'esecuzione di un
gestore in occasione dell'arrivo di un segnale, e rilevarne l'avvenuta
ricezione leggendone la notifica tramite l'uso di uno speciale file
descriptor. Trattandosi di un file descriptor questo potrà essere tenuto sotto
osservazione con le ordinarie funzioni dell'\textit{I/O multiplexing} (vale a
dire con le solite \func{select}, \func{poll} e \func{epoll\_wait}) allo
stesso modo di quelli associati a file o socket, per cui alla fine si potrà
attendere in contemporanea sia l'arrivo del segnale che la disponibilità di
accesso ai dati relativi a questi ultimi.

La funzione di sistema che permette di abilitare la ricezione dei segnali
tramite file descriptor è \funcd{signalfd},\footnote{in realtà quella
  riportata è l'interfaccia alla funzione fornita dalla \acr{glibc}, esistono
  infatti due versioni diverse della \textit{system call}; una prima versione,
  \func{signalfd}, introdotta nel kernel 2.6.22 e disponibile con la
  \acr{glibc} 2.8 che non supporta l'argomento \texttt{flags}, ed una seconda
  versione, \funcm{signalfd4}, introdotta con il kernel 2.6.27 e che è quella
  che viene sempre usata a partire dalla \acr{glibc} 2.9, che prende un
  argomento aggiuntivo \code{size\_t sizemask} che indica la dimensione della
  maschera dei segnali, il cui valore viene impostato automaticamente dalla
  \acr{glibc}.}  il cui prototipo è:

\begin{funcproto}{
\fhead{sys/signalfd.h}
\fdecl{int signalfd(int fd, const sigset\_t *mask, int flags)}

\fdesc{Crea o modifica un file descriptor per la ricezione dei segnali.}
}

{La funzione ritorna un numero di file descriptor in caso di successo e $-1$
  per un errore, nel qual caso \var{errno} assumerà uno dei valori:
  \begin{errlist}
  \item[\errcode{EBADF}] il valore \param{fd} non indica un file descriptor.
  \item[\errcode{EINVAL}] il file descriptor \param{fd} non è stato ottenuto
    con \func{signalfd} o il valore di \param{flags} non è valido.
  \item[\errcode{ENODEV}] il kernel non può montare internamente il
    dispositivo per la gestione anonima degli \textit{inode}
    associati al file descriptor.
  \item[\errcode{ENOMEM}] non c'è memoria sufficiente per creare un nuovo file
    descriptor di \func{signalfd}.
  \end{errlist}
  ed inoltre \errval{EMFILE} e \errval{ENFILE} nel loro significato generico.
  
}  
\end{funcproto}

La funzione consente di creare o modificare le caratteristiche di un file
descriptor speciale su cui ricevere le notifiche della ricezione di
segnali. Per creare ex-novo uno di questi file descriptor è necessario passare
$-1$ come valore per l'argomento \param{fd}, ogni altro valore positivo verrà
invece interpretato come il numero del file descriptor (che deve esser stato
precedentemente creato sempre con \func{signalfd}) di cui si vogliono
modificare le caratteristiche. Nel primo caso la funzione ritornerà il valore
del nuovo file descriptor e nel secondo caso il valore indicato
con \param{fd}, in caso di errore invece verrà restituito $-1$.

L'elenco dei segnali che si vogliono gestire con \func{signalfd} deve essere
specificato tramite l'argomento \param{mask}. Questo deve essere passato come
puntatore ad una maschera di segnali creata con l'uso delle apposite macro già
illustrate in sez.~\ref{sec:sig_sigset}. La maschera deve indicare su quali
segnali si intende operare con \func{signalfd}; l'elenco può essere modificato
con una successiva chiamata a \func{signalfd}. Dato che \signal{SIGKILL} e
\signal{SIGSTOP} non possono essere intercettati (e non prevedono neanche la
possibilità di un gestore) un loro inserimento nella maschera verrà ignorato
senza generare errori.

L'argomento \param{flags} consente di impostare direttamente in fase di
creazione due flag per il file descriptor analoghi a quelli che si possono
impostare con una creazione ordinaria con \func{open}, evitando una
impostazione successiva con \func{fcntl} (si ricordi che questo è un argomento
aggiuntivo, introdotto con la versione fornita a partire dal kernel 2.6.27,
per kernel precedenti il valore deve essere nullo).  L'argomento deve essere
specificato come maschera binaria dei valori riportati in
tab.~\ref{tab:signalfd_flags}.

\begin{table}[htb]
  \centering
  \footnotesize
  \begin{tabular}[c]{|l|p{8cm}|}
    \hline
    \textbf{Valore}  & \textbf{Significato} \\
    \hline
    \hline
    \constd{SFD\_NONBLOCK}&imposta sul file descriptor il flag di
                           \const{O\_NONBLOCK} per renderlo non bloccante.\\ 
    \constd{SFD\_CLOEXEC}& imposta il flag di \const{O\_CLOEXEC} per la
                           chiusura automatica del file descriptor nella
                           esecuzione di \func{exec}.\\
    \hline    
  \end{tabular}
  \caption{Valori dell'argomento \param{flags} per la funzione \func{signalfd}
    che consentono di impostare i flag del file descriptor.} 
  \label{tab:signalfd_flags}
\end{table}

Si tenga presente che la chiamata a \func{signalfd} non disabilita la gestione
ordinaria dei segnali indicati da \param{mask}; questa, se si vuole effettuare
la ricezione tramite il file descriptor, dovrà essere disabilitata
esplicitamente bloccando gli stessi segnali con \func{sigprocmask}, altrimenti
verranno comunque eseguite le azioni di default (o un eventuale gestore
installato in precedenza). Il blocco non ha invece nessun effetto sul file
descriptor restituito da \func{signalfd}, dal quale sarà possibile pertanto
ricevere qualunque segnale, anche se questo risultasse bloccato.

Si tenga presente inoltre che la lettura di una struttura
\struct{signalfd\_siginfo} relativa ad un segnale pendente è equivalente alla
esecuzione di un gestore, vale a dire che una volta letta il segnale non sarà
più pendente e non potrà essere ricevuto, qualora si ripristino le normali
condizioni di gestione, né da un gestore, né dalla funzione \func{sigwaitinfo}.

Come anticipato, essendo questo lo scopo principale della nuova interfaccia,
il file descriptor può essere tenuto sotto osservazione tramite le funzioni
dell'\textit{I/O multiplexing} (vale a dire con le solite \func{select},
\func{poll} e \func{epoll\_wait}), e risulterà accessibile in lettura quando
uno o più dei segnali indicati tramite \param{mask} sarà pendente.

La funzione può essere chiamata più volte dallo stesso processo, consentendo
così di tenere sotto osservazione segnali diversi tramite file descriptor
diversi. Inoltre è anche possibile tenere sotto osservazione lo stesso segnale
con più file descriptor, anche se la pratica è sconsigliata; in tal caso la
ricezione del segnale potrà essere effettuata con una lettura da uno qualunque
dei file descriptor a cui è associato, ma questa potrà essere eseguita
soltanto una volta. Questo significa che tutti i file descriptor su cui è
presente lo stesso segnale risulteranno pronti in lettura per le funzioni di
\textit{I/O multiplexing}, ma una volta eseguita la lettura su uno di essi il
segnale sarà considerato ricevuto ed i relativi dati non saranno più
disponibili sugli altri file descriptor, che (a meno di una ulteriore
occorrenza del segnale nel frattempo) di non saranno più pronti.

Quando il file descriptor per la ricezione dei segnali non serve più potrà
essere chiuso con \func{close} liberando tutte le risorse da esso allocate. In
tal caso qualora vi fossero segnali pendenti questi resteranno tali, e
potranno essere ricevuti normalmente una volta che si rimuova il blocco
imposto con \func{sigprocmask}.

Oltre a poter essere usato con le funzioni dell'\textit{I/O multiplexing}, il
file descriptor restituito da \func{signalfd} cerca di seguire la semantica di
un sistema unix-like anche con altre \textit{system call}; in particolare esso
resta aperto (come ogni altro file descriptor) attraverso una chiamata ad
\func{exec}, a meno che non lo si sia creato con il flag di
\const{SFD\_CLOEXEC} o si sia successivamente impostato il
\textit{close-on-exec} con \func{fcntl}. Questo comportamento corrisponde
anche alla ordinaria semantica relativa ai segnali bloccati, che restano
pendenti attraverso una \func{exec}.

Analogamente il file descriptor resta sempre disponibile attraverso una
\func{fork} per il processo figlio, che ne riceve una copia; in tal caso però
il figlio potrà leggere dallo stesso soltanto i dati relativi ai segnali
ricevuti da lui stesso. Nel caso di \textit{thread} viene nuovamente seguita
la semantica ordinaria dei segnali, che prevede che un singolo \textit{thread}
possa ricevere dal file descriptor solo le notifiche di segnali inviati
direttamente a lui o al processo in generale, e non quelli relativi ad altri
\textit{thread} appartenenti allo stesso processo.

L'interfaccia fornita da \func{signalfd} prevede che la ricezione dei segnali
sia eseguita leggendo i dati relativi ai segnali pendenti dal file descriptor
restituito dalla funzione con una normalissima \func{read}.  Qualora non vi
siano segnali pendenti la \func{read} si bloccherà a meno di non aver
impostato la modalità di I/O non bloccante sul file descriptor, o direttamente
in fase di creazione con il flag \const{SFD\_NONBLOCK}, o in un momento
successivo con \func{fcntl}.  

\begin{figure}[!htb]
  \footnotesize \centering
  \begin{minipage}[c]{0.95\textwidth}
    \includestruct{listati/signalfd_siginfo.h}
  \end{minipage} 
  \normalsize 
  \caption{La struttura \structd{signalfd\_siginfo}, restituita in lettura da
    un file descriptor creato con \func{signalfd}.}
  \label{fig:signalfd_siginfo}
\end{figure}

I dati letti dal file descriptor vengono scritti sul buffer indicato come
secondo argomento di \func{read} nella forma di una sequenza di una o più
strutture \struct{signalfd\_siginfo} (la cui definizione si è riportata in
fig.~\ref{fig:signalfd_siginfo}) a seconda sia della dimensione del buffer che
del numero di segnali pendenti. Per questo motivo il buffer deve essere almeno
di dimensione pari a quella di \struct{signalfd\_siginfo}, qualora sia di
dimensione maggiore potranno essere letti in unica soluzione i dati relativi
ad eventuali più segnali pendenti, fino al numero massimo di strutture
\struct{signalfd\_siginfo} che possono rientrare nel buffer.

\begin{figure}[!htb]
  \footnotesize \centering
  \begin{minipage}[c]{\codesamplewidth}
    \includecodesample{listati/FifoReporter-init.c}
  \end{minipage} 
  \normalsize 
  \caption{Sezione di inizializzazione del codice del programma
    \file{FifoReporter.c}.}
  \label{fig:fiforeporter_code_init}
\end{figure}

Il contenuto di \struct{signalfd\_siginfo} ricalca da vicino quella
dell'analoga struttura \struct{siginfo\_t} (illustrata in
fig.~\ref{fig:sig_siginfo_t}) usata dall'interfaccia ordinaria dei segnali, e
restituisce dati simili. Come per \struct{siginfo\_t} i campi che vengono
avvalorati dipendono dal tipo di segnale e ricalcano i valori che abbiamo già
illustrato in sez.~\ref{sec:sig_sigaction}.\footnote{si tenga presente però
  che per un bug i kernel fino al 2.6.25 non avvalorano correttamente i campi
  \var{ssi\_ptr} e \var{ssi\_int} per segnali inviati con \func{sigqueue}.}

Come esempio di questa nuova interfaccia ed anche come esempio di applicazione
della interfaccia di \textit{epoll}, si è scritto un programma elementare che
stampi sullo \textit{standard output} sia quanto viene scritto da terzi su una
\textit{named fifo}, che l'avvenuta ricezione di alcuni segnali.  Il codice
completo si trova al solito nei sorgenti allegati alla guida (nel file
\texttt{FifoReporter.c}).

In fig.~\ref{fig:fiforeporter_code_init} si è riportata la parte iniziale del
programma in cui vengono effettuate le varie inizializzazioni necessarie per
l'uso di \textit{epoll} e \func{signalfd}, a partire (\texttt{\small 12-16})
dalla definizione delle varie variabili e strutture necessarie. Al solito si è
tralasciata la parte dedicata alla decodifica delle opzioni che consentono ad
esempio di cambiare il nome del file associato alla \textit{fifo}.

Il primo passo (\texttt{\small 19-20}) è la creazione di un file descriptor
\texttt{epfd} di \textit{epoll} con \func{epoll\_create} che è quello che
useremo per il controllo degli altri.  É poi necessario disabilitare la
ricezione dei segnali (nel caso \signal{SIGINT}, \signal{SIGQUIT} e
\signal{SIGTERM}) per i quali si vuole la notifica tramite file
descriptor. Per questo prima li si inseriscono (\texttt{\small 22-25}) in una
maschera di segnali \texttt{sigmask} che useremo con (\texttt{\small 26})
\func{sigprocmask} per disabilitarli.  Con la stessa maschera si potrà per
passare all'uso (\texttt{\small 28-29}) di \func{signalfd} per abilitare la
notifica sul file descriptor \var{sigfd}. Questo poi (\texttt{\small 30-33})
dovrà essere aggiunto con \func{epoll\_ctl} all'elenco di file descriptor
controllati con \texttt{epfd}.

Occorrerà infine (\texttt{\small 35-38}) creare la \textit{named fifo} se
questa non esiste ed aprirla per la lettura (\texttt{\small 39-40}); una volta
fatto questo sarà necessario aggiungere il relativo file descriptor
(\var{fifofd}) a quelli osservati da \textit{epoll} in maniera del tutto
analoga a quanto fatto con quello relativo alla notifica dei segnali.

\begin{figure}[!htb]
  \footnotesize \centering
  \begin{minipage}[c]{\codesamplewidth}
    \includecodesample{listati/FifoReporter-main.c}
  \end{minipage} 
  \normalsize 
  \caption{Ciclo principale del codice del programma \file{FifoReporter.c}.}
  \label{fig:fiforeporter_code_body}
\end{figure}

Una volta completata l'inizializzazione verrà eseguito indefinitamente il
ciclo principale del programma (\texttt{\small 2-45}) che si è riportato in
fig.~\ref{fig:fiforeporter_code_body}, fintanto che questo non riceva un
segnale di \signal{SIGINT} (ad esempio con la pressione di \texttt{C-c}). Il
ciclo prevede che si attenda (\texttt{\small 2-3}) la presenza di un file
descriptor pronto in lettura con \func{epoll\_wait} (si ricordi che entrambi i
file descriptor \var{fifofd} e \var{sigfd} sono stati posti in osservazioni
per eventi di tipo \const{EPOLLIN}) che si bloccherà fintanto che non siano
stati scritti dati sulla \textit{fifo} o che non sia arrivato un
segnale.\footnote{per semplificare il codice non si è trattato il caso in cui
  \func{epoll\_wait} viene interrotta da un segnale, assumendo che tutti
  quelli che possano interessare siano stati predisposti per la notifica
  tramite file descriptor, per gli altri si otterrà semplicemente l'uscita dal
  programma.}

Anche se in questo caso i file descriptor pronti possono essere al più due, si
è comunque adottato un approccio generico in cui questi verranno letti
all'interno di un opportuno ciclo (\texttt{\small 5-44}) sul numero
restituito da \func{epoll\_wait}, esaminando i risultati presenti nel vettore
\var{events} all'interno di una catena di condizionali alternativi sul valore
del file descriptor riconosciuto come pronto, controllando cioè a quale dei
due file descriptor possibili corrisponde il campo relativo,
\var{events[i].data.fd}.

Il primo condizionale (\texttt{\small 6-24}) è relativo al caso che si sia
ricevuto un segnale e che il file descriptor pronto corrisponda
(\texttt{\small 6}) a \var{sigfd}. Dato che in generale si possono ricevere
anche notifiche relativi a più di un singolo segnale, si è scelto di leggere
una struttura \struct{signalfd\_siginfo} alla volta, eseguendo la lettura
all'interno di un ciclo (\texttt{\small 8-24}) che prosegue fintanto che vi
siano dati da leggere.

Per questo ad ogni lettura si esamina (\texttt{\small 9-14}) se il valore di
ritorno della funzione \func{read} è negativo, uscendo dal programma
(\texttt{\small 11}) in caso di errore reale, o terminando il ciclo
(\texttt{\small 13}) con un \texttt{break} qualora si ottenga un errore di
\errcode{EAGAIN} per via dell'esaurimento dei dati. Si ricordi infatti come
sia la \textit{fifo} che il file descriptor per i segnali siano stati aperti in
modalità non-bloccante, come previsto per l’\textit{I/O multiplexing},
pertanto ci si aspetta di ricevere un errore di \errcode{EAGAIN} quando non vi
saranno più dati da leggere.

In presenza di dati invece il programma proseguirà l'esecuzione stampando
(\texttt{\small 19-20}) il nome del segnale ottenuto all'interno della
struttura \struct{signalfd\_siginfo} letta in \var{siginf} ed il \textit{pid}
del processo da cui lo ha ricevuto;\footnote{per la stampa si è usato il
  vettore \var{sig\_names} a ciascun elemento del quale corrisponde il nome
  del segnale avente il numero corrispondente, la cui definizione si è omessa
  dal codice di fig.~\ref{fig:fiforeporter_code_init} per brevità.} inoltre
(\texttt{\small 21-24}) si controllerà anche se il segnale ricevuto è
\signal{SIGINT}, che si è preso come segnale da utilizzare per la terminazione
del programma, che verrà eseguita dopo aver rimosso il file della \textit{name
  fifo}.
 
Il secondo condizionale (\texttt{\small 26-39}) è invece relativo al caso in
cui ci siano dati pronti in lettura sulla \textit{fifo} e che il file
descriptor pronto corrisponda (\texttt{\small 26}) a \var{fifofd}. Di nuovo si
effettueranno le letture in un ciclo (\texttt{\small 28-39}) ripetendole fin
tanto che la funzione \func{read} non restituisce un errore di
\errcode{EAGAIN} (\texttt{\small 29-35}). Il procedimento è lo stesso adottato
per il file descriptor associato al segnale, in cui si esce dal programma in
caso di errore reale, in questo caso però alla fine dei dati prima di uscire
si stampa anche (\texttt{\small 32}) un messaggio di chiusura.

Se invece vi sono dati validi letti dalla \textit{fifo} si inserirà
(\texttt{\small 36}) una terminazione di stringa sul buffer e si stamperà il
tutto (\texttt{\small 37-38}) sullo \textit{standard output}. L'ultimo
condizionale (\texttt{\small 40-44}) è semplicemente una condizione di cattura
per una eventualità che comunque non dovrebbe mai verificarsi, e che porta
alla uscita dal programma con una opportuna segnalazione di errore.

A questo punto si potrà eseguire il comando lanciandolo su un terminale, ed
osservarne le reazioni agli eventi generati da un altro terminale; lanciando
il programma otterremo qualcosa del tipo:
\begin{Console}
piccardi@hain:~/gapil/sources$ \textbf{./a.out} 
FifoReporter starting, pid 4568
\end{Console}
%$
e scrivendo qualcosa sull'altro terminale con:
\begin{Console}
root@hain:~# \textbf{echo prova > /tmp/reporter.fifo}  
\end{Console}
si otterrà:
\begin{Console}
Message from fifo:
prova
end message
\end{Console}
mentre inviando un segnale:
\begin{Console}
root@hain:~# \textbf{kill 4568}
\end{Console}
si avrà:
\begin{Console}
Signal received:
Got SIGTERM       
From pid 3361
\end{Console}
ed infine premendo \texttt{C-\bslash} sul terminale in cui è in esecuzione si
vedrà:
\begin{Console}
^\\Signal received:
Got SIGQUIT       
From pid 0
\end{Console}
e si potrà far uscire il programma con \texttt{C-c} ottenendo:
\begin{Console}
^CSignal received:
Got SIGINT        
From pid 0
SIGINT means exit
\end{Console}

Lo stesso paradigma di notifica tramite file descriptor usato per i segnali è
stato adottato anche per i timer. In questo caso, rispetto a quanto visto in
sez.~\ref{sec:sig_timer_adv}, la scadenza di un timer potrà essere letta da un
file descriptor senza dover ricorrere ad altri meccanismi di notifica come un
segnale o un \textit{thread}. Di nuovo questo ha il vantaggio di poter
utilizzare le funzioni dell'\textit{I/O multiplexing} per attendere allo
stesso tempo la disponibilità di dati o la ricezione della scadenza di un
timer. In realtà per questo sarebbe già sufficiente \func{signalfd} per
ricevere i segnali associati ai timer, ma la nuova interfaccia semplifica
notevolmente la gestione e consente di fare tutto con una sola \textit{system
  call}.

Le funzioni di questa nuova interfaccia ricalcano da vicino la struttura delle
analoghe versioni ordinarie introdotte con lo standard POSIX.1-2001, che
abbiamo già illustrato in sez.~\ref{sec:sig_timer_adv}.\footnote{questa
  interfaccia è stata introdotta in forma considerata difettosa con il kernel
  2.6.22, per cui è stata immediatamente tolta nel successivo 2.6.23 e
  reintrodotta in una forma considerata adeguata nel kernel 2.6.25, il
  supporto nella \acr{glibc} è stato introdotto a partire dalla versione
  2.8.6, la versione del kernel 2.6.22, presente solo su questo kernel, non è
  supportata e non deve essere usata.} La prima funzione di sistema prevista,
quella che consente di creare un timer, è \funcd{timerfd\_create}, il cui
prototipo è:

\begin{funcproto}{
\fhead{sys/timerfd.h}
\fdecl{int timerfd\_create(int clockid, int flags)}

\fdesc{Crea un timer associato ad un file descriptor di notifica.}
}

{La funzione ritorna un numero di file descriptor in caso di successo e $-1$
  per un errore, nel qual caso \var{errno} assumerà uno dei valori:
  \begin{errlist}
  \item[\errcode{EINVAL}] l'argomento \param{clockid} non è
    \const{CLOCK\_MONOTONIC} o \const{CLOCK\_REALTIME}, o
    l'argomento \param{flag} non è valido, o è diverso da zero per kernel
    precedenti il 2.6.27.
  \item[\errcode{ENODEV}] il kernel non può montare internamente il
    dispositivo per la gestione anonima degli \textit{inode} associati al file
    descriptor.
  \item[\errcode{ENOMEM}] non c'è memoria sufficiente per creare un nuovo file
    descriptor di \func{signalfd}.
  \end{errlist}
  ed inoltre \errval{EMFILE} e \errval{ENFILE} nel loro significato generico.
}  
\end{funcproto}

La funzione prende come primo argomento un intero che indica il tipo di
orologio a cui il timer deve fare riferimento, i valori sono gli stessi delle
funzioni dello standard POSIX-1.2001 già illustrati in
tab.~\ref{tab:sig_timer_clockid_types}, ma al momento i soli utilizzabili sono
\const{CLOCK\_REALTIME} e \const{CLOCK\_MONOTONIC}. L'argomento \param{flags},
come l'analogo di \func{signalfd}, consente di impostare i flag per l'I/O non
bloccante ed il \textit{close-on-exec} sul file descriptor
restituito,\footnote{il flag è stato introdotto a partire dal kernel 2.6.27,
  per le versioni precedenti deve essere passato un valore nullo.} e deve
essere specificato come una maschera binaria delle costanti riportate in
tab.~\ref{tab:timerfd_flags}.

\begin{table}[htb]
  \centering
  \footnotesize
  \begin{tabular}[c]{|l|p{8cm}|}
    \hline
    \textbf{Valore}  & \textbf{Significato} \\
    \hline
    \hline
    \constd{TFD\_NONBLOCK}& imposta sul file descriptor il flag di
                            \const{O\_NONBLOCK} per renderlo non bloccante.\\ 
    \constd{TFD\_CLOEXEC} & imposta il flag di \const{O\_CLOEXEC} per la
                            chiusura automatica del file descriptor nella
                            esecuzione di \func{exec}.\\
    \hline    
  \end{tabular}
  \caption{Valori dell'argomento \param{flags} per la funzione
    \func{timerfd\_create} che consentono di impostare i flag del file
    descriptor.}  
  \label{tab:timerfd_flags}
\end{table}

In caso di successo la funzione restituisce un file descriptor sul quale
verranno notificate le scadenze dei timer. Come per quelli restituiti da
\func{signalfd} anche questo file descriptor segue la semantica dei sistemi
unix-like, in particolare resta aperto attraverso una \func{exec} (a meno che
non si sia impostato il flag di \textit{close-on exec} con
\const{TFD\_CLOEXEC}) e viene duplicato attraverso una \func{fork}; questa
ultima caratteristica comporta però che anche il figlio può utilizzare i dati
di un timer creato nel padre, a differenza di quanto avviene invece con i
timer impostati con le funzioni ordinarie. Si ricordi infatti che, come
illustrato in sez.~\ref{sec:proc_fork}, allarmi, timer e segnali pendenti nel
padre vengono cancellati per il figlio dopo una \func{fork}.

Una volta creato il timer con \func{timerfd\_create} per poterlo utilizzare
occorre \textsl{armarlo} impostandone un tempo di scadenza ed una eventuale
periodicità di ripetizione, per farlo si usa una funzione di sistema omologa
di \func{timer\_settime} per la nuova interfaccia; questa è
\funcd{timerfd\_settime} ed il suo prototipo è:

\begin{funcproto}{
\fhead{sys/timerfd.h}
\fdecl{int timerfd\_settime(int fd, int flags,
                           const struct itimerspec *new\_value,\\
\phantom{int timerfd\_settime(}struct itimerspec *old\_value)}

\fdesc{Arma un timer associato ad un file descriptor di notifica.}
}

{La funzione ritorna un numero di file descriptor in caso di successo e $-1$
  per un errore, nel qual caso \var{errno} assumerà uno dei valori:
  \begin{errlist}
  \item[\errcode{EBADF}] l'argomento \param{fd} non corrisponde ad un file
    descriptor. 
  \item[\errcode{EFAULT}] o \param{new\_value} o \param{old\_value} non sono
    puntatori validi.
  \item[\errcode{EINVAL}] il file descriptor \param{fd} non è stato ottenuto
    con \func{timerfd\_create}, o i valori di \param{flag} o dei campi
    \var{tv\_nsec} in \param{new\_value} non sono validi.
  \end{errlist}
}  
\end{funcproto}

In questo caso occorre indicare su quale timer si intende operare specificando
come primo argomento il file descriptor ad esso associato, che deve essere
stato ottenuto da una precedente chiamata a \func{timerfd\_create}. I restanti
argomenti sono del tutto analoghi a quelli della omologa funzione
\func{timer\_settime}, e prevedono l'uso di strutture \struct{itimerspec}
(vedi fig.~\ref{fig:struct_itimerspec}) per le indicazioni di temporizzazione.

I valori ed il significato di questi argomenti sono gli stessi che sono già
stati illustrati in dettaglio in sez.~\ref{sec:sig_timer_adv} e non staremo a
ripetere quanto detto in quell'occasione; per brevità si ricordi che
con \param{new\_value.it\_value} si indica la prima scadenza del timer e
con \param{new\_value.it\_interval} la sua periodicità.  L'unica differenza
riguarda l'argomento \param{flags} che serve sempre ad indicare se il tempo di
scadenza del timer è da considerarsi relativo o assoluto rispetto al valore
corrente dell'orologio associato al timer, ma che in questo caso ha come
valori possibili rispettivamente soltanto $0$ e \constd{TFD\_TIMER\_ABSTIME}
(l'analogo di \const{TIMER\_ABSTIME}).

L'ultima funzione di sistema prevista dalla nuova interfaccia è
\funcd{timerfd\_gettime}, che è l'analoga di \func{timer\_gettime}, il suo
prototipo è:

\begin{funcproto}{
\fhead{sys/timerfd.h}
\fdecl{int timerfd\_gettime(int fd, struct itimerspec *curr\_value)}

\fdesc{Legge l'impostazione di un timer associato ad un file descriptor di
  notifica.} 
}

{La funzione ritorna un numero di file descriptor in caso di successo e $-1$
  per un errore, nel qual caso \var{errno} assumerà uno dei valori:
  \begin{errlist}
  \item[\errcode{EBADF}] l'argomento \param{fd} non corrisponde ad un file
    descriptor. 
  \item[\errcode{EINVAL}] il file descriptor \param{fd} non è stato ottenuto
    con \func{timerfd\_create}.
  \item[\errcode{EFAULT}] o \param{curr\_value} non è un puntatore valido.
  \end{errlist}
}  
\end{funcproto}

La funzione consente di rileggere le impostazioni del timer associato al file
descriptor \param{fd} nella struttura \struct{itimerspec} puntata
da \param{curr\_value}. Il campo \var{it\_value} riporta il tempo rimanente
alla prossima scadenza del timer, che viene sempre espresso in forma relativa,
anche se lo si è armato specificando \const{TFD\_TIMER\_ABSTIME}. Un valore
nullo (di entrambi i campi di \var{it\_value}) indica invece che il timer non
è stato ancora armato. Il campo \var{it\_interval} riporta la durata
dell'intervallo di ripetizione del timer, ed un valore nullo (di entrambi i
campi) indica che il timer è stato impostato per scadere una sola volta.

Il timer creato con \func{timerfd\_create} notificherà la sua scadenza
rendendo pronto per la lettura il file descriptor ad esso associato, che
pertanto potrà essere messo sotto controllo con una qualunque delle varie
funzioni dell'I/O multiplexing viste in precedenza. Una volta che il file
descriptor risulta pronto sarà possibile leggere il numero di volte che il
timer è scaduto con una ordinaria \func{read}. 

La funzione legge il valore in un dato di tipo \typed{uint64\_t}, e necessita
pertanto che le si passi un buffer di almeno 8 byte, fallendo con
\errval{EINVAL} in caso contrario, in sostanza la lettura deve essere
effettuata con una istruzione del tipo:
\includecodesnip{listati/readtimerfd.c} 

Il valore viene restituito da \func{read} seguendo l'ordinamento dei bit
(\textit{big-endian} o \textit{little-endian}) nativo della macchina in uso,
ed indica il numero di volte che il timer è scaduto dall'ultima lettura
eseguita con successo, o, se lo si legge per la prima volta, da quando lo si è
impostato con \func{timerfd\_settime}. Se il timer non è scaduto la funzione
si blocca fino alla prima scadenza, a meno di non aver creato il file
descriptor in modalità non bloccante con \const{TFD\_NONBLOCK} o aver
impostato la stessa con \func{fcntl}, nel qual caso fallisce con l'errore di
\errval{EAGAIN}.


% TODO trattare qui eventfd introdotto con il 2.6.22 


\section{L'accesso \textsl{asincrono} ai file}
\label{sec:file_asyncronous_operation}

Benché l'\textit{I/O multiplexing} sia stata la prima, e sia tutt'ora una fra
le più diffuse modalità di gestire l'I/O in situazioni complesse in cui si
debba operare su più file contemporaneamente, esistono altre modalità di
gestione delle stesse problematiche. In particolare sono importanti in questo
contesto le modalità di accesso ai file eseguibili in maniera
\textsl{asincrona}, quelle cioè in cui un processo non deve bloccarsi in
attesa della disponibilità dell'accesso al file, ma può proseguire
nell'esecuzione utilizzando invece un meccanismo di notifica asincrono (di
norma un segnale, ma esistono anche altre interfacce, come \textit{inotify}),
per essere avvisato della possibilità di eseguire le operazioni di I/O volute.


\subsection{Il \textit{Signal driven I/O}}
\label{sec:signal_driven_io}

\itindbeg{signal~driven~I/O}

Abbiamo accennato in sez.~\ref{sec:file_open_close} che è definito un flag
\const{O\_ASYNC}, che consentirebbe di aprire un file in modalità asincrona,
anche se in realtà è opportuno attivare in un secondo tempo questa modalità
impostando questo flag attraverso l'uso di \func{fcntl} con il comando
\const{F\_SETFL} (vedi sez.~\ref{sec:file_fcntl_ioctl}).\footnote{l'uso del
  flag di \const{O\_ASYNC} e dei comandi \const{F\_SETOWN} e \const{F\_GETOWN}
  per \func{fcntl} è specifico di Linux e BSD.}  In realtà parlare di apertura
in modalità asincrona non significa che le operazioni di lettura o scrittura
del file vengono eseguite in modo asincrono (tratteremo questo, che è ciò che
più propriamente viene chiamato \textsl{I/O asincrono}, in
sez.~\ref{sec:file_asyncronous_io}), quanto dell'attivazione un meccanismo di
notifica asincrona delle variazione dello stato del file descriptor aperto in
questo modo.

Quello che succede è che per tutti i file posti in questa modalità il sistema
genera un apposito segnale, \signal{SIGIO}, tutte le volte che diventa
possibile leggere o scrivere dal file descriptor; si tenga presente però che
essa non è utilizzabile con i file ordinari ma solo con socket, file di
terminale o pseudo terminale, ed anche, a partire dal kernel 2.6, per
\textit{fifo} e \textit{pipe}. Inoltre è possibile, come illustrato in
sez.~\ref{sec:file_fcntl_ioctl}, selezionare con il comando \const{F\_SETOWN}
di \func{fcntl} quale processo o quale gruppo di processi dovrà ricevere il
segnale. In questo modo diventa possibile effettuare le operazioni di I/O in
risposta alla ricezione del segnale, e non ci sarà più la necessità di restare
bloccati in attesa della disponibilità di accesso ai file.

% TODO: per i thread l'uso di F_SETOWN ha un significato diverso

Per questo motivo Stevens, ed anche le pagine di manuale di Linux, chiamano
questa modalità ``\textit{Signal driven I/O}''.  Si tratta di un'altra
modalità di gestione dell'I/O, alternativa all'uso di
\textit{epoll},\footnote{anche se le prestazioni ottenute con questa tecnica
  sono inferiori, il vantaggio è che questa modalità è utilizzabile anche con
  kernel che non supportano \textit{epoll}, come quelli della serie 2.4,
  ottenendo comunque prestazioni superiori a quelle che si hanno con
  \func{poll} e \func{select}.} che consente di evitare l'uso delle funzioni
\func{poll} o \func{select} che, come illustrato in sez.~\ref{sec:file_epoll},
quando vengono usate con un numero molto grande di file descriptor, non hanno
buone prestazioni.

Tuttavia con l'implementazione classica dei segnali questa modalità di I/O
presenta notevoli problemi, dato che non è possibile determinare, quando i
file descriptor sono più di uno, qual è quello responsabile dell'emissione del
segnale. Inoltre dato che i segnali normali non si accodano (si ricordi quanto
illustrato in sez.~\ref{sec:sig_notification}), in presenza di più file
descriptor attivi contemporaneamente, più segnali emessi nello stesso momento
verrebbero notificati una volta sola.

Linux però supporta le estensioni POSIX.1b dei segnali \textit{real-time}, che
vengono accodati e che permettono di riconoscere il file descriptor che li ha
emessi.  In questo caso infatti si può fare ricorso alle informazioni
aggiuntive restituite attraverso la struttura \struct{siginfo\_t}, utilizzando
la forma estesa \var{sa\_sigaction} del gestore installata con il flag
\const{SA\_SIGINFO} (si riveda quanto illustrato in
sez.~\ref{sec:sig_sigaction}).

Per far questo però occorre utilizzare le funzionalità dei segnali
\textit{real-time} (vedi sez.~\ref{sec:sig_real_time}) impostando
esplicitamente con il comando \const{F\_SETSIG} di \func{fcntl} un segnale
\textit{real-time} da inviare in caso di I/O asincrono (il segnale predefinito
è \signal{SIGIO}). In questo caso il gestore, tutte le volte che riceverà
\const{SI\_SIGIO} come valore del campo \var{si\_code} di \struct{siginfo\_t},
troverà nel campo \var{si\_fd} il valore del file descriptor che ha generato
il segnale. Si noti che il valore di\var{si\_code} resta \const{SI\_SIGIO}
qualunque sia il segnale che si è associato all'I/O, in quanto indica che il
segnale è stato generato a causa di attività di I/O.

Un secondo vantaggio dell'uso dei segnali \textit{real-time} è che essendo
questi ultimi dotati di una coda di consegna ogni segnale sarà associato ad
uno solo file descriptor; inoltre sarà possibile stabilire delle priorità
nella risposta a seconda del segnale usato, dato che i segnali
\textit{real-time} supportano anche questa funzionalità. In questo modo si può
identificare immediatamente un file su cui l'accesso è diventato possibile
evitando completamente l'uso di funzioni come \func{poll} e \func{select},
almeno fintanto che non si satura la coda.

Se infatti si eccedono le dimensioni di quest'ultima, il kernel, non potendo
più assicurare il comportamento corretto per un segnale \textit{real-time},
invierà al suo posto un solo \signal{SIGIO}, su cui si saranno accumulati
tutti i segnali in eccesso, e si dovrà allora determinare con un ciclo quali
sono i file diventati attivi. L'unico modo per essere sicuri che questo non
avvenga è di impostare la lunghezza della coda dei segnali \textit{real-time}
ad una dimensione identica al valore massimo del numero di file descriptor
utilizzabili, vale a dire impostare il contenuto di
\sysctlfile{kernel/rtsig-max} allo stesso valore del contenuto di
\sysctlfile{fs/file-max}.

% TODO fare esempio che usa O_ASYNC

\itindend{signal~driven~I/O}



\subsection{I meccanismi di notifica asincrona.}
\label{sec:file_asyncronous_lease}

Una delle domande più frequenti nella programmazione in ambiente unix-like è
quella di come fare a sapere quando un file viene modificato. La risposta, o
meglio la non risposta, tanto che questa nelle Unix FAQ \cite{UnixFAQ} viene
anche chiamata una \textit{Frequently Unanswered Question}, è che
nell'architettura classica di Unix questo non è possibile. Al contrario di
altri sistemi operativi infatti un kernel unix-like classico non prevedeva
alcun meccanismo per cui un processo possa essere \textsl{notificato} di
eventuali modifiche avvenute su un file. 

Questo è il motivo per cui i demoni devono essere \textsl{avvisati} in qualche
modo se il loro file di configurazione è stato modificato, perché possano
rileggerlo e riconoscere le modifiche; in genere questo vien fatto inviandogli
un segnale di \signal{SIGHUP} che, per una convenzione adottata dalla gran
parte di detti programmi, causa la rilettura della configurazione.

Questa scelta è stata fatta perché provvedere un simile meccanismo a livello
generico per qualunque file comporterebbe un notevole aumento di complessità
dell'architettura della gestione dei file, il tutto per fornire una
funzionalità che serve soltanto in alcuni casi particolari. Dato che
all'origine di Unix i soli programmi che potevano avere una tale esigenza
erano i demoni, attenendosi a uno dei criteri base della progettazione, che
era di far fare al kernel solo le operazioni strettamente necessarie e
lasciare tutto il resto a processi in \textit{user space}, non era stata
prevista nessuna funzionalità di notifica.

Visto però il crescente interesse nei confronti di una funzionalità di questo
tipo, che è molto richiesta specialmente nello sviluppo dei programmi ad
interfaccia grafica quando si deve presentare all'utente lo stato del
filesystem, sono state successivamente introdotte delle estensioni che
permettessero la creazione di meccanismi di notifica più efficienti dell'unica
soluzione disponibile con l'interfaccia tradizionale, che è quella del
\textit{polling}.

Queste nuove funzionalità sono delle estensioni specifiche, non
standardizzate, che sono disponibili soltanto su Linux (anche se altri kernel
supportano meccanismi simili). Alcune di esse sono realizzate, e solo a
partire dalla versione 2.4 del kernel, attraverso l'uso di alcuni
\textsl{comandi} aggiuntivi per la funzione \func{fcntl} (vedi
sez.~\ref{sec:file_fcntl_ioctl}), che divengono disponibili soltanto se si è
definita la macro \macro{\_GNU\_SOURCE} prima di includere \headfile{fcntl.h}.

\itindbeg{file~lease} 

La prima di queste funzionalità è quella del cosiddetto \textit{file lease};
questo è un meccanismo che consente ad un processo, detto \textit{lease
  holder}, di essere notificato quando un altro processo, chiamato a sua volta
\textit{lease breaker}, cerca di eseguire una \func{open} o una
\func{truncate} sul file del quale l'\textit{holder} detiene il
\textit{lease}.  La notifica avviene in maniera analoga a come illustrato in
precedenza per l'uso di \const{O\_ASYNC}: di default viene inviato al
\textit{lease holder} il segnale \signal{SIGIO}, ma questo segnale può essere
modificato usando il comando \const{F\_SETSIG} di \func{fcntl} (anche in
questo caso si può rispecificare lo stesso \signal{SIGIO}).

Se si è fatto questo (ed in genere è opportuno farlo, come in precedenza, per
utilizzare segnali \textit{real-time}) e se inoltre si è installato il gestore
del segnale con \const{SA\_SIGINFO} si riceverà nel campo \var{si\_fd} della
struttura \struct{siginfo\_t} il valore del file descriptor del file sul quale
è stato compiuto l'accesso; in questo modo un processo può mantenere anche più
di un \textit{file lease}.

Esistono due tipi di \textit{file lease}: di lettura (\textit{read lease}) e
di scrittura (\textit{write lease}). Nel primo caso la notifica avviene quando
un altro processo esegue l'apertura del file in scrittura o usa
\func{truncate} per troncarlo. Nel secondo caso la notifica avviene anche se
il file viene aperto in lettura; in quest'ultimo caso però il \textit{lease}
può essere ottenuto solo se nessun altro processo ha aperto lo stesso file.

Come accennato in sez.~\ref{sec:file_fcntl_ioctl} il comando di \func{fcntl}
che consente di acquisire un \textit{file lease} è \const{F\_SETLEASE}, che
viene utilizzato anche per rilasciarlo. In tal caso il file
descriptor \param{fd} passato a \func{fcntl} servirà come riferimento per il
file su cui si vuole operare, mentre per indicare il tipo di operazione
(acquisizione o rilascio) occorrerà specificare come valore
dell'argomento \param{arg} di \func{fcntl} uno dei tre valori di
tab.~\ref{tab:file_lease_fctnl}.

\begin{table}[htb]
  \centering
  \footnotesize
  \begin{tabular}[c]{|l|l|}
    \hline
    \textbf{Valore}  & \textbf{Significato} \\
    \hline
    \hline
    \constd{F\_RDLCK} & Richiede un \textit{read lease}.\\
    \constd{F\_WRLCK} & Richiede un \textit{write lease}.\\
    \constd{F\_UNLCK} & Rilascia un \textit{file lease}.\\
    \hline    
  \end{tabular}
  \caption{Costanti per i tre possibili valori dell'argomento \param{arg} di
    \func{fcntl} quando usata con i comandi \const{F\_SETLEASE} e
    \const{F\_GETLEASE}.} 
  \label{tab:file_lease_fctnl}
\end{table}

Se invece si vuole conoscere lo stato di eventuali \textit{file lease}
occorrerà chiamare \func{fcntl} sul relativo file descriptor \param{fd} con il
comando \const{F\_GETLEASE}, e si otterrà indietro nell'argomento \param{arg}
uno dei valori di tab.~\ref{tab:file_lease_fctnl}, che indicheranno la
presenza del rispettivo tipo di \textit{lease}, o, nel caso di
\const{F\_UNLCK}, l'assenza di qualunque \textit{file lease}.

Si tenga presente che un processo può mantenere solo un tipo di \textit{lease}
su un file, e che un \textit{lease} può essere ottenuto solo su file di dati
(\textit{pipe} e dispositivi sono quindi esclusi). Inoltre un processo non
privilegiato può ottenere un \textit{lease} soltanto per un file appartenente
ad un \ids{UID} corrispondente a quello del processo. Soltanto un processo con
privilegi di amministratore (cioè con la capacità \const{CAP\_LEASE}, vedi
sez.~\ref{sec:proc_capabilities}) può acquisire \textit{lease} su qualunque
file.

Se su un file è presente un \textit{lease} quando il \textit{lease breaker}
esegue una \func{truncate} o una \func{open} che confligge con
esso,\footnote{in realtà \func{truncate} confligge sempre, mentre \func{open},
  se eseguita in sola lettura, non confligge se si tratta di un \textit{read
    lease}.} la funzione si blocca (a meno di non avere aperto il file con
\const{O\_NONBLOCK}, nel qual caso \func{open} fallirebbe con un errore di
\errcode{EWOULDBLOCK}) e viene eseguita la notifica al \textit{lease holder},
così che questo possa completare le sue operazioni sul file e rilasciare il
\textit{lease}.  In sostanza con un \textit{read lease} si rilevano i
tentativi di accedere al file per modificarne i dati da parte di un altro
processo, mentre con un \textit{write lease} si rilevano anche i tentativi di
accesso in lettura.  Si noti comunque che le operazioni di notifica avvengono
solo in fase di apertura del file e non sulle singole operazioni di lettura e
scrittura.

L'utilizzo dei \textit{file lease} consente al \textit{lease holder} di
assicurare la consistenza di un file, a seconda dei due casi, prima che un
altro processo inizi con le sue operazioni di scrittura o di lettura su di
esso. In genere un \textit{lease holder} che riceve una notifica deve
provvedere a completare le necessarie operazioni (ad esempio scaricare
eventuali buffer), per poi rilasciare il \textit{lease} così che il
\textit{lease breaker} possa eseguire le sue operazioni. Questo si fa con il
comando \const{F\_SETLEASE}, o rimuovendo il \textit{lease} con
\const{F\_UNLCK}, o, nel caso di \textit{write lease} che confligge con una
operazione di lettura, declassando il \textit{lease} a lettura con
\const{F\_RDLCK}.

Se il \textit{lease holder} non provvede a rilasciare il \textit{lease} entro
il numero di secondi specificato dal parametro di sistema mantenuto in
\sysctlfiled{fs/lease-break-time} sarà il kernel stesso a rimuoverlo o
declassarlo automaticamente (questa è una misura di sicurezza per evitare che
un processo blocchi indefinitamente l'accesso ad un file acquisendo un
\textit{lease}). Una volta che un \textit{lease} è stato rilasciato o
declassato (che questo sia fatto dal \textit{lease holder} o dal kernel è lo
stesso) le chiamate a \func{open} o \func{truncate} eseguite dal \textit{lease
  breaker} rimaste bloccate proseguono automaticamente.

Benché possa risultare utile per sincronizzare l'accesso ad uno stesso file da
parte di più processi, l'uso dei \textit{file lease} non consente comunque di
risolvere il problema di rilevare automaticamente quando un file o una
directory vengono modificati,\footnote{questa funzionalità venne aggiunta
  principalmente ad uso di Samba per poter facilitare l'emulazione del
  comportamento di Windows sui file, ma ad oggi viene considerata una
  interfaccia mal progettata ed il suo uso è fortemente sconsigliato a favore
  di \textit{inotify}.} che è quanto necessario ad esempio ai programma di
gestione dei file dei vari desktop grafici.

\itindbeg{dnotify}

Per risolvere questo problema a partire dal kernel 2.4 è stata allora creata
un'altra interfaccia,\footnote{si ricordi che anche questa è una interfaccia
  specifica di Linux che deve essere evitata se si vogliono scrivere programmi
  portabili, e che le funzionalità illustrate sono disponibili soltanto se è
  stata definita la macro \macro{\_GNU\_SOURCE}.} chiamata \textit{dnotify},
che consente di richiedere una notifica quando una directory, o uno qualunque
dei file in essa contenuti, viene modificato.  Come per i \textit{file lease}
la notifica avviene di default attraverso il segnale \signal{SIGIO}, ma se ne
può utilizzare un altro, e di nuovo, per le ragioni già esposte in precedenza,
è opportuno che si utilizzino dei segnali \textit{real-time}.  Inoltre, come
in precedenza, si potrà ottenere nel gestore del segnale il file descriptor
che è stato modificato tramite il contenuto della struttura
\struct{siginfo\_t}.

\itindend{file~lease}

\begin{table}[htb]
  \centering
  \footnotesize
  \begin{tabular}[c]{|l|p{8cm}|}
    \hline
    \textbf{Valore}  & \textbf{Significato} \\
    \hline
    \hline
    \constd{DN\_ACCESS} & Un file è stato acceduto, con l'esecuzione di una fra
                          \func{read}, \func{pread}, \func{readv}.\\ 
    \constd{DN\_MODIFY} & Un file è stato modificato, con l'esecuzione di una
                          fra \func{write}, \func{pwrite}, \func{writev}, 
                          \func{truncate}, \func{ftruncate}.\\ 
    \constd{DN\_CREATE} & È stato creato un file nella directory, con
                          l'esecuzione di una fra \func{open}, \func{creat},
                          \func{mknod}, \func{mkdir}, \func{link},
                          \func{symlink}, \func{rename} (da un'altra
                          directory).\\
    \constd{DN\_DELETE} & È stato cancellato un file dalla directory con
                          l'esecuzione di una fra \func{unlink}, \func{rename}
                          (su un'altra directory), \func{rmdir}.\\
    \constd{DN\_RENAME} & È stato rinominato un file all'interno della
                          directory (con \func{rename}).\\
    \constd{DN\_ATTRIB} & È stato modificato un attributo di un file con
                          l'esecuzione di una fra \func{chown}, \func{chmod},
                          \func{utime}.\\ 
    \constd{DN\_MULTISHOT}& Richiede una notifica permanente di tutti gli
                            eventi.\\ 
    \hline    
  \end{tabular}
  \caption{Le costanti che identificano le varie classi di eventi per i quali
    si richiede la notifica con il comando \const{F\_NOTIFY} di \func{fcntl}.} 
  \label{tab:file_notify}
\end{table}

Ci si può registrare per le notifiche dei cambiamenti al contenuto di una
certa directory eseguendo la funzione \func{fcntl} su un file descriptor
associato alla stessa con il comando \const{F\_NOTIFY}. In questo caso
l'argomento \param{arg} di \func{fcntl} serve ad indicare per quali classi
eventi si vuole ricevere la notifica, e prende come valore una maschera
binaria composta dall'OR aritmetico di una o più delle costanti riportate in
tab.~\ref{tab:file_notify}.

A meno di non impostare in maniera esplicita una notifica permanente usando il
valore \const{DN\_MULTISHOT}, la notifica è singola: viene cioè inviata una
sola volta quando si verifica uno qualunque fra gli eventi per i quali la si è
richiesta. Questo significa che un programma deve registrarsi un'altra volta
se desidera essere notificato di ulteriori cambiamenti. Se si eseguono diverse
chiamate con \const{F\_NOTIFY} e con valori diversi per \param{arg} questi
ultimi si \textsl{accumulano}; cioè eventuali nuovi classi di eventi
specificate in chiamate successive vengono aggiunte a quelle già impostate
nelle precedenti.  Se si vuole rimuovere la notifica si deve invece
specificare un valore nullo.

\itindbeg{inotify}

Il maggiore problema di \textit{dnotify} è quello della scalabilità: si deve
usare un file descriptor per ciascuna directory che si vuole tenere sotto
controllo, il che porta facilmente ad avere un eccesso di file aperti. Inoltre
quando la directory che si controlla è all'interno di un dispositivo
rimovibile, mantenere il relativo file descriptor aperto comporta
l'impossibilità di smontare il dispositivo e di rimuoverlo, il che in genere
complica notevolmente la gestione dell'uso di questi dispositivi.

Un altro problema è che l'interfaccia di \textit{dnotify} consente solo di
tenere sotto controllo il contenuto di una directory; la modifica di un file
viene segnalata, ma poi è necessario verificare di quale file si tratta
(operazione che può essere molto onerosa quando una directory contiene un gran
numero di file).  Infine l'uso dei segnali come interfaccia di notifica
comporta tutti i problemi di gestione visti in sez.~\ref{sec:sig_management} e
sez.~\ref{sec:sig_adv_control}.  Per tutta questa serie di motivi in generale
quella di \textit{dnotify} viene considerata una interfaccia di usabilità
problematica ed il suo uso oggi è fortemente sconsigliato.

\itindend{dnotify}

Per risolvere i problemi appena illustrati è stata introdotta una nuova
interfaccia per l'osservazione delle modifiche a file o directory, chiamata
\textit{inotify}.\footnote{l'interfaccia è disponibile a partire dal kernel
  2.6.13, le relative funzioni sono state introdotte nelle glibc 2.4.}  Anche
questa è una interfaccia specifica di Linux (pertanto non deve essere usata se
si devono scrivere programmi portabili), ed è basata sull'uso di una coda di
notifica degli eventi associata ad un singolo file descriptor, il che permette
di risolvere il principale problema di \textit{dnotify}.  La coda viene creata
attraverso la funzione di sistema \funcd{inotify\_init}, il cui prototipo è:

\begin{funcproto}{
\fhead{sys/inotify.h}
\fdecl{int inotify\_init(void)}
\fdesc{Inizializza una istanza di \textit{inotify}.}
}

{La funzione ritornaun file descriptor in caso di successo, o $-1$ in caso di
  errore, nel qual caso \var{errno} assumerà uno dei valori:
  \begin{errlist}
  \item[\errcode{EMFILE}] si è raggiunto il numero massimo di istanze di
    \textit{inotify} consentite all'utente.
  \item[\errcode{ENFILE}] si è raggiunto il massimo di file descriptor aperti
    nel sistema.
  \item[\errcode{ENOMEM}] non c'è sufficiente memoria nel kernel per creare
    l'istanza.
  \end{errlist}
}
\end{funcproto}

La funzione non prende alcun argomento; inizializza una istanza di
\textit{inotify} e restituisce un file descriptor attraverso il quale verranno
effettuate le operazioni di notifica; si tratta di un file descriptor speciale
che non è associato a nessun file su disco, e che viene utilizzato solo per
notificare gli eventi che sono stati posti in osservazione. Per evitare abusi
delle risorse di sistema è previsto che un utente possa utilizzare un numero
limitato di istanze di \textit{inotify}; il valore di default del limite è di
128, ma questo valore può essere cambiato con \func{sysctl} o usando il file
\sysctlfiled{fs/inotify/max\_user\_instances}.

Dato che questo file descriptor non è associato a nessun file o directory
reale, l'inconveniente di non poter smontare un filesystem i cui file sono
tenuti sotto osservazione viene completamente eliminato; anzi, una delle
capacità dell'interfaccia di \textit{inotify} è proprio quella di notificare
il fatto che il filesystem su cui si trova il file o la directory osservata è
stato smontato.

Inoltre trattandosi di un file descriptor a tutti gli effetti, esso potrà
essere utilizzato come argomento per le funzioni \func{select} e \func{poll} e
con l'interfaccia di \textit{epoll}, ed a partire dal kernel 2.6.25 è stato
introdotto anche il supporto per il \texttt{signal-driven I/O}.  Siccome gli
eventi vengono notificati come dati disponibili in lettura, dette funzioni
ritorneranno tutte le volte che si avrà un evento di notifica.

Così, invece di dover utilizzare i segnali, considerati una pessima scelta dal
punto di vista dell'interfaccia utente, si potrà gestire l'osservazione degli
eventi con una qualunque delle modalità di \textit{I/O multiplexing}
illustrate in sez.~\ref{sec:file_multiplexing}. Qualora si voglia cessare
l'osservazione, sarà sufficiente chiudere il file descriptor e tutte le
risorse allocate saranno automaticamente rilasciate. Infine l'interfaccia di
\textit{inotify} consente di mettere sotto osservazione, oltre che una
directory, anche singoli file.

Una volta creata la coda di notifica si devono definire gli eventi da tenere
sotto osservazione; questo viene fatto attraverso una \textsl{lista di
  osservazione} (o \textit{watch list}) che è associata alla coda. Per gestire
la lista di osservazione l'interfaccia fornisce due funzioni di sistema, la
prima di queste è \funcd{inotify\_add\_watch}, il cui prototipo è:

\begin{funcproto}{
\fhead{sys/inotify.h}
\fdecl{int inotify\_add\_watch(int fd, const char *pathname, uint32\_t mask)}
\fdesc{Aggiunge un evento di osservazione a una lista di osservazione.} 
}

{La funzione ritorna un valore positivo in caso di successo, o $-1$ per un
  errore, nel qual caso \var{errno} assumerà uno dei valori:
  \begin{errlist}
  \item[\errcode{EACCES}] non si ha accesso in lettura al file indicato.
  \item[\errcode{EINVAL}] \param{mask} non contiene eventi legali o \param{fd}
    non è un file descriptor di \textit{inotify}.
  \item[\errcode{ENOSPC}] si è raggiunto il numero massimo di voci di
    osservazione o il kernel non ha potuto allocare una risorsa necessaria.
  \end{errlist}
  ed inoltre \errval{EFAULT}, \errval{ENOMEM} e \errval{EBADF} nel loro
  significato generico.}
\end{funcproto}

La funzione consente di creare un ``\textsl{osservatore}'' (il cosiddetto
``\textit{watch}'') nella lista di osservazione di una coda di notifica, che
deve essere indicata specificando il file descriptor ad essa associato
nell'argomento \param{fd}, che ovviamente dovrà essere un file descriptor
creato con \func{inotify\_init}.  Il file o la directory da porre sotto
osservazione vengono invece indicati per nome, da passare
nell'argomento \param{pathname}.  Infine il terzo argomento, \param{mask},
indica che tipo di eventi devono essere tenuti sotto osservazione e le
modalità della stessa.  L'operazione può essere ripetuta per tutti i file e le
directory che si vogliono tenere sotto osservazione,\footnote{anche in questo
  caso c'è un limite massimo che di default è pari a 8192, ed anche questo
  valore può essere cambiato con \func{sysctl} o usando il file
  \sysctlfiled{fs/inotify/max\_user\_watches}.} e si utilizzerà sempre un solo
file descriptor.

Il tipo di evento che si vuole osservare deve essere specificato
nell'argomento \param{mask} come maschera binaria, combinando i valori delle
costanti riportate in tab.~\ref{tab:inotify_event_watch} che identificano i
singoli bit della maschera ed il relativo significato. In essa si sono marcati
con un ``$\bullet$'' gli eventi che, quando specificati per una directory,
vengono osservati anche su tutti i file che essa contiene.  Nella seconda
parte della tabella si sono poi indicate alcune combinazioni predefinite dei
flag della prima parte.

\begin{table}[htb]
  \centering
  \footnotesize
  \begin{tabular}[c]{|l|c|p{8cm}|}
    \hline
    \textbf{Valore}  & & \textbf{Significato} \\
    \hline
    \hline
    \constd{IN\_ACCESS}        &$\bullet$& C'è stato accesso al file in
                                           lettura.\\  
    \constd{IN\_ATTRIB}        &$\bullet$& Ci sono stati cambiamenti sui dati
                                           dell'\textit{inode}
                                           (o sugli attributi estesi, vedi
                                           sez.~\ref{sec:file_xattr}).\\ 
    \constd{IN\_CLOSE\_WRITE}  &$\bullet$& È stato chiuso un file aperto in
                                           scrittura.\\  
    \constd{IN\_CLOSE\_NOWRITE}&$\bullet$& È stato chiuso un file aperto in
                                           sola lettura.\\
    \constd{IN\_CREATE}        &$\bullet$& È stato creato un file o una
                                           directory in una directory sotto
                                           osservazione.\\  
    \constd{IN\_DELETE}        &$\bullet$& È stato cancellato un file o una
                                           directory in una directory sotto
                                           osservazione.\\ 
    \constd{IN\_DELETE\_SELF}  & --      & È stato cancellato il file (o la
                                          directory) sotto osservazione.\\ 
    \constd{IN\_MODIFY}        &$\bullet$& È stato modificato il file.\\ 
    \constd{IN\_MOVE\_SELF}    &         & È stato rinominato il file (o la
                                           directory) sotto osservazione.\\ 
    \constd{IN\_MOVED\_FROM}   &$\bullet$& Un file è stato spostato fuori dalla
                                           directory sotto osservazione.\\ 
    \constd{IN\_MOVED\_TO}     &$\bullet$& Un file è stato spostato nella
                                           directory sotto osservazione.\\ 
    \constd{IN\_OPEN}          &$\bullet$& Un file è stato aperto.\\ 
    \hline    
    \constd{IN\_CLOSE}         &         & Combinazione di
                                           \const{IN\_CLOSE\_WRITE} e
                                           \const{IN\_CLOSE\_NOWRITE}.\\  
    \constd{IN\_MOVE}          &         & Combinazione di
                                           \const{IN\_MOVED\_FROM} e
                                           \const{IN\_MOVED\_TO}.\\
    \constd{IN\_ALL\_EVENTS}   &         & Combinazione di tutti i flag
                                           possibili.\\
    \hline    
  \end{tabular}
  \caption{Le costanti che identificano i bit della maschera binaria
    dell'argomento \param{mask} di \func{inotify\_add\_watch} che indicano il
    tipo di evento da tenere sotto osservazione.} 
  \label{tab:inotify_event_watch}
\end{table}

Oltre ai flag di tab.~\ref{tab:inotify_event_watch}, che indicano il tipo di
evento da osservare e che vengono utilizzati anche in uscita per indicare il
tipo di evento avvenuto, \func{inotify\_add\_watch} supporta ulteriori
flag,\footnote{i flag \const{IN\_DONT\_FOLLOW}, \const{IN\_MASK\_ADD} e
  \const{IN\_ONLYDIR} sono stati introdotti a partire dalle glibc 2.5, se si
  usa la versione 2.4 è necessario definirli a mano.}  riportati in
tab.~\ref{tab:inotify_add_watch_flag}, che indicano le modalità di
osservazione (da passare sempre nell'argomento \param{mask}) e che al
contrario dei precedenti non vengono mai impostati nei risultati in uscita.

\begin{table}[htb]
  \centering
  \footnotesize
  \begin{tabular}[c]{|l|p{8cm}|}
    \hline
    \textbf{Valore}  & \textbf{Significato} \\
    \hline
    \hline
    \constd{IN\_DONT\_FOLLOW}& Non dereferenzia \param{pathname} se questo è un
                               link simbolico.\\
    \constd{IN\_MASK\_ADD}   & Aggiunge a quelli già impostati i flag indicati
                               nell'argomento \param{mask}, invece di
                               sovrascriverli.\\
    \constd{IN\_ONESHOT}     & Esegue l'osservazione su \param{pathname} per
                               una sola volta, rimuovendolo poi dalla
                               \textit{watch list}.\\ 
    \constd{IN\_ONLYDIR}     & Se \param{pathname} è una directory riporta
                               soltanto gli eventi ad essa relativi e non
                               quelli per i file che contiene.\\ 
    \hline    
  \end{tabular}
  \caption{Le costanti che identificano i bit della maschera binaria
    dell'argomento \param{mask} di \func{inotify\_add\_watch} che indicano le
    modalità di osservazione.} 
  \label{tab:inotify_add_watch_flag}
\end{table}

Se non esiste nessun \textit{watch} per il file o la directory specificata
questo verrà creato per gli eventi specificati dall'argomento \param{mask},
altrimenti la funzione sovrascriverà le impostazioni precedenti, a meno che
non si sia usato il flag \const{IN\_MASK\_ADD}, nel qual caso gli eventi
specificati saranno aggiunti a quelli già presenti.

Come accennato quando si tiene sotto osservazione una directory vengono
restituite le informazioni sia riguardo alla directory stessa che ai file che
essa contiene; questo comportamento può essere disabilitato utilizzando il
flag \const{IN\_ONLYDIR}, che richiede di riportare soltanto gli eventi
relativi alla directory stessa. Si tenga presente inoltre che quando si
osserva una directory vengono riportati solo gli eventi sui file che essa
contiene direttamente, non quelli relativi a file contenuti in eventuali
sottodirectory; se si vogliono osservare anche questi sarà necessario creare
ulteriori \textit{watch} per ciascuna sottodirectory.

Infine usando il flag \const{IN\_ONESHOT} è possibile richiedere una notifica
singola;\footnote{questa funzionalità però è disponibile soltanto a partire dal
  kernel 2.6.16.} una volta verificatosi uno qualunque fra gli eventi
richiesti con \func{inotify\_add\_watch} l'\textsl{osservatore} verrà
automaticamente rimosso dalla lista di osservazione e nessun ulteriore evento
sarà più notificato.

In caso di successo \func{inotify\_add\_watch} ritorna un intero positivo,
detto \textit{watch descriptor}, che identifica univocamente un
\textsl{osservatore} su una coda di notifica; esso viene usato per farvi
riferimento sia riguardo i risultati restituiti da \textit{inotify}, che per
la eventuale rimozione dello stesso. 

La seconda funzione di sistema per la gestione delle code di notifica, che
permette di rimuovere un \textsl{osservatore}, è \funcd{inotify\_rm\_watch},
ed il suo prototipo è:

\begin{funcproto}{
\fhead{sys/inotify.h}
\fdecl{int inotify\_rm\_watch(int fd, uint32\_t wd)}
\fdesc{Rimuove un \textsl{osservatore} da una coda di notifica.} 
}

{La funzione ritorna $0$ in caso di successo e $-1$ per un errore, nel qual
  caso \var{errno} assumerà uno dei valori: 
  \begin{errlist}
  \item[\errcode{EBADF}] non si è specificato in \param{fd} un file descriptor
    valido.
  \item[\errcode{EINVAL}] il valore di \param{wd} non è corretto, o \param{fd}
    non è associato ad una coda di notifica.
  \end{errlist}
}
\end{funcproto}

La funzione rimuove dalla coda di notifica identificata dall'argomento
\param{fd} l'osservatore identificato dal \textit{watch descriptor}
\param{wd}; ovviamente deve essere usato per questo argomento un valore
ritornato da \func{inotify\_add\_watch}, altrimenti si avrà un errore di
\errval{EINVAL}. In caso di successo della rimozione, contemporaneamente alla
cancellazione dell'osservatore, sulla coda di notifica verrà generato un
evento di tipo \const{IN\_IGNORED} (vedi
tab.~\ref{tab:inotify_read_event_flag}). Si tenga presente che se un file
viene cancellato o un filesystem viene smontato i relativi osservatori vengono
rimossi automaticamente e non è necessario utilizzare
\func{inotify\_rm\_watch}.

Come accennato l'interfaccia di \textit{inotify} prevede che gli eventi siano
notificati come dati presenti in lettura sul file descriptor associato alla
coda di notifica. Una applicazione pertanto dovrà leggere i dati da detto file
con una \func{read}, che ritornerà sul buffer i dati presenti nella forma di
una o più strutture di tipo \struct{inotify\_event} (la cui definizione è
riportata in fig.~\ref{fig:inotify_event}). Qualora non siano presenti dati la
\func{read} si bloccherà (a meno di non aver impostato il file descriptor in
modalità non bloccante) fino all'arrivo di almeno un evento.

\begin{figure}[!htb]
  \footnotesize \centering
  \begin{minipage}[c]{0.90\textwidth}
    \includestruct{listati/inotify_event.h}
  \end{minipage} 
  \normalsize 
  \caption{La struttura \structd{inotify\_event} usata dall'interfaccia di
    \textit{inotify} per riportare gli eventi.}
  \label{fig:inotify_event}
\end{figure}

Una ulteriore caratteristica dell'interfaccia di \textit{inotify} è che essa
permette di ottenere con \func{ioctl}, come per i file descriptor associati ai
socket (si veda sez.~\ref{sec:sock_ioctl_IP}), il numero di byte disponibili
in lettura sul file descriptor, utilizzando su di esso l'operazione
\const{FIONREAD}.\footnote{questa è una delle operazioni speciali per i file
  (vedi sez.~\ref{sec:file_fcntl_ioctl}), che è disponibile solo per i socket
  e per i file descriptor creati con \func{inotify\_init}.} Si può così
utilizzare questa operazione, oltre che per predisporre una operazione di
lettura con un buffer di dimensioni adeguate, anche per ottenere rapidamente
il numero di file che sono cambiati.

Una volta effettuata la lettura con \func{read} a ciascun evento sarà
associata una struttura \struct{inotify\_event} contenente i rispettivi dati.
Per identificare a quale file o directory l'evento corrisponde viene
restituito nel campo \var{wd} il \textit{watch descriptor} con cui il relativo
osservatore è stato registrato. Il campo \var{mask} contiene invece una
maschera di bit che identifica il tipo di evento verificatosi; in essa
compariranno sia i bit elencati nella prima parte di
tab.~\ref{tab:inotify_event_watch}, che gli eventuali valori aggiuntivi di
tab.~\ref{tab:inotify_read_event_flag} (questi compaiono solo nel campo
\var{mask} di \struct{inotify\_event}, e non sono utilizzabili in fase di
registrazione dell'osservatore).

\begin{table}[htb]
  \centering
  \footnotesize
  \begin{tabular}[c]{|l|p{10cm}|}
    \hline
    \textbf{Valore}  & \textbf{Significato} \\
    \hline
    \hline
    \constd{IN\_IGNORED}    & L'osservatore è stato rimosso, sia in maniera 
                              esplicita con l'uso di \func{inotify\_rm\_watch}, 
                              che in maniera implicita per la rimozione 
                              dell'oggetto osservato o per lo smontaggio del
                              filesystem su cui questo si trova.\\
    \constd{IN\_ISDIR}      & L'evento avvenuto fa riferimento ad una directory
                              (consente così di distinguere, quando si pone
                              sotto osservazione una directory, fra gli eventi
                              relativi ad essa e quelli relativi ai file che
                              essa contiene).\\
    \constd{IN\_Q\_OVERFLOW}& Si sono eccedute le dimensioni della coda degli
                              eventi (\textit{overflow} della coda); in questo
                              caso il valore di \var{wd} è $-1$.\footnotemark\\
    \constd{IN\_UNMOUNT}    & Il filesystem contenente l'oggetto posto sotto
                              osservazione è stato smontato.\\
    \hline    
  \end{tabular}
  \caption{Le costanti che identificano i bit aggiuntivi usati nella maschera
    binaria del campo \var{mask} di \struct{inotify\_event}.} 
  \label{tab:inotify_read_event_flag}
\end{table}

\footnotetext{la coda di notifica ha una dimensione massima che viene
  controllata dal parametro di sistema
  \sysctlfiled{fs/inotify/max\_queued\_events}, che indica il numero massimo di
  eventi che possono essere mantenuti sulla stessa; quando detto valore viene
  ecceduto gli ulteriori eventi vengono scartati, ma viene comunque generato
  un evento di tipo \const{IN\_Q\_OVERFLOW}.}

Il campo \var{cookie} contiene invece un intero univoco che permette di
identificare eventi correlati (per i quali avrà lo stesso valore), al momento
viene utilizzato soltanto per rilevare lo spostamento di un file, consentendo
così all'applicazione di collegare la corrispondente coppia di eventi
\const{IN\_MOVED\_TO} e \const{IN\_MOVED\_FROM}.

Infine due campi \var{name} e \var{len} sono utilizzati soltanto quando
l'evento è relativo ad un file presente in una directory posta sotto
osservazione, in tal caso essi contengono rispettivamente il nome del file
(come \textit{pathname} relativo alla directory osservata) e la relativa
dimensione in byte. Il campo \var{name} viene sempre restituito come stringa
terminata da NUL, con uno o più zeri di terminazione, a seconda di eventuali
necessità di allineamento del risultato, ed il valore di \var{len} corrisponde
al totale della dimensione di \var{name}, zeri aggiuntivi compresi. La stringa
con il nome del file viene restituita nella lettura subito dopo la struttura
\struct{inotify\_event}; questo significa che le dimensioni di ciascun evento
di \textit{inotify} saranno pari a \code{sizeof(\struct{inotify\_event}) +
  len}.

Vediamo allora un esempio dell'uso dell'interfaccia di \textit{inotify} con un
semplice programma che permette di mettere sotto osservazione uno o più file e
directory. Il programma si chiama \texttt{inotify\_monitor.c} ed il codice
completo è disponibile coi sorgenti allegati alla guida, il corpo principale
del programma, che non contiene la sezione di gestione delle opzioni e le
funzioni di ausilio è riportato in fig.~\ref{fig:inotify_monitor_example}.

\begin{figure}[!htbp]
  \footnotesize \centering
  \begin{minipage}[c]{\codesamplewidth}
    \includecodesample{listati/inotify_monitor.c}
  \end{minipage}
  \normalsize
  \caption{Esempio di codice che usa l'interfaccia di \textit{inotify}.}
  \label{fig:inotify_monitor_example}
\end{figure}

Una volta completata la scansione delle opzioni il corpo del programma inizia
controllando (\texttt{\small 11-15}) che sia rimasto almeno un argomento che
indichi quale file o directory mettere sotto osservazione (e qualora questo
non avvenga esce stampando la pagina di aiuto); dopo di che passa
(\texttt{\small 16-20}) all'inizializzazione di \textit{inotify} ottenendo con
\func{inotify\_init} il relativo file descriptor (o si esce in caso di
errore).

Il passo successivo è aggiungere (\texttt{\small 21-30}) alla coda di
notifica gli opportuni osservatori per ciascuno dei file o directory indicati
all'invocazione del comando; questo viene fatto eseguendo un ciclo
(\texttt{\small 22-29}) fintanto che la variabile \var{i}, inizializzata a
zero (\texttt{\small 21}) all'inizio del ciclo, è minore del numero totale di
argomenti rimasti. All'interno del ciclo si invoca (\texttt{\small 23})
\func{inotify\_add\_watch} per ciascuno degli argomenti, usando la maschera
degli eventi data dalla variabile \var{mask} (il cui valore viene impostato
nella scansione delle opzioni), in caso di errore si esce dal programma
altrimenti si incrementa l'indice (\texttt{\small 29}).

Completa l'inizializzazione di \textit{inotify} inizia il ciclo principale
(\texttt{\small 32-56}) del programma, nel quale si resta in attesa degli
eventi che si intendono osservare. Questo viene fatto eseguendo all'inizio del
ciclo (\texttt{\small 33}) una \func{read} che si bloccherà fintanto che non
si saranno verificati eventi.

Dato che l'interfaccia di \textit{inotify} può riportare anche più eventi in
una sola lettura, si è avuto cura di passare alla \func{read} un buffer di
dimensioni adeguate, inizializzato in (\texttt{\small 7}) ad un valore di
approssimativamente 512 eventi (si ricordi che la quantità di dati restituita
da \textit{inotify} è variabile a causa della diversa lunghezza del nome del
file restituito insieme a \struct{inotify\_event}). In caso di errore di
lettura (\texttt{\small 35-40}) il programma esce con un messaggio di errore
(\texttt{\small 37-39}), a meno che non si tratti di una interruzione della
\textit{system call}, nel qual caso (\texttt{\small 36}) si ripete la lettura.

Se la lettura è andata a buon fine invece si esegue un ciclo (\texttt{\small
  43-52}) per leggere tutti gli eventi restituiti, al solito si inizializza
l'indice \var{i} a zero (\texttt{\small 42}) e si ripetono le operazioni
(\texttt{\small 43}) fintanto che esso non supera il numero di byte restituiti
in lettura. Per ciascun evento all'interno del ciclo si assegna alla variabile
\var{event} (si noti come si sia eseguito un opportuno \textit{casting} del
puntatore) l'indirizzo nel buffer della corrispondente struttura
\struct{inotify\_event} (\texttt{\small 44}), e poi si stampano il numero di
\textit{watch descriptor} (\texttt{\small 45}) ed il file a cui questo fa
riferimento (\texttt{\small 46}), ricavato dagli argomenti passati a riga di
comando sfruttando il fatto che i \textit{watch descriptor} vengono assegnati
in ordine progressivo crescente a partire da 1.

Qualora sia presente il riferimento ad un nome di file associato all'evento lo
si stampa (\texttt{\small 47-49}); si noti come in questo caso si sia
controllato il valore del campo \var{event->len} e non il fatto che
\var{event->name} riporti o meno un puntatore nullo. L'interfaccia infatti,
qualora il nome non sia presente, non tocca il campo \var{event->name}, che
si troverà pertanto a contenere quello che era precedentemente presente nella
rispettiva locazione di memoria, nel caso più comune il puntatore al nome di
un file osservato in precedenza.

Si utilizza poi (\texttt{\small 50}) la funzione \code{printevent}, che
interpreta il valore del campo \var{event->mask}, per stampare il tipo di
eventi accaduti.\footnote{per il relativo codice, che non riportiamo in quanto
  non essenziale alla comprensione dell'esempio, si possono utilizzare
  direttamente i sorgenti allegati alla guida.} Infine (\texttt{\small 51}) si
provvede ad aggiornare l'indice \var{i} per farlo puntare all'evento
successivo.

Se adesso usiamo il programma per mettere sotto osservazione una directory, e
da un altro terminale eseguiamo il comando \texttt{ls} otterremo qualcosa del
tipo di:
\begin{Console}
piccardi@gethen:~/gapil/sources$ \textbf{./inotify_monitor -a /home/piccardi/gapil/}
Watch descriptor 1
Observed event on /home/piccardi/gapil/
IN_OPEN, 
Watch descriptor 1
Observed event on /home/piccardi/gapil/
IN_CLOSE_NOWRITE, 
\end{Console}
%$

I lettori più accorti si saranno resi conto che nel ciclo di lettura degli
eventi appena illustrato non viene trattato il caso particolare in cui la
funzione \func{read} restituisce in \var{nread} un valore nullo. Lo si è fatto
perché con \textit{inotify} il ritorno di una \func{read} con un valore nullo
avviene soltanto, come forma di avviso, quando si sia eseguita la funzione
specificando un buffer di dimensione insufficiente a contenere anche un solo
evento. Nel nostro caso le dimensioni erano senz'altro sufficienti, per cui
tale evenienza non si verificherà mai.

Ci si potrà però chiedere cosa succede se il buffer è sufficiente per un
evento, ma non per tutti gli eventi verificatisi. Come si potrà notare nel
codice illustrato in precedenza non si è presa nessuna precauzione per
verificare che non ci fossero stati troncamenti dei dati. Anche in questo caso
il comportamento scelto è corretto, perché l'interfaccia di \textit{inotify}
garantisce automaticamente, anche quando ne sono presenti in numero maggiore,
di restituire soltanto il numero di eventi che possono rientrare completamente
nelle dimensioni del buffer specificato.\footnote{si avrà cioè, facendo
  riferimento sempre al codice di fig.~\ref{fig:inotify_monitor_example}, che
  \var{read} sarà in genere minore delle dimensioni di \var{buffer} ed uguale
  soltanto qualora gli eventi corrispondano esattamente alle dimensioni di
  quest'ultimo.} Se gli eventi sono di più saranno restituiti solo quelli che
entrano interamente nel buffer e gli altri saranno restituiti alla successiva
chiamata di \func{read}.

Infine un'ultima caratteristica dell'interfaccia di \textit{inotify} è che gli
eventi restituiti nella lettura formano una sequenza ordinata, è cioè
garantito che se si esegue uno spostamento di un file gli eventi vengano
generati nella sequenza corretta. L'interfaccia garantisce anche che se si
verificano più eventi consecutivi identici (vale a dire con gli stessi valori
dei campi \var{wd}, \var{mask}, \var{cookie}, e \var{name}) questi vengono
raggruppati in un solo evento.

\itindend{inotify}

% TODO trattare fanotify, vedi http://lwn.net/Articles/339399/ e 
% http://lwn.net/Articles/343346/ (incluso nel 2.6.36)
% fanotify_mark() ha FAN_MARK_FILESYSTEM dal 4.20
% fanotify() ha FAN_OPEN_EXEC dal 4.21/5.0


\subsection{L'interfaccia POSIX per l'I/O asincrono}
\label{sec:file_asyncronous_io}

Una modalità alternativa all'uso dell'\textit{I/O multiplexing} per gestione
dell'I/O simultaneo su molti file è costituita dal cosiddetto \textsl{I/O
  asincrono} o ``AIO''. Il concetto base dell'\textsl{I/O asincrono} è che le
funzioni di I/O non attendono il completamento delle operazioni prima di
ritornare, così che il processo non viene bloccato.  In questo modo diventa ad
esempio possibile effettuare una richiesta preventiva di dati, in modo da
poter effettuare in contemporanea le operazioni di calcolo e quelle di I/O.

Benché la modalità di apertura asincrona di un file vista in
sez.~\ref{sec:signal_driven_io} possa risultare utile in varie occasioni (in
particolar modo con i socket e gli altri file per i quali le funzioni di I/O
sono \textit{system call} lente), essa è comunque limitata alla notifica della
disponibilità del file descriptor per le operazioni di I/O, e non ad uno
svolgimento asincrono delle medesime.  Lo standard POSIX.1b definisce una
interfaccia apposita per l'I/O asincrono vero e proprio,\footnote{questa è
  stata ulteriormente perfezionata nelle successive versioni POSIX.1-2001 e
  POSIX.1-2008.} che prevede un insieme di funzioni dedicate per la lettura e
la scrittura dei file, completamente separate rispetto a quelle usate
normalmente.

In generale questa interfaccia è completamente astratta e può essere
implementata sia direttamente nel kernel che in \textit{user space} attraverso
l'uso di \textit{thread}. Per le versioni del kernel meno recenti esiste una
implementazione di questa interfaccia fornita completamente dalla \acr{glibc}
a partire dalla versione 2.1, che è realizzata completamente in \textit{user
  space}, ed è accessibile linkando i programmi con la libreria
\file{librt}. A partire dalla versione 2.5.32 è stato introdotto nel kernel
una nuova infrastruttura per l'I/O asincrono, ma ancora il supporto è parziale
ed insufficiente ad implementare tutto l'AIO POSIX.

Lo standard POSIX prevede che tutte le operazioni di I/O asincrono siano
controllate attraverso l'uso di una apposita struttura \struct{aiocb} (il cui
nome sta per \textit{asyncronous I/O control block}), che viene passata come
argomento a tutte le funzioni dell'interfaccia. La sua definizione, come
effettuata in \headfiled{aio.h}, è riportata in
fig.~\ref{fig:file_aiocb}. Nello steso file è definita la macro
\macrod{\_POSIX\_ASYNCHRONOUS\_IO}, che dichiara la disponibilità
dell'interfaccia per l'I/O asincrono.

\begin{figure}[!htb]
  \footnotesize \centering
  \begin{minipage}[c]{0.90\textwidth}
    \includestruct{listati/aiocb.h}
  \end{minipage}
  \normalsize 
  \caption{La struttura \structd{aiocb}, usata per il controllo dell'I/O
    asincrono.}
  \label{fig:file_aiocb}
\end{figure}

Le operazioni di I/O asincrono possono essere effettuate solo su un file già
aperto; il file deve inoltre supportare la funzione \func{lseek}, pertanto
terminali e \textit{pipe} sono esclusi. Non c'è limite al numero di operazioni
contemporanee effettuabili su un singolo file.  Ogni operazione deve
inizializzare opportunamente un \textit{control block}.  Il file descriptor su
cui operare deve essere specificato tramite il campo \var{aio\_fildes}; dato
che più operazioni possono essere eseguita in maniera asincrona, il concetto
di posizione corrente sul file viene a mancare; pertanto si deve sempre
specificare nel campo \var{aio\_offset} la posizione sul file da cui i dati
saranno letti o scritti.  Nel campo \var{aio\_buf} deve essere specificato
l'indirizzo del buffer usato per l'I/O, ed in \var{aio\_nbytes} la lunghezza
del blocco di dati da trasferire.

Il campo \var{aio\_reqprio} permette di impostare la priorità delle operazioni
di I/O, in generale perché ciò sia possibile occorre che la piattaforma
supporti questa caratteristica, questo viene indicato dal fatto che le macro
\macrod{\_POSIX\_PRIORITIZED\_IO}, e \macrod{\_POSIX\_PRIORITY\_SCHEDULING}
sono definite. La priorità viene impostata a partire da quella del processo
chiamante (vedi sez.~\ref{sec:proc_priority}), cui viene sottratto il valore
di questo campo.  Il campo \var{aio\_lio\_opcode} è usato solo dalla funzione
\func{lio\_listio}, che, come vedremo, permette di eseguire con una sola
chiamata una serie di operazioni, usando un vettore di \textit{control
  block}. Tramite questo campo si specifica quale è la natura di ciascuna di
esse.

Infine il campo \var{aio\_sigevent} è una struttura di tipo \struct{sigevent}
(illustrata in fig.~\ref{fig:struct_sigevent}) che serve a specificare il modo
in cui si vuole che venga effettuata la notifica del completamento delle
operazioni richieste; per la trattazione delle modalità di utilizzo della
stessa si veda quanto già visto in proposito in sez.~\ref{sec:sig_timer_adv}.

Le due funzioni base dell'interfaccia per l'I/O asincrono sono
\funcd{aio\_read} ed \funcd{aio\_write}.  Esse permettono di richiedere una
lettura od una scrittura asincrona di dati usando la struttura \struct{aiocb}
appena descritta; i rispettivi prototipi sono:

\begin{funcproto}{
\fhead{aio.h}
\fdecl{int aio\_read(struct aiocb *aiocbp)}
\fdesc{Richiede una lettura asincrona.} 
\fdecl{int aio\_write(struct aiocb *aiocbp)}
\fdesc{Richiede una scrittura asincrona.} 
}

{Le funzioni ritornano $0$ in caso di successo e $-1$ per un errore, nel qual
  caso \var{errno} assumerà uno dei valori: 
  \begin{errlist}
  \item[\errcode{EAGAIN}] la coda delle richieste è momentaneamente piena.
  \item[\errcode{EBADF}] si è specificato un file descriptor sbagliato.
  \item[\errcode{EINVAL}] si è specificato un valore non valido per i campi
    \var{aio\_offset} o \var{aio\_reqprio} di \param{aiocbp}.
  \item[\errcode{ENOSYS}] la funzione non è implementata.
  \end{errlist}
}
\end{funcproto}


Entrambe le funzioni ritornano immediatamente dopo aver messo in coda la
richiesta, o in caso di errore. Non è detto che gli errori \errcode{EBADF} ed
\errcode{EINVAL} siano rilevati immediatamente al momento della chiamata,
potrebbero anche emergere nelle fasi successive delle operazioni. Lettura e
scrittura avvengono alla posizione indicata da \var{aio\_offset}, a meno che
il file non sia stato aperto in \textit{append mode} (vedi
sez.~\ref{sec:file_open_close}), nel qual caso le scritture vengono effettuate
comunque alla fine del file, nell'ordine delle chiamate a \func{aio\_write}.

Si tenga inoltre presente che deallocare la memoria indirizzata da
\param{aiocbp} o modificarne i valori prima della conclusione di una
operazione può dar luogo a risultati impredicibili, perché l'accesso ai vari
campi per eseguire l'operazione può avvenire in un momento qualsiasi dopo la
richiesta. Questo comporta che non si devono usare per \param{aiocbp}
variabili automatiche e che non si deve riutilizzare la stessa struttura per
un'altra operazione fintanto che la precedente non sia stata ultimata. In
generale per ogni operazione si deve utilizzare una diversa struttura
\struct{aiocb}.

Dato che si opera in modalità asincrona, il successo di \func{aio\_read} o
\func{aio\_write} non implica che le operazioni siano state effettivamente
eseguite in maniera corretta; per verificarne l'esito l'interfaccia prevede
altre due funzioni, che permettono di controllare lo stato di esecuzione. La
prima è \funcd{aio\_error}, che serve a determinare un eventuale stato di
errore; il suo prototipo è:

\begin{funcproto}{
\fhead{aio.h}
\fdecl{int aio\_error(const struct aiocb *aiocbp)} 
\fdesc{Determina lo stato di errore di una operazione di I/O asincrono.} 
}

{La funzione ritorna $0$ se le operazioni si sono concluse con successo,
  altrimenti restituisce \errval{EINPROGRESS} se non sono concluse,
  \errcode{ECANCELED} se sono state cancellate o il relativo codice di errore
  se sono fallite.}
\end{funcproto}

Se l'operazione non si è ancora completata viene sempre restituito l'errore di
\errcode{EINPROGRESS}, mentre se è stata cancellata ritorna
\errcode{ECANCELED}. La funzione ritorna zero quando l'operazione si è
conclusa con successo, altrimenti restituisce il codice dell'errore
verificatosi, ed esegue la corrispondente impostazione di \var{errno}. Il
codice può essere sia \errcode{EINVAL} ed \errcode{EBADF}, dovuti ad un valore
errato per \param{aiocbp}, che uno degli errori possibili durante l'esecuzione
dell'operazione di I/O richiesta, nel qual caso saranno restituiti, a seconda
del caso, i codici di errore delle \textit{system call} \func{read},
\func{write}, \func{fsync} e \func{fdatasync}.

Una volta che si sia certi che le operazioni siano state concluse (cioè dopo
che una chiamata ad \func{aio\_error} non ha restituito
\errcode{EINPROGRESS}), si potrà usare la funzione \funcd{aio\_return}, che
permette di verificare il completamento delle operazioni di I/O asincrono; il
suo prototipo è:

\begin{funcproto}{
\fhead{aio.h}
\fdecl{ssize\_t aio\_return(const struct aiocb *aiocbp)}
\fdesc{Ottiene lo stato dei risultati di una operazione di I/O asincrono.} 
}

{La funzione ritorna lo stato di uscita dell'operazione eseguita (il valore
  che avrebbero restituito le equivalenti funzioni eseguite in maniera
  sincrona).}
\end{funcproto}

La funzione recupera il valore dello stato di ritorno delle operazioni di I/O
associate a \param{aiocbp} e deve essere chiamata una sola volta per ciascuna
operazione asincrona, essa infatti fa sì che il sistema rilasci le risorse ad
essa associate. É per questo motivo che occorre chiamare la funzione solo dopo
che l'operazione cui \param{aiocbp} fa riferimento si è completata
verificandolo con \func{aio\_error}, ed usarla una sola volta. Una chiamata
precedente il completamento delle operazioni darebbe risultati indeterminati,
così come chiamarla più di una volta.

La funzione restituisce il valore di ritorno relativo all'operazione eseguita,
così come ricavato dalla sottostante \textit{system call} (il numero di byte
letti, scritti o il valore di ritorno di \func{fsync} o \func{fdatasync}).  É
importante chiamare sempre questa funzione, altrimenti le risorse disponibili
per le operazioni di I/O asincrono non verrebbero liberate, rischiando di
arrivare ad un loro esaurimento.

Oltre alle operazioni di lettura e scrittura l'interfaccia POSIX.1b mette a
disposizione un'altra operazione, quella di sincronizzazione dell'I/O,
compiuta dalla funzione \funcd{aio\_fsync}, che ha lo stesso effetto della
analoga \func{fsync}, ma viene eseguita in maniera asincrona; il suo prototipo
è:

\begin{funcproto}{
\fhead{aio.h}
\fdecl{int aio\_fsync(int op, struct aiocb *aiocbp)} 
\fdesc{Richiede la sincronizzazione dei dati su disco.} 
}

{La funzione ritorna $0$ in caso di successo e $-1$ per un errore, nel qual
  caso \var{errno} assumerà gli stessi valori visti \func{aio\_read} con lo
  stesso significato.
}
\end{funcproto}

La funzione richiede la sincronizzazione dei dati delle operazioni di I/O
relative al file descriptor indicato in \texttt{aiocbp->aio\_fildes},
ritornando immediatamente. Si tenga presente che la funzione mette
semplicemente in coda la richiesta, l'esecuzione effettiva della
sincronizzazione dovrà essere verificata con \func{aio\_error} e
\func{aio\_return} come per le operazioni di lettura e
scrittura. L'argomento \param{op} permette di indicare la modalità di
esecuzione, se si specifica il valore \const{O\_DSYNC} le operazioni saranno
completate con una chiamata a \func{fdatasync}, se si specifica
\const{O\_SYNC} con una chiamata a \func{fsync} (per i dettagli vedi
sez.~\ref{sec:file_sync}).

Il successo della chiamata assicura la richiesta di sincronizzazione dei dati
relativi operazioni di I/O asincrono richieste fino a quel momento, niente è
garantito riguardo la sincronizzazione dei dati relativi ad eventuali
operazioni richieste successivamente. Se si è specificato un meccanismo di
notifica questo sarà innescato una volta che le operazioni di sincronizzazione
dei dati saranno completate (\texttt{aio\_sigevent} è l'unico altro campo
di \param{aiocbp} che viene usato.

In alcuni casi può essere necessario interrompere le operazioni di I/O (in
genere quando viene richiesta un'uscita immediata dal programma), per questo
lo standard POSIX.1b prevede una funzione apposita, \funcd{aio\_cancel}, che
permette di cancellare una operazione richiesta in precedenza; il suo
prototipo è:

\begin{funcproto}{
\fhead{aio.h}
\fdecl{int aio\_cancel(int fd, struct aiocb *aiocbp)}
\fdesc{Richiede la cancellazione delle operazioni di I/O asincrono.} 
}

{La funzione ritorna un intero positivo che indica il risultato
  dell'operazione in caso di successo e $-1$ per un errore, nel qual caso
  \var{errno} assumerà uno dei valori:
  \begin{errlist}
  \item[\errcode{EBADF}] \param{fd} non è un file descriptor valido.
  \item[\errcode{ENOSYS}] la funzione non è implementata.
  \end{errlist}
}
\end{funcproto}

La funzione permette di cancellare una operazione specifica sul file
\param{fd}, idicata con \param{aiocbp}, o tutte le operazioni pendenti,
specificando \val{NULL} come valore di \param{aiocbp}. Quando una operazione
viene cancellata una successiva chiamata ad \func{aio\_error} riporterà
\errcode{ECANCELED} come codice di errore, ed mentre il valore di ritorno per
\func{aio\_return} sarà $-1$, inoltre il meccanismo di notifica non verrà
invocato. Se con \param{aiocbp} si specifica una operazione relativa ad un
file descriptor diverso da \param{fd} il risultato è indeterminato.  In caso
di successo, i possibili valori di ritorno per \func{aio\_cancel} (anch'essi
definiti in \headfile{aio.h}) sono tre:
\begin{basedescript}{\desclabelwidth{3.0cm}}
\item[\constd{AIO\_ALLDONE}] indica che le operazioni di cui si è richiesta la
  cancellazione sono state già completate,
  
\item[\constd{AIO\_CANCELED}] indica che tutte le operazioni richieste sono
  state cancellate,  
  
\item[\constd{AIO\_NOTCANCELED}] indica che alcune delle operazioni erano in
  corso e non sono state cancellate.
\end{basedescript}

Nel caso si abbia \const{AIO\_NOTCANCELED} occorrerà chiamare
\func{aio\_error} per determinare quali sono le operazioni effettivamente
cancellate. Le operazioni che non sono state cancellate proseguiranno il loro
corso normale, compreso quanto richiesto riguardo al meccanismo di notifica
del loro avvenuto completamento.

Benché l'I/O asincrono preveda un meccanismo di notifica, l'interfaccia
fornisce anche una apposita funzione, \funcd{aio\_suspend}, che permette di
sospendere l'esecuzione del processo chiamante fino al completamento di una
specifica operazione; il suo prototipo è:

\begin{funcproto}{
\fhead{aio.h}
\fdecl{int aio\_suspend(const struct aiocb * const list[], int nent, \\
\phantom{int aio\_suspend(}const struct timespec *timeout)}
\fdesc{Attende il completamento di una operazione di I/O asincrono.} 
}

{La funzione ritorna $0$ se una (o più) operazioni sono state completate e
  $-1$ per un errore, nel qual caso \var{errno} assumerà uno dei valori:
  \begin{errlist}
    \item[\errcode{EAGAIN}] nessuna operazione è stata completata entro
      \param{timeout}.
    \item[\errcode{EINTR}] la funzione è stata interrotta da un segnale.
    \item[\errcode{ENOSYS}] la funzione non è implementata.
  \end{errlist}
}
\end{funcproto}
  
La funzione permette di bloccare il processo fintanto che almeno una delle
\param{nent} operazioni specificate nella lista \param{list} è completata, per
un tempo massimo specificato dalla struttura \struct{timespec} puntata
da \param{timout}, o fintanto che non arrivi un segnale (si tenga conto che
questo segnale potrebbe essere anche quello utilizzato come meccanismo di
notifica). La lista deve essere inizializzata con delle strutture
\struct{aiocb} relative ad operazioni effettivamente richieste, ma può
contenere puntatori nulli, che saranno ignorati. In caso si siano specificati
valori non validi l'effetto è indefinito.  
Un valore \val{NULL} per \param{timout} comporta l'assenza di timeout, mentre
se si vuole effettuare un \textit{polling} sulle operazioni occorrerà
specificare un puntatore valido ad una struttura \texttt{timespec} (vedi
fig.~\ref{fig:sys_timespec_struct}) contenente valori nulli, e verificare poi
con \func{aio\_error} quale delle operazioni della lista \param{list} è stata
completata.

Lo standard POSIX.1b infine ha previsto pure una funzione, \funcd{lio\_listio},
che permette di effettuare la richiesta di una intera lista di operazioni di
lettura o scrittura; il suo prototipo è:


\begin{funcproto}{
\fhead{aio.h}
\fdecl{int lio\_listio(int mode, struct aiocb * const list[], int nent, struct
    sigevent *sig)}

\fdesc{Richiede l'esecuzione di una serie di operazioni di I/O.} 
}

{La funzione ritorna $0$ in caso di successo e $-1$ per un errore, nel qual
  caso \var{errno} assumerà uno dei valori: 
  \begin{errlist}
    \item[\errcode{EAGAIN}] nessuna operazione è stata completata entro
      \param{timeout}.
    \item[\errcode{EINTR}] la funzione è stata interrotta da un segnale.
    \item[\errcode{EINVAL}] si è passato un valore di \param{mode} non valido
      o un numero di operazioni \param{nent} maggiore di
      \const{AIO\_LISTIO\_MAX}.
    \item[\errcode{ENOSYS}] la funzione non è implementata.
  \end{errlist}
}
\end{funcproto}

La funzione esegue la richiesta delle \param{nent} operazioni indicate nella
lista \param{list} un vettore di puntatori a strutture \struct{aiocb}
indicanti le operazioni da compiere (che verranno eseguite senza un ordine
particolare). La lista può contenere anche puntatori nulli, che saranno
ignorati (si possono così eliminare facilmente componenti della lista senza
doverla rigenerare).

Ciascuna struttura \struct{aiocb} della lista deve contenere un
\textit{control block} opportunamente inizializzato; in particolare per
ognuna di esse dovrà essere specificato il tipo di operazione con il campo
\var{aio\_lio\_opcode}, che può prendere i valori:
\begin{basedescript}{\desclabelwidth{2.0cm}}
\item[\constd{LIO\_READ}]  si richiede una operazione di lettura.
\item[\constd{LIO\_WRITE}] si richiede una operazione di scrittura.
\item[\constd{LIO\_NOP}] non si effettua nessuna operazione.
\end{basedescript}
dove \const{LIO\_NOP} viene usato quando si ha a che fare con un vettore di
dimensione fissa, per poter specificare solo alcune operazioni, o quando si
sono dovute cancellare delle operazioni e si deve ripetere la richiesta per
quelle non completate. 

L'argomento \param{mode} controlla il comportamento della funzione, se viene
usato il valore \constd{LIO\_WAIT} la funzione si blocca fino al completamento
di tutte le operazioni richieste; se si usa \constd{LIO\_NOWAIT} la funzione
ritorna immediatamente dopo aver messo in coda tutte le richieste. In tal caso
il chiamante può richiedere la notifica del completamento di tutte le
richieste, impostando l'argomento \param{sig} in maniera analoga a come si fa
per il campo \var{aio\_sigevent} di \struct{aiocb}.

% TODO: trattare libaio e le system call del kernel per l'I/O asincrono, vedi
% http://lse.sourceforge.net/io/aio.html,
% http://webfiveoh.com/content/guides/2012/aug/mon-13th/linux-asynchronous-io-and-libaio.html, 
% https://code.google.com/p/kernel/wiki/AIOUserGuide,
% http://bert-hubert.blogspot.de/2012/05/on-linux-asynchronous-file-io.html 
% https://www.fsl.cs.sunysb.edu/~vass/linux-aio.txt

% TODO trattare la poll API basata sull'I/O asicrono, introdotta con il kernel
% 4.18, vedi  https://lwn.net/Articles/743714/,
% https://lwn.net/Articles/742978/, https://lwn.net/Articles/758324/
% http://git.infradead.org/users/hch/vfs.git/commit/d2d9e26c7cb6d95d521153897910080cf56c7fad
% Reverted

% TODO trattare la nuova API per l'I/O asincrono (io_uring), introdotta con il
% kernel 5.1, vedi https://lwn.net/Articles/776703/,
% https://lwn.net/ml/linux-fsdevel/20190112213011.1439-1-axboe@kernel.dk/ 

\section{Altre modalità di I/O avanzato}
\label{sec:file_advanced_io}

Oltre alle precedenti modalità di \textit{I/O multiplexing} e \textsl{I/O
  asincrono}, esistono altre funzioni che implementano delle modalità di
accesso ai file più evolute rispetto alle normali funzioni di lettura e
scrittura che abbiamo esaminato in sez.~\ref{sec:file_unix_interface}. In
questa sezione allora prenderemo in esame le interfacce per l'\textsl{I/O
  mappato in memoria}, per l'\textsl{I/O vettorizzato} e altre funzioni di I/O
avanzato.


\subsection{File mappati in memoria}
\label{sec:file_memory_map}

\itindbeg{memory~mapping}

Una modalità alternativa di I/O, che usa una interfaccia completamente diversa
rispetto a quella classica vista in sez.~\ref{sec:file_unix_interface}, è il
cosiddetto \textit{memory-mapped I/O}, che attraverso il meccanismo della
\textsl{paginazione}  usato dalla memoria virtuale (vedi
sez.~\ref{sec:proc_mem_gen}) permette di \textsl{mappare} il contenuto di un
file in una sezione dello spazio di indirizzi del processo che lo ha allocato.

\begin{figure}[htb]
  \centering
  \includegraphics[width=12cm]{img/mmap_layout}
  \caption{Disposizione della memoria di un processo quando si esegue la
  mappatura in memoria di un file.}
  \label{fig:file_mmap_layout}
\end{figure}

Il meccanismo è illustrato in fig.~\ref{fig:file_mmap_layout}, una sezione del
file viene \textsl{mappata} direttamente nello spazio degli indirizzi del
programma.  Tutte le operazioni di lettura e scrittura su variabili contenute
in questa zona di memoria verranno eseguite leggendo e scrivendo dal contenuto
del file attraverso il sistema della memoria virtuale illustrato in
sez.~\ref{sec:proc_mem_gen} che in maniera analoga a quanto avviene per le
pagine che vengono salvate e rilette nella \textit{swap}, si incaricherà di
sincronizzare il contenuto di quel segmento di memoria con quello del file
mappato su di esso.  Per questo motivo si può parlare tanto di \textsl{file
  mappato in memoria}, quanto di \textsl{memoria mappata su file}.

L'uso del \textit{memory-mapping} comporta una notevole semplificazione delle
operazioni di I/O, in quanto non sarà più necessario utilizzare dei buffer
intermedi su cui appoggiare i dati da traferire, poiché questi potranno essere
acceduti direttamente nella sezione di memoria mappata; inoltre questa
interfaccia è più efficiente delle usuali funzioni di I/O, in quanto permette
di caricare in memoria solo le parti del file che sono effettivamente usate ad
un dato istante.

Infatti, dato che l'accesso è fatto direttamente attraverso la memoria
virtuale, la sezione di memoria mappata su cui si opera sarà a sua volta letta
o scritta sul file una pagina alla volta e solo per le parti effettivamente
usate, il tutto in maniera completamente trasparente al processo; l'accesso
alle pagine non ancora caricate avverrà allo stesso modo con cui vengono
caricate in memoria le pagine che sono state salvate sullo \textit{swap}.

Infine in situazioni in cui la memoria è scarsa, le pagine che mappano un file
vengono salvate automaticamente, così come le pagine dei programmi vengono
scritte sulla \textit{swap}; questo consente di accedere ai file su dimensioni
il cui solo limite è quello dello spazio di indirizzi disponibile, e non della
memoria su cui possono esserne lette delle porzioni.

L'interfaccia POSIX implementata da Linux prevede varie funzioni di sistema
per la gestione del \textit{memory mapped I/O}, la prima di queste, che serve
ad eseguire la mappatura in memoria di un file, è \funcd{mmap}; il suo
prototipo è:

\begin{funcproto}{
%\fhead{unistd.h}
\fhead{sys/mman.h} 
\fdecl{void * mmap(void * start, size\_t length, int prot, int flags, int
    fd, off\_t offset)}
\fdesc{Esegue la mappatura in memoria di una sezione di un file.} 
}

{La funzione ritorna il puntatore alla zona di memoria mappata in caso di
  successo, e \const{MAP\_FAILED} (\texttt{(void *) -1}) per un errore, nel
  qual caso \var{errno} assumerà uno dei valori:
  \begin{errlist}
    \item[\errcode{EACCES}] o \param{fd} non si riferisce ad un file regolare,
      o si è usato \const{MAP\_PRIVATE} ma \param{fd} non è aperto in lettura,
      o si è usato \const{MAP\_SHARED} e impostato \const{PROT\_WRITE} ed
      \param{fd} non è aperto in lettura/scrittura, o si è impostato
      \const{PROT\_WRITE} ed \param{fd} è in \textit{append-only}.
    \item[\errcode{EAGAIN}] il file è bloccato, o si è bloccata troppa memoria
      rispetto a quanto consentito dai limiti di sistema (vedi
      sez.~\ref{sec:sys_resource_limit}).
    \item[\errcode{EBADF}] il file descriptor non è valido, e non si è usato
      \const{MAP\_ANONYMOUS}.
    \item[\errcode{EINVAL}] i valori di \param{start}, \param{length} o
      \param{offset} non sono validi (o troppo grandi o non allineati sulla
      dimensione delle pagine), o \param{lengh} è zero (solo dal 2.6.12)
      o \param{flags} contiene sia \const{MAP\_PRIVATE} che
      \const{MAP\_SHARED} o nessuno dei due.
    \item[\errcode{ENFILE}] si è superato il limite del sistema sul numero di
      file aperti (vedi sez.~\ref{sec:sys_resource_limit}).
    \item[\errcode{ENODEV}] il filesystem di \param{fd} non supporta il memory
      mapping.
    \item[\errcode{ENOMEM}] non c'è memoria o si è superato il limite sul
      numero di mappature possibili.
    \item[\errcode{EOVERFLOW}] su architettura a 32 bit con il supporto per i
      \textit{large file} (che hanno una dimensione a 64 bit) il numero di
      pagine usato per \param{lenght} aggiunto a quello usato
      per \param{offset} eccede i 32 bit (\texttt{unsigned long}).
    \item[\errcode{EPERM}] l'argomento \param{prot} ha richiesto
      \const{PROT\_EXEC}, ma il filesystem di \param{fd} è montato con
      l'opzione \texttt{noexec}.
    \item[\errcode{ETXTBSY}] si è impostato \const{MAP\_DENYWRITE} ma
      \param{fd} è aperto in scrittura.
  \end{errlist}
}
\end{funcproto}

La funzione richiede di mappare in memoria la sezione del file \param{fd} a
partire da \param{offset} per \param{length} byte, preferibilmente
all'indirizzo \param{start}. Il valore \param{start} viene normalmente
considerato come un suggerimento, ma l'uso di un qualunque valore diverso da
\val{NULL}, in cui si rimette completamente al kernel la scelta
dell'indirizzo, viene sconsigliato per ragioni di portabilità. Il valore
di \param{offset} deve essere un multiplo della dimensione di una pagina di
memoria.

\begin{table}[htb]
  \centering
  \footnotesize
  \begin{tabular}[c]{|l|l|}
    \hline
    \textbf{Valore} & \textbf{Significato} \\
    \hline
    \hline
    \constd{PROT\_EXEC}  & Le pagine possono essere eseguite.\\
    \constd{PROT\_READ}  & Le pagine possono essere lette.\\
    \constd{PROT\_WRITE} & Le pagine possono essere scritte.\\
    \constd{PROT\_NONE}  & L'accesso alle pagine è vietato.\\
    \hline    
  \end{tabular}
  \caption{Valori dell'argomento \param{prot} di \func{mmap}, relativi alla
    protezione applicate alle pagine del file mappate in memoria.}
  \label{tab:file_mmap_prot}
\end{table}

Il valore dell'argomento \param{prot} indica la protezione\footnote{come
  accennato in sez.~\ref{sec:proc_memory} in Linux la memoria reale è divisa
  in pagine, ogni processo vede la sua memoria attraverso uno o più segmenti
  lineari di memoria virtuale; per ciascuno di questi segmenti il kernel
  mantiene nella \textit{page table} la mappatura sulle pagine di memoria
  reale, ed le modalità di accesso (lettura, esecuzione, scrittura); una loro
  violazione causa quella una \textit{segment violation}, e la relativa
  emissione del segnale \signal{SIGSEGV}.} da applicare al segmento di memoria
e deve essere specificato come maschera binaria ottenuta dall'OR di uno o più
dei valori riportati in tab.~\ref{tab:file_mmap_prot}; il valore specificato
deve essere compatibile con la modalità di accesso con cui si è aperto il
file.

\begin{table}[!htb]
  \centering
  \footnotesize
  \begin{tabular}[c]{|l|p{11cm}|}
    \hline
    \textbf{Valore} & \textbf{Significato} \\
    \hline
    \hline
    \constd{MAP\_32BIT}    & Esegue la mappatura sui primi 2Gb dello spazio
                             degli indirizzi, viene supportato solo sulle
                             piattaforme \texttt{x86-64} per compatibilità con
                             le applicazioni a 32 bit. Viene ignorato se si è
                             richiesto \const{MAP\_FIXED} (dal kernel 2.4.20).\\
    \constd{MAP\_ANON}     & Sinonimo di \const{MAP\_ANONYMOUS}, deprecato.\\
    \constd{MAP\_ANONYMOUS}& La mappatura non è associata a nessun file. Gli
                             argomenti \param{fd} e \param{offset} sono
                             ignorati. L'uso di questo flag con
                             \const{MAP\_SHARED} è stato implementato in Linux
                             a partire dai kernel della serie 2.4.x.\\
    \constd{MAP\_DENYWRITE}& In Linux viene ignorato per evitare
                             \textit{DoS}
                             (veniva usato per segnalare che tentativi di
                             scrittura sul file dovevano fallire con
                             \errcode{ETXTBSY}).\\ 
    \constd{MAP\_EXECUTABLE}& Ignorato.\\
    \constd{MAP\_FILE}     & Valore di compatibilità, ignorato.\\
    \constd{MAP\_FIXED}    & Non permette di restituire un indirizzo diverso
                             da \param{start}, se questo non può essere usato
                             \func{mmap} fallisce. Se si imposta questo flag il
                             valore di \param{start} deve essere allineato
                             alle dimensioni di una pagina.\\
    \constd{MAP\_GROWSDOWN}& Usato per gli \textit{stack}. 
                             Indica che la mappatura deve essere effettuata 
                             con gli indirizzi crescenti verso il basso.\\
    \constd{MAP\_HUGETLB}  & Esegue la mappatura usando le cosiddette
                             ``\textit{huge pages}'' (dal kernel 2.6.32).\\
    \constd{MAP\_LOCKED}   & Se impostato impedisce lo \textit{swapping} delle
                             pagine mappate (dal kernel 2.5.37).\\
    \constd{MAP\_NONBLOCK} & Esegue un \textit{prefaulting} più limitato che
                             non causa I/O (dal kernel 2.5.46).\\
    \constd{MAP\_NORESERVE}& Si usa con \const{MAP\_PRIVATE}. Non riserva
                             delle pagine di \textit{swap} ad uso del meccanismo
                             del \textit{copy on write} 
                             per mantenere le modifiche fatte alla regione
                             mappata, in questo caso dopo una scrittura, se
                             non c'è più memoria disponibile, si ha
                             l'emissione di un \signal{SIGSEGV}.\\
    \constd{MAP\_POPULATE} & Esegue il \textit{prefaulting} delle pagine di
                             memoria necessarie alla mappatura (dal kernel
                             2.5.46).\\ 
    \constd{MAP\_PRIVATE}  & I cambiamenti sulla memoria mappata non vengono
                             riportati sul file. Ne viene fatta una copia
                             privata cui solo il processo chiamante ha
                             accesso.  Incompatibile con \const{MAP\_SHARED}.\\
    \constd{MAP\_SHARED}   & I cambiamenti sulla memoria mappata vengono
                             riportati sul file e saranno immediatamente
                             visibili agli altri processi che mappano lo stesso
                             file. Incompatibile
                             con \const{MAP\_PRIVATE}.\\ 
    \const{MAP\_STACK}     & Al momento è ignorato, è stato fornito (dal kernel
                             2.6.27) a supporto della implementazione dei
                             \textit{thread} nella \acr{glibc}, per allocare
                             memoria in uno spazio utilizzabile come
                             \textit{stack} per le architetture hardware che
                             richiedono un trattamento speciale di
                             quest'ultimo.\\ 
    \constd{MAP\_UNINITIALIZED}& Specifico per i sistemi embedded ed
                             utilizzabile dal kernel 2.6.33 solo se è stata
                             abilitata in fase di compilazione dello stesso
                             l'opzione
                             \texttt{CONFIG\_MMAP\_ALLOW\_UNINITIALIZED}. Se
                             usato le pagine di memoria usate nella mappatura
                             anonima non vengono cancellate; questo migliora
                             le prestazioni sui sistemi con risorse minime, ma
                             comporta la possibilità di rileggere i dati di
                             altri processi che han chiuso una mappatura, per
                             cui viene usato solo quando (come si suppone sia
                             per i sistemi embedded) si ha il completo
                             controllo dell'uso della memoria da parte degli
                             utenti.\\ 
%    \constd{MAP\_DONTEXPAND}& Non consente una successiva espansione dell'area
%                              mappata con \func{mremap}, proposto ma pare non
%                              implementato.\\
    \hline
  \end{tabular}
  \caption{Valori possibili dell'argomento \param{flag} di \func{mmap}.}
  \label{tab:file_mmap_flag}
\end{table}

% TODO trattare MAP_HUGETLB introdotto con il kernel 2.6.32, e modifiche
% introdotte con il 3.8 per le dimensioni variabili delle huge pages

% TODO trattare  MAP_FIXED_NOREPLACE vedi https://lwn.net/Articles/751651/ e
% https://lwn.net/Articles/741369/ 

% TODO: verificare MAP_SYNC e MAP_SHARED_VALIDATE, vedi
% https://lwn.net/Articles/731706/, https://lwn.net/Articles/758594/ incluse
% con il 4.15


L'argomento \param{flags} specifica infine qual è il tipo di oggetto mappato,
le opzioni relative alle modalità con cui è effettuata la mappatura e alle
modalità con cui le modifiche alla memoria mappata vengono condivise o
mantenute private al processo che le ha effettuate. Deve essere specificato
come maschera binaria ottenuta dall'OR di uno o più dei valori riportati in
tab.~\ref{tab:file_mmap_flag}. Fra questi comunque deve sempre essere
specificato o \const{MAP\_PRIVATE} o \const{MAP\_SHARED} per indicare la
modalità con cui viene effettuata la mappatura.

Esistono infatti due modalità alternative di eseguire la mappatura di un file;
la più comune è \const{MAP\_SHARED} in cui la memoria è condivisa e le
modifiche effettuate su di essa sono visibili a tutti i processi che hanno
mappato lo stesso file. In questo caso le modifiche vengono anche riportate su
disco, anche se questo può non essere immediato a causa della bufferizzazione:
si potrà essere sicuri dell'aggiornamento solo in seguito alla chiamata di
\func{msync} o \func{munmap}, e solo allora le modifiche saranno visibili sul
file con l'I/O convenzionale.

Con \const{MAP\_PRIVATE} invece viene creata una copia privata del file,
questo non viene mai modificato e solo il processo chiamante ha accesso alla
mappatura. Le modifiche eseguite dal processo sulla mappatura vengono
effettuate utilizzando il meccanismo del \textit{copy on write}, mentenute in
memoria e salvate su \textit{swap} in caso di necessità.  Non è specificato se
i cambiamenti sul file originale vengano riportati sulla regione mappata.

Gli altri valori di \func{flag} modificano le caratteristiche della
mappatura. Fra questi il più rilevante è probabilmente \const{MAP\_ANONYMOUS}
che consente di creare segmenti di memoria condivisa fra processi diversi
senza appoggiarsi a nessun file (torneremo sul suo utilizzo in
sez.~\ref{sec:ipc_mmap_anonymous}). In tal caso gli argomenti \param{fd}
e \param{offset} vangono ignorati, anche se alcune implementazioni richiedono
che invece \param{fd} sia $-1$, convenzione che è opportuno seguire se si ha a
cuore la portabilità dei programmi.

Gli effetti dell'accesso ad una zona di memoria mappata su file possono essere
piuttosto complessi, essi si possono comprendere solo tenendo presente che
tutto quanto è comunque basato sul meccanismo della memoria virtuale. Questo
comporta allora una serie di conseguenze. La più ovvia è che se si cerca di
scrivere su una zona mappata in sola lettura si avrà l'emissione di un segnale
di violazione di accesso (\signal{SIGSEGV}), dato che i permessi sul segmento
di memoria relativo non consentono questo tipo di accesso.

È invece assai diversa la questione relativa agli accessi al di fuori della
regione di cui si è richiesta la mappatura. A prima vista infatti si potrebbe
ritenere che anch'essi debbano generare un segnale di violazione di accesso;
questo però non tiene conto del fatto che, essendo basata sul meccanismo della
paginazione, la mappatura in memoria non può che essere eseguita su un
segmento di dimensioni rigorosamente multiple di quelle di una pagina, ed in
generale queste potranno non corrispondere alle dimensioni effettive del file
o della sezione che si vuole mappare.

\begin{figure}[!htb] 
  \centering
  \includegraphics[height=6cm]{img/mmap_boundary}
  \caption{Schema della mappatura in memoria di una sezione di file di
    dimensioni non corrispondenti al bordo di una pagina.}
  \label{fig:file_mmap_boundary}
\end{figure}

Il caso più comune è quello illustrato in fig.~\ref{fig:file_mmap_boundary},
in cui la sezione di file non rientra nei confini di una pagina: in tal caso
il file sarà mappato su un segmento di memoria che si estende fino al
bordo della pagina successiva.  In questo caso è possibile accedere a quella
zona di memoria che eccede le dimensioni specificate da \param{length}, senza
ottenere un \signal{SIGSEGV} poiché essa è presente nello spazio di indirizzi
del processo, anche se non è mappata sul file. Il comportamento del sistema è
quello di restituire un valore nullo per quanto viene letto, e di non
riportare su file quanto viene scritto.

Un caso più complesso è quello che si viene a creare quando le dimensioni del
file mappato sono più corte delle dimensioni della mappatura, oppure quando il
file è stato troncato, dopo che è stato mappato, ad una dimensione inferiore a
quella della mappatura in memoria.  In questa situazione, per la sezione di
pagina parzialmente coperta dal contenuto del file, vale esattamente quanto
visto in precedenza; invece per la parte che eccede, fino alle dimensioni date
da \param{length}, l'accesso non sarà più possibile, ma il segnale emesso non
sarà \signal{SIGSEGV}, ma \signal{SIGBUS}, come illustrato in
fig.~\ref{fig:file_mmap_exceed}.

Non tutti i file possono venire mappati in memoria, dato che, come illustrato
in fig.~\ref{fig:file_mmap_layout}, la mappatura introduce una corrispondenza
biunivoca fra una sezione di un file ed una sezione di memoria. Questo
comporta che ad esempio non è possibile mappare in memoria file descriptor
relativi a \textit{pipe}, socket e \textit{fifo}, per i quali non ha senso
parlare di \textsl{sezione}. Lo stesso vale anche per alcuni file di
dispositivo, che non dispongono della relativa operazione \func{mmap} (si
ricordi quanto esposto in sez.~\ref{sec:file_vfs_work}). Si tenga presente
però che esistono anche casi di dispositivi (un esempio è l'interfaccia al
ponte PCI-VME del chip Universe) che sono utilizzabili solo con questa
interfaccia.

\begin{figure}[htb]
  \centering
  \includegraphics[height=6cm]{img/mmap_exceed}
  \caption{Schema della mappatura in memoria di file di dimensioni inferiori
    alla lunghezza richiesta.}
  \label{fig:file_mmap_exceed}
\end{figure}

Dato che passando attraverso una \func{fork} lo spazio di indirizzi viene
copiato integralmente, i file mappati in memoria verranno ereditati in maniera
trasparente dal processo figlio, mantenendo gli stessi attributi avuti nel
padre; così se si è usato \const{MAP\_SHARED} padre e figlio accederanno allo
stesso file in maniera condivisa, mentre se si è usato \const{MAP\_PRIVATE}
ciascuno di essi manterrà una sua versione privata indipendente. Non c'è
invece nessun passaggio attraverso una \func{exec}, dato che quest'ultima
sostituisce tutto lo spazio degli indirizzi di un processo con quello di un
nuovo programma.

Quando si effettua la mappatura di un file vengono pure modificati i tempi ad
esso associati (di cui si è trattato in sez.~\ref{sec:file_file_times}). Il
valore di \var{st\_atime} può venir cambiato in qualunque istante a partire
dal momento in cui la mappatura è stata effettuata: il primo riferimento ad
una pagina mappata su un file aggiorna questo tempo.  I valori di
\var{st\_ctime} e \var{st\_mtime} possono venir cambiati solo quando si è
consentita la scrittura sul file (cioè per un file mappato con
\const{PROT\_WRITE} e \const{MAP\_SHARED}) e sono aggiornati dopo la scrittura
o in corrispondenza di una eventuale \func{msync}.

Dato per i file mappati in memoria le operazioni di I/O sono gestite
direttamente dalla memoria virtuale, occorre essere consapevoli delle
interazioni che possono esserci con operazioni effettuate con l'interfaccia
dei file ordinaria illustrata in sez.~\ref{sec:file_unix_interface}. Il
problema è che una volta che si è mappato un file, le operazioni di lettura e
scrittura saranno eseguite sulla memoria, e riportate su disco in maniera
autonoma dal sistema della memoria virtuale.

Pertanto se si modifica un file con l'interfaccia ordinaria queste modifiche
potranno essere visibili o meno a seconda del momento in cui la memoria
virtuale trasporterà dal disco in memoria quella sezione del file, perciò è
del tutto imprevedibile il risultato della modifica di un file nei confronti
del contenuto della memoria su cui è mappato.

Per questo è sempre sconsigliabile eseguire scritture su un file attraverso
l'interfaccia ordinaria quando lo si è mappato in memoria, è invece possibile
usare l'interfaccia ordinaria per leggere un file mappato in memoria, purché
si abbia una certa cura; infatti l'interfaccia dell'I/O mappato in memoria
mette a disposizione la funzione \funcd{msync} per sincronizzare il contenuto
della memoria mappata con il file su disco; il suo prototipo è:

\begin{funcproto}{
%\fhead{unistd.h}
\fhead{sys/mman.h}
\fdecl{int msync(const void *start, size\_t length, int flags)}
\fdesc{Sincronizza i contenuti di una sezione di un file mappato in memoria.} 
}

{La funzione ritorna $0$ in caso di successo e $-1$ per un errore, nel qual
  caso \var{errno} assumerà uno dei valori: 
  \begin{errlist}
    \item[\errcode{EBUSY}] si è indicato \const{MS\_INVALIDATE} ma
      nell'intervallo di memoria specificato è presente un \textit{memory lock}.
    \item[\errcode{EFAULT}] l'intervallo indicato, o parte di esso, non
      risulta mappato (prima del kernel 2.4.19).
    \item[\errcode{EINVAL}] o \param{start} non è multiplo di
      \const{PAGE\_SIZE}, o si è specificato un valore non valido per
      \param{flags}.
    \item[\errcode{ENOMEM}] l'intervallo indicato, o parte di esso, non
      risulta mappato (dal kernel 2.4.19).
  \end{errlist}
}
\end{funcproto}

La funzione esegue la sincronizzazione di quanto scritto nella sezione di
memoria indicata da \param{start} e \param{offset}, scrivendo le modifiche sul
file (qualora questo non sia già stato fatto).  Provvede anche ad aggiornare i
relativi tempi di modifica. In questo modo si è sicuri che dopo l'esecuzione
di \func{msync} le funzioni dell'interfaccia ordinaria troveranno un contenuto
del file aggiornato.

\begin{table}[htb]
  \centering
  \footnotesize
  \begin{tabular}[c]{|l|p{11cm}|}
    \hline
    \textbf{Valore} & \textbf{Significato} \\
    \hline
    \hline
    \constd{MS\_SYNC}      & richiede una sincronizzazione e ritorna soltanto
                             quando questa è stata completata.\\
    \constd{MS\_ASYNC}     & richiede una sincronizzazione, ma ritorna subito 
                             non attendendo che questa sia finita.\\
    \constd{MS\_INVALIDATE}& invalida le pagine per tutte le mappature
                             in memoria così da rendere necessaria una
                             rilettura immediata delle stesse.\\
    \hline
  \end{tabular}
  \caption{Valori possibili dell'argomento \param{flag} di \func{msync}.}
  \label{tab:file_mmap_msync}
\end{table}

L'argomento \param{flag} è specificato come maschera binaria composta da un OR
dei valori riportati in tab.~\ref{tab:file_mmap_msync}, di questi però
\const{MS\_ASYNC} e \const{MS\_SYNC} sono incompatibili; con il primo valore
infatti la funzione si limita ad inoltrare la richiesta di sincronizzazione al
meccanismo della memoria virtuale, ritornando subito, mentre con il secondo
attende che la sincronizzazione sia stata effettivamente eseguita. Il terzo
valore fa sì che vengano invalidate, per tutte le mappature dello stesso file,
le pagine di cui si è richiesta la sincronizzazione, così che esse possano
essere immediatamente aggiornate con i nuovi valori.

Una volta che si sono completate le operazioni di I/O si può eliminare la
mappatura della memoria usando la funzione \funcd{munmap}, il suo prototipo è:

\begin{funcproto}{
%\fhead{unistd.h}
\fhead{sys/mman.h}
\fdecl{int munmap(void *start, size\_t length)}
\fdesc{Rilascia la mappatura sulla sezione di memoria specificata.} 
}

{La funzione ritorna $0$ in caso di successo e $-1$ per un errore, nel qual
  caso \var{errno} assumerà uno dei valori: 
  \begin{errlist}
    \item[\errcode{EINVAL}] l'intervallo specificato non ricade in una zona
      precedentemente mappata.
  \end{errlist}
}
\end{funcproto}

La funzione cancella la mappatura per l'intervallo specificato con
\param{start} e \param{length}; ogni successivo accesso a tale regione causerà
un errore di accesso in memoria. L'argomento \param{start} deve essere
allineato alle dimensioni di una pagina, e la mappatura di tutte le pagine
contenute anche parzialmente nell'intervallo indicato, verrà rimossa.
Indicare un intervallo che non contiene mappature non è un errore.  Si tenga
presente inoltre che alla conclusione di un processo ogni pagina mappata verrà
automaticamente rilasciata, mentre la chiusura del file descriptor usato per
il \textit{memory mapping} non ha alcun effetto su di esso.

Lo standard POSIX prevede anche una funzione che permetta di cambiare le
protezioni delle pagine di memoria; lo standard prevede che essa si applichi
solo ai \textit{memory mapping} creati con \func{mmap}, ma nel caso di Linux
la funzione può essere usata con qualunque pagina valida nella memoria
virtuale. Questa funzione di sistema è \funcd{mprotect} ed il suo prototipo è:

\begin{funcproto}{
\fhead{sys/mman.h} 
\fdecl{int mprotect(const void *addr, size\_t len, int prot)}
\fdesc{Modifica le protezioni delle pagine di memoria.} 
}

{La funzione ritorna $0$ in caso di successo e $-1$ per un errore, nel qual
  caso \var{errno} assumerà uno dei valori: 
  \begin{errlist}
    \item[\errcode{EINVAL}] il valore di \param{addr} non è valido o non è un
      multiplo di \const{PAGE\_SIZE}.
    \item[\errcode{EACCES}] l'operazione non è consentita, ad esempio si è
      cercato di marcare con \const{PROT\_WRITE} un segmento di memoria cui si
      ha solo accesso in lettura.
    \item[\errcode{ENOMEM}] non è stato possibile allocare le risorse
      necessarie all'interno del kernel o si è specificato un indirizzo di
      memoria non valido del processo o non corrispondente a pagine mappate
      (negli ultimi due casi prima del kernel 2.4.19 veniva prodotto,
      erroneamente, \errcode{EFAULT}).
  \end{errlist}
}
\end{funcproto}

La funzione prende come argomenti un indirizzo di partenza in \param{addr},
allineato alle dimensioni delle pagine di memoria, ed una dimensione
\param{size}. La nuova protezione deve essere specificata in \param{prot} con
una combinazione dei valori di tab.~\ref{tab:file_mmap_prot}.  La nuova
protezione verrà applicata a tutte le pagine contenute, anche parzialmente,
dall'intervallo fra \param{addr} e \param{addr}+\param{size}-1.

Infine Linux supporta alcune operazioni specifiche non disponibili su altri
kernel unix-like per poter usare le quali occorre però dichiarare
\macro{\_GNU\_SOURCE} prima dell'inclusione di \texttt{sys/mman.h}. La prima
di queste è la possibilità di modificare un precedente \textit{memory
  mapping}, ad esempio per espanderlo o restringerlo.  Questo è realizzato
dalla funzione di sistema \funcd{mremap}, il cui prototipo è:

\begin{funcproto}{
\fhead{sys/mman.h} 
\fdecl{void * mremap(void *old\_address, size\_t old\_size , size\_t
    new\_size, unsigned long flags)}
\fdesc{Restringe o allarga una mappatura in memoria.} 
}

{La funzione ritorna l'indirizzo alla nuova area di memoria in caso di
  successo o il valore \const{MAP\_FAILED} (pari a \texttt{(void *) -1}), nel
  qual caso \var{errno} assumerà uno dei valori:
  \begin{errlist}
    \item[\errcode{EINVAL}] il valore di \param{old\_address} non è un
      puntatore valido.
    \item[\errcode{EFAULT}] ci sono indirizzi non validi nell'intervallo
      specificato da \param{old\_address} e \param{old\_size}, o ci sono altre
      mappature di tipo non corrispondente a quella richiesta.
    \item[\errcode{ENOMEM}] non c'è memoria sufficiente oppure l'area di
      memoria non può essere espansa all'indirizzo virtuale corrente, e non si
      è specificato \const{MREMAP\_MAYMOVE} nei flag.
    \item[\errcode{EAGAIN}] il segmento di memoria scelto è bloccato e non può
      essere rimappato.
  \end{errlist}
}
\end{funcproto}

La funzione richiede come argomenti \param{old\_address} (che deve essere
allineato alle dimensioni di una pagina di memoria) che specifica il
precedente indirizzo del \textit{memory mapping} e \param{old\_size}, che ne
indica la dimensione. Con \param{new\_size} si specifica invece la nuova
dimensione che si vuole ottenere. Infine l'argomento \param{flags} è una
maschera binaria per i flag che controllano il comportamento della funzione.
Il solo valore utilizzato è \constd{MREMAP\_MAYMOVE} che consente di eseguire
l'espansione anche quando non è possibile utilizzare il precedente
indirizzo. Per questo motivo, se si è usato questo flag, la funzione può
restituire un indirizzo della nuova zona di memoria che non è detto coincida
con \param{old\_address}.

La funzione si appoggia al sistema della memoria virtuale per modificare
l'associazione fra gli indirizzi virtuali del processo e le pagine di memoria,
modificando i dati direttamente nella \textit{page table} del processo. Come
per \func{mprotect} la funzione può essere usata in generale, anche per pagine
di memoria non corrispondenti ad un \textit{memory mapping}, e consente così
di implementare la funzione \func{realloc} in maniera molto efficiente.

Una caratteristica comune a tutti i sistemi unix-like è che la mappatura in
memoria di un file viene eseguita in maniera lineare, cioè parti successive di
un file vengono mappate linearmente su indirizzi successivi in memoria.
Esistono però delle applicazioni (in particolare la tecnica è usata dai
database o dai programmi che realizzano macchine virtuali) in cui è utile
poter mappare sezioni diverse di un file su diverse zone di memoria.

Questo è ovviamente sempre possibile eseguendo ripetutamente la funzione
\func{mmap} per ciascuna delle diverse aree del file che si vogliono mappare
in sequenza non lineare (ed in effetti è quello che veniva fatto anche con
Linux prima che fossero introdotte queste estensioni) ma questo approccio ha
delle conseguenze molto pesanti in termini di prestazioni.  Infatti per
ciascuna mappatura in memoria deve essere definita nella \textit{page table}
del processo una nuova area di memoria virtuale, quella che nel gergo del
kernel viene chiamata VMA (\textit{virtual memory area}, che corrisponda alla
mappatura, in modo che questa diventi visibile nello spazio degli indirizzi
come illustrato in fig.~\ref{fig:file_mmap_layout}.

Quando un processo esegue un gran numero di mappature diverse (si può arrivare
anche a centinaia di migliaia) per realizzare a mano una mappatura non-lineare
esso vedrà un accrescimento eccessivo della sua \textit{page table}, e lo
stesso accadrà per tutti gli altri processi che utilizzano questa tecnica. In
situazioni in cui le applicazioni hanno queste esigenze si avranno delle
prestazioni ridotte, dato che il kernel dovrà impiegare molte risorse per
mantenere i dati relativi al \textit{memory mapping}, sia in termini di
memoria interna per i dati delle \textit{page table}, che di CPU per il loro
aggiornamento.

Per questo motivo con il kernel 2.5.46 è stato introdotto, ad opera di Ingo
Molnar, un meccanismo che consente la mappatura non-lineare. Anche questa è
una caratteristica specifica di Linux, non presente in altri sistemi
unix-like.  Diventa così possibile utilizzare una sola mappatura iniziale, e
quindi una sola \textit{virtual memory area} nella \textit{page table} del
processo, e poi rimappare a piacere all'interno di questa i dati del file. Ciò
è possibile grazie ad una nuova \textit{system call},
\funcd{remap\_file\_pages}, il cui prototipo è:

\begin{funcproto}{
\fhead{sys/mman.h} 
\fdecl{int remap\_file\_pages(void *start, size\_t size, int prot,
    ssize\_t pgoff, int flags)}
\fdesc{Rimappa non linearmente un \textit{memory mapping}.} 
}

{La funzione ritorna $0$ in caso di successo e $-1$ per un errore, nel qual
  caso \var{errno} assumerà uno dei valori: 
  \begin{errlist}
    \item[\errcode{EINVAL}] si è usato un valore non valido per uno degli
      argomenti o \param{start} non fa riferimento ad un \textit{memory
        mapping} valido creato con \const{MAP\_SHARED}.
  \end{errlist}
  ed inoltre 
 nel loro significato generico.}
\end{funcproto}

Per poter utilizzare questa funzione occorre anzitutto effettuare
preliminarmente una chiamata a \func{mmap} con \const{MAP\_SHARED} per
definire l'area di memoria che poi sarà rimappata non linearmente. Poi si
chiamerà questa funzione per modificare le corrispondenze fra pagine di
memoria e pagine del file; si tenga presente che \func{remap\_file\_pages}
permette anche di mappare la stessa pagina di un file in più pagine della
regione mappata.

La funzione richiede che si identifichi la sezione del file che si vuole
riposizionare all'interno del \textit{memory mapping} con gli argomenti
\param{pgoff} e \param{size}; l'argomento \param{start} invece deve indicare
un indirizzo all'interno dell'area definita dall'\func{mmap} iniziale, a
partire dal quale la sezione di file indicata verrà rimappata. L'argomento
\param{prot} deve essere sempre nullo, mentre \param{flags} prende gli stessi
valori di \func{mmap} (quelli di tab.~\ref{tab:file_mmap_prot}) ma di tutti i
flag solo \const{MAP\_NONBLOCK} non viene ignorato.

\itindbeg{prefaulting} 

Insieme alla funzione \func{remap\_file\_pages} nel kernel 2.5.46 con sono
stati introdotti anche due nuovi flag per \func{mmap}: \const{MAP\_POPULATE} e
\const{MAP\_NONBLOCK}.  Il primo dei due consente di abilitare il meccanismo
del \textit{prefaulting}. Questo viene di nuovo in aiuto per migliorare le
prestazioni in certe condizioni di utilizzo del \textit{memory mapping}.

Il problema si pone tutte le volte che si vuole mappare in memoria un file di
grosse dimensioni. Il comportamento normale del sistema della memoria virtuale
è quello per cui la regione mappata viene aggiunta alla \textit{page table}
del processo, ma i dati verranno effettivamente utilizzati (si avrà cioè un
\textit{page fault} che li trasferisce dal disco alla memoria) soltanto in
corrispondenza dell'accesso a ciascuna delle pagine interessate dal
\textit{memory mapping}.

Questo vuol dire che il passaggio dei dati dal disco alla memoria avverrà una
pagina alla volta con un gran numero di \textit{page fault}, chiaramente se si
sa in anticipo che il file verrà utilizzato immediatamente, è molto più
efficiente eseguire un \textit{prefaulting} in cui tutte le pagine di memoria
interessate alla mappatura vengono ``\textsl{popolate}'' in una sola volta,
questo comportamento viene abilitato quando si usa con \func{mmap} il flag
\const{MAP\_POPULATE}.

Dato che l'uso di \const{MAP\_POPULATE} comporta dell'I/O su disco che può
rallentare l'esecuzione di \func{mmap} è stato introdotto anche un secondo
flag, \const{MAP\_NONBLOCK}, che esegue un \textit{prefaulting} più limitato
in cui vengono popolate solo le pagine della mappatura che già si trovano
nella cache del kernel.\footnote{questo può essere utile per il linker
  dinamico, in particolare quando viene effettuato il \textit{prelink} delle
  applicazioni.}

\itindend{prefaulting}

Per i vantaggi illustrati all'inizio del paragrafo l'interfaccia del
\textit{memory mapped I/O} viene usata da una grande varietà di programmi,
spesso con esigenze molto diverse fra di loro riguardo le modalità con cui
verranno eseguiti gli accessi ad un file; è ad esempio molto comune per i
database effettuare accessi ai dati in maniera pressoché casuale, mentre un
riproduttore audio o video eseguirà per lo più letture sequenziali.

\itindend{memory~mapping}

Per migliorare le prestazioni a seconda di queste modalità di accesso è
disponibile una apposita funzione, \funcd{madvise},\footnote{tratteremo in
  sez.~\ref{sec:file_fadvise} le funzioni che consentono di ottimizzare
  l'accesso ai file con l'interfaccia classica.} che consente di fornire al
kernel delle indicazioni su come un processo intende accedere ad un segmento
di memoria, anche al di là delle mappature dei file, così che possano essere
adottate le opportune strategie di ottimizzazione. Il suo prototipo è:

\begin{funcproto}{
\fhead{sys/mman.h}
\fdecl{int madvise(void *start, size\_t length, int advice)}
\fdesc{Fornisce indicazioni sull'uso previsto di un segmento di memoria.} 
}

{La funzione ritorna $0$ in caso di successo e $-1$ per un errore, nel qual
  caso \var{errno} assumerà uno dei valori: 
  \begin{errlist}
    \item[\errcode{EBADF}] la mappatura esiste ma non corrisponde ad un file.
    \item[\errcode{EINVAL}] \param{start} non è allineato alla dimensione di
      una pagina, \param{length} ha un valore negativo, o \param{advice} non è
      un valore valido, o si è richiesto il rilascio (con
      \const{MADV\_DONTNEED}) di pagine bloccate o condivise o si è usato
      \const{MADV\_MERGEABLE} o \const{MADV\_UNMERGEABLE} ma il kernel non è
      stato compilato per il relativo supporto.
    \item[\errcode{EIO}] la paginazione richiesta eccederebbe i limiti (vedi
      sez.~\ref{sec:sys_resource_limit}) sulle pagine residenti in memoria del
      processo (solo in caso di \const{MADV\_WILLNEED}).
    \item[\errcode{ENOMEM}] gli indirizzi specificati non sono mappati, o, in
      caso \const{MADV\_WILLNEED}, non c'è sufficiente memoria per soddisfare
      la richiesta.
  \end{errlist}
  ed inoltre \errval{EAGAIN} e \errval{ENOSYS} nel loro significato generico.}
\end{funcproto}

La sezione di memoria sulla quale si intendono fornire le indicazioni deve
essere indicata con l'indirizzo iniziale \param{start} e l'estensione
\param{length}, il valore di \param{start} deve essere allineato,
mentre \param{length} deve essere un numero positivo; la versione di Linux
consente anche un valore nullo per \param{length}, inoltre se una parte
dell'intervallo non è mappato in memoria l'indicazione viene comunque
applicata alle restanti parti, anche se la funzione ritorna un errore di
\errval{ENOMEM}.

L'indicazione viene espressa dall'argomento \param{advice} che deve essere
specificato con uno dei valori riportati in
tab.~\ref{tab:madvise_advice_values}; si tenga presente che i valori indicati
nella seconda parte della tabella sono specifici di Linux e non sono previsti
dallo standard POSIX.1b.  La funzione non ha, tranne il caso di
\const{MADV\_DONTFORK}, nessun effetto sul comportamento di un programma, ma
può influenzarne le prestazioni fornendo al kernel indicazioni sulle esigenze
dello stesso, così che sia possibile scegliere le opportune strategie per la
gestione del \textit{read-ahead} (vedi sez.~\ref{sec:file_fadvise}) e del
caching dei dati.

\begin{table}[!htb]
  \centering
  \footnotesize
  \begin{tabular}[c]{|l|p{10 cm}|}
    \hline
    \textbf{Valore} & \textbf{Significato} \\
    \hline
    \hline
    \constd{MADV\_DONTNEED}& non ci si aspetta nessun accesso nell'immediato
                             futuro, pertanto le pagine possono essere
                             liberate dal kernel non appena necessario; l'area
                             di memoria resterà accessibile, ma un accesso
                             richiederà che i dati vengano ricaricati dal file
                             a cui la mappatura fa riferimento.\\
    \constd{MADV\_NORMAL}  & nessuna indicazione specifica, questo è il valore
                             di default usato quando non si è chiamato
                             \func{madvise}.\\
    \constd{MADV\_RANDOM}  & ci si aspetta un accesso casuale all'area
                             indicata, pertanto l'applicazione di una lettura
                             anticipata con il meccanismo del
                             \textit{read-ahead} (vedi 
                             sez.~\ref{sec:file_fadvise}) è di
                             scarsa utilità e verrà disabilitata.\\
    \constd{MADV\_SEQUENTIAL}& ci si aspetta un accesso sequenziale al file,
                               quindi da una parte sarà opportuno eseguire una
                               lettura anticipata, e dall'altra si potranno
                               scartare immediatamente le pagine una volta che
                               queste siano state lette.\\
    \const{MADV\_WILLNEED}& ci si aspetta un accesso nell'immediato futuro,
                            pertanto l'applicazione del \textit{read-ahead}
                            deve essere incentivata.\\
    \hline
    \constd{MADV\_DONTDUMP}& esclude da un \textit{core dump} (vedi
                             sez.~\ref{sec:sig_standard}) le pagine 
                             specificate, viene usato per evitare di scrivere
                             su disco dati relativi a zone di memoria che si sa
                             non essere utili in un \textit{core dump}.\\
    \constd{MADV\_DODUMP}  & rimuove l'effetto della precedente
                             \const{MADV\_DONTDUMP} (dal kernel 3.4).\\ 
    \constd{MADV\_DONTFORK}& impedisce che l'intervallo specificato venga
                             ereditato dal processo figlio dopo una
                             \func{fork}; questo consente di evitare che il
                             meccanismo del \textit{copy on write} effettui la
                             rilocazione delle pagine quando il padre scrive
                             sull'area di memoria dopo la \func{fork}, cosa che
                             può causare problemi per l'hardware che esegue
                             operazioni in DMA su quelle pagine (dal kernel
                             2.6.16).\\
    \constd{MADV\_DOFORK}  & rimuove l'effetto della precedente
                             \const{MADV\_DONTFORK} (dal kernel 2.6.16).\\ 
    \constd{MADV\_HUGEPAGE}& abilita il meccanismo delle \textit{Transparent
                             Huge Page} (vedi sez.~\ref{sec:huge_pages})
                             sulla regione indicata; se questa è allineata
                             alle relative dimensioni il kernel alloca
                             direttamente delle \textit{huge page}; è
                             utilizzabile solo con mappature anomime private
                             (dal kernel 2.6.38).\\
    \constd{MADV\_NOHUGEPAGE}& impedisce che la regione indicata venga
                               collassata in eventuali \textit{huge page} (dal
                               kernel 2.6.38).\\
    \constd{MADV\_HWPOISON} &opzione ad uso di debug per verificare codice
                              che debba gestire errori nella gestione della
                              memoria; richiede una apposita opzione di
                              compilazione del kernel, privilegi amministrativi
                              (la capacità \const{CAP\_SYS\_ADMIN}) e provoca
                              l'emissione di un segnale di \const{SIGBUS} dal
                              programma chiamante e rimozione della mappatura
                              (dal kernel 2.6.32).\\
    \constd{MADV\_SOFT\_OFFLINE}&opzione utilizzata per il debug del
                              codice di verifica degli errori di gestione
                              memoria, richiede una apposita opzione di
                              compilazione (dal kernel 2.6.33).\\
    \constd{MADV\_MERGEABLE}& marca la pagina come accorpabile, indicazione
                              principalmente ad uso dei sistemi di
                              virtualizzazione\footnotemark (dal kernel
                              2.6.32).\\ 
    \constd{MADV\_REMOVE}  & libera un intervallo di pagine di memoria ed il
                             relativo supporto sottostante; è supportato
                             soltanto sui filesystem in RAM \textit{tmpfs} e
                             \textit{shmfs} se usato su altri tipi di
                             filesystem causa un errore di \errcode{ENOSYS}
                             (dal kernel 2.6.16).\\
    \constd{MADV\_UNMERGEABLE}& rimuove l'effetto della precedente
                                \const{MADV\_MERGEABLE} (dal kernel 2.6.32). \\
    \hline
  \end{tabular}
  \caption{Valori dell'argomento \param{advice} di \func{madvise}.}
  \label{tab:madvise_advice_values}
\end{table}

% TODO aggiunta MADV_FREE dal kernel 4.5 (vedi http://lwn.net/Articles/590991/)
% TODO aggiunta MADV_WIPEONFORK dal kernel 4.14 that causes the affected memory
% region to appear to be full of zeros in the child process after a fork. It
% differs from the existing MADV_DONTFORK in that the address range will
% remain valid in the child (dalla notizia in https://lwn.net/Articles/733256/).  

\footnotetext{a partire dal kernel 2.6.32 è stato introdotto un meccanismo che
  identifica pagine di memoria identiche e le accorpa in una unica pagina
  (soggetta al \textit{copy-on-write} per successive modifiche); per evitare
  di controllare tutte le pagine solo quelle marcate con questo flag vengono
  prese in considerazione per l'accorpamento; in questo modo si possono
  migliorare le prestazioni nella gestione delle macchine virtuali diminuendo
  la loro occupazione di memoria, ma il meccanismo può essere usato anche in
  altre applicazioni in cui sian presenti numerosi processi che usano gli
  stessi dati; per maggiori dettagli si veda
  \href{http://kernelnewbies.org/Linux_2_6_32\#head-d3f32e41df508090810388a57efce73f52660ccb}{\texttt{http://kernelnewbies.org/Linux\_2\_6\_32}}
  e la documentazione nei sorgenti del kernel
  (\texttt{Documentation/vm/ksm.txt}).} 


A differenza da quanto specificato nello standard POSIX.1b, per il quale l'uso
di \func{madvise} è a scopo puramente indicativo, Linux considera queste
richieste come imperative, per cui ritorna un errore qualora non possa
soddisfarle; questo comportamento differisce da quanto specificato nello
standard.

Nello standard POSIX.1-2001 è prevista una ulteriore funzione
\funcd{posix\_madvise} che su Linux viene reimplementata utilizzando
\func{madvise}; il suo prototipo è:

\begin{funcproto}{
\fhead{sys/mman.h}
\fdecl{int posix\_madvise(void *start, size\_t lenght, int advice)}
\fdesc{Fornisce indicazioni sull'uso previsto di un segmento di memoria.} 
}

{La funzione ritorna $0$ in caso di successo ed un valore positivo per un
  errore, nel qual caso \var{errno} assumerà uno dei valori:
  \begin{errlist}
    \item[\errcode{EINVAL}] \param{start} non è allineato alla dimensione di
      una pagina, \param{length} ha un valore negativo, o \param{advice} non è
      un valore valido.
    \item[\errcode{ENOMEM}] gli indirizzi specificati non sono nello spazio di
      indirizzi del processo.
  \end{errlist}
}
\end{funcproto}

Gli argomenti \param{start} e \param{lenght} hanno lo stesso identico
significato degli analoghi di \func{madvise}, a cui si rimanda per la loro
descrizione ma a differenza di quanto indicato dallo standard per questa
funzione, su Linux un valore nullo di \param{len} è consentito.

\begin{table}[!htb]
  \centering
  \footnotesize
  \begin{tabular}[c]{|l|l|}
    \hline
    \textbf{Valore} & \textbf{Significato} \\
    \hline
    \hline
    \constd{POSIX\_MADV\_DONTNEED}& analogo a \const{MADV\_DONTNEED}.\\
    \constd{POSIX\_MADV\_NORMAL}  & identico a \const{MADV\_NORMAL}.\\
    \constd{POSIX\_MADV\_RANDOM}  & identico a \const{MADV\_RANDOM}.\\
    \constd{POSIX\_MADV\_SEQUENTIAL}& identico a \const{MADV\_SEQUENTIAL}.\\
    \constd{POSIX\_MADV\_WILLNEED}& identico a \const{MADV\_WILLNEED}.\\
     \hline
  \end{tabular}
  \caption{Valori dell'argomento \param{advice} di \func{posix\_madvise}.}
  \label{tab:posix_madvise_advice_values}
\end{table}


L'argomento \param{advice} invece può assumere solo i valori indicati in
tab.~\ref{tab:posix_madvise_advice_values}, che riflettono gli analoghi di
\func{madvise}, con lo stesso effetto per tutti tranne
\const{POSIX\_MADV\_DONTNEED}.  Infatti a partire dalla \acr{glibc} 2.6
\const{POSIX\_MADV\_DONTNEED} viene ignorato, in quanto l'uso del
corrispondente \const{MADV\_DONTNEED} di \func{madvise} ha, per la semantica
imperativa, l'effetto immediato di far liberare le pagine da parte del kernel,
che viene considerato distruttivo.



\subsection{I/O vettorizzato: \func{readv} e \func{writev}}
\label{sec:file_multiple_io}

Una seconda modalità di I/O diversa da quella ordinaria è il cosiddetto
\textsl{I/O vettorizzato}, che nasce per rispondere al caso abbastanza comune
in cui ci si trova nell'esigenza di dover eseguire una serie multipla di
operazioni di I/O, come una serie di letture o scritture di vari buffer. Un
esempio tipico è quando i dati sono strutturati nei campi di una struttura ed
essi devono essere caricati o salvati su un file.  Benché l'operazione sia
facilmente eseguibile attraverso una serie multipla di chiamate a \func{read}
e \func{write}, ci sono casi in cui si vuole poter contare sulla atomicità
delle operazioni di lettura e scrittura rispetto all'esecuzione del programma.

Per questo motivo fino da BSD 4.2 vennero introdotte delle nuove
\textit{system call} che permettessero di effettuare con una sola chiamata una
serie di letture da, o scritture su, una serie di buffer, quello che poi venne
chiamato \textsl{I/O vettorizzato}. Queste funzioni di sistema sono
\funcd{readv} e \funcd{writev},\footnote{in Linux le due funzioni sono riprese
  da BSD4.4, esse sono previste anche dallo standard POSIX.1-2001.} ed i
relativi prototipi sono:


\begin{funcproto}{
\fhead{sys/uio.h}
\fdecl{int readv(int fd, const struct iovec *vector, int count)}
\fdecl{int writev(int fd, const struct iovec *vector, int count)}
\fdesc{Eseguono rispettivamente una lettura o una scrittura vettorizzata.} 
}

{Le funzioni ritornano il numero di byte letti o scritti in caso di successo e
  $-1$ per un errore, nel qual caso \var{errno} assumerà uno dei valori:
  \begin{errlist}
    \item[\errcode{EINVAL}] si è specificato un valore non valido per uno degli
    argomenti (ad esempio \param{count} è maggiore di \const{IOV\_MAX}).
  \end{errlist}
  più tutti i valori, con lo stesso significato, che possono risultare
  dalle condizioni di errore di \func{read} e \func{write}.
 }
\end{funcproto}


Entrambe le funzioni usano una struttura \struct{iovec}, la cui definizione è
riportata in fig.~\ref{fig:file_iovec}, che definisce dove i dati devono
essere letti o scritti ed in che quantità. Il primo campo della struttura,
\var{iov\_base}, contiene l'indirizzo del buffer ed il secondo,
\var{iov\_len}, la dimensione dello stesso.

\begin{figure}[!htb]
  \footnotesize \centering
  \begin{minipage}[c]{\textwidth}
    \includestruct{listati/iovec.h}
  \end{minipage} 
  \normalsize 
  \caption{La struttura \structd{iovec}, usata dalle operazioni di I/O
    vettorizzato.} 
  \label{fig:file_iovec}
\end{figure}

La lista dei buffer da utilizzare viene indicata attraverso l'argomento
\param{vector} che è un vettore di strutture \struct{iovec}, la cui lunghezza
è specificata dall'argomento \param{count}.\footnote{fino alle libc5, Linux
  usava \type{size\_t} come tipo dell'argomento \param{count}, una scelta
  logica, che però è stata dismessa per restare aderenti allo standard
  POSIX.1-2001.}  Ciascuna struttura dovrà essere inizializzata opportunamente
per indicare i vari buffer da e verso i quali verrà eseguito il trasferimento
dei dati. Essi verranno letti (o scritti) nell'ordine in cui li si sono
specificati nel vettore \param{vector}.

La standardizzazione delle due funzioni all'interno della revisione
POSIX.1-2001 prevede anche che sia possibile avere un limite al numero di
elementi del vettore \param{vector}. Qualora questo sussista, esso deve essere
indicato dal valore dalla costante \const{IOV\_MAX}, definita come le altre
costanti analoghe (vedi sez.~\ref{sec:sys_limits}) in \headfile{limits.h}; lo
stesso valore deve essere ottenibile in esecuzione tramite la funzione
\func{sysconf} richiedendo l'argomento \const{\_SC\_IOV\_MAX} (vedi
sez.~\ref{sec:sys_limits}).

Nel caso di Linux il limite di sistema è di 1024, però se si usa la
\acr{glibc} essa fornisce un \textit{wrapper} per le \textit{system call}
che si accorge se una operazione supererà il precedente limite, in tal caso i
dati verranno letti o scritti con le usuali \func{read} e \func{write} usando
un buffer di dimensioni sufficienti appositamente allocato in grado di
contenere tutti i dati indicati da \param{vector}. L'operazione avrà successo
ma si perderà l'atomicità del trasferimento da e verso la destinazione finale.

Si tenga presente infine che queste funzioni operano sui file con
l'interfaccia dei file descriptor, e non è consigliabile mescolarle con
l'interfaccia classica dei \textit{file stream} di
sez.~\ref{sec:files_std_interface}; a causa delle bufferizzazioni interne di
quest'ultima infatti si potrebbero avere risultati indefiniti e non
corrispondenti a quanto aspettato.

Come per le normali operazioni di lettura e scrittura, anche per l'\textsl{I/O
  vettorizzato} si pone il problema di poter effettuare le operazioni in
maniera atomica a partire da un certa posizione sul file. Per questo motivo a
partire dal kernel 2.6.30 sono state introdotte anche per l'\textsl{I/O
  vettorizzato} le analoghe delle funzioni \func{pread} e \func{pwrite} (vedi
sez.~\ref{sec:file_read} e \ref{sec:file_write}); le due funzioni sono
\funcd{preadv} e \funcd{pwritev} ed i rispettivi prototipi sono:\footnote{le
  due funzioni sono analoghe alle omonime presenti in BSD; le \textit{system
    call} usate da Linux (introdotte a partire dalla versione 2.6.30)
  utilizzano degli argomenti diversi per problemi collegati al formato a 64
  bit dell'argomento \param{offset}, che varia a seconda delle architetture,
  ma queste differenze vengono gestite dalle funzioni di librerie di libreria
  che mantengono l'interfaccia delle analoghe tratte da BSD.}


\begin{funcproto}{
\fhead{sys/uio.h}
\fdecl{int preadv(int fd, const struct iovec *vector, int count, off\_t
    offset)}
\fdecl{int pwritev(int fd, const struct iovec *vector, int count, off\_t
    offset)}
\fdesc{Eseguono una lettura o una scrittura vettorizzata a partire da una data
  posizione sul file.} 
}

{ Le funzioni hanno gli stessi valori di ritorno delle corrispondenti
  \func{readv} e \func{writev} ed anche gli eventuali errori sono gli stessi,
  con in più quelli che si possono ottenere dalle possibili condizioni di
  errore di \func{lseek}.
}
\end{funcproto}

Le due funzioni eseguono rispettivamente una lettura o una scrittura
vettorizzata a partire dalla posizione \param{offset} sul file indicato
da \param{fd}, la posizione corrente sul file, come vista da eventuali altri
processi che vi facciano riferimento, non viene alterata. A parte la presenza
dell'ulteriore argomento il comportamento delle funzioni è identico alle
precedenti \func{readv} e \func{writev}. 

Con l'uso di queste funzioni si possono evitare eventuali \textit{race
  condition} quando si deve eseguire la una operazione di lettura e scrittura
vettorizzata a partire da una certa posizione su un file, mentre al contempo
si possono avere in concorrenza processi che utilizzano lo stesso file
descriptor (si ricordi quanto visto in sez.~\ref{sec:file_adv_func}) con delle
chiamate a \func{lseek}.

% TODO trattare preadv2() e pwritev2(), introdotte con il kernel 4.6, vedi
% http://lwn.net/Articles/670231/ ed il flag RWF_HIPRI, anche l'aggiunta del
% flag RWF_APPEND a pwritev2 con il kernel 4.16, vedi
% https://lwn.net/Articles/746129/ 


\subsection{L'I/O diretto fra file descriptor: \func{sendfile} e
  \func{splice}} 
\label{sec:file_sendfile_splice}

Uno dei problemi che si presentano nella gestione dell'I/O è quello in cui si
devono trasferire grandi quantità di dati da un file descriptor ed un altro;
questo usualmente comporta la lettura dei dati dal primo file descriptor in un
buffer in memoria, da cui essi vengono poi scritti sul secondo.

Benché il kernel ottimizzi la gestione di questo processo quando si ha a che
fare con file normali, in generale quando i dati da trasferire sono molti si
pone il problema di effettuare trasferimenti di grandi quantità di dati da
\textit{kernel space} a \textit{user space} e all'indietro, quando in realtà
potrebbe essere più efficiente mantenere tutto in \textit{kernel
  space}. Tratteremo in questa sezione alcune funzioni specialistiche che
permettono di ottimizzare le prestazioni in questo tipo di situazioni.

La prima funzione che è stata ideata per ottimizzare il trasferimento dei dati
fra due file descriptor è \func{sendfile}.\footnote{la funzione è stata
  introdotta con i kernel della serie 2.2, e disponibile dalla \acr{glibc}
  2.1.} La funzione è presente in diverse versioni di Unix (la si ritrova ad
esempio in FreeBSD, HPUX ed altri Unix) ma non è presente né in POSIX.1-2001
né in altri standard (pertanto si eviti di utilizzarla se si devono scrivere
programmi portabili) per cui per essa vengono utilizzati prototipi e
semantiche differenti. Nel caso di Linux il prototipo di \funcd{sendfile} è:


\begin{funcproto}{
\fhead{sys/sendfile.h}
\fdecl{ssize\_t sendfile(int out\_fd, int in\_fd, off\_t *offset, size\_t
    count)}
\fdesc{Copia dei dati da un file descriptor ad un altro.} 
}

{La funzione ritorna il numero di byte trasferiti in caso di successo e $-1$
  per un errore, nel qual caso \var{errno} assumerà uno dei valori:
  \begin{errlist}
    \item[\errcode{EAGAIN}] si è impostata la modalità non bloccante su
      \param{out\_fd} e la scrittura si bloccherebbe.
    \item[\errcode{EINVAL}] i file descriptor non sono validi, o sono bloccati
      (vedi sez.~\ref{sec:file_locking}), o \func{mmap} non è disponibile per
      \param{in\_fd}.
    \item[\errcode{EIO}] si è avuto un errore di lettura da \param{in\_fd}.
    \item[\errcode{ENOMEM}] non c'è memoria sufficiente per la lettura da
      \param{in\_fd}.
  \end{errlist}
  ed inoltre \errcode{EBADF} e \errcode{EFAULT} nel loro significato
  generico.}
\end{funcproto}

La funzione copia direttamente \param{count} byte dal file descriptor
\param{in\_fd} al file descriptor \param{out\_fd}. In caso di successo la
funzione ritorna il numero di byte effettivamente copiati da \param{in\_fd} a
\param{out\_fd} e come per le ordinarie \func{read} e \func{write} questo
valore può essere inferiore a quanto richiesto con \param{count}.

Se il puntatore \param{offset} è nullo la funzione legge i dati a partire
dalla posizione corrente su \param{in\_fd}, altrimenti verrà usata la
posizione indicata dal valore puntato da \param{offset}; in questo caso detto
valore sarà aggiornato, come \textit{value result argument}, per indicare la
posizione del byte successivo all'ultimo che è stato letto, mentre la
posizione corrente sul file non sarà modificata. Se invece \param{offset} è
nullo la posizione corrente sul file sarà aggiornata tenendo conto dei byte
letti da \param{in\_fd}.

Fino ai kernel della serie 2.4 la funzione era utilizzabile su un qualunque
file descriptor, e permetteva di sostituire la invocazione successiva di una
\func{read} e una \func{write} (e l'allocazione del relativo buffer) con una
sola chiamata a \funcd{sendfile}. In questo modo si poteva diminuire il numero
di chiamate al sistema e risparmiare in trasferimenti di dati da
\textit{kernel space} a \textit{user space} e viceversa.  La massima utilità
della funzione si ottiene comunque per il trasferimento di dati da un file su
disco ad un socket di rete,\footnote{questo è il caso classico del lavoro
  eseguito da un server web, ed infatti Apache ha una opzione per il supporto
  esplicito di questa funzione.} dato che in questo caso diventa possibile
effettuare il trasferimento diretto via DMA dal controller del disco alla
scheda di rete, senza neanche allocare un buffer nel kernel (il meccanismo è
detto \textit{zerocopy} in quanto i dati non vengono mai copiati dal kernel,
che si limita a programmare solo le operazioni di lettura e scrittura via DMA)
ottenendo la massima efficienza possibile senza pesare neanche sul processore.

In seguito però ci si accorse che, fatta eccezione per il trasferimento
diretto da file a socket, non sempre \func{sendfile} comportava miglioramenti
significativi delle prestazioni rispetto all'uso in sequenza di \func{read} e
\func{write}. Nel caso generico infatti il kernel deve comunque allocare un
buffer ed effettuare la copia dei dati, e in tal caso spesso il guadagno
ottenibile nel ridurre il numero di chiamate al sistema non compensa le
ottimizzazioni che possono essere fatte da una applicazione in \textit{user
  space} che ha una conoscenza diretta su come questi sono strutturati, per
cui in certi casi si potevano avere anche dei peggioramenti.  Questo ha
portato, per i kernel della serie 2.6,\footnote{per alcune motivazioni di
  questa scelta si può fare riferimento a quanto illustrato da Linus Torvalds
  in \url{http://www.cs.helsinki.fi/linux/linux-kernel/2001-03/0200.html}.}
alla decisione di consentire l'uso della funzione soltanto quando il file da
cui si legge supporta le operazioni di \textit{memory mapping} (vale a dire
non è un socket) e quello su cui si scrive è un socket; in tutti gli altri
casi l'uso di \func{sendfile} da luogo ad un errore di \errcode{EINVAL}.

Nonostante ci possano essere casi in cui \func{sendfile} non migliora le
prestazioni, resta il dubbio se la scelta di disabilitarla sempre per il
trasferimento fra file di dati sia davvero corretta. Se ci sono peggioramenti
di prestazioni infatti si può sempre fare ricorso al metodo ordinario, ma
lasciare a disposizione la funzione consentirebbe se non altro di semplificare
la gestione della copia dei dati fra file, evitando di dover gestire
l'allocazione di un buffer temporaneo per il loro trasferimento. Comunque a
partire dal kernel 2.6.33 la restrizione su \param{out\_fd} è stata rimossa e
questo può essere un file qualunque, rimane però quella di non poter usare un
socket per \param{in\_fd}.

A partire dal kernel 2.6.17 come alternativa a \func{sendfile} è disponibile
la nuova \textit{system call} \func{splice}. Lo scopo di questa funzione è
quello di fornire un meccanismo generico per il trasferimento di dati da o
verso un file, utilizzando un buffer gestito internamente dal
kernel. Descritta in questi termini \func{splice} sembra semplicemente un
``\textsl{dimezzamento}'' di \func{sendfile}, nel senso che un trasferimento
di dati fra due file con \func{sendfile} non sarebbe altro che la lettura
degli stessi su un buffer seguita dalla relativa scrittura, cosa che in questo
caso si dovrebbe eseguire con due chiamate a \func{splice}.

In realtà le due \textit{system call} sono profondamente diverse nel loro
meccanismo di funzionamento;\footnote{questo fino al kernel 2.6.23, dove
  \func{sendfile} è stata reimplementata in termini di \func{splice}, pur
  mantenendo disponibile la stessa interfaccia verso l'\textit{user space}.}
\func{sendfile} infatti, come accennato, non necessita di avere a disposizione
un buffer interno, perché esegue un trasferimento diretto di dati; questo la
rende in generale più efficiente, ma anche limitata nelle sue applicazioni,
dato che questo tipo di trasferimento è possibile solo in casi specifici che
nel caso di Linux questi sono anche solo quelli in cui essa può essere
effettivamente utilizzata.

Il concetto che sta dietro a \func{splice} invece è diverso,\footnote{in
  realtà la proposta originale di Larry Mc Voy non differisce poi tanto negli
  scopi da \func{sendfile}, quello che rende \func{splice} davvero diversa è
  stata la reinterpretazione che ne è stata fatta nell'implementazione su
  Linux realizzata da Jens Anxboe, concetti che sono esposti sinteticamente
  dallo stesso Linus Torvalds in \url{http://kerneltrap.org/node/6505}.} si
tratta semplicemente di una funzione che consente di fare in maniera del tutto
generica delle operazioni di trasferimento di dati fra un file e un buffer
gestito interamente in \textit{kernel space}. In questo caso il cuore della
funzione (e delle affini \func{vmsplice} e \func{tee}, che tratteremo più
avanti) è appunto l'uso di un buffer in \textit{kernel space}, e questo è
anche quello che ne ha semplificato l'adozione, perché l'infrastruttura per la
gestione di un tale buffer è presente fin dagli albori di Unix per la
realizzazione delle \textit{pipe} (vedi sez.~\ref{sec:ipc_unix}). Dal punto di
vista concettuale allora \func{splice} non è altro che una diversa interfaccia
(rispetto alle \textit{pipe}) con cui utilizzare in \textit{user space}
l'oggetto ``\textsl{buffer in kernel space}''.

Così se per una \textit{pipe} o una \textit{fifo} il buffer viene utilizzato
come area di memoria (vedi fig.~\ref{fig:ipc_pipe_singular}) dove appoggiare i
dati che vengono trasferiti da un capo all'altro della stessa per creare un
meccanismo di comunicazione fra processi, nel caso di \func{splice} il buffer
viene usato o come fonte dei dati che saranno scritti su un file, o come
destinazione dei dati che vengono letti da un file. La funzione fornisce
quindi una interfaccia generica che consente di trasferire dati da un buffer
ad un file o viceversa; il prototipo di \funcd{splice}, accessibile solo dopo
aver definito la macro \macro{\_GNU\_SOURCE},\footnote{si ricordi che questa
  funzione non è contemplata da nessuno standard, è presente solo su Linux, e
  pertanto deve essere evitata se si vogliono scrivere programmi portabili.}
è il seguente:

\begin{funcproto}{
\fhead{fcntl.h} 
\fdecl{long splice(int fd\_in, off\_t *off\_in, int fd\_out, off\_t
    *off\_out, size\_t len, \\
\phantom{long splice(}unsigned int flags)}
\fdesc{Trasferisce dati da un file verso una \textit{pipe} o viceversa.} 
}

{La funzione ritorna il numero di byte trasferiti in caso di successo e $-1$
  per un errore, nel qual caso \var{errno} assumerà uno dei valori:
  \begin{errlist}
    \item[\errcode{EBADF}] uno o entrambi fra \param{fd\_in} e \param{fd\_out}
      non sono file descriptor validi o, rispettivamente, non sono stati
      aperti in lettura o scrittura.
    \item[\errcode{EINVAL}] il filesystem su cui si opera non supporta
      \func{splice}, oppure nessuno dei file descriptor è una \textit{pipe},
      oppure si 
      è dato un valore a \param{off\_in} o \param{off\_out} ma il
      corrispondente file è un dispositivo che non supporta la funzione
      \func{lseek}.
    \item[\errcode{ENOMEM}] non c'è memoria sufficiente per l'operazione
      richiesta.
    \item[\errcode{ESPIPE}] o \param{off\_in} o \param{off\_out} non sono
      \val{NULL} ma il corrispondente file descriptor è una \textit{pipe}.
  \end{errlist}
}
\end{funcproto}


La funzione esegue un trasferimento di \param{len} byte dal file descriptor
\param{fd\_in} al file descriptor \param{fd\_out}, uno dei quali deve essere
una \textit{pipe}; l'altro file descriptor può essere qualunque, questo
significa che può essere, oltre che un file di dati, anche un altra
\textit{pipe}, o un socket.  Come accennato una \textit{pipe} non è altro che
un buffer in \textit{kernel space}, per cui a seconda che essa sia usata
per \param{fd\_in} o \param{fd\_out} si avrà rispettivamente la copia dei dati
dal buffer al file o viceversa.

In caso di successo la funzione ritorna il numero di byte trasferiti, che può
essere, come per le normali funzioni di lettura e scrittura su file, inferiore
a quelli richiesti; un valore negativo indicherà un errore mentre un valore
nullo indicherà che non ci sono dati da trasferire (ad esempio si è giunti
alla fine del file in lettura). Si tenga presente che, a seconda del verso del
trasferimento dei dati, la funzione si comporta nei confronti del file
descriptor che fa riferimento al file ordinario, come \func{read} o
\func{write}, e pertanto potrà anche bloccarsi (a meno che non si sia aperto
il suddetto file in modalità non bloccante).

I due argomenti \param{off\_in} e \param{off\_out} consentono di specificare,
come per l'analogo \param{offset} di \func{sendfile}, la posizione all'interno
del file da cui partire per il trasferimento dei dati. Come per
\func{sendfile} un valore nullo indica di usare la posizione corrente sul
file, ed essa sarà aggiornata automaticamente secondo il numero di byte
trasferiti. Un valore non nullo invece deve essere un puntatore ad una
variabile intera che indica la posizione da usare; questa verrà aggiornata, al
ritorno della funzione, al byte successivo all'ultimo byte trasferito.
Ovviamente soltanto uno di questi due argomenti, e più precisamente quello che
fa riferimento al file descriptor non associato alla \textit{pipe}, può essere
specificato come valore non nullo.

Infine l'argomento \param{flags} consente di controllare alcune
caratteristiche del funzionamento della funzione; il contenuto è una maschera
binaria e deve essere specificato come OR aritmetico dei valori riportati in
tab.~\ref{tab:splice_flag}. Alcuni di questi valori vengono utilizzati anche
dalle funzioni \func{vmsplice} e \func{tee} per cui la tabella riporta le
descrizioni complete di tutti i valori possibili anche quando, come per
\const{SPLICE\_F\_GIFT}, questi non hanno effetto su \func{splice}.

\begin{table}[htb]
  \centering
  \footnotesize
  \begin{tabular}[c]{|l|p{10cm}|}
    \hline
    \textbf{Valore} & \textbf{Significato} \\
    \hline
    \hline
    \constd{SPLICE\_F\_MOVE} & Suggerisce al kernel di spostare le pagine
                               di memoria contenenti i dati invece di
                               copiarle: per una maggiore efficienza
                               \func{splice} usa quando possibile i
                               meccanismi della memoria virtuale per
                               eseguire i trasferimenti di dati. In maniera
                               analoga a \func{mmap}), qualora le pagine non
                               possano essere spostate dalla \textit{pipe} o
                               il buffer non corrisponda a pagine intere
                               esse saranno comunque copiate. Viene usato
                               soltanto da \func{splice}.\\ 
    \constd{SPLICE\_F\_NONBLOCK}& Richiede di operare in modalità non
                                  bloccante; questo flag influisce solo sulle
                                  operazioni che riguardano l'I/O da e verso la
                                  \textit{pipe}. Nel caso di \func{splice}
                                  questo significa che la funzione potrà
                                  comunque bloccarsi nell'accesso agli altri
                                  file descriptor (a meno che anch'essi non
                                  siano stati aperti in modalità non
                                  bloccante).\\
    \constd{SPLICE\_F\_MORE} & Indica al kernel che ci sarà l'invio di
                               ulteriori dati in una \func{splice}
                               successiva, questo è un suggerimento utile
                               che viene usato quando \param{fd\_out} è un
                               socket. Questa opzione consente di utilizzare
                               delle opzioni di gestione dei socket che
                               permettono di ottimizzare le trasmissioni via
                               rete (si veda la descrizione di
                               \const{TCP\_CORK} in
                               sez.~\ref{sec:sock_tcp_udp_options} e quella
                               di \const{MSG\_MORE} in
                               sez.~\ref{sec:net_sendmsg}).  Attualmente
                               viene usato solo da \func{splice}, potrà essere
                               implementato in futuro anche per
                               \func{vmsplice} e \func{tee}.\\
    \constd{SPLICE\_F\_GIFT} & Le pagine di memoria utente sono
                               ``\textsl{donate}'' al kernel; questo
                               significa che la cache delle pagine e i dati
                               su disco potranno differire, e che
                               l'applicazione non potrà modificare
                               quest'area di memoria. 
                               Se impostato una seguente \func{splice} che
                               usa \const{SPLICE\_F\_MOVE} potrà spostare le 
                               pagine con successo, altrimenti esse dovranno
                               essere copiate; per usare questa opzione i
                               dati dovranno essere opportunamente allineati
                               in posizione ed in dimensione alle pagine di
                               memoria. Viene usato soltanto da
                               \func{vmsplice}.\\
    \hline
  \end{tabular}
  \caption{Le costanti che identificano i bit della maschera binaria
    dell'argomento \param{flags} di \func{splice}, \func{vmsplice} e
    \func{tee}.} 
  \label{tab:splice_flag}
\end{table}


Per capire meglio il funzionamento di \func{splice} vediamo un esempio con un
semplice programma che usa questa funzione per effettuare la copia di un file
su un altro senza utilizzare buffer in \textit{user space}. Lo scopo del
programma è quello di eseguire la copia dei dati con \func{splice}, questo
significa che si dovrà usare la funzione due volte, prima per leggere i dati
dal file di ingresso e poi per scriverli su quello di uscita, appoggiandosi ad
una \textit{pipe}: lo schema del flusso dei dati è illustrato in
fig.~\ref{fig:splicecp_data_flux}.

\begin{figure}[htb]
  \centering
  \includegraphics[height=3.5cm]{img/splice_copy}
  \caption{Struttura del flusso di dati usato dal programma \texttt{splicecp}.}
  \label{fig:splicecp_data_flux}
\end{figure}

Il programma si chiama \texttt{splicecp.c} ed il codice completo è disponibile
coi sorgenti allegati alla guida, il corpo principale del programma, che non
contiene la sezione di gestione delle opzioni, le funzioni di ausilio, le
aperture dei file di ingresso e di uscita passati come argomenti e quella
della \textit{pipe} intermedia, è riportato in fig.~\ref{fig:splice_example}.

\begin{figure}[!htb]
  \footnotesize \centering
  \begin{minipage}[c]{\codesamplewidth}
    \includecodesample{listati/splicecp.c}
  \end{minipage}
  \normalsize
  \caption{Esempio di codice che usa \func{splice} per effettuare la copia di
    un file.}
  \label{fig:splice_example}
\end{figure}

Il ciclo principale (\texttt{\small 13-38}) inizia con la lettura dal file
sorgente tramite la prima \func{splice} (\texttt{\small 14-15}), in questo
caso si è usato come primo argomento il file descriptor del file sorgente e
come terzo quello del capo in scrittura della \textit{pipe}. Il funzionamento
delle \textit{pipe} e l'uso della coppia di file descriptor ad esse associati
è trattato in dettaglio in sez.~\ref{sec:ipc_unix}; non ne parleremo qui dato
che nell'ottica dell'uso di \func{splice} questa operazione corrisponde
semplicemente al trasferimento dei dati dal file al buffer in \textit{kernel
  space}.

La lettura viene eseguita in blocchi pari alla dimensione specificata
dall'opzione \texttt{-s} (il default è 4096); essendo in questo caso
\func{splice} equivalente ad una \func{read} sul file, se ne controlla il
valore di uscita in \var{nread} che indica quanti byte sono stati letti, se
detto valore è nullo (\texttt{\small 16}) questo significa che si è giunti
alla fine del file sorgente e pertanto l'operazione di copia è conclusa e si
può uscire dal ciclo arrivando alla conclusione del programma (\texttt{\small
  59}). In caso di valore negativo (\texttt{\small 17-24}) c'è stato un
errore ed allora si ripete la lettura (\texttt{\small 16}) se questo è dovuto
ad una interruzione, o altrimenti si esce con un messaggio di errore
(\texttt{\small 21-23}).

Una volta completata con successo la lettura si avvia il ciclo di scrittura
(\texttt{\small 25-37}); questo inizia (\texttt{\small 26-27}) con la
seconda \func{splice} che cerca di scrivere gli \var{nread} byte letti, si
noti come in questo caso il primo argomento faccia di nuovo riferimento alla
\textit{pipe} (in questo caso si usa il capo in lettura, per i dettagli si
veda al solito sez.~\ref{sec:ipc_unix}) mentre il terzo sia il file descriptor
del file di destinazione.

Di nuovo si controlla il numero di byte effettivamente scritti restituito in
\var{nwrite} e in caso di errore al solito si ripete la scrittura se questo è
dovuto a una interruzione o si esce con un messaggio negli altri casi
(\texttt{\small 28-35}). Infine si chiude il ciclo di scrittura sottraendo
(\texttt{\small 37}) il numero di byte scritti a quelli di cui è richiesta la
scrittura,\footnote{in questa parte del ciclo \var{nread}, il cui valore
  iniziale è dato dai byte letti dalla precedente chiamata a \func{splice},
  viene ad assumere il significato di byte da scrivere.} così che il ciclo di
scrittura venga ripetuto fintanto che il valore risultante sia maggiore di
zero, indice che la chiamata a \func{splice} non ha esaurito tutti i dati
presenti sul buffer.

Si noti come il programma sia concettualmente identico a quello che si sarebbe
scritto usando \func{read} al posto della prima \func{splice} e \func{write}
al posto della seconda, utilizzando un buffer in \textit{user space} per
eseguire la copia dei dati, solo che in questo caso non è stato necessario
allocare nessun buffer e non si è trasferito nessun dato in \textit{user
  space}.  Si noti anche come si sia usata la combinazione
\texttt{SPLICE\_F\_MOVE | SPLICE\_F\_MORE } per l'argomento \param{flags} di
\func{splice}, infatti anche se un valore nullo avrebbe dato gli stessi
risultati, l'uso di questi flag, che si ricordi servono solo a dare
suggerimenti al kernel, permette in genere di migliorare le prestazioni.

Come accennato con l'introduzione di \func{splice} sono state realizzate anche
altre due \textit{system call}, \func{vmsplice} e \func{tee}, che utilizzano
la stessa infrastruttura e si basano sullo stesso concetto di manipolazione e
trasferimento di dati attraverso un buffer in \textit{kernel space}; benché
queste non attengono strettamente ad operazioni di trasferimento dati fra file
descriptor, le tratteremo qui, essendo strettamente correlate fra loro.

La prima funzione, \funcd{vmsplice}, è la più simile a \func{splice} e come
indica il suo nome consente di trasferire i dati dalla memoria virtuale di un
processo (ad esempio per un file mappato in memoria) verso una \textit{pipe};
il suo prototipo è:

\begin{funcproto}{
\fhead{fcntl.h} 
\fhead{sys/uio.h}
\fdecl{long vmsplice(int fd, const struct iovec *iov, unsigned long nr\_segs,\\
\phantom{long vmsplice(}unsigned int flags)}
\fdesc{Trasferisce dati dalla memoria di un processo verso una \textit{pipe}.} 
}

{La funzione ritorna il numero di byte trasferiti in caso di successo e $-1$
  per un errore, nel qual caso \var{errno} assumerà uno dei valori:
  \begin{errlist}
    \item[\errcode{EBADF}] o \param{fd} non è un file descriptor valido o non
      fa riferimento ad una \textit{pipe}.
    \item[\errcode{EINVAL}] si è usato un valore nullo per \param{nr\_segs}
      oppure si è usato \const{SPLICE\_F\_GIFT} ma la memoria non è allineata.
    \item[\errcode{ENOMEM}] non c'è memoria sufficiente per l'operazione
      richiesta.
  \end{errlist}
}
\end{funcproto}

La \textit{pipe} indicata da \param{fd} dovrà essere specificata tramite il
file descriptor corrispondente al suo capo aperto in scrittura (di nuovo si
faccia riferimento a sez.~\ref{sec:ipc_unix}), mentre per indicare quali
segmenti della memoria del processo devono essere trasferiti verso di essa si
dovrà utilizzare un vettore di strutture \struct{iovec} (vedi
fig.~\ref{fig:file_iovec}), esattamente con gli stessi criteri con cui le si
usano per l'I/O vettorizzato, indicando gli indirizzi e le dimensioni di
ciascun segmento di memoria su cui si vuole operare; le dimensioni del
suddetto vettore devono essere passate nell'argomento \param{nr\_segs} che
indica il numero di segmenti di memoria da trasferire.  Sia per il vettore che
per il valore massimo di \param{nr\_segs} valgono le stesse limitazioni
illustrate in sez.~\ref{sec:file_multiple_io}.

In caso di successo la funzione ritorna il numero di byte trasferiti sulla
\textit{pipe}. In generale, se i dati una volta creati non devono essere
riutilizzati (se cioè l'applicazione che chiama \func{vmsplice} non
modificherà più la memoria trasferita), è opportuno utilizzare
per \param{flag} il valore \const{SPLICE\_F\_GIFT}; questo fa sì che il kernel
possa rimuovere le relative pagine dalla cache della memoria virtuale, così
che queste possono essere utilizzate immediatamente senza necessità di
eseguire una copia dei dati che contengono.

La seconda funzione aggiunta insieme a \func{splice} è \func{tee}, che deve il
suo nome all'omonimo comando in \textit{user space}, perché in analogia con
questo permette di duplicare i dati in ingresso su una \textit{pipe} su
un'altra \textit{pipe}. In sostanza, sempre nell'ottica della manipolazione
dei dati su dei buffer in \textit{kernel space}, la funzione consente di
eseguire una copia del contenuto del buffer stesso. Il prototipo di
\funcd{tee} è il seguente:

\begin{funcproto}{
\fhead{fcntl.h}
\fdecl{long tee(int fd\_in, int fd\_out, size\_t len, unsigned int
    flags)}
\fdesc{Duplica i dati da una \textit{pipe} ad un'altra.} 
}

{La funzione ritorna restituisce il numero di byte copiati in caso di successo
  e $-1$ per un errore, nel qual caso \var{errno} assumerà uno dei valori:
  \begin{errlist}
    \item[\errcode{EINVAL}] o uno fra \param{fd\_in} e \param{fd\_out} non fa
      riferimento ad una \textit{pipe} o entrambi fanno riferimento alla
      stessa \textit{pipe}.
    \item[\errcode{ENOMEM}] non c'è memoria sufficiente per l'operazione
      richiesta.
  \end{errlist}
}
\end{funcproto}

La funzione copia \param{len} byte del contenuto di una \textit{pipe} su di
un'altra; \param{fd\_in} deve essere il capo in lettura della \textit{pipe}
sorgente e \param{fd\_out} il capo in scrittura della \textit{pipe}
destinazione; a differenza di quanto avviene con \func{read} i dati letti con
\func{tee} da \param{fd\_in} non vengono \textsl{consumati} e restano
disponibili sulla \textit{pipe} per una successiva lettura (di nuovo per il
comportamento delle \textit{pipe} si veda sez.~\ref{sec:ipc_unix}). Al
momento\footnote{quello della stesura di questo paragrafo, avvenuta il Gennaio
  2010, in futuro potrebbe essere implementato anche \const{SPLICE\_F\_MORE}.}
il solo valore utilizzabile per \param{flag}, fra quelli elencati in
tab.~\ref{tab:splice_flag}, è \const{SPLICE\_F\_NONBLOCK} che rende la
funzione non bloccante.

La funzione restituisce il numero di byte copiati da una \textit{pipe}
all'altra (o $-1$ in caso di errore), un valore nullo indica che non ci sono
byte disponibili da copiare e che il capo in scrittura della \textit{pipe} è
stato chiuso; si tenga presente però che questo non avviene se si è impostato
il flag \const{SPLICE\_F\_NONBLOCK}, in tal caso infatti si avrebbe un errore
di \errcode{EAGAIN}. Un esempio di realizzazione del comando \texttt{tee}
usando questa funzione, ripreso da quello fornito nella pagina di manuale e
dall'esempio allegato al patch originale, è riportato in
fig.~\ref{fig:tee_example}. Il programma consente di copiare il contenuto
dello \textit{standard input} sullo \textit{standard output} e su un file
specificato come argomento, il codice completo si trova nel file
\texttt{tee.c} dei sorgenti allegati alla guida.

\begin{figure}[!htb]
  \footnotesize \centering
  \begin{minipage}[c]{\codesamplewidth}
    \includecodesample{listati/tee.c}
  \end{minipage}
  \normalsize
  \caption{Esempio di codice che usa \func{tee} per copiare i dati dello
    standard input sullo standard output e su un file.}
  \label{fig:tee_example}
\end{figure}

La prima parte del programma, che si è omessa per brevità, si cura
semplicemente di controllare che sia stato fornito almeno un argomento (il
nome del file su cui scrivere), di aprirlo e che sia lo standard input che lo
standard output corrispondano ad una \textit{pipe}.

Il ciclo principale (\texttt{\small 11-32}) inizia con la chiamata a
\func{tee} che duplica il contenuto dello standard input sullo standard output
(\texttt{\small 13}), questa parte è del tutto analoga ad una lettura ed
infatti come nell'esempio di fig.~\ref{fig:splice_example} si controlla il
valore di ritorno della funzione in \var{len}; se questo è nullo significa che
non ci sono più dati da leggere e si chiude il ciclo (\texttt{\small 14}), se
è negativo c'è stato un errore, ed allora si ripete la chiamata se questo è
dovuto ad una interruzione (\texttt{\small 15-48}) o si stampa un messaggio
di errore e si esce negli altri casi (\texttt{\small 18-21}).

Una volta completata la copia dei dati sullo \textit{standard output} si
possono estrarre dallo \textit{standard input} e scrivere sul file, di nuovo
su usa un ciclo di scrittura (\texttt{\small 24-31}) in cui si ripete una
chiamata a \func{splice} (\texttt{\small 25}) fintanto che non si sono scritti
tutti i \var{len} byte copiati in precedenza con \func{tee} (il funzionamento
è identico all'analogo ciclo di scrittura del precedente esempio di
fig.~\ref{fig:splice_example}).

Infine una nota finale riguardo \func{splice}, \func{vmsplice} e \func{tee}:
occorre sottolineare che benché finora si sia parlato di trasferimenti o copie
di dati in realtà nella implementazione di queste \textit{system call} non è
affatto detto che i dati vengono effettivamente spostati o copiati, il kernel
infatti realizza le \textit{pipe} come un insieme di puntatori\footnote{per
  essere precisi si tratta di un semplice buffer circolare, un buon articolo
  sul tema si trova su \url{http://lwn.net/Articles/118750/}.}  alle pagine di
memoria interna che contengono i dati, per questo una volta che i dati sono
presenti nella memoria del kernel tutto quello che viene fatto è creare i
suddetti puntatori ed aumentare il numero di referenze; questo significa che
anche con \func{tee} non viene mai copiato nessun byte, vengono semplicemente
copiati i puntatori.

% TODO?? dal 2.6.25 splice ha ottenuto il supporto per la ricezione su rete


% TODO trattare qui copy_file_range (vedi http://lwn.net/Articles/659523/),
% introdotta nel kernel 4.5

\subsection{Gestione avanzata dell'accesso ai dati dei file}
\label{sec:file_fadvise}

Nell'uso generico dell'interfaccia per l'accesso al contenuto dei file le
operazioni di lettura e scrittura non necessitano di nessun intervento di
supervisione da parte dei programmi, si eseguirà una \func{read} o una
\func{write}, i dati verranno passati al kernel che provvederà ad effettuare
tutte le operazioni (e a gestire il \textit{caching} dei dati) per portarle a
termine in quello che ritiene essere il modo più efficiente.

Il problema è che il concetto di migliore efficienza impiegato dal kernel è
relativo all'uso generico, mentre esistono molti casi in cui ci sono esigenze
specifiche dei singoli programmi, che avendo una conoscenza diretta di come
verranno usati i file, possono necessitare di effettuare delle ottimizzazioni
specifiche, relative alle proprie modalità di I/O sugli stessi. Tratteremo in
questa sezione una serie funzioni che consentono ai programmi di ottimizzare
il loro accesso ai dati dei file e controllare la gestione del relativo
\textit{caching}.

\itindbeg{read-ahead}

Una prima funzione che può essere utilizzata per modificare la gestione
ordinaria dell'I/O su un file è \funcd{readahead} (questa è una funzione
specifica di Linux, introdotta con il kernel 2.4.13, e non deve essere usata
se si vogliono scrivere programmi portabili), che consente di richiedere una
lettura anticipata del contenuto dello stesso in cache, così che le seguenti
operazioni di lettura non debbano subire il ritardo dovuto all'accesso al
disco; il suo prototipo è:

\begin{funcproto}{
\fhead{fcntl.h}
\fdecl{ssize\_t readahead(int fd, off64\_t *offset, size\_t count)}
\fdesc{Esegue una lettura preventiva del contenuto di un file in cache.} 
}

{La funzione ritorna $0$ in caso di successo e $-1$ per un errore, nel qual
  caso \var{errno} assumerà uno dei valori: 
  \begin{errlist}
    \item[\errcode{EBADF}] l'argomento \param{fd} non è un file descriptor
      valido o non è aperto in lettura.
    \item[\errcode{EINVAL}] l'argomento \param{fd} si riferisce ad un tipo di
      file che non supporta l'operazione (come una \textit{pipe} o un socket).
  \end{errlist}
}
\end{funcproto}

La funzione richiede che venga letto in anticipo il contenuto del file
\param{fd} a partire dalla posizione \param{offset} e per un ammontare di
\param{count} byte, in modo da portarlo in cache.  La funzione usa la memoria
virtuale ed il meccanismo della paginazione per cui la lettura viene eseguita
in blocchi corrispondenti alle dimensioni delle pagine di memoria, ed i valori
di \param{offset} e \param{count} vengono arrotondati di conseguenza.

La funzione estende quello che è un comportamento normale del kernel che,
quando si legge un file, aspettandosi che l'accesso prosegua, esegue sempre
una lettura preventiva di una certa quantità di dati; questo meccanismo di
lettura anticipata viene chiamato \textit{read-ahead}, da cui deriva il nome
della funzione. La funzione \func{readahead}, per ottimizzare gli accessi a
disco, effettua la lettura in cache della sezione richiesta e si blocca
fintanto che questa non viene completata.  La posizione corrente sul file non
viene modificata ed indipendentemente da quanto indicato con \param{count} la
lettura dei dati si interrompe una volta raggiunta la fine del file.

Si può utilizzare questa funzione per velocizzare le operazioni di lettura
all'interno di un programma tutte le volte che si conosce in anticipo quanti
dati saranno necessari nelle elaborazioni successive. Si potrà così
concentrare in un unico momento (ad esempio in fase di inizializzazione) la
lettura dei dati da disco, così da ottenere una migliore velocità di risposta
nelle operazioni successive.

\itindend{read-ahead}

Il concetto di \func{readahead} viene generalizzato nello standard
POSIX.1-2001 dalla funzione \func{posix\_fadvise} (anche se
l'argomento \param{len} è stato modificato da \type{size\_t} a \type{off\_t}
nella revisione POSIX.1-2003 TC1) che consente di ``\textsl{avvisare}'' il
kernel sulle modalità con cui si intende accedere nel futuro ad una certa
porzione di un file, così che esso possa provvedere le opportune
ottimizzazioni; il prototipo di \funcd{posix\_fadvise}\footnote{la funzione è
  stata introdotta su Linux solo a partire dal kernel 2.5.60, ed è disponibile
  soltanto se è stata definita la macro \macro{\_XOPEN\_SOURCE} ad valore di
  almeno \texttt{600} o la macro \macro{\_POSIX\_C\_SOURCE} ad valore di
  almeno \texttt{200112L}.} è:


\begin{funcproto}{
\fhead{fcntl.h}
\fdecl{int posix\_fadvise(int fd, off\_t offset, off\_t len, int advice)}
\fdesc{Dichiara al kernel le future modalità di accesso ad un file.}
}

{La funzione ritorna $0$ in caso di successo e $-1$ per un errore, nel qual
  caso \var{errno} assumerà uno dei valori: 
  \begin{errlist}
    \item[\errcode{EBADF}] l'argomento \param{fd} non è un file descriptor
      valido.
    \item[\errcode{EINVAL}] il valore di \param{advice} non è valido o
      \param{fd} si riferisce ad un tipo di file che non supporta l'operazione
      (come una \textit{pipe} o un socket).
    \item[\errcode{ESPIPE}] previsto dallo standard se \param{fd} è una
      \textit{pipe} o un socket (ma su Linux viene restituito
      \errcode{EINVAL}).
  \end{errlist}
}
\end{funcproto}

La funzione dichiara al kernel le modalità con cui intende accedere alla
regione del file indicato da \param{fd} che inizia alla posizione
\param{offset} e si estende per \param{len} byte. Se per \param{len} si usa un
valore nullo la regione coperta sarà da \param{offset} alla fine del file, ma
questo è vero solo per le versioni più recenti, fino al kernel 2.6.6 il valore
nullo veniva interpretato letteralmente. Le modalità sono indicate
dall'argomento \param{advice} che è una maschera binaria dei valori illustrati
in tab.~\ref{tab:posix_fadvise_flag}, che riprendono il significato degli
analoghi già visti in sez.~\ref{sec:file_memory_map} per
\func{madvise}.\footnote{dato che si tratta dello stesso tipo di funzionalità,
  in questo caso applicata direttamente al sistema ai contenuti di un file
  invece che alla sua mappatura in memoria.} Si tenga presente comunque che la
funzione dà soltanto un avvertimento, non esiste nessun vincolo per il kernel,
che utilizza semplicemente l'informazione.

\begin{table}[htb]
  \centering
  \footnotesize
  \begin{tabular}[c]{|l|p{10cm}|}
    \hline
    \textbf{Valore} & \textbf{Significato} \\
    \hline
    \hline
    \constd{POSIX\_FADV\_NORMAL}  & Non ci sono avvisi specifici da fare
                                   riguardo le modalità di accesso, il
                                   comportamento sarà identico a quello che si
                                   avrebbe senza nessun avviso.\\ 
    \constd{POSIX\_FADV\_SEQUENTIAL}& L'applicazione si aspetta di accedere di
                                   accedere ai dati specificati in maniera
                                   sequenziale, a partire dalle posizioni più
                                   basse.\\ 
    \constd{POSIX\_FADV\_RANDOM}  & I dati saranno letti in maniera
                                   completamente causale.\\
    \constd{POSIX\_FADV\_NOREUSE} & I dati saranno acceduti una sola volta.\\ 
    \constd{POSIX\_FADV\_WILLNEED}& I dati saranno acceduti a breve.\\ 
    \constd{POSIX\_FADV\_DONTNEED}& I dati non saranno acceduti a breve.\\ 
    \hline
  \end{tabular}
  \caption{Valori delle costanti usabili per l'argomento \param{advice} di
    \func{posix\_fadvise}, che indicano la modalità con cui si intende accedere
    ad un file.}
  \label{tab:posix_fadvise_flag}
\end{table}

Come \func{madvise} anche \func{posix\_fadvise} si appoggia al sistema della
memoria virtuale ed al meccanismo standard del \textit{read-ahead} utilizzato
dal kernel; in particolare utilizzando il valore
\const{POSIX\_FADV\_SEQUENTIAL} si raddoppia la dimensione dell'ammontare di
dati letti preventivamente rispetto al default, aspettandosi appunto una
lettura sequenziale che li utilizzerà, mentre con \const{POSIX\_FADV\_RANDOM}
si disabilita del tutto il suddetto meccanismo, dato che con un accesso del
tutto casuale è inutile mettersi a leggere i dati immediatamente successivi
gli attuali; infine l'uso di \const{POSIX\_FADV\_NORMAL} consente di
riportarsi al comportamento di default.

Le due modalità \const{POSIX\_FADV\_NOREUSE} e \const{POSIX\_FADV\_WILLNEED}
fino al kernel 2.6.18 erano equivalenti, a partire da questo kernel la prima
viene non ha più alcun effetto, mentre la seconda dà inizio ad una lettura in
cache della regione del file indicata.  La quantità di dati che verranno letti
è ovviamente limitata in base al carico che si viene a creare sul sistema
della memoria virtuale, ma in genere una lettura di qualche megabyte viene
sempre soddisfatta (ed un valore superiore è solo raramente di qualche
utilità). In particolare l'uso di \const{POSIX\_FADV\_WILLNEED} si può
considerare l'equivalente POSIX di \func{readahead}.

Infine con \const{POSIX\_FADV\_DONTNEED} si dice al kernel di liberare le
pagine di cache occupate dai dati presenti nella regione di file indicata.
Questa è una indicazione utile che permette di alleggerire il carico sulla
cache, ed un programma può utilizzare periodicamente questa funzione per
liberare pagine di memoria da dati che non sono più utilizzati per far posto a
nuovi dati utili; la pagina di manuale riporta l'esempio dello streaming di
file di grosse dimensioni, dove le pagine occupate dai dati già inviati
possono essere tranquillamente scartate.

Sia \func{posix\_fadvise} che \func{readahead} attengono alla ottimizzazione
dell'accesso in lettura; lo standard POSIX.1-2001 prevede anche una funzione
specifica per le operazioni di scrittura, \funcd{posix\_fallocate} (la
funzione è stata introdotta a partire dalle glibc 2.1.94), che consente di
preallocare dello spazio disco per assicurarsi che una seguente scrittura non
fallisca, il suo prototipo, anch'esso disponibile solo se si definisce la
macro \macro{\_XOPEN\_SOURCE} ad almeno 600, è:

\begin{funcproto}{
\fhead{fcntl.h}
\fdecl{int posix\_fallocate(int fd, off\_t offset, off\_t len)}
\fdesc{Richiede la allocazione di spazio disco per un file.} 
}

{La funzione ritorna $0$ in caso di successo e direttamente un codice di
  errore altrimenti, in tal caso \var{errno} non viene impostato, e si otterrà
  direttamente uno dei valori:
  \begin{errlist}
    \item[\errcode{EBADF}] l'argomento \param{fd} non è un file descriptor
      valido o non è aperto in scrittura.
    \item[\errcode{EINVAL}] o \param{offset} o \param{len} sono minori di
      zero.
    \item[\errcode{EFBIG}] il valore di (\param{offset} + \param{len}) eccede
      la dimensione massima consentita per un file.
    \item[\errcode{ENODEV}] l'argomento \param{fd} non fa riferimento ad un
      file regolare.
    \item[\errcode{ENOSPC}] non c'è sufficiente spazio disco per eseguire
      l'operazione. 
    \item[\errcode{ESPIPE}] l'argomento \param{fd} è una \textit{pipe}.
  \end{errlist}
}
\end{funcproto}

La funzione assicura che venga allocato sufficiente spazio disco perché sia
possibile scrivere sul file indicato dall'argomento \param{fd} nella regione
che inizia dalla posizione \param{offset} e si estende per \param{len} byte;
se questa regione si estende oltre la fine del file le dimensioni di
quest'ultimo saranno incrementate di conseguenza. Dopo aver eseguito con
successo la funzione è garantito che una successiva scrittura nella regione
indicata non fallirà per mancanza di spazio disco. La funzione non ha nessun
effetto né sul contenuto, né sulla posizione corrente del file.

Ci si può chiedere a cosa possa servire una funzione come
\func{posix\_fallocate} dato che è sempre possibile ottenere l'effetto voluto
eseguendo esplicitamente sul file la scrittura di una serie di zeri (usando
\funcd{pwrite} per evitare spostamenti della posizione corrente sul file) per
l'estensione di spazio necessaria qualora il file debba essere esteso o abbia
dei buchi.\footnote{si ricordi che occorre scrivere per avere l'allocazione e
  che l'uso di \func{truncate} per estendere un file creerebbe soltanto uno
  \textit{sparse file} (vedi sez.~\ref{sec:file_lseek}) senza una effettiva
  allocazione dello spazio disco.}  In realtà questa è la modalità con cui la
funzione veniva realizzata nella prima versione fornita dalla \acr{glibc}, per
cui la funzione costituiva in sostanza soltanto una standardizzazione delle
modalità di esecuzione di questo tipo di allocazioni.

Questo metodo, anche se funzionante, comporta però l'effettiva esecuzione una
scrittura su tutto lo spazio disco necessario, da fare al momento della
richiesta di allocazione, pagandone il conseguente prezzo in termini di
prestazioni; il tutto quando in realtà servirebbe solo poter riservare lo
spazio per poi andarci a scrivere, una sola volta, quando il contenuto finale
diventa effettivamente disponibile.  Per poter fare tutto questo è però
necessario il supporto da parte del kernel, e questo è divenuto disponibile
solo a partire dal kernel 2.6.23 in cui è stata introdotta la nuova
\textit{system call} \func{fallocate},\footnote{non è detto che la funzione
  sia disponibile per tutti i filesystem, ad esempio per XFS il supporto è
  stato introdotto solo a partire dal kernel 2.6.25.}  che consente di
realizzare direttamente all'interno del kernel l'allocazione dello spazio
disco così da poter realizzare una versione di \func{posix\_fallocate} con
prestazioni molto più elevate; nella \acr{glibc} la nuova \textit{system call}
viene sfruttata per la realizzazione di \func{posix\_fallocate} a partire
dalla versione 2.10.

Trattandosi di una funzione di servizio, ed ovviamente disponibile
esclusivamente su Linux, inizialmente \funcd{fallocate} non era stata definita
come funzione di libreria,\footnote{pertanto poteva essere invocata soltanto
  in maniera indiretta con l'ausilio di \func{syscall}, vedi
  sez.~\ref{sec:proc_syscall}, come \code{long fallocate(int fd, int mode,
      loff\_t offset, loff\_t len)}.} ma a partire dalla \acr{glibc} 2.10 è
  stato fornito un supporto esplicito; il suo prototipo è:

\begin{funcproto}{
\fhead{fcntl.h} 
\fdecl{int fallocate(int fd, int mode, off\_t offset, off\_t len)}
\fdesc{Prealloca dello spazio disco per un file.} 
}

{La funzione ritorna $0$ in caso di successo e $-1$ per un errore, nel qual
  caso \var{errno} assumerà uno dei valori: 
  \begin{errlist}
    \item[\errcode{EBADF}] \param{fd} non fa riferimento ad un file descriptor
      valido aperto in scrittura.
    \item[\errcode{EFBIG}] la somma di \param{offset} e \param{len} eccede le
      dimensioni massime di un file. 
    \item[\errcode{EINVAL}] \param{offset} è minore di zero o \param{len} è
      minore o uguale a zero. 
    \item[\errcode{ENODEV}] \param{fd} non fa riferimento ad un file ordinario
      o a una directory. 
    \item[\errcode{EPERM}] il file è immutabile o \textit{append-only} (vedi
      sez.~\ref{sec:file_perm_overview}).
    \item[\errcode{ENOSYS}] il filesystem contenente il file associato
      a \param{fd} non supporta \func{fallocate}.
    \item[\errcode{EOPNOTSUPP}] il filesystem contenente il file associato
      a \param{fd} non supporta l'operazione \param{mode}.
  \end{errlist}
  ed inoltre \errval{EINTR}, \errval{EIO} e \errval{ENOSPC} nel loro significato
  generico.}
\end{funcproto}

La funzione prende gli stessi argomenti di \func{posix\_fallocate} con lo
stesso significato, a cui si aggiunge l'argomento \param{mode} che indica le
modalità di allocazione; se questo è nullo il comportamento è identico a
quello di \func{posix\_fallocate} e si può considerare \func{fallocate} come
l'implementazione ottimale della stessa a livello di kernel.

Inizialmente l'unico altro valore possibile per \param{mode} era
\const{FALLOC\_FL\_KEEP\_SIZE} che richiede che la dimensione del file
(quella ottenuta nel campo \var{st\_size} di una struttura \struct{stat} dopo
una chiamata a \texttt{fstat}) non venga modificata anche quando la somma
di \param{offset} e \param{len} eccede la dimensione corrente, che serve
quando si deve comunque preallocare dello spazio per scritture in append. In
seguito sono stati introdotti altri valori, riassunti in
tab.\ref{tab:fallocate_mode}, per compiere altre operazioni relative alla
allocazione dello spazio disco dei file.

\begin{table}[htb]
  \centering
  \footnotesize
  \begin{tabular}[c]{|l|p{10cm}|}
    \hline
    \textbf{Valore} & \textbf{Significato} \\
    \hline
    \hline
    \constd{FALLOC\_FL\_INSERT}     & .\\
    \constd{FALLOC\_FL\_COLLAPSE\_RANGE}& .\\ 
    \constd{FALLOC\_FL\_KEEP\_SIZE} & Mantiene invariata la dimensione del
                                     file, pur allocando lo spazio disco anche
                                     oltre la dimensione corrente del file.\\
    \constd{FALLOC\_FL\_PUNCH\_HOLE}& Crea un \textsl{buco} nel file (vedi
                                     sez.~\ref{sec:file_lseek}) rendendolo una
                                     \textit{sparse file} (dal kernel
                                     2.6.38).\\  
    \constd{FALLOC\_FL\_ZERO\_RANGE}& .\\ 
    \hline
  \end{tabular}
  \caption{Valori delle costanti usabili per l'argomento \param{mode} di
    \func{fallocate}.}
  \label{tab:fallocate_mode}
\end{table}

In particolare con \const{FALLOC\_FL\_PUNCH\_HOLE} è possibile scartare il
contenuto della sezione di file indicata da \param{offser} e \param{len},
creando un \textsl{buco} (si ricordi quanto detto in
sez.~\ref{sec:file_lseek}); i blocchi del file interamente contenuti
nell'intervallo verranno disallocati, la parte di intervallo contenuta
parzialmente in altri blocchi verrà riempita con zeri e la lettura dal file
restituirà degli zeri per tutto l'intervallo indicato. In sostanza si rende il
file uno \textit{sparse file} a posteriori.

% vedi http://lwn.net/Articles/226710/ e http://lwn.net/Articles/240571/
% http://kernelnewbies.org/Linux_2_6_23

% TODO aggiungere FALLOC_FL_ZERO_RANGE e FALLOC_FL_COLLAPSE_RANGE, inseriti
% nel kernel 3.15 (sul secondo vedi http://lwn.net/Articles/589260/), vedi
% anche http://lwn.net/Articles/629965/

% TODO aggiungere FALLOC_FL_INSERT vedi  http://lwn.net/Articles/629965/


% TODO non so dove trattarli, ma dal 2.6.39 ci sono i file handle, vedi
% http://lwn.net/Articles/432757/ (probabilmente da associare alle
% at-functions) 


% LocalWords:  dell'I locking multiplexing cap sez system call socket BSD GID
% LocalWords:  descriptor client deadlock NONBLOCK EAGAIN polling select kernel
% LocalWords:  pselect like sys unistd int fd readfds writefds exceptfds struct
% LocalWords:  timeval errno EBADF EINTR EINVAL ENOMEM sleep tab signal void of
% LocalWords:  CLR ISSET SETSIZE POSIX read NULL nell'header l'header glibc fig
% LocalWords:  libc header psignal sigmask SOURCE XOPEN timespec sigset race DN
% LocalWords:  condition sigprocmask tut self trick oldmask poll XPG pollfd l'I
% LocalWords:  ufds unsigned nfds RLIMIT NOFILE EFAULT ndfs events revents hung
% LocalWords:  POLLIN POLLRDNORM POLLRDBAND POLLPRI POLLOUT POLLWRNORM POLLERR
% LocalWords:  POLLWRBAND POLLHUP POLLNVAL POLLMSG SysV stream ASYNC SETOWN FAQ
% LocalWords:  GETOWN fcntl SETFL SIGIO SETSIG Stevens driven siginfo sigaction
% LocalWords:  all'I nell'I Frequently Unanswered Question SIGHUP lease holder
% LocalWords:  breaker truncate write SETLEASE arg RDLCK WRLCK UNLCK GETLEASE
% LocalWords:  uid capabilities capability EWOULDBLOCK notify dall'OR ACCESS st
% LocalWords:  pread readv MODIFY pwrite writev ftruncate creat mknod mkdir buf
% LocalWords:  symlink rename DELETE unlink rmdir ATTRIB chown chmod utime lio
% LocalWords:  MULTISHOT thread linkando librt layer aiocb asyncronous control
% LocalWords:  block ASYNCHRONOUS lseek fildes nbytes reqprio PRIORITIZED sigev
% LocalWords:  PRIORITY SCHEDULING opcode listio sigevent signo value function
% LocalWords:  aiocbp ENOSYS append error const EINPROGRESS fsync return ssize
% LocalWords:  DSYNC fdatasync SYNC cancel ECANCELED ALLDONE CANCELED suspend
% LocalWords:  NOTCANCELED list nent timout sig NOP WAIT NOWAIT size count iov
% LocalWords:  iovec vector EOPNOTSUPP EISDIR len memory mapping mapped swap NB
% LocalWords:  mmap length prot flags off MAP FAILED ANONYMOUS EACCES SHARED SH
% LocalWords:  only ETXTBSY DENYWRITE ENODEV filesystem EPERM EXEC noexec table
% LocalWords:  ENFILE lenght segment violation SIGSEGV FIXED msync munmap copy
% LocalWords:  DoS Denial Service EXECUTABLE NORESERVE LOCKED swapping stack fs
% LocalWords:  GROWSDOWN ANON POPULATE prefaulting SIGBUS fifo VME fork old SFD
% LocalWords:  exec atime ctime mtime mprotect addr mremap address new Failed
% LocalWords:  long MAYMOVE realloc VMA virtual Ingo Molnar remap pages pgoff
% LocalWords:  dall' fault cache linker prelink advisory discrectionary lock fl
% LocalWords:  flock shared exclusive operation dup inode linked NFS cmd ENOLCK
% LocalWords:  EDEADLK whence SEEK CUR type pid GETLK SETLK SETLKW HP EACCESS
% LocalWords:  switch bsd lockf mandatory SVr sgid group root mount mand TRUNC
% LocalWords:  SVID UX Documentation sendfile dnotify inotify NdA ppoll fds add
% LocalWords:  init EMFILE FIONREAD ioctl watch char pathname uint mask ENOSPC
% LocalWords:  CLOSE NOWRITE MOVE MOVED FROM TO rm wd event page ctl acquired
% LocalWords:  attribute Universe epoll Solaris kqueue level triggered Jonathan
% LocalWords:  Lemon BSDCON edge Libenzi kevent backporting epfd EEXIST ENOENT
% LocalWords:  MOD wait EPOLLIN EPOLLOUT EPOLLRDHUP SOCK EPOLLPRI EPOLLERR one
% LocalWords:  EPOLLHUP EPOLLET EPOLLONESHOT shot maxevents ctlv ALL DONT HPUX
% LocalWords:  FOLLOW ONESHOT ONLYDIR FreeBSD EIO caching sysctl instances name
% LocalWords:  watches IGNORED ISDIR OVERFLOW overflow UNMOUNT queued cookie ls
% LocalWords:  NUL sizeof casting printevent nread limits sysconf SC wrapper Di
% LocalWords:  splice result argument DMA controller zerocopy Linus Larry Voy
% LocalWords:  Jens Anxboe vmsplice seek ESPIPE GIFT TCP CORK MSG splicecp nr
% LocalWords:  nwrite segs patch readahead posix fadvise TC advice FADV NORMAL
% LocalWords:  SEQUENTIAL NOREUSE WILLNEED DONTNEED streaming fallocate EFBIG
% LocalWords:  POLLRDHUP half close pwait Gb madvise MADV ahead REMOVE tmpfs it
% LocalWords:  DONTFORK DOFORK shmfs preadv pwritev syscall linux loff head XFS
% LocalWords:  MERGEABLE EOVERFLOW prealloca hole FALLOC KEEP stat fstat union
% LocalWords:  conditions sigwait CLOEXEC signalfd sizemask SIGKILL SIGSTOP ssi
% LocalWords:  sigwaitinfo FifoReporter Windows ptr sigqueue named timerfd TFD
% LocalWords:  clockid CLOCK MONOTONIC REALTIME itimerspec interval Resource
% LocalWords:  ABSTIME gettime temporarily unavailable SIGINT SIGQUIT SIGTERM
% LocalWords:  sigfd fifofd break siginf names starting echo Message from Got
% LocalWords:  message kill received means exit TLOCK ULOCK EPOLLWAKEUP


%%% Local Variables: 
%%% mode: latex
%%% TeX-master: "gapil"
%%% End: 


%% ipc.tex
%%
%% Copyright (C) 2000-2018 Simone Piccardi.  Permission is granted to
%% copy, distribute and/or modify this document under the terms of the GNU Free
%% Documentation License, Version 1.1 or any later version published by the
%% Free Software Foundation; with the Invariant Sections being "Un preambolo",
%% with no Front-Cover Texts, and with no Back-Cover Texts.  A copy of the
%% license is included in the section entitled "GNU Free Documentation
%% License".
%%

\chapter{L'intercomunicazione fra processi}
\label{cha:IPC}


Uno degli aspetti fondamentali della programmazione in un sistema unix-like è
la comunicazione fra processi. In questo capitolo affronteremo solo i
meccanismi più elementari che permettono di mettere in comunicazione processi
diversi, come quelli tradizionali che coinvolgono \textit{pipe} e
\textit{fifo} e i meccanismi di intercomunicazione di System V e quelli POSIX.

Tralasceremo invece tutte le problematiche relative alla comunicazione
attraverso la rete (e le relative interfacce) che saranno affrontate in
dettaglio in un secondo tempo.  Non affronteremo neanche meccanismi più
complessi ed evoluti come le RPC (\textit{Remote Procedure Calls}) che in
genere sono implementati da un ulteriore livello di librerie sopra i
meccanismi elementari.


\section{L'intercomunicazione fra processi tradizionale}
\label{sec:ipc_unix}

Il primo meccanismo di comunicazione fra processi introdotto nei sistemi Unix,
è quello delle cosiddette \textit{pipe}; esse costituiscono una delle
caratteristiche peculiari del sistema, in particolar modo dell'interfaccia a
linea di comando. In questa sezione descriveremo le sue basi, le funzioni che
ne gestiscono l'uso e le varie forme in cui si è evoluto.


\subsection{Le \textit{pipe} standard}
\label{sec:ipc_pipes}

Le \textit{pipe} nascono sostanzialmente con Unix, e sono il primo, e tuttora
uno dei più usati, meccanismi di comunicazione fra processi. Si tratta in
sostanza di una coppia di file descriptor connessi fra di loro in modo che
quanto scrive su di uno si può rileggere dall'altro.  Si viene così a
costituire un canale di comunicazione realizzato tramite i due file
descriptor, che costituisce appunto una sorta di \textsl{tubo} (che appunto il
significato del termine inglese \textit{pipe}) attraverso cui si possono far
passare i dati.

In pratica si tratta di un buffer circolare in memoria in cui il kernel
appoggia i dati immessi nel file descriptor su cui si scrive per farli poi
riemergere dal file descriptor da cui si legge. Si tenga ben presente che in
questo passaggio di dati non è previsto nessun tipo di accesso al disco e che
nonostante l'uso dei file descriptor le \textit{pipe} non han nulla a che fare
con i file di dati di cui si è parlato al cap.~\ref{cha:file_IO_interface}.

La funzione di sistema che permette di creare questa speciale coppia di file
descriptor associati ad una \textit{pipe} è appunto \funcd{pipe}, ed il suo
prototipo è:

\begin{funcproto}{
\fhead{unistd.h}
\fdecl{int pipe(int filedes[2])}
\fdesc{Crea la coppia di file descriptor di una \textit{pipe}.} 
}

{La funzione ritorna $0$ in caso di successo e $-1$ per un errore, nel qual
  caso \var{errno} assumerà uno dei valori: 
  \begin{errlist}
  \item[\errcode{EFAULT}] \param{filedes} non è un indirizzo valido.
  \end{errlist}
  ed inoltre \errval{EMFILE} e \errval{ENFILE} nel loro significato generico.}
\end{funcproto}

La funzione restituisce la coppia di file descriptor nel vettore
\param{filedes}, il primo è aperto in lettura ed il secondo in scrittura. Come
accennato concetto di funzionamento di una \textit{pipe} è semplice: quello
che si scrive nel file descriptor aperto in scrittura viene ripresentato tale
e quale nel file descriptor aperto in lettura. 

I file descriptor infatti non sono connessi a nessun file reale, ma, come
accennato, ad un buffer nel kernel la cui dimensione è specificata dal
parametro di sistema \const{PIPE\_BUF}, (vedi
sez.~\ref{sec:sys_file_limits}). Lo schema di funzionamento di una
\textit{pipe} è illustrato in fig.~\ref{fig:ipc_pipe_singular}, in cui sono
indicati i due capi della \textit{pipe}, associati a ciascun file descriptor,
con le frecce che indicano la direzione del flusso dei dati.

\begin{figure}[!htb]
  \centering
  \includegraphics[height=5cm]{img/pipe}
  \caption{Schema della struttura di una \textit{pipe}.}
  \label{fig:ipc_pipe_singular}
\end{figure}

Della funzione di sistema esiste una seconda versione, \funcd{pipe2},
introdotta con il kernel 2.6.27 e la \acr{glibc} 2.9 e specifica di Linux
(utilizzabile solo definendo la macro \macro{\_GNU\_SOURCE}), che consente di
impostare atomicamente le caratteristiche dei file descriptor restituiti, il
suo prototipo è:

\begin{funcproto}{
\fhead{unistd.h}
\fhead{fcntl.h}
\fdecl{int pipe2(int pipefd[2], int flags)}
\fdesc{Crea la coppia di file descriptor di una \textit{pipe}.} 
}

{La funzione ritorna $0$ in caso di successo e $-1$ per un errore, nel qual
  caso \var{errno} assumerà uno dei valori: 
  \begin{errlist}
  \item[\errcode{EINVAL}] il valore di \param{flags} non valido.
  \end{errlist}
  e gli altri già visti per \func{pipe} con lo stesso significato.}
\end{funcproto}

Utilizzando un valore nullo per \param{flags} la funzione è identica a
\func{pipe}, si può però passare come valore l'OR aritmetico di uno qualunque
fra \const{O\_NONBLOCK} o \const{O\_CLOEXEC} che hanno l'effetto di impostare
su entrambi i file descriptor restituiti dalla funzione i relativi flag, già
descritti per \func{open} in tab.~\ref{tab:open_operation_flag}, che attivano
rispettivamente la modalità di accesso \textsl{non-bloccante} ed il
\textit{close-on-exec}.

Chiaramente creare una \textit{pipe} all'interno di un singolo processo non
serve a niente; se però ricordiamo quanto esposto in
sez.~\ref{sec:file_shared_access} riguardo al comportamento dei file
descriptor nei processi figli, è immediato capire come una \textit{pipe} possa
diventare un meccanismo di intercomunicazione. Un processo figlio infatti
condivide gli stessi file descriptor del padre, compresi quelli associati ad
una \textit{pipe} (secondo la situazione illustrata in
fig.~\ref{fig:ipc_pipe_fork}). In questo modo se uno dei processi scrive su un
capo della \textit{pipe}, l'altro può leggere.

\begin{figure}[!htb]
  \centering
  \includegraphics[height=5cm]{img/pipefork}
  \caption{Schema dei collegamenti ad una \textit{pipe}, condivisi fra
    processo padre e figlio dopo l'esecuzione \func{fork}.}
  \label{fig:ipc_pipe_fork}
\end{figure}

Tutto ciò ci mostra come sia immediato realizzare un meccanismo di
comunicazione fra processi attraverso una \textit{pipe}, utilizzando le
proprietà ordinarie dei file, ma ci mostra anche qual è il principale limite
nell'uso delle \textit{pipe}.\footnote{Stevens in \cite{APUE} riporta come
  limite anche il fatto che la comunicazione è unidirezionale, ma in realtà
  questo è un limite superabile usando una coppia di \textit{pipe}, anche se
  al costo di una maggiore complessità di gestione.}  È necessario infatti che
i processi possano condividere i file descriptor della \textit{pipe}, e per
questo essi devono comunque essere \textsl{parenti} (dall'inglese
\textit{siblings}), cioè o derivare da uno stesso processo padre in cui è
avvenuta la creazione della \textit{pipe}, o, più comunemente, essere nella
relazione padre/figlio.

A differenza di quanto avviene con i file normali, la lettura da una
\textit{pipe} può essere bloccante (qualora non siano presenti dati), inoltre
se si legge da una \textit{pipe} il cui capo in scrittura è stato chiuso, si
avrà la ricezione di un EOF (vale a dire che la funzione \func{read} ritornerà
restituendo 0).  Se invece si esegue una scrittura su una \textit{pipe} il cui
capo in lettura non è aperto il processo riceverà il segnale \signal{SIGPIPE},
e la funzione di scrittura restituirà un errore di \errcode{EPIPE} (al ritorno
del gestore, o qualora il segnale sia ignorato o bloccato).

La dimensione del buffer della \textit{pipe} (\const{PIPE\_BUF}) ci dà inoltre
un'altra importante informazione riguardo il comportamento delle operazioni di
lettura e scrittura su di una \textit{pipe}; esse infatti sono atomiche
fintanto che la quantità di dati da scrivere non supera questa
dimensione. Qualora ad esempio si effettui una scrittura di una quantità di
dati superiore l'operazione verrà effettuata in più riprese, consentendo
l'intromissione di scritture effettuate da altri processi.

La dimensione originale del buffer era di 4096 byte (uguale ad una pagina di
memoria) fino al kernel 2.6.11, ed è stata portata in seguito a 64kb; ma a
partire dal kernel 2.6.35 è stata resa disponibile l'operazione di controllo
\const{F\_SETPIPE\_SZ} (vedi sez.~\ref{sec:ipc_pipes}) che consente di
modificarne la dimensione.



\subsection{Un esempio dell'uso delle \textit{pipe}}
\label{sec:ipc_pipe_use}

Per capire meglio il funzionamento delle \textit{pipe} faremo un esempio di
quello che è il loro uso più comune, analogo a quello effettuato della shell,
e che consiste nell'inviare l'output di un processo (lo standard output)
sull'input di un altro. Realizzeremo il programma di esempio nella forma di un
\textit{CGI}\footnote{quella dei CGI (\textit{Common Gateway Interface}) è una
  interfaccia che consente ad un server web di eseguire un programma il cui
  output (che deve essere opportunamente formattato seguendo le specifiche
  dell'interfaccia) può essere presentato come risposta ad una richiesta HTTP
  al posto del contenuto di un file, e che ha costituito probabilmente la
  prima modalità con cui sono state create pagine HTML dinamiche.}  che genera
una immagine JPEG di un codice a barre, specificato come argomento in
ingresso.

Un programma che deve essere eseguito come \textit{CGI} deve rispondere a
delle caratteristiche specifiche, esso infatti non viene lanciato da una
shell, ma dallo stesso web server, alla richiesta di una specifica URL, che di
solito ha la forma:
\begin{Example}
http://www.sito.it/cgi-bin/programma?argomento
\end{Example}
ed il risultato dell'elaborazione deve essere presentato (con una intestazione
che ne descrive il \textit{mime-type}) sullo \textit{standard output}, in modo
che il server web possa reinviarlo al browser che ha effettuato la richiesta,
che in questo modo è in grado di visualizzarlo opportunamente.

\begin{figure}[!htb]
  \centering
  \includegraphics[height=5cm]{img/pipeuse}
  \caption{Schema dell'uso di una \textit{pipe} come mezzo di comunicazione fra
    due processi attraverso l'esecuzione una \func{fork} e la chiusura dei
    capi non utilizzati.}
  \label{fig:ipc_pipe_use}
\end{figure}

Per realizzare quanto voluto useremo in sequenza i programmi \cmd{barcode} e
\cmd{gs}, il primo infatti è in grado di generare immagini PostScript di
codici a barre corrispondenti ad una qualunque stringa, mentre il secondo
serve per poter effettuare la conversione della stessa immagine in formato
JPEG. Usando una \textit{pipe} potremo inviare l'output del primo sull'input
del secondo, secondo lo schema mostrato in fig.~\ref{fig:ipc_pipe_use}, in cui
la direzione del flusso dei dati è data dalle frecce continue.

Si potrebbe obiettare che sarebbe molto più semplice salvare il risultato
intermedio su un file temporaneo. Questo però non tiene conto del fatto che un
\textit{CGI} può essere eseguito più volte in contemporanea, e si avrebbe una
evidente \textit{race condition} in caso di accesso simultaneo a detto file da
istanze diverse. Il problema potrebbe essere superato utilizzando un sempre
diverso per il file temporaneo, che verrebbe creato all'avvio di ogni istanza,
utilizzato dai sottoprocessi, e cancellato alla fine della sua esecuzione; ma
a questo punto le cose non sarebbero più tanto semplici.  L'uso di una
\textit{pipe} invece permette di risolvere il problema in maniera semplice ed
elegante, oltre ad essere molto più efficiente, dato che non si deve scrivere
su disco.

Il programma ci servirà anche come esempio dell'uso delle funzioni di
duplicazione dei file descriptor che abbiamo trattato in
sez.~\ref{sec:file_dup}, in particolare di \func{dup2}. È attraverso queste
funzioni infatti che è possibile dirottare gli stream standard dei processi
(che abbiamo visto in tab.~\ref{tab:file_std_files} e
sez.~\ref{sec:file_stream}) sulla \textit{pipe}. In
fig.~\ref{fig:ipc_barcodepage_code} abbiamo riportato il corpo del programma,
il cui codice completo è disponibile nel file \file{BarCodePage.c} che si
trova nella directory dei sorgenti.

\begin{figure}[!htb]
  \footnotesize \centering
  \begin{minipage}[c]{\codesamplewidth}
    \includecodesample{listati/BarCodePage.c}
  \end{minipage} 
  \normalsize 
  \caption{Sezione principale del codice del \textit{CGI} 
    \file{BarCodePage.c}.}
  \label{fig:ipc_barcodepage_code}
\end{figure}

La prima operazione del programma (\texttt{\small 4-12}) è quella di creare
le due \textit{pipe} che serviranno per la comunicazione fra i due comandi
utilizzati per produrre il codice a barre; si ha cura di controllare la
riuscita della chiamata, inviando in caso di errore un messaggio invece
dell'immagine richiesta. La funzione \func{WriteMess} non è riportata in
fig.~\ref{fig:ipc_barcodepage_code}; essa si incarica semplicemente di
formattare l'uscita alla maniera dei CGI, aggiungendo l'opportuno
\textit{mime-type}, e formattando il messaggio in HTML, in modo che
quest'ultimo possa essere visualizzato correttamente da un browser.

Una volta create le \textit{pipe}, il programma può creare (\texttt{\small
  13-17}) il primo processo figlio, che si incaricherà (\texttt{\small
  19-25}) di eseguire \cmd{barcode}. Quest'ultimo legge dallo standard input
una stringa di caratteri, la converte nell'immagine PostScript del codice a
barre ad essa corrispondente, e poi scrive il risultato direttamente sullo
standard output.

Per poter utilizzare queste caratteristiche prima di eseguire \cmd{barcode} si
chiude (\texttt{\small 20}) il capo aperto in scrittura della prima
\textit{pipe}, e se ne collega (\texttt{\small 21}) il capo in lettura allo
\textit{standard input} usando \func{dup2}. Si ricordi che invocando
\func{dup2} il secondo file, qualora risulti aperto, viene, come nel caso
corrente, chiuso prima di effettuare la duplicazione. Allo stesso modo, dato
che \cmd{barcode} scrive l'immagine PostScript del codice a barre sullo
standard output, per poter effettuare una ulteriore redirezione il capo in
lettura della seconda \textit{pipe} viene chiuso (\texttt{\small 22}) mentre
il capo in scrittura viene collegato allo standard output (\texttt{\small
  23}).

In questo modo all'esecuzione (\texttt{\small 25}) di \cmd{barcode} (cui si
passa in \var{size} la dimensione della pagina per l'immagine) quest'ultimo
leggerà dalla prima \textit{pipe} la stringa da codificare che gli sarà
inviata dal padre, e scriverà l'immagine PostScript del codice a barre sulla
seconda.

Al contempo una volta lanciato il primo figlio, il processo padre prima chiude
(\texttt{\small 26}) il capo inutilizzato della prima \textit{pipe} (quello in
ingresso) e poi scrive (\texttt{\small 27}) la stringa da convertire sul capo
in uscita, così che \cmd{barcode} possa riceverla dallo \textit{standard
  input}. A questo punto l'uso della prima \textit{pipe} da parte del padre è
finito ed essa può essere definitivamente chiusa (\texttt{\small 28}), si
attende poi (\texttt{\small 29}) che l'esecuzione di \cmd{barcode} sia
completata.

Alla conclusione della sua esecuzione \cmd{barcode} avrà inviato l'immagine
PostScript del codice a barre sul capo in scrittura della seconda
\textit{pipe}; a questo punto si può eseguire la seconda conversione, da PS a
JPEG, usando il programma \cmd{gs}. Per questo si crea (\texttt{\small
  30-34}) un secondo processo figlio, che poi (\texttt{\small 35-42})
eseguirà questo programma leggendo l'immagine PostScript creata da
\cmd{barcode} dallo \textit{standard input}, per convertirla in JPEG.

Per fare tutto ciò anzitutto si chiude (\texttt{\small 37}) il capo in
scrittura della seconda \textit{pipe}, e se ne collega (\texttt{\small 38}) il
capo in lettura allo \textit{standard input}. Per poter formattare l'output
del programma in maniera utilizzabile da un browser, si provvede anche
\texttt{\small 40}) alla scrittura dell'apposita stringa di identificazione
del \textit{mime-type} in testa allo \textit{standard output}. A questo punto
si può invocare \texttt{\small 41}) \cmd{gs}, provvedendo le opportune opzioni
del comando che consentono di leggere il file da convertire dallo
\textit{standard input} e di inviare la conversione sullo \textit{standard
  output}.

Per completare le operazioni il processo padre chiude (\texttt{\small 44}) il
capo in scrittura della seconda \textit{pipe}, e attende la conclusione del
figlio (\texttt{\small 45}); a questo punto può (\texttt{\small 46})
uscire. Si tenga conto che l'operazione di chiudere il capo in scrittura della
seconda \textit{pipe} è necessaria, infatti, se non venisse chiusa, \cmd{gs},
che legge il suo \textit{standard input} da detta \textit{pipe}, resterebbe
bloccato in attesa di ulteriori dati in ingresso (l'unico modo che un
programma ha per sapere che i dati in ingresso sono terminati è rilevare che
lo \textit{standard input} è stato chiuso), e la \func{wait} non ritornerebbe.


\subsection{Le funzioni \func{popen} e \func{pclose}}
\label{sec:ipc_popen}

Come si è visto la modalità più comune di utilizzo di una \textit{pipe} è
quella di utilizzarla per fare da tramite fra output ed input di due programmi
invocati in sequenza; per questo motivo lo standard POSIX.2 ha introdotto due
funzioni che permettono di sintetizzare queste operazioni. La prima di esse si
chiama \funcd{popen} ed il suo prototipo è:


\begin{funcproto}{
\fhead{stdio.h}
\fdecl{FILE *popen(const char *command, const char *type)}
\fdesc{Esegue un programma dirottando l'uscita su una \textit{pipe}.} 
}

{La funzione ritorna l'indirizzo dello stream associato alla \textit{pipe} in
  caso di successo e \val{NULL} per un errore, nel qual caso \var{errno} potrà
  assumere i valori relativi alle sottostanti invocazioni di \func{pipe} e
  \func{fork} o \errcode{EINVAL} se \param{type} non è valido.}
\end{funcproto}

La funzione crea una \textit{pipe}, esegue una \func{fork} creando un nuovo
processo nel quale invoca il programma \param{command} attraverso la shell (in
sostanza esegue \file{/bin/sh} con il flag \code{-c}).
L'argomento \param{type} deve essere una delle due stringhe \verb|"w"| o
\verb|"r"|, per richiedere che la \textit{pipe} restituita come valore di
ritorno sia collegata allo \textit{standard input} o allo \textit{standard
  output} del comando invocato.

La funzione restituisce il puntatore ad uno stream associato alla
\textit{pipe} creata, che sarà aperto in sola lettura (e quindi associato allo
\textit{standard output} del programma indicato) in caso si sia indicato
\code{r}, o in sola scrittura (e quindi associato allo \textit{standard
  input}) in caso di \code{w}. A partire dalla versione 2.9 della \acr{glibc}
(questa è una estensione specifica di Linux) all'argomento \param{type} può
essere aggiunta la lettera ``\texttt{e}'' per impostare automaticamente il
flag di \textit{close-on-exec} sul file descriptor sottostante (si ricordi
quanto spiegato in sez.~\ref{sec:file_open_close}).

Lo \textit{stream} restituito da \func{popen} è identico a tutti gli effetti
ai \textit{file stream} visti in sez.~\ref{sec:files_std_interface}, anche se
è collegato ad una \textit{pipe} e non ad un file, e viene sempre aperto in
modalità \textit{fully-buffered} (vedi sez.~\ref{sec:file_buffering}); l'unica
differenza con gli usuali stream è che dovrà essere chiuso dalla seconda delle
due nuove funzioni, \funcd{pclose}, il cui prototipo è:

\begin{funcproto}{
\fhead{stdio.h}
\fdecl{int pclose(FILE *stream)}
\fdesc{Chiude una \textit{pipe} creata con \func{popen}.} 
}

{La funzione ritorna lo stato del processo creato da \func{popen} in caso di
  successo e $-1$ per un errore, nel qual caso \var{errno} assumerà i valori
  derivanti dalle sottostanti funzioni \func{fclose} e \func{wait4}.}
\end{funcproto}

La funzione chiude il file \param{stream} associato ad una \textit{pipe}
creato da una precedente \func{popen}, ed oltre alla chiusura dello stream si
incarica anche di attendere (tramite \func{wait4}) la conclusione del processo
creato dalla precedente \func{popen}. Se lo stato di uscita non può essere
letto la funzione restituirà per \var{errno} un errore di \errval{ECHILD}.

Per illustrare l'uso di queste due funzioni riprendiamo il problema
precedente: il programma mostrato in fig.~\ref{fig:ipc_barcodepage_code} per
quanto funzionante, è volutamente codificato in maniera piuttosto complessa,
inoltre doveva scontare un problema di \cmd{gs} che non era in grado di
riconoscere correttamente l'Encapsulated PostScript,\footnote{si fa
  riferimento alla versione di GNU Ghostscript 6.53 (2002-02-13), usata quando
  l'esempio venne scritto per la prima volta.} per cui si era utilizzato il
PostScript semplice, generando una pagina intera invece che una immagine delle
dimensioni corrispondenti al codice a barre.

Se si vuole generare una immagine di dimensioni appropriate si deve usare un
approccio diverso. Una possibilità sarebbe quella di ricorrere ad ulteriore
programma, \cmd{epstopsf}, per convertire in PDF un file EPS (che può essere
generato da \cmd{barcode} utilizzando l'opzione \cmd{-E}).  Utilizzando un PDF
al posto di un EPS \cmd{gs} esegue la conversione rispettando le dimensioni
originarie del codice a barre e produce un JPEG di dimensioni corrette.

Questo approccio però non può funzionare per via di una delle caratteristiche
principali delle \textit{pipe}. Per poter effettuare la conversione di un PDF
infatti è necessario, per la struttura del formato, potersi spostare (con
\func{lseek}) all'interno del file da convertire. Se si esegue la conversione
con \cmd{gs} su un file regolare non ci sono problemi, una \textit{pipe} però
è rigidamente sequenziale, e l'uso di \func{lseek} su di essa fallisce sempre
con un errore di \errcode{ESPIPE}, rendendo impossibile la conversione.
Questo ci dice che in generale la concatenazione di vari programmi funzionerà
soltanto quando tutti prevedono una lettura sequenziale del loro input.

Per questo motivo si è dovuto utilizzare un procedimento diverso, eseguendo
prima la conversione (sempre con \cmd{gs}) del PS in un altro formato
intermedio, il PPM,\footnote{il \textit{Portable PixMap file format} è un
  formato usato spesso come formato intermedio per effettuare conversioni, è
  infatti molto facile da manipolare, dato che usa caratteri ASCII per
  memorizzare le immagini, anche se per questo è estremamente inefficiente
  come formato di archiviazione.}  dal quale poi si può ottenere un'immagine
di dimensioni corrette attraverso vari programmi di manipolazione
(\cmd{pnmcrop}, \cmd{pnmmargin}) che può essere infine trasformata in PNG (con
\cmd{pnm2png}).

In questo caso però occorre eseguire in sequenza ben quattro comandi diversi,
inviando l'uscita di ciascuno all'ingresso del successivo, per poi ottenere il
risultato finale sullo \textit{standard output}: un caso classico di
utilizzazione delle \textit{pipe}, in cui l'uso di \func{popen} e \func{pclose}
permette di semplificare notevolmente la stesura del codice.

Nel nostro caso, dato che ciascun processo deve scrivere la sua uscita sullo
\textit{standard input} del successivo, occorrerà usare \func{popen} aprendo
la \textit{pipe} in scrittura. Il codice del nuovo programma è riportato in
fig.~\ref{fig:ipc_barcode_code}.  Come si può notare l'ordine di invocazione
dei programmi è l'inverso di quello in cui ci si aspetta che vengano
effettivamente eseguiti. Questo non comporta nessun problema dato che la
lettura su una \textit{pipe} è bloccante, per cui un processo, anche se
lanciato per primo, se non ottiene i dati che gli servono si bloccherà in
attesa sullo \textit{standard input} finché non otterrà il risultato
dell'elaborazione del processo che li deve creare, che pur essendo logicamente
precedente, viene lanciato dopo di lui.

\begin{figure}[!htb]
  \footnotesize \centering
  \begin{minipage}[c]{\codesamplewidth}
    \includecodesample{listati/BarCode.c}
  \end{minipage} 
  \normalsize 
  \caption{Codice completo del \textit{CGI} \file{BarCode.c}.}
  \label{fig:ipc_barcode_code}
\end{figure}

Nel nostro caso il primo passo (\texttt{\small 14}) è scrivere il
\textit{mime-type} sullo \textit{standard output}; a questo punto il processo
padre non necessita più di eseguire ulteriori operazioni sullo
\textit{standard output} e può tranquillamente provvedere alla redirezione.

Dato che i vari programmi devono essere lanciati in successione, si è
approntato un ciclo (\texttt{\small 15-19}) che esegue le operazioni in
sequenza: prima crea una \textit{pipe} (\texttt{\small 17}) per la scrittura
eseguendo il programma con \func{popen}, in modo che essa sia collegata allo
\textit{standard input}, e poi redirige (\texttt{\small 18}) lo
\textit{standard output} su detta \textit{pipe}.

In questo modo il primo processo ad essere invocato (che è l'ultimo della
catena) scriverà ancora sullo \textit{standard output} del processo padre, ma
i successivi, a causa di questa redirezione, scriveranno sulla \textit{pipe}
associata allo \textit{standard input} del processo invocato nel ciclo
precedente.

Alla fine tutto quello che resta da fare è lanciare (\texttt{\small 21}) il
primo processo della catena, che nel caso è \cmd{barcode}, e scrivere
(\texttt{\small 23}) la stringa del codice a barre sulla \textit{pipe}, che è
collegata al suo \textit{standard input}, infine si può eseguire
(\texttt{\small 24-27}) un ciclo che chiuda con \func{pclose}, nell'ordine
inverso rispetto a quello in cui le si sono create, tutte le \textit{pipe}
create in precedenza.


\subsection{Le \textit{pipe} con nome, o \textit{fifo}}
\label{sec:ipc_named_pipe}

Come accennato in sez.~\ref{sec:ipc_pipes} il problema delle \textit{pipe} è
che esse possono essere utilizzate solo da processi con un progenitore comune
o nella relazione padre/figlio. Per superare questo problema lo standard
POSIX.1 ha introdotto le \textit{fifo}, che hanno le stesse caratteristiche
delle \textit{pipe}, ma che invece di essere visibili solo attraverso un file
descriptor creato all'interno di un processo da una \textit{system call}
apposita, costituiscono un oggetto che risiede sul filesystem (si rammenti
quanto detto in sez.~\ref{sec:file_file_types}) che può essere aperto come un
qualunque file, così che i processi le possono usare senza dovere per forza
essere in una relazione di \textsl{parentela}.

Utilizzando una \textit{fifo} tutti i dati passeranno, come per le
\textit{pipe}, attraverso un buffer nel kernel, senza transitare dal
filesystem. Il fatto che siano associate ad un \textit{inode} presente sul
filesystem serve infatti solo a fornire un punto di accesso per i processi,
che permetta a questi ultimi di accedere alla stessa \textit{fifo} senza avere
nessuna relazione, con una semplice \func{open}. Il comportamento delle
funzioni di lettura e scrittura è identico a quello illustrato per le
\textit{pipe} in sez.~\ref{sec:ipc_pipes}.

Abbiamo già trattato in sez.~\ref{sec:file_mknod} le funzioni \func{mknod} e
\func{mkfifo} che permettono di creare una \textit{fifo}. Per utilizzarne una
un processo non avrà che da aprire il relativo file speciale o in lettura o
scrittura; nel primo caso il processo sarà collegato al capo di uscita della
\textit{fifo}, e dovrà leggere, nel secondo al capo di ingresso, e dovrà
scrivere.

Il kernel alloca un singolo buffer per ciascuna \textit{fifo} che sia stata
aperta, e questa potrà essere acceduta contemporaneamente da più processi, sia
in lettura che in scrittura. Dato che per funzionare deve essere aperta in
entrambe le direzioni, per una \textit{fifo} la funzione \func{open} di norma
si blocca se viene eseguita quando l'altro capo non è aperto.

Le \textit{fifo} però possono essere anche aperte in modalità
\textsl{non-bloccante}, nel qual caso l'apertura del capo in lettura avrà
successo solo quando anche l'altro capo è aperto, mentre l'apertura del capo
in scrittura restituirà l'errore di \errcode{ENXIO} fintanto che non verrà
aperto il capo in lettura.

In Linux è possibile aprire le \textit{fifo} anche in lettura/scrittura (lo
standard POSIX lascia indefinito il comportamento in questo caso) operazione
che avrà sempre successo immediato qualunque sia la modalità di apertura,
bloccante e non bloccante.  Questo può essere utilizzato per aprire comunque
una \textit{fifo} in scrittura anche se non ci sono ancora processi il
lettura. Infine è possibile anche usare la \textit{fifo} all'interno di un
solo processo, nel qual caso però occorre stare molto attenti alla possibili
situazioni di stallo: se si cerca di leggere da una \textit{fifo} che non
contiene dati si avrà infatti un \textit{deadlock} immediato, dato che il
processo si blocca e quindi non potrà mai eseguire le funzioni di scrittura.

Per la loro caratteristica di essere accessibili attraverso il filesystem, è
piuttosto frequente l'utilizzo di una \textit{fifo} come canale di
comunicazione nelle situazioni un processo deve ricevere informazioni da
altri. In questo caso è fondamentale che le operazioni di scrittura siano
atomiche; per questo si deve sempre tenere presente che questo è vero soltanto
fintanto che non si supera il limite delle dimensioni di \const{PIPE\_BUF} (si
ricordi quanto detto in sez.~\ref{sec:ipc_pipes}).

A parte il caso precedente, che resta probabilmente il più comune, Stevens
riporta in \cite{APUE} altre due casistiche principali per l'uso delle
\textit{fifo}:
\begin{itemize*}
\item Da parte dei comandi di shell, per evitare la creazione di file
  temporanei quando si devono inviare i dati di uscita di un processo
  sull'input di parecchi altri (attraverso l'uso del comando \cmd{tee}).  
\item Come canale di comunicazione fra un client ed un
  server (il modello \textit{client-server} è illustrato in
  sez.~\ref{sec:net_cliserv}).
\end{itemize*}

Nel primo caso quello che si fa è creare tante \textit{fifo} da usare come
\textit{standard input} quanti sono i processi a cui i vogliono inviare i
dati; questi ultimi saranno stati posti in esecuzione ridirigendo lo
\textit{standard input} dalle \textit{fifo}, si potrà poi eseguire il processo
che fornisce l'output replicando quest'ultimo, con il comando \cmd{tee}, sulle
varie \textit{fifo}.

Il secondo caso è relativamente semplice qualora si debba comunicare con un
processo alla volta, nel qual caso basta usare due \textit{fifo}, una per
leggere ed una per scrivere. Le cose diventano invece molto più complesse
quando si vuole effettuare una comunicazione fra un server ed un numero
imprecisato di client. Se il primo infatti può ricevere le richieste
attraverso una \textit{fifo} ``\textsl{nota}'', per le risposte non si può
fare altrettanto, dato che, per la struttura sequenziale delle \textit{fifo},
i client dovrebbero sapere prima di leggerli quando i dati inviati sono
destinati a loro.

Per risolvere questo problema, si può usare un'architettura come quella
illustrata in fig.~\ref{fig:ipc_fifo_server_arch} in cui i client
inviano le richieste al server su una \textit{fifo} nota mentre le
risposte vengono reinviate dal server a ciascuno di essi su una
\textit{fifo} temporanea creata per l'occasione.

\begin{figure}[!htb]
  \centering
  \includegraphics[height=9cm]{img/fifoserver}
  \caption{Schema dell'utilizzo delle \textit{fifo} nella realizzazione di una
    architettura di comunicazione client/server.}
  \label{fig:ipc_fifo_server_arch}
\end{figure}

Come esempio di uso questa architettura e dell'uso delle \textit{fifo},
abbiamo scritto un server di \textit{fortunes}, che restituisce, alle
richieste di un client, un detto a caso estratto da un insieme di frasi. Sia
il numero delle frasi dell'insieme, che i file da cui esse vengono lette
all'avvio, sono impostabili da riga di comando. Il corpo principale del
server è riportato in fig.~\ref{fig:ipc_fifo_server}, dove si è
tralasciata la parte che tratta la gestione delle opzioni a riga di comando,
che effettua l'impostazione delle variabili \var{fortunefilename}, che indica
il file da cui leggere le frasi, ed \var{n}, che indica il numero di frasi
tenute in memoria, ad un valore diverso da quelli preimpostati. Il codice
completo è nel file \file{FortuneServer.c}.

\begin{figure}[!htbp]
  \footnotesize \centering
  \begin{minipage}[c]{\codesamplewidth}
    \includecodesample{listati/FortuneServer.c}
  \end{minipage} 
  \normalsize 
  \caption{Sezione principale del codice del server di \textit{fortunes}
    basato sulle \textit{fifo}.}
  \label{fig:ipc_fifo_server}
\end{figure}

Il server richiede (\texttt{\small 12}) che sia stata impostata una dimensione
dell'insieme delle frasi non nulla, dato che l'inizializzazione del vettore
\var{fortune} avviene solo quando questa dimensione viene specificata, la
presenza di un valore nullo provoca l'uscita dal programma attraverso la
funzione (non riportata) che ne stampa le modalità d'uso.  Dopo di che
installa (\texttt{\small 13-15}) la funzione che gestisce i segnali di
interruzione (anche questa non è riportata in fig.~\ref{fig:ipc_fifo_server})
che si limita a rimuovere dal filesystem la \textit{fifo} usata dal server per
comunicare.

Terminata l'inizializzazione (\texttt{\small 16}) si effettua la chiamata alla
funzione \code{FortuneParse} che legge dal file specificato in
\var{fortunefilename} le prime \var{n} frasi e le memorizza (allocando
dinamicamente la memoria necessaria) nel vettore di puntatori \var{fortune}.
Anche il codice della funzione non è riportato, in quanto non direttamente
attinente allo scopo dell'esempio.

Il passo successivo (\texttt{\small 17-22}) è quello di creare con
\func{mkfifo} la \textit{fifo} nota sulla quale il server ascolterà le
richieste, qualora si riscontri un errore il server uscirà (escludendo
ovviamente il caso in cui la funzione \func{mkfifo} fallisce per la precedente
esistenza della \textit{fifo}).

Una volta che si è certi che la \textit{fifo} di ascolto esiste la procedura
di inizializzazione è completata. A questo punto (\texttt{\small 23}) si può
chiamare la funzione \func{daemon} per far proseguire l'esecuzione del
programma in background come demone.  Si può quindi procedere (\texttt{\small
  24-33}) alla apertura della \textit{fifo}: si noti che questo viene fatto
due volte, prima in lettura e poi in scrittura, per evitare di dover gestire
all'interno del ciclo principale il caso in cui il server è in ascolto ma non
ci sono client che effettuano richieste.  Si ricordi infatti che quando una
\textit{fifo} è aperta solo dal capo in lettura, l'esecuzione di \func{read}
ritorna con zero byte (si ha cioè una condizione di \textit{end-of-file}).

Nel nostro caso la prima apertura si bloccherà fintanto che un qualunque
client non apre a sua volta la \textit{fifo} nota in scrittura per effettuare
la sua richiesta. Pertanto all'inizio non ci sono problemi, il client però,
una volta ricevuta la risposta, uscirà, chiudendo tutti i file aperti,
compresa la \textit{fifo}.  A questo punto il server resta (se non ci sono
altri client che stanno effettuando richieste) con la \textit{fifo} chiusa sul
lato in lettura, ed in questo stato la funzione \func{read} non si bloccherà
in attesa di dati in ingresso, ma ritornerà in continuazione, restituendo una
condizione di \textit{end-of-file}.

Si è usata questa tecnica per compatibilità, Linux infatti supporta l'apertura
delle \textit{fifo} in lettura/scrittura, per cui si sarebbe potuto effettuare
una singola apertura con \const{O\_RDWR}; la doppia apertura comunque ha il
vantaggio che non si può scrivere per errore sul capo aperto in sola lettura.

Per questo motivo, dopo aver eseguito l'apertura in lettura (\texttt{\small
  24-28}),\footnote{di solito si effettua l'apertura del capo in lettura di
  una \textit{fifo} in modalità non bloccante, per evitare il rischio di uno
  stallo: se infatti nessuno apre la \textit{fifo} in scrittura il processo
  non ritornerà mai dalla \func{open}. Nel nostro caso questo rischio non
  esiste, mentre è necessario potersi bloccare in lettura in attesa di una
  richiesta.}  si esegue una seconda apertura in scrittura (\texttt{\small
  29-32}), scartando il relativo file descriptor, che non sarà mai usato, in
questo modo però la \textit{fifo} resta comunque aperta anche in scrittura,
cosicché le successive chiamate a \func{read} possono bloccarsi.

A questo punto si può entrare nel ciclo principale del programma che fornisce
le risposte ai client (\texttt{\small 34-50}); questo viene eseguito
indefinitamente (l'uscita del server viene effettuata inviando un segnale, in
modo da passare attraverso la funzione di chiusura che cancella la
\textit{fifo}).

Il server è progettato per accettare come richieste dai client delle stringhe
che contengono il nome della \textit{fifo} sulla quale deve essere inviata la
risposta.  Per cui prima (\texttt{\small 35-39}) si esegue la lettura dalla
stringa di richiesta dalla \textit{fifo} nota (che a questo punto si bloccherà
tutte le volte che non ci sono richieste). Dopo di che, una volta terminata la
stringa (\texttt{\small 40}) e selezionato (\texttt{\small 41}) un numero
casuale per ricavare la frase da inviare, si procederà (\texttt{\small
  42-46}) all'apertura della \textit{fifo} per la risposta, che poi 
(\texttt{\small 47-48}) vi sarà scritta. Infine (\texttt{\small 49}) si
chiude la \textit{fifo} di risposta che non serve più.

Il codice del client è invece riportato in fig.~\ref{fig:ipc_fifo_client},
anche in questo caso si è omessa la gestione delle opzioni e la funzione che
stampa a video le informazioni di utilizzo ed esce, riportando solo la sezione
principale del programma e le definizioni delle variabili. Il codice completo
è nel file \file{FortuneClient.c} dei sorgenti allegati.

\begin{figure}[!htbp]
  \footnotesize \centering
  \begin{minipage}[c]{\codesamplewidth}
    \includecodesample{listati/FortuneClient.c}
  \end{minipage} 
  \normalsize 
  \caption{Sezione principale del codice del client di \textit{fortunes}
    basato sulle \textit{fifo}.}
  \label{fig:ipc_fifo_client}
\end{figure}

La prima istruzione (\texttt{\small 12}) compone il nome della \textit{fifo}
che dovrà essere utilizzata per ricevere la risposta dal server.  Si usa il
\ids{PID} del processo per essere sicuri di avere un nome univoco; dopo di che
(\texttt{\small 13-18}) si procede alla creazione del relativo file, uscendo
in caso di errore (a meno che il file non sia già presente sul filesystem).

A questo punto il client può effettuare l'interrogazione del server, per
questo prima si apre la \textit{fifo} nota (\texttt{\small 19-23}), e poi ci
si scrive (\texttt{\small 24}) la stringa composta in precedenza, che contiene
il nome della \textit{fifo} da utilizzare per la risposta. Infine si richiude
la \textit{fifo} del server che a questo punto non serve più (\texttt{\small
  25}). 

Inoltrata la richiesta si può passare alla lettura della risposta; anzitutto
si apre (\texttt{\small 26-30}) la \textit{fifo} appena creata, da cui si
deve riceverla, dopo di che si effettua una lettura (\texttt{\small 31})
nell'apposito buffer; si è supposto, come è ragionevole, che le frasi inviate
dal server siano sempre di dimensioni inferiori a \const{PIPE\_BUF},
tralasciamo la gestione del caso in cui questo non è vero. Infine si stampa
(\texttt{\small 32}) a video la risposta, si chiude (\texttt{\small 33}) la
\textit{fifo} e si cancella (\texttt{\small 34}) il relativo file.  Si noti
come la \textit{fifo} per la risposta sia stata aperta solo dopo aver inviato
la richiesta, se non si fosse fatto così si avrebbe avuto uno stallo, in
quanto senza la richiesta, il server non avrebbe potuto aprirne il capo in
scrittura e l'apertura si sarebbe bloccata indefinitamente.

Verifichiamo allora il comportamento dei nostri programmi, in questo, come in
altri esempi precedenti, si fa uso delle varie funzioni di servizio, che sono
state raccolte nella libreria \file{libgapil.so}, e per poterla usare
occorrerà definire la variabile di ambiente \envvar{LD\_LIBRARY\_PATH} in modo
che il linker dinamico possa accedervi.

In generale questa variabile indica il \textit{pathname} della directory
contenente la libreria. Nell'ipotesi (che daremo sempre per verificata) che si
facciano le prove direttamente nella directory dei sorgenti (dove di norma
vengono creati sia i programmi che la libreria), il comando da dare sarà
\code{export LD\_LIBRARY\_PATH=./}; a questo punto potremo lanciare il server,
facendogli leggere una decina di frasi, con:
\begin{Console}
[piccardi@gont sources]$ \textbf{./fortuned -n10}
\end{Console}
%$

Avendo usato \func{daemon} per eseguire il server in background il comando
ritornerà immediatamente, ma potremo verificare con \cmd{ps} che in effetti il
programma resta un esecuzione in background, e senza avere associato un
terminale di controllo (si ricordi quanto detto in sez.~\ref{sec:sess_daemon}):
\begin{Console}
[piccardi@gont sources]$ \textbf{ps aux}
...
piccardi 27489  0.0  0.0  1204  356 ?        S    01:06   0:00 ./fortuned -n10
piccardi 27492  3.0  0.1  2492  764 pts/2    R    01:08   0:00 ps aux
\end{Console}
%$
e si potrà verificare anche che in \file{/tmp} è stata creata la \textit{fifo}
di ascolto \file{fortune.fifo}. A questo punto potremo interrogare il server
con il programma client; otterremo così:
\begin{Console}
[piccardi@gont sources]$ \textbf{./fortune}
Linux ext2fs has been stable for a long time, now it's time to break it
        -- Linuxkongreß '95 in Berlin
[piccardi@gont sources]$ \textbf{./fortune}
Let's call it an accidental feature.
        --Larry Wall
[piccardi@gont sources]$ \textbf{./fortune}
.........    Escape the 'Gates' of Hell
  `:::'                  .......  ......
   :::  *                  `::.    ::'
   ::: .::  .:.::.  .:: .::  `::. :'
   :::  ::   ::  ::  ::  ::    :::.
   ::: .::. .::  ::.  `::::. .:'  ::.
...:::.....................::'   .::::..
        -- William E. Roadcap
[piccardi@gont sources]$ ./fortune
Linux ext2fs has been stable for a long time, now it's time to break it
        -- Linuxkongreß '95 in Berlin
\end{Console}
%$
e ripetendo varie volte il comando otterremo, in ordine casuale, le dieci
frasi tenute in memoria dal server.

Infine per chiudere il server basterà inviargli un segnale di terminazione (ad
esempio con \cmd{killall fortuned}) e potremo verificare che il gestore del
segnale ha anche correttamente cancellato la \textit{fifo} di ascolto da
\file{/tmp}.

Benché il nostro sistema client-server funzioni, la sua struttura è piuttosto
complessa e continua ad avere vari inconvenienti\footnote{lo stesso Stevens,
  che esamina questa architettura in \cite{APUE}, nota come sia impossibile
  per il server sapere se un client è andato in crash, con la possibilità di
  far restare le \textit{fifo} temporanee sul filesystem, di come sia
  necessario intercettare \signal{SIGPIPE} dato che un client può terminare
  dopo aver fatto una richiesta, ma prima che la risposta sia inviata (cosa
  che nel nostro esempio non è stata fatta).}; in generale infatti
l'interfaccia delle \textit{fifo} non è adatta a risolvere questo tipo di
problemi, che possono essere affrontati in maniera più semplice ed efficace o
usando i socket (che tratteremo in dettaglio a partire da
cap.~\ref{cha:socket_intro}) o ricorrendo a meccanismi di comunicazione
diversi, come quelli che esamineremo in seguito.



\subsection{La funzione \func{socketpair}}
\label{sec:ipc_socketpair}

Un meccanismo di comunicazione molto simile alle \textit{pipe}, ma che non
presenta il problema della unidirezionalità del flusso dei dati, è quello dei
cosiddetti \textsl{socket locali} (o \textit{Unix domain socket}).  Tratteremo
in generale i socket in cap.~\ref{cha:socket_intro}, nell'ambito
dell'interfaccia che essi forniscono per la programmazione di rete, e vedremo
anche (in~sez.~\ref{sec:sock_sa_local}) come si possono utilizzare i file
speciali di tipo socket, analoghi a quelli associati alle \textit{fifo} (si
rammenti sez.~\ref{sec:file_file_types}) cui si accede però attraverso quella
medesima interfaccia; vale però la pena esaminare qui una modalità di uso dei
socket locali che li rende sostanzialmente identici ad una \textit{pipe}
bidirezionale.

La funzione di sistema \funcd{socketpair}, introdotta da BSD ma supportata in
genere da qualunque sistema che fornisca l'interfaccia dei socket ed inclusa
in POSIX.1-2001, consente infatti di creare una coppia di file descriptor
connessi fra loro (tramite un socket, appunto) senza dover ricorrere ad un
file speciale sul filesystem. I descrittori sono del tutto analoghi a quelli
che si avrebbero con una chiamata a \func{pipe}, con la sola differenza è che
in questo caso il flusso dei dati può essere effettuato in entrambe le
direzioni. Il prototipo della funzione è:

\begin{funcproto}{
\fhead{sys/types.h} 
\fhead{sys/socket.h}
\fdecl{int socketpair(int domain, int type, int protocol, int sv[2])}
\fdesc{Crea una coppia di socket connessi fra loro.} 
}

{La funzione ritorna $0$ in caso di successo e $-1$ per un errore, nel qual
  caso \var{errno} assumerà uno dei valori: 
  \begin{errlist}
  \item[\errcode{EAFNOSUPPORT}] i socket locali non sono supportati.
  \item[\errcode{EOPNOTSUPP}] il protocollo specificato non supporta la
  creazione di coppie di socket.
  \item[\errcode{EPROTONOSUPPORT}] il protocollo specificato non è supportato.
  \end{errlist}
  ed inoltre \errval{EFAULT}, \errval{EMFILE} e \errval{ENFILE} nel loro
  significato generico.}
\end{funcproto}

La funzione restituisce in \param{sv} la coppia di descrittori connessi fra di
loro: quello che si scrive su uno di essi sarà ripresentato in input
sull'altro e viceversa. Gli argomenti \param{domain}, \param{type} e
\param{protocol} derivano dall'interfaccia dei socket (vedi
sez.~\ref{sec:sock_creation}) che è quella che fornisce il substrato per
connettere i due descrittori, ma in questo caso i soli valori validi che
possono essere specificati sono rispettivamente \const{AF\_UNIX},
\const{SOCK\_STREAM} e \val{0}.  

A partire dal kernel 2.6.27 la funzione supporta nell'indicazione del tipo di
socket anche i due flag \const{SOCK\_NONBLOCK} e \const{SOCK\_CLOEXEC}
(trattati in sez.~\ref{sec:sock_type}), con effetto identico agli analoghi
\const{O\_CLOEXEC} e \const{O\_NONBLOCK} di una \func{open} (vedi
tab.~\ref{tab:open_operation_flag}).

L'utilità di chiamare questa funzione per evitare due chiamate a \func{pipe}
può sembrare limitata; in realtà l'utilizzo di questa funzione (e dei socket
locali in generale) permette di trasmettere attraverso le linea non solo dei
dati, ma anche dei file descriptor: si può cioè passare da un processo ad un
altro un file descriptor, con una sorta di duplicazione dello stesso non
all'interno di uno stesso processo, ma fra processi distinti (torneremo su
questa funzionalità in sez.~\ref{sec:sock_fd_passing}).


\section{L'intercomunicazione fra processi di System V}
\label{sec:ipc_sysv}

Benché le \textit{pipe} e le \textit{fifo} siano ancora ampiamente usate, esse
scontano il limite fondamentale che il meccanismo di comunicazione che
forniscono è rigidamente sequenziale: una situazione in cui un processo scrive
qualcosa che molti altri devono poter leggere non può essere implementata con
una \textit{pipe}.

Per questo nello sviluppo di System V vennero introdotti una serie di nuovi
oggetti per la comunicazione fra processi ed una nuova interfaccia di
programmazione, poi inclusa anche in POSIX.1-2001, che fossero in grado di
garantire una maggiore flessibilità.  In questa sezione esamineremo come Linux
supporta quello che viene chiamato il \textsl{Sistema di comunicazione fra
  processi} di System V, cui da qui in avanti faremo riferimento come
\textit{SysV-IPC} (dove IPC è la sigla di \textit{Inter-Process
  Comunication}).



\subsection{Considerazioni generali}
\label{sec:ipc_sysv_generic}

La principale caratteristica del \textit{SysV-IPC} è quella di essere basato
su oggetti permanenti che risiedono nel kernel. Questi, a differenza di quanto
avviene per i file descriptor, non mantengono un contatore dei riferimenti, e
non vengono cancellati dal sistema una volta che non sono più in uso.  Questo
comporta due problemi: il primo è che, al contrario di quanto avviene per
\textit{pipe} e \textit{fifo}, la memoria allocata per questi oggetti non
viene rilasciata automaticamente quando non c'è più nessuno che li utilizzi ed
essi devono essere cancellati esplicitamente, se non si vuole che restino
attivi fino al riavvio del sistema. Il secondo problema è, dato che non c'è
come per i file un contatore del numero di riferimenti che ne indichi l'essere
in uso, che essi possono essere cancellati anche se ci sono dei processi che
li stanno utilizzando, con tutte le conseguenze (ovviamente assai sgradevoli)
del caso.

Un'ulteriore caratteristica negativa è che gli oggetti usati nel
\textit{SysV-IPC} vengono creati direttamente dal kernel, e sono accessibili
solo specificando il relativo \textsl{identificatore}. Questo è un numero
progressivo (un po' come il \ids{PID} dei processi) che il kernel assegna a
ciascuno di essi quanto vengono creati (sul procedimento di assegnazione
torneremo in sez.~\ref{sec:ipc_sysv_id_use}). L'identificatore viene
restituito dalle funzioni che creano l'oggetto, ed è quindi locale al processo
che le ha eseguite. Dato che l'identificatore viene assegnato dinamicamente
dal kernel non è possibile prevedere quale sarà, né utilizzare un qualche
valore statico, si pone perciò il problema di come processi diversi possono
accedere allo stesso oggetto.

Per risolvere il problema nella struttura \struct{ipc\_perm} che il kernel
associa a ciascun oggetto, viene mantenuto anche un campo apposito che
contiene anche una \textsl{chiave}, identificata da una variabile del tipo
primitivo \type{key\_t}, da specificare in fase di creazione dell'oggetto, e
tramite la quale è possibile ricavare l'identificatore.\footnote{in sostanza
  si sposta il problema dell'accesso dalla classificazione in base
  all'identificatore alla classificazione in base alla chiave, una delle tante
  complicazioni inutili presenti nel \textit{SysV-IPC}.} Oltre la chiave, la
struttura, la cui definizione è riportata in fig.~\ref{fig:ipc_ipc_perm},
mantiene varie proprietà ed informazioni associate all'oggetto.

\begin{figure}[!htb]
  \footnotesize \centering
  \begin{minipage}[c]{0.80\textwidth}
    \includestruct{listati/ipc_perm.h}
  \end{minipage} 
  \normalsize 
  \caption{La struttura \structd{ipc\_perm}, come definita in
    \headfiled{sys/ipc.h}.}
  \label{fig:ipc_ipc_perm}
\end{figure}

Usando la stessa chiave due processi diversi possono ricavare l'identificatore
associato ad un oggetto ed accedervi. Il problema che sorge a questo punto è
come devono fare per accordarsi sull'uso di una stessa chiave. Se i processi
sono \textsl{imparentati} la soluzione è relativamente semplice, in tal caso
infatti si può usare il valore speciale \constd{IPC\_PRIVATE} per creare un
nuovo oggetto nel processo padre, l'identificatore così ottenuto sarà
disponibile in tutti i figli, e potrà essere passato come argomento attraverso
una \func{exec}.

Però quando i processi non sono \textsl{imparentati} (come capita tutte le
volte che si ha a che fare con un sistema client-server) tutto questo non è
possibile; si potrebbe comunque salvare l'identificatore su un file noto, ma
questo ovviamente comporta lo svantaggio di doverselo andare a rileggere.  Una
alternativa più efficace è quella che i programmi usino un valore comune per
la chiave (che ad esempio può essere dichiarato in un header comune), ma c'è
sempre il rischio che questa chiave possa essere stata già utilizzata da
qualcun altro.  Dato che non esiste una convenzione su come assegnare queste
chiavi in maniera univoca l'interfaccia mette a disposizione una funzione
apposita, \funcd{ftok}, che permette di ottenere una chiave specificando il
nome di un file ed un numero di versione; il suo prototipo è:

\begin{funcproto}{
\fhead{sys/types.h} 
\fhead{sys/ipc.h} 
\fdecl{key\_t ftok(const char *pathname, int proj\_id)}
\fdesc{Restituisce una chiave per identificare un oggetto del \textit{SysV
    IPC}.}}

{La funzione ritorna la chiave in caso di successo e $-1$ per un errore, nel
  qual caso \var{errno} assumerà uno dei possibili codici di errore di
  \func{stat}.}
\end{funcproto}

La funzione determina un valore della chiave sulla base di \param{pathname},
che deve specificare il \textit{pathname} di un file effettivamente esistente
e di un numero di progetto \param{proj\_id)}, che di norma viene specificato
come carattere, dato che ne vengono utilizzati solo gli 8 bit meno
significativi. Nelle \acr{libc4} e \acr{libc5}, come avviene in SunOS,
l'argomento \param{proj\_id} è dichiarato tipo \ctyp{char}, la \acr{glibc} usa
il prototipo specificato da XPG4, ma vengono lo stesso utilizzati gli 8 bit
meno significativi.

Il problema è che anche così non c'è la sicurezza che il valore della chiave
sia univoco, infatti esso è costruito combinando il byte di \param{proj\_id)}
con i 16 bit meno significativi dell'inode del file \param{pathname} (che
vengono ottenuti attraverso \func{stat}, da cui derivano i possibili errori),
e gli 8 bit meno significativi del numero del dispositivo su cui è il file.
Diventa perciò relativamente facile ottenere delle collisioni, specie se i
file sono su dispositivi con lo stesso \textit{minor number}, come
\file{/dev/hda1} e \file{/dev/sda1}.

In genere quello che si fa è utilizzare un file comune usato dai programmi che
devono comunicare (ad esempio un header comune, o uno dei programmi che devono
usare l'oggetto in questione), utilizzando il numero di progetto per ottenere
le chiavi che interessano. In ogni caso occorre sempre controllare, prima di
creare un oggetto, che la chiave non sia già stata utilizzata. Se questo va
bene in fase di creazione, le cose possono complicarsi per i programmi che
devono solo accedere, in quanto, a parte gli eventuali controlli sugli altri
attributi di \struct{ipc\_perm}, non esiste una modalità semplice per essere
sicuri che l'oggetto associato ad una certa chiave sia stato effettivamente
creato da chi ci si aspetta.

Questo è, insieme al fatto che gli oggetti sono permanenti e non mantengono un
contatore di riferimenti per la cancellazione automatica, il principale
problema del \textit{SysV-IPC}. Non esiste infatti una modalità chiara per
identificare un oggetto, come sarebbe stato se lo si fosse associato ad in
file, e tutta l'interfaccia è inutilmente complessa.  Per questo se ne
sconsiglia assolutamente l'uso nei nuovi programmi, considerato che è ormai
disponibile una revisione completa dei meccanismi di IPC fatta secondo quanto
indicato dallo standard POSIX.1b, che presenta una realizzazione più sicura ed
una interfaccia più semplice, che tratteremo in sez.~\ref{sec:ipc_posix}.


\subsection{Il controllo di accesso}
\label{sec:ipc_sysv_access_control}

Oltre alle chiavi, abbiamo visto in fig.~\ref{fig:ipc_ipc_perm} che ad ogni
oggetto sono associate in \struct{ipc\_perm} ulteriori informazioni, come gli
identificatori del creatore (nei campi \var{cuid} e \var{cgid}) e del
proprietario (nei campi \var{uid} e \var{gid}) dello stesso, e un insieme di
permessi (nel campo \var{mode}). In questo modo è possibile definire un
controllo di accesso sugli oggetti di IPC, simile a quello che si ha per i
file (vedi sez.~\ref{sec:file_perm_overview}).

Benché questo controllo di accesso sia molto simile a quello dei file, restano
delle importanti differenze. La prima è che il permesso di esecuzione non
esiste (e se specificato viene ignorato), per cui si può parlare solo di
permessi di lettura e scrittura (nel caso dei semafori poi quest'ultimo è più
propriamente un permesso di modifica). I valori di \var{mode} sono gli stessi
ed hanno lo stesso significato di quelli riportati in
tab.~\ref{tab:file_mode_flags} e come per i file definiscono gli accessi per
il proprietario, il suo gruppo e tutti gli altri. 

Se però si vogliono usare le costanti simboliche di
tab.~\ref{tab:file_mode_flags} occorrerà includere anche il file
\headfile{sys/stat.h}; alcuni sistemi definiscono le costanti \constd{MSG\_R}
(il valore ottale \texttt{0400}) e \constd{MSG\_W} (il valore ottale
\texttt{0200}) per indicare i permessi base di lettura e scrittura per il
proprietario, da utilizzare, con gli opportuni shift, pure per il gruppo e gli
altri. In Linux, visto la loro scarsa utilità, queste costanti non sono
definite.

Quando l'oggetto viene creato i campi \var{cuid} e \var{uid} di
\struct{ipc\_perm} ed i campi \var{cgid} e \var{gid} vengono impostati
rispettivamente al valore dell'\ids{UID} e del \ids{GID} effettivo del processo
che ha chiamato la funzione, ma, mentre i campi \var{uid} e \var{gid} possono
essere cambiati, i campi \var{cuid} e \var{cgid} restano sempre gli stessi.

Il controllo di accesso è effettuato a due livelli. Il primo livello è nelle
funzioni che richiedono l'identificatore di un oggetto data la chiave. Queste
specificano tutte un argomento \param{flag}, in tal caso quando viene
effettuata la ricerca di una chiave, qualora \param{flag} specifichi dei
permessi, questi vengono controllati e l'identificatore viene restituito solo
se corrispondono a quelli dell'oggetto. Se ci sono dei permessi non presenti
in \var{mode} l'accesso sarà negato. Questo controllo però è di utilità
indicativa, dato che è sempre possibile specificare per \param{flag} un valore
nullo, nel qual caso l'identificatore sarà restituito comunque.

Il secondo livello di controllo è quello delle varie funzioni che accedono
direttamente (in lettura o scrittura) all'oggetto. In tal caso lo schema dei
controlli è simile a quello dei file, ed avviene secondo questa sequenza:
\begin{itemize*}
\item se il processo ha i privilegi di amministratore (più precisamente
  \const{CAP\_IPC\_OWNER}) l'accesso è sempre consentito.
\item se l'\ids{UID} effettivo del processo corrisponde o al valore del campo
  \var{cuid} o a quello del campo \var{uid} ed il permesso per il proprietario
  in \var{mode} è appropriato\footnote{per appropriato si intende che è
    impostato il permesso di scrittura per le operazioni di scrittura e quello
    di lettura per le operazioni di lettura.} l'accesso è consentito.
\item se il \ids{GID} effettivo del processo corrisponde o al
  valore del campo \var{cgid} o a quello del campo \var{gid} ed il permesso
  per il gruppo in \var{mode} è appropriato l'accesso è consentito.
\item se il permesso per gli altri è appropriato l'accesso è consentito.
\end{itemize*}
solo se tutti i controlli elencati falliscono l'accesso è negato. Si noti che
a differenza di quanto avviene per i permessi dei file, fallire in uno dei
passi elencati non comporta il fallimento dell'accesso. Un'ulteriore
differenza rispetto a quanto avviene per i file è che per gli oggetti di IPC
il valore di \textit{umask} (si ricordi quanto esposto in
sez.~\ref{sec:file_perm_management}) non ha alcun significato.


\subsection{Gli identificatori ed il loro utilizzo}
\label{sec:ipc_sysv_id_use}

L'unico campo di \struct{ipc\_perm} del quale non abbiamo ancora parlato è
\var{seq}, che in fig.~\ref{fig:ipc_ipc_perm} è qualificato con un criptico
``\textsl{numero di sequenza}'', ne parliamo adesso dato che esso è
strettamente attinente alle modalità con cui il kernel assegna gli
identificatori degli oggetti del sistema di IPC.

Quando il sistema si avvia, alla creazione di ogni nuovo oggetto di IPC viene
assegnato un numero progressivo, pari al numero di oggetti di quel tipo
esistenti. Se il comportamento fosse sempre questo sarebbe identico a quello
usato nell'assegnazione dei file descriptor nei processi, ed i valori degli
identificatori tenderebbero ad essere riutilizzati spesso e restare di piccole
dimensioni (inferiori al numero massimo di oggetti disponibili).

Questo va benissimo nel caso dei file descriptor, che sono locali ad un
processo, ma qui il comportamento varrebbe per tutto il sistema, e per
processi del tutto scorrelati fra loro. Così si potrebbero avere situazioni
come quella in cui un server esce e cancella le sue code di messaggi, ed il
relativo identificatore viene immediatamente assegnato a quelle di un altro
server partito subito dopo, con la possibilità che i client del primo non
facciano in tempo ad accorgersi dell'avvenuto, e finiscano con l'interagire
con gli oggetti del secondo, con conseguenze imprevedibili.

Proprio per evitare questo tipo di situazioni il sistema usa il valore di
\var{seq} per provvedere un meccanismo che porti gli identificatori ad
assumere tutti i valori possibili, rendendo molto più lungo il periodo in cui
un identificatore può venire riutilizzato.

Il sistema dispone sempre di un numero fisso di oggetti di IPC, fino al kernel
2.2.x questi erano definiti dalle costanti \const{MSGMNI}, \const{SEMMNI} e
\const{SHMMNI}, e potevano essere cambiati (come tutti gli altri limiti
relativi al \textit{SysV-IPC}) solo con una ricompilazione del kernel.  A
partire dal kernel 2.4.x è possibile cambiare questi valori a sistema attivo
scrivendo sui file \sysctlrelfile{kernel}{shmmni},
\sysctlrelfile{kernel}{msgmni} e \sysctlrelfiled{kernel}{sem} di
\file{/proc/sys/kernel} o con l'uso di \func{sysctl}.

\begin{figure}[!htb]
  \footnotesize \centering
  \begin{minipage}[c]{\codesamplewidth}
    \includecodesample{listati/IPCTestId.c}
  \end{minipage} 
  \normalsize 
  \caption{Sezione principale del programma di test per l'assegnazione degli
    identificatori degli oggetti di IPC \file{IPCTestId.c}.}
  \label{fig:ipc_sysv_idtest}
\end{figure}

Per ciascun tipo di oggetto di IPC viene mantenuto in \var{seq} un numero di
sequenza progressivo che viene incrementato di uno ogni volta che l'oggetto
viene cancellato. Quando l'oggetto viene creato usando uno spazio che era già
stato utilizzato in precedenza, per restituire il nuovo identificatore al
numero di oggetti presenti viene sommato il valore corrente del campo
\var{seq}, moltiplicato per il numero massimo di oggetti di quel tipo.

Questo in realtà è quanto avveniva fino ai kernel della serie 2.2, dalla serie
2.4 viene usato lo stesso fattore di moltiplicazione per qualunque tipo di
oggetto, utilizzando il valore dalla costante \constd{IPCMNI} (definita in
\file{include/linux/ipc.h}), che indica il limite massimo complessivo per il
numero di tutti gli oggetti presenti nel \textit{SysV-IPC}, ed il cui default
è 32768.  Si evita così il riutilizzo degli stessi numeri, e si fa sì che
l'identificatore assuma tutti i valori possibili.

In fig.~\ref{fig:ipc_sysv_idtest} è riportato il codice di un semplice
programma di test che si limita a creare un oggetto di IPC (specificato con
una opzione a riga di comando), stamparne il numero di identificatore, e
cancellarlo, il tutto un numero di volte specificato tramite una seconda
opzione.  La figura non riporta il codice di selezione delle opzioni, che
permette di inizializzare i valori delle variabili \var{type} al tipo di
oggetto voluto, e \var{n} al numero di volte che si vuole effettuare il ciclo
di creazione, stampa, cancellazione.

I valori di default sono per l'uso delle code di messaggi e per 5 ripetizioni
del ciclo. Per questo motivo se non si utilizzano opzioni verrà eseguito per
cinque volte il ciclo (\texttt{\small 7-11}), in cui si crea una coda di
messaggi (\texttt{\small 8}), se ne stampa l'identificativo (\texttt{\small
  9}) e la si rimuove (\texttt{\small 10}). Non stiamo ad approfondire adesso
il significato delle funzioni utilizzate, che verranno esaminate nelle
prossime sezioni.

Quello che ci interessa infatti è verificare l'allocazione degli
identificativi associati agli oggetti; lanciando il comando si otterrà
pertanto qualcosa del tipo:
\begin{Console}
piccardi@gont sources]$ \textbf{./ipctestid}
Identifier Value 0 
Identifier Value 32768 
Identifier Value 65536 
Identifier Value 98304 
Identifier Value 131072
\end{Console}
%$
il che ci mostra che stiamo lavorando con un kernel posteriore alla serie 2.2
nel quale non avevamo ancora usato nessuna coda di messaggi (il valore nullo
del primo identificativo indica che il campo \var{seq} era zero). Ripetendo il
comando, e quindi eseguendolo in un processo diverso, in cui non può esistere
nessuna traccia di quanto avvenuto in precedenza, otterremo come nuovo
risultato:
\begin{Console}
[piccardi@gont sources]$ \textbf{./ipctestid}
Identifier Value 163840
Identifier Value 196608 
Identifier Value 229376 
Identifier Value 262144 
Identifier Value 294912 
\end{Console}
%$
in cui la sequenza numerica prosegue, cosa che ci mostra come il valore di
\var{seq} continui ad essere incrementato e costituisca in effetti una
quantità mantenuta all'interno del sistema ed indipendente dai processi.


\subsection{Code di messaggi}
\label{sec:ipc_sysv_mq}

Il primo oggetto introdotto dal \textit{SysV-IPC} è quello delle code di
messaggi.  Le code di messaggi sono oggetti analoghi alle \textit{pipe} o alle
\textit{fifo} ed il loro scopo principale è quello di fornire a processi
diversi un meccanismo con cui scambiarsi dei dati in forma di messaggio. Dato
che le \textit{pipe} e le \textit{fifo} costituiscono una ottima alternativa,
ed in genere sono molto più semplici da usare, le code di messaggi sono il
meno utilizzato degli oggetti introdotti dal \textit{SysV-IPC}.

La funzione di sistema che permette di ottenere l'identificativo di una coda
di messaggi esistente per potervi accedere, oppure di creare una nuova coda
qualora quella indicata non esista ancora, è \funcd{msgget}, e il suo
prototipo è:

\begin{funcproto}{
\fhead{sys/types.h}
\fhead{sys/ipc.h} 
\fhead{sys/msg.h} 
\fdecl{int msgget(key\_t key, int flag)}
\fdesc{Ottiene o crea una coda di messaggi.} 
}

{La funzione ritorna l'identificatore (un intero positivo) in caso di successo
  e $-1$ per un errore, nel qual caso \var{errno} assumerà uno dei valori:
  \begin{errlist}
  \item[\errcode{EACCES}] il processo chiamante non ha i privilegi per
    accedere alla coda richiesta.
  \item[\errcode{EEXIST}] si è richiesta la creazione di una coda che già
    esiste, ma erano specificati sia \const{IPC\_CREAT} che \const{IPC\_EXCL}.
  \item[\errcode{EIDRM}] la coda richiesta è marcata per essere cancellata
    (solo fino al kernel 2.3.20).
  \item[\errcode{ENOENT}] si è cercato di ottenere l'identificatore di una coda
    di messaggi specificando una chiave che non esiste e \const{IPC\_CREAT}
    non era specificato.
  \item[\errcode{ENOSPC}] si è cercato di creare una coda di messaggi quando è
    stato superato il limite massimo di code (\const{MSGMNI}).
 \end{errlist}
 ed inoltre \errval{ENOMEM} nel suo significato generico.}
\end{funcproto}

Le funzione (come le analoghe che si usano per gli altri oggetti) serve sia a
ottenere l'identificatore di una coda di messaggi esistente, che a crearne una
nuova. L'argomento \param{key} specifica la chiave che è associata
all'oggetto, eccetto il caso in cui si specifichi il valore
\const{IPC\_PRIVATE}, nel qual caso la coda è creata ex-novo e non vi è
associata alcuna chiave (per questo viene detta \textsl{privata}), ed il
processo e i suoi eventuali figli potranno farvi riferimento solo attraverso
l'identificatore.

Se invece si specifica un valore diverso da \const{IPC\_PRIVATE} (in Linux
questo significa un valore diverso da zero) l'effetto della funzione dipende
dal valore di \param{flag}, se questo è nullo la funzione si limita ad
effettuare una ricerca sugli oggetti esistenti, restituendo l'identificatore
se trova una corrispondenza, o fallendo con un errore di \errcode{ENOENT} se
non esiste o di \errcode{EACCES} se si sono specificati dei permessi non
validi.

Se invece si vuole creare una nuova coda di messaggi \param{flag} non può
essere nullo e deve essere fornito come maschera binaria, impostando il bit
corrispondente al valore \constd{IPC\_CREAT}. In questo caso i nove bit meno
significativi di \param{flag} saranno usati come permessi per il nuovo
oggetto, secondo quanto illustrato in sez.~\ref{sec:ipc_sysv_access_control}.
Se si imposta anche il bit corrispondente a \constd{IPC\_EXCL} la funzione avrà
successo solo se l'oggetto non esiste già, fallendo con un errore di
\errcode{EEXIST} altrimenti.

Si tenga conto che l'uso di \const{IPC\_PRIVATE} non impedisce ad altri
processi di accedere alla coda, se hanno privilegi sufficienti, una volta che
questi possano indovinare o ricavare, ad esempio per tentativi,
l'identificatore ad essa associato. Per come sono implementati gli oggetti di
IPC infatti non esiste alcun modo in cui si possa garantire l'accesso
esclusivo ad una coda di messaggi.  Usare \const{IPC\_PRIVATE} o
\const{IPC\_CREAT} e \const{IPC\_EXCL} per \param{flag} comporta solo la
creazione di una nuova coda.

\begin{table}[htb]
  \footnotesize
  \centering
  \begin{tabular}[c]{|c|r|l|l|}
    \hline
    \textbf{Costante} & \textbf{Valore} & \textbf{File in \file{/proc}}
    & \textbf{Significato} \\
    \hline
    \hline
    \constd{MSGMNI}&   16& \file{msgmni} & Numero massimo di code di
                                          messaggi.\\
    \constd{MSGMAX}& 8192& \file{msgmax} & Dimensione massima di un singolo
                                          messaggio.\\
    \constd{MSGMNB}&16384& \file{msgmnb} & Dimensione massima del contenuto di 
                                          una coda.\\
    \hline
  \end{tabular}
  \caption{Valori delle costanti associate ai limiti delle code di messaggi.}
  \label{tab:ipc_msg_limits}
\end{table}

Le code di messaggi sono caratterizzate da tre limiti fondamentali, un tempo
definiti staticamente e corrispondenti alle prime tre costanti riportate in
tab.~\ref{tab:ipc_msg_limits}.  Come accennato però con tutte le versioni più
recenti del kernel con Linux è possibile modificare questi limiti attraverso
l'uso di \func{sysctl} o scrivendo nei file \sysctlrelfiled{kernel}{msgmax},
\sysctlrelfiled{kernel}{msgmnb} e \sysctlrelfiled{kernel}{msgmni} di
\file{/proc/sys/kernel/}.

\itindbeg{linked~list}

Una coda di messaggi è costituita da una \textit{linked list}.\footnote{una
  \textit{linked list} è una tipica struttura di dati, organizzati in una
  lista in cui ciascun elemento contiene un puntatore al successivo. In questo
  modo la struttura è veloce nell'estrazione ed immissione dei dati dalle
  estremità dalla lista (basta aggiungere un elemento in testa o in coda ed
  aggiornare un puntatore), e relativamente veloce da attraversare in ordine
  sequenziale (seguendo i puntatori), è invece relativamente lenta
  nell'accesso casuale e nella ricerca.}  I nuovi messaggi vengono inseriti in
coda alla lista e vengono letti dalla cima, in fig.~\ref{fig:ipc_mq_schema} si
è riportato uno schema semplificato con cui queste strutture vengono mantenute
dal kernel. Lo schema illustrato in realtà è una semplificazione di quello
usato fino ai kernel della serie 2.2. A partire della serie 2.4 la gestione
delle code di messaggi è effettuata in maniera diversa (e non esiste una
struttura \kstruct{msqid\_ds} nel kernel), ma abbiamo mantenuto lo schema
precedente dato che illustra in maniera più che adeguata i principi di
funzionamento delle code di messaggi.

\itindend{linked~list}

\begin{figure}[!htb]
  \centering \includegraphics[width=13cm]{img/mqstruct}
  \caption{Schema delle strutture di una coda di messaggi
    (\kstructd{msqid\_ds} e \kstructd{msg}).}
  \label{fig:ipc_mq_schema}
\end{figure}


A ciascuna coda è associata una struttura \kstruct{msqid\_ds} la cui
definizione è riportata in fig.~\ref{fig:ipc_msqid_ds} ed a cui si accede
includendo \headfiled{sys/msg.h};
%
% INFO: sotto materiale obsoleto e non interessante
% In questa struttura il
% kernel mantiene le principali informazioni riguardo lo stato corrente della
% coda. Come accennato questo vale fino ai kernel della serie 2.2, essa viene
% usata nei kernel della serie 2.4 solo per compatibilità in quanto è quella
% restituita dalle funzioni dell'interfaccia; si noti come ci sia una differenza
% con i campi mostrati nello schema di fig.~\ref{fig:ipc_mq_schema} che sono
% presi dalla definizione di \file{include/linux/msg.h}, e fanno riferimento
% alla definizione della omonima struttura usata nel kernel. 
%In fig.~\ref{fig:ipc_msqid_ds} sono elencati i campi definiti in
%\headfile{sys/msg.h};  
si tenga presente che il campo \var{\_\_msg\_cbytes} non è previsto dallo
standard POSIX.1-2001 e che alcuni campi fino al kernel 2.2 erano definiti
come \ctyp{short}.

\begin{figure}[!htb]
  \footnotesize \centering
  \begin{minipage}[c]{.91\textwidth}
    \includestruct{listati/msqid_ds.h}
  \end{minipage} 
  \normalsize 
  \caption{La struttura \structd{msqid\_ds}, associata a ciascuna coda di
    messaggi.}
  \label{fig:ipc_msqid_ds}
\end{figure}

Quando si crea una nuova coda con \func{msgget} questa struttura viene
inizializzata,\footnote{in realtà viene inizializzata una struttura interna al
  kernel, ma i dati citati sono gli stessi.} in particolare il campo
\var{msg\_perm} che esprime i permessi di accesso viene inizializzato nella
modalità illustrata in sez.~\ref{sec:ipc_sysv_access_control}. Per quanto
riguarda gli altri campi invece:
\begin{itemize*}
\item il campo \var{msg\_qnum}, che esprime il numero di messaggi presenti
  sulla coda, viene inizializzato a 0.
\item i campi \var{msg\_lspid} e \var{msg\_lrpid}, che esprimono
  rispettivamente il \ids{PID} dell'ultimo processo che ha inviato o ricevuto
  un messaggio sulla coda, sono inizializzati a 0.
\item i campi \var{msg\_stime} e \var{msg\_rtime}, che esprimono
  rispettivamente il tempo in cui è stato inviato o ricevuto l'ultimo
  messaggio sulla coda, sono inizializzati a 0.
\item il campo \var{msg\_ctime}, che esprime il tempo di ultima modifica della
  coda, viene inizializzato al tempo corrente.
\item il campo \var{msg\_qbytes}, che esprime la dimensione massima del
  contenuto della coda (in byte) viene inizializzato al valore preimpostato
  del sistema (\const{MSGMNB}).
\item il campo \var{\_\_msg\_cbytes}, che esprime la dimensione in byte dei
  messaggi presenti sulla coda, viene inizializzato a zero.
% i campi \var{msg\_first} e \var{msg\_last} che esprimono l'indirizzo del
%   primo e ultimo messaggio sono inizializzati a \val{NULL} e
%   \var{msg\_cbytes}, che esprime la dimensione in byte dei messaggi presenti è
%   inizializzato a zero. Questi campi sono ad uso interno dell'implementazione
%   e non devono essere utilizzati da programmi in \textit{user space}).
\end{itemize*}

Una volta creata una coda di messaggi le operazioni di controllo vengono
effettuate con la funzione di sistema \funcd{msgctl}, che, come le analoghe
\func{semctl} e \func{shmctl}, fa le veci di quello che \func{ioctl} è per i
file; il suo prototipo è:

\begin{funcproto}{
\fhead{sys/types.h}
\fhead{sys/ipc.h}
\fhead{sys/msg.h}
\fdecl{int msgctl(int msqid, int cmd, struct msqid\_ds *buf)}
\fdesc{Esegue una operazione su una coda.} 
}

{La funzione ritorna $0$ in caso di successo e $-1$ per un errore, nel qual
  caso \var{errno} assumerà uno dei valori: 
  \begin{errlist}
  \item[\errcode{EACCES}] si è richiesto \const{IPC\_STAT} ma processo
    chiamante non ha i privilegi di lettura sulla coda.
  \item[\errcode{EIDRM}] la coda richiesta è stata cancellata.
  \item[\errcode{EPERM}] si è richiesto \const{IPC\_SET} o \const{IPC\_RMID} ma
    il processo non ha i privilegi, o si è richiesto di aumentare il valore di
    \var{msg\_qbytes} oltre il limite \const{MSGMNB} senza essere
    amministratore.
  \end{errlist}
  ed inoltre \errval{EFAULT} ed \errval{EINVAL} nel loro significato
  generico.}
\end{funcproto}

La funzione permette di eseguire una operazione di controllo per la coda
specificata dall'identificatore \param{msqid}, utilizzando i valori della
struttura \struct{msqid\_ds}, mantenuta all'indirizzo \param{buf}. Il
comportamento della funzione dipende dal valore dell'argomento \param{cmd},
che specifica il tipo di azione da eseguire. I valori possibili
per \param{cmd} sono:
\begin{basedescript}{\desclabelwidth{1.6cm}\desclabelstyle{\nextlinelabel}}
\item[\constd{IPC\_STAT}] Legge le informazioni riguardo la coda nella
  struttura \struct{msqid\_ds} indicata da \param{buf}. Occorre avere il
  permesso di lettura sulla coda.
\item[\constd{IPC\_RMID}] Rimuove la coda, cancellando tutti i dati, con
  effetto immediato. Tutti i processi che cercheranno di accedere alla coda
  riceveranno un errore di \errcode{EIDRM}, e tutti processi in attesa su
  funzioni di lettura o di scrittura sulla coda saranno svegliati ricevendo
  il medesimo errore. Questo comando può essere eseguito solo da un processo
  con \ids{UID} effettivo corrispondente al creatore o al proprietario della
  coda, o all'amministratore.
\item[\constd{IPC\_SET}] Permette di modificare i permessi ed il proprietario
  della coda, ed il limite massimo sulle dimensioni del totale dei messaggi in
  essa contenuti (\var{msg\_qbytes}). I valori devono essere passati in una
  struttura \struct{msqid\_ds} puntata da \param{buf}.  Per modificare i
  valori di \var{msg\_perm.mode}, \var{msg\_perm.uid} e \var{msg\_perm.gid}
  occorre essere il proprietario o il creatore della coda, oppure
  l'amministratore e lo stesso vale per \var{msg\_qbytes}. Infine solo
  l'amministratore (più precisamente un processo con la capacità
  \const{CAP\_IPC\_RESOURCE}) ha la facoltà di incrementarne il valore a
  limiti superiori a \const{MSGMNB}. Se eseguita con successo la funzione
  aggiorna anche il campo \var{msg\_ctime}.
\end{basedescript}

A questi tre valori, che sono quelli previsti dallo standard, su Linux se ne
affiancano altri tre (\constd{IPC\_INFO}, \constd{MSG\_STAT} e
\constd{MSG\_INFO}) introdotti ad uso del programma \cmd{ipcs} per ottenere le
informazioni generali relative alle risorse usate dalle code di
messaggi. Questi potranno essere modificati o rimossi in favore dell'uso di
\texttt{/proc}, per cui non devono essere usati e non li tratteremo.

Una volta che si abbia a disposizione l'identificatore, per inviare un
messaggio su una coda si utilizza la funzione di sistema \funcd{msgsnd}, il
cui prototipo è:

\begin{funcproto}{
\fhead{sys/types.h}
\fhead{sys/ipc.h}
\fhead{sys/msg.h}
\fdecl{int msgsnd(int msqid, struct msgbuf *msgp, size\_t msgsz, int msgflg)}
\fdesc{Invia un messaggio su una coda.}
}

{La funzione ritorna $0$ in caso di successo e $-1$ per un errore, nel qual
  caso \var{errno} assumerà uno dei valori: 
  \begin{errlist}
  \item[\errcode{EACCES}] non si hanno i privilegi di accesso sulla coda.
  \item[\errcode{EAGAIN}] il messaggio non può essere inviato perché si è
    superato il limite \var{msg\_qbytes} sul numero massimo di byte presenti
    sulla coda, e si è richiesto \const{IPC\_NOWAIT} in \param{flag}.
  \item[\errcode{EIDRM}] la coda è stata cancellata.
  \item[\errcode{EINVAL}] si è specificato un \param{msgid} invalido, o un
    valore non positivo per \param{mtype}, o un valore di \param{msgsz}
    maggiore di \const{MSGMAX}.
  \end{errlist}
  ed inoltre \errval{EFAULT}, \errval{EINTR} e \errval{ENOMEM} nel loro
  significato generico.}
\end{funcproto}

La funzione inserisce il messaggio sulla coda specificata da \param{msqid}; il
messaggio ha lunghezza specificata da \param{msgsz} ed è passato attraverso il
l'argomento \param{msgp}.  Quest'ultimo deve venire passato sempre come
puntatore ad una struttura \struct{msgbuf} analoga a quella riportata in
fig.~\ref{fig:ipc_msbuf} che è quella che deve contenere effettivamente il
messaggio.  La dimensione massima per il testo di un messaggio non può
comunque superare il limite \const{MSGMAX}.

\begin{figure}[!htb]
  \footnotesize \centering
  \begin{minipage}[c]{0.8\textwidth}
    \includestruct{listati/msgbuf.h}
  \end{minipage} 
  \normalsize 
  \caption{Schema della struttura \structd{msgbuf}, da utilizzare come
    argomento per inviare/ricevere messaggi.}
  \label{fig:ipc_msbuf}
\end{figure}

La struttura di fig.~\ref{fig:ipc_msbuf} è comunque solo un modello, tanto che
la definizione contenuta in \headfile{sys/msg.h} usa esplicitamente per il
secondo campo il valore \code{mtext[1]}, che non è di nessuna utilità ai fini
pratici.  La sola cosa che conta è che la struttura abbia come primo membro un
campo \var{mtype} come nell'esempio; esso infatti serve ad identificare il
tipo di messaggio e deve essere sempre specificato come intero positivo di
tipo \ctyp{long}.  Il campo \var{mtext} invece può essere di qualsiasi tipo e
dimensione, e serve a contenere il testo del messaggio.

In generale pertanto per inviare un messaggio con \func{msgsnd} si usa
ridefinire una struttura simile a quella di fig.~\ref{fig:ipc_msbuf}, adattando
alle proprie esigenze il campo \var{mtype}, (o ridefinendo come si vuole il
corpo del messaggio, anche con più campi o con strutture più complesse) avendo
però la cura di mantenere nel primo campo un valore di tipo \ctyp{long} che ne
indica il tipo.

Si tenga presente che la lunghezza che deve essere indicata in questo
argomento è solo quella del messaggio, non quella di tutta la struttura, se
cioè \var{message} è una propria struttura che si passa alla funzione,
\param{msgsz} dovrà essere uguale a \code{sizeof(message)-sizeof(long)}, (se
consideriamo il caso dell'esempio in fig.~\ref{fig:ipc_msbuf}, \param{msgsz}
dovrà essere pari a \var{LENGTH}).

Per capire meglio il funzionamento della funzione riprendiamo in
considerazione la struttura della coda illustrata in
fig.~\ref{fig:ipc_mq_schema}. Alla chiamata di \func{msgsnd} il nuovo
messaggio sarà aggiunto in fondo alla lista inserendo una nuova struttura
\kstruct{msg}, il puntatore \var{msg\_last} di \kstruct{msqid\_ds} verrà
aggiornato, come pure il puntatore al messaggio successivo per quello che era
il precedente ultimo messaggio; il valore di \var{mtype} verrà mantenuto in
\var{msg\_type} ed il valore di \param{msgsz} in \var{msg\_ts}; il testo del
messaggio sarà copiato all'indirizzo specificato da \var{msg\_spot}.

Il valore dell'argomento \param{flag} permette di specificare il comportamento
della funzione. Di norma, quando si specifica un valore nullo, la funzione
ritorna immediatamente a meno che si sia ecceduto il valore di
\var{msg\_qbytes}, o il limite di sistema sul numero di messaggi, nel qual
caso si blocca.  Se si specifica per \param{flag} il valore
\constd{IPC\_NOWAIT} la funzione opera in modalità non-bloccante, ed in questi
casi ritorna immediatamente con un errore di \errcode{EAGAIN}.

Se non si specifica \const{IPC\_NOWAIT} la funzione resterà bloccata fintanto
che non si liberano risorse sufficienti per poter inserire nella coda il
messaggio, nel qual caso ritornerà normalmente. La funzione può ritornare con
una condizione di errore anche in due altri casi: quando la coda viene rimossa
(nel qual caso si ha un errore di \errcode{EIDRM}) o quando la funzione viene
interrotta da un segnale (nel qual caso si ha un errore di \errcode{EINTR}).

Una volta completato con successo l'invio del messaggio sulla coda, la
funzione aggiorna i dati mantenuti in \struct{msqid\_ds}, in particolare
vengono modificati:
\begin{itemize*}
\item Il valore di \var{msg\_lspid}, che viene impostato al \ids{PID} del
  processo chiamante.
\item Il valore di \var{msg\_qnum}, che viene incrementato di uno.
\item Il valore \var{msg\_stime}, che viene impostato al tempo corrente.
\end{itemize*}

La funzione di sistema che viene utilizzata per estrarre un messaggio da una
coda è \funcd{msgrcv}, ed il suo prototipo è:

\begin{funcproto}{
\fhead{sys/types.h}
\fhead{sys/ipc.h} 
\fhead{sys/msg.h}
\fdecl{ssize\_t msgrcv(int msqid, struct msgbuf *msgp, size\_t msgsz,
    long msgtyp, int msgflg)}
\fdesc{Legge un messaggio da una coda.} 
}

{La funzione ritorna il numero di byte letti in caso di successo e $-1$ per un
  errore, nel qual caso \var{errno} assumerà uno dei valori:
  \begin{errlist}
  \item[\errcode{E2BIG}] il testo del messaggio è più lungo di \param{msgsz} e
    non si è specificato \const{MSG\_NOERROR} in \param{msgflg}.
  \item[\errcode{EACCES}] non si hanno i privilegi di accesso sulla coda.
  \item[\errcode{EIDRM}] la coda è stata cancellata.
  \item[\errcode{EINTR}] la funzione è stata interrotta da un segnale mentre
    era in attesa di ricevere un messaggio.
  \item[\errcode{EINVAL}] si è specificato un \param{msgid} invalido o un
    valore di \param{msgsz} negativo.
  \end{errlist}
  ed inoltre \errval{EFAULT} nel suo significato generico.}
\end{funcproto}

La funzione legge un messaggio dalla coda specificata da \param{msqid},
scrivendolo sulla struttura puntata da \param{msgp}, che dovrà avere un
formato analogo a quello di fig.~\ref{fig:ipc_msbuf}.  Una volta estratto, il
messaggio sarà rimosso dalla coda.  L'argomento \param{msgsz} indica la
lunghezza massima del testo del messaggio (equivalente al valore del parametro
\var{LENGTH} nell'esempio di fig.~\ref{fig:ipc_msbuf}).

Se il testo del messaggio ha lunghezza inferiore a \param{msgsz} esso viene
rimosso dalla coda; in caso contrario, se \param{msgflg} è impostato a
\constd{MSG\_NOERROR}, il messaggio viene troncato e la parte in eccesso viene
perduta, altrimenti il messaggio non viene estratto e la funzione ritorna con
un errore di \errcode{E2BIG}.

L'argomento \param{msgtyp} permette di restringere la ricerca ad un
sottoinsieme dei messaggi presenti sulla coda; la ricerca infatti è fatta con
una scansione della struttura mostrata in fig.~\ref{fig:ipc_mq_schema},
restituendo il primo messaggio incontrato che corrisponde ai criteri
specificati (che quindi, visto come i messaggi vengono sempre inseriti dalla
coda, è quello meno recente); in particolare:
\begin{itemize*}
\item se \param{msgtyp} è 0 viene estratto il messaggio in cima alla coda, cioè
  quello fra i presenti che è stato inserito per primo. 
\item se \param{msgtyp} è positivo viene estratto il primo messaggio il cui
  tipo (il valore del campo \var{mtype}) corrisponde al valore di
  \param{msgtyp}.
\item se \param{msgtyp} è negativo viene estratto il primo fra i messaggi con
  il valore più basso del tipo, fra tutti quelli il cui tipo ha un valore
  inferiore al valore assoluto di \param{msgtyp}.
\end{itemize*}

Il valore di \param{msgflg} permette di controllare il comportamento della
funzione, esso può essere nullo o una maschera binaria composta da uno o più
valori.  Oltre al precedente \const{MSG\_NOERROR}, sono possibili altri due
valori: \constd{MSG\_EXCEPT}, che permette, quando \param{msgtyp} è positivo,
di leggere il primo messaggio nella coda con tipo diverso da \param{msgtyp}, e
\const{IPC\_NOWAIT} che causa il ritorno immediato della funzione quando non
ci sono messaggi sulla coda.

Il comportamento usuale della funzione infatti, se non ci sono messaggi
disponibili per la lettura, è di bloccare il processo. Nel caso però si sia
specificato \const{IPC\_NOWAIT} la funzione ritorna immediatamente con un
errore \errcode{ENOMSG}. Altrimenti la funzione ritorna normalmente non appena
viene inserito un messaggio del tipo desiderato, oppure ritorna con errore
qualora la coda sia rimossa (con \var{errno} impostata a \errcode{EIDRM}) o se
il processo viene interrotto da un segnale (con \var{errno} impostata a
\errcode{EINTR}).

Una volta completata con successo l'estrazione del messaggio dalla coda, la
funzione aggiorna i dati mantenuti in \struct{msqid\_ds}, in particolare
vengono modificati:
\begin{itemize*}
\item Il valore di \var{msg\_lrpid}, che viene impostato al \ids{PID} del
  processo chiamante.
\item Il valore di \var{msg\_qnum}, che viene decrementato di uno.
\item Il valore \var{msg\_rtime}, che viene impostato al tempo corrente.
\end{itemize*}

Le code di messaggi presentano il solito problema di tutti gli oggetti del
\textit{SysV-IPC} che essendo questi permanenti restano nel sistema occupando
risorse anche quando un processo è terminato, al contrario delle \textit{pipe}
per le quali tutte le risorse occupate vengono rilasciate quanto l'ultimo
processo che le utilizzava termina. Questo comporta che in caso di errori si
può saturare il sistema, e che devono comunque essere esplicitamente previste
delle funzioni di rimozione in caso di interruzioni o uscite dal programma
(come vedremo in fig.~\ref{fig:ipc_mq_fortune_server}).

L'altro problema è che non facendo uso di file descriptor le tecniche di
\textit{I/O multiplexing} descritte in sez.~\ref{sec:file_multiplexing} non
possono essere utilizzate, e non si ha a disposizione niente di analogo alle
funzioni \func{select} e \func{poll}. Questo rende molto scomodo usare più di
una di queste strutture alla volta; ad esempio non si può scrivere un server
che aspetti un messaggio su più di una coda senza fare ricorso ad una tecnica
di \textit{polling} che esegua un ciclo di attesa su ciascuna di esse.

Come esempio dell'uso delle code di messaggi possiamo riscrivere il nostro
server di \textit{fortunes} usando queste al posto delle \textit{fifo}. In
questo caso useremo una sola coda di messaggi, usando il tipo di messaggio per
comunicare in maniera indipendente con client diversi.

\begin{figure}[!htbp]
  \footnotesize \centering
  \begin{minipage}[c]{\codesamplewidth}
    \includecodesample{listati/MQFortuneServer.c}
  \end{minipage} 
  \normalsize 
  \caption{Sezione principale del codice del server di \textit{fortunes}
    basato sulle \textit{message queue}.}
  \label{fig:ipc_mq_fortune_server}
\end{figure}

In fig.~\ref{fig:ipc_mq_fortune_server} si è riportato un estratto delle parti
principali del codice del nuovo server (il codice completo è nel file
\file{MQFortuneServer.c} nei sorgenti allegati). Il programma è basato su un
uso accorto della caratteristica di poter associate un ``tipo'' ai messaggi
per permettere una comunicazione indipendente fra il server ed i vari client,
usando il \ids{PID} di questi ultimi come identificativo. Questo è possibile
in quanto, al contrario di una \textit{fifo}, la lettura di una coda di
messaggi può non essere sequenziale, proprio grazie alla classificazione dei
messaggi sulla base del loro tipo.

Il programma, oltre alle solite variabili per il nome del file da cui leggere
le \textit{fortunes} e per il vettore di stringhe che contiene le frasi,
definisce due strutture appositamente per la comunicazione; con
\var{msgbuf\_read} vengono passate (\texttt{\small 8-11}) le richieste mentre
con \var{msgbuf\_write} vengono restituite (\texttt{\small 12-15}) le frasi.

La gestione delle opzioni si è al solito omessa, essa si curerà di impostare
nella variabile \var{n} il numero di frasi da leggere specificato a linea di
comando ed in \var{fortunefilename} il file da cui leggerle. Dopo aver
installato (\texttt{\small 19-21}) i gestori dei segnali per trattare
l'uscita dal server, viene prima controllato (\texttt{\small 22}) il numero di
frasi richieste abbia senso (cioè sia maggiore di zero), le quali poi vengono
lette (\texttt{\small 23}) nel vettore in memoria con la stessa funzione
\code{FortuneParse} usata anche per il server basato sulle \textit{fifo}.

Una volta inizializzato il vettore di stringhe coi messaggi presi dal file
delle \textit{fortune} si procede (\texttt{\small 25}) con la generazione di
una chiave per identificare la coda di messaggi (si usa il nome del file dei
sorgenti del server) con la quale poi si esegue (\texttt{\small 26}) la
creazione della stessa (si noti come si sia chiamata \func{msgget} con un
valore opportuno per l'argomento \param{flag}), avendo cura di abortire il
programma (\texttt{\small 27-29}) in caso di errore.

Finita la fase di inizializzazione il server prima (\texttt{\small 32}) chiama
la funzione \func{daemon} per andare in background e poi esegue in permanenza
il ciclo principale (\texttt{\small 33-40}). Questo inizia (\texttt{\small
  34}) con il porsi in attesa di un messaggio di richiesta da parte di un
client. Si noti infatti come \func{msgrcv} richieda un messaggio con
\var{mtype} uguale a 1, questo è il valore usato per le richieste dato che
corrisponde al \ids{PID} di \cmd{init}, che non può essere un client. L'uso
del flag \const{MSG\_NOERROR} è solo per sicurezza, dato che i messaggi di
richiesta sono di dimensione fissa (e contengono solo il \ids{PID} del
client).

Se non sono presenti messaggi di richiesta \func{msgrcv} si bloccherà,
ritornando soltanto in corrispondenza dell'arrivo sulla coda di un messaggio
di richiesta da parte di un client, in tal caso il ciclo prosegue
(\texttt{\small 35}) selezionando una frase a caso, copiandola (\texttt{\small
  36}) nella struttura \var{msgbuf\_write} usata per la risposta e
calcolandone (\texttt{\small 37}) la dimensione.

Per poter permettere a ciascun client di ricevere solo la risposta indirizzata
a lui il tipo del messaggio in uscita viene inizializzato (\texttt{\small 38})
al valore del \ids{PID} del client ricevuto nel messaggio di richiesta.
L'ultimo passo del ciclo (\texttt{\small 39}) è inviare sulla coda il
messaggio di risposta. Si tenga conto che se la coda è piena anche questa
funzione potrà bloccarsi fintanto che non venga liberato dello spazio.

Si noti che il programma può terminare solo grazie ad una interruzione da
parte di un segnale; in tal caso verrà eseguito (\texttt{\small 45-48}) il
gestore \code{HandSIGTERM}, che semplicemente si limita a cancellare la coda
(\texttt{\small 46}) ed ad uscire (\texttt{\small 47}).

\begin{figure}[!htbp]
  \footnotesize \centering
  \begin{minipage}[c]{\codesamplewidth}
    \includecodesample{listati/MQFortuneClient.c}
  \end{minipage} 
  \normalsize 
  \caption{Sezione principale del codice del client di \textit{fortunes}
    basato sulle \textit{message queue}.}
  \label{fig:ipc_mq_fortune_client}
\end{figure}

In fig.~\ref{fig:ipc_mq_fortune_client} si è riportato un estratto il codice
del programma client.  Al solito il codice completo è con i sorgenti allegati,
nel file \file{MQFortuneClient.c}.  Come sempre si sono rimosse le parti
relative alla gestione delle opzioni, ed in questo caso, anche la
dichiarazione delle variabili, che, per la parte relative alle strutture usate
per la comunicazione tramite le code, sono le stesse viste in
fig.~\ref{fig:ipc_mq_fortune_server}.

Il client in questo caso è molto semplice; la prima parte del programma
(\texttt{\small 4-9}) si occupa di accedere alla coda di messaggi, ed è
identica a quanto visto per il server, solo che in questo caso \func{msgget}
non viene chiamata con il flag di creazione in quanto la coda deve essere
preesistente. In caso di errore (ad esempio se il server non è stato avviato)
il programma termina immediatamente. 

Una volta acquisito l'identificatore della coda il client compone
(\texttt{\small 12-13}) il messaggio di richiesta in \var{msg\_read}, usando
1 per il tipo ed inserendo il proprio \ids{PID} come dato da passare al
server.  Calcolata (\texttt{\small 14}) la dimensione, provvede
(\texttt{\small 15}) ad immettere la richiesta sulla coda.

A questo punto non resta che rileggere la risposta (\texttt{\small 16}) dalla
coda del server richiedendo a \func{msgrcv} di selezionare i messaggi di tipo
corrispondente al valore del \ids{PID} inviato nella richiesta. L'ultimo passo
(\texttt{\small 17}) prima di uscire è quello di stampare a video il messaggio
ricevuto.
 
Proviamo allora il nostro nuovo sistema, al solito occorre definire
\code{LD\_LIBRARY\_PATH} per accedere alla libreria \file{libgapil.so}, dopo
di che, in maniera del tutto analoga a quanto fatto con il programma che usa
le \textit{fifo}, potremo far partire il server con:
\begin{Console}
[piccardi@gont sources]$ \textbf{./mqfortuned -n10}
\end{Console}
%$
come nel caso precedente, avendo eseguito il server in background, il comando
ritornerà immediatamente; potremo però verificare con \cmd{ps} che il
programma è effettivamente in esecuzione, e che ha creato una coda di
messaggi:
\begin{Console}
[piccardi@gont sources]$ \textbf{ipcs}

------ Shared Memory Segments --------
key        shmid      owner      perms      bytes      nattch     status      

------ Semaphore Arrays --------
key        semid      owner      perms      nsems     

------ Message Queues --------
key        msqid      owner      perms      used-bytes   messages    
0x0102dc6a 0          piccardi   666        0            0           
\end{Console}
%$
a questo punto potremo usare il client per ottenere le nostre frasi:
\begin{Console}
[piccardi@gont sources]$ \textbf{./mqfortune}
Linux ext2fs has been stable for a long time, now it's time to break it
        -- Linuxkongreß '95 in Berlin
[piccardi@gont sources]$ \textbf{./mqfortune}
Let's call it an accidental feature.
        --Larry Wall
\end{Console} 
con un risultato del tutto equivalente al precedente. Infine potremo chiudere
il server inviando il segnale di terminazione con il comando \code{killall
  mqfortuned}, verificando che effettivamente la coda di messaggi venga
rimossa.

Benché funzionante questa architettura risente dello stesso inconveniente
visto anche nel caso del precedente server basato sulle \textit{fifo}; se il
client viene interrotto dopo l'invio del messaggio di richiesta e prima della
lettura della risposta, quest'ultima resta nella coda (così come per le
\textit{fifo} si aveva il problema delle \textit{fifo} che restavano nel
filesystem). In questo caso però il problemi sono maggiori, sia perché è molto
più facile esaurire la memoria dedicata ad una coda di messaggi che gli
\textit{inode} di un filesystem, sia perché, con il riutilizzo dei \ids{PID}
da parte dei processi, un client eseguito in un momento successivo potrebbe
ricevere un messaggio non indirizzato a lui.


\subsection{I semafori}
\label{sec:ipc_sysv_sem}

I semafori non sono propriamente meccanismi di intercomunicazione come
\textit{pipe}, \textit{fifo} e code di messaggi, poiché non consentono di
scambiare dati fra processi, ma servono piuttosto come meccanismi di
sincronizzazione o di protezione per le \textsl{sezioni critiche} del codice
(si ricordi quanto detto in sez.~\ref{sec:proc_race_cond}).  Un semaforo
infatti non è altro che un contatore mantenuto nel kernel che determina se
consentire o meno la prosecuzione dell'esecuzione di un programma. In questo
modo si può controllare l'accesso ad una risorsa condivisa da più processi,
associandovi un semaforo che assicuri che non possa essere usata da più di un
processo alla volta.

Il concetto di semaforo è uno dei concetti base nella programmazione ed è
assolutamente generico, così come del tutto generali sono modalità con cui lo
si utilizza. Un processo che deve accedere ad una risorsa condivisa eseguirà
un controllo del semaforo: se questo è positivo il suo valore sarà
decrementato, indicando che si è consumato una unità della risorsa, ed il
processo potrà proseguire nell'utilizzo di quest'ultima, provvedendo a
rilasciarla, una volta completate le operazioni volute, reincrementando il
semaforo.

Se al momento del controllo il valore del semaforo è nullo la risorsa viene
considerata non disponibile, ed il processo si bloccherà fin quando chi la sta
utilizzando non la rilascerà, incrementando il valore del semaforo. Non appena
il semaforo diventa positivo, indicando che la risorsa è tornata disponibile,
il processo bloccato in attesa riprenderà l'esecuzione, e potrà operare come
nel caso precedente (decremento del semaforo, accesso alla risorsa, incremento
del semaforo).

Per poter implementare questo tipo di logica le operazioni di controllo e
decremento del contatore associato al semaforo devono essere atomiche,
pertanto una realizzazione di un oggetto di questo tipo è necessariamente
demandata al kernel. La forma più semplice di semaforo è quella del
\textsl{semaforo binario}, o \textit{mutex}, in cui un valore diverso da zero
(normalmente 1) indica la libertà di accesso, e un valore nullo l'occupazione
della risorsa. In generale però si possono usare semafori con valori interi,
utilizzando il valore del contatore come indicatore del ``numero di risorse''
ancora disponibili.

Il sistema di intercomunicazione di \textit{SysV-IPC} prevede anche una
implementazione dei semafori, ma gli oggetti utilizzati sono tuttavia non
semafori singoli, ma gruppi (più propriamente \textsl{insiemi}) di semafori
detti ``\textit{semaphore set}''. La funzione di sistema che permette di
creare o ottenere l'identificatore di un insieme di semafori è \funcd{semget},
ed il suo prototipo è:

\begin{funcproto}{
\fhead{sys/types.h}
\fhead{sys/ipc.h}
\fhead{sys/sem.h}
\fdecl{int semget(key\_t key, int nsems, int flag)}
\fdesc{Restituisce l'identificatore di un insieme di semafori.} 
}

{La funzione ritorna l'identificatore (un intero positivo) in caso di successo
  e $-1$ per un errore, nel qual caso \var{errno} assumerà uno dei valori:
  \begin{errlist}
  \item[\errcode{ENOSPC}] si è superato il limite di sistema per il numero
    totale di semafori (\const{SEMMNS}) o di insiemi (\const{SEMMNI}).
  \item[\errcode{EINVAL}] \param{nsems} è minore di zero o maggiore del limite
    sul numero di semafori di un insieme (\const{SEMMSL}), o se l'insieme già
    esiste, maggiore del numero di semafori che contiene.
  \item[\errcode{ENOMEM}] il sistema non ha abbastanza memoria per poter
    contenere le strutture per un nuovo insieme di semafori.
  \end{errlist}
  ed inoltre \errval{EACCES}, \errval{EEXIST}, \errval{EIDRM} e
  \errval{ENOENT} con lo stesso significato che hanno per \func{msgget}.}
\end{funcproto}

La funzione è del tutto analoga a \func{msgget}, solo che in questo caso
restituisce l'identificatore di un insieme di semafori, in particolare è
identico l'uso degli argomenti \param{key} e \param{flag}, per cui non
ripeteremo quanto detto al proposito in sez.~\ref{sec:ipc_sysv_mq}. L'argomento
\param{nsems} permette di specificare quanti semafori deve contenere l'insieme
quando se ne richieda la creazione, e deve essere nullo quando si effettua una
richiesta dell'identificatore di un insieme già esistente.

Purtroppo questa implementazione complica inutilmente lo schema elementare che
abbiamo descritto, dato che non è possibile definire un singolo semaforo, ma
se ne deve creare per forza un insieme.  Ma questa in definitiva è solo una
complicazione inutile dell'interfaccia, il problema è che i semafori forniti
dal \textit{SysV-IPC} soffrono di altri due difetti progettuali molto più
gravi.

Il primo difetto è che non esiste una funzione che permetta di creare ed
inizializzare un semaforo in un'unica chiamata; occorre prima creare l'insieme
dei semafori con \func{semget} e poi inizializzarlo con \func{semctl}, si
perde così ogni possibilità di eseguire l'operazione atomicamente. Eventuali
accessi che possono avvenire fra la creazione e l'inizializzazione potranno
avere effetti imprevisti.

Il secondo difetto deriva dalla caratteristica generale degli oggetti del
\textit{SysV-IPC} di essere risorse globali di sistema, che non vengono
cancellate quando nessuno le usa più. In questo caso il problema è più grave
perché ci si a trova a dover affrontare esplicitamente il caso in cui un
processo termina per un qualche errore lasciando un semaforo occupato, che
resterà tale fino al successivo riavvio del sistema. Come vedremo esistono
delle modalità per evitare tutto ciò, ma diventa necessario indicare
esplicitamente che si vuole il ripristino del semaforo all'uscita del
processo, e la gestione diventa più complicata.

\begin{figure}[!htb]
  \footnotesize \centering
  \begin{minipage}[c]{.85\textwidth}
    \includestruct{listati/semid_ds.h}
  \end{minipage} 
  \normalsize 
  \caption{La struttura \structd{semid\_ds}, associata a ciascun insieme di
    semafori.}
  \label{fig:ipc_semid_ds}
\end{figure}

A ciascun insieme di semafori è associata una struttura \struct{semid\_ds},
riportata in fig.~\ref{fig:ipc_semid_ds}.\footnote{anche in questo caso in
  realtà il kernel usa una sua specifica struttura interna, ma i dati
  significativi sono sempre quelli citati.}  Come nel caso delle code di
messaggi quando si crea un nuovo insieme di semafori con \func{semget} questa
struttura viene inizializzata. In particolare il campo \var{sem\_perm}, che
esprime i permessi di accesso, viene inizializzato come illustrato in
sez.~\ref{sec:ipc_sysv_access_control} (si ricordi che in questo caso il
permesso di scrittura è in realtà permesso di alterare il semaforo), per
quanto riguarda gli altri campi invece:
\begin{itemize*}
\item il campo \var{sem\_nsems}, che esprime il numero di semafori
  nell'insieme, viene inizializzato al valore di \param{nsems}.
\item il campo \var{sem\_ctime}, che esprime il tempo di ultimo cambiamento
  dell'insieme, viene inizializzato al tempo corrente.
\item il campo \var{sem\_otime}, che esprime il tempo dell'ultima operazione
  effettuata, viene inizializzato a zero.
\end{itemize*}

\begin{figure}[!htb]
  \footnotesize \centering
  \begin{minipage}[c]{.85\textwidth}
    \includestruct{listati/sem.h}
  \end{minipage} 
  \normalsize 
  \caption{La struttura \structd{sem}, che contiene i dati di un singolo
    semaforo.} 
  \label{fig:ipc_sem}
\end{figure}

Ciascun semaforo dell'insieme è realizzato come una struttura di tipo
\struct{sem} che ne contiene i dati essenziali, la cui definizione è riportata
in fig.~\ref{fig:ipc_sem}.\footnote{in realtà in fig~\ref{fig:ipc_sem} si è
  riportata la definizione originaria del kernel 1.0, che contiene la prima
  realizzazione del \textit{SysV-IPC} in Linux; ormai questa struttura è
  ridotta ai soli due primi membri, e gli altri vengono calcolati
  dinamicamente, la si è usata solo a scopo di esempio, perché indica tutti i
  valori associati ad un semaforo, restituiti dalle funzioni di controllo, e
  citati dalle pagine di manuale.}  Questa struttura non è accessibile
direttamente dallo \textit{user space}, ma i valori in essa specificati
possono essere letti in maniera indiretta, attraverso l'uso delle opportune
funzioni di controllo.  I dati mantenuti nella struttura, ed elencati in
fig.~\ref{fig:ipc_sem}, indicano rispettivamente:
\begin{basedescript}{\desclabelwidth{2cm}\desclabelstyle{\nextlinelabel}}
\item[\var{semval}] il valore numerico del semaforo.
\item[\var{sempid}] il \ids{PID} dell'ultimo processo che ha eseguito una
  operazione sul semaforo.
\item[\var{semncnt}] il numero di processi in attesa che esso venga
  incrementato.
\item[\var{semzcnt}] il numero di processi in attesa che esso si annulli.
\end{basedescript}

Come per le code di messaggi anche per gli insiemi di semafori esistono una
serie di limiti, i cui valori sono associati ad altrettante costanti, che si
sono riportate in tab.~\ref{tab:ipc_sem_limits}. Alcuni di questi limiti sono
al solito accessibili e modificabili attraverso \func{sysctl} o scrivendo
direttamente nel file \sysctlfile{kernel/sem}.

\begin{table}[htb]
  \footnotesize
  \centering
  \begin{tabular}[c]{|c|r|p{8cm}|}
    \hline
    \textbf{Costante} & \textbf{Valore} & \textbf{Significato} \\
    \hline
    \hline
    \constd{SEMMNI}&          128 & Numero massimo di insiemi di semafori.\\
    \constd{SEMMSL}&          250 & Numero massimo di semafori per insieme.\\
    \constd{SEMMNS}&\const{SEMMNI}*\const{SEMMSL}& Numero massimo di semafori
                                                   nel sistema.\\
    \constd{SEMVMX}&        32767 & Massimo valore per un semaforo.\\
    \constd{SEMOPM}&           32 & Massimo numero di operazioni per chiamata a
                                    \func{semop}. \\
    \constd{SEMMNU}&\const{SEMMNS}& Massimo numero di strutture di ripristino.\\
    \constd{SEMUME}&\const{SEMOPM}& Massimo numero di voci di ripristino.\\
    \constd{SEMAEM}&\const{SEMVMX}& Valore massimo per l'aggiustamento
                                    all'uscita. \\
    \hline
  \end{tabular}
  \caption{Valori delle costanti associate ai limiti degli insiemi di
    semafori, definite in \file{linux/sem.h}.} 
  \label{tab:ipc_sem_limits}
\end{table}


La funzione di sistema che permette di effettuare le varie operazioni di
controllo sui semafori fra le quali, come accennato, è impropriamente compresa
anche la loro inizializzazione, è \funcd{semctl}; il suo prototipo è:

\begin{funcproto}{
\fhead{sys/types.h}
\fhead{sys/ipc.h}
\fhead{sys/sem.h}
\fdecl{int semctl(int semid, int semnum, int cmd)}
\fdecl{int semctl(int semid, int semnum, int cmd, union semun arg)}
\fdesc{Esegue le operazioni di controllo su un semaforo o un insieme di
  semafori.}
}

{La funzione ritorna in caso di successo un valore positivo quanto usata con
  tre argomenti ed un valore nullo quando usata con quattro e $-1$ per un
  errore, nel qual caso \var{errno} assumerà uno dei valori:
  \begin{errlist}
  \item[\errcode{EACCES}] i permessi assegnati al semaforo non consentono
    l'operazione di lettura o scrittura richiesta e non si hanno i privilegi
    di amministratore.
    \item[\errcode{EIDRM}] l'insieme di semafori è stato cancellato.
    \item[\errcode{EPERM}] si è richiesto \const{IPC\_SET} o \const{IPC\_RMID}
      ma il processo non è né il creatore né il proprietario del semaforo e
      non ha i privilegi di amministratore.
    \item[\errcode{ERANGE}] si è richiesto \const{SETALL} \const{SETVAL} ma il
      valore a cui si vuole impostare il semaforo è minore di zero o maggiore
      di \const{SEMVMX}.
   \end{errlist}
   ed inoltre \errval{EFAULT} ed \errval{EINVAL} nel loro significato
   generico.}
\end{funcproto}

La funzione può avere tre o quattro argomenti, a seconda dell'operazione
specificata con \param{cmd}, ed opera o sull'intero insieme specificato da
\param{semid} o sul singolo semaforo di un insieme, specificato da
\param{semnum}. 

\begin{figure}[!htb]
  \footnotesize \centering
  \begin{minipage}[c]{0.80\textwidth}
    \includestruct{listati/semun.h}
  \end{minipage} 
  \normalsize 
  \caption{La definizione dei possibili valori di una \dirct{union}
    \structd{semun}, usata come quarto argomento della funzione
    \func{semctl}.}
  \label{fig:ipc_semun}
\end{figure}

Qualora la funzione operi con quattro argomenti \param{arg} è un argomento
generico, che conterrà un dato diverso a seconda dell'azione richiesta; per
unificare detto argomento esso deve essere passato come una unione
\struct{semun}, la cui definizione, con i possibili valori che può assumere, è
riportata in fig.~\ref{fig:ipc_semun}.

Nelle versioni più vecchie della \acr{glibc} questa unione veniva definita in
\file{sys/sem.h}, ma nelle versioni più recenti questo non avviene più in
quanto lo standard POSIX.1-2001 richiede che sia sempre definita a cura del
chiamante. In questa seconda evenienza la \acr{glibc} definisce però la
macro \macrod{\_SEM\_SEMUN\_UNDEFINED} che può essere usata per controllare la
situazione.

Come già accennato sia il comportamento della funzione che il numero di
argomenti con cui deve essere invocata dipendono dal valore dell'argomento
\param{cmd}, che specifica l'azione da intraprendere. Per questo argomento i
valori validi, quelli cioè che non causano un errore di \errcode{EINVAL}, sono
i seguenti:
\begin{basedescript}{\desclabelwidth{1.6cm}\desclabelstyle{\nextlinelabel}}
\item[\const{IPC\_STAT}] Legge i dati dell'insieme di semafori, copiandone i
  valori nella struttura \struct{semid\_ds} posta all'indirizzo specificato
  con \var{arg.buf}. Occorre avere il permesso di lettura.
  L'argomento \param{semnum} viene ignorato.
\item[\const{IPC\_RMID}] Rimuove l'insieme di semafori e le relative strutture
  dati, con effetto immediato. Tutti i processi che erano bloccati in attesa
  vengono svegliati, ritornando con un errore di \errcode{EIDRM}.  L'\ids{UID}
  effettivo del processo deve corrispondere o al creatore o al proprietario
  dell'insieme, o all'amministratore. L'argomento
  \param{semnum} viene ignorato.
\item[\const{IPC\_SET}] Permette di modificare i permessi ed il proprietario
  dell'insieme. I valori devono essere passati in una struttura
  \struct{semid\_ds} puntata da \param{arg.buf} di cui saranno usati soltanto
  i campi \var{sem\_perm.uid}, \var{sem\_perm.gid} e i nove bit meno
  significativi di \var{sem\_perm.mode}. La funziona aggiorna anche il campo
  \var{sem\_ctime}.  L'\ids{UID} effettivo del processo deve corrispondere o
  al creatore o al proprietario dell'insieme, o all'amministratore.
  L'argomento \param{semnum} viene ignorato.
\item[\constd{GETALL}] Restituisce il valore corrente di ciascun semaforo
  dell'insieme (corrispondente al campo \var{semval} di \struct{sem}) nel
  vettore indicato da \param{arg.array}. Occorre avere il permesso di lettura.
  L'argomento \param{semnum} viene ignorato.
\item[\constd{GETNCNT}] Restituisce come valore di ritorno della funzione il
  numero di processi in attesa che il semaforo \param{semnum} dell'insieme
  \param{semid} venga incrementato (corrispondente al campo \var{semncnt} di
  \struct{sem}). Va invocata con tre argomenti.  Occorre avere il permesso di
  lettura.
\item[\constd{GETPID}] Restituisce come valore di ritorno della funzione il
  \ids{PID} dell'ultimo processo che ha compiuto una operazione sul semaforo
  \param{semnum} dell'insieme \param{semid} (corrispondente al campo
  \var{sempid} di \struct{sem}). Va invocata con tre argomenti.  Occorre avere
  il permesso di lettura.
\item[\constd{GETVAL}] Restituisce come valore di ritorno della funzione il 
  valore corrente del semaforo \param{semnum} dell'insieme \param{semid}
  (corrispondente al campo \var{semval} di \struct{sem}). Va invocata con tre
  argomenti.  Occorre avere il permesso di lettura.
\item[\constd{GETZCNT}] Restituisce come valore di ritorno della funzione il
  numero di processi in attesa che il valore del semaforo \param{semnum}
  dell'insieme \param{semid} diventi nullo (corrispondente al campo
  \var{semncnt} di \struct{sem}). Va invocata con tre argomenti.  Occorre
  avere il permesso di lettura.
\item[\constd{SETALL}] Inizializza il valore di tutti i semafori dell'insieme,
  aggiornando il campo \var{sem\_ctime} di \struct{semid\_ds}. I valori devono
  essere passati nel vettore indicato da \param{arg.array}.  Si devono avere i
  privilegi di scrittura.  L'argomento \param{semnum} viene ignorato.
\item[\constd{SETVAL}] Inizializza il semaforo \param{semnum} al valore passato
  dall'argomento \param{arg.val}, aggiornando il campo \var{sem\_ctime} di
  \struct{semid\_ds}.  Si devono avere i privilegi di scrittura.
\end{basedescript}

Come per \func{msgctl} esistono tre ulteriori valori, \const{IPC\_INFO},
\constd{SEM\_STAT} e \constd{SEM\_INFO}, specifici di Linux e fuori da ogni
standard, creati specificamente ad uso del comando \cmd{ipcs}. Dato che anche
questi potranno essere modificati o rimossi, non devono essere utilizzati e
pertanto non li tratteremo.

Quando si imposta il valore di un semaforo (sia che lo si faccia per tutto
l'insieme con \const{SETALL}, che per un solo semaforo con \const{SETVAL}), i
processi in attesa su di esso reagiscono di conseguenza al cambiamento di
valore.  Inoltre la coda delle operazioni di ripristino viene cancellata per
tutti i semafori il cui valore viene modificato.

\begin{table}[htb]
  \footnotesize
  \centering
  \begin{tabular}[c]{|c|l|}
    \hline
    \textbf{Operazione}  & \textbf{Valore restituito} \\
    \hline
    \hline
    \const{GETNCNT}& Valore di \var{semncnt}.\\
    \const{GETPID} & Valore di \var{sempid}.\\
    \const{GETVAL} & Valore di \var{semval}.\\
    \const{GETZCNT}& Valore di \var{semzcnt}.\\
    \hline
  \end{tabular}
  \caption{Valori di ritorno della funzione \func{semctl}.} 
  \label{tab:ipc_semctl_returns}
\end{table}

Il valore di ritorno della funzione in caso di successo dipende
dall'operazione richiesta; per tutte le operazioni che richiedono quattro
argomenti esso è sempre nullo, per le altre operazioni, elencate in
tab.~\ref{tab:ipc_semctl_returns} viene invece restituito il valore richiesto,
corrispondente al campo della struttura \struct{sem} indicato nella seconda
colonna della tabella.

Le operazioni ordinarie sui semafori, come l'acquisizione o il rilascio degli
stessi (in sostanza tutte quelle non comprese nell'uso di \func{semctl})
vengono effettuate con la funzione di sistema \funcd{semop}, il cui prototipo
è:

\begin{funcproto}{
\fhead{sys/types.h}
\fhead{sys/ipc.h}
\fhead{sys/sem.h}
\fdecl{int semop(int semid, struct sembuf *sops, unsigned nsops)}
\fdesc{Esegue operazioni ordinarie su un semaforo o un insieme di semafori.} 
}

{La funzione ritorna $0$ in caso di successo e $-1$ per un errore, nel qual
  caso \var{errno} assumerà uno dei valori: 
  \begin{errlist}
    \item[\errcode{E2BIG}] l'argomento \param{nsops} è maggiore del numero
      massimo di operazioni \const{SEMOPM}.
    \item[\errcode{EACCES}] il processo non ha i permessi per eseguire
      l'operazione richiesta e non ha i privilegi di amministratore.
    \item[\errcode{EAGAIN}] un'operazione comporterebbe il blocco del processo,
      ma si è specificato \const{IPC\_NOWAIT} in \var{sem\_flg}.
    \item[\errcode{EFBIG}] il valore del campo \var{sem\_num} è negativo o
      maggiore o uguale al numero di semafori dell'insieme.
    \item[\errcode{EIDRM}] l'insieme di semafori è stato cancellato.
    \item[\errcode{EINTR}] la funzione, bloccata in attesa dell'esecuzione
      dell'operazione, viene interrotta da un segnale.
    \item[\errcode{ENOMEM}] si è richiesto un \const{SEM\_UNDO} ma il sistema
      non ha le risorse per allocare la struttura di ripristino.
    \item[\errcode{ERANGE}] per alcune operazioni il valore risultante del
      semaforo viene a superare il limite massimo \const{SEMVMX}.
  \end{errlist}
  ed inoltre \errval{EFAULT} ed \errval{EINVAL} nel loro significato
  generico.}
\end{funcproto}

La funzione permette di eseguire operazioni multiple sui singoli semafori di
un insieme. La funzione richiede come primo argomento l'identificatore
\param{semid} dell'insieme su cui si vuole operare, il numero di operazioni da
effettuare viene specificato con l'argomento \param{nsop}, mentre il loro
contenuto viene passato con un puntatore ad un vettore di strutture
\struct{sembuf} nell'argomento \param{sops}. Le operazioni richieste vengono
effettivamente eseguite se e soltanto se è possibile effettuarle tutte quante,
ed in tal caso vengono eseguite nella sequenza passata nel
vettore \param{sops}.

Con lo standard POSIX.1-2001 è stata introdotta una variante di \func{semop}
che consente di specificare anche un tempo massimo di attesa. La nuova
funzione di sistema, disponibile a partire dal kernel 2.4.22 e dalla
\acr{glibc} 2.3.3, ed utilizzabile solo dopo aver definito la macro
\macro{\_GNU\_SOURCE}, è \funcd{semtimedop}, ed il suo prototipo è:

\begin{funcproto}{
\fhead{sys/types.h}
\fhead{sys/ipc.h}
\fhead{sys/sem.h}
\fdecl{int semtimedop(int semid, struct sembuf *sops, unsigned nsops,
                      struct timespec *timeout)}
\fdesc{Esegue operazioni ordinarie su un semaforo o un insieme di semafori.} 
}

{La funzione ritorna $0$ in caso di successo e $-1$ per un errore, nel qual
  caso \var{errno} assumerà uno dei valori: 
  \begin{errlist}
  \item[\errcode{EAGAIN}] l'operazione comporterebbe il blocco del processo,
    ma si è specificato \const{IPC\_NOWAIT} in \var{sem\_flg} oppure si è
    atteso oltre quanto indicato da \param{timeout}.
  \end{errlist}
  e gli altri valori già visti per \func{semop}, con lo stesso significato.}
\end{funcproto}

Rispetto a \func{semop} la funzione consente di specificare un tempo massimo
di attesa, indicato con una struttura \struct{timespec} (vedi
fig.~\ref{fig:sys_timespec_struct}), per le operazioni che verrebbero
bloccate. Alla scadenza di detto tempo la funzione ritorna comunque con un
errore di \errval{EAGAIN} senza che nessuna delle operazioni richieste venga
eseguita. 

Si tenga presente che la precisione della temporizzazione è comunque limitata
dalla risoluzione dell'orologio di sistema, per cui il tempo di attesa verrà
arrotondato per eccesso. In caso si passi un valore \val{NULL}
per \param{timeout} il comportamento di \func{semtimedop} è identico a quello
di \func{semop}.


\begin{figure}[!htb]
  \footnotesize \centering
  \begin{minipage}[c]{.80\textwidth}
    \includestruct{listati/sembuf.h}
  \end{minipage} 
  \normalsize 
  \caption{La struttura \structd{sembuf}, usata per le operazioni sui
    semafori.}
  \label{fig:ipc_sembuf}
\end{figure}

Come indicato il contenuto di ciascuna operazione deve essere specificato
attraverso una struttura \struct{sembuf} (la cui definizione è riportata in
fig.~\ref{fig:ipc_sembuf}) che il programma chiamante deve avere cura di
allocare in un opportuno vettore. La struttura permette di indicare il
semaforo su cui operare, il tipo di operazione, ed un flag di controllo.  

Il campo \var{sem\_num} serve per indicare a quale semaforo dell'insieme fa
riferimento l'operazione. Si ricordi che i semafori sono numerati come gli
elementi di un vettore, per cui il primo semaforo di un insieme corrisponde ad
un valore nullo di \var{sem\_num}.

Il campo \var{sem\_flg} è un flag, mantenuto come maschera binaria, per il
quale possono essere impostati i due valori \const{IPC\_NOWAIT} e
\constd{SEM\_UNDO}.  Impostando \const{IPC\_NOWAIT} si fa sì che in tutti quei
casi in cui l'esecuzione di una operazione richiederebbe di porre il processo
vada nello stato di \textit{sleep}, invece di bloccarsi \func{semop} ritorni
immediatamente (abortendo così le eventuali operazioni restanti) con un errore
di \errcode{EAGAIN}.  Impostando \const{SEM\_UNDO} si richiede invece che
l'operazione in questione venga registrata, in modo che il valore del semaforo
possa essere ripristinato all'uscita del processo.

Infine \var{sem\_op} è il campo che controlla qual'è l'operazione che viene
eseguita e determina in generale il comportamento della chiamata a
\func{semop}. I casi possibili per il valore di questo campo sono tre:
\begin{basedescript}{\desclabelwidth{1.8cm}}
\item[\var{sem\_op} $>0$] In questo caso il valore viene aggiunto al valore
  corrente di \var{semval} per il semaforo indicato. Questa operazione non
  causa mai un blocco del processo, ed eventualmente \func{semop} ritorna
  immediatamente con un errore di \errcode{ERANGE} qualora si sia superato il
  limite \const{SEMVMX}. Se l'operazione ha successo si passa immediatamente
  alla successiva.  Specificando \const{SEM\_UNDO} si aggiorna il contatore
  per il ripristino del valore del semaforo. Al processo chiamante è richiesto
  il privilegio di alterazione (scrittura) sull'insieme di semafori.
  
\item[\var{sem\_op} $=0$] Nel caso \var{semval} sia zero l'operazione ha
  successo immediato, e o si passa alla successiva o \func{semop} ritorna con
  successo se questa era l'ultima. Se \var{semval} è diverso da zero il
  comportamento è controllato da \var{sem\_flg}, se è stato impostato
  \const{IPC\_NOWAIT} \func{semop} ritorna immediatamente abortendo tutte le
  operazioni con un errore di \errcode{EAGAIN}, altrimenti viene incrementato
  \var{semzcnt} di uno ed il processo viene bloccato fintanto che non si
  verifica una delle condizioni seguenti:
  \begin{itemize*}
  \item \var{semval} diventa zero, nel qual caso \var{semzcnt} viene
    decrementato di uno, l'operazione ha successo e si passa alla successiva,
    oppure \func{semop} ritorna con successo se questa era l'ultima.
  \item l'insieme di semafori viene rimosso, nel qual caso \func{semop}
    ritorna abortendo tutte le operazioni con un errore di \errcode{EIDRM}.
  \item il processo chiamante riceve un segnale, nel qual caso \var{semzcnt}
    viene decrementato di uno e \func{semop} ritorna abortendo tutte le
    operazioni con un errore di \errcode{EINTR}.
  \end{itemize*}
  Al processo chiamante è richiesto soltanto il privilegio di lettura
  dell'insieme dei semafori.
  
\item[\var{sem\_op} $<0$] Nel caso in cui \var{semval} è maggiore o uguale del
  valore assoluto di \var{sem\_op} (se cioè la somma dei due valori resta
  positiva o nulla) i valori vengono sommati e l'operazione ha successo e si
  passa alla successiva, oppure \func{semop} ritorna con successo se questa
  era l'ultima. Qualora si sia impostato \const{SEM\_UNDO} viene anche
  aggiornato il contatore per il ripristino del valore del semaforo. In caso
  contrario (quando cioè la somma darebbe luogo ad un valore di \var{semval}
  negativo) se si è impostato \const{IPC\_NOWAIT} \func{semop} ritorna
  immediatamente abortendo tutte le operazioni con un errore di
  \errcode{EAGAIN}, altrimenti viene incrementato di uno \var{semncnt} ed il
  processo resta in stato di \textit{sleep} fintanto che non si ha una delle
  condizioni seguenti:
  \begin{itemize*}
  \item \var{semval} diventa maggiore o uguale del valore assoluto di
    \var{sem\_op}, nel qual caso \var{semncnt} viene decrementato di uno, il
    valore di \var{sem\_op} viene sommato a \var{semval}, e se era stato
    impostato \const{SEM\_UNDO} viene aggiornato il contatore per il
    ripristino del valore del semaforo.
  \item l'insieme di semafori viene rimosso, nel qual caso \func{semop}
    ritorna abortendo tutte le operazioni con un errore di \errcode{EIDRM}.
  \item il processo chiamante riceve un segnale, nel qual caso \var{semncnt}
    viene decrementato di uno e \func{semop} ritorna abortendo tutte le
    operazioni con un errore di \errcode{EINTR}.
  \end{itemize*}    
  Al processo chiamante è richiesto il privilegio di alterazione (scrittura)
  sull'insieme di semafori.
\end{basedescript}

Qualora si sia usato \func{semtimedop} alle condizioni di errore precedenti si
aggiunge anche quella di scadenza del tempo di attesa indicato
con \param{timeout} che farà abortire la funzione, qualora resti bloccata
troppo a lungo nell'esecuzione delle operazioni richieste, con un errore di
\errcode{EAGAIN}.

In caso di successo (sia per \func{semop} che per \func{semtimedop}) per ogni
semaforo modificato verrà aggiornato il campo \var{sempid} al valore del
\ids{PID} del processo chiamante; inoltre verranno pure aggiornati al tempo
corrente i campi \var{sem\_otime} e \var{sem\_ctime}.

Dato che, come già accennato in precedenza, in caso di uscita inaspettata i
semafori possono restare occupati, abbiamo visto come \func{semop} (e
\func{semtimedop}) permetta di attivare un meccanismo di ripristino attraverso
l'uso del flag \const{SEM\_UNDO}. Il meccanismo è implementato tramite una
apposita struttura \kstruct{sem\_undo}, associata ad ogni processo per ciascun
semaforo che esso ha modificato; all'uscita i semafori modificati vengono
ripristinati, e le strutture disallocate.  Per mantenere coerente il
comportamento queste strutture non vengono ereditate attraverso una
\func{fork} (altrimenti si avrebbe un doppio ripristino), mentre passano
inalterate nell'esecuzione di una \func{exec} (altrimenti non si avrebbe
ripristino).

Tutto questo però ha un problema di fondo. Per capire di cosa si tratta
occorre fare riferimento all'implementazione usata in Linux, che è riportata
in maniera semplificata nello schema di fig.~\ref{fig:ipc_sem_schema}.  Si è
presa come riferimento l'architettura usata fino al kernel 2.2.x che è più
semplice (ed illustrata in dettaglio in \cite{tlk}). Nel kernel 2.4.x la
struttura del \textit{SysV-IPC} è stata modificata, ma le definizioni relative
a queste strutture restano per compatibilità (in particolare con le vecchie
versioni delle librerie del C, come le \acr{libc5}).

\begin{figure}[!htb]
  \centering \includegraphics[width=12cm]{img/semtruct}
  \caption{Schema delle varie strutture di un insieme di semafori
    (\kstructd{semid\_ds}, \kstructd{sem}, \kstructd{sem\_queue} e
    \kstructd{sem\_undo}).}
  \label{fig:ipc_sem_schema}
\end{figure}

Alla creazione di un nuovo insieme viene allocata una nuova strutture
\kstruct{semid\_ds} ed il relativo vettore di strutture \kstruct{sem}. Quando
si richiede una operazione viene anzitutto verificato che tutte le operazioni
possono avere successo; se una di esse comporta il blocco del processo il
kernel crea una struttura \kstruct{sem\_queue} che viene aggiunta in fondo
alla coda di attesa associata a ciascun insieme di semafori, che viene
referenziata tramite i campi \var{sem\_pending} e \var{sem\_pending\_last} di
\kstruct{semid\_ds}.  Nella struttura viene memorizzato il riferimento alle
operazioni richieste (nel campo \var{sops}, che è un puntatore ad una
struttura \struct{sembuf}) e al processo corrente (nel campo \var{sleeper})
poi quest'ultimo viene messo stato di attesa e viene invocato lo
\textit{scheduler} per passare all'esecuzione di un altro processo.

Se invece tutte le operazioni possono avere successo vengono eseguite
immediatamente, dopo di che il kernel esegue una scansione della coda di
attesa (a partire da \var{sem\_pending}) per verificare se qualcuna delle
operazioni sospese in precedenza può essere eseguita, nel qual caso la
struttura \kstruct{sem\_queue} viene rimossa e lo stato del processo associato
all'operazione (\var{sleeper}) viene riportato a \textit{running}; il tutto
viene ripetuto fin quando non ci sono più operazioni eseguibili o si è
svuotata la coda.  Per gestire il meccanismo del ripristino tutte le volte che
per un'operazione si è specificato il flag \const{SEM\_UNDO} viene mantenuta
per ciascun insieme di semafori una apposita struttura \kstruct{sem\_undo} che
contiene (nel vettore puntato dal campo \var{semadj}) un valore di
aggiustamento per ogni semaforo cui viene sommato l'opposto del valore usato
per l'operazione.

Queste strutture sono mantenute in due liste (rispettivamente attraverso i due
campi \var{id\_next} e \var{proc\_next}) una associata all'insieme di cui fa
parte il semaforo, che viene usata per invalidare le strutture se questo viene
cancellato o per azzerarle se si è eseguita una operazione con \func{semctl},
l'altra associata al processo che ha eseguito l'operazione, attraverso il
campo \var{semundo} di \kstruct{task\_struct}, come mostrato in
\ref{fig:ipc_sem_schema}. Quando un processo termina, la lista ad esso
associata viene scandita e le operazioni applicate al semaforo.  Siccome un
processo può accumulare delle richieste di ripristino per semafori differenti
attraverso diverse chiamate a \func{semop}, si pone il problema di come
eseguire il ripristino dei semafori all'uscita del processo, ed in particolare
se questo può essere fatto atomicamente.

Il punto è cosa succede quando una delle operazioni previste per il ripristino
non può essere eseguita immediatamente perché ad esempio il semaforo è
occupato; in tal caso infatti, se si pone il processo in stato di
\textit{sleep} aspettando la disponibilità del semaforo (come faceva
l'implementazione originaria) si perde l'atomicità dell'operazione. La scelta
fatta dal kernel è pertanto quella di effettuare subito le operazioni che non
prevedono un blocco del processo e di ignorare silenziosamente le altre;
questo però comporta il fatto che il ripristino non è comunque garantito in
tutte le occasioni.

Come esempio di uso dell'interfaccia dei semafori vediamo come implementare
con essa dei semplici \textit{mutex} (cioè semafori binari), tutto il codice
in questione, contenuto nel file \file{Mutex.c} allegato ai sorgenti, è
riportato in fig.~\ref{fig:ipc_mutex_create}. Utilizzeremo l'interfaccia per
creare un insieme contenente un singolo semaforo, per il quale poi useremo un
valore unitario per segnalare la disponibilità della risorsa, ed un valore
nullo per segnalarne l'indisponibilità. 

\begin{figure}[!htbp]
  \footnotesize \centering
  \begin{minipage}[c]{\codesamplewidth}
    \includecodesample{listati/Mutex.c}
  \end{minipage} 
  \normalsize 
  \caption{Il codice delle funzioni che permettono di creare o recuperare
    l'identificatore di un semaforo da utilizzare come \textit{mutex}.}
  \label{fig:ipc_mutex_create}
\end{figure}

La prima funzione (\texttt{\small 2-15}) è \func{MutexCreate} che data una
chiave crea il semaforo usato per il mutex e lo inizializza, restituendone
l'identificatore. Il primo passo (\texttt{\small 6}) è chiamare \func{semget}
con \const{IPC\_CREATE} per creare il semaforo qualora non esista,
assegnandogli i privilegi di lettura e scrittura per tutti. In caso di errore
(\texttt{\small 7-9}) si ritorna subito il risultato di \func{semget},
altrimenti (\texttt{\small 10}) si inizializza il semaforo chiamando
\func{semctl} con il comando \const{SETVAL}, utilizzando l'unione
\struct{semunion} dichiarata ed avvalorata in precedenza (\texttt{\small 4})
ad 1 per significare che risorsa è libera. In caso di errore (\texttt{\small
  11-13}) si restituisce il valore di ritorno di \func{semctl}, altrimenti
(\texttt{\small 14}) si ritorna l'identificatore del semaforo.

La seconda funzione (\texttt{\small 17-20}) è \func{MutexFind}, che, data una
chiave, restituisce l'identificatore del semaforo ad essa associato. La
comprensione del suo funzionamento è immediata in quanto essa è soltanto un
\textit{wrapper}\footnote{si chiama così una funzione usata per fare da
  \textsl{involucro} alla chiamata di un altra, usata in genere per
  semplificare un'interfaccia (come in questo caso) o per utilizzare con la
  stessa funzione diversi substrati (librerie, ecc.)  che possono fornire le
  stesse funzionalità.} di una chiamata a \func{semget} per cercare
l'identificatore associato alla chiave, il valore di ritorno di quest'ultima
viene passato all'indietro al chiamante.

La terza funzione (\texttt{\small 22-25}) è \func{MutexRead} che, dato un
identificatore, restituisce il valore del semaforo associato al mutex. Anche
in questo caso la funzione è un \textit{wrapper} per una chiamata a
\func{semctl} con il comando \const{GETVAL}, che permette di restituire il
valore del semaforo.

La quarta e la quinta funzione (\texttt{\small 36-44}) sono \func{MutexLock},
e \func{MutexUnlock}, che permettono rispettivamente di bloccare e sbloccare
il mutex. Entrambe fanno da wrapper per \func{semop}, utilizzando le due
strutture \var{sem\_lock} e \var{sem\_unlock} definite in precedenza
(\texttt{\small 27-34}). Si noti come per queste ultime si sia fatto uso
dell'opzione \const{SEM\_UNDO} per evitare che il semaforo resti bloccato in
caso di terminazione imprevista del processo.

L'ultima funzione (\texttt{\small 46-49}) della serie, è \func{MutexRemove},
che rimuove il mutex. Anche in questo caso si ha un wrapper per una chiamata a
\func{semctl} con il comando \const{IPC\_RMID}, che permette di cancellare il
semaforo; il valore di ritorno di quest'ultima viene passato all'indietro.

Chiamare \func{MutexLock} decrementa il valore del semaforo: se questo è
libero (ha già valore 1) sarà bloccato (valore nullo), se è bloccato la
chiamata a \func{semop} si bloccherà fintanto che la risorsa non venga
rilasciata. Chiamando \func{MutexUnlock} il valore del semaforo sarà
incrementato di uno, sbloccandolo qualora fosse bloccato.  

Si noti che occorre eseguire sempre prima \func{MutexLock} e poi
\func{MutexUnlock}, perché se per un qualche errore si esegue più volte
quest'ultima il valore del semaforo crescerebbe oltre 1, e \func{MutexLock}
non avrebbe più l'effetto aspettato (bloccare la risorsa quando questa è
considerata libera).  Infine si tenga presente che usare \func{MutexRead} per
controllare il valore dei mutex prima di proseguire in una operazione di
sblocco non servirebbe comunque, dato che l'operazione non sarebbe atomica.
Vedremo in sez.~\ref{sec:ipc_lock_file} come sia possibile ottenere
un'interfaccia analoga a quella appena illustrata, senza incorrere in questi
problemi, usando il \textit{file locking}.


\subsection{Memoria condivisa}
\label{sec:ipc_sysv_shm}

Il terzo oggetto introdotto dal \textit{SysV-IPC} è quello dei segmenti di
memoria condivisa. La funzione di sistema che permette di ottenerne uno è
\funcd{shmget}, ed il suo prototipo è:

\begin{funcproto}{
\fhead{sys/types.h}
\fhead{sys/ipc.h}
\fhead{sys/shm.h}
\fdecl{int shmget(key\_t key, int size, int flag)}
\fdesc{Ottiene o crea una memoria condivisa.} 
}

{La funzione ritorna l'identificatore (un intero positivo) in caso di successo
  e $-1$ per un errore, nel qual caso \var{errno} assumerà uno dei valori:
  \begin{errlist}
    \item[\errcode{ENOSPC}] si è superato il limite (\const{SHMMNI}) sul numero
      di segmenti di memoria nel sistema, o cercato di allocare un segmento le
      cui dimensioni fanno superare il limite di sistema (\const{SHMALL}) per
      la memoria ad essi riservata.
    \item[\errcode{EINVAL}] si è richiesta una dimensione per un nuovo segmento
      maggiore di \const{SHMMAX} o minore di \const{SHMMIN}, o se il segmento
      già esiste \param{size} è maggiore delle sue dimensioni.
    \item[\errcode{ENOMEM}] il sistema non ha abbastanza memoria per poter
      contenere le strutture per un nuovo segmento di memoria condivisa.
    \item[\errcode{ENOMEM}] si è specificato \const{IPC\_HUGETLB} ma non si
      hanno i privilegi di amministratore.
   \end{errlist}
   ed inoltre \errval{EACCES}, \errval{ENOENT}, \errval{EEXIST},
   \errval{EIDRM}, con lo stesso significato che hanno per \func{msgget}.}
\end{funcproto}


La funzione, come \func{semget}, è analoga a \func{msgget}, ed identico è
l'uso degli argomenti \param{key} e \param{flag} per cui non ripeteremo quanto
detto al proposito in sez.~\ref{sec:ipc_sysv_mq}.  A partire dal kernel 2.6
però sono stati introdotti degli ulteriori bit di controllo per
l'argomento \param{flag}, specifici di \func{shmget}, attinenti alle modalità
di gestione del segmento di memoria condivisa in relazione al sistema della
memoria virtuale.

Il primo dei due flag è \constd{SHM\_HUGETLB} che consente di richiedere la
creazione del segmento usando una \textit{huge page}, le pagine di memoria di
grandi dimensioni introdotte con il kernel 2.6 per ottimizzare le prestazioni
nei sistemi più recenti che hanno grandi quantità di memoria. L'operazione è
privilegiata e richiede che il processo abbia la \textit{capability}
\const{CAP\_IPC\_LOCK}. Questa funzionalità è specifica di Linux e non è
portabile.

Il secondo flag aggiuntivo, introdotto a partire dal kernel 2.6.15, è
\constd{SHM\_NORESERVE}, ed ha lo stesso scopo del flag \const{MAP\_NORESERVE}
di \func{mmap} (vedi sez.~\ref{sec:file_memory_map}): non vengono riservate
delle pagine di swap ad uso del meccanismo del \textit{copy on write} per
mantenere le modifiche fatte sul segmento. Questo significa che caso di
scrittura sul segmento quando non c'è più memoria disponibile, si avrà
l'emissione di un \signal{SIGSEGV}.

Infine l'argomento \param{size} specifica la dimensione del segmento di
memoria condivisa; il valore deve essere specificato in byte, ma verrà
comunque arrotondato al multiplo superiore di \const{PAGE\_SIZE}. Il valore
deve essere specificato quando si crea un nuovo segmento di memoria con
\const{IPC\_CREAT} o \const{IPC\_PRIVATE}, se invece si accede ad un segmento
di memoria condivisa esistente non può essere maggiore del valore con cui esso
è stato creato.

La memoria condivisa è la forma più veloce di comunicazione fra due processi,
in quanto permette agli stessi di vedere nel loro spazio di indirizzi una
stessa sezione di memoria.  Pertanto non è necessaria nessuna operazione di
copia per trasmettere i dati da un processo all'altro, in quanto ciascuno può
accedervi direttamente con le normali operazioni di lettura e scrittura dei
dati in memoria.

Ovviamente tutto questo ha un prezzo, ed il problema fondamentale della
memoria condivisa è la sincronizzazione degli accessi. È evidente infatti che
se un processo deve scambiare dei dati con un altro, si deve essere sicuri che
quest'ultimo non acceda al segmento di memoria condivisa prima che il primo
non abbia completato le operazioni di scrittura, inoltre nel corso di una
lettura si deve essere sicuri che i dati restano coerenti e non vengono
sovrascritti da un accesso in scrittura sullo stesso segmento da parte di un
altro processo. Per questo in genere la memoria condivisa viene sempre
utilizzata in abbinamento ad un meccanismo di sincronizzazione, il che, di
norma, significa insieme a dei semafori.

\begin{figure}[!htb]
  \footnotesize \centering
  \begin{minipage}[c]{0.80\textwidth}
    \includestruct{listati/shmid_ds.h}
  \end{minipage} 
  \normalsize 
  \caption{La struttura \structd{shmid\_ds}, associata a ciascun segmento di
    memoria condivisa.}
  \label{fig:ipc_shmid_ds}
\end{figure}

A ciascun segmento di memoria condivisa è associata una struttura
\struct{shmid\_ds}, riportata in fig.~\ref{fig:ipc_shmid_ds}.  Come nel caso
delle code di messaggi quando si crea un nuovo segmento di memoria condivisa
con \func{shmget} questa struttura viene inizializzata, in particolare il
campo \var{shm\_perm} viene inizializzato come illustrato in
sez.~\ref{sec:ipc_sysv_access_control}, e valgono le considerazioni ivi fatte
relativamente ai permessi di accesso; per quanto riguarda gli altri campi
invece:
\begin{itemize*}
\item il campo \var{shm\_segsz}, che esprime la dimensione del segmento, viene
  inizializzato al valore di \param{size}.
\item il campo \var{shm\_ctime}, che esprime il tempo di creazione del
  segmento, viene inizializzato al tempo corrente.
\item i campi \var{shm\_atime} e \var{shm\_dtime}, che esprimono
  rispettivamente il tempo dell'ultima volta che il segmento è stato
  agganciato o sganciato da un processo, vengono inizializzati a zero.
\item il campo \var{shm\_lpid}, che esprime il \ids{PID} del processo che ha
  eseguito l'ultima operazione, viene inizializzato a zero.
\item il campo \var{shm\_cpid}, che esprime il \ids{PID} del processo che ha
  creato il segmento, viene inizializzato al \ids{PID} del processo chiamante.
\item il campo \var{shm\_nattac}, che esprime il numero di processi agganciati
  al segmento viene inizializzato a zero.
\end{itemize*}

Come per le code di messaggi e gli insiemi di semafori, anche per i segmenti
di memoria condivisa esistono una serie di limiti imposti dal sistema.  Alcuni
di questi limiti sono al solito accessibili e modificabili attraverso
\func{sysctl} o scrivendo direttamente nei rispettivi file di
\file{/proc/sys/kernel/}. 

In tab.~\ref{tab:ipc_shm_limits} si sono riportate le
costanti simboliche associate a ciascuno di essi, il loro significato, i
valori preimpostati, e, quando presente, il file in \file{/proc/sys/kernel/}
che permettono di cambiarne il valore. 


\begin{table}[htb]
  \footnotesize
  \centering
  \begin{tabular}[c]{|c|r|c|p{7cm}|}
    \hline
    \textbf{Costante} & \textbf{Valore} & \textbf{File in \texttt{proc}}
    & \textbf{Significato} \\
    \hline
    \hline
    \constd{SHMALL}& 0x200000&\sysctlrelfiled{kernel}{shmall}
                             & Numero massimo di pagine che 
                               possono essere usate per i segmenti di
                               memoria condivisa.\\
    \constd{SHMMAX}&0x2000000&\sysctlrelfiled{kernel}{shmmax} 
                             & Dimensione massima di un segmento di memoria
                               condivisa.\\ 
    \constd{SHMMNI}&     4096&\sysctlrelfiled{kernel}{shmmni}
                             & Numero massimo di segmenti di memoria condivisa
                              presenti nel kernel.\\ 
    \constd{SHMMIN}&        1& ---         & Dimensione minima di un segmento di
                                             memoria condivisa.\\
    \constd{SHMLBA}&\const{PAGE\_SIZE}&--- & Limite inferiore per le dimensioni
                                             minime di un segmento (deve essere
                                             allineato alle dimensioni di una
                                             pagina di memoria).\\
    \constd{SHMSEG}&   ---   &     ---     & Numero massimo di segmenti di
                                             memoria condivisa per ciascun
                                             processo (l'implementazione non
                                             prevede l'esistenza di questo
                                             limite).\\


    \hline
  \end{tabular}
  \caption{Valori delle costanti associate ai limiti dei segmenti di memoria
    condivisa, insieme al relativo file in \file{/proc/sys/kernel/} ed al
    valore preimpostato presente nel sistema.} 
  \label{tab:ipc_shm_limits}
\end{table}

Al solito la funzione di sistema che permette di effettuare le operazioni di
controllo su un segmento di memoria condivisa è \funcd{shmctl}; il suo
prototipo è:

\begin{funcproto}{
\fhead{sys/ipc.h}
\fhead{sys/shm.h}
\fdecl{int shmctl(int shmid, int cmd, struct shmid\_ds *buf)}

\fdesc{Esegue le operazioni di controllo su un segmento di memoria condivisa.}
}

{La funzione ritorna $0$ in caso di successo e $-1$ per un errore, nel qual
  caso \var{errno} assumerà uno dei valori:
  \begin{errlist}
    \item[\errcode{EACCES}] si è richiesto \const{IPC\_STAT} ma i permessi non
      consentono l'accesso in lettura al segmento.
    \item[\errcode{EFAULT}] l'indirizzo specificato con \param{buf} non è
      valido.
    \item[\errcode{EIDRM}] l'argomento \param{shmid} fa riferimento ad un
      segmento che è stato cancellato.
    \item[\errcode{EINVAL}] o \param{shmid} non è un identificatore valido o
      \param{cmd} non è un comando valido.
    \item[\errcode{ENOMEM}] si è richiesto un \textit{memory lock} di
      dimensioni superiori al massimo consentito.
    \item[\errcode{EOVERFLOW}] si è tentato il comando \const{IPC\_STAT} ma il
      valore del \ids{GID} o dell'\ids{UID} è troppo grande per essere
      memorizzato nella struttura puntata da \param{buf}.
    \item[\errcode{EPERM}] si è specificato un comando con \const{IPC\_SET} o
      \const{IPC\_RMID} senza i permessi necessari.
  \end{errlist}
}  
\end{funcproto}

Il comando specificato attraverso l'argomento \param{cmd} determina i diversi
effetti della funzione. Nello standard POSIX.1-2001 i valori che esso può
assumere, ed il corrispondente comportamento della funzione, sono i seguenti:

\begin{basedescript}{\desclabelwidth{2.2cm}\desclabelstyle{\nextlinelabel}}
\item[\const{IPC\_STAT}] Legge le informazioni riguardo il segmento di memoria
  condivisa nella struttura \struct{shmid\_ds} puntata da \param{buf}. Occorre
  che il processo chiamante abbia il permesso di lettura sulla segmento.
\item[\const{IPC\_RMID}] Marca il segmento di memoria condivisa per la
  rimozione, questo verrà cancellato effettivamente solo quando l'ultimo
  processo ad esso agganciato si sarà staccato. Questo comando può essere
  eseguito solo da un processo con \ids{UID} effettivo corrispondente o al
  creatore del segmento, o al proprietario del segmento, o all'amministratore.
\item[\const{IPC\_SET}] Permette di modificare i permessi ed il proprietario
  del segmento.  Per modificare i valori di \var{shm\_perm.mode},
  \var{shm\_perm.uid} e \var{shm\_perm.gid} occorre essere il proprietario o
  il creatore del segmento, oppure l'amministratore. Compiuta l'operazione
  aggiorna anche il valore del campo \var{shm\_ctime}.
\end{basedescript}

Oltre ai precedenti su Linux sono definiti anche degli ulteriori comandi, che
consentono di estendere le funzionalità, ovviamente non devono essere usati se
si ha a cuore la portabilità. Questi comandi aggiuntivi sono:

\begin{basedescript}{\desclabelwidth{2.2cm}\desclabelstyle{\nextlinelabel}}
\item[\constd{SHM\_LOCK}] Abilita il \textit{memory locking} sul segmento di
  memoria condivisa, impedendo che la memoria usata per il segmento venga
  salvata su disco dal meccanismo della memoria virtuale. Come illustrato in
  sez.~\ref{sec:proc_mem_lock} fino al kernel 2.6.9 solo l'amministratore
  poteva utilizzare questa capacità,\footnote{che richiedeva la
    \textit{capability} \const{CAP\_IPC\_LOCK}.} a partire dal kernel 2.6.10
  anche gli utenti normali possono farlo fino al limite massimo determinato da
  \const{RLIMIT\_MEMLOCK} (vedi sez.~\ref{sec:sys_resource_limit}).
\item[\constd{SHM\_UNLOCK}] Disabilita il \textit{memory locking} sul segmento
  di memoria condivisa.  Fino al kernel 2.6.9 solo l'amministratore poteva
  utilizzare questo comando in corrispondenza di un segmento da lui bloccato.
\end{basedescript}

A questi due, come per \func{msgctl} e \func{semctl}, si aggiungono tre
ulteriori valori, \const{IPC\_INFO}, \constd{SHM\_STAT} e \constd{SHM\_INFO},
introdotti ad uso del programma \cmd{ipcs} per ottenere le informazioni
generali relative alle risorse usate dai segmenti di memoria condivisa. Dato
che potranno essere modificati o rimossi in favore dell'uso di \texttt{/proc},
non devono essere usati e non li tratteremo.

L'argomento \param{buf} viene utilizzato solo con i comandi \const{IPC\_STAT}
e \const{IPC\_SET} nel qual caso esso dovrà puntare ad una struttura
\struct{shmid\_ds} precedentemente allocata, in cui nel primo caso saranno
scritti i dati del segmento di memoria restituiti dalla funzione e da cui, nel
secondo caso, verranno letti i dati da impostare sul segmento.

Una volta che lo si è creato, per utilizzare un segmento di memoria condivisa
l'interfaccia prevede due funzioni, \funcd{shmat} e \func{shmdt}. La prima di
queste serve ad agganciare un segmento al processo chiamante, in modo che
quest'ultimo possa inserirlo nel suo spazio di indirizzi per potervi accedere;
il suo prototipo è:

\begin{funcproto}{
\fhead{sys/types.h} 
\fhead{sys/shm.h}
\fdecl{void *shmat(int shmid, const void *shmaddr, int shmflg)}

\fdesc{Aggancia un segmento di memoria condivisa al processo chiamante.}
}

{La funzione ritorna l'indirizzo del segmento in caso di successo e $-1$ (in
  un cast a \ctyp{void *}) per un errore, nel qual caso \var{errno} assumerà
  uno dei valori:
  \begin{errlist}
    \item[\errcode{EACCES}] il processo non ha i privilegi per accedere al
      segmento nella modalità richiesta.
    \item[\errcode{EINVAL}] si è specificato un identificatore invalido per
      \param{shmid}, o un indirizzo non allineato sul confine di una pagina
      per \param{shmaddr} o il valore \val{NULL} indicando \const{SHM\_REMAP}.
  \end{errlist}
  ed inoltre \errval{ENOMEM} nel suo significato generico.
}  
\end{funcproto}

La funzione inserisce un segmento di memoria condivisa all'interno dello
spazio di indirizzi del processo, in modo che questo possa accedervi
direttamente, la situazione dopo l'esecuzione di \func{shmat} è illustrata in
fig.~\ref{fig:ipc_shmem_layout} (per la comprensione del resto dello schema si
ricordi quanto illustrato al proposito in sez.~\ref{sec:proc_mem_layout}). In
particolare l'indirizzo finale del segmento dati (quello impostato da
\func{brk}, vedi sez.~\ref{sec:proc_mem_alloc}) non viene influenzato.
Si tenga presente infine che la funzione ha successo anche se il segmento è
stato marcato per la cancellazione.

\begin{figure}[!htb]
  \centering \includegraphics[height=10cm]{img/sh_memory_layout}
  \caption{Disposizione dei segmenti di memoria di un processo quando si è
    agganciato un segmento di memoria condivisa.}
  \label{fig:ipc_shmem_layout}
\end{figure}

L'argomento \param{shmaddr} specifica a quale indirizzo\footnote{lo standard
  SVID prevede che l'argomento \param{shmaddr} sia di tipo \ctyp{char *}, così
  come il valore di ritorno della funzione; in Linux è stato così con la
  \acr{libc4} e la \acr{libc5}, con il passaggio alla \acr{glibc} il tipo di
  \param{shmaddr} è divenuto un \ctyp{const void *} e quello del valore di
  ritorno un \ctyp{void *} seguendo POSIX.1-2001.} deve essere associato il
segmento, se il valore specificato è \val{NULL} è il sistema a scegliere
opportunamente un'area di memoria libera (questo è il modo più portabile e
sicuro di usare la funzione).  Altrimenti il kernel aggancia il segmento
all'indirizzo specificato da \param{shmaddr}; questo però può avvenire solo se
l'indirizzo coincide con il limite di una pagina, cioè se è un multiplo esatto
del parametro di sistema \const{SHMLBA}, che in Linux è sempre uguale
\const{PAGE\_SIZE}.

Si tenga presente però che quando si usa \val{NULL} come valore di
\param{shmaddr}, l'indirizzo restituito da \func{shmat} può cambiare da
processo a processo; pertanto se nell'area di memoria condivisa si salvano
anche degli indirizzi, si deve avere cura di usare valori relativi (in genere
riferiti all'indirizzo di partenza del segmento).

L'argomento \param{shmflg} permette di cambiare il comportamento della
funzione; esso va specificato come maschera binaria, i bit utilizzati al
momento sono tre e sono identificati dalle costanti \const{SHM\_RND},
\const{SHM\_RDONLY} e \const{SHM\_REMAP} che vanno combinate con un OR
aritmetico.

Specificando \constd{SHM\_RND} si evita che \func{shmat} ritorni un errore
quando \param{shmaddr} non è allineato ai confini di una pagina. Si può quindi
usare un valore qualunque per \param{shmaddr}, e il segmento verrà comunque
agganciato, ma al più vicino multiplo di \const{SHMLBA}; il nome della
costante sta infatti per \textit{rounded}, e serve per specificare un
indirizzo come arrotondamento.

L'uso di \constd{SHM\_RDONLY} permette di agganciare il segmento in sola
lettura (si ricordi che anche le pagine di memoria hanno dei permessi), in tal
caso un tentativo di scrivere sul segmento comporterà una violazione di
accesso con l'emissione di un segnale di \signal{SIGSEGV}. Il comportamento
usuale di \func{shmat} è quello di agganciare il segmento con l'accesso in
lettura e scrittura (ed il processo deve aver questi permessi in
\var{shm\_perm}), non è prevista la possibilità di agganciare un segmento in
sola scrittura.

Infine \constd{SHM\_REMAP} è una estensione specifica di Linux (quindi non
portabile) che indica che la mappatura del segmento deve rimpiazzare ogni
precedente mappatura esistente nell'intervallo iniziante
all'indirizzo \param{shmaddr} e di dimensione pari alla lunghezza del
segmento. In condizioni normali questo tipo di richiesta fallirebbe con un
errore di \errval{EINVAL}. Ovviamente usando \const{SHM\_REMAP}
l'argomento \param{shmaddr} non può essere nullo.

In caso di successo la funzione \func{shmat} aggiorna anche i seguenti campi
della struttura \struct{shmid\_ds}:
\begin{itemize*}
\item il tempo \var{shm\_atime} dell'ultima operazione di aggancio viene
  impostato al tempo corrente.
\item il \ids{PID} \var{shm\_lpid} dell'ultimo processo che ha operato sul
  segmento viene impostato a quello del processo corrente.
\item il numero \var{shm\_nattch} di processi agganciati al segmento viene
  aumentato di uno.
\end{itemize*}

Come accennato in sez.~\ref{sec:proc_fork} un segmento di memoria condivisa
agganciato ad un processo viene ereditato da un figlio attraverso una
\func{fork}, dato che quest'ultimo riceve una copia dello spazio degli
indirizzi del padre. Invece, dato che attraverso una \func{exec} viene
eseguito un diverso programma con uno spazio di indirizzi completamente
diverso, tutti i segmenti agganciati al processo originario vengono
automaticamente sganciati. Lo stesso avviene all'uscita del processo
attraverso una \func{exit}.

Una volta che un segmento di memoria condivisa non serve più, si può
sganciarlo esplicitamente dal processo usando l'altra funzione
dell'interfaccia, \funcd{shmdt}, il cui prototipo è:

\begin{funcproto}{
\fhead{sys/types.h} 
\fhead{sys/shm.h}
\fdecl{int shmdt(const void *shmaddr)}

\fdesc{Sgancia dal processo un segmento di memoria condivisa.}
}

{La funzione ritorna $0$ in caso di successo e $-1$ per un errore, la funzione
  fallisce solo quando non c'è un segmento agganciato
  all'indirizzo \param{shmaddr}, con \var{errno} che assume il valore
  \errval{EINVAL}.  
}  
\end{funcproto}

La funzione sgancia dallo spazio degli indirizzi del processo un segmento di
memoria condivisa; questo viene identificato con l'indirizzo \param{shmaddr}
restituito dalla precedente chiamata a \func{shmat} con il quale era stato
agganciato al processo.

In caso di successo la funzione aggiorna anche i seguenti campi di
\struct{shmid\_ds}:
\begin{itemize*}
\item il tempo \var{shm\_dtime} dell'ultima operazione di sganciamento viene
  impostato al tempo corrente.
\item il \ids{PID} \var{shm\_lpid} dell'ultimo processo che ha operato sul
  segmento viene impostato a quello del processo corrente.
\item il numero \var{shm\_nattch} di processi agganciati al segmento viene
  decrementato di uno.
\end{itemize*} 
inoltre la regione di indirizzi usata per il segmento di memoria condivisa
viene tolta dallo spazio di indirizzi del processo.

\begin{figure}[!htbp]
  \footnotesize \centering
  \begin{minipage}[c]{\codesamplewidth}
    \includecodesample{listati/SharedMem.c}
  \end{minipage} 
  \normalsize 
  \caption{Il codice delle funzioni che permettono di creare, trovare e
    rimuovere un segmento di memoria condivisa.}
  \label{fig:ipc_sysv_shm_func}
\end{figure}

Come esempio di uso di queste funzioni vediamo come implementare una serie di
funzioni di libreria che ne semplifichino l'uso, automatizzando le operazioni
più comuni; il codice, contenuto nel file \file{SharedMem.c}, è riportato in
fig.~\ref{fig:ipc_sysv_shm_func}.

La prima funzione (\texttt{\small 1-16}) è \func{ShmCreate} che, data una
chiave, crea il segmento di memoria condivisa restituendo il puntatore allo
stesso. La funzione comincia (\texttt{\small 6}) con il chiamare
\func{shmget}, usando il flag \const{IPC\_CREATE} per creare il segmento
qualora non esista, ed assegnandogli i privilegi specificati dall'argomento
\var{perm} e la dimensione specificata dall'argomento \var{shm\_size}.  In
caso di errore (\texttt{\small 7-9}) si ritorna immediatamente un puntatore
nullo, altrimenti (\texttt{\small 10}) si prosegue agganciando il segmento di
memoria condivisa al processo con \func{shmat}. In caso di errore
(\texttt{\small 11-13}) si restituisce di nuovo un puntatore nullo, infine
(\texttt{\small 14}) si inizializza con \func{memset} il contenuto del
segmento al valore costante specificato dall'argomento \var{fill}, e poi si
ritorna il puntatore al segmento stesso.

La seconda funzione (\texttt{\small 17-31}) è \func{ShmFind}, che, data una
chiave, restituisce l'indirizzo del segmento ad essa associato. Anzitutto
(\texttt{\small 22}) si richiede l'identificatore del segmento con
\func{shmget}, ritornando (\texttt{\small 23-25}) un puntatore nullo in caso
di errore. Poi si prosegue (\texttt{\small 26}) agganciando il segmento al
processo con \func{shmat}, restituendo (\texttt{\small 27-29}) di nuovo un
puntatore nullo in caso di errore, se invece non ci sono errori si restituisce
il puntatore ottenuto da \func{shmat}.

La terza funzione (\texttt{\small 32-51}) è \func{ShmRemove} che, data la
chiave ed il puntatore associati al segmento di memoria condivisa, prima lo
sgancia dal processo e poi lo rimuove. Il primo passo (\texttt{\small 37}) è
la chiamata a \func{shmdt} per sganciare il segmento, restituendo
(\texttt{\small 38-39}) un valore -1 in caso di errore. Il passo successivo
(\texttt{\small 41}) è utilizzare \func{shmget} per ottenere l'identificatore
associato al segmento data la chiave \var{key}. Al solito si restituisce un
valore di -1 (\texttt{\small 42-45}) in caso di errore, mentre se tutto va
bene si conclude restituendo un valore nullo.

Benché la memoria condivisa costituisca il meccanismo di intercomunicazione
fra processi più veloce, essa non è sempre il più appropriato, dato che, come
abbiamo visto, si avrà comunque la necessità di una sincronizzazione degli
accessi.  Per questo motivo, quando la comunicazione fra processi è
sequenziale, altri meccanismi come le \textit{pipe}, le \textit{fifo} o i
socket, che non necessitano di sincronizzazione esplicita, sono da
preferire. Essa diventa l'unico meccanismo possibile quando la comunicazione
non è sequenziale\footnote{come accennato in sez.~\ref{sec:ipc_sysv_mq} per la
  comunicazione non sequenziale si possono usare le code di messaggi,
  attraverso l'uso del campo \var{mtype}, ma solo se quest'ultima può essere
  effettuata in forma di messaggio.} o quando non può avvenire secondo una
modalità predefinita.

Un esempio classico di uso della memoria condivisa è quello del
``\textit{monitor}'', in cui viene per scambiare informazioni fra un processo
server, che vi scrive dei dati di interesse generale che ha ottenuto, e i
processi client interessati agli stessi dati che così possono leggerli in
maniera completamente asincrona.  Con questo schema di funzionamento da una
parte si evita che ciascun processo client debba compiere l'operazione,
potenzialmente onerosa, di ricavare e trattare i dati, e dall'altra si evita
al processo server di dover gestire l'invio a tutti i client di tutti i dati
(non potendo il server sapere quali di essi servono effettivamente al singolo
client).

Nel nostro caso implementeremo un ``\textsl{monitor}'' di una directory: un
processo si incaricherà di tenere sotto controllo alcuni parametri relativi ad
una directory (il numero dei file contenuti, la dimensione totale, quante
directory, link simbolici, file normali, ecc.) che saranno salvati in un
segmento di memoria condivisa cui altri processi potranno accedere per
ricavare la parte di informazione che interessa.

In fig.~\ref{fig:ipc_dirmonitor_main} si è riportata la sezione principale del
corpo del programma server, insieme alle definizioni delle altre funzioni
usate nel programma e delle variabili globali, omettendo tutto quello che
riguarda la gestione delle opzioni e la stampa delle istruzioni di uso a
video; al solito il codice completo si trova con i sorgenti allegati nel file
\file{DirMonitor.c}.

\begin{figure}[!htbp]
  \footnotesize \centering
  \begin{minipage}[c]{\codesamplewidth}
    \includecodesample{listati/DirMonitor.c}
  \end{minipage} 
  \normalsize 
  \caption{Codice della funzione principale del programma \file{DirMonitor.c}.}
  \label{fig:ipc_dirmonitor_main}
\end{figure}

Il programma usa delle variabili globali (\texttt{\small 2-14}) per mantenere
i valori relativi agli oggetti usati per la comunicazione inter-processo; si è
definita inoltre una apposita struttura \struct{DirProp} che contiene i dati
relativi alle proprietà che si vogliono mantenere nella memoria condivisa, per
l'accesso da parte dei client.

Il programma, dopo la sezione, omessa, relativa alla gestione delle opzioni da
riga di comando (che si limitano alla eventuale stampa di un messaggio di
aiuto a video ed all'impostazione della durata dell'intervallo con cui viene
ripetuto il calcolo delle proprietà della directory) controlla (\texttt{\small
  20-23}) che sia stato specificato l'argomento necessario contenente il nome
della directory da tenere sotto controllo, senza il quale esce immediatamente
con un messaggio di errore.

Poi, per verificare che l'argomento specifichi effettivamente una directory,
si esegue (\texttt{\small 24-26}) su di esso una \func{chdir}, uscendo
immediatamente in caso di errore.  Questa funzione serve anche per impostare
la directory di lavoro del programma nella directory da tenere sotto
controllo, in vista del successivo uso della funzione \func{daemon}. Si noti
come si è potuta fare questa scelta, nonostante le indicazioni illustrate in
sez.~\ref{sec:sess_daemon}, per il particolare scopo del programma, che
necessita comunque di restare all'interno di una directory.

Infine (\texttt{\small 27-29}) si installano i gestori per i vari segnali di
terminazione che, avendo a che fare con un programma che deve essere eseguito
come server, sono il solo strumento disponibile per concluderne l'esecuzione.

Il passo successivo (\texttt{\small 30-39}) è quello di creare gli oggetti di
intercomunicazione necessari. Si inizia costruendo (\texttt{\small 30}) la
chiave da usare come riferimento con il nome del programma,\footnote{si è
  usato un riferimento relativo alla home dell'utente, supposto che i sorgenti
  di GaPiL siano stati installati direttamente in essa; qualora si effettui
  una installazione diversa si dovrà correggere il programma.} dopo di che si
richiede (\texttt{\small 31}) la creazione di un segmento di memoria condivisa
con usando la funzione \func{ShmCreate} illustrata in precedenza (una pagina
di memoria è sufficiente per i dati che useremo), uscendo (\texttt{\small
  32-35}) qualora la creazione ed il successivo agganciamento al processo non
abbia successo. Con l'indirizzo \var{shmptr} così ottenuto potremo poi
accedere alla memoria condivisa, che, per come abbiamo lo abbiamo definito,
sarà vista nella forma data da \struct{DirProp}. Infine (\texttt{\small
  36-39}) utilizzando sempre la stessa chiave, si crea, tramite le funzioni
di interfaccia già descritte in sez.~\ref{sec:ipc_sysv_sem}, anche un mutex,
che utilizzeremo per regolare l'accesso alla memoria condivisa.

Completata l'inizializzazione e la creazione degli oggetti di
intercomunicazione il programma entra nel ciclo principale (\texttt{\small
  40-49}) dove vengono eseguite indefinitamente le attività di monitoraggio.
Il primo passo (\texttt{\small 41}) è eseguire \func{daemon} per proseguire
con l'esecuzione in background come si conviene ad un programma demone; si
noti che si è mantenuta, usando un valore non nullo del primo argomento, la
directory di lavoro corrente.  Una volta che il programma è andato in
background l'esecuzione prosegue all'interno di un ciclo infinito
(\texttt{\small 42-48}).

Si inizia (\texttt{\small 43}) bloccando il mutex con \func{MutexLock} per
poter accedere alla memoria condivisa (la funzione si bloccherà
automaticamente se qualche client sta leggendo), poi (\texttt{\small 44}) si
cancellano i valori precedentemente immagazzinati nella memoria condivisa con
\func{memset}, e si esegue (\texttt{\small 45}) un nuovo calcolo degli stessi
utilizzando la funzione \myfunc{dir\_scan}; infine (\texttt{\small 46}) si
sblocca il mutex con \func{MutexUnlock}, e si attende (\texttt{\small 47}) per
il periodo di tempo specificato a riga di comando con l'opzione \code{-p}
usando una \func{sleep}.

Si noti come per il calcolo dei valori da mantenere nella memoria condivisa si
sia usata ancora una volta la funzione \myfunc{dir\_scan}, già utilizzata (e
descritta in dettaglio) in sez.~\ref{sec:file_dir_read}, che ci permette di
effettuare la scansione delle voci della directory, chiamando per ciascuna di
esse la funzione \func{ComputeValues}, che esegue tutti i calcoli necessari.

\begin{figure}[!htbp]
  \footnotesize \centering
  \begin{minipage}[c]{\codesamplewidth}
    \includecodesample{listati/ComputeValues.c}
  \end{minipage} 
  \normalsize 
  \caption{Codice delle funzioni ausiliarie usate da \file{DirMonitor.c}.}
  \label{fig:ipc_dirmonitor_sub}
\end{figure}


Il codice di quest'ultima è riportato in fig.~\ref{fig:ipc_dirmonitor_sub}.
Come si vede la funzione (\texttt{\small 2-16}) è molto semplice e si limita a
chiamare (\texttt{\small 5}) la funzione \func{stat} sul file indicato da
ciascuna voce, per ottenerne i dati, che poi utilizza per incrementare i vari
contatori nella memoria condivisa, cui accede grazie alla variabile globale
\var{shmptr}.

Dato che la funzione è chiamata da \myfunc{dir\_scan}, si è all'interno del
ciclo principale del programma, con un mutex acquisito, perciò non è
necessario effettuare nessun controllo e si può accedere direttamente alla
memoria condivisa usando \var{shmptr} per riempire i campi della struttura
\struct{DirProp}; così prima (\texttt{\small 6-7}) si sommano le dimensioni
dei file ed il loro numero, poi, utilizzando le macro di
tab.~\ref{tab:file_type_macro}, si contano (\texttt{\small 8-14}) quanti ce
ne sono per ciascun tipo.

In fig.~\ref{fig:ipc_dirmonitor_sub} è riportato anche il codice
(\texttt{\small 17-23}) del gestore dei segnali di terminazione, usato per
chiudere il programma. Esso, oltre a provocare l'uscita del programma, si
incarica anche di cancellare tutti gli oggetti di intercomunicazione non più
necessari.  Per questo anzitutto (\texttt{\small 19}) acquisisce il mutex con
\func{MutexLock}, per evitare di operare mentre un client sta ancora leggendo
i dati, dopo di che (\texttt{\small 20}) distacca e rimuove il segmento di
memoria condivisa usando \func{ShmRemove}.  Infine (\texttt{\small 21})
rimuove il mutex con \func{MutexRemove} ed esce (\texttt{\small 22}).

\begin{figure}[!htbp]
  \footnotesize \centering
  \begin{minipage}[c]{\codesamplewidth}
    \includecodesample{listati/ReadMonitor.c}
  \end{minipage} 
  \normalsize 
  \caption{Codice del programma client del monitor delle proprietà di una
    directory, \file{ReadMonitor.c}.}
  \label{fig:ipc_dirmonitor_client}
\end{figure}

Il codice del client usato per leggere le informazioni mantenute nella memoria
condivisa è riportato in fig.~\ref{fig:ipc_dirmonitor_client}. Al solito si è
omessa la sezione di gestione delle opzioni e la funzione che stampa a video
le istruzioni; il codice completo è nei sorgenti allegati, nel file
\file{ReadMonitor.c}.

Una volta conclusa la gestione delle opzioni a riga di comando il programma
rigenera (\texttt{\small 7}) con \func{ftok} la stessa chiave usata dal server
per identificare il segmento di memoria condivisa ed il mutex, poi
(\texttt{\small 8}) richiede con \func{ShmFind} l'indirizzo della memoria
condivisa agganciando al contempo il segmento al processo, Infine
(\texttt{\small 17-20}) con \func{MutexFind} si richiede l'identificatore del
mutex.  Completata l'inizializzazione ed ottenuti i riferimenti agli oggetti
di intercomunicazione necessari viene eseguito il corpo principale del
programma (\texttt{\small 21-33}); si comincia (\texttt{\small 22})
acquisendo il mutex con \func{MutexLock}; qui avviene il blocco del processo
se la memoria condivisa non è disponibile.  Poi (\texttt{\small 23-31}) si
stampano i vari valori mantenuti nella memoria condivisa attraverso l'uso di
\var{shmptr}.  Infine (\texttt{\small 41}) con \func{MutexUnlock} si rilascia
il mutex, prima di uscire.

Verifichiamo allora il funzionamento dei nostri programmi; al solito, usando
le funzioni di libreria occorre definire opportunamente
\code{LD\_LIBRARY\_PATH}; poi si potrà lanciare il server con:
\begin{Console}
[piccardi@gont sources]$ \textbf{./dirmonitor ./}
\end{Console}
%$
ed avendo usato \func{daemon} il comando ritornerà immediatamente. Una volta
che il server è in esecuzione, possiamo passare ad invocare il client per
verificarne i risultati, in tal caso otterremo:
\begin{Console}
[piccardi@gont sources]$ \textbf{./readmon} 
Ci sono 68 file dati
Ci sono 3 directory
Ci sono 0 link
Ci sono 0 fifo
Ci sono 0 socket
Ci sono 0 device a caratteri
Ci sono 0 device a blocchi
Totale  71 file, per 489831 byte
\end{Console}
%$
ed un rapido calcolo (ad esempio con \code{ls -a | wc} per contare i file) ci
permette di verificare che il totale dei file è giusto. Un controllo con
\cmd{ipcs} ci permette inoltre di verificare la presenza di un segmento di
memoria condivisa e di un semaforo:
\begin{Console}
[piccardi@gont sources]$ \textbf{ipcs}
------ Shared Memory Segments --------
key        shmid      owner      perms      bytes      nattch     status      
0xffffffff 54067205   piccardi  666        4096       1                       

------ Semaphore Arrays --------
key        semid      owner      perms      nsems     
0xffffffff 229376     piccardi  666        1         

------ Message Queues --------
key        msqid      owner      perms      used-bytes   messages    
\end{Console}
%$

Se a questo punto aggiungiamo un file, ad esempio con \code{touch prova},
potremo verificare che, passati nel peggiore dei casi almeno 10 secondi (o
l'eventuale altro intervallo impostato per la rilettura dei dati) avremo:
\begin{Console}
[piccardi@gont sources]$ \textbf{./readmon} 
Ci sono 69 file dati
Ci sono 3 directory
Ci sono 0 link
Ci sono 0 fifo
Ci sono 0 socket
Ci sono 0 device a caratteri
Ci sono 0 device a blocchi
Totale  72 file, per 489887 byte
\end{Console}
%$

A questo punto possiamo far uscire il server inviandogli un segnale di
\signal{SIGTERM} con il comando \code{killall dirmonitor}, a questo punto
ripetendo la lettura, otterremo un errore:
\begin{Console}
[piccardi@gont sources]$ \textbf{./readmon} 
Cannot find shared memory: No such file or directory
\end{Console}
%$
e inoltre potremo anche verificare che anche gli oggetti di intercomunicazione
visti in precedenza sono stati regolarmente  cancellati:
\begin{Console}
[piccardi@gont sources]$ \textbf{ipcs}
------ Shared Memory Segments --------
key        shmid      owner      perms      bytes      nattch     status      

------ Semaphore Arrays --------
key        semid      owner      perms      nsems     

------ Message Queues --------
key        msqid      owner      perms      used-bytes   messages    
\end{Console}
%$


%% Per capire meglio il funzionamento delle funzioni facciamo ancora una volta
%% riferimento alle strutture con cui il kernel implementa i segmenti di memoria
%% condivisa; uno schema semplificato della struttura è illustrato in
%% fig.~\ref{fig:ipc_shm_struct}. 

%% \begin{figure}[!htb]
%%   \centering
%%   \includegraphics[width=10cm]{img/shmstruct}
%%    \caption{Schema dell'implementazione dei segmenti di memoria condivisa in
%%     Linux.}
%%   \label{fig:ipc_shm_struct}
%% \end{figure}




\section{Tecniche alternative}
\label{sec:ipc_alternatives}

Come abbiamo detto in sez.~\ref{sec:ipc_sysv_generic}, e ripreso nella
descrizione dei singoli oggetti che ne fan parte, il \textit{SysV-IPC}
presenta numerosi problemi; in \cite{APUE}\footnote{in particolare nel
  capitolo 14.}  Stevens ne effettua una accurata analisi (alcuni dei concetti
sono già stati accennati in precedenza) ed elenca alcune possibili tecniche
alternative, che vogliamo riprendere in questa sezione.


\subsection{Alternative alle code di messaggi}
\label{sec:ipc_mq_alternative}
 
Le code di messaggi sono probabilmente il meno usato degli oggetti del
\textit{SysV-IPC}; esse infatti nacquero principalmente come meccanismo di
comunicazione bidirezionale quando ancora le \textit{pipe} erano
unidirezionali; con la disponibilità di \func{socketpair} (vedi
sez.~\ref{sec:ipc_socketpair}) o utilizzando una coppia di \textit{pipe}, si
può ottenere questo risultato senza incorrere nelle complicazioni introdotte
dal \textit{SysV-IPC}.

In realtà, grazie alla presenza del campo \var{mtype}, le code di messaggi
hanno delle caratteristiche ulteriori, consentendo una classificazione dei
messaggi ed un accesso non rigidamente sequenziale; due caratteristiche che
sono impossibili da ottenere con le \textit{pipe} e i socket di
\func{socketpair}.  A queste esigenze però si può comunque ovviare in maniera
diversa con un uso combinato della memoria condivisa e dei meccanismi di
sincronizzazione, per cui alla fine l'uso delle code di messaggi classiche è
relativamente poco diffuso.

% TODO: trattare qui, se non si trova posto migliore, copy_from_process e
% copy_to_process, introdotte con il kernel 3.2. Vedi
% http://lwn.net/Articles/405346/ e
% http://ozlabs.org/~cyeoh/cma/process_vm_readv.txt 


\subsection{I \textsl{file di lock}}
\label{sec:ipc_file_lock}

\index{file!di~lock|(}

Come illustrato in sez.~\ref{sec:ipc_sysv_sem} i semafori del \textit{SysV-IPC}
presentano una interfaccia inutilmente complessa e con alcuni difetti
strutturali, per questo quando si ha una semplice esigenza di sincronizzazione
per la quale basterebbe un semaforo binario (quello che abbiamo definito come
\textit{mutex}), per indicare la disponibilità o meno di una risorsa, senza la
necessità di un contatore come i semafori, si possono utilizzare metodi
alternativi.

La prima possibilità, utilizzata fin dalle origini di Unix, è quella di usare
dei \textsl{file di lock} (per i quali è stata anche riservata una opportuna
directory, \file{/var/lock}, nella standardizzazione del \textit{Filesystem
  Hierarchy Standard}). Per questo si usa la caratteristica della funzione
\func{open} (illustrata in sez.~\ref{sec:file_open_close}) che
prevede\footnote{questo è quanto dettato dallo standard POSIX.1, ciò non
  toglie che in alcune implementazioni questa tecnica possa non funzionare; in
  particolare per Linux, nel caso di NFS, si è comunque soggetti alla
  possibilità di una \textit{race condition}.} che essa ritorni un errore
quando usata con i flag di \const{O\_CREAT} e \const{O\_EXCL}. In tal modo la
creazione di un \textsl{file di lock} può essere eseguita atomicamente, il
processo che crea il file con successo si può considerare come titolare del
lock (e della risorsa ad esso associata) mentre il rilascio si può eseguire
con una chiamata ad \func{unlink}.

Un esempio dell'uso di questa funzione è mostrato dalle funzioni
\func{LockFile} ed \func{UnlockFile} riportate in fig.~\ref{fig:ipc_file_lock}
(sono contenute in \file{LockFile.c}, un altro dei sorgenti allegati alla
guida) che permettono rispettivamente di creare e rimuovere un \textsl{file di
  lock}. Come si può notare entrambe le funzioni sono elementari; la prima
(\texttt{\small 4-10}) si limita ad aprire il file di lock (\texttt{\small
  9}) nella modalità descritta, mentre la seconda (\texttt{\small 11-17}) lo
cancella con \func{unlink}.

\begin{figure}[!htbp]
  \footnotesize \centering
  \begin{minipage}[c]{\codesamplewidth}
    \includecodesample{listati/LockFile.c}
  \end{minipage} 
  \normalsize 
  \caption{Il codice delle funzioni \func{LockFile} e \func{UnlockFile} che
    permettono di creare e rimuovere un \textsl{file di lock}.}
  \label{fig:ipc_file_lock}
\end{figure}

Uno dei limiti di questa tecnica è che, come abbiamo già accennato in
sez.~\ref{sec:file_open_close}, questo comportamento di \func{open} può non
funzionare (la funzione viene eseguita, ma non è garantita l'atomicità
dell'operazione) se il filesystem su cui si va ad operare è su NFS; in tal
caso si può adottare una tecnica alternativa che prevede l'uso della
\func{link} per creare come \textsl{file di lock} un hard link ad un file
esistente; se il link esiste già e la funzione fallisce, significa che la
risorsa è bloccata e potrà essere sbloccata solo con un \func{unlink},
altrimenti il link è creato ed il lock acquisito; il controllo e l'eventuale
acquisizione sono atomici; la soluzione funziona anche su NFS, ma ha un altro
difetto è che è quello di poterla usare solo se si opera all'interno di uno
stesso filesystem.

In generale comunque l'uso di un \textsl{file di lock} presenta parecchi
problemi che non lo rendono una alternativa praticabile per la
sincronizzazione: anzitutto in caso di terminazione imprevista del processo,
si lascia allocata la risorsa (il \textsl{file di lock}) e questa deve essere
sempre cancellata esplicitamente.  Inoltre il controllo della disponibilità
può essere eseguito solo con una tecnica di \textit{polling}, ed è quindi
molto inefficiente.

La tecnica dei file di lock ha comunque una sua utilità, e può essere usata
con successo quando l'esigenza è solo quella di segnalare l'occupazione di una
risorsa, senza necessità di attendere che questa si liberi; ad esempio la si
usa spesso per evitare interferenze sull'uso delle porte seriali da parte di
più programmi: qualora si trovi un file di lock il programma che cerca di
accedere alla seriale si limita a segnalare che la risorsa non è disponibile.

\index{file!di~lock|)}


\subsection{La sincronizzazione con il \textit{file locking}}
\label{sec:ipc_lock_file}

Dato che i file di lock presentano gli inconvenienti illustrati in precedenza,
la tecnica alternativa di sincronizzazione più comune è quella di fare ricorso
al \textit{file locking} (trattato in sez.~\ref{sec:file_locking}) usando
\func{fcntl} su un file creato per l'occasione per ottenere un write lock. In
questo modo potremo usare il lock come un \textit{mutex}: per bloccare la
risorsa basterà acquisire il lock, per sbloccarla basterà rilasciare il
lock. Una richiesta fatta con un write lock metterà automaticamente il
processo in stato di attesa, senza necessità di ricorrere al \textit{polling}
per determinare la disponibilità della risorsa, e al rilascio della stessa da
parte del processo che la occupava si otterrà il nuovo lock atomicamente.

Questo approccio presenta il notevole vantaggio che alla terminazione di un
processo tutti i lock acquisiti vengono rilasciati automaticamente (alla
chiusura dei relativi file) e non ci si deve preoccupare di niente; inoltre
non consuma risorse permanentemente allocate nel sistema. Lo svantaggio è che,
dovendo fare ricorso a delle operazioni sul filesystem, esso è in genere
leggermente più lento.

\begin{figure}[!htbp]
  \footnotesize \centering
  \begin{minipage}[c]{\codesamplewidth}
    \includecodesample{listati/MutexLocking.c}
  \end{minipage} 
  \normalsize 
  \caption{Il codice delle funzioni che permettono per la gestione dei
    \textit{mutex} con il \textit{file locking}.}
  \label{fig:ipc_flock_mutex}
\end{figure}

Il codice delle varie funzioni usate per implementare un mutex utilizzando il
\textit{file locking} è riportato in fig.~\ref{fig:ipc_flock_mutex}; si è
mantenuta volutamente una struttura analoga alle precedenti funzioni che usano
i semafori, anche se le due interfacce non possono essere completamente
equivalenti, specie per quanto riguarda la rimozione del mutex.

La prima funzione (\texttt{\small 1-5}) è \func{CreateMutex}, e serve a
creare il mutex; la funzione è estremamente semplice, e si limita
(\texttt{\small 4}) a creare, con una opportuna chiamata ad \func{open}, il
file che sarà usato per il successivo \textit{file locking}, assicurandosi che
non esista già (nel qual caso segnala un errore); poi restituisce il file
descriptor che sarà usato dalle altre funzioni per acquisire e rilasciare il
mutex.

La seconda funzione (\texttt{\small 6-10}) è \func{FindMutex}, che, come la
precedente, è stata definita per mantenere una analogia con la corrispondente
funzione basata sui semafori. Anch'essa si limita (\texttt{\small 9}) ad
aprire il file da usare per il \textit{file locking}, solo che in questo caso
le opzioni di \func{open} sono tali che il file in questione deve esistere di
già.

La terza funzione (\texttt{\small 11-22}) è \func{LockMutex} e serve per
acquisire il mutex. La funzione definisce (\texttt{\small 14}) e inizializza
(\texttt{\small 16-19}) la struttura \var{lock} da usare per acquisire un
write lock sul file, che poi (\texttt{\small 21}) viene richiesto con
\func{fcntl}, restituendo il valore di ritorno di quest'ultima. Se il file è
libero il lock viene acquisito e la funzione ritorna immediatamente;
altrimenti \func{fcntl} si bloccherà (si noti che la si è chiamata con
\const{F\_SETLKW}) fino al rilascio del lock.

La quarta funzione (\texttt{\small 24-34}) è \func{UnlockMutex} e serve a
rilasciare il mutex. La funzione è analoga alla precedente, solo che in questo
caso si inizializza (\texttt{\small 28-31}) la struttura \var{lock} per il
rilascio del lock, che viene effettuato (\texttt{\small 33}) con la opportuna
chiamata a \func{fcntl}. Avendo usato il \textit{file locking} in semantica
POSIX (si riveda quanto detto sez.~\ref{sec:file_posix_lock}) solo il processo
che ha precedentemente eseguito il lock può sbloccare il mutex.

La quinta funzione (\texttt{\small 36-39}) è \func{RemoveMutex} e serve a
cancellare il mutex. Anche questa funzione è stata definita per mantenere una
analogia con le funzioni basate sui semafori, e si limita a cancellare
(\texttt{\small 38}) il file con una chiamata ad \func{unlink}. Si noti che in
questo caso la funzione non ha effetto sui mutex già ottenuti con precedenti
chiamate a \func{FindMutex} o \func{CreateMutex}, che continueranno ad essere
disponibili fintanto che i relativi file descriptor restano aperti. Pertanto
per rilasciare un mutex occorrerà prima chiamare \func{UnlockMutex} oppure
chiudere il file usato per il lock.

La sesta funzione (\texttt{\small 41-55}) è \func{ReadMutex} e serve a
leggere lo stato del mutex. In questo caso si prepara (\texttt{\small 46-49})
la solita struttura \var{lock} come l'acquisizione del lock, ma si effettua
(\texttt{\small 51}) la chiamata a \func{fcntl} usando il comando
\const{F\_GETLK} per ottenere lo stato del lock, e si restituisce
(\texttt{\small 52}) il valore di ritorno in caso di errore, ed il valore del
campo \var{l\_type} (che descrive lo stato del lock) altrimenti
(\texttt{\small 54}). Per questo motivo la funzione restituirà -1 in caso di
errore e uno dei due valori \const{F\_UNLCK} o \const{F\_WRLCK}\footnote{non
  si dovrebbe mai avere il terzo valore possibile, \const{F\_RDLCK}, dato che
  la nostra interfaccia usa solo i write lock. Però è sempre possibile che
  siano richiesti altri lock sul file al di fuori dell'interfaccia, nel qual
  caso si potranno avere, ovviamente, interferenze indesiderate.} in caso di
successo, ad indicare che il mutex è, rispettivamente, libero o occupato.

Basandosi sulla semantica dei file lock POSIX valgono tutte le considerazioni
relative al comportamento di questi ultimi fatte in
sez.~\ref{sec:file_posix_lock}; questo significa ad esempio che, al contrario
di quanto avveniva con l'interfaccia basata sui semafori, chiamate multiple a
\func{UnlockMutex} o \func{LockMutex} non si cumulano e non danno perciò
nessun inconveniente.


\subsection{Il \textit{memory mapping} anonimo}
\label{sec:ipc_mmap_anonymous}

\itindbeg{memory~mapping} 

Abbiamo già visto che quando i processi sono \textsl{correlati}, se cioè hanno
almeno un progenitore comune, l'uso delle \textit{pipe} può costituire una
valida alternativa alle code di messaggi; nella stessa situazione si può
evitare l'uso di una memoria condivisa facendo ricorso al cosiddetto
\textit{memory mapping} anonimo.

In sez.~\ref{sec:file_memory_map} abbiamo visto come sia possibile mappare il
contenuto di un file nella memoria di un processo, e che, quando viene usato
il flag \const{MAP\_SHARED}, le modifiche effettuate al contenuto del file
vengono viste da tutti i processi che lo hanno mappato. Utilizzare questa
tecnica per creare una memoria condivisa fra processi diversi è estremamente
inefficiente, in quanto occorre passare attraverso il disco. 

Però abbiamo visto anche che se si esegue la mappatura con il flag
\const{MAP\_ANONYMOUS} la regione mappata non viene associata a nessun file,
anche se quanto scritto rimane in memoria e può essere riletto; allora, dato
che un processo figlio mantiene nel suo spazio degli indirizzi anche le
regioni mappate, esso sarà anche in grado di accedere a quanto in esse è
contenuto.

In questo modo diventa possibile creare una memoria condivisa fra processi
diversi, purché questi abbiano almeno un progenitore comune che ha effettuato
il \textit{memory mapping} anonimo.\footnote{nei sistemi derivati da SysV una
  funzionalità simile a questa viene implementata mappando il file speciale
  \file{/dev/zero}. In tal caso i valori scritti nella regione mappata non
  vengono ignorati (come accade qualora si scriva direttamente sul file), ma
  restano in memoria e possono essere riletti secondo le stesse modalità usate
  nel \textit{memory mapping} anonimo.} Vedremo come utilizzare questa tecnica
più avanti, quando realizzeremo una nuova versione del monitor visto in
sez.~\ref{sec:ipc_sysv_shm} che possa restituisca i risultati via rete.

\itindend{memory~mapping}

% TODO: fare esempio di mmap anonima

% TODO: con il kernel 3.2 è stata introdotta un nuovo meccanismo di
% intercomunicazione veloce chiamato Cross Memory Attach, da capire se e come
% trattarlo qui, vedi http://lwn.net/Articles/405346/
% https://git.kernel.org/?p=linux/kernel/git/torvalds/linux-2.6.git;a=commitdiff;h=fcf634098c00dd9cd247447368495f0b79be12d1

% TODO: con il kernel 3.17 è stata introdotta una fuunzionalità di
% sigillatura dei file mappati in memoria e la system call memfd
% (capire se va messo qui o altrove) vedi: http://lwn.net/Articles/593918/
% col 5.1 aggiunta a memfd F_SEAL_FUTURE_WRITE, vedi 
% https://git.kernel.org/linus/ab3948f58ff8 e https://lwn.net/Articles/782511/


\section{L'intercomunicazione fra processi di POSIX}
\label{sec:ipc_posix}

Per superare i numerosi problemi del \textit{SysV-IPC}, evidenziati per i suoi
aspetti generali in coda a sez.~\ref{sec:ipc_sysv_generic} e per i singoli
oggetti nei paragrafi successivi, lo standard POSIX.1b ha introdotto dei nuovi
meccanismi di comunicazione, che vanno sotto il nome di POSIX IPC, definendo
una interfaccia completamente nuova, che tratteremo in questa sezione.


\subsection{Considerazioni generali}
\label{sec:ipc_posix_generic}

Oggi Linux supporta tutti gli oggetti definito nello standard POSIX per l'IPC,
ma a lungo non è stato così; la memoria condivisa è presente a partire dal
kernel 2.4.x, i semafori sono forniti dalla \acr{glibc} nella sezione che
implementa i \textit{thread} POSIX di nuova generazione che richiedono il
kernel 2.6, le code di messaggi sono supportate a partire dal kernel 2.6.6.

La caratteristica fondamentale dell'interfaccia POSIX è l'abbandono dell'uso
degli identificatori e delle chiavi visti nel \textit{SysV-IPC}, per passare ai
\itindex{POSIX~IPC~names} \textit{POSIX IPC names}, che sono sostanzialmente
equivalenti ai nomi dei file. Tutte le funzioni che creano un oggetto di IPC
POSIX prendono come primo argomento una stringa che indica uno di questi nomi;
lo standard è molto generico riguardo l'implementazione, ed i nomi stessi
possono avere o meno una corrispondenza sul filesystem; tutto quello che è
richiesto è che:
\begin{itemize*}
\item i nomi devono essere conformi alle regole che caratterizzano i
  \textit{pathname}, in particolare non essere più lunghi di \const{PATH\_MAX}
  byte e terminati da un carattere nullo.
\item se il nome inizia per una \texttt{/} chiamate differenti allo stesso
  nome fanno riferimento allo stesso oggetto, altrimenti l'interpretazione del
  nome dipende dall'implementazione.
\item l'interpretazione di ulteriori \texttt{/} presenti nel nome dipende
  dall'implementazione.
\end{itemize*}

Data la assoluta genericità delle specifiche, il comportamento delle funzioni
è subordinato in maniera quasi completa alla relativa implementazione, tanto
che Stevens in \cite{UNP2} cita questo caso come un esempio della maniera
standard usata dallo standard POSIX per consentire implementazioni non
standardizzabili. 

Nel caso di Linux, sia per quanto riguarda la memoria condivisa ed i semafori,
che per le code di messaggi, tutto viene creato usando come radici delle
opportune directory (rispettivamente \file{/dev/shm} e \file{/dev/mqueue}, per
i dettagli si faccia riferimento a sez.~\ref{sec:ipc_posix_shm},
sez.~\ref{sec:ipc_posix_sem} e sez.~\ref{sec:ipc_posix_mq}).  I nomi
specificati nelle relative funzioni devono essere nella forma di un
\textit{pathname} assoluto (devono cioè iniziare con ``\texttt{/}'') e
corrisponderanno ad altrettanti file creati all'interno di queste directory;
per questo motivo detti nomi non possono contenere altre ``\texttt{/}'' oltre
quella iniziale.

Il vantaggio degli oggetti di IPC POSIX è comunque che essi vengono inseriti
nell'albero dei file, e possono essere maneggiati con le usuali funzioni e
comandi di accesso ai file, che funzionano come su dei file normali. Questo
però è vero nel caso di Linux, che usa una implementazione che lo consente,
non è detto che altrettanto valga per altri kernel. In particolare, come si
può facilmente verificare con il comando \cmd{strace}, sia per la memoria
condivisa che per le code di messaggi varie \textit{system call} utilizzate da
Linux corrispondono in realtà a quelle ordinarie dei file, essendo detti
oggetti realizzati come tali usando degli specifici filesystem.

In particolare i permessi associati agli oggetti di IPC POSIX sono identici ai
permessi dei file, ed il controllo di accesso segue esattamente la stessa
semantica (quella illustrata in sez.~\ref{sec:file_access_control}), e non
quella particolare (si ricordi quanto visto in
sez.~\ref{sec:ipc_sysv_access_control}) che viene usata per gli oggetti del
SysV-IPC. Per quanto riguarda l'attribuzione dell'utente e del gruppo
proprietari dell'oggetto alla creazione di quest'ultimo essa viene effettuata
secondo la semantica SysV: corrispondono cioè a \ids{UID} e \ids{GID} effettivi
del processo che esegue la creazione.


\subsection{Code di messaggi Posix}
\label{sec:ipc_posix_mq}

Le code di messaggi POSIX sono supportate da Linux a partire dalla versione
2.6.6 del kernel. In generale, come le corrispettive del \textit{SysV-IPC}, le
code di messaggi sono poco usate, dato che i socket, nei casi in cui sono
sufficienti, sono più comodi, e che in casi più complessi la comunicazione può
essere gestita direttamente con mutex (o semafori) e memoria condivisa con
tutta la flessibilità che occorre.

Per poter utilizzare le code di messaggi, oltre ad utilizzare un kernel
superiore al 2.6.6 occorre utilizzare la libreria \file{librt} che contiene le
funzioni dell'interfaccia POSIX ed i programmi che usano le code di messaggi
devono essere compilati aggiungendo l'opzione \code{-lrt} al comando
\cmd{gcc}. In corrispondenza all'inclusione del supporto nel kernel ufficiale
le funzioni di libreria sono state inserite nella \acr{glibc}, e sono
disponibili a partire dalla versione 2.3.4 delle medesime.

La libreria inoltre richiede la presenza dell'apposito filesystem di tipo
\texttt{mqueue} montato sulla directory \file{/dev/mqueue}; questo può essere
fatto aggiungendo ad \conffile{/etc/fstab} una riga come:
\begin{FileExample}[label=/etc/fstab]
mqueue   /dev/mqueue       mqueue    defaults        0      0
\end{FileExample}
ed esso sarà utilizzato come radice sulla quale vengono risolti i nomi delle
code di messaggi che iniziano con una ``\texttt{/}''. Le opzioni di mount
accettate sono \texttt{uid}, \texttt{gid} e \texttt{mode} che permettono
rispettivamente di impostare l'utente, il gruppo ed i permessi associati al
filesystem.


La funzione di sistema che permette di aprire (e crearla se non esiste ancora)
una coda di messaggi POSIX è \funcd{mq\_open}, ed il suo prototipo è:

\begin{funcproto}{
\fhead{fcntl.h}
\fhead{sys/stat.h}
\fhead{mqueue.h}
\fdecl{mqd\_t mq\_open(const char *name, int oflag)}
\fdecl{mqd\_t mq\_open(const char *name, int oflag, unsigned long mode,
    struct mq\_attr *attr)}

\fdesc{Apre una coda di messaggi POSIX impostandone le caratteristiche.}
}

{La funzione ritorna il descrittore associato alla coda in caso di successo e
  $-1$ per un errore, nel qual caso \var{errno} assumerà uno dei valori:
  \begin{errlist}
  \item[\errcode{EACCES}] il processo non ha i privilegi per accedere alla
    coda secondo quanto specificato da \param{oflag} oppure \const{name}
    contiene più di una ``\texttt{/}''.
  \item[\errcode{EEXIST}] si è specificato \const{O\_CREAT} e \const{O\_EXCL}
    ma la coda già esiste.
  \item[\errcode{EINVAL}] il file non supporta la funzione, o si è specificato
    \const{O\_CREAT} con una valore non nullo di \param{attr} e valori non
    validi dei campi \var{mq\_maxmsg} e \var{mq\_msgsize}; questi valori
    devono essere positivi ed inferiori ai limiti di sistema se il processo
    non ha privilegi amministrativi, inoltre \var{mq\_maxmsg} non può comunque
    superare \const{HARD\_MAX}.
  \item[\errcode{ENOENT}] non si è specificato \const{O\_CREAT} ma la coda non
    esiste o si è usato il nome ``\texttt{/}''.
  \item[\errcode{ENOSPC}] lo spazio è insufficiente, probabilmente per aver
    superato il limite di \texttt{queues\_max}.
  \end{errlist}
  ed inoltre \errval{EMFILE}, \errval{ENAMETOOLONG}, \errval{ENFILE},
  \errval{ENOMEM} ed nel loro significato generico.  }
\end{funcproto}

La funzione apre la coda di messaggi identificata dall'argomento \param{name}
restituendo il descrittore ad essa associato, del tutto analogo ad un file
descriptor, con l'unica differenza che lo standard prevede un apposito tipo
\typed{mqd\_t}. Nel caso di Linux si tratta in effetti proprio di un normale
file descriptor; pertanto, anche se questo comportamento non è portabile, lo
si può tenere sotto osservazione con le funzioni dell'I/O multiplexing (vedi
sez.~\ref{sec:file_multiplexing}) come possibile alternativa all'uso
dell'interfaccia di notifica di \func{mq\_notify} (che vedremo a breve).

Se il nome indicato fa riferimento ad una coda di messaggi già esistente, il
descrittore ottenuto farà riferimento allo stesso oggetto, pertanto tutti i
processi che hanno usato \func{mq\_open} su quel nome otterranno un
riferimento alla stessa coda. Diventa così immediato costruire un canale di
comunicazione fra detti processi.

La funzione è del tutto analoga ad \func{open} ed analoghi sono i valori che
possono essere specificati per \param{oflag}, che deve essere specificato come
maschera binaria; i valori possibili per i vari bit sono quelli visti in
sez.~\ref{sec:file_open_close} (per questo occorre includere \texttt{fcntl.h})
dei quali però \func{mq\_open} riconosce solo i seguenti:
\begin{basedescript}{\desclabelwidth{2.2cm}\desclabelstyle{\nextlinelabel}}
\item[\const{O\_RDONLY}] Apre la coda solo per la ricezione di messaggi. Il
  processo potrà usare il descrittore con \func{mq\_receive} ma non con
  \func{mq\_send}.
\item[\const{O\_WRONLY}] Apre la coda solo per la trasmissione di messaggi. Il
  processo potrà usare il descrittore con \func{mq\_send} ma non con
  \func{mq\_receive}.
\item[\const{O\_RDWR}] Apre la coda solo sia per la trasmissione che per la
  ricezione. 
\item[\const{O\_CREAT}] Necessario qualora si debba creare la coda; la
  presenza di questo bit richiede la presenza degli ulteriori argomenti
  \param{mode} e \param{attr}.
\item[\const{O\_EXCL}] Se usato insieme a \const{O\_CREAT} fa fallire la
  chiamata se la coda esiste già, altrimenti esegue la creazione atomicamente.
\item[\const{O\_NONBLOCK}] Imposta la coda in modalità non bloccante, le
  funzioni di ricezione e trasmissione non si bloccano quando non ci sono le
  risorse richieste, ma ritornano immediatamente con un errore di
  \errcode{EAGAIN}.
\end{basedescript}

I primi tre bit specificano la modalità di apertura della coda, e sono fra
loro esclusivi. Ma qualunque sia la modalità in cui si è aperta una coda,
questa potrà essere riaperta più volte in una modalità diversa, e vi si potrà
sempre accedere attraverso descrittori diversi, esattamente come si può fare
per i file normali.

Se la coda non esiste e la si vuole creare si deve specificare
\const{O\_CREAT}, in tal caso occorre anche specificare i permessi di
creazione con l'argomento \param{mode};\footnote{fino al 2.6.14 per un bug i
  valori della \textit{umask} del processo non venivano applicati a questi
  permessi.} i valori di quest'ultimo sono identici a quelli usati per
\func{open} (per questo occorre includere \texttt{sys/stat.h}), anche se per
le code di messaggi han senso solo i permessi di lettura e scrittura.

Oltre ai permessi di creazione possono essere specificati anche gli attributi
specifici della coda tramite l'argomento \param{attr}; quest'ultimo è un
puntatore ad una apposita struttura \struct{mq\_attr}, la cui definizione è
riportata in fig.~\ref{fig:ipc_mq_attr}.

\begin{figure}[!htb]
  \footnotesize \centering
  \begin{minipage}[c]{0.90\textwidth}
    \includestruct{listati/mq_attr.h}
  \end{minipage} 
  \normalsize
  \caption{La struttura \structd{mq\_attr}, contenente gli attributi di una
    coda di messaggi POSIX.}
  \label{fig:ipc_mq_attr}
\end{figure}

Per la creazione della coda i campi della struttura che devono essere
specificati sono \var{mq\_maxmsg} e \var{mq\_msgsize}, che indicano
rispettivamente il numero massimo di messaggi che può contenere e la
dimensione massima di un messaggio. Il valore dovrà essere positivo e minore
dei rispettivi limiti di sistema altrimenti la funzione fallirà con un errore
di \errcode{EINVAL}.  Se \param{attr} è un puntatore nullo gli attributi della
coda saranno impostati ai valori predefiniti.

I suddetti limiti di sistema sono impostati attraverso una serie di file
presenti sotto \texttt{/proc/sys/fs/mqueue}, in particolare i file che
controllano i valori dei limiti sono:
\begin{basedescript}{\desclabelwidth{1.5cm}\desclabelstyle{\nextlinelabel}}
\item[\sysctlfiled{fs/mqueue/msg\_max}] Indica il valore massimo del numero di
  messaggi in una coda e agisce come limite superiore per il valore di
  \var{attr->mq\_maxmsg} in \func{mq\_open}. Il suo valore di default è 10. Il
  valore massimo è \constd{HARD\_MAX} che vale \code{(131072/sizeof(void *))},
  ed il valore minimo 1 (ma era 10 per i kernel precedenti il 2.6.28). Questo
  limite viene ignorato per i processi con privilegi amministrativi (più
  precisamente con la \textit{capability} \const{CAP\_SYS\_RESOURCE}) ma
  \const{HARD\_MAX} resta comunque non superabile.

\item[\sysctlfiled{fs/mqueue/msgsize\_max}] Indica il valore massimo della
  dimensione in byte di un messaggio sulla coda ed agisce come limite
  superiore per il valore di \var{attr->mq\_msgsize} in \func{mq\_open}. Il
  suo valore di default è 8192.  Il valore massimo è 1048576 ed il valore
  minimo 128 (ma per i kernel precedenti il 2.6.28 detti limiti erano
  rispettivamente \const{INT\_MAX} e 8192). Questo limite viene ignorato dai
  processi con privilegi amministrativi (con la \textit{capability}
  \const{CAP\_SYS\_RESOURCE}).

\item[\sysctlfiled{fs/mqueue/queues\_max}] Indica il numero massimo di code di
  messaggi creabili in totale sul sistema, il valore di default è 256 ma si
  può usare un valore qualunque fra $0$ e \const{INT\_MAX}. Il limite non
  viene applicato ai processi con privilegi amministrativi (cioè con la
  \textit{capability} \const{CAP\_SYS\_RESOURCE}).

\end{basedescript}

Infine sulle code di messaggi si applica il limite imposto sulla risorsa
\const{RLIMIT\_MSGQUEUE} (vedi sez.~\ref{sec:sys_resource_limit}) che indica
lo spazio massimo (in byte) occupabile dall'insieme di tutte le code di
messaggi appartenenti ai processi di uno stesso utente, che viene identificato
in base al \textit{real user ID} degli stessi.

Quando l'accesso alla coda non è più necessario si può chiudere il relativo
descrittore con la funzione \funcd{mq\_close}, il cui prototipo è:

\begin{funcproto}{
\fhead{mqueue.h}
\fdecl{int mq\_close(mqd\_t mqdes)}

\fdesc{Chiude una coda di messaggi.}
}

{La funzione ritorna $0$ in caso di successo e $-1$ per un errore, nel qual
  caso \var{errno} assumerà uno dei valori \errval{EBADF} o \errval{EINTR} nel
  loro significato generico.
}  
\end{funcproto}

La funzione è analoga a \func{close},\footnote{su Linux, dove le code sono
  implementate come file su un filesystem dedicato, è esattamente la stessa
  funzione, per cui non esiste una \textit{system call} autonoma e la funzione
  viene rimappata su \func{close} dalla \acr{glibc}.}  dopo la sua esecuzione
il processo non sarà più in grado di usare il descrittore della coda, ma
quest'ultima continuerà ad esistere nel sistema e potrà essere acceduta con
un'altra chiamata a \func{mq\_open}. All'uscita di un processo tutte le code
aperte, così come i file, vengono chiuse automaticamente. Inoltre se il
processo aveva agganciato una richiesta di notifica sul descrittore che viene
chiuso, questa sarà rilasciata e potrà essere richiesta da qualche altro
processo.

Quando si vuole effettivamente rimuovere una coda dal sistema occorre usare la
funzione di sistema \funcd{mq\_unlink}, il cui prototipo è:

\begin{funcproto}{
\fhead{mqueue.h}
\fdecl{int mq\_unlink(const char *name)}

\fdesc{Rimuove una coda di messaggi.}
}

{La funzione ritorna $0$ in caso di successo e $-1$ per un errore, nel qual
  caso \var{errno} assumerà gli uno dei valori:
  \begin{errlist}
    \item[\errcode{EACCES}] non si hanno i permessi per cancellare la coda.
    \item[\errcode{ENAMETOOLONG}] il nome indicato è troppo lungo.
    \item[\errcode{ENOENT}] non esiste una coda con il nome indicato.
  \end{errlist}
}  
\end{funcproto}

Anche in questo caso il comportamento della funzione è analogo a quello di
\func{unlink} per i file, la funzione rimuove la coda \param{name} (ed il
relativo file sotto \texttt{/dev/mqueue}), così che una successiva chiamata a
\func{mq\_open} fallisce o crea una coda diversa.

% TODO, verificare se mq_unlink è davvero una system call indipendente.

Come per i file ogni coda di messaggi ha un contatore di riferimenti, per cui
la coda non viene effettivamente rimossa dal sistema fin quando questo non si
annulla. Pertanto anche dopo aver eseguito con successo \func{mq\_unlink} la
coda resterà accessibile a tutti i processi che hanno un descrittore aperto su
di essa.  Allo stesso modo una coda ed i suoi contenuti resteranno disponibili
all'interno del sistema anche quando quest'ultima non è aperta da nessun
processo (questa è una delle differenze più rilevanti nei confronti di
\textit{pipe} e \textit{fifo}).  La sola differenza fra code di messaggi POSIX
e file normali è che, essendo il filesystem delle code di messaggi virtuale, e
basato su oggetti interni al kernel, il suo contenuto viene perduto con il
riavvio del sistema.

Come accennato ad ogni coda di messaggi è associata una struttura
\struct{mq\_attr}, che può essere letta e modificata attraverso le due
funzioni \funcd{mq\_getattr} e \funcd{mq\_setattr}, i cui prototipi sono:

\begin{funcproto}{
\fhead{mqueue.h}
\fdecl{int mq\_getattr(mqd\_t mqdes, struct mq\_attr *mqstat)}
\fdesc{Legge gli attributi di una coda di messaggi POSIX.}
\fdecl{int mq\_setattr(mqd\_t mqdes, const struct mq\_attr *mqstat,
    struct mq\_attr *omqstat)}
\fdesc{Modifica gli attributi di una coda di messaggi POSIX.}
}
{
Entrambe le funzioni ritornano $0$ in caso di successo e $-1$ per un errore,
  nel qual caso \var{errno} assumerà i valori \errval{EBADF}
    o \errval{EINVAL} nel loro significato generico.
}  
\end{funcproto}

La funzione \func{mq\_getattr} legge i valori correnti degli attributi della
coda \param{mqdes} nella struttura \struct{mq\_attr} puntata
da \param{mqstat}; di questi l'unico relativo allo stato corrente della coda è
\var{mq\_curmsgs} che indica il numero di messaggi da essa contenuti, gli
altri indicano le caratteristiche generali della stessa impostate in fase di
apertura.

La funzione \func{mq\_setattr} permette di modificare gli attributi di una
coda (indicata da \param{mqdes}) tramite i valori contenuti nella struttura
\struct{mq\_attr} puntata da \param{mqstat}, ma può essere modificato solo il
campo \var{mq\_flags}, gli altri campi vengono comunque ignorati.

In particolare i valori di \var{mq\_maxmsg} e \var{mq\_msgsize} possono essere
specificati solo in fase ci creazione della coda.  Inoltre i soli valori
possibili per \var{mq\_flags} sono 0 e \const{O\_NONBLOCK}, per cui alla fine
la funzione può essere utilizzata solo per abilitare o disabilitare la
modalità non bloccante. L'argomento \param{omqstat} viene usato, quando
diverso da \val{NULL}, per specificare l'indirizzo di una struttura su cui
salvare i valori degli attributi precedenti alla chiamata della funzione.

Per inserire messaggi su di una coda sono previste due funzioni di sistema,
\funcd{mq\_send} e \funcd{mq\_timedsend}. In realtà su Linux la \textit{system
  call} è soltanto \func{mq\_timedsend}, mentre \func{mq\_send} viene
implementata come funzione di libreria che si appoggia alla
precedente. Inoltre \func{mq\_timedsend} richiede che sia definita la macro
\macro{\_XOPEN\_SOURCE} ad un valore pari ad almeno \texttt{600} o la macro
\macro{\_POSIX\_C\_SOURCE} ad un valore uguale o maggiore di \texttt{200112L}.
I rispettivi prototipi sono:

\begin{funcproto}{
\fhead{mqueue.h}
\fdecl{int mq\_send(mqd\_t mqdes, const char *msg\_ptr, size\_t msg\_len,
    unsigned int msg\_prio)}
\fdesc{Esegue l'inserimento di un messaggio su una coda.}
\fhead{mqueue.h}
\fhead{time.h}
\fdecl{int mq\_timedsend(mqd\_t mqdes, const char *msg\_ptr, size\_t
    msg\_len, \\ 
\phantom{int mq\_timedsend(}unsigned int msg\_prio, const struct timespec
*abs\_timeout)} 
\fdesc{Esegue l'inserimento di un messaggio su una coda entro un tempo
  specificato}
}

{Entrambe le funzioni ritornano $0$ in caso di successo e $-1$ per un errore,
  nel qual caso \var{errno} assumerà uno dei valori:
  \begin{errlist}
    \item[\errcode{EAGAIN}] si è aperta la coda con \const{O\_NONBLOCK}, e la
      coda è piena.
    \item[\errcode{EBADF}] si specificato un file descriptor non valido.
    \item[\errcode{EINVAL}] si è specificato un valore nullo per
      \param{msg\_len}, o un valore di \param{msg\_prio} fuori dai limiti, o
      un valore non valido per \param{abs\_timeout}.
    \item[\errcode{EMSGSIZE}] la lunghezza del messaggio \param{msg\_len}
      eccede il limite impostato per la coda.
    \item[\errcode{ETIMEDOUT}] l'inserimento del messaggio non è stato
      effettuato entro il tempo stabilito (solo \func{mq\_timedsend}).
  \end{errlist}
  ed inoltre \errval{EBADF} e \errval{EINTR} nel loro significato generico.
}
\end{funcproto}

Entrambe le funzioni richiedono un puntatore ad un buffer in memoria
contenente il testo del messaggio da inserire nella coda \param{mqdes}
nell'argomento \param{msg\_ptr}, e la relativa lunghezza in \param{msg\_len}.
Se quest'ultima eccede la dimensione massima specificata da \var{mq\_msgsize}
le funzioni ritornano immediatamente con un errore di \errcode{EMSGSIZE}.

L'argomento \param{msg\_prio} indica la priorità dell'argomento che essendo
definito come \ctyp{unsigned int} è sempre un intero positivo. I messaggi di
priorità maggiore vengono inseriti davanti a quelli di priorità inferiore, e
quindi saranno riletti per primi. A parità del valore della priorità il
messaggio sarà inserito in coda a tutti quelli che hanno la stessa priorità
che quindi saranno letti con la politica di una \textit{fifo}. Il valore della
priorità non può eccedere il limite di sistema \constd{MQ\_PRIO\_MAX}, che al
momento è pari a 32768.

Qualora la coda sia piena, entrambe le funzioni si bloccano, a meno che non
sia stata selezionata in fase di apertura della stessa la modalità non
bloccante o non si sia impostato il flag \const{O\_NONBLOCK} sul file
descriptor della coda, nel qual caso entrambe ritornano con un codice di
errore di \errcode{EAGAIN}.

La sola differenza fra le due funzioni è che \func{mq\_timedsend}, passato il
tempo massimo impostato con l'argomento \param{abs\_timeout}, ritorna con un
errore di \errcode{ETIMEDOUT}, se invece il tempo è già scaduto al momento
della chiamata e la coda è piena la funzione ritorna immediatamente. Il valore
di \param{abs\_timeout} deve essere specificato come tempo assoluto tramite
una struttura \struct{timespec} (vedi fig.~\ref{fig:sys_timespec_struct})
indicato in numero di secondi e nanosecondi a partire dal 1 gennaio 1970.

Come per l'inserimento, anche per l'estrazione dei messaggi da una coda sono
previste due funzioni di sistema, \funcd{mq\_receive} e
\funcd{mq\_timedreceive}. Anche in questo caso su Linux soltanto
\func{mq\_timedreceive} è effettivamente, una \textit{system call} e per
usarla devono essere definite le opportune macro come per
\func{mq\_timedsend}. I rispettivi prototipi sono:

\begin{funcproto}{
\fhead{mqueue.h} 
\fdecl{ssize\_t mq\_receive(mqd\_t mqdes, char *msg\_ptr, size\_t
    msg\_len, unsigned int *msg\_prio)} 
\fdesc{Effettua la ricezione di un messaggio da una coda.}
\fhead{mqueue.h} 
\fhead{time.h} 
\fdecl{ssize\_t mq\_timedreceive(mqd\_t mqdes, char *msg\_ptr, size\_t
    msg\_len,\\ 
\phantom{ssize\_t  mq\_timedreceive(}unsigned int *msg\_prio, const struct timespec
*abs\_timeout)} 
\fdesc{Riceve un messaggio da una coda entro un limite di tempo.}
}
{Entrambe le funzioni ritornano il numero di byte del messaggio in caso di
  successo e $-1$ per un errore, nel qual caso \var{errno} assumerà uno dei
  valori:
  \begin{errlist}
    \item[\errcode{EAGAIN}] si è aperta la coda con \const{O\_NONBLOCK}, e la
      coda è vuota.
    \item[\errcode{EINVAL}] si è specificato un valore nullo per
      \param{msg\_ptr}, o un valore non valido per \param{abs\_timeout}.
    \item[\errcode{EMSGSIZE}] la lunghezza del messaggio sulla coda eccede il
      valore \param{msg\_len} specificato per la ricezione.
    \item[\errcode{ETIMEDOUT}] la ricezione del messaggio non è stata
      effettuata entro il tempo stabilito.
  \end{errlist}
  ed inoltre \errval{EBADF} o \errval{EINTR} nel loro significato generico.  }
\end{funcproto}

La funzione estrae dalla coda \param{mqdes} il messaggio a priorità più alta,
o il più vecchio fra quelli della stessa priorità. Una volta ricevuto il
messaggio viene tolto dalla coda e la sua dimensione viene restituita come
valore di ritorno; si tenga presente che 0 è una dimensione valida e che la
condizione di errore è indicata soltanto da un valore di
$-1$.\footnote{Stevens in \cite{UNP2} fa notare che questo è uno dei casi in
  cui vale ciò che lo standard \textsl{non} dice, una dimensione nulla
  infatti, pur non essendo citata, non viene proibita.}

Se la dimensione specificata da \param{msg\_len} non è sufficiente a contenere
il messaggio, entrambe le funzioni, al contrario di quanto avveniva nelle code
di messaggi di SysV, ritornano un errore di \errcode{EMSGSIZE} senza estrarre
il messaggio.  È pertanto opportuno eseguire sempre una chiamata a
\func{mq\_getattr} prima di eseguire una ricezione, in modo da ottenere la
dimensione massima dei messaggi sulla coda, per poter essere in grado di
allocare dei buffer sufficientemente ampi per la lettura.

Se si specifica un puntatore per l'argomento \param{msg\_prio} il valore della
priorità del messaggio viene memorizzato all'indirizzo da esso indicato.
Qualora non interessi usare la priorità dei messaggi si può specificare
\var{NULL}, ed usare un valore nullo della priorità nelle chiamate a
\func{mq\_send}.

Si noti che con le code di messaggi POSIX non si ha la possibilità di
selezionare quale messaggio estrarre con delle condizioni sulla priorità, a
differenza di quanto avveniva con le code di messaggi di SysV che permettono
invece la selezione in base al valore del campo \var{mtype}. 

Qualora la coda sia vuota entrambe le funzioni si bloccano, a meno che non si
sia selezionata la modalità non bloccante; in tal caso entrambe ritornano
immediatamente con l'errore \errcode{EAGAIN}. Anche in questo caso la sola
differenza fra le due funzioni è che la seconda non attende indefinitamente e
passato il tempo massimo \param{abs\_timeout} ritorna comunque con un errore
di \errcode{ETIMEDOUT}.

Uno dei problemi sottolineati da Stevens in \cite{UNP2}, comuni ad entrambe le
tipologie di code messaggi, è che non è possibile per chi riceve identificare
chi è che ha inviato il messaggio, in particolare non è possibile sapere da
quale utente esso provenga. Infatti, in mancanza di un meccanismo interno al
kernel, anche se si possono inserire delle informazioni nel messaggio, queste
non possono essere credute, essendo completamente dipendenti da chi lo invia.
Vedremo però come, attraverso l'uso del meccanismo di notifica, sia possibile
superare in parte questo problema.

Una caratteristica specifica delle code di messaggi POSIX è la possibilità di
usufruire di un meccanismo di notifica asincrono; questo può essere attivato
usando la funzione \funcd{mq\_notify}, il cui prototipo è:

\begin{funcproto}{
\fhead{mqueue.h}
\fdecl{int mq\_notify(mqd\_t mqdes, const struct sigevent *notification)}

\fdesc{Attiva il meccanismo di notifica per una coda.}
}

{La funzione ritorna $0$ in caso di successo e $-1$ per un errore, nel qual
  caso \var{errno} assumerà uno dei valori:
  \begin{errlist}
    \item[\errcode{EBADF}] il descrittore non fa riferimento ad una coda di
      messaggi.
    \item[\errcode{EBUSY}] c'è già un processo registrato per la notifica.
    \item[\errcode{EINVAL}] si è richiesto un meccanismo di notifica invalido
      o specificato nella notifica con i segnali il valore di un segnale non
      esistente.
  \end{errlist}
  ed inoltre \errval{ENOMEM} nel suo significato generico.
}  
\end{funcproto}

Il meccanismo di notifica permette di segnalare in maniera asincrona ad un
processo la presenza di dati sulla coda indicata da \param{mqdes}, in modo da
evitare la necessità di bloccarsi nell'attesa. Per far questo un processo deve
registrarsi con la funzione \func{mq\_notify}, ed il meccanismo è disponibile
per un solo processo alla volta per ciascuna coda.

Il comportamento di \func{mq\_notify} dipende dai valori passati con
l'argomento \param{notification}, che è un puntatore ad una apposita struttura
\struct{sigevent}, (definita in fig.~\ref{fig:struct_sigevent}) introdotta
dallo standard POSIX.1b per gestire la notifica di eventi; per altri dettagli
su di essa si può rivedere quanto detto in sez.~\ref{sec:sig_timer_adv} a
proposito dell'uso della stessa struttura per la notifica delle scadenze dei
\textit{timer}.

Attraverso questa struttura si possono impostare le modalità con cui viene
effettuata la notifica nel campo \var{sigev\_notify}, che può assumere i
valori di tab.~\ref{tab:sigevent_sigev_notify}; fra questi la pagina di
manuale riporta soltanto i primi tre, ed inizialmente era possibile solo
\const{SIGEV\_SIGNAL}. Il metodo consigliato è quello di usare
\const{SIGEV\_SIGNAL} usando il campo \var{sigev\_signo} per indicare il quale
segnale deve essere inviato al processo. Inoltre il campo \var{sigev\_value} è
un puntatore ad una struttura \struct{sigval} (definita in
fig.~\ref{fig:sig_sigval}) che permette di restituire al gestore del segnale
un valore numerico o un indirizzo,\footnote{per il suo uso si riveda la
  trattazione fatta in sez.~\ref{sec:sig_real_time} a proposito dei segnali
  \textit{real-time}.} posto che questo sia installato nella forma estesa
vista in sez.~\ref{sec:sig_sigaction}.

La funzione registra il processo chiamante per la notifica se
\param{notification} punta ad una struttura \struct{sigevent} opportunamente
inizializzata, o cancella una precedente registrazione se è \val{NULL}. Dato
che un solo processo alla volta può essere registrato, la funzione fallisce
con \errcode{EBUSY} se c'è un altro processo già registrato. Questo significa
anche che se si registra una notifica con \const{SIGEV\_NONE} il processo non
la riceverà, ma impedirà anche che altri possano registrarsi per poterlo fare.
Si tenga presente inoltre che alla chiusura del descrittore associato alla
coda (e quindi anche all'uscita del processo) ogni eventuale registrazione di
notifica presente viene cancellata.

La notifica del segnale avviene all'arrivo di un messaggio in una coda vuota
(cioè solo se sulla coda non ci sono messaggi) e se non c'è nessun processo
bloccato in una chiamata a \func{mq\_receive}, in questo caso infatti il
processo bloccato ha la precedenza ed il messaggio gli viene immediatamente
inviato, mentre per il meccanismo di notifica tutto funziona come se la coda
fosse rimasta vuota.

Quando un messaggio arriva su una coda vuota al processo che si era registrato
viene inviato il segnale specificato da \code{notification->sigev\_signo}, e
la coda diventa disponibile per una ulteriore registrazione.  Questo comporta
che se si vuole mantenere il meccanismo di notifica occorre ripetere la
registrazione chiamando nuovamente \func{mq\_notify} all'interno del gestore
del segnale di notifica. A differenza della situazione simile che si aveva con
i segnali non affidabili (l'argomento è stato affrontato in
\ref{sec:sig_semantics}) questa caratteristica non configura una \textit{race
  condition} perché l'invio di un segnale avviene solo se la coda è vuota;
pertanto se si vuole evitare di correre il rischio di perdere eventuali
ulteriori segnali inviati nel lasso di tempo che occorre per ripetere la
richiesta di notifica basta avere cura di eseguire questa operazione prima di
estrarre i messaggi presenti dalla coda.

L'invio del segnale di notifica avvalora alcuni campi di informazione
restituiti al gestore attraverso la struttura \struct{siginfo\_t} (definita in
fig.~\ref{fig:sig_siginfo_t}). In particolare \var{si\_pid} viene impostato al
valore del \ids{PID} del processo che ha emesso il segnale, \var{si\_uid}
all'\textsl{user-ID} effettivo, \var{si\_code} a \const{SI\_MESGQ}, e
\var{si\_errno} a 0. Questo ci dice che, se si effettua la ricezione dei
messaggi usando esclusivamente il meccanismo di notifica, è possibile ottenere
le informazioni sul processo che ha inserito un messaggio usando un gestore
per il segnale in forma estesa, di nuovo si faccia riferimento a quanto detto
al proposito in sez.~\ref{sec:sig_sigaction} e sez.~\ref{sec:sig_real_time}.


\subsection{Memoria condivisa}
\label{sec:ipc_posix_shm}

La memoria condivisa è stato il primo degli oggetti di IPC POSIX inserito nel
kernel ufficiale; il supporto a questo tipo di oggetti è realizzato attraverso
il filesystem \texttt{tmpfs}, uno speciale filesystem che mantiene tutti i
suoi contenuti in memoria, che viene attivato abilitando l'opzione
\texttt{CONFIG\_TMPFS} in fase di compilazione del kernel.

Per potere utilizzare l'interfaccia POSIX per la memoria condivisa la
\acr{glibc} (le funzioni sono state introdotte con la versione 2.2) richiede
di compilare i programmi con l'opzione \code{-lrt}; inoltre è necessario che
in \file{/dev/shm} sia montato un filesystem \texttt{tmpfs}; questo di norma
viene fatto aggiungendo una riga del tipo di:
\begin{FileExample}[label=/etc/fstab]
tmpfs   /dev/shm        tmpfs   defaults        0      0
\end{FileExample}
ad \conffile{/etc/fstab}. In realtà si può montare un filesystem
\texttt{tmpfs} dove si vuole, per usarlo come RAM disk, con un comando del
tipo:
\begin{Example}
mount -t tmpfs -o size=128M,nr_inodes=10k,mode=700 tmpfs /mytmpfs
\end{Example}

Il filesystem riconosce, oltre quelle mostrate, le opzioni \texttt{uid} e
\texttt{gid} che identificano rispettivamente utente e gruppo cui assegnarne
la titolarità, e \texttt{nr\_blocks} che permette di specificarne la
dimensione in blocchi, cioè in multipli di \constd{PAGECACHE\_SIZE} che in
questo caso è l'unità di allocazione elementare.

La funzione che permette di aprire un segmento di memoria condivisa POSIX, ed
eventualmente di crearlo se non esiste ancora, è \funcd{shm\_open}; il suo
prototipo è:

\begin{funcproto}{
\fhead{sys/mman.h}
\fhead{sys/stat.h}
\fhead{fcntl.h}
\fdecl{int shm\_open(const char *name, int oflag, mode\_t mode)}

\fdesc{Apre un segmento di memoria condivisa.}
}

{La funzione ritorna un file descriptor in caso di successo e $-1$ per un
  errore, nel qual caso \var{errno} assumerà uno dei valori:
  \begin{errlist}
  \item[\errcode{EACCES}] non si hanno i permessi di aprire il segmento nella
    modalità scelta o si richiesto \const{O\_TRUNC} per un segmento su cui non
    si ha il permesso di scrittura.
  \item[\errcode{EINVAL}] si è utilizzato un nome non valido.
  \end{errlist}
  ed inoltre \errval{EEXIST}, \errval{EMFILE}, \errval{ENAMETOOLONG},
  \errval{ENFILE} e \errval{ENOENT} nello stesso significato che hanno per
  \func{open}.
}  
\end{funcproto}


La funzione apre un segmento di memoria condivisa identificato dal nome
\param{name}. Come già spiegato in sez.~\ref{sec:ipc_posix_generic} questo
nome può essere specificato in forma standard solo facendolo iniziare per
``\texttt{/}'' e senza ulteriori ``\texttt{/}''. Linux supporta comunque nomi
generici, che verranno interpretati prendendo come radice
\file{/dev/shm}.\footnote{occorre pertanto evitare di specificare qualcosa del
  tipo \file{/dev/shm/nome} all'interno di \param{name}, perché questo
  comporta, da parte delle funzioni di libreria, il tentativo di accedere a
  \file{/dev/shm/dev/shm/nome}.}

La funzione è del tutto analoga ad \func{open} ed analoghi sono i valori che
possono essere specificati per \param{oflag}, che deve essere specificato come
maschera binaria comprendente almeno uno dei due valori \const{O\_RDONLY} e
\const{O\_RDWR}; i valori possibili per i vari bit sono quelli visti in
sez.~\ref{sec:file_open_close} dei quali però \func{shm\_open} riconosce solo
i seguenti:
\begin{basedescript}{\desclabelwidth{2.0cm}\desclabelstyle{\nextlinelabel}}
\item[\const{O\_RDONLY}] Apre il file descriptor associato al segmento di
  memoria condivisa per l'accesso in sola lettura.
\item[\const{O\_RDWR}] Apre il file descriptor associato al segmento di
  memoria condivisa per l'accesso in lettura e scrittura.
\item[\const{O\_CREAT}] Necessario qualora si debba creare il segmento di
  memoria condivisa se esso non esiste; in questo caso viene usato il valore
  di \param{mode} per impostare i permessi, che devono essere compatibili con
  le modalità con cui si è aperto il file.
\item[\const{O\_EXCL}] Se usato insieme a \const{O\_CREAT} fa fallire la
  chiamata a \func{shm\_open} se il segmento esiste già, altrimenti esegue la
  creazione atomicamente.
\item[\const{O\_TRUNC}] Se il segmento di memoria condivisa esiste già, ne
  tronca le dimensioni a 0 byte.
\end{basedescript}

In caso di successo la funzione restituisce un file descriptor associato al
segmento di memoria condiviso con le stesse modalità di \func{open} viste in
sez.~\ref{sec:file_open_close}. Inoltre sul file descriptor viene sempre
impostato il flag \const{FD\_CLOEXEC}.  Chiamate effettuate da diversi
processi usando lo stesso nome restituiranno file descriptor associati allo
stesso segmento, così come, nel caso di file ordinari, essi sono associati
allo stesso inode. In questo modo è possibile effettuare una chiamata ad
\func{mmap} sul file descriptor restituito da \func{shm\_open} ed i processi
vedranno lo stesso segmento di memoria condivisa.

Quando il nome non esiste si può creare un nuovo segmento specificando
\const{O\_CREAT}; in tal caso il segmento avrà (così come i nuovi file)
lunghezza nulla. Il nuovo segmento verrà creato con i permessi indicati
da \param{mode} (di cui vengono usati solo i 9 bit meno significativi, non si
applicano pertanto i permessi speciali di sez.~\ref{sec:file_special_perm})
filtrati dal valore dell'\textit{umask} del processo. Come gruppo ed utente
proprietario del segmento saranno presi quelli facenti parte del gruppo
\textit{effective} del processo chiamante.

Dato che un segmento di lunghezza nulla è di scarsa utilità, una vola che lo
si è creato per impostarne la dimensione si dovrà poi usare \func{ftruncate}
(vedi sez.~\ref{sec:file_file_size}) prima di mapparlo in memoria con
\func{mmap}.  Si tenga presente che una volta chiamata \func{mmap} si può
chiudere il file descriptor ad esso associato (semplicemente con
\func{close}), senza che la mappatura ne risenta, e che questa può essere
rimossa usando \func{munmap}.

Come per i file, quando si vuole rimuovere completamente un segmento di
memoria condivisa occorre usare la funzione \funcd{shm\_unlink}, il cui
prototipo è:

\begin{funcproto}{
\fhead{sys/mman.h}
\fdecl{int shm\_unlink(const char *name)}

\fdesc{Rimuove un segmento di memoria condivisa.}
}

{La funzione ritorna $0$ in caso di successo e $-1$ per un errore, nel qual
  caso \var{errno} assumerà uno dei valori:
  \begin{errlist}
  \item[\errcode{EACCES}] non si è proprietari del segmento.
  \end{errlist}
  ed inoltre \errval{ENAMETOOLONG} e \errval{ENOENT}, nel loro significato
  generico.
}  
\end{funcproto}

La funzione è del tutto analoga ad \func{unlink}, e si limita a cancellare il
nome del segmento da \file{/dev/shm}, senza nessun effetto né sui file
descriptor precedentemente aperti con \func{shm\_open}, né sui segmenti già
mappati in memoria; questi verranno cancellati automaticamente dal sistema
solo con le rispettive chiamate a \func{close} e \func{munmap}.  Una volta
eseguita questa funzione però, qualora si richieda l'apertura di un segmento
con lo stesso nome, la chiamata a \func{shm\_open} fallirà, a meno di non aver
usato \const{O\_CREAT}, in quest'ultimo caso comunque si otterrà un file
descriptor che fa riferimento ad un segmento distinto da eventuali precedenti.

Dato che i segmenti di memoria condivisa sono trattati come file del
filesystem \texttt{tmpfs}, si possono usare su di essi, con lo stesso
significato che assumono sui file ordinari, anche funzioni come quelle delle
famiglie \func{fstat}, \func{fchown} e \func{fchmod}. Inoltre a partire dal
kernel 2.6.19 per i permessi sono supportate anche le ACL illustrate in
sez.~\ref{sec:file_ACL}.

Come esempio dell'uso delle funzioni attinenti ai segmenti di memoria
condivisa POSIX, vediamo come è possibile riscrivere una interfaccia
semplificata analoga a quella vista in fig.~\ref{fig:ipc_sysv_shm_func} per la
memoria condivisa in stile SysV. Il codice completo, di cui si sono riportate
le parti essenziali in fig.~\ref{fig:ipc_posix_shmmem}, è contenuto nel file
\file{SharedMem.c} dei sorgenti allegati.

\begin{figure}[!htb]
  \footnotesize \centering
  \begin{minipage}[c]{\codesamplewidth}
    \includecodesample{listati/MemShared.c}
  \end{minipage} 
  \normalsize 
  \caption{Il codice delle funzioni di gestione dei segmenti di memoria
    condivisa POSIX.}
  \label{fig:ipc_posix_shmmem}
\end{figure}

La prima funzione (\texttt{\small 1-24}) è \func{CreateShm} che, dato un nome
nell'argomento \var{name} crea un nuovo segmento di memoria condivisa,
accessibile in lettura e scrittura, e ne restituisce l'indirizzo. Anzitutto si
definiscono (\texttt{\small 8}) i flag per la successiva (\texttt{\small 9})
chiamata a \func{shm\_open}, che apre il segmento in lettura e scrittura
(creandolo se non esiste, ed uscendo in caso contrario) assegnandogli sul
filesystem i permessi specificati dall'argomento \var{perm}. 

In caso di errore (\texttt{\small 10-12}) si restituisce un puntatore nullo,
altrimenti si prosegue impostando (\texttt{\small 14}) la dimensione del
segmento con \func{ftruncate}. Di nuovo (\texttt{\small 15-16}) si esce
immediatamente restituendo un puntatore nullo in caso di errore. Poi si passa
(\texttt{\small 18}) a mappare in memoria il segmento con \func{mmap}
specificando dei diritti di accesso corrispondenti alla modalità di apertura.
Di nuovo si restituisce (\texttt{\small 19-21}) un puntatore nullo in caso di
errore, altrimenti si inizializza (\texttt{\small 22}) il contenuto del
segmento al valore specificato dall'argomento \var{fill} con \func{memset}, e
se ne restituisce (\texttt{\small 23}) l'indirizzo.

La seconda funzione (\texttt{\small 25-40}) è \func{FindShm} che trova un
segmento di memoria condiviso esistente, restituendone l'indirizzo. In questo
caso si apre (\texttt{\small 31}) il segmento con \func{shm\_open} richiedendo
che il segmento sia già esistente, in caso di errore (\texttt{\small 31-33})
si ritorna immediatamente un puntatore nullo.  Ottenuto il file descriptor del
segmento lo si mappa (\texttt{\small 35}) in memoria con \func{mmap},
restituendo (\texttt{\small 36-38}) un puntatore nullo in caso di errore, o
l'indirizzo (\texttt{\small 39}) dello stesso in caso di successo.

La terza funzione (\texttt{\small 40-45}) è \func{RemoveShm}, e serve a
cancellare un segmento di memoria condivisa. Dato che al contrario di quanto
avveniva con i segmenti del \textit{SysV-IPC} gli oggetti allocati nel kernel
vengono rilasciati automaticamente quando nessuna li usa più, tutto quello che
c'è da fare (\texttt{\small 44}) in questo caso è chiamare \func{shm\_unlink},
restituendo al chiamante il valore di ritorno.




\subsection{Semafori}
\label{sec:ipc_posix_sem}

Fino alla serie 2.4.x del kernel esisteva solo una implementazione parziale
dei semafori POSIX che li realizzava solo a livello di \textit{thread} e non
di processi,\footnote{questo significava che i semafori erano visibili solo
  all'interno dei \textit{thread} creati da un singolo processo, e non
  potevano essere usati come meccanismo di sincronizzazione fra processi
  diversi.} fornita attraverso la sezione delle estensioni \textit{real-time}
della \acr{glibc} (quelle che si accedono collegandosi alla libreria
\texttt{librt}). Esisteva inoltre una libreria che realizzava (parzialmente)
l'interfaccia POSIX usando le funzioni dei semafori di \textit{SysV-IPC}
(mantenendo così tutti i problemi sottolineati in
sez.~\ref{sec:ipc_sysv_sem}).

A partire dal kernel 2.5.7 è stato introdotto un meccanismo di
sincronizzazione completamente nuovo, basato sui cosiddetti \textit{futex} (la
sigla sta per \textit{fast user mode mutex}) con il quale è stato possibile
implementare una versione nativa dei semafori POSIX.  Grazie a questo con i
kernel della serie 2.6 e le nuove versioni della \acr{glibc} che usano questa
nuova infrastruttura per quella che viene chiamata \textit{New Posix Thread
  Library}, sono state implementate anche tutte le funzioni dell'interfaccia
dei semafori POSIX.

Anche in questo caso è necessario appoggiarsi alla libreria per le estensioni
\textit{real-time} \texttt{librt}, questo significa che se si vuole utilizzare
questa interfaccia, oltre ad utilizzare gli opportuni file di definizione,
occorrerà compilare i programmi con l'opzione \texttt{-lrt} o con
\texttt{-lpthread} se si usano questi ultimi. 

La funzione che permette di creare un nuovo semaforo POSIX, creando il
relativo file, o di accedere ad uno esistente, è \funcd{sem\_open}, questa
prevede due forme diverse a seconda che sia utilizzata per aprire un semaforo
esistente o per crearne uno nuovi, i relativi prototipi sono:

\begin{funcproto}{
\fhead{semaphore.h}
\fhead{sys/stat.h}
\fhead{fcntl.h}
\fdecl{sem\_t *sem\_open(const char *name, int oflag)}
\fdecl{sem\_t *sem\_open(const char *name, int oflag, mode\_t mode,
    unsigned int value)} 

\fdesc{Crea un semaforo o ne apre uno esistente.}
}
{La funzione ritorna l'indirizzo del semaforo in caso di successo e
  \constd{SEM\_FAILED} per un errore, nel qual caso \var{errno} assumerà uno
  dei valori:
  \begin{errlist}
    \item[\errcode{EACCES}] il semaforo esiste ma non si hanno permessi
      sufficienti per accedervi.
    \item[\errcode{EEXIST}] si sono specificati \const{O\_CREAT} e
      \const{O\_EXCL} ma il semaforo esiste.
    \item[\errcode{EINVAL}] il valore di \param{value} eccede
      \constd{SEM\_VALUE\_MAX} o il nome è solo ``\texttt{/}''.
    \item[\errcode{ENAMETOOLONG}] si è utilizzato un nome troppo lungo.
    \item[\errcode{ENOENT}] non si è usato \const{O\_CREAT} ed il nome
      specificato non esiste.
  \end{errlist}
  ed inoltre \errval{EMFILE}, \errval{ENFILE} ed \errval{ENOMEM} nel loro
  significato generico.

}
\end{funcproto}

L'argomento \param{name} definisce il nome del semaforo che si vuole
utilizzare, ed è quello che permette a processi diversi di accedere allo
stesso semaforo. Questo deve essere specificato nella stessa forma utilizzata
per i segmenti di memoria condivisa, con un nome che inizia con ``\texttt{/}''
e senza ulteriori ``\texttt{/}'', vale a dire nella forma
\texttt{/nome-semaforo}.

Con Linux i file associati ai semafori sono mantenuti nel filesystem virtuale
\texttt{/dev/shm}, e gli viene assegnato automaticamente un nome nella forma
\texttt{sem.nome-semaforo}, si ha cioè una corrispondenza per cui
\texttt{/nome-semaforo} viene rimappato, nella creazione tramite
\func{sem\_open}, su \texttt{/dev/shm/sem.nome-semaforo}. Per questo motivo la
dimensione massima per il nome di un semaforo, a differenza di quanto avviene
per i segmenti di memoria condivisa, è pari a \const{NAME\_MAX}$ - 4$.

L'argomento \param{oflag} è quello che controlla le modalità con cui opera la
funzione, ed è passato come maschera binaria; i bit corrispondono a quelli
utilizzati per l'analogo argomento di \func{open}, anche se dei possibili
valori visti in sez.~\ref{sec:file_open_close} sono utilizzati soltanto
\const{O\_CREAT} e \const{O\_EXCL}.

Se si usa \const{O\_CREAT} si richiede la creazione del semaforo qualora
questo non esista, ed in tal caso occorre utilizzare la seconda forma della
funzione, in cui si devono specificare sia un valore iniziale con l'argomento
\param{value},\footnote{e si noti come così diventa possibile, differenza di
  quanto avviene per i semafori del \textit{SysV-IPC}, effettuare in maniera
  atomica creazione ed inizializzazione di un semaforo usando una unica
  funzione.} che una maschera dei permessi con l'argomento
\param{mode}; se il semaforo esiste già questi saranno semplicemente
ignorati. Usando il flag \const{O\_EXCL} si richiede invece la verifica che il
semaforo non esista, ed usandolo insieme ad \const{O\_CREAT} la funzione
fallisce qualora un semaforo con lo stesso nome sia già presente.

Si tenga presente che, come accennato in sez.~\ref{sec:ipc_posix_generic}, i
semafori usano la semantica standard dei file per quanto riguarda i controlli
di accesso, questo significa che un nuovo semaforo viene sempre creato con
l'\ids{UID} ed il \ids{GID} effettivo del processo chiamante, e che i permessi
indicati con \param{mode} vengono filtrati dal valore della \textit{umask} del
processo.  Inoltre per poter aprire un semaforo è necessario avere su di esso
sia il permesso di lettura che quello di scrittura.

La funzione restituisce in caso di successo un puntatore all'indirizzo del
semaforo con un valore di tipo \ctyp{sem\_t *}, è questo valore che dovrà
essere passato alle altre funzioni per operare sul semaforo stesso, e non sarà
più necessario fare riferimento al nome, che potrebbe anche essere rimosso con
\func{sem\_unlink}.

Una volta che si sia ottenuto l'indirizzo di un semaforo, sarà possibile
utilizzarlo; se si ricorda quanto detto all'inizio di
sez.~\ref{sec:ipc_sysv_sem}, dove si sono introdotti i concetti generali
relativi ai semafori, le operazioni principali sono due, quella che richiede
l'uso di una risorsa bloccando il semaforo e quella che rilascia la risorsa
liberando il semaforo. La prima operazione è effettuata dalla funzione
\funcd{sem\_wait}, il cui prototipo è:

\begin{funcproto}{
\fhead{semaphore.h}
\fdecl{int sem\_wait(sem\_t *sem)}

\fdesc{Blocca un semaforo.}
}

{La funzione ritorna $0$ in caso di successo e $-1$ per un errore, nel qual
  caso \var{errno} assumerà uno dei valori:
  \begin{errlist}
    \item[\errcode{EINTR}] la funzione è stata interrotta da un segnale.
    \item[\errcode{EINVAL}] il semaforo \param{sem} non esiste.
  \end{errlist}
}  
\end{funcproto}

La funzione cerca di decrementare il valore del semaforo indicato dal
puntatore \param{sem}, se questo ha un valore positivo, cosa che significa che
la risorsa è disponibile, la funzione ha successo, il valore del semaforo
viene diminuito di 1 ed essa ritorna immediatamente consentendo la
prosecuzione del processo.

Se invece il valore è nullo la funzione si blocca (fermando l'esecuzione del
processo) fintanto che il valore del semaforo non ritorna positivo (cosa che a
questo punto può avvenire solo per opera di altro processo che rilascia il
semaforo con una chiamata a \func{sem\_post}) così che poi essa possa
decrementarlo con successo e proseguire.

Si tenga presente che la funzione può sempre essere interrotta da un segnale,
nel qual caso si avrà un errore di \errval{EINTR}; inoltre questo avverrà
comunque, anche qualora si fosse richiesta la gestione con la semantica BSD,
installando il gestore del suddetto segnale con l'opzione \const{SA\_RESTART}
(vedi sez.~\ref{sec:sig_sigaction}) per riavviare le \textit{system call}
interrotte.

Della funzione \func{sem\_wait} esistono due varianti che consentono di
gestire diversamente le modalità di attesa in caso di risorsa occupata, la
prima di queste è \funcd{sem\_trywait}, che serve ad effettuare un tentativo
di acquisizione senza bloccarsi; il suo prototipo è:

\begin{funcproto}{
\fhead{semaphore.h} 
\fdecl{int sem\_trywait(sem\_t *sem)}

\fdesc{Tenta di bloccare un semaforo.}
}
{La funzione ritorna $0$ in caso di successo e $-1$ per un errore, nel qual
  caso \var{errno} assumerà uno dei valori:
  \begin{errlist}
    \item[\errcode{EAGAIN}] il semaforo non può essere acquisito senza
      bloccarsi. 
    \item[\errcode{EINVAL}] l'argomento \param{sem} non indica un semaforo
      valido.
  \end{errlist}
}
\end{funcproto}

La funzione è identica a \func{sem\_wait} ed se la risorsa è libera ha lo
stesso effetto, vale a dire che in caso di semaforo diverso da zero la
funzione lo decrementa e ritorna immediatamente; la differenza è che nel caso
in cui il semaforo è occupato essa non si blocca e di nuovo ritorna
immediatamente, restituendo però un errore di \errval{EAGAIN}, così che il
programma possa proseguire.

La seconda variante di \func{sem\_wait} è una estensione specifica che può
essere utilizzata soltanto se viene definita la macro \macro{\_XOPEN\_SOURCE}
ad un valore di almeno 600 o la macro \macro{\_POSIX\_C\_SOURCE} ad un valore
uguale o maggiore di \texttt{200112L} prima di includere
\headfiled{semaphore.h}, la funzione è \funcd{sem\_timedwait}, ed il suo
prototipo è:

\begin{funcproto}{
\fhead{semaphore.h}
\fdecl{int sem\_timedwait(sem\_t *sem, const struct timespec
    *abs\_timeout)}

\fdesc{Blocca un semaforo.}
}

{La funzione ritorna $0$ in caso di successo e $-1$ per un errore, nel qual
  caso \var{errno} assumerà uno dei valori:
  \begin{errlist}
    \item[\errcode{EINTR}] la funzione è stata interrotta da un segnale.
    \item[\errcode{EINVAL}] l'argomento \param{sem} non indica un semaforo
      valido.
    \item[\errcode{ETIMEDOUT}] è scaduto il tempo massimo di attesa. 
  \end{errlist}
}  
\end{funcproto}

Anche in questo caso il comportamento della funzione è identico a quello di
\func{sem\_wait}, ma è possibile impostare un tempo limite per l'attesa
tramite la struttura \struct{timespec} (vedi
fig.~\ref{fig:sys_timespec_struct}) puntata
dall'argomento \param{abs\_timeout}, indicato in secondi e nanosecondi a
partire dalla cosiddetta \textit{Epoch} (00:00:00, 1 January 1970
UTC). Scaduto il limite la funzione ritorna anche se non è possibile acquisire
il semaforo fallendo con un errore di \errval{ETIMEDOUT}.

La seconda funzione principale utilizzata per l'uso dei semafori è quella che
viene usata per rilasciare un semaforo occupato o, in generale, per aumentare
di una unità il valore dello stesso anche qualora non fosse occupato (si
ricordi che in generale un semaforo viene usato come indicatore di un numero
di risorse disponibili). Detta funzione è \funcd{sem\_post} ed il suo
prototipo è:

\begin{funcproto}{
\fhead{semaphore.h}
\fdecl{int sem\_post(sem\_t *sem)}

\fdesc{Rilascia un semaforo.}
}

{La funzione ritorna $0$ in caso di successo e $-1$ per un errore, nel qual
  caso \var{errno} assumerà uno dei valori:
  \begin{errlist}
    \item[\errcode{EINVAL}] l'argomento \param{sem} non indica un semaforo
      valido.
    \item[\errcode{EOVERFLOW}] si superato il massimo valore di un semaforo.
  \end{errlist}
}  
\end{funcproto}

La funzione incrementa di uno il valore corrente del semaforo indicato
dall'argomento \param{sem}, se questo era nullo la relativa risorsa risulterà
sbloccata, cosicché un altro processo (o \textit{thread}) eventualmente
bloccato in una \func{sem\_wait} sul semaforo possa essere svegliato e rimesso
in esecuzione.  Si tenga presente che la funzione è sicura per l'uso
all'interno di un gestore di segnali (si ricordi quanto detto in
sez.~\ref{sec:sig_signal_handler}).

Se invece di operare su un semaforo se ne volesse semplicemente leggere il
valore, si potrà usare la funzione \funcd{sem\_getvalue}, il cui prototipo è:

\begin{funcproto}{
\fhead{semaphore.h}
\fdecl{int sem\_getvalue(sem\_t *sem, int *sval)}

\fdesc{Richiede il valore di un semaforo.}
}
{La funzione ritorna $0$ in caso di successo e $-1$ per un errore, nel qual
  caso \var{errno} assumerà uno dei valori:
  \begin{errlist}
    \item[\errcode{EINVAL}] l'argomento \param{sem} non indica un semaforo
      valido.
  \end{errlist}
}  
\end{funcproto}

La funzione legge il valore del semaforo indicato dall'argomento \param{sem} e
lo restituisce nella variabile intera puntata dall'argomento
\param{sval}. Qualora ci siano uno o più processi bloccati in attesa sul
semaforo lo standard prevede che la funzione possa restituire un valore nullo
oppure il numero di processi bloccati in una \func{sem\_wait} sul suddetto
semaforo; nel caso di Linux vale la prima opzione.

Questa funzione può essere utilizzata per avere un suggerimento sullo stato di
un semaforo, ovviamente non si può prendere il risultato riportato in
\param{sval} che come indicazione, il valore del semaforo infatti potrebbe
essere già stato modificato al ritorno della funzione.

% TODO verificare comportamento sem_getvalue

Una volta che non ci sia più la necessità di operare su un semaforo se ne può
terminare l'uso con la funzione \funcd{sem\_close}, il cui prototipo è:

\begin{funcproto}{
\fhead{semaphore.h}
\fdecl{int sem\_close(sem\_t *sem)}

\fdesc{Chiude un semaforo.}
}

{La funzione ritorna $0$ in caso di successo e $-1$ per un errore, nel qual
  caso \var{errno} assumerà uno dei valori:
  \begin{errlist}
    \item[\errcode{EINVAL}] l'argomento \param{sem} non indica un semaforo
      valido.
  \end{errlist}
}  
\end{funcproto}

La funzione chiude il semaforo indicato dall'argomento \param{sem}, che non
potrà più essere utilizzato nelle altre funzioni. La chiusura comporta anche
che tutte le risorse che il sistema poteva avere assegnato al processo
nell'uso del semaforo vengono immediatamente rilasciate. Questo significa che
un eventuale altro processo bloccato sul semaforo a causa della acquisizione
dello stesso da parte del processo che chiama \func{sem\_close} potrà essere
immediatamente riavviato.

Si tenga presente poi che come avviene per i file, all'uscita di un processo
anche tutti i semafori che questo aveva aperto vengono automaticamente chiusi.
Questo comportamento risolve il problema che si aveva con i semafori del
\textit{SysV IPC} (di cui si è parlato in sez.~\ref{sec:ipc_sysv_sem}) per i
quali le risorse possono restare bloccate. Si tenga infine presente che, a
differenza di quanto avviene per i file, in caso di una chiamata ad
\func{execve} tutti i semafori vengono chiusi automaticamente.

Come per i semafori del \textit{SysV-IPC} anche quelli POSIX hanno una
persistenza di sistema; questo significa che una volta che si è creato un
semaforo con \func{sem\_open} questo continuerà ad esistere fintanto che il
kernel resta attivo (vale a dire fino ad un successivo riavvio) a meno che non
lo si cancelli esplicitamente. Per far questo si può utilizzare la funzione
\funcd{sem\_unlink}, il cui prototipo è:

\begin{funcproto}{
\fhead{semaphore.h}
\fdecl{int sem\_unlink(const char *name)}

\fdesc{Rimuove un semaforo.}
}

{La funzione ritorna $0$ in caso di successo e $-1$ per un errore, nel qual
  caso \var{errno} assumerà uno dei valori:
  \begin{errlist}
    \item[\errcode{EACCES}] non si hanno i permessi necessari a cancellare il
      semaforo.
    \item[\errcode{ENAMETOOLONG}] il nome indicato è troppo lungo.
    \item[\errcode{ENOENT}] il semaforo \param{name} non esiste.
  \end{errlist}
}  
\end{funcproto}

La funzione rimuove il semaforo indicato dall'argomento \param{name}, che
prende un valore identico a quello usato per creare il semaforo stesso con
\func{sem\_open}. Il semaforo viene rimosso dal filesystem immediatamente; ma
il semaforo viene effettivamente cancellato dal sistema soltanto quando tutti
i processi che lo avevano aperto lo chiudono. Si segue cioè la stessa
semantica usata con \func{unlink} per i file, trattata in dettaglio in
sez.~\ref{sec:link_symlink_rename}.

Una delle caratteristiche peculiari dei semafori POSIX è che questi possono
anche essere utilizzati anche in forma anonima, senza necessità di fare
ricorso ad un nome sul filesystem o ad altri indicativi.  In questo caso si
dovrà porre la variabile che contiene l'indirizzo del semaforo in un tratto di
memoria che sia accessibile a tutti i processi in gioco.  La funzione che
consente di inizializzare un semaforo anonimo è \funcd{sem\_init}, il cui
prototipo è:

\begin{funcproto}{
\fhead{semaphore.h}
\fdecl{int sem\_init(sem\_t *sem, int pshared, unsigned int value)}
\fdesc{Inizializza un semaforo anonimo.}
}

{La funzione ritorna $0$ in caso di successo e $-1$ per un errore, nel qual
  caso \var{errno} assumerà uno dei valori:
  \begin{errlist}
    \item[\errcode{EINVAL}] il valore di \param{value} eccede
      \const{SEM\_VALUE\_MAX}.
    \item[\errcode{ENOSYS}] il valore di \param{pshared} non è nullo ed il
      sistema non supporta i semafori per i processi.
  \end{errlist}
}  
\end{funcproto}

La funzione inizializza un semaforo all'indirizzo puntato dall'argomento
\param{sem}, e come per \func{sem\_open} consente di impostare un valore
iniziale con \param{value}. L'argomento \param{pshared} serve ad indicare se
il semaforo deve essere utilizzato dai \textit{thread} di uno stesso processo
(con un valore nullo) o condiviso fra processi diversi (con un valore non
nullo).

Qualora il semaforo debba essere condiviso dai \textit{thread} di uno stesso
processo (nel qual caso si parla di \textit{thread-shared semaphore}),
occorrerà che \param{sem} sia l'indirizzo di una variabile visibile da tutti i
\textit{thread}, si dovrà usare cioè una variabile globale o una variabile
allocata dinamicamente nello \textit{heap}.

Qualora il semaforo debba essere condiviso fra più processi (nel qual caso si
parla di \textit{process-shared semaphore}) la sola scelta possibile per
renderlo visibile a tutti è di porlo in un tratto di memoria condivisa. Questo
potrà essere ottenuto direttamente sia con \func{shmget} (vedi
sez.~\ref{sec:ipc_sysv_shm}) che con \func{shm\_open} (vedi
sez.~\ref{sec:ipc_posix_shm}), oppure, nel caso che tutti i processi in gioco
abbiano un genitore comune, con una mappatura anonima con \func{mmap} (vedi
sez.~\ref{sec:file_memory_map}) a cui essi poi potranno accedere (si ricordi
che i tratti di memoria condivisa vengono mantenuti nei processi figli
attraverso la funzione \func{fork}).

Una volta inizializzato il semaforo anonimo con \func{sem\_init} lo si potrà
utilizzare nello stesso modo dei semafori normali con \func{sem\_wait} e
\func{sem\_post}. Si tenga presente però che inizializzare due volte lo stesso
semaforo può dar luogo ad un comportamento indefinito. 

Qualora non si intenda più utilizzare un semaforo anonimo questo può essere
eliminato dal sistema; per far questo di deve utilizzare una apposita
funzione, \funcd{sem\_destroy}, il cui prototipo è:

\begin{funcproto}{
\fhead{semaphore.h}
\fdecl{int sem\_destroy(sem\_t *sem)}
\fdesc{Elimina un semaforo anonimo.}
}
{La funzione ritorna $0$ in caso di successo e $-1$ per un errore, nel qual
  caso \var{errno} assumerà uno dei valori:
  \begin{errlist}
    \item[\errcode{EINVAL}] l'argomento \param{sem} non indica un semaforo
      valido.
  \end{errlist}
}  
\end{funcproto}

La funzione prende come unico argomento l'indirizzo di un semaforo che deve
essere stato inizializzato con \func{sem\_init}; non deve quindi essere
applicata a semafori creati con \func{sem\_open}. Inoltre si deve essere
sicuri che il semaforo sia effettivamente inutilizzato, la distruzione di un
semaforo su cui sono presenti processi (o \textit{thread}) in attesa (cioè
bloccati in una \func{sem\_wait}) provoca un comportamento indefinito.

Si tenga presente infine che utilizzare un semaforo che è stato distrutto con
\func{sem\_destroy} di nuovo può dare esito a comportamenti indefiniti.  Nel
caso ci si trovi in una tale evenienza occorre reinizializzare il semaforo una
seconda volta con \func{sem\_init}.

Come esempio di uso sia della memoria condivisa che dei semafori POSIX si sono
scritti due semplici programmi con i quali è possibile rispettivamente
monitorare il contenuto di un segmento di memoria condivisa e modificarne il
contenuto. 

\begin{figure}[!htbp]
  \footnotesize \centering
  \begin{minipage}[c]{\codesamplewidth}
    \includecodesample{listati/message_getter.c}
  \end{minipage} 
  \normalsize 
  \caption{Sezione principale del codice del programma
    \file{message\_getter.c}.}
  \label{fig:ipc_posix_sem_shm_message_server}
\end{figure}

Il corpo principale del primo dei due, il cui codice completo è nel file
\file{message\_getter.c} dei sorgenti allegati, è riportato in
fig.~\ref{fig:ipc_posix_sem_shm_message_server}; si è tralasciata la parte che
tratta la gestione delle opzioni a riga di comando (che consentono di
impostare un nome diverso per il semaforo e il segmento di memoria condivisa)
ed il controllo che al programma venga fornito almeno un argomento, contenente
la stringa iniziale da inserire nel segmento di memoria condivisa.

Lo scopo del programma è quello di creare un segmento di memoria condivisa su
cui registrare una stringa, e tenerlo sotto osservazione stampando la stessa
una volta al secondo. Si utilizzerà un semaforo per proteggere l'accesso in
lettura alla stringa, in modo che questa non possa essere modificata
dall'altro programma prima di averla finita di stampare.

La parte iniziale del programma contiene le definizioni (\texttt{\small 1-8})
del gestore del segnale usato per liberare le risorse utilizzate, delle
variabili globali contenenti i nomi di default del segmento di memoria
condivisa e del semaforo (il default scelto è \texttt{messages}), e delle
altre variabili utilizzate dal programma.

Come prima istruzione (\texttt{\small 10}) si è provveduto ad installare un
gestore di segnale che consentirà di effettuare le operazioni di pulizia
(usando la funzione \func{Signal} illustrata in
fig.~\ref{fig:sig_Signal_code}), dopo di che (\texttt{\small 12-16}) si è
creato il segmento di memoria condivisa con la funzione \func{CreateShm} che
abbiamo appena trattato in sez.~\ref{sec:ipc_posix_shm}, uscendo con un
messaggio in caso di errore. 

Si tenga presente che la funzione \func{CreateShm} richiede che il segmento
non sia già presente e fallirà qualora un'altra istanza, o un altro programma
abbia già allocato un segmento con quello stesso nome. Per semplicità di
gestione si è usata una dimensione fissa pari a 256 byte, definita tramite la
costante \texttt{MSGMAXSIZE}.

Il passo successivo (\texttt{\small 17-21}) è quello della creazione del
semaforo che regola l'accesso al segmento di memoria condivisa con
\func{sem\_open}; anche in questo caso si gestisce l'uscita con stampa di un
messaggio in caso di errore. Anche per il semaforo, avendo specificato la
combinazione di flag \code{O\_CREAT|O\_EXCL} come secondo argomento, si esce
qualora fosse già esistente; altrimenti esso verrà creato con gli opportuni
permessi specificati dal terzo argomento, (indicante lettura e scrittura in
notazione ottale). Infine il semaforo verrà inizializzato ad un valore nullo
(il quarto argomento), corrispondete allo stato in cui risulta bloccato.

A questo punto (\texttt{\small 22}) si potrà inizializzare il messaggio posto
nel segmento di memoria condivisa usando la stringa passata come argomento al
programma. Essendo il semaforo stato creato già bloccato non ci si dovrà
preoccupare di eventuali \textit{race condition} qualora il programma di
modifica del messaggio venisse lanciato proprio in questo momento.  Una volta
inizializzato il messaggio occorrerà però rilasciare il semaforo
(\texttt{\small 24-27}) per consentirne l'uso; in tutte queste operazioni si
provvederà ad uscire dal programma con un opportuno messaggio in caso di
errore.

Una volta completate le inizializzazioni il ciclo principale del programma
(\texttt{\small 29-47}) viene ripetuto indefinitamente (\texttt{\small 29})
per stampare sia il contenuto del messaggio che una serie di informazioni di
controllo. Il primo passo (\texttt{\small 30-34}) è quello di acquisire (con
\func{sem\_getvalue}, con uscita in caso di errore) e stampare il valore del
semaforo ad inizio del ciclo; seguito (\texttt{\small 35-36}) dal tempo
corrente.

\begin{figure}[!htb]
  \footnotesize \centering
  \begin{minipage}[c]{\codesamplewidth}
    \includecodesample{listati/HandSigInt.c}
  \end{minipage} 
  \normalsize 
  \caption{Codice del gestore di segnale del programma
    \file{message\_getter.c}.}
  \label{fig:ipc_posix_sem_shm_message_server_handler}
\end{figure}

Prima della stampa del messaggio invece si deve acquisire il semaforo
(\texttt{\small 30-33}) per evitare accessi concorrenti alla stringa da parte
del programma di modifica. Una volta eseguita la stampa (\texttt{\small 41})
il semaforo dovrà essere rilasciato (\texttt{\small 42-45}). Il passo finale
(\texttt{\small 46}) è attendere per un secondo prima di eseguire da capo il
ciclo. 

Per uscire in maniera corretta dal programma sarà necessario fermarlo con una
interruzione da tastiera (\texttt{C-c}), che corrisponde all'invio del segnale
\signal{SIGINT}, per il quale si è installato (\texttt{\small 10}) una
opportuna funzione di gestione, riportata in
fig.~\ref{fig:ipc_posix_sem_shm_message_server_handler}. La funzione è molto
semplice e richiama le funzioni di rimozione sia per il segmento di memoria
condivisa che per il semaforo, garantendo così che possa essere riaperto
ex-novo senza errori in un futuro riutilizzo del comando.

\begin{figure}[!htb]
  \footnotesize \centering
  \begin{minipage}[c]{\codesamplewidth}
    \includecodesample{listati/message_setter.c}
  \end{minipage} 
  \normalsize 
  \caption{Sezione principale del codice del programma
    \file{message\_setter.c}.}
  \label{fig:ipc_posix_sem_shm_message_setter}
\end{figure}

Il secondo programma di esempio è \file{message\_setter.c}, di cui si è
riportato il corpo principale in
fig.~\ref{fig:ipc_posix_sem_shm_message_setter},\footnote{al solito il codice
  completo è nel file dei sorgenti allegati.} dove si è tralasciata, non
essendo significativa per quanto si sta trattando, la parte relativa alla
gestione delle opzioni a riga di comando e degli argomenti, che sono identici
a quelli usati da \file{message\_getter}, con l'unica aggiunta di un'opzione
``\texttt{-t}'' che consente di indicare un tempo di attesa (in secondi) in
cui il programma si ferma tenendo bloccato il semaforo.

Una volta completata la gestione delle opzioni e degli argomenti (ne deve
essere presente uno solo, contenente la nuova stringa da usare come
messaggio), il programma procede (\texttt{\small 10-14}) con l'acquisizione
del segmento di memoria condivisa usando la funzione \func{FindShm} (trattata
in sez.~\ref{sec:ipc_posix_shm}) che stavolta deve già esistere.  Il passo
successivo (\texttt{\small 16-19}) è quello di aprire il semaforo, e a
differenza di \file{message\_getter}, in questo caso si richiede a
\func{sem\_open} che questo esista, passando uno zero come secondo ed unico
argomento.

Una volta completate con successo le precedenti inizializzazioni, il passo
seguente (\texttt{\small 21-24}) è quello di acquisire il semaforo, dopo di
che sarà possibile eseguire la sostituzione del messaggio (\texttt{\small 25})
senza incorrere in possibili \textit{race condition} con la stampa dello
stesso da parte di \file{message\_getter}.

Una volta effettuata la modifica viene stampato (\texttt{\small 26}) il tempo
di attesa impostato con l'opzione ``\texttt{-t}'' dopo di che (\texttt{\small
  27}) viene eseguita la stessa, senza rilasciare il semaforo che resterà
quindi bloccato (causando a questo punto una interruzione delle stampe
eseguite da \file{message\_getter}). Terminato il tempo di attesa si rilascerà
(\texttt{\small 29-32}) il semaforo per poi uscire.

Per verificare il funzionamento dei programmi occorrerà lanciare prima
\file{message\_getter} (lanciare per primo \file{message\_setter} darebbe
luogo ad un errore, non essendo stati creati il semaforo ed il segmento di
memoria condivisa) che inizierà a stampare una volta al secondo il contenuto
del messaggio ed i suoi dati, con qualcosa del tipo:
\begin{Console}
piccardi@hain:~/gapil/sources$  \textbf{./message_getter messaggio}
sem=1, Fri Dec 31 14:12:41 2010
message: messaggio
sem=1, Fri Dec 31 14:12:42 2010
message: messaggio
...
\end{Console}
%$
proseguendo indefinitamente fintanto che non si prema \texttt{C-c} per farlo
uscire. Si noti come il valore del semaforo risulti sempre pari ad 1 (in
quanto al momento esso sarà sempre libero). 

A questo punto si potrà lanciare \file{message\_setter} per cambiare il
messaggio, nel nostro caso per rendere evidente il funzionamento del blocco
richiederemo anche una attesa di 3 secondi, ed otterremo qualcosa del tipo:
\begin{Console}
piccardi@hain:~/gapil/sources$ \textbf{./message_setter -t 3 ciao}
Sleeping for 3 seconds
\end{Console}
%$
dove il programma si fermerà per 3 secondi prima di rilasciare il semaforo e
terminare. L'effetto di tutto ciò si potrà vedere nell'\textit{output} di
\file{message\_getter}, che verrà interrotto per questo stesso tempo, prima di
ricominciare con il nuovo testo:
\begin{Console}
...
sem=1, Fri Dec 31 14:16:27 2010
message: messaggio
sem=1, Fri Dec 31 14:16:28 2010
message: messaggio
sem=0, Fri Dec 31 14:16:29 2010
message: ciao
sem=1, Fri Dec 31 14:16:32 2010
message: ciao
sem=1, Fri Dec 31 14:16:33 2010
message: ciao
...
\end{Console}
%$

E si noterà come nel momento in cui si lancia \file{message\_setter} le stampe
di \file{message\_getter} si bloccheranno, come corretto, dopo aver registrato
un valore nullo per il semaforo.  Il programma infatti resterà bloccato nella
\func{sem\_wait} (quella di riga (\texttt{\small 37}) in
fig.~\ref{fig:ipc_posix_sem_shm_message_server}) fino alla scadenza
dell'attesa di \file{message\_setter} (con l'esecuzione della \func{sem\_post}
della riga (\texttt{\small 29}) di
fig.~\ref{fig:ipc_posix_sem_shm_message_setter}), e riprenderanno con il nuovo
testo alla terminazione di quest'ultimo.


% LocalWords:  like fifo System POSIX RPC Calls Common Object Request Brocker
% LocalWords:  Architecture descriptor kernel unistd int filedes errno EMFILE
% LocalWords:  ENFILE EFAULT BUF sez fig fork Stevens siblings EOF read SIGPIPE
% LocalWords:  EPIPE shell CGI Gateway Interface HTML JPEG URL mime type gs dup
% LocalWords:  barcode PostScript race condition stream BarCodePage WriteMess
% LocalWords:  size PS switch wait popen pclose stdio const char command NULL
% LocalWords:  EINVAL cap fully buffered Ghostscript l'Encapsulated epstopsf of
% LocalWords:  PDF EPS lseek ESPIPE PPM Portable PixMap format pnmcrop PNG pnm
% LocalWords:  pnmmargin png BarCode inode filesystem l'inode mknod mkfifo RDWR
% LocalWords:  ENXIO deadlock client reinviate fortunes fortunefilename daemon
% LocalWords:  FortuneServer FortuneParse FortuneClient pid libgapil  LD librt
% LocalWords:  PATH linker pathname ps tmp killall fortuned crash socket domain
% LocalWords:  socketpair BSD sys protocol sv EAFNOSUPPORT EPROTONOSUPPORT AF
% LocalWords:  EOPNOTSUPP SOCK Process Comunication ipc perm key exec pipefd SZ
% LocalWords:  header ftok proj stat libc SunOS glibc XPG dell'inode number uid
% LocalWords:  cuid cgid gid tab MSG shift group umask seq MSGMNI SEMMNI SHMMNI
% LocalWords:  shmmni msgmni sem sysctl IPCMNI IPCTestId msgget EACCES EEXIST
% LocalWords:  CREAT EXCL EIDRM ENOENT ENOSPC ENOMEM novo proc MSGMAX msgmax ds
% LocalWords:  MSGMNB msgmnb linked list msqid msgid linux msg qnum lspid lrpid
% LocalWords:  rtime ctime qbytes first last cbytes msgctl semctl shmctl ioctl
% LocalWords:  cmd struct buf EPERM RMID msgsnd msgbuf msgp msgsz msgflg EAGAIN
% LocalWords:  NOWAIT EINTR mtype mtext long message sizeof LENGTH ts sleep BIG
% LocalWords:  msgrcv ssize msgtyp NOERROR EXCEPT ENOMSG multiplexing select ls
% LocalWords:  poll polling queue MQFortuneServer write init HandSIGTERM  l'IPC
% LocalWords:  MQFortuneClient mqfortuned mutex risorse' inter semaphore semget
% LocalWords:  nsems SEMMNS SEMMSL semid otime semval sempid semncnt semzcnt nr
% LocalWords:  SEMVMX SEMOPM semop SEMMNU SEMUME SEMAEM semnum union semun arg
% LocalWords:  ERANGE SETALL SETVAL GETALL array GETNCNT GETPID GETVAL GETZCNT
% LocalWords:  sembuf sops unsigned nsops UNDO flg nsop num undo pending semadj
% LocalWords:  sleeper scheduler running next semundo MutexCreate semunion lock
% LocalWords:  MutexFind wrapper MutexRead MutexLock MutexUnlock unlock locking
% LocalWords:  MutexRemove shmget SHMALL SHMMAX SHMMIN shmid shm segsz atime FD
% LocalWords:  dtime lpid cpid nattac shmall shmmax SHMLBA SHMSEG EOVERFLOW brk
% LocalWords:  memory shmat shmdt void shmaddr shmflg SVID RND RDONLY rounded
% LocalWords:  SIGSEGV nattch exit SharedMem ShmCreate memset fill ShmFind home
% LocalWords:  ShmRemove DirMonitor DirProp chdir GaPiL shmptr ipcs NFS SETPIPE
% LocalWords:  ComputeValues ReadMonitor touch SIGTERM dirmonitor unlink fcntl
% LocalWords:  LockFile UnlockFile CreateMutex FindMutex LockMutex SETLKW GETLK
% LocalWords:  UnlockMutex RemoveMutex ReadMutex UNLCK WRLCK RDLCK mapping MAP
% LocalWords:  SHARED ANONYMOUS thread patch names strace system call userid Di
% LocalWords:  groupid Michal Wronski Krzysztof Benedyczak wrona posix mqueue
% LocalWords:  lmqueue gcc mount mqd name oflag attr maxmsg msgsize receive ptr
% LocalWords:  send WRONLY NONBLOCK close mqdes EBADF getattr setattr mqstat to
% LocalWords:  omqstat curmsgs flags timedsend len prio timespec abs EMSGSIZE
% LocalWords:  ETIMEDOUT timedreceive getaddr notify sigevent notification l'I
% LocalWords:  EBUSY sigev SIGNAL signo value sigval siginfo all'userid MESGQ
% LocalWords:  Konstantin Knizhnik futex tmpfs ramfs cache shared swap CONFIG
% LocalWords:  lrt blocks PAGECACHE TRUNC CLOEXEC mmap ftruncate munmap FindShm
% LocalWords:  CreateShm RemoveShm LIBRARY Library libmqueue FAILED has fclose
% LocalWords:  ENAMETOOLONG qualchenome RESTART trywait XOPEN SOURCE timedwait
% LocalWords:  process getvalue sval execve pshared ENOSYS heap PAGE destroy it
% LocalWords:  xffffffff Arrays owner perms Queues used bytes messages device
% LocalWords:  Cannot find such Segments getter Signal MSGMAXSIZE been stable
% LocalWords:  for now it's break Berlin sources Let's an accidental feature fs
% LocalWords:  Larry Wall Escape the Hell William ipctestid Identifier segment
% LocalWords:  violation dell'I SIGINT setter Fri Dec Sleeping seconds ECHILD
% LocalWords:  SysV capability short RESOURCE INFO UNDEFINED EFBIG semtimedop
% LocalWords:  scan HUGETLB huge page NORESERVE copy RLIMIT MEMLOCK REMAP UTC
% LocalWords:  readmon Hierarchy defaults queues MSGQUEUE effective fstat
% LocalWords:  fchown fchmod Epoch January


%%% Local Variables: 
%%% mode: latex
%%% TeX-master: "gapil"
%%% End: 

%% thread.tex
%%
%% Copyright (C) 2007-2018 Simone Piccardi.  Permission is granted to
%% copy, distribute and/or modify this document under the terms of the GNU Free
%% Documentation License, Version 1.1 or any later version published by the
%% Free Software Foundation; with the Invariant Sections being "Un preambolo",
%% with no Front-Cover Texts, and with no Back-Cover Texts.  A copy of the
%% license is included in the section entitled "GNU Free Documentation
%% License".
%%

\chapter{I thread}
\label{cha:threads}


\itindbeg{thread} 

Tratteremo in questo capitolo un modello di programmazione multitasking,
quello dei \textit{thread}, alternativo al modello classico dei processi,
tipico di Unix. Ne esamineremo le caratteristiche, vantaggi e svantaggi, e le
diverse realizzazioni che sono disponibili per Linux; nella seconda parte
tratteremo in dettaglio quella che è l'implementazione principale, che fa
riferimento all'interfaccia standardizzata da POSIX.1e.


\section{Introduzione ai \textit{thread}}
\label{sec:thread_intro}

Questa prima sezione costituisce una introduzione ai \textit{thread} e
tratterà i concetti principali del relativo modello di programmazione,
esamineremo anche quali modelli sono disponibili per Linux, dando una breve
panoramica sulle implementazioni alternative.


\subsection{Una panoramica}
\label{sec:thread_overview}

% riferimenti
% http://vergil.chemistry.gatech.edu/resources/programming/threads.html
% http://math.arizona.edu/~swig/documentation/pthreads/
% http://www.humanfactor.com/pthreads/

Il modello classico dell'esecuzione dei programmi nei sistemi Unix, illustrato
in sez.~\ref{cha:process_interface}, è fondato sui processi. Il modello nasce
per assicurare la massima stabilità al sistema e prevede una rigida
separazione fra i diversi processi, in modo che questi non possano disturbarsi
a vicenda. 

Le applicazioni moderne però sono altamente concorrenti, e necessitano quindi
di un gran numero di processi; questo ha portato a scontrarsi con alcuni
limiti dell'architettura precedente. In genere i fautori del modello di
programmazione a \texttt{thread} sottolineano due problemi connessi all'uso
dei processi:
\begin{itemize}
\item
\item 
\end{itemize}



\subsection{\textit{Thread} e processi}
\label{sec:thread_process}

Per un utilizzo effettivo dei \textit{thread} è sempre opportuno capire se
questi sono davvero adatti allo scopo che ci si pone.


\subsection{Implementazioni alternative}
\label{sec:thread_other}

Vedremo nella prossima sezione le caratteristiche del supporto per i
\textit{thread} fornita dal kernel, ma esistono diversi possibili approcci
alle modalità in cui questi possono essere realizzati. 

% TODO cenni su pth e le implementazioni in userspace


\section{I \textit{thread} e Linux}
\label{sec:linux_thread}

In questa sezione tratteremo le implementazioni dei \textit{thread}
disponibili con Linux che ha visto un radicale cambiamento nel passaggio dalla
serie 2.4 alla serie 2.6, che ha portato alla versione attuale. 

\subsection{I \textit{LinuxThread}}
\label{sec:linux_old_thread}


\subsection{La \textit{Native Thread Posix Library}}
\label{sec:linux_ntpl}





% http://www.gnu.org/software/pth/


\section{Posix \textit{thread}}
\label{sec:thread_posix_intro}


Tratteremo in questa sezione l'interfaccia di programmazione con i
\textit{thread} standardizzata dallo standard POSIX 1.c, che è quella che è
stata seguita anche dalle varie implementazioni dei \textit{thread} realizzate
su Linux, ed in particolare dalla \textit{Native Thread Posix Library} che è
stata integrata con i kernel della serie 2.6 e che fa parte a pieno titolo
della \acr{glibc}.


\subsection{Una panoramica}
\label{sec:pthread_overview}


\subsection{La gestione dei \textit{thread}}
\label{sec:pthread_management}


Benché la funzione sia utilizzabile anche con i processi, tanto che a partire
dalla versione 2.3 della \acr{glibc} viene a sostituire \func{\_exit} (tramite
un \textit{wrapper} che la utilizza al suo posto) per la terminazione di tutti
i \textit{thread} di un processo si deve usare la funzione di sistema
\func{exit\_group}, il cui prototipo è:

\begin{funcproto}{
\fhead{linux/unistd.h}
\fdecl{void exit\_group(int status)}
\fdesc{Termina tutti i \textit{thread} di un processo.} 
}
{La funzione non ha errori e pertanto non ritorna.}
\end{funcproto}

La funzione è sostanzialmente identica alla \textit{system call} \func{\_exit}
ma a differenza di quest'ultima, che termina solo il \textit{thread}
chiamante, termina tutti \textit{thread} del processo. 


\section{La sincronizzazione dei \textit{thread}}
\label{sec:pthread_sync}

\subsection{I \textit{mutex}}
\label{sec:pthread_mutex}


\subsection{Le variabili di condizione}
\label{sec:pthread_cond}


\itindend{thread} 


\subsection{I \textit{thread} e i segnali.}
\label{sec:thread_signal}

% TODO trattare tkill e tgkill per l'invio di segnali a thread, fare un
% capitolo apposito su thread e segnali




% TODO troppe cose, ma segue list di notizie correlate
% aggiunta rt_tgsigqueueinfo con il kernel 2.6.31



% LocalWords:  thread multitasking POSIX sez Posix Library kernel glibc mutex


%%% Local Variables: 
%%% mode: latex
%%% TeX-master: "gapil"
%%% End: 



% Commentare sotto se si genera la prima parte
\part{Programmazione di rete}
\label{part:progr-di-rete}

%% network.tex
%%
%% Copyright (C) 2000-2018 Simone Piccardi.  Permission is granted to
%% copy, distribute and/or modify this document under the terms of the GNU Free
%% Documentation License, Version 1.1 or any later version published by the
%% Free Software Foundation; with the Invariant Sections being "Un preambolo",
%% with no Front-Cover Texts, and with no Back-Cover Texts.  A copy of the
%% license is included in the section entitled "GNU Free Documentation
%% License".
%%

\chapter{Introduzione alla programmazione di rete}
\label{cha:network}

In questo capitolo sarà fatta un'introduzione ai concetti generali che servono
come prerequisiti per capire la programmazione di rete, non tratteremo quindi
aspetti specifici ma faremo una breve introduzione ai modelli più comuni usati
nella programmazione di rete, per poi passare ad un esame a grandi linee dei
protocolli di rete e di come questi sono organizzati e interagiscono.

In particolare, avendo assunto l'ottica di un'introduzione mirata alla
programmazione, ci concentreremo sul gruppo di protocolli più diffuso, il
TCP/IP, che è quello che sta alla base di Internet, avendo cura di
sottolineare i concetti più importanti da conoscere per la scrittura dei
programmi.



\section{Modelli di programmazione}
\label{sec:net_prog_model}

La differenza principale fra un'applicazione di rete e un programma normale è
che quest'ultima per definizione concerne la comunicazione fra processi
diversi, che in generale non girano neanche sulla stessa macchina. Questo già
prefigura un cambiamento completo rispetto all'ottica del programma monolitico
all'interno del quale vengono eseguite tutte le istruzioni, e chiaramente
presuppone un sistema operativo multitasking in grado di eseguire più processi
contemporaneamente.

In questa prima sezione esamineremo brevemente i principali modelli di
programmazione in uso. Ne daremo una descrizione assolutamente generica e
superficiale, che ne illustri le caratteristiche principali, non essendo fra
gli scopi del testo approfondire questi argomenti.

\subsection{Il modello \textit{client-server}}
\label{sec:net_cliserv}

L'architettura fondamentale su cui si basa gran parte della programmazione di
rete sotto Linux (e sotto Unix in generale) è il modello
\textit{client-server} caratterizzato dalla presenza di due categorie di
soggetti, i programmi di servizio, chiamati \textit{server}, che ricevono le
richieste e forniscono le risposte, ed i programmi di utilizzo, detti
\textit{client}.

In generale un server può (di norma deve) essere in grado di rispondere a più
di un client, per cui è possibile che molti programmi possano interagire
contemporaneamente, quello che contraddistingue il modello però è che
l'architettura dell'interazione è sempre nei termini di molti verso uno, il
server, che viene ad assumere un ruolo privilegiato.

Seguono questo modello tutti i servizi fondamentali di Internet, come le
pagine web, la posta elettronica, ftp, telnet, ssh e praticamente ogni
servizio che viene fornito tramite la rete, anche se, come abbiamo visto, il
modello è utilizzato in generale anche per programmi che non fanno
necessariamente uso della rete, come gli esempi che abbiamo usato in
cap.~\ref{cha:IPC} a proposito della comunicazione fra processi nello stesso
sistema.

Normalmente si dividono i server in due categorie principali, e vengono detti
\textsl{concorrenti} o \textsl{iterativi}, sulla base del loro comportamento.
Un \textsl{server iterativo} risponde alla richiesta inviando i dati e resta
occupato e non rispondendo ad ulteriori richieste fintanto che non ha fornito
una risposta alla richiesta. Una volta completata la risposta il server
diventa di nuovo disponibile.

Un \textsl{server concorrente} al momento di trattare la richiesta crea un
processo figlio (o un \textit{thread}) incaricato di fornire i servizi
richiesti, per porsi immediatamente in attesa di ulteriori richieste. In
questo modo, con sistemi multitasking, più richieste possono essere
soddisfatte contemporaneamente. Una volta che il processo figlio ha concluso
il suo lavoro esso di norma viene terminato, mentre il server originale resta
sempre attivo.


\subsection{Il modello \textit{peer-to-peer}}
\label{sec:net_peertopeer}

Come abbiamo visto il tratto saliente dell'architettura \textit{client-server}
è quello della preminenza del server rispetto ai client, le architetture
\textit{peer-to-peer} si basano su un approccio completamente opposto che è
quello di non avere nessun programma che svolga un ruolo preminente.

Questo vuol dire che in generale ciascun programma viene ad agire come un nodo
in una rete potenzialmente paritetica; ciascun programma si trova pertanto a
ricevere ed inviare richieste ed a ricevere ed inviare risposte, e non c'è più
la separazione netta dei compiti che si ritrova nelle architetture
\textit{client-server}.

Le architetture \textit{peer-to-peer} sono salite alla ribalta con
l'esplosione del fenomeno Napster, ma gli stessi protocolli di routing sono un
buon esempio di architetture \textit{peer-to-peer}, in cui ciascun nodo,
tramite il demone che gestisce il routing, richiede ed invia informazioni ad
altri nodi.

In realtà in molti casi di architetture classificate come
\textit{peer-to-peer} non è detto che la struttura sia totalmente paritetica e
ci sono parecchi esempi in cui alcuni servizi vengono centralizzati o
distribuiti gerarchicamente, come avveniva per lo stesso Napster, in cui le
ricerche erano effettuate su un server centrale.



\subsection{Il modello \textit{three-tier}}
\label{sec:net_three_tier}

Benché qui sia trattato a parte, il modello \textit{three-tier} in realtà è
una estensione del modello \textit{client-server}. Con il crescere della
quantità dei servizi forniti in rete (in particolare su Internet) ed al numero
di accessi richiesto. Si è così assistito anche ad una notevole crescita di
complessità, in cui diversi servizi venivano ad essere integrati fra di loro.

In particolare sempre più spesso si assiste ad una integrazione di servizi di
database con servizi di web, in cui le pagine vengono costruite dinamicamente
sulla base dei dati contenuti nel database. In tutti questi casi il problema
fondamentale di una architettura \textit{client-server} è che la richiesta di
un servizio da parte di un gran numero di client si scontra con il collo di
bottiglia dell'accesso diretto ad un unico server, con gravi problemi di
scalabilità.

Rispondere a queste esigenze di scalabilità il modello più semplice (chiamato
talvolta \textit{two-tier}) da adottare è stata quello di distribuire il
carico delle richieste su più server identici, mantenendo quindi
sostanzialmente inalterata l'architettura \textit{client-server} originale.

Nel far questo ci si scontra però con gravi problemi di manutenibilità dei
servizi, in particolare per quanto riguarda la sincronizzazione dei dati, e di
inefficienza dell'uso delle risorse. Il problema è particolarmente grave ad
esempio per i database che non possono essere replicati e sincronizzati
facilmente, e che sono molto onerosi, la loro replicazione è costosa e
complessa.

È a partire da queste problematiche che nasce il modello \textit{three-tier},
che si struttura, come dice il nome, su tre livelli. Il primo livello, quello
dei client che eseguono le richieste e gestiscono l'interfaccia con l'utente,
resta sostanzialmente lo stesso del modello \textit{client-server}, ma la
parte server viene suddivisa in due livelli, introducendo un
\textit{middle-tier}, su cui deve appoggiarsi tutta la logica di analisi delle
richieste dei client per ottimizzare l'accesso al terzo livello, che è quello
che si limita a fornire i dati dinamici che verranno usati dalla logica
implementata nel \textit{middle-tier} per eseguire le operazioni richieste dai
client.

In questo modo si può disaccoppiare la logica dai dati, replicando la prima,
che è molto meno soggetta a cambiamenti ed evoluzione, e non soffre di
problemi di sincronizzazione, e centralizzando opportunamente i secondi. In
questo modo si può distribuire il carico ed accedere in maniera efficiente i
dati.


\subsection{Il modello \textit{broadcast}}
\label{sec:net_broadcast}

Uno specifico modello relativo alla programmazione di rete è poi quello in cui
è possibile, invece della classica comunicazione uno ad uno comunque usata in
tutti i modelli precedenti (anche nel \textit{peer-to-peer} la comunicazione è
comunque fra singoli ``\textit{peer}''), una comunicazione da uno a molti.

\itindbeg{broadcast}

Questo modello nasce dal fatto che molte tecnologie di rete (ed in particolare
Ethernet, che è probabilmente la più diffusa) hanno il supporto per effettuare
una comunicazione in cui un nodo qualunque della rete più inviare informazioni
in contemporanea a tutti gli altri. In questo caso si parla di
\textit{broadcast}, utilizzando la nomenclatura usata per le trasmissioni
radio, anche se in realtà questo tipo di comunicazione è eseguibile da un nodo
qualunque per cui tutti quanti possono ricoprire sia il ruolo di trasmettitore
che quello di ricevitore.

\itindbeg{multicast}

In genere si parla di \textit{broadcast} quando la trasmissione uno a molti è
possibile fra qualunque nodo di una rete e gli altri, ed è supportata
direttamente dalla tecnologia di collegamento utilizzata. L'utilizzo di questa
forma di comunicazione da uno a molti però può risultare molto utile anche
quando questo tipo di supporto non è disponibile (come ad esempio su Internet,
dove non si possono contattare tutti i nodi presenti). 

\itindend{broadcast}

In tal caso alcuni protocolli di rete (e quelli usati per Internet sono fra
questi) supportano una variante del\textit{broadcast}, detta
\textit{multicast}, in cui resta possibile fare una comunicazione uno a molti,
in cui una applicazione invia i pacchetti a molte altre, in genere passando
attraverso un opportuno supporto degli apparati ed una qualche forma di
registrazione che consente la distribuzione della cominicazione ai nodi
interessati. 

\itindend{multicast}

Ovviamente i programmi che devono realizzare un tipo di comunicazione di
questo tipo (come ad esempio potrebbero essere quelli che effettuano uno
\textit{streaming} di informazioni) devono rispondere a delle problematiche
del tutto diverse da quelle classiche illustrate nei modelli precedenti, e
costituiscono pertanto un'altra classe completamente a parte.


\section{I protocolli di rete}
\label{sec:net_protocols}

Parlando di reti di computer si parla in genere di un insieme molto vasto ed
eterogeneo di mezzi di comunicazione che vanno dal cavo telefonico, alla fibra
ottica, alle comunicazioni via satellite o via radio; per rendere possibile la
comunicazione attraverso un così variegato insieme di mezzi sono stati
adottati molti protocolli, il più famoso dei quali, quello alla base del
funzionamento di Internet, è il gruppo di protocolli comunemente chiamato
TCP/IP.

\subsection{Il modello ISO/OSI}
\label{sec:net_iso_osi}

Una caratteristica comune dei protocolli di rete è il loro essere strutturati
in livelli sovrapposti; in questo modo ogni protocollo di un certo livello
realizza le sue funzionalità basandosi su un protocollo del livello
sottostante.  Questo modello di funzionamento è stato standardizzato dalla
\textit{International Standards Organization} (ISO) che ha preparato fin dal
1984 il Modello di Riferimento \textit{Open Systems Interconnection} (OSI),
strutturato in sette livelli, secondo quanto riportato in
tab.~\ref{tab:net_osilayers}.

\begin{table}[htb]
  \centering
  \begin{tabular}{|l|c|c|} 
    \hline
    \textbf{Livello} & \multicolumn{2}{|c|}{\textbf{Nome}} \\
    \hline
    \hline
    Livello 7&\textit{Application}  &\textsl{Applicazione}\\ 
    Livello 6&\textit{Presentation} &\textsl{Presentazione} \\ 
    Livello 5&\textit{Session}      &\textsl{Sessione} \\ 
    Livello 4&\textit{Transport}    &\textsl{Trasporto} \\ 
    Livello 3&\textit{Network}      &\textsl{Rete}\\ 
    Livello 2&\textit{DataLink}     &\textsl{Collegamento Dati} \\
    Livello 1&\textit{Physical}   &\textsl{Connessione Fisica} \\
    \hline
\end{tabular}
\caption{I sette livelli del protocollo ISO/OSI.}
\label{tab:net_osilayers}
\end{table}

Il modello ISO/OSI è stato sviluppato in corrispondenza alla definizione della
serie di protocolli X.25 per la commutazione di pacchetto; come si vede è un
modello abbastanza complesso\footnote{infatti per memorizzarne i vari livelli
  è stata creata la frase \textit{All people seem to need data processing}, in
  cui ciascuna parola corrisponde all'iniziale di uno dei livelli.}, tanto che
usualmente si tende a suddividerlo in due parti, secondo lo schema mostrato in
fig.~\ref{fig:net_osi_tcpip_comp}, con un \textit{upper layer} che riguarda
solo le applicazioni, che viene realizzato in \textit{user space}, ed un
\textit{lower layer} in cui si mescolano la gestione fatta dal kernel e le
funzionalità fornite dall'hardware.

Il modello ISO/OSI mira ad effettuare una classificazione completamente
generale di ogni tipo di protocollo di rete; nel frattempo però era stato
sviluppato anche un altro modello, relativo al protocollo TCP/IP, che è quello
su cui è basata Internet, che è diventato uno standard de facto.  Questo
modello viene talvolta chiamato anche modello \textit{DoD} (sigla che sta per
\textit{Department of Defense}), dato che fu sviluppato dall'agenzia ARPA per
il Dipartimento della Difesa Americano.

\begin{figure}[!htb]
  \centering
  \includegraphics[width=12cm]{img/iso_tcp_comp}
  \caption{Struttura a livelli dei protocolli OSI e TCP/IP, con la relative
    corrispondenze e la divisione fra \textit{kernel space} e \textit{user
      space}.}
  \label{fig:net_osi_tcpip_comp}
\end{figure}

La scelta fra quale dei due modelli utilizzare dipende per lo più dai gusti
personali. Come caratteristiche generali il modello ISO/OSI è più teorico e
generico, basato separazioni funzionali, mentre il modello TCP/IP è più vicino
alla separazione concreta dei vari strati del sistema operativo; useremo
pertanto quest'ultimo, anche per la sua maggiore semplicità. Questa semplicità
ha un costo quando si fa riferimento agli strati più bassi, che sono in
effetti descritti meglio dal modello ISO/OSI, in quanto gran parte dei
protocolli di trasmissione hardware sono appunto strutturati sui due livelli
di \textit{Data Link} e \textit{Connection}.


\subsection{Il modello TCP/IP (o DoD)}
\label{sec:net_tcpip_overview}

Così come ISO/OSI anche il modello del TCP/IP è stato strutturato in livelli
(riassunti in tab.~\ref{tab:net_layers}); un confronto fra i due è riportato
in fig.~\ref{fig:net_osi_tcpip_comp} dove viene evidenziata anche la
corrispondenza fra i rispettivi livelli (che comunque è approssimativa) e su
come essi vanno ad inserirsi all'interno del sistema rispetto alla divisione
fra \textit{user space} e \textit{kernel space} spiegata in
sez.~\ref{sec:intro_unix_struct}.\footnote{in realtà è sempre possibile
  accedere dallo \textit{user space}, attraverso una opportuna interfaccia
  (come vedremo in sez.~\ref{sec:sock_sa_packet}), ai livelli inferiori del
  protocollo.}

\begin{table}[htb]
  \centering
  \begin{tabular}{|l|c|c|l|} 
    \hline
    \textbf{Livello} & \multicolumn{2}{|c|}{\textbf{Nome}} & \textbf{Esempi} \\
    \hline
    \hline
    Livello 4 & \textit{Application} & \textsl{Applicazione}& 
                                       Telnet, FTP, ecc. \\ 
    Livello 3 & \textit{Transport}   & \textsl{Trasporto} & TCP, UDP\\ 
    Livello 2 & \textit{Network}     & \textsl{Rete}      & IP, (ICMP, IGMP)\\ 
    Livello 1 & \textit{Link}        & \textsl{Collegamento}& 
                                       Device driver \& scheda di interfaccia\\
    \hline
\end{tabular}
\caption{I quattro livelli del protocollo TCP/IP.}
\label{tab:net_layers}
\end{table}

Come si può notare come il modello TCP/IP è più semplice del modello ISO/OSI
ed è strutturato in soli quattro livelli. Il suo nome deriva dai due
principali protocolli che lo compongono, il TCP (\textit{Trasmission Control
  Protocol}) che copre il livello 3 e l'IP (\textit{Internet Protocol}) che
copre il livello 2. Le funzioni dei vari livelli sono le seguenti:

\begin{basedescript}{\desclabelwidth{2.5cm}\desclabelstyle{\nextlinelabel}}
\item[\textbf{Applicazione}] É relativo ai programmi di interfaccia con la
  rete, in genere questi vengono realizzati secondo il modello client-server
  (vedi sez.~\ref{sec:net_cliserv}), realizzando una comunicazione secondo un
  protocollo che è specifico di ciascuna applicazione.
\item[\textbf{Trasporto}] Fornisce la comunicazione tra le due stazioni
  terminali su cui girano gli applicativi, regola il flusso delle
  informazioni, può fornire un trasporto affidabile, cioè con recupero degli
  errori o inaffidabile. I protocolli principali di questo livello sono il TCP
  e l'UDP.
\item[\textbf{Rete}] Si occupa dello smistamento dei singoli pacchetti su una
  rete complessa e interconnessa, a questo stesso livello operano i protocolli
  per il reperimento delle informazioni necessarie allo smistamento, per lo
  scambio di messaggi di controllo e per il monitoraggio della rete. Il
  protocollo su cui si basa questo livello è IP (sia nella attuale versione,
  IPv4, che nella nuova versione, IPv6).
\item[\textbf{Collegamento}] È responsabile per l'interfacciamento al
  dispositivo elettronico che effettua la comunicazione fisica, gestendo
  l'invio e la ricezione dei pacchetti da e verso l'hardware.
\end{basedescript}

La comunicazione fra due stazioni remote avviene secondo le modalità
illustrate in fig.~\ref{fig:net_tcpip_data_flux}, dove si è riportato il flusso
dei dati reali e i protocolli usati per lo scambio di informazione su ciascun
livello. Si è genericamente indicato \textit{ethernet} per il livello 1, anche
se in realtà i protocolli di trasmissione usati possono essere molti altri.

\begin{figure}[!htb]
  \centering \includegraphics[width=13cm]{img/tcp_data_flux}
  \caption{Strutturazione del flusso dei dati nella comunicazione fra due
    applicazioni attraverso i protocolli della suite TCP/IP.}
  \label{fig:net_tcpip_data_flux}
\end{figure}

Per chiarire meglio la struttura della comunicazione attraverso i vari
protocolli mostrata in fig.~\ref{fig:net_tcpip_data_flux}, conviene prendere in
esame i singoli passaggi fatti per passare da un livello al sottostante,
la procedura si può riassumere nei seguenti passi:
\begin{itemize}
\item Le singole applicazioni comunicano scambiandosi i dati ciascuna secondo
  un suo specifico formato. Per applicazioni generiche, come la posta o le
  pagine web, viene di solito definito ed implementato quello che viene
  chiamato un protocollo di applicazione (esempi possono essere HTTP, POP,
  SMTP, ecc.), ciascuno dei quali è descritto in un opportuno standard, di
  solito attraverso un RFC (l'acronimo RFC sta per
  \itindex{Request~For~Comment~(RFC)} \textit{Request For Comment} ed è la
  procedura attraverso la quale vengono proposti gli standard per Internet).
\item I dati delle applicazioni vengono inviati al livello di trasporto usando
  un'interfaccia opportuna (i \textit{socket}, che esamineremo in dettaglio in
  cap.~\ref{cha:socket_intro}). Qui verranno spezzati in pacchetti di
  dimensione opportuna e inseriti nel protocollo di trasporto, aggiungendo ad
  ogni pacchetto le informazioni necessarie per la sua gestione. Questo
  processo viene svolto direttamente nel kernel, ad esempio dallo stack TCP,
  nel caso il protocollo di trasporto usato sia questo.
\item Una volta composto il pacchetto nel formato adatto al protocollo di
  trasporto usato questo sarà passato al successivo livello, quello di rete,
  che si occupa di inserire le opportune informazioni per poter effettuare
  l'instradamento nella rete ed il recapito alla destinazione finale. In
  genere questo è il livello di IP (Internet Protocol), a cui vengono inseriti
  i numeri IP che identificano i computer su Internet.
\item L'ultimo passo è il trasferimento del pacchetto al driver della
  interfaccia di trasmissione, che si incarica di incapsularlo nel relativo
  protocollo di trasmissione. Questo può avvenire sia in maniera diretta, come
  nel caso di ethernet, in cui i pacchetti vengono inviati sulla linea
  attraverso le schede di rete, che in maniera indiretta con protocolli come
  PPP o SLIP, che vengono usati come interfaccia per far passare i dati su
  altri dispositivi di comunicazione (come la seriale o la parallela).
\end{itemize}


\subsection{Criteri generali dell'architettura del TCP/IP}
\label{sec:net_tcpip_design}

La filosofia architetturale del TCP/IP è semplice: costruire una rete che
possa sopportare il carico in transito, ma permettere ai singoli nodi di
scartare pacchetti se il carico è temporaneamente eccessivo, o se risultano
errati o non recapitabili.

L'incarico di rendere il recapito pacchetti affidabile non spetta al livello
di rete, ma ai livelli superiori. Pertanto il protocollo IP è per sua natura
inaffidabile, in quanto non è assicurata né una percentuale di successo né un
limite sui tempi di consegna dei pacchetti.

È il livello di trasporto che si deve occupare (qualora necessiti) del
controllo del flusso dei dati e del recupero degli errori; questo è realizzato
dal protocollo TCP. La sede principale di "\textit{intelligenza}" della rete è
pertanto al livello di trasporto o ai livelli superiori.

Infine le singole stazioni collegate alla rete non fungono soltanto da punti
terminali di comunicazione, ma possono anche assumere il ruolo di
\textit{router} (\textsl{instradatori}), per l'interscambio di pacchetti da
una rete ad un'altra. Questo rende possibile la flessibilità della rete che è
in grado di adattarsi ai mutamenti delle interconnessioni.

La caratteristica essenziale che rende tutto ciò possibile è la strutturazione
a livelli tramite l'incapsulamento. Ogni pacchetto di dati viene incapsulato
nel formato del livello successivo, fino al livello del collegamento fisico.
In questo modo il pacchetto ricevuto ad un livello \textit{n} dalla stazione
di destinazione è esattamente lo stesso spedito dal livello \textit{n} dalla
sorgente.  Questo rende facile il progettare il software facendo riferimento
unicamente a quanto necessario ad un singolo livello, con la confidenza che
questo poi sarà trattato uniformemente da tutti i nodi della rete.


\section{La struttura del TCP/IP}
\label{sec:net_tpcip}

Come accennato in sez.~\ref{sec:net_protocols} il TCP/IP è un insieme di
protocolli diversi, che operano su 4 livelli diversi. Per gli interessi della
programmazione di rete però sono importanti principalmente i due livelli
centrali, e soprattutto quello di trasporto.

La principale interfaccia usata nella programmazione di rete, quella dei
socket (che vedremo in sez.~\ref{cha:socket_intro}), è infatti un'interfaccia
nei confronti di quest'ultimo.  Questo avviene perché al di sopra del livello
di trasporto i programmi hanno a che fare solo con dettagli specifici delle
applicazioni, mentre al di sotto vengono curati tutti i dettagli relativi alla
comunicazione. È pertanto naturale definire una interfaccia di programmazione
su questo confine, tanto più che è proprio lì (come evidenziato in
fig.~\ref{fig:net_osi_tcpip_comp}) che nei sistemi Unix (e non solo) viene
inserita la divisione fra \textit{kernel space} e \textit{user space}.

In realtà in un sistema Unix è possibile accedere anche agli altri livelli (e
non solo a quello di trasporto) con opportune interfacce di programmazione
(vedi sez.~\ref{sec:sock_sa_packet}), ma queste vengono usate solo quando si
debbano fare applicazioni di sistema per il controllo della rete a basso
livello, di uso quindi molto specialistico.

In questa sezione daremo una descrizione sommaria dei vari protocolli del
TCP/IP, concentrandoci, per le ragioni appena esposte, sul livello di
trasporto.  All'interno di quest'ultimo privilegeremo poi il protocollo TCP,
per il ruolo centrale che svolge nella maggior parte delle applicazioni.


\subsection{Il quadro generale}
\label{sec:net_tcpip_general}

Benché si parli di TCP/IP questa famiglia di protocolli è composta anche da
molti membri. In fig.~\ref{fig:net_tcpip_overview} si è riportato uno schema
che mostra un panorama sui principali protocolli della famiglia, e delle loro
relazioni reciproche e con alcune dalle principali applicazioni che li usano.

\begin{figure}[!htb]
  \centering
  \includegraphics[width=13cm]{img/tcpip_overview}  
  \caption{Panoramica sui vari protocolli che compongono la suite TCP/IP.}
  \label{fig:net_tcpip_overview}
\end{figure}

I vari protocolli riportati in fig.~\ref{fig:net_tcpip_overview} sono i
seguenti:
\begin{basedescript}{\desclabelwidth{1.7cm}\desclabelstyle{\nextlinelabel}}
\item[\textsl{IPv4}] \textit{Internet Protocol version 4}. È quello che
  comunemente si chiama IP. Ha origine negli anni '80 e da allora è la base su
  cui è costruita Internet. Usa indirizzi a 32 bit, e mantiene tutte le
  informazioni di instradamento e controllo per la trasmissione dei pacchetti
  sulla rete; tutti gli altri protocolli della suite (eccetto ARP e RARP, e
  quelli specifici di IPv6) vengono trasmessi attraverso di esso.
\item[\textsl{IPv6}] \textit{Internet Protocol version 6}. È stato progettato
  a metà degli anni '90 per rimpiazzare IPv4. Ha uno spazio di indirizzi
  ampliato 128 bit che consente più gerarchie di indirizzi,
  l'auto-configurazione, ed un nuovo tipo di indirizzi, gli \textit{anycast},
  che consentono di inviare un pacchetto ad una stazione su un certo gruppo.
  Effettua lo stesso servizio di trasmissione dei pacchetti di IPv4 di cui
  vuole essere un sostituto.
\item[\textsl{TCP}] \textit{Trasmission Control Protocol}. È un protocollo
  orientato alla connessione che provvede un trasporto affidabile per un
  flusso di dati bidirezionale fra due stazioni remote. Il protocollo ha cura
  di tutti gli aspetti del trasporto dei dati, come l'\textit{acknowledgment}
  (il ricevuto), i timeout, la ritrasmissione, ecc. È usato dalla maggior
  parte delle applicazioni.
\item[\textsl{UDP}] \textit{User Datagram Protocol}. È un protocollo senza
  connessione, per l'invio di dati a pacchetti. Contrariamente al TCP il
  protocollo non è affidabile e non c'è garanzia che i pacchetti raggiungano
  la loro destinazione, si perdano, vengano duplicati, o abbiano un
  particolare ordine di arrivo.
\item[\textsl{ICMP}] \textit{Internet Control Message Protocol}. È il
  protocollo usato a livello 2 per gestire gli errori e trasportare le
  informazioni di controllo fra stazioni remote e instradatori (cioè fra
  \textit{host} e \textit{router}). I messaggi sono normalmente generati dal
  software del kernel che gestisce la comunicazione TCP/IP, anche se ICMP può
  venire usato direttamente da alcuni programmi come \cmd{ping}. A volte ci
  si riferisce ad esso come ICPMv4 per distinguerlo da ICMPv6.
\item[\textsl{IGMP}] \textit{Internet Group Management Protocol}. É un
  protocollo di livello 2 usato per il \textit{multicast} (vedi
  sez.~\ref{sec:xxx_multicast}).  Permette alle stazioni remote di notificare
  ai router che supportano questa comunicazione a quale gruppo esse
  appartengono.  Come ICMP viene implementato direttamente sopra IP.
\item[\textsl{ARP}] \textit{Address Resolution Protocol}. È il protocollo che
  mappa un indirizzo IP in un indirizzo hardware sulla rete locale. È usato in
  reti di tipo \textit{broadcast} come Ethernet, Token Ring o FDDI che hanno
  associato un indirizzo fisico (il \textit{MAC address}) alla interfaccia, ma
  non serve in connessioni punto-punto.
\item[\textsl{RARP}] \textit{Reverse Address Resolution Protocol}. È il
  protocollo che esegue l'operazione inversa rispetto ad ARP (da cui il nome)
  mappando un indirizzo hardware in un indirizzo IP. Viene usato a volte per
  durante l'avvio per assegnare un indirizzo IP ad una macchina.
\item[\textsl{ICMPv6}] \textit{Internet Control Message Protocol, version 6}.
  Combina per IPv6 le funzionalità di ICMPv4, IGMP e ARP.
\item[\textsl{EGP}] \textit{Exterior Gateway Protocol}. È un protocollo di
  routing usato per comunicare lo stato fra gateway vicini a livello di
  \textsl{sistemi autonomi} (vengono chiamati \textit{autonomous
      systems} i raggruppamenti al livello più alto della rete), con
  meccanismi che permettono di identificare i vicini, controllarne la
  raggiungibilità e scambiare informazioni sullo stato della rete. Viene
  implementato direttamente sopra IP. 
\item[\textsl{OSPF}] \textit{Open Shortest Path First}. È in protocollo di
  routing per router su reti interne, che permette a questi ultimi di
  scambiarsi informazioni sullo stato delle connessioni e dei legami che
  ciascuno ha con gli altri. Viene implementato direttamente sopra IP.
\item[\textsl{GRE}] \textit{Generic Routing Encapsulation}. È un protocollo
  generico di incapsulamento che permette di incapsulare un qualunque altro
  protocollo all'interno di IP. 
\item[\textsl{AH}] \textit{Authentication Header}. Provvede l'autenticazione
  dell'integrità e dell'origine di un pacchetto. È una opzione nativa in IPv6
  e viene implementato come protocollo a sé su IPv4. Fa parte della suite di
  IPSEC che provvede la trasmissione cifrata ed autenticata a livello IP.
\item[\textsl{ESP}] \textit{Encapsulating Security Payload}. Provvede la
  cifratura insieme all'autenticazione dell'integrità e dell'origine di un
  pacchetto. Come per AH è opzione nativa in IPv6 e viene implementato come
  protocollo a sé su IPv4.
\item[\textsl{PPP}] \textit{Point-to-Point Protocol}. È un protocollo a
  livello 1 progettato per lo scambio di pacchetti su connessioni punto punto.
  Viene usato per configurare i collegamenti, definire i protocolli di rete
  usati ed incapsulare i pacchetti di dati. È un protocollo complesso con
  varie componenti.
\item[\textsl{SLIP}] \textit{Serial Line over IP}. È un protocollo di livello
  1 che permette di trasmettere un pacchetto IP attraverso una linea seriale.
\end{basedescript}

Gran parte delle applicazioni comunicano usando TCP o UDP, solo alcune, e per
scopi particolari si rifanno direttamente ad IP (ed i suoi correlati ICMP e
IGMP); benché sia TCP che UDP siano basati su IP e sia possibile intervenire a
questo livello con i \textit{raw socket} questa tecnica è molto meno diffusa e
a parte applicazioni particolari si preferisce sempre usare i servizi messi a
disposizione dai due protocolli precedenti.  Per questo, motivo a parte alcuni
brevi accenni su IP in questa sezione, ci concentreremo sul livello di
trasporto.

\subsection{Internet Protocol (IP)}
\label{sec:net_ip}

Quando si parla di \textit{Internet Protocol} (IP) si fa in genere riferimento
ad una versione (la quarta, da cui il nome IPv4) che è quella più usata
comunemente, anche se ormai si sta diffondendo sempre di più la nuova versione
IPv6. Il protocollo IPv4 venne standardizzato nel 1981
dall'\href{http://www.ietf.org/rfc/rfc0719.txt}{RFC~719}.

Il protocollo IP (indipendentemente dalla versione) nasce per disaccoppiare le
applicazioni della struttura hardware delle reti di trasmissione, e creare una
interfaccia di trasmissione dei dati indipendente dal sottostante substrato di
interconnessione fisica, che può essere realizzato con le tecnologie più
disparate (Ethernet, Token Ring, FDDI, ecc.).  Il compito di IP è pertanto
quello di trasmettere i pacchetti da un computer all'altro della rete; le
caratteristiche essenziali con cui questo viene realizzato in sono due:

\begin{itemize}
\item \textit{Universal addressing} la comunicazione avviene fra due stazioni
  remote identificate univocamente con un indirizzo a 32 bit che può
  appartenere ad una sola interfaccia di rete.
\item \textit{Best effort} viene assicurato il massimo impegno nella
  trasmissione, ma non c'è nessuna garanzia per i livelli superiori né sulla
  percentuale di successo né sul tempo di consegna dei pacchetti di dati.
\end{itemize}

Negli anni '90 la crescita vertiginosa del numero di macchine connesse a
Internet ha iniziato a far emergere i vari limiti di IPv4, per risolverne i
problemi si è perciò definita una nuova versione del protocollo, che (saltando
un numero) è diventata la versione 6. IPv6 nasce quindi come evoluzione di
IPv4, mantenendone inalterate le funzioni che si sono dimostrate valide,
eliminando quelle inutili e aggiungendone poche altre per mantenere il
protocollo il più snello e veloce possibile.

I cambiamenti apportati sono comunque notevoli e si possono essere riassunti a
grandi linee nei seguenti punti:
\begin{itemize}
\item l'espansione delle capacità di indirizzamento e instradamento, per
  supportare una gerarchia con più livelli di indirizzamento, un numero di
  nodi indirizzabili molto maggiore e una auto-configurazione degli indirizzi.
\item l'introduzione un nuovo tipo di indirizzamento, l'\textit{anycast} che
  si aggiunge agli usuali \textit{unicast} e \textit{multicast}.
\item la semplificazione del formato dell'intestazione (\textit{header}) dei
  pacchetti, eliminando o rendendo opzionali alcuni dei campi di IPv4, per
  eliminare la necessità di rielaborazione della stessa da parte dei router e
  contenere l'aumento di dimensione dovuto all'ampliamento degli indirizzi.
\item un supporto per le opzioni migliorato, per garantire una trasmissione
  più efficiente del traffico normale, limiti meno stringenti sulle dimensioni
  delle opzioni, e la flessibilità necessaria per introdurne di nuove in
  futuro.
\item il supporto per delle capacità di \textsl{qualità di servizio} (QoS) che
  permettano di identificare gruppi di dati per i quali si può provvedere un
  trattamento speciale (in vista dell'uso di Internet per applicazioni
  multimediali e/o ``real-time'').
\end{itemize}

Maggiori dettagli riguardo a caratteristiche, notazioni e funzionamento del
protocollo IP sono forniti nell'appendice sez.~\ref{sec:ip_protocol}.

 
\subsection{User Datagram Protocol (UDP)}
\label{sec:net_udp}

Il protocollo UDP è un protocollo di trasporto molto semplice; la sua
descrizione completa è contenuta
dell'\href{http://www.ietf.org/rfc/rfc0768.txt}{RFC~768}, ma in sostanza esso
è una semplice interfaccia al protocollo IP dal livello di trasporto. Quando
un'applicazione usa UDP essa scrive un pacchetto di dati (il cosiddetto
\textit{datagram} che da il nome al protocollo) su un socket, al pacchetto
viene aggiunto un header molto semplice (per una descrizione più accurata vedi
sez.~\ref{sec:udp_protocol}), e poi viene passato al livello superiore (IPv4 o
IPv6 che sia) che lo spedisce verso la destinazione.  Dato che né IPv4 né IPv6
garantiscono l'affidabilità niente assicura che il pacchetto arrivi a
destinazione, né che più pacchetti arrivino nello stesso ordine in cui sono
stati spediti.

Pertanto il problema principale che si affronta quando si usa UDP è la
mancanza di affidabilità, se si vuole essere sicuri che i pacchetti arrivino a
destinazione occorrerà provvedere con l'applicazione, all'interno della quale
si dovrà inserire tutto quanto necessario a gestire la notifica di
ricevimento, la ritrasmissione, il timeout. 

Si tenga conto poi che in UDP niente garantisce che i pacchetti arrivino nello
stesso ordine in cui sono stati trasmessi, e può anche accadere che i
pacchetti vengano duplicati nella trasmissione, e non solo perduti. Di tutto
questo di nuovo deve tenere conto l'applicazione.

Un altro aspetto di UDP è che se un pacchetto raggiunge correttamente la
destinazione esso viene passato all'applicazione ricevente in tutta la sua
lunghezza, la trasmissione avviene perciò per \textit{record} la cui lunghezza
viene anche essa trasmessa all'applicazione all'atto del ricevimento.

Infine UDP è un protocollo che opera senza connessione
(\textit{connectionless}) in quanto non è necessario stabilire nessun tipo di
relazione tra origine e destinazione dei pacchetti. Si hanno così situazioni
in cui un client può scrivere su uno stesso socket pacchetti destinati a
server diversi, o un server ricevere su un socket pacchetti provenienti da
client diversi.  Il modo più semplice di immaginarsi il funzionamento di UDP è
quello della radio, in cui si può \textsl{trasmettere} e \textsl{ricevere} da
più stazioni usando la stessa frequenza.

Nonostante gli evidenti svantaggi comportati dall'inaffidabilità UDP ha il
grande pregio della velocità, che in certi casi è essenziale; inoltre si
presta bene per le applicazioni in cui la connessione non è necessaria, e
costituirebbe solo un peso in termini di prestazioni, mentre una perdita di
pacchetti può essere tollerata: ad esempio le applicazioni di streaming e
quelle che usano il \textit{multicast}.

\subsection{Transport Control Protocol (TCP)}
\label{sec:net_tcp}

Il TCP è un protocollo molto complesso, definito
nell'\href{http://www.ietf.org/rfc/rfc0739.txt}{RFC~739} e completamente
diverso da UDP; alla base della sua progettazione infatti non stanno
semplicità e velocità, ma la ricerca della massima affidabilità possibile
nella trasmissione dei dati.

La prima differenza con UDP è che TCP provvede sempre una connessione diretta
fra un client e un server, attraverso la quale essi possono comunicare; per
questo il paragone più appropriato per questo protocollo è quello del
collegamento telefonico, in quanto prima viene stabilita una connessione fra
due i due capi della comunicazione su cui poi effettuare quest'ultima.

Caratteristica fondamentale di TCP è l'affidabilità; quando i dati vengono
inviati attraverso una connessione ne viene richiesto un ``\textsl{ricevuto}''
(il cosiddetto \textit{acknowlegment}), se questo non arriva essi verranno
ritrasmessi per un determinato numero di tentativi, intervallati da un periodo
di tempo crescente, fino a che sarà considerata fallita o caduta la
connessione (e sarà generato un errore di \textit{timeout}); il periodo di
tempo dipende dall'implementazione e può variare far i quattro e i dieci
minuti.

Inoltre, per tenere conto delle diverse condizioni in cui può trovarsi la
linea di comunicazione, TCP comprende anche un algoritmo di calcolo dinamico
del tempo di andata e ritorno dei pacchetti fra un client e un server (il
cosiddetto RTT, \textit{Round Trip Time}), che lo rende in grado di adattarsi
alle condizioni della rete per non generare inutili ritrasmissioni o cadere
facilmente in timeout.

Inoltre TCP è in grado di preservare l'ordine dei dati assegnando un numero di
sequenza ad ogni byte che trasmette. Ad esempio se un'applicazione scrive 3000
byte su un socket TCP, questi potranno essere spezzati dal protocollo in due
segmenti (le unità di dati passate da TCP a IP vengono chiamate
\textit{segment}) di 1500 byte, di cui il primo conterrà il numero di sequenza
$1-1500$ e il secondo il numero $1501-3000$. In questo modo anche se i
segmenti arrivano a destinazione in un ordine diverso, o se alcuni arrivano
più volte a causa di ritrasmissioni dovute alla perdita degli
\textit{acknowlegment}, all'arrivo sarà comunque possibile riordinare i dati e
scartare i duplicati.

\itindbeg{advertised~window}

Il protocollo provvede anche un controllo di flusso (\textit{flow control}),
cioè specifica sempre all'altro capo della trasmissione quanti dati può
ricevere tramite una \textit{advertised window} (letteralmente
``\textsl{finestra annunciata}''), che indica lo spazio disponibile nel buffer
di ricezione, cosicché nella trasmissione non vengano inviati più dati di
quelli che possono essere ricevuti.

Questa finestra cambia dinamicamente diminuendo con la ricezione dei dati dal
socket ed aumentando con la lettura di quest'ultimo da parte
dell'applicazione, se diventa nulla il buffer di ricezione è pieno e non
verranno accettati altri dati.  Si noti che UDP non provvede niente di tutto
ciò per cui nulla impedisce che vengano trasmessi pacchetti ad un ritmo che il
ricevente non può sostenere.

\itindend{advertised~window}

Infine attraverso TCP la trasmissione è sempre bidirezionale (in inglese si
dice che è \textit{full-duplex}). È cioè possibile sia trasmettere che
ricevere allo stesso tempo, il che comporta che quanto dicevamo a proposito
del controllo di flusso e della gestione della sequenzialità dei dati viene
effettuato per entrambe le direzioni di comunicazione.

% TODO mettere riferimento alla appendice su TCP quando ci sarà
%% Una descrizione più accurata del protocollo è fornita in appendice
%% sez.~\ref{sec:tcp_protocol}.

\subsection{Limiti e dimensioni riguardanti la trasmissione dei dati}
\label{sec:net_lim_dim}

Un aspetto di cui bisogna tenere conto nella programmazione di rete, e che
ritornerà in seguito, quando tratteremo gli aspetti più avanzati, è che ci sono
una serie di limiti a cui la trasmissione dei dati attraverso i vari livelli
del protocollo deve sottostare; limiti che è opportuno tenere presente perché
in certi casi si possono avere delle conseguenze sul comportamento delle
applicazioni.

Un elenco di questi limiti, insieme ad un breve accenno alle loro origini ed
alle eventuali implicazioni che possono avere, è il seguente:
\begin{itemize}
\item La dimensione massima di un pacchetto IP è di 65535 byte, compresa
  l'intestazione. Questo è dovuto al fatto che la dimensione è indicata da un
  campo apposito nell'header di IP che è lungo 16 bit (vedi
  fig.~\ref{fig:IP_ipv4_head}).
\item La dimensione massima di un pacchetto normale di IPv6 è di 65575 byte;
  il campo apposito nell'header infatti è sempre a 16 bit, ma la dimensione
  dell'header è fissa e di 40 byte e non è compresa nel valore indicato dal
  suddetto campo. Inoltre IPv6 ha la possibilità di estendere la dimensione di
  un pacchetto usando la \textit{jumbo payload option}.
\itindbeg{Maximum~Transfer~Unit~(MTU)}
\item Molte reti fisiche hanno una MTU (\textit{Maximum Transfer Unit}) che
  dipende dal protocollo specifico usato al livello di connessione fisica. Il
  più comune è quello di ethernet che è pari a 1500 byte, una serie di altri
  valori possibili sono riportati in tab.~\ref{tab:net_mtu_values}.
\end{itemize}

Quando un pacchetto IP viene inviato su una interfaccia di rete e le sue
dimensioni eccedono la MTU viene eseguita la cosiddetta
\textit{frammentazione}, i pacchetti cioè vengono suddivisi in blocchi più
piccoli che possono essere trasmessi attraverso l'interfaccia.\footnote{questo
  accade sia per IPv4 che per IPv6, anche se i pacchetti frammentati sono
  gestiti con modalità diverse, IPv4 usa un flag nell'header, IPv6 una
  opportuna opzione, si veda sez.~\ref{sec:ipv6_protocol}.}

\begin{table}[!htb]
  \centering
  \begin{tabular}[c]{|l|c|}
    \hline
    \textbf{Rete} & \textbf{MTU} \\
    \hline
    \hline
    Hyperlink & 65535 \\
    Token Ring IBM (16 Mbit/sec) & 17914 \\
    Token Ring IEEE 802.5 (4 Mbit/sec) & 4464 \\
    FDDI & 4532 \\
    Ethernet & 1500 \\
    X.25 & 576 \\
    \hline
  \end{tabular}
  \caption{Valori della MTU (\textit{Maximum Transfer Unit}) per una serie di
    diverse tecnologie di rete.} 
  \label{tab:net_mtu_values}
\end{table}

%TODO aggiornare la tabella con dati più recenti

\itindbeg{Path~MTU}

La MTU più piccola fra due stazioni viene in genere chiamata \textit{path
  MTU}, che dice qual è la lunghezza massima oltre la quale un pacchetto
inviato da una stazione ad un'altra verrebbe senz'altro frammentato. Si tenga
conto che non è affatto detto che la \textit{path MTU} sia la stessa in
entrambe le direzioni, perché l'instradamento può essere diverso nei due
sensi, con diverse tipologie di rete coinvolte.

Una delle differenze fra IPv4 e IPv6 é che per IPv6 la frammentazione può
essere eseguita solo alla sorgente, questo vuol dire che i router IPv6 non
frammentano i pacchetti che ritrasmettono (anche se possono frammentare i
pacchetti che generano loro stessi), al contrario di quanto fanno i router
IPv4. In ogni caso una volta frammentati i pacchetti possono essere
riassemblati solo alla destinazione.

Nell'header di IPv4 è previsto il flag \texttt{DF} che specifica che il
pacchetto non deve essere frammentato; un router che riceva un pacchetto le
cui dimensioni eccedano quelle dell'MTU della rete di destinazione genererà un
messaggio di errore ICMPv4 di tipo \textit{destination unreachable,
  fragmentation needed but DF bit set}.  Dato che i router IPv6 non possono
effettuare la frammentazione la ricezione di un pacchetto di dimensione
eccessiva per la ritrasmissione genererà sempre un messaggio di errore ICMPv6
di tipo \textit{packet too big}.

Dato che il meccanismo di frammentazione e riassemblaggio dei pacchetti
comporta inefficienza, normalmente viene utilizzato un procedimento, detto
\textit{path MTU discovery} che permette di determinare il \textit{path MTU}
fra due stazioni; per la realizzazione del procedimento si usa il flag
\texttt{DF} di IPv4 e il comportamento normale di IPv6 inviando delle
opportune serie di pacchetti (per i dettagli vedere
l'\href{http://www.ietf.org/rfc/rfc1191.txt}{RFC~1191} per IPv4 e
l'\href{http://www.ietf.org/rfc/rfc1981.txt}{RFC~1981} per IPv6) fintanto che
non si hanno più errori.

Il TCP usa sempre questo meccanismo, che per le implementazioni di IPv4 è
opzionale, mentre diventa obbligatorio per IPv6.  Per IPv6 infatti, non
potendo i router frammentare i pacchetti, è necessario, per poter comunicare,
conoscere da subito il \textit{path MTU}.

\itindend{Path~MTU}

Infine il TCP definisce una \textit{Maximum Segment Size} o MSS (vedi
sez.~\ref{sec:tcp_protocol}) che annuncia all'altro capo della connessione la
dimensione massima del segmento di dati che può essere ricevuto, così da
evitare la frammentazione. Di norma viene impostato alla dimensione della MTU
dell'interfaccia meno la lunghezza delle intestazioni di IP e TCP, in Linux il
default, mantenuto nella costante \constd{TCP\_MSS} è 512.

\itindend{Maximum~Transfer~Unit~(MTU)}


%%% Local Variables: 
%%% mode: latex
%%% TeX-master: "gapil"
%%% End: 

% LocalWords:  TCP multitasking client ftp telnet ssh cap thread peer to three
% LocalWords:  Napster routing tier two middle International Standards Systems
% LocalWords:  Organization Interconnection tab Application Presentation All of
% LocalWords:  Session Transport DataLink Physical people seem need processing
% LocalWords:  fig upper layer lower kernel DoD Department Defense Connection
% LocalWords:  sez UDP ICMP IGMP device Trasmission Control Protocol l'IP l'UDP
% LocalWords:  IPv ethernet SMTP RFC Request For Comment socket stack PPP ARP
% LocalWords:  router instradatori version RARP anycast Di
% LocalWords:  l'acknoweledgment Datagram Message host ping ICPMv ICMPv Group
% LocalWords:  multicast Address Resolution broadcast Token FDDI MAC address DF
% LocalWords:  Reverse EGP Exterior Gateway gateway autonomous systems OSPF GRE
% LocalWords:  Shortest Path First Generic Encapsulation Authentication Header
% LocalWords:  IPSEC ESP Encapsulating Security Payload Point Line over raw QoS
% LocalWords:  dall' Universal addressing Best effort unicast header dell' RTT
% LocalWords:  datagram connectionless streaming nell' acknowlegment trip flow
% LocalWords:  segment control advertised window nell'header dell'header option
% LocalWords:  payload MTU Transfer Unit Hyperlink IBM Mbit sec IEEE path but
% LocalWords:  dell'MTU destination unreachable fragmentation needed packet too
% LocalWords:  big discovery MSS Size

%% socket.tex
%%
%% Copyright (C) 2000-2018 Simone Piccardi.  Permission is granted to
%% copy, distribute and/or modify this document under the terms of the GNU Free
%% Documentation License, Version 1.1 or any later version published by the
%% Free Software Foundation; with the Invariant Sections being "Un preambolo",
%% with no Front-Cover Texts, and with no Back-Cover Texts.  A copy of the
%% license is included in the section entitled "GNU Free Documentation
%% License".
%%

\chapter{I socket}
\label{cha:socket_intro}

In questo capitolo inizieremo a spiegare le caratteristiche salienti della
principale interfaccia per la programmazione di rete, quella dei
\textit{socket}, che, pur essendo nata in ambiente Unix, è usata ormai da
tutti i sistemi operativi.

Dopo una breve panoramica sulle caratteristiche di questa interfaccia vedremo
come creare un socket e come collegarlo allo specifico protocollo di rete che
si utilizzerà per la comunicazione. Per evitare un'introduzione puramente
teorica concluderemo il capitolo con un primo esempio di applicazione.

\section{Introduzione ai socket}
\label{sec:sock_overview}

In questa sezione daremo descrizione essenziale di cosa sono i \textit{socket}
e di quali sono i concetti fondamentali da tenere presente quando si ha a che
fare con essi; ne illustreremo poi le caratteristiche e le differenti
tipologie presenti ed infine tratteremo le modalità con cui possono essere
creati.

\index{socket!definizione|(}

\subsection{Cosa sono i \textit{socket}}
\label{sec:sock_socket_def}

I \textit{socket} (una traduzione letterale potrebbe essere \textsl{presa}, ma
essendo universalmente noti come \textit{socket} utilizzeremo sempre la parola
inglese) sono uno dei principali meccanismi di comunicazione utilizzato in
ambito Unix, e li abbiamo brevemente incontrati in
sez.~\ref{sec:ipc_socketpair}, fra i vari meccanismi di intercomunicazione fra
processi. 

Un socket costituisce in sostanza un canale di comunicazione fra due processi
su cui si possono leggere e scrivere dati analogo a quello di una
\textit{pipe} (vedi sez.~\ref{sec:ipc_pipes}) ma, a differenza di questa e
degli altri meccanismi esaminati nel capitolo cap.~\ref{cha:IPC}, i socket non
sono limitati alla comunicazione fra processi che girano sulla stessa
macchina, ma possono realizzare la comunicazione anche attraverso la rete.

Quella dei socket costituisce infatti la principale interfaccia usata nella
programmazione di rete.  La loro origine risale al 1983, quando furono
introdotti in BSD 4.2; l'interfaccia è rimasta sostanzialmente la stessa, con
piccole modifiche, negli anni successivi. Benché siano state sviluppate
interfacce alternative, originate dai sistemi SVr4 come la XTI (\textit{X/Open
  Transport Interface}) nessuna ha mai raggiunto la diffusione e la popolarità
di quella dei socket (né tantomeno la stessa usabilità e flessibilità) ed oggi
sono praticamente dimenticate.

La flessibilità e la genericità dell'interfaccia inoltre consente di
utilizzare i socket con i più disparati meccanismi di comunicazione, e non
solo con l'insieme dei protocolli TCP/IP, anche se questa sarà comunque quella
di cui tratteremo in maniera più estesa.

Per capire il funzionamento dei socket occorre avere presente il funzionamento
dei protocolli di rete che su utilizzeranno (ed in particolare quelli del
TCP/IP già illustrati in sez.~\ref{sec:net_tpcip}), ma l'interfaccia è del
tutto generale e benché le problematiche, e quindi le modalità di risolvere i
problemi, siano diverse a seconda del tipo di protocollo di comunicazione
usato, le funzioni da usare nella gestione dei socket restano le stesse.

Per questo motivo una semplice descrizione dell'interfaccia è assolutamente
inutile, in quanto il comportamento di quest'ultima e le problematiche da
affrontare cambiano radicalmente a seconda del tipo di comunicazione usato.
La scelta di questo tipo di comunicazione (sovente anche detto \textsl{stile})
va infatti ad incidere sulla semantica che verrà utilizzata a livello utente
per gestire la comunicazione cioè su come inviare e ricevere i dati e sul
comportamento effettivo delle funzioni utilizzate.

La scelta di uno \textsl{stile} dipende sia dai meccanismi disponibili, sia
dal tipo di comunicazione che si vuole effettuare. Ad esempio alcuni tipi di
comunicazione considerano i dati come una sequenza continua di byte, in quello
che viene chiamato un \textsl{flusso} (in inglese \textit{stream}), mentre
altri invece li raggruppano in \textsl{pacchetti} (in inglese
\textit{datagram}) che vengono sempre inviati in blocchi separati e non
divisibili.

Un altro esempio delle differenze fra i diversi tipi di comunicazione concerne
la possibilità che essa possa o meno perdere dati nella trasmissione, che
possa o meno rispettare l'ordine in cui i dati inviati e ricevuti, o che possa
accadere di inviare dei pacchetti di dati più volte (differenze che ad esempio
sono presenti nel caso di utilizzo dei protocolli TCP o UDP).

Un terzo esempio di differenza nel tipo di comunicazione concerne il modo in
cui essa avviene nei confronti dei corrispondenti, in certi casi essa può
essere condotta con una connessione diretta con un solo corrispondente, come
per una telefonata; altri casi possono prevedere una comunicazione come per
lettera, in cui si scrive l'indirizzo su ogni pacchetto, altri ancora una
comunicazione uno a molti come il \textit{broadcast} ed il \textit{multicast},
in cui i pacchetti possono venire emessi su appositi ``\textsl{canali}'' dove
chiunque si collega possa riceverli.

É chiaro che ciascuno di questi diversi aspetti è associato ad un tipo di
comunicazione che comporta una modalità diversa di gestire la stessa, ad
esempio se la comunicazione è inaffidabile occorrerà essere in grado di
gestire la perdita o il rimescolamento dei dati, se è a pacchetti questi
dovranno essere opportunamente trattati, se è uno a molti occorrerà tener
conto della eventuale unidirezionalità della stessa, ecc.

\index{socket!definizione|)}


\subsection{La creazione di un socket}
\label{sec:sock_creation}

Come accennato l'interfaccia dei socket è estremamente flessibile e permette
di interagire con protocolli di comunicazione anche molto diversi fra di loro;
in questa sezione vedremo come è possibile creare un socket e come specificare
il tipo di comunicazione che esso deve utilizzare.

La creazione di un socket avviene attraverso l'uso della funzione di sistema 
\funcd{socket}; essa restituisce un \textit{file descriptor} (del tutto
analogo a quelli che si ottengono per i file di dati e le \textit{pipe},
descritti in sez.~\ref{sec:file_fd}) che serve come riferimento al socket; il
suo prototipo è:

\begin{funcproto}{
\fhead{sys/socket.h}
\fdecl{int socket(int domain, int type, int protocol)}
\fdesc{Apre un socket.} 
}

{La funzione ritorna un valore positivo in caso di successo e $-1$ per un
  errore, nel qual caso \var{errno} assumerà uno dei valori:
  \begin{errlist}
  \item[\errcode{EACCES}] non si hanno privilegi per creare un socket nel
    dominio o con il protocollo specificato.
  \item[\errcode{EAFNOSUPPORT}] famiglia di indirizzi non supportata.
  \item[\errcode{EINVAL}] argomento \param{type} invalido.
  \item[\errcode{EMFILE}] si è ecceduta la tabella dei file.
  \item[\errcode{ENFILE}] si è raggiunto il limite massimo di file aperti.
  \item[\errcode{ENOBUFS}] non c'è sufficiente memoria per creare il socket
    (può essere anche \errval{ENOMEM}).
  \item[\errcode{EPROTONOSUPPORT}] il tipo di socket o il protocollo scelto
    non sono supportati nel dominio.
  \end{errlist}
  ed inoltre a seconda del protocollo usato, potranno essere generati altri
  errori, che sono riportati nelle pagine di manuale relative al protocollo.}
\end{funcproto}


La funzione ha tre argomenti, \param{domain} specifica il dominio del socket
(definisce cioè, come vedremo in sez.~\ref{sec:sock_domain}, la famiglia di
protocolli usata), \param{type} specifica il tipo di socket (definisce cioè,
come vedremo in sez.~\ref{sec:sock_type}, lo stile di comunicazione) e
\param{protocol} il protocollo; in genere quest'ultimo è indicato
implicitamente dal tipo di socket, per cui di norma questo valore viene messo
a zero (con l'eccezione dei \textit{raw socket}).

% TODO: l'ultimo argomento viene usato anche dai nuovi ping socket introdotti
% con il kernel 3.0, vedi anche http://lwn.net/Articles/420799/ e
% http://git.kernel.org/?p=linux/kernel/git/torvalds/linux-2.6.git;a=commitdiff;h=c319b4d76b9e583a5d88d6bf190e079c4e43213d 

Si noti che la creazione del socket si limita ad allocare le opportune
strutture nel kernel (sostanzialmente una voce nella \textit{file table}) e
non comporta nulla riguardo all'indicazione degli indirizzi remoti o locali
attraverso i quali si vuole effettuare la comunicazione. Questo significa che
la funzione da sola non è in grado di fornire alcun tipo di comunicazione. 


\subsection{Il dominio dei socket}
\label{sec:sock_domain}

Dati i tanti e diversi protocolli di comunicazione disponibili, esistono vari
tipi di socket, che vengono classificati raggruppandoli in quelli che si
chiamano \textsl{domini}.  La scelta di un dominio equivale in sostanza alla
scelta di una famiglia di protocolli, e viene effettuata attraverso
l'argomento \param{domain} della funzione \func{socket}. Ciascun dominio ha un
suo nome simbolico che convenzionalmente è indicato da una costante che inizia
per \texttt{PF\_}, sigla che sta per \textit{protocol family}, altro nome con
cui si indicano i domini.

A ciascun tipo di dominio corrisponde un analogo nome simbolico, anch'esso
associato ad una costante, che inizia invece per \texttt{AF\_} (da
\textit{address family}) che identifica il formato degli indirizzi usati in
quel dominio. Le pagine di manuale di Linux si riferiscono a questi indirizzi
anche come \textit{name space}, (nome che invece il manuale della \acr{glibc}
riserva a quello che noi abbiamo chiamato domini) dato che identificano il
formato degli indirizzi usati in quel dominio per identificare i capi della
comunicazione.

\begin{table}[!htb]
  \footnotesize
  \centering
  \begin{tabular}[c]{|l|l|l|l|}
       \hline
       \textbf{Nome}&\textbf{Valore}&\textbf{Utilizzo}&\textbf{Man page} \\
       \hline
       \hline
       \constd{AF\_UNSPEC}   & 0& Non specificato               &            \\
       \constd{AF\_LOCAL}    & 1& Local communication           & unix(7)    \\
       \constd{AF\_UNIX}, \constd{AF\_FILE}&1&Sinonimi di \const{AF\_LOCAL}& \\
       \constd{AF\_INET}     & 2& IPv4 Internet protocols       & ip(7)      \\
       \constd{AF\_AX25}     & 3& Amateur radio AX.25 protocol  &            \\
       \constd{AF\_IPX}      & 4& IPX - Novell protocols        &            \\
       \constd{AF\_APPLETALK}& 5& Appletalk                     & ddp(7)     \\
       \constd{AF\_NETROM}   & 6& Amateur radio NetROM          &            \\
       \constd{AF\_BRIDGE}   & 7& Multiprotocol bridge          &            \\
       \constd{AF\_ATMPVC}   & 8& Access to raw ATM PVCs        &            \\
       \constd{AF\_X25}      & 9& ITU-T X.25 / ISO-8208 protocol& x25(7)     \\
       \constd{AF\_INET6}    &10& IPv6 Internet protocols       & ipv6(7)    \\
       \constd{AF\_ROSE}     &11& Amateur Radio X.25 PLP        &            \\
       \constd{AF\_DECnet}   &12& Reserved for DECnet project   &            \\
       \constd{AF\_NETBEUI}  &13& Reserved for 802.2LLC project &            \\
       \constd{AF\_SECURITY} &14& Security callback pseudo AF   &            \\
       \constd{AF\_KEY}      &15& AF\_KEY key management API    &            \\
       \constd{AF\_NETLINK}  &16& Kernel user interface device  & netlink(7) \\
       \constd{AF\_ROUTE}    &16& Sinonimo di \const{AF\_NETLINK} emula BSD.&\\
       \constd{AF\_PACKET}   &17& Low level packet interface    & packet(7)  \\
       \constd{AF\_ASH}      &18& Ash                           &    \\
       \constd{AF\_ECONET}   &19& Acorn Econet                  &    \\
       \constd{AF\_ATMSVC}   &20& ATM SVCs                      &    \\
       \constd{AF\_RDS}      &21& RDS Sockets                   &    \\
       \constd{AF\_SNA}      &22& Linux SNA Project             &    \\
       \constd{AF\_IRDA}     &23& IRDA socket (infrarossi)      & irda(7)    \\
       \constd{AF\_PPPOX}    &24& PPPoX socket                  &    \\
       \constd{AF\_WANPIPE}  &25& Wanpipe API socket            &    \\
       \constd{AF\_LLC}      &26& Linux LLC                     &    \\
       \constd{AF\_IB}       &27&  Native InfiniBand address    &    \\
       \constd{AF\_MPLS}     &28& MPSL                          &    \\
       \constd{AF\_CAN}      &29& Controller Are Network        &    \\
       \constd{AF\_TIPC}     &30& TIPC sockets                  &    \\
       \constd{AF\_BLUETOOTH}&31& Bluetooth socket              &    \\
       \constd{AF\_IUCV}     &32& IUCV sockets                  &    \\
       \constd{AF\_RXRPC}    &33& RxRPC sockets                 &    \\
       \constd{AF\_ISDN}     &34& mISDN sockets                 &    \\
       \constd{AF\_PHONET}   &35& Phonet sockets                &    \\
       \constd{AF\_IEEE802154}&36& IEEE802154 sockets           &    \\
       \constd{AF\_CAIF}     &37& CAIF sockets                  &    \\
       \constd{AF\_ALG}      &38& Algorithm sockets             &    \\
       \constd{AF\_NFC}      &39& NFC sockets                   &    \\
       \constd{AF\_VSOCK}    &40& vSockets                      &    \\
       \constd{AF\_KCM}      &41& Kernel Connection Multiplexor &    \\
       \constd{AF\_QIPCRTR}  &42& Qualcomm IPC Router           &    \\
       \constd{AF\_SMC}      &43& smc sockets                   &    \\
       \hline
  \end{tabular}
  \caption{Famiglie di protocolli definiti in Linux.} 
  \label{tab:net_pf_names}
\end{table}

L'idea alla base della distinzione fra questi due insiemi di costanti era che
una famiglia di protocolli potesse supportare vari tipi di indirizzi, per cui
il prefisso \texttt{PF\_} si sarebbe dovuto usare nella creazione dei socket e
il prefisso \texttt{AF\_} in quello delle strutture degli indirizzi. Questo è
quanto specificato anche dallo standard POSIX.1g, ma non esistono a tuttora
famiglie di protocolli che supportino diverse strutture di indirizzi, per cui
nella pratica questi due nomi sono equivalenti e corrispondono agli stessi
valori numerici.\footnote{in Linux, come si può verificare andando a guardare
  il contenuto di \file{bits/socket.h}, le costanti sono esattamente le stesse
  e ciascuna \texttt{AF\_} è definita alla corrispondente \texttt{PF\_} e con
  lo stesso nome.} Qui si sono indicati i nomi con il prefisso \texttt{AF\_}
seguendo la convenzione usata nelle pagine di manuale.

I domini (e i relativi nomi simbolici), così come i nomi delle famiglie di
indirizzi, sono definiti dall'\textit{header file} \headfiled{socket.h}. Un
elenco, aggiornato alla versione 4.15, delle famiglie di protocolli
disponibili in Linux è riportato in tab.~\ref{tab:net_pf_names}.  L'elenco
indica tutti i protocolli definiti; fra questi però saranno utilizzabili solo
quelli per i quali si è compilato il supporto nel kernel (o si sono caricati
gli opportuni moduli), viene definita anche una costante \constd{AF\_MAX} che
indica il valore massimo associabile ad un dominio.

Si tenga presente che non tutte le famiglie di protocolli sono utilizzabili
dall'utente generico, ad esempio in generale tutti i socket di tipo
\const{SOCK\_RAW} possono essere creati solo da processi che hanno i privilegi
di amministratore (cioè con \ids{UID} effettivo uguale a zero) o dotati della
\textit{capability} \const{CAP\_NET\_RAW}.


\subsection{Il tipo di socket}
\label{sec:sock_type}

La scelta di un dominio non comporta però la scelta dello stile di
comunicazione, questo infatti viene a dipendere dal protocollo che si andrà ad
utilizzare fra quelli disponibili nella famiglia scelta. L'interfaccia dei
socket permette di scegliere lo stile di comunicazione indicando il tipo di
socket con l'argomento \param{type} di \func{socket}. Linux mette a
disposizione vari tipi di socket (che corrispondono a quelli che il manuale
della \acr{glibc} \cite{GlibcMan} chiama \textit{styles}) identificati dalle
seguenti costanti:\footnote{le pagine di manuale POSIX riportano solo i primi
  tre tipi, Linux supporta anche gli altri, come si può verificare nel file
  \texttt{include/linux/net.h} dei sorgenti del kernel.}

\begin{basedescript}{\desclabelwidth{2.0cm}\desclabelstyle{\nextlinelabel}}
\item[\constd{SOCK\_STREAM}] Provvede un canale di trasmissione dati
  bidirezionale, sequenziale e affidabile. Opera su una connessione con un
  altro socket. I dati vengono ricevuti e trasmessi come un flusso continuo di
  byte (da cui il nome \textit{stream}) e possono essere letti in blocchi di
  dimensioni qualunque. Può supportare la trasmissione dei cosiddetti dati
  urgenti (o \textit{out-of-band}, vedi sez.~\ref{sec:TCP_urgent_data}).
\item[\constd{SOCK\_DGRAM}] Viene usato per trasmettere pacchetti di dati
  (\textit{datagram}) di lunghezza massima prefissata, indirizzati
  singolarmente. Non esiste una connessione e la trasmissione è effettuata in
  maniera non affidabile.
\item[\constd{SOCK\_SEQPACKET}] Provvede un canale di trasmissione di dati
  bidirezionale, sequenziale e affidabile. Opera su una connessione con un
  altro socket. I dati possono vengono trasmessi per pacchetti di dimensione
  massima fissata, e devono essere letti integralmente da ciascuna chiamata a
  \func{read}.
\item[\constd{SOCK\_RAW}] Provvede l'accesso a basso livello ai protocolli di
  rete e alle varie interfacce. I normali programmi di comunicazione non
  devono usarlo, è riservato all'uso di sistema.
\item[\constd{SOCK\_RDM}] Provvede un canale di trasmissione di dati
  affidabile, ma in cui non è garantito l'ordine di arrivo dei pacchetti.
\item[\constd{SOCK\_PACKET}] Obsoleto, non deve essere più usato (e pertanto
  non ne parleremo ulteriormente).
\end{basedescript}

A partire dal kernel 2.6.27 l'argomento \param{type} della funzione
\func{socket} assume un significato ulteriore perché può essere utlizzato per
impostare dei flag relativi alle caratteristiche generali del \textit{socket}
non strettamente attinenti all'indicazione del tipo secondo i valori appena
illustrati. Essi infatti possono essere combinati con un OR aritmetico delle
ulteriori costanti:
\begin{basedescript}{\desclabelwidth{2.5cm}\desclabelstyle{\nextlinelabel}}
\item[\constd{SOCK\_CLOEXEC}] imposta il flag di \textit{close-on-exec} sul
  file descriptor del socket, ottenendo lo stesso effetto del flag
  \const{O\_CLOEXEC} di \func{open} (vedi tab.~\ref{tab:open_operation_flag}),
  di cui costituisce l'analogo.

\item[\constd{SOCK\_NONBLOCK}] crea il socket in modalità non-bloccante, con
  effetti identici ad una successiva chiamata a \func{fcntl} per impostare il
  flag di \const{O\_NONBLOCK} sul file descriptor (si faccia di nuovo
  riferimenti al significato di quest'ultimo come spiegato in
  tab.~\ref{tab:open_operation_flag}).
\end{basedescript}

Si tenga presente inoltre che non tutte le combinazioni fra una famiglia di
protocolli e un tipo di socket sono valide, in quanto non è detto che in una
famiglia esista un protocollo per ciascuno dei diversi stili di comunicazione
appena elencati.

\begin{table}[htb]
  \footnotesize
  \centering
  \begin{tabular}{|l|c|c|c|c|c|}
    \hline
    \multicolumn{1}{|c|}{\textbf{Famiglia}}&
    \multicolumn{5}{|c|}{\textbf{Tipo}}\\
    \hline
    \hline
    &\const{SOCK\_STREAM} &\const{SOCK\_DGRAM}     &\const{SOCK\_RAW}& 
      \const{SOCK\_RDM}&\const{SOCK\_SEQPACKET} \\
     \hline
    \const{AF\_UNIX}     &  si & si  &  --  & -- &  si\footnotemark \\
     \hline
    \const{AF\_LOCAL}&\multicolumn{5}{|c|}{sinonimo di \const{AF\_UNIX}}\\
     \hline
    \const{AF\_INET}      & TCP & UDP & IPv4 & --  & --  \\
     \hline
    \const{AF\_INET6}     & TCP & UDP & IPv6 & --  & -- \\
     \hline
    \const{AF\_IPX}       & --  & si  &  --  & --  & -- \\
     \hline
    \const{AF\_NETLINK}   & --  & si  &  si  & --  & -- \\
     \hline
    \const{AF\_X25}       & --  & --  &  --  & --  & si \\
     \hline
    \const{AF\_AX25}      & --  & si  &  si  & --  & si \\
     \hline
%    \const{AF\_ATMPVC}    & --  & --  &  --  & --  & -- \\
%     \hline
    \const{AF\_APPLETALK} & --  & si  &  si  & --  & -- \\
     \hline
    \const{AF\_PACKET}    & --  & si  &  si  & --  & -- \\
     \hline
    \const{AF\_KEY}       & --  & --  &  si  & --  & -- \\
     \hline
    \const{AF\_IRDA}      & si  & si  &  si  & --  & si \\
     \hline
    \const{AF\_NETROM}    & --  & --  &  --  & --  & si \\
     \hline
    \const{AF\_ROSE}      & --  & --  &  --  & --  & si \\
     \hline
    \const{AF\_RDS}       & --  & --  &  --  & --  & si \\
     \hline
    \const{AF\_ECONET}    & --  & si  &  --  & --  & -- \\
     \hline
  \end{tabular}
  \caption{Combinazioni valide di dominio e tipo di protocollo per la 
    funzione \func{socket}.}
  \label{tab:sock_sock_valid_combinations}
\end{table}

\footnotetext{supportati a partire dal kernel 2.6.4 per socket che conservano
  i limiti dei messaggi e li consegnano in sequenza ordinata.}

In tab.~\ref{tab:sock_sock_valid_combinations} sono mostrate le combinazioni
valide possibili per le principali famiglie di protocolli. Per ogni
combinazione valida si è indicato il tipo di protocollo, o la parola
\textsl{si} qualora il protocollo non abbia un nome definito, mentre si sono
lasciate vuote le caselle per le combinazioni non supportate.


\section{Le strutture degli indirizzi dei socket}
\label{sec:sock_sockaddr}

Come si è visto nella creazione di un socket non si specifica nulla oltre al
tipo di famiglia di protocolli che si vuole utilizzare, in particolare nessun
indirizzo che identifichi i due capi della comunicazione. La funzione infatti
si limita ad allocare nel kernel quanto necessario per poter poi realizzare la
comunicazione.

Gli indirizzi infatti vengono specificati attraverso apposite strutture che
vengono utilizzate dalle altre funzioni della interfaccia dei socket, quando
la comunicazione viene effettivamente realizzata.  Ogni famiglia di protocolli
ha ovviamente una sua forma di indirizzamento e in corrispondenza a questa una
sua peculiare struttura degli indirizzi. I nomi di tutte queste strutture
iniziano per \var{sockaddr\_}; quelli propri di ciascuna famiglia vengono
identificati dal suffisso finale, aggiunto al nome precedente.


\subsection{La struttura generica}
\label{sec:sock_sa_gen}

Le strutture degli indirizzi vengono sempre passate alle varie funzioni
attraverso puntatori (cioè \textit{by reference}), ma le funzioni devono poter
maneggiare puntatori a strutture relative a tutti gli indirizzi possibili
nelle varie famiglie di protocolli; questo pone il problema di come passare
questi puntatori, il C moderno risolve questo problema coi i puntatori
generici (i \ctyp{void *}), ma l'interfaccia dei socket è antecedente alla
definizione dello standard ANSI C, e per questo nel 1982 fu scelto di definire
una struttura generica per gli indirizzi dei socket, \struct{sockaddr}, che si
è riportata in fig.~\ref{fig:sock_sa_gen_struct}.

\begin{figure}[!htb]
  \footnotesize \centering
  \begin{minipage}[c]{0.80\textwidth}
    \includestruct{listati/sockaddr.h}
  \end{minipage} 
  \caption{La struttura generica degli indirizzi dei socket
    \structd{sockaddr}.} 
  \label{fig:sock_sa_gen_struct}
\end{figure}

Tutte le funzioni dei socket che usano gli indirizzi sono definite usando nel
prototipo un puntatore a questa struttura; per questo motivo quando si
invocano dette funzioni passando l'indirizzo di un protocollo specifico
occorrerà eseguire una conversione del relativo puntatore.

I tipi di dati che compongono la struttura sono stabiliti dallo standard
POSIX.1g e li abbiamo riassunti in tab.~\ref{tab:sock_data_types} con i
rispettivi file di include in cui sono definiti; la struttura è invece
definita nell'include file \headfile{sys/socket.h}.

\begin{table}[!htb]
  \centering
  \footnotesize
  \begin{tabular}{|l|l|l|}
    \hline
    \multicolumn{1}{|c|}{\textbf{Tipo}}& 
    \multicolumn{1}{|c|}{\textbf{Descrizione}}& 
    \multicolumn{1}{|c|}{\textbf{Header}} \\
    \hline
    \hline
    \typed{int8\_t}   & intero a 8 bit con segno   & \headfile{sys/types.h}\\
    \typed{uint8\_t}  & intero a 8 bit senza segno & \headfile{sys/types.h}\\
    \typed{int16\_t}  & intero a 16 bit con segno  & \headfile{sys/types.h}\\
    \typed{uint16\_t} & intero a 16 bit senza segno& \headfile{sys/types.h}\\
    \typed{int32\_t}  & intero a 32 bit con segno  & \headfile{sys/types.h}\\
    \typed{uint32\_t} & intero a 32 bit senza segno& \headfile{sys/types.h}\\
    \hline
    \typed{sa\_family\_t} & famiglia degli indirizzi&\headfile{sys/socket.h}\\
    \typed{socklen\_t} & lunghezza (\type{uint32\_t}) dell'indirizzo di
    un socket& \headfile{sys/socket.h}\\
    \hline
    \typed{in\_addr\_t} & indirizzo IPv4 (\type{uint32\_t}) & 
    \headfile{netinet/in.h}\\
    \typed{in\_port\_t} & porta TCP o UDP (\type{uint16\_t})& 
    \headfile{netinet/in.h}\\
    \hline
  \end{tabular}
  \caption{Tipi di dati usati nelle strutture degli indirizzi, secondo quanto 
    stabilito dallo standard POSIX.1g.}
  \label{tab:sock_data_types}
\end{table}

In alcuni sistemi la struttura è leggermente diversa e prevede un primo membro
aggiuntivo \code{uint8\_t sin\_len} (come riportato da R. Stevens in
\cite{UNP1}). Questo campo non verrebbe usato direttamente dal programmatore e
non è richiesto dallo standard POSIX.1g, in Linux pertanto non esiste. Il
campo \type{sa\_family\_t} era storicamente un \ctyp{unsigned short}.

Dal punto di vista del programmatore l'unico uso di questa struttura è quello
di fare da riferimento per il casting, per il kernel le cose sono un po'
diverse, in quanto esso usa il puntatore per recuperare il campo
\var{sa\_family}, comune a tutte le famiglie, con cui determinare il tipo di
indirizzo; per questo motivo, anche se l'uso di un puntatore \ctyp{void *}
sarebbe più immediato per l'utente (che non dovrebbe più eseguire il casting),
è stato mantenuto l'uso di questa struttura.

Se si usa una struttura \struct{sockaddr} per allocare delle variabili
generiche da usare in seguito per degli indirizzi si pone il problema che
niente assicura che i dati necessari per le varie famiglie di indirizzi
possano rientrare nella dimensione del campo \var{sa\_data} indicata in
fig.~\ref{fig:sock_sa_gen_struct}, anzi, come vedremo in
sez.~\ref{sec:sock_sa_ipv6}, nel caso di indirizzi IPv6 questa non è proprio
sufficiente. 

\begin{figure}[!htb]
  \footnotesize \centering
  \begin{minipage}[c]{0.95\textwidth}
    \includestruct{listati/sockaddr_storage.h}
  \end{minipage} 
  \caption{La struttura generica degli indirizzi dei socket
    \structd{sockaddr\_storage}.} 
  \label{fig:sock_sa_storage_struct}
\end{figure}

Per questo l'interfaccia di programmazione dei socket prevede la defizione di
una speciale struttura \struct{sockaddr\_storage} illustrata in
fig.~\ref{fig:sock_sa_storage_struct}, in cui di nuovo si usa il primo campo
(\var{ss\_family}) per indicare il tipo di indirizzo, ed in cui i campi
successivi sono utilizzati per allineare i dati al tipo di architettura
hardware utilizzata, e per allocare uno spazio sufficiente ampio per contenere
qualunque tipo di indirizzo supportato. Allocando questa struttura si ha la
certezza di non eccedere le dimensioni qualunque sia il tipo di indirizzi che
si useranno, pertanto risulta utile tutte le volte che si devono gestire in
maniera generica tipi di indirizzi diversi (ad esempio IPv4 ed IPv6).


\subsection{La struttura degli indirizzi IPv4}
\label{sec:sock_sa_ipv4}

I socket di tipo \const{AF\_INET} vengono usati per la comunicazione
attraverso Internet; la struttura per gli indirizzi per un socket Internet (se
si usa IPv4) è definita come \struct{sockaddr\_in} nell'header file
\headfiled{netinet/in.h} ed ha la forma mostrata in
fig.~\ref{fig:sock_sa_ipv4_struct}, conforme allo standard POSIX.1g.

\begin{figure}[!htb]
  \footnotesize\centering
  \begin{minipage}[c]{0.80\textwidth}
    \includestruct{listati/sockaddr_in.h}
  \end{minipage} 
  \caption{La struttura \structd{sockaddr\_in} degli indirizzi dei socket
    Internet (IPv4) e la struttura \structd{in\_addr} degli indirizzi IPv4.}
  \label{fig:sock_sa_ipv4_struct}
\end{figure}

L'indirizzo di un socket Internet (secondo IPv4) comprende l'indirizzo
Internet di un'interfaccia più un \textsl{numero di porta} (affronteremo in
dettaglio il significato di questi numeri in sez.~\ref{sec:TCP_port_num}).  Il
protocollo IP di per sé non prevede numeri di porta, questi sono utilizzati
solo dai protocolli di livello superiore come TCP e UDP, ma devono essere
indicati qui. Inoltre questa struttura viene usata anche per i socket RAW che
accedono direttamente al livello di IP, in questo caso il numero della porta
deve essere impostato al numero di protocollo.

Il membro \var{sin\_family} deve essere sempre impostato a \constd{AF\_INET},
altrimenti si avrà un errore di \errcode{EINVAL}; il membro \var{sin\_port}
specifica il \textsl{numero di porta}. I numeri di porta sotto il 1024 sono
chiamati \textsl{riservati} in quanto utilizzati da servizi standard e
soltanto processi con i privilegi di amministratore (con \ids{UID} effettivo
uguale a zero) o con la \textit{capability} \const{CAP\_NET\_BIND\_SERVICE}
possono usare la funzione \func{bind} (che vedremo in
sez.~\ref{sec:TCP_func_bind}) su queste porte.

Il membro \var{sin\_addr} contiene un indirizzo Internet, e viene acceduto sia
come struttura (un resto di una implementazione precedente in cui questa era
una \dirct{union} usata per accedere alle diverse classi di indirizzi) che
direttamente come intero. In \headfile{netinet/in.h} vengono definite anche
alcune costanti che identificano alcuni indirizzi speciali, riportati in
tab.~\ref{tab:TCP_ipv4_addr}, che rincontreremo più avanti.

Infine occorre sottolineare che sia gli indirizzi che i numeri di porta devono
essere specificati in quello che viene chiamato \textit{network order}, cioè
con i bit ordinati in formato \textit{big endian} (vedi
sez.~\ref{sec:endianness}), questo comporta la necessità di usare apposite
funzioni di conversione per mantenere la portabilità del codice (vedi
sez.~\ref{sec:sock_addr_func} per i dettagli del problema e le relative
soluzioni).


\subsection{La struttura degli indirizzi IPv6}
\label{sec:sock_sa_ipv6}

Essendo IPv6 un'estensione di IPv4, i socket di tipo \const{AF\_INET6} sono
sostanzialmente identici ai precedenti; la parte in cui si trovano
praticamente tutte le differenze fra i due socket è quella della struttura
degli indirizzi; la sua definizione, presa da \headfile{netinet/in.h}, è
riportata in fig.~\ref{fig:sock_sa_ipv6_struct}.

\begin{figure}[!htb]
  \footnotesize \centering
  \begin{minipage}[c]{0.80\textwidth}
    \includestruct{listati/sockaddr_in6.h}
  \end{minipage} 
  \caption{La struttura \structd{sockaddr\_in6} degli indirizzi dei socket
    IPv6 e la struttura \structd{in6\_addr} degli indirizzi IPv6.}
  \label{fig:sock_sa_ipv6_struct}
\end{figure}

Il campo \var{sin6\_family} deve essere sempre impostato ad \constd{AF\_INET6},
il campo \var{sin6\_port} è analogo a quello di IPv4 e segue le stesse regole;
il campo \var{sin6\_flowinfo} è a sua volta diviso in tre parti di cui i 24
bit inferiori indicano l'etichetta di flusso, i successivi 4 bit la priorità e
gli ultimi 4 sono riservati. Questi valori fanno riferimento ad alcuni campi
specifici dell'header dei pacchetti IPv6 (vedi sez.~\ref{sec:IP_ipv6head}) ed
il loro uso è sperimentale.

Il campo \var{sin6\_addr} contiene l'indirizzo a 128 bit usato da IPv6,
espresso da un vettore di 16 byte; anche in questo caso esistono alcuni valori
predediniti, ma essendo il campo un vettore di byte non è possibile assegnarli
con il calore di una costante. Esistono però le variabili predefinite
\var{in6addr\_any} (che indica l'indirizzo generico) e \var{in6addr\_loopback}
(che indica l'indirizzo di loopback) il cui valore può essere copiato in
questo campo. A queste due variabili si aggiungono le macro
\macrod{IN6ADDR\_ANY\_INIT} e \macrod{IN6ADDR\_LOOPBACK\_INIT} per effettuare
delle assegnazioni statiche.

Infine il campo \var{sin6\_scope\_id} è un campo introdotto in Linux con il
kernel 2.4, per gestire alcune operazioni riguardanti il
\textit{multicasting}, è supportato solo per gli indirizzi di tipo
\textit{link-local} (vedi sez.~\ref{sec:IP_ipv6_unicast}) e deve contenere
l'\textit{interface index} (vedi sez.~\ref{sec:sock_ioctl_netdevice}) della
scheda di rete.  Si noti infine che \struct{sockaddr\_in6} ha una dimensione
maggiore della struttura \struct{sockaddr} generica di
fig.~\ref{fig:sock_sa_gen_struct}, quindi occorre stare attenti a non avere
fatto assunzioni riguardo alla possibilità di contenere i dati nelle
dimensioni di quest'ultima (per questo se necessario è opportuno usare
\struct{sockaddr\_storage}).


\subsection{La struttura degli indirizzi locali}
\label{sec:sock_sa_local}

I socket di tipo \const{AF\_UNIX} o \const{AF\_LOCAL} vengono usati per una
comunicazione fra processi che stanno sulla stessa macchina (per questo
vengono chiamati \textit{local domain} o anche \textit{Unix domain}); essi
hanno la caratteristica ulteriore di poter essere creati anche in maniera
anonima attraverso la funzione \func{socketpair} (che abbiamo trattato in
sez.~\ref{sec:ipc_socketpair}).  Quando però si vuole fare riferimento
esplicito ad uno di questi socket si deve usare una struttura degli indirizzi
di tipo \struct{sockaddr\_un}, la cui definizione si è riportata in
fig.~\ref{fig:sock_sa_local_struct}.

\begin{figure}[!htb]
  \footnotesize \centering
  \begin{minipage}[c]{0.80\textwidth}
    \includestruct{listati/sockaddr_un.h}
  \end{minipage} 
  \caption{La struttura \structd{sockaddr\_un} degli indirizzi dei socket
    locali (detti anche \textit{unix domain}) definita in
    \headfiled{sys/un.h}.}
  \label{fig:sock_sa_local_struct}
\end{figure}

In questo caso il campo \var{sun\_family} deve essere \constd{AF\_UNIX},
mentre il campo \var{sun\_path} deve specificare un indirizzo. Questo ha due
forme; può essere ``\textit{named}'' ed in tal caso deve corrispondere ad un
file (di tipo socket) presente nel filesystem o essere ``\textit{abstract}''
nel qual caso viene identificato da una stringa univoca in uno spazio di nomi
astratto. 

Nel primo caso l'indirizzo viene specificato in \var{sun\_path} come una
stringa (terminata da uno zero) corrispondente al \textit{pathname} del file;
nel secondo caso (che è specifico di Linux e non portabile) \var{sun\_path}
deve iniziare con uno zero ed il nome verrà costituito dai restanti byte che
verranno interpretati come stringa senza terminazione (un byte nullo non ha in
questo caso nessun significato).

In realtà esiste una terza forma, \textit{unnamed}, che non è possibile
indicare in fase di scrittura, ma che è quella che viene usata quando si legge
l'indirizzo di un socket anonimo creato con \texttt{socketpair}; in tal caso
la struttura restituita è di dimensione \code{sizeof(sa\_family\_t)}, quindi
\var{sun\_path} non esiste e non deve essere referenziato.

\subsection{La struttura degli indirizzi AppleTalk}
\label{sec:sock_sa_appletalk}

I socket di tipo \const{AF\_APPLETALK} sono usati dalla libreria
\file{netatalk} per implementare la comunicazione secondo il protocollo
AppleTalk, uno dei primi protocolli di rete usato nel mondo dei personal
computer, usato dalla Apple per connettere fra loro computer e stampanti. Il
kernel supporta solo due strati del protocollo, DDP e AARP, e di norma è
opportuno usare le funzioni della libreria \texttt{netatalk}, tratteremo qui
questo argomento principalmente per mostrare l'uso di un protocollo
alternativo.

I socket AppleTalk permettono di usare il protocollo DDP, che è un protocollo
a pacchetto, di tipo \const{SOCK\_DGRAM}; l'argomento \param{protocol} di
\func{socket} deve essere nullo. È altresì possibile usare i socket raw
specificando un tipo \const{SOCK\_RAW}, nel qual caso l'unico valore valido
per \param{protocol} è \constd{ATPROTO\_DDP}.

Gli indirizzi AppleTalk devono essere specificati tramite una struttura
\struct{sockaddr\_atalk}, la cui definizione è riportata in
fig.~\ref{fig:sock_sa_atalk_struct}; la struttura viene dichiarata includendo
il file \headfiled{netatalk/at.h}.

\begin{figure}[!htb]
  \footnotesize \centering
  \begin{minipage}[c]{0.80\textwidth}
    \includestruct{listati/sockaddr_atalk.h}
  \end{minipage} 
  \caption{La struttura \structd{sockaddr\_atalk} degli indirizzi dei socket
    AppleTalk, e la struttura \structd{at\_addr} degli indirizzi AppleTalk.}
  \label{fig:sock_sa_atalk_struct}
\end{figure}

Il campo \var{sat\_family} deve essere sempre \constd{AF\_APPLETALK}, mentre
il campo \var{sat\_port} specifica la porta che identifica i vari
servizi. Valori inferiori a 129 sono usati per le \textsl{porte riservate}, e
possono essere usati solo da processi con i privilegi di amministratore o con
la \textit{capability} \const{CAP\_NET\_BIND\_SERVICE}.  

L'indirizzo remoto è specificato nella struttura \var{sat\_addr}, e deve
essere in \textit{network order} (vedi sez.~\ref{sec:endianness}); esso è
composto da un parte di rete data dal campo \var{s\_net}, che può assumere il
valore \constd{AT\_ANYNET}, che indica una rete generica e vale anche per
indicare la rete su cui si è, il singolo nodo è indicato da \var{s\_node}, e
può prendere il valore generico \constd{AT\_ANYNODE} che indica anche il nodo
corrente, ed il valore \constd{ATADDR\_BCAST} che indica tutti i nodi della
rete.


\subsection{La struttura degli indirizzi dei \textit{packet socket}}
\label{sec:sock_sa_packet}

I \textit{packet socket}, identificati dal dominio \const{AF\_PACKET}, sono
un'interfaccia specifica di Linux per inviare e ricevere pacchetti
direttamente su un'interfaccia di rete, senza passare per le funzioni di
gestione dei protocolli di livello superiore. In questo modo è possibile
implementare dei protocolli in \textit{user space}, agendo direttamente sul
livello fisico. In genere comunque si preferisce usare la libreria
\file{pcap},\footnote{la libreria è mantenuta insieme al comando
  \cmd{tcpdump}, informazioni e documentazione si possono trovare sul sito del
  progetto \url{http://www.tcpdump.org/}.}  che assicura la portabilità su
altre piattaforme, anche se con funzionalità ridotte.

Questi socket possono essere di tipo \const{SOCK\_RAW} o \const{SOCK\_DGRAM}.
Con socket di tipo \const{SOCK\_RAW} si può operare sul livello di
collegamento, ed i pacchetti vengono passati direttamente dal socket al driver
del dispositivo e viceversa.  In questo modo, in fase di trasmissione, il
contenuto completo dei pacchetti, comprese le varie intestazioni, deve essere
fornito dall'utente. In fase di ricezione invece tutto il contenuto del
pacchetto viene passato inalterato sul socket, anche se il kernel analizza
comunque il pacchetto, riempiendo gli opportuni campi della struttura
\struct{sockaddr\_ll} ad esso associata.

Si usano invece socket di tipo \const{SOCK\_DGRAM} quando si vuole operare a
livello di rete. 

In questo caso in fase di ricezione l'intestazione del protocollo di
collegamento viene rimossa prima di passare il resto del pacchetto all'utente,
mentre in fase di trasmissione viene creata una opportuna intestazione per il
protocollo a livello di collegamento utilizzato, usando le informazioni
necessarie che devono essere specificate sempre con una struttura
\struct{sockaddr\_ll}.

Nella creazione di un \textit{packet socket} il valore dell'argomento
\param{protocol} di \func{socket} serve a specificare, in \textit{network
  order}, il numero identificativo del protocollo di collegamento si vuole
utilizzare. I valori possibili sono definiti secondo lo standard IEEE 802.3, e
quelli disponibili in Linux sono accessibili attraverso opportune costanti
simboliche definite nel file \file{linux/if\_ether.h}. Se si usa il valore
speciale \constd{ETH\_P\_ALL} passeranno sul \textit{packet socket} tutti i
pacchetti, qualunque sia il loro protocollo di collegamento. Ovviamente l'uso
di questi socket è una operazione privilegiata e può essere effettuati solo da
un processo con i privilegi di amministratore (\ids{UID} effettivo nullo) o
con la \textit{capability} \const{CAP\_NET\_RAW}.

Una volta aperto un \textit{packet socket}, tutti i pacchetti del protocollo
specificato passeranno attraverso di esso, qualunque sia l'interfaccia da cui
provengono; se si vuole limitare il passaggio ad una interfaccia specifica
occorre usare la funzione \func{bind} (vedi sez.~\ref{sec:TCP_func_bind}) per
agganciare il socket a quest'ultima.

\begin{figure}[!htb]
  \footnotesize \centering
  \begin{minipage}[c]{\textwidth}
    \includestruct{listati/sockaddr_ll.h}
  \end{minipage} 
  \caption{La struttura \structd{sockaddr\_ll} degli indirizzi dei
    \textit{packet socket}.}
  \label{fig:sock_sa_packet_struct}
\end{figure}

Nel caso dei \textit{packet socket} la struttura degli indirizzi è di tipo
\struct{sockaddr\_ll}, e la sua definizione è riportata in
fig.~\ref{fig:sock_sa_packet_struct}; essa però viene ad assumere un ruolo
leggermente diverso rispetto a quanto visto finora per gli altri tipi di
socket.  Infatti se il socket è di tipo \const{SOCK\_RAW} si deve comunque
scrivere tutto direttamente nel pacchetto, quindi la struttura non serve più a
specificare gli indirizzi. Essa mantiene questo ruolo solo per i socket di
tipo \const{SOCK\_DGRAM}, per i quali permette di specificare i dati necessari
al protocollo di collegamento, mentre viene sempre utilizzata in lettura (per
entrambi i tipi di socket), per la ricezione dei dati relativi a ciascun
pacchetto.

Al solito il campo \var{sll\_family} deve essere sempre impostato al valore
\constd{AF\_PACKET}. Il campo \var{sll\_protocol} indica il protocollo scelto,
e deve essere indicato in \textit{network order}, facendo uso delle costanti
simboliche definite in \file{linux/if\_ether.h}. Il campo \var{sll\_ifindex} è
l'indice dell'interfaccia (l'\textit{inxterface index} (vedi
sez.~\ref{sec:sock_ioctl_netdevice}) che in caso di presenza di più
interfacce dello stesso tipo (se ad esempio si hanno più schede Ethernet),
permette di selezionare quella con cui si vuole operare (un valore nullo
indica qualunque interfaccia).  Questi sono i due soli campi che devono essere
specificati quando si vuole selezionare una interfaccia specifica, usando
questa struttura con la funzione \func{bind}.

I campi \var{sll\_halen} e \var{sll\_addr} indicano rispettivamente
l'indirizzo associato all'interfaccia sul protocollo di collegamento e la
relativa lunghezza; ovviamente questi valori cambiano a seconda del tipo di
collegamento che si usa, ad esempio, nel caso di Ethernet, questi saranno il
MAC address della scheda e la relativa lunghezza. Essi vengono usati, insieme
ai campi \var{sll\_family} e \var{sll\_ifindex} quando si inviano dei
pacchetti, in questo caso tutti gli altri campi devono essere nulli.

Il campo \var{sll\_hatype} indica il tipo ARP, come definito in
\file{linux/if\_arp.h}, mentre il campo \var{sll\_pkttype} indica il tipo di
pacchetto; entrambi vengono impostati alla ricezione di un pacchetto ed han
senso solo in questo caso. In particolare \var{sll\_pkttype} può assumere i
seguenti valori: \constd{PACKET\_HOST} per un pacchetto indirizzato alla
macchina ricevente, \constd{PACKET\_BROADCAST} per un pacchetto di
\textit{broadcast}, \constd{PACKET\_MULTICAST} per un pacchetto inviato ad un
indirizzo fisico di \textit{multicast}, \constd{PACKET\_OTHERHOST} per un
pacchetto inviato ad un'altra stazione (e ricevuto su un'interfaccia in modo
promiscuo), \constd{PACKET\_OUTGOING} per un pacchetto originato dalla propria
macchina che torna indietro sul socket.

Si tenga presente infine che in fase di ricezione, anche se si richiede il
troncamento del pacchetto, le funzioni \func{recv}, \func{recvfrom} e
\func{recvmsg} (vedi sez.~\ref{sec:net_sendmsg}) restituiranno comunque la
lunghezza effettiva del pacchetto così come arrivato sulla linea.

%% \subsection{La struttura degli indirizzi DECnet}
%% \label{sec:sock_sa_decnet}
 
%% I socket di tipo \const{AF\_DECnet} usano il protocollo DECnet, usato dai VAX
%% Digital sotto VMS quando ancora il TCP/IP non era diventato lo standard di
%% fatto. Il protocollo è un protocollo chiuso, ed il suo uso attuale è limitato
%% alla comunicazione con macchine che stanno comunque scomparendo. Lo si riporta
%% solo come esempio 


% TODO aggiungere AF_CAN, vedi http://lwn.net/Articles/253425, dal 2.6.25 ?

% TODO: trattare i socket RDS, vedi documentazione del kernel, file 
% Documentation/networking/rds.txt



\section{Le funzioni di conversione degli indirizzi}
\label{sec:sock_addr_func}

In questa sezione tratteremo delle varie funzioni usate per manipolare gli
indirizzi, limitandoci però agli indirizzi Internet.  Come accennato gli
indirizzi e i numeri di porta usati nella rete devono essere forniti nel
cosiddetto \textit{network order}, che corrisponde al formato \textit{big
  endian} (vedi sez.~\ref{sec:endianness}), anche quando la proprio macchina
non usa questo formato, cosa che può comportare la necessità di eseguire delle
conversioni.


\subsection{Le funzioni per il riordinamento}
\label{sec:sock_func_ord}

Come già visto in sez.~\ref{sec:endianness} il problema connesso
all'\textit{endianness} è che quando si passano dei dati da un tipo di
architettura all'altra i dati vengono interpretati in maniera diversa, e ad
esempio nel caso dell'intero a 16 bit ci si ritroverà con i due byte in cui è
suddiviso scambiati di posto.  

Per questo motivo si usano delle funzioni di conversione che servono a tener
conto automaticamente della possibile differenza fra l'ordinamento usato sul
computer e quello che viene usato nelle trasmissione sulla rete; queste
funzioni sono \funcd{htonl}, \funcd{htons}, \funcd{ntohl} e \funcd{ntohs} ed i
rispettivi prototipi sono:

\begin{funcproto}{
\fhead{arpa/inet.h}
\fdecl{unsigned long int htonl(unsigned long int hostlong)}
\fdesc{Converte l'intero a 32 bit \param{hostlong} dal formato della macchina a
  quello della rete.} 
\fdecl{unsigned short int htons(unsigned short int hostshort)}
\fdesc{Converte l'intero a 16 bit \param{hostshort} dal formato della macchina a
  quello della rete.}
\fdecl{unsigned long int ntohl(unsigned long int netlong)}
\fdesc{Converte l'intero a 32 bit \param{netlong} dal formato della rete a
  quello della macchina.}
\fdecl{unsigned sort int ntohs(unsigned short int netshort)}
\fdesc{Converte l'intero a 16 bit \param{netshort} dal formato della rete a
  quello della macchina.}
}

{Tutte le funzioni restituiscono il valore convertito, e non prevedono
  errori.}
\end{funcproto}

I nomi sono assegnati usando la lettera \texttt{n} come mnemonico per indicare
l'ordinamento usato sulla rete (da \textit{network order}) e la lettera
\texttt{h} come mnemonico per l'ordinamento usato sulla macchina locale (da
\textit{host order}), mentre le lettere \texttt{s} e \texttt{l} stanno ad
indicare i tipi di dato (\ctyp{long} o \ctyp{short}, riportati anche dai
prototipi).

Usando queste funzioni si ha la conversione automatica: nel caso in cui la
macchina che si sta usando abbia una architettura \textit{big endian} queste
funzioni sono definite come macro che non fanno nulla. Per questo motivo vanno
sempre utilizzate, anche quando potrebbero non essere necessarie, in modo da
assicurare la portabilità del codice su tutte le architetture.


\subsection{Le funzioni di conversione per gli indirizzi IPv4}
\label{sec:sock_func_ipv4}

Un secondo insieme di funzioni di manipolazione è quello che serve per passare
dalla rappresentazione simbolica degli indirizzi IP al formato binario
previsto dalla struttura degli indirizzi di
fig.~\ref{fig:sock_sa_ipv4_struct}, e viceversa. La notazione più comune è la
cosiddetta notazione \itindex{dotted-decimal} \textit{dotted-decimal}, che
prevede che gli indirizzi IPv4 siano indicati con l'espressione del valore
numerico decimale di ciascuno dei 4 byte che li costituiscono separati da un
punto (ad esempio \texttt{192.168.0.1}).

In realtà le funzioni che illustreremo supportano una notazione che più
propriamente dovrebbe esser chiamata \textit{numbers-and-dot} in quanto il
valore può essere indicato con numeri espressi sia in decimale, che in ottale
(se indicati apponendo uno zero) che in esadecimale (se indicati apponendo
\texttt{0x}). Inoltre per la parte meno significativa dell'espressione, quella
che riguarda l'indirizzo locale, si può usare, eliminando altrettanti punti,
valori a 16 o a 24 bit, e togliendo tutti i punti, si può usare anche
direttamente un valore numerico a 32 bit.\footnote{la funzionalità si trova
  anche in gran parte dei programmi che usano indirizzi di rete, e deriva
  direttamente da queste funzioni.}

Tradizionalmente la conversione di un indirizzo \textit{dotted-decimal} al
valore numerico veniva eseguita dalla funzione \funcd{inet\_addr} (prevista
fin dalle origini in BSD e inclusa in POSIX.1-2001) il cui prototipo è:

\begin{funcproto}{
\fhead{arpa/inet.h}
\fdecl{in\_addr\_t inet\_addr(const char *strptr)}
\fdesc{Converte la stringa dell'indirizzo \textit{dotted decimal} in nel
  numero IP in network order.} 
}

{La funzione ritorna il valore dell'indirizzo in caso di successo e
  \const{INADDR\_NONE} per un errore e non genera codici di errore.}
\end{funcproto}

La prima funzione, \func{inet\_addr}, restituisce l'indirizzo a 32 bit in
\textit{network order} (del tipo \type{in\_addr\_t}) a partire dalla stringa
passata nell'argomento \param{strptr}. In caso di errore (quando la stringa
non esprime un indirizzo valido) restituisce invece il valore
\const{INADDR\_NONE}, che tipicamente sono trentadue bit a uno.  Questo però
comporta che la stringa \texttt{255.255.255.255}, che pure è un indirizzo
valido, non può essere usata con questa funzione dato che genererebe comunque
un errore; per questo motivo essa è generalmente deprecata in favore di
\func{inet\_aton}.

Per effettuare la conversione inversa la funzione usata tradizionalmente è
\funcd{inet\_ntoa}, anch'essa presente fin da BSD 4.3, in cui si riprende la
notazione già vista in sez.~\ref{sec:sock_func_ord} che usa la lettera
\texttt{n} come mnemonico per indicare la rete ed \texttt{a} (per ASCII) come
mnemonico per indicare la stringa corrispodente all'indirizzo; il suo
prototipo è:

\begin{funcproto}{
\fhead{arpa/inet.h}
\fdecl{char *inet\_ntoa(struct in\_addr addrptr)}
\fdesc{Converte un indirizzo IP in una stringa \textit{dotted decimal}.} 
}

{La funzione l'indirizzo della stringa con il valore dell'indirizzo convertito
  e non prevede errori.}
\end{funcproto}

La funzione converte il valore a 32 bit dell'indirizzo, espresso in
\textit{network order}, e preso direttamente con un puntatore al relativo
campo della struttura degli indirizzi, restituendo il puntatore alla stringa
che contiene l'espressione in formato \textit{dotted-decimal}. Si deve tenere
presente che la stringa risiede in un segmento di memoria statica, per cui
viene riscritta ad ogni chiamata e la funzione non è rientrante.

Per rimediare ai problemi di \funcd{inet\_addr} è stata sostituita da
\funcd{inet\_aton}, che però non è stata standardizzata e non è presente in
POSIX.1-2001, anche se è definita sulla gran parte dei sistemi Unix; il suo
prototipo è:

\begin{funcproto}{
\fhead{arpa/inet.h}
\fdecl{int inet\_aton(const char *src, struct in\_addr *dest)}
\fdesc{Converte la stringa dell'indirizzo \textit{dotted decimal} in un
  indirizzo IP.}
}

{La funzione ritorna un valore non nullo in caso di successo e $0$ per un
  errore e non genera codici di errore.}
\end{funcproto}

La funzione converte la stringa puntata da \param{src} nell'indirizzo binario
che viene memorizzato nell'opportuna struttura \struct{in\_addr} (si veda
fig.~\ref{fig:sock_sa_ipv4_struct}) situata all'indirizzo dato
dall'argomento \param{dest} (è espressa in questa forma in modo da poterla
usare direttamente con il puntatore usato per passare la struttura degli
indirizzi). La funzione restituisce un valore diverso da zero se l'indirizzo è
valido e la conversione ha successo e 0 in caso contrario. Se usata
con \param{dest} inizializzato a \val{NULL} può essere usata per effettuare la
validazione dell'indirizzo espresso da \param{src}.

Oltre a queste tre funzioni esistono le ulteriori \funcm{inet\_lnaof},
\funcm{inet\_netof} e \funcm{inet\_makeaddr} che assumono la ormai obsoleta e
deprecata suddivisione in classi degli indirizzi IP per fornire la parte di
rete e quella di indirizzo locale. Ad oggi il loro uso non ha più alcun senso
per ciò non le tratteremo.


\subsection{Le funzioni di conversione per indirizzi IP generici}
\label{sec:sock_conv_func_gen}

Le tre funzioni precedenti sono limitate solo ad indirizzi IPv4, per questo
motivo è preferibile usare le due nuove funzioni \func{inet\_pton} e
\func{inet\_ntop} che possono convertire anche gli indirizzi IPv6. Anche in
questo caso le lettere \texttt{n} e \texttt{p} sono degli mnemonici per
ricordare il tipo di conversione effettuata e stanno per \textit{presentation}
e \textit{numeric}.

Entrambe le funzioni accettano l'argomento \param{af} che indica il tipo di
indirizzo, e che può essere soltanto \const{AF\_INET} o \const{AF\_INET6}. La
prima funzione, \funcd{inet\_pton}, serve a convertire una stringa in un
indirizzo; il suo prototipo è:

\begin{funcproto}{
\fhead{sys/socket.h}
\fdecl{int inet\_pton(int af, const char *src, void *addr\_ptr)} 
\fdesc{Converte l'indirizzo espresso tramite una stringa nel valore numerico.} 
}

{La funzione ritorna $1$ in caso di successo, $0$ se \param{src} non contiene
  una rappresentazione valida per la famiglia di indirizzi indicati
  da \param{af} e $-1$ se \param{af} specifica una famiglia di indirizzi non
  valida, e solo in quest'ultimo caso  \var{errno} assumerà il valore
  \errcode{EAFNOSUPPORT}.
}
\end{funcproto}

La funzione converte la stringa indicata tramite \param{src} nel valore
numerico dell'indirizzo IP del tipo specificato da \param{af} che viene
memorizzato all'indirizzo puntato da \param{addr\_ptr}. La funzione supporta
per IPv4 la sola notazione \textit{dotted-decimal}, e non quella più completa
\textit{number-and-dot} che abbiamo visto per \func{inet\_aton}. Per IPv6 la
notazione prevede la suddivisione dei 128 bit dell'indirizzo in 16 parti di 16
bit espresse con valori esadecimali separati dal carattere ``\texttt{:}'' ed
una serie di valori nulli possono essere sostituiti (una sola volta, sempre a
partire dalla sinistra) con la notazione ``\texttt{::}'', un esempio di
indirizzo in questa forma potrebbe essere \texttt{2001:db8::8:ba98:2078:e3e3},
per una descrizione più completa si veda sez.~\ref{sec:IP_ipv6_notation}.

La seconda funzione di conversione è \funcd{inet\_ntop} che converte un
indirizzo in una stringa; il suo prototipo è:

\begin{funcproto}{
\fhead{sys/socket.h}
\fdecl{char *inet\_ntop(int af, const void *addr\_ptr, char *dest, size\_t len)}
\fdesc{Converte l'indirizzo dalla relativa struttura in una stringa simbolica.} 
}

{La funzione ritorna un puntatore non nullo alla stringa convertita in caso di
  successo e \val{NULL} per un errore, nel qual caso \var{errno} assumerà uno
  dei valori:
  \begin{errlist}
    \item[\errcode{ENOSPC}] le dimensioni della stringa con la conversione
      dell'indirizzo eccedono la lunghezza specificata da \param{len}.
    \item[\errcode{ENOAFSUPPORT}] la famiglia di indirizzi \param{af} non è
      una valida.
  \end{errlist}
}
\end{funcproto}


La funzione converte la struttura dell'indirizzo puntata da \param{addr\_ptr}
in una stringa che viene copiata nel buffer puntato dall'indirizzo
\param{dest}; questo deve essere preallocato dall'utente e la lunghezza deve
essere almeno \constd{INET\_ADDRSTRLEN} in caso di indirizzi IPv4 e
\constd{INET6\_ADDRSTRLEN} per indirizzi IPv6; la lunghezza del buffer deve
comunque essere specificata con il parametro \param{len}.

Gli indirizzi vengono convertiti da/alle rispettive strutture di indirizzo
(una struttura \struct{in\_addr} per IPv4, e una struttura \struct{in6\_addr}
per IPv6), che devono essere precedentemente allocate e passate attraverso il
puntatore \param{addr\_ptr}; l'argomento \param{dest} di \func{inet\_ntop} non
può essere nullo e deve essere allocato precedentemente.





% LocalWords:  socket sez cap BSD SVr XTI Transport Interface TCP stream UDP PF
% LocalWords:  datagram broadcast descriptor sys int domain type protocol errno
% LocalWords:  EPROTONOSUPPORT ENFILE kernel EMFILE EACCES EINVAL ENOBUFS raw
% LocalWords:  ENOMEM table family AF address name glibc UNSPEC LOCAL Local IPv
% LocalWords:  communication INET protocols ip AX Amateur IPX Novell APPLETALK
% LocalWords:  Appletalk ddp NETROM NetROM Multiprotocol ATMPVC Access to ATM
% LocalWords:  PVCs ITU ipv PLP DECnet Reserved for project NETBEUI LLC KEY key
% LocalWords:  SECURITY Security callback NETLINK interface device netlink Low
% LocalWords:  PACKET level packet ASH Ash ECONET Acorn Econet ATMSVC SVCs SNA
% LocalWords:  IRDA PPPOX PPPoX WANPIPE Wanpipe BLUETOOTH Bluetooth POSIX bits
% LocalWords:  dall'header tab SOCK capabilities capability styles DGRAM read
% LocalWords:  SEQPACKET RDM sockaddr reference void fig Header uint socklen at
% LocalWords:  addr netinet port len Stevens unsigned short casting nell'header
% LocalWords:  BIND SERVICE bind union order big endian flowinfo dell'header ll
% LocalWords:  multicast multicasting local socketpair sun path filesystem AARP
% LocalWords:  pathname AppleTalk netatalk personal Apple ATPROTO atalk sat if
% LocalWords:  ANYNET node ANYNODE ATADDR BCAST pcap IEEE linux ether ETH ALL
% LocalWords:  sll ifindex ethernet halen MAC hatype ARP arp pkttype HOST recv
% LocalWords:  OTHERHOST OUTGOING recvfrom recvmsg endianness little endtest Mac
% LocalWords:  Intel Digital Motorola IBM VME PowerPC l'Intel xABCD ptr htonl
% LocalWords:  htons ntohl ntohs long hostlong hostshort netlong
% LocalWords:  sort netshort host inet aton ntoa dotted decimal const char src
% LocalWords:  strptr struct dest addrptr INADDR NULL pton ntop presentation af
% LocalWords:  numeric EAFNOSUPPORT size ENOSPC ENOAFSUPPORT ADDRSTRLEN ROUTE
% LocalWords:  of tcpdump page


%%% Local Variables: 
%%% mode: latex
%%% TeX-master: "gapil"
%%% End: 

%% tcpsock.tex
%%
%% Copyright (C) 2000-2018 Simone Piccardi.  Permission is granted to
%% copy, distribute and/or modify this document under the terms of the GNU Free
%% Documentation License, Version 1.1 or any later version published by the
%% Free Software Foundation; with the Invariant Sections being "Un preambolo",
%% with no Front-Cover Texts, and with no Back-Cover Texts.  A copy of the
%% license is included in the section entitled "GNU Free Documentation
%% License".
%%

\chapter{I socket TCP}
\label{cha:TCP_socket}

In questo capitolo tratteremo le basi dei socket TCP, iniziando con una
descrizione delle principali caratteristiche del funzionamento di una
connessione TCP; vedremo poi le varie funzioni che servono alla creazione di
una connessione fra client e server, fornendo alcuni esempi elementari, e
finiremo prendendo in esame l'uso dell'\textit{I/O multiplexing}.


\section{Il funzionamento di una connessione TCP}
\label{sec:TCP_connession}

Prima di entrare nei dettagli delle singole funzioni usate nelle applicazioni
che utilizzano i socket TCP, è fondamentale spiegare alcune delle basi del
funzionamento del protocollo, poiché questa conoscenza è essenziale per
comprendere il comportamento di dette funzioni per questo tipo di socket, ed
il relativo modello di programmazione.

Si ricordi che il protocollo TCP serve a creare degli \textit{stream socket},
cioè una forma di canale di comunicazione che stabilisce una connessione
stabile fra due macchine in rete, in modo che queste possano scambiarsi dei
dati. In questa sezione ci concentreremo sulle modalità con le quali il
protocollo dà inizio e conclude una connessione e faremo inoltre un breve
accenno al significato di alcuni dei vari \textsl{stati} ad essa associati.


\subsection{La creazione della connessione: il \textit{three way handshake}}
\label{sec:TCP_conn_cre}

\itindbeg{three~way~handshake} 

Il processo che porta a creare una connessione TCP viene chiamato
\textit{three way handshake}; la successione tipica degli eventi (e dei
\textsl{segmenti}\footnote{si ricordi che il \textsl{segmento} è l'unità
  elementare di dati trasmessa dal protocollo TCP al livello successivo; tutti
  i segmenti hanno un'intestazione che contiene le informazioni che servono
  allo \textit{stack TCP} (così viene di solito chiamata la parte del kernel
  che realizza il protocollo) per effettuare la comunicazione, fra questi dati
  ci sono una serie di flag usati per gestire la connessione, come SYN, ACK,
  URG, FIN, alcuni di essi, come SYN (che sta per \textit{syncronize})
  corrispondono a funzioni particolari del protocollo e danno il nome al
  segmento, (per maggiori dettagli si veda sez.~\ref{sec:tcp_protocol}).}  di
dati che vengono scambiati) che porta alla creazione di una connessione è la
seguente:
 
\begin{enumerate*}
\item Il server deve essere preparato per accettare le connessioni in arrivo;
  il procedimento si chiama \textsl{apertura passiva} del socket (in inglese
  \textit{passive open}). Questo viene fatto chiamando la sequenza di funzioni
  \func{socket}, \func{bind} e \func{listen}. Completata l'apertura passiva il
  server chiama la funzione \func{accept} e il processo si blocca in attesa di
  connessioni.
  
\item Il client richiede l'inizio della connessione usando la funzione
  \func{connect}, attraverso un procedimento che viene chiamato
  \textsl{apertura attiva}, dall'inglese \textit{active open}. La chiamata di
  \func{connect} blocca il processo e causa l'invio da parte del client di un
  segmento SYN, in sostanza viene inviato al server un pacchetto IP che
  contiene solo le intestazioni di IP e TCP (con il numero di sequenza
  iniziale e il flag SYN) e le opzioni di TCP.
  
\item Il server deve dare ricevuto (l'\textit{acknowledge}) del SYN del
  client, inoltre anche il server deve inviare il suo SYN al client (e
  trasmettere il suo numero di sequenza iniziale) questo viene fatto
  ritrasmettendo un singolo segmento in cui sono impostati entrambi i flag SYN
  e ACK.
  
\item Una volta che il client ha ricevuto l'\textit{acknowledge} dal server la
  funzione \func{connect} ritorna, l'ultimo passo è dare il ricevuto del SYN
  del server inviando un ACK. Alla ricezione di quest'ultimo la funzione
  \func{accept} del server ritorna e la connessione è stabilita.
\end{enumerate*}

Il procedimento viene chiamato \textit{three way handshake} dato che per
realizzarlo devono essere scambiati tre segmenti. In fig.~\ref{fig:TCP_TWH}
si è rappresentata graficamente la sequenza di scambio dei segmenti che
stabilisce la connessione.

% Una analogia citata da R. Stevens per la connessione TCP è quella con il
% sistema del telefono. La funzione \func{socket} può essere considerata
% l'equivalente di avere un telefono. La funzione \func{bind} è analoga al
% dire alle altre persone qual è il proprio numero di telefono perché possano
% chiamare. La funzione \func{listen} è accendere il campanello del telefono
% per sentire le chiamate in arrivo.  La funzione \func{connect} richiede di
% conoscere il numero di chi si vuole chiamare. La funzione \func{accept} è
% quando si risponde al telefono.

\begin{figure}[!htb]
  \centering \includegraphics[width=10cm]{img/three_way_handshake}  
  \caption{Il \textit{three way handshake} del TCP.}
  \label{fig:TCP_TWH}
\end{figure}

\index{numeri~di~sequenza|(}

Si è accennato in precedenza ai \textsl{numeri di sequenza}, che sono anche
riportati in fig.~\ref{fig:TCP_TWH}: per gestire una connessione affidabile
infatti il protocollo TCP prevede nell'intestazione la presenza di un numero a
32 bit (chiamato appunto \textit{sequence number}) che identifica a quale byte
nella sequenza del flusso corrisponde il primo byte della sezione dati
contenuta nel segmento.

Il numero di sequenza di ciascun segmento viene calcolato a partire da un
\textsl{numero di sequenza iniziale} generato in maniera casuale del kernel
all'inizio della connessione e trasmesso con il SYN;
l'\textit{acknowledgement} di ciascun segmento viene effettuato dall'altro
capo della connessione impostando il flag ACK e restituendo nell'apposito
campo dell'intestazione un \textit{acknowledge number}) pari al numero di
sequenza che il ricevente si aspetta di ricevere con il pacchetto successivo;
dato che il primo pacchetto SYN consuma un byte, nel \textit{three way
  handshake} il numero di \textit{acknowledge} è sempre pari al numero di
sequenza iniziale incrementato di uno; lo stesso varrà anche (vedi
fig.~\ref{fig:TCP_close}) per l'\textit{acknowledgement} di un segmento FIN.

\index{numeri~di~sequenza|)}
\itindend{three~way~handshake}


\subsection{Le opzioni TCP.}
\label{sec:TCP_TCP_opt}

Ciascun segmento SYN contiene in genere delle opzioni per il protocollo TCP,
le cosiddette \textit{TCP options}, da non confondere con le opzioni dei
socket TCP che tratteremo in sez.~\ref{sec:sock_tcp_udp_options}. In questo
caso infatti si tratta delle opzioni che vengono trasmesse come parte di un
pacchetto TCP, e non delle funzioni che consentono di impostare i relativi
valori. Queste opzioni vengono inserite fra l'intestazione ed i dati, e
servono a comunicare all'altro capo una serie di parametri utili a regolare la
connessione. Normalmente vengono usate le seguenti opzioni:

\begin{itemize*}
\item \textit{MSS option}, con questa opzione ciascun capo della connessione
  annuncia all'altro il massimo ammontare di dati (MMS sta appunto per
  \textit{Maximum Segment Size}, vedi sez.~\ref{sec:tcp_protocol}) che
  vorrebbe accettare per ciascun segmento nella connessione corrente. È
  possibile leggere e scrivere questo valore attraverso l'opzione del socket
  \const{TCP\_MAXSEG} (vedi sez.~\ref{sec:sock_tcp_udp_options}).
  
\item \textit{window scale option}, il protocollo TCP realizza il controllo di
  flusso attraverso una \textit{advertised window} (la ``\textsl{finestra
    annunciata}'', vedi sez.~\ref{sec:tcp_protocol}) con la quale ciascun capo
  della comunicazione dichiara quanto spazio disponibile ha in memoria per i
  dati. Questo è un numero a 16 bit dell'header, che così può indicare un
  massimo di 65535 byte\footnote{in Linux il massimo è 32767 per evitare
    problemi con alcune realizzazione dello \textit{stack TCP} che usano
    l'aritmetica con segno.} ma alcuni tipi di connessione come quelle ad alta
  velocità (sopra i 45Mbit/sec) e quelle che hanno grandi ritardi nel cammino
  dei pacchetti (come i satelliti) richiedono una finestra più grande per
  poter ottenere il massimo dalla trasmissione.

  Per questo esiste un'altra opzione che indica un fattore di scala da
  applicare al valore della finestra annunciata per la connessione corrente, 
  espresso come numero di bit cui spostare a sinistra il valore della
  finestra annunciata inserito nel pacchetto. Essendo una nuova opzione per
  garantire la compatibilità con delle vecchie realizzazione del protocollo la
  procedura che la attiva prevede come negoziazione che l'altro capo della
  connessione riconosca esplicitamente l'opzione inserendola anche lui nel suo
  SYN di risposta dell'apertura della connessione.

  Con Linux è possibile indicare al kernel di far negoziare il fattore di
  scala in fase di creazione di una connessione tramite la \textit{sysctl}
  \texttt{tcp\_window\_scaling} (vedi sez.~\ref{sec:sock_ipv4_sysctl}). Per
  poter usare questa funzionalità è comunque necessario ampliare le dimensioni
  dei buffer di ricezione e spedizione, cosa che può essere fatta sia a
  livello di sistema con le opportune \textit{sysctl} (vedi
  sez.~\ref{sec:sock_ipv4_sysctl}) che a livello di singoli socket con le
  relative opzioni (vedi sez.~\ref{sec:sock_tcp_udp_options}).

\item \textit{timestamp option}, è anche questa una nuova opzione necessaria
  per le connessioni ad alta velocità per evitare possibili corruzioni di dati
  dovute a pacchetti perduti che riappaiono; anche questa viene negoziata
  all'inizio della connessione come la precedente.
\end{itemize*}

La \textit{MSS option} è generalmente supportata da quasi tutte le
realizzazioni del protocollo, le ultime due opzioni (trattate
nell'\href{http://www.ietf.org/rfc/rfc1323.txt}{RFC~1323}) sono meno comuni;
vengono anche dette \textit{long fat pipe options} dato che questo è il nome
che viene dato alle connessioni caratterizzate da alta velocità o da ritardi
elevati. In ogni caso Linux supporta pienamente entrambe queste opzioni
aggiuntive.


\subsection{La terminazione della connessione}
\label{sec:TCP_conn_term}

Mentre per la creazione di una connessione occorre un interscambio di tre
segmenti, la procedura di chiusura ne richiede normalmente quattro. In questo
caso la successione degli eventi è la seguente:

\begin{enumerate*}
\item Un processo ad uno dei due capi chiama la funzione \func{close}, dando
  l'avvio a quella che viene chiamata \textsl{chiusura attiva} (o
  \textit{active close}). Questo comporta l'emissione di un segmento FIN, che
  serve ad indicare che si è finito con l'invio dei dati sulla connessione.
  
\item L'altro capo della connessione riceve il segmento FIN e dovrà eseguire
  la \textsl{chiusura passiva} (o \textit{passive close}). Al FIN, come ad
  ogni altro pacchetto, viene risposto con un ACK, inoltre il ricevimento del
  FIN viene segnalato al processo che ha aperto il socket (dopo che ogni altro
  eventuale dato rimasto in coda è stato ricevuto) come un
  \textit{end-of-file} sulla lettura: questo perché il ricevimento di un
  segmento FIN significa che non si riceveranno altri dati sulla connessione.
  
\item Una volta rilevata l'\textit{end-of-file} anche il secondo processo
  chiamerà la funzione \func{close} sul proprio socket, causando l'emissione
  di un altro segmento FIN.

\item L'altro capo della connessione riceverà il segmento FIN conclusivo e
  risponderà con un ACK.
\end{enumerate*}

Dato che in questo caso sono richiesti un FIN ed un ACK per ciascuna direzione
normalmente i segmenti scambiati sono quattro.  Questo non è vero sempre
giacché in alcune situazioni il FIN del passo 1 è inviato insieme a dei dati.
Inoltre è possibile che i segmenti inviati nei passi 2 e 3 dal capo che
effettua la chiusura passiva, siano accorpati in un singolo segmento. Come per
il SYN anche il FIN occupa un byte nel numero di sequenza, per cui l'ACK
riporterà un \textit{acknowledge number} incrementato di uno. In
fig.~\ref{fig:TCP_close} si è rappresentata graficamente la sequenza di
scambio dei segmenti che conclude la connessione.

\begin{figure}[!htb]
  \centering \includegraphics[width=10cm]{img/tcp_close}  
  \caption{La chiusura di una connessione TCP.}
  \label{fig:TCP_close}
\end{figure}

Si noti che, nella sequenza di chiusura, fra i passi 2 e 3, è in teoria
possibile che si mantenga un flusso di dati dal capo della connessione che
deve ancora eseguire la chiusura passiva a quello che sta eseguendo la
chiusura attiva.  Nella sequenza indicata i dati verrebbero persi, dato che si
è chiuso il socket dal lato che esegue la chiusura attiva; esistono tuttavia
situazioni in cui si vuole poter sfruttare questa possibilità, usando una
procedura che è chiamata \textit{half-close}; torneremo su questo aspetto e su
come utilizzarlo in sez.~\ref{sec:TCP_shutdown}, quando parleremo della
funzione \func{shutdown}.

La emissione del segmento FIN avviene quando il socket viene chiuso, questo
però non avviene solo per la chiamata esplicita della funzione \func{close},
ma anche alla terminazione di un processo, quando tutti i file vengono chiusi.
Questo comporta ad esempio che se un processo viene terminato da un segnale
tutte le connessioni aperte verranno chiuse.

Infine occorre sottolineare che, benché nella figura (e nell'esempio che
vedremo più avanti in sez.~\ref{sec:TCP_echo}) sia stato il client ad eseguire
la chiusura attiva, nella realtà questa può essere eseguita da uno qualunque
dei due capi della comunicazione (come nell'esempio di
fig.~\ref{fig:TCP_daytime_iter_server_code}), e che anche se il caso più
comune resta quello del client, ci sono alcuni servizi, il più noto dei quali
è l'HTTP 1.0 (con le versioni successive il default è cambiato) per i quali è
il server ad effettuare la chiusura attiva.


\subsection{Un esempio di connessione}
\label{sec:TCP_conn_dia}

Come abbiamo visto le operazioni del TCP nella creazione e conclusione di una
connessione sono piuttosto complesse, ed abbiamo esaminato soltanto quelle
relative ad un andamento normale.  In sez.~\ref{sec:TCP_states} vedremo con
maggiori dettagli che una connessione può assumere vari stati, che ne
caratterizzano il funzionamento, e che sono quelli che vengono riportati dal
comando \cmd{netstat}, per ciascun socket TCP aperto, nel campo
\textit{State}.

Non possiamo affrontare qui una descrizione completa del funzionamento del
protocollo; un approfondimento sugli aspetti principali si trova in
sez.~\ref{sec:tcp_protocol}, ma per una trattazione completa il miglior
riferimento resta \cite{TCPIll1}. Qui ci limiteremo a descrivere brevemente un
semplice esempio di connessione e le transizioni che avvengono nei due casi
appena citati (creazione e terminazione della connessione).

In assenza di connessione lo stato del TCP è \texttt{CLOSED}; quando una
applicazione esegue una apertura attiva il TCP emette un SYN e lo stato
diventa \texttt{SYN\_SENT}; quando il TCP riceve la risposta del SYN$+$ACK
emette un ACK e passa allo stato \texttt{ESTABLISHED}; questo è lo stato
finale in cui avviene la gran parte del trasferimento dei dati.

Dal lato server in genere invece il passaggio che si opera con l'apertura
passiva è quello di portare il socket dallo stato \texttt{CLOSED} allo
stato \texttt{LISTEN} in cui vengono accettate le connessioni.

Dallo stato \texttt{ESTABLISHED} si può uscire in due modi; se un'applicazione
chiama la funzione \func{close} prima di aver ricevuto un
\textit{end-of-file} (chiusura attiva) la transizione è verso lo stato
\texttt{FIN\_WAIT\_1}; se invece l'applicazione riceve un FIN nello stato
\texttt{ESTABLISHED} (chiusura passiva) la transizione è verso lo stato
\texttt{CLOSE\_WAIT}.

In fig.~\ref{fig:TCP_conn_example} è riportato lo schema dello scambio dei
pacchetti che avviene per una un esempio di connessione, insieme ai vari stati
che il protocollo viene ad assumere per i due lati, server e client.

\begin{figure}[!htb]
  \centering \includegraphics[width=9cm]{img/tcp_connection}  
  \caption{Schema dello scambio di pacchetti per un esempio di connessione.}
  \label{fig:TCP_conn_example}
\end{figure}

La connessione viene iniziata dal client che annuncia una \textit{Maximum
  Segment Size} di 1460, un valore tipico con Linux per IPv4 su Ethernet, il
server risponde con lo stesso valore (ma potrebbe essere anche un valore
diverso).

Una volta che la connessione è stabilita il client scrive al server una
richiesta (che assumiamo stare in un singolo segmento, cioè essere minore dei
1460 byte annunciati dal server), quest'ultimo riceve la richiesta e
restituisce una risposta (che di nuovo supponiamo stare in un singolo
segmento). Si noti che l'\textit{acknowledge} della richiesta è mandato
insieme alla risposta: questo viene chiamato \textit{piggybacking} ed avviene
tutte le volte che il server è sufficientemente veloce a costruire la
risposta; in caso contrario si avrebbe prima l'emissione di un ACK e poi
l'invio della risposta.

Infine si ha lo scambio dei quattro segmenti che terminano la connessione
secondo quanto visto in sez.~\ref{sec:TCP_conn_term}; si noti che il capo della
connessione che esegue la chiusura attiva entra nello stato
\texttt{TIME\_WAIT}, sul cui significato torneremo fra poco.

È da notare come per effettuare uno scambio di due pacchetti (uno di richiesta
e uno di risposta) TCP necessiti di ulteriori otto segmenti, se invece si
fosse usato UDP sarebbero stati sufficienti due soli pacchetti. Questo è il
costo che occorre pagare per avere l'affidabilità garantita da TCP, se si
fosse usato UDP si sarebbe dovuto trasferire la gestione di tutta una serie di
dettagli (come la verifica della ricezione dei pacchetti) dal livello del
trasporto all'interno dell'applicazione.

Quello che è bene sempre tenere presente è allora quali sono le esigenze che
si hanno in una applicazione di rete, perché non è detto che TCP sia la
miglior scelta in tutti i casi (ad esempio se si devono solo scambiare dati
già organizzati in piccoli pacchetti l'overhead aggiunto può essere eccessivo)
per questo esistono applicazioni che usano UDP e lo fanno perché nel caso
specifico le sue caratteristiche di velocità e compattezza nello scambio dei
dati rispondono meglio alle esigenze che devono essere affrontate.


\subsection{Lo stato \texttt{TIME\_WAIT}}
\label{sec:TCP_time_wait}

Come riportato da Stevens in \cite{UNP1} lo stato \texttt{TIME\_WAIT} è
probabilmente uno degli aspetti meno compresi del protocollo TCP, è infatti
comune trovare domande su come sia possibile evitare che un'applicazione resti
in questo stato lasciando attiva una connessione ormai conclusa; la risposta è
che non deve essere fatto, ed il motivo cercheremo di spiegarlo adesso.

\itindbeg{Maximum~Segment~Lifetime~(MSL)}
Come si è visto nell'esempio precedente (vedi fig.~\ref{fig:TCP_conn_example})
\texttt{TIME\_WAIT} è lo stato finale in cui il capo di una connessione che
esegue la chiusura attiva resta prima di passare alla chiusura definitiva
della connessione. Il tempo in cui l'applicazione resta in questo stato deve
essere due volte la \textit{Maximum Segment Lifetime} (da qui in avanti
abbreviata in MSL).

La \textit{Maximum Segment Lifetime} è la stima del massimo periodo di tempo
in secondi che un pacchetto IP può vivere sulla rete. Questo tempo è limitato
perché ogni pacchetto IP può essere ritrasmesso dai router un numero massimo
di volte (detto \textit{hop limit}).  Il numero di ritrasmissioni consentito è
indicato dal campo TTL dell'intestazione di IP (per maggiori dettagli vedi
sez.~\ref{sec:ip_protocol}), e viene decrementato ad ogni passaggio da un
router; quando si annulla il pacchetto viene scartato.  Siccome il numero è ad
8 bit il numero massimo di ``\textsl{salti}'' è di 255, pertanto anche se il
TTL (da \textit{time to live}) non è propriamente un limite sul tempo, sulla
sua base si può stimare che un pacchetto IP non possa restare nella rete per
più un certo numero di secondi, che costituisce la MSL.
\itindend{Maximum~Segment~Lifetime~(MSL)}

Ogni realizzazione del TCP deve scegliere un valore per la MSL;
l'\href{http://www.ietf.org/rfc/rfc1122.txt}{RFC~1122} raccomanda 2 minuti,
mentre Linux usa 30 secondi, questo comporta una durata dello stato
\texttt{TIME\_WAIT} che a seconda delle realizzazioni può variare fra 1 a 4
minuti.  Lo stato \texttt{TIME\_WAIT} viene utilizzato dal protocollo per due
motivi principali:
\begin{enumerate*}
\item effettuare in maniera affidabile la terminazione della connessione
  in entrambe le direzioni.
\item consentire l'eliminazione dei segmenti duplicati dalla rete. 
\end{enumerate*}

Il punto è che entrambe le ragioni sono importanti, anche se spesso si fa
riferimento solo alla prima; ma è solo se si tiene conto della seconda che si
capisce il perché della scelta di un tempo pari al doppio della MSL come
durata di questo stato.

Il primo dei due motivi precedenti si può capire tornando a
fig.~\ref{fig:TCP_conn_example}: assumendo che l'ultimo ACK della sequenza
(quello del capo che ha eseguito la chiusura attiva) venga perso, chi esegue
la chiusura passiva non ricevendo risposta rimanderà un ulteriore FIN, per
questo motivo chi esegue la chiusura attiva deve mantenere lo stato della
connessione per essere in grado di reinviare l'ACK e chiuderla correttamente.
Se non fosse così la risposta sarebbe un RST (un altro tipo si segmento) che
verrebbe interpretato come un errore.

Se il TCP deve poter chiudere in maniera pulita entrambe le direzioni della
connessione allora deve essere in grado di affrontare la perdita di uno
qualunque dei quattro segmenti che costituiscono la chiusura. Per questo
motivo un socket deve rimanere attivo nello stato \texttt{TIME\_WAIT} anche
dopo l'invio dell'ultimo ACK, per potere essere in grado di gestirne
l'eventuale ritrasmissione, in caso esso venga perduto.

Il secondo motivo è più complesso da capire, e necessita di una spiegazione
degli scenari in cui può accadere che i pacchetti TCP si possano perdere nella
rete o restare intrappolati, per poi riemergere in un secondo tempo. 

Il caso più comune in cui questo avviene è quello di anomalie
nell'instradamento; può accadere cioè che un router smetta di funzionare o che
una connessione fra due router si interrompa. In questo caso i protocolli di
instradamento dei pacchetti possono impiegare diverso tempo (anche dell'ordine
dei minuti) prima di trovare e stabilire un percorso alternativo per i
pacchetti. Nel frattempo possono accadere casi in cui un router manda i
pacchetti verso un altro e quest'ultimo li rispedisce indietro, o li manda ad
un terzo router che li rispedisce al primo, si creano cioè dei circoli (i
cosiddetti \textit{routing loop}) in cui restano intrappolati i pacchetti.

Se uno di questi pacchetti intrappolati è un segmento TCP, chi l'ha inviato,
non ricevendo un ACK in risposta, provvederà alla ritrasmissione e se nel
frattempo sarà stata stabilita una strada alternativa il pacchetto ritrasmesso
giungerà a destinazione.

Ma se dopo un po' di tempo (che non supera il limite dell'MSL, dato che
altrimenti verrebbe ecceduto il TTL) l'anomalia viene a cessare, il circolo di
instradamento viene spezzato i pacchetti intrappolati potranno essere inviati
alla destinazione finale, con la conseguenza di avere dei pacchetti duplicati;
questo è un caso che il TCP deve essere in grado di gestire.

Allora per capire la seconda ragione per l'esistenza dello stato
\texttt{TIME\_WAIT} si consideri il caso seguente: si supponga di avere una
connessione fra l'IP \texttt{195.110.112.236} porta 1550 e l'IP
\texttt{192.84.145.100} porta 22 (affronteremo il significato delle porte
nella prossima sezione), che questa venga chiusa e che poco dopo si
ristabilisca la stessa connessione fra gli stessi IP sulle stesse porte
(quella che viene detta, essendo gli stessi porte e numeri IP, una nuova
\textsl{incarnazione} della connessione precedente); in questo caso ci si
potrebbe trovare con dei pacchetti duplicati relativi alla precedente
connessione che riappaiono nella nuova.

Ma fintanto che il socket non è chiuso una nuova incarnazione non può essere
creata: per questo un socket TCP resta sempre nello stato \texttt{TIME\_WAIT}
per un periodo di 2 volte la MSL, prima per attendere MSL secondi per essere
sicuri che tutti i pacchetti duplicati in arrivo siano stati ricevuti e
scartati, o che nel frattempo siano stati eliminati dalla rete, poi altri MSL
secondi per essere sicuri che lo stesso avvenga per le risposte nella
direzione opposta.

In questo modo, prima che venga creata una nuova connessione, il protocollo
TCP si assicura che tutti gli eventuali segmenti residui di una precedente
connessione, che potrebbero causare disturbi, siano stati eliminati dalla
rete.


\subsection{I numeri di porta}
\label{sec:TCP_port_num}

In un ambiente multitasking in un dato momento più processi devono poter usare
sia UDP che TCP, e ci devono poter essere più connessioni in contemporanea.
Per poter tenere distinte le diverse connessioni entrambi i protocolli usano i
\textsl{numeri di porta}, che fanno parte, come si può vedere in
sez.~\ref{sec:sock_sa_ipv4} e sez.~\ref{sec:sock_sa_ipv6} pure delle strutture
degli indirizzi del socket.

Quando un client contatta una macchina server deve poter identificare con
quale dei vari possibili programmi server attivi intende parlare. Sia TCP che
UDP definiscono un gruppo di \textsl{porte conosciute} (le cosiddette
\textit{well-known port}) che identificano una serie di servizi noti (ad
esempio la porta 22 identifica il servizio SSH) effettuati da appositi
programmi server che rispondono alle connessioni verso tali porte.

D'altra parte un client non ha necessità di usare dalla sua parte un numero di
porta specifico, per cui in genere vengono usate le cosiddette \textsl{porte
  effimere} (o \textit{ephemeral ports}) cioè porte a cui non è assegnato
nessun servizio noto e che vengono assegnate automaticamente dal kernel alla
creazione della connessione. Queste sono dette effimere in quanto vengono
usate solo per la durata della connessione, e l'unico requisito che deve
essere soddisfatto è che ognuna di esse sia assegnata in maniera univoca.

La lista delle porte conosciute è definita
dall'\href{http://www.ietf.org/rfc/rfc1700.txt}{RFC~1700} che contiene
l'elenco delle porte assegnate dalla IANA (la \textit{Internet Assigned Number
  Authority}) ma l'elenco viene costantemente aggiornato e pubblicato su
Internet (una versione aggiornata si può trovare all'indirizzo
\url{http://www.iana.org/assignments/port-numbers}); inoltre in un sistema
unix-like un analogo elenco viene mantenuto nel file \conffile{/etc/services},
con la corrispondenza fra i vari numeri di porta ed il nome simbolico del
servizio.  I numeri sono divisi in tre intervalli:

\begin{enumerate*}
\item \textsl{le porte note}. I numeri da 0 a 1023. Queste sono controllate e
  assegnate dalla IANA. Se è possibile la stessa porta è assegnata allo stesso
  servizio sia su UDP che su TCP (ad esempio la porta 22 è assegnata a SSH su
  entrambi i protocolli, anche se viene usata solo dal TCP).
  
\item \textsl{le porte registrate}. I numeri da 1024 a 49151. Queste porte non
  sono controllate dalla IANA, che però registra ed elenca chi usa queste
  porte come servizio agli utenti. Come per le precedenti si assegna una porta
  ad un servizio sia per TCP che UDP anche se poi il servizio è effettuato
  solo su TCP. Ad esempio \textit{X Window} usa le porte TCP e UDP dal 6000 al
  6063 anche se il protocollo viene usato solo con TCP.
  
\item \textsl{le porte private} o \textsl{dinamiche}. I numeri da 49152 a
  65535. La IANA non dice nulla riguardo a queste porte che pertanto
  sono i candidati naturali ad essere usate come porte effimere.
\end{enumerate*}

In realtà rispetto a quanto indicato
nell'\href{http://www.ietf.org/rfc/rfc1700.txt}{RFC~1700} i vari sistemi hanno
fatto scelte diverse per le porte effimere, in particolare in
fig.~\ref{fig:TCP_port_alloc} sono riportate quelle di BSD e Linux.

\begin{figure}[!htb]
  \centering \includegraphics[width=13cm]{img/port_alloc}  
  \caption{Allocazione dei numeri di porta.}
  \label{fig:TCP_port_alloc}
\end{figure}

I sistemi Unix hanno inoltre il concetto di \textsl{porte riservate}, che
corrispondono alle porte con numero minore di 1024 e coincidono quindi con le
\textsl{porte note}. La loro caratteristica è che possono essere assegnate a
un socket solo da un processo con i privilegi di amministratore, per far sì
che solo l'amministratore possa allocare queste porte per far partire i
relativi servizi.

Le \textsl{glibc} definiscono in \headfile{netinet/in.h} le costanti
\constd{IPPORT\_RESERVED} e \constd{IPPORT\_USERRESERVED}, in cui la prima
(che vale 1024) indica il limite superiore delle porte riservate, e la seconda
(che vale 5000) il limite inferiore delle porte a disposizione degli utenti.
La convenzione vorrebbe che le porte \textsl{effimere} siano allocate fra
questi due valori. Nel caso di Linux questo è vero solo in uno dei due casi di
fig.~\ref{fig:TCP_port_alloc}, e la scelta fra i due possibili intervalli
viene fatta dinamicamente dal kernel a seconda della memoria disponibile per
la gestione delle relative tabelle.

Si tenga conto poi che ci sono alcuni client, in particolare \cmd{rsh} e
\cmd{rlogin}, che richiedono una connessione su una porta riservata anche dal
lato client come parte dell'autenticazione, contando appunto sul fatto che
solo l'amministratore può usare queste porte. Data l'assoluta inconsistenza in
termini di sicurezza di un tale metodo, al giorno d'oggi esso è in completo
disuso.

Data una connessione TCP, ma la cosa vale anche per altri protocolli del
livello di trasporto come UDP, si suole chiamare \itindex{socket~pair}
\textit{socket pair}\footnote{da non confondere con la coppia di socket della
  omonima funzione \func{socketpair} di sez.~\ref{sec:ipc_socketpair} che
  fanno riferimento ad una coppia di socket sulla stessa macchina, non ai capi
  di una connessione TCP.} la combinazione dei quattro numeri che definiscono
i due capi della connessione e cioè l'indirizzo IP locale e la porta TCP
locale, e l'indirizzo IP remoto e la porta TCP remota.  Questa combinazione,
che scriveremo usando una notazione del tipo (\texttt{195.110.112.152:22},
\texttt{192.84.146.100:20100}), identifica univocamente una connessione su
Internet.  Questo concetto viene di solito esteso anche a UDP, benché in
questo caso non abbia senso parlare di connessione. L'utilizzo del programma
\cmd{netstat} permette di visualizzare queste informazioni nei campi
\textit{Local Address} e \textit{Foreing Address}.

Per capire meglio l'uso delle porte e come vengono utilizzate quando si ha a
che fare con un'applicazione client/server (come quelle che descriveremo in
sez.~\ref{sec:TCP_daytime_application} e sez.~\ref{sec:TCP_echo_application})
esamineremo cosa accade con le connessioni nel caso di un server TCP che deve
gestire connessioni multiple.

Se eseguiamo un \cmd{netstat} su una macchina di prova (il cui indirizzo sia
\texttt{195.110.112.152}) potremo avere un risultato del tipo:
\begin{Terminal}
Active Internet connections (servers and established)
Proto Recv-Q Send-Q Local Address           Foreign Address         State      
tcp        0      0 0.0.0.0:22              0.0.0.0:*               LISTEN
tcp        0      0 0.0.0.0:25              0.0.0.0:*               LISTEN
tcp        0      0 127.0.0.1:53            0.0.0.0:*               LISTEN
\end{Terminal}
essendo presenti e attivi un server SSH, un server di posta e un DNS per il
caching locale.

Questo ci mostra ad esempio che il server SSH ha compiuto un'apertura passiva,
mettendosi in ascolto sulla porta 22 riservata a questo servizio, e che si è
posto in ascolto per connessioni provenienti da uno qualunque degli indirizzi
associati alle interfacce locali. La notazione \texttt{0.0.0.0} usata da
\cmd{netstat} è equivalente all'asterisco utilizzato per il numero di porta,
indica il valore generico, e corrisponde al valore \const{INADDR\_ANY}
definito in \headfiled{arpa/inet.h} (vedi \ref{tab:TCP_ipv4_addr}).

Inoltre si noti come la porta e l'indirizzo di ogni eventuale connessione
esterna non sono specificati; in questo caso la \textit{socket pair} associata
al socket potrebbe essere indicata come (\texttt{*:22}, \texttt{*:*}), usando
anche per gli indirizzi l'asterisco come carattere che indica il valore
generico.

Dato che in genere una macchina è associata ad un solo indirizzo IP, ci si può
chiedere che senso abbia l'utilizzo dell'indirizzo generico per specificare
l'indirizzo locale; ma a parte il caso di macchine che hanno più di un
indirizzo IP (il cosiddetto \textit{multihoming}) esiste sempre anche
l'indirizzo di \textit{loopback}, per cui con l'uso dell'indirizzo generico si
possono accettare connessioni indirizzate verso uno qualunque degli indirizzi
IP presenti. Ma, come si può vedere nell'esempio con il DNS che è in ascolto
sulla porta 53, è possibile anche restringere l'accesso ad uno specifico
indirizzo, cosa che nel caso è fatta accettando solo connessioni che arrivino
sull'interfaccia di \textit{loopback}.

Una volta che ci si vorrà collegare a questa macchina da un'altra, per esempio
quella con l'indirizzo \texttt{192.84.146.100}, si dovrà lanciare su
quest'ultima un client \cmd{ssh} per creare una connessione, e il kernel gli
assocerà una porta effimera (ad esempio la 21100), per cui la connessione sarà
espressa dalla \textit{socket pair} (\texttt{192.84.146.100:21100},
\texttt{195.110.112.152:22}).

Alla ricezione della richiesta dal client il server creerà un processo figlio
per gestire la connessione, se a questo punto eseguiamo nuovamente il
programma \cmd{netstat} otteniamo come risultato:
\begin{Terminal}
Active Internet connections (servers and established)
Proto Recv-Q Send-Q Local Address           Foreign Address         State      
tcp        0      0 0.0.0.0:22              0.0.0.0:*               LISTEN
tcp        0      0 0.0.0.0:25              0.0.0.0:*               LISTEN
tcp        0      0 127.0.0.1:53            0.0.0.0:*               LISTEN
tcp        0      0 195.110.112.152:22      192.84.146.100:21100    ESTABLISHED
\end{Terminal}

Come si può notare il server è ancora in ascolto sulla porta 22, però adesso
c'è un nuovo socket (con lo stato \texttt{ESTABLISHED}) che utilizza anch'esso
la porta 22, ed ha specificato l'indirizzo locale, questo è il socket con cui
il processo figlio gestisce la connessione mentre il padre resta in ascolto
sul socket originale.

Se a questo punto lanciamo un'altra volta il client \cmd{ssh} per una seconda
connessione quello che otterremo usando \cmd{netstat} sarà qualcosa del
genere:
\begin{Terminal}
Active Internet connections (servers and established)
Proto Recv-Q Send-Q Local Address           Foreign Address         State      
tcp        0      0 0.0.0.0:22              0.0.0.0:*               LISTEN
tcp        0      0 0.0.0.0:25              0.0.0.0:*               LISTEN
tcp        0      0 127.0.0.1:53            0.0.0.0:*               LISTEN
tcp        0      0 195.110.112.152:22      192.84.146.100:21100    ESTABLISHED
tcp        0      0 195.110.112.152:22      192.84.146.100:21101    ESTABLISHED
\end{Terminal}
cioè il client effettuerà la connessione usando un'altra porta effimera: con
questa sarà aperta la connessione, ed il server creerà un altro processo
figlio per gestirla.

Tutto ciò mostra come il TCP, per poter gestire le connessioni con un server
concorrente, non può suddividere i pacchetti solo sulla base della porta di
destinazione, ma deve usare tutta l'informazione contenuta nella
\textit{socket pair}, compresa la porta dell'indirizzo remoto.  E se andassimo
a vedere quali sono i processi (ad esempio con il comando \cmd{fuser}, o con
\cmd{lsof}, o usando l'opzione \texttt{-p}) a cui fanno riferimento i vari
socket vedremmo che i pacchetti che arrivano dalla porta remota 21100 vanno al
primo figlio e quelli che arrivano alla porta 21101 al secondo.


\section{Le funzioni di base per la gestione dei socket}
\label{sec:TCP_functions}

In questa sezione descriveremo in maggior dettaglio le varie funzioni che
vengono usate per la gestione di base dei socket TCP, non torneremo però sulla
funzione \func{socket}, che è già stata esaminata accuratamente nel capitolo
precedente in sez.~\ref{sec:sock_creation}. Pur trattandole principalmente dal
punto di vista dei socket TCP, daremo brevemente conto del loro comportamento
anche per altri tipi di socket.


\subsection{La funzione \func{bind}}
\label{sec:TCP_func_bind}

La funzione \funcd{bind} assegna un indirizzo locale ad un
socket,\footnote{nel nostro caso la utilizzeremo per socket TCP, ma la
  funzione è generica e deve essere usata per qualunque tipo di socket
  \const{SOCK\_STREAM} prima che questo possa accettare connessioni.} ed è
usata cioè per specificare la prima parte dalla \textit{socket pair}.  Viene
usata sul lato server per specificare la porta (e gli eventuali indirizzi
locali) su cui poi ci si porrà in ascolto. Il prototipo della funzione è il
seguente:


\begin{funcproto}{
\fhead{sys/socket.h}
\fdecl{int bind(int sockfd, const struct sockaddr *serv\_addr, socklen\_t
  addrlen)} 
\fdesc{Assegna un indirizzo ad un socket.} 
}

{La funzione ritorna $0$ in caso di successo e $-1$ per un errore, nel qual
  caso \var{errno} assumerà uno dei valori: 
  \begin{errlist}
  \item[\errcode{EACCES}] si è cercato di usare una porta riservata senza
    sufficienti privilegi.
  \item[\errcode{EADDRINUSE}] qualche altro socket sta già usando l'indirizzo
    richiesto oppure quando non si è richiesta un porta specifica per avere
    una porta effimera non ve ne sono di disponibili nell'intervallo ad esse
    riservato correntemente in uso.
  \item[\errcode{EBADF}] il file descriptor non è valido.
  \item[\errcode{EINVAL}] il socket ha già un indirizzo assegnato,
    o \param{addrlen} è sbagliata, o \param{serv\_addr} non è valido per il
    dominio.
  \item[\errcode{ENOTSOCK}] il file descriptor non è associato ad un socket.
  \end{errlist}
  ed \errval{EFAULT} nel suo significato generico; inoltre per i socket di
  tipo \const{AF\_UNIX}:
  \begin{errlist}
  \item[\errcode{EADDRNOTAVAIL}] il tipo di indirizzo specificato non è
    disponibile.
  \end{errlist}
  ed \errval{ELOOP}, \errval{ENAMETOOLONG}, \errval{ENOENT}, \errval{ENOMEM}, 
  \errval{ENOTDIR} e \errval{EROFS} nel loro significato generico.}
\end{funcproto}

Il primo argomento è un file descriptor ottenuto da una precedente chiamata a
\func{socket}, mentre il secondo e terzo argomento sono rispettivamente
l'indirizzo (locale) del socket e la dimensione della struttura che lo
contiene, secondo quanto già trattato in sez.~\ref{sec:sock_sockaddr}. 

Con i socket TCP la chiamata \func{bind} permette di specificare l'indirizzo,
la porta, entrambi o nessuno dei due. In genere i server utilizzano una porta
nota che assegnano all'avvio, se questo non viene fatto è il kernel a
scegliere una porta effimera quando vengono eseguite la funzioni
\func{connect} o \func{listen}, ma se questo è normale per il client non lo è
per il server\footnote{un'eccezione a tutto ciò sono i server che usano RPC.
  In questo caso viene fatta assegnare dal kernel una porta effimera che poi
  viene registrata presso il \textit{portmapper}; quest'ultimo è un altro
  demone che deve essere contattato dai client per ottenere la porta effimera
  su cui si trova il server.} che in genere viene identificato dalla porta su
cui risponde (l'elenco di queste porte, e dei relativi servizi, è in
\conffile{/etc/services}).

Con \func{bind} si può assegnare un indirizzo IP specifico ad un socket,
purché questo appartenga ad una interfaccia della macchina. Per un client TCP
questo diventerà l'indirizzo sorgente usato per i tutti i pacchetti inviati
sul socket, mentre per un server TCP questo restringerà l'accesso al socket
solo alle connessioni che arrivano verso tale indirizzo.

Normalmente un client non specifica mai l'indirizzo di un socket, ed il kernel
sceglie l'indirizzo di origine quando viene effettuata la connessione, sulla
base dell'interfaccia usata per trasmettere i pacchetti, che dipenderà dalle
regole di instradamento usate per raggiungere il server (è comunque possibile
impostarlo in maniera specifica con i comandi di gestione avanzata del
routing, vedi sez.~7.3.4 di \cite{AGL}). Se un server non specifica il suo
indirizzo locale il kernel userà come indirizzo di origine l'indirizzo di
destinazione specificato dal SYN del client.

Per specificare un indirizzo generico, con IPv4 si usa il valore
\const{INADDR\_ANY}, il cui valore, come accennato in
sez.~\ref{sec:sock_sa_ipv4}, è pari a zero; nell'esempio
fig.~\ref{fig:TCP_daytime_iter_server_code} si è usata un'assegnazione
immediata del tipo: \includecodesnip{listati/serv_addr_sin_addr.c}

Si noti che si è usato \func{htonl} per assegnare il valore
\const{INADDR\_ANY}, anche se, essendo questo nullo, il riordinamento è
inutile.  Si tenga presente comunque che tutte le costanti \val{INADDR\_}
(riportate in tab.~\ref{tab:TCP_ipv4_addr}) sono definite secondo
l'\textit{endianness} della macchina, ed anche se esse possono essere
invarianti rispetto all'ordinamento dei bit, è comunque buona norma usare
sempre la funzione \func{htonl}.

\begin{table}[htb]
  \centering
  \footnotesize
  \begin{tabular}[c]{|l|l|}
    \hline
    \textbf{Costante} & \textbf{Significato} \\
    \hline
    \hline
    \constd{INADDR\_ANY}      & Indirizzo generico (\texttt{0.0.0.0})\\
    \constd{INADDR\_BROADCAST}& Indirizzo di \textit{broadcast}.\\ 
    \constd{INADDR\_LOOPBACK} & Indirizzo di \textit{loopback}
                                (\texttt{127.0.0.1}).\\ 
    \constd{INADDR\_NONE}     & Indirizzo errato.\\
    \hline    
  \end{tabular}
  \caption{Costanti di definizione di alcuni indirizzi generici per IPv4.}
  \label{tab:TCP_ipv4_addr}
\end{table}

L'esempio precedente funziona correttamente con IPv4 poiché che l'indirizzo è
rappresentabile anche con un intero a 32 bit; non si può usare lo stesso
metodo con IPv6, in cui l'indirizzo deve necessariamente essere specificato
con una struttura, perché il linguaggio C non consente l'uso di una struttura
costante come operando a destra in una assegnazione.

Per questo motivo nell'header \headfile{netinet/in.h} è definita una variabile
\var{in6addr\_any} (dichiarata come \dirct{extern}, ed inizializzata dal
sistema al valore \constd{IN6ADRR\_ANY\_INIT}) che permette di effettuare una
assegnazione del tipo: \includecodesnip{listati/serv_addr_sin6_addr.c} in
maniera analoga si può utilizzare la variabile \var{in6addr\_loopback} per
indicare l'indirizzo di \textit{loopback}, che a sua volta viene inizializzata
staticamente a \constd{IN6ADRR\_LOOPBACK\_INIT}.


\subsection{La funzione \func{connect}}
\label{sec:TCP_func_connect}

La funzione \funcd{connect} è usata da un client TCP per stabilire la
connessione con un server TCP, ma la funzione è generica e supporta vari tipi
di socket.  La differenza è che per socket senza connessione come quelli di
tipo \const{SOCK\_DGRAM} la sua chiamata si limiterà ad impostare l'indirizzo
dal quale e verso il quale saranno inviati e ricevuti i pacchetti, mentre per
socket di tipo \const{SOCK\_STREAM} o \const{SOCK\_SEQPACKET} essa attiverà
la procedura di avvio della connessione (nel caso del TCP il \textit{three way
  handshake}). Il suo prototipo è il seguente:


\begin{funcproto}{
\fhead{sys/socket.h}
\fdecl{int connect(int sockfd, const struct sockaddr *servaddr, socklen\_t
    addrlen)}
\fdesc{Stabilisce una connessione fra due socket.} 
}

{La funzione ritorna $0$ in caso di successo e $-1$ per un errore, nel qual
  caso \var{errno} assumerà uno dei valori: 
  \begin{errlist}
  \item[\errcode{EACCES}, \errcode{EPERM}] si è tentato di eseguire una
    connessione ad un indirizzo \textit{broadcast} senza che il socket fosse
    stato abilitato per il \textit{broadcast}.
  \item[\errcode{EAFNOSUPPORT}] l'indirizzo non ha una famiglia di indirizzi
    corretta nel relativo campo.
  \item[\errcode{EAGAIN}] non ci sono più porte locali libere. 
  \item[\errcode{EALREADY}] il socket è non bloccante (vedi
    sez.~\ref{sec:file_noblocking}) e un tentativo precedente di connessione
    non si è ancora concluso.
  \item[\errcode{ECONNREFUSED}] non c'è nessuno in ascolto sull'indirizzo
    remoto.
  \item[\errcode{EINPROGRESS}] il socket è non bloccante (vedi
    sez.~\ref{sec:file_noblocking}) e la connessione non può essere conclusa
    immediatamente.
  \item[\errcode{ENETUNREACH}] la rete non è raggiungibile.
  \item[\errcode{ETIMEDOUT}] si è avuto timeout durante il tentativo di
    connessione.
  \end{errlist}
  ed inoltre \errval{EADDRINUSE}, \errval{EBADF}, \errval{EFAULT},
  \errval{EINTR}, \errval{EISCONN} e \errval{ENOTSOCK} e nel loro significato
  generico.}
\end{funcproto}

Il primo argomento è un file descriptor ottenuto da una precedente chiamata a
\func{socket}, mentre il secondo e terzo argomento sono rispettivamente
l'indirizzo e la dimensione della struttura che contiene l'indirizzo del
socket, già descritta in sez.~\ref{sec:sock_sockaddr}.

La struttura dell'indirizzo deve essere inizializzata con l'indirizzo IP e il
numero di porta del server a cui ci si vuole connettere usando le funzioni
illustrate in sez.~\ref{sec:sock_addr_func} come mostrato nell'esempio che
vedremo in sez.~\ref{sec:TCP_daytime_client}.

Nel caso di socket TCP la funzione \func{connect} avvia il \textit{three way
  handshake}, e ritorna solo quando la connessione è stabilita o si è
verificato un errore. Le possibili cause di errore sono molteplici (ed i
relativi codici riportati sopra), quelle che però dipendono dalla situazione
della rete e non da errori o problemi nella chiamata della funzione sono le
seguenti:
\begin{enumerate}
\item Il client non riceve risposta al SYN: l'errore restituito è
  \errcode{ETIMEDOUT}. Stevens riporta che BSD invia un primo SYN alla
  chiamata di \func{connect}, un altro dopo 6 secondi, un terzo dopo 24
  secondi, se dopo 75 secondi non ha ricevuto risposta viene ritornato
  l'errore. Linux invece ripete l'emissione del SYN ad intervalli di 30
  secondi per un numero di volte che può essere stabilito dall'utente. Questo
  può essere fatto a livello globale con una opportuna \func{sysctl} (o più
  semplicemente scrivendo il valore voluto in
  \sysctlfile{net/ipv4/tcp\_syn\_retries}, vedi
  sez.~\ref{sec:sock_ipv4_sysctl}) e a livello di singolo socket con l'opzione
  \const{TCP\_SYNCNT} (vedi sez.~\ref{sec:sock_tcp_udp_options}). Il valore
  predefinito per la ripetizione dell'invio è di 5 volte, che comporta un
  timeout dopo circa 180 secondi.

\item Il client riceve come risposta al SYN un RST significa che non c'è
  nessun programma in ascolto per la connessione sulla porta specificata (il
  che vuol dire probabilmente che o si è sbagliato il numero della porta o che
  non è stato avviato il server), questo è un errore fatale e la funzione
  ritorna non appena il RST viene ricevuto riportando un errore
  \errcode{ECONNREFUSED}.
  
  Il flag RST sta per \textit{reset} ed è un segmento inviato direttamente
  dal TCP quando qualcosa non va. Tre condizioni che generano un RST sono:
  quando arriva un SYN per una porta che non ha nessun server in ascolto,
  quando il TCP abortisce una connessione in corso, quando TCP riceve un
  segmento per una connessione che non esiste.
  
\item Il SYN del client provoca l'emissione di un messaggio ICMP di
  destinazione non raggiungibile. In questo caso dato che il messaggio può
  essere dovuto ad una condizione transitoria si ripete l'emissione dei SYN
  come nel caso precedente, fino al timeout, e solo allora si restituisce il
  codice di errore dovuto al messaggio ICMP, che da luogo ad un
  \errcode{ENETUNREACH}.
   
\end{enumerate}

Se si fa riferimento al diagramma degli stati del TCP riportato in
fig.~\ref{fig:TCP_state_diag} la funzione \func{connect} porta un socket dallo
stato \texttt{CLOSED} (lo stato iniziale in cui si trova un socket appena
creato) prima allo stato \texttt{SYN\_SENT} e poi, al ricevimento dell'ACK,
nello stato \texttt{ESTABLISHED}. Se invece la connessione fallisce il socket
non è più utilizzabile e deve essere chiuso.

Si noti infine che con la funzione \func{connect} si è specificato solo
indirizzo e porta del server, quindi solo una metà della \textit{socket pair};
essendo questa funzione usata nei client l'altra metà contenente indirizzo e
porta locale viene lasciata all'assegnazione automatica del kernel, e non è
necessario effettuare una \func{bind}.


\subsection{La funzione \func{listen}}
\label{sec:TCP_func_listen}

La funzione \funcd{listen} serve ad usare un socket in modalità passiva, cioè,
come dice il nome, per metterlo in ascolto di eventuali
connessioni;\footnote{questa funzione può essere usata con socket che
  supportino le connessioni, cioè di tipo \const{SOCK\_STREAM} o
  \const{SOCK\_SEQPACKET}.} in sostanza l'effetto della funzione è di portare
il socket dallo stato \texttt{CLOSED} a quello \texttt{LISTEN}. In genere si
chiama la funzione in un server dopo le chiamate a \func{socket} e \func{bind}
e prima della chiamata ad \func{accept}. Il prototipo della funzione, come
definito dalla pagina di manuale, è:

\begin{funcproto}{
\fhead{sys/socket.h}
\fdecl{int listen(int sockfd, int backlog)}
\fdesc{Pone un socket in attesa di una connessione.} 
}

{La funzione ritorna $0$ in caso di successo e $-1$ per un errore, nel qual
  caso \var{errno} assumerà uno dei valori: 
  \begin{errlist}
  \item[\errcode{EADDRINUSE}] qualche altro socket sta già usando l'indirizzo.
  \item[\errcode{EBADF}] l'argomento \param{sockfd} non è un file descriptor
    valido.
  \item[\errcode{ENOTSOCK}] l'argomento \param{sockfd} non è un socket.
  \item[\errcode{EOPNOTSUPP}] il socket è di un tipo che non supporta questa
    operazione.
  \end{errlist}
}
\end{funcproto}

La funzione pone il socket specificato da \param{sockfd} in modalità passiva e
predispone una coda per le connessioni in arrivo di lunghezza pari a
\param{backlog}. La funzione si può applicare solo a socket di tipo
\const{SOCK\_STREAM} o \const{SOCK\_SEQPACKET}. L'argomento \param{backlog}
indica il numero massimo di connessioni pendenti accettate; se esso viene
ecceduto il client al momento della richiesta della connessione riceverà un
errore di tipo \errcode{ECONNREFUSED}, o se il protocollo, come accade nel
caso del TCP, supporta la ritrasmissione, la richiesta sarà ignorata in modo
che la connessione possa venire ritentata.

Per capire meglio il significato di tutto ciò occorre approfondire la modalità
con cui il kernel tratta le connessioni in arrivo. Per ogni socket in ascolto
infatti vengono mantenute due code:
\begin{enumerate}
\item La coda delle connessioni incomplete (\textit{incomplete connection
    queue}) che contiene una voce per ciascun socket per il quale è arrivato
  un SYN ma il \textit{three way handshake} non si è ancora concluso.  Questi
  socket sono tutti nello stato \texttt{SYN\_RECV}.
\item La coda delle connessioni complete (\textit{complete connection queue})
  che contiene una voce per ciascun socket per il quale il \textit{three way
    handshake} è stato completato ma ancora \func{accept} non è ritornata.
  Questi socket sono tutti nello stato \texttt{ESTABLISHED}.
\end{enumerate}

Lo schema di funzionamento è descritto in fig.~\ref{fig:TCP_listen_backlog}:
quando arriva un segmento SYN da un client il kernel crea una voce nella coda
delle connessioni incomplete e risponde con il segmento SYN$+$ACK. La voce
resterà nella coda delle connessioni incomplete fino al ricevimento del
segmento ACK dal client o fino ad un timeout.

Nel caso di completamento del \textit{three way handshake} la voce viene
spostata nella coda delle connessioni complete.  Quando il processo chiama la
funzione \func{accept} (vedi sez.~\ref{sec:TCP_func_accept}) gli viene passata
la prima voce nella coda delle connessioni complete, oppure, se la coda è
vuota, il processo viene posto in attesa in stato di \textit{sleep} e
risvegliato all'arrivo della prima connessione completa.

\begin{figure}[!htb]
  \centering \includegraphics[width=12cm]{img/tcp_listen_backlog}  
  \caption{Schema di funzionamento delle code delle connessioni complete ed
    incomplete.}
  \label{fig:TCP_listen_backlog}
\end{figure}

Storicamente il valore dell'argomento \param{backlog} era corrispondente al
massimo valore della somma del numero di voci possibili per ciascuna delle due
code. Stevens in \cite{UNP1} riporta che BSD ha sempre applicato un fattore di
1.5 a detto valore, e fornisce una tabella con i risultati ottenuti con vari
kernel, compreso anche il vecchio Linux 2.0, che mostrano le differenze fra
diverse realizzazioni.

In Linux il significato di questo valore è cambiato a partire dal kernel 2.2
per prevenire il \textit{denial of service} chiamato \itindex{SYN~flood}
\textit{SYN flood}. Questo attacco si basa sull'emissione da parte
dell'attaccante di un grande numero di segmenti SYN indirizzati verso una
porta, forgiati con indirizzo IP fasullo (con la tecnica che viene detta
\textit{ip spoofing}); in questo modo i segmenti SYN$+$ACK di risposta vanno
perduti e la coda delle connessioni incomplete viene saturata, impedendo di
fatto ulteriori connessioni.

Per ovviare a questo problema il significato del \param{backlog} è stato
cambiato e adesso indica la lunghezza della coda delle connessioni
complete. La lunghezza della coda delle connessioni incomplete può essere
ancora controllata ma occorre usare esplicitamente la funzione \func{sysctl}
con il parametro \constd{NET\_TCP\_MAX\_SYN\_BACKLOG} o scrivere il valore 
direttamente sul file \sysctlfile{net/ipv4/tcp\_max\_syn\_backlog}.

Quando si attiva la protezione dei \textit{syncookies} però (con l'opzione da
compilare nel kernel e da attivare usando \func{sysctl} o scrivendo nel file
\sysctlfile{net/ipv4/tcp\_syncookies}, vedi sez.~\ref{sec:sock_ipv4_sysctl})
questo valore viene ignorato e non esiste più un valore massimo.  In ogni caso
in Linux il valore di \param{backlog} viene sempre troncato ad un massimo di
\const{SOMAXCONN} se è superiore a detta costante (che di default vale 128);
per i kernel precedenti il 2.4.25 questo valore era fisso e non modificabile,
nelle versioni successive può essere controllato con un parametro di
\func{sysctl}, o scrivendo nel file \sysctlfile{net/core/somaxconn} 
(vedi sez.~\ref{sec:sock_ioctl_IP}).

La scelta storica per il valore assegnato a questo argomento era di 5, e
alcuni vecchi kernel non supportavano neanche valori superiori, ma la
situazione corrente è molto cambiata per via della presenza di server web che
devono gestire un gran numero di connessioni per cui un tale valore non è più
adeguato. Non esiste comunque una risposta univoca per la scelta del valore,
per questo non conviene specificarlo con una costante (il cui cambiamento
richiederebbe la ricompilazione del server) ma usare piuttosto una variabile
di ambiente (vedi sez.~\ref{sec:proc_environ}). 

Stevens tratta accuratamente questo argomento in \cite{UNP1}, con esempi presi
da casistiche reali trovate su dei server web, ed in particolare evidenzia
come non sia più vero che il compito principale della coda sia quello di
gestire il caso in cui il server è occupato fra chiamate successive alla
\func{accept} (per cui la coda più occupata sarebbe quella delle connessioni
completate), ma piuttosto quello di gestire la presenza di un gran numero di
SYN in attesa di concludere il \textit{three way handshake}.

Infine va messo in evidenza che, nel caso di socket TCP, quando un segmento
SYN arriva con tutte le code piene, il pacchetto verrà semplicemente
ignorato. Questo avviene perché la condizione in cui le code sono piene è
ovviamente transitoria, per cui se il client ritrasmette in seguito un
segmento SYN, come previsto dal protocollo, è probabile che essendo passato un
po' di tempo esso possa trovare nella coda lo spazio per una nuova
connessione.

Se al contrario si rispondesse immediatamente con un segmento RST, per
indicare che è impossibile effettuare la connessione, la chiamata a
\func{connect} eseguita dal client fallirebbe ritornando una condizione di
errore. In questo modo si sarebbe costretti ad inserire nell'applicazione la
gestione dei tentativi di riconnessione, che invece grazie a questa modalità
di funzionamento viene effettuata in maniera trasparente dal protocollo
TCP.


\subsection{La funzione \func{accept}}
\label{sec:TCP_func_accept}

La funzione \funcd{accept} è chiamata da un server per gestire la connessione,
nel caso di TCP una volta che sia stato completato il \textit{three way
  handshake};\footnote{come le precedenti, la funzione è generica ed è
  utilizzabile su socket di tipo \const{SOCK\_STREAM}, \const{SOCK\_SEQPACKET}
  e \const{SOCK\_RDM}, ma qui la tratteremo solo per gli aspetti riguardanti
  le connessioni con TCP.} la funzione restituisce un nuovo socket descriptor
su cui si potrà operare per effettuare la comunicazione. Se non ci sono
connessioni completate il processo viene messo in attesa. Il prototipo della
funzione è il seguente:

\begin{funcproto}{
\fhead{sys/socket.h}
\fdecl{int accept(int sockfd, struct sockaddr *addr, socklen\_t *addrlen)} 
\fdesc{Accetta una connessione sul socket specificato.} 
}

{La funzione ritorna un numero di socket descriptor positivo in caso di
  successo e $-1$ per un errore, nel qual caso \var{errno} assumerà uno dei
  valori:
  \begin{errlist}
  \item[\errcode{EAGAIN} o \errcode{EWOULDBLOCK}] il socket è stato impostato
    come non bloccante (vedi sez.~\ref{sec:file_noblocking}), e non ci sono
    connessioni in attesa di essere accettate. In generale possono essere
    restituiti entrambi i valori, per cui se si ha a cuore la portabilità
    occorre controllare entrambi.
  \item[\errcode{EBADF}] l'argomento \param{sockfd} non è un file descriptor
    valido.
  \item[\errcode{ECONNABORTED}] la connessione è stata abortita.
  \item[\errcode{EINTR}] la funzione è stata interrotta da un segnale.
  \item[\errcode{EINVAL}] il socket non è in ascolto o \param{addrlen} non ha
    un valore valido. 
  \item[\errcode{ENOBUFS}, \errcode{ENOMEM}] questo spesso significa che
    l'allocazione della memoria è limitata dai limiti sui buffer dei socket,
    non dalla memoria di sistema.
  \item[\errcode{ENOTSOCK}] l'argomento \param{sockfd} non è un socket.
  \item[\errcode{EOPNOTSUPP}] il socket è di un tipo che non supporta questa
    operazione.
  \item[\errcode{EPERM}] le regole del firewall non consentono la connessione.
  \end{errlist}
  ed inoltre nel loro significato generico: \errval{EFAULT}, \errval{EMFILE},
  \errval{ENFILE}; infine a seconda del protocollo e del kernel possono essere
  restituiti errori di rete relativi al nuovo socket come: \errval{ENOSR},
  \errval{ESOCKTNOSUPPORT}, \errval{EPROTONOSUPPORT}, \errval{ETIMEDOUT},
  \errval{ERESTARTSYS}.}
\end{funcproto}

La funzione estrae la prima connessione relativa al socket \param{sockfd} in
attesa sulla coda delle connessioni complete, che associa ad nuovo socket con
le stesse caratteristiche di \param{sockfd}.  Il socket originale non viene
toccato e resta nello stato di \texttt{LISTEN}, mentre il nuovo socket viene
posto nello stato \texttt{ESTABLISHED}. 

I due argomenti \param{addr} e \param{addrlen} (si noti che quest'ultimo è un
\textit{valure-result argument} passato con un puntatore per riavere indietro
il valore) sono usati rispettivamente per ottenere l'indirizzo del client da
cui proviene la connessione e la lunghezza dello stesso; la dimensione dipende
da quale famiglia di indirizzi si sta utilizzando.  Prima della
chiamata \param{addrlen} deve essere inizializzato alle dimensioni della
struttura degli indirizzi cui punta \param{addr} (un numero positivo); al
ritorno della funzione \param{addrlen} conterrà il numero di byte scritti
dentro \param{addr}. Se questa informazione non interessa basterà
inizializzare a \val{NULL} detti puntatori.

Se la funzione ha successo restituisce il descrittore di un nuovo socket
creato dal kernel (detto \textit{connected socket}) a cui viene associata la
prima connessione completa (estratta dalla relativa coda, vedi
sez.~\ref{sec:TCP_func_listen}) che il client ha effettuato verso il socket
\param{sockfd}. Quest'ultimo (detto \textit{listening socket}) è quello creato
all'inizio e messo in ascolto con \func{listen}, e non viene toccato dalla
funzione.  Se non ci sono connessioni pendenti da accettare la funzione mette
in attesa il processo\footnote{a meno che non si sia impostato il socket per
  essere non bloccante (vedi sez.~\ref{sec:file_noblocking}), nel qual caso
  ritorna con l'errore \errcode{EAGAIN}.  Torneremo su questa modalità di
  operazione in sez.~\ref{sec:TCP_sock_multiplexing}.}  fintanto che non ne
arriva una.

La funzione può essere usata solo con socket che supportino la connessione
(cioè di tipo \const{SOCK\_STREAM}, \const{SOCK\_SEQPACKET} o
\const{SOCK\_RDM}). Per alcuni protocolli che richiedono una conferma
esplicita della connessione,\footnote{attualmente in Linux solo DECnet ha
  questo comportamento.} la funzione opera solo l'estrazione dalla coda delle
connessioni, la conferma della connessione viene eseguita implicitamente dalla
prima chiamata ad una \func{read} o una \func{write}, mentre il rifiuto della
connessione viene eseguito con la funzione \func{close}. 

Si tenga presente che con Linux, seguendo POSIX.1, è sufficiente includere
\texttt{sys/socket.h}, ma alcune implementazioni di altri sistemi possono
richiedere l'inclusione di \texttt{sys/types.h}, per cui dovendo curare la
portabilità può essere il caso di includere anche questo file.  Inoltre Linux
presenta un comportamento diverso nella gestione degli errori rispetto ad
altre realizzazioni dei socket BSD, infatti la funzione \func{accept} passa
gli errori di rete pendenti sul nuovo socket come codici di errore per
\func{accept}, per cui l'applicazione deve tenerne conto ed eventualmente
ripetere la chiamata alla funzione come per l'errore di \errcode{EAGAIN}
(torneremo su questo in sez.~\ref{sec:TCP_echo_critical}).

Un'altra differenza con BSD è che la funzione non fa ereditare al nuovo socket
i flag del socket originale, come \const{O\_NONBLOCK},\footnote{ed in generale
  tutti quelli che si possono impostare con \func{fcntl}, vedi
  sez.~\ref{sec:file_fcntl_ioctl}.} che devono essere rispecificati ogni
volta. Tutto questo deve essere tenuto in conto se si devono scrivere
programmi portabili. Per poter effettuare questa impostazione in maniera
atomica, senza dover ricorrere ad ulteriori chiamate a \func{fcntl} su Linux è
disponibile anche la funzione \funcd{accept4}, il cui prototipo è:\footnote{la
  funzione è utilizzabile solo se si è definito la macro \macro{\_GNU\_SOURCE}
  ed ovviamente non è portabile.}

\begin{funcproto}{
\fhead{sys/socket.h}
\fdecl{int accept4(int sockfd, struct sockaddr *addr, socklen\_t *addrlen, int
  flags)} 
\fdesc{Accetta una connessione sul socket specificato.} 
}

{La funzione ritorna $0$ in caso di successo e $-1$ per un errore, nel qual
  caso \var{errno} assumerà gli stessi valori di \func{accept}.}
\end{funcproto}

La funzione aggiunge un quarto argomento \param{flags} usato come maschera
binaria, e se questo è nullo il suo comportamento è identico a quello di
\func{accept}. Con \param{flags} si possono impostare contestualmente
all'esecuzione sul file descriptor restituito i due flag di
\const{O\_NONBLOCK} e \const{O\_CLOEXEC}, fornendo un valore che sia un OR
aritmetico delle costanti in tab.\ref{tab:accept4_flags_arg}.

\begin{table}[htb]
  \centering
  \footnotesize
  \begin{tabular}[c]{|l|l|}
    \hline
    \textbf{Costante} & \textbf{Significato} \\
    \hline
    \hline
    \constd{SOCK\_NONBLOCK} & imposta sul file descriptor restituito il flag
                              di \const{O\_NONBLOCK}\\
    \constd{SOCK\_NOXEC}    & imposta sul file descriptor restituito il flag
                              di \const{O\_CLOEXEC}\\ 
    \hline    
  \end{tabular}
  \caption{Costanti per i possibili valori dell'argomento \param{flags} di
    \func{accept4}.}
  \label{tab:accept4_flags_arg}
\end{table}

Il meccanismo di funzionamento di \func{accept} è essenziale per capire il
funzionamento di un server: in generale infatti c'è sempre un solo socket in
ascolto, detto per questo \textit{listening socket}, che resta per tutto il
tempo nello stato \texttt{LISTEN}, mentre le connessioni vengono gestite dai
nuovi socket, detti \textit{connected socket}, ritornati da \func{accept}, che
si trovano automaticamente nello stato \texttt{ESTABLISHED}, e vengono
utilizzati per lo scambio dei dati, che avviene su di essi, fino alla chiusura
della connessione.  Si può riconoscere questo schema anche nell'esempio
elementare di fig.~\ref{fig:TCP_daytime_iter_server_code}, dove per ogni
connessione il socket creato da \func{accept} viene chiuso dopo l'invio dei
dati.


\subsection{Le funzioni \func{getsockname} e \func{getpeername}}
\label{sec:TCP_get_names}

Oltre a tutte quelle viste finora, dedicate all'utilizzo dei socket, esistono
alcune funzioni ausiliarie che possono essere usate per recuperare alcune
informazioni relative ai socket ed alle connessioni ad essi associate. Le due
più elementari sono le seguenti, usate per ottenere i dati relativi alla
\textit{socket pair} associata ad un certo socket.  La prima è
\funcd{getsockname} e serve ad ottenere l'indirizzo locale associato ad un
socket; il suo prototipo è:

\begin{funcproto}{
\fhead{sys/socket.h}
\fdecl{int getsockname(int sockfd, struct sockaddr *name, socklen\_t *namelen)}
\fdesc{Legge l'indirizzo locale di un socket.} 
}

{La funzione ritorna $0$ in caso di successo e $-1$ per un errore, nel qual
  caso \var{errno} assumerà uno dei valori: 
  \begin{errlist}
  \item[\errcode{EBADF}] l'argomento \param{sockfd} non è un file descriptor
    valido.
  \item[\errcode{ENOBUFS}] non ci sono risorse sufficienti nel sistema per
  \item[\errcode{ENOTSOCK}] l'argomento \param{sockfd} non è un socket.
    eseguire l'operazione.
  \end{errlist}
  ed \errval{EFAULT} nel suo significato generico.}
\end{funcproto}

La funzione restituisce la struttura degli indirizzi del socket \param{sockfd}
nella struttura indicata dal puntatore \param{name} la cui lunghezza è
specificata tramite l'argomento \param{namlen}. Quest'ultimo è un
\textit{value result argument} e pertanto viene passato come indirizzo per
avere indietro anche il numero di byte effettivamente scritti nella struttura
puntata da \param{name}. Si tenga presente che se si è utilizzato un buffer
troppo piccolo per \param{name} l'indirizzo risulterà troncato.

La funzione si usa tutte le volte che si vuole avere l'indirizzo locale di un
socket; ad esempio può essere usata da un client (che usualmente non chiama
\func{bind}) per ottenere l'indirizzo IP e la porta locale associati al socket
restituito da una \func{connect}, o da un server che ha chiamato \func{bind}
su un socket usando $0$ come porta locale per ottenere il numero di porta
effimera assegnato dal kernel.

Inoltre quando un server esegue una \func{bind} su un indirizzo generico, se
chiamata dopo il completamento di una connessione sul socket restituito da
\func{accept}, restituisce l'indirizzo locale che il kernel ha assegnato a
quella connessione. Invece tutte le volte che si vuole avere l'indirizzo
remoto di un socket si usa la funzione \funcd{getpeername}, il cui prototipo
è:

\begin{funcproto}{
\fhead{sys/socket.h}
\fdecl{int getpeername(int sockfd, struct sockaddr * name, socklen\_t *
  namelen)} 
\fdesc{Legge l'indirizzo remoto di un socket.} 
}

{La funzione ritorna $0$ in caso di successo e $-1$ per un errore, nel qual
  caso \var{errno} assumerà uno dei valori: 
  \begin{errlist}
  \item[\errcode{EBADF}] l'argomento \param{sockfd} non è un file descriptor
    valido.
  \item[\errcode{ENOBUFS}] non ci sono risorse sufficienti nel sistema per
    eseguire l'operazione.
  \item[\errcode{ENOTCONN}] il socket non è connesso.
  \item[\errcode{ENOTSOCK}] l'argomento \param{sockfd} non è un socket.
  \end{errlist}
  ed \errval{EFAULT} nel suo significato generico.}
\end{funcproto}


La funzione è identica a \func{getsockname}, ed usa la stessa sintassi, ma
restituisce l'indirizzo remoto del socket, cioè quello associato all'altro
capo della connessione.  Ci si può chiedere a cosa serva questa funzione dato
che dal lato client l'indirizzo remoto è sempre noto quando si esegue la
\func{connect} mentre dal lato server si possono usare, come vedremo in
fig.~\ref{fig:TCP_daytime_cunc_server_code}, i valori di ritorno di
\func{accept}.

Il fatto è che in generale quest'ultimo caso non è sempre possibile.  In
particolare questo avviene quando il server, invece di gestire la connessione
direttamente in un processo figlio, come vedremo nell'esempio di server
concorrente di sez.~\ref{sec:TCP_daytime_cunc_server}, lancia per ciascuna
connessione un altro programma, usando \func{exec}. Questa ad esempio è la
modalità con cui opera il \textsl{super-server} \cmd{xinetd}, che può gestire
tutta una serie di servizi diversi, eseguendo su ogni connessione ricevuta
sulle porte tenute sotto controllo, il relativo server.

In questo caso benché il processo figlio abbia una immagine della memoria che
è copia di quella del processo padre (e contiene quindi anche la struttura
ritornata da \func{accept}), all'esecuzione di \func{exec} verrà caricata in
memoria l'immagine del programma eseguito, che a questo punto perde ogni
riferimento ai valori tornati da \func{accept}.  Il socket descriptor però
resta aperto, e se si è seguita una opportuna convenzione per rendere noto al
programma eseguito qual è il socket connesso (ad esempio il solito
\cmd{xinetd} fa sempre in modo che i file descriptor 0, 1 e 2 corrispondano al
socket connesso) quest'ultimo potrà usare la funzione \func{getpeername} per
determinare l'indirizzo remoto del client.

Infine è da chiarire (si legga la pagina di manuale) che, come per
\func{accept}, il terzo argomento, che è specificato dallo standard POSIX.1g
come di tipo \code{socklen\_t *} in realtà deve sempre corrispondere ad un
\ctyp{int *} come prima dello standard, perché tutte le realizzazioni dei
socket BSD fanno questa assunzione.


\subsection{La funzione \func{close}}
\label{sec:TCP_func_close}

La funzione standard Unix \func{close} (vedi sez.~\ref{sec:file_open_close})
che si usa sui file può essere usata con lo stesso effetto anche sui file
descriptor associati ad un socket.

L'azione di questa funzione quando applicata a socket è di marcarlo come
chiuso e ritornare immediatamente al processo. Una volta chiamata il socket
descriptor non è più utilizzabile dal processo e non può essere usato come
argomento per una \func{write} o una \func{read} (anche se l'altro capo della
connessione non avesse chiuso la sua parte).  Il kernel invierà comunque tutti
i dati che ha in coda prima di iniziare la sequenza di chiusura.  Vedremo più
avanti in sez.~\ref{sec:sock_generic_options} come sia possibile cambiare
questo comportamento, e cosa può essere fatto perché il processo possa
assicurarsi che l'altro capo abbia ricevuto tutti i dati.

Come per tutti i file descriptor anche per i socket viene mantenuto un numero
di riferimenti, per cui se più di un processo ha lo stesso socket aperto
l'emissione del FIN e la sequenza di chiusura di TCP non viene innescata
fintanto che il numero di riferimenti non si annulla, questo si applica, come
visto in sez.~\ref{sec:file_shared_access}, sia ai file descriptor duplicati
che a quelli ereditati dagli eventuali processi figli, ed è il comportamento
che ci si aspetta in una qualunque applicazione client/server.

Per attivare immediatamente l'emissione del FIN e la sequenza di chiusura
descritta in sez.~\ref{sec:TCP_conn_term}, si può invece usare la funzione
\func{shutdown} su cui torneremo in seguito (vedi
sez.~\ref{sec:TCP_shutdown}).


\section{Un esempio elementare: il servizio \textit{daytime}}
\label{sec:TCP_daytime_application}

Avendo introdotto le funzioni di base per la gestione dei socket, potremo
vedere in questa sezione un primo esempio di applicazione elementare che
realizza il servizio \textit{daytime} su TCP, secondo quanto specificato
dall'\href{http://www.ietf.org/rfc/rfc867.txt}{RFC~867}.  Prima di passare
agli esempi del client e del server, inizieremo riesaminando con maggiori
dettagli una peculiarità delle funzioni di I/O, già accennata in
sez.~\ref{sec:file_read} e sez.~\ref{sec:file_write}, che nel caso dei socket
è particolarmente rilevante.  Passeremo poi ad illustrare gli esempi della
realizzazione, sia dal lato client, che dal lato server, che si è effettuata
sia in forma iterativa che concorrente.


\subsection{Il comportamento delle funzioni di I/O}
\label{sec:sock_io_behav}

Una cosa che si tende a dimenticare quando si ha a che fare con i socket è che
le funzioni di input/output non sempre hanno lo stesso comportamento che
avrebbero con i normali file di dati (in particolare questo accade per i
socket di tipo stream).

Infatti con i socket è comune che funzioni come \func{read} o \func{write}
possano restituire in input o scrivere in output un numero di byte minore di
quello richiesto. Come già accennato in sez.~\ref{sec:file_read} questo è un
comportamento normale per le funzioni di I/O, ma con i normali file di dati il
problema si avverte solo in lettura, quando si incontra la fine del file. In
generale non è così, e con i socket questo è particolarmente evidente.


\begin{figure}[!htbp]
  \footnotesize \centering
  \begin{minipage}[c]{\codesamplewidth}
    \includecodesample{listati/FullRead.c}
  \end{minipage} 
  \normalsize
  \caption{La funzione \func{FullRead}, che legge esattamente \var{count} byte
    da un file descriptor, iterando opportunamente le letture.}
  \label{fig:sock_FullRead_code}
\end{figure}

Quando ci si trova ad affrontare questo comportamento tutto quello che si deve
fare è semplicemente ripetere la lettura (o la scrittura) per la quantità di
byte restanti, tenendo conto che le funzioni si possono bloccare se i dati non
sono disponibili: è lo stesso comportamento che si può avere scrivendo più di
\const{PIPE\_BUF} byte in una \textit{pipe} (si riveda quanto detto in
sez.~\ref{sec:ipc_pipes}).

Per questo motivo, seguendo l'esempio di R. W. Stevens in \cite{UNP1}, si sono
definite due funzioni, \func{FullRead} e \func{FullWrite}, che eseguono
lettura e scrittura tenendo conto di questa caratteristica, ed in grado di
ritornare solo dopo avere letto o scritto esattamente il numero di byte
specificato; il sorgente è riportato rispettivamente in
fig.~\ref{fig:sock_FullRead_code} e fig.~\ref{fig:sock_FullWrite_code} ed è
disponibile fra i sorgenti allegati alla guida nei file \file{FullRead.c} e
\file{FullWrite.c}.

\begin{figure}[!htbp]
  \centering
  \footnotesize \centering
  \begin{minipage}[c]{\codesamplewidth}
    \includecodesample{listati/FullWrite.c}
  \end{minipage} 
  \normalsize
  \caption{La funzione \func{FullWrite}, che scrive esattamente \var{count}
    byte su un file descriptor, iterando opportunamente le scritture.}
  \label{fig:sock_FullWrite_code}
\end{figure}

Come si può notare le due funzioni ripetono la lettura/scrittura in un ciclo
fino all'esaurimento del numero di byte richiesti, in caso di errore viene
controllato se questo è \errcode{EINTR} (cioè un'interruzione della
\textit{system call} dovuta ad un segnale), nel qual caso l'accesso viene
ripetuto, altrimenti l'errore viene ritornato al programma chiamante,
interrompendo il ciclo.

Nel caso della lettura, se il numero di byte letti è zero, significa che si è
arrivati alla fine del file (per i socket questo significa in genere che
l'altro capo è stato chiuso, e quindi non sarà più possibile leggere niente) e
pertanto si ritorna senza aver concluso la lettura di tutti i byte
richiesti. Entrambe le funzioni restituiscono 0 in caso di successo, ed un
valore negativo in caso di errore, \func{FullRead} restituisce il numero di
byte non letti in caso di \textit{end-of-file} prematuro.


\subsection{Il client \textit{daytime}}
\label{sec:TCP_daytime_client}

Il primo esempio di applicazione delle funzioni di base illustrate in
sez.~\ref{sec:TCP_functions} è relativo alla creazione di un client elementare
per il servizio \textit{daytime}, un servizio elementare, definito
nell'\href{http://www.ietf.org/rfc/rfc867.txt}{RFC~867}, che restituisce
l'ora locale della macchina a cui si effettua la richiesta, e che è assegnato
alla porta 13.

In fig.~\ref{fig:TCP_daytime_client_code} è riportata la sezione principale
del codice del nostro client. Il sorgente completo del programma
(\texttt{TCP\_daytime.c}, che comprende il trattamento delle opzioni ed una
funzione per stampare un messaggio di aiuto) è allegato alla guida nella
sezione dei codici sorgente e può essere compilato su una qualunque macchina
GNU/Linux.

\begin{figure}[!htbp]
  \footnotesize \centering
  \begin{minipage}[c]{\codesamplewidth}
    \includecodesample{listati/TCP_daytime.c}
  \end{minipage} 
  \normalsize
  \caption{Esempio di codice di un client elementare per il servizio
    \textit{daytime}.} 
  \label{fig:TCP_daytime_client_code}
\end{figure}

Il programma anzitutto (\texttt{\small 1--5}) include gli header necessari;
dopo la dichiarazione delle variabili (\texttt{\small 9--12}) si è omessa
tutta la parte relativa al trattamento degli argomenti passati dalla linea di
comando (effettuata con le apposite funzioni illustrate in
sez.~\ref{sec:proc_opt_handling}).

Il primo passo (\texttt{\small 14--18}) è creare un socket TCP (quindi di tipo
\const{SOCK\_STREAM} e di famiglia \const{AF\_INET}). La funzione
\func{socket} ritorna il descrittore che viene usato per identificare il
socket in tutte le chiamate successive. Nel caso la chiamata fallisca si
stampa un errore (\texttt{\small 16}) con la funzione \func{perror} e si esce
(\texttt{\small 17}) con un codice di errore.

Il passo seguente (\texttt{\small 19--27}) è quello di costruire un'apposita
struttura \struct{sockaddr\_in} in cui sarà inserito l'indirizzo del server ed
il numero della porta del servizio. Il primo passo (\texttt{\small 20}) è
inizializzare tutto a zero, per poi inserire il tipo di indirizzo
(\texttt{\small 21}) e la porta (\texttt{\small 22}), usando per quest'ultima
la funzione \func{htons} per convertire il formato dell'intero usato dal
computer a quello usato nella rete, infine (\texttt{\small 23--27}) si può
utilizzare la funzione \func{inet\_pton} per convertire l'indirizzo numerico
passato dalla linea di comando.

A questo punto (\texttt{\small 28--32}) usando la funzione \func{connect} sul
socket creato in precedenza (\texttt{\small 29}) si può stabilire la
connessione con il server. Per questo si deve utilizzare come secondo
argomento la struttura preparata in precedenza con il relativo indirizzo; si
noti come, esistendo diversi tipi di socket, si sia dovuto effettuare un cast.
Un valore di ritorno della funzione negativo implica il fallimento della
connessione, nel qual caso si stampa un errore (\texttt{\small 30}) e si
ritorna (\texttt{\small 31}).

Completata con successo la connessione il passo successivo (\texttt{\small
  34--40}) è leggere la data dal socket; il protocollo prevede che il server
invii sempre una stringa alfanumerica, il formato della stringa non è
specificato dallo standard, per cui noi useremo il formato usato dalla
funzione \func{ctime}, seguito dai caratteri di terminazione \verb|\r\n|, cioè
qualcosa del tipo:
\begin{Verbatim}
Wed Apr 4 00:53:00 2001\r\n
\end{Verbatim}
questa viene letta dal socket (\texttt{\small 34}) con la funzione \func{read}
in un buffer temporaneo; la stringa poi deve essere terminata (\texttt{\small
  35}) con il solito carattere nullo per poter essere stampata (\texttt{\small
  36}) sullo \textit{standard output} con l'uso di \func{fputs}.

Come si è già spiegato in sez.~\ref{sec:sock_io_behav} la risposta dal socket
potrà arrivare in un unico pacchetto di 26 byte (come avverrà senz'altro nel
caso in questione) ma potrebbe anche arrivare in 26 pacchetti di un byte.  Per
questo nel caso generale non si può mai assumere che tutti i dati arrivino con
una singola lettura, pertanto quest'ultima deve essere effettuata in un ciclo
in cui si continui a leggere fintanto che la funzione \func{read} non ritorni
uno zero (che significa che l'altro capo ha chiuso la connessione) o un numero
minore di zero (che significa un errore nella connessione).

Si noti come in questo caso la fine dei dati sia specificata dal server che
chiude la connessione (anche questo è quanto richiesto dal protocollo); questa
è una delle tecniche possibili (è quella usata pure dal protocollo HTTP), ma
ce ne possono essere altre, ad esempio FTP marca la conclusione di un blocco
di dati con la sequenza ASCII \verb|\r\n| (carriage return e line feed),
mentre il DNS mette la lunghezza in testa ad ogni blocco che trasmette. Il
punto essenziale è che TCP non provvede nessuna indicazione che permetta di
marcare dei blocchi di dati, per cui se questo è necessario deve provvedere il
programma stesso.

Se abilitiamo il servizio \textit{daytime}\footnote{in genere questo viene
  fornito direttamente dal \textsl{superdemone} \cmd{xinetd}, pertanto basta
  assicurarsi che esso sia abilitato nel relativo file di configurazione.}
possiamo verificare il funzionamento del nostro client, avremo allora:
\begin{Console}
[piccardi@gont sources]$ \textbf{./daytime 127.0.0.1}
Mon Apr 21 20:46:11 2003
\end{Console}
%$
e come si vede tutto funziona regolarmente.


\subsection{Un server \textit{daytime} iterativo}
\label{sec:TCP_daytime_iter_server}

Dopo aver illustrato il client daremo anche un esempio di un server
elementare, che sia anche in grado di rispondere al precedente client. Come
primo esempio realizzeremo un server iterativo, in grado di fornire una sola
risposta alla volta. Il codice del programma è nuovamente mostrato in
fig.~\ref{fig:TCP_daytime_iter_server_code}, il sorgente completo
(\texttt{TCP\_iter\_daytimed.c}) è allegato insieme agli altri file degli
esempi.

\begin{figure}[!htbp]
  \footnotesize \centering
  \begin{minipage}[c]{\codesamplewidth}
    \includecodesample{listati/TCP_iter_daytimed.c}
  \end{minipage} 
  \normalsize
  \caption{Esempio di codice di un semplice server per il servizio daytime.}
  \label{fig:TCP_daytime_iter_server_code}
\end{figure}

Come per il client si includono (\texttt{\small 1--9}) gli header necessari a
cui è aggiunto quello per trattare i tempi, e si definiscono (\texttt{\small
  14--18}) alcune costanti e le variabili necessarie in seguito. Come nel caso
precedente si sono omesse le parti relative al trattamento delle opzioni da
riga di comando.

La creazione del socket (\texttt{\small 20--24}) è analoga al caso precedente,
come pure l'inizializzazione (\texttt{\small 25--29}) della struttura
\struct{sockaddr\_in}.  Anche in questo caso (\texttt{\small 28}) si usa la
porta standard del servizio \textit{daytime}, ma come indirizzo IP si usa
(\texttt{\small 29}) il valore predefinito \const{INET\_ANY}, che corrisponde
all'indirizzo generico.

Si effettua poi (\texttt{\small 30--34}) la chiamata alla funzione \func{bind}
che permette di associare la precedente struttura al socket, in modo che
quest'ultimo possa essere usato per accettare connessioni su una qualunque
delle interfacce di rete locali. In caso di errore si stampa (\texttt{\small
  31}) un messaggio, e si termina (\texttt{\small 32}) immediatamente il
programma.

Il passo successivo (\texttt{\small 35--39}) è quello di mettere ``\textsl{in
  ascolto}'' il socket; questo viene fatto (\texttt{\small 36}) con la
funzione \func{listen} che dice al kernel di accettare connessioni per il
socket che abbiamo creato; la funzione indica inoltre, con il secondo
argomento, il numero massimo di connessioni che il kernel accetterà di mettere
in coda per il suddetto socket. Di nuovo in caso di errore si stampa
(\texttt{\small 37}) un messaggio, e si esce (\texttt{\small 38})
immediatamente.

La chiamata a \func{listen} completa la preparazione del socket per l'ascolto
(che viene chiamato anche \textit{listening descriptor}) a questo punto si può
procedere con il ciclo principale (\texttt{\small 40--53}) che viene eseguito
indefinitamente. Il primo passo (\texttt{\small 42}) è porsi in attesa di
connessioni con la chiamata alla funzione \func{accept}, come in precedenza in
caso di errore si stampa (\texttt{\small 43}) un messaggio, e si esce
(\texttt{\small 44}).

Il processo resterà in stato di \textit{sleep} fin quando non arriva e viene
accettata una connessione da un client; quando questo avviene \func{accept}
ritorna, restituendo un secondo descrittore, che viene chiamato
\textit{connected descriptor}, e che è quello che verrà usato dalla successiva
chiamata alla \func{write} per scrivere la risposta al client.

Il ciclo quindi proseguirà determinando (\texttt{\small 46}) il tempo corrente
con una chiamata a \texttt{time}, con il quale si potrà opportunamente
costruire (\texttt{\small 47}) la stringa con la data da trasmettere
(\texttt{\small 48}) con la chiamata a \func{write}. Completata la
trasmissione il nuovo socket viene chiuso (\texttt{\small 52}).  A questo
punto il ciclo si chiude ricominciando da capo in modo da poter ripetere
l'invio della data in risposta ad una successiva connessione.

È importante notare che questo server è estremamente elementare, infatti, a
parte il fatto di poter essere usato solo con indirizzi IPv4, esso è in grado
di rispondere ad un solo un client alla volta: è cioè, come dicevamo, un
\textsl{server iterativo}. Inoltre è scritto per essere lanciato da linea di
comando, se lo si volesse utilizzare come demone occorrerebbero le opportune
modifiche (come una chiamata a \func{daemon} prima dell'inizio del ciclo
principale) per tener conto di quanto illustrato in
sez.~\ref{sec:sess_daemon}. Si noti anche che non si è inserita nessuna forma
di gestione della terminazione del processo, dato che tutti i file descriptor
vengono chiusi automaticamente alla sua uscita, e che, non generando figli,
non è necessario preoccuparsi di gestire la loro terminazione.


\subsection{Un server \textit{daytime} concorrente}
\label{sec:TCP_daytime_cunc_server}

Il server \texttt{daytime} dell'esempio in
sez.~\ref{sec:TCP_daytime_iter_server} è un tipico esempio di server iterativo,
in cui viene servita una richiesta alla volta; in generale però, specie se il
servizio è più complesso e comporta uno scambio di dati più sostanzioso di
quello in questione, non è opportuno bloccare un server nel servizio di un
client per volta; per questo si ricorre alle capacità di multitasking del
sistema.

Come spiegato in sez.~\ref{sec:proc_fork} una delle modalità più comuni di
funzionamento da parte dei server è quella di usare la funzione \func{fork}
per creare, ad ogni richiesta da parte di un client, un processo figlio che si
incarichi della gestione della comunicazione.  Si è allora riscritto il server
\textit{daytime} dell'esempio precedente in forma concorrente, inserendo anche
una opzione per la stampa degli indirizzi delle connessioni ricevute.

In fig.~\ref{fig:TCP_daytime_cunc_server_code} è mostrato un estratto del
codice, in cui si sono tralasciati il trattamento delle opzioni e le parti
rimaste invariate rispetto al precedente esempio (cioè tutta la parte
riguardante l'apertura passiva del socket). Al solito il sorgente completo del
server, nel file \texttt{TCP\_cunc\_daytimed.c}, è allegato insieme ai
sorgenti degli altri esempi.

\begin{figure}[!htbp]
  \footnotesize \centering
  \begin{minipage}[c]{\codesamplewidth}
    \includecodesample{listati/TCP_cunc_daytimed.c}
  \end{minipage} 
  \normalsize
  \caption{Esempio di codice di un server concorrente elementare per il 
    servizio daytime.}
  \label{fig:TCP_daytime_cunc_server_code}
\end{figure}

Stavolta (\texttt{\small 21--26}) la funzione \func{accept} è chiamata
fornendo una struttura di indirizzi in cui saranno ritornati l'indirizzo IP e
la porta da cui il client effettua la connessione, che in un secondo tempo,
(\texttt{\small 40--44}), se il logging è abilitato, stamperemo sullo
\textit{standard output}.

Quando \func{accept} ritorna il server chiama la funzione \func{fork}
(\texttt{\small 27--31}) per creare il processo figlio che effettuerà
(\texttt{\small 32--46}) tutte le operazioni relative a quella connessione,
mentre il padre proseguirà l'esecuzione del ciclo principale in attesa di
ulteriori connessioni.

Si noti come il figlio operi solo sul socket connesso, chiudendo
immediatamente (\texttt{\small 33}) il socket \var{list\_fd}; mentre il padre
continua ad operare solo sul socket in ascolto chiudendo (\texttt{\small 48})
\var{conn\_fd} al ritorno dalla \func{fork}. Per quanto abbiamo detto in
sez.~\ref{sec:TCP_func_close} nessuna delle due chiamate a \func{close} causa
l'innesco della sequenza di chiusura perché il numero di riferimenti al file
descriptor non si è annullato.

Infatti subito dopo la creazione del socket \var{list\_fd} ha una referenza, e
lo stesso vale per \var{conn\_fd} dopo il ritorno di \func{accept}, ma dopo la
\func{fork} i descrittori vengono duplicati nel padre e nel figlio per cui
entrambi i socket si trovano con due referenze. Questo fa sì che quando il
padre chiude \var{sock\_fd} esso resta con una referenza da parte del figlio,
e sarà definitivamente chiuso solo quando quest'ultimo, dopo aver completato
le sue operazioni, chiamerà (\texttt{\small 45}) la funzione \func{close}.

In realtà per il figlio non sarebbe necessaria nessuna chiamata a
\func{close}, in quanto con la \func{exit} finale (\texttt{\small 45}) tutti i
file descriptor, quindi anche quelli associati ai socket, vengono
automaticamente chiusi.  Tuttavia si è preferito effettuare esplicitamente le
chiusure per avere una maggiore chiarezza del codice, e per evitare eventuali
errori, prevenendo ad esempio un uso involontario del \textit{listening
  descriptor}.

Si noti invece come sia essenziale che il padre chiuda ogni volta il socket
connesso dopo la \func{fork}; se così non fosse nessuno di questi socket
sarebbe effettivamente chiuso dato che alla chiusura da parte del figlio
resterebbe ancora un riferimento nel padre. Si avrebbero così due effetti: il
padre potrebbe esaurire i descrittori disponibili (che sono un numero limitato
per ogni processo) e soprattutto nessuna delle connessioni con i client
verrebbe chiusa.

Come per ogni server iterativo il lavoro di risposta viene eseguito
interamente dal processo figlio. Questo si incarica (\texttt{\small 34}) di
chiamare \func{time} per leggere il tempo corrente, e di stamparlo
(\texttt{\small 35}) sulla stringa contenuta in \var{buffer} con l'uso di
\func{snprintf} e \func{ctime}. Poi la stringa viene scritta (\texttt{\small
  36--39}) sul socket, controllando che non ci siano errori. Anche in questo
caso si è evitato il ricorso a \func{FullWrite} in quanto la stringa è
estremamente breve e verrà senz'altro scritta in un singolo segmento.

Inoltre nel caso sia stato abilitato il \textit{logging} delle connessioni, si
provvede anche (\texttt{\small 40--43}) a stampare sullo \textit{standard
  output} l'indirizzo e la porta da cui il client ha effettuato la
connessione, usando i valori contenuti nelle strutture restituite da
\func{accept}, eseguendo le opportune conversioni con \func{inet\_ntop} e
\func{ntohs}.

Ancora una volta l'esempio è estremamente semplificato, si noti come di nuovo
non si sia gestita né la terminazione del processo né il suo uso come demone,
che tra l'altro sarebbe stato incompatibile con l'uso della opzione di logging
che stampa gli indirizzi delle connessioni sullo standard output. Un altro
aspetto tralasciato è la gestione della terminazione dei processi figli,
torneremo su questo più avanti quando tratteremo alcuni esempi di server più
complessi.


\section{Un esempio più completo: il servizio \textit{echo}}
\label{sec:TCP_echo_application}

L'esempio precedente, basato sul servizio \textit{daytime}, è un esempio molto
elementare, in cui il flusso dei dati va solo nella direzione dal server al
client. In questa sezione esamineremo un esempio di applicazione client/server
un po' più complessa, che usi i socket TCP per una comunicazione in entrambe
le direzioni.

Ci limiteremo a fornire una realizzazione elementare, che usi solo le funzioni
di base viste finora, ma prenderemo in esame, oltre al comportamento in
condizioni normali, anche tutti i possibili scenari particolari (errori,
sconnessione della rete, crash del client o del server durante la connessione)
che possono avere luogo durante l'impiego di un'applicazione di rete, partendo
da una versione primitiva che dovrà essere rimaneggiata di volta in volta per
poter tenere conto di tutte le evenienze che si possono manifestare nella vita
reale di un'applicazione di rete, fino ad arrivare ad una realizzazione
completa.


\subsection{Il servizio \textit{echo}}
\label{sec:TCP_echo}

Nella ricerca di un servizio che potesse fare da esempio per una comunicazione
bidirezionale, si è deciso, seguendo la scelta di Stevens in \cite{UNP1}, di
usare il servizio \textit{echo}, che si limita a restituire in uscita quanto
immesso in ingresso. Infatti, nonostante la sua estrema semplicità, questo
servizio costituisce il prototipo ideale per una generica applicazione di rete
in cui un server risponde alle richieste di un client.  Nel caso di una
applicazione più complessa quello che si potrà avere in più è una elaborazione
dell'input del client, che in molti casi viene interpretato come un comando,
da parte di un server che risponde fornendo altri dati in uscita.

Il servizio \textit{echo} è uno dei servizi standard solitamente provvisti
direttamente dal superserver \cmd{inetd}, ed è definito
dall'\href{http://www.ietf.org/rfc/rfc862.txt}{RFC~862}. Come dice il nome il
servizio deve riscrivere indietro sul socket i dati che gli vengono inviati in
ingresso. L'RFC descrive le specifiche del servizio sia per TCP che UDP, e per
qil primo stabilisce che una volta stabilita la connessione ogni dato in
ingresso deve essere rimandato in uscita fintanto che il chiamante non ha
chiude la connessione. Al servizio è assegnata la porta riservata 7.

Nel nostro caso l'esempio sarà costituito da un client che legge una linea di
caratteri dallo \textit{standard input} e la scrive sul server. A sua volta il
server leggerà la linea dalla connessione e la riscriverà immutata
all'indietro. Sarà compito del client leggere la risposta del server e
stamparla sullo \textit{standard output}.


\subsection{Il client \textit{echo}: prima versione}
\label{sec:TCP_echo_client}

Il codice della prima versione del client per il servizio \textit{echo} (file
\texttt{TCP\_echo\_first.c}) dei sorgenti allegati alla guida) è riportato in
fig.~\ref{fig:TCP_echo_client_1}. Esso ricalca la struttura del precedente
client per il servizio \textit{daytime} (vedi
sez.~\ref{sec:TCP_daytime_client}), e la prima parte (\texttt{\small 10--27})
è sostanzialmente identica, a parte l'uso di una porta diversa.

\begin{figure}[!htb]
  \footnotesize \centering
  \begin{minipage}[c]{\codesamplewidth}
    \includecodesample{listati/TCP_echo_first.c}
  \end{minipage} 
  \normalsize
  \caption{Codice della prima versione del client \textit{echo}.}
  \label{fig:TCP_echo_client_1}
\end{figure}

Al solito si è tralasciata la sezione relativa alla gestione delle opzioni a
riga di comando.  Una volta dichiarate le variabili, si prosegue
(\texttt{\small 10--13}) con la creazione del socket con l'usuale controllo
degli errori, alla preparazione (\texttt{\small 14--17}) della struttura
dell'indirizzo, che stavolta usa la porta 7 riservata al servizio
\textit{echo}, infine si converte (\texttt{\small 18--22}) l'indirizzo
specificato a riga di comando.  A questo punto (\texttt{\small 23--27}) si può
eseguire la connessione al server secondo la stessa modalità usata in
sez.~\ref{sec:TCP_daytime_client}.

Completata la connessione, per gestire il funzionamento del protocollo si usa
la funzione \code{ClientEcho}, il cui codice si è riportato a parte in
fig.~\ref{fig:TCP_client_echo_sub}. Questa si preoccupa di gestire tutta la
comunicazione, leggendo una riga alla volta dallo \textit{standard input}
\file{stdin}, scrivendola sul socket e ristampando su \file{stdout} quanto
ricevuto in risposta dal server. Al ritorno dalla funzione (\texttt{\small
  30--31}) anche il programma termina.

La funzione \code{ClientEcho} utilizza due buffer (\texttt{\small 3}) per
gestire i dati inviati e letti sul socket.  La comunicazione viene gestita
all'interno di un ciclo (\texttt{\small 5--10}), i dati da inviare sulla
connessione vengono presi dallo \file{stdin} usando la funzione \func{fgets},
che legge una linea di testo (terminata da un \texttt{CR} e fino al massimo di
\const{MAXLINE} caratteri) e la salva sul buffer di invio.

Si usa poi (\texttt{\small 6}) la funzione \func{FullWrite}, vista in
sez.~\ref{sec:sock_io_behav}, per scrivere i dati sul socket, gestendo
automaticamente l'invio multiplo qualora una singola \func{write} non sia
sufficiente.  I dati vengono riletti indietro (\texttt{\small 7}) con una
\func{read}\footnote{si è fatta l'assunzione implicita che i dati siano
  contenuti tutti in un solo segmento, così che la chiamata a \func{read} li
  restituisca sempre tutti; avendo scelto una dimensione ridotta per il buffer
  questo sarà sempre vero, vedremo più avanti come superare il problema di
  rileggere indietro tutti e soli i dati disponibili, senza bloccarsi.} sul
buffer di ricezione e viene inserita (\texttt{\small 8}) la terminazione della
stringa e per poter usare (\texttt{\small 9}) la funzione \func{fputs} per
scriverli su \file{stdout}.

\begin{figure}[!htbp]
  \footnotesize \centering
  \begin{minipage}[c]{\codesamplewidth}
    \includecodesample{listati/ClientEcho_first.c}
  \end{minipage} 
  \normalsize
  \caption{Codice della prima versione della funzione \texttt{ClientEcho} per 
    la gestione del servizio \textit{echo}.}
  \label{fig:TCP_client_echo_sub}
\end{figure}

Quando si concluderà l'invio di dati mandando un \textit{end-of-file} sullo
\textit{standard input} si avrà il ritorno di \func{fgets} con un puntatore
nullo (si riveda quanto spiegato in sez.~\ref{sec:file_line_io}) e la
conseguente uscita dal ciclo; al che la subroutine ritorna ed il nostro
programma client termina.

Si può effettuare una verifica del funzionamento del client abilitando il
servizio \textit{echo} nella configurazione di \cmd{xinetd} sulla propria
macchina ed usandolo direttamente verso di esso in locale, vedremo in
dettaglio più avanti (in sez.~\ref{sec:TCP_echo_startup}) il funzionamento del
programma, usato però con la nostra versione del server \textit{echo}, che
illustriamo immediatamente.


\subsection{Il server \textit{echo}: prima versione}
\label{sec:TCPsimp_server_main}

La prima versione del server, contenuta nel file \texttt{TCP\_echod\_first.c},
è riportata in fig.~\ref{fig:TCP_echo_server_first_code} e
fig.~\ref{fig:TCP_echo_server_first_main}. Come abbiamo fatto per il client
anche il server è stato diviso in un corpo principale, costituito dalla
funzione \code{main}, che è molto simile a quello visto nel precedente esempio
per il server del servizio \textit{daytime} di
sez.~\ref{sec:TCP_daytime_cunc_server}, e da una funzione ausiliaria
\code{ServEcho} che si cura della gestione del servizio.

\begin{figure}[!htb]
  \footnotesize \centering
  \begin{minipage}[c]{\codesamplewidth}
    \includecodesample{listati/TCP_echod_first_init.c}
  \end{minipage} 
  \normalsize
  \caption{Codice di inizializzatione della prima versione del server
    per il servizio \textit{echo}.}
  \label{fig:TCP_echo_server_first_code}
\end{figure}

In questo caso però, rispetto a quanto visto nell'esempio di
fig.~\ref{fig:TCP_daytime_cunc_server_code} si è preferito scrivere il server
curando maggiormente alcuni dettagli, per tenere conto anche di alcune
esigenze generali (che non riguardano direttamente la rete), come la
possibilità di lanciare il server anche in modalità interattiva e la cessione
dei privilegi di amministratore non appena questi non sono più necessari.

La sezione iniziale del programma (\texttt{\small 8--21}) è la stessa del
server di sez.~\ref{sec:TCP_daytime_cunc_server}, ed ivi descritta in
dettaglio: crea il socket, inizializza l'indirizzo e esegue \func{bind}; dato
che quest'ultima funzione viene usata su una porta riservata, il server dovrà
essere eseguito da un processo con i privilegi di amministratore, pena il
fallimento della chiamata.

Una volta eseguita la funzione \func{bind} però i privilegi di amministratore
non sono più necessari, per questo è sempre opportuno rilasciarli, in modo da
evitare problemi in caso di eventuali vulnerabilità del server.  Per questo
prima (\texttt{\small 22--26}) si esegue \func{setgid} per assegnare il
processo ad un gruppo senza privilegi,\footnote{si è usato il valore 65534,
  ovvero -1 per il formato \ctyp{short}, che di norma in tutte le
  distribuzioni viene usato per identificare il gruppo \texttt{nogroup} e
  l'utente \texttt{nobody}, usati appunto per eseguire programmi che non
  richiedono nessun privilegio particolare.} e poi si ripete (\texttt{\small
  27--30}) l'operazione usando \func{setuid} per cambiare anche
l'utente.\footnote{si tenga presente che l'ordine in cui si eseguono queste
  due operazioni è importante, infatti solo avendo i privilegi di
  amministratore si può cambiare il gruppo di un processo ad un altro di cui
  non si fa parte, per cui chiamare prima \func{setuid} farebbe fallire una
  successiva chiamata a \func{setgid}.  Inoltre si ricordi (si riveda quanto
  esposto in sez.~\ref{sec:proc_perms}) che usando queste due funzioni il
  rilascio dei privilegi è irreversibile.}  Infine (\texttt{\small 30--36}),
qualora sia impostata la variabile \var{demonize}, prima (\texttt{\small 31})
si apre il sistema di logging per la stampa degli errori, e poi
(\texttt{\small 32--35}) si invoca \func{daemon} per eseguire in background il
processo come demone.


\begin{figure}[!htb]
  \footnotesize \centering
  \begin{minipage}[c]{\codesamplewidth}
    \includecodesample{listati/TCP_echod_first_main.c}
  \end{minipage} 
  \normalsize
  \caption{Codice di inizializzatione della prima versione del server
    per il servizio \textit{echo}.}
  \label{fig:TCP_echo_server_first_main}
\end{figure}

A questo punto il programma prosegue nel ciclo principale, illustrato in
fig.~\ref{fig:TCP_echo_server_first_main}, usando lo schema già visto in
precedenza per server per il servizio \textit{daytime}, con l'unica differenza
della chiamata alla funzione \code{PrintErr}, riportata in
fig.~\ref{fig:TCP_PrintErr}, al posto di \func{perror} per la stampa degli
errori.

Si inizia con il porre (\texttt{\small 3--6}) in ascolto il socket, e poi si
esegue indefinitamente il ciclo principale (\texttt{\small 7--26}).
All'interno di questo si ricevono (\texttt{\small 9--12}) le connessioni,
creando (\texttt{\small 13--16}) un processo figlio per ciascuna di esse.
Quest'ultimo (\texttt{\small 17--21}), chiuso (\texttt{\small 18}) il
\textit{listening socket}, esegue (\texttt{\small 19}) la funzione di gestione
del servizio \code{ServEcho}, ed al ritorno di questa esce (\texttt{\small
  20}).

Il padre invece si limita (\texttt{\small 22}) a chiudere il \textit{connected
  socket} per ricominciare da capo il ciclo in attesa di nuove connessioni. In
questo modo si ha un server concorrente. La terminazione del padre non è
gestita esplicitamente, e deve essere effettuata inviando un segnale al
processo.

Avendo trattato direttamente la gestione del programma come demone, si è
dovuto anche provvedere alla necessità di poter stampare eventuali messaggi di
errore attraverso il sistema del \textit{syslog} trattato in
sez.~\ref{sec:sess_daemon}. Come accennato questo è stato fatto utilizzando
come \textit{wrapper} la funzione \code{PrintErr}, il cui codice è riportato
in fig.~\ref{fig:TCP_PrintErr}.

\begin{figure}[!htb]
  \footnotesize \centering
  \begin{minipage}[c]{\codesamplewidth}
    \includecodesample{listati/PrintErr.c}
  \end{minipage} 
  \normalsize
  \caption{Codice della funzione \code{PrintErr} per la generalizzazione della
    stampa degli errori sullo \textit{standard input} o attraverso il
    \texttt{syslog}.}
  \label{fig:TCP_PrintErr}
\end{figure}


In essa ci si limita a controllare (\texttt{\small 2}) se è stato impostato
(valore attivo per default) l'uso come demone, nel qual caso (\texttt{\small
  3}) si usa \func{syslog} (vedi sez.~\ref{sec:sess_daemon}) per stampare il
messaggio di errore fornito come argomento sui log di sistema. Se invece si è
in modalità interattiva (attivabile con l'opzione \texttt{-i}) si usa
(\texttt{\small 5}) semplicemente la funzione \func{perror} per stampare sullo
\textit{standard error}.

La gestione del servizio \textit{echo} viene effettuata interamente nella
funzione \code{ServEcho}, il cui codice è mostrato in
fig.~\ref{fig:TCP_ServEcho_first}, e la comunicazione viene gestita all'interno
di un ciclo (\texttt{\small 6--13}).  I dati inviati dal client vengono letti
(\texttt{\small 6}) dal socket con una semplice \func{read}, di cui non si
controlla lo stato di uscita, assumendo che ritorni solo in presenza di dati
in arrivo. 

La riscrittura (\texttt{\small 7}) viene invece gestita dalla funzione
\func{FullWrite} (descritta in fig.~\ref{fig:sock_FullWrite_code}) che si
incarica di tenere conto automaticamente della possibilità che non tutti i
dati di cui è richiesta la scrittura vengano trasmessi con una singola
\func{write}.

\begin{figure}[!htb] 
  \footnotesize \centering
  \begin{minipage}[c]{\codesamplewidth}
    \includecodesample{listati/ServEcho_first.c}
  \end{minipage} 
  \normalsize
  \caption{Codice della prima versione della funzione \code{ServEcho} per la
    gestione del servizio \textit{echo}.}
  \label{fig:TCP_ServEcho_first}
\end{figure}

In caso di errore di scrittura (si ricordi che \func{FullWrite} restituisce un
valore nullo in caso di successo) si provvede (\texttt{\small 8--10}) a
stampare il relativo messaggio con \func{PrintErr}.  Quando il client chiude
la connessione il ricevimento del segmento FIN fa ritornare la \func{read} con
un numero di byte letti pari a zero, il che causa l'uscita dal ciclo e il
ritorno (\texttt{\small 12}) della funzione, che a sua volta causa la
terminazione del processo figlio.


\subsection{L'avvio e il funzionamento ordinario}
\label{sec:TCP_echo_startup}

Benché il codice dell'esempio precedente sia molto ridotto, esso ci permetterà
di considerare in dettaglio le varie problematiche che si possono incontrare
nello scrivere un'applicazione di rete. Infatti attraverso l'esame delle sue
modalità di funzionamento ordinarie, all'avvio e alla terminazione, e di
quello che avviene nelle varie situazioni limite, da una parte potremo
approfondire la comprensione del protocollo TCP/IP e dall'altra ricavare le
indicazioni necessarie per essere in grado di scrivere applicazioni robuste,
in grado di gestire anche i casi limite.

Il primo passo è compilare e lanciare il server (usando l'utente di
amministrazione, per poter usare la porta 7 che è riservata). All'avvio esso
eseguirà l'apertura passiva con la sequenza delle chiamate a \func{socket},
\func{bind}, \func{listen} e poi si bloccherà nella \func{accept}. A questo
punto si potrà controllarne lo stato con \cmd{netstat}:
\begin{Console}
[piccardi@roke piccardi]$ \textbf{netstat -at}
Active Internet connections (servers and established)
Proto Recv-Q Send-Q Local Address           Foreign Address         State 
...
tcp        0      0 *:echo                  *:*                     LISTEN
...
\end{Console}
%$
che ci mostra come il socket sia in ascolto sulla porta richiesta, accettando
connessioni da qualunque indirizzo e da qualunque porta e su qualunque
interfaccia locale.

A questo punto si può lanciare il client, esso chiamerà \func{socket} e
\func{connect}; una volta completato il \textit{three way handshake} la
connessione è stabilita; la \func{connect} ritornerà nel client\footnote{si
  noti che è sempre la \func{connect} del client a ritornare per prima, in
  quanto questo avviene alla ricezione del secondo segmento (l'ACK del server)
  del \textit{three way handshake}, la \func{accept} del server ritorna solo
  dopo un altro mezzo RTT quando il terzo segmento (l'ACK del client) viene
  ricevuto.}  e la \func{accept} nel server, ed usando di nuovo \cmd{netstat}
otterremmo che:
\begin{Terminal}
Active Internet connections (servers and established)
Proto Recv-Q Send-Q Local Address           Foreign Address         State
tcp        0      0 *:echo                  *:*                     LISTEN
tcp        0      0 roke:echo               gont:32981              ESTABLISHED
\end{Terminal}
mentre per quanto riguarda l'esecuzione dei programmi avremo che:
\begin{itemize}
\item il client chiama la funzione \code{ClientEcho} che si blocca sulla
  \func{fgets} dato che non si è ancora scritto nulla sul terminale.
\item il server eseguirà una \func{fork} facendo chiamare al processo figlio
  la funzione \code{ServEcho}, quest'ultima si bloccherà sulla \func{read}
  dal socket sul quale ancora non sono presenti dati.
\item il processo padre del server chiamerà di nuovo \func{accept}
  bloccandosi fino all'arrivo di un'altra connessione.
\end{itemize}
e se usiamo il comando \cmd{ps} per esaminare lo stato dei processi otterremo
un risultato del tipo:
\begin{Console}
[piccardi@roke piccardi]$ \textbf{ps ax}
  PID TTY      STAT   TIME COMMAND
 ...  ...      ...    ...  ...
 2356 pts/0    S      0:00 ./echod
 2358 pts/1    S      0:00 ./echo 127.0.0.1
 2359 pts/0    S      0:00 ./echod
\end{Console} 
%$
(dove si sono cancellate le righe inutili) da cui si evidenzia la presenza di
tre processi, tutti in stato di \textit{sleep} (vedi
tab.~\ref{tab:proc_proc_states}).

Se a questo punto si inizia a scrivere qualcosa sul client non sarà trasmesso
niente fintanto che non si preme il tasto di a capo (si ricordi quanto detto
in sez.~\ref{sec:file_line_io} a proposito dell'I/O su terminale).  Solo
allora \func{fgets} ritornerà ed il client scriverà quanto immesso dal
terminale sul socket, per poi passare a rileggere quanto gli viene inviato
all'indietro dal server, che a sua volta sarà inviato sullo \textit{standard
  output}, che nel caso ne provoca l'immediata stampa a video.


\subsection{La conclusione normale}
\label{sec:TCP_echo_conclusion}

Tutto quello che scriveremo sul client sarà rimandato indietro dal server e
ristampato a video fintanto che non concluderemo l'immissione dei dati; una
sessione tipica sarà allora del tipo: 
\begin{Console}
[piccardi@roke sources]$ \textbf{./echo 127.0.0.1}
Questa e` una prova
Questa e` una prova
Ho finito
Ho finito
\end{Console} 
%$
che termineremo inviando un EOF dal terminale (usando la combinazione di tasti
ctrl-D, che non compare a schermo); se eseguiamo un \cmd{netstat} a questo
punto avremo:
\begin{Console}
[piccardi@roke piccardi]$ \textbf{netstat -at} 
tcp        0      0 *:echo                  *:*                     LISTEN
tcp        0      0 localhost:33032         localhost:echo          TIME_WAIT
\end{Console} 
%$
con il client che entra in \texttt{TIME\_WAIT}.

Esaminiamo allora in dettaglio la sequenza di eventi che porta alla
terminazione normale della connessione, che ci servirà poi da riferimento
quando affronteremo il comportamento in caso di conclusioni anomale:

\begin{enumerate}
\item inviando un carattere di EOF da terminale la \func{fgets} ritorna
  restituendo un puntatore nullo che causa l'uscita dal ciclo di \code{while},
  così la funzione \code{ClientEcho} ritorna.
\item al ritorno di \code{ClientEcho} ritorna anche la funzione \code{main}, e
  come parte del processo di terminazione tutti i file descriptor vengono chiusi
  (si ricordi quanto detto in sez.~\ref{sec:proc_term_conclusion}); questo
  causa la chiusura del socket di comunicazione; il client allora invierà un
  FIN al server a cui questo risponderà con un ACK.  A questo punto il client
  verrà a trovarsi nello stato \texttt{FIN\_WAIT\_2} ed il server nello stato
  \texttt{CLOSE\_WAIT} (si riveda quanto spiegato in
  sez.~\ref{sec:TCP_conn_term}).
\item quando il server riceve il FIN la \func{read} del processo figlio che
  gestisce la connessione ritorna restituendo 0 causando così l'uscita dal
  ciclo e il ritorno di \code{ServEcho}, a questo punto il processo figlio
  termina chiamando \func{exit}.
\item all'uscita del figlio tutti i file descriptor vengono chiusi, la
  chiusura del socket connesso fa sì che venga effettuata la sequenza finale
  di chiusura della connessione, viene emesso un FIN dal server che riceverà
  un ACK dal client, a questo punto la connessione è conclusa e il client
  resta nello stato \texttt{TIME\_WAIT}.
\end{enumerate}


\subsection{La gestione dei processi figli}
\label{sec:TCP_child_hand}

Tutto questo riguarda la connessione, c'è però da tenere conto dell'effetto
del procedimento di chiusura del processo figlio nel server (si veda quanto
esaminato in sez.~\ref{sec:proc_termination}). In questo caso avremo l'invio
del segnale \signal{SIGCHLD} al padre, ma dato che non si è installato un
gestore e che l'azione predefinita per questo segnale è quella di essere
ignorato, non avendo predisposto la ricezione dello stato di terminazione,
otterremo che il processo figlio entrerà nello stato di \textit{zombie} (si
riveda quanto illustrato in sez.~\ref{sec:sig_sigchld}), come risulterà
ripetendo il comando \cmd{ps}:
\begin{Terminal}
 2356 pts/0    S      0:00 ./echod
 2359 pts/0    Z      0:00 [echod <defunct>]
\end{Terminal}

Dato che non è il caso di lasciare processi \textit{zombie}, occorrerà
ricevere opportunamente lo stato di terminazione del processo (si veda
sez.~\ref{sec:proc_wait}), cosa che faremo utilizzando \signal{SIGCHLD}
secondo quanto illustrato in sez.~\ref{sec:sig_sigchld}. Una prima modifica al
nostro server è pertanto quella di inserire la gestione della terminazione dei
processi figli attraverso l'uso di un gestore.  Per questo useremo la funzione
\code{Signal} (che abbiamo illustrato in fig.~\ref{fig:sig_Signal_code}), per
installare il gestore che riceve i segnali dei processi figli terminati già
visto in fig.~\ref{fig:sig_sigchld_handl}.  Basterà allora aggiungere il
seguente codice: \includecodesnip{listati/sigchildhand.c}
\noindent
all'esempio illustrato in fig.~\ref{fig:TCP_echo_server_first_code}.

In questo modo però si introduce un altro problema. Si ricordi infatti che,
come spiegato in sez.~\ref{sec:sig_gen_beha}, quando un programma si trova in
stato di \texttt{sleep} durante l'esecuzione di una \textit{system call},
questa viene interrotta alla ricezione di un segnale. Per questo motivo, alla
fine dell'esecuzione del gestore del segnale, se questo ritorna, il programma
riprenderà l'esecuzione ritornando dalla \textit{system call} interrotta con
un errore di \errcode{EINTR}.

Vediamo allora cosa comporta tutto questo nel nostro caso: quando si chiude il
client, il processo figlio che gestisce la connessione terminerà, ed il padre,
per evitare la creazione di \textit{zombie}, riceverà il segnale
\signal{SIGCHLD} eseguendo il relativo gestore. Al ritorno del gestore però
l'esecuzione nel padre ripartirà subito con il ritorno della funzione
\func{accept} (a meno di un caso fortuito in cui il segnale arriva durante
l'esecuzione del programma in risposta ad una connessione) con un errore di
\errcode{EINTR}. Non avendo previsto questa eventualità il programma considera
questo un errore fatale terminando a sua volta con un messaggio del tipo:
\begin{Console}
[root@gont sources]# \textbf{./echod -i}
accept error: Interrupted system call
\end{Console}

Come accennato in sez.~\ref{sec:sig_gen_beha} le conseguenze di questo
comportamento delle \textit{system call} possono essere superate in due modi
diversi, il più semplice è quello di modificare il codice di \func{Signal} per
richiedere il riavvio automatico delle \textit{system call} interrotte secondo
la semantica di BSD, usando l'opzione \const{SA\_RESTART} di \func{sigaction};
rispetto a quanto visto in fig.~\ref{fig:sig_Signal_code}. Definiremo allora
la nuova funzione \func{SignalRestart}\footnote{anche questa è definita,
  insieme alle altre funzioni riguardanti la gestione dei segnali, nel file
  \file{SigHand.c}, il cui contento completo può essere trovato negli esempi
  allegati.} come mostrato in fig.~\ref{fig:sig_SignalRestart_code}, ed
installeremo il gestore usando quest'ultima.

\begin{figure}[!htbp]
  \footnotesize \centering
  \begin{minipage}[c]{\codesamplewidth}
    \includecodesample{listati/SignalRestart.c}
  \end{minipage}  
  \normalsize 
  \caption{La funzione \func{SignalRestart}, che installa un gestore di
    segnali in semantica BSD per il riavvio automatico delle \textit{system
      call} interrotte.}
  \label{fig:sig_SignalRestart_code}
\end{figure}

Come si può notare questa funzione è identica alla precedente \func{Signal},
illustrata in fig.~\ref{fig:sig_Signal_code}, solo che in questo caso invece
di inizializzare a zero il campo \var{sa\_flags} di \struct{sigaction}, lo si
inizializza (\texttt{\small 5}) al valore \const{SA\_RESTART}. Usando questa
funzione al posto di \func{Signal} nel server non è necessaria nessuna altra
modifica: le \textit{system call} interrotte saranno automaticamente
riavviate, e l'errore \errcode{EINTR} non si manifesterà più.

La seconda soluzione è più invasiva e richiede di controllare tutte le volte
l'errore restituito dalle varie \textit{system call}, ripetendo la chiamata
qualora questo corrisponda ad \errcode{EINTR}. Questa soluzione ha però il
pregio della portabilità, infatti lo standard POSIX dice che la funzionalità
di riavvio automatico delle \textit{system call}, fornita da
\const{SA\_RESTART}, è opzionale, per cui non è detto che essa sia disponibile
su qualunque sistema.  Inoltre in certi casi,\footnote{Stevens in \cite{UNP1}
  accenna che la maggior parte degli Unix derivati da BSD non fanno ripartire
  \func{select}; altri non riavviano neanche \func{accept} e \func{recvfrom},
  cosa che invece nel caso di Linux viene sempre fatta.} anche quando questa è
presente, non è detto possa essere usata con \func{accept}.

La portabilità nella gestione dei segnali però viene al costo di una
riscrittura parziale del server, la nuova versione di questo, in cui si sono
introdotte una serie di nuove opzioni che ci saranno utili per il debug, è
mostrata in fig.~\ref{fig:TCP_echo_server_code_second}, dove si sono riportate
la sezioni di codice modificate nella seconda versione del programma, il
codice completo di quest'ultimo si trova nel file
\texttt{TCP\_echod\_second.c} dei sorgenti allegati alla guida.

La prima modifica effettuata è stata quella di introdurre una nuova opzione a
riga di comando, \texttt{-c}, che permette di richiedere il comportamento
compatibile nella gestione di \signal{SIGCHLD} al posto della semantica BSD
impostando la variabile \var{compat} ad un valore non nullo. Questa è
preimpostata al valore nullo, cosicché se non si usa questa opzione il
comportamento di default del server è di usare la semantica BSD. 

Una seconda opzione aggiunta è quella di inserire un tempo di attesa fisso
specificato in secondi fra il ritorno della funzione \func{listen} e la
chiamata di \func{accept}, specificabile con l'opzione \texttt{-w}, che
permette di impostare la variabile \var{waiting}.  Infine si è introdotta una
opzione \texttt{-d} per abilitare il debugging che imposta ad un valore non
nullo la variabile \var{debugging}. Al solito si è omessa da
fig.~\ref{fig:TCP_echo_server_code_second} la sezione di codice relativa alla
gestione di tutte queste opzioni, che può essere trovata nel sorgente del
programma.

\begin{figure}[!htb ]
  \footnotesize \centering
  \begin{minipage}[c]{\codesamplewidth}
    \includecodesample{listati/TCP_echod_second.c}
  \end{minipage} 
  \normalsize
  \caption{La sezione nel codice della seconda versione del server
    per il servizio \textit{echo} modificata per tener conto dell'interruzione
    delle \textit{system call}.}
  \label{fig:TCP_echo_server_code_second}
\end{figure}

Vediamo allora come è cambiato il nostro server; una volta definite le
variabili e trattate le opzioni il primo passo (\texttt{\small 9--13}) è
verificare la semantica scelta per la gestione di \signal{SIGCHLD}, a seconda
del valore di \var{compat} (\texttt{\small 9}) si installa il gestore con la
funzione \func{Signal} (\texttt{\small 10}) o con \texttt{SignalRestart}
(\texttt{\small 12}), essendo quest'ultimo il valore di default.

Tutta la sezione seguente, che crea il socket, cede i privilegi di
amministratore ed eventualmente lancia il programma come demone, è rimasta
invariata e pertanto è stata omessa in
fig.~\ref{fig:TCP_echo_server_code_second}; l'unica modifica effettuata prima
dell'entrata nel ciclo principale è stata quella di aver introdotto, subito
dopo la chiamata (\texttt{\small 17--20}) alla funzione \func{listen}, una
eventuale pausa con una condizione (\texttt{\small 21}) sulla variabile
\var{waiting}, che viene inizializzata, con l'opzione \texttt{-w Nsec}, al
numero di secondi da aspettare (il valore preimpostato è nullo).

Si è potuto lasciare inalterata tutta la sezione di creazione del socket
perché nel server l'unica chiamata ad una \textit{system call} lenta, che può
essere interrotta dall'arrivo di \signal{SIGCHLD}, è quella ad \func{accept},
che è l'unica funzione che può mettere il processo padre in stato di sleep nel
periodo in cui un figlio può terminare; si noti infatti come le altre
\textit{system call} lente (si ricordi la distinzione fatta in
sez.~\ref{sec:sig_gen_beha}) o sono chiamate prima di entrare nel ciclo
principale, quando ancora non esistono processi figli, o sono chiamate dai
figli stessi e non risentono di \signal{SIGCHLD}.

Per questo l'unica modifica sostanziale nel ciclo principale (\texttt{\small
  23--42}), rispetto precedente versione di fig.~\ref{fig:TCP_ServEcho_first},
è nella sezione (\texttt{\small 25--31}) in cui si effettua la chiamata di
\func{accept}.  Quest'ultima viene effettuata (\texttt{\small 26--27})
all'interno di un ciclo di \code{while}\footnote{la sintassi del C relativa a
  questo ciclo può non essere del tutto chiara. In questo caso infatti si è
  usato un ciclo vuoto che non esegue nessuna istruzione, in questo modo
  quello che viene ripetuto con il ciclo è soltanto il codice che esprime la
  condizione all'interno del \code{while}.}  che la ripete indefinitamente
qualora in caso di errore il valore di \var{errno} sia \errcode{EINTR}. Negli
altri casi si esce in caso di errore effettivo (\texttt{\small 27--29}),
altrimenti il programma prosegue.

\begin{figure}[!htb]
  \footnotesize \centering
  \begin{minipage}[c]{\codesamplewidth}
    \includecodesample{listati/ServEcho_second.c}
  \end{minipage} 
  \normalsize
  \caption{Codice della seconda versione della funzione \code{ServEcho} per la
    gestione del servizio \textit{echo}.}
  \label{fig:TCP_ServEcho_second}
\end{figure}

Si noti che in questa nuova versione si è aggiunta una ulteriore sezione
(\texttt{\small 32--40}) di aiuto per il debug del programma, che eseguita con
un controllo (\texttt{\small 33}) sul valore della variabile \var{debugging}
impostato dall'opzione \texttt{-d}. Qualora questo sia nullo, come
preimpostato, non accade nulla, altrimenti (\texttt{\small 33}) l'indirizzo
ricevuto da \var{accept} viene convertito in una stringa che poi
(\texttt{\small 34--39}) viene opportunamente stampata o sullo schermo o nei
log.

Infine come ulteriore miglioria si è perfezionata la funzione \code{ServEcho},
sia per tenere conto della nuova funzionalità di debugging, che per effettuare
un controllo in caso di errore; il codice della nuova versione è mostrato in
fig.~\ref{fig:TCP_ServEcho_second}.

 Rispetto alla precedente versione di fig.~\ref{fig:TCP_ServEcho_first} in
questo caso si è provveduto a controllare (\texttt{\small 7--10}) il valore di
ritorno di \func{read} per rilevare un eventuale errore, in modo da stampare
(\texttt{\small 8}) un messaggio di errore e ritornare (\texttt{\small 9})
concludendo la connessione.

Inoltre qualora sia stata attivata la funzionalità di debug (avvalorando
\var{debugging} tramite l'apposita opzione \texttt{-d}) si provvederà a
stampare (\texttt{\small 16--24}) il numero di byte e la stringa letta dal
client, tenendo conto della modalità di invocazione del server, se interattiva
o in forma di demone.


\section{I vari scenari critici}
\label{sec:TCP_echo_critical}

Con le modifiche viste in sez.~\ref{sec:TCP_child_hand} il nostro esempio
diventa in grado di affrontare la gestione ordinaria delle connessioni, ma un
server di rete deve tenere conto che, al contrario di quanto avviene per i
server che operano nei confronti di processi presenti sulla stessa macchina,
la rete è di sua natura inaffidabile, per cui è necessario essere in grado di
gestire tutta una serie di situazioni critiche che non esistono per i processi
locali.


\subsection{La terminazione precoce della connessione}
\label{sec:TCP_conn_early_abort}

La prima situazione critica è quella della terminazione precoce, causata da un
qualche errore sulla rete, della connessione effettuata da un client. Come
accennato in sez.~\ref{sec:TCP_func_accept} la funzione \func{accept} riporta
tutti gli eventuali errori di rete pendenti su una connessione sul
\textit{connected socket}. Di norma questo non è un problema, in quanto non
appena completata la connessione, \func{accept} ritorna, e l'errore sarà
rilevato in seguito dal processo che gestisce la connessione, alla prima
chiamata di una funzione che opera sul socket.

È però possibile, dal punto di vista teorico, incorrere anche in uno scenario
del tipo di quello mostrato in fig.~\ref{fig:TCP_early_abort}, in cui la
connessione viene abortita sul lato client per un qualche errore di rete con
l'invio di un segmento RST, prima che nel server sia stata chiamata la
funzione \func{accept}.

\begin{figure}[!htb]
  \centering \includegraphics[width=10cm]{img/tcp_client_early_abort}  
  \caption{Un possibile caso di terminazione precoce della connessione.}
  \label{fig:TCP_early_abort}
\end{figure}

Benché questo non sia un fatto comune, un evento simile può essere osservato
con dei server molto occupati. In tal caso, con una struttura del server
simile a quella del nostro esempio, in cui la gestione delle singole
connessioni è demandata a processi figli, può accadere che il \textit{three
  way handshake} venga completato e la relativa connessione abortita subito
dopo, prima che il padre, per via del carico della macchina, abbia fatto in
tempo ad eseguire la chiamata ad \func{accept}. Di nuovo si ha una situazione
analoga a quella illustrata in fig.~\ref{fig:TCP_early_abort}, in cui la
connessione viene stabilita, ma subito dopo si ha una condizione di errore che
la chiude prima che essa sia stata accettata dal programma.

Questo significa che oltre alla interruzione da parte di un segnale, che
abbiamo trattato in sez.~\ref{sec:TCP_child_hand} nel caso particolare di
\signal{SIGCHLD}, si possono ricevere altri errori non fatali all'uscita di
\func{accept}; questi, come nel caso precedente, necessitano semplicemente la
ripetizione della chiamata senza uscire dal programma. In questo caso anche la
versione modificata del nostro server non sarebbe adatta, in quanto uno di
questi errori causerebbe la terminazione dello stesso. In Linux i possibili
errori di rete non fatali, riportati sul socket connesso al ritorno di
\func{accept}, sono \errcode{ENETDOWN}, \errcode{EPROTO},
\errcode{ENOPROTOOPT}, \errcode{EHOSTDOWN}, \errcode{ENONET},
\errcode{EHOSTUNREACH}, \errcode{EOPNOTSUPP} e \errcode{ENETUNREACH}.

Si tenga presente che questo tipo di terminazione non è riproducibile
terminando il client prima della chiamata ad \func{accept}, come si potrebbe
fare usando l'opzione \texttt{-w} per introdurre una pausa dopo il lancio del
demone, in modo da poter avere il tempo per lanciare e terminare una
connessione usando il programma client. In tal caso infatti, alla terminazione
del client, il socket associato alla connessione viene semplicemente chiuso,
attraverso la sequenza vista in sez.~\ref{sec:TCP_conn_term}, per cui la
\func{accept} ritornerà senza errori, e si avrà semplicemente un
\textit{end-of-file} al primo accesso al socket. Nel caso di Linux inoltre,
anche qualora si modifichi il client per fargli gestire l'invio di un segmento
RST alla chiusura dal socket (usando l'opzione \const{SO\_LINGER}, vedi
sez.~\ref{sec:sock_options_main}), non si ha comunque nessun errore al ritorno
di \func{accept}, quanto un errore di \errcode{ECONNRESET} al primo tentativo
di accesso al socket.



\subsection{La terminazione precoce del server}
\label{sec:TCP_server_crash}

Un secondo caso critico è quello in cui si ha una terminazione precoce del
server, ad esempio perché il programma ha un crash. In tal caso si suppone che
il processo termini per un errore fatale, cosa che potremo simulare
inviandogli un segnale di terminazione. La conclusione del processo comporta
la chiusura di tutti i file descriptor aperti, compresi tutti i socket
relativi a connessioni stabilite; questo significa che al momento del crollo
del servizio il client riceverà un FIN dal server in corrispondenza della
chiusura del socket.

Vediamo allora cosa succede nel nostro caso, facciamo partire una connessione
con il server e scriviamo una prima riga, poi terminiamo il server con un
\texttt{C-c}. A questo punto scriviamo una seconda riga e poi un'altra riga
ancora. Il risultato finale della sessione è il seguente:
\begin{Console}
[piccardi@gont sources]$ \textbf{./echo 192.168.1.141}
Prima riga
Prima riga
Seconda riga dopo il C-c
Altra riga
[piccardi@gont sources]$
\end{Console}

Come si vede il nostro client, nonostante la connessione sia stata interrotta
prima dell'invio della seconda riga, non solo accetta di inviarla, ma prende
anche un'altra riga prima di terminare senza riportare nessun
errore. 

Per capire meglio cosa è successo conviene analizzare il flusso dei pacchetti
utilizzando un analizzatore di traffico come \cmd{tcpdump}. Il comando
permette di selezionare, nel traffico di rete generato su una macchina, i
pacchetti che interessano, stampando a video (o salvando su disco) il loro
contenuto. Non staremo qui ad entrare nei dettagli dell'uso del programma, che
sono spiegati dalla pagina di manuale;\footnote{per una trattazione di base si
  può consultare sez.~5.2.2 di \cite{SGL}.} per l'uso che vogliamo farne
quello che ci interessa è, posizionandosi sulla macchina che fa da client,
selezionare tutti i pacchetti che sono diretti o provengono dalla macchina che
fa da server. In questo modo (posto che non ci siano altre connessioni col
server, cosa che avremo cura di evitare) tutti i pacchetti rilevati
apparterranno alla nostra sessione di interrogazione del servizio.

Il comando \cmd{tcpdump} permette selezioni molto complesse, basate sulle
interfacce su cui passano i pacchetti, sugli indirizzi IP, sulle porte, sulle
caratteristiche ed il contenuto dei pacchetti stessi, inoltre permette di
combinare fra loro diversi criteri di selezione con degli operatori logici;
quando un pacchetto che corrisponde ai criteri di selezione scelti viene
rilevato i suoi dati vengono stampati sullo schermo (anche questi secondo un
formato configurabile in maniera molto precisa).  

Lanciando il comando prima di ripetere la sessione di lavoro mostrata
nell'esempio precedente potremo allora catturare tutti pacchetti scambiati fra
il client ed il server; i risultati (in realtà si è ridotta la lunghezza
dell'output rispetto al reale tagliando alcuni dati non necessari alla
comprensione del flusso) prodotti in questa occasione da \cmd{tcpdump} sono
allora i seguenti:
\begin{Console}
[root@gont gapil]# \textbf{tcpdump src 192.168.1.141 or dst 192.168.1.141 -N -t}
tcpdump: listening on eth0
gont.34559 > anarres.echo: S 800922320:800922320(0) win 5840 
anarres.echo > gont.34559: S 511689719:511689719(0) ack 800922321 win 5792 
gont.34559 > anarres.echo: . ack 1 win 5840 
gont.34559 > anarres.echo: P 1:12(11) ack 1 win 5840 
anarres.echo > gont.34559: . ack 12 win 5792 
anarres.echo > gont.34559: P 1:12(11) ack 12 win 5792 
gont.34559 > anarres.echo: . ack 12 win 5840 
anarres.echo > gont.34559: F 12:12(0) ack 12 win 5792 
gont.34559 > anarres.echo: . ack 13 win 5840 
gont.34559 > anarres.echo: P 12:37(25) ack 13 win 5840 
anarres.echo > gont.34559: R 511689732:511689732(0) win 0 
\end{Console}

Le prime tre righe vengono prodotte al momento in cui lanciamo il nostro
client, e corrispondono ai tre pacchetti del \textit{three way handshake}.
L'output del comando riporta anche i numeri di sequenza iniziali, mentre la
lettera \texttt{S} indica che per quel pacchetto si aveva il SYN flag attivo.
Si noti come a partire dal secondo pacchetto sia sempre attivo il campo
\texttt{ack}, seguito dal numero di sequenza per il quale si da il ricevuto;
quest'ultimo, a partire dal terzo pacchetto, viene espresso in forma relativa
per maggiore compattezza.  Il campo \texttt{win} in ogni riga indica la
\textit{advertised window} di cui parlavamo in sez.~\ref{sec:TCP_TCP_opt}.

Allora si può verificare dall'output del comando come venga appunto realizzata
la sequenza di pacchetti descritta in sez.~\ref{sec:TCP_conn_cre}: prima viene
inviato dal client un primo pacchetto con il SYN che inizia la connessione, a
cui il server risponde dando il ricevuto con un secondo pacchetto, che a sua
volta porta un SYN, cui il client risponde con un il terzo pacchetto di
ricevuto.

Ritorniamo allora alla nostra sessione con il servizio echo: dopo le tre righe
del \textit{three way handshake} non avremo nulla fin tanto che non scriveremo
una prima riga sul client; al momento in cui facciamo questo si genera una
sequenza di altri quattro pacchetti. Il primo, dal client al server,
contraddistinto da una lettera \texttt{P} che significa che il flag PSH è
impostato, contiene la nostra riga (che è appunto di 11 caratteri), e ad esso
il server risponde immediatamente con un pacchetto vuoto di ricevuto.

Poi tocca al server riscrivere indietro quanto gli è stato inviato, per cui
sarà lui a mandare indietro un terzo pacchetto con lo stesso contenuto appena
ricevuto, e a sua volta riceverà dal client un ACK nel quarto pacchetto.
Questo causerà la ricezione dell'eco nel client che lo stamperà a video.

A questo punto noi procediamo ad interrompere l'esecuzione del server con un
\texttt{C-c} (cioè con l'invio di \signal{SIGTERM}): nel momento in cui
facciamo questo vengono immediatamente generati altri due pacchetti. La
terminazione del processo infatti comporta la chiusura di tutti i suoi file
descriptor, il che comporta, per il socket che avevamo aperto, l'inizio della
sequenza di chiusura illustrata in sez.~\ref{sec:TCP_conn_term}.  Questo
significa che dal server partirà un FIN, che è appunto il primo dei due
pacchetti, contraddistinto dalla lettera \texttt{F}, cui seguirà al solito un
ACK da parte del client.  

A questo punto la connessione dalla parte del server è chiusa, ed infatti se
usiamo \cmd{netstat} per controllarne lo stato otterremo che sul server si ha:
\begin{Console}
anarres:/home/piccardi# \textbf{netstat -ant}
Active Internet connections (servers and established)
Proto Recv-Q Send-Q Local Address           Foreign Address         State
...      ...    ... ...                     ...                     ...
tcp        0      0 192.168.1.141:7         192.168.1.2:34626       FIN_WAIT2
\end{Console}
cioè essa è andata nello stato \texttt{FIN\_WAIT2}, che indica l'avvenuta
emissione del segmento FIN, mentre sul client otterremo che essa è andata
nello stato \texttt{CLOSE\_WAIT}:
\begin{Console}
[root@gont gapil]# \textbf{netstat -ant}
Active Internet connections (servers and established)
Proto Recv-Q Send-Q Local Address           Foreign Address         State
...      ...    ... ...                     ...                     ...
tcp        1      0 192.168.1.2:34582       192.168.1.141:7         CLOSE_WAIT 
\end{Console}

Il problema è che in questo momento il client è bloccato dentro la funzione
\texttt{ClientEcho} nella chiamata a \func{fgets}, e sta attendendo dell'input
dal terminale, per cui non è in grado di accorgersi di nulla. Solo quando
inseriremo la seconda riga il comando uscirà da \func{fgets} e proverà a
scriverla sul socket. Questo comporta la generazione degli ultimi due
pacchetti riportati da \cmd{tcpdump}: il primo, inviato dal client contenente
i 25 caratteri della riga appena letta, e ad esso la macchina server
risponderà, non essendoci più niente in ascolto sulla porta 7, con un segmento
di RST, contraddistinto dalla lettera \texttt{R}, che causa la conclusione
definitiva della connessione anche nel client, dove non comparirà più
nell'output di \cmd{netstat}.

Come abbiamo accennato in sez.~\ref{sec:TCP_conn_term} e come vedremo più
avanti in sez.~\ref{sec:TCP_shutdown} la chiusura di un solo capo di un socket
è una operazione lecita, per cui la nostra scrittura avrà comunque successo
(come si può constatare lanciando usando \cmd{strace}\footnote{il comando
  \cmd{strace} è un comando di debug molto utile che prende come argomento un
  altro comando e ne stampa a video tutte le invocazioni di una \textit{system
    call}, coi relativi argomenti e valori di ritorno, per cui usandolo in
  questo contesto potremo verificare che effettivamente la \func{write} ha
  scritto la riga, che in effetti è stata pure trasmessa via rete.}), in
quanto il nostro programma non ha a questo punto alcun modo di sapere che
dall'altra parte non c'è più nessuno processo in grado di leggere quanto
scriverà. Questo sarà chiaro solo dopo il tentativo di scrittura, e la
ricezione del segmento RST di risposta che indica che dall'altra parte non si
è semplicemente chiuso un capo del socket, ma è completamente terminato il
programma.

Per questo motivo il nostro client proseguirà leggendo dal socket, e dato che
questo è stato chiuso avremo che, come spiegato in
sez.~\ref{sec:TCP_conn_term}, la funzione \func{read} ritorna normalmente con
un valore nullo. Questo comporta che la seguente chiamata a \func{fputs} non
ha effetto (viene stampata una stringa nulla) ed il client si blocca di nuovo
nella successiva chiamata a \func{fgets}. Per questo diventa possibile
inserire una terza riga e solo dopo averlo fatto si avrà la terminazione del
programma.

Per capire come questa avvenga comunque, non avendo inserito nel codice nessun
controllo di errore, occorre ricordare che, a parte la bidirezionalità del
flusso dei dati, dal punto di vista del funzionamento nei confronti delle
funzioni di lettura e scrittura, i socket sono del tutto analoghi a delle
\textit{pipe}. Allora, da quanto illustrato in sez.~\ref{sec:ipc_pipes},
sappiamo che tutte le volte che si cerca di scrivere su una \textit{pipe} il
cui altro capo non è aperto il lettura il processo riceve un segnale di
\signal{SIGPIPE}, e questo è esattamente quello che avviene in questo caso, e
siccome non abbiamo un gestore per questo segnale, viene eseguita l'azione
preimpostata, che è quella di terminare il processo.

Per gestire in maniera più corretta questo tipo di evento dovremo allora
modificare il nostro client perché sia in grado di trattare le varie tipologie
di errore, per questo dovremo riscrivere la funzione \func{ClientEcho}, in
modo da controllare gli stati di uscita delle varie chiamate. Si è riportata
la nuova versione della funzione in fig.~\ref{fig:TCP_ClientEcho_second}.

\begin{figure}[!htb]
  \footnotesize \centering
  \begin{minipage}[c]{\codesamplewidth}
    \includecodesample{listati/ClientEcho_second.c}
  \end{minipage} 
  \normalsize
  \caption{La sezione nel codice della seconda versione della funzione
    \func{ClientEcho} usata dal client per il servizio \textit{echo}
    modificata per tener conto degli eventuali errori.}
  \label{fig:TCP_ClientEcho_second}
\end{figure}

Come si può vedere in questo caso si controlla il valore di ritorno di tutte
le funzioni, ed inoltre si verifica la presenza di un eventuale
\textit{end-of-file} in caso di lettura. Con questa modifica il nostro client
\cmd{echo} diventa in grado di accorgersi della chiusura del socket da parte
del server, per cui ripetendo la sequenza di operazioni precedenti stavolta
otterremo che:
\begin{Console}
[piccardi@gont sources]$ \textbf{./echo 192.168.1.141}
Prima riga
Prima riga
Seconda riga dopo il C-c
EOF sul socket
\end{Console}
%$
ma di nuovo si tenga presente che non c'è modo di accorgersi della chiusura
del socket fin quando non si esegue la scrittura della seconda riga; il
protocollo infatti prevede che ci debba essere una scrittura prima di ricevere
un RST che confermi la chiusura del file, e solo alle successive scritture si
potrà ottenere un errore.

Questa caratteristica dei socket ci mette di fronte ad un altro problema
relativo al nostro client, e che cioè esso non è in grado di accorgersi di
nulla fintanto che è bloccato nella lettura del terminale fatta con
\func{gets}. In questo caso il problema è minimo, ma esso riemergerà più
avanti, ed è quello che si deve affrontare tutte le volte quando si ha a che
fare con la necessità di lavorare con più descrittori, nel qual caso si pone
la questione di come fare a non restare bloccati su un socket quando altri
potrebbero essere liberi. Vedremo come affrontare questa problematica in
sez.~\ref{sec:TCP_sock_multiplexing}.
 

\subsection{Altri scenari di terminazione della connessione}
\label{sec:TCP_conn_crash}

La terminazione del server è solo uno dei possibili scenari di terminazione
della connessione, un altro caso è ad esempio quello in cui si ha
un'interruzione sulla rete, cosa che potremo simulare facilmente staccando il
cavo di rete.  Un'altra condizione è quella di un blocco completo della
macchina su cui gira il server che deve essere riavviata, cosa che potremo
simulare sia eseguendo un reset fisico (un normale shutdown non va bene; in
tal caso infatti il sistema provvede a terminare tutti i processi, per cui la
situazione sarebbe sostanzialmente identica alla precedente) oppure, in
maniera più gentile, riavviando la macchina dopo aver interrotto la
connessione di rete.

Cominciamo ad analizzare il primo caso, l'interruzione del collegamento di
rete. Ripetiamo la nostra sessione di lavoro precedente, lanciamo il client,
scriviamo una prima riga, poi stacchiamo il cavo e scriviamo una seconda
riga. Il risultato che otterremo è:
\begin{Console}
[piccardi@gont sources]$ \textbf{./echo 192.168.1.141}
Prima riga
Prima riga
Seconda riga dopo l'interruzione
Errore in lettura: No route to host
\end{Console}
%$

Quello che succede in questo è che il programma, dopo aver scritto la seconda
riga, resta bloccato per un tempo molto lungo, prima di dare l'errore
\errcode{EHOSTUNREACH}.  Se andiamo ad osservare con \cmd{strace} cosa accade
nel periodo in cui il programma è bloccato vedremo che stavolta, a differenza
del caso precedente, il programma è bloccato nella lettura dal socket.

Se poi, come nel caso precedente, usiamo l'accortezza di analizzare il
traffico di rete fra client e server con \cmd{tcpdump}, otterremo il seguente
risultato:
\begin{Console}
[root@gont sources]# \textbf{tcpdump src 192.168.1.141 or dst 192.168.1.141 -N -t}
tcpdump: listening on eth0
gont.34685 > anarres.echo: S 1943495663:1943495663(0) win 5840
anarres.echo > gont.34685: S 1215783131:1215783131(0) ack 1943495664 win 5792 
gont.34685 > anarres.echo: . ack 1 win 5840
gont.34685 > anarres.echo: P 1:12(11) ack 1 win 5840 
anarres.echo > gont.34685: . ack 12 win 5792 
anarres.echo > gont.34685: P 1:12(11) ack 12 win 5792 
gont.34685 > anarres.echo: . ack 12 win 5840 
gont.34685 > anarres.echo: P 12:45(33) ack 12 win 5840 
gont.34685 > anarres.echo: P 12:45(33) ack 12 win 5840 
gont.34685 > anarres.echo: P 12:45(33) ack 12 win 5840 
gont.34685 > anarres.echo: P 12:45(33) ack 12 win 5840 
gont.34685 > anarres.echo: P 12:45(33) ack 12 win 5840 
gont.34685 > anarres.echo: P 12:45(33) ack 12 win 5840 
gont.34685 > anarres.echo: P 12:45(33) ack 12 win 5840 
gont.34685 > anarres.echo: P 12:45(33) ack 12 win 5840 
gont.34685 > anarres.echo: P 12:45(33) ack 12 win 5840 
arp who-has anarres tell gont
arp who-has anarres tell gont
arp who-has anarres tell gont
arp who-has anarres tell gont
arp who-has anarres tell gont
arp who-has anarres tell gont
...
\end{Console}

In questo caso l'andamento dei primi sette pacchetti è esattamente lo stesso
di prima. Solo che stavolta, non appena inviata la seconda riga, il programma
si bloccherà nella successiva chiamata a \func{read}, non ottenendo nessuna
risposta. Quello che succede è che nel frattempo il kernel provvede, come
richiesto dal protocollo TCP, a tentare la ritrasmissione della nostra riga un
certo numero di volte, con tempi di attesa crescente fra un tentativo ed il
successivo, per tentare di ristabilire la connessione.

Il risultato finale qui dipende dalla realizzazione dello \textit{stack TCP},
e nel caso di Linux anche dall'impostazione di alcuni dei parametri di sistema
che si trovano in \file{/proc/sys/net/ipv4}, che ne controllano il
comportamento: in questo caso in particolare da
\sysctlrelfile{net/ipv4}{tcp\_retries2} (vedi
sez.~\ref{sec:sock_ipv4_sysctl}). Questo parametro infatti specifica il numero
di volte che deve essere ritentata la ritrasmissione di un pacchetto nel mezzo
di una connessione prima di riportare un errore di timeout.  Il valore
preimpostato è pari a 15, il che comporterebbe 15 tentativi di ritrasmissione,
ma nel nostro caso le cose sono andate diversamente, dato che le
ritrasmissioni registrate da \cmd{tcpdump} sono solo 8; inoltre l'errore
riportato all'uscita del client non è stato \errcode{ETIMEDOUT}, come dovrebbe
essere in questo caso, ma \errcode{EHOSTUNREACH}.

Per capire l'accaduto continuiamo ad analizzare l'output di \cmd{tcpdump}:
esso ci mostra che a un certo punto i tentativi di ritrasmissione del
pacchetto sono cessati, per essere sostituiti da una serie di richieste di
protocollo ARP in cui il client richiede l'indirizzo del server.

Come abbiamo accennato in sez.~\ref{sec:net_tcpip_general} ARP è il protocollo
che si incarica di trovare le corrispondenze fra indirizzo IP e indirizzo
hardware sulla scheda di rete. È evidente allora che nel nostro caso, essendo
client e server sulla stessa rete, è scaduta la voce nella \textit{ARP
  cache}\footnote{la \textit{ARP cache} è una tabella mantenuta internamente
  dal kernel che contiene tutte le corrispondenze fra indirizzi IP e indirizzi
  fisici, ottenute appunto attraverso il protocollo ARP; le voci della tabella
  hanno un tempo di vita limitato, passato il quale scadono e devono essere
  nuovamente richieste.} relativa ad \texttt{anarres}, ed il nostro client ha
iniziato ad effettuare richieste ARP sulla rete per sapere l'IP di
quest'ultimo, che essendo scollegato non poteva rispondere. Anche per questo
tipo di richieste esiste un timeout, per cui dopo un certo numero di tentativi
il meccanismo si è interrotto, e l'errore riportato al programma a questo
punto è stato \errcode{EHOSTUNREACH}, in quanto non si era più in grado di
contattare il server.

Un altro errore possibile in questo tipo di situazione, che si può avere
quando la macchina è su una rete remota, è \errcode{ENETUNREACH}; esso viene
riportato alla ricezione di un pacchetto ICMP di \textit{destination
  unreachable} da parte del router che individua l'interruzione della
connessione. Di nuovo anche qui il risultato finale dipende da quale è il
meccanismo più veloce che porta ad accorgersi del problema.

Se però agiamo sui parametri del kernel, e scriviamo in \file{tcp\_retries2}
un valore di tentativi più basso, possiamo evitare la scadenza della
\textit{ARP cache} e vedere cosa succede. Così se ad esempio richiediamo 4
tentativi di ritrasmissione, l'analisi di \cmd{tcpdump} ci riporterà il
seguente scambio di pacchetti:
\begin{Console}
[root@gont gapil]# \textbf{tcpdump src 192.168.1.141 or dst 192.168.1.141 -N -t}
tcpdump: listening on eth0
gont.34752 > anarres.echo: S 3646972152:3646972152(0) win 5840
anarres.echo > gont.34752: S 2735190336:2735190336(0) ack 3646972153 win 5792 
gont.34752 > anarres.echo: . ack 1 win 5840 
gont.34752 > anarres.echo: P 1:12(11) ack 1 win 5840
anarres.echo > gont.34752: . ack 12 win 5792 
anarres.echo > gont.34752: P 1:12(11) ack 12 win 5792 
gont.34752 > anarres.echo: . ack 12 win 5840 
gont.34752 > anarres.echo: P 12:45(33) ack 12 win 5840 
gont.34752 > anarres.echo: P 12:45(33) ack 12 win 5840 
gont.34752 > anarres.echo: P 12:45(33) ack 12 win 5840 
gont.34752 > anarres.echo: P 12:45(33) ack 12 win 5840 
gont.34752 > anarres.echo: P 12:45(33) ack 12 win 5840 
\end{Console}
e come si vede in questo caso i tentativi di ritrasmissione del pacchetto
iniziale sono proprio 4 (per un totale di 5 voci con quello trasmesso la prima
volta), ed in effetti, dopo un tempo molto più breve rispetto a prima ed in
corrispondenza dell'invio dell'ultimo tentativo, quello che otterremo come
errore all'uscita del client sarà diverso, e cioè:
\begin{Console}
[piccardi@gont sources]$ \textbf{./echo 192.168.1.141}
Prima riga
Prima riga
Seconda riga dopo l'interruzione
Errore in lettura: Connection timed out
\end{Console}
%$
che corrisponde appunto, come ci aspettavamo, alla ricezione di un
\errcode{ETIMEDOUT}.

Analizziamo ora il secondo scenario, in cui si ha un crollo della macchina che
fa da server. Al solito lanciamo il nostro client, scriviamo una prima riga
per verificare che sia tutto a posto, poi stacchiamo il cavo e riavviamo il
server. A questo punto, ritornato attivo il server, scriviamo una seconda
riga. Quello che otterremo in questo caso è:
\begin{Console}
[piccardi@gont sources]$ \textbf{./echo 192.168.1.141}
Prima riga
Prima riga
Seconda riga dopo l'interruzione
Errore in lettura Connection reset by peer
\end{Console}
%$
e l'errore ricevuti da \func{read} stavolta è \errcode{ECONNRESET}. Se al
solito riportiamo l'analisi dei pacchetti effettuata con \cmd{tcpdump},
avremo:
\begin{Console}
[root@gont gapil]# \textbf{tcpdump src 192.168.1.141 or dst 192.168.1.141 -N -t}
tcpdump: listening on eth0
gont.34756 > anarres.echo: S 904864257:904864257(0) win 5840 
anarres.echo > gont.34756: S 4254564871:4254564871(0) ack 904864258 win 5792
gont.34756 > anarres.echo: . ack 1 win 5840 
gont.34756 > anarres.echo: P 1:12(11) ack 1 win 5840
anarres.echo > gont.34756: . ack 12 win 5792 
anarres.echo > gont.34756: P 1:12(11) ack 12 win 5792
gont.34756 > anarres.echo: . ack 12 win 5840 
gont.34756 > anarres.echo: P 12:45(33) ack 12 win 5840
anarres.echo > gont.34756: R 4254564883:4254564883(0) win 0 
\end{Console}

Ancora una volta i primi sette pacchetti sono gli stessi; ma in questo caso
quello che succede dopo lo scambio iniziale è che, non avendo inviato nulla
durante il periodo in cui si è riavviato il server, il client è del tutto
ignaro dell'accaduto per cui quando effettuerà una scrittura, dato che la
macchina server è stata riavviata e che tutti gli stati relativi alle
precedenti connessioni sono completamente persi, anche in presenza di una
nuova istanza del server echo non sarà possibile consegnare i dati in arrivo,
per cui alla loro ricezione il kernel risponderà con un segmento di RST.

Il client da parte sua, dato che neanche in questo caso non è stato emesso un
FIN, dopo aver scritto verrà bloccato nella successiva chiamata a \func{read},
che però adesso ritornerà immediatamente alla ricezione del segmento RST,
riportando appunto come errore \errcode{ECONNRESET}. Occorre precisare che se
si vuole che il client sia in grado di accorgersi del crollo del server anche
quando non sta effettuando uno scambio di dati, è possibile usare una
impostazione speciale del socket (ci torneremo in
sez.~\ref{sec:sock_generic_options}) che provvede all'esecuzione di questo
controllo.


\section{L'uso dell'\textit{I/O multiplexing}}
\label{sec:TCP_sock_multiplexing}

Affronteremo in questa sezione l'utilizzo dell'\textit{I/O multiplexing},
affrontato in sez.~\ref{sec:file_multiplexing}, nell'ambito delle applicazioni
di rete. Già in sez.~\ref{sec:TCP_server_crash} era emerso il problema
relativo al client del servizio \textit{echo} che non era in grado di
accorgersi della terminazione precoce del server, essendo bloccato nella
lettura dei dati immessi da tastiera.

Abbiamo visto in sez.~\ref{sec:file_multiplexing} quali sono le funzionalità
del sistema che ci permettono di tenere sotto controllo più file descriptor in
contemporanea; in quella occasione non abbiamo fatto esempi, in quanto quando
si tratta con file normali questa tipologia di I/O normalmente non viene
usata, è invece un caso tipico delle applicazioni di rete quello di dover
gestire varie connessioni da cui possono arrivare dati comuni in maniera
asincrona, per cui riprenderemo l'argomento in questa sezione.


\subsection{Il comportamento della funzione \func{select} con i socket.}
\label{sec:TCP_sock_select}

Iniziamo con la prima delle funzioni usate per l'\textit{I/O multiplexing},
\func{select}; il suo funzionamento è già stato descritto in dettaglio in
sez.~\ref{sec:file_multiplexing} e non staremo a ripetere quanto detto lì;
sappiamo che la funzione ritorna quando uno o più dei file descriptor messi
sotto controllo è pronto per la relativa operazione.

In quell'occasione non abbiamo però definito cosa si intende per pronto,
infatti per dei normali file, o anche per delle \textit{pipe}, la condizione
di essere pronti per la lettura o la scrittura è ovvia; invece lo è molto meno
nel caso dei socket, visto che possono intervenire tutte una serie di
possibili condizioni di errore dovute alla rete. Occorre allora specificare
chiaramente quali sono le condizioni per cui un socket risulta essere
``\textsl{pronto}'' quando viene passato come membro di uno dei tre
\textit{file descriptor set} usati da \func{select}.

Le condizioni che fanno si che la funzione \func{select} ritorni segnalando
che un socket (che sarà riportato nel primo insieme di file descriptor) è
pronto per la lettura sono le seguenti:
\begin{itemize*}
\item nel buffer di ricezione del socket sono arrivati dei dati in quantità
  sufficiente a superare il valore di una \textsl{soglia di basso livello} (il
  cosiddetto \textit{low watermark}). Questo valore è espresso in numero di
  byte e può essere impostato con l'opzione del socket \const{SO\_RCVLOWAT}
  (tratteremo l'uso di questa opzione in sez.~\ref{sec:sock_generic_options});
  il suo valore di default è 1 per i socket TCP e UDP. In questo caso una
  operazione di lettura avrà successo e leggerà un numero di byte maggiore di
  zero.
\item il lato in lettura della connessione è stato chiuso; si è cioè ricevuto
  un segmento FIN (si ricordi quanto illustrato in
  sez.~\ref{sec:TCP_conn_term}) sulla connessione. In questo caso una
  operazione di lettura avrà successo, ma non risulteranno presenti dati (in
  sostanza \func{read} ritornerà con un valore nullo) per indicare la
  condizione di \textit{end-of-file}.
\item c'è stato un errore sul socket. In questo caso una operazione di lettura
  non si bloccherà ma restituirà una condizione di errore (ad esempio
  \func{read} restituirà $-1$) e imposterà la variabile \var{errno} al relativo
  valore. Vedremo in sez.~\ref{sec:sock_generic_options} come sia possibile
  estrarre e cancellare gli errori pendenti su un socket senza usare
  \func{read} usando l'opzione \const{SO\_ERROR}.
\item quando si sta utilizzando un \textit{listening socket} ed ci sono delle
  connessioni completate. In questo caso la funzione \func{accept} non si
  bloccherà.\footnote{in realtà questo non è sempre vero, come accennato in
    sez.~\ref{sec:TCP_conn_early_abort} una connessione può essere abortita
    dalla ricezione di un segmento RST una volta che è stata completata,
    allora se questo avviene dopo che \func{select} è ritornata, ma prima
    della chiamata ad \func{accept}, quest'ultima, in assenza di altre
    connessioni, potrà bloccarsi.}
\end{itemize*}

Le condizioni che fanno si che la funzione \func{select} ritorni segnalando
che un socket (che sarà riportato nel secondo insieme di file descriptor) è
pronto per la scrittura sono le seguenti:
\begin{itemize*}
\item nel buffer di invio è disponibile una quantità di spazio superiore al
  valore della \textsl{soglia di basso livello} in scrittura ed inoltre o il
  socket è già connesso o non necessita (ad esempio è UDP) di connessione.  Il
  valore della soglia è espresso in numero di byte e può essere impostato con
  l'opzione del socket \const{SO\_SNDLOWAT} (trattata in 
  sez.~\ref{sec:sock_generic_options}); il suo valore di default è 2048 per i
  socket TCP e UDP. In questo caso una operazione di scrittura non si
  bloccherà e restituirà un valore positivo pari al numero di byte accettati
  dal livello di trasporto.
\item il lato in scrittura della connessione è stato chiuso. In questo caso
  una operazione di scrittura sul socket genererà il segnale \signal{SIGPIPE}.
\item c'è stato un errore sul socket. In questo caso una operazione di
  scrittura non si bloccherà ma restituirà una condizione di errore ed
  imposterà opportunamente la variabile \var{errno}. Vedremo in
  sez.~\ref{sec:sock_generic_options} come sia possibile estrarre e cancellare
  errori pendenti su un socket usando l'opzione \const{SO\_ERROR}.
\end{itemize*}

Infine c'è una sola condizione che fa sì che \func{select} ritorni segnalando
che un socket (che sarà riportato nel terzo insieme di file descriptor) ha una
condizione di eccezione pendente, e cioè la ricezione sul socket di
\textsl{dati urgenti} (o \textit{out-of-band}), una caratteristica specifica
dei socket TCP su cui torneremo in sez.~\ref{sec:TCP_urgent_data}.

Si noti come nel caso della lettura \func{select} si applichi anche ad
operazioni che non hanno nulla a che fare con l'I/O di dati come il
riconoscimento della presenza di connessioni pronte, in modo da consentire
anche l'utilizzo di \func{accept} in modalità non bloccante. Si noti infine
come in caso di errore un socket venga sempre riportato come pronto sia per la
lettura che per la scrittura.

Lo scopo dei due valori di soglia per i buffer di ricezione e di invio è
quello di consentire maggiore flessibilità nell'uso di \func{select} da parte
dei programmi, se infatti si sa che una applicazione non è in grado di fare
niente fintanto che non può ricevere o inviare una certa quantità di dati, si
possono utilizzare questi valori per far sì che \func{select} ritorni solo
quando c'è la certezza di avere dati a sufficienza.\footnote{questo tipo di
  controllo è utile di norma solo per la lettura, in quanto in genere le
  operazioni di scrittura sono già controllate dall'applicazione, che sa
  sempre quanti dati invia, mentre non è detto possa conoscere la quantità di
  dati in ricezione; per cui, nella situazione in cui si conosce almeno un
  valore minimo, per evitare la penalizzazione dovuta alla ripetizione delle
  operazioni di lettura per accumulare dati sufficienti, si può lasciare al
  kernel il compito di impostare un minimo al di sotto del quale il socket,
  pur avendo disponibili dei dati, non viene dato per pronto in lettura.}



\subsection{Un esempio di \textit{I/O multiplexing}}
\label{sec:TCP_multiplex_example}

Abbiamo incontrato la problematica tipica che conduce all'uso dell'\textit{I/O
  multiplexing} nella nostra analisi degli errori in
sez.~\ref{sec:TCP_conn_early_abort}, quando il nostro client non era in grado
di rendersi conto di errori sulla connessione essendo impegnato nella attesa
di dati in ingresso dallo \textit{standard input}.

In questo caso il problema è quello di dover tenere sotto controllo due
diversi file descriptor, lo \textit{standard input}, da cui viene letto il
testo che vogliamo inviare al server, e il socket connesso con il server su
cui detto testo sarà scritto e dal quale poi si vorrà ricevere la
risposta. L'uso dell'\textit{I/O multiplexing} consente di tenere sotto controllo
entrambi, senza restare bloccati.

Nel nostro caso quello che ci interessa è non essere bloccati in lettura sullo
\textit{standard input} in caso di errori sulla connessione o chiusura della
stessa da parte del server. Entrambi questi casi possono essere rilevati
usando \func{select}, per quanto detto in sez.~\ref{sec:TCP_sock_select},
mettendo sotto osservazione i file descriptor per la condizione di essere
pronti in lettura: sia infatti che si ricevano dati, che la connessione sia
chiusa regolarmente (con la ricezione di un segmento FIN) che si riceva una
condizione di errore (con un segmento RST) il socket connesso sarà pronto in
lettura (nell'ultimo caso anche in scrittura, ma questo non è necessario ai
nostri scopi).

\begin{figure}[!htb]
  \footnotesize \centering
  \begin{minipage}[c]{\codesamplewidth}
    \includecodesample{listati/ClientEcho_third.c}
  \end{minipage} 
  \normalsize
  \caption{La sezione nel codice della terza versione della funzione
    \func{ClientEcho} usata dal client per il servizio \textit{echo}
    modificata per l'uso di \func{select}.}
  \label{fig:TCP_ClientEcho_third}
\end{figure}

Riprendiamo allora il codice del client, modificandolo per l'uso di
\func{select}. Quello che dobbiamo modificare è la funzione \func{ClientEcho}
di fig.~\ref{fig:TCP_ClientEcho_second}, dato che tutto il resto, che riguarda
le modalità in cui viene stabilita la connessione con il server, resta
assolutamente identico. La nostra nuova versione di \func{ClientEcho}, la
terza della serie, è riportata in fig.~\ref{fig:TCP_ClientEcho_third}, il
codice completo si trova nel file \file{TCP\_echo\_third.c} dei sorgenti
allegati alla guida.

In questo caso la funzione comincia (\texttt{\small 8--9}) con l'azzeramento
del \textit{file descriptor set} \var{fset} e l'impostazione del valore
\var{maxfd}, da passare a \func{select} come massimo per il numero di file
descriptor. Per determinare quest'ultimo si usa la macro \code{max} definita
nel nostro file \file{macro.h} che raccoglie una collezione di macro di
preprocessore di varia utilità.

La funzione prosegue poi (\texttt{\small 10--41}) con il ciclo principale, che
viene ripetuto indefinitamente. Per ogni ciclo si reinizializza
(\texttt{\small 11--12}) il \textit{file descriptor set}, impostando i valori
per il file descriptor associato al socket \var{socket} e per lo
\textit{standard input} (il cui valore si recupera con la funzione
\func{fileno}). Questo è necessario in quanto la successiva (\texttt{\small
  13}) chiamata a \func{select} comporta una modifica dei due bit relativia
questi file, che quindi devono essere reimpostati all'inizio di ogni ciclo.

Si noti come la chiamata a \func{select} venga eseguita usando come primo
argomento il valore di \var{maxfd}, precedentemente calcolato, e passando poi
il solo \textit{file descriptor set} per il controllo dell'attività in
lettura, negli altri argomenti vengono passati tutti puntatori nulli, non
interessando in questo caso né il controllo delle altre attività, né
l'impostazione di un valore di timeout.

Al ritorno di \func{select} si provvede a controllare quale dei due file
descriptor presenta attività in lettura, cominciando (\texttt{\small 14--24})
con il file descriptor associato allo \textit{standard input}. In caso di
attività (quando cioè \macro{FD\_ISSET} ritorna una valore diverso da zero) si
esegue (\texttt{\small 15}) una \func{fgets} per leggere gli eventuali dati
presenti; se non ve ne sono (e la funzione restituisce pertanto un puntatore
nullo) si ritorna immediatamente (\texttt{\small 16}) dato che questo
significa che si è chiuso lo \textit{standard input} e quindi concluso l'utilizzo del
client; altrimenti (\texttt{\small 18--22}) si scrivono i dati appena letti
sul socket, prevedendo una uscita immediata in caso di errore di scrittura.

Controllato lo \textit{standard input} si passa a controllare (\texttt{\small
  25--40}) il socket connesso, in caso di attività (\texttt{\small 26}) si
esegue subito una \func{read} di cui si controlla il valore di ritorno; se
questo è negativo si è avuto un errore e pertanto si esce immediatamente
segnalandolo (\texttt{\small 27--30}), se è nullo significa che il server ha
chiuso la connessione, e di nuovo si esce con stampando prima un messaggio di
avviso (\texttt{\small 31--34}), altrimenti (\texttt{\small 35--39}) si
effettua la terminazione della stringa e la si stampa a sullo \textit{standard
  output}, uscendo in caso di errore, per ripetere il ciclo da capo.

Con questo meccanismo il programma invece di essere bloccato in lettura sullo
\textit{standard input} resta bloccato sulla \func{select}, che ritorna
soltanto quando viene rilevata attività su uno dei due file descriptor posti
sotto controllo.  Questo di norma avviene solo quando si è scritto qualcosa
sullo \textit{standard input}, o quando si riceve dal socket la risposta a quanto si
era appena scritto. 

Ma adesso il client diventa capace di accorgersi immediatamente della
terminazione del server; in tal caso infatti il server chiuderà il socket
connesso, ed alla ricezione del FIN la funzione \func{select} ritornerà (come
illustrato in sez.~\ref{sec:TCP_sock_select}) segnalando una condizione di end
of file, per cui il nostro client potrà uscire immediatamente.

Riprendiamo la situazione affrontata in sez.~\ref{sec:TCP_server_crash},
terminando il server durante una connessione, in questo caso quello che
otterremo, una volta scritta una prima riga ed interrotto il server con un
\texttt{C-c}, sarà:
\begin{Console}
[piccardi@gont sources]$ \textbf{./echo 192.168.1.1}
Prima riga
Prima riga
EOF sul socket
\end{Console}
%$
dove l'ultima riga compare immediatamente dopo aver interrotto il server. Il
nostro client infatti è in grado di accorgersi immediatamente che il socket
connesso è stato chiuso ed uscire immediatamente.

Veniamo allora agli altri scenari di terminazione anomala visti in
sez.~\ref{sec:TCP_conn_crash}. Il primo di questi è l'interruzione fisica della
connessione; in questo caso avremo un comportamento analogo al precedente, in
cui si scrive una riga e non si riceve risposta dal server e non succede
niente fino a quando non si riceve un errore di \errcode{EHOSTUNREACH} o
\errcode{ETIMEDOUT} a seconda dei casi.

La differenza è che stavolta potremo scrivere più righe dopo l'interruzione,
in quanto il nostro client dopo aver inviato i dati non si bloccherà più nella
lettura dal socket, ma nella \func{select}; per questo potrà accettare
ulteriore dati che scriverà di nuovo sul socket, fintanto che c'è spazio sul
buffer di uscita (ecceduto il quale si bloccherà in scrittura). 

Si ricordi infatti che il client non ha modo di determinare se la connessione
è attiva o meno (dato che in molte situazioni reali l'inattività può essere
temporanea).  Tra l'altro se si ricollega la rete prima della scadenza del
timeout, potremo anche verificare come tutto quello che si era scritto viene
poi effettivamente trasmesso non appena la connessione ridiventa attiva, per
cui otterremo qualcosa del tipo:
\begin{Console}
[piccardi@gont sources]$ \textbf{./echo 192.168.1.1}
Prima riga
Prima riga
Seconda riga dopo l'interruzione
Terza riga
Quarta riga
Seconda riga dopo l'interruzione
Terza riga
Quarta riga
\end{Console}
%$
in cui, una volta riconnessa la rete, tutto quello che abbiamo scritto durante
il periodo di disconnessione restituito indietro e stampato immediatamente.

Lo stesso comportamento visto in sez.~\ref{sec:TCP_server_crash} si riottiene
nel caso di un crollo completo della macchina su cui sta il server. In questo
caso di nuovo il client non è in grado di accorgersi di niente dato che si
suppone che il programma server non venga terminato correttamente, ma si
blocchi tutto senza la possibilità di avere l'emissione di un segmento FIN che
segnala la terminazione della connessione. 

Di nuovo fintanto che la connessione non si riattiva (con il riavvio della
macchina del server) il client non è in grado di fare altro che accettare
dell'input e tentare di inviarlo. La differenza in questo caso è che non
appena la connessione ridiventa attiva i dati verranno sì trasmessi, ma
essendo state perse tutte le informazioni relative alle precedenti connessioni
ai tentativi di scrittura del client sarà risposto con un segmento RST che
provocherà il ritorno di \func{select} per la ricezione di un errore di
\errcode{ECONNRESET}.


\subsection{La funzione \func{shutdown}}
\label{sec:TCP_shutdown}

Come spiegato in sez.~\ref{sec:TCP_conn_term} il procedimento di chiusura di
un socket TCP prevede che da entrambe le parti venga emesso un segmento FIN. È
pertanto del tutto normale dal punto di vista del protocollo che uno dei due
capi chiuda la connessione quando l'altro capo la lascia aperta; abbiamo
incontrato questa situazione nei vari scenari critici di
sez.~\ref{sec:TCP_echo_critical}.

\itindbeg{half-close}

È pertanto possibile avere una situazione in cui un capo della connessione,
non avendo più nulla da scrivere, possa chiudere il socket, segnalando così
l'avvenuta terminazione della trasmissione (l'altro capo riceverà infatti un
\textit{end-of-file} in lettura) mentre dall'altra parte si potrà proseguire
la trasmissione dei dati scrivendo sul socket che da quel lato è ancora
aperto.  Questa è quella situazione in cui si dice che il socket è
\textsl{mezzo chiuso} (``\textit{half closed}'').

Il problema che si pone è che se la chiusura del socket è effettuata con la
funzione \func{close}, come spiegato in sez.~\ref{sec:TCP_func_close}, si
perde ogni possibilità di poter leggere quanto l'altro capo può star
continuando a scrivere. Per permettere di segnalare che si è finito con la
scrittura, continuando al contempo a leggere quanto può provenire dall'altro
capo del socket, si può usare la funzione \funcd{shutdown}, il cui prototipo
è:

\begin{funcproto}{
\fhead{sys/socket.h}
\fdecl{int shutdown(int sockfd, int how)}
\fdesc{Chiude un lato della connessione fra due socket.} 
}

{La funzione ritorna $0$ in caso di successo e $-1$ per un errore, nel qual
  caso \var{errno} assumerà uno dei valori: 
  \begin{errlist}
  \item[\errcode{EBADF}] \param{sockfd} non è un file descriptor valido.
  \item[\errcode{EINVAL}] il valore di \param{how} non è valido.
  \item[\errcode{ENOTCONN}] il socket non è connesso.
  \item[\errcode{ENOTSOCK}] il file descriptor non corrisponde a un socket.
  \end{errlist}}
\end{funcproto}

La funzione prende come primo argomento il socket \param{sockfd} su cui si
vuole operare e come secondo argomento un valore intero \param{how} che indica
la modalità di chiusura del socket, quest'ultima può prendere soltanto tre
valori: 
\begin{basedescript}{\desclabelwidth{2.2cm}\desclabelstyle{\nextlinelabel}}
\item[\constd{SHUT\_RD}] chiude il lato in lettura del socket, non sarà più
  possibile leggere dati da esso, tutti gli eventuali dati trasmessi
  dall'altro capo del socket saranno automaticamente scartati dal kernel, che,
  in caso di socket TCP, provvederà comunque ad inviare i relativi segmenti di
  ACK.
\item[\constd{SHUT\_WR}] chiude il lato in scrittura del socket, non sarà più
  possibile scrivere dati su di esso. Nel caso di socket TCP la chiamata causa
  l'emissione di un segmento FIN, secondo la procedura chiamata
  \textit{half-close}. Tutti i dati presenti nel buffer di scrittura prima
  della chiamata saranno inviati, seguiti dalla sequenza di chiusura
  illustrata in sez.~\ref{sec:TCP_conn_term}.
\item[\constd{SHUT\_RDWR}] chiude sia il lato in lettura che quello in
  scrittura del socket. È equivalente alla chiamata in sequenza con
  \const{SHUT\_RD} e \const{SHUT\_WR}.
\end{basedescript}

\itindend{half-close}

Ci si può chiedere quale sia l'utilità di avere introdotto \const{SHUT\_RDWR}
quando questa sembra rendere \func{shutdown} del tutto equivalente ad una
\func{close}. In realtà non è così, esiste infatti un'altra differenza con
\func{close}, più sottile. Finora infatti non ci siamo presi la briga di
sottolineare in maniera esplicita che, come per i file e le \textit{fifo},
anche per i socket possono esserci più riferimenti contemporanei ad uno stesso
socket. 

Per cui si avrebbe potuto avere l'impressione che sia una corrispondenza
univoca fra un socket ed il file descriptor con cui vi si accede. Questo non è
assolutamente vero, (e lo abbiamo già visto nel codice del server di
fig.~\ref{fig:TCP_echo_server_first_code}), ed è invece assolutamente normale
che, come per gli altri oggetti, ci possano essere più file descriptor che
fanno riferimento allo stesso socket.

Allora se avviene uno di questi casi quello che succederà è che la chiamata a
\func{close} darà effettivamente avvio alla sequenza di chiusura di un socket
soltanto quando il numero di riferimenti a quest'ultimo diventerà nullo.
Fintanto che ci sono file descriptor che fanno riferimento ad un socket l'uso
di \func{close} si limiterà a deallocare nel processo corrente il file
descriptor utilizzato, ma il socket resterà pienamente accessibile attraverso
tutti gli altri riferimenti. 

Se torniamo all'esempio originale del server di
fig.~\ref{fig:TCP_echo_server_first_code} abbiamo infatti che ci sono due
\func{close}, una sul socket connesso nel padre, ed una sul socket in ascolto
nel figlio, ma queste non effettuano nessuna chiusura reale di detti socket,
dato che restano altri riferimenti attivi, uno al socket connesso nel figlio
ed uno a quello in ascolto nel padre.

Questo non avviene affatto se si usa \func{shutdown} con argomento
\const{SHUT\_RDWR} al posto di \func{close}; in questo caso infatti la
chiusura del socket viene effettuata immediatamente, indipendentemente dalla
presenza di altri riferimenti attivi, e pertanto sarà efficace anche per tutti
gli altri file descriptor con cui, nello stesso o in altri processi, si fa
riferimento allo stesso socket.

Il caso più comune di uso di \func{shutdown} è comunque quello della chiusura
del lato in scrittura, per segnalare all'altro capo della connessione che si è
concluso l'invio dei dati, restando comunque in grado di ricevere quanto
questi potrà ancora inviarci. Questo è ad esempio l'uso che ci serve per
rendere finalmente completo il nostro esempio sul servizio \textit{echo}. 

Il nostro client infatti presenta ancora un problema, che nell'uso che finora
ne abbiamo fatto non è emerso, ma che ci aspetta dietro l'angolo non appena
usciamo dall'uso interattivo e proviamo ad eseguirlo redirigendo
\textit{standard input} e \textit{standard output}. Così se eseguiamo:
\begin{Console}
[piccardi@gont sources]$ \textbf{./echo 192.168.1.1 < ../fileadv.tex  > copia}
\end{Console}
%$
vedremo che il file \texttt{copia} risulta mancare della parte finale.

Per capire cosa avviene in questo caso occorre tenere presente come avviene la
comunicazione via rete; quando redirigiamo lo \textit{standard input} il
nostro client inizierà a leggere il contenuto del file \texttt{../fileadv.tex}
a blocchi di dimensione massima pari a \texttt{MAXLINE} per poi scriverlo,
alla massima velocità consentitagli dalla rete, sul socket. Dato che la
connessione è con una macchina remota, occorre un certo tempo perché i
pacchetti vi arrivino, vengano processati, e poi tornino
indietro. Considerando trascurabile il tempo di processo, questo tempo è
quello impiegato nella trasmissione via rete, che viene detto RTT (dalla
denominazione inglese \itindex{Round~Trip~Time~(RTT)} \textit{Round Trip
  Time}) ed è quello che viene stimato con l'uso del comando \cmd{ping}.

A questo punto, se torniamo al codice mostrato in
fig.~\ref{fig:TCP_ClientEcho_third}, possiamo vedere che mentre i pacchetti
sono in transito sulla rete il client continua a leggere e a scrivere fintanto
che il file in ingresso finisce. Però non appena viene ricevuto un
\textit{end-of-file} in ingresso il nostro client termina. Nel caso
interattivo, in cui si inviavano brevi stringhe una alla volta, c'era sempre
il tempo di eseguire la lettura completa di quanto il server rimandava
indietro. In questo caso invece, quando il client termina, essendo la
comunicazione saturata e a piena velocità, ci saranno ancora pacchetti in
transito sulla rete che devono arrivare al server e poi tornare indietro, ma
siccome il client esce immediatamente dopo la fine del file in ingresso,
questi non faranno a tempo a completare il percorso e verranno persi.

Per evitare questo tipo di problema, invece di uscire una volta completata la
lettura del file in ingresso, occorre usare \func{shutdown} per effettuare la
chiusura del lato in scrittura del socket. In questo modo il client segnalerà
al server la chiusura del flusso dei dati, ma potrà continuare a leggere
quanto il server gli sta ancora inviando indietro, fino a quando anch'esso,
riconosciuta la chiusura del socket in scrittura da parte del client,
effettuerà la chiusura dalla sua parte. Solo alla ricezione della chiusura del
socket da parte del server il client potrà essere sicuro della ricezione di
tutti i dati e della terminazione effettiva della connessione.

\begin{figure}[!htbp]
  \footnotesize \centering
  \begin{minipage}[c]{\codesamplewidth}
    \includecodesample{listati/ClientEcho.c}
  \end{minipage} 
  \normalsize
  \caption{La sezione nel codice della versione finale della funzione
    \func{ClientEcho}, che usa \func{shutdown} per una conclusione corretta
    della connessione.}
  \label{fig:TCP_ClientEcho}
\end{figure}

Si è allora riportato in fig.~\ref{fig:TCP_ClientEcho} la versione finale
della nostra funzione \func{ClientEcho}, in grado di gestire correttamente
l'intero flusso di dati fra client e server. Il codice completo del client,
comprendente la gestione delle opzioni a riga di comando e le istruzioni per
la creazione della connessione, si trova nel file
\texttt{TCP\_echo\_fourth.c}, distribuito coi sorgenti allegati alla guida.

La nuova versione è molto simile alla precedente di
fig.~\ref{fig:TCP_ClientEcho_third}; la prima differenza è l'introduzione
(\texttt{\small 7}) della variabile \var{eof}, inizializzata ad un valore
nullo, che serve a mantenere traccia dell'avvenuta conclusione della lettura
del file in ingresso.

La seconda modifica (\texttt{\small 12--15}) è stata quella di rendere
subordinata ad un valore nullo di \var{eof} l'impostazione del file descriptor
set per l'osservazione dello \textit{standard input}. Se infatti il valore di
\var{eof} è non nullo significa che si è già raggiunta la fine del file in
ingresso ed è pertanto inutile continuare a tenere sotto controllo lo
\textit{standard input} nella successiva (\texttt{\small 16}) chiamata a
\func{select}.

Le maggiori modifiche rispetto alla precedente versione sono invece nella
gestione (\texttt{\small 18--22}) del caso in cui la lettura con \func{fgets}
restituisce un valore nullo, indice della fine del file. Questa nella
precedente versione causava l'immediato ritorno della funzione; in questo caso
prima (\texttt{\small 19}) si imposta opportunamente \var{eof} ad un valore
non nullo, dopo di che (\texttt{\small 20}) si effettua la chiusura del lato
in scrittura del socket con \func{shutdown}. Infine (\texttt{\small 21}) si
usa la macro \macro{FD\_CLR} per togliere lo \textit{standard input} dal
\textit{file descriptor set}.

In questo modo anche se la lettura del file in ingresso è conclusa, la
funzione non esce dal ciclo principale (\texttt{\small 11--50}), ma continua
ad eseguirlo ripetendo la chiamata a \func{select} per tenere sotto controllo
soltanto il socket connesso, dal quale possono arrivare altri dati, che
saranno letti (\texttt{\small 31}) ed opportunamente trascritti
(\texttt{\small 44--48}) sullo \textit{standard output}.

Il ritorno della funzione, e la conseguente terminazione normale del client,
viene invece adesso gestito all'interno (\texttt{\small 30--49}) della lettura
dei dati dal socket; se infatti dalla lettura del socket si riceve una
condizione di \textit{end-of-file}, la si tratterà (\texttt{\small 36--43}) in
maniera diversa a seconda del valore di \var{eof}. Se infatti questa è diversa
da zero (\texttt{\small 37--39}), essendo stata completata la lettura del file
in ingresso, vorrà dire che anche il server ha concluso la trasmissione dei
dati restanti, e si potrà uscire senza errori, altrimenti si stamperà
(\texttt{\small 40--42}) un messaggio di errore per la chiusura precoce della
connessione.


\subsection{Un server basato sull'\textit{I/O multiplexing}}
\label{sec:TCP_serv_select}

Seguendo di nuovo le orme di Stevens in \cite{UNP1} vediamo ora come con
l'utilizzo dell'\textit{I/O multiplexing} diventi possibile riscrivere
completamente il nostro server \textit{echo} con una architettura
completamente diversa, in modo da evitare di dover creare un nuovo processo
tutte le volte che si ha una connessione.\footnote{ne faremo comunque una
  realizzazione diversa rispetto a quella presentata da Stevens in
  \cite{UNP1}.}

La struttura del nuovo server è illustrata in
fig.~\ref{fig:TCP_echo_multiplex}, in questo caso avremo un solo processo che
ad ogni nuova connessione da parte di un client sul socket in ascolto si
limiterà a registrare l'entrata in uso di un nuovo file descriptor ed
utilizzerà \func{select} per rilevare la presenza di dati in arrivo su tutti i
file descriptor attivi, operando direttamente su ciascuno di essi.

\begin{figure}[!htb]
  \centering \includegraphics[width=13cm]{img/TCPechoMult}
  \caption{Schema del nuovo server echo basato sull'\textit{I/O multiplexing}.}
  \label{fig:TCP_echo_multiplex}
\end{figure}

La sezione principale del codice del nuovo server è illustrata in
fig.~\ref{fig:TCP_SelectEchod}. Si è tralasciata al solito la gestione delle
opzioni, che è identica alla versione precedente. Resta invariata anche tutta
la parte relativa alla gestione dei segnali, degli errori, e della cessione
dei privilegi, così come è identica la gestione della creazione del socket (si
può fare riferimento al codice già illustrato in
sez.~\ref{sec:TCPsimp_server_main}); al solito il codice completo del server è
disponibile coi sorgenti allegati nel file \texttt{select\_echod.c}.

\begin{figure}[!htbp]
  \footnotesize \centering
  \begin{minipage}[c]{\codesamplewidth}
    \includecodesample{listati/select_echod.c}
  \end{minipage} 
  \normalsize
  \caption{La sezione principale della nuova versione di server
    \textit{echo} basato sull'uso della funzione \func{select}.}
  \label{fig:TCP_SelectEchod}
\end{figure}

In questo caso, una volta aperto e messo in ascolto il socket, tutto quello
che ci servirà sarà chiamare \func{select} per rilevare la presenza di nuove
connessioni o di dati in arrivo, e processarli immediatamente. Per realizzare
lo schema mostrato in fig.~\ref{fig:TCP_echo_multiplex} il programma usa una
tabella dei socket connessi mantenuta nel vettore \var{fd\_open} dimensionato
al valore di \const{FD\_SETSIZE}, ed una variabile \var{max\_fd} per
registrare il valore più alto dei file descriptor aperti.

Prima di entrare nel ciclo principale (\texttt{\small 5--53}) la nostra
tabella viene inizializzata (\texttt{\small 2}) a zero (valore che
utilizzeremo come indicazione del fatto che il relativo file descriptor non è
aperto), mentre il valore massimo (\texttt{\small 3}) per i file descriptor
aperti viene impostato a quello del socket in ascolto, in quanto esso è
l'unico file aperto, oltre i tre standard, e pertanto avrà il valore più alto,
che verrà anche (\texttt{\small 4}) inserito nella tabella.

La prima sezione (\texttt{\small 6--8}) del ciclo principale esegue la
costruzione del \textit{file descriptor set} \var{fset} in base ai socket
connessi in un certo momento; all'inizio ci sarà soltanto il socket in ascolto
ma nel prosieguo delle operazioni verranno utilizzati anche tutti i socket
connessi registrati nella tabella \var{fd\_open}.  Dato che la chiamata di
\func{select} modifica il valore del \textit{file descriptor set} è necessario
ripetere (\texttt{\small 6}) ogni volta il suo azzeramento per poi procedere
con il ciclo (\texttt{\small 7--8}) in cui si impostano i socket trovati
attivi.

Per far questo si usa la caratteristica dei file descriptor, descritta in
sez.~\ref{sec:file_open_close}, per cui il kernel associa sempre ad ogni nuovo
file il file descriptor con il valore più basso disponibile. Questo fa sì che
si possa eseguire il ciclo (\texttt{\small 7}) a partire da un valore minimo,
che sarà sempre quello del socket in ascolto, mantenuto in \var{list\_fd},
fino al valore massimo di \var{max\_fd} che dovremo aver cura di tenere
aggiornato.  Dopo di che basterà controllare (\texttt{\small 8}) nella nostra
tabella se il file descriptor è in uso o meno,\footnote{si tenga presente che
  benché il kernel assegni sempre il primo valore libero, si potranno sempre
  avere dei \textsl{buchi} nella nostra tabella dato che nelle operazioni i
  socket saranno aperti e chiusi in corrispondenza della creazione e
  conclusione delle connessioni.} e impostare \var{fset} di conseguenza.

Una volta inizializzato con i socket aperti il nostro \textit{file descriptor
  set} potremo chiamare \func{select} per fargli osservare lo stato degli
stessi (in lettura, presumendo che la scrittura sia sempre consentita). Come
per il precedente esempio di sez.~\ref{sec:TCP_child_hand}, essendo questa
l'unica funzione che può bloccarsi ed essere interrotta da un segnale, la
eseguiremo (\texttt{\small 9--10}) all'interno di un ciclo di \code{while},
che la ripete indefinitamente qualora esca con un errore di \errcode{EINTR}.
Nel caso invece di un errore normale si provvede (\texttt{\small 11--14}) ad
uscire dal programma stampando un messaggio di errore.

Infine quando la funzione ritorna normalmente avremo in \var{n} il numero di
socket da controllare. Nello specifico si danno due casi per cui \func{select}
può ritornare: o si è ricevuta una nuova connessione ed è pronto il socket in
ascolto, sul quale si può eseguire \func{accept}, o c'è attività su uno dei
socket connessi, sui quali si può eseguire \func{read}.

Il primo caso viene trattato immediatamente (\texttt{\small 15--24}): si
controlla (\texttt{\small 15}) che il socket in ascolto sia fra quelli attivi,
nel qual caso anzitutto (\texttt{\small 16}) se ne decrementa il numero
mantenuto nella variabile \var{n}. Poi, inizializzata (\texttt{\small 17}) la
lunghezza della struttura degli indirizzi, si esegue \func{accept} per
ottenere il nuovo socket connesso, controllando che non ci siano errori
(\texttt{\small 18--21}). In questo caso non c'è più la necessità di
controllare per interruzioni dovute a segnali, in quanto siamo sicuri che
\func{accept} non si bloccherà. Per completare la trattazione occorre a questo
punto aggiungere (\texttt{\small 22}) il nuovo file descriptor alla tabella di
quelli connessi, ed inoltre, se è il caso, aggiornare (\texttt{\small 23}) il
valore massimo in \var{max\_fd}.

Una volta controllato l'arrivo di nuove connessioni si passa a verificare se
ci sono dati sui socket connessi, per questo si ripete un ciclo
(\texttt{\small 26--52}) fintanto che il numero di socket attivi indicato
dalla variabile \var{n} resta diverso da zero. In questo modo, se l'unico
socket con attività era quello connesso, avendola opportunamente decrementata
in precedenza, essa risulterà nulla, pertanto il ciclo di verifica verrà
saltato e si ritornerà all'inzizio del ciclo principale, ripetendo, dopo
l'inizializzazione del \textit{file descriptor set} con i nuovi valori nella
tabella, la chiamata di \func{select}.

Se il socket attivo non è quello in ascolto, o ce ne sono comunque anche
altri, il valore di \var{n} non sarà nullo ed il controllo sarà
eseguito. Prima di entrare nel ciclo di veridica comunque si inizializza
(\texttt{\small 25}) il valore della variabile \var{i}, che useremo come
indice nella tabella \var{fd\_open}, al valore minimo, corrispondente al file
descriptor del socket in ascolto.

Il primo passo (\texttt{\small 27}) nella verifica è incrementare il valore
dell'indice \var{i} per posizionarsi sul primo valore possibile per un file
descriptor associato ad un eventuale socket connesso, dopo di che si controlla
(\texttt{\small 28}) se questo è nella tabella dei socket connessi, chiedendo
la ripetizione del ciclo in caso contrario. Altrimenti si passa a verificare
(\texttt{\small 29}) se il file descriptor corrisponde ad uno di quelli
attivi, e nel caso si esegue (\texttt{\small 30}) una lettura, uscendo con un
messaggio in caso di errore (\texttt{\small 31--35}).

Se (\texttt{\small 36}) il numero di byte letti \var{nread} è nullo si è in
presenza di una \textit{end-of-file}, indice che una connessione che si è
chiusa, che deve essere trattata (\texttt{\small 36--45}) opportunamente.  Il
primo passo è chiudere (\texttt{\small 37}) anche il proprio capo del socket e
rimuovere (\texttt{\small 38}) il file descriptor dalla tabella di quelli
aperti, inoltre occorre verificare (\texttt{\small 39}) se il file descriptor
chiuso è quello con il valore più alto, nel qual caso occorre trovare
(\texttt{\small 39--43}) il nuovo massimo, altrimenti (\texttt{\small 44}) si
può ripetere il ciclo da capo per esaminare (se ne restano) ulteriori file
descriptor attivi.

Se però è stato chiuso il file descriptor più alto, dato che la scansione dei
file descriptor attivi viene fatta a partire dal valore più basso, questo
significa che siamo anche arrivati alla fine della scansione, per questo
possiamo utilizzare direttamente il valore dell'indice \var{i} con un ciclo
all'indietro (\texttt{\small 40}) che trova il primo valore per cui la tabella
presenta un file descriptor aperto, e lo imposta (\texttt{\small 41}) come
nuovo massimo, per poi tornare (\texttt{\small 42}) al ciclo principale con un
\code{break}, e rieseguire \func{select}.

Se infine si sono effettivamente letti dei dati dal socket (ultimo caso
rimasto) si potrà invocare immediatamente (\texttt{\small 46})
\func{FullWrite} per riscriverli indietro sul socket stesso, avendo cura di
uscire con un messaggio in caso di errore (\texttt{\small 47--50}). Si noti
che nel ciclo si esegue una sola lettura, contrariamente a quanto fatto con la
precedente versione (si riveda il codice di fig.~\ref{fig:TCP_ServEcho_second})
in cui si continuava a leggere fintanto che non si riceveva un
\textit{end-of-file}, questo perché usando l'\textit{I/O multiplexing} non si
vuole essere bloccati in lettura.  

L'uso di \func{select} ci permette di trattare automaticamente anche il caso
in cui la \func{read} non è stata in grado di leggere tutti i dati presenti
sul socket, dato che alla iterazione successiva \func{select} ritornerà
immediatamente segnalando l'ulteriore disponibilità.

Il nostro server comunque soffre di una vulnerabilità per un attacco di tipo
\textit{Denial of Service}. Il problema è che in caso di blocco di una
qualunque delle funzioni di I/O, non avendo usato processi separati, tutto il
server si ferma e non risponde più a nessuna richiesta. Abbiamo scongiurato
questa evenienza per l'I/O in ingresso con l'uso di \func{select}, ma non vale
altrettanto per l'I/O in uscita. Il problema pertanto può sorgere qualora una
delle chiamate a \func{write} effettuate da \func{FullWrite} si blocchi. 

Con il funzionamento normale questo non accade in quanto il server si limita a
scrivere quanto riceve in ingresso, ma qualora venga utilizzato un client
malevolo che esegua solo scritture e non legga mai indietro l'\textsl{eco} del
server, si potrebbe giungere alla saturazione del buffer di scrittura, ed al
conseguente blocco del server su di una \func{write}.

Le possibili soluzioni in questo caso sono quelle di ritornare ad eseguire il
ciclo di risposta alle richieste all'interno di processi separati, utilizzare
un timeout per le operazioni di scrittura, o eseguire queste ultime in
modalità non bloccante, concludendo le operazioni qualora non vadano a buon
fine.


\subsection{\textit{I/O multiplexing} con \func{poll}}
\label{sec:TCP_serv_poll}

Finora abbiamo trattato le problematiche risolubili con l'\textit{I/O
  multiplexing} impiegando la funzione \func{select}. Questo è quello che
avviene nella maggior parte dei casi, in quanto essa è nata sotto BSD proprio
per affrontare queste problematiche con i socket.  Abbiamo però visto in
sez.~\ref{sec:file_multiplexing} come la funzione \func{poll} possa costituire
una alternativa a \func{select}, con alcuni vantaggi, non soffrendo delle
limitazioni dovute all'uso dei \textit{file descriptor set}.

Ancora una volta in sez.~\ref{sec:file_poll} abbiamo trattato la funzione in
maniera generica, parlando di file descriptor, ma come per \func{select}
quando si ha a che fare con dei socket, il concetto di essere \textsl{pronti}
per l'I/O deve essere specificato nei dettagli per tener conto delle
condizioni della rete. Inoltre deve essere specificato come viene classificato
il traffico nella suddivisione fra dati normali e prioritari. In generale
pertanto:
\begin{itemize}
\item i dati inviati su un socket vengono considerati traffico normale,
  pertanto vengono rilevati alla loro ricezione sull'altro capo da una
  selezione effettuata con \const{POLLIN} o \const{POLLRDNORM};
\item i dati urgenti \textit{out-of-band} (vedi
  sez.~\ref{sec:TCP_urgent_data}) su un socket TCP vengono considerati
  traffico prioritario e vengono rilevati da una condizione \const{POLLIN},
  \const{POLLPRI} o \const{POLLRDBAND}.
\item la chiusura di una connessione (cioè la ricezione di un segmento FIN)
  viene considerato traffico normale, pertanto viene rilevato da una
  condizione \const{POLLIN} o \const{POLLRDNORM}, ma una conseguente chiamata
  a \func{read} restituirà 0.
\item la disponibilità di spazio sul socket per la scrittura di dati viene
  segnalata con una condizione \const{POLLOUT}.
\item quando uno dei due capi del socket chiude un suo lato della connessione
  con \func{shutdown} si riceve una condizione di \const{POLLHUP}.
\item la presenza di un errore sul socket (sia dovuta ad un segmento RST che a
  timeout) viene considerata traffico normale, ma viene segnalata anche dalla
  condizione \const{POLLERR}.
\item la presenza di una nuova connessione su un socket in ascolto può essere
  considerata sia traffico normale che prioritario, nel caso di Linux la
  realizzazione dello \textit{stack TCP} la classifica come normale.
\end{itemize}

Come esempio dell'uso di \func{poll} proviamo allora a riscrivere il server
\textit{echo} secondo lo schema di fig.~\ref{fig:TCP_echo_multiplex} usando
\func{poll} al posto di \func{select}. In questo caso dovremo fare qualche
modifica, per tenere conto della diversa sintassi delle due funzioni, ma la
struttura del programma resta sostanzialmente la stessa.


\begin{figure}[!htbp]
  \footnotesize \centering
  \begin{minipage}[c]{\codesamplewidth}
    \includecodesample{listati/poll_echod.c}
  \end{minipage} 
  \normalsize
  \caption{La sezione principale della nuova versione di server
    \textit{echo} basato sull'uso della funzione \func{poll}.}
  \label{fig:TCP_PollEchod}
\end{figure}

In fig.~\ref{fig:TCP_PollEchod} è riportata la sezione principale della nuova
versione del server, la versione completa del codice è riportata nel file
\texttt{poll\_echod.c} dei sorgenti allegati alla guida. Al solito nella
figura si sono tralasciate la gestione delle opzioni, la creazione del socket
in ascolto, la cessione dei privilegi e le operazioni necessarie a far
funzionare il programma come demone, privilegiando la sezione principale del
programma.

Come per il precedente server basato su \func{select} il primo passo
(\texttt{\small 2--8}) è quello di inizializzare le variabili necessarie. Dato
che in questo caso dovremo usare un vettore di strutture occorre anzitutto
(\texttt{\small 2}) allocare la memoria necessaria utilizzando il numero
massimo \var{n} di socket osservabili, che viene impostato attraverso
l'opzione \texttt{-n} ed ha un valore di default di 256. 

Dopo di che si preimposta (\texttt{\small 3}) il valore \var{max\_fd} del file
descriptor aperto con valore più alto a quello del socket in ascolto (al
momento l'unico), e si provvede (\texttt{\small 4--7}) ad inizializzare le
strutture, disabilitando l'osservazione (\texttt{\small 5}) con un valore
negativo del campo \var{fd}, ma predisponendo (\texttt{\small 6}) il campo
\var{events} per l'osservazione dei dati normali con \const{POLLRDNORM}.
Infine (\texttt{\small 8}) si attiva l'osservazione del socket in ascolto
inizializzando la corrispondente struttura. Questo metodo comporta, in
modalità interattiva, lo spreco di tre strutture (quelle relative a
\textit{standard input}, \textit{standard output} e \textit{standard error})
che non vengono mai utilizzate in quanto la prima è sempre quella relativa al
socket in ascolto.

Una volta completata l'inizializzazione tutto il lavoro viene svolto
all'interno del ciclo principale \texttt{\small 9--53}) che ha una struttura
sostanzialmente identica a quello usato per il precedente esempio basato su
\func{select}. La prima istruzione (\texttt{\small 10}) è quella di eseguire
\func{poll} all'interno di un ciclo che la ripete qualora venisse interrotta
da un segnale, da cui si esce soltanto quando la funzione ritorna restituendo
nella variabile \var{n} il numero di file descriptor trovati attivi.  Qualora
invece si sia ottenuto un errore si procede (\texttt{\small 11--14}) alla
terminazione immediata del processo provvedendo a stampare una descrizione
dello stesso.

Una volta ottenuta dell'attività su un file descriptor si hanno di nuovo due
possibilità. La prima è che ci sia attività sul socket in ascolto, indice di
una nuova connessione, nel qual caso si controlla (\texttt{\small 17}) se il
campo \var{revents} della relativa struttura è attivo; se è così si provvede
(\texttt{\small 16}) a decrementare la variabile \var{n} (che assume il
significato di numero di file descriptor attivi rimasti da controllare) per
poi (\texttt{\small 17--21}) effettuare la chiamata ad \func{accept},
terminando il processo in caso di errore. Se la chiamata ad \func{accept} ha
successo si procede attivando (\texttt{\small 22}) la struttura relativa al
nuovo file descriptor da essa ottenuto, modificando (\texttt{\small 23})
infine quando necessario il valore massimo dei file descriptor aperti
mantenuto in \var{max\_fd}.

La seconda possibilità è che vi sia dell'attività su uno dei socket aperti in
precedenza, nel qual caso si inizializza (\texttt{\small 25}) l'indice \var{i}
del vettore delle strutture \struct{pollfd} al valore del socket in ascolto,
dato che gli ulteriori socket aperti avranno comunque un valore superiore.  Il
ciclo (\texttt{\small 26--52}) prosegue fintanto che il numero di file
descriptor attivi, mantenuto nella variabile \var{n}, è diverso da zero. Se
pertanto ci sono ancora socket attivi da individuare si comincia con
l'incrementare (\texttt{\small 27}) l'indice e controllare (\texttt{\small
  28}) se corrisponde ad un file descriptor in uso analizzando il valore del
campo \var{fd} della relativa struttura e chiudendo immediatamente il ciclo
qualora non lo sia. Se invece il file descriptor è in uso si verifica
(\texttt{\small 29}) se c'è stata attività controllando il campo
\var{revents}.

Di nuovo se non si verifica la presenza di attività il ciclo si chiude subito,
altrimenti si provvederà (\texttt{\small 30}) a decrementare il numero \var{n}
di file descriptor attivi da controllare e ad eseguire (\texttt{\small 31}) la
lettura, ed in caso di errore (\texttt{\small 32--35}) al solito lo si
notificherà uscendo immediatamente. Qualora invece si ottenga una condizione
di \textit{end-of-file} (\texttt{\small 36--45}) si provvederà a chiudere
(\texttt{\small 37}) anche il nostro capo del socket e a marcarlo
(\texttt{\small 38}) come inutilizzato nella struttura ad esso associata.
Infine dovrà essere ricalcolato (\texttt{\small 39--43}) un eventuale nuovo
valore di \var{max\_fd}. L'ultimo passo è chiudere (\texttt{\small 44}) il
ciclo in quanto in questo caso non c'è più niente da riscrivere all'indietro
sul socket.

Se invece si sono letti dei dati si provvede (\texttt{\small 46}) ad
effettuarne la riscrittura all'indietro, con il solito controllo ed eventuale
uscita e notifica in caso di errore (\texttt{\small 47--51}).

Come si può notare la logica del programma è identica a quella vista in
fig.~\ref{fig:TCP_SelectEchod} per l'analogo server basato su \func{select};
la sola differenza significativa è che in questo caso non c'è bisogno di
rigenerare i \textit{file descriptor set} in quanto l'uscita è indipendente
dai dati in ingresso. Si applicano comunque anche a questo server le
considerazioni finali di sez.~\ref{sec:TCP_serv_select}.


%\subsection{\textit{I/O multiplexing} con \textit{epoll}}
%\label{sec:TCP_serv_epoll}

%Da fare.

% TODO fare esempio con epoll



% LocalWords:  socket TCP client dell'I multiplexing stream three way handshake
% LocalWords:  header stack kernel SYN ACK URG syncronize sez bind listen fig
% LocalWords:  accept connect active acknowledge l'acknowledge nell'header MSS
% LocalWords:  sequence number l'acknowledgement dell'header options l'header
% LocalWords:  option MMS segment size MAXSEG window advertised Mbit sec nell'
% LocalWords:  timestamp RFC long fat close of l'end l'ACK half shutdown CLOSED
% LocalWords:  netstat SENT ESTABLISHED WAIT IPv Ethernet piggybacking UDP MSL
% LocalWords:  l'overhead Stevens Lifetime router hop limit TTL to live RST SSH
% LocalWords:  routing dell'MSL l'IP multitasking well known port ephemeral BSD
% LocalWords:  ports dall' IANA Assigned Authority like glibc netinet IPPORT AF
% LocalWords:  RESERVED USERRESERVED rsh rlogin pair socketpair Local Address
% LocalWords:  Foreing DNS caching INADDR ANY multihoming loopback ssh fuser ip
% LocalWords:  lsof SOCK sys int sockfd const struct sockaddr serv addr socklen
% LocalWords:  addrlen errno EBADF descriptor EINVAL ENOTSOCK EACCES EADDRINUSE
% LocalWords:  EADDRNOTAVAIL EFAULT ENOTDIR ENOENT ENOMEM ELOOP ENOSR EROFS RPC
% LocalWords:  portmapper htonl tab endianness BROADCAST broadcast any extern fd
% LocalWords:  ADRR INIT DGRAM SEQPACKET servaddr ECONNREFUSED ETIMEDOUT EAGAIN
% LocalWords:  ENETUNREACH EINPROGRESS EALREADY EAFNOSUPPORT EPERM EISCONN proc
% LocalWords:  sysctl filesystem syn retries reset ICMP backlog EOPNOTSUPP RECV
% LocalWords:  connection queue dell'ACK flood spoofing syncookies SOMAXCONN CR
% LocalWords:  RDM EWOULDBLOCK firewall ENOBUFS EINTR EMFILE ECONNABORTED NULL
% LocalWords:  ESOCKTNOSUPPORT EPROTONOSUPPORT ERESTARTSYS connected listening
% LocalWords:  DECnet read write NONBLOCK fcntl getsockname getpeername name ps
% LocalWords:  namelen namlen ENOTCONN exec inetd POSIX daytime FullRead count
% LocalWords:  BUF FullWrite system call INET perror htons inet pton ctime FTP
% LocalWords:  fputs carriage return line feed superdemone daytimed sleep fork
% LocalWords:  daemon cunc logging list conn sock exit snprintf ntop ntohs echo
% LocalWords:  crash superserver L'RFC first ClientEcho stdin stdout fgets main
% LocalWords:  MAXLINE initd echod ServEcho setgid short nogroup nobody setuid
% LocalWords:  demonize PrintErr syslog wrapper log error root RTT EOF ctrl ack
% LocalWords:  while SIGCHLD Signal RESTART sigaction SignalRestart SigHand win
% LocalWords:  flags select recvfrom debug second compat waiting Nsec ENETDOWN
% LocalWords:  EPROTO ENOPROTOOPT EHOSTDOWN ENONET EHOSTUNREACH LINGER tcpdump
% LocalWords:  ECONNRESET advertising PSH SIGTERM strace SIGPIPE gets tcp ARP
% LocalWords:  cache anarres destination unreachable l'I low watermark RCVLOWAT
% LocalWords:  SNDLOWAT third fset maxfd fileno ISSET closed how SHUT RD WR eof
% LocalWords:  RDWR fifo Trip ping fourth CLR sull'I SETSIZE nread break Denial
% LocalWords:  Service poll POLLIN POLLRDNORM POLLPRI POLLRDBAND POLLOUT events
% LocalWords:  POLLHUP POLLERR revents pollfd Di scaling SYNCNT DoS

%%% Local Variables: 
%%% mode: latex
%%% TeX-master: "gapil"
%%% End: 


%% sockctrl.tex
%%
%% Copyright (C) 2004-2018 Simone Piccardi.  Permission is granted to
%% copy, distribute and/or modify this document under the terms of the GNU Free
%% Documentation License, Version 1.1 or any later version published by the
%% Free Software Foundation; with the Invariant Sections being "Prefazione",
%% with no Front-Cover Texts, and with no Back-Cover Texts.  A copy of the
%% license is included in the section entitled "GNU Free Documentation
%% License".
%%

\chapter{La gestione dei socket}
\label{cha:sock_generic_management}

Esamineremo in questo capitolo una serie di funzionalità aggiuntive relative
alla gestione dei socket, come la gestione della risoluzione di nomi e
indirizzi, le impostazioni delle varie proprietà ed opzioni relative ai
socket, e le funzioni di controllo che permettono di modificarne il
comportamento.


\section{La risoluzione dei nomi}
\label{sec:sock_name_resolution}

Negli esempi dei capitoli precedenti abbiamo sempre identificato le singole
macchine attraverso indirizzi numerici, sfruttando al più le funzioni di
conversione elementare illustrate in sez.~\ref{sec:sock_addr_func} che
permettono di passare da un indirizzo espresso in forma \textit{dotted
  decimal} ad un numero. Vedremo in questa sezione le funzioni utilizzate per
poter utilizzare dei nomi simbolici al posto dei valori numerici, e viceversa
quelle che permettono di ottenere i nomi simbolici associati ad indirizzi,
porte o altre proprietà del sistema.


\subsection{La struttura del \textit{resolver}}
\label{sec:sock_resolver}

\itindbeg{resolver} La risoluzione dei nomi è associata tradizionalmente al
servizio del \itindex{Domain~Name~Service~(DNS)} \textit{Domain Name Service}
che permette di identificare le macchine su internet invece che per numero IP
attraverso il relativo \textsl{nome a dominio}.\footnote{non staremo ad
  entrare nei dettagli della definizione di cosa è un nome a dominio, dandolo
  per noto, una introduzione alla problematica si trova in \cite{AGL} (cap.~9)
  mentre per una trattazione approfondita di tutte le problematiche relative
  al DNS si può fare riferimento a \cite{DNSbind}.} In realtà per DNS si
intendono spesso i server che forniscono su internet questo servizio, mentre
nel nostro caso affronteremo la problematica dal lato client, di un qualunque
programma che necessita di compiere questa operazione.

\begin{figure}[!htb]
  \centering \includegraphics[width=11cm]{img/resolver}
  \caption{Schema di funzionamento delle funzioni del \textit{resolver}.}
  \label{fig:sock_resolver_schema}
\end{figure}

Inoltre quella fra nomi a dominio e indirizzi IP non è l'unica corrispondenza
possibile fra nomi simbolici e valori numerici, come abbiamo visto anche in
sez.~\ref{sec:sys_user_group} per le corrispondenze fra nomi di utenti e
gruppi e relativi identificatori numerici; per quanto riguarda però tutti i
nomi associati a identificativi o servizi relativi alla rete il servizio di
risoluzione è gestito in maniera unificata da un insieme di funzioni fornite
con le librerie del C, detto appunto \textit{resolver}.

Lo schema di funzionamento del \textit{resolver} è illustrato in
fig.~\ref{fig:sock_resolver_schema}; in sostanza i programmi hanno a
disposizione un insieme di funzioni di libreria con cui chiamano il
\textit{resolver}, indicate con le frecce nere. Ricevuta la richiesta è
quest'ultimo che, sulla base della sua configurazione, esegue le operazioni
necessarie a fornire la risposta, che possono essere la lettura delle
informazioni mantenute nei relativi dei file statici presenti sulla macchina,
una interrogazione ad un DNS (che a sua volta, per il funzionamento del
protocollo, può interrogarne altri) o la richiesta ad altri server per i quali
sia fornito il supporto, come LDAP.\footnote{la sigla LDAP fa riferimento ad
  un protocollo, il \textit{Lightweight Directory Access Protocol}, che
  prevede un meccanismo per la gestione di \textsl{elenchi} di informazioni
  via rete; il contenuto di un elenco può essere assolutamente generico, e
  questo permette il mantenimento dei più vari tipi di informazioni su una
  infrastruttura di questo tipo.}

La configurazione del \textit{resolver} attiene più alla amministrazione di
sistema che alla programmazione, ciò non di meno, prima di trattare le varie
funzioni di librerie utilizzate dai programmi, vale la pena fare una
panoramica generale.  Originariamente la configurazione del \textit{resolver}
riguardava esclusivamente le questioni relative alla gestione dei nomi a
dominio, e prevedeva solo l'utilizzo del DNS e del file statico
\conffile{/etc/hosts}.

Per questo aspetto il file di configurazione principale del sistema è
\conffile{/etc/resolv.conf} che contiene in sostanza l'elenco degli indirizzi
IP dei server DNS da contattare; a questo si affiancava (fino alla \acr{glibc}
2.4) il file \conffile{/etc/host.conf} il cui scopo principale era indicare
l'ordine in cui eseguire la risoluzione dei nomi (se usare prima i valori di
\conffile{/etc/hosts} o quelli del DNS). Tralasciamo i dettagli relativi alle
varie direttive che possono essere usate in questi file, che si trovano nelle
rispettive pagine di manuale.

Con il tempo però è divenuto possibile fornire diversi sostituti per
l'utilizzo delle associazione statiche in \conffile{/etc/hosts}, inoltre oltre
alla risoluzione dei nomi a dominio ci sono anche altri nomi da risolvere,
come quelli che possono essere associati ad una rete (invece che ad una
singola macchina) o ai gruppi di macchine definiti dal servizio
NIS,\footnote{il \textit{Network Information Service} è un servizio, creato da
  Sun, e poi diffuso su tutte le piattaforme unix-like, che permette di
  raggruppare all'interno di una rete (in quelli che appunto vengono chiamati
  \textit{netgroup}) varie macchine, centralizzando i servizi di definizione
  di utenti e gruppi e di autenticazione, oggi è sempre più spesso sostituito
  da LDAP.} o come quelli dei protocolli e dei servizi che sono mantenuti nei
file statici \conffile{/etc/protocols} e \conffile{/etc/services}.  

Molte di queste informazioni non si trovano su un DNS, ma in una rete locale
può essere molto utile centralizzare il mantenimento di alcune di esse su
opportuni server.  Inoltre l'uso di diversi supporti possibili per le stesse
informazioni (ad esempio il nome delle macchine può essere mantenuto sia
tramite \conffile{/etc/hosts}, che con il DNS, che con NIS) comporta il
problema dell'ordine in cui questi vengono interrogati. Con le implementazioni
classiche i vari supporti erano introdotti modificando direttamente le
funzioni di libreria, prevedendo un ordine di interrogazione predefinito e non
modificabile (a meno di una ricompilazione delle librerie stesse).

\itindbeg{Name~Service~Switch~(NSS)} 

Per risolvere questa serie di problemi la risoluzione dei nomi a dominio
eseguità dal \textit{resolver} è stata inclusa all'interno di un meccanismo
generico per la risoluzione di corrispondenze fra nomi ed informazioni ad essi
associate chiamato \textit{Name Service Switch}, cui abbiamo accennato anche in
sez.~\ref{sec:sys_user_group} per quanto riguarda la gestione dei dati
associati a utenti e gruppi. Il sistema è stato introdotto la prima volta
nella libreria standard di Solaris e la \acr{glibc} ha ripreso lo stesso
schema; si tenga presente che questo sistema non esiste per altre librerie
standard come la \acr{libc5} o la \acr{uClib}.

Il \textit{Name Service Switch} (cui spesso si fa riferimento con l'acronimo
NSS) è un sistema di librerie dinamiche che permette di definire in maniera
generica sia i supporti su cui mantenere i dati di corrispondenza fra nomi e
valori numerici, sia l'ordine in cui effettuare le ricerche sui vari supporti
disponibili. Il sistema prevede una serie di possibili classi di
corrispondenza, quelle attualmente definite sono riportate in
tab.~\ref{tab:sys_NSS_classes}.

\begin{table}[htb]
  \footnotesize
  \centering
  \begin{tabular}[c]{|l|p{8cm}|}
    \hline
    \textbf{Classe} & \textbf{Tipo di corrispondenza}\\
    \hline
    \hline
    \texttt{passwd}   & Corrispondenze fra nome dell'utente e relative
                        proprietà (\ids{UID}, gruppo principale, ecc.).\\  
    \texttt{shadow}   & Corrispondenze fra username e password dell'utente
                        (e altre informazioni relative alle password).\\  
    \texttt{group}    & Corrispondenze fra nome del gruppo e proprietà dello 
                        stesso.\\  
    \texttt{aliases}  & Alias per la posta elettronica.\\ 
    \texttt{ethers}   & Corrispondenze fra numero IP e MAC address della
                        scheda di rete.\\ 
    \texttt{hosts}    & Corrispondenze fra nome a dominio e numero IP.\\ 
    \texttt{netgroup} & Corrispondenze fra gruppo di rete e macchine che lo
                        compongono.\\  
    \texttt{networks} & Corrispondenze fra nome di una rete e suo indirizzo
                        IP.\\  
    \texttt{protocols}& Corrispondenze fra nome di un protocollo e relativo
                        numero identificativo.\\ 
    \texttt{rpc}      & Corrispondenze fra nome di un servizio RPC e relativo 
                        numero identificativo.\\ 
    \texttt{publickey}& Chiavi pubbliche e private usate per gli RFC sicuri,
                        utilizzate da NFS e NIS+. \\ 
    \texttt{services} & Corrispondenze fra nome di un servizio e numero di
                        porta. \\ 
    \hline
  \end{tabular}
  \caption{Le diverse classi di corrispondenze definite
    all'interno del \textit{Name Service Switch}.} 
  \label{tab:sys_NSS_classes}
\end{table}

% TODO rivedere meglio la tabella

Il sistema del \textit{Name Service Switch} è controllato dal contenuto del
file \conffile{/etc/nsswitch.conf}; questo contiene una riga di configurazione
per ciascuna di queste classi, che viene inizia col nome di
tab.~\ref{tab:sys_NSS_classes} seguito da un carattere ``\texttt{:}'' e
prosegue con la lista dei \textsl{servizi} su cui le relative informazioni
sono raggiungibili, scritti nell'ordine in cui si vuole siano interrogati.
Pertanto nelle versioni recenti delle librerie è questo file e non
\conffile{/etc/host.conf}  a indicare l'ordine con cui si esegue la
risoluzione dei nomi.

Ogni servizio è specificato a sua volta da un nome, come \texttt{file},
\texttt{dns}, \texttt{db}, ecc.  che identifica la libreria dinamica che
realizza l'interfaccia con esso. Per ciascun servizio se \texttt{NAME} è il
nome utilizzato dentro \conffile{/etc/nsswitch.conf}, dovrà essere presente
(usualmente in \file{/lib}) una libreria \texttt{libnss\_NAME} che ne
implementa le funzioni.

In ogni caso, qualunque sia la modalità con cui ricevono i dati o il supporto
su cui vengono mantenuti, e che si usino o meno funzionalità aggiuntive
fornite dal sistema del \textit{Name Service Switch}, dal punto di vista di un
programma che deve effettuare la risoluzione di un nome a dominio, tutto
quello che conta sono le funzioni classiche che il \textit{resolver} mette a
disposizione (è cura della \acr{glibc} tenere conto della presenza del
\textit{Name Service Switch}) e sono queste quelle che tratteremo nelle
sezioni successive.  

\itindend{Name~Service~Switch~(NSS)}


\subsection{Le funzioni di interrogazione del DNS}
\label{sec:sock_resolver_functions}

Prima di trattare le funzioni usate normalmente nella risoluzione dei nomi a
dominio conviene trattare in maniera più dettagliata il servizio DNS. Come
accennato questo, benché esso in teoria sia solo uno dei possibili supporti su
cui mantenere le informazioni, in pratica costituisce il meccanismo principale
con cui vengono risolti i nomi a dominio. Inolte esso può fornire anche
ulteriori informazioni oltre relative alla risoluzione dei nomi a dominio.
Per questo motivo esistono una serie di funzioni di libreria che servono
specificamente ad eseguire delle interrogazioni verso un server DNS, funzioni
che poi vengono utilizzate anche per realizzare le funzioni generiche di
libreria usate dal sistema del \textit{resolver}.

Il sistema del DNS è in sostanza di un database distribuito organizzato in
maniera gerarchica, i dati vengono mantenuti in tanti server distinti ciascuno
dei quali si occupa della risoluzione del proprio \textsl{dominio}; i nomi a
dominio sono organizzati in una struttura ad albero analoga a quella
dell'albero dei file, con domini di primo livello (come i \texttt{.org}),
secondo livello (come \texttt{.truelite.it}), ecc.  In questo caso le
separazioni sono fra i vari livelli sono definite dal carattere ``\texttt{.}''
ed i nomi devono essere risolti da destra verso sinistra.\footnote{per chi si
  stia chiedendo quale sia la radice di questo albero, cioè l'equivalente di
  ``\texttt{/}'', la risposta è il dominio speciale ``\texttt{.}'', che in
  genere non viene mai scritto esplicitamente, ma che, come chiunque abbia
  configurato un server DNS sa bene, esiste ed è gestito dai cosiddetti
  \textit{root DNS} che risolvono i domini di primo livello.} Il meccanismo
funziona con il criterio della \textsl{delegazione}, un server responsabile
per un dominio di primo livello può delegare la risoluzione degli indirizzi
per un suo dominio di secondo livello ad un altro server, il quale a sua volta
potrà delegare la risoluzione di un eventuale sotto-dominio di terzo livello ad
un altro server ancora.

Come accennato un server DNS è in grado di fare molto altro rispetto alla
risoluzione di un nome a dominio in un indirizzo IP: ciascuna voce nel
database viene chiamata \textit{resource record}, e può contenere diverse
informazioni. In genere i \textit{resource record} vengono classificati per la
\textsl{classe di indirizzi} cui i dati contenuti fanno riferimento, e per il
\textsl{tipo} di questi ultimi (ritroveremo classi di indirizzi e tipi di
record più avanti in tab.~\ref{tab:DNS_address_class} e
tab.~\ref{tab:DNS_record_type}).  Oggigiorno i dati mantenuti nei server DNS
sono quasi esclusivamente relativi ad indirizzi internet, per cui in pratica
viene utilizzata soltanto una classe di indirizzi; invece le corrispondenze
fra un nome a dominio ed un indirizzo IP sono solo uno fra i vari tipi di
informazione che un server DNS fornisce normalmente.

L'esistenza di vari tipi di informazioni è un altro dei motivi per cui il
\textit{resolver} prevede, oltre a quelle relative alla semplice risoluzione
dei nomi, un insieme di funzioni specifiche dedicate all'interrogazione di un
server DNS, tutte nella forma \texttt{res\_}\textsl{\texttt{nome}}. La prima
di queste funzioni è \funcd{res\_init}, il cui prototipo è:

\begin{funcproto}{
\fhead{netinet/in.h} 
\fhead{arpa/nameser.h} 
\fhead{resolv.h}
\fdecl{int res\_init(void)}
\fdesc{Inizializza il sistema del \textit{resolver}.} 
}
{La funzione ritorna $0$ in caso di successo e $-1$ per un errore.
}
\end{funcproto}

La funzione legge il contenuto dei file di configurazione per impostare il
dominio di default, gli indirizzi dei server DNS da contattare e l'ordine
delle ricerche; se non sono specificati server verrà utilizzato l'indirizzo
locale, e se non è definito un dominio di default sarà usato quello associato
con l'indirizzo locale (ma questo può essere sovrascritto con l'uso della
variabile di ambiente \envvar{LOCALDOMAIN}). In genere non è necessario
eseguire questa funzione esplicitamente, in quanto viene automaticamente
chiamata la prima volta che si esegue una qualunque delle altre.

Le impostazioni e lo stato del \textit{resolver} inizializzati da
\func{res\_init} vengono mantenuti in una serie di variabili raggruppate nei
campi di una apposita struttura. Questa struttura viene definita in
\headfiled{resolv.h} e mantenuta nella variabile globale \var{\_res}, che
viene utilizzata internamente da tutte le funzioni dell'interfaccia. Questo
consente anche di accedere direttamente al contenuto della variabile
all'interno di un qualunque programma, una volta che la sia opportunamente
dichiarata con:
\includecodesnip{listati/resolv_option.c}

Dato che l'uso di una variabile globale rende tutte le funzioni
dell'interfaccia classica non rientranti, queste sono state deprecate in
favore di una nuova interfaccia in cui esse sono state sostituite da
altrettante nuove funzioni, il cui nome è ottenuto apponendo una
``\texttt{n}'' al nome di quella tradizionale (cioè nella forma
\texttt{res\_n\textsl{nome}}). Tutte le nuove funzioni sono identiche alle
precedenti, ma hanno un primo argomento aggiuntivo, \param{statep}, puntatore
ad una struttura dello stesso tipo di \var{\_res}. Questo consente di usare
una variabile locale per mantenere lo stato del \textit{resolver}, rendendo le
nuove funzioni rientranti.  In questo caso per poter utilizzare il nuovo
argomento occorrerà una opportuna dichiarazione del relativo tipo di dato con:
\includecodesnip{listati/resolv_newoption.c}

Così la nuova funzione utilizzata per inizializzare il \textit{resolver} (che
come la precedente viene chiamata automaticamente da tutte altre funzioni) è
\funcd{res\_ninit}, ed il suo prototipo è:

\begin{funcproto}{
\fhead{netinet/in.h} 
\fhead{arpa/nameser.h} 
\fhead{resolv.h}
\fdecl{int res\_ninit(res\_state statep)}
\fdesc{Inizializza il sistema del \textit{resolver}.} 
}
{La funzione ritorna $0$ in caso di successo e $-1$ per un errore.
}
\end{funcproto}

Indipendentemente da quale versione delle funzioni si usino, tutti i campi
della struttura (\var{\_res} o la variabile puntata da \param{statep}) sono ad
uso interno, e vengono usualmente inizializzate da \func{res\_init} o
\func{res\_ninit} in base al contenuto dei file di configurazione e ad una
serie di valori di default. L'unico campo che può essere utile modificare è
\var{\_res.options} (o l'equivalente della variabile puntata
da\param{statep}), una maschera binaria che contiene una serie di bit che
esprimono le opzioni che permettono di controllare il comportamento del
\textit{resolver}.

\begin{table}[htb]
  \centering
  \footnotesize
  \begin{tabular}[c]{|l|p{8cm}|}
    \hline
    \textbf{Costante} & \textbf{Significato} \\
    \hline
    \hline
    \constd{RES\_INIT}       & Viene attivato se è stata chiamata
                               \func{res\_init}. \\
    \constd{RES\_DEBUG}      & Stampa dei messaggi di debug.\\
    \constd{RES\_AAONLY}     & Accetta solo risposte autoritative.\\
    \constd{RES\_USEVC}      & Usa connessioni TCP per contattare i server 
                               invece che l'usuale UDP.\\
    \constd{RES\_PRIMARY}    & Interroga soltanto server DNS primari.\\
    \constd{RES\_IGNTC}      & Ignora gli errori di troncamento, non ritenta la
                               richiesta con una connessione TCP.\\
    \constd{RES\_RECURSE}    & Imposta il bit che indica che si desidera
                               eseguire una interrogazione ricorsiva.\\
    \constd{RES\_DEFNAMES}   & Se attivo \func{res\_search} aggiunge il nome
                               del dominio di default ai nomi singoli (che non
                               contengono cioè un ``\texttt{.}'').\\
    \constd{RES\_STAYOPEN}   & Usato con \const{RES\_USEVC} per mantenere
                               aperte le connessioni TCP fra interrogazioni
                               diverse.\\
    \constd{RES\_DNSRCH}     & Se attivo \func{res\_search} esegue le ricerche
                               di nomi di macchine nel dominio corrente o nei
                               domini ad esso sovrastanti.\\
    \constd{RES\_INSECURE1}  & Blocca i controlli di sicurezza di tipo 1.\\
    \constd{RES\_INSECURE2}  & Blocca i controlli di sicurezza di tipo 2.\\
    \constd{RES\_NOALIASES}  & Blocca l'uso della variabile di ambiente
                               \envvar{HOSTALIASES}.\\ 
    \constd{RES\_USE\_INET6} & Restituisce indirizzi IPv6 con
                               \func{gethostbyname}. \\
    \constd{RES\_ROTATE}     & Ruota la lista dei server DNS dopo ogni
                               interrogazione.\\
    \constd{RES\_NOCHECKNAME}& Non controlla i nomi per verificarne la
                               correttezza sintattica. \\
    \constd{RES\_KEEPTSIG}   & Non elimina i record di tipo \texttt{TSIG}.\\
    \constd{RES\_BLAST}      & Effettua un ``\textit{blast}'' inviando
                               simultaneamente le richieste a tutti i server;
                               non ancora implementata. \\
    \constd{RES\_DEFAULT}    & Combinazione di \const{RES\_RECURSE},
                               \const{RES\_DEFNAMES} e \const{RES\_DNSRCH}.\\
    \hline
  \end{tabular}
  \caption{Costanti utilizzabili come valori per \var{\_res.options}.}
  \label{tab:resolver_option}
\end{table}

Per utilizzare questa funzionalità per modificare le impostazioni direttamente
da programma occorrerà impostare un opportuno valore per questo campo ed
invocare esplicitamente \func{res\_init} o \func{res\_ninit}, dopo di che le
altre funzioni prenderanno le nuove impostazioni. Le costanti che definiscono
i vari bit di questo campo, ed il relativo significato sono illustrate in
tab.~\ref{tab:resolver_option}; trattandosi di una maschera binaria un valore
deve essere espresso con un opportuno OR aritmetico di dette costanti; ad
esempio il valore di default delle opzioni, espresso dalla costante
\const{RES\_DEFAULT}, è definito come:
\includecodesnip{listati/resolv_option_def.c}

Non tratteremo il significato degli altri campi non essendovi necessità di
modificarli direttamente; gran parte di essi sono infatti impostati dal
contenuto dei file di configurazione, mentre le funzionalità controllate da
alcuni di esse possono essere modificate con l'uso delle opportune variabili
di ambiente come abbiamo visto per \envvar{LOCALDOMAIN}. In particolare con
\envvar{RES\_RETRY} si soprassiede il valore del campo \var{retry} che
controlla quante volte viene ripetuto il tentativo di connettersi ad un server
DNS prima di dichiarare fallimento; il valore di default è 4, un valore nullo
significa bloccare l'uso del DNS. Infine con \envvar{RES\_TIMEOUT} si
soprassiede il valore del campo \var{retrans} (preimpostato al valore della
omonima costante \const{RES\_TIMEOUT} di \headfile{resolv.h}) che è il valore
preso come base (in numero di secondi) per definire la scadenza di una
richiesta, ciascun tentativo di richiesta fallito viene ripetuto raddoppiando
il tempo di scadenza per il numero massimo di volte stabilito da
\texttt{RES\_RETRY}.

La funzione di interrogazione principale è \funcd{res\_query}
(\funcd{res\_nquery} per la nuova interfaccia), che serve ad eseguire una
richiesta ad un server DNS per un nome a dominio \textsl{completamente
  specificato} (quello che si chiama
\itindex{Fully~Qualified~Domain~Name~(FQDN)} FQDN, \textit{Fully Qualified 
  Domain Name}); il loro prototipo è:

\begin{funcproto}{
\fhead{netinet/in.h} 
\fhead{arpa/nameser.h} 
\fhead{resolv.h}
\fdecl{int res\_query(const char *dname, int class, int type,
              unsigned char *answer, int anslen)}
\fdecl{int res\_nquery(res\_state statep, const char *dname, int class, int
  type, \\
  \phantom{int res\_nquery(}unsigned char *answer, int anslen)}
\fdesc{Esegue una interrogazione al DNS.} 
}
{Le funzioni ritornano un valore positivo pari alla lunghezza dei dati scritti
  nel buffer \param{answer} in caso di successo e $-1$ per un errore.
}
\end{funcproto}

Le funzioni eseguono una interrogazione ad un server DNS relativa al nome da
risolvere passato nella stringa indirizzata da \param{dname}, inoltre deve
essere specificata la classe di indirizzi in cui eseguire la ricerca con
\param{class}, ed il tipo di \textit{resource record} che si vuole ottenere
con \param{type}. Il risultato della ricerca verrà scritto nel buffer di
lunghezza \param{anslen} puntato da \param{answer} che si sarà opportunamente
allocato in precedenza.

Una seconda funzione di ricerca analoga a \func{res\_query}, che prende gli
stessi argomenti ma che esegue l'interrogazione con le funzionalità
addizionali previste dalle due opzioni \const{RES\_DEFNAMES} e
\const{RES\_DNSRCH}, è \funcd{res\_search} (\funcd{res\_nsearch} per la nuova
interfaccia), il cui prototipo è:

\begin{funcproto}{
\fhead{netinet/in.h} 
\fhead{arpa/nameser.h} 
\fhead{resolv.h}
\fdecl{int res\_search(const char *dname, int class, int type,
  unsigned char *answer, \\
  \phantom{int res\_search}int anslen)}
\fdecl{int res\_nsearch(res\_state statep, const char *dname, int class, 
  int type, \\
  \phantom{int res\_nsearch(}unsigned char *answer, int anslen)}
\fdesc{Esegue una interrogazione al DNS.} 
}
{Le funzioni ritornano un valore positivo pari alla lunghezza dei dati scritti
  nel buffer \param{answer} in caso di successo e $-1$ per un errore.
}
\end{funcproto}

In sostanza la funzione ripete una serie di chiamate a \func{res\_query}
(\func{res\_nquery}) aggiungendo al nome contenuto nella stringa \param{dname}
il dominio di default da cercare, fermandosi non appena trova un risultato.
Il risultato di entrambe le funzioni viene scritto nel formato opportuno (che
sarà diverso a seconda del tipo di record richiesto) nel buffer di ritorno;
sarà compito del programma (o di altre funzioni) estrarre i relativi dati,
esistono una serie di funzioni interne usate per la scansione di questi dati,
per chi fosse interessato una trattazione dettagliata è riportata nel
quattordicesimo capitolo di \cite{DNSbind}.

Le classi di indirizzi supportate da un server DNS sono tre, ma di queste in
pratica oggi viene utilizzata soltanto quella degli indirizzi internet; le
costanti che identificano dette classi, da usare come valore per l'argomento
\param{class} delle precedenti funzioni, sono riportate in
tab.~\ref{tab:DNS_address_class} (esisteva in realtà anche una classe
\constd{C\_CSNET} per la omonima rete, ma è stata dichiarata obsoleta).

\begin{table}[htb]
  \centering
  \footnotesize
  \begin{tabular}[c]{|l|p{8cm}|}
    \hline
    \textbf{Costante} & \textbf{Significato} \\
    \hline
    \hline
    \constd{C\_IN}   & Indirizzi internet, in pratica i soli utilizzati oggi.\\
    \constd{C\_HS}   & Indirizzi \textit{Hesiod}, utilizzati solo al MIT, oggi
                       completamente estinti. \\
    \constd{C\_CHAOS}& Indirizzi per la rete \textit{Chaosnet}, un'altra rete
                       sperimentale nata al MIT. \\
    \constd{C\_ANY}  & Indica un indirizzo di classe qualunque.\\
    \hline
  \end{tabular}
  \caption{Costanti identificative delle classi di indirizzi per l'argomento
    \param{class} di \func{res\_query}.}
  \label{tab:DNS_address_class}
\end{table}

Come accennato le tipologie di dati che sono mantenibili su un server DNS sono
diverse, ed a ciascuna di essa corrisponde un diverso tipo di \textit{resource
  record}. L'elenco delle costanti, ripreso dai file di dichiarazione
\headfiled{arpa/nameser.h} e \headfiled{arpa/nameser\_compat.h}, che
definiscono i valori che si possono usare per l'argomento \param{type} per
specificare il tipo di \textit{resource record} da richiedere è riportato in
tab.~\ref{tab:DNS_record_type}; le costanti (tolto il \texttt{T\_} iniziale)
hanno gli stessi nomi usati per identificare i record nei file di zona di
BIND,\footnote{BIND, acronimo di \textit{Berkley Internet Name Domain}, è una
  implementazione di un server DNS, ed, essendo utilizzata nella stragrande
  maggioranza dei casi, fa da riferimento; i dati relativi ad un certo dominio
  (cioè i suoi \textit{resource record} vengono mantenuti in quelli che sono
  usualmente chiamati \textsl{file di zona}, e in essi ciascun tipo di dominio
  è identificato da un nome che è appunto identico a quello delle costanti di
  tab.~\ref{tab:DNS_record_type} senza il \texttt{T\_} iniziale.} e che
normalmente sono anche usati come nomi per indicare i record.

\begin{table}[!htb]
  \centering
  \footnotesize
  \begin{tabular}[c]{|l|l|}
    \hline
    \textbf{Costante} & \textbf{Significato} \\
    \hline
    \hline
    \constd{T\_A}     & Indirizzo di una stazione.\\
    \constd{T\_NS}    & Server DNS autoritativo per il dominio richiesto.\\
    \constd{T\_MD}    & Destinazione per la posta elettronica.\\
    \constd{T\_MF}    & Redistributore per la posta elettronica.\\
    \constd{T\_CNAME} & Nome canonico.\\
    \constd{T\_SOA}   & Inizio di una zona di autorità.\\
    \constd{T\_MB}    & Nome a dominio di una casella di posta.\\
    \constd{T\_MG}    & Nome di un membro di un gruppo di posta.\\
    \constd{T\_MR}    & Nome di un cambiamento di nome per la posta.\\
    \constd{T\_NULL}  & Record nullo.\\
    \constd{T\_WKS}   & Servizio noto.\\
    \constd{T\_PTR}   & Risoluzione inversa di un indirizzo numerico.\\
    \constd{T\_HINFO} & Informazione sulla stazione.\\
    \constd{T\_MINFO} & Informazione sulla casella di posta.\\
    \constd{T\_MX}    & Server cui instradare la posta per il dominio.\\
    \constd{T\_TXT}   & Stringhe di testo (libere).\\
    \constd{T\_RP}    & Nome di un responsabile (\textit{responsible person}).\\
    \constd{T\_AFSDB} & Database per una cella AFS.\\
    \constd{T\_X25}   & Indirizzo di chiamata per X.25.\\
    \constd{T\_ISDN}  & Indirizzo di chiamata per ISDN.\\
    \constd{T\_RT}    & Router.\\
    \constd{T\_NSAP}  & Indirizzo NSAP.\\
    \constd{T\_NSAP\_PTR}& Risoluzione inversa per NSAP (deprecato).\\
    \constd{T\_SIG}   & Firma digitale di sicurezza.\\
    \constd{T\_KEY}   & Chiave per firma.\\
    \constd{T\_PX}    & Corrispondenza per la posta X.400.\\
    \constd{T\_GPOS}  & Posizione geografica.\\
    \constd{T\_AAAA}  & Indirizzo IPv6.\\
    \constd{T\_LOC}   & Informazione di collocazione.\\
    \constd{T\_NXT}   & Dominio successivo.\\
    \constd{T\_EID}   & Identificatore di punto conclusivo.\\
    \constd{T\_NIMLOC}& Posizionatore \textit{nimrod}.\\
    \constd{T\_SRV}   & Servizio.\\
    \constd{T\_ATMA}  & Indirizzo ATM.\\
    \constd{T\_NAPTR} & Puntatore ad una \textit{naming authority}.\\
    \constd{T\_TSIG}  & Firma di transazione.\\
    \constd{T\_IXFR}  & Trasferimento di zona incrementale.\\
    \constd{T\_AXFR}  & Trasferimento di zona di autorità.\\
    \constd{T\_MAILB} & Trasferimento di record di caselle di posta.\\
    \constd{T\_MAILA} & Trasferimento di record di server di posta.\\
    \constd{T\_ANY}   & Valore generico.\\
    \hline
  \end{tabular}
  \caption{Costanti identificative del tipo di record per l'argomento
    \param{type} di \func{res\_query}.}
  \label{tab:DNS_record_type}
\end{table}


L'elenco di tab.~\ref{tab:DNS_record_type} è quello di \textsl{tutti} i
\textit{resource record} definiti, con una breve descrizione del relativo
significato.  Di tutti questi però viene impiegato correntemente solo un
piccolo sottoinsieme, alcuni sono obsoleti ed altri fanno riferimento a dati
applicativi che non ci interessano non avendo nulla a che fare con la
risoluzione degli indirizzi IP, pertanto non entreremo nei dettagli del
significato di tutti i \textit{resource record}, ma solo di quelli usati dalle
funzioni del \textit{resolver}. Questi sono sostanzialmente i seguenti (per
indicarli si è usata la notazione dei file di zona di BIND):
\begin{basedescript}{\desclabelwidth{1.2cm}\desclabelstyle{\nextlinelabel}}
\item[\texttt{A}] viene usato per indicare la corrispondenza fra un nome a
  dominio ed un indirizzo IPv4; ad esempio la corrispondenza fra
  \texttt{jojo.truelite.it} e l'indirizzo IP \texttt{62.48.34.25}.
\item[\texttt{AAAA}] viene usato per indicare la corrispondenza fra un nome a
  dominio ed un indirizzo IPv6; è chiamato in questo modo dato che la
  dimensione di un indirizzo IPv6 è quattro volte quella di un indirizzo IPv4.
\item[\texttt{PTR}] per fornire la corrispondenza inversa fra un indirizzo IP
  ed un nome a dominio ad esso associato si utilizza questo tipo di record (il
  cui nome sta per \textit{pointer}).
\item[\texttt{CNAME}] qualora si abbiamo più nomi che corrispondono allo
  stesso indirizzo (come ad esempio \texttt{www.truelite.it} e
  \texttt{sources.truelite.it}, che fanno entrambi riferimento alla stessa
  macchina (nel caso \texttt{dodds.truelite.it}) si può usare questo tipo di
  record per creare degli \textit{alias} in modo da associare un qualunque
  altro nome al \textsl{nome canonico} della macchina (si chiama così quello
  associato al record \texttt{A}).
\end{basedescript}

Come accennato in caso di successo le due funzioni di richiesta restituiscono
il risultato della interrogazione al server, in caso di insuccesso l'errore
invece viene segnalato da un valore di ritorno pari a $-1$, ma in questo caso,
non può essere utilizzata la variabile \var{errno} per riportare un codice di
errore, in quanto questo viene impostato per ciascuna delle chiamate al
sistema utilizzate dalle funzioni del \textit{resolver}, non avrà alcun
significato nell'indicare quale parte del procedimento di risoluzione è
fallita.

Per questo motivo è stata definita una variabile di errore separata,
\var{h\_errno}, che viene utilizzata dalle funzioni del \textit{resolver} per
indicare quale problema ha causato il fallimento della risoluzione del nome.
Ad essa si può accedere una volta che la si dichiara con:
\includecodesnip{listati/herrno.c} 
ed i valori che può assumere, con il relativo significato, sono riportati in
tab.~\ref{tab:h_errno_values}.

\begin{table}[!htb]
  \centering
  \footnotesize
  \begin{tabular}[c]{|l|p{9cm}|}
    \hline
    \textbf{Costante} & \textbf{Significato} \\
    \hline
    \hline
    \constd{HOST\_NOT\_FOUND}& L'indirizzo richiesto non è valido e la
                               macchina indicata è sconosciuta.\\
    \constd{NO\_ADDRESS}     & Il nome a dominio richiesto è valido, ma non ha
                               un indirizzo associato ad esso
                               (alternativamente può essere indicato come 
                               \constd{NO\_DATA}).\\
    \constd{NO\_RECOVERY}    & Si è avuto un errore non recuperabile
                               nell'interrogazione di un server DNS.\\
    \constd{TRY\_AGAIN}      & Si è avuto un errore temporaneo
                               nell'interrogazione di un server DNS, si può
                               ritentare l'interrogazione in un secondo
                               tempo.\\
    \hline
  \end{tabular}
  \caption{Valori possibili della variabile \var{h\_errno}.}
  \label{tab:h_errno_values}
\end{table}

Insieme alla nuova variabile vengono definite anche delle nuove funzioni per
stampare l'errore a video, analoghe a quelle di sez.~\ref{sec:sys_strerror}
per \var{errno}, ma che usano il valore di \var{h\_errno}; la prima è
\funcd{herror} ed il suo prototipo è:

{\centering
\vspace{3pt}
\begin{funcbox}{
\fhead{netdb.h}
\fdecl{void herror(const char *string)}
\fdesc{Stampa un errore di risoluzione.} 
}
\end{funcbox}}

La funzione è l'analoga di \func{perror} e stampa sullo \textit{standard error} un
messaggio di errore corrispondente al valore corrente di \var{h\_errno}, a cui
viene anteposta la stringa \param{string} passata come argomento.  La seconda
funzione è \funcd{hstrerror} ed il suo prototipo è:

{\centering
\vspace{3pt}
\begin{funcbox}{
\fhead{netdb.h}
\fdecl{const char *hstrerror(int err)}
\fdesc{Restituisce una stringa corrispondente ad un errore di risoluzione.}
}
\end{funcbox}}

\noindent che, come  l'analoga \func{strerror}, restituisce una stringa con un
messaggio di errore già formattato, corrispondente al codice passato come
argomento (che si presume sia dato da \var{h\_errno}).

\itindend{resolver}


\subsection{La vecchia interfaccia per la risoluzione dei nomi a dominio}
\label{sec:sock_name_services}

La principale funzionalità del \textit{resolver} resta quella di risolvere i
nomi a dominio in indirizzi IP, per cui non ci dedicheremo oltre alle funzioni
per effetture delle richieste generiche al DNS ed esamineremo invece le
funzioni del \textit{resolver} dedicate specificamente a questo. Tratteremo in
questa sezione l'interfaccia tradizionale, che ormai è deprecata, mentre
vedremo nella sezione seguente la nuova interfaccia.

La prima funzione dell'interfaccia tradizionale è \funcd{gethostbyname} il cui
scopo è ottenere l'indirizzo di una stazione noto il suo nome a dominio, il
suo prototipo è:

\begin{funcproto}{
\fhead{netdb.h}
\fdecl{struct hostent *gethostbyname(const char *name)}
\fdesc{Determina l'indirizzo associato ad un nome a dominio.} 
}

{La funzione ritorna il puntatore ad una
  struttura di tipo \struct{hostent} contenente i dati associati al nome a
  dominio in caso di successo o un puntatore nullo per un errore.}
\end{funcproto}

La funzione prende come argomento una stringa \param{name} contenente il nome
a dominio che si vuole risolvere, in caso di successo i dati ad esso relativi
vengono memorizzati in una opportuna struttura \struct{hostent} la cui
definizione è riportata in fig.~\ref{fig:sock_hostent_struct}. 

\begin{figure}[!htb]
  \footnotesize \centering
  \begin{minipage}[c]{0.80\textwidth}
    \includestruct{listati/hostent.h}
  \end{minipage}
  \caption{La struttura \structd{hostent} per la risoluzione dei nomi a
    dominio e degli indirizzi IP.}
  \label{fig:sock_hostent_struct}
\end{figure}

Quando un programma chiama \func{gethostbyname} e questa usa il DNS per
effettuare la risoluzione del nome, è con i valori contenuti nei relativi
record che vengono riempite le varie parti della struttura \struct{hostent}.
Il primo campo della struttura, \var{h\_name} contiene sempre il \textsl{nome
  canonico}, che nel caso del DNS è appunto il nome associato ad un record
\texttt{A}. Il secondo campo della struttura, \var{h\_aliases}, invece è un
puntatore ad vettore di puntatori, terminato da un puntatore nullo. Ciascun
puntatore del vettore punta ad una stringa contenente uno degli altri
possibili nomi associati allo stesso \textsl{nome canonico} (quelli che nel
DNS vengono inseriti come record di tipo \texttt{CNAME}).

Il terzo campo della struttura, \var{h\_addrtype}, indica il tipo di indirizzo
che è stato restituito, e può assumere soltanto i valori \const{AF\_INET} o
\const{AF\_INET6}, mentre il quarto campo, \var{h\_length}, indica la
lunghezza dell'indirizzo stesso in byte. 

Infine il campo \var{h\_addr\_list} è il puntatore ad un vettore di puntatori
ai singoli indirizzi; il vettore è terminato da un puntatore nullo.  Inoltre,
come illustrato in fig.~\ref{fig:sock_hostent_struct}, viene definito il campo
\var{h\_addr} come sinonimo di \code{h\_addr\_list[0]}, cioè un riferimento
diretto al primo indirizzo della lista.

Oltre ai normali nomi a dominio la funzione accetta come argomento
\param{name} anche indirizzi numerici, in formato dotted decimal per IPv4 o
con la notazione illustrata in sez.~\ref{sec:IP_ipv6_notation} per IPv6. In
tal caso \func{gethostbyname} non eseguirà nessuna interrogazione remota, ma
si limiterà a copiare la stringa nel campo \var{h\_name} ed a creare la
corrispondente struttura \var{in\_addr} da indirizzare con
\code{h\_addr\_list[0]}.

Con l'uso di \func{gethostbyname} si ottengono solo gli indirizzi IPv4, se si
vogliono ottenere degli indirizzi IPv6 occorrerà prima impostare l'opzione
\const{RES\_USE\_INET6} nel campo \texttt{\_res.options} e poi chiamare
\func{res\_init} (vedi sez.~\ref{sec:sock_resolver_functions}) per modificare
le opzioni del \textit{resolver}; dato che questo non è molto comodo è stata
definita (è una estensione fornita dalla \acr{glibc}, disponibile anche in
altri sistemi unix-like) un'altra funzione, \funcd{gethostbyname2}, il cui
prototipo è:

\begin{funcproto}{
\fhead{netdb.h}
\fhead{sys/socket.h}
\fdecl{struct hostent *gethostbyname2(const char *name, int af)}
\fdesc{Determina l'indirizzo del tipo scelto associato ad un nome a dominio.} 
}

{La funzione ritorna il puntatore ad una
  struttura di tipo \struct{hostent} contenente i dati associati al nome a
  dominio in caso di successo e  un puntatore nullo per un errore.}
\end{funcproto}

In questo caso la funzione prende un secondo argomento \param{af} che indica
quale famiglia di indirizzi è quella che dovrà essere utilizzata per
selezionare i risultati restituiti dalla funzione; i soli valori consentiti
sono \const{AF\_INET} o \const{AF\_INET6} per indicare rispettivamente IPv4 e
IPv6 (per questo è necessario l'uso di \headfile{sys/socket.h}). Per tutto il
resto la funzione è identica a \func{gethostbyname}, ed identici sono i suoi
risultati.

\begin{figure}[!htb]
  \footnotesize \centering
  \begin{minipage}[c]{\codesamplewidth}
    \includecodesample{listati/mygethost.c}
  \end{minipage}
  \normalsize
  \caption{Esempio di codice per la risoluzione di un indirizzo.}
  \label{fig:mygethost_example}
\end{figure}

Vediamo allora un primo esempio dell'uso delle funzioni di risoluzione, in
fig.~\ref{fig:mygethost_example} è riportato un estratto del codice di un
programma che esegue una semplice interrogazione al \textit{resolver} usando
\func{gethostbyname} e poi ne stampa a video i risultati. Al solito il
sorgente completo, che comprende il trattamento delle opzioni ed una funzione
per stampare un messaggio di aiuto, è nel file \texttt{mygethost.c} dei
sorgenti allegati alla guida.

Il programma richiede un solo argomento che specifichi il nome da cercare,
senza il quale (\texttt{\small 15--18}) esce con un errore. Dopo di che
(\texttt{\small 20}) si limita a chiamare \func{gethostbyname}, ricevendo il
risultato nel puntatore \var{data}. Questo (\texttt{\small 21--24}) viene
controllato per rilevare eventuali errori, nel qual caso il programma esce
dopo aver stampato un messaggio con \func{herror}. 

Se invece la risoluzione è andata a buon fine si inizia (\texttt{\small 25})
con lo stampare il nome canonico, dopo di che (\texttt{\small 26--30}) si
stampano eventuali altri nomi. Per questo prima (\texttt{\small 26}) si prende
il puntatore alla cima della lista che contiene i nomi e poi (\texttt{\small
  27--30}) si esegue un ciclo che sarà ripetuto fin tanto che nella lista si
troveranno dei puntatori validi per le stringhe dei nomi (si ricordi che la
lista viene terminata da un puntatore nullo); prima (\texttt{\small 28}) si
stamperà la stringa e poi (\texttt{\small 29}) si provvederà ad incrementare
il puntatore per passare al successivo elemento della lista.

Una volta stampati i nomi si passerà a stampare gli indirizzi, il primo passo
(\texttt{\small 31--38}) è allora quello di riconoscere il tipo di indirizzo
sulla base del valore del campo \var{h\_addrtype}, stampandolo a video. Si è
anche previsto di stampare un errore nel caso (che non dovrebbe mai accadere)
di un indirizzo non valido.

Infine (\texttt{\small 39--44}) si stamperanno i valori degli indirizzi, di
nuovo (\texttt{\small 39}) si inizializzerà un puntatore alla cima della lista
e si eseguirà un ciclo fintanto che questo punterà ad indirizzi validi in
maniera analoga a quanto fatto in precedenza per i nomi a dominio. Si noti
come, essendo il campo \var{h\_addr\_list} un puntatore ad strutture di
indirizzi generiche, questo sia ancora di tipo \texttt{char **} e si possa
riutilizzare lo stesso puntatore usato per i nomi.

Per ciascun indirizzo valido si provvederà (\texttt{\small 41}) a una
conversione con la funzione \func{inet\_ntop} (vedi
sez.~\ref{sec:sock_addr_func}) passandole gli opportuni argomenti, questa
restituirà la stringa da stampare (\texttt{\small 42}) con il valore
dell'indirizzo in \var{buffer}, che si è avuto la cura di dichiarare
inizialmente (\texttt{\small 10}) con dimensioni adeguate; dato che la
funzione è in grado di tenere conto automaticamente del tipo di indirizzo non
ci sono precauzioni particolari da prendere.\footnote{volendo essere pignoli
  si dovrebbe controllarne lo stato di uscita, lo si è tralasciato per non
  appesantire il codice, dato che in caso di indirizzi non validi si sarebbe
  avuto un errore con \func{gethostbyname}, ma si ricordi che la sicurezza non
  è mai troppa.}

Le funzioni illustrate finora hanno un difetto: utilizzano tutte una area di
memoria interna per allocare i contenuti della struttura \struct{hostent} e
per questo non possono essere rientranti. L'uso della memoria interna inoltre
comporta anche che in due successive chiamate i dati potranno essere
sovrascritti.

Si tenga presente poi che copiare il contenuto della sola struttura non è
sufficiente per salvare tutti i dati, in quanto questa contiene puntatori ad
altri dati, che pure possono essere sovrascritti; per questo motivo, se si
vuole salvare il risultato di una chiamata, occorrerà eseguire quella che si
chiama una \itindex{deep~copy} \textit{deep copy}.\footnote{si chiama così
  quella tecnica per cui, quando si deve copiare il contenuto di una struttura
  complessa (con puntatori che puntano ad altri dati, che a loro volta possono
  essere puntatori ad altri dati) si deve copiare non solo il contenuto della
  struttura, ma eseguire una scansione per risolvere anche tutti i puntatori
  contenuti in essa (e così via se vi sono altre sotto-strutture con altri
  puntatori) e copiare anche i dati da questi referenziati.}

Per ovviare a questi problemi nella \acr{glibc} sono definite anche delle
versioni rientranti delle precedenti funzioni, al solito queste sono
caratterizzate dall'avere un suffisso \texttt{\_r}, pertanto avremo le due
funzioni \funcd{gethostbyname\_r} e \funcd{gethostbyname2\_r} i cui prototipi
sono:

\begin{funcproto}{
\fhead{netdb.h}
\fhead{sys/socket.h}
\fdecl{int gethostbyname\_r(const char *name, struct hostent *ret, 
    char *buf, size\_t buflen,\\
\phantom{int gethostbyname\_r(}struct hostent **result, int *h\_errnop)}
\fdecl{int gethostbyname2\_r(const char *name, int af,
         struct hostent *ret, char *buf,\\
\phantom{int gethostbyname2\_r(}size\_t buflen, struct hostent **result, int *h\_errnop)}

\fdesc{Versioni rientranti delle funzioni \func{gethostbyname} e
  \func{gethostbyname2}.} 
}

{Le funzioni ritornano $0$ in caso di successo ed un valore diverso da zero per
  un errore.}
\end{funcproto}

Gli argomenti \param{name} (e \param{af} per \func{gethostbyname2\_r}) hanno
lo stesso significato visto in precedenza. Tutti gli altri argomenti hanno lo
stesso significato per entrambe le funzioni. Per evitare l'uso di variabili
globali si dovrà allocare preventivamente una struttura \struct{hostent} in
cui ricevere il risultato, passandone l'indirizzo alla funzione nell'argomento
\param{ret}.  Inoltre, dato che \struct{hostent} contiene dei puntatori, dovrà
essere allocato anche un buffer in cui le funzioni possano scrivere tutti i
dati del risultato dell'interrogazione da questi puntati; l'indirizzo e la
lunghezza di questo buffer devono essere indicati con gli argomenti
\param{buf} e \param{buflen}.

Gli ultimi due argomenti vengono utilizzati per avere indietro i risultati
come \textit{value result argument}, si deve specificare l'indirizzo della
variabile su cui la funzione dovrà salvare il codice di errore
con \param{h\_errnop} e quello su cui dovrà salvare il puntatore che si userà
per accedere i dati con \param{result}.

In caso di successo entrambe le funzioni restituiscono un valore nullo,
altrimenti restituiscono un valore non nulla e all'indirizzo puntato
da \param{result} sarà salvato un puntatore nullo, mentre a quello puntato da
\param{h\_errnop} sarà salvato il valore del codice di errore, dato che per
essere rientrante la funzione non può la variabile globale \var{h\_errno}. In
questo caso il codice di errore, oltre ai valori di
tab.~\ref{tab:h_errno_values}, può avere anche quello di \errcode{ERANGE}
qualora il buffer allocato su \param{buf} non sia sufficiente a contenere i
dati, in tal caso si dovrà semplicemente ripetere l'esecuzione della funzione
con un buffer di dimensione maggiore.

Una delle caratteristiche delle interrogazioni al servizio DNS è che queste
sono normalmente eseguite con il protocollo UDP, ci sono casi in cui si
preferisce che vengano usate connessioni permanenti con il protocollo TCP. Per
ottenere questo sono previste delle funzioni apposite (si potrebbero impostare
direttamente le opzioni di \var{\_\_res.options}, ma queste funzioni
permettono di semplificare la procedura); la prime sono \funcd{sethostent} e
\func{endhostent}, il cui prototipo è:

\begin{funcproto}{
\fhead{netdb.h}
\fdecl{void sethostent(int stayopen)}
\fdesc{Richiede l'uso di connessioni TCP per le interrogazioni ad un server DNS.} 
\fdecl{void endhostent(void)}
\fdesc{Disattiva l'uso di connessioni TCP per le interrogazioni ad un server DNS.} 
}

{Le funzioni non restituiscono nulla, e non danno errori.}
\end{funcproto}

La funzione \func{sethostent} permette di richiedere l'uso di connessioni TCP
per la richiesta dei dati, e che queste restino aperte per successive
richieste; il valore dell'argomento \param{stayopen} indica se attivare questa
funzionalità, un valore diverso da zero, che indica una condizione vera in C,
attiva la funzionalità.  Per disattivare l'uso delle connessioni TCP si può
invece usare \func{endhostent}, e come si vede la funzione è estremamente
semplice, non richiedendo nessun argomento.

Infine si può richiedere la risoluzione inversa di un indirizzo IP od IPv6,
per ottenerne il nome a dominio ad esso associato, per fare questo si può
usare la funzione \funcd{gethostbyaddr}, il cui prototipo è:

\begin{funcproto}{
\fhead{netdb.h}
\fhead{sys/socket.h} 
\fdecl{struct hostent *gethostbyaddr(const char *addr, int len, int type)}
\fdesc{Richiede la risoluzione inversa di un indirizzo IP.} 
}

{La funzione ritorna l'indirizzo ad una struttura \struct{hostent} in caso di
  successo e \val{NULL} per un errore.}
\end{funcproto}

In questo caso l'argomento \param{addr} dovrà essere il puntatore ad una
appropriata struttura contenente il valore dell'indirizzo IP (o IPv6) che si
vuole risolvere. L'uso del tipo \texttt{char *} per questo argomento è
storico, il dato dovrà essere fornito in una struttura \struct{in\_addr} per
un indirizzo IPv4 ed una struttura \struct{in6\_addr} per un indirizzo IPv6.

Si ricordi inoltre, come illustrato in fig.~\ref{fig:sock_sa_ipv4_struct}, che
mentre \struct{in\_addr} corrisponde in realtà ad un oridinario numero intero
a 32 bit  (da esprimere comunque in \textit{network order}) non altrettanto
avviene per \struct{in6\_addr}, pertanto è sempre opportuno inizializzare
questi indirizzi con \func{inet\_pton} (vedi
sez.~\ref{sec:sock_conv_func_gen}).

Nell'argomento \param{len} se ne dovrà poi specificare la dimensione
(rispettivamente 4 o 16), infine l'argomento
\param{type} deve indicare il tipo di indirizzo, e dovrà essere o
\const{AF\_INET} o \const{AF\_INET6}.

La funzione restituisce, in caso di successo, un puntatore ad una struttura
\struct{hostent}, solo che in questo caso la ricerca viene eseguita
richiedendo al DNS un record di tipo \texttt{PTR} corrispondente all'indirizzo
specificato. In caso di errore al solito viene usata la variabile
\var{h\_errno} per restituire un opportuno codice. In questo caso l'unico
campo del risultato che interessa è \var{h\_name} che conterrà il nome a
dominio, la funzione comunque inizializza anche il primo campo della lista
\var{h\_addr\_list} col valore dell'indirizzo passato come argomento.

Dato che \func{gethostbyaddr} usa un buffer statico, anche di questa funzione
esiste una versione rientrante \funcd{gethostbyaddr\_r} fornita come
estensione dalla \acr{glibc}, il cui prototipo è:

\begin{funcproto}{
\fhead{netdb.h}
\fhead{sys/socket.h} 
\fdecl{struct hostent *gethostbyaddr\_r(const void *addr, socklen\_t len, int type,\\
\phantom{struct hostent *gethostbyaddr\_r(}struct hostent *ret, char *buf, size\_t buflen,\\
\phantom{struct hostent *gethostbyaddr\_r(}struct hostent **result, int *h\_errnop);}
\fdesc{Richiede la risoluzione inversa di un indirizzo IP.} 
}

{La funzione ritorna $0$ in caso di successo e un valore non nullo per un errore.}
\end{funcproto}

La funzione prende per gli argomenti \param{addr}, \param{len} e \param{type}
gli stessi valori di \func{gethostbyaddr} con lo stesso significato, gli
argomenti successivi vengono utilizzati per restituire i dati, sono identici a
quelli già illustrati in per \func{gethostbyname\_r} e
\func{gethostbyname2\_r} e devono essere usati allo stesso modo.

Infine lo standard POSIX prevede la presenza della funzione
\funcd{gethostent}, il cui prototipo è:

\begin{funcproto}{
\fhead{netdb.h}
\fdecl{struct hostent *gethostent(void)}
\fdesc{Ottiene la voce successiva nel database dei nomi a dominio.} 
}

{La funzione ritorna l'indirizzo ad una struttura \struct{hostent} in caso di
  successo e \val{NULL} per un errore.}
\end{funcproto}

La funzione dovrebbe ritornare (come puntatore alla solita struttura
\struct{hostent} allocata internamente) la voce successiva nel database dei
nomi a dominio, ma questo ha un significato soltato quando è relativo alla
lettura dei dati da un file come \conffile{/etc/hosts} e non per i risultati
del DNS. Nel caso della \acr{glibc} questa viene usata allora solo per la
lettura delle voci presenti in quest'ultimo, come avviene anche in altri
sistemi unix-like, ed inoltre ignora le voci relative ad indirizzi IPv6.

Della stessa funzione la \acr{glibc} fornisce anche una versione rientrante
\funcd{gethostent\_r}, il cui prototipo è:

\begin{funcproto}{
\fhead{netdb.h}
\fdecl{struct hostent *gethostent\_r(struct hostent *ret, char *buf, size\_t buflen,\\
\phantom{struct hostent *gethostent\_r(}struct hostent **result, int *h\_errnop);}
\fdesc{Ottiene la voce successiva nel database dei nomi a dominio.} 
}

{La funzione ritorna $0$ in caso di successo e un valore non nullo per un errore.}
\end{funcproto}

La funzione ha lo stesso effetto di \func{gethostent}; gli argomenti servono a
restituire i risultati in maniera rientrante e vanno usati secondo le modalità
già illustrate per \func{gethostbyname\_r} e \func{gethostbyname2\_r}.

Dati i limiti delle funzioni \func{gethostbyname} e \func{gethostbyaddr} con
l'uso di memoria statica che può essere sovrascritta fra due chiamate
successive, e per avere sempre la possibilità di indicare esplicitamente il
tipo di indirizzi voluto (cosa che non è possibile con \func{gethostbyname}),
è stata successivamente proposta,
nell'\href{http://www.ietf.org/rfc/rfc2553.txt}{RFC~2553} un diversa
interfaccia con l'introduzione due nuove funzioni di
risoluzione,\footnote{dette funzioni sono presenti nella \acr{glibc} 2.1.96,
  ma essendo considerate deprecate (vedi
  sez.~\ref{sec:sock_advanced_name_services}) sono state rimosse nelle
  versioni successive.} \funcd{getipnodebyname} e \funcd{getipnodebyaddr}, i
cui prototipi sono:

\begin{funcproto}{
\fhead{netdb.h}
\fhead{sys/types.h}
\fhead{sys/socket.h}
\fdecl{struct hostent *getipnodebyname(const char *name, int af, int
    flags, int *error\_num)} 
\fdesc{Richiede la risoluzione di un nome a dominio.} 
\fdecl{struct hostent *getipnodebyaddr(const void *addr, size\_t len,
    int af, int *error\_num)}
\fdesc{Richiede la risoluzione inversa di un indirizzo IP.} 
}

{Le funzioni ritornano l'indirizzo ad una struttura \struct{hostent} in caso
  di successo e \val{NULL} per un errore.}
\end{funcproto}

Entrambe le funzioni supportano esplicitamente la scelta di una famiglia di
indirizzi con l'argomento \param{af} (che può assumere i valori
\const{AF\_INET} o \const{AF\_INET6}), e restituiscono un codice di errore
(con valori identici a quelli precedentemente illustrati in
tab.~\ref{tab:h_errno_values}) nella variabile puntata da \param{error\_num}.
La funzione \func{getipnodebyaddr} richiede poi che si specifichi l'indirizzo
come per \func{gethostbyaddr} passando anche la lunghezza dello stesso
nell'argomento \param{len}.

La funzione \func{getipnodebyname} prende come primo argomento il nome da
risolvere, inoltre prevede un apposito argomento \param{flags}, da usare come
maschera binaria, che permette di specificarne il comportamento nella
risoluzione dei diversi tipi di indirizzi (IPv4 e IPv6); ciascun bit
dell'argomento esprime una diversa opzione, e queste possono essere specificate
con un OR aritmetico delle costanti riportate in
tab.~\ref{tab:sock_getipnodebyname_flags}.

\begin{table}[!htb]
  \centering
  \footnotesize
  \begin{tabular}[c]{|l|p{8cm}|}
    \hline
    \textbf{Costante} & \textbf{Significato} \\
    \hline
    \hline
    \constd{AI\_V4MAPPED}  & Usato con \const{AF\_INET6} per richiedere una
                             ricerca su un indirizzo IPv4 invece che IPv6; gli
                             eventuali risultati saranno rimappati su indirizzi 
                             IPv6.\\
    \constd{AI\_ALL}       & Usato con \const{AI\_V4MAPPED}; richiede sia
                             indirizzi IPv4 che IPv6, e gli indirizzi IPv4
                             saranno rimappati in IPv6.\\
    \constd{AI\_ADDRCONFIG}& Richiede che una richiesta IPv4 o IPv6 venga
                             eseguita solo se almeno una interfaccia del
                             sistema è associata ad un indirizzo di tale tipo.\\
    \constd{AI\_DEFAULT}   & Il valore di default, è equivalente alla
                             combinazione di \const{AI\_ADDRCONFIG} e di
                             \const{AI\_V4MAPPED}.\\  
    \hline
  \end{tabular}
  \caption{Valori possibili per i bit dell'argomento \param{flags} della
    funzione \func{getipnodebyname}.}
  \label{tab:sock_getipnodebyname_flags}
\end{table}

Entrambe le funzioni restituiscono un puntatore ad una struttura \var{hostent}
che contiene i risultati della ricerca, che viene allocata dinamicamente
insieme a tutto lo spazio necessario a contenere i dati in essa referenziati;
per questo motivo queste funzioni non soffrono dei problemi dovuti all'uso di
una sezione statica di memoria presenti con le precedenti \func{gethostbyname}
e \func{gethostbyaddr}.  L'uso di una allocazione dinamica però comporta anche
la necessità di disallocare esplicitamente la memoria occupata dai risultati
una volta che questi non siano più necessari; a tale scopo viene fornita la
funzione \funcd{freehostent}, il cui prototipo è:

\begin{funcproto}{
\fhead{netdb.h}
\fhead{sys/types.h}
\fhead{sys/socket.h}
\fdecl{void freehostent(struct hostent *ip)}
\fdesc{Disalloca una struttura \var{hostent}.} 
}

{La funzione non ritorna nulla, e non da errori.}
\end{funcproto}

La funzione permette di disallocare una struttura \var{hostent}
precedentemente allocata in una chiamata di \func{getipnodebyname} o
\func{getipnodebyaddr}, e prende come argomento l'indirizzo restituito da una
di queste funzioni.

Infine per concludere la nostra panoramica sulle funzioni di risoluzione dei
nomi dobbiamo citare le funzioni che permettono di interrogare gli altri
servizi di risoluzione dei nomi illustrati in sez.~\ref{sec:sock_resolver}; in
generale infatti ci sono una serie di funzioni nella forma
\texttt{get\textsl{XXX}byname} e \texttt{get\textsl{XXX}byaddr} (dove
\texttt{\textsl{XXX}} indica il servizio) per ciascuna delle informazioni di
rete mantenute dal \textit{Name Service Switch} che permettono rispettivamente
di trovare una corrispondenza cercando per nome o per numero.

L'elenco di queste funzioni è riportato nelle colonne finali di
tab.~\ref{tab:name_resolution_functions}, dove le si sono suddivise rispetto
al tipo di informazione che forniscono (riportato in prima colonna). Nella
tabella si è anche riportato il file su cui vengono ordinariamente mantenute
queste informazioni, che però può essere sostituito da un qualunque supporto
interno al \textit{Name Service Switch} (anche se usualmente questo avviene
solo per la risoluzione degli indirizzi).  Ciascuna funzione fa riferimento ad
una sua apposita struttura che contiene i relativi dati, riportata in terza
colonna.

\begin{table}[!htb]
  \centering
  \footnotesize
  \begin{tabular}[c]{|l|l|l|l|l|}
    \hline
    \textbf{Informazione}&\textbf{File}&\textbf{Struttura}&
    \multicolumn{2}{|c|}{\textbf{Funzioni}}\\
    \hline
    \hline
    indirizzo &\conffile{/etc/hosts}&\struct{hostent}&\func{gethostbyname}&
               \func{gethostbyaddr}\\ 
    servizio  &\conffile{/etc/services}&\struct{servent}&\func{getservbyname}&
               \func{getservbyport}\\ 
    rete      &\conffile{/etc/networks}&\struct{netent}&\funcm{getnetbyname}&
               \funcm{getnetbyaddr}\\ 
    protocollo&\conffile{/etc/protocols}&\struct{protoent}&
               \funcm{getprotobyname}&\funcm{getprotobyaddr}\\ 
    \hline
  \end{tabular}
  \caption{Funzioni di risoluzione dei nomi per i vari servizi del
    \textit{Name Service Switch} riguardanti la rete.}
  \label{tab:name_resolution_functions}
\end{table}

Delle funzioni di tab.~\ref{tab:name_resolution_functions} abbiamo trattato
finora soltanto quelle relative alla risoluzione dei nomi, dato che sono le
più usate, e prevedono praticamente da sempre la necessità di rivolgersi ad
una entità esterna; per le altre invece, estensioni fornite dal \textit{Name
  Service Switch} a parte, si fa sempre riferimento ai dati mantenuti nei
rispettivi file.

Dopo la risoluzione dei nomi a dominio una delle ricerche più comuni è quella
sui nomi dei servizi di rete più comuni (cioè \texttt{http}, \texttt{smtp},
ecc.) da associare alle rispettive porte. Le due funzioni da utilizzare per
questo sono \funcd{getservbyname} e \funcd{getservbyport}, che permettono
rispettivamente di ottenere il numero di porta associato ad un servizio dato
il nome e viceversa; i loro prototipi sono:



\begin{funcproto}{
\fhead{netdb.h} 
\fdecl{struct servent *getservbyname(const char *name, const char *proto)}
\fdecl{struct servent *getservbyport(int port, const char *proto)} 
\fdesc{Risolvono il nome di un servizio nel rispettivo numero di porta e viceversa.} 
}

{Le funzioni ritornano il puntatore ad una struttura \struct{servent} con i
    risultati in caso di successo e \val{NULL} per un errore.}
\end{funcproto}

Entrambe le funzioni prendono come ultimo argomento una stringa \param{proto}
che indica il protocollo per il quale si intende effettuare la ricerca (le
informazioni mantenute in \conffile{/etc/services} infatti sono relative sia
alle porte usate su UDP che su TCP, occorre quindi specificare a quale dei due
protocolli si fa riferimento) che nel caso di IP può avere come valori
possibili solo \texttt{udp} o \texttt{tcp};\footnote{in teoria si potrebbe
  avere un qualunque protocollo fra quelli citati in
  \conffile{/etc/protocols}, posto che lo stesso supporti il concetto di
  \textsl{porta}, in pratica questi due sono gli unici presenti.}  se si
specifica un puntatore nullo la ricerca sarà eseguita su un protocollo
qualsiasi.

Il primo argomento è il nome del servizio per \func{getservbyname},
specificato tramite la stringa \param{name}, mentre \func{getservbyport}
richiede il numero di porta in \param{port}. Entrambe le funzioni eseguono una
ricerca sul file \conffile{/etc/services}\footnote{il \textit{Name Service
    Switch} astrae il concetto a qualunque supporto su cui si possano
  mantenere i suddetti dati.}  ed estraggono i dati dalla prima riga che
corrisponde agli argomenti specificati; se la risoluzione ha successo viene
restituito un puntatore ad una apposita struttura \struct{servent} contenente
tutti i risultati, altrimenti viene restituito un puntatore nullo.  Si tenga
presente che anche in questo caso i dati vengono mantenuti in una area di
memoria statica e che quindi la funzione non è rientrante.

\begin{figure}[!htb]
  \footnotesize \centering
  \begin{minipage}[c]{0.80\textwidth}
    \includestruct{listati/servent.h}
  \end{minipage}
  \caption{La struttura \structd{servent} per la risoluzione dei nomi dei
    servizi e dei numeri di porta.}
  \label{fig:sock_servent_struct}
\end{figure}

La definizione della struttura \struct{servent} è riportata in
fig.~\ref{fig:sock_servent_struct}, il primo campo, \var{s\_name} contiene
sempre il nome canonico del servizio, mentre \var{s\_aliases} è un puntatore
ad un vettore di stringhe contenenti gli eventuali nomi alternativi
utilizzabili per identificare lo stesso servizio. Infine \var{s\_port}
contiene il numero di porta e \var{s\_proto} il nome del protocollo.

Come riportato in tab.~\ref{tab:name_resolution_functions} ci sono analoghe
funzioni per la risoluzione del nome dei protocolli e delle reti; non staremo
a descriverle nei dettagli, in quanto il loro uso è molto limitato, esse
comunque utilizzano una loro struttura dedicata del tutto analoga alle
precedenti: tutti i dettagli relativi al loro funzionamento possono essere
trovati nelle rispettive pagine di manuale.

Oltre alle funzioni di ricerca esistono delle ulteriori funzioni che prevedono
una lettura sequenziale delle informazioni mantenute nel \textit{Name Service
  Switch} (in sostanza permettono di leggere i file contenenti le informazioni
riga per riga), che sono analoghe a quelle elencate in
tab.~\ref{tab:sys_passwd_func} per le informazioni relative ai dati degli
utenti e dei gruppi. Nel caso specifico dei servizi avremo allora le tre
funzioni \funcd{setservent}, \funcd{getservent} e \funcd{endservent} i cui
prototipi sono:

\begin{funcproto}{
\fhead{netdb.h} 
\fdecl{struct servent *getservent(void)}
\fdesc{Legge la voce successiva nel file \conffile{/etc/services}.} 
\fdecl{void setservent(int stayopen)} 
\fdesc{Apre il file \conffile{/etc/services} e si posiziona al suo inizio.} 
\fdecl{void endservent(void)}
\fdesc{Chiude il file \conffile{/etc/services}.} 
}

{Le due funzioni \func{setservent} e \func{endservent} non ritornano nulla,
  \func{getservent} restituisce il puntatore ad una struttura \struct{servent}
  in caso di successo e \val{NULL} per un errore o fine del file.}
\end{funcproto}

La prima funzione, \func{getservent}, legge una singola voce a partire dalla
posizione corrente in \conffile{/etc/services}, pertanto si può eseguire una
lettura sequenziale dello stesso invocandola più volte. Se il file non è
aperto provvede automaticamente ad aprirlo, nel qual caso leggerà la prima
voce. La seconda funzione, \func{setservent}, permette di aprire il file
\conffile{/etc/services} per una successiva lettura, ma se il file è già stato
aperto riporta la posizione di lettura alla prima voce del file, in questo
modo si può far ricominciare da capo una lettura sequenziale. 

\begin{table}[!htb]
  \centering
  \footnotesize
  \begin{tabular}[c]{|l|l|l|l|}
    \hline
    \textbf{Informazione}&\multicolumn{3}{|c|}{\textbf{Funzioni}}\\
    \hline
    \hline
    indirizzo &\func{sethostent} &\func{gethostent} &\func{endhostent} \\
    servizio  &\func{setservent} &\func{getservent} &\func{endservent}\\ 
    rete      &\funcm{setnetent}  &\funcm{getnetent}  &\funcm{endnetent}\\ 
    protocollo&\funcm{setprotoent}&\funcm{getprotoent}&\funcm{endprotoent}\\ 
    \hline
  \end{tabular}
  \caption{Funzioni lettura sequenziale dei dati del
    \textit{Name Service Switch}.} 
  \label{tab:name_sequential_read}
\end{table}

L'argomento \param{stayopen} di \func{setservent}, se diverso da zero, fa sì
che il file resti aperto anche fra diverse chiamate a \func{getservbyname} e
\func{getservbyport}; di default dopo una chiamata a queste funzioni il file
viene chiuso, cosicché una successiva chiamata a \func{getservent} riparte
dall'inizio.  La terza funzione, \func{endservent}, provvede semplicemente a
chiudere il file.

Queste tre funzioni per la lettura sequenziale di nuovo sono presenti per
ciascuno dei vari tipi di informazione relative alle reti di
tab.~\ref{tab:name_resolution_functions}; questo significa che esistono
altrettante funzioni nella forma \texttt{setXXXent}, \texttt{getXXXent} e
\texttt{endXXXent}, analoghe alle precedenti per la risoluzione dei servizi,
che abbiamo riportato in tab.~\ref{tab:name_sequential_read}.  Essendo, a
parte il tipo di informazione che viene trattato, sostanzialmente identiche
nel funzionamento e di scarso utilizzo, non staremo a trattarle una per una,
rimandando alle rispettive pagine di manuale.


\subsection{Le funzioni avanzate per la risoluzione dei nomi}
\label{sec:sock_advanced_name_services}

Quelle illustrate nella sezione precedente sono le funzioni classiche per la
risoluzione di nomi ed indirizzi IP, ma abbiamo già visto come esse soffrano
di vari inconvenienti come il fatto che usano informazioni statiche, e non
prevedono la possibilità di avere diverse classi di indirizzi. Anche se sono
state create delle estensioni o metodi diversi che permettono di risolvere
alcuni di questi inconvenienti, comunque esse non forniscono una interfaccia
sufficientemente generica.\footnote{rimane ad esempio il problema generico che
  si deve sapere in anticipo quale tipo di indirizzi IP (IPv4 o IPv6)
  corrispondono ad un certo nome a dominio.}

Inoltre in genere quando si ha a che fare con i socket non esiste soltanto il
problema della risoluzione del nome che identifica la macchina, ma anche
quello del servizio a cui ci si vuole rivolgere.  Per questo motivo con lo
standard POSIX 1003.1-2001 sono state indicate come deprecate le varie
funzioni \func{gethostbyaddr}, \func{gethostbyname}, \var{getipnodebyname} e
\var{getipnodebyaddr} ed è stata introdotta una interfaccia completamente
nuova.

La prima funzione di questa interfaccia è \funcd{getaddrinfo}, che combina le
funzionalità delle precedenti \func{getipnodebyname}, \func{getipnodebyaddr},
\func{getservbyname} e \func{getservbyport}, consentendo di ottenere
contemporaneamente sia la risoluzione di un indirizzo simbolico che del nome
di un servizio; la funzione è stata introdotta, insieme a \func{getnameinfo}
che vedremo più avanti,
nell'\href{http://www.ietf.org/rfc/rfc2553.txt}{RFC~2553} ed il suo prototipo è:


\begin{funcproto}{
\fhead{netdb.h}
\fhead{sys/socket.h} 
\fdecl{int getaddrinfo(const char *node, const char *service, const
    struct addrinfo *hints, \\
\phantom{int getaddrinfo(}struct addrinfo **res)}
\fdesc{Esegue una risoluzione di un nome a dominio e di un nome di servizio.} 
}

{La funzione ritorna $0$ in caso di successo e un codice di errore diverso da
  zero per un errore.}
\end{funcproto}

La funzione prende come primo argomento il nome della macchina che si vuole
risolvere, specificato tramite la stringa \param{node}. Questo argomento,
oltre ad un comune nome a dominio, può indicare anche un indirizzo numerico in
forma \textit{dotted-decimal} per IPv4 o in formato esadecimale per IPv6.  Si
può anche specificare il nome di una rete invece che di una singola macchina.
Il secondo argomento, \param{service}, specifica invece il nome del servizio
che si intende risolvere. Per uno dei due argomenti si può anche usare il
valore \val{NULL}, nel qual caso la risoluzione verrà effettuata soltanto
sulla base del valore dell'altro.

Il terzo argomento, \param{hints}, deve essere invece un puntatore ad una
struttura \struct{addrinfo} usata per dare dei \textsl{suggerimenti} al
procedimento di risoluzione riguardo al protocollo o del tipo di socket che si
intenderà utilizzare; \func{getaddrinfo} infatti permette di effettuare
ricerche generiche sugli indirizzi, usando sia IPv4 che IPv6, e richiedere
risoluzioni sui nomi dei servizi indipendentemente dal protocollo (ad esempio
TCP o UDP) che questi possono utilizzare.

Come ultimo argomento in \param{res} deve essere passato un puntatore ad una
variabile (di tipo puntatore ad una struttura \struct{addrinfo}) che verrà
utilizzata dalla funzione per riportare (come \textit{value result argument})
i propri risultati. La funzione infatti è rientrante, ed alloca autonomamente
tutta la memoria necessaria in cui verranno riportati i risultati della
risoluzione.  La funzione scriverà all'indirizzo puntato da \param{res} il
puntatore iniziale ad una \textit{linked list} di strutture di tipo
\struct{addrinfo} contenenti tutte le informazioni ottenute.

\begin{figure}[!htb]
  \footnotesize \centering
  \begin{minipage}[c]{0.90\textwidth}
    \includestruct{listati/addrinfo.h}
  \end{minipage}
  \caption{La struttura \structd{addrinfo} usata nella nuova interfaccia POSIX
    per la risoluzione di nomi a dominio e servizi.}
  \label{fig:sock_addrinfo_struct}
\end{figure}

Come illustrato la struttura \struct{addrinfo}, la cui definizione è riportata
in fig.~\ref{fig:sock_addrinfo_struct}, viene usata sia in ingresso, per
passare dei valori di controllo alla funzione, che in uscita, per ricevere i
risultati. La definizione è ripresa direttamente dal file \headfiled{netdb.h}
in cui questa struttura viene dichiarata, la pagina di manuale riporta
\type{size\_t} come tipo di dato per il campo \var{ai\_addrlen}, qui viene
usata quanto previsto dallo standard POSIX, in cui viene utilizzato
\type{socklen\_t}; i due tipi di dati sono comunque equivalenti.

Il primo campo, \var{ai\_flags}, è una maschera binaria di bit che
permettono di controllare le varie modalità di risoluzione degli indirizzi,
che viene usato soltanto in ingresso. I tre campi successivi \var{ai\_family},
\var{ai\_socktype}, e \var{ai\_protocol} contengono rispettivamente la
famiglia di indirizzi, il tipo di socket e il protocollo, in ingresso vengono
usati per impostare una selezione (impostandone il valore nella struttura
puntata da \param{hints}), mentre in uscita indicano il tipo di risultato
contenuto nella struttura.

Tutti i campi seguenti vengono usati soltanto in uscita e devono essere nulli
o \val{NULL} in ingresso; il campo \var{ai\_addrlen} indica la dimensione
della struttura degli indirizzi ottenuta come risultato, il cui contenuto sarà
memorizzato nella struttura \struct{sockaddr} posta all'indirizzo puntato dal
campo \var{ai\_addr}. Il campo \var{ai\_canonname} è un puntatore alla stringa
contenente il nome canonico della macchina, ed infine, quando la funzione
restituisce più di un risultato, \var{ai\_next} è un puntatore alla successiva
struttura \struct{addrinfo} della lista.

Ovviamente non è necessario dare dei suggerimenti in ingresso, ed usando
\val{NULL} come valore per l'argomento \param{hints} si possono compiere
ricerche generiche.  Se però si specifica un valore non nullo questo deve
puntare ad una struttura \struct{addrinfo} precedentemente allocata nella
quale siano stati opportunamente impostati i valori dei campi
\var{ai\_family}, \var{ai\_socktype}, \var{ai\_protocol} ed \var{ai\_flags}.

I due campi \var{ai\_family} e \var{ai\_socktype} prendono gli stessi valori
degli analoghi argomenti della funzione \func{socket}; in particolare per
\var{ai\_family} si possono usare i valori di tab.~\ref{tab:net_pf_names} ma
sono presi in considerazione solo \const{AF\_INET} e \const{AF\_INET6}, mentre
se non si vuole specificare nessuna famiglia di indirizzi si può usare il
valore \const{AF\_UNSPEC}.  Allo stesso modo per \var{ai\_socktype} si possono
usare i valori illustrati in sez.~\ref{sec:sock_type} per indicare per quale
tipo di socket si vuole risolvere il servizio indicato, anche se i soli
significativi sono \const{SOCK\_STREAM} e \const{SOCK\_DGRAM}; in questo caso,
se non si vuole effettuare nessuna risoluzione specifica, si potrà usare un
valore nullo.

Il campo \var{ai\_protocol} permette invece di effettuare la selezione dei
risultati per il nome del servizio usando il numero identificativo del
rispettivo protocollo di trasporto (i cui valori possibili sono riportati in
\conffile{/etc/protocols}); di nuovo i due soli valori utilizzabili sono quelli
relativi a UDP e TCP, o il valore nullo che indica di ignorare questo campo
nella selezione.

\begin{table}[!htb]
  \centering
  \footnotesize
  \begin{tabular}[c]{|l|p{8cm}|}
    \hline
    \textbf{Costante} & \textbf{Significato} \\
    \hline
    \hline
    \const{AI\_ADDRCONFIG} & Stesso significato dell'analoga di
                             tab.~\ref{tab:sock_getipnodebyname_flags}.\\ 
    \const{AI\_ALL}        & Stesso significato dell'analoga di
                             tab.~\ref{tab:sock_getipnodebyname_flags}.\\ 
    \constd{AI\_CANONNAME}  & Richiede la restituzione del nome canonico della
                              macchina, che verrà salvato in una stringa il cui
                              indirizzo sarà restituito nel campo
                              \var{ai\_canonname} della prima struttura
                              \struct{addrinfo} dei risultati. Se il nome
                              canonico non è disponibile al suo posto
                              viene restituita una copia di \param{node}. \\ 
    \constd{AI\_NUMERICHOST}& Se impostato il nome della macchina specificato
                              con \param{node} deve essere espresso in forma
                              numerica, altrimenti sarà restituito un errore
                              \const{EAI\_NONAME} (vedi
                              tab.~\ref{tab:addrinfo_error_code}), in questo
                              modo si evita ogni chiamata alle funzioni di
                              risoluzione.\\ 
    \constd{AI\_NUMERICSERVICE}& Analogo di \const{AI\_NUMERICHOST} per la
                                 risoluzione di un servizio, con 
                                 \param{service} che deve essere espresso in forma
                                 numerica.\\ 
    \constd{AI\_PASSIVE}    & Viene utilizzato per ottenere un indirizzo in
                              formato adatto per una successiva chiamata a
                              \func{bind}. Se specificato quando si è usato 
                              \val{NULL} come valore per \param{node} gli
                              indirizzi restituiti saranno inizializzati al
                              valore generico (\const{INADDR\_ANY} per IPv4 e
                              \const{IN6ADDR\_ANY\_INIT} per IPv6), altrimenti
                              verrà usato l'indirizzo dell'interfaccia di
                              \textit{loopback}. Se invece non è impostato gli
                              indirizzi verranno restituiti in formato adatto ad
                              una chiamata a \func{connect} o \func{sendto}.\\
    \const{AI\_V4MAPPED}   & Stesso significato dell'analoga di
                             tab.~\ref{tab:sock_getipnodebyname_flags}.\\  
    \hline
    \const{AI\_CANONIDN}   & Se il nome canonico richiesto con
                             \const{AI\_CANONNAME} è codificato con questo
                             flag la codifica viene convertita in forma
                             leggibile nella localizzazione corrente.\\ 
    \const{AI\_IDN}        & Se specificato il nome viene convertito, se
                             necessario, nella codifica IDN, usando la
                             localizzazione corrente.\\ 
    \const{AI\_IDN\_ALLOW\_UNASSIGNED} & attiva il controllo 
                                         \texttt{IDNA\_ALLOW\_UNASSIGNED}.\\ 
    \const{AI\_AI\_IDN\_USE\_STD3\_ASCII\_RULES} & attiva il controllo
                                            \texttt{IDNA\_USE\_STD3\_ASCII\_RULES}\\ 
    \hline
  \end{tabular}
  \caption{Costanti associate ai bit del campo \var{ai\_flags} della struttura 
    \struct{addrinfo}.} 
  \label{tab:ai_flags_values}
\end{table}

% TODO mettere riferimento a IDNA_ALLOW_UNASSIGNED e IDNA_USE_STD3_ASCII_RULES

Infine gli ultimi dettagli si controllano con il campo \var{ai\_flags}; che
deve essere impostato come una maschera binaria; i bit di questa variabile
infatti vengono usati per dare delle indicazioni sul tipo di risoluzione
voluta, ed hanno valori analoghi a quelli visti in
sez.~\ref{sec:sock_name_services} per \func{getipnodebyname}; il valore di
\var{ai\_flags} può essere impostata con un OR aritmetico delle costanti di
tab.~\ref{tab:ai_flags_values}, ciascuna delle quali identifica un bit della
maschera. 

Nella seconda parte della tabella si sono riportati i valori delle costanti
aggiunte a partire dalla \acr{glibc} 2.3.4 per gestire la
internazionalizazione dei nomi a dominio (IDN o \textit{Internationalized
  Domain Names}) secondo quanto specificato
nell'\href{http://www.ietf.org/rfc/rfc3490.txt}{RFC~3490} (potendo cioè usare
codifiche di caratteri che consentono l'espressione di nomi a dominio in
qualunque lingua).

Come accennato passando un valore \val{NULL} per l'argomento \param{hints} si
effettua una risuluzione generica, equivalente ad aver impostato un valore
nullo per \var{ai\_family} e \var{ai\_socktype}, un valore \const{AF\_UNSPEC}
per \var{ai\_family} e il valore \code{(AI\_V4MAPPED|AI\_ADDRCONFIG)} per
\var{ai\_flags}. 

La funzione restituisce un valore nullo in caso di successo, o un codice in
caso di errore. I valori usati come codice di errore sono riportati in
tab.~\ref{tab:addrinfo_error_code}; dato che la funzione utilizza altre
funzioni e chiamate al sistema per ottenere il suo risultato in generale il
valore di \var{errno} non è significativo, eccetto il caso in cui si sia
ricevuto un errore di \const{EAI\_SYSTEM}, nel qual caso l'errore
corrispondente è riportato tramite \var{errno}.

\begin{table}[!htb]
  \centering
  \footnotesize
  \begin{tabular}[c]{|l|p{10cm}|}
    \hline
    \textbf{Costante} & \textbf{Significato} \\
    \hline
    \hline
    \constd{EAI\_ADDRFAMILY}& La richiesta non ha nessun indirizzo di rete
                              per la famiglia di indirizzi specificata. \\
    \constd{EAI\_AGAIN}   & Il DNS ha restituito un errore di risoluzione  
                            temporaneo, si può ritentare in seguito. \\
    \constd{EAI\_BADFLAGS}& Il campo \var{ai\_flags} contiene dei valori non
                            validi per i flag o si è richiesto
                            \const{AI\_CANONNAME} con \param{name} nullo. \\
    \constd{EAI\_FAIL}    & Il DNS ha restituito un errore di risoluzione  
                            permanente. \\
    \constd{EAI\_FAMILY}  & La famiglia di indirizzi richiesta non è
                            supportata. \\ 
    \constd{EAI\_MEMORY}  & È stato impossibile allocare la memoria necessaria
                            alle operazioni. \\
    \constd{EAI\_NODATA}  & La macchina specificata esiste, ma non ha nessun
                            indirizzo di rete definito. \\
    \constd{EAI\_NONAME}  & Il nome a dominio o il servizio non sono noti,
                            viene usato questo errore anche quando si specifica
                            il valore \val{NULL} per entrambi gli argomenti
                            \param{node} e \param{service}. \\
    \constd{EAI\_SERVICE} & Il servizio richiesto non è disponibile per il tipo
                            di socket richiesto, anche se può esistere per
                            altri tipi di socket. \\
    \constd{EAI\_SOCKTYPE}& Il tipo di socket richiesto non è supportato. \\
    \constd{EAI\_SYSTEM}  & C'è stato un errore di sistema, si può controllare
                            \var{errno} per i dettagli. \\
%    \hline
% TODO estensioni GNU, trovarne la documentazione
%    \constd{EAI\_INPROGRESS}& Richiesta in corso. \\
%    \constd{EAI\_CANCELED}& La richiesta è stata cancellata.\\
%    \constd{EAI\_NOTCANCELED}& La richiesta non è stata cancellata. \\
%    \constd{EAI\_ALLDONE} & Tutte le richieste sono complete. \\
%    \constd{EAI\_INTR}    & Richiesta interrotta. \\
    \hline
  \end{tabular}
  \caption{Costanti associate ai valori dei codici di errore della funzione
    \func{getaddrinfo}.} 
  \label{tab:addrinfo_error_code}
\end{table}

Come per i codici di errore di \func{gethostbyname} anche in questo caso è
fornita una apposita funzione, simile a \func{strerror}, che consente di
utilizzare direttamente il codice restituito dalla funzione per stampare a
video un messaggio esplicativo; la funzione è \funcd{gai\_strerror} ed il suo
prototipo è:

\begin{funcproto}{
\fhead{netdb.h} 
\fdecl{const char *gai\_strerror(int errcode)}
\fdesc{Fornisce il messaggio corrispondente ad un errore di \func{getaddrinfo}.} 
}

{La funzione ritorna il puntatore alla stringa contenente il messaggio di
  errore.}
\end{funcproto}

La funzione restituisce un puntatore alla stringa contenente il messaggio
corrispondente dal codice di errore \param{errcode} ottenuto come valore di
ritorno di \func{getaddrinfo}.  La stringa è allocata staticamente, ma essendo
costante ed accessibile in sola lettura, la funzione è rientrante.

Dato che ad un certo nome a dominio possono corrispondere più indirizzi IP
(sia IPv4 che IPv6), e che un certo servizio può essere fornito su protocolli
e tipi di socket diversi, in generale, a meno di non aver eseguito una
selezione specifica attraverso l'uso di \param{hints}, si otterrà una diversa
struttura \struct{addrinfo} per ciascuna possibilità.  

Ad esempio se si richiede la risoluzione del servizio \textit{echo} per
l'indirizzo \texttt{www.truelite.it}, e si imposta \const{AI\_CANONNAME} per
avere anche la risoluzione del nome canonico, si avrà come risposta della
funzione la lista illustrata in fig.~\ref{fig:sock_addrinfo_list}.

\begin{figure}[!htb]
  \centering
  \includegraphics[width=10cm]{img/addrinfo_list}
  \caption{La \textit{linked list} delle strutture \struct{addrinfo}
    restituite da \func{getaddrinfo}.}
  \label{fig:sock_addrinfo_list}
\end{figure}

Come primo esempio di uso di \func{getaddrinfo} vediamo un programma
elementare di interrogazione del \textit{resolver} basato questa funzione, il
cui corpo principale è riportato in fig.~\ref{fig:mygetaddr_example}. Il
codice completo del programma, compresa la gestione delle opzioni in cui è
gestita l'eventuale inizializzazione dell'argomento \var{hints} per
restringere le ricerche su protocolli, tipi di socket o famiglie di indirizzi,
è disponibile nel file \texttt{mygetaddr.c} dei sorgenti allegati alla guida.

\begin{figure}[!htb]
  \footnotesize \centering
  \begin{minipage}[c]{\codesamplewidth}
    \includecodesample{listati/mygetaddr.c}
  \end{minipage}
  \normalsize
  \caption{Esempio di codice per la risoluzione di un indirizzo.}
  \label{fig:mygetaddr_example}
\end{figure}

Il corpo principale inizia controllando (\texttt{\small 1--5}) il numero di
argomenti passati, che devono essere sempre due, e corrispondere
rispettivamente all'indirizzo ed al nome del servizio da risolvere. A questo
segue la chiamata (\texttt{\small 7}) alla funzione \func{getaddrinfo}, ed il
successivo controllo (\texttt{\small 8--11}) del suo corretto funzionamento,
senza il quale si esce immediatamente stampando il relativo codice di errore.

Se la funzione ha restituito un valore nullo il programma prosegue
inizializzando (\texttt{\small 12}) il puntatore \var{ptr} che sarà usato nel
successivo ciclo (\texttt{\small 14--35}) di scansione della lista delle
strutture \struct{addrinfo} restituite dalla funzione. Prima di eseguire
questa scansione (\texttt{\small 12}) viene stampato il valore del nome
canonico che è presente solo nella prima struttura.

La scansione viene ripetuta (\texttt{\small 14}) fintanto che si ha un
puntatore valido. La selezione principale è fatta sul campo \var{ai\_family},
che stabilisce a quale famiglia di indirizzi fa riferimento la struttura in
esame. Le possibilità sono due, un indirizzo IPv4 o IPv6, se nessuna delle due
si verifica si provvede (\texttt{\small 27--30}) a stampare un messaggio di
errore ed uscire (questa eventualità non dovrebbe comunque mai verificarsi,
almeno fintanto che la funzione \func{getaddrinfo} lavora correttamente).

Per ciascuno delle due possibili famiglie di indirizzi si estraggono le
informazioni che poi verranno stampate alla fine del ciclo (\texttt{\small
  31--34}). Il primo caso esaminato (\texttt{\small 15--21}) è quello degli
indirizzi IPv4, nel qual caso prima se ne stampa l'identificazione
(\texttt{\small 16}) poi si provvede a ricavare la struttura degli indirizzi
(\texttt{\small 17}) indirizzata dal campo \var{ai\_addr}, eseguendo un
opportuno casting del puntatore per poter estrarre da questa la porta
(\texttt{\small 18}) e poi l'indirizzo (\texttt{\small 19}) che verrà
convertito con una chiamata ad \func{inet\_ntop}.

La stessa operazione (\texttt{\small 21--27}) viene ripetuta per gli indirizzi
IPv6, usando la rispettiva struttura degli indirizzi. Si noti anche come in
entrambi i casi per la chiamata a \func{inet\_ntop} si sia dovuto passare il
puntatore al campo contenente l'indirizzo IP nella struttura puntata dal campo
\var{ai\_addr}.\footnote{il meccanismo è complesso a causa del fatto che al
  contrario di IPv4, in cui l'indirizzo IP può essere espresso con un semplice
  numero intero, in IPv6 questo deve essere necessariamente fornito come
  struttura, e pertanto anche se nella struttura puntata da \var{ai\_addr}
  sono presenti direttamente i valori finali, per l'uso con \func{inet\_ntop}
  occorre comunque passare un puntatore agli stessi (ed il costrutto
  \code{\&addr6->sin6\_addr} è corretto in quanto l'operatore \texttt{->} ha
  in questo caso precedenza su \texttt{\&}).}

Una volta estratte dalla struttura \struct{addrinfo} tutte le informazioni
relative alla risoluzione richiesta e stampati i relativi valori, l'ultimo
passo (\texttt{\small 34}) è di estrarre da \var{ai\_next} l'indirizzo della
eventuale successiva struttura presente nella lista e ripetere il ciclo, fin
tanto che, completata la scansione, questo avrà un valore nullo e si potrà
terminare (\texttt{\small 36}) il programma.

Si tenga presente che \func{getaddrinfo} non garantisce nessun particolare
ordinamento della lista delle strutture \struct{addrinfo} restituite, anche se
usualmente i vari indirizzi IP (se ne è presente più di uno) sono forniti
nello stesso ordine in cui vengono inviati dal server DNS. In particolare
nulla garantisce che vengano forniti prima i dati relativi ai servizi di un
determinato protocollo o tipo di socket, se ne sono presenti di diversi.  Se
allora utilizziamo il nostro programma potremo verificare il risultato:
\begin{Console}
[piccardi@gont sources]$ \textbf{./mygetaddr -c  gapil.truelite.it echo}
Canonical name sources2.truelite.it
IPv4 address:
        Indirizzo 62.48.34.25
        Protocollo 6
        Porta 7
IPv4 address:
        Indirizzo 62.48.34.25
        Protocollo 17
        Porta 7
\end{Console}
%$

Una volta estratti i risultati dalla \textit{linked list} puntata
da \param{res} se questa non viene più utilizzata si dovrà avere cura di
disallocare opportunamente tutta la memoria, per questo viene fornita
l'apposita funzione \funcd{freeaddrinfo}, il cui prototipo è:


\begin{funcproto}{
\fhead{netdb.h}
\fdecl{void freeaddrinfo(struct addrinfo *res)}
\fdesc{Libera la memoria allocata da una precedente chiamata a \func{getaddrinfo}.} 
}

{La funzione non restituisce nessun codice di errore.}
\end{funcproto}

La funzione prende come unico argomento il puntatore \param{res}, ottenuto da
una precedente chiamata a \func{getaddrinfo}, e scandisce la lista delle
strutture per liberare tutta la memoria allocata. Dato che la funzione non ha
valori di ritorno deve essere posta molta cura nel passare un valore valido
per \param{res} ed usare un indirizzo non valido o già liberato può avere
conseguenze non prevedibili.

Si tenga presente infine che se si copiano i risultati da una delle strutture
\struct{addrinfo} restituite nella lista indicizzata da \param{res}, occorre
avere cura di eseguire una \textit{deep copy} in cui si copiano anche tutti i
dati presenti agli indirizzi contenuti nella struttura \struct{addrinfo},
perché una volta disallocati i dati con \func{freeaddrinfo} questi non
sarebbero più disponibili.

Anche la nuova interfaccia definita da POSIX prevede una nuova funzione per
eseguire la risoluzione inversa e determinare nomi di servizi e di dominio
dati i rispettivi valori numerici. La funzione che sostituisce le varie
\func{gethostbyname}, \func{getipnodebyname} e \func{getservbyname} è
\funcd{getnameinfo}, ed il suo prototipo è:

\begin{funcproto}{
\fhead{netdb.h}
\fhead{sys/socket.h}
\fdecl{int getnameinfo(const struct sockaddr *sa, socklen\_t salen, \\
\phantom{int getnameinfo(}char *host, size\_t hostlen, char *serv, size\_t
servlen, int flags)} 
\fdesc{Effettua una risoluzione di un indirizzo di rete in maniera
  indipendente dal protocollo.} 
}

{La funzione ritorna $0$ in caso di successo e un codice di errore diverso da
  zero per un errore.}
\end{funcproto}

La principale caratteristica di \func{getnameinfo} è che la funzione è in
grado di eseguire una risoluzione inversa in maniera indipendente dal
protocollo; il suo primo argomento \param{sa} infatti è il puntatore ad una
struttura degli indirizzi generica, che può contenere sia indirizzi IPv4 che
IPv6, la cui dimensione deve comunque essere specificata con l'argomento
\param{salen}. 

I risultati della funzione saranno restituiti nelle due stringhe puntate da
\param{host} e \param{serv}, che dovranno essere state precedentemente
allocate per una lunghezza massima che deve essere specificata con gli altri
due argomenti \param{hostlen} e \param{servlen}. Quando non si è
interessati ad uno dei due, si può passare il valore \val{NULL} come argomento,
così che la corrispondente informazione non venga richiesta. Infine l'ultimo
argomento \param{flags} è una maschera binaria i cui bit consentono di
impostare le modalità con cui viene eseguita la ricerca, e deve essere
specificato attraverso l'OR aritmetico dei valori illustrati in
tab.~\ref{tab:getnameinfo_flags}, nella seconda parte della tabella si sono
aggiunti i valori introdotto con la \acr{glibc} 2.3.4 per gestire la
internazionalizzione dei nomi a dominio.

\begin{table}[!htb]
  \centering
  \footnotesize
  \begin{tabular}[c]{|l|p{8cm}|}
    \hline
    \textbf{Costante} & \textbf{Significato} \\
    \hline
    \hline
    \constd{NI\_DGRAM}      & Richiede che venga restituito il nome del
                              servizio su UDP invece che quello su TCP per quei
                              pichi servizi (porte 512-214) che soni diversi
                              nei due protocolli.\\
    \constd{NI\_NOFQDN}     & Richiede che venga restituita solo il nome della
                              macchina all'interno del dominio al posto del
                              nome completo (FQDN).\\
    \constd{NI\_NAMEREQD}   & Richiede la restituzione di un errore se il nome
                              non può essere risolto.\\
    \constd{NI\_NUMERICHOST}& Richiede che venga restituita la forma numerica
                              dell'indirizzo (questo succede sempre se il nome
                              non può essere ottenuto).\\ 
    \constd{NI\_NUMERICSERV}& Richiede che il servizio venga restituito in
                              forma numerica (attraverso il numero di
                              porta).\\
    \hline
    \const{NI\_IDN}        & Se specificato il nome restituito viene convertito usando la
                             localizzazione corrente, se necessario, nella
                             codifica IDN.\\ 
    \const{NI\_IDN\_ALLOW\_UNASSIGNED} & attiva il controllo 
                                       \texttt{IDNA\_ALLOW\_UNASSIGNED}.\\ 
    \const{NI\_AI\_IDN\_USE\_STD3\_ASCII\_RULES} & attiva il controllo
                                       \texttt{IDNA\_USE\_STD3\_ASCII\_RULES}\\ 
    \hline
  \end{tabular}
  \caption{Costanti associate ai bit dell'argomento \param{flags} della  
    funzione \func{getnameinfo}.} 
  \label{tab:getnameinfo_flags}
\end{table}

La funzione ritorna zero in caso di successo, e scrive i propri risultati agli
indirizzi indicati dagli argomenti \param{host} e \param{serv} come stringhe
terminate dal carattere NUL, a meno che queste non debbano essere troncate
qualora la loro dimensione ecceda quelle specificate dagli argomenti
\param{hostlen} e \param{servlen}. Sono comunque definite le due costanti
\constd{NI\_MAXHOST} e \constd{NI\_MAXSERV}\footnote{in Linux le due costanti
  sono definite in \headfile{netdb.h} ed hanno rispettivamente il valore 1024
  e 12.}  che possono essere utilizzate come limiti massimi.  In caso di
errore viene restituito invece un codice che assume gli stessi valori
illustrati in tab.~\ref{tab:addrinfo_error_code}.

A questo punto possiamo fornire degli esempi di utilizzo della nuova
interfaccia, adottandola per le precedenti implementazioni del client e del
server per il servizio \textit{echo}; dato che l'uso delle funzioni appena
illustrate (in particolare di \func{getaddrinfo}) è piuttosto complesso,
essendo necessaria anche una impostazione diretta dei campi dell'argomento
\param{hints}, provvederemo una interfaccia semplificata per i due casi visti
finora, quello in cui si specifica nel client un indirizzo remoto per la
connessione al server, e quello in cui si specifica nel server un indirizzo
locale su cui porsi in ascolto.

La prima funzione della nostra interfaccia semplificata è \texttt{sockconn}
che permette di ottenere un socket, connesso all'indirizzo ed al servizio
specificati. Il corpo della funzione è riportato in
fig.~\ref{fig:sockconn_code}, il codice completo è nel file \file{SockUtil.c}
dei sorgenti allegati alla guida, che contiene varie funzioni di utilità per
l'uso dei socket.

\begin{figure}[!htb]
  \footnotesize \centering
  \begin{minipage}[c]{\codesamplewidth}
    \includecodesample{listati/sockconn.c}
  \end{minipage}
  \normalsize
  \caption{Il codice della funzione \texttt{sockconn}.}
  \label{fig:sockconn_code}
\end{figure}

La funzione prende quattro argomenti, i primi due sono le stringhe che
indicano il nome della macchina a cui collegarsi ed il relativo servizio su
cui sarà effettuata la risoluzione; seguono il protocollo da usare (da
specificare con il valore numerico di \conffile{/etc/protocols}) ed il tipo di
socket (al solito specificato con i valori illustrati in
sez.~\ref{sec:sock_type}).  La funzione ritorna il valore del file descriptor
associato al socket (un numero positivo) in caso di successo, o $-1$ in caso
di errore.

Per risolvere il problema di non poter passare indietro i valori di ritorno di
\func{getaddrinfo} contenenti i relativi codici di errore si sono stampati i
messaggi d'errore direttamente nella funzione; infatti non si può avere
nessuna certezza che detti valori siano negativi e per cui stampare subito
l'errore diventa necessario per evitare ogni possibile ambiguità nei confronti
del valore di ritorno in caso di successo.

Una volta definite le variabili occorrenti (\texttt{\small 3--5}) la funzione
prima (\texttt{\small 6}) azzera il contenuto della struttura \var{hint} e poi
provvede (\texttt{\small 7--9}) ad inizializzarne i valori necessari per la
chiamata (\texttt{\small 10}) a \func{getaddrinfo}. Di quest'ultima si
controlla (\texttt{\small 12--16}) il codice di ritorno, in modo da stampare
un avviso di errore, azzerare \var{errno} ed uscire in caso di errore.  

Dato che ad una macchina possono corrispondere più indirizzi IP, e di tipo
diverso (sia IPv4 che IPv6), mentre il servizio può essere in ascolto soltanto
su uno solo di questi, si provvede a tentare la connessione per ciascun
indirizzo restituito all'interno di un ciclo (\texttt{\small 18--40}) di
scansione della lista restituita da \func{getaddrinfo}, ma prima
(\texttt{\small 17}) si salva il valore del puntatore per poterlo riutilizzare
alla fine per disallocare la lista.

Il ciclo viene ripetuto (\texttt{\small 18}) fintanto che si hanno indirizzi
validi, ed inizia (\texttt{\small 19}) con l'apertura del socket; se questa
fallisce si controlla (\texttt{\small 20}) se sono disponibili altri
indirizzi, nel qual caso si passa al successivo (\texttt{\small 21}) e si
riprende (\texttt{\small 22}) il ciclo da capo; se non ve ne sono si stampa
l'errore ritornando immediatamente (\texttt{\small 24--27}). 

Quando la creazione del socket ha avuto successo si procede (\texttt{\small
  29}) direttamente con la connessione, di nuovo in caso di fallimento viene
ripetuto (\texttt{\small 30--38}) il controllo se vi sono o no altri indirizzi
da provare nella stessa modalità fatta in precedenza, aggiungendovi però in
entrambi i casi (\texttt{\small 32} e (\texttt{\small 36}) la chiusura del
socket precedentemente aperto, che non è più utilizzabile.

Se la connessione ha avuto successo invece si termina (\texttt{\small 39})
direttamente il ciclo, e prima di ritornare (\texttt{\small 31}) il valore del
file descriptor del socket si provvede (\texttt{\small 30}) a liberare le
strutture \struct{addrinfo} allocate da \func{getaddrinfo} utilizzando il
valore del relativo puntatore precedentemente (\texttt{\small 17}) salvato.
Si noti come per la funzione sia del tutto irrilevante se la struttura
ritornata contiene indirizzi IPv6 o IPv4, in quanto si fa uso direttamente dei
dati relativi alle strutture degli indirizzi di \struct{addrinfo} che sono
opachi rispetto all'uso della funzione \func{connect}.

Per usare questa funzione possiamo allora modificare ulteriormente il nostro
programma client per il servizio \textit{echo}; in questo caso rispetto al
codice usato finora per collegarsi (vedi fig.~\ref{fig:TCP_echo_client_1})
avremo una semplificazione per cui il corpo principale del nostro client
diventerà quello illustrato in fig.~\ref{fig:TCP_echo_fifth}, in cui le
chiamate a \func{socket}, \func{inet\_pton} e \func{connect} sono sostituite
da una singola chiamata a \texttt{sockconn}. Inoltre il nuovo client (il cui
codice completo è nel file \file{TCP\_echo\_fifth.c} dei sorgenti allegati)
consente di utilizzare come argomento del programma un nome a dominio al posto
dell'indirizzo numerico, e può utilizzare sia indirizzi IPv4 che IPv6.

\begin{figure}[!htb]
  \footnotesize \centering
  \begin{minipage}[c]{\codesamplewidth}
    \includecodesample{listati/TCP_echo_fifth.c}
  \end{minipage}
  \normalsize
  \caption{Il nuovo codice per la connessione del client \textit{echo}.}
  \label{fig:TCP_echo_fifth}
\end{figure}

La seconda funzione di ausilio che abbiamo creato è \texttt{sockbind}, il cui
corpo principale è riportato in fig.~\ref{fig:sockbind_code} (al solito il
sorgente completo è nel file \file{sockbind.c} dei sorgenti allegati alla
guida). Come si può notare la funzione è del tutto analoga alla precedente
\texttt{sockconn}, e prende gli stessi argomenti, però invece di eseguire una
connessione con \func{connect} si limita a chiamare \func{bind} per collegare
il socket ad una porta.

\begin{figure}[!htb]
  \footnotesize \centering
  \begin{minipage}[c]{\codesamplewidth}
    \includecodesample{listati/sockbind.c}
  \end{minipage}
  \normalsize
  \caption{Il codice della funzione \texttt{sockbind}.}
  \label{fig:sockbind_code}
\end{figure}

Dato che la funzione è pensata per essere utilizzata da un server ci si può
chiedere a quale scopo mantenere l'argomento \param{host} quando l'indirizzo
di questo è usualmente noto. Si ricordi però quanto detto in
sez.~\ref{sec:TCP_func_bind}, relativamente al significato della scelta di un
indirizzo specifico come argomento di \func{bind}, che consente di porre il
server in ascolto su uno solo dei possibili diversi indirizzi presenti su di
una macchina.  Se non si vuole che la funzione esegua \func{bind} su un
indirizzo specifico, ma utilizzi l'indirizzo generico, occorrerà avere cura di
passare un valore \val{NULL} come valore per l'argomento \var{host}; l'uso
del valore \const{AI\_PASSIVE} serve ad ottenere il valore generico nella
rispettiva struttura degli indirizzi.

Come già detto la funzione è analoga a \texttt{sockconn} ed inizia azzerando
ed inizializzando (\texttt{\small 6--11}) opportunamente la struttura
\var{hint} con i valori ricevuti come argomenti, soltanto che in questo caso
si è usata (\texttt{\small 8}) una impostazione specifica dei flag di
\var{hint} usando \const{AI\_PASSIVE} per indicare che il socket sarà usato
per una apertura passiva. Per il resto la chiamata (\texttt{\small 12--18}) a
\func{getaddrinfo} e ed il ciclo principale (\texttt{\small 20--42}) sono
identici, solo che si è sostituita (\texttt{\small 31}) la chiamata a
\func{connect} con una chiamata a \func{bind}. Anche la conclusione
(\texttt{\small 43--44}) della funzione è identica.

Si noti come anche in questo caso si siano inserite le stampe degli errori
sullo \textit{standard error}, nonostante la funzione possa essere invocata da
un demone. Nel nostro caso questo non è un problema in quanto se la funzione
non ha successo il programma deve uscire immediatamente prima di essere posto
in background, e può quindi scrivere gli errori direttamente sullo
\textit{standard error}.

\begin{figure}[!htbp]
  \footnotesize \centering
  \begin{minipage}[c]{\codesamplewidth}
    \includecodesample{listati/TCP_echod_third.c}
  \end{minipage}
  \normalsize
  \caption{Nuovo codice per l'apertura passiva del server \textit{echo}.}
  \label{fig:TCP_echod_third}
\end{figure}

Con l'uso di questa funzione si può modificare anche il codice del nostro
server \textit{echo}, che rispetto a quanto illustrato nella versione iniziale
di fig.~\ref{fig:TCP_echo_server_first_code} viene modificato nella forma
riportata in fig.~\ref{fig:TCP_echod_third}. In questo caso il socket su cui
porsi in ascolto viene ottenuto (\texttt{\small 15--18}) da \texttt{sockbind}
che si cura anche della eventuale risoluzione di un indirizzo specifico sul
quale si voglia far ascoltare il server.


\section{Le opzioni dei socket}
\label{sec:sock_options}

Benché dal punto di vista del loro uso come canali di trasmissione di dati i
socket vengano trattati allo stesso modo dei file, acceduti tramite i file
descriptor, e gestiti con le ordinarie funzioni di lettura e scrittura dei
file, l'interfaccia standard usata per la gestione dei file generici non è
comunque sufficiente a controllare la moltitudine di caratteristiche
specifiche che li contraddistinguono, considerato tra l'altro che queste
possono essere completamente diverse fra loro a seconda del tipo di socket e
della relativa forma di comunicazione sottostante.

In questa sezione vedremo allora quali sono le funzioni dedicate alla gestione
delle caratteristiche specifiche dei vari tipi di socket, che vengono
raggruppate sotto il nome generico di ``\textit{socket options}'', ma
soprattutto analizzaremo quali sono queste opzioni e quali caretteristiche e
comportamenti dei socket permettono di controllare.


\subsection{Le funzioni di gestione delle opzioni dei socket}
\label{sec:sock_setsockopt}

La modalità principale con cui si possono gestire le caratteristiche dei
socket (ne vedremo delle ulteriori nelle prossime sezioni) è quella che passa
attraverso l'uso di due funzioni di sistema generiche che permettono
rispettivamente di impostarle e di recuperarne il valore corrente. La prima di
queste due funzioni, quella usata per impostare le \textit{socket options}, è
\funcd{setsockopt}, ed il suo prototipo è:

\begin{funcproto}{
\fhead{sys/socket.h}
\fhead{sys/types.h}
\fdecl{int setsockopt(int sock, int level, int optname, const void
    *optval, socklen\_t optlen)}
\fdesc{Imposta le opzioni di un socket.} 
}

{La funzione ritorna $0$ in caso di successo e $-1$ per un errore, nel qual
  caso \var{errno} assumerà uno dei valori: 
  \begin{errlist}
  \item[\errcode{EBADF}]  il file descriptor \param{sock} non è valido.
  \item[\errcode{EFAULT}] l'indirizzo \param{optval} non è valido.
  \item[\errcode{EINVAL}] il valore di \param{optlen} non è valido.
  \item[\errcode{ENOPROTOOPT}] l'opzione scelta non esiste per il livello
    indicato. 
  \item[\errcode{ENOTSOCK}] il file descriptor \param{sock} non corrisponde ad
    un socket.
  \end{errlist}
}
\end{funcproto}

Il primo argomento della funzione, \param{sock}, indica il socket su cui si
intende operare; per indicare l'opzione da impostare si devono usare i due
argomenti successivi, \param{level} e \param{optname}.  Come abbiamo visto in
sez.~\ref{sec:net_protocols} i protocolli di rete sono strutturati su vari
livelli, ed l'interfaccia dei socket può usarne più di uno. Si avranno allora
funzionalità e caratteristiche diverse per ciascun protocollo usato da un
socket, e quindi saranno anche diverse le opzioni che si potranno impostare
per ciascun socket, a seconda del \textsl{livello} (trasporto, rete, ecc.) su
cui si vuole andare ad operare.

Il valore di \param{level} seleziona allora il protocollo su cui vuole
intervenire, mentre \param{optname} permette di scegliere su quale delle
opzioni che sono definite per quel protocollo si vuole operare. In sostanza la
selezione di una specifica opzione viene fatta attraverso una coppia di valori
\param{level} e \param{optname} e chiaramente la funzione avrà successo
soltanto se il protocollo in questione prevede quella opzione ed è utilizzato
dal socket.  Infine \param{level} prevede anche il valore speciale
\const{SOL\_SOCKET} usato per le opzioni generiche che sono disponibili per
qualunque tipo di socket.

\begin{table}[!htb]
  \centering
  \footnotesize
  \begin{tabular}[c]{|l|l|}
    \hline
    \textbf{Livello} & \textbf{Significato} \\
    \hline
    \hline
    \constd{SOL\_SOCKET}& Opzioni generiche dei socket.\\
    \constd{SOL\_IP}    & Opzioni specifiche per i socket che usano IPv4.\\
    \constd{SOL\_TCP}   & Opzioni per i socket che usano TCP.\\
    \constd{SOL\_IPV6}  & Opzioni specifiche per i socket che usano IPv6.\\
    \constd{SOL\_ICMPV6}& Opzioni specifiche per i socket che usano ICMPv6.\\
    \hline
  \end{tabular}
  \caption{Possibili valori dell'argomento \param{level} delle 
    funzioni \func{setsockopt} e \func{getsockopt}.} 
  \label{tab:sock_option_levels}
\end{table}

I valori usati per \param{level}, corrispondenti ad un dato protocollo usato
da un socket, sono quelli corrispondenti al valore numerico che identifica il
suddetto protocollo in \conffile{/etc/protocols}; dato che la leggibilità di un
programma non trarrebbe certo beneficio dall'uso diretto dei valori numerici,
più comunemente si indica il protocollo tramite le apposite costanti
\texttt{SOL\_*} riportate in tab.~\ref{tab:sock_option_levels}, dove si sono
riassunti i valori che possono essere usati per l'argomento
\param{level}. 

La notazione in questo caso è, purtroppo, abbastanza confusa: infatti in Linux
il valore si può impostare sia usando le costanti \texttt{SOL\_*}, che le
analoghe \texttt{IPPROTO\_*} (citate anche da Stevens in \cite{UNP1});
entrambe hanno gli stessi valori che sono equivalenti ai numeri di protocollo
di \conffile{/etc/protocols}, con una eccezione specifica, che è quella del
protocollo ICMP, per la quale non esiste una costante, il che è comprensibile
dato che il suo valore, 1, è quello che viene assegnato a \const{SOL\_SOCKET}.

Il quarto argomento, \param{optval} è un puntatore ad una zona di memoria che
contiene i dati che specificano il valore dell'opzione che si vuole passare al
socket, mentre l'ultimo argomento \param{optlen},\footnote{questo argomento è
  in realtà sempre di tipo \ctyp{int}, come era nelle \acr{libc4} e
  \acr{libc5}; l'uso di \type{socklen\_t} è stato introdotto da POSIX (valgono
  le stesse considerazioni per l'uso di questo tipo di dato fatte in
  sez.~\ref{sec:TCP_func_accept}) ed adottato dalla \acr{glibc}.} è la
dimensione in byte dei dati presenti all'indirizzo indicato da \param{optval}.
Dato che il tipo di dati varia a seconda dell'opzione scelta, occorrerà
individuare qual è quello che deve essere usato, ed utilizzare le opportune
variabili.

La gran parte delle opzioni utilizzano per \param{optval} un valore intero, se
poi l'opzione esprime una condizione logica, il valore è sempre un intero, ma
si dovrà usare un valore non nullo per abilitarla ed un valore nullo per
disabilitarla.  Se invece l'opzione non prevede di dover ricevere nessun tipo
di valore si deve impostare \param{optval} a \val{NULL}. Un piccolo numero
di opzioni però usano dei tipi di dati peculiari, è questo il motivo per cui
\param{optval} è stato definito come puntatore generico.

La seconda funzione usata per controllare le proprietà dei socket è
\funcd{getsockopt}, che serve a leggere i valori delle opzioni dei socket ed a
farsi restituire i dati relativi al loro funzionamento; il suo prototipo è:


\begin{funcproto}{
\fhead{sys/socket.h}
\fhead{sys/types.h}
\fdecl{int getsockopt(int s, int level, int optname, void *optval,
    socklen\_t *optlen)}
\fdesc{Legge le opzioni di un socket.} 
}

{La funzione ritorna $0$ in caso di successo e $-1$ per un errore, nel qual
  caso \var{errno} assumerà uno dei valori: 
  \begin{errlist}
  \item[\errcode{EBADF}] il file descriptor \param{sock} non è valido.
  \item[\errcode{EFAULT}] l'indirizzo \param{optval} o quello di
    \param{optlen} non è valido.
  \item[\errcode{ENOPROTOOPT}] l'opzione scelta non esiste per il livello
    indicato. 
  \item[\errcode{ENOTSOCK}] il file descriptor \param{sock} non corrisponde ad
    un socket.
  \end{errlist}
}
\end{funcproto}

I primi tre argomenti sono identici ed hanno lo stesso significato di quelli
di \func{setsockopt}, anche se non è detto che tutte le opzioni siano definite
per entrambe le funzioni. In questo caso \param{optval} viene usato per
ricevere le informazioni ed indica l'indirizzo a cui andranno scritti i dati
letti dal socket, infine \param{optlen} diventa un puntatore ad una variabile
che viene usata come \textit{value result argument} per indicare, prima della
chiamata della funzione, la lunghezza del buffer allocato per \param{optval} e
per ricevere indietro, dopo la chiamata della funzione, la dimensione
effettiva dei dati scritti su di esso.  Se la dimensione del buffer allocato
per \param{optval} non è sufficiente si avrà un errore.



\subsection{Le opzioni generiche}
\label{sec:sock_generic_options}

Come accennato esiste un insieme generico di opzioni dei socket che possono
applicarsi a qualunque tipo di socket, indipendentemente da quale protocollo
venga poi utilizzato. Una descrizione di queste opzioni è generalmente
disponibile nella settima sezione delle pagine di manuale; nel caso specifico
la si può consultare con \texttt{man 7 socket}. Se si vuole operare su queste
opzioni generiche il livello da utilizzare è \const{SOL\_SOCKET}; si è
riportato un elenco di queste opzioni in tab.~\ref{tab:sock_opt_socklevel}.


\begin{table}[!htb]
  \centering
  \footnotesize
  \begin{tabular}[c]{|l|c|c|c|l|l|}
    \hline
    \textbf{Opzione}&\texttt{get}&\texttt{set}&\textbf{flag}&\textbf{Tipo}&
                    \textbf{Descrizione}\\
    \hline
    \hline
    \const{SO\_ACCEPTCONN}&$\bullet$&        &         &\texttt{int}& 
                          Indica se il socket è in ascolto.\\
    \const{SO\_ATTACH\_FILTER}&     &$\bullet$&         &\texttt{sock\_fprog}& 
                          Aggancia un filtro BPF al socket.\\
    \const{SO\_BINDTODEVICE}&$\bullet$&$\bullet$&         &\texttt{char *}& 
                          Lega il socket ad un dispositivo.\\
    \const{SO\_BROADCAST}&$\bullet$&$\bullet$&$\bullet$&\texttt{int}& 
                          Attiva o disattiva il \textit{broadcast}.\\ 
    \const{SO\_BSDCOMPAT}&$\bullet$&$\bullet$&$\bullet$&\texttt{int}& 
                          Abilita la compatibilità con BSD.\\
    \const{SO\_BUSY\_POLL}&$\bullet$&$\bullet$&$\bullet$&\texttt{int}& 
                          Attiva il ``\textit{busy poll}'' sul socket.\\
    \const{SO\_DEBUG}    &$\bullet$&$\bullet$&$\bullet$&\texttt{int}& 
                          Abilita il debugging sul socket.\\
    \const{SO\_DETACH\_FILTER}&    &$\bullet$&$\bullet$&\texttt{int}& 
                          Rimuove il filtro BPF agganciato al socket.\\
    \const{SO\_DOMAIN}   &$\bullet$&         &         &\texttt{int}& 
                          Legge il tipo di socket.\\
    \const{SO\_DONTROUTE}&$\bullet$&$\bullet$&$\bullet$&\texttt{int}& 
                          Non invia attraverso un gateway.\\
    \const{SO\_ERROR}    &$\bullet$&         &         &\texttt{int}& 
                          Riceve e cancella gli errori pendenti.\\
    \const{SO\_KEEPALIVE}&$\bullet$&$\bullet$&$\bullet$&\texttt{int}& 
                          Controlla l'attività della connessione.\\
    \const{SO\_LINGER}   &$\bullet$&$\bullet$&         &\texttt{linger}&
                          Indugia nella chiusura con dati da spedire.\\
    \const{SO\_LOCK\_FILTER}&    &$\bullet$&$\bullet$&\texttt{int}& 
                          Blocca il filtro BPF agganciato al socket.\\
    \const{SO\_MARK}     &$\bullet$&$\bullet$&         &\texttt{int}&
                          Imposta un ``\textit{firewall mark}'' sul socket.\\
    \const{SO\_OOBINLINE}&$\bullet$&$\bullet$&$\bullet$&\texttt{int}& 
                          Lascia in linea i dati \textit{out-of-band}.\\
    \const{SO\_PASSCRED} &$\bullet$&$\bullet$&$\bullet$&\texttt{int}& 
                          Abilita la ricezione di credenziali.\\
    \const{SO\_PEEK\_OFF} &$\bullet$&$\bullet$&$\bullet$&\texttt{int}& 
                          Imposta il valore del ``\textit{peek offset}''.\\
    \const{SO\_PEERCRED} &$\bullet$&         &         &\texttt{ucred}& 
                          Restituisce le credenziali del processo remoto.\\
    \const{SO\_PRIORITY} &$\bullet$&$\bullet$&         &\texttt{int}& 
                          Imposta la priorità del socket.\\
    \const{SO\_PROTOCOL} &$\bullet$&         &         &\texttt{int}& 
                          Ottiene il protocollo usato dal socket.\\
    \const{SO\_RCVBUF}   &$\bullet$&$\bullet$&         &\texttt{int}& 
                          Imposta dimensione del buffer di ricezione.\\
    \const{SO\_RCVBUFFORCE}&$\bullet$&$\bullet$&         &\texttt{int}& 
                          Forza dimensione del buffer di ricezione.\\
    \const{SO\_RCVLOWAT} &$\bullet$&$\bullet$&$\bullet$&\texttt{int}& 
                          Basso livello sul buffer di ricezione.\\
    \const{SO\_RCVTIMEO} &$\bullet$&$\bullet$&         &\texttt{timeval}& 
                          Timeout in ricezione.\\
    \const{SO\_REUSEADDR}&$\bullet$&$\bullet$&$\bullet$&\texttt{int}& 
                          Consente il riutilizzo di un indirizzo locale.\\
    \const{SO\_REUSEPORT}&$\bullet$&$\bullet$&$\bullet$&\texttt{int}& 
                          Consente il riutilizzo di una porta.\\
    \const{SO\_RXQ\_OVFL} &$\bullet$&$\bullet$&$\bullet$&\texttt{int}& 
                          Richiede messaggio ancillare con pacchetti persi.\\
    \const{SO\_SNDBUF}   &$\bullet$&$\bullet$&         &\texttt{int}& 
                          Imposta dimensione del buffer di trasmissione.\\
    \const{SO\_SNDBUFFORCE}&$\bullet$&$\bullet$&         &\texttt{int}& 
                          Forza dimensione del buffer di trasmissione.\\
    \const{SO\_SNDLOWAT} &$\bullet$&$\bullet$&         &\texttt{int}&
                          Basso livello sul buffer di trasmissione.\\
    \const{SO\_SNDTIMEO} &$\bullet$&$\bullet$&         &\texttt{timeval}& 
                          Timeout in trasmissione.\\
    \const{SO\_TIMESTAMP}&$\bullet$&$\bullet$&$\bullet$&\texttt{int}&
                          Abilita/disabilita la ricezione dei \textit{timestamp}.\\
    \const{SO\_TYPE}     &$\bullet$&         &         &\texttt{int}& 
                          Restituisce il tipo di socket.\\
   \hline
  \end{tabular}
  \caption{Le opzioni disponibili al livello \const{SOL\_SOCKET}.} 
  \label{tab:sock_opt_socklevel}
\end{table}

% TODO aggiungere e documentare SO_ATTACH_BPF, introdotta con il kernel 3.19,
% vedi http://lwn.net/Articles/625224/
% TODO aggiungere e documentare SO_INCOMING_CPU, introdotta con il kernel 3.19,
% TODO documentare SO_PEERGROUPS introdotta con il kernel 4.13, citata
% in https://lwn.net/Articles/727385/
% TODO documentare SO_ZEROCOPY, vedi https://lwn.net/Articles/726917/ (e il
% resto pure) introdotta con il kernel 4.14

La tabella elenca le costanti che identificano le singole opzioni da usare
come valore per \param{optname}; le due colonne seguenti indicano per quali
delle due funzioni (\func{getsockopt} o \func{setsockopt}) l'opzione è
disponibile, mentre la colonna successiva indica, quando si ha a che fare con
un valore di \param{optval} intero, se l'opzione è da considerare un numero o
un valore logico. Si è inoltre riportato sulla quinta colonna il tipo di dato
usato per \param{optval} ed una breve descrizione del significato delle
singole opzioni sulla sesta.

Le descrizioni delle opzioni presenti in tab.~\ref{tab:sock_opt_socklevel}
sono estremamente sommarie, è perciò necessario fornire un po' più di
informazioni. Alcune opzioni inoltre hanno una notevole rilevanza nella
gestione dei socket, e pertanto il loro utilizzo sarà approfondito
separatamente in sez.~\ref{sec:sock_options_main}. Quello che segue è quindi
soltanto un elenco più dettagliato della breve descrizione di
tab.~\ref{tab:sock_opt_socklevel} sul significato delle varie opzioni:

\begin{basedescript}{\desclabelwidth{2.2cm}\desclabelstyle{\nextlinelabel}}
\item[\constd{SO\_ACCEPTCONN}] questa opzione permette di rilevare se il socket
  su cui opera è stato posto in modalità di ricezione di eventuali connessioni
  con una chiamata a \func{listen}. L'opzione può essere usata soltanto con
  \func{getsockopt} ed utilizza per \param{optval} un intero in cui viene
  restituito 1 se il socket è in ascolto e 0 altrimenti.

\item[\constd{SO\_ATTACH\_FILTER}] questa opzione permette di agganciare ad un
  socket un filtro di selezione dei pacchetti con la stessa sintassi del BPF
  (\textit{Berkley Packet Filter}) di BSD, che consente di selezionare, fra
  tutti quelli ricevuti, verranno letti. Può essere usata solo con
  \func{setsockopt} ed utilizza per \param{optval} un puntatore ad una
  struttura \struct{sock\_fprog} (definita in
  \headfile{linux/filter.h}). Questa opzione viene usata principalmente con i
  socket di tipo \const{AF\_PACKET} (torneremo su questo in
  sez.~\ref{sec:packet_socket}) dalla libreria \texttt{libpcap} per
  implementare programmi di cattura dei pacchetti, e per questo tipo di
  applicazione è opportuno usare sempre quest'ultima.\footnote{la trattazione
    del BPF va al di là dell'argomento di questa sezione per la documentazione
    si consulti il file \texttt{networking/filter.txt} nella documentazione
    del kernel.}

\item[\constd{SO\_BINDTODEVICE}] questa opzione permette di \textsl{legare} il
  socket ad una particolare interfaccia, in modo che esso possa ricevere ed
  inviare pacchetti solo su quella. L'opzione richiede per \param{optval} il
  puntatore ad una stringa contenente il nome dell'interfaccia (ad esempio
  \texttt{eth0}); utilizzando una stringa nulla o un valore nullo per
  \param{optlen} si può rimuovere un precedente collegamento.

  Il nome della interfaccia deve essere specificato con una stringa terminata
  da uno zero e di lunghezza massima pari a \constd{IFNAMSIZ}; l'opzione è
  effettiva solo per alcuni tipi di socket, ed in particolare per quelli della
  famiglia \const{AF\_INET}; non è invece supportata per i \textit{packet
    socket} (vedi sez.~\ref{sec:socket_raw}). 

\item[\constd{SO\_BROADCAST}] questa opzione abilita il \textit{broadcast}:
  quando abilitata i socket di tipo \const{SOCK\_DGRAM} riceveranno i
  pacchetti inviati all'indirizzo di \textit{broadcast} e potranno scrivere
  pacchetti su tale indirizzo.  Prende per \param{optval} un intero usato come
  valore logico. L'opzione non ha effetti su un socket di tipo
  \const{SOCK\_STREAM}.

\item[\constd{SO\_BSDCOMPAT}] questa opzione abilita la compatibilità con il
  comportamento di BSD (in particolare ne riproduce alcuni bug).  Attualmente
  è una opzione usata solo per il protocollo UDP e ne è prevista la rimozione.
  L'opzione utilizza per \param{optval} un intero usato come valore logico.

  Quando viene abilitata gli errori riportati da messaggi ICMP per un socket
  UDP non vengono passati al programma in \textit{user space}. Con le versioni
  2.0.x del kernel erano anche abilitate altre opzioni di compatibilità per i
  socket raw (modifiche casuali agli header, perdita del flag di
  \textit{broadcast}) che sono state rimosse con il passaggio al 2.2; è
  consigliato correggere i programmi piuttosto che usare questa funzione. Dal
  kernel 2.4 viene ignorata, e dal 2.6 genera un messaggio di log del kernel.

\item[\constd{SO\_BUSY\_POLL}] questa opzione, presente dal kernel 3.11,
  imposta un tempo approssimato in microsecondi, per cui in caso di ricezione
  bloccante verrà eseguito un \itindex{busy~poll} ``\textit{busy
    poll}'',\footnote{con \textit{busy poll} si fa riferimento al
    \textit{polling} su una risorsa occupata; si continuerà cioè a tentare di
    leggere anche quando non ci sono dati senza portare il processo stato di
    \textit{sleep}, in alcuni casi, quando ci si aspetta che i dati arrivino a
    breve, questa tecnica può dare un miglioramento delle prestazioni.}  da
  indicare in \param{optval} con un valore intero. Si tratta di una opzione
  utilizzabile solo con socket che ricevono dati da un dispositivo di rete che
  la supporti, e che consente di ridurre la latenza per alcune applicazioni,
  ma che comporta un maggiore utilizzo della CPU (e quindi di energia); per
  questo il valore può essere aumentato solo da processi con i privilegi di
  amministratore (in particolare con la \textit{capability}
  \const{CAP\_NET\_ADMIN}).  Il valore di default viene controllato dal file
  \sysctlrelfile{net/core}{busy\_read} per le funzioni di lettura mentre il
  file \sysctlrelfile{net/core}{busy\_poll} controlla il ``\textit{busy
    poll}'' per \func{select} e \func{poll}.

\item[\constd{SO\_DEBUG}] questa opzione abilita il debugging delle operazioni
  dei socket; l'opzione utilizza per \param{optval} un intero usato come
  valore logico, e può essere utilizzata solo da un processo con i privilegi
  di amministratore (in particolare con la \textit{capability}
  \const{CAP\_NET\_ADMIN}).  L'opzione necessita inoltre dell'opportuno
  supporto nel kernel;\footnote{deve cioè essere definita la macro di
    preprocessore \macrod{SOCK\_DEBUGGING} nel file \file{include/net/sock.h}
    dei sorgenti del kernel, questo è sempre vero nei kernel delle serie
    superiori alla 2.3, per i kernel delle serie precedenti invece è
    necessario aggiungere a mano detta definizione; è inoltre possibile
    abilitare anche il tracciamento degli stati del TCP definendo la macro
    \macrod{STATE\_TRACE} in \file{include/net/tcp.h}.}  quando viene
  abilitata una serie di messaggi con le informazioni di debug vengono inviati
  direttamente al sistema del kernel log.\footnote{si tenga presente che il
    comportamento è diverso da quanto avviene con BSD, dove l'opzione opera
    solo sui socket TCP, causando la scrittura di tutti i pacchetti inviati
    sulla rete su un buffer circolare che viene letto da un apposito
    programma, \cmd{trpt}.}

\item[\constd{SO\_DETACH\_FILTER}] consente di distaccare un filtro
  precedentemente aggiunto ad un socket con l'opzione
  \const{SO\_ATTACH\_FILTER}, in genere non viene usata direttamente in quanto
  i filtri BPF vengono automaticamente rimossi alla chiusura del socket, il
  suo utilizzo è pertanto limitato ai rari casi in cui si vuole rimuovere un
  precedente filtro per inserirne uno diverso. Come \const{SO\_ATTACH\_FILTER}
  può essere usato solo \func{setsockopt} e prende per \param{optval} un
  intero usato come valore logico.

\item[\constd{SO\_DOMAIN}] questa opzione, presente dal kernel 2.6.32, legge
  il ``\textsl{dominio}'' (la famiglia di indirizzi) del socket.
. Funziona solo con \func{getsockopt}, ed
  utilizza per \param{optval} un intero in cui verrà restituito il valore
  numerico che lo identifica (ad esempio \texttt{AF\_INET}).

\item[\constd{SO\_DONTROUTE}] questa opzione richiede che l'invio dei
  pacchetti del socket sia eseguito soltanto verso una destinazione
  direttamente connessa, impedendo l'uso di un \textit{gateway} e saltando
  ogni processo relativo all'uso della tabella di routing del kernel. Prende
  per \param{optval} un intero usato come valore logico. È equivalente all'uso
  del flag \const{MSG\_DONTROUTE} su una \func{send} (vedi
  sez.~\ref{sec:net_send_recv}).

\item[\constd{SO\_ERROR}] questa opzione riceve un errore presente sul socket;
  può essere utilizzata soltanto con \func{getsockopt} e prende per
  \param{optval} un valore intero, nel quale viene restituito il codice di
  errore, e la condizione di errore sul socket viene cancellata. Viene
  usualmente utilizzata per ricevere il codice di errore, come accennato in
  sez.~\ref{sec:TCP_sock_select}, quando si sta osservando il socket con una
  \func{select} che ritorna a causa dello stesso.

\item[\const{SO\_KEEPALIVE}] questa opzione abilita un meccanismo di verifica
  della persistenza di una connessione associata al socket; è pertanto
  effettiva solo sui socket che supportano le connessioni, ed è usata
  principalmente con il TCP. L'opzione utilizza per \param{optval} un intero
  usato come valore logico. Maggiori dettagli sul suo funzionamento sono
  forniti in sez.~\ref{sec:sock_options_main}.

\item[\const{SO\_LINGER}] questa opzione controlla le modalità con cui viene
  chiuso un socket quando si utilizza un protocollo che supporta le
  connessioni (è pertanto usata con i socket TCP ed ignorata per UDP) e
  modifica il comportamento delle funzioni \func{close} e \func{shutdown}.
  L'opzione richiede che l'argomento \param{optval} sia una struttura di tipo
  \struct{linger}, definita in \headfile{sys/socket.h} ed illustrata in
  fig.~\ref{fig:sock_linger_struct}.  Maggiori dettagli sul suo funzionamento
  sono forniti in sez.~\ref{sec:sock_options_main}.

\item[\constd{SO\_LOCK\_FILTER}] consente di bloccare un filtro
  precedentemente aggiunto ad un socket con l'opzione
  \const{SO\_ATTACH\_FILTER}, in modo che non possa essere né rimosso né
  modificato, questo consente di impostare un filtro su un socket, bloccarlo e
  poi cedere i privilegi con la sicurezza che il filtro permarrà fino alla
  chiusura del socket. Come \const{SO\_ATTACH\_FILTER} può essere usato solo
  \func{setsockopt} e prende per \param{optval} un intero usato come valore
  logico.

\item[\constd{SO\_MARK}] questa opzione, presente dal kernel 2.6.25, imposta
  un valore di marcatura sui pacchetti del socket. Questa è una funzionalità
  specifica di Linux, ottenibile anche con l'uso del target \texttt{MARK}
  del comando \texttt{iptables} (l'argomento è trattato in sez.~3.3.5 di
  \cite{SGL}). Il valore di marcatura viene mantenuto all'interno dello stack
  di rete del kernel e può essere usato sia dal \itindex{netfilter}
  \textit{netfilter}\footnote{il \textit{netfilter} è l'infrastruttura usata
    per il filtraggio dei pacchetti del kernel, per maggiori dettagli si
    consulti il cap.~2 di \cite{FwGL}.} di Linux che per impostare politiche
  di routing avanzato. Il valore deve essere specificato in \param{optval}
  come un intero. L'opzione richiede i privilegi di amministratore con la
  \textit{capability} \const{CAP\_NET\_ADMIN}.

\item[\constd{SO\_OOBINLINE}] se questa opzione viene abilitata i dati
  \textit{out-of-band} vengono inviati direttamente nel flusso di dati del
  socket (e sono quindi letti con una normale \func{read}) invece che restare
  disponibili solo per l'accesso con l'uso del flag \const{MSG\_OOB} di
  \func{recvmsg}. L'argomento è trattato in dettaglio in
  sez.~\ref{sec:TCP_urgent_data}. L'opzione funziona soltanto con socket che
  supportino i dati \textit{out-of-band} (non ha senso per socket UDP ad
  esempio), ed utilizza per \param{optval} un intero usato come valore logico.

\item[\constd{SO\_PASSCRED}] questa opzione abilita sui socket unix-domain
  (vedi sez.~\ref{sec:unix_socket}) la ricezione dei messaggi di controllo di
  tipo \const{SCM\_CREDENTIALS}. Prende come \param{optval} un intero usato
  come valore logico.

\item[\constd{SO\_PEEK\_OFF}] questa opzione, disponibile a partire dal kernel
  3.4, imposta un ``\textit{peek offset}'' sul socket (attualmente disponibile
  solo per i socket unix-domain (vedi sez.~\ref{sec:unix_socket}). La funzione
  serve come ausilio per l'uso del flag \const{MSG\_PEEK} di \func{recv} che
  consente di ``\textsl{sbirciare}'' nei dati di un socket, cioè di leggerli
  senza rimuoverli dalla coda in cui sono mantenuti, così che vengano
  restituiti anche da una successiva lettura ordinaria.
 
  Un valore negativo (il default è $-1$) riporta alla situazione ordinaria in
  cui si ``\textsl{sbircia}'' a partire dalla testa della coda, un valore
  positivo consente di leggere a partire dalla posizione indicata nella coda e
  tutte le volte che si sbirciano dei dati il valore dell'offset viene
  automaticamente aumentato della quantità di dati sbirciati, in modo che si
  possa proseguire da dove si era arrivati. Il valore deve essere specificato
  in \param{optval} come intero.

\item[\constd{SO\_PEERCRED}] questa opzione restituisce le credenziali del
  processo remoto connesso al socket; l'opzione è disponibile solo per socket
  unix-domain di tipo \textit{stream} e anche per quelli di tipo
  \textit{datagram} quando se ne crea una coppia con \func{socketpair}, e può
  essere usata solo con \func{getsockopt}.  Utilizza per
  \param{optval} una apposita struttura \struct{ucred} (vedi
  sez.~\ref{sec:unix_socket}). 

\item[\constd{SO\_PRIORITY}] questa opzione permette di impostare le priorità
  per tutti i pacchetti che sono inviati sul socket, prende per \param{optval}
  un valore intero. Con questa opzione il kernel usa il valore per ordinare le
  priorità sulle code di rete,\footnote{questo richiede che sia abilitato il
    sistema di \textit{Quality of Service} disponibile con le opzioni di
    routing avanzato.} i pacchetti con priorità più alta vengono processati
  per primi, in modalità che dipendono dalla disciplina di gestione della
  coda. Nel caso di protocollo IP questa opzione permette anche di impostare i
  valori del campo \textit{type of service} (noto come TOS, vedi
  sez.~\ref{sec:IP_header}) per i pacchetti uscenti. Per impostare una
  priorità al di fuori dell'intervallo di valori fra 0 e 6 sono richiesti i
  privilegi di amministratore con la \textit{capability}
  \const{CAP\_NET\_ADMIN}.

\item[\constd{SO\_PROTOCOL}] questa opzione, presente dal kernel 2.6.32, legge
  il protocollo usato dal socket. Funziona solo con \func{getsockopt}, ed
  utilizza per \param{optval} un intero in cui verrà restituito il valore
  numerico che lo identifica (ad esempio \texttt{IPPROTO\_TCP}).

\item[\constd{SO\_RCVBUF}] questa opzione imposta la dimensione del buffer di
  ricezione del socket. Prende per \param{optval} un intero indicante il
  numero di byte. Il valore di default ed il valore massimo che si può
  specificare come argomento per questa opzione sono impostabili tramiti gli
  opportuni valori di \func{sysctl} (vedi sez.~\ref{sec:sock_sysctl}).

  Si tenga presente che nel caso di socket TCP il kernel alloca effettivamente
  una quantità di memoria doppia rispetto a quanto richiesto con
  \func{setsockopt} per entrambe le opzioni \const{SO\_RCVBUF} e
  \const{SO\_SNDBUF}. Questo comporta che una successiva lettura con
  \func{getsockopt} riporterà un valore diverso da quello impostato con
  \func{setsockopt}.  Questo avviene perché TCP necessita dello spazio in più
  per mantenere dati amministrativi e strutture interne, e solo una parte
  viene usata come buffer per i dati, mentre il valore letto da
  \func{getsockopt} e quello riportato nei vari parametri di
  \textit{sysctl}\footnote{cioè \sysctlrelfile{net/core}{wmem\_max} e
    \sysctlrelfile{net/core}{rmem\_max} in \texttt{/proc/sys/net/core} e
    \sysctlrelfile{net/ipv4}{tcp\_wmem} e \sysctlrelfile{net/ipv4}{tcp\_rmem}
    in \texttt{/proc/sys/net/ipv4}, vedi sez.~\ref{sec:sock_sysctl}.} indica
  la memoria effettivamente impiegata.  Si tenga presente inoltre che le
  modifiche alle dimensioni dei buffer di ricezione e trasmissione, per poter
  essere effettive, devono essere impostate prima della chiamata alle funzioni
  \func{listen} o \func{connect}.

\item[\constd{SO\_RCVBUFFORCE}] questa opzione, presente dal kernel 2.6.14, è
  identica a \const{SO\_RCVBUF} ma consente ad un processo con i privilegi di
  amministratore (per la precisione con la \textit{capability}
  \const{CAP\_NET\_ADMIN}) di impostare in valore maggiore del limite di
  \sysctlrelfile{net/core}{rmem\_max}.

\item[\constd{SO\_RCVLOWAT}] questa opzione imposta il valore che indica il
  numero minimo di byte che devono essere presenti nel buffer di ricezione
  perché il kernel passi i dati all'utente, restituendoli ad una \func{read} o
  segnalando ad una \func{select} (vedi sez.~\ref{sec:TCP_sock_select}) che ci
  sono dati in ingresso. L'opzione utilizza per \param{optval} un intero che
  specifica il numero di byte; con Linux questo valore è sempre 1 e può essere
  cambiato solo con i kernel a partire dal 2.4. Si tenga presente però che per
  i kernel prima del 2.6.28 sia \func{poll} che \func{select} non supportano
  questa funzionalità e ritornano comunque, indicando il socket come
  leggibile, non appena almeno un byte è presente, con una successiva lettura
  con \func{read} che si blocca fintanto che non diventa disponibile la
  quantità di byte indicati. Con \func{getsockopt} si può leggere questo
  valore mentre \func{setsockopt} darà un errore di \errcode{ENOPROTOOPT}
  quando il cambiamento non è supportato.

\item[\constd{SO\_RCVTIMEO}] l'opzione permette di impostare un tempo massimo
  sulle operazioni di lettura da un socket, e prende per \param{optval} una
  struttura di tipo \struct{timeval} (vedi fig.~\ref{fig:sys_timeval_struct})
  identica a quella usata con \func{select}. Con \func{getsockopt} si può
  leggere il valore attuale, mentre con \func{setsockopt} si imposta il tempo
  voluto, usando un valore nullo per \struct{timeval} il timeout viene
  rimosso. 

  Se l'opzione viene attivata tutte le volte che una delle funzioni di lettura
  (\func{read}, \func{readv}, \func{recv}, \func{recvfrom} e \func{recvmsg})
  si blocca in attesa di dati per un tempo maggiore di quello impostato, essa
  ritornerà un valore -1 e la variabile \var{errno} sarà impostata con un
  errore di \errcode{EAGAIN} e \errcode{EWOULDBLOCK}, così come sarebbe
  avvenuto se si fosse aperto il socket in modalità non bloccante.\footnote{in
    teoria, se il numero di byte presenti nel buffer di ricezione fosse
    inferiore a quello specificato da \const{SO\_RCVLOWAT}, l'effetto potrebbe
    essere semplicemente quello di provocare l'uscita delle funzioni di
    lettura restituendo il numero di byte fino ad allora ricevuti.}

  In genere questa opzione non è molto utilizzata se si ha a che fare con la
  lettura dei dati, in quanto è sempre possibile usare una \func{select} che
  consente di specificare un \textit{timeout}; l'uso di \func{select} non
  consente però di impostare il timeout per l'uso di \func{connect}, per avere
  il quale si può ricorrere a questa opzione (nel qual caso il raggiungimento
  del \textit{timeout} restituisce l'errore \errcode{EINPROGRESS}). 

\item[\const{SO\_REUSEADDR}] questa opzione permette di eseguire la funzione
  \func{bind} su indirizzi locali che siano già in uso da altri socket;
  l'opzione utilizza per \param{optval} un intero usato come valore logico.
  Questa opzione modifica il comportamento normale dell'interfaccia dei socket
  che fa fallire l'esecuzione della funzione \func{bind} con un errore di
  \errcode{EADDRINUSE} quando l'indirizzo locale (più propriamente il
  controllo viene eseguito sulla porta) è già in uso da parte di un altro
  socket.  Maggiori dettagli sul suo funzionamento sono forniti in
  sez.~\ref{sec:sock_options_main}.

\item[\const{SO\_REUSEPORT}] questa opzione, presente a partire dal kernel
  3.9, permette di far usare a più socket di tipo \const{AF\_INET} o
  \const{AF\_INET6} lo stesso indirizzo locale, e costituisce una estensione
  della precedente \const{SO\_REUSEADDR}.  Maggiori dettagli sul suo
  funzionamento sono forniti in sez.~\ref{sec:sock_options_main}.

\item[\constd{SO\_RXQ\_OVFL}] questa opzione, presente dal kernel 2.6.33,
  permette di abilitare o disabilitare sul socket la ricezione di un messaggio
  ancillare (tratteremo l'argomento in sez.~\ref{sec:net_ancillary_data})
  contenente un intero senza segno a 32 bit che indica il numero di pacchetti
  scartati sul socket fra l'ultimo pacchetto ricevuto e questo.  L'opzione
  utilizza per \param{optval} un intero usato come valore logico.

% TODO documentare SO_RXQ_OVFL, vedi https://lwn.net/Articles/626150/ capire
% che significa 'questo'

\item[\constd{SO\_SNDLOWAT}] questa opzione imposta il valore che indica il
  numero minimo di byte che devono essere presenti nel buffer di trasmissione
  perché il kernel li invii al protocollo successivo, consentendo ad una
  \func{write} di ritornare o segnalando ad una \func{select} (vedi
  sez.~\ref{sec:TCP_sock_select}) che è possibile eseguire una scrittura.
  L'opzione utilizza per \param{optval} un intero che specifica il numero di
  byte; con Linux questo valore è sempre 1 e non può essere cambiato;
  \func{getsockopt} leggerà questo valore mentre \func{setsockopt} darà un
  errore di \errcode{ENOPROTOOPT}.

\item[\constd{SO\_SNDBUF}] questa opzione imposta la dimensione del buffer di
  trasmissione del socket. Prende per \param{optval} un intero indicante il
  numero di byte. Il valore di default ed il valore massimo che si possono
  specificare come argomento per questa opzione sono impostabili
  rispettivamente tramite gli opportuni valori di \func{sysctl} (vedi
  sez.~\ref{sec:sock_sysctl}).

\item[\constd{SO\_SNDBUFFORCE}] questa opzione, presente dal kernel 2.6.14, è
  identica a \const{SO\_SNDBUF} ma consente ad un processo con i privilegi di
  amministratore (per la precisione con la \textit{capability}
  \const{CAP\_NET\_ADMIN}) di impostare in valore maggiore del limite di
  \sysctlrelfile{net/core}{wmem\_max}.

% TODO verificare il timeout con un programma di test

\item[\constd{SO\_SNDTIMEO}] l'opzione permette di impostare un tempo massimo
  sulle operazioni di scrittura su un socket, ed usa gli stessi valori di
  \const{SO\_RCVTIMEO}.  In questo caso però si avrà un errore di
  \errcode{EAGAIN} o \errcode{EWOULDBLOCK} per le funzioni di scrittura
  \func{write}, \func{writev}, \func{send}, \func{sendto} e \func{sendmsg}
  qualora queste restino bloccate per un tempo maggiore di quello specificato. 

\item[\constd{SO\_TIMESTAMP}] questa opzione, presente dal kernel 2.6.30,
  permette di abilitare o disabilitare sul socket la ricezione di un messaggio
  ancillare di tipo \const{SO\_TIMESTAMP}. Il messaggio viene inviato con
  livello \const{SOL\_SOCKET} ed nel campo \var{cmsg\_data} (per i dettagli si
  veda sez.~\ref{sec:net_ancillary_data}) viene impostata una struttura
  \struct{timeval} che indica il tempo di ricezione dell'ultimo pacchetto
  passato all'utente.  L'opzione utilizza per \param{optval} un intero usato
  come valore logico.

% TODO capire meglio la tempistica di questi messaggi ancillari
% TODO documentare SO_TIMESTAMP e le altre opzioni di timestamping dei 
% pacchetti, introdotte nel 2.6.30, vedi nei sorgenti del kernel:
% Documentation/networking/timestamping.txt

\item[\constd{SO\_TYPE}] questa opzione permette di leggere il tipo di socket
  su cui si opera; funziona solo con \func{getsockopt}, ed utilizza per
  \param{optval} un intero in cui verrà restituito il valore numerico che lo
  identifica (ad esempio \const{SOCK\_STREAM}). 


\end{basedescript}


\subsection{L'uso delle principali opzioni dei socket}
\label{sec:sock_options_main}

La descrizione sintetica del significato delle opzioni generiche dei socket,
riportata nell'elenco in sez.~\ref{sec:sock_generic_options}, è
necessariamente sintetica, alcune di queste però possono essere utilizzate
per controllare delle funzionalità che hanno una notevole rilevanza nella
programmazione dei socket.  Per questo motivo faremo in questa sezione un
approfondimento sul significato delle opzioni generiche più importanti.


\constbeg{SO\_KEEPALIVE}
\subsubsection{L'opzione \const{SO\_KEEPALIVE}}

La prima opzione da approfondire è \const{SO\_KEEPALIVE} che permette di
tenere sotto controllo lo stato di una connessione. Una connessione infatti
resta attiva anche quando non viene effettuato alcun traffico su di essa; è
allora possibile, in caso di una interruzione completa della rete, che la
caduta della connessione non venga rilevata, dato che sulla stessa non passa
comunque alcun traffico.

Se si imposta questa opzione, è invece cura del kernel inviare degli appositi
messaggi sulla rete, detti appunto \textit{keep-alive}, per verificare se la
connessione è attiva.  L'opzione funziona soltanto con i socket che supportano
le connessioni (non ha senso per socket UDP ad esempio) e si applica
principalmente ai socket TCP.

Con le impostazioni di default (che sono riprese da BSD) Linux emette un
messaggio di \textit{keep-alive} (in sostanza un segmento ACK vuoto, cui sarà
risposto con un altro segmento ACK vuoto) verso l'altro capo della connessione
se questa è rimasta senza traffico per più di due ore.  Se è tutto a posto il
messaggio viene ricevuto e verrà emesso un segmento ACK di risposta, alla cui
ricezione ripartirà un altro ciclo di attesa per altre due ore di inattività;
il tutto avviene all'interno del kernel e le applicazioni non riceveranno
nessun dato.

Qualora ci siano dei problemi di rete si possono invece verificare i due casi
di terminazione precoce del server già illustrati in
sez.~\ref{sec:TCP_conn_crash}. Il primo è quello in cui la macchina remota ha
avuto un crollo del sistema ed è stata riavviata, per cui dopo il riavvio la
connessione non esiste più.\footnote{si ricordi che un normale riavvio o il
  crollo dell'applicazione non ha questo effetto, in quanto in tal caso si
  passa sempre per la chiusura del processo, e questo, come illustrato in
  sez.~\ref{sec:file_open_close}, comporta anche la regolare chiusura del
  socket con l'invio di un segmento FIN all'altro capo della connessione.} In
questo caso all'invio del messaggio di \textit{keep-alive} si otterrà come
risposta un segmento RST che indica che l'altro capo non riconosce più
l'esistenza della connessione ed il socket verrà chiuso riportando un errore
di \errcode{ECONNRESET}.

Se invece non viene ricevuta nessuna risposta (indice che la macchina non è
più raggiungibile) l'emissione dei messaggi viene ripetuta ad intervalli di 75
secondi per un massimo di 9 volte\footnote{entrambi questi valori possono
  essere modificati a livello di sistema (cioè per tutti i socket) con gli
  opportuni parametri illustrati in sez.~\ref{sec:sock_sysctl} ed a livello di
  singolo socket con le opzioni \texttt{TCP\_KEEP*} di
  sez.~\ref{sec:sock_tcp_udp_options}.}  (per un totale di 11 minuti e 15
secondi) dopo di che, se non si è ricevuta nessuna risposta, il socket viene
chiuso dopo aver impostato un errore di \errcode{ETIMEDOUT}. Qualora la
connessione si sia ristabilita e si riceva un successivo messaggio di risposta
il ciclo riparte come se niente fosse avvenuto.  Infine se si riceve come
risposta un pacchetto ICMP di destinazione irraggiungibile (vedi
sez.~\ref{sec:ICMP_protocol}), verrà restituito l'errore corrispondente.

In generale questa opzione serve per individuare una caduta della connessione
anche quando non si sta facendo traffico su di essa.  Viene usata
principalmente sui server per evitare di mantenere impegnate le risorse che
verrebbero dedicate a trattare delle connessioni che in realtà sono già
terminate (quelle che vengono anche chiamate connessioni
\textsl{semi-aperte}); in tutti quei casi cioè in cui il server si trova in
attesa di dati in ingresso su una connessione che non arriveranno mai o perché
il client sull'altro capo non è più attivo o perché non è più in grado di
comunicare con il server via rete.

\begin{figure}[!htbp]
  \footnotesize \centering
  \begin{minipage}[c]{\codesamplewidth}
    \includecodesample{listati/TCP_echod_fourth.c}
  \end{minipage}
  \normalsize
  \caption{La sezione della nuova versione del server del servizio
    \textit{echo} che prevede l'attivazione del \textit{keepalive} sui
    socket.}
  \label{fig:echod_keepalive_code}
\end{figure}

Abilitandola dopo un certo tempo le connessioni effettivamente terminate
verranno comunque chiuse per cui, utilizzando ad esempio una \func{select}, se
ne potrà rilevare la conclusione e ricevere il relativo errore. Si tenga
presente però che non può avere la certezza assoluta che un errore di
\errcode{ETIMEDOUT} ottenuto dopo aver abilitato questa opzione corrisponda
necessariamente ad una reale conclusione della connessione, il problema
potrebbe anche essere dovuto ad un problema di routing che perduri per un
tempo maggiore di quello impiegato nei vari tentativi di ritrasmissione del
\textit{keep-alive} (anche se questa non è una condizione molto probabile).

Come esempio dell'utilizzo di questa opzione introduciamo all'interno del
nostro server per il servizio \textit{echo} la nuova opzione \texttt{-k} che
permette di attivare il \textit{keep-alive} sui socket; tralasciando la parte
relativa alla gestione di detta opzione (che si limita ad assegnare ad 1 la
variabile \var{keepalive}) tutte le modifiche al server sono riportate in
fig.~\ref{fig:echod_keepalive_code}. Al solito il codice completo è contenuto
nel file \texttt{TCP\_echod\_fourth.c} dei sorgenti allegati alla guida.

Come si può notare la variabile \var{keepalive} è preimpostata (\texttt{\small
  8}) ad un valore nullo; essa viene utilizzata sia come variabile logica per
la condizione (\texttt{\small 14}) che controlla l'attivazione del
\textit{keep-alive} che come valore dell'argomento \param{optval} della
chiamata a \func{setsockopt} (\texttt{\small 16}).  A seconda del suo valore
tutte le volte che un processo figlio viene eseguito in risposta ad una
connessione verrà pertanto eseguita o meno la sezione (\texttt{\small 14--17})
che esegue l'impostazione di \const{SO\_KEEPALIVE} sul socket connesso,
attivando il relativo comportamento.
\constend{SO\_KEEPALIVE}



\constbeg{SO\_REUSEADDR}
\subsubsection{Le opzioni \const{SO\_REUSEADDR} e \const{SO\_REUSEPORT}}

La seconda opzione da approfondire è \constd{SO\_REUSEADDR}, che consente di
eseguire \func{bind} su un socket anche quando la porta specificata è già in
uso da parte di un altro socket. Si ricordi infatti che, come accennato in
sez.~\ref{sec:TCP_func_bind}, normalmente la funzione \func{bind} fallisce con
un errore di \errcode{EADDRINUSE} se la porta scelta è già utilizzata da un
altro socket, proprio per evitare che possano essere lanciati due server sullo
stesso indirizzo e la stessa porta, che verrebbero a contendersi i pacchetti
aventi quella destinazione.

Esistono però situazioni ed esigenze particolari in cui non si vuole che
questo comportamento di salvaguardia accada, ed allora si può fare ricorso a
questa opzione.  La questione è comunque abbastanza complessa in quanto, come
sottolinea Stevens in \cite{UNP1}, si distinguono ben quattro casi diversi in
cui è prevista la possibilità di un utilizzo di questa opzione, il che la
rende una delle più difficili da capire.

Il primo caso in cui si fa ricorso a \const{SO\_REUSEADDR}, che è anche il più
comune, è quello in cui un server è terminato ma esistono ancora dei processi
figli che mantengono attiva almeno una connessione remota che utilizza
l'indirizzo locale, mantenendo occupata la porta. Quando si riesegue il server
allora questo riceve un errore sulla chiamata a \func{bind} dato che la porta
è ancora utilizzata in una connessione esistente.\footnote{questa è una delle
  domande più frequenti relative allo sviluppo, in quanto è piuttosto comune
  trovarsi in questa situazione quando si sta sviluppando un server che si
  ferma e si riavvia in continuazione dopo aver fatto modifiche.}  Inoltre se
si usa il protocollo TCP questo può avvenire anche dopo tutti i processi figli
sono terminati, dato che una connessione può restare attiva anche dopo la
chiusura del socket, mantenendosi nello stato \texttt{TIME\_WAIT} (vedi
sez.~\ref{sec:TCP_time_wait}).

\begin{figure}[!htbp]
  \footnotesize \centering
  \begin{minipage}[c]{\codesamplewidth}
    \includecodesample{listati/sockbindopt.c}
  \end{minipage}
  \normalsize
  \caption{Le sezioni della funzione \texttt{sockbindopt} modificate rispetto al
    codice della precedente \texttt{sockbind}.} 
  \label{fig:sockbindopt_code}
\end{figure}

Usando \const{SO\_REUSEADDR} fra la chiamata a \func{socket} e quella a
\func{bind} si consente a quest'ultima di avere comunque successo anche se la
connessione è attiva (o nello stato \texttt{TIME\_WAIT}). È bene però
ricordare (si riveda quanto detto in sez.~\ref{sec:TCP_time_wait}) che la
presenza dello stato \texttt{TIME\_WAIT} ha una ragione, ed infatti se si usa
questa opzione esiste sempre una probabilità, anche se estremamente
remota,\footnote{perché ciò avvenga infatti non solo devono coincidere gli
  indirizzi IP e le porte degli estremi della nuova connessione, ma anche i
  numeri di sequenza dei pacchetti, e questo è estremamente improbabile.}  che
eventuali pacchetti rimasti intrappolati in una precedente connessione possano
finire fra quelli di una nuova.

Come esempio di uso di questa connessione abbiamo predisposto una nuova
versione della funzione \texttt{sockbind} (vedi fig.~\ref{fig:sockbind_code})
che consenta l'impostazione di questa opzione. La nuova funzione è
\texttt{sockbindopt}, e le principali differenze rispetto alla precedente sono
illustrate in fig.~\ref{fig:sockbindopt_code}, dove si sono riportate le
sezioni di codice modificate rispetto alla versione precedente. Il codice
completo della funzione si trova, insieme alle altre funzioni di servizio dei
socket, all'interno del file \texttt{SockUtils.c} dei sorgenti allegati alla
guida.

In realtà tutto quello che si è fatto è stato introdurre nella nuova funzione
(\texttt{\small 1}) un nuovo argomento intero, \param{reuse}, che conterrà il
valore logico da usare nella successiva chiamata (\texttt{\small 14}) a
\func{setsockopt}. Si è poi aggiunta una sezione (\texttt{\small 13--17}) che
esegue l'impostazione dell'opzione fra la chiamata a \func{socket} e quella a
\func{bind}.

\begin{figure}[!htbp] 
  \footnotesize \centering
  \begin{minipage}[c]{\codesamplewidth}
    \includecodesample{listati/TCP_echod_fifth.c}
  \end{minipage}
  \normalsize
  \caption{Il nuovo codice per l'apertura passiva del server \textit{echo} che
    usa la nuova funzione \texttt{sockbindopt}.}
  \label{fig:TCP_echod_fifth}
\end{figure}

A questo punto basterà modificare il  server per utilizzare la nuova
funzione; in fig.~\ref{fig:TCP_echod_fifth} abbiamo riportato le sezioni
modificate rispetto alla precedente versione di
fig.~\ref{fig:TCP_echod_third}. Al solito il codice completo è coi sorgenti
allegati alla guida, nel file \texttt{TCP\_echod\_fifth.c}.

Anche in questo caso si è introdotta (\texttt{\small 8}) una nuova variabile
\var{reuse} che consente di controllare l'uso dell'opzione e che poi sarà
usata (\texttt{\small 14}) come ultimo argomento di \func{setsockopt}. Il
valore di default di questa variabile è nullo, ma usando l'opzione \texttt{-r}
nell'invocazione del server (al solito la gestione delle opzioni non è
riportata in fig.~\ref{fig:TCP_echod_fifth}) se ne potrà impostare ad 1 il
valore, per cui in tal caso la successiva chiamata (\texttt{\small 13--17}) a
\func{setsockopt} attiverà l'opzione \const{SO\_REUSEADDR}.

Il secondo caso in cui viene usata \const{SO\_REUSEADDR} è quando si ha una
macchina cui sono assegnati diversi indirizzi IP (o come suol dirsi
\textit{multi-homed}) e si vuole porre in ascolto sulla stessa porta un
programma diverso (o una istanza diversa dello stesso programma) per indirizzi
IP diversi. Si ricordi infatti che è sempre possibile indicare a \func{bind}
di collegarsi solo su di un indirizzo specifico; in tal caso se un altro
programma cerca di riutilizzare la stessa porta (anche specificando un
indirizzo diverso) otterrà un errore, a meno di non aver preventivamente
impostato \const{SO\_REUSEADDR}.

Usando questa opzione diventa anche possibile eseguire \func{bind}
sull'indirizzo generico, e questo permetterà il collegamento per tutti gli
indirizzi (di quelli presenti) per i quali la porta non risulti occupata da
una precedente chiamata più specifica. Si tenga presente infatti che con il
protocollo TCP non è in genere possibile far partire più server che eseguano
\func{bind} sullo stesso indirizzo e la stessa porta se su di esso c'è già un
socket in ascolto, cioè ottenere quello che viene chiamato un
\itindex{completely~duplicate~binding} \textit{completely duplicate binding}
(per questo è stata introdotta \const{SO\_REUSEPORT}).

Il terzo impiego è simile al precedente e prevede l'uso di \func{bind}
all'interno dello stesso programma per associare indirizzi locali diversi a
socket diversi. In genere questo viene fatto per i socket UDP quando è
necessario ottenere l'indirizzo a cui sono rivolte le richieste del client ed
il sistema non supporta l'opzione \constd{IP\_RECVDSTADDR};\footnote{nel caso
  di Linux questa opzione è stata supportata per in certo periodo nello
  sviluppo del kernel 2.1.x, ma è in seguito stata soppiantata dall'uso di
  \const{IP\_PKTINFO} (vedi sez.~\ref{sec:sock_ipv4_options}).} in tale modo
si può sapere a quale socket corrisponde un certo indirizzo.  Non ha senso
fare questa operazione per un socket TCP dato che su di essi si può sempre
invocare \func{getsockname} una volta che si è completata la connessione.

Infine il quarto caso è quello in cui si vuole effettivamente ottenere un
\textit{completely duplicate binding}, quando cioè si vuole eseguire
\func{bind} su un indirizzo ed una porta che sono già ``\textsl{legati}'' ad
un altro socket.  Come accennato questo con TCP non è possibile, ed ha anche
poco senso pensare di mettere in ascolto due server sulla stessa porta. Se
però si prende in considerazione il traffico in \textit{multicast}, diventa
assolutamente normale che i pacchetti provenienti dal traffico in
\textit{multicast} possano essere ricevuti da più applicazioni o da diverse
istanze della stessa applicazione sulla stessa porte di un indirizzo di
\textit{multicast}.

Un esempio classico è quello di uno streaming di dati (audio, video, ecc.) in
cui l'uso del \textit{multicast} consente di trasmettere un solo pacchetto da
recapitare a tutti i possibili destinatari (invece di inviarne un duplicato a
ciascuno); in questo caso è perfettamente logico aspettarsi che sulla stessa
macchina più utenti possano lanciare un programma che permetta loro di
ricevere gli stessi dati, e quindi effettuare un \textit{completely duplicate
  binding}.

In questo caso utilizzando \const{SO\_REUSEADDR} si può consentire ad una
applicazione eseguire \func{bind} sulla stessa porta ed indirizzo usata da
un'altra, così che anche essa possa ricevere gli stessi pacchetti. Come detto
la cosa non è possibile con i socket TCP (per i quali il \textit{multicast}
comunque non è applicabile), ma lo è per quelli UDP che è il protocollo
normalmente in uso da parte di queste applicazioni. La regola è che quando si
hanno più applicazioni che hanno eseguito \func{bind} sulla stessa porta, di
tutti pacchetti destinati ad un indirizzo di \textit{broadcast} o di
\textit{multicast} viene inviata una copia a ciascuna applicazione. Non è
definito invece cosa accade qualora il pacchetto sia destinato ad un indirizzo
normale (\textit{unicast}).

% TODO documentare SO_REUSEPORT, vedi https://lwn.net/Articles/542260/

\constbeg{SO\_REUSEPORT}

Esistono però dei casi, in particolare per l'uso di programmi
\textit{multithreaded}, in cui può servire un \textit{completely duplicate
  binding} anche per delle ordinarie connessioni TCP.  Per supportare queste
esigenze a partire dal kernel 3.9 è stata introdotta un'altra opzione,
\const{SO\_REUSEPORT} (già presente in altri sistemi come BSD), che consente
di eseguire il \textit{completely duplicate binding}, fintanto che essa venga
specificata per tutti i socket interessati. Come per \const{SO\_REUSEADDR}
sarà possibile usare l'opzione su un socket legato allo stesso indirizzo e
porta solo se il programma che ha eseguito il primo \func{bind} ha impostato
questa opzione.

Nel caso di \const{SO\_REUSEPORT} oltre al fatto che l'opzione deve essere
attivata sul socket prima di chiamare \func{bind} ed attiva su tutti i socket
con \textit{completely duplicate binding}, è richiesto pure che tutti i
processi che si mettono in ascolto sullo stesso indirizzo e porta abbiano lo
stesso \ids{UID} effettivo, per evitare che un altro utente possa ottenere il
relativo traffico (eseguendo quello che viene definito \textit{port hijacking}
o \textit{port stealing}). L'opzione utilizza per \param{optval} un intero
usato come valore logico.

L'opzione si usa sia per socket TCP che UDP, nel primo caso consente un uso
distribuito di \func{accept} in una applicazione \textit{multithreaded}
passando un diverso \textit{listening socket} ad ogni thread, cosa che
migliora le prestazioni rispetto all'approccio tradizionale di usare un thread
per usare \func{accept} e distribuire le connessioni agli altri o avere più
thread che competono per usare \func{accept} sul socket. Nel caso di UDP
l'opzione consente di distribuire meglio i pacchetti su più processi o thread,
rispetto all'approccio tradizionale di far competere gli stessi per l'accesso
in ricezione al socket.

% TODO documentare SO_REUSEPORT introdotta con il kernel 3.9, vedi
% http://git.kernel.org/linus/c617f398edd4db2b8567a28e899a88f8f574798d TODO:
% in cosa differisce da REUSEADDR?

\constend{SO\_REUSEPORT}
\constend{SO\_REUSEADDR}


\constbeg{SO\_LINGER}
\subsubsection{L'opzione \const{SO\_LINGER}}

La terza opzione da approfondire è \const{SO\_LINGER}; essa, come il nome
suggerisce, consente di ``\textsl{indugiare}'' nella chiusura di un socket. Il
comportamento standard sia di \func{close} che \func{shutdown} è infatti
quello di terminare immediatamente dopo la chiamata, mentre il procedimento di
chiusura della connessione (o di un lato di essa) ed il rispettivo invio sulla
rete di tutti i dati ancora presenti nei buffer, viene gestito in sottofondo
dal kernel.

\begin{figure}[!htb]
  \footnotesize \centering
  \begin{minipage}[c]{0.80\textwidth}
    \includestruct{listati/linger.h}
  \end{minipage}
  \caption{La struttura \structd{linger} richiesta come valore dell'argomento
    \param{optval} per l'impostazione dell'opzione dei socket
    \const{SO\_LINGER}.}
  \label{fig:sock_linger_struct}
\end{figure}

L'uso di \const{SO\_LINGER} con \func{setsockopt} permette di modificare (ed
eventualmente ripristinare) questo comportamento in base ai valori passati nei
campi della struttura \struct{linger}, illustrata in
fig.~\ref{fig:sock_linger_struct}.  Fintanto che il valore del campo
\var{l\_onoff} di \struct{linger} è nullo la modalità che viene impostata
(qualunque sia il valore di \var{l\_linger}) è quella standard appena
illustrata; questa combinazione viene utilizzata per riportarsi al
comportamento normale qualora esso sia stato cambiato da una precedente
chiamata.

Se si utilizza un valore di \var{l\_onoff} diverso da zero, il comportamento
alla chiusura viene a dipendere dal valore specificato per il campo
\var{l\_linger}; se quest'ultimo è nullo l'uso delle funzioni \func{close} e
\func{shutdown} provoca la terminazione immediata della connessione: nel caso
di TCP cioè non viene eseguito il procedimento di chiusura illustrato in
sez.~\ref{sec:TCP_conn_term}, ma tutti i dati ancora presenti nel buffer
vengono immediatamente scartati e sulla rete viene inviato un segmento di RST
che termina immediatamente la connessione.

Un esempio di questo comportamento si può abilitare nel nostro client del
servizio \textit{echo} utilizzando l'opzione \texttt{-r}; riportiamo in
fig.~\ref{fig:TCP_echo_sixth} la sezione di codice che permette di introdurre
questa funzionalità; al solito il codice completo è disponibile nei sorgenti
allegati.

\begin{figure}[!htbp] 
  \footnotesize \centering
  \begin{minipage}[c]{\codesamplewidth}
    \includecodesample{listati/TCP_echo_sixth.c}
  \end{minipage}
  \normalsize
  \caption{La sezione del codice del client \textit{echo} che imposta la
    terminazione immediata della connessione in caso di chiusura.}
  \label{fig:TCP_echo_sixth}
\end{figure}

La sezione indicata viene eseguita dopo aver effettuato la connessione e prima
di chiamare la funzione di gestione, cioè fra le righe (\texttt{\small 12}) e
(\texttt{\small 13}) del precedente esempio di fig.~\ref{fig:TCP_echo_fifth}.
Il codice si limita semplicemente a controllare (\texttt{\small 3}) il
valore della variabile \var{reset} che assegnata nella gestione delle opzioni
in corrispondenza all'uso di \texttt{-r} nella chiamata del client. Nel caso
questa sia diversa da zero vengono impostati (\texttt{\small 5--6}) i valori
della struttura \var{ling} che permettono una terminazione immediata della
connessione. Questa viene poi usata nella successiva (\texttt{\small 7})
chiamata a \func{setsockopt}. Al solito si controlla (\texttt{\small 7--10})
il valore di ritorno e si termina il programma in caso di errore, stampandone
il valore.

Infine l'ultima possibilità, quella in cui si utilizza effettivamente
\const{SO\_LINGER} per \textsl{indugiare} nella chiusura, è quella in cui sia
\var{l\_onoff} che \var{l\_linger} hanno un valore diverso da zero. Se si
esegue l'impostazione con questi valori sia \func{close} che \func{shutdown}
si bloccano, nel frattempo viene eseguita la normale procedura di conclusione
della connessione (quella di sez.~\ref{sec:TCP_conn_term}) ma entrambe le
funzioni non ritornano fintanto che non si sia concluso il procedimento di
chiusura della connessione, o non sia passato un numero di
secondi\footnote{questa è l'unità di misura indicata da POSIX ed adottata da
  Linux, altri kernel possono usare unità di misura diverse, oppure usare il
  campo \var{l\_linger} come valore logico (ignorandone il valore) per rendere
  (quando diverso da zero) \func{close} e \func{shutdown} bloccanti fino al
  completamento della trasmissione dei dati sul buffer.}  pari al valore
specificato in \var{l\_linger}.

\constend{SO\_LINGER}



\subsection{Le opzioni per il protocollo IPv4}
\label{sec:sock_ipv4_options}

Il secondo insieme di opzioni dei socket che tratteremo è quello relativo ai
socket che usano il protocollo IPv4; come per le precedenti opzioni generiche
una descrizione di esse è disponibile nella settima sezione delle pagine di
manuale, nel caso specifico la documentazione si può consultare con
\texttt{man 7 ip}.  

\begin{table}[!htb]
  \centering
  \footnotesize
  \begin{tabular}[c]{|l|c|c|c|l|p{4.8cm}|}
    \hline
    \textbf{Opzione}&\texttt{get}&\texttt{set}&\textbf{flag}&\textbf{Tipo}&
                    \textbf{Descrizione}\\
    \hline
    \hline
    \const{IP\_ADD\_MEMBERSHIP} &         &$\bullet$&   &\struct{ip\_mreqn}& 
      Si unisce a un gruppo di \textit{multicast}.\\
    \const{IP\_ADD\_SOURCE\_MEMBERSHIP} & &$\bullet$&   &\struct{ip\_mreq\_source}& 
      Si unisce a un gruppo di \textit{multicast} per una sorgente.\\
    \const{IP\_BLOCK\_SOURCE} & &$\bullet$&   &\struct{ip\_mreq\_source}& 
      Smette di ricevere dati di \textit{multicast} per una sorgente.\\
    \const{IP\_DROP\_MEMBERSHIP}&         &$\bullet$&   &\struct{ip\_mreqn}& 
      Si sgancia da un gruppo di \textit{multicast}.\\
    \const{IP\_DROP\_SOURCE\_MEMBERSHIP}& &$\bullet$&   &\struct{ip\_mreq\_source}& 
      Si sgancia da un gruppo di \textit{multicast} per una sorgente.\\
    \const{IP\_FREEBIND}       &$\bullet$&$\bullet$&$\bullet$&\texttt{int}& 
      Abilita il \textit{binding} a un indirizzo IP non locale ancora non assegnato.\\
    \const{IP\_HDRINCL}         &$\bullet$&$\bullet$&$\bullet$&\texttt{int}& 
      Passa l'intestazione di IP nei dati.\\
    \const{IP\_MINTTL}          &$\bullet$&$\bullet$&   &\texttt{int}& 
      Imposta il valore minimo del TTL per i pacchetti accettati.\\ 
    \const{IP\_MSFILTER}        &$\bullet$&$\bullet$&   &\struct{ip\_msfilter}& 
      Accesso completo all'interfaccia per il filtraggio delle sorgenti
      \textit{multicast}.\\  
    \const{IP\_MTU}             &$\bullet$&         &         &\texttt{int}& 
      Legge il valore attuale della MTU.\\
    \const{IP\_MTU\_DISCOVER}   &$\bullet$&$\bullet$&         &\texttt{int}& 
      Imposta il \textit{Path MTU Discovery}.\\
    \const{IP\_MULTICAST\_ALL}  &$\bullet$&$\bullet$&$\bullet$&\texttt{int}& 
      Imposta l'interfaccia locale di un socket \textit{multicast}.\\ 
    \const{IP\_MULTICAST\_IF}   &         &$\bullet$&   &\struct{ip\_mreqn}& 
      Imposta l'interfaccia locale di un socket \textit{multicast}.\\ 
    \const{IP\_MULTICAST\_LOOP} &$\bullet$&$\bullet$&$\bullet$&\texttt{int}& 
      Controlla il reinvio a se stessi dei dati di \textit{multicast}.\\ 
    \const{IP\_MULTICAST\_TTL}  &$\bullet$&$\bullet$&         &\texttt{int}& 
      Imposta il TTL per i pacchetti \textit{multicast}.\\
    \const{IP\_NODEFRAG}        &$\bullet$&$\bullet$&$\bullet$&\texttt{int}&
      Disabilita il riassemblaggio di pacchetti frammentati.\\
    \const{IP\_OPTIONS}         &$\bullet$&$\bullet$&&\texttt{void *}& %??? 
      Imposta o riceve le opzioni di IP.\\
    \const{IP\_PKTINFO}         &$\bullet$&$\bullet$&$\bullet$&\texttt{int}& 
      Abilita i messaggi di informazione.\\
    \const{IP\_RECVERR}         &$\bullet$&$\bullet$&$\bullet$&\texttt{int}& 
      Abilita i messaggi di errore affidabili.\\
    \const{IP\_RECVOPTS}        &$\bullet$&$\bullet$&$\bullet$&\texttt{int}& 
      Passa un messaggio con le opzioni IP.\\
    \const{IP\_RECVORIGSTADDR}  &$\bullet$&$\bullet$&$\bullet$&\texttt{int}& 
      Abilita i messaggi con l'indirizzo di destinazione originale.\\
    \const{IP\_RECVTOS}         &$\bullet$&$\bullet$&$\bullet$&\texttt{int}& 
      Passa un messaggio col campo TOS.\\
    \const{IP\_RECVTTL}         &$\bullet$&$\bullet$&$\bullet$&\texttt{int}& 
      Abilita i messaggi col campo TTL.\\
    \const{IP\_RETOPTS}         &$\bullet$&$\bullet$&$\bullet$&\texttt{int}& 
      Abilita i messaggi con le opzioni IP non trattate.\\
    \const{IP\_ROUTER\_ALERT}   &$\bullet$&$\bullet$&$\bullet$&\texttt{int}& 
      Imposta l'opzione \textit{IP router alert} sui pacchetti.\\
    \const{IP\_TOS}             &$\bullet$&$\bullet$&         &\texttt{int}& 
      Imposta il valore del campo TOS.\\
    \const{IP\_TRANSPARENT}     &$\bullet$&$\bullet$&$\bullet$&\texttt{int}& 
      Abilita un \textit{proxing} trasparente sul socket.\\
    \const{IP\_TTL}             &$\bullet$&$\bullet$&         &\texttt{int}& 
      Imposta il valore del campo TTL.\\
    \const{IP\_UNBLOCK\_SOURCE} & &$\bullet$&   &\struct{ip\_mreq\_source}& 
      Ricomincia a ricevere dati di \textit{multicast} per una sorgente.\\
   \hline
  \end{tabular}
  \caption{Le opzioni disponibili al livello \const{IPPROTO\_IP}.} 
  \label{tab:sock_opt_iplevel}
\end{table}

Se si vuole operare su queste opzioni generiche il livello da utilizzare è
\constd{IPPROTO\_IP} (o l'equivalente \const{SOL\_IP}); si è riportato un
elenco di queste opzioni in tab.~\ref{tab:sock_opt_iplevel}.  Le costanti
indicanti le opzioni, e tutte le altre costanti ad esse collegate, sono
definite in \headfiled{netinet/ip.h}, ed accessibili includendo detto file.

Le descrizioni riportate in tab.~\ref{tab:sock_opt_iplevel} sono estremamente
succinte, una maggiore quantità di dettagli sulle varie opzioni è fornita nel
seguente elenco:
\begin{basedescript}{\desclabelwidth{2.0cm}\desclabelstyle{\nextlinelabel}}
\item[\constd{IP\_ADD\_MEMBERSHIP}] L'opzione consente di unirsi ad gruppo di
  \textit{multicast}, e può essere usata solo con \func{setsockopt}.
  L'argomento \param{optval} in questo caso deve essere una struttura di tipo
  \struct{ip\_mreqn}, illustrata in fig.~\ref{fig:ip_mreqn_struct}, che
  permette di indicare, con il campo \var{imr\_multiaddr} l'indirizzo del
  gruppo di \textit{multicast} a cui ci si vuole unire, con il campo
  \var{imr\_address} l'indirizzo dell'interfaccia locale con cui unirsi al
  gruppo di \textit{multicast} e con \var{imr\_ifindex} l'indice
  dell'interfaccia da utilizzare (un valore nullo indica una interfaccia
  qualunque).

\begin{figure}[!htb]
  \footnotesize \centering
  \begin{minipage}[c]{0.90\textwidth}
    \includestruct{listati/ip_mreqn.h}
  \end{minipage}
  \caption{La struttura \structd{ip\_mreqn} utilizzata per unirsi a un gruppo di
    \textit{multicast}.}
  \label{fig:ip_mreqn_struct}
\end{figure}

Questa struttura è presente a partire dal kernel 2.2, per compatibilità è
possibile utilizzare anche un argomento di tipo \struct{ip\_mreq}, presente
fino dal kernel 1.2, che differisce da essa soltanto per l'assenza del campo
\var{imr\_ifindex}; il kernel riconosce il tipo di struttura in base alla
differente dimensione passata in \param{optlen}. 

\item[\constd{IP\_ADD\_SOURCE\_MEMBERSHIP}] L'opzione consente di unirsi ad
  gruppo di \textit{multicast}, ricevendo i dati solo da una sorgente
  specifica; come \const{IP\_ADD\_MEMBERSHIP} può essere usata solo con
  \func{setsockopt}.

\begin{figure}[!htb]
  \footnotesize \centering
  \begin{minipage}[c]{0.90\textwidth}
    \includestruct{listati/ip_mreq_source.h}
  \end{minipage}
  \caption{La struttura \structd{ip\_mreqn} utilizzata per unirsi a un gruppo di
    \textit{multicast} per una specifica sorgente.}
  \label{fig:ip_mreq_source_struct}
\end{figure}

L'argomento \param{optval} in questo caso è una struttura
\struct{ip\_mreq\_source} illustrata in fig.~\ref{fig:ip_mreq_source_struct},
dove il campo \var{imr\_multiaddr} è l'indirizzo del gruppo di
\textit{multicast}, il campo \var{imr\_interface} l'indirizzo dell'interfaccia
locale che deve essere usata per aggiungersi al gruppo di \textit{multicast} e
il campo \var{imr\_sourceaddr} l'indirizzo della sorgente da cui
l'applicazione vuole ricevere i dati. L'opzione può essere ripetuta più volte
per collegarsi a diverse sorgenti.

\item[\constd{IP\_BLOCK\_SOURCE}] Questa opzione, disponibile dal kernel
  2.4.22, consente di smettere di ricevere dati di \textit{multicast} dalla
  sorgente (e relativo gruppo) specificati dalla struttura
  \struct{ip\_mreq\_source} (vedi fig.~\ref{fig:ip_mreq_source_struct})
  passata come argomento \param{optval}. L'opzione è utilizzabile solo se si è
  già registrati nel gruppo di \textit{multicast} indicato con un precedente
  utilizzo di \const{IP\_ADD\_MEMBERSHIP} o
  \const{IP\_ADD\_SOURCE\_MEMBERSHIP}.

\item[\constd{IP\_DROP\_MEMBERSHIP}] Lascia un gruppo di \textit{multicast},
  prende per \param{optval} la stessa struttura \struct{ip\_mreqn} (o
  \struct{ip\_mreq}) usata per \const{IP\_ADD\_MEMBERSHIP} (vedi
  fig.~\ref{fig:ip_mreqn_struct}).

\item[\constd{IP\_DROP\_SOURCE\_MEMBERSHIP}] Lascia un gruppo di
  \textit{multicast} per una specifica sorgente, prende per \param{optval} la
  stessa struttura \struct{ip\_mreq\_source} usata per
  \const{IP\_ADD\_SOURCE\_MEMBERSHIP} (vedi
  fig.~\ref{fig:ip_mreq_source_struct}). Se ci si è registrati per più
  sorgenti nello stesso gruppo, si continuerà a ricevere dati sulle altre. Per
  smettere di ricevere dati da tutte le sorgenti occorre usare l'opzione
  \const{IP\_LEAVE\_GROUP}.

\item[\constd{IP\_FREEBIND}] Se abilitata questa opzione, disponibile dal
  kernel 2.4, consente di usare \func{bind} anche su un indirizzo IP non
  locale o che ancora non è stato assegnato. Questo permette ad una
  applicazione di mettersi in ascolto su un socket prima che l'interfaccia
  sottostante o l'indirizzo che questa deve usare sia stato configurato. È
  l'equivalente a livello di singolo socket dell'uso della \textit{sysctl}
  \texttt{ip\_nonlocal\_bind} che vedremo in
  sez.~\ref{sec:sock_ipv4_sysctl}. Prende per
  \param{optval} un intero usato come valore logico.

\item[\constd{IP\_HDRINCL}] Se viene abilitata questa opzione, presente dal
  kernel 2.0, l'utente deve fornire lui stesso l'intestazione del protocollo
  IP in testa ai propri dati. L'opzione è valida soltanto per socket di tipo
  \const{SOCK\_RAW}, e quando utilizzata eventuali valori impostati con
  \const{IP\_OPTIONS}, \const{IP\_TOS} o \const{IP\_TTL} sono ignorati. In
  ogni caso prima della spedizione alcuni campi dell'intestazione vengono
  comunque modificati dal kernel, torneremo sull'argomento in
  sez.~\ref{sec:socket_raw}.  Prende per \param{optval} un intero usato come
  valore logico.

\item[\constd{IP\_MSFILTER}] L'opzione, introdotta con il kernel 2.4.22,
  fornisce accesso completo all'interfaccia per il filtraggio delle sorgenti
  di \textit{multicast} (il cosiddetto \textit{multicast source filtering},
  definito dall'\href{http://www.ietf.org/rfc/rfc3376.txt}{RFC~3376}).

\begin{figure}[!htb]
  \footnotesize \centering
  \begin{minipage}[c]{0.90\textwidth}
    \includestruct{listati/ip_msfilter.h}
  \end{minipage}
  \caption{La struttura \structd{ip\_msfilter} utilizzata per il
    \textit{multicast source filtering}.}
  \label{fig:ip_msfilter_struct}
\end{figure}

L'argomento \param{optval} deve essere una struttura di tipo
\struct{ip\_msfilter} (illustrata in fig.~\ref{fig:ip_msfilter_struct}); il
campo \var{imsf\_multiaddr} indica l'indirizzo del gruppo di
\textit{multicast}, il campo \var{imsf\_interface} l'indirizzo
dell'interfaccia locale, il campo \var{imsf\_mode} indica la modalità di
filtraggio e con i campi \var{imsf\_numsrc} e \var{imsf\_slist}
rispettivamente la lunghezza della lista, e la lista stessa, degli indirizzi
sorgente.

Come ausilio all'uso di questa opzione sono disponibili le macro
\macro{MCAST\_INCLUDE} e \macro{MCAST\_EXCLUDE} che si possono usare per
\var{imsf\_mode}. Inoltre si può usare la macro
\macro{IP\_MSFILTER\_SIZE}\texttt{(\textsf{n})} per determinare il valore
di \param{optlen} con una struttura \struct{ip\_msfilter} contenente
\texttt{\textsf{n}} sorgenti in \var{imsf\_slist}.

\item[\constd{IP\_MINTTL}] L'opzione, introdotta con il kernel 2.6.34, imposta
  un valore minimo per il campo \textit{Time to Live} dei pacchetti associati
  al socket su cui è attivata, che se non rispettato ne causa lo scarto
  automatico. L'opzione è nata per implementare
  l'\href{http://www.ietf.org/rfc/rfc5082.txt}{RFC~5082} che la prevede come
  forma di protezione per i router che usano il protocollo BGP poiché questi,
  essendo in genere adiacenti, possono, impostando un valore di 255, scartare
  automaticamente tutti gli eventuali pacchetti falsi creati da un attacco a
  questo protocollo, senza doversi curare di verificarne la
  validità.\footnote{l'attacco viene in genere portato per causare un
    \textit{Denial of Service} aumentando il consumo di CPU del router nella
    verifica dell'autenticità di un gran numero di pacchetti falsi; questi,
    arrivando da sorgenti diverse da un router adiacente, non potrebbero più
    avere un TTL di 255 anche qualora questo fosse stato il valore di
    partenza, e l'impostazione dell'opzione consente di scartarli senza carico
    aggiuntivo sulla CPU (che altrimenti dovrebbe calcolare una checksum).}

\itindbeg{Path~MTU}
\item[\constd{IP\_MTU}] Permette di leggere il valore della \textit{Path MTU}
  del socket.  L'opzione richiede per \param{optval} un intero che conterrà il
  valore della \textit{Path MTU} in byte.  Questa è una opzione introdotta con
  i kernel della serie 2.2.x, ed è specifica di Linux.

  È tramite questa opzione che un programma può leggere, quando si è avuto un
  errore di \errval{EMSGSIZE}, il valore della MTU corrente del socket. Si
  tenga presente che per poter usare questa opzione, oltre ad avere abilitato
  la scoperta della \textit{Path MTU}, occorre che il socket sia stato
  esplicitamente connesso con \func{connect}. 

  Ad esempio con i socket UDP si può ottenere una stima iniziale della
  \textit{Path MTU} eseguendo prima una \func{connect} verso la destinazione,
  e poi usando \func{getsockopt} con questa opzione. Si può anche avviare
  esplicitamente il procedimento di scoperta inviando un pacchetto di grosse
  dimensioni (che verrà scartato) e ripetendo l'invio coi dati aggiornati. Si
  tenga infine conto che durante il procedimento i pacchetti iniziali possono
  essere perduti, ed è compito dell'applicazione gestirne una eventuale
  ritrasmissione.

\item[\constd{IP\_MTU\_DISCOVER}] Questa è una opzione introdotta con i kernel
  della serie 2.2.x, ed è specifica di Linux.  L'opzione permette di scrivere
  o leggere le impostazioni della modalità usata per la determinazione della
  \textit{Path MTU} (vedi sez.~\ref{sec:net_lim_dim}) del
  socket. L'opzione prende per \param{optval} un valore intero che indica la
  modalità usata, da specificare con una delle costanti riportate in
  tab.~\ref{tab:sock_ip_mtu_discover}.

  \begin{table}[!htb]
    \centering
    \footnotesize
    \begin{tabular}[c]{|l|r|p{7cm}|}
      \hline
      \multicolumn{2}{|c|}{\textbf{Valore}}&\textbf{Significato} \\
      \hline
      \hline
      \constd{IP\_PMTUDISC\_DONT}&0& Non effettua la ricerca dalla \textit{Path
                                     MTU}.\\
      \constd{IP\_PMTUDISC\_WANT}&1& Utilizza il valore impostato per la rotta
                                     utilizzata dai pacchetti (dal comando
                                     \texttt{route}).\\ 
      \constd{IP\_PMTUDISC\_DO}  &2& Esegue la procedura di determinazione
                                     della \textit{Path MTU} come richiesto
                                     dall'\href{http://www.ietf.org/rfc/rfc1191.txt}{RFC~1191}.\\ 
      \constd{IP\_PMTUDISC\_PROBE}&?& .\\
      \hline
    \end{tabular}
    \caption{Valori possibili per l'argomento \param{optval} di
      \const{IP\_MTU\_DISCOVER}.} 
    \label{tab:sock_ip_mtu_discover}
  \end{table}

  Il valore di default applicato ai socket di tipo \const{SOCK\_STREAM} è
  determinato dal parametro \texttt{ip\_no\_pmtu\_disc} (vedi
  sez.~\ref{sec:sock_sysctl}), mentre per tutti gli altri socket di default la
  ricerca è disabilitata ed è responsabilità del programma creare pacchetti di
  dimensioni appropriate e ritrasmettere eventuali pacchetti persi. Se
  l'opzione viene abilitata, il kernel si incaricherà di tenere traccia
  automaticamente della \textit{Path MTU} verso ciascuna destinazione, e
  rifiuterà immediatamente la trasmissione di pacchetti di dimensioni maggiori
  della MTU con un errore di \errval{EMSGSIZE}.\footnote{in caso contrario la
    trasmissione del pacchetto sarebbe effettuata, ottenendo o un fallimento
    successivo della trasmissione, o la frammentazione dello stesso.}
\itindend{Path~MTU}

\item[\constd{IP\_MULTICAST\_IF}] Imposta l'interfaccia locale per l'utilizzo
  del \textit{multicast}, ed utilizza come \param{optval} le stesse strutture
  \struct{ip\_mreqn} o \struct{ip\_mreq} delle due precedenti opzioni.

\item[\constd{IP\_MULTICAST\_LOOP}] L'opzione consente di decidere se i dati
  che si inviano su un socket usato con il \textit{multicast} vengano ricevuti
  anche sulla stessa macchina da cui li si stanno inviando.  Prende per
  \param{optval} un intero usato come valore logico.

  In generale se si vuole che eventuali client possano ricevere i dati che si
  inviano occorre che questa funzionalità sia abilitata (come avviene di
  default). Qualora però non si voglia generare traffico per dati che già sono
  disponibili in locale l'uso di questa opzione permette di disabilitare
  questo tipo di traffico.

\item[\constd{IP\_MULTICAST\_TTL}] L'opzione permette di impostare o leggere il
  valore del campo TTL per i pacchetti \textit{multicast} in uscita associati
  al socket. È importante che questo valore sia il più basso possibile, ed il
  default è 1, che significa che i pacchetti non potranno uscire dalla rete
  locale. Questa opzione consente ai programmi che lo richiedono di superare
  questo limite.  L'opzione richiede per
  \param{optval} un intero che conterrà il valore del TTL.
% TODO chiarire quale è la struttura \struct{ip\_mreq}

\item[\constd{IP\_OPTIONS}] l'opzione permette di impostare o leggere le
  opzioni del protocollo IP (si veda sez.~\ref{sec:IP_options}). L'opzione
  prende come valore dell'argomento \param{optval} un puntatore ad un buffer
  dove sono mantenute le opzioni, mentre \param{optlen} indica la dimensione
  di quest'ultimo. Quando la si usa con \func{getsockopt} vengono lette le
  opzioni IP utilizzate per la spedizione, quando la si usa con
  \func{setsockopt} vengono impostate le opzioni specificate. L'uso di questa
  opzione richiede una profonda conoscenza del funzionamento del protocollo,
  torneremo in parte sull'argomento in sez.~\ref{sec:sock_IP_options}.

\item[\constd{IP\_PKTINFO}] Quando abilitata l'opzione permette di ricevere
  insieme ai pacchetti un messaggio ancillare (vedi
  sez.~\ref{sec:net_ancillary_data}) di tipo \const{IP\_PKTINFO} contenente
  una struttura \struct{pktinfo} (vedi fig.~\ref{fig:sock_pktinfo_struct}) che
  mantiene una serie di informazioni riguardo i pacchetti in arrivo. In
  particolare è possibile conoscere l'interfaccia su cui è stato ricevuto un
  pacchetto (nel campo \var{ipi\_ifindex}),\footnote{in questo campo viene
    restituito il valore numerico dell'indice dell'interfaccia,
    sez.~\ref{sec:sock_ioctl_netdevice}.} l'indirizzo locale da esso
  utilizzato (nel campo \var{ipi\_spec\_dst}) e l'indirizzo remoto dello
  stesso (nel campo \var{ipi\_addr}).

\begin{figure}[!htb]
  \footnotesize \centering
  \begin{minipage}[c]{0.80\textwidth}
    \includestruct{listati/pktinfo.h}
  \end{minipage}
  \caption{La struttura \structd{pktinfo} usata dall'opzione
    \const{IP\_PKTINFO} per ricavare informazioni sui pacchetti di un socket
    di tipo \const{SOCK\_DGRAM}.}
  \label{fig:sock_pktinfo_struct}
\end{figure}


L'opzione è utilizzabile solo per socket di tipo \const{SOCK\_DGRAM}. Questa è
una opzione introdotta con i kernel della serie 2.2.x, ed è specifica di
Linux;\footnote{non dovrebbe pertanto essere utilizzata se si ha a cuore la
  portabilità.} essa permette di sostituire le opzioni \const{IP\_RECVDSTADDR}
e \const{IP\_RECVIF} presenti in altri Unix (la relativa informazione è quella
ottenibile rispettivamente dai campi \var{ipi\_addr} e \var{ipi\_ifindex} di
\struct{pktinfo}). 

L'opzione prende per \param{optval} un intero usato come valore logico, che
specifica soltanto se insieme al pacchetto deve anche essere inviato o
ricevuto il messaggio \const{IP\_PKTINFO} (vedi
sez.~\ref{sec:net_ancillary_data}); il messaggio stesso dovrà poi essere
letto o scritto direttamente con \func{recvmsg} e \func{sendmsg} (vedi
sez.~\ref{sec:net_sendmsg}).

\item[\constd{IP\_RECVERR}] Questa è una opzione introdotta con i kernel della
  serie 2.2.x, ed è specifica di Linux. Essa permette di usufruire di un
  meccanismo affidabile per ottenere un maggior numero di informazioni in caso
  di errori. Se l'opzione è abilitata tutti gli errori generati su un socket
  vengono memorizzati su una coda, dalla quale poi possono essere letti con
  \func{recvmsg} (vedi sez.~\ref{sec:net_sendmsg}) come messaggi ancillari
  (torneremo su questo in sez.~\ref{sec:net_ancillary_data}) di tipo
  \const{IP\_RECVERR}.  L'opzione richiede per \param{optval} un intero usato
  come valore logico e non è applicabile a socket di tipo
  \const{SOCK\_STREAM}.

\item[\constd{IP\_RECVOPTS}] Quando abilitata l'opzione permette di ricevere
  insieme ai pacchetti un messaggio ancillare (vedi
  sez.~\ref{sec:net_ancillary_data}) di tipo \const{IP\_OPTIONS}, contenente
  le opzioni IP del protocollo (vedi sez.~\ref{sec:IP_options}). Le
  intestazioni di instradamento e le altre opzioni sono già riempite con i
  dati locali. L'opzione richiede per \param{optval} un intero usato come
  valore logico.  L'opzione non è supportata per socket di tipo
  \const{SOCK\_STREAM}.

\item[\constd{IP\_RECVTOS}] Quando abilitata l'opzione permette di ricevere
  insieme ai pacchetti un messaggio ancillare (vedi
  sez.~\ref{sec:net_ancillary_data}) di tipo \const{IP\_TOS}, che contiene un
  byte con il valore del campo \textit{Type of Service} dell'intestazione IP
  del pacchetto stesso (vedi sez.~\ref{sec:IP_header}).  Prende per
  \param{optval} un intero usato come valore logico.

\item[\constd{IP\_RECVTTL}] Quando abilitata l'opzione permette di ricevere
  insieme ai pacchetti un messaggio ancillare (vedi
  sez.~\ref{sec:net_ancillary_data}) di tipo \const{IP\_RECVTTL}, contenente
  un byte con il valore del campo \textit{Time to Live} dell'intestazione IP
  (vedi sez.~\ref{sec:IP_header}).  L'opzione richiede per \param{optval} un
  intero usato come valore logico. L'opzione non è supportata per socket di
  tipo \const{SOCK\_STREAM}.

\item[\constd{IP\_RETOPTS}] Identica alla precedente \const{IP\_RECVOPTS}, ma
  in questo caso restituisce i dati grezzi delle opzioni, senza che siano
  riempiti i capi di instradamento e le marche temporali.  L'opzione richiede
  per \param{optval} un intero usato come valore logico.  L'opzione non è
  supportata per socket di tipo \const{SOCK\_STREAM}.

\item[\constd{IP\_ROUTER\_ALERT}] Questa è una opzione introdotta con i
  kernel della serie 2.2.x, ed è specifica di Linux. Prende per
  \param{optval} un intero usato come valore logico. Se abilitata
  passa tutti i pacchetti con l'opzione \textit{IP Router Alert} (vedi
  sez.~\ref{sec:IP_options}) che devono essere inoltrati al socket
  corrente. Può essere usata soltanto per socket di tipo raw.

\item[\constd{IP\_TOS}] L'opzione consente di leggere o impostare il campo
  \textit{Type of Service} dell'intestazione IP (per una trattazione più
  dettagliata, che riporta anche i valori possibili e le relative costanti di
  definizione si veda sez.~\ref{sec:IP_header}) che permette di indicare le
  priorità dei pacchetti. Se impostato il valore verrà mantenuto per tutti i
  pacchetti del socket; alcuni valori (quelli che aumentano la priorità)
  richiedono i privilegi di amministrazione con la \textit{capability}
  \const{CAP\_NET\_ADMIN}.

  Il campo TOS è di 8 bit e l'opzione richiede per \param{optval} un intero
  che ne contenga il valore. Sono definite anche alcune costanti che
  definiscono alcuni valori standardizzati per il \textit{Type of Service},
  riportate in tab.~\ref{tab:IP_TOS_values}, il valore di default usato da
  Linux è \const{IPTOS\_LOWDELAY}, ma esso può essere modificato con le
  funzionalità del cosiddetto \textit{Advanced Routing}. Si ricordi che la
  priorità dei pacchetti può essere impostata anche in maniera indipendente
  dal protocollo utilizzando l'opzione \const{SO\_PRIORITY} illustrata in
  sez.~\ref{sec:sock_generic_options}.

\item[\constd{IP\_TTL}] L'opzione consente di leggere o impostare per tutti i
  pacchetti associati al socket il campo \textit{Time to Live}
  dell'intestazione IP che indica il numero massimo di \textit{hop} (passaggi
  da un router ad un altro) restanti al paccheto (per una trattazione più
  estesa si veda sez.~\ref{sec:IP_header}).  Il campo TTL è di 8 bit e
  l'opzione richiede che \param{optval} sia un intero, che ne conterrà il
  valore.
\end{basedescript}



\subsection{Le opzioni per i protocolli TCP e UDP}
\label{sec:sock_tcp_udp_options}

In questa sezione tratteremo le varie opzioni disponibili per i socket che
usano i due principali protocolli di comunicazione del livello di trasporto;
UDP e TCP.\footnote{come per le precedenti, una descrizione di queste opzioni
  è disponibile nella settima sezione delle pagine di manuale, che si può
  consultare rispettivamente con \texttt{man 7 tcp} e \texttt{man 7 udp}; le
  pagine di manuale però, alla stesura di questa sezione (Agosto 2006) sono
  alquanto incomplete.}  Dato che questi due protocolli sono entrambi
trasportati su IP,\footnote{qui si sottintende IPv4, ma le opzioni per TCP e
  UDP sono le stesse anche quando si usa IPv6.} oltre alle opzioni generiche
di sez.~\ref{sec:sock_generic_options} saranno comunque disponibili anche le
precedenti opzioni di sez.~\ref{sec:sock_ipv4_options}.\footnote{in realtà in
  sez.~\ref{sec:sock_ipv4_options} si sono riportate le opzioni per IPv4, al
  solito, qualora si stesse utilizzando IPv6, si potrebbero utilizzare le
  opzioni di quest'ultimo.}

Il protocollo che supporta il maggior numero di opzioni è TCP; per poterle
utilizzare occorre specificare \const{SOL\_TCP} (o l'equivalente
\constd{IPPROTO\_TCP}) come valore per l'argomento \param{level}. Si sono
riportate le varie opzioni disponibili in tab.~\ref{tab:sock_opt_tcplevel},
dove sono elencate le rispettive costanti da utilizzare come valore per
l'argomento \param{optname}. Dette costanti e tutte le altre costanti e
strutture collegate all'uso delle opzioni TCP sono definite in
\headfiled{netinet/tcp.h}, ed accessibili includendo detto file.\footnote{in
  realtà questo è il file usato dalle librerie; la definizione delle opzioni
  effettivamente supportate da Linux si trova nel file
  \texttt{include/linux/tcp.h} dei sorgenti del kernel, dal quale si sono
  estratte le costanti di tab.~\ref{tab:sock_opt_tcplevel}.}

\begin{table}[!htb]
  \centering
  \footnotesize
  \begin{tabular}[c]{|l|c|c|c|l|l|}
    \hline
    \textbf{Opzione}&\texttt{get}&\texttt{set}&\textbf{flag}&\textbf{Tipo}&
                    \textbf{Descrizione}\\
    \hline
    \hline
    \const{TCP\_NODELAY}      &$\bullet$&$\bullet$&$\bullet$&\texttt{int}& 
      Spedisce immediatamente i dati in segmenti singoli.\\
    \const{TCP\_MAXSEG}       &$\bullet$&$\bullet$&         &\texttt{int}&
      Valore della MSS per i segmenti in uscita.\\  
    \const{TCP\_CORK}         &$\bullet$&$\bullet$&$\bullet$&\texttt{int}&
      Accumula i dati in un unico segmento.\\
    \const{TCP\_KEEPIDLE}     &$\bullet$&$\bullet$&         &\texttt{int}& 
      Tempo in secondi prima di inviare un \textit{keepalive}.\\
    \const{TCP\_KEEPINTVL}    &$\bullet$&$\bullet$&         &\texttt{int}&
      Tempo in secondi prima fra \textit{keepalive} successivi.\\
    \const{TCP\_KEEPCNT}      &$\bullet$&$\bullet$&         &\texttt{int}& 
      Numero massimo di \textit{keepalive} inviati.\\
    \const{TCP\_SYNCNT}       &$\bullet$&$\bullet$&         &\texttt{int}& 
      Numero massimo di ritrasmissioni di un SYN.\\
    \const{TCP\_LINGER2}      &$\bullet$&$\bullet$&         &\texttt{int}&
      Tempo di vita in stato \texttt{FIN\_WAIT2}.\\
    \const{TCP\_DEFER\_ACCEPT}&$\bullet$&$\bullet$&         &\texttt{int}&
      Ritorna da \func{accept} solo in presenza di dati.\\
    \const{TCP\_WINDOW\_CLAMP}&$\bullet$&$\bullet$&         &\texttt{int}&
      Valore della \textit{advertised window}.\\
    \const{TCP\_INFO}         &$\bullet$&        &       &\struct{tcp\_info}& 
      Restituisce informazioni sul socket.\\
    \const{TCP\_QUICKACK}     &$\bullet$&$\bullet$&$\bullet$&\texttt{int}&
      Abilita la modalità \textit{quickack}.\\
    \const{TCP\_CONGESTION}   &$\bullet$&$\bullet$&         &\texttt{char *}&
      Imposta l'algoritmo per il controllo della congestione.\\
   \hline
  \end{tabular}
  \caption{Le opzioni per i socket TCP disponibili al livello
    \const{SOL\_TCP}.}
  \label{tab:sock_opt_tcplevel}
\end{table}

Le descrizioni delle varie opzioni riportate in
tab.~\ref{tab:sock_opt_tcplevel} sono estremamente sintetiche ed indicative,
la spiegazione del funzionamento delle singole opzioni con una maggiore
quantità di dettagli è fornita nel seguente elenco:
\begin{basedescript}{\desclabelwidth{1.5cm}\desclabelstyle{\nextlinelabel}}

\item[\constd{TCP\_NODELAY}] il protocollo TCP utilizza un meccanismo di
  bufferizzazione dei dati uscenti, per evitare la trasmissione di tanti
  piccoli segmenti con un utilizzo non ottimale della banda
  disponibile.\footnote{il problema è chiamato anche
    \itindex{silly~window~syndrome} \textit{silly window syndrome}, per averne
    un'idea si pensi al risultato che si ottiene quando un programma di
    terminale invia un segmento TCP per ogni tasto premuto, 40 byte di
    intestazione di protocollo con 1 byte di dati trasmessi; per evitare
    situazioni del genere è stato introdotto \index{algoritmo~di~Nagle}
    l'\textsl{algoritmo di Nagle}.}  Questo meccanismo è controllato da un
  apposito algoritmo (detto \textsl{algoritmo di Nagle}, vedi
  sez.~\ref{sec:tcp_protocol_xxx}).  Il comportamento normale del protocollo
  prevede che i dati siano accumulati fintanto che non si raggiunge una
  quantità considerata adeguata per eseguire la trasmissione di un singolo
  segmento.

  Ci sono però delle situazioni in cui questo comportamento può non essere
  desiderabile, ad esempio quando si sa in anticipo che l'applicazione invierà
  soltanto un piccolo quantitativo di dati;\footnote{è il caso classico di una
    richiesta HTTP.} in tal caso l'attesa introdotta dall'algoritmo di
  bufferizzazione non soltanto è inutile, ma peggiora le prestazioni
  introducendo un ritardo.  Impostando questa opzione si disabilita l'uso
  dell'\textsl{algoritmo di Nagle} ed i dati vengono inviati immediatamente in
  singoli segmenti, qualunque sia la loro dimensione.  Ovviamente l'uso di
  questa opzione è dedicato a chi ha esigenze particolari come quella
  illustrata, che possono essere stabilite solo per la singola applicazione.

  Si tenga conto che questa opzione viene sovrascritta dall'eventuale
  impostazione dell'opzione \const{TCP\_CORK} (il cui scopo è sostanzialmente
  l'opposto) che blocca l'invio immediato. Tuttavia quando la si abilita viene
  sempre forzato lo scaricamento della coda di invio (con conseguente
  trasmissione di tutti i dati pendenti), anche qualora si fosse già abilitata
  \const{TCP\_CORK}.\footnote{si tenga presente però che \const{TCP\_CORK} può
    essere specificata insieme a \const{TCP\_NODELAY} soltanto a partire dal
    kernel 2.5.71.}

\item[\constd{TCP\_MAXSEG}] con questa opzione si legge o si imposta il valore
  della MSS (\textit{Maximum Segment Size}, vedi sez.~\ref{sec:net_lim_dim} e
  sez.~\ref{sec:tcp_protocol_xxx}) dei segmenti TCP uscenti. Se l'opzione è
  impostata prima di stabilire la connessione, si cambia anche il valore della
  MSS annunciata all'altro capo della connessione. Se si specificano valori
  maggiori della MTU questi verranno ignorati, inoltre TCP imporrà anche i
  suoi limiti massimo e minimo per questo valore.

\item[\constd{TCP\_CORK}] questa opzione è il complemento naturale di
  \const{TCP\_NODELAY} e serve a gestire a livello applicativo la situazione
  opposta, cioè quella in cui si sa fin dal principio che si dovranno inviare
  grosse quantità di dati. Anche in questo caso l'\textsl{algoritmo di Nagle}
  tenderà a suddividerli in dimensioni da lui ritenute
  opportune,\footnote{l'algoritmo cerca di tenere conto di queste situazioni,
    ma essendo un algoritmo generico tenderà comunque ad introdurre delle
    suddivisioni in segmenti diversi, anche quando potrebbero non essere
    necessarie, con conseguente spreco di banda.}  ma sapendo fin dall'inizio
  quale è la dimensione dei dati si potranno di nuovo ottenere delle migliori
  prestazioni disabilitandolo, e gestendo direttamente l'invio del nostro
  blocco di dati in soluzione unica.

  Quando questa opzione viene abilitata non vengono inviati segmenti di dati
  fintanto che essa non venga disabilitata; a quel punto tutti i dati rimasti
  in coda saranno inviati in un solo segmento TCP. In sostanza con questa
  opzione si può controllare il flusso dei dati mettendo una sorta di
  ``\textsl{tappo}'' (da cui il nome in inglese) al flusso di uscita, in modo
  ottimizzare a mano l'uso della banda. Si tenga presente che per l'effettivo
  funzionamento ci si deve ricordare di disattivare l'opzione al termine
  dell'invio del blocco dei dati.

  Si usa molto spesso \const{TCP\_CORK} quando si effettua il trasferimento
  diretto di un blocco di dati da un file ad un socket con \func{sendfile}
  (vedi sez.~\ref{sec:file_sendfile_splice}), per inserire una intestazione
  prima della chiamata a questa funzione; senza di essa l'intestazione
  potrebbe venire spedita in un segmento a parte, che a seconda delle
  condizioni potrebbe richiedere anche una risposta di ACK, portando ad una
  notevole penalizzazione delle prestazioni.

  Si tenga presente che l'implementazione corrente di \const{TCP\_CORK} non
  consente di bloccare l'invio dei dati per più di 200 millisecondi, passati i
  quali i dati accumulati in coda sanno inviati comunque.  Questa opzione è
  tipica di Linux\footnote{l'opzione è stata introdotta con i kernel della
    serie 2.4.x.} e non è disponibile su tutti i kernel unix-like, pertanto
  deve essere evitata se si vuole scrivere codice portabile.

\item[\constd{TCP\_KEEPIDLE}] con questa opzione si legge o si imposta
  l'intervallo di tempo, in secondi, che deve trascorrere senza traffico sul
  socket prima che vengano inviati, qualora si sia attivata su di esso
  l'opzione \const{SO\_KEEPALIVE}, i messaggi di \textit{keep-alive} (si veda
  la trattazione relativa al \textit{keep-alive} in
  sez.~\ref{sec:sock_options_main}).  Anche questa opzione non è disponibile
  su tutti i kernel unix-like e deve essere evitata se si vuole scrivere
  codice portabile.

\item[\constd{TCP\_KEEPINTVL}] con questa opzione si legge o si imposta
  l'intervallo di tempo, in secondi, fra due messaggi di \textit{keep-alive}
  successivi (si veda sempre quanto illustrato in
  sez.~\ref{sec:sock_options_main}). Come la precedente non è disponibile su
  tutti i kernel unix-like e deve essere evitata se si vuole scrivere codice
  portabile.

\item[\constd{TCP\_KEEPCNT}] con questa opzione si legge o si imposta il numero
  totale di messaggi di \textit{keep-alive} da inviare prima di concludere che
  la connessione è caduta per assenza di risposte ad un messaggio di
  \textit{keep-alive} (di nuovo vedi sez.~\ref{sec:sock_options_main}). Come
  la precedente non è disponibile su tutti i kernel unix-like e deve essere
  evitata se si vuole scrivere codice portabile.

\item[\constd{TCP\_SYNCNT}] con questa opzione si legge o si imposta il numero
  di tentativi di ritrasmissione dei segmenti SYN usati nel \textit{three way
    handshake} prima che il tentativo di connessione venga abortito (si
  ricordi quanto accennato in sez.~\ref{sec:TCP_func_connect}). Sovrascrive
  per il singolo socket il valore globale impostato con la \textit{sysctl}
  \texttt{tcp\_syn\_retries} (vedi sez.~\ref{sec:sock_ipv4_sysctl}).  Non
  vengono accettati valori maggiori di 255; anche questa opzione non è
  standard e deve essere evitata se si vuole scrivere codice portabile.

\item[\constd{TCP\_LINGER2}] con questa opzione si legge o si imposta, in
  numero di secondi, il tempo di sussistenza dei socket terminati nello stato
  \texttt{FIN\_WAIT2} (si ricordi quanto visto in
  sez.~\ref{sec:TCP_conn_term}).\footnote{si tenga ben presente che questa
    opzione non ha nulla a che fare con l'opzione \const{SO\_LINGER} che
    abbiamo visto in sez.~\ref{sec:sock_options_main}.}  Questa opzione
  consente di sovrascrivere per il singolo socket il valore globale impostato
  con la \textit{sysctl} \texttt{tcp\_fin\_timeout} (vedi
  sez.~\ref{sec:sock_ipv4_sysctl}).  Anche questa opzione è da evitare se si
  ha a cuore la portabilità del codice.

\item[\constd{TCP\_DEFER\_ACCEPT}] questa opzione consente di modificare il
  comportamento standard del protocollo TCP nello stabilirsi di una
  connessione; se ricordiamo il meccanismo del \textit{three way handshake}
  illustrato in fig.~\ref{fig:TCP_TWH} possiamo vedere che in genere un client
  inizierà ad inviare i dati ad un server solo dopo l'emissione dell'ultimo
  segmento di ACK.

  Di nuovo esistono situazioni (e la più tipica è quella di una richiesta
  HTTP) in cui sarebbe utile inviare immediatamente la richiesta all'interno
  del segmento con l'ultimo ACK del \textit{three way handshake}; si potrebbe
  così risparmiare l'invio di un segmento successivo per la richiesta e il
  ritardo sul server fra la ricezione dell'ACK e quello della richiesta.

  Se si invoca \const{TCP\_DEFER\_ACCEPT} su un socket dal lato client (cioè
  dal lato da cui si invoca \func{connect}) si istruisce il kernel a non
  inviare immediatamente l'ACK finale del \textit{three way handshake},
  attendendo per un po' di tempo la prima scrittura, in modo da inviare i dati
  di questa insieme col segmento ACK.  Chiaramente la correttezza di questo
  comportamento dipende in maniera diretta dal tipo di applicazione che usa il
  socket; con HTTP, che invia una breve richiesta, permette di risparmiare un
  segmento, con FTP, in cui invece si attende la ricezione del prompt del
  server, introduce un inutile ritardo.

  Allo stesso tempo il protocollo TCP prevede che sul lato del server la
  funzione \func{accept} ritorni dopo la ricezione dell'ACK finale, in tal
  caso quello che si fa usualmente è lanciare un nuovo processo per leggere i
  successivi dati, che si bloccherà su una \func{read} se questi non sono
  disponibili; in questo modo si saranno impiegate delle risorse (per la
  creazione del nuovo processo) che non vengono usate immediatamente.  L'uso
  di \const{TCP\_DEFER\_ACCEPT} consente di intervenire anche in questa
  situazione; quando la si invoca sul lato server (vale a dire su un socket in
  ascolto) l'opzione fa sì che \func{accept} ritorni soltanto quando sono
  presenti dei dati sul socket, e non alla ricezione dell'ACK conclusivo del
  \textit{three way handshake}.

  L'opzione prende un valore intero che indica il numero massimo di secondi
  per cui mantenere il ritardo, sia per quanto riguarda il ritorno di
  \func{accept} su un server, che per l'invio dell'ACK finale insieme ai dati
  su un client. L'opzione è specifica di Linux non deve essere utilizzata in
  codice che vuole essere portabile.\footnote{su FreeBSD è presente una
    opzione \texttt{SO\_ACCEPTFILTER} che consente di ottenere lo stesso
    comportamento di \const{TCP\_DEFER\_ACCEPT} per quanto riguarda il lato
    server.}

\item[\constd{TCP\_WINDOW\_CLAMP}] con questa opzione si legge o si imposta
  alla dimensione specificata, in byte, il valore dichiarato della
  \textit{advertised window} (vedi sez.~\ref{sec:tcp_protocol_xxx}). Il kernel
  impone comunque una dimensione minima pari a \texttt{SOCK\_MIN\_RCVBUF/2}.
  Questa opzione non deve essere utilizzata in codice che vuole essere
  portabile.

\begin{figure}[!htb]
  \footnotesize \centering
  \begin{minipage}[c]{0.80\textwidth}
    \includestruct{listati/tcp_info.h}
  \end{minipage}
  \caption{La struttura \structd{tcp\_info} contenente le informazioni sul
    socket restituita dall'opzione \const{TCP\_INFO}.}
  \label{fig:tcp_info_struct}
\end{figure}

\item[\constd{TCP\_INFO}] questa opzione, specifica di Linux, ma introdotta
  anche in altri kernel (ad esempio FreeBSD) permette di controllare lo stato
  interno di un socket TCP direttamente da un programma in \textit{user space}.
  L'opzione restituisce in una speciale struttura \struct{tcp\_info}, la cui
  definizione è riportata in fig.~\ref{fig:tcp_info_struct}, tutta una serie
  di dati che il kernel mantiene, relativi al socket.  Anche questa opzione
  deve essere evitata se si vuole scrivere codice portabile.

  Con questa opzione diventa possibile ricevere una serie di informazioni
  relative ad un socket TCP così da poter effettuare dei controlli senza dover
  passare attraverso delle operazioni di lettura. Ad esempio si può verificare
  se un socket è stato chiuso usando una funzione analoga a quella illustrata
  in fig.~\ref{fig:is_closing}, in cui si utilizza il valore del campo
  \var{tcpi\_state} di \struct{tcp\_info} per controllare lo stato del socket.

\begin{figure}[!htbp]
  \footnotesize \centering
  \begin{minipage}[c]{\codesamplewidth}
    \includecodesample{listati/is_closing.c}
  \end{minipage}
  \caption{Codice della funzione \texttt{is\_closing.c}, che controlla lo stato
    di un socket TCP per verificare se si sta chiudendo.}
  \label{fig:is_closing}
\end{figure}

%Si noti come nell'esempio si sia (


\item[\constd{TCP\_QUICKACK}] con questa opzione è possibile eseguire una forma
  di controllo sull'invio dei segmenti ACK all'interno di in flusso di dati su
  TCP. In genere questo invio viene gestito direttamente dal kernel, il
  comportamento standard, corrispondente la valore logico di vero (in genere
  1) per questa opzione, è quello di inviare immediatamente i segmenti ACK, in
  quanto normalmente questo significa che si è ricevuto un blocco di dati e si
  può passare all'elaborazione del blocco successivo.

  Qualora però la nostra applicazione sappia in anticipo che alla ricezione di
  un blocco di dati seguirà immediatamente l'invio di un altro
  blocco,\footnote{caso tipico ad esempio delle risposte alle richieste HTTP.}
  poter accorpare quest'ultimo al segmento ACK permette di risparmiare sia in
  termini di dati inviati che di velocità di risposta. Per far questo si può
  utilizzare \const{TCP\_QUICKACK} impostando un valore logico falso (cioè 0),
  in questo modo il kernel attenderà così da inviare il prossimo segmento di
  ACK insieme ai primi dati disponibili.

  Si tenga presente che l'opzione non è permanente, vale a dire che una volta
  che la si sia impostata a 0 il kernel la riporterà al valore di default dopo
  il suo primo utilizzo. Sul lato server la si può impostare anche una volta
  sola su un socket in ascolto, ed essa verrà ereditata da tutti i socket che
  si otterranno da esso al ritorno di \func{accept}.

% TODO trattare con gli esempi di apache

\item[\constd{TCP\_CONGESTION}] questa opzione permette di impostare quale
  algoritmo per il controllo della congestione\footnote{il controllo della
    congestione è un meccanismo previsto dal protocollo TCP (vedi
    sez.~\ref{sec:tcp_protocol_xxx}) per evitare di trasmettere inutilmente
    dati quando una connessione è congestionata; un buon algoritmo è
    fondamentale per il funzionamento del protocollo, dato che i pacchetti
    persi andrebbero ritrasmessi, per cui inviare un pacchetto su una linea
    congestionata potrebbe causare facilmente un peggioramento della
    situazione.} utilizzare per il singolo socket. L'opzione è stata
  introdotta con il kernel 2.6.13,\footnote{alla data di stesura di queste
    note (Set. 2006) è pure scarsamente documentata, tanto che non è neanche
    definita nelle intestazioni della \acr{glibc} per cui occorre definirla a
    mano al suo valore che è 13.} e prende come per \param{optval} il
  puntatore ad un buffer contenente il nome dell'algoritmo di controllo che
  si vuole usare. 

  L'uso di un nome anziché di un valore numerico è dovuto al fatto che gli
  algoritmi di controllo della congestione sono realizzati attraverso
  altrettanti moduli del kernel, e possono pertanto essere attivati a
  richiesta; il nome consente di caricare il rispettivo modulo e di introdurre
  moduli aggiuntivi che implementino altri meccanismi. 

  Per poter disporre di questa funzionalità occorre aver compilato il kernel
  attivando l'opzione di configurazione generale
  \texttt{TCP\_CONG\_ADVANCED},\footnote{disponibile come \textit{TCP:
      advanced congestion control} nel menù \textit{Network->Networking
      options}, che a sua volta renderà disponibile un ulteriore menù con gli
    algoritmi presenti.} e poi abilitare i singoli moduli voluti con le varie
  \texttt{TCP\_CONG\_*} presenti per i vari algoritmi disponibili; un elenco
  di quelli attualmente supportati nella versione ufficiale del kernel è
  riportato in tab.~\ref{tab:sock_tcp_congestion_algo}.\footnote{la lista è
    presa dalla versione 2.6.17.}


  Si tenga presente che prima della implementazione modulare alcuni di questi
  algoritmi erano disponibili soltanto come caratteristiche generali del
  sistema, attivabili per tutti i socket, questo è ancora possibile con la
  \textit{sysctl} \texttt{tcp\_congestion\_control} (vedi
  sez.~\ref{sec:sock_ipv4_sysctl}) che ha sostituito le precedenti
  \textit{sysctl}.\footnote{riportate anche, alla data di stesura di queste
    pagine (Set. 2006) nelle pagine di manuale, ma non più presenti.}

  \begin{table}[!htb]
    \centering
    \footnotesize
    \begin{tabular}[c]{|l|l|l|}
      \hline
      \textbf{Nome}&\textbf{Configurazione}&\textbf{Riferimento} \\
      \hline
      \hline
      reno& -- &Algoritmo tradizionale, usato in caso di assenza degli altri.\\
      \texttt{bic}     &\texttt{TCP\_CONG\_BIC}     & 
      \url{http://www.csc.ncsu.edu/faculty/rhee/export/bitcp/index.htm}.\\
      \texttt{cubic}   &\texttt{TCP\_CONG\_CUBIC}     & 
      \url{http://www.csc.ncsu.edu/faculty/rhee/export/bitcp/index.htm}.\\
      \texttt{highspeed}&\texttt{TCP\_CONG\_HSTCP}  & 
      \url{http://www.icir.org/floyd/hstcp.html}.\\
      \texttt{htcp}    &\texttt{TCP\_CONG\_HTCP}    & 
      \url{http://www.hamilton.ie/net/htcp/}.\\
      \texttt{hybla}   &\texttt{TCP\_CONG\_HYBLA}   &       
      \url{http://www.danielinux.net/projects.html}.\\
      \texttt{scalable}&\texttt{TCP\_CONG\_SCALABLE}&  
      \url{http://www.deneholme.net/tom/scalable/}.\\
      \texttt{vegas}   &\texttt{TCP\_CONG\_VEGAS}   &  
      \url{http://www.cs.arizona.edu/protocols/}.\\
      \texttt{westwood}&\texttt{TCP\_CONG\_WESTWOOD}& 
      \url{http://www.cs.ucla.edu/NRL/hpi/tcpw/}.\\
%      \texttt{}&\texttt{}& .\\
      \hline
    \end{tabular}
    \caption{Gli algoritmi per il controllo della congestione disponibili con
      Linux con le relative opzioni di configurazione da attivare.}  
    \label{tab:sock_tcp_congestion_algo}
  \end{table}

\end{basedescript}


Il protocollo UDP, anche per la sua maggiore semplicità, supporta un numero
ridotto di opzioni, riportate in tab.~\ref{tab:sock_opt_udplevel}; anche in
questo caso per poterle utilizzare occorrerà impostare l'opportuno valore per
l'argomento \param{level}, che è \const{SOL\_UDP} (o l'equivalente
\constd{IPPROTO\_UDP}).  Le costanti che identificano dette opzioni sono
definite in \headfiled{netinet/udp.h}, ed accessibili includendo detto
file.\footnote{come per TCP, la definizione delle opzioni effettivamente
  supportate dal kernel si trova in realtà nel file
  \texttt{include/linux/udp.h}, dal quale si sono estratte le costanti di
  tab.~\ref{tab:sock_opt_udplevel}.}

\begin{table}[!htb]
  \centering
  \footnotesize
  \begin{tabular}[c]{|l|c|c|c|l|l|}
    \hline
    \textbf{Opzione}&\texttt{get}&\texttt{set}&\textbf{flag}&\textbf{Tipo}&
                    \textbf{Descrizione}\\
    \hline
    \hline
    \constd{UDP\_CORK}  &$\bullet$&$\bullet$&$\bullet$&\texttt{int}& %??? 
      Accumula tutti i dati su un unico pacchetto.\\
    \constd{UDP\_ENCAP} &$\bullet$&$\bullet$&$\bullet$&\texttt{int}& %??? 
      Non documentata.\\
   \hline
  \end{tabular}
  \caption{Le opzioni per i socket UDP disponibili al livello
    \const{SOL\_UDP}.}
  \label{tab:sock_opt_udplevel}
\end{table}

% TODO documentare \const{UDP\_ENCAP} 

Ancora una volta le descrizioni contenute tab.~\ref{tab:sock_opt_udplevel}
sono un semplice riferimento, una maggiore quantità di dettagli sulle
caratteristiche delle opzioni citate è quello dell'elenco seguente:
\begin{basedescript}{\desclabelwidth{1.5cm}\desclabelstyle{\nextlinelabel}}

\item[\constd{UDP\_CORK}] questa opzione ha l'identico effetto dell'analoga
  \const{TCP\_CORK} vista in precedenza per il protocollo TCP, e quando
  abilitata consente di accumulare i dati in uscita su un solo pacchetto che
  verrà inviato una volta che la si disabiliti. L'opzione è stata introdotta
  con il kernel 2.5.44, e non deve essere utilizzata in codice che vuole
  essere portabile.

\item[\constd{UDP\_ENCAP}] Questa opzione permette di gestire l'incapsulazione
  dei dati nel protocollo UDP. L'opzione è stata introdotta con il kernel
  2.5.67, e non è documentata. Come la precedente è specifica di Linux e non
  deve essere utilizzata in codice portabile.

\end{basedescript}




\section{La gestione attraverso le funzioni di controllo}
\label{sec:sock_ctrl_func}

Benché la maggior parte delle caratteristiche dei socket sia gestibile con le
funzioni \func{setsockopt} e \func{getsockopt}, alcune proprietà possono
essere impostate attraverso le funzioni \func{fcntl} e \func{ioctl} già
trattate in sez.~\ref{sec:file_fcntl_ioctl}; in quell'occasione abbiamo
parlato di queste funzioni esclusivamente nell'ambito della loro applicazione
a file descriptor associati a dei file normali; qui tratteremo invece i
dettagli del loro utilizzo con file descriptor associati a dei socket.


\subsection{L'uso di \func{ioctl} e \func{fcntl} per i socket generici}
\label{sec:sock_ioctl}

Tratteremo in questa sezione le caratteristiche specifiche delle funzioni
\func{ioctl} e \func{fcntl} quando esse vengono utilizzate con dei socket
generici. Quanto già detto in precedenza sez.~\ref{sec:file_fcntl_ioctl}
continua a valere; quello che tratteremo qui sono le operazioni ed i comandi
che sono validi, o che hanno significati peculiari, quando queste funzioni
vengono applicate a dei socket generici.

Nell'elenco seguente si riportano i valori specifici che può assumere il
secondo argomento della funzione \func{ioctl} (\param{request}, che indica il
tipo di operazione da effettuare) quando essa viene applicata ad un socket
generico. Nell'elenco si illustrerà anche, per ciascuna operazione, il tipo di
dato usato come terzo argomento della funzione ed il significato che esso
viene ad assumere.  Dato che in caso di lettura questi dati vengono restituiti
come \textit{value result argument}, con queste operazioni il terzo argomento
deve sempre essere passato come puntatore ad una variabile (o struttura)
precedentemente allocata. Le costanti che identificano le operazioni sono le
seguenti:
\begin{basedescript}{\desclabelwidth{1.5cm}\desclabelstyle{\nextlinelabel}}
\item[\constd{SIOCGSTAMP}] restituisce il contenuto di una struttura
  \struct{timeval} con la marca temporale dell'ultimo pacchetto ricevuto sul
  socket, questa operazione può essere utilizzata per effettuare delle
  misurazioni precise del tempo di andata e ritorno\footnote{il \textit{Round
      Trip Time} cui abbiamo già accennato in sez.~\ref{sec:net_tcp}.} dei
  pacchetti sulla rete.

\item[\constd{SIOCSPGRP}] imposta il processo o il \textit{process group} a
  cui inviare i segnali \signal{SIGIO} e \signal{SIGURG} quando viene
  completata una operazione di I/O asincrono o arrivano dei dati urgenti
  (\texttt{out-of-band}). Il terzo argomento deve essere un puntatore ad una
  variabile di tipo \type{pid\_t}; un valore positivo indica direttamente il
  \ids{PID} del processo, mentre un valore negativo indica (col valore
  assoluto) il \textit{process group}. Senza privilegi di amministratore o la
  \textit{capability} \const{CAP\_KILL} si può impostare solo se stessi o il
  proprio \textit{process group}.

\item[\constd{SIOCGPGRP}] legge le impostazioni presenti sul socket
  relativamente all'eventuale processo o \textit{process group} cui devono
  essere inviati i segnali \signal{SIGIO} e \signal{SIGURG}. Come per
  \const{SIOCSPGRP} l'argomento passato deve un puntatore ad una variabile di
  tipo \type{pid\_t}, con lo stesso significato.  Qualora non sia presente
  nessuna impostazione verrà restituito un valore nullo.

\item[\const{FIOASYNC}] Abilita o disabilita la modalità di I/O asincrono sul
  socket. Questo significa (vedi sez.~\ref{sec:signal_driven_io}) che verrà
  inviato il segnale di \signal{SIGIO} (o quanto impostato con
  \const{F\_SETSIG}, vedi sez.~\ref{sec:file_fcntl_ioctl}) in caso di eventi
  di I/O sul socket.
\end{basedescript}

Nel caso dei socket generici anche \func{fcntl} prevede un paio di comandi
specifici; in questo caso il secondo argomento (\param{cmd}, che indica il
comando) può assumere i due valori \const{FIOGETOWN} e \const{FIOSETOWN},
mentre il terzo argomento dovrà essere un puntatore ad una variabile di tipo
\type{pid\_t}. Questi due comandi sono una modalità alternativa di eseguire le
stesse operazioni (lettura o impostazione del processo o del gruppo di
processo che riceve i segnali) che si effettuano chiamando \func{ioctl} con
\const{SIOCGPGRP} e \const{SIOCSPGRP}.


\subsection{L'uso di \func{ioctl} per l'accesso ai dispositivi di rete}
\label{sec:sock_ioctl_netdevice}

Benché non strettamente attinenti alla gestione dei socket, vale la pena di
trattare qui l'interfaccia di accesso a basso livello ai dispositivi di rete
che viene appunto fornita attraverso la funzione \texttt{ioctl}. Questa non è
attinente a caratteristiche specifiche di un qualche protocollo, ma si applica
a tutti i socket, indipendentemente da tipo e famiglia degli stessi, e
permette di impostare e rilevare le funzionalità delle interfacce di rete.

\begin{figure}[!htb]
  \footnotesize \centering
  \begin{minipage}[c]{0.80\textwidth}
    \includestruct{listati/ifreq.h}
  \end{minipage}
  \caption{La struttura \structd{ifreq} utilizzata dalle \func{ioctl} per le
    operazioni di controllo sui dispositivi di rete.}
  \label{fig:netdevice_ifreq_struct}
\end{figure}

Tutte le operazioni di questo tipo utilizzano come terzo argomento di
\func{ioctl} il puntatore ad una struttura \struct{ifreq}, la cui definizione
è illustrata in fig.~\ref{fig:netdevice_ifreq_struct}. Normalmente si utilizza
il primo campo della struttura, \var{ifr\_name} per specificare il nome
dell'interfaccia su cui si vuole operare (ad esempio \texttt{eth0},
\texttt{ppp0}, ecc.), e si inseriscono (o ricevono) i valori relativi alle
diversa caratteristiche e funzionalità nel secondo campo, che come si può
notare è definito come una \dirct{union} proprio in quanto il suo significato
varia a secondo dell'operazione scelta.

Si tenga inoltre presente che alcune di queste operazioni (in particolare
quelle che modificano le caratteristiche dell'interfaccia) sono privilegiate e
richiedono i privilegi di amministratore o la \textit{capability}
\const{CAP\_NET\_ADMIN}, altrimenti si otterrà un errore di \errval{EPERM}.
Le costanti che identificano le operazioni disponibili sono le seguenti:
\begin{basedescript}{\desclabelwidth{1.5cm}\desclabelstyle{\nextlinelabel}}
\item[\constd{SIOCGIFNAME}] questa è l'unica operazione che usa il campo
  \var{ifr\_name} per restituire un risultato, tutte le altre lo utilizzano
  per indicare l'interfaccia sulla quale operare. L'operazione richiede che si
  indichi nel campo \var{ifr\_ifindex} il valore numerico dell'\textsl{indice}
  dell'interfaccia, e restituisce il relativo nome in \var{ifr\_name}.

  Il kernel infatti assegna ad ogni interfaccia un numero progressivo, detto
  appunto \itindex{interface~index} \textit{interface index}, che è quello che
  effettivamente la identifica nelle operazioni a basso livello, il nome
  dell'interfaccia è soltanto una etichetta associata a detto \textsl{indice},
  che permette di rendere più comprensibile l'indicazione dell'interfaccia
  all'interno dei comandi. Una modalità per ottenere questo valore è usare il
  comando \cmd{ip link}, che fornisce un elenco delle interfacce presenti
  ordinato in base a tale valore (riportato come primo campo).
  

\item[\constd{SIOCGIFINDEX}] restituisce nel campo \var{ifr\_ifindex} il valore
  numerico dell'indice dell'interfaccia specificata con \var{ifr\_name}, è in
  sostanza l'operazione inversa di \const{SIOCGIFNAME}.

\item[\constd{SIOCGIFFLAGS}] permette di ottenere nel campo \var{ifr\_flags} il
  valore corrente dei flag dell'interfaccia specificata (con \var{ifr\_name}).
  Il valore restituito è una maschera binaria i cui bit sono identificabili
  attraverso le varie costanti di tab.~\ref{tab:netdevice_iface_flag}.

\begin{table}[htb]
  \centering
  \footnotesize
  \begin{tabular}[c]{|l|p{8cm}|}
    \hline
    \textbf{Flag} & \textbf{Significato} \\
    \hline
    \hline
    \constd{IFF\_UP}        & L'interfaccia è attiva.\\
    \constd{IFF\_BROADCAST} & L'interfaccia ha impostato un indirizzo di
                             \textit{broadcast} valido.\\
    \constd{IFF\_DEBUG}     & È attivo il flag interno di debug.\\
    \constd{IFF\_LOOPBACK}  & L'interfaccia è una interfaccia di
                             \textit{loopback}.\\ 
    \constd{IFF\_POINTOPOINT}&L'interfaccia è associata ad un collegamento
                             \textsl{punto-punto}.\\ 
    \constd{IFF\_RUNNING}   & L'interfaccia ha delle risorse allocate (non può
                             quindi essere disattivata).\\
    \constd{IFF\_NOARP}     & L'interfaccia ha il protocollo ARP disabilitato o
                             l'indirizzo del livello di rete non è impostato.\\
    \constd{IFF\_PROMISC}   & L'interfaccia è nel cosiddetto
                             \index{modo~promiscuo} \textsl{modo promiscuo},
                             riceve cioè tutti i pacchetti che vede passare,
                             compresi quelli non direttamente indirizzati a
                             lei.\\
    \constd{IFF\_NOTRAILERS}& Evita l'uso di \textit{trailer} nei pacchetti.\\
    \constd{IFF\_ALLMULTI}  & Riceve tutti i pacchetti di \textit{multicast}.\\
    \constd{IFF\_MASTER}    & L'interfaccia è il master di un bundle per il
                             bilanciamento di carico.\\
    \constd{IFF\_SLAVE}     & L'interfaccia è uno slave di un bundle per il
                             bilanciamento di carico.\\
    \constd{IFF\_MULTICAST} & L'interfaccia ha il supporto per il
                             \textit{multicast} attivo.\\
    \constd{IFF\_PORTSEL}   & L'interfaccia può impostare i suoi parametri
                             hardware (con l'uso di \struct{ifmap}).\\
    \constd{IFF\_AUTOMEDIA} & L'interfaccia è in grado di selezionare
                             automaticamente il tipo di collegamento.\\
    \constd{IFF\_DYNAMIC}   & Gli indirizzi assegnati all'interfaccia vengono
                             persi quando questa viene disattivata.\\
%    \const{IFF\_}      & .\\
    \hline
  \end{tabular}
  \caption{Le costanti che identificano i vari bit della maschera binaria
    \var{ifr\_flags} che esprime i flag di una interfaccia di rete.}
  \label{tab:netdevice_iface_flag}
\end{table}


\item[\constd{SIOCSIFFLAGS}] permette di impostare il valore dei flag
  dell'interfaccia specificata (sempre con \var{ifr\_name}, non staremo a
  ripeterlo oltre) attraverso il valore della maschera binaria da passare nel
  campo \var{ifr\_flags}, che può essere ottenuta con l'OR aritmetico delle
  costanti di tab.~\ref{tab:netdevice_iface_flag}; questa operazione è
  privilegiata.

\item[\constd{SIOCGIFMETRIC}] permette di leggere il valore della metrica del
  dispositivo associato all'interfaccia specificata nel campo
  \var{ifr\_metric}.  Attualmente non è implementato, e l'operazione
  restituisce sempre un valore nullo.

\item[\constd{SIOCSIFMETRIC}] permette di impostare il valore della metrica del
  dispositivo al valore specificato nel campo \var{ifr\_metric}, attualmente
  non ancora implementato, restituisce un errore di \errval{EOPNOTSUPP}.

\item[\constd{SIOCGIFMTU}] permette di leggere il valore della \textit{Maximum
    Transfer Unit} del dispositivo nel campo \var{ifr\_mtu}.

\item[\constd{SIOCSIFMTU}] permette di impostare il valore della
  \textit{Maximum Transfer Unit} del dispositivo al valore specificato campo
  \var{ifr\_mtu}. L'operazione è privilegiata, e si tenga presente che
  impostare un valore troppo basso può causare un blocco del kernel.

\item[\constd{SIOCGIFHWADDR}] permette di leggere il valore dell'indirizzo
  hardware del dispositivo associato all'interfaccia nel campo
  \var{ifr\_hwaddr}; questo viene restituito come struttura \struct{sockaddr}
  in cui il campo \var{sa\_family} contiene un valore \texttt{ARPHRD\_*}
  indicante il tipo di indirizzo ed il campo \var{sa\_data} il valore binario
  dell'indirizzo hardware a partire dal byte 0.

\item[\constd{SIOCSIFHWADDR}] permette di impostare il valore dell'indirizzo
  hardware del dispositivo associato all'interfaccia attraverso il valore
  della struttura \struct{sockaddr} (con lo stesso formato illustrato per
  \const{SIOCGIFHWADDR}) passata nel campo \var{ifr\_hwaddr}. L'operazione è
  privilegiata.

\item[\constd{SIOCSIFHWBROADCAST}] imposta l'indirizzo \textit{broadcast}
  hardware dell'interfaccia al valore specificato dal campo
  \var{ifr\_hwaddr}. L'operazione è privilegiata.

\item[\constd{SIOCGIFMAP}] legge alcuni parametri hardware (memoria, interrupt,
  canali di DMA) del driver dell'interfaccia specificata, restituendo i
  relativi valori nel campo \var{ifr\_map}; quest'ultimo contiene una
  struttura di tipo \struct{ifmap}, la cui definizione è illustrata in
  fig.~\ref{fig:netdevice_ifmap_struct}.

\begin{figure}[!htb]
  \footnotesize \centering
  \begin{minipage}[c]{0.80\textwidth}
    \includestruct{listati/ifmap.h}
  \end{minipage}
  \caption{La struttura \structd{ifmap} utilizzata per leggere ed impostare i
    valori dei parametri hardware di un driver di una interfaccia.}
  \label{fig:netdevice_ifmap_struct}
\end{figure}

\item[\constd{SIOCSIFMAP}] imposta i parametri hardware del driver
  dell'interfaccia specificata, restituendo i relativi valori nel campo
  \var{ifr\_map}. Come per \const{SIOCGIFMAP} questo deve essere passato come
  struttura \struct{ifmap}, secondo la definizione di
  fig.~\ref{fig:netdevice_ifmap_struct}.

\item[\constd{SIOCADDMULTI}] aggiunge un indirizzo di \textit{multicast} ai
  filtri del livello di collegamento associati dell'interfaccia. Si deve usare
  un indirizzo hardware da specificare attraverso il campo \var{ifr\_hwaddr},
  che conterrà l'opportuna struttura \struct{sockaddr}; l'operazione è
  privilegiata. Per una modalità alternativa per eseguire la stessa operazione
  si possono usare i \textit{packet socket}, vedi
  sez.~\ref{sec:packet_socket}.

\item[\constd{SIOCDELMULTI}] rimuove un indirizzo di \textit{multicast} ai
  filtri del livello di collegamento dell'interfaccia, vuole un indirizzo
  hardware specificato come per \const{SIOCADDMULTI}. Anche questa operazione
  è privilegiata e può essere eseguita in forma alternativa con i
  \textit{packet socket}.

\item[\constd{SIOCGIFTXQLEN}] permette di leggere la lunghezza della coda di
  trasmissione del dispositivo associato all'interfaccia specificata nel campo
  \var{ifr\_qlen}.

\item[\constd{SIOCSIFTXQLEN}] permette di impostare il valore della lunghezza
  della coda di trasmissione del dispositivo associato all'interfaccia, questo
  deve essere specificato nel campo \var{ifr\_qlen}. L'operazione è
  privilegiata. 

\item[\constd{SIOCSIFNAME}] consente di cambiare il nome dell'interfaccia
  indicata da \var{ifr\_name} utilizzando il nuovo nome specificato nel campo
  \var{ifr\_rename}.

\end{basedescript}


% TODO aggiunta con il kernel 3.14 SIOCGHWTSTAMP per ottenere il timestamp
% hardware senza modificarlo

Una ulteriore operazione, che consente di ricavare le caratteristiche delle
interfacce di rete, è \constd{SIOCGIFCONF}; però per ragioni di compatibilità
questa operazione è disponibile soltanto per i socket della famiglia
\const{AF\_INET} (vale ad dire per socket IPv4). In questo caso l'utente dovrà
passare come argomento una struttura \struct{ifconf}, definita in
fig.~\ref{fig:netdevice_ifconf_struct}.

\begin{figure}[!htb]
  \footnotesize \centering
  \begin{minipage}[c]{0.80\textwidth}
    \includestruct{listati/ifconf.h}
  \end{minipage}
  \caption{La struttura \structd{ifconf}.}
  \label{fig:netdevice_ifconf_struct}
\end{figure}

Per eseguire questa operazione occorrerà allocare preventivamente un buffer di
contenente un vettore di strutture \struct{ifreq}. La dimensione (in byte) di
questo buffer deve essere specificata nel campo \var{ifc\_len} di
\struct{ifconf}, mentre il suo indirizzo andrà specificato nel campo
\var{ifc\_req}. Qualora il buffer sia stato allocato come una stringa, il suo
indirizzo potrà essere fornito usando il campo \var{ifc\_buf}.\footnote{si
  noti che l'indirizzo del buffer è definito in \struct{ifconf} con una
  \dirct{union}, questo consente di utilizzare una delle due forme a piacere.}

La funzione restituisce nel buffer indicato una serie di strutture
\struct{ifreq} contenenti nel campo \var{ifr\_name} il nome dell'interfaccia e
nel campo \var{ifr\_addr} il relativo indirizzo IP.  Se lo spazio allocato nel
buffer è sufficiente il kernel scriverà una struttura \struct{ifreq} per
ciascuna interfaccia attiva, restituendo nel campo \var{ifc\_len} il totale
dei byte effettivamente scritti. Il valore di ritorno è 0 se l'operazione ha
avuto successo e negativo in caso contrario. 

Si tenga presente che il kernel non scriverà mai sul buffer di uscita dati
eccedenti numero di byte specificato col valore di \var{ifc\_len} impostato
alla chiamata della funzione, troncando il risultato se questi non dovessero
essere sufficienti. Questa condizione non viene segnalata come errore per cui
occorre controllare il valore di \var{ifc\_len} all'uscita della funzione, e
verificare che esso sia inferiore a quello di ingresso. In caso contrario si è
probabilmente\footnote{probabilmente perché si potrebbe essere nella
  condizione in cui sono stati usati esattamente quel numero di byte.} avuta
una situazione di troncamento dei dati.

\begin{figure}[!htbp]
  \footnotesize \centering
  \begin{minipage}[c]{\codesamplewidth}
    \includecodesample{listati/iflist.c}
  \end{minipage}
  \caption{Il corpo principale del programma \texttt{iflist.c}.}
  \label{fig:netdevice_iflist}
\end{figure}

Come esempio dell'uso di queste funzioni si è riportato in
fig.~\ref{fig:netdevice_iflist} il corpo principale del programma
\texttt{iflist} in cui si utilizza l'operazione \const{SIOCGIFCONF} per
ottenere una lista delle interfacce attive e dei relativi indirizzi. Al solito
il codice completo è fornito nei sorgenti allegati alla guida.

Il programma inizia (\texttt{\small 7--11}) con la creazione del socket
necessario ad eseguire l'operazione, dopo di che si inizializzano
opportunamente (\texttt{\small 13--14}) i valori della struttura
\struct{ifconf} indicando la dimensione del buffer ed il suo
indirizzo;\footnote{si noti come in questo caso si sia specificato l'indirizzo
  usando il campo \var{ifc\_buf}, mentre nel seguito del programma si accederà
  ai valori contenuti nel buffer usando \var{ifc\_req}.} si esegue poi
l'operazione invocando \func{ioctl}, controllando come sempre la corretta
esecuzione, ed uscendo in caso di errore (\texttt{\small 15--19}).

Si esegue poi un controllo sulla quantità di dati restituiti segnalando un
eventuale overflow del buffer (\texttt{\small 21--23}); se invece è tutto a
posto (\texttt{\small 24--27}) si calcola e si stampa a video il numero di
interfacce attive trovate.  L'ultima parte del programma (\texttt{\small
  28--33}) è il ciclo sul contenuto delle varie strutture \struct{ifreq}
restituite in cui si estrae (\texttt{\small 30}) l'indirizzo ad esse
assegnato\footnote{si è definito \var{access} come puntatore ad una struttura
  di tipo \struct{sockaddr\_in} per poter eseguire un \textit{casting}
  dell'indirizzo del valore restituito nei vari campi \var{ifr\_addr}, così
  poi da poterlo poi usare come argomento di \func{inet\_ntoa}.}  e lo si
stampa (\texttt{\small 31--32}) insieme al nome dell'interfaccia.



\subsection{L'uso di \func{ioctl} per i socket TCP e UDP}
\label{sec:sock_ioctl_IP}

Non esistono operazioni specifiche per i socket IP in quanto tali,\footnote{a
  parte forse \const{SIOCGIFCONF}, che però resta attinente alle proprietà
  delle interfacce di rete, per cui l'abbiamo trattata in
  sez.~\ref{sec:sock_ioctl_netdevice} insieme alle altre che comunque si
  applicano anche ai socket IP.} mentre per i pacchetti di altri protocolli
trasportati su IP, qualora li si gestisca attraverso dei socket, si dovrà fare
riferimento direttamente all'eventuale supporto presente per il tipo di socket
usato: ad esempio si possono ricevere pacchetti ICMP con socket di tipo
\texttt{raw}, nel qual caso si dovrà fare riferimento alle operazioni di
quest'ultimo.

Tuttavia la gran parte dei socket utilizzati nella programmazione di rete
utilizza proprio il protocollo IP, e quello che succede è che in realtà la
funzione \func{ioctl} consente di effettuare alcune operazioni specifiche per
i socket che usano questo protocollo, ma queste vendono eseguite, invece che a
livello di IP, al successivo livello di trasporto, vale a dire in maniera
specifica per i socket TCP e UDP.

Le operazioni di controllo disponibili per i socket TCP sono illustrate dalla
relativa pagina di manuale, accessibile con \texttt{man 7 tcp}, e prevedono
come possibile valore per il secondo argomento della funzione le costanti
illustrate nell'elenco seguente; il terzo argomento della funzione, gestito
come \textit{value result argument}, deve essere sempre il puntatore ad una
variabile di tipo \ctyp{int}:
\begin{basedescript}{\desclabelwidth{1.5cm}\desclabelstyle{\nextlinelabel}}
\item[\constd{SIOCINQ}] restituisce la quantità di dati non ancora letti
  presenti nel buffer di ricezione; il socket non deve essere in stato
  \texttt{LISTEN}, altrimenti si avrà un errore di \errval{EINVAL}.
\item[\constd{SIOCATMARK}] ritorna un intero non nullo, da intendere come
  valore logico, se il flusso di dati letti sul socket è arrivato sulla
  posizione (detta anche \textit{urgent mark}) in cui sono stati ricevuti dati
  urgenti (vedi sez.~\ref{sec:TCP_urgent_data}).  Una operazione di lettura da
  un socket non attraversa mai questa posizione, per cui è possibile
  controllare se la si è raggiunta o meno con questa operazione.

  Questo è utile quando si attiva l'opzione \const{SO\_OOBINLINE} (vedi
  sez.~\ref{sec:sock_generic_options}) per ricevere i dati urgenti all'interno
  del flusso dei dati ordinari del socket;\footnote{vedremo in
    sez.~\ref{sec:TCP_urgent_data} che in genere i dati urgenti presenti su un
    socket si leggono \textit{out-of-band} usando un opportuno flag per
    \func{recvmsg}.}  in tal caso quando \const{SIOCATMARK} restituisce un
  valore non nullo si saprà che la successiva lettura dal socket restituirà i
  dati urgenti e non il normale traffico; torneremo su questo in maggior
  dettaglio in sez.~\ref{sec:TCP_urgent_data}.

\item[\constd{SIOCOUTQ}] restituisce la quantità di dati non ancora inviati
  presenti nel buffer di spedizione; come per \const{SIOCINQ} il socket non
  deve essere in stato \texttt{LISTEN}, altrimenti si avrà un errore di
  \errval{EINVAL}.
\end{basedescript}

Le operazioni di controllo disponibili per i socket UDP, anch'esse illustrate
dalla relativa pagina di manuale accessibile con \texttt{man 7 udp}, sono
quelle indicate dalle costanti del seguente elenco; come per i socket TCP il
terzo argomento viene gestito come \textit{value result argument} e deve
essere un puntatore ad una variabile di tipo \ctyp{int}:
\begin{basedescript}{\desclabelwidth{1.5cm}\desclabelstyle{\nextlinelabel}}
\item[\const{FIONREAD}] restituisce la dimensione in byte del primo pacchetto
  in attesa di ricezione, o 0 qualora non ci sia nessun pacchetto.
\item[\const{TIOCOUTQ}] restituisce il numero di byte presenti nella coda di
  invio locale; questa opzione è supportata soltanto a partire dal kernel 2.4
\end{basedescript}



\section{La gestione con \func{sysctl} ed il filesystem \texttt{/proc}}
\label{sec:sock_sysctl_proc}

Come ultimo argomento di questo capitolo tratteremo l'uso della funzione
\func{sysctl} (che è stata introdotta nelle sue funzionalità generiche in
sez.~\ref{sec:sys_sysctl}) per quanto riguarda le sue capacità di effettuare
impostazioni relative alle proprietà dei socket.  Dato che le stesse
funzionalità sono controllabili direttamente attraverso il filesystem
\texttt{/proc}, le tratteremo attraverso i file presenti in quest'ultimo.


\subsection{L'uso di \func{sysctl} e \texttt{/proc} per le proprietà della
  rete}
\label{sec:sock_sysctl}

La differenza nell'uso di \func{sysctl} e del filesystem \texttt{/proc}
rispetto a quello delle funzioni \func{ioctl} e \func{fcntl} visto in
sez.~\ref{sec:sock_ctrl_func} o all'uso di \func{getsockopt} e
\func{setsockopt} è che queste funzioni consentono di controllare le proprietà
di un singolo socket, mentre con \func{sysctl} e con \texttt{/proc} si
impostano proprietà (o valori di default) validi a livello dell'intero
sistema, e cioè per tutti i socket.

Le opzioni disponibili per le proprietà della rete, nella gerarchia dei valori
impostabili con \func{sysctl}, sono riportate sotto il nodo \texttt{net}, o,
se acceduti tramite l'interfaccia del filesystem \texttt{/proc}, sotto
\texttt{/proc/sys/net}. In genere sotto questa directory compaiono le
sottodirectory (corrispondenti ad altrettanti sotto-nodi per \func{sysctl})
relative ai vari protocolli e tipi di interfacce su cui è possibile
intervenire per effettuare impostazioni; un contenuto tipico di questa
directory è il seguente:
\begin{verbatim}
/proc/sys/net/
|-- core
|-- ethernet
|-- ipv4
|-- ipv6
|-- irda
|-- token-ring
`-- unix
\end{verbatim}
e sono presenti varie centinaia di parametri, molti dei quali non sono neanche
documentati; nel nostro caso ci limiteremo ad illustrare quelli più
significativi.

Si tenga presente infine che se è sempre possibile utilizzare il filesystem
\texttt{/proc} come sostituto di \func{sysctl}, dato che i valori di nodi e
sotto-nodi di quest'ultima sono mappati come file e directory sotto
\texttt{/proc/sys/}, non è vero il contrario, ed in particolare Linux consente
di impostare alcuni parametri o leggere lo stato della rete a livello di
sistema sotto \texttt{/proc/net}, dove sono presenti dei file che non
corrispondono a nessun nodo di \func{sysctl}.


\subsection{I valori di controllo per i socket generici}
\label{sec:sock_gen_sysctl}

Nella directory \texttt{/proc/sys/net/core/} sono presenti i file
corrispondenti ai parametri generici di \textit{sysctl} validi per tutti i
socket.  Quelli descritti anche nella pagina di manuale, accessibile con
\texttt{man 7 socket} sono i seguenti:

\begin{basedescript}{\desclabelwidth{2.2cm}\desclabelstyle{\nextlinelabel}}
\item[\sysctlrelfiled{net/core}{rmem\_default}] imposta la dimensione
  di default del buffer di ricezione (cioè per i dati in ingresso) dei socket.
\item[\sysctlrelfiled{net/core}{rmem\_max}] imposta la dimensione
  massima che si può assegnare al buffer di ricezione dei socket attraverso
  l'uso dell'opzione \const{SO\_RCVBUF}.
\item[\sysctlrelfiled{net/core}{wmem\_default}] imposta la dimensione
  di default del buffer di trasmissione (cioè per i dati in uscita) dei
  socket.
\item[\sysctlrelfiled{net/core}{wmem\_max}] imposta la dimensione
  massima che si può assegnare al buffer di trasmissione dei socket attraverso
  l'uso dell'opzione \const{SO\_SNDBUF}.
\item[\sysctlrelfiled{net/core}{message\_cost},
  \sysctlrelfiled{net/core}{message\_burst}] contengono le impostazioni del
  \textit{bucket filter} che controlla l'emissione di
  messaggi di avviso da parte del kernel per eventi relativi a problemi sulla
  rete, imponendo un limite che consente di prevenire eventuali attacchi di
  \textit{Denial of Service} usando i log.\footnote{senza questo limite un
    attaccante potrebbe inviare ad arte un traffico che generi
    intenzionalmente messaggi di errore, per saturare il sistema dei log.}

  \itindbeg{bucket~filter} 

  Il \textit{bucket filter} è un algoritmo generico che permette di impostare
  dei limiti di flusso su una quantità\footnote{uno analogo viene usato per
    imporre dei limiti sul flusso dei pacchetti nel \textit{netfilter} di
    Linux.}  senza dovere eseguire medie temporali, che verrebbero a dipendere
  in misura non controllabile dalla dimensione dell'intervallo su cui si media
  e dalla distribuzione degli eventi;\footnote{in caso di un picco di flusso
    (il cosiddetto \textit{burst}) il flusso medio verrebbe a dipendere in
    maniera esclusiva dalla dimensione dell'intervallo di tempo su cui calcola
    la media.} in questo caso si definisce la dimensione di un
  ``\textsl{bidone}'' (il \textit{bucket}) e del flusso che da esso può
  uscire, la presenza di una dimensione iniziale consente di assorbire
  eventuali picchi di emissione, l'aver fissato un flusso di uscita garantisce
  che a regime questo sarà il valore medio del flusso ottenibile dal
  \textit{bucket}.

  \itindend{bucket~filter} 

  I due valori indicano rispettivamente il flusso a regime (non sarà inviato
  più di un messaggio per il numero di secondi specificato da
  \texttt{message\_cost}) e la dimensione iniziale per in caso di picco di
  emissione (verranno accettati inizialmente fino ad un massimo di
  \texttt{message\_cost/message\_burst} messaggi).

\item[\sysctlrelfiled{net/core}{netdev\_max\_backlog}] numero massimo
  di pacchetti che possono essere contenuti nella coda di ingresso generale.

\item[\sysctlrelfiled{net/core}{optmem\_max}] lunghezza massima dei
  dati ancillari e di controllo (vedi sez.~\ref{sec:net_ancillary_data}).
\end{basedescript}

Oltre a questi, nella directory \texttt{/proc/sys/net/core} si trovano diversi
altri file, la cui documentazione, come per gli altri, dovrebbe essere
mantenuta nei sorgenti del kernel nel file
\texttt{Documentation/networking/ip-sysctl.txt}; la maggior parte di questi
però non è documentato:
\begin{basedescript}{\desclabelwidth{2.2cm}\desclabelstyle{\nextlinelabel}}
\item[\sysctlrelfiled{net/core}{dev\_weight}] blocco di lavoro (\textit{work
    quantum}) dello \textit{scheduler} di processo dei pacchetti.

% TODO da documentare meglio

\item[\sysctlrelfiled{net/core}{lo\_cong}] valore per l'occupazione
  della coda di ricezione sotto la quale si considera di avere una bassa
  congestione.

\item[\sysctlrelfiled{net/core}{mod\_cong}] valore per l'occupazione della
  coda di ricezione sotto la quale si considera di avere una congestione
  moderata.

\item[\sysctlrelfiled{net/core}{no\_cong}] valore per l'occupazione della coda
  di ricezione sotto la quale si considera di non avere congestione.

\item[\sysctlrelfiled{net/core}{no\_cong\_thresh}] valore minimo (\textit{low
    water mark}) per il riavvio dei dispositivi congestionati.

  % \item[\sysctlrelfiled{net/core}{netdev\_fastroute}] è presente
  %   soltanto quando si è compilato il kernel con l'apposita opzione di
  %   ottimizzazione per l'uso come router.

\item[\sysctlrelfiled{net/core}{somaxconn}] imposta la dimensione massima
  utilizzabile per il \textit{backlog} della funzione \func{listen} (vedi
  sez.~\ref{sec:TCP_func_listen}), e corrisponde al valore della costante
  \constd{SOMAXCONN}; il suo valore di default è 128.

\end{basedescript}


\subsection{I valori di controllo per il protocollo IPv4}
\label{sec:sock_ipv4_sysctl}

Nella directory \texttt{/proc/sys/net/ipv4} sono presenti i file che
corrispondono ai parametri dei socket che usano il protocollo IPv4, relativi
quindi sia alle caratteristiche di IP, che a quelle degli altri protocolli che
vengono usati all'interno di quest'ultimo (come ICMP, TCP e UDP) o a fianco
dello stesso (come ARP).

I file che consentono di controllare le caratteristiche specifiche del
protocollo IP in quanto tale, che sono descritti anche nella relativa pagina
di manuale accessibile con \texttt{man 7 ip}, sono i seguenti:
\begin{basedescript}{\desclabelwidth{2.2cm}\desclabelstyle{\nextlinelabel}}

\item[\sysctlrelfiled{net/ipv4}{ip\_default\_ttl}] imposta il valore di
  default per il campo TTL (vedi sez.~\ref{sec:IP_header}) di tutti i
  pacchetti uscenti, stabilendo così il numero massimo di router che i
  pacchetti possono attraversare. Il valore può essere modificato anche per il
  singolo socket con l'opzione \const{IP\_TTL}.  Prende un valore intero, ma
  dato che il campo citato è di 8 bit hanno senso solo valori fra 0 e 255. Il
  valore di default è 64, e normalmente non c'è nessuna necessità di
  modificarlo.\footnote{l'unico motivo sarebbe per raggiungere macchine
    estremamente ``{lontane}'' in termini di \textit{hop}, ma è praticamente
    impossibile trovarne.} Aumentare il valore è una pratica poco gentile, in
  quanto in caso di problemi di routing si allunga inutilmente il numero di
  ritrasmissioni.

\item[\sysctlrelfiled{net/ipv4}{ip\_forward}] abilita l'inoltro dei
  pacchetti da una interfaccia ad un altra, e può essere impostato anche per
  la singola interfaccia. Prende un valore logico (0 disabilita, diverso da
  zero abilita), di default è disabilitato.

\item[\sysctlrelfiled{net/ipv4}{ip\_dynaddr}] abilita la riscrittura
  automatica degli indirizzi associati ad un socket quando una interfaccia
  cambia indirizzo. Viene usato per le interfacce usate nei collegamenti in
  dial-up, il cui indirizzo IP viene assegnato dinamicamente dal provider, e
  può essere modificato. Prende un valore intero, con 0 si disabilita la
  funzionalità, con 1 la si abilita, con 2 (o con qualunque altro valore
  diverso dai precedenti) la si abilità in modalità \textsl{prolissa}; di
  default la funzionalità è disabilitata.

\item[\sysctlrelfiled{net/ipv4}{ip\_autoconfig}] specifica se
  l'indirizzo IP è stato configurato automaticamente dal kernel all'avvio
  attraverso DHCP, BOOTP o RARP. Riporta un valore logico (0 falso, 1 vero)
  accessibile solo in lettura, è inutilizzato nei kernel recenti ed eliminato
  a partire dal kernel 2.6.18.

\item[\sysctlrelfiled{net/ipv4}{ip\_local\_port\_range}] imposta l'intervallo
  dei valori usati per l'assegnazione delle porte effimere, permette cioè di
  modificare i valori illustrati in fig.~\ref{fig:TCP_port_alloc}; prende due
  valori interi separati da spazi, che indicano gli estremi
  dell'intervallo. Si abbia cura di non definire un intervallo che si
  sovrappone a quello delle porte usate per il \textit{masquerading}, il
  kernel può gestire la sovrapposizione, ma si avrà una perdita di
  prestazioni. Si imposti sempre un valore iniziale maggiore di 1024 (o meglio
  ancora di 4096) per evitare conflitti con le porte usate dai servizi noti.

\item[\sysctlrelfiled{net/ipv4}{ip\_no\_pmtu\_disc}] permette di disabilitare
  per i socket \const{SOCK\_STREAM} la ricerca automatica della \textit{Path
    MTU} (vedi sez.~\ref{sec:net_lim_dim} e
  sez.~\ref{sec:sock_ipv4_options}). Prende un valore logico, e di default è
  disabilitato (cioè la ricerca viene eseguita).

  In genere si abilita questo parametro quando per qualche motivo il
  procedimento del \textit{Path MTU discovery} fallisce; dato che questo può
  avvenire a causa di router\footnote{ad esempio se si scartano tutti i
    pacchetti ICMP, il problema è affrontato anche in sez.~3.4.4 di
    \cite{SGL}.} o interfacce\footnote{ad esempio se i due capi di un
    collegamento \textit{point-to-point} non si accordano sulla stessa MTU.}
  mal configurati è opportuno correggere le configurazioni, perché
  disabilitare globalmente il procedimento con questo parametro ha pesanti
  ripercussioni in termini di prestazioni di rete.

\item[\sysctlrelfiled{net/ipv4}{ip\_always\_defrag}] fa sì che tutti i
  pacchetti IP frammentati siano riassemblati, anche in caso in successivo
  immediato inoltro.\footnote{introdotto con il kernel 2.2.13, nelle versioni
    precedenti questo comportamento poteva essere solo stabilito un volta per
    tutte in fase di compilazione del kernel con l'opzione
    \texttt{CONFIG\_IP\_ALWAYS\_DEFRAG}.} Prende un valore logico e di default
  è disabilitato. Con i kernel dalla serie 2.4 in poi la deframmentazione
  viene attivata automaticamente quando si utilizza il sistema del
  \textit{netfilter}, e questo parametro non è più presente.

\item[\sysctlrelfiled{net/ipv4}{ipfrag\_high\_thresh}] indica il limite
  massimo (espresso in numero di byte) sui pacchetti IP frammentati presenti
  in coda; quando questo valore viene raggiunto la coda viene ripulita fino al
  valore \texttt{ipfrag\_low\_thresh}. Prende un valore intero.

\item[\sysctlrelfiled{net/ipv4}{ipfrag\_low\_thresh}] indica la dimensione
  (specificata in byte) della soglia inferiore a cui viene riportata la coda
  dei pacchetti IP frammentati quando si raggiunge il valore massimo dato da
  \texttt{ipfrag\_high\_thresh}. Prende un valore intero.

\item[\sysctlrelfiled{net/ipv4}{ip\_nonlocal\_bind}] se abilitato rende
  possibile ad una applicazione eseguire \func{bind} anche su un indirizzo che
  non è presente su nessuna interfaccia locale. Prende un valore logico e di
  default è disabilitato. La funzionalità può essere abilitata a livello di
  singolo socket con l'opzione \const{IP\_FREEBIND} illustrata in
  sez.~\ref{sec:sock_ipv4_options}. 

  Questo può risultare utile per applicazioni particolari (come gli
  \textit{sniffer}) che hanno la necessità di ricevere pacchetti anche non
  diretti agli indirizzi presenti sulla macchina, ad esempio per intercettare
  il traffico per uno specifico indirizzo che si vuole tenere sotto
  controllo. Il suo uso però può creare problemi ad alcune applicazioni.

% \item[\texttt{neigh/*}] La directory contiene i valori 
% TODO trattare neigh/* nella parte su arp, da capire dove sarà.

\end{basedescript}

I file di \texttt{/proc/sys/net/ipv4} che invece fanno riferimento alle
caratteristiche specifiche del protocollo TCP, elencati anche nella rispettiva
pagina di manuale (accessibile con \texttt{man 7 tcp}), sono i seguenti:
\begin{basedescript}{\desclabelwidth{2.2cm}\desclabelstyle{\nextlinelabel}}

\item[\sysctlrelfiled{net/ipv4}{tcp\_abort\_on\_overflow}] indica al
  kernel di azzerare le connessioni quando il programma che le riceve è troppo
  lento ed incapace di accettarle. Prende un valore logico ed è disabilitato
  di default.  Questo consente di recuperare le connessioni se si è avuto un
  eccesso dovuto ad un qualche picco di traffico, ma ovviamente va a discapito
  dei client che interrogano il server. Pertanto è da abilitare soltanto
  quando si è sicuri che non è possibile ottimizzare il server in modo che sia
  in grado di accettare connessioni più rapidamente.

\item[\sysctlrelfiled{net/ipv4}{tcp\_adv\_win\_scale}] indica al kernel quale
  frazione del buffer associato ad un socket\footnote{quello impostato con
    \sysctlrelfile{net/ipv4}{tcp\_rmem}.} deve essere utilizzata per la
  finestra del protocollo TCP\footnote{in sostanza il valore che costituisce
    la \textit{advertised window} annunciata all'altro capo del socket.} e
  quale come buffer applicativo per isolare la rete dalle latenze
  dell'applicazione.  Prende un valore intero che determina la suddetta
  frazione secondo la formula
  $\texttt{buffer}/2^\texttt{tcp\_adv\_win\_scale}$ se positivo o con
  $\texttt{buffer}-\texttt{buffer}/2^\texttt{tcp\_adv\_win\_scale}$ se
  negativo.  Il default è 2 che significa che al buffer dell'applicazione
  viene riservato un quarto del totale.

\item[\sysctlrelfiled{net/ipv4}{tcp\_app\_win}] indica la frazione della
  finestra TCP che viene riservata per gestire l'overhaed dovuto alla
  bufferizzazione. Prende un valore intero che consente di calcolare la
  dimensione in byte come il massimo fra la MSS e
  $\texttt{window}/2^\texttt{tcp\_app\_win}$. Un valore nullo significa che
  non viene riservato nessuno spazio; il valore di default è 31.

% vecchi, presumibilmente usati quando gli algoritmi di congestione non erano
% modularizzabili 
% \item[\texttt{tcp\_bic}]
% \item[\texttt{tcp\_bic\_low\_window}] 
% \item[\texttt{tcp\_bic\_fast\_convergence}] 

\item[\sysctlrelfiled{net/ipv4}{tcp\_dsack}] abilita il supporto,
  definito nell'\href{http://www.ietf.org/rfc/rfc2884.txt}{RFC~2884}, per il
  cosiddetto \textit{Duplicate SACK}.\footnote{si indica con SACK
    (\textit{Selective Acknowledgement}) un'opzione TCP, definita
    nell'\href{http://www.ietf.org/rfc/rfc2018.txt}{RFC~2018}, usata per dare
    un \textit{acknowledgement} unico su blocchi di pacchetti non contigui,
    che consente di diminuire il numero di pacchetti scambiati.} Prende un
  valore logico e di default è abilitato.
% TODO documentare o descrivere che cos'è il Duplicate SACK o
% mettere riferimento nelle appendici


\item[\sysctlrelfiled{net/ipv4}{tcp\_ecn}] abilita il meccanismo della
  \textit{Explicit Congestion Notification} (in breve ECN) nelle connessioni
  TCP. Prende valore logico che di default è disabilitato. La \textit{Explicit
    Congestion Notification} \itindex{Explicit~Congestion~Notification~(ECN)}
  è un meccanismo che consente di notificare quando una rotta o una rete è
  congestionata da un eccesso di traffico,\footnote{il meccanismo è descritto
    in dettaglio nell'\href{http://www.ietf.org/rfc/rfc3168.txt}{RFC~3168}
    mentre gli effetti sulle prestazioni del suo utilizzo sono documentate
    nell'\href{http://www.ietf.org/rfc/rfc2884.txt}{RFC~2884}.} si può così
  essere avvisati e cercare rotte alternative oppure diminuire l'emissione di
  pacchetti (in modo da non aumentare la congestione).

  Si tenga presente che se si abilita questa opzione si possono avere dei
  malfunzionamenti apparentemente casuali dipendenti dalla destinazione,
  dovuti al fatto che alcuni vecchi router non supportano il meccanismo ed
  alla sua attivazione scartano i relativi pacchetti, bloccando completamente
  il traffico.
% TODO documentare o descrivere che cos'è l'Explicit Congestion Notification o
% mettere riferimento nelle appendici


\item[\sysctlrelfiled{net/ipv4}{tcp\_fack}] abilita il supporto per il
  \textit{TCP Forward Acknowledgement}, un algoritmo per il controllo della
  congestione del traffico. Prende un valore logico e di default è abilitato.

% TODO documentare o descrivere che cos'è il TCP Forward Acknowledgement o
% mettere riferimento nelle appendici

\item[\sysctlrelfiled{net/ipv4}{tcp\_fin\_timeout}] specifica il numero di
  secondi da passare in stato \texttt{FIN\_WAIT2} nell'attesa delle ricezione
  del pacchetto FIN conclusivo, passati quali il socket viene comunque chiuso
  forzatamente.  Prende un valore intero che indica i secondi e di default è
  60.\footnote{nei kernel della serie 2.2.x era il valore utilizzato era
    invece di 120 secondi.} L'uso di questa opzione realizza quella che in
  sostanza è una violazione delle specifiche del protocollo TCP, ma è utile
  per fronteggiare alcuni attacchi di \textit{Denial of Service}.

\item[\sysctlrelfiled{net/ipv4}{tcp\_frto}] abilita il supporto per
  l'algoritmo F-RTO, un algoritmo usato per la ritrasmissione dei timeout del
  protocollo TCP, che diventa molto utile per le reti wireless dove la perdita
  di pacchetti è usualmente dovuta a delle interferenze radio, piuttosto che
  alla congestione dei router. Prende un valore logico e di default è
  disabilitato.

\item[\sysctlrelfiled{net/ipv4}{tcp\_keepalive\_intvl}] indica il
  numero di secondi che deve trascorrere fra l'emissione di due successivi
  pacchetti di test quando è abilitata la funzionalità del \textit{keepalive}
  (vedi sez.~\ref{sec:sock_options_main}). Prende un valore intero che di
  default è 75.

\item[\sysctlrelfiled{net/ipv4}{tcp\_keepalive\_probes}] indica il
  massimo numero pacchetti di \textit{keepalive} (vedi
  sez.~\ref{sec:sock_options_main}) che devono essere inviati senza ricevere
  risposta prima che il kernel decida che la connessione è caduta e la
  termini. Prende un valore intero che di default è 9.

\item[\sysctlrelfiled{net/ipv4}{tcp\_keepalive\_time}] indica il numero di
  secondi che devono passare senza traffico sulla connessione prima che il
  kernel inizi ad inviare pacchetti di \textit{keepalive}.\footnote{ha effetto
    solo per i socket per cui si è impostata l'opzione \const{SO\_KEEPALIVE}
    (vedi sez.~\ref{sec:sock_options_main}.}  Prende un valore intero che di
  default è 7200, pari a due ore.

\item[\sysctlrelfiled{net/ipv4}{tcp\_low\_latency}] indica allo stack
  TCP del kernel di ottimizzare il comportamento per ottenere tempi di latenza
  più bassi a scapito di valori più alti per l'utilizzo della banda. Prende un
  valore logico che di default è disabilitato in quanto un maggior utilizzo
  della banda è preferito, ma esistono applicazioni particolari in cui la
  riduzione della latenza è più importante (ad esempio per i cluster di
  calcolo parallelo) nelle quali lo si può abilitare.

\item[\sysctlrelfiled{net/ipv4}{tcp\_max\_orphans}] indica il numero
  massimo di socket TCP ``\textsl{orfani}'' (vale a dire non associati a
  nessun file descriptor) consentito nel sistema.\footnote{trattasi in genere
    delle connessioni relative a socket chiusi che non hanno completato il
    processo di chiusura.}  Quando il limite viene ecceduto la connessione
  orfana viene resettata e viene stampato un avvertimento. Questo limite viene
  usato per contrastare alcuni elementari attacchi di \textit{denial of
    service}.  Diminuire il valore non è mai raccomandato, in certe condizioni
  di rete può essere opportuno aumentarlo, ma si deve tenere conto del fatto
  che ciascuna connessione orfana può consumare fino a 64K di memoria del
  kernel. Prende un valore intero, il valore di default viene impostato
  inizialmente al valore del parametro del kernel \texttt{NR\_FILE}, e viene
  aggiustato a seconda della memoria disponibile.

% TODO verificare la spiegazione di connessione orfana.

\item[\sysctlrelfiled{net/ipv4}{tcp\_max\_syn\_backlog}] indica la lunghezza
  della coda delle connessioni incomplete, cioè delle connessioni per le quali
  si è ricevuto un SYN di richiesta ma non l'ACK finale del \textit{three way
    handshake} (si riveda quanto illustrato in
  sez.~\ref{sec:TCP_func_listen}).

  Quando questo valore è superato il kernel scarterà immediatamente ogni
  ulteriore richiesta di connessione. Prende un valore intero; il default, che
  è 256, viene automaticamente portato a 1024 qualora nel sistema ci sia
  sufficiente memoria (se maggiore di 128Mb) e ridotto a 128 qualora la
  memoria sia poca (inferiore a 32Mb).\footnote{si raccomanda, qualora si
    voglia aumentare il valore oltre 1024, di seguire la procedura citata
    nella pagina di manuale di TCP, e modificare il valore della costante
    \texttt{TCP\_SYNQ\_HSIZE} nel file \texttt{include/net/tcp.h} dei sorgenti
    del kernel, in modo che sia $\mathtt{tcp\_max\_syn\_backlog} \ge
    \mathtt{16*TCP\_SYNQ\_HSIZE}$, per poi ricompilare il kernel.}

\item[\sysctlrelfiled{net/ipv4}{tcp\_max\_tw\_buckets}] indica il
  numero massimo di socket in stato \texttt{TIME\_WAIT} consentito nel
  sistema. Prende un valore intero di default è impostato al doppio del valore
  del parametro \texttt{NR\_FILE}, ma che viene aggiustato automaticamente a
  seconda della memoria presente. Se il valore viene superato il socket viene
  chiuso con la stampa di un avviso; l'uso di questa funzionalità consente di
  prevenire alcuni semplici attacchi di \textit{denial of service}.
  

\item[\sysctlrelfiled{net/ipv4}{tcp\_mem}] viene usato dallo stack TCP
  per gestire le modalità con cui esso utilizzerà la memoria. Prende una
  tripletta di valori interi, che indicano un numero di pagine:

  \begin{itemize*}
  \item il primo valore, chiamato \textit{low} nelle pagine di manuale, indica
    il numero di pagine allocate sotto il quale non viene usato nessun
    meccanismo di regolazione dell'uso della memoria.

  \item il secondo valore, chiamato \textit{pressure} indica il numero di
    pagine allocate passato il quale lo stack TCP inizia a moderare il suo
    consumo di memoria; si esce da questo stato di \textsl{pressione} sulla
    memoria quando il numero di pagine scende sotto il precedente valore
    \textit{low}.

  \item il terzo valore, chiamato \textit{high} indica il numero massimo di
    pagine che possono essere utilizzate dallo stack TCP/IP, e soprassiede
    ogni altro valore specificato dagli altri limiti del kernel.
  \end{itemize*}

\item[\sysctlrelfiled{net/ipv4}{tcp\_orphan\_retries}] indica il numero
  massimo di volte che si esegue un tentativo di controllo sull'altro capo di
  una connessione che è stata già chiusa dalla nostra parte. Prende un valore
  intero che di default è 8.

\item[\sysctlrelfiled{net/ipv4}{tcp\_reordering}] indica il numero
  massimo di volte che un pacchetto può essere riordinato nel flusso di dati,
  prima che lo stack TCP assuma che è andato perso e si ponga nello stato di
  \textit{slow start} (si veda sez.~\ref{sec:tcp_protocol_xxx}) viene usata
  questa metrica di riconoscimento dei riordinamenti per evitare inutili
  ritrasmissioni provocate dal riordinamento. Prende un valore intero che di
  default che è 3, e che non è opportuno modificare.

\item[\sysctlrelfiled{net/ipv4}{tcp\_retrans\_collapse}] in caso di
  pacchetti persi durante una connessione, per ottimizzare l'uso della banda
  il kernel cerca di eseguire la ritrasmissione inviando pacchetti della
  massima dimensione possibile; in sostanza dati che in precedenza erano stati
  trasmessi su pacchetti diversi possono essere ritrasmessi riuniti su un solo
  pacchetto (o su un numero minore di pacchetti di dimensione
  maggiore). Prende un valore logico e di default è abilitato.

\item[\sysctlrelfiled{net/ipv4}{tcp\_retries1}] imposta il massimo
  numero di volte che protocollo tenterà la ritrasmissione si un pacchetto su
  una connessione stabilita prima di fare ricorso ad ulteriori sforzi che
  coinvolgano anche il livello di rete. Passato questo numero di
  ritrasmissioni verrà fatto eseguire al livello di rete un tentativo di
  aggiornamento della rotta verso la destinazione prima di eseguire ogni
  successiva ritrasmissione. Prende un valore intero che di default è 3.

\item[\sysctlrelfiled{net/ipv4}{tcp\_retries2}] imposta il numero di
  tentativi di ritrasmissione di un pacchetto inviato su una connessione già
  stabilita per il quale non si sia ricevuto una risposta di ACK (si veda
  anche quanto illustrato in sez.~\ref{sec:TCP_server_crash}). Prende un
  valore intero che di default è 15, il che comporta un tempo variabile fra 13
  e 30 minuti; questo non corrisponde a quanto richiesto
  nell'\href{http://www.ietf.org/rfc/rfc1122.txt}{RFC~1122} dove è indicato un
  massimo di 100 secondi, che però è un valore considerato troppo basso.

\item[\sysctlrelfiled{net/ipv4}{tcp\_rfc1337}] indica al kernel di
  abilitare il comportamento richiesto
  nell'\href{http://www.ietf.org/rfc/rfc1337.txt}{RFC~1337}. Prende un valore
  logico e di default è disabilitato, il che significa che alla ricezione di
  un segmento RST in stato \texttt{TIME\_WAIT} il socket viene chiuso
  immediatamente senza attendere la conclusione del periodo di
  \texttt{TIME\_WAIT}.

\item[\sysctlrelfiled{net/ipv4}{tcp\_rmem}] viene usato dallo stack TCP
  per controllare dinamicamente le dimensioni dei propri buffer di ricezione,
  anche in rapporto alla memoria disponibile.  Prende una tripletta di valori
  interi separati da spazi che indicano delle dimensioni in byte:

  \begin{itemize}
  \item il primo valore, chiamato \textit{min} nelle pagine di manuale, indica
    la dimensione minima in byte del buffer di ricezione; il default è 4Kb, ma
    in sistemi con poca memoria viene automaticamente ridotto a
    \const{PAGE\_SIZE}.  Questo valore viene usato per assicurare che anche in
    situazioni di pressione sulla memoria (vedi quanto detto per
    \sysctlrelfile{net/ipv4}{tcp\_rmem}) le allocazioni al di sotto di
    questo limite abbiamo comunque successo.  Questo valore non viene comunque
    ad incidere sulla dimensione del buffer di ricezione di un singolo socket
    dichiarata con l'opzione \const{SO\_RCVBUF}.

  \item il secondo valore, denominato \textit{default} nelle pagine di
    manuale, indica la dimensione di default, in byte, del buffer di ricezione
    di un socket TCP.  Questo valore sovrascrive il default iniziale impostato
    per tutti i socket con \sysctlfile{net/core/mem\_default} che vale per
    qualunque protocollo. Il default è 87380 byte, ridotto a 43689 per sistemi
    con poca memoria. Se si desiderano dimensioni più ampie per tutti i socket
    si può aumentare questo valore, ma se si vuole che in corrispondenza
    aumentino anche le dimensioni usate per la finestra TCP si deve abilitare
    il \textit{TCP window scaling} (di default è abilitato, vedi più avanti
    \sysctlrelfile{net/ipv4}{tcp\_window\_scaling}).

  \item il terzo valore, denominato \textit{max} nelle pagine di manuale,
    indica la dimensione massima in byte del buffer di ricezione di un socket
    TCP; il default è 174760 byte, che viene ridotto automaticamente a 87380
    per sistemi con poca memoria. Il valore non può comunque eccedere il
    limite generale per tutti i socket posto con
    \sysctlfile{net/core/rmem\_max}. Questo valore non viene ad
    incidere sulla dimensione del buffer di ricezione di un singolo socket
    dichiarata con l'opzione \const{SO\_RCVBUF}.
  \end{itemize}

\item[\sysctlrelfiled{net/ipv4}{tcp\_sack}] indica al kernel di
  utilizzare il meccanismo del \textit{TCP selective acknowledgement} definito
  nell'\href{http://www.ietf.org/rfc/rfc2018.txt}{RFC~2018}. Prende un valore
  logico e di default è abilitato.

\item[\sysctlrelfiled{net/ipv4}{tcp\_stdurg}] indica al kernel di
  utilizzare l'interpretazione che viene data
  dall'\href{http://www.ietf.org/rfc/rfc1122.txt}{RFC~1122} del puntatore dei
  \textit{dati urgenti} (vedi sez.~\ref{sec:TCP_urgent_data}) in cui questo
  punta all'ultimo byte degli stessi; se disabilitato viene usata
  l'interpretazione usata da BSD per cui esso punta al primo byte successivo.
  Prende un valore logico e di default è disabilitato, perché abilitarlo può
  dar luogo a problemi di interoperabilità.

\item[\sysctlrelfiled{net/ipv4}{tcp\_synack\_retries}] indica il numero
  massimo di volte che verrà ritrasmesso il segmento SYN/ACK nella creazione di
  una connessione (vedi sez.~\ref{sec:TCP_conn_cre}). Prende un valore intero
  ed il valore di default è 5; non si deve superare il valore massimo di 255.

\item[\sysctlrelfiled{net/ipv4}{tcp\_syncookies}] abilita i \textit{TCP
    syncookies}.\footnote{per poter usare questa funzionalità è necessario
    avere abilitato l'opzione \texttt{CONFIG\_SYN\_COOKIES} nella compilazione
    del kernel.} Prende un valore logico, e di default è disabilitato. Questa
  funzionalità serve a fornire una protezione in caso di un attacco di tipo
  \textit{SYN flood}, e deve essere utilizzato come ultima risorsa dato che
  costituisce una violazione del protocollo TCP e confligge con altre
  funzionalità come le estensioni e può causare problemi per i client ed il
  reinoltro dei pacchetti.

\item[\sysctlrelfiled{net/ipv4}{tcp\_syn\_retries}] imposta il numero di
  tentativi di ritrasmissione dei pacchetti SYN di inizio connessione del
  \textit{three way handshake} (si ricordi quanto illustrato in
  sez.~\ref{sec:TCP_func_connect}). Prende un valore intero che di default è
  5; non si deve superare il valore massimo di 255.

\item[\sysctlrelfiled{net/ipv4}{tcp\_timestamps}] abilita l'uso dei
  \textit{TCP timestamps}, come definiti
  nell'\href{http://www.ietf.org/rfc/rfc1323.txt}{RFC~1323}. Prende un valore
  logico e di default è abilitato.

\item[\sysctlrelfiled{net/ipv4}{tcp\_tw\_recycle}] abilita il riutilizzo
  rapido dei socket in stato \texttt{TIME\_WAIT}. Prende un valore logico e di
  default è disabilitato. Non è opportuno abilitare questa opzione che può
  causare problemi con il NAT.\footnote{la
    \itindex{Network~Address~Translation~(NAT)} \textit{Network Address
      Translation} (abbreviato in NAT) è una tecnica, impiegata nei firewall e
    nei router, che consente di modificare al volo gli indirizzi dei pacchetti
    che transitano per una macchina, Linux la supporta con il
    \textit{netfilter}.}

\item[\sysctlrelfiled{net/ipv4}{tcp\_tw\_reuse}] abilita il riutilizzo
  dello stato \texttt{TIME\_WAIT} quando questo è sicuro dal punto di vista
  del protocollo. Prende un valore logico e di default è disabilitato. 

\item[\sysctlrelfiled{net/ipv4}{tcp\_window\_scaling}] un valore
  logico, attivo di default, che abilita la funzionalità del
  \itindex{TCP~window~scaling} \textit{TCP window scaling} definita
  dall'\href{http://www.ietf.org/rfc/rfc1323.txt}{RFC~1323}. Prende un valore
  logico e di default è abilitato. Come accennato in
  sez.~\ref{sec:TCP_TCP_opt} i 16 bit della finestra TCP comportano un limite
  massimo di dimensione di 64Kb, ma esiste una opportuna opzione del
  protocollo che permette di applicare un fattore di scale che consente di
  aumentarne le dimensioni. Questa è pienamente supportata dallo stack TCP di
  Linux, ma se lo si disabilita la negoziazione del
  \itindex{TCP~window~scaling} \textit{TCP window scaling} con l'altro capo
  della connessione non viene effettuata.

%\item[\sysctlrelfiled{net/ipv4}{tcp\_vegas\_cong\_avoid}] 
% TODO: controllare su internet

%\item[\sysctlrelfiled{net/ipv4}{tcp\_westwood}] 
% TODO: controllare su internet

\item[\sysctlrelfiled{net/ipv4}{tcp\_wmem}] viene usato dallo stack TCP
  per controllare dinamicamente le dimensioni dei propri buffer di spedizione,
  adeguandole in rapporto alla memoria disponibile.  Prende una tripletta di
  valori interi separati da spazi che indicano delle dimensioni in byte:

  \begin{itemize}
  \item il primo valore, chiamato \textit{min}, indica la dimensione minima in
    byte del buffer di spedizione; il default è 4Kb. Come per l'analogo di
    \sysctlrelfile{net/ipv4}{tcp\_rmem}) viene usato per assicurare
    che anche in situazioni di pressione sulla memoria (vedi
    \sysctlrelfile{net/ipv4}{tcp\_mem}) le allocazioni al di sotto di
    questo limite abbiamo comunque successo.  Di nuovo questo valore non viene
    ad incidere sulla dimensione del buffer di trasmissione di un singolo
    socket dichiarata con l'opzione \const{SO\_SNDBUF}.

  \item il secondo valore, denominato \textit{default}, indica la dimensione
    di default in byte del buffer di spedizione di un socket TCP.  Questo
    valore sovrascrive il default iniziale impostato per tutti i tipi di
    socket sul file \sysctlfile{net/core/wmem\_default}. Il default è 87380
    byte, ridotto a 43689 per sistemi con poca memoria. Si può aumentare
    questo valore quando si desiderano dimensioni più ampie del buffer di
    trasmissione per i socket TCP, ma come per il precedente
    \sysctlrelfile{net/ipv4}{tcp\_rmem}) se si vuole che in corrispondenza
    aumentino anche le dimensioni usate per la finestra TCP si deve abilitare
    il \textit{TCP window scaling} con
    \sysctlrelfile{net/ipv4}{tcp\_window\_scaling}.

  \item il terzo valore, denominato \textit{max}, indica la dimensione massima
    in byte del buffer di spedizione di un socket TCP; il default è 128Kb, che
    viene ridotto automaticamente a 64Kb per sistemi con poca memoria. Il
    valore non può comunque eccedere il limite generale per tutti i socket
    posto con \sysctlfile{net/core/wmem\_max}. Questo valore non viene
    ad incidere sulla dimensione del buffer di trasmissione di un singolo
    socket dichiarata con l'opzione \const{SO\_SNDBUF}.
  \end{itemize}

\end{basedescript}



% LocalWords:  socket sez dotted decimal resolver Domain Name Service cap DNS
% LocalWords:  client fig LDAP Lightweight Access Protocol NIS Information Sun
% LocalWords:  like netgroup Switch Solaris glibc libc uclib NSS tab shadow uid
% LocalWords:  username group aliases ethers MAC address hosts networks rpc RPC
% LocalWords:  protocols services dns db lib libnss org truelite it root res HS
% LocalWords:  resource init netinet resolv int void conf host LOCALDOMAIN TCP
% LocalWords:  options DEBUG debug AAONLY USEVC UDP PRIMARY IGNTC RECURSE INET
% LocalWords:  DEFNAMES search STAYOPEN DNSRCH INSECURE NOALIASES HOSTALIASES
% LocalWords:  IPv gethostbyname NOCHECKNAME KEEPTSIG TSIG BLAST RETRY retry NS
% LocalWords:  retrans query FQDN Fully Qualified const char dname class type
% LocalWords:  unsigned answer anslen CSNET Hesiod MIT CHAOS Chaosnet ANY BIND
% LocalWords:  nameser compat Berkley MF CNAME SOA MB MR NULL WKS PTR HINFO TXT
% LocalWords:  MINFO RP responsible person AFSDB AFS RT router NSAP SIG KEY PX
% LocalWords:  GPOS AAAA LOC NXT EID NIMLOC nimrod SRV ATMA ATM NAPTR naming AF
% LocalWords:  authority IXFR AXFR MAILB MAILA errno NOT FOUND RECOVERY TRY err
% LocalWords:  AGAIN herror netdb string perror error hstrerror strerror struct
% LocalWords:  hostent name addrtype length addr list sys af mygethost inet ret
% LocalWords:  ntop deep copy buf size buflen result errnop value argument len
% LocalWords:  ERANGE sethostent stayopen endhostent gethostbyaddr order pton
% LocalWords:  getipnodebyname getipnodebyaddr flags num MAPPED ALL ADDRCONFIG
% LocalWords:  freehostent ip getXXXbyname getXXXbyaddr servent getservbyname
% LocalWords:  netent getnetbyname getnetbyaddr protoent smtp udp
% LocalWords:  getprotobyname getprotobyaddr getservbyport port tcp setservent
% LocalWords:  getservent endservent setXXXent getXXXent endXXXent gethostent
% LocalWords:  setnetent getnetent endnetent setprotoent getprotoent POSIX RFC
% LocalWords:  endprotoent getaddrinfo getnameinfo nell' node service addrinfo
% LocalWords:  hints linked addrlen socklen family socktype protocol sockaddr
% LocalWords:  canonname next PF UNSPEC SOCK STREAM DGRAM bind INADDR loopback
% LocalWords:  connect sendto NUMERICHOST EAI NONAME SYSTEM BADFLAGS ADDRFAMILY
% LocalWords:  NODATA MEMORY FAIL errcode echo mygetaddr ptr casting Canonical
% LocalWords:  freeaddrinfo getservname salen hostlen serv servlen l'OR NI NUL
% LocalWords:  NOFQDN NAMEREQD NUMERICSERV MAXHOST MAXSERV sockconn SockUtil of
% LocalWords:  descriptor hint fifth sockbind setsockopt getsockopt sock level
% LocalWords:  optname optval optlen EBADF EFAULT EINVAL ENOPROTOOPT ENOTSOCK
% LocalWords:  IPPROTO Stevens ICMP ICMPV ICMPv get KEEPALIVE OOBINLINE timeval
% LocalWords:  RCVLOWAT SNDLOWAT RCVTIMEO SNDTIMEO BSDCOMPAT BSD PASSCRED ucred
% LocalWords:  PEERCRED BINDTODEVICE REUSEADDR ACCEPTCONN DONTROUTE gateway MSG
% LocalWords:  BROADCAST broadcast SNDBUF RCVBUF LINGER linger PRIORITY read IF
% LocalWords:  OOB recvmsg kernel select write readv recv recvfrom EAGAIN send
% LocalWords:  EWOULDBLOCK writev sendmsg raw domain SCM CREDENTIALS eth packet
% LocalWords:  IFNAMSIZ capabilities capability ADMIN log trpt EADDRINUSE close
% LocalWords:  listen routing sysctl shutdown Quality TOS keep alive ACK RST to
% LocalWords:  ECONNRESET ETIMEDOUT keepalive echod fourth newsgroup WAIT reuse
% LocalWords:  sockbindopt SockUtils homed completely binding RECVDSTADDR onoff
% LocalWords:  PKTINFO getsockname multicast streaming unicast REUSEPORT reset
% LocalWords:  stealing ling RECVTOS RECVTTL TTL RECVOPTS RETOPTS HDRINCL MTU
% LocalWords:  RECVERR DISCOVER Path Discovery ALERT alert ADD MEMBERSHIP mreqn
% LocalWords:  pktinfo ipi ifindex spec dst RECVIF Live IPTOS LOWDELAY Advanced
% LocalWords:  Transfer Unit PMTUDISC DONT WANT route dall' pmtu EMSGSIZE imr
% LocalWords:  multiaddr mreq fcntl ioctl request SIOCGSTAMP trip SIOCSPGRP pid
% LocalWords:  process SIGIO SIGURG KILL FIOASYNC SIOCGPGRP filesystem proc ttl
% LocalWords:  rmem wmem message cost burst bucket filter netdev backlog optmem
% LocalWords:  forward dynaddr dial autoconfig local masquerading ipfrag high
% LocalWords:  thresh low always defrag CONFIG SETSIG cmd FIOGETOWN FIOSETOWN
% LocalWords:  quest'ultime neigh dev weight cong mod somaxconn Di SIOCINQ DoS
% LocalWords:  Documentation SIOCATMARK SIOCOUTQ FIONREAD TIOCOUTQ Denial work
% LocalWords:  netfilter scheduler mark ARP DHCP BOOTP RARP nonlocal sniffer is
% LocalWords:  linux NODELAY MAXSEG CORK KEEPIDLE KEEPINTVL KEEPCNT SYNCNT INFO
% LocalWords:  DEFER ACCEPT WINDOW CLAMP QUICKACK CONGESTION ENCAP urgent MSS
% LocalWords:  Segment SYN accept advertised window info quickack Nagle ifreq
% LocalWords:  ifr ppp union EPERM SIOCGIFNAME dell' interface index IFF NOARP
% LocalWords:  SIOCGIFINDEX SIOCGIFFLAGS POINTOPOINT RUNNING PROMISC NOTRAILERS
% LocalWords:  ALLMULTI bundle PORTSEL ifmap AUTOMEDIA DYNAMIC SIOCSIFFLAGS way
% LocalWords:  SIOCGIFMETRIC SIOCSIFMETRIC SIOCGIFMTU SIOCSIFMTU SIOCGIFHWADDR
% LocalWords:  SIOCSIFHWADDR SIOCSIFHWBROADCAST SIOCGIFMAP SIOCSIFMAP sendfile
% LocalWords:  SIOCADDMULTI SIOCDELMULTI SIOCGIFTXQLEN SIOCSIFTXQLEN three syn
% LocalWords:  SIOCSIFNAME SIOCGIFCONF handshake retries MIN FreeBSD closing Mb
% LocalWords:  abort overflow adv win app bic convergence dsack ecn fack frto
% LocalWords:  intvl probes latency orphans l'ACK SYNQ HSIZE tw buckets mem rfc
% LocalWords:  orphan reordering collapse sack stdurg synack syncookies recycle
% LocalWords:  timestamps scaling vegas avoid westwood tcpi l'incapsulazione NR
% LocalWords:  metric EOPNOTSUPP mtu hwaddr ARPHRD interrupt DMA map qlen silly
% LocalWords:  rename ifconf syndrome dell'ACK FTP ACCEPTFILTER advanced reno
% LocalWords:  congestion control Networking cubic CUBIC highspeed HSTCP htcp
% LocalWords:  HTCP hybla HYBLA scalable SCALABLE ifc req iflist access ntoa Kb
% LocalWords:  hop Selective acknowledgement Explicit RTO stack firewall passwd
% LocalWords:  Notification wireless denial pressure ATTACH DETACH publickey
% LocalWords:  libpcap discovery point l'overhaed min PAGE flood NFS blast
% LocalWords:  selective COOKIES NAT Translation

%%% Local Variables: 
%%% mode: latex
%%% TeX-master: "gapil"
%%% End: 


%% othersock.tex
%%
%% Copyright (C) 2004-2018 Simone Piccardi.  Permission is granted to
%% copy, distribute and/or modify this document under the terms of the GNU Free
%% Documentation License, Version 1.1 or any later version published by the
%% Free Software Foundation; with the Invariant Sections being "Un preambolo",
%% with no Front-Cover Texts, and with no Back-Cover Texts.  A copy of the
%% license is included in the section entitled "GNU Free Documentation
%% License".
%%

\chapter{Gli altri tipi di socket}
\label{cha:other_socket}

Dopo aver trattato in cap.~\ref{cha:TCP_socket} i socket TCP, che
costituiscono l'esempio più comune dell'interfaccia dei socket, esamineremo in
questo capitolo gli altri tipi di socket, a partire dai socket UDP, e i socket
\textit{Unix domain} già incontrati in sez.~\ref{sec:ipc_socketpair}.


\section{I socket UDP}
\label{sec:UDP_socket}

Dopo i socket TCP i socket più utilizzati nella programmazione di rete sono i
socket UDP: protocolli diffusi come NFS o il DNS usano principalmente questo
tipo di socket. Tratteremo in questa sezione le loro caratteristiche
principali e le modalità per il loro utilizzo.


\subsection{Le caratteristiche di un socket UDP}
\label{sec:UDP_characteristics}

Come illustrato in sez.\ref{sec:net_udp} UDP è un protocollo molto semplice che
non supporta le connessioni e non è affidabile: esso si appoggia direttamente
sopra IP (per i dettagli sul protocollo si veda sez.~\ref{sec:udp_protocol}).
I dati vengono inviati in forma di pacchetti, e non ne è assicurata né la
effettiva ricezione né l'arrivo nell'ordine in cui vengono inviati. Il
vantaggio del protocollo è la velocità, non è necessario trasmettere le
informazioni di controllo ed il risultato è una trasmissione di dati più
veloce ed immediata.

Questo significa che a differenza dei socket TCP i socket UDP non supportano
una comunicazione di tipo \textit{stream} in cui si ha a disposizione un
flusso continuo di dati che può essere letto un po' alla volta, ma piuttosto
una comunicazione di tipo \textit{datagram}, in cui i dati arrivano in singoli
blocchi che devono essere letti integralmente.

Questo diverso comportamento significa anche che i socket UDP, pur
appartenendo alla famiglia \const{PF\_INET}\footnote{o \const{PF\_INET6}
  qualora si usasse invece il protocollo IPv6, che pure supporta UDP.} devono
essere aperti quando si usa la funzione \func{socket} (si riveda quanto
illustrato a suo tempo in tab.~\ref{tab:sock_sock_valid_combinations})
utilizzando per il tipo di socket il valore \const{SOCK\_DGRAM}.

Questa differenza comporta ovviamente che anche le modalità con cui si usano i
socket UDP sono completamente diverse rispetto ai socket TCP, ed in
particolare non esistendo il concetto di connessione non esiste il meccanismo
del \textit{three way handshake} né quello degli stati del protocollo. In
realtà tutto quello che avviene nella comunicazione attraverso dei socket UDP
è la trasmissione di un pacchetto da un client ad un server o viceversa,
secondo lo schema illustrato in fig.~\ref{fig:UDP_packet-exchange}.

\begin{figure}[!htb]
  \centering \includegraphics[width=10cm]{img/udp_connection}  
  \caption{Lo schema di interscambio dei pacchetti per una comunicazione via
     UDP.}
  \label{fig:UDP_packet-exchange}
\end{figure}

Come illustrato in fig.~\ref{fig:UDP_packet-exchange} la struttura generica di
un server UDP prevede, una volta creato il socket, la chiamata a \func{bind}
per mettersi in ascolto dei dati, questa è l'unica parte comune con un server
TCP. Non essendovi il concetto di connessione le funzioni \func{listen} ed
\func{accept} non sono mai utilizzate nel caso di server UDP. La ricezione dei
dati dal client avviene attraverso la funzione \func{recvfrom}, mentre una
eventuale risposta sarà inviata con la funzione \func{sendto}.

Da parte del client invece, una volta creato il socket non sarà necessario
connettersi con \func{connect} (anche se, come vedremo in
sez.~\ref{sec:UDP_connect}, è possibile usare questa funzione, con un
significato comunque diverso) ma si potrà effettuare direttamente una
richiesta inviando un pacchetto con la funzione \func{sendto} e si potrà
leggere una eventuale risposta con la funzione \func{recvfrom}.

Anche se UDP è completamente diverso rispetto a TCP resta identica la
possibilità di gestire più canali di comunicazione fra due macchine
utilizzando le porte. In questo caso il server dovrà usare comunque la
funzione \func{bind} per scegliere la porta su cui ricevere i dati, e come nel
caso dei socket TCP si potrà usare il comando \cmd{netstat} per
verificare quali socket sono in ascolto:
\begin{verbatim}
[piccardi@gont gapil]# netstat -anu
Active Internet connections (servers and established)
Proto Recv-Q Send-Q Local Address           Foreign Address         State
udp        0      0 0.0.0.0:32768           0.0.0.0:*
udp        0      0 192.168.1.2:53          0.0.0.0:*
udp        0      0 127.0.0.1:53            0.0.0.0:*
udp        0      0 0.0.0.0:67              0.0.0.0:*
\end{verbatim}
in questo caso abbiamo attivi il DNS (sulla porta 53, e sulla 32768 per la
connessione di controllo del server \cmd{named}) ed un server DHCP (sulla
porta 67).

Si noti però come in questo caso la colonna che indica lo stato sia vuota. I
socket UDP infatti non hanno uno stato. Inoltre anche in presenza di traffico
non si avranno indicazioni delle connessioni attive, proprio perché questo
concetto non esiste per i socket UDP, il kernel si limita infatti a ricevere i
pacchetti ed inviarli al processo in ascolto sulla porta cui essi sono
destinati, oppure a scartarli inviando un messaggio \textit{ICMP port
  unreachable} qualora non vi sia nessun processo in ascolto.


\subsection{Le funzioni \func{sendto} e \func{recvfrom}}
\label{sec:UDP_sendto_recvfrom}

Come accennato in sez.~\ref{sec:UDP_characteristics} le due funzioni
principali usate per la trasmissione di dati attraverso i socket UDP (ed in
generale per i socket di tipo \textit{datagram}) sono \func{sendto} e
\func{recvfrom}. La necessità di usare queste funzioni è dovuta al fatto che
non esistendo con UDP il concetto di connessione, non si può stabilire (come
avviene con i socket TCP grazie alla chiamata ad \func{accept} che li associa
ad una connessione) quali sono la sorgente e la destinazione dei dati che
passano sul socket.\footnote{anche se in alcuni casi, come quello di un client
  che contatta un server, è possibile connettere il socket (vedi
  sez.~\ref{sec:UDP_connect}), la necessità di \func{sendto} e \func{recvfrom}
  resta, dato che questo è possibile solo sul lato del client.}

Per questo anche se in generale si possono comunque leggere i dati con
\func{read}, usando questa funzione non si sarà in grado di determinare da
quale fra i possibili corrispondenti (se ve ne sono più di uno, come avviene
sul lato del server) questi arrivino. E non sarà comunque possibile usare
\func{write} (che fallisce un errore di \errval{EDESTADDRREQ}) in quanto non
è determinato la destinazione che i dati avrebbero.

Per questo motivo nel caso di UDP diventa essenziale utilizzare queste due
funzioni, che sono comunque utilizzabili in generale per la trasmissione di
dati attraverso qualunque tipo di socket. Esse hanno la caratteristica di
prevedere tre argomenti aggiuntivi attraverso i quali è possibile specificare
la destinazione dei dati trasmessi o ottenere l'origine dei dati ricevuti. La
prima di queste funzioni è \funcd{sendto} ed il suo prototipo\footnote{il
  prototipo illustrato è quello utilizzato dalla \acr{glibc}, che segue le
  \textit{Single Unix Specification}, l'argomento \param{flags} era di tipo
  \ctyp{int} nei vari BSD4.*, mentre nelle \acr{libc4} e \acr{libc5} veniva
  usato un \texttt{unsigned int}; l'argomento \param{len} era \ctyp{int} nei
  vari BSD4.* e nella \acr{libc4}, ma \type{size\_t} nella \acr{libc5}; infine
  l'argomento \param{tolen} era \ctyp{int} nei vari BSD4.*, nella \acr{libc4} e
  nella \acr{libc5}.} è:
\begin{functions}
  \headdecl{sys/types.h}
  \headdecl{sys/socket.h}
  
  \funcdecl{ssize\_t sendto(int sockfd, const void *buf, size\_t len, int
    flags, const struct sockaddr *to, socklen\_t tolen)}
  
  Trasmette un messaggio ad un altro socket.
  
  \bodydesc{La funzione restituisce il numero di caratteri inviati in caso di
    successo e -1 per un errore; nel qual caso \var{errno} viene impostata al
    rispettivo codice di errore:
  \begin{errlist}
  \item[\errcode{EAGAIN}] il socket è in modalità non bloccante, ma
    l'operazione richiede che la funzione si blocchi.
  \item[\errcode{ECONNRESET}] l'altro capo della comunicazione ha resettato la
    connessione.
  \item[\errcode{EDESTADDRREQ}] il socket non è di tipo connesso, e non si è
    specificato un indirizzo di destinazione.
  \item[\errcode{EISCONN}] il socket è già connesso, ma si è specificato un
    destinatario.
  \item[\errcode{EMSGSIZE}] il tipo di socket richiede l'invio dei dati in un
    blocco unico, ma la dimensione del messaggio lo rende impossibile.
  \item[\errcode{ENOBUFS}] la coda di uscita dell'interfaccia è già piena (di
    norma Linux non usa questo messaggio ma scarta silenziosamente i
    pacchetti).
  \item[\errcode{ENOTCONN}] il socket non è connesso e non si è specificata
    una destinazione.
  \item[\errcode{EOPNOTSUPP}] il valore di \param{flag} non è appropriato per
    il tipo di socket usato.
  \item[\errcode{EPIPE}] il capo locale della connessione è stato chiuso, si
    riceverà anche un segnale di \signal{SIGPIPE}, a meno di non aver impostato
    \const{MSG\_NOSIGNAL} in \param{flags}.
  \end{errlist}
  ed anche \errval{EFAULT}, \errval{EBADF}, \errval{EINVAL}, \errval{EINTR},
  \errval{ENOMEM}, \errval{ENOTSOCK} più gli eventuali altri errori relativi
  ai protocolli utilizzati.}
\end{functions}

I primi tre argomenti sono identici a quelli della funzione \func{write} e
specificano il socket \param{sockfd} a cui si fa riferimento, il buffer
\param{buf} che contiene i dati da inviare e la relativa lunghezza
\param{len}.  Come per \func{write} la funzione ritorna il numero di byte
inviati; nel caso di UDP però questo deve sempre corrispondere alla dimensione
totale specificata da \param{len} in quanto i dati vengono sempre inviati in
forma di pacchetto e non possono essere spezzati in invii successivi.  Qualora
non ci sia spazio nel buffer di uscita la funzione si blocca (a meno di non
avere aperto il socket in modalità non bloccante), se invece non è possibile
inviare il messaggio all'interno di un unico pacchetto (ad esempio perché
eccede le dimensioni massime del protocollo sottostante utilizzato) essa
fallisce con l'errore di \errcode{EMSGSIZE}. 

I due argomenti \param{to} e \param{tolen} servono a specificare la
destinazione del messaggio da inviare, e indicano rispettivamente la struttura
contenente l'indirizzo di quest'ultima e la sua dimensione; questi argomenti
vanno specificati stessa forma in cui li si sarebbero usati con
\func{connect}. Nel nostro caso \param{to} dovrà puntare alla struttura
contenente l'indirizzo IP e la porta di destinazione verso cui si vogliono
inviare i dati (questo è indifferente rispetto all'uso di TCP o UDP, usando
socket diversi si sarebbero dovute utilizzare le rispettive strutture degli
indirizzi).

Se il socket è di un tipo che prevede le connessioni (ad esempio un socket
TCP), questo deve essere già connesso prima di poter eseguire la funzione, in
caso contrario si riceverà un errore di \errcode{ENOTCONN}. In questo
specifico caso in cui gli argomenti \param{to} e \param{tolen} non servono
essi dovranno essere inizializzati rispettivamente a \val{NULL} e 0;
normalmente quando si opera su un socket connesso essi vengono ignorati, ma
qualora si sia specificato un indirizzo è possibile ricevere un errore di
\errcode{EISCONN}.

Finora abbiamo tralasciato l'argomento \param{flags}; questo è un intero usato
come maschera binaria che permette di impostare una serie di modalità di
funzionamento della comunicazione attraverso il socket (come
\constd{MSG\_NOSIGNAL} che impedisce l'invio del segnale \signal{SIGPIPE}
quando si è già chiuso il capo locale della connessione). Torneremo con
maggiori dettagli sul significato di questo argomento in
sez.~\ref{sec:net_sendmsg}, dove tratteremo le funzioni avanzate dei socket,
per il momento ci si può limitare ad usare sempre un valore nullo.

La seconda funzione utilizzata nella comunicazione fra socket UDP è
\funcd{recvfrom}, che serve a ricevere i dati inviati da un altro socket; il
suo prototipo\footnote{il prototipo è quello della \acr{glibc} che segue le
  \textit{Single Unix Specification}, i vari BSD4.*, le \acr{libc4} e la
  \acr{libc5} usano un \ctyp{int} come valore di ritorno; per gli argomenti
  \param{flags} e \param{len} vale quanto detto a proposito di \func{sendto};
  infine l'argomento \param{fromlen} è \ctyp{int} per i vari BSD4.*, la
  \acr{libc4} e la \acr{libc5}.} è:
\begin{functions}
  \headdecl{sys/types.h}
  \headdecl{sys/socket.h}
  
  \funcdecl{ssize\_t recvfrom(int sockfd, const void *buf, size\_t len, int
    flags, const struct sockaddr *from, socklen\_t *fromlen)}
  
  Riceve un messaggio da un socket.
  
  \bodydesc{La funzione restituisce il numero di byte ricevuti in caso di
    successo e -1 in caso di errore; nel qual caso \var{errno} assumerà il
    valore:
  \begin{errlist}
  \item[\errcode{EAGAIN}] il socket è in modalità non bloccante, ma
    l'operazione richiede che la funzione si blocchi, oppure si è impostato un
    timeout in ricezione e questo è scaduto.
  \item[\errcode{ECONNREFUSED}] l'altro capo della comunicazione ha rifiutato
    la connessione (in genere perché il relativo servizio non è disponibile).
  \item[\errcode{ENOTCONN}] il socket è di tipo connesso, ma non si è eseguita
    la connessione.
  \end{errlist}
  ed anche \errval{EFAULT}, \errval{EBADF}, \errval{EINVAL}, \errval{EINTR},
  \errval{ENOMEM}, \errval{ENOTSOCK} più gli eventuali altri errori relativi
  ai protocolli utilizzati.}
\end{functions}

Come per \func{sendto} i primi tre argomenti sono identici agli analoghi di
\func{read}: dal socket vengono letti \param{len} byte che vengono salvati nel
buffer \param{buf}. A seconda del tipo di socket (se di tipo \textit{datagram}
o di tipo \textit{stream}) i byte in eccesso che non sono stati letti possono
rispettivamente andare persi o restare disponibili per una lettura successiva.
Se non sono disponibili dati la funzione si blocca, a meno di non aver aperto
il socket in modalità non bloccante, nel qual caso si avrà il solito errore di
\errcode{EAGAIN}.  Qualora \param{len} ecceda la dimensione del pacchetto la
funzione legge comunque i dati disponibili, ed il suo valore di ritorno è
comunque il numero di byte letti.

I due argomenti \param{from} e \param{fromlen} sono utilizzati per ottenere
l'indirizzo del mittente del pacchetto che è stato ricevuto, e devono essere
opportunamente inizializzati; il primo deve contenere il puntatore alla
struttura (di tipo \ctyp{sockaddr}) che conterrà l'indirizzo e il secondo il
puntatore alla variabile con la dimensione di detta struttura. Si tenga
presente che mentre il contenuto della struttura \ctyp{sockaddr} cui punta
\param{from} può essere qualunque, la variabile puntata da \param{fromlen}
deve essere opportunamente inizializzata a \code{sizeof(sockaddr)},
assicurandosi che la dimensione sia sufficiente a contenere tutti i dati
dell'indirizzo.\footnote{si ricordi che \ctyp{sockaddr} è un tipo generico che
  serve ad indicare la struttura corrispondente allo specifico tipo di
  indirizzo richiesto, il valore di \param{fromlen} pone un limite alla
  quantità di dati che verranno scritti sulla struttura puntata da
  \param{from} e se è insufficiente l'indirizzo risulterà corrotto.}  Al
ritorno della funzione si otterranno i dati dell'indirizzo e la sua effettiva
lunghezza, (si noti che \param{fromlen} è un valore intero ottenuto come
\textit{value result argument}).  Se non si è interessati a questa
informazione, entrambi gli argomenti devono essere inizializzati al valore
\val{NULL}.

Una differenza fondamentale del comportamento di queste funzioni rispetto alle
usuali \func{read} e \func{write} che abbiamo usato con i socket TCP è che in
questo caso è perfettamente legale inviare con \func{sendto} un pacchetto
vuoto (che nel caso conterrà solo le intestazioni di IP e di UDP),
specificando un valore nullo per \param{len}. Allo stesso modo è possibile
ricevere con \func{recvfrom} un valore di ritorno di 0 byte, senza che questo
possa configurarsi come una chiusura della connessione\footnote{dato che la
  connessione non esiste, non ha senso parlare di chiusura della connessione,
  questo significa anche che con i socket UDP non è necessario usare
  \func{close} o \func{shutdown} per terminare la comunicazione.} o come una
cessazione delle comunicazioni.



\subsection{Un client UDP elementare}
\label{sec:UDP_daytime_client}

Vediamo allora come implementare un primo client elementare con dei socket
UDP.  Ricalcando quanto fatto nel caso dei socket TCP prenderemo come primo
esempio l'uso del servizio \textit{daytime}, utilizzando questa volta UDP. Il
servizio è definito nell'\href{http://www.ietf.org/rfc/rfc0862.txt}{RFC~867},
che nel caso di uso di UDP prescrive che il client debba inviare un pacchetto
UDP al server (di contenuto non specificato), il quale risponderà a inviando a
sua volta un pacchetto UDP contenente la data.

\begin{figure}[!htbp] 
  \footnotesize \centering
  \begin{minipage}[c]{\codesamplewidth}
    \includecodesample{listati/UDP_daytime.c}
  \end{minipage} 
  \normalsize
  \caption{Sezione principale del client per il servizio \textit{daytime} su
    UDP.}
  \label{fig:UDP_daytime_client}
\end{figure}

In fig.~\ref{fig:UDP_daytime_client} è riportato la sezione principale del
codice del nostro client, il sorgente completo si trova nel file
\file{UDP\_daytime.c} distribuito con gli esempi allegati alla guida; al
solito si è tralasciato di riportare in figura la sezione relativa alla
gestione delle opzioni a riga di comando (nel caso praticamente assenti).

Il programma inizia (\texttt{\small 9--12}) con la creazione del socket, al
solito uscendo dopo aver stampato un messaggio in caso errore. Si noti come in
questo caso, rispetto all'analogo client basato su socket TCP di
fig.~\ref{fig:TCP_daytime_client_code} si sia usato per il tipo di socket il
valore \const{SOCK\_DGRAM}, pur mantenendosi nella stessa famiglia data da
\const{AF\_INET}.

Il passo successivo (\texttt{\small 13--21}) è l'inizializzazione della
struttura degli indirizzi; prima (\texttt{\small 14}) si cancella
completamente la stessa con \func{memset}, (\texttt{\small 15}) poi si imposta
la famiglia dell'indirizzo ed infine (\texttt{\small 16} la porta. Infine
(\texttt{\small 18--21}) si ricava l'indirizzo del server da contattare
dall'argomento passato a riga di comando, convertendolo con \func{inet\_pton}.
Si noti come questa sezione sia identica a quella del client TCP di
fig.~\ref{fig:TCP_daytime_client_code}, in quanto la determinazione dell'uso
di UDP al posto di TCP è stata effettuata quando si è creato il socket.

Una volta completate le inizializzazioni inizia il corpo principale del
programma, il primo passo è inviare, come richiesto dal protocollo, un
pacchetto al server. Questo lo si fa (\texttt{\small 16}) inviando un
pacchetto vuoto (si ricordi quanto detto in
sez.~\ref{sec:UDP_sendto_recvfrom}) con \func{sendto}, avendo cura di passare
un valore nullo per il puntatore al buffer e la lunghezza del messaggio. In
realtà il protocollo non richiede che il pacchetto sia vuoto, ma dato che il
server comunque ne ignorerà il contenuto, è inutile inviare dei dati.

Verificato (\texttt{\small 24--27}) che non ci siano stati errori nell'invio
si provvede (\texttt{\small 28}) ad invocare \func{recvfrom} per ricevere la
risposta del server. Si controlla poi (\texttt{\small 29--32}) che non vi
siano stati errori in ricezione (uscendo con un messaggio in caso contrario);
se è tutto a posto la variabile \var{nread} conterrà la dimensione del
messaggio di risposta inviato dal server che è stato memorizzato su
\var{buffer}, se (\texttt{\small 34}) pertanto il valore è positivo si
provvederà (\texttt{\small 35}) a terminare la stringa contenuta nel buffer di
lettura\footnote{si ricordi che, come illustrato in
  sez.~\ref{sec:TCP_daytime_client}, il server invia in risposta una stringa
  contenente la data, terminata dai due caratteri CR e LF, che pertanto prima
  di essere stampata deve essere opportunamente terminata con un NUL.} e a
stamparla (\texttt{\small 36}) sullo standard output, controllando anche in
questo caso (\texttt{\small 36--38}) l'esito dell'operazione, ed uscendo con
un messaggio in caso di errore.

Se pertanto si è avuto cura di attivare il server del servizio
\textit{daytime}\footnote{di norma questo è un servizio standard fornito dal
  \textsl{superdemone} \cmd{inetd}, per cui basta abilitarlo nel file di
  configurazione di quest'ultimo, avendo cura di predisporre il servizio su
  UDP.} potremo verificare il funzionamento del nostro client interrogando
quest'ultimo con:
\begin{verbatim}
[piccardi@gont sources]$ ./daytime 127.0.0.1
Sat Mar 20 23:17:13 2004
\end{verbatim}%$
ed osservando il traffico con uno sniffer potremo effettivamente vedere lo
scambio dei due pacchetti, quello vuoto di richiesta, e la risposta del
server:
\begin{verbatim}
[root@gont gapil]# tcpdump -i lo
tcpdump: listening on lo
23:41:21.645579 localhost.32780 > localhost.daytime: udp 0 (DF)
23:41:21.645710 localhost.daytime > localhost.32780: udp 26 (DF)
\end{verbatim}

Una differenza fondamentale del nostro client è che in questo caso, non
disponendo di una connessione, è per lui impossibile riconoscere errori di
invio relativi alla rete. La funzione \func{sendto} infatti riporta solo
errori locali, i dati vengono comunque scritti e la funzione ritorna senza
errori anche se il server non è raggiungibile o non esiste un server in
ascolto sull'indirizzo di destinazione. Questo comporta ad esempio che se si
usa il nostro programma interrogando un server inesistente questo resterà
perennemente bloccato nella chiamata a \func{recvfrom}, fin quando non lo
interromperemo. Vedremo in sez.~\ref{sec:UDP_connect} come si può porre rimedio
a questa problematica.


\subsection{Un server UDP elementare}
\label{sec:UDP_daytime_server}

Nella sezione precedente abbiamo visto come scrivere un client elementare per
servizio \textit{daytime}, vediamo in questa come deve essere scritto un
server.  Si ricordi che il compito di quest'ultimo è quello di ricevere un
pacchetto di richiesta ed inviare in risposta un pacchetto contenente una
stringa con la data corrente. 

\begin{figure}[!htbp] 
  \footnotesize \centering
  \begin{minipage}[c]{\codesamplewidth}
    \includecodesample{listati/UDP_daytimed.c}
  \end{minipage} 
  \normalsize
  \caption{Sezione principale del server per il servizio \textit{daytime} su
    UDP.}
  \label{fig:UDP_daytime_server}
\end{figure}

In fig.~\ref{fig:UDP_daytime_server} è riportato la sezione principale del
codice del nostro client, il sorgente completo si trova nel file
\file{UDP\_daytimed.c} distribuito con gli esempi allegati alla guida; anche
in questo caso si è omessa la sezione relativa alla gestione delle opzioni a
riga di comando (la sola presente è \texttt{-v} che permette di stampare a
video l'indirizzo associato ad ogni richiesta).

Anche in questo caso la prima parte del server (\texttt{\small 9--23}) è
sostanzialmente identica a quella dell'analogo server per TCP illustrato in
fig.~\ref{fig:TCP_daytime_cunc_server_code}; si inizia (\texttt{\small 10})
con il creare il socket, uscendo con un messaggio in caso di errore
(\texttt{\small 10--13}), e di nuovo la sola differenza con il caso precedente
è il diverso tipo di socket utilizzato. Dopo di che (\texttt{\small 14--18})
si inizializza la struttura degli indirizzi che poi (\texttt{\small 20}) verrà
usata da \func{bind}; si cancella (\texttt{\small 15}) preventivamente il
contenuto, si imposta (\texttt{\small 16}) la famiglia dell'indirizzo, la
porta (\texttt{\small 17}) e l'indirizzo (\texttt{\small 18}) su cui si
riceveranno i pacchetti.  Si noti come in quest'ultimo sia l'indirizzo
generico \const{INADDR\_ANY}; questo significa (si ricordi quanto illustrato
in sez.~\ref{sec:TCP_func_bind}) che il server accetterà pacchetti su uno
qualunque degli indirizzi presenti sulle interfacce di rete della macchina.

Completata l'inizializzazione tutto quello che resta da fare è eseguire
(\texttt{\small 20--23}) la chiamata a \func{bind}, controllando la presenza
di eventuali errori, ed uscendo con un avviso qualora questo fosse il caso.
Nel caso di socket UDP questo è tutto quello che serve per consentire al
server di ricevere i pacchetti a lui indirizzati, e non è più necessario
chiamare successivamente \func{listen}. In questo caso infatti non esiste il
concetto di connessione, e quindi non deve essere predisposta una coda delle
connessioni entranti. Nel caso di UDP i pacchetti arrivano al kernel con un
certo indirizzo ed una certa porta di destinazione, il kernel controlla se
corrispondono ad un socket che è stato \textsl{legato} ad essi con
\func{bind}, qualora questo sia il caso scriverà il contenuto all'interno del
socket, così che il programma possa leggerlo, altrimenti risponderà alla
macchina che ha inviato il pacchetto con un messaggio ICMP di tipo
\textit{port unreachable}.

Una volta completata la fase di inizializzazione inizia il corpo principale
(\texttt{\small 24--44}) del server, mantenuto all'interno di un ciclo
infinito in cui si trattano le richieste. Il ciclo inizia (\texttt{\small 26})
con una chiamata a \func{recvfrom}, che si bloccherà in attesa di pacchetti
inviati dai client. Lo scopo della funzione è quello di ritornare tutte le
volte che un pacchetto viene inviato al server, in modo da poter ricavare da
esso l'indirizzo del client a cui inviare la risposta in \var{addr}. Per
questo motivo in questo caso (al contrario di quanto fatto in
fig.~\ref{fig:UDP_daytime_client}) si è avuto cura di passare gli argomenti
\var{addr} e \var{len} alla funzione.  Dopo aver controllato (\texttt{\small
  27--30}) la presenza di eventuali errori (uscendo con un messaggio di errore
qualora ve ne siano) si verifica (\texttt{\small 31}) se è stata attivata
l'opzione \texttt{-v} (che imposta la variabile \var{verbose}) stampando nel
caso (\texttt{\small 32--35}) l'indirizzo da cui si è appena ricevuto una
richiesta (questa sezione è identica a quella del server TCP illustrato in
fig.~\ref{fig:TCP_daytime_cunc_server_code}).

Una volta ricevuta la richiesta resta solo da ottenere il tempo corrente
(\texttt{\small 36}) e costruire (\texttt{\small 37}) la stringa di risposta,
che poi verrà inviata (\texttt{\small 38}) al client usando \func{sendto},
avendo al solito cura di controllare (\texttt{\small 40--42}) lo stato di
uscita della funzione e trattando opportunamente la condizione di errore.

Si noti come per le peculiarità del protocollo si sia utilizzato un server
iterativo, che processa le richieste una alla volta via via che gli arrivano.
Questa è una caratteristica comune dei server UDP, conseguenza diretta del
fatto che non esiste il concetto di connessione, per cui non c'è la necessità
di trattare separatamente le singole connessioni. Questo significa anche che è
il kernel a gestire la possibilità di richieste multiple in contemporanea;
quello che succede è semplicemente che il kernel accumula in un buffer in
ingresso i pacchetti UDP che arrivano e li restituisce al processo uno alla
volta per ciascuna chiamata di \func{recvfrom}; nel nostro caso sarà poi
compito del server distribuire le risposte sulla base dell'indirizzo da cui
provengono le richieste.


\subsection{Le problematiche dei socket UDP}
\label{sec:UDP_problems}

L'esempio del servizio \textit{daytime} illustrato nelle precedenti sezioni
è in realtà piuttosto particolare, e non evidenzia quali possono essere i
problemi collegati alla mancanza di affidabilità e all'assenza del concetto di
connessione che sono tipiche dei socket UDP. In tal caso infatti il protocollo
è estremamente semplice, dato che la comunicazione consiste sempre in una
richiesta seguita da una risposta, per uno scambio di dati effettuabile con un
singolo pacchetto, per cui tutti gli eventuali problemi sarebbero assai più
complessi da rilevare.

Anche qui però possiamo notare che se il pacchetto di richiesta del client, o
la risposta del server si perdono, il client resterà permanentemente bloccato
nella chiamata a \func{recvfrom}. Per evidenziare meglio quali problemi si
possono avere proviamo allora con un servizio leggermente più complesso come
\textit{echo}. 

\begin{figure}[!htbp] 
  \footnotesize \centering
  \begin{minipage}[c]{\codesamplewidth}
    \includecodesample{listati/UDP_echo_first.c}
  \end{minipage} 
  \normalsize
  \caption{Sezione principale della prima versione client per il servizio
    \textit{echo} su UDP.}
  \label{fig:UDP_echo_client}
\end{figure}

In fig.~\ref{fig:UDP_echo_client} è riportato un estratto del corpo principale
del nostro client elementare per il servizio \textit{echo} (al solito il
codice completo è con i sorgenti allegati). Le uniche differenze con l'analogo
client visto in fig.~\ref{fig:TCP_echo_client_1} sono che al solito si crea
(\texttt{\small 14}) un socket di tipo \const{SOCK\_DGRAM}, e che non è
presente nessuna chiamata a \func{connect}. Per il resto il funzionamento del
programma è identico, e tutto il lavoro viene effettuato attraverso la
chiamata (\texttt{\small 28}) alla funzione \func{ClientEcho} che stavolta
però prende un argomento in più, che è l'indirizzo del socket.

\begin{figure}[!htbp] 
  \footnotesize \centering
  \begin{minipage}[c]{\codesamplewidth}
    \includecodesample{listati/UDP_ClientEcho_first.c}
  \end{minipage}
  \normalsize
  \caption{Codice della funzione \func{ClientEcho} usata dal client per il
    servizio \textit{echo} su UDP.}
  \label{fig:UDP_echo_client_echo}
\end{figure}

Ovviamente in questo caso il funzionamento della funzione, il cui codice è
riportato in fig.~\ref{fig:UDP_echo_client_echo}, è completamente diverso
rispetto alla analoga del server TCP, e dato che non esiste una connessione
questa necessita anche di un terzo argomento, che è l'indirizzo del server cui
inviare i pacchetti.

Data l'assenza di una connessione come nel caso di TCP il meccanismo è molto
più semplice da gestire. Al solito si esegue un ciclo infinito (\texttt{\small
  6--30}) che parte dalla lettura (\texttt{\small 7}) sul buffer di invio
\var{sendbuff} di una stringa dallo standard input, se la stringa è vuota
(\texttt{\small 7--9}), indicando che l'input è terminato, si ritorna
immediatamente causando anche la susseguente terminazione del programma.

Altrimenti si procede (\texttt{\small 10--11}) all'invio della stringa al
destinatario invocando \func{sendto}, utilizzando, oltre alla stringa appena
letta, gli argomenti passati nella chiamata a \func{ClientEcho}, ed in
particolare l'indirizzo del server che si è posto in \var{serv\_addr}; qualora
(\texttt{\small 12}) si riscontrasse un errore si provvederà al solito
(\texttt{\small 13--14}) ad uscire con un messaggio di errore.

Il passo immediatamente seguente (\texttt{\small 17}) l'invio è quello di
leggere l'eventuale risposta del server con \func{recvfrom}; si noti come in
questo caso si sia scelto di ignorare l'indirizzo dell'eventuale pacchetto di
risposta, controllando (\texttt{\small 18--21}) soltanto la presenza di un
errore (nel qual caso al solito si ritorna dopo la stampa di un adeguato
messaggio). Si noti anche come, rispetto all'analoga funzione
\func{ClientEcho} utilizzata nel client TCP illustrato in
sez.~\ref{sec:TCP_echo_client} non si sia controllato il caso di un messaggio
nullo, dato che, nel caso di socket UDP, questo non significa la terminazione
della comunicazione.

L'ultimo passo (\texttt{\small 17}) è quello di terminare opportunamente la
stringa di risposta nel relativo buffer per poi provvedere alla sua stampa
sullo standard output, eseguendo il solito controllo (ed eventuale uscita con
adeguato messaggio informativo) in caso di errore.

In genere fintanto che si esegue il nostro client in locale non sorgerà nessun
problema, se però si proverà ad eseguirlo attraverso un collegamento remoto
(nel caso dell'esempio seguente su una VPN, attraverso una ADSL abbastanza
congestionata) e in modalità non interattiva, la probabilità di perdere
qualche pacchetto aumenta, ed infatti, eseguendo il comando come:
\begin{verbatim}
[piccardi@gont sources]$ cat UDP_echo.c | ./echo 192.168.1.120
/* UDP_echo.c
 *
 * Copyright (C) 2004 Simone Piccardi
...
...
/*
 * Include needed headers

\end{verbatim}%$
si otterrà che, dopo aver correttamente stampato alcune righe, il programma si
blocca completamente senza stampare più niente. Se al contempo si fosse tenuto
sotto controllo il traffico UDP diretto o proveniente dal servizio
\textit{echo} con \cmd{tcpdump} si sarebbe ottenuto:
\begin{verbatim}
[root@gont gapil]# tcpdump  \( dst port 7 or src port 7 \)
...
...
18:48:16.390255 gont.earthsea.ea.32788 > 192.168.1.120.echo: udp 4 (DF)
18:48:17.177613 192.168.1.120.echo > gont.earthsea.ea.32788: udp 4 (DF)
18:48:17.177790 gont.earthsea.ea.32788 > 192.168.1.120.echo: udp 26 (DF)
18:48:17.964917 192.168.1.120.echo > gont.earthsea.ea.32788: udp 26 (DF)
18:48:17.965408 gont.earthsea.ea.32788 > 192.168.1.120.echo: udp 4 (DF)
\end{verbatim}
che come si vede il traffico fra client e server si interrompe dopo l'invio di
un pacchetto UDP per il quale non si è ricevuto risposta.

Il problema è che in tutti i casi in cui un pacchetto di risposta si perde, o
una richiesta non arriva a destinazione, il nostro programma si bloccherà
nell'esecuzione di \func{recvfrom}. Lo stesso avviene anche se il server non è
in ascolto, in questo caso però, almeno dal punto di vista dello scambio di
pacchetti, il risultato è diverso, se si lancia al solito il programma e si
prova a scrivere qualcosa si avrà ugualmente un blocco su \func{recvfrom} ma
se si osserva il traffico con \cmd{tcpdump} si vedrà qualcosa del tipo:
\begin{verbatim}
[root@gont gapil]# tcpdump  \( dst 192.168.0.2 and src 192.168.1.120 \) \
   or \( src 192.168.0.2 and dst 192.168.1.120 \)
tcpdump: listening on eth0
00:43:27.606944 gont.earthsea.ea.32789 > 192.168.1.120.echo: udp 6 (DF)
00:43:27.990560 192.168.1.120 > gont.earthsea.ea: icmp: 192.168.1.120 udp port 
echo unreachable [tos 0xc0]
\end{verbatim}
cioè in questo caso si avrà in risposta un pacchetto ICMP di destinazione
irraggiungibile che ci segnala che la porta in questione non risponde. 

Ci si può chiedere allora perché, benché la situazione di errore sia
rilevabile, questa non venga segnalata. Il luogo più naturale in cui
riportarla sarebbe la chiamata di \func{sendto}, in quanto è a causa dell'uso
di un indirizzo sbagliato che il pacchetto non può essere inviato; farlo in
questo punto però è impossibile, dato che l'interfaccia di programmazione
richiede che la funzione ritorni non appena il kernel invia il
pacchetto,\footnote{questo è il classico caso di \textsl{errore asincrono},
  una situazione cioè in cui la condizione di errore viene rilevata in maniera
  asincrona rispetto all'operazione che l'ha causata, una eventualità
  piuttosto comune quando si ha a che fare con la rete, tutti i pacchetti ICMP
  che segnalano errori rientrano in questa tipologia.} e non può bloccarsi in
una attesa di una risposta che potrebbe essere molto lunga (si noti infatti
che il pacchetto ICMP arriva qualche decimo di secondo più tardi) o non
esserci affatto.

Si potrebbe allora pensare di riportare l'errore nella \func{recvfrom} che è
comunque bloccata in attesa di una risposta che nel caso non arriverà mai.  La
ragione per cui non viene fatto è piuttosto sottile e viene spiegata da
Stevens in \cite{UNP2} con il seguente esempio: si consideri un client che
invia tre pacchetti a tre diverse macchine, due dei quali vengono regolarmente
ricevuti, mentre al terzo, non essendo presente un server sulla relativa
macchina, viene risposto con un messaggio ICMP come il precedente. Detto
messaggio conterrà anche le informazioni relative ad indirizzo e porta del
pacchetto che ha fallito, però tutto quello che il kernel può restituire al
programma è un codice di errore in \var{errno}, con il quale è impossibile di
distinguere per quale dei pacchetti inviati si è avuto l'errore; per questo è
stata fatta la scelta di non riportare un errore su un socket UDP, a meno che,
come vedremo in sez.~\ref{sec:UDP_connect}, questo non sia connesso.



\subsection{L'uso della funzione \func{connect} con i socket UDP}
\label{sec:UDP_connect}

Come illustrato in sez.~\ref{sec:UDP_characteristics} essendo i socket UDP
privi di connessione non è necessario per i client usare \func{connect} prima
di iniziare una comunicazione con un server. Ciò non di meno abbiamo accennato
come questa possa essere utilizzata per gestire la presenza di errori
asincroni.

Quando si chiama \func{connect} su di un socket UDP tutto quello che succede è
che l'indirizzo passato alla funzione viene registrato come indirizzo di
destinazione del socket. A differenza di quanto avviene con TCP non viene
scambiato nessun pacchetto, tutto quello che succede è che da quel momento in
qualunque cosa si scriva sul socket sarà inviata a quell'indirizzo; non sarà
più necessario usare l'argomento \param{to} di \func{sendto} per specificare
la destinazione dei pacchetti, che potranno essere inviati e ricevuti usando
le normali funzioni \func{read} e \func{write} che a questo punto non
falliranno più con l'errore di \errval{EDESTADDRREQ}.\footnote{in realtà si
  può anche continuare ad usare la funzione \func{sendto}, ma in tal caso
  l'argomento \param{to} deve essere inizializzato a \val{NULL}, e
  \param{tolen} deve essere inizializzato a zero, pena un errore di
  \errval{EISCONN}.}

Una volta che il socket è connesso cambia però anche il comportamento in
ricezione; prima infatti il kernel avrebbe restituito al socket qualunque
pacchetto ricevuto con un indirizzo di destinazione corrispondente a quello
del socket, senza nessun controllo sulla sorgente; una volta che il socket
viene connesso saranno riportati su di esso solo i pacchetti con un indirizzo
sorgente corrispondente a quello a cui ci si è connessi.

Infine quando si usa un socket connesso, venendo meno l'ambiguità segnalata
alla fine di sez.~\ref{sec:UDP_problems}, tutti gli eventuali errori asincroni
vengono riportati alle funzioni che operano su di esso; pertanto potremo
riscrivere il nostro client per il servizio \textit{echo} con le modifiche
illustrate in fig.~\ref{fig:UDP_echo_conn_cli}.

\begin{figure}[!htbp] 
  \footnotesize \centering
  \begin{minipage}[c]{\codesamplewidth}
    \includecodesample{listati/UDP_echo.c}
  \end{minipage}
  \normalsize
  \caption{Seconda versione del client del servizio \textit{echo} che utilizza
    socket UDP connessi.}
  \label{fig:UDP_echo_conn_cli}
\end{figure}

Ed in questo caso rispetto alla precedente versione, il solo cambiamento è
l'utilizzo (\texttt{\small 17}) della funzione \func{connect} prima della
chiamata alla funzione di gestione del protocollo, che a sua volta è stata
modificata eliminando l'indirizzo passato come argomento e sostituendo le
chiamata a \func{sendto} e \func{recvfrom} con chiamate a \func{read} e
\func{write} come illustrato dal nuovo codice riportato in
fig.~\ref{fig:UDP_echo_conn_echo_client}.

\begin{figure}[!htbp] 
  \footnotesize \centering
  \begin{minipage}[c]{\codesamplewidth}
    \includecodesample{listati/UDP_ClientEcho.c}
  \end{minipage}
  \normalsize
  \caption{Seconda versione della funzione \func{ClientEcho}.}
  \label{fig:UDP_echo_conn_echo_client}
\end{figure}

Utilizzando questa nuova versione del client si può verificare che quando ci
si rivolge verso un indirizzo inesistente o su cui non è in ascolto un server
si è in grado rilevare l'errore, se infatti eseguiamo il nuovo programma
otterremo un qualcosa del tipo:
\begin{verbatim}
[piccardi@gont sources]$ ./echo 192.168.1.1
prova
Errore in lettura: Connection refused
\end{verbatim}%$

Ma si noti che a differenza di quanto avveniva con il client TCP qui l'errore
viene rilevato soltanto dopo che si è tentato di inviare qualcosa, ed in
corrispondenza al tentativo di lettura della risposta. Questo avviene perché
con UDP non esiste una connessione, e fintanto che non si invia un pacchetto
non c'è traffico sulla rete. In questo caso l'errore sarà rilevato alla
ricezione del pacchetto ICMP \textit{destination unreachable} emesso dalla
macchina cui ci si è rivolti, e questa volta, essendo il socket UDP connesso,
il kernel potrà riportare detto errore in \textit{user space} in maniera non
ambigua, ed esso apparirà alla successiva lettura sul socket.

Si tenga presente infine che l'uso dei socket connessi non risolve l'altro
problema del client, e cioè il fatto che in caso di perdita di un pacchetto
questo resterà bloccato permanentemente in attesa di una risposta. Per
risolvere questo problema l'unico modo sarebbe quello di impostare un
\textit{timeout} o riscrivere il client in modo da usare l'I/O non bloccante.



\index{socket!locali|(}


\section{I socket \textit{Unix domain}}
\label{sec:unix_socket}

Benché i socket Unix domain, come meccanismo di comunicazione fra processi che
girano sulla stessa macchina, non siano strettamente attinenti alla rete, li
tratteremo comunque in questa sezione. Nonostante le loro peculiarità infatti,
l'interfaccia di programmazione che serve ad utilizzarli resta sempre quella
dei socket.



\subsection{Il passaggio di file descriptor}
\label{sec:sock_fd_passing}



\index{socket!locali|)}


\section{Altri socket}
\label{sec:socket_other}

Tratteremo in questa sezione gli altri tipi particolari di socket supportati
da Linux, come quelli relativi a particolare protocolli di trasmissione, i
socket \textit{netlink} che definiscono una interfaccia di comunicazione con
il kernel, ed i \textit{packet socket} che consentono di inviare pacchetti
direttamente a livello delle interfacce di rete. 

\subsection{I socket \textit{raw}}
\label{sec:socket_raw}

Tratteremo in questa sezione i cosiddetti \textit{raw socket}, con i quali si
possono \textsl{forgiare} direttamente i pacchetti a tutti i livelli dello
stack dei protocolli. 

\subsection{I socket \textit{netlink}}
\label{sec:socket_netlink}


\subsection{I \textit{packet socket}}
\label{sec:packet_socket}


% articoli interessanti:
% http://www.linuxjournal.com/article/5617
% http://www.linuxjournal.com/article/4659
% 


% LocalWords:  socket cap TCP UDP domain sez NFS DNS stream datagram PF INET to
% LocalWords:  IPv tab SOCK DGRAM three way handshake client fig bind listen AF
% LocalWords:  accept recvfrom sendto connect netstat named DHCP kernel ICMP CR
% LocalWords:  port unreachable read write glibc Specification flags int BSD LF
% LocalWords:  libc unsigned len size tolen sys ssize sockfd const void buf MSG
% LocalWords:  struct sockaddr socklen errno EAGAIN ECONNRESET EDESTADDRREQ RFC
% LocalWords:  EISCONN EMSGSIZE ENOBUFS ENOTCONN EOPNOTSUPP EPIPE SIGPIPE EBADF
% LocalWords:  NOSIGNAL EFAULT EINVAL EINTR ENOMEM ENOTSOCK NULL fromlen from
% LocalWords:  ECONNREFUSED value result argument close shutdown daytime nell'
% LocalWords:  memset inet pton nread NUL superdemone inetd sniffer daytimed
% LocalWords:  INADDR ANY addr echo ClientEcho sendbuff serv VPN tcpdump l'I
% LocalWords:  Stevens destination descriptor raw stack netlink packet



%%% Local Variables: 
%%% mode: latex
%%% TeX-master: "gapil"
%%% End: 

%% sockadv.tex
%%
%% Copyright (C) 2004-2018 Simone Piccardi.  Permission is granted to
%% copy, distribute and/or modify this document under the terms of the GNU Free
%% Documentation License, Version 1.1 or any later version published by the
%% Free Software Foundation; with the Invariant Sections being "Un preambolo",
%% with no Front-Cover Texts, and with no Back-Cover Texts.  A copy of the
%% license is included in the section entitled "GNU Free Documentation
%% License".
%%

\chapter{Socket avanzati}
\label{cha:advanced_socket}


Esamineremo in questo capitolo le funzionalità più evolute della gestione dei
socket, le funzioni avanzate, la gestione dei dati urgenti e
\textit{out-of-band} e dei messaggi ancillari, come l'uso come l'uso del I/O
multiplexing (vedi sez.~\ref{sec:file_multiplexing}) con i socket.


\section{Le funzioni di I/O avanzate}
\label{sec:sock_advanced_IO}

Tratteremo in questa sezione le funzioni di I/O più avanzate che permettono di
controllare le funzionalità specifiche della comunicazione dei dati che sono
disponibili con i vari tipi di socket.

\subsection{La funzioni \func{send} e \func{recv}}
\label{sec:net_send_recv}

Da fare


% TODO: note su MSG_ZEROCOPY/SOCK_ZEROCOPY, aggiunte con il kernel 4.14 (e per
% la ricezione con il kernel 4.18, vedi https://lwn.net/Articles/726917/ e
% https://lwn.net/Articles/752300/ 

\subsection{La funzioni \func{sendmsg} e \func{recvmsg}}
\label{sec:net_sendmsg}

Finora abbiamo trattato delle funzioni che permettono di inviare dati sul
socket in forma semplificata. Se infatti si devono semplicemente ...

% TODO trattare anche recvmmsg, introdotta con il kernel 2.3.33, vedi
% http://kernelnewbies.org/Linux_2_6_33 

% TODO trattare anche sendmmsg, introdotta con il kernel 3.0, vedi
% 


\subsection{I messaggi ancillari}
\label{sec:net_ancillary_data}

Quanto è stata attivata l'opzione \const{IP\_RECVERR} il kernel attiva per il
socket una speciale coda su cui vengono inviati tutti gli errori riscontrati.
Questi possono essere riletti usando il flag \const{MSG\_ERRQUEUE}, nel qual
caso sarà passato come messaggio ancillare una struttura di tipo
\struct{sock\_extended\_err} illustrata in
fig.~\ref{fig:sock_extended_err_struct}.

\begin{figure}[!htb]
  \footnotesize \centering
  \begin{minipage}[c]{\textwidth}
    \includestruct{listati/sock_extended_err.h}
  \end{minipage}
  \caption{La struttura \structd{sock\_extended\_err} usata dall'opzione
    \const{IP\_RECVERR} per ottenere le informazioni relative agli errori su
    un socket.}
  \label{fig:sock_extended_err_struct}
\end{figure}

% TODO vedi man cmsg

\subsection{I \textsl{dati urgenti} o \textit{out-of-band}}
\label{sec:TCP_urgent_data}

\itindbeg{out-of-band} 

Una caratteristica particolare dei socket TCP è quella che consente di inviare
all'altro capo della comunicazione una sorta di messaggio privilegiato, che si
richiede che sia trattato il prima possibile. Si fa riferimento a questa
funzionalità come all'invio dei cosiddetti \textsl{dati urgenti} (o
\textit{urgent data}); talvolta essi chiamati anche dati \textit{out-of-band}
poiché, come vedremo più avanti, possono essere letti anche al di fuori del
flusso di dati normale.

Come già accennato in sez.~\ref{sec:file_multiplexing} la presenza di dati
urgenti viene rilevata in maniera specifica sia di \func{select} (con il
\textit{file descriptor set} \param{exceptfds}) che da \func{poll} (con la
condizione \const{POLLRDBAND}).


Le modalità di lettura dei dati urgenti sono due, la prima e più comune
prevede l'uso di \func{recvmsg} con 


% TODO aggiungere pezzo di codice per inviare dati urgenti all'echo server

La seconda modalità di lettura prevede invece l'uso dell'opzione dei socket
\const{SO\_OOBINLINE} (vedi sez.~\ref{sec:sock_generic_options}) che consente
di ricevere i dati urgenti direttamente nel flusso dei dati del socket; in tal
caso però si pone il problema di come distinguere i dati normali da quelli
urgenti. Come già accennato in sez.~\ref{sec:sock_ioctl_IP} a questo scopo si
può usare \func{ioctl} con l'operazione \const{SIOCATMARK}, che consente di
sapere se si è arrivati o meno all'\textit{urgent mark}. 

La procedura allora prevede che, una volta che si sia rilevata la presenza di
dati urgenti, si ripeta la lettura ordinaria dal socket fintanto che
\const{SIOCATMARK} non restituisce un valore diverso da zero; la successiva
lettura restituirà i dati urgenti.


\itindend{out-of-band} 


% TODO trattare sendmmsg aggiunta con il kernel 3.0.
% e gli ICMP socket, vedi  http://lwn.net/Articles/420799/


\section{L'uso dell'I/O non bloccante}
\label{sec:sock_noblok_IO}

Tratteremo in questa sezione le modalità avanzate che permettono di utilizzare
i socket con una comunicazione non bloccante, in modo da 





\subsection{La gestione delle opzioni IP}  
\label{sec:sock_IP_options}  

Abbiamo visto in sez.~\ref{sec:sock_ipv4_options} come di possa usare  
\func{setsockopt} con l'opzione \const{IP\_OPTIONS} per impostare le opzioni  
IP associate per i pacchetti associati ad un socket.  Vedremo qui il
significato di tali opzioni e le modalità con cui esse possono essere
utilizzate ed impostate.

% LocalWords:  socket of multiplexing sez sendmsg recvmsg RECVERR kernel MSG
% LocalWords:  ERRQUEUE sock err fig TCP dell'I setsockopt OPTIONS urgent poll
% LocalWords:  select descriptor exceptfds POLLRDBAND OOBINLINE ioctl all' mark
% LocalWords:  SIOCATMARK

%%% Local Variables: 
%%% mode: latex
%%% TeX-master: "gapil"
%%% End: 

\appendix

\part{Appendici}
\label{part:appendici}
%% errors.tex
%%
%% Copyright (C) 2000-2018 Simone Piccardi.  Permission is granted to
%% copy, distribute and/or modify this document under the terms of the GNU Free
%% Documentation License, Version 1.1 or any later version published by the
%% Free Software Foundation; with the Invariant Sections being "Un preambolo",
%% with no Front-Cover Texts, and with no Back-Cover Texts.  A copy of the
%% license is included in the section entitled "GNU Free Documentation
%% License".
%%

\chapter{I codici di errore}
\label{cha:errors}

Si riportano in questa appendice tutti i codici di errore. Essi sono
accessibili attraverso l'inclusione del file di header \headfile{errno.h}, che
definisce anche la variabile globale \var{errno}. Per ogni errore definito
riporteremo la stringa stampata da \func{perror} ed una breve spiegazione. Si
tenga presente che spiegazioni più particolareggiate del significato
dell'errore, qualora necessarie per casi specifici, possono essere trovate
nella descrizione del prototipo della funzione per cui detto errore si è
verificato.

I codici di errore sono riportati come costanti di tipo \ctyp{int}, i valori
delle costanti sono definiti da macro di preprocessore nel file citato, e
possono variare da architettura a architettura; è pertanto necessario
riferirsi ad essi tramite i nomi simbolici. Le funzioni \func{perror} e
\func{strerror} (vedi sez.~\ref{sec:sys_strerror}) possono essere usate per
ottenere dei messaggi di errore più espliciti.


\section{Gli errori dei file}
\label{sec:err_file_errors}

In questa sezione sono raccolti i codici restituiti dalle \textit{system call}
attinenti ad errori che riguardano operazioni specifiche relative alla
gestione dei file.

\begin{basedescript}{\desclabelwidth{1.5cm}\desclabelstyle{\nextlinelabel}}
\item[\errcode{EACCES} \textit{Permission denied}.] Permesso negato; l'accesso
  al file o alla directory non è consentito: i permessi del file o della
  directory o quelli necessari ad attraversare un \textit{pathname} non
  consentono l'operazione richiesta. 
\item[\errcode{EBADF} \textit{Bad file descriptor}.] File descriptor non
  valido: si è usato un file descriptor inesistente, o aperto in sola lettura
  per scrivere, o viceversa, o si è cercato di eseguire un'operazione non
  consentita per quel tipo di file descriptor.
\item[\errcode{EBUSY} \textit{Resource busy}.] Una risorsa di sistema che non
  può essere condivisa è occupata. Ad esempio si è tentato di cancellare la
  directory su cui si è montato un filesystem.
\item[\errcode{EDQUOT} \textit{Quota exceeded}.] Si è ecceduta la quota di
  disco dell'utente, usato sia per lo spazio disco che per il numero di
  \textit{inode}.
\item[\errcode{EEXIST} \textit{File exists}.] Si è specificato un file
  esistente in un contesto in cui ha senso solo specificare un nuovo file.
\item[\errcode{EFBIG} \textit{File too big}.] Si è ecceduto il limite imposto
  dal sistema sulla dimensione massima che un file può avere.
\item[\errcode{EFTYPE} \textit{Inappropriate file type or format}.] Il file è
  di tipo sbagliato rispetto all'operazione richiesta o un file di dati ha un
  formato sbagliato. Alcuni sistemi restituiscono questo errore quando si
  cerca di impostare lo \textit{sticky bit} su un file che non è una
  directory.
\item[\errcode{EIO} \textit{Input/output error}.] Errore di input/output:
  usato per riportare errori hardware in lettura/scrittura su un dispositivo.
\item[\errcode{EISDIR} \textit{Is a directory}.] Il file specificato è una
  directory; non può essere aperto in scrittura, né si possono creare o
  rimuovere link diretti ad essa.
\item[\errcode{ELOOP} \textit{Too many symbolic links encountered}.] Ci sono
  troppi link simbolici nella risoluzione di un \textit{pathname}.
\item[\errcode{EMFILE} \textit{Too many open files}.] Il processo corrente ha
  troppi file aperti e non può aprirne altri. Anche i descrittori duplicati ed
  i socket vengono tenuti in conto.\footnote{il numero massimo di file aperti
    è controllabile dal sistema; in Linux si può impostare usando il comando
    \cmd{ulimit}, esso è in genere indicato dalla costante \const{OPEN\_MAX},
    vedi sez.~\ref{sec:sys_limits}.}
\item[\errcode{EMLINK} \textit{Too many links}.] Ci sono già troppi link al
  file, il numero massimo è specificato dalla variabile \const{LINK\_MAX},
  vedi sez.~\ref{sec:sys_limits}.
\item[\errcode{ENAMETOOLONG} \textit{File name too long}.] Si è indicato un
  \textit{pathname} troppo lungo per un file o una directory.
\item[\errcode{ENFILE} \textit{File table overflow}.] Il sistema ha troppi
  file aperti in contemporanea. Si tenga presente che anche i socket contano
  come file. Questa è una condizione temporanea, ed è molto difficile che si
  verifichi nei sistemi moderni.
\item[\errcode{ENODEV} \textit{No such device}.] Si è indicato un tipo di
  device sbagliato ad una funzione che ne richiede uno specifico.
\item[\errcode{ENOENT} \textit{No such file or directory}.] Il file indicato
  dal \textit{pathname} non esiste: o una delle componenti non esiste o il
  \textit{pathname} contiene un link simbolico spezzato.  Errore tipico di un
  riferimento ad un file che si suppone erroneamente essere esistente.
\item[\errcode{ENOEXEC} \textit{Invalid executable file format}.] Il file non
  ha un formato eseguibile, è un errore riscontrato dalle funzioni
  \func{exec}.
\item[\errcode{ENOLCK} \textit{No locks available}.] È usato dalle utilità per
  la gestione del file locking; non viene generato da un sistema GNU, ma può
  risultare da un'operazione su un server NFS di un altro sistema.
\item[\errcode{ENOSPC} \textit{No space left on device}.] La directory in cui
  si vuole creare il link non ha spazio per ulteriori voci, o si è cercato di
  scrivere o di creare un nuovo file su un dispositivo che è già pieno.
\item[\errcode{ENOTBLK} \textit{Block device required}.] Si è specificato un
  file che non è un \textit{block device} in un contesto in cui era necessario
  specificare un \textit{block device} (ad esempio si è tentato di montare un
  file ordinario).
\item[\errcode{ENOTDIR} \textit{Not a directory}.] Si è specificato un file
  che non è una directory in una operazione che richiede una directory.
\item[\errcode{ENOTEMPTY} \textit{Directory not empty}.] La directory non è
  vuota quando l'operazione richiede che lo sia. È l'errore tipico che si ha
  quando si cerca di cancellare una directory contenente dei file.
\item[\errcode{ENOTTY} \textit{Not a terminal}.] Si è tentata una operazione
  di controllo relativa ad un terminale su un file che non lo è.
\item[\errcode{ENXIO} \textit{No such device or address}.] Dispositivo
  inesistente: il sistema ha tentato di usare un dispositivo attraverso il
  file specificato, ma non lo ha trovato. Può significare che il file di
  dispositivo non è corretto, che il modulo relativo non è stato caricato nel
  kernel, o che il dispositivo è fisicamente assente o non funzionante.
\item[\errcode{EPERM} \textit{Operation not permitted}.] L'operazione non è
  permessa: solo il proprietario del file o un processo con sufficienti
  privilegi può eseguire l'operazione.
\item[\errcode{EPIPE} \textit{Broken pipe}.] Non c'è un processo che stia
  leggendo l'altro capo della \textit{pipe}. Ogni funzione che restituisce
  questo errore genera anche un segnale \signal{SIGPIPE}, la cui azione
  predefinita è terminare il programma; pertanto non si potrà vedere questo
  errore fintanto che \signal{SIGPIPE} non viene gestito o bloccato.
\item[\errcode{EREMOTE} \textit{Object is remote}.] Si è fatto un tentativo di
  montare via NFS un filesystem remoto con un nome  che già specifica un
  filesystem montato via NFS. 
\item[\errcode{EROFS} \textit{Read-only file system}.]  Si è cercato di
  eseguire una operazione di scrittura su un file o una directory che risiede
  su un filesystem montato un sola lettura.
\item[\errcode{ESPIPE} \textit{Invalid seek operation}.] Si cercato di
  eseguire una \func{lseek} su un file che non supporta questa operazione (ad
  esempio su una \textit{pipe}, da cui il nome).
\item[\errcode{ESTALE} \textit{Stale NFS file handle}.] Indica un problema
  interno a NFS causato da cambiamenti del filesystem del sistema remoto. Per
  recuperare questa condizione in genere è necessario smontare e rimontare il
  filesystem NFS.
\item[\errcode{ETXTBSY} \textit{Text file busy}.] Si è cercato di eseguire un
  file che è aperto in scrittura, o di scrivere su un file che è in
  esecuzione.
\item[\errcode{EUSERS} \textit{Too many users}.] Troppi utenti, il sistema delle
  quote rileva troppi utenti nel sistema.
\item[\errcode{EXDEV} \textit{Cross-device link}.] Si è tentato di creare un
  link diretto che attraversa due filesystem differenti.
\end{basedescript}

\section{Gli errori dei processi}
\label{sec:err_proc_errors}

In questa sezione sono raccolti i codici restituiti dalle \textit{system call}
attinenti ad errori che riguardano operazioni specifiche relative alla
gestione dei processi.

\begin{basedescript}{\desclabelwidth{1.5cm}\desclabelstyle{\nextlinelabel}}
\item[\errcode{E2BIG} \textit{Argument list too long}.] La lista degli
  argomenti passati è troppo lunga: è una condizione prevista da POSIX quando
  la lista degli argomenti passata ad una delle funzioni \func{exec} occupa
  troppa memoria.
\item[\errcode{ECHILD} \textit{There are no child processes}.] Non esistono
  processi figli di cui attendere la terminazione. Viene rilevato dalle
  funzioni \func{wait} e \func{waitpid} (vedi sez.~\ref{sec:proc_wait}).
\item[\errcode{EPROCLIM} \textit{Too many processes}.]  Il limite dell'utente
  per nuovi processi (vedi sez.~\ref{sec:sys_resource_limit}) sarà ecceduto
  alla prossima \func{fork}; è un codice di errore di BSD, che non viene
  utilizzato al momento su Linux. 
\item[\errcode{ESRCH} \textit{No process matches the specified process ID}.]
  Non esiste un processo o un \textit{process group} corrispondenti al valore
  dell'identificativo specificato.
\end{basedescript}


\section{Gli errori di rete}
\label{sec:err_network}

In questa sezione sono raccolti i codici restituiti dalle \textit{system call}
attinenti ad errori che riguardano operazioni specifiche relative alla
gestione dei socket e delle connessioni di rete.

\begin{basedescript}{\desclabelwidth{1.5cm}\desclabelstyle{\nextlinelabel}}
\item[\errcode{EADDRINUSE} \textit{Address already in use}.] L'indirizzo del
  socket richiesto è già utilizzato (ad esempio si è eseguita \func{bind}
  su una porta già in uso).
\item[\errcode{EADDRNOTAVAIL} \textit{Cannot assign requested
    address}.]  L'indirizzo richiesto non è disponibile (ad esempio si
  è cercato di dare al socket un nome che non corrisponde al nome
  della stazione locale), o l'interfaccia richiesta non esiste.
\item[\errcode{EAFNOSUPPORT} \textit{Address family not supported by
    protocol}.]  Famiglia di indirizzi non supportata. La famiglia di
  indirizzi richiesta non è supportata, o è inconsistente con il protocollo
  usato dal socket.
\item[\errcode{ECONNABORTED} \textit{Software caused connection abort}.] Una
  connessione è stata abortita localmente. 
\item[\errcode{ECONNREFUSED} \textit{Connection refused}.] Un host remoto ha
  rifiutato la connessione (in genere dipende dal fatto che non c'è un server
  per soddisfare il servizio richiesto).  
\item[\errcode{ECONNRESET} \textit{Connection reset by peer}.] Una connessione
  è stata chiusa per ragioni fuori dal controllo dell'host locale, come il
  riavvio di una macchina remota o un qualche errore non recuperabile sul
  protocollo.
\item[\errcode{EDESTADDRREQ} \textit{Destination address required}.] Non c'è
  un indirizzo di destinazione predefinito per il socket. Si ottiene questo
  errore mandando dato su un socket senza connessione senza averne prima
  specificato una destinazione.
\item[\errcode{EISCONN} \textit{Transport endpoint is already connected}.] Si
  è tentato di connettere un socket che è già connesso.
\item[\errcode{EMSGSIZE} \textit{Message too long}.] Le dimensioni di un
  messaggio inviato su un socket sono eccedono la massima lunghezza supportata.
\item[\errcode{ENETDOWN} \textit{Network is down}.] L'operazione sul socket è
  fallita perché la rete è sconnessa.
\item[\errcode{ENETRESET} \textit{Network dropped connection because of
    reset}.]  Una connessione è stata cancellata perché l'host remoto è
  caduto.
\item[\errcode{ENETUNREACH} \textit{Network is unreachable}.] L'operazione è
  fallita perché l'indirizzo richiesto è irraggiungibile (ad esempio la
  sottorete della stazione remota è irraggiungibile).
\item[\errcode{ENOBUFS} \textit{No buffer space available}.] Tutti i buffer per
  le operazioni di I/O del kernel sono occupati. In generale questo errore è
  sinonimo di \errcode{ENOMEM}, ma attiene alle funzioni di input/output. In
  caso di operazioni sulla rete si può ottenere questo errore invece
  dell'altro.
\item[\errcode{ENOPROTOOPT} \textit{Protocol not available}.] Protocollo non
  disponibile. Si è richiesta un'opzione per il socket non disponibile con il
  protocollo usato.
\item[\errcode{ENOTCONN} \textit{Transport endpoint is not connected}.] Il
  socket non è connesso a niente. Si ottiene questo errore quando si cerca di
  trasmettere dati su un socket senza avere specificato in precedenza la loro
  destinazione. Nel caso di socket senza connessione (ad esempio socket UDP)
  l'errore che si ottiene è \errcode{EDESTADDRREQ}.
\item[\errcode{ENOTSOCK} \textit{Socket operation on non-socket}.] Si è
  tentata un'operazione su un file descriptor che non è un socket quando
  invece era richiesto un socket.
\item[\errcode{EOPNOTSUPP} \textit{Operation not supported on transport
    endpoint}.] L'operazione richiesta non è supportata. Alcune funzioni non
  hanno senso per tutti i tipi di socket, ed altre non sono implementate per
  tutti i protocolli di trasmissione. Questo errore quando un socket non
  supporta una particolare operazione, e costituisce una indicazione generica
  che il server non sa cosa fare per la chiamata effettuata.
\item[\errcode{EPFNOSUPPORT} \textit{Protocol family not supported}.] Famiglia
  di protocolli non supportata. La famiglia di protocolli richiesta non è
  supportata.
\item[\errcode{EPROTONOSUPPORT} \textit{Protocol not supported}.] Protocollo
  non supportato. Il tipo di socket non supporta il protocollo richiesto (un
  probabile errore nella specificazione del protocollo).
\item[\errcode{EPROTOTYPE} \textit{Protocol wrong type for socket}.]
  Protocollo sbagliato per il socket. Il socket usato non supporta il
  protocollo di comunicazione richiesto.
\item[\errcode{ESHUTDOWN} \textit{Cannot send after transport endpoint
    shutdown}.] Il socket su cui si cerca di inviare dei dati ha avuto uno
  shutdown.
\item[\errcode{ESOCKTNOSUPPORT} \textit{Socket type not supported}.] Socket
  non supportato. Il tipo di socket scelto non è supportato.
\item[\errcode{ETIMEDOUT} \textit{Connection timed out}.] Un'operazione sul
  socket non ha avuto risposta entro il periodo di timeout.
\item[\errcode{ETOOMANYREFS} \textit{Too many references: cannot splice}.] La
  \acr{glibc} dice ???
\end{basedescript}


\section{Errori generici}

In questa sezione sono raccolti i codici restituiti dalle \textit{system call}
attinenti ad errori generici, si trovano qui tutti i codici di errore non
specificati nelle sezioni precedenti.

\begin{basedescript}{\desclabelwidth{1.5cm}\desclabelstyle{\nextlinelabel}}
\item[\errcode{EAGAIN} \textit{Resource temporarily unavailable}.] La funzione è
  fallita ma potrebbe funzionare se la chiamata fosse ripetuta. Questo errore
  accade in due tipologie di situazioni:
  \begin{itemize}
  \item Si è effettuata un'operazione che si sarebbe bloccata su un oggetto
    che è stato posto in modalità non bloccante. Nei vecchi sistemi questo era
    un codice diverso, \errcode{EWOULDBLOCK}. In genere questo ha a che fare
    con file o socket, per i quali si può usare la funzione \func{select} per
    vedere quando l'operazione richiesta (lettura, scrittura o connessione)
    diventa possibile.
  \item Indica la carenza di una risorsa di sistema che non è al momento
    disponibile (ad esempio \func{fork} può fallire con questo errore se si è
    esaurito il numero di processi contemporanei disponibili). La ripetizione
    della chiamata in un periodo successivo, in cui la carenza della risorsa
    richiesta può essersi attenuata, può avere successo. Questo tipo di
    carenza è spesso indice di qualcosa che non va nel sistema, è pertanto
    opportuno segnalare esplicitamente questo tipo di errori.
  \end{itemize}
\item[\errcode{EALREADY} \textit{Operation already in progress}.] L'operazione è
  già in corso. Si è tentata un'operazione già in corso su un oggetto posto in
  modalità non-bloccante.
\item[\errcode{EDEADLK} \textit{Deadlock avoided}.] L'allocazione di una
  risorsa avrebbe causato un \textit{deadlock}. Non sempre il sistema è in
  grado di riconoscere queste situazioni, nel qual caso si avrebbe il blocco.
\item[\errcode{EFAULT} \textit{Bad address}.] Una stringa passata come
  argomento è fuori dello spazio di indirizzi del processo, in genere questa
  situazione provoca direttamente l'emissione di un segnale di \textit{segment
    violation} (\signal{SIGSEGV}).
\item[\errcode{EDOM} \textit{Domain error}.] È usato dalle funzioni matematiche
  quando il valore di un argomento è al di fuori dell'intervallo in cui esse
  sono definite.
\item[\errcode{EILSEQ} \textit{Illegal byte sequence}.] Nella decodifica di un
  carattere esteso si è avuta una sequenza errata o incompleta o si è
  specificato un valore non valido.
\item[\errcode{EINPROGRESS} \textit{Operation now in progress}.] Operazione in
  corso. Un'operazione che non può essere completata immediatamente è stata
  avviata su un oggetto posto in modalità non-bloccante. Questo errore viene
  riportato per operazioni che si dovrebbero sempre bloccare (come per una
  \func{connect}) e che pertanto non possono riportare \errcode{EAGAIN},
  l'errore indica che l'operazione è stata avviata correttamente e occorrerà
  del tempo perché si possa completare. La ripetizione della chiamata darebbe
  luogo ad un errore \errcode{EALREADY}.
\item[\errcode{EINTR} \textit{Interrupted function call}.] Una funzione di
  libreria è stata interrotta. In genere questo avviene causa di un segnale
  asincrono al processo che impedisce la conclusione della chiamata, la
  funzione ritorna con questo errore una volta che si sia correttamente
  eseguito il gestore del segnale. In questo caso è necessario ripetere la
  chiamata alla funzione.
\item[\errcode{EINVAL} \textit{Invalid argument}.] Errore utilizzato per
  segnalare vari tipi di problemi dovuti all'aver passato un argomento
  sbagliato ad una funzione di libreria.
\item[\errcode{ENOMEM} \textit{No memory available}.] Il kernel non è in grado
  di allocare ulteriore memoria per completare l'operazione richiesta.
\item[\errcode{ENOSYS} \textit{Function not implemented}.] Indica che la
  funzione non è supportata o nelle librerie del C o nel kernel. Può dipendere
  sia dalla mancanza di una implementazione, che dal fatto che non si è
  abilitato l'opportuno supporto nel kernel; nel caso di Linux questo può
  voler dire anche che un modulo necessario non è stato caricato nel sistema.
\item[\errcode{ENOTSUP} \textit{Not supported}.] Una funzione ritorna questo
  errore quando gli argomenti sono validi ma l'operazione richiesta non è
  supportata. Questo significa che la funzione non implementa quel particolare
  comando o opzione o che, in caso di oggetti specifici (file descriptor o
  altro) non è in grado di supportare i parametri richiesti.
\item[\errcode{ERANGE} \textit{Range error}.] È usato dalle funzioni
  matematiche quando il risultato dell'operazione non è rappresentabile nel
  valore di ritorno a causa di un overflow o di un underflow.
\item[\errcode{EWOULDBLOCK} \textit{Operation would block}.] Indica che
  l'operazione richiesta si bloccherebbe, ad esempio se si apre un file in
  modalità non bloccante, una \func{read} restituirebbe questo errore per
  indicare che non ci sono dati; in Linux è identico a \errcode{EAGAIN}, ma in
  altri sistemi può essere specificato un valore diverso.
\end{basedescript}


\begin{basedescript}{\desclabelwidth{1.5cm}\desclabelstyle{\nextlinelabel}}
% definiti nel manuale delle glibc ma inesistenti in linux/errno.h
%\item[\errcode{EBADRPC} \textit{}.] 
%\item[\errcode{ERPCMISMATCH} \textit{}.] 
%\item[\errcode{EPROGUNAVAIL} \textit{}.] 
%\item[\errcode{EPROGMISMATCH} \textit{}.] 
%\item[\errcode{EPROCUNAVAIL} \textit{}.] 
%\item[\errcode{EAUTH} \textit{}.] 
%\item[\errcode{ENEEDAUTH} \textit{}.] 
%\item[\errcode{EBACKGROUND} \textit{}.] 
%\item[\errcode{EDIED} \textit{}.] 
% questi sembrano scherzi, sempre dal manuale delle glibc...
%\item[\errcode{ED} \textit{}.] 
%\item[\errcode{EGREGIOUS} \textit{}.] 
%\item[\errcode{EIEIO} \textit{}.] 
%\item[\errcode{EGRATUITOUS} \textit{} roba di Hurd, pare. 


\item[\errcode{EBADMSG} \textit{Not a data message}.] Definito da POSIX come
errore che arriva ad una funzione di lettura che opera su uno stream. Non
essendo gli stream definiti su Linux il kernel non genera mai questo tipo di
messaggio. 

\item[\errcode{EIDRM} \textit{Identifier removed}.] Indica che l'oggetto del
  \textit{SysV IPC} a cui si fa riferimento è stato cancellato.

\item[\errcode{EMULTIHOP} \textit{Multihop attempted}.] Definito da POSIX come
  errore dovuto all'accesso a file remoti attraverso più macchine, quando ciò
  non è consentito. Non viene mai generato su Linux.

\item[\errcode{ENOATTR} \textit{No such attribute}.] È un codice di errore
  specifico di Linux utilizzato dalle funzioni per la gestione degli attributi
  estesi dei file (vedi sez.~\ref{sec:file_xattr}) quando il nome
  dell'attributo richiesto non viene trovato.

\item[\errcode{ENODATA} \textit{No data available}.] Viene indicato da POSIX
  come restituito da una \func{read} eseguita su un file descriptor in
  modalità non bloccante quando non ci sono dati. In realtà in questo caso su
  Linux viene utilizzato \errcode{EAGAIN}. Lo stesso valore però viene usato
  come sinonimo di \errcode{ENOATTR}.

\item[\errcode{ENOLINK} \textit{Link has been severed}.] È un errore il cui
  valore è indicato come \textsl{riservato} nelle \textit{Single Unix
    Specification}. Dovrebbe indicare l'impossibilità di accedere ad un file a
  causa di un errore sul collegamento di rete, ma non ci sono indicazioni
  precise del suo utilizzo. Per quanto riguarda Linux viene riportato nei
  sorgenti del kernel in alcune operazioni relative ad operazioni di rete. 

\item[\errcode{ENOMSG} \textit{No message of desired type}.] Indica che in una
  coda di messaggi del \textit{SysV IPC} non è presente nessun messaggio del
  tipo desiderato.

\item[\errcode{ENOSR} \textit{Out of streams resources}.] Errore relativo agli
  \textit{STREAMS}, che indica l'assenza di risorse sufficienti a completare
  l'operazione richiesta. Quella degli \textit{STREAMS}\footnote{che non vanno
    confusi con gli \textit{stream} di sez.~\ref{sec:files_std_interface}.}  è
  interfaccia di programmazione originaria di System V, che non è implementata
  da Linux, per cui questo errore non viene utilizzato.

\item[\errcode{ENOSTR} \textit{Device not a stream}.] Altro errore relativo
  agli \textit{STREAMS}, anch'esso non utilizzato da Linux.

\item[\errcode{EOVERFLOW} \textit{Value too large for defined data type}.] Si è
  chiesta la lettura di un dato dal \textit{SysV IPC} con \const{IPC\_STAT} ma
  il valore eccede la dimensione usata nel buffer di lettura.

\item[\errcode{EPROTO} \textit{Protocol error}.] Indica che c'è stato un errore
  nel protocollo di rete usato dal socket; viene usato come errore generico
  dall'interfaccia degli \textit{STREAMS} quando non si è in grado di
  specificare un altro codice di errore che esprima più accuratamente la
  situazione.

\item[\errcode{ETIME} \textit{Timer expired}.] Indica che è avvenuto un timeout
  nell'accesso ad una risorsa (ad esempio un semaforo). Compare nei sorgenti
  del kernel (in particolare per le funzioni relativa al bus USB) come
  indicazione di una mancata risposta di un dispositivo, con una descrizione
  alternativa di \textit{Device did not respond}.
\end{basedescript}



% \section{Errori del kernel}
% \label{sec:err_kernel_err}

% In questa sezione sono raccolti i codici di errore interni del kernel. Non
% sono usati dalle funzioni di libreria, ma vengono riportati da alcune system
% call 
% TODO verificare i dettagli degli errori del kernel, eventualmente cassare.

% \begin{description}
% \item[\errcode{ERESTART} \textit{Interrupted system call should be restarted}.] 
% \item[\errcode{ECHRNG} \textit{Channel number out of range}.] 
% \item[\errcode{EL2NSYNC} \textit{Level 2 not synchronized}.] 
% \item[\errcode{EL3HLT} \textit{Level 3 halted}.] 
% \item[\errcode{EL3RST} \textit{Level 3 reset}.] 
% \item[\errcode{ELNRNG} \textit{Link number out of range}.] 
% \item[\errcode{EUNATCH} \textit{Protocol driver not attached}.] 
% \item[\errcode{ENOCSI} \textit{No CSI structure available}.] 
% \item[\errcode{EL2HLT} \textit{Level 2 halted}.] 
% \item[\errcode{EBADE} \textit{Invalid exchange}.] 
% \item[\errcode{EBADR} \textit{Invalid request descriptor}.] 
% \item[\errcode{EXFULL} \textit{Exchange full}.] 
% \item[\errcode{ENOANO} \textit{No anode}.] 
% \item[\errcode{EBADRQC} \textit{Invalid request code}.] 
% \item[\errcode{EBADSLT} \textit{Invalid slot}.] 
% \item[\errcode{EDEADLOCK} Identico a \errcode{EDEADLK}.] 
% \item[\errcode{EBFONT} \textit{Bad font file format}.] 
% \item[\errcode{ENONET} \textit{Machine is not on the network}.] 
% \item[\errcode{ENOPKG} \textit{Package not installed}.] 
% \item[\errcode{EADV} \textit{Advertise error}.] 
% \item[\errcode{ESRMNT} \textit{Srmount error}.] 
% \item[\errcode{ECOMM} \textit{Communication error on send}.] 
% \item[\errcode{EDOTDOT} \textit{RFS specific error}.] 
% \item[\errcode{ENOTUNIQ} \textit{Name not unique on network}.] 
% \item[\errcode{EBADFD} \textit{File descriptor in bad state}.] 
% \item[\errcode{EREMCHG} \textit{Remote address changed}.] 
% \item[\errcode{ELIBACC} \textit{Can not access a needed shared library}.] 
% \item[\errcode{ELIBBAD} \textit{Accessing a corrupted shared library}.] 
% \item[\errcode{ELIBSCN} \textit{.lib section in a.out corrupted}.] 
% \item[\errcode{ELIBMAX} \textit{Attempting to link in too many shared
%     libraries}.]
% \item[\errcode{ELIBEXEC} \textit{Cannot exec a shared library directly}.] 
% \item[\errcode{ESTRPIPE} \textit{Streams pipe error}.] 
% \item[\errcode{EUCLEAN} \textit{Structure needs cleaning}.] 
% \item[\errcode{ENAVAIL} \textit{No XENIX semaphores available}.] 
% \item[\errcode{EISNAM} \textit{Is a named type file}.] 
% \item[\errcode{EREMOTEIO} \textit{Remote I/O error}.] 
% \item[\errcode{ENOMEDIUM} \textit{No medium found}.] 
% \item[\errcode{EMEDIUMTYPE} \textit{Wrong medium type}.] 
% \end{description}


% LocalWords:  header errno perror int strerror sez EPERM Operation not ENOENT
% LocalWords:  permitted such pathname EIO error ENXIO device address kernel Is
% LocalWords:  ENOEXEC Invalid executable format exec EBADF Bad descriptor Too
% LocalWords:  EACCES Permission denied ELOOP many symbolic links encountered
% LocalWords:  ENAMETOOLONG name too long ENOTBLK Block required block EEXIST
% LocalWords:  exists EBUSY Resource busy filesystem EXDEV ENODEV ENOTDIR files
% LocalWords:  EISDIR EMFILE ulimit ENFILE table overflow ENOTTY ETXTBSY Text
% LocalWords:  EFBIG big ENOSPC left ESPIPE seek operation EROFS Read only read
% LocalWords:  system EMLINK EPIPE Broken SIGPIPE ENOTEMPTY empty EUSERS users
% LocalWords:  EDQUOT exceeded ESTALE NFS handle EREMOTE Object is ENOLCK locks
% LocalWords:  available locking EFTYPE type sticky ESRCH process matches the
% LocalWords:  specified pid Argument list POSIX ECHILD There child processes
% LocalWords:  socket ENOTSOCK EMSGSIZE Message EPROTOTYPE Protocol wrong for
% LocalWords:  ENOPROTOOPT EPROTONOSUPPORT supported ESOCKTNOSUPPORT EOPNOTSUPP
% LocalWords:  transport endpoint EPFNOSUPPORT family EAFNOSUPPORT protocol of
% LocalWords:  EADDRINUSE already bind EADDRNOTAVAIL Cannot assign requested to
% LocalWords:  ENETDOWN ENETUNREACH unreachable ENETRESET dropped connection
% LocalWords:  because reset l'host ECONNABORTED caused abort ECONNRESET peer
% LocalWords:  dell'host ENOBUFS ENOMEM EISCONN connected ENOTCONN UDP send now
% LocalWords:  EDESTADDRREQ Destination ESHUTDOWN after shutdown ETOOMANYREFS
% LocalWords:  references cannot splice glibc ETIMEDOUT timed ECONNREFUSED host
% LocalWords:  refused EHOSTDOWN EHOSTUNREACH route EINTR Interrupted function
% LocalWords:  call memory EDEADLK Deadlock avoided deadlock EFAULT segment IPC
% LocalWords:  violation SIGSEGV EINVAL argument EDOM Domain ERANGE underflow
% LocalWords:  EAGAIN temporarily unavailable EWOULDBLOCK select fork would has
% LocalWords:  EINPROGRESS progress connect EALREADY ENOSYS implemented ENOTSUP
% LocalWords:  EILSEQ Illegal sequence EBADMSG message EIDRM Identifier removed
% LocalWords:  SysV EMULTIHOP Multihop attempted ENODATA ENOLINK been severed
% LocalWords:  ENOMSG desired ENOSR streams resources ENOSTR stream EOVERFLOW
% LocalWords:  Value large defined STAT EPROTO ETIME Timer expired group wait
% LocalWords:  waitpid Specification cap USB did respond Stale


%%% Local Variables: 
%%% mode: latex
%%% TeX-master: "gapil"
%%% End: 

%% netlayer.tex
%%
%% Copyright (C) 2000-2018 Simone Piccardi.  Permission is granted to
%% copy, distribute and/or modify this document under the terms of the GNU Free
%% Documentation License, Version 1.1 or any later version published by the
%% Free Software Foundation; with the Invariant Sections being "Un preambolo",
%% with no Front-Cover Texts, and with no Back-Cover Texts.  A copy of the
%% license is included in the section entitled "GNU Free Documentation
%% License".
%% 

\chapter{Il livello di rete}
\label{cha:network_layer}

In questa appendice prenderemo in esame i vari protocolli disponibili a
livello di rete.\footnote{per la spiegazione della suddivisione in livelli dei
  protocolli di rete, si faccia riferimento a quanto illustrato in
  sez.~\ref{sec:net_protocols}.} Per ciascuno di essi forniremo una descrizione
generica delle principali caratteristiche, del formato di dati usato e quanto
possa essere necessario per capirne meglio il funzionamento dal punto di vista
della programmazione.

Data la loro prevalenza il capitolo sarà sostanzialmente incentrato sui due
protocolli principali esistenti su questo livello: il protocollo IP, sigla che
sta per \textit{Internet Protocol}, (ma che più propriamente si dovrebbe
chiamare IPv4) ed la nuova versione di questo stesso protocollo, denominata
IPv6. Tratteremo comunque anche il protocollo ICMP e la sua versione
modificata per IPv6 (cioè ICMPv6).


\section{Il protocollo IP}
\label{sec:ip_protocol}

L'attuale \textit{Internet Protocol} (IPv4) viene standardizzato nel 1981
dall'\href{http://www.ietf.org/rfc/rfc791.txt}{RFC~791}; esso nasce per
disaccoppiare le applicazioni della struttura hardware delle reti di
trasmissione, e creare una interfaccia di trasmissione dei dati indipendente
dal sottostante substrato di rete, che può essere realizzato con le tecnologie
più disparate (Ethernet, Token Ring, FDDI, ecc.).


\subsection{Introduzione}
\label{sec:IP_intro}

Il compito principale di IP è quello di trasmettere i pacchetti da un computer
all'altro della rete; le caratteristiche essenziali con cui questo viene
realizzato in IPv4 sono due:
\begin{itemize}
\item \textit{Universal addressing} la comunicazione avviene fra due host
  identificati univocamente con un indirizzo a 32 bit che può appartenere ad
  una sola interfaccia di rete.
\item \textit{Best effort} viene assicurato il massimo impegno nella
  trasmissione, ma non c'è nessuna garanzia per i livelli superiori né
  sulla percentuale di successo né sul tempo di consegna dei pacchetti di
  dati, né sull'ordine in cui vengono consegnati.
\end{itemize}

Per effettuare la comunicazione e l'instradamento dei pacchetti fra le varie
reti di cui è composta Internet IPv4 organizza gli indirizzi in una
gerarchia a due livelli, in cui una parte dei 32 bit dell'indirizzo indica il
numero di rete, e un'altra l'host al suo interno.  Il numero di rete serve
ai router per stabilire a quale rete il pacchetto deve essere inviato, il
numero di host indica la macchina di destinazione finale all'interno di detta
rete.

Per garantire l'unicità dell'indirizzo Internet esiste un'autorità centrale
(la IANA, \textit{Internet Assigned Number Authority}) che assegna i numeri di
rete alle organizzazioni che ne fanno richiesta; è poi compito di quest'ultime
assegnare i numeri dei singoli host all'interno della propria rete.

Per venire incontro alle richieste dei vari enti e organizzazioni che volevano
utilizzare questo protocollo di comunicazione, originariamente gli indirizzi
di rete erano stati suddivisi all'interno delle cosiddette \textit{classi},
(rappresentate in tab.~\ref{tab:IP_ipv4class}), in modo da consentire
dispiegamenti di reti di varie dimensioni a seconda delle diverse esigenze.

\begin{table}[htb]
  \centering
  \footnotesize
  \begin{usepicture} 
  \begin{tabular} {c@{\hspace{1mm}\vrule}
      p{3mm}@{\vrule}p{3mm}@{\vrule}p{3mm}@{\vrule}p{3mm}@{\vrule}
      p{3mm}@{\vrule}p{3mm}@{\vrule}p{3mm}@{\vrule}p{3mm}@{\vrule}
      p{3mm}@{\vrule}p{3mm}@{\vrule}p{3mm}@{\vrule}p{3mm}@{\vrule}
      p{3mm}@{\vrule}p{3mm}@{\vrule}p{3mm}@{\vrule}p{3mm}@{\vrule}
      p{3mm}@{\vrule}p{3mm}@{\vrule}p{3mm}@{\vrule}p{3mm}@{\vrule}
      p{3mm}@{\vrule}p{3mm}@{\vrule}p{3mm}@{\vrule}p{3mm}@{\vrule}
      p{3mm}@{\vrule}p{3mm}@{\vrule}p{3mm}@{\vrule}p{3mm}@{\vrule}
      p{3mm}@{\vrule}p{3mm}@{\vrule}p{3mm}@{\vrule}p{3mm}@{\vrule}}
    \omit&\omit& \multicolumn{7}{c}{7 bit}&\multicolumn{24}{c}{24 bit} \\
    \cline{2-33}
    \omit\hfill\vrule &&&&&&&& &&&&&&&& &&&&&&&& &&&&&&&& \\
    classe A &\centering 0&
    \multicolumn{7}{@{}c@{\vrule}}{\parbox[c]{21mm}{\centering net Id}} &
    \multicolumn{24}{@{}c@{\vrule}}{\parbox[c]{72mm}{\centering host Id}} \\
    \omit\hfill\vrule &&&&&&&& &&&&&&&& &&&&&&&& &&&&&&&& \\
    \cline{2-33}
    \multicolumn{33}{c}{ } \\
    \omit&\omit&\omit& 
    \multicolumn{14}{c}{14 bit}&\multicolumn{16}{c}{16 bit} \\
    \cline{2-33}
    \omit\hfill\vrule &&&&&&&& &&&&&&&& &&&&&&&& &&&&&&&& \\
    classe B&\centering 1&\centering 0& 
    \multicolumn{14}{@{}c@{\vrule}}{\parbox{42mm}{\centering net Id}} &
    \multicolumn{16}{@{}c@{\vrule}}{\parbox{48mm}{\centering host Id}} \\
    \omit\hfill\vrule &&&&&&&& &&&&&&&& &&&&&&&& &&&&&&&& \\
    \cline{2-33}
   
    \multicolumn{33}{c}{ } \\
    \omit&\omit&\omit& 
    \multicolumn{21}{c}{21 bit}&\multicolumn{8}{c}{8 bit} \\
    \cline{2-33}
    \omit\hfill\vrule &&&&&&&& &&&&&&&& &&&&&&&& &&&&&&&& \\
    classe C&\centering 1&\centering 1&\centering 0&
    \multicolumn{21}{@{}c@{\vrule}}{\parbox{63mm}{\centering net Id}} &
    \multicolumn{8}{@{}c@{\vrule}}{\parbox{24mm}{\centering host Id}} \\
    \omit\hfill\vrule &&&&&&&& &&&&&&&& &&&&&&&& &&&&&&&& \\
    \cline{2-33}


    \multicolumn{33}{c}{ } \\
    \omit&\omit&\omit&\omit& 
    \multicolumn{28}{c}{28 bit} \\
    \cline{2-33}
    \omit\hfill\vrule &&&&&&&& &&&&&&&& &&&&&&&& &&&&&&&& \\
    classe D&\centering 1&\centering 1&\centering 1&\centering 0&
    \multicolumn{28}{@{}c@{\vrule}}{\parbox{63mm}{\centering 
        multicast group Id}} \\
    \omit\hfill\vrule &&&&&&&& &&&&&&&& &&&&&&&& &&&&&&&& \\
    \cline{2-33}

    \multicolumn{33}{c}{ } \\
    \omit&\omit&\omit&\omit&\omit&
    \multicolumn{27}{c}{27 bit} \\
    \cline{2-33}
    \omit\hfill\vrule &&&&&&&& &&&&&&&& &&&&&&&& &&&&&&&& \\
    classe E&\centering 1&\centering 1&\centering 1&\centering 1&\centering 0&
    \multicolumn{27}{@{}c@{\vrule}}{\parbox{59mm}{\centering 
        reserved for future use}} \\
    \omit\hfill\vrule &&&&&&&& &&&&&&&& &&&&&&&& &&&&&&&& \\
    \cline{2-33}

  \end{tabular}
  \end{usepicture} 
\caption{Le classi di indirizzi secondo IPv4.}
\label{tab:IP_ipv4class}
\end{table}

Le classi di indirizzi usate per il dispiegamento delle reti su quella che
comunemente viene chiamata \textit{Internet} sono le prime tre; la classe D è
destinata al \textit{multicast} mentre la classe E è riservata per usi
sperimentali e non viene impiegata.

Come si può notare però la suddivisione riportata in
tab.~\ref{tab:IP_ipv4class} è largamente inefficiente in quanto se ad un
utente necessita anche solo un indirizzo in più dei 256 disponibili con una
classe A occorre passare a una classe B, che ne prevede 65536,\footnote{in
  realtà i valori esatti sarebbero 254 e 65534, una rete con a disposizione
  $N$ bit dell'indirizzo IP, ha disponibili per le singole macchine soltanto
  $@^N-2$ numeri, dato che uno deve essere utilizzato come indirizzo di rete e
  uno per l'indirizzo di \textit{broadcast}.} con un conseguente spreco di
numeri.

Inoltre, in particolare per le reti di classe C, la presenza di tanti
indirizzi di rete diversi comporta una crescita enorme delle tabelle di
instradamento che ciascun router dovrebbe tenere in memoria per sapere dove
inviare il pacchetto, con conseguente crescita dei tempi di elaborazione da
parte di questi ultimi ed inefficienza nel trasporto.

\begin{table}[htb]
  \centering
  \footnotesize
  \begin{usepicture} 
  \begin{tabular} {c@{\hspace{1mm}\vrule}
      p{3mm}@{\vrule}p{3mm}@{\vrule}p{3mm}@{\vrule}p{3mm}@{\vrule}
      p{3mm}@{\vrule}p{3mm}@{\vrule}p{3mm}@{\vrule}p{3mm}@{\vrule}
      p{3mm}@{\vrule}p{3mm}@{\vrule}p{3mm}@{\vrule}p{3mm}@{\vrule}
      p{3mm}@{\vrule}p{3mm}@{\vrule}p{3mm}@{\vrule}p{3mm}@{\vrule}
      p{3mm}@{\vrule}p{3mm}@{\vrule}p{3mm}@{\vrule}p{3mm}@{\vrule}
      p{3mm}@{\vrule}p{3mm}@{\vrule}p{3mm}@{\vrule}p{3mm}@{\vrule}
      p{3mm}@{\vrule}p{3mm}@{\vrule}p{3mm}@{\vrule}p{3mm}@{\vrule}
      p{3mm}@{\vrule}p{3mm}@{\vrule}p{3mm}@{\vrule}p{3mm}@{\vrule}}
    \omit&
    \multicolumn{12}{c}{$n$ bit}&\multicolumn{20}{c}{$32-n$ bit} \\
    \cline{2-33}
    \omit\hfill\vrule &&&&&&&& &&&&&&&& &&&&&&&& &&&&&&&& \\
    CIDR &
    \multicolumn{12}{@{}c@{\vrule}}{\parbox[c]{36mm}{\centering net Id}} &
    \multicolumn{20}{@{}c@{\vrule}}{\parbox[c]{60mm}{\centering host Id}} \\
    \omit\hfill\vrule &&&&&&&& &&&&&&&& &&&&&&&& &&&&&&&& \\
    \cline{2-33}
  \end{tabular}
  \end{usepicture} 
\caption{Uno esempio di indirizzamento CIDR.}
\label{tab:IP_ipv4cidr}
\end{table}

Per questo nel 1992 è stato introdotto un indirizzamento senza classi (il
CIDR, \textit{Classless Inter-Domain Routing}) in cui il limite fra i bit
destinati a indicare il numero di rete e quello destinati a indicare l'host
finale può essere piazzato in qualunque punto (vedi
tab.~\ref{tab:IP_ipv4cidr}), permettendo di accorpare più classi A su un'unica
rete o suddividere una classe B e diminuendo al contempo il numero di
indirizzi di rete da inserire nelle tabelle di instradamento dei router.


\subsection{L'intestazione di IP}
\label{sec:IP_header}

Come illustrato in fig.~\ref{fig:net_tcpip_data_flux} (si ricordi quanto detto
in sez.~\ref{sec:net_tcpip_overview} riguardo al funzionamento generale del
TCP/IP), per eseguire il suo compito il protocollo IP inserisce (come
praticamente ogni protocollo di rete) una opportuna intestazione in cima ai
dati che deve trasmettere, la cui schematizzazione è riportata in
fig.~\ref{fig:IP_ipv4_head}.

\begin{figure}[!htb]
  \centering
  \includegraphics[width=10cm]{img/ipv4_head}
  \caption{L'intestazione o \textit{header} di IPv4.}
  \label{fig:IP_ipv4_head}
\end{figure}

Ciascuno dei campi illustrati in fig.~\ref{fig:IP_ipv4_head} ha un suo preciso
scopo e significato, che si è riportato brevemente in
tab.~\ref{tab:IP_ipv4field}; si noti come l'intestazione riporti sempre due
indirizzi IP, quello \textsl{sorgente}, che indica l'IP da cui è partito il
pacchetto (cioè l'indirizzo assegnato alla macchina che lo spedisce) e quello
\textsl{destinazione} che indica l'indirizzo a cui deve essere inviato il
pacchetto (cioè l'indirizzo assegnato alla macchina che lo riceverà).

\begin{table}[!htb]
  \footnotesize
  \begin{center}
    \begin{tabular}{|l|c|p{10cm}|}
      \hline
      \textbf{Nome} & \textbf{Bit} & \textbf{Significato} \\
      \hline
      \hline
      \textit{version}       & 4& Numero di \textsl{versione}, nel caso 
                                  specifico vale sempre 4.\\
      \textit{head length}   & 4& Lunghezza dell'intestazione,
                                  in multipli di 32 bit.\\
      \textit{type of service}&8& Il ``\textsl{tipo di servizio}'', è suddiviso
                                  in: 3 bit di precedenza, che nelle attuali
                                  implementazioni del protocollo non vengono
                                  comunque utilizzati; un bit riservato che
                                  deve essere mantenuto a 0; 4 bit che
                                  identificano il tipo di servizio
                                  richiesto, uno solo dei quali può essere
                                  attivo.\\ 
      \textit{total length}  &16& La \textsl{lunghezza totale}, indica 
                                  la dimensione del carico di dati del
                                  pacchetto IP in byte.\\ 
      \textit{identification}&16& L'\textsl{identificazione}, assegnato alla
                                  creazione, è aumentato di uno all'origine
                                  della trasmissione di ciascun pacchetto, ma
                                  resta lo stesso per i pacchetti
                                  frammentati, consentendo così di
                                  identificare quelli che derivano dallo
                                  stesso pacchetto originario.\\
      \textit{flag}          & 3& I \textsl{flag} di controllo nell'ordine: il 
                                  primo è riservato e sempre nullo, il secondo
                                  indica se il pacchetto non può essere 
                                  frammentato, il terzo se ci sono ulteriori
                                  frammenti.\\  
      \textit{fragmentation offset}&13& L'\textsl{offset di frammento}, indica
                                  la posizione del frammento rispetto al
                                  pacchetto originale.\\
      \textit{time to live}  &16& Il \textsl{tempo di vita}, è decrementato di
                                  uno ogni volta che un router ritrasmette il
                                  pacchetto, se arriva a zero il pacchetto
                                  viene scartato.\\ 
      \textit{protocol}      & 8& Il \textsl{protocollo}, identifica il tipo di
                                  pacchetto che segue l'intestazione di IPv4.\\
      \textit{header checksum}&16&La \textsl{checksum di intestazione}, somma 
                                  di controllo per l'intestazione.\\ 
      \textit{source IP}     &32& L'\textsl{indirizzo di origine}.\\
      \textit{destination IP}&32& L'\textsl{indirizzo di destinazione}.\\
      \hline
    \end{tabular}
    \caption{Legenda per il significato dei campi dell'intestazione di IPv4}
    \label{tab:IP_ipv4field}
  \end{center}
\end{table}


Il campo TOS definisce il cosiddetto \textit{Type of Service}; questo permette
di definire il tipo di traffico contenuto nei pacchetti, e può essere
utilizzato dai router per dare diverse priorità in base al valore assunto da
questo campo. Abbiamo già visto come il valore di questo campo può essere
impostato sul singolo socket con l'opzione \const{IP\_TOS} (vedi
sez.~\ref{sec:sock_ipv4_options}), esso inoltre può essere manipolato sia dal
sistema del \textit{netfilter} di Linux con il comando \texttt{iptables} che
dal sistema del routing avanzato del comando \texttt{ip route} per consentire
un controllo più dettagliato dell'instradamento dei pacchetti e l'uso di
priorità e politiche di distribuzione degli stessi.

\begin{table}[!htb]
  \centering
  \footnotesize
  \begin{tabular}{|l|l|p{8cm}|}
    \hline
    \multicolumn{2}{|c|}{\textbf{Valore}}&\textbf{Significato}\\
    \hline
    \hline
    \constd{IPTOS\_LOWDELAY}    &\texttt{0x10}& Minimizza i ritardi
                                                per rendere più veloce
                                                possibile la ritrasmissione
                                                dei pacchetti (usato per
                                                traffico interattivo di
                                                controllo come SSH).\\
    \constd{IPTOS\_THROUGHPUT}  &\texttt{0x8} & Ottimizza la trasmissione
                                                per rendere il più elevato
                                                possibile il flusso netto di
                                                dati (usato su traffico dati,
                                                come quello di FTP).\\ 
    \constd{IPTOS\_RELIABILITY} &\texttt{0x4} & Ottimizza la trasmissione
                                                per ridurre al massimo le
                                                perdite di pacchetti (usato su
                                                traffico soggetto a rischio di
                                                perdita di pacchetti come TFTP
                                                o DHCP).\\
    \constd{IPTOS\_MINCOST}     &\texttt{0x2} & Indica i dati di riempimento,
                                                dove non interessa se si ha
                                                una bassa velocità di
                                                trasmissione, da utilizzare
                                                per i collegamenti con minor
                                                costo (usato per i protocolli
                                                di streaming).\\ 
    \textit{Normal-Service}&\texttt{0x0}      & Nessuna richiesta specifica.\\
    \hline

    \hline
  \end{tabular}
  \caption{Le costanti che definiscono alcuni valori standard per il campo TOS
    da usare come argomento \param{optval} per l'opzione \const{IP\_TOS}.} 
  \label{tab:IP_TOS_values}
\end{table}

I possibili valori del campo TOS, insieme al relativo significato ed alle
costanti numeriche ad esso associati, sono riportati in
tab.~\ref{tab:IP_TOS_values}. Per il valore nullo, usato di default per tutti
i pacchetti, e relativo al traffico normale, non esiste nessuna costante
associata. 

Il campo TTL, acromino di \textit{Time To Live}, viene utilizzato per
stabilire una sorta di tempo di vita massimo dei pacchetti sulla rete. In
realtà più che di un tempo, il campo serve a limitare il numero massimo di
salti (i cosiddetti \textit{hop}) che un pacchetto IP può compiere nel passare
da un router ad un altro nel suo attraversamento della rete verso la
destinazione.

Il protocollo IP prevede infatti che il valore di questo campo venga
decrementato di uno da ciascun router che ritrasmette il pacchetto verso la
sua destinazione, e che quando questo diventa nullo il router lo debba
scartare, inviando all'indirizzo sorgente un pacchetto ICMP di tipo
\textit{time-exceeded} con un codice \textit{ttl-zero-during-transit} se
questo avviene durante il transito sulla rete o
\textit{ttl-zero-during-reassembly} se questo avviene alla destinazione finale
(vedi sez.~\ref{sec:ICMP_protocol}).

In sostanza grazie all'uso di questo accorgimento un pacchetto non può
continuare a vagare indefinitamente sulla rete, e viene comunque scartato dopo
un certo tempo, o meglio, dopo che ha attraversato in certo numero di
router. Nel caso di Linux il valore iniziale utilizzato normalmente è 64 (vedi
sez.~\ref{sec:sock_ipv4_sysctl}).



\subsection{Le opzioni di IP}
\label{sec:IP_options}

Da fare ...


\section{Il protocollo IPv6}
\label{sec:ipv6_protocol}

Negli anni '90 con la crescita del numero di macchine connesse su Internet si
arrivò a temere l'esaurimento dello spazio degli indirizzi disponibili, specie
in vista di una prospettiva (per ora rivelatasi prematura) in cui ogni
apparecchio elettronico sarebbe stato inserito all'interno della rete. 

Per questo motivo si iniziò a progettare una nuova versione del protocollo 

L'attuale Internet Protocol (IPv4) viene standardizzato nel 1981
dall'\href{http://www.ietf.org/rfc/rfc719.txt}{RFC~719}; esso nasce per
disaccoppiare le applicazioni della struttura hardware delle reti di
trasmissione, e creare una interfaccia di trasmissione dei dati indipendente
dal sottostante substrato di rete, che può essere realizzato con le tecnologie
più disparate (Ethernet, Token Ring, FDDI, ecc.).


\subsection{I motivi della transizione}
\label{sec:IP_whyipv6}

Negli ultimi anni la crescita vertiginosa del numero di macchine connesse a
internet ha iniziato a far emergere i vari limiti di IPv4; in particolare si
è iniziata a delineare la possibilità di arrivare a una carenza di
indirizzi disponibili.

In realtà il problema non è propriamente legato al numero di indirizzi
disponibili; infatti con 32 bit si hanno $2^{32}$, cioè circa 4 miliardi,
numeri diversi possibili, che sono molti di più dei computer attualmente
esistenti.

Il punto è che la suddivisione di questi numeri nei due livelli rete/host e
l'utilizzo delle classi di indirizzamento mostrate in precedenza, ha
comportato che, nella sua evoluzione storica, il dispiegamento delle reti e
l'allocazione degli indirizzi siano stati inefficienti; neanche l'uso del CIDR
ha permesso di eliminare le inefficienze che si erano formate, dato che il
ridispiegamento degli indirizzi comporta cambiamenti complessi a tutti i
livelli e la riassegnazione di tutti gli indirizzi dei computer di ogni
sottorete.

Diventava perciò necessario progettare un nuovo protocollo che permettesse
di risolvere questi problemi, e garantisse flessibilità sufficiente per
poter continuare a funzionare a lungo termine; in particolare necessitava un
nuovo schema di indirizzamento che potesse rispondere alle seguenti
necessità:

\begin{itemize}
\item un maggior numero di numeri disponibili che consentisse di non restare
  più a corto di indirizzi
\item un'organizzazione gerarchica più flessibile dell'attuale 
\item uno schema di assegnazione degli indirizzi in grado di minimizzare le
  dimensioni delle tabelle di instradamento
\item uno spazio di indirizzi che consentisse un passaggio automatico dalle
  reti locali a internet
\end{itemize}


\subsection{Principali caratteristiche di IPv6}
\label{sec:IP_ipv6over}

Per rispondere alle esigenze descritte in sez.~\ref{sec:IP_whyipv6} IPv6 nasce
come evoluzione di IPv4, mantenendone inalterate le funzioni che si sono
dimostrate valide, eliminando quelle inutili e aggiungendone poche altre
ponendo al contempo una grande attenzione a mantenere il protocollo il più
snello e veloce possibile.

I cambiamenti apportati sono comunque notevoli e possono essere riassunti a
grandi linee nei seguenti punti:
\begin{itemize}
\item l'espansione delle capacità di indirizzamento e instradamento, per
  supportare una gerarchia con più livelli di indirizzamento, un numero di
  nodi indirizzabili molto maggiore e una auto-configurazione degli indirizzi
\item l'introduzione un nuovo tipo di indirizzamento, l'\textit{anycast} che
  si aggiungono agli usuali \textit{unicast} e \textit{multicast}
\item la semplificazione del formato dell'intestazione, eliminando o rendendo
  opzionali alcuni dei campi di IPv4, per eliminare la necessità di
  riprocessare la stessa da parte dei router e contenere l'aumento di
  dimensione dovuto ai nuovi indirizzi
\item un supporto per le opzioni migliorato, per garantire una trasmissione
  più efficiente del traffico normale, limiti meno stringenti sulle
  dimensioni delle opzioni, e la flessibilità necessaria per introdurne di
  nuove in futuro
\item il supporto per delle capacità di qualità di servizio (QoS) che
  permetta di identificare gruppi di dati per i quali si può provvedere un
  trattamento speciale (in vista dell'uso di internet per applicazioni
  multimediali e/o ``real-time'')
\end{itemize}


\subsection{L'intestazione di IPv6}
\label{sec:IP_ipv6head}

Per capire le caratteristiche di IPv6 partiamo dall'intestazione usata dal
protocollo per gestire la trasmissione dei pacchetti; in
fig.~\ref{fig:IP_ipv6head} è riportato il formato dell'intestazione di IPv6 da
confrontare con quella di IPv4 in fig.~\ref{fig:IP_ipv4_head}. La spiegazione
del significato dei vari campi delle due intestazioni è riportato
rispettivamente in tab.~\ref{tab:IP_ipv6field} e tab.~\ref{tab:IP_ipv4field})

% \begin{table}[htb]
%   \footnotesize
%   \begin{center}
%     \begin{tabular}{@{\vrule}p{16mm}@{\vrule}p{16mm}@{\vrule}p{16mm}
%         @{\vrule}p{16mm}@{\vrule}p{16mm}@{\vrule}p{16mm}
%         @{\vrule}p{16mm}@{\vrule}p{16mm}@{\vrule} }
%     \multicolumn{8}{@{}c@{}}{0\hfill 15 16\hfill 31}\\
%     \hline
%     \centering version&\centering priority& 
%     \multicolumn{6}{@{}p{96mm}@{\vrule}}{\centering flow label} \\
%     \hline
%     \multicolumn{4}{@{\vrule}p{64mm}@{\vrule}}{\centering payload length} & 
%     \multicolumn{2}{@{}p{32mm}@{\vrule}}{\centering next header} & 
%     \multicolumn{2}{@{}p{32mm}@{\vrule}}{\centering hop limit}\\
%     \hline
%     \multicolumn{8}{@{\vrule}c@{\vrule}}{} \\
%     \multicolumn{8}{@{\vrule}c@{\vrule}}{
%       source} \\
%     \multicolumn{8}{@{\vrule}c@{\vrule}}{
%       IP address} \\
%     \multicolumn{8}{@{\vrule}c@{\vrule}}{} \\
%     \hline
%     \multicolumn{8}{@{\vrule}c@{\vrule}}{} \\
%     \multicolumn{8}{@{\vrule}c@{\vrule}}{
%       destination} \\
%     \multicolumn{8}{@{\vrule}c@{\vrule}}{
%      IP address} \\
%     \multicolumn{8}{@{\vrule}c@{\vrule}}{} \\
%     \hline
%     \end{tabular}
%     \caption{L'intestazione o \textit{header} di IPv6}
%     \label{tab:IP_ipv6head}
%   \end{center}
% \end{table}

\begin{figure}[!htb]
  \centering
  \includegraphics[width=10cm]{img/ipv6_head}
  \caption{L'intestazione o \textit{header} di IPv6.}
  \label{fig:IP_ipv6head}
\end{figure}


Come si può notare l'intestazione di IPv6 diventa di dimensione fissa, pari a
40 byte, contro una dimensione (minima, in assenza di opzioni) di 20 byte per
IPv4; un semplice raddoppio nonostante lo spazio destinato agli indirizzi sia
quadruplicato, questo grazie a una notevole semplificazione che ha ridotto il
numero dei campi da 12 a 8.

\begin{table}[htb]
  \begin{center}
  \footnotesize
    \begin{tabular}{|l|c|p{8cm}|}
      \hline
      \textbf{Nome} & \textbf{Bit} & \textbf{Significato} \\
      \hline
      \hline
      \textit{version}       & 4& La \textsl{versione}, nel caso specifico vale
                                  sempre 6.\\ 
      \textit{priority}      & 4& La \textsl{priorità}, vedi
                                  sez.~\ref{sec:IPv6_prio}.\\
      \textit{flow label}    &24& L'\textsl{etichetta di flusso}, vedi
                                  sez.~\ref{sec:IP_ipv6_flow}.\\ 
      \textit{payload length}&16& La \textsl{lunghezza del carico}, cioè del
                                  corpo dei dati che segue l'intestazione, in
                                  byte. \\ 
      \textit{next header}   & 8& L'\textsl{intestazione successiva}, 
                                  identifica il tipo di pacchetto che segue
                                  l'intestazione di IPv6, ed usa gli stessi
                                  valori del campo protocollo
                                  nell'intestazione di IPv4.\\ 
      \textit{hop limit}     & 8& Il \textsl{limite di salti}, ha lo stesso 
                                  significato del \textit{time to live} 
                                  nell'intestazione di IPv4.\\ 
      \textit{source IP}     &128&L'\textsl{indirizzo di origine}.\\
      \textit{destination IP}&128&L'\textsl{indirizzo di destinazione}.\\
      \hline
    \end{tabular}
    \caption{Legenda per il significato dei campi dell'intestazione di IPv6}
    \label{tab:IP_ipv6field}
  \end{center}
\end{table}

Abbiamo già anticipato in sez.~\ref{sec:IP_ipv6over} uno dei criteri
principali nella progettazione di IPv6 è stato quello di ridurre al minimo il
tempo di elaborazione dei pacchetti da parte dei router, un confronto con
l'intestazione di IPv4 (vedi fig.~\ref{fig:IP_ipv4_head}) mostra le seguenti
differenze:

\begin{itemize}
\item è stato eliminato il campo \textit{header length} in quanto le opzioni
  sono state tolte dall'intestazione che ha così dimensione fissa; ci possono
  essere più intestazioni opzionali (\textsl{intestazioni di estensione}, vedi
  sez.~\ref{sec:IP_ipv6_extens}), ciascuna delle quali avrà un suo campo di
  lunghezza all'interno.
\item l'intestazione e gli indirizzi sono allineati a 64 bit, questo rende più
  veloce il processo da parte di computer con processori a 64 bit.
\item i campi per gestire la frammentazione (\textit{identification},
  \textit{flag} e \textit{fragment offset}) sono stati eliminati; questo
  perché la frammentazione è un'eccezione che non deve rallentare
  l'elaborazione dei pacchetti nel caso normale.
\item è stato eliminato il campo \textit{checksum} in quanto tutti i
  protocolli di livello superiore (TCP, UDP e ICMPv6) hanno un campo di
  checksum che include, oltre alla loro intestazione e ai dati, pure i campi
  \textit{payload length}, \textit{next header}, e gli indirizzi di origine e
  di destinazione; una checksum esiste anche per la gran parte protocolli di
  livello inferiore (anche se quelli che non lo hanno, come SLIP, non possono
  essere usati con grande affidabilità); con questa scelta si è ridotto di
  molto il tempo di elaborazione dato che i router non hanno più la necessità
  di ricalcolare la checksum ad ogni passaggio di un pacchetto per il
  cambiamento del campo \textit{hop limit}.
\item è stato eliminato il campo \textit{type of service}, che praticamente
  non è mai stato utilizzato; una parte delle funzionalità ad esso delegate
  sono state reimplementate (vedi il campo \textit{priority} al prossimo
  punto) con altri metodi.
\item è stato introdotto un nuovo campo \textit{flow label}, che viene usato,
  insieme al campo \textit{priority} (che recupera i bit di precedenza del
  campo \textit{type of service}) per implementare la gestione di una
  ``\textsl{qualità di servizio}'' (vedi sez.~\ref{sec:IP_ipv6_qos}) che
  permette di identificare i pacchetti appartenenti a un ``\textsl{flusso}''
  di dati per i quali si può provvedere un trattamento speciale.
\end{itemize}

Oltre alle differenze precedenti, relative ai singoli campi nell'intestazione,
ulteriori caratteristiche che diversificano il comportamento di IPv4 da
quello di IPv6 sono le seguenti:

\begin{itemize}
\item il \textit{broadcasting} non è previsto in IPv6, le applicazioni che lo
  usano dovono essere reimplementate usando il \textit{multicasting} (vedi
  sez.~\ref{sec:IP_ipv6_multicast}), che da opzionale diventa obbligatorio.
\item è stato introdotto un nuovo tipo di indirizzi, gli \textit{anycast}.
\item i router non possono più frammentare i pacchetti lungo il cammino, la
  frammentazione di pacchetti troppo grandi potrà essere gestita solo ai
  capi della comunicazione (usando un'apposita estensione vedi
  sez.~\ref{sec:IP_ipv6_extens}).
\item IPv6 richiede il supporto per il \textit{path MTU discovery} (cioè il
  protocollo per la selezione della massima lunghezza del pacchetto); seppure
  questo sia in teoria opzionale, senza di esso non sarà possibile inviare
  pacchetti più larghi della dimensione minima (576 byte).
\end{itemize}

\subsection{Gli indirizzi di IPv6}
\label{sec:IP_ipv6_addr}

Come già abbondantemente anticipato la principale novità di IPv6 è
costituita dall'ampliamento dello spazio degli indirizzi, che consente di avere
indirizzi disponibili in un numero dell'ordine di quello degli atomi che
costituiscono la terra. 

In realtà l'allocazione di questi indirizzi deve tenere conto della
necessità di costruire delle gerarchie che consentano un instradamento
rapido ed efficiente dei pacchetti, e flessibilità nel dispiegamento delle
reti, il che comporta una riduzione drastica dei numeri utilizzabili; uno
studio sull'efficienza dei vari sistemi di allocazione usati in altre
architetture (come i sistemi telefonici) è comunque giunto alla conclusione
che anche nella peggiore delle ipotesi IPv6 dovrebbe essere in grado di
fornire più di un migliaio di indirizzi per ogni metro quadro della
superficie terrestre.


\subsection{La notazione}
\label{sec:IP_ipv6_notation}
Con un numero di bit quadruplicato non è più possibile usare la notazione
coi numeri decimali di IPv4 per rappresentare un numero IP. Per questo gli
indirizzi di IPv6 sono in genere scritti come sequenze di otto numeri
esadecimali di 4 cifre (cioè a gruppi di 16 bit) usando i due punti come
separatore; cioè qualcosa del tipo
\texttt{1080:0000:0000:0008:0800:ba98:2078:e3e3}.

Visto che la notazione resta comunque piuttosto pesante esistono alcune
abbreviazioni: si può evitare di scrivere gli zeri iniziali delle singole
cifre, abbreviando l'indirizzo precedente in
\texttt{1080:0:0:8:800:ba98:2078:e3e3}; se poi un intero è zero si può
omettere del tutto, così come un insieme di zeri (ma questo solo una volta per
non generare ambiguità) per cui il precedente indirizzo si può scrivere anche
come \texttt{1080::8:800:ba98:2078:e3e3}.

Infine per scrivere un indirizzo IPv4 all'interno di un indirizzo IPv6 si
può usare la vecchia notazione con i punti, per esempio
\texttt{::192.84.145.138}.

\begin{table}[htb]
  \centering 
  \footnotesize
  \begin{tabular}{|l|l|l|}
    \hline
    \centering \textbf{Tipo di indirizzo}
    & \centering \textbf{Prefisso} & {\centering \textbf{Frazione}} \\
    \hline
    \hline
    riservato & \texttt{0000 0000} & 1/256 \\
    non assegnato  & \texttt{0000 0001} & 1/256 \\
    \hline
    riservato per NSAP & \texttt{0000 001} & 1/128\\
    riservato per IPX & \texttt{0000 010} & 1/128\\
    \hline
    non assegnato  & \texttt{0000 011} & 1/128 \\
    non assegnato  & \texttt{0000 1} & 1/32 \\
    non assegnato  & \texttt{0001} & 1/16 \\
    \hline
    provider-based & \texttt{001} & 1/8\\
    \hline
    non assegnato  & \texttt{010} & 1/8 \\
    non assegnato  & \texttt{011} & 1/8 \\
    geografic-based& \texttt{100} & 1/8 \\
    non assegnato  & \texttt{101} & 1/8 \\
    non assegnato  & \texttt{110} & 1/8 \\
    non assegnato  & \texttt{1110} & 1/16 \\
    non assegnato  & \texttt{1111 0} & 1/32 \\
    non assegnato  & \texttt{1111 10} & 1/64 \\
    non assegnato  & \texttt{1111 110} & 1/128 \\
    non assegnato  & \texttt{1111 1110 0} & 1/512 \\
    \hline
    unicast link-local & \texttt{1111 1110 10} & 1/1024 \\
    unicast site-local & \texttt{1111 1110 11} & 1/1024 \\
    \hline
    \hline
    \textit{multicast} & \texttt{1111 1111} & 1/256 \\
    \hline
  \end{tabular}
  \caption{Classificazione degli indirizzi IPv6 a seconda dei bit più 
    significativi}
  \label{tab:IP_ipv6addr}
\end{table}


\subsection{La architettura degli indirizzi di IPv6}
\label{sec:IP_ipv6_addr_arch}

Come per IPv4 gli indirizzi sono identificatori per una singola (indirizzi
\textit{unicast}) o per un insieme (indirizzi \textit{multicast} e
\textit{anycast}) di interfacce di rete.

Gli indirizzi sono sempre assegnati all'interfaccia, non al nodo che la
ospita; dato che ogni interfaccia appartiene ad un nodo quest'ultimo può
essere identificato attraverso uno qualunque degli indirizzi unicast delle sue
interfacce. A una interfaccia possono essere associati anche più indirizzi.

IPv6 presenta tre tipi diversi di indirizzi: due di questi, gli indirizzi
\textit{unicast} e \textit{multicast} hanno le stesse caratteristiche che in
IPv4, un terzo tipo, gli indirizzi \textit{anycast} è completamente nuovo.  In
IPv6 non esistono più gli indirizzi \textit{broadcast}, la funzione di questi
ultimi deve essere reimplementata con gli indirizzi \textit{multicast}.

Gli indirizzi \textit{unicast} identificano una singola interfaccia: i
pacchetti mandati ad un tale indirizzo verranno inviati a quella interfaccia,
gli indirizzi \textit{anycast} identificano un gruppo di interfacce tale che
un pacchetto mandato a uno di questi indirizzi viene inviato alla più vicina
(nel senso di distanza di routing) delle interfacce del gruppo, gli indirizzi
\textit{multicast} identificano un gruppo di interfacce tale che un pacchetto
mandato a uno di questi indirizzi viene inviato a tutte le interfacce del
gruppo.

In IPv6 non ci sono più le classi ma i bit più significativi indicano il tipo
di indirizzo; in tab.~\ref{tab:IP_ipv6addr} sono riportati i valori di detti
bit e il tipo di indirizzo che loro corrispondente.  I bit più significativi
costituiscono quello che viene chiamato il \textit{format prefix} ed è sulla
base di questo che i vari tipi di indirizzi vengono identificati.  Come si
vede questa architettura di allocazione supporta l'allocazione di indirizzi
per i provider, per uso locale e per il \textit{multicast}; inoltre è stato
riservato lo spazio per indirizzi NSAP, IPX e per le connessioni; gran parte
dello spazio (più del 70\%) è riservato per usi futuri.

Si noti infine che gli indirizzi \textit{anycast} non sono riportati in
tab.~\ref{tab:IP_ipv6addr} in quanto allocati al di fuori dello spazio di
allocazione degli indirizzi unicast.

\subsection{Indirizzi unicast \textit{provider-based}}
\label{sec:IP_ipv6_unicast}

Gli indirizzi \textit{provider-based} sono gli indirizzi usati per le
comunicazioni globali, questi sono definiti
nell'\href{http://www.ietf.org/rfc/rfc2073.txt}{RFC~2073} e sono gli
equivalenti degli attuali indirizzi delle classi da A a C.

L'autorità che presiede all'allocazione di questi indirizzi è la IANA; per
evitare i problemi di crescita delle tabelle di instradamento e una procedura
efficiente di allocazione la struttura di questi indirizzi è organizzata fin
dall'inizio in maniera gerarchica; pertanto lo spazio di questi indirizzi è
stato suddiviso in una serie di campi secondo lo schema riportato in
tab.~\ref{tab:IP_ipv6_unicast}.

\begin{table}[htb]
  \centering
  \footnotesize
  \begin{tabular} {@{\vrule}p{6mm}
      @{\vrule}p{16mm}@{\vrule}p{24mm}
      @{\vrule}p{30mm}@{\vrule}c@{\vrule}}
    \multicolumn{1}{@{}c@{}}{3}&\multicolumn{1}{c}{5 bit}&
    \multicolumn{1}{c}{$n$ bit}&\multicolumn{1}{c}{$56-n$ bit}&
    \multicolumn{1}{c}{64 bit} \\
    \hline
    \omit\vrule\hfill\vrule&\hspace{16mm} & & &\omit\hspace{76mm}\hfill\vrule\\
    \centering 010&
    \centering \textsl{Registry Id}&
    \centering \textsl{Provider Id}& 
    \centering \textsl{Subscriber Id}& 
    \textsl{Intra-Subscriber} \\
    \omit\vrule\hfill\vrule& & & &\omit\hspace{6mm}\hfill\vrule\\ 
    \hline
  \end{tabular}
\caption{Formato di un indirizzo unicast \textit{provider-based}.}
\label{tab:IP_ipv6_unicast}
\end{table}

Al livello più alto la IANA può delegare l'allocazione a delle autorità
regionali (i Regional Register) assegnando ad esse dei blocchi di indirizzi; a
queste autorità regionali è assegnato un Registry Id che deve seguire
immediatamente il prefisso di formato. Al momento sono definite tre registri
regionali (INTERNIC, RIPE NCC e APNIC), inoltre la IANA si è riservata la
possibilità di allocare indirizzi su base regionale; pertanto sono previsti
i seguenti possibili valori per il \textsl{Registry Id};
gli altri valori restano riservati per la IANA.
\begin{table}[htb]
  \centering 
  \footnotesize
    \begin{tabular}{|l|l|l|}
      \hline
      \textbf{Regione} & \textbf{Registro} & \textbf{Id} \\
      \hline
      \hline
      Nord America &INTERNIC & \texttt{11000} \\
      Europa & RIPE NCC & \texttt{01000} \\
      Asia & APNIC & \texttt{00100} \\
      Multi-regionale & IANA &\texttt{10000} \\
      \hline
    \end{tabular}
    \caption{Valori dell'identificativo dei 
      Regional Register allocati ad oggi.}
    \label{tab:IP_ipv6_regid}
\end{table}

L'organizzazione degli indirizzi prevede poi che i due livelli successivi, di
suddivisione fra \textit{Provider Id}, che identifica i grandi fornitori di
servizi, e \textit{Subscriber Id}, che identifica i fruitori, sia gestita dai
singoli registri regionali. Questi ultimi dovranno definire come dividere lo
spazio di indirizzi assegnato a questi due campi (che ammonta a un totale di
56~bit), definendo lo spazio da assegnare al \textit{Provider Id} e
al \textit{Subscriber Id}, ad essi spetterà inoltre anche l'allocazione dei
numeri di \textit{Provider Id} ai singoli fornitori, ai quali sarà delegata
l'autorità di allocare i \textit{Subscriber Id} al loro interno.

L'ultimo livello è quello \textit{Intra-subscriber} che è lasciato alla
gestione dei singoli fruitori finali, gli indirizzi \textit{provider-based}
lasciano normalmente gli ultimi 64~bit a disposizione per questo livello, la
modalità più immediata è quella di usare uno schema del tipo mostrato in
tab.~\ref{tab:IP_ipv6_uninterf} dove l'\textit{Interface Id} è dato dal
MAC-address a 48~bit dello standard Ethernet, scritto in genere nell'hardware
delle scheda di rete, e si usano i restanti 16~bit per indicare la sottorete.

\begin{table}[htb]
  \centering
  \footnotesize
  \begin{tabular} {@{\vrule}p{64mm}@{\vrule}p{16mm}@{\vrule}c@{\vrule}}
    \multicolumn{1}{c}{64 bit}&\multicolumn{1}{c}{16 bit}&
    \multicolumn{1}{c}{48 bit}\\
    \hline
    \omit\vrule\hfill\vrule&\hspace{16mm}&\omit\hspace{48mm}\hfill\vrule\\ 
    \centering \textsl{Subscriber Prefix}& 
    \centering \textsl{Subnet Id}&
    \textsl{Interface Id}\\
    \omit\vrule\hfill\vrule& &\omit\hspace{6mm}\hfill\vrule\\ 
    \hline
  \end{tabular}
\caption{Formato del campo \textit{Intra-subscriber} per un indirizzo unicast
  \textit{provider-based}.}
\label{tab:IP_ipv6_uninterf}
\end{table}

Qualora si dovesse avere a che fare con una necessità di un numero più
elevato di sotto-reti, il precedente schema andrebbe modificato, per evitare
l'enorme spreco dovuto all'uso dei MAC-address, a questo scopo si possono
usare le capacità di auto-configurazione di IPv6 per assegnare indirizzi
generici con ulteriori gerarchie per sfruttare efficacemente tutto lo spazio
di indirizzi.

Un registro regionale può introdurre un ulteriore livello nella gerarchia
degli indirizzi, allocando dei blocchi per i quali delegare l'autorità a dei
registri nazionali, quest'ultimi poi avranno il compito di gestire la
attribuzione degli indirizzi per i fornitori di servizi nell'ambito del/i
paese coperto dal registro nazionale con le modalità viste in precedenza.
Una tale ripartizione andrà effettuata all'interno dei soliti 56~bit come
mostrato in tab.~\ref{tab:IP_ipv6_uninaz}.

\begin{table}[htb]
  \centering
  \footnotesize
  \begin{tabular} {@{\vrule}p{3mm}
      @{\vrule}p{10mm}@{\vrule}p{12mm}@{\vrule}p{18mm}
      @{\vrule}p{18mm}@{\vrule}c@{\vrule}}
    \multicolumn{1}{@{}c@{}}{3}&\multicolumn{1}{c}{5 bit}&
    \multicolumn{1}{c}{n bit}&\multicolumn{1}{c}{m bit}&
    \multicolumn{1}{c}{56-n-m bit}&\multicolumn{1}{c}{64 bit} \\
    \hline
    \omit\vrule\hfill\vrule& & & & &\omit\hspace{64mm}\hfill\vrule\\
    \centering \texttt{3}&
    \centering \textsl{Reg.}&
    \centering \textsl{Naz.}&
    \centering \textsl{Prov.}& 
    \centering \textsl{Subscr.}& 
    \textsl{Intra-Subscriber} \\
    \omit\vrule\hfill\vrule &&&&&\\ 
    \hline
  \end{tabular}
\caption{Formato di un indirizzo unicast \textit{provider-based} che prevede
      un registro nazionale.}
\label{tab:IP_ipv6_uninaz}
\end{table}


\subsection{Indirizzi ad uso locale}
\label{sec:IP_ipv6_linksite}

Gli indirizzi ad uso locale sono indirizzi unicast che sono instradabili solo
localmente (all'interno di un sito o di una sottorete), e possono avere una
unicità locale o globale.

Questi indirizzi sono pensati per l'uso all'interno di un sito per mettere su
una comunicazione locale immediata, o durante le fasi di auto-configurazione
prima di avere un indirizzo globale.

\begin{table}[htb]
  \centering
  \footnotesize
  \begin{tabular} {@{\vrule}p{10mm}@{\vrule}p{54mm}@{\vrule}c@{\vrule}}
    \multicolumn{1}{c}{10} &\multicolumn{1}{c}{54 bit} & 
    \multicolumn{1}{c}{64 bit} \\
    \hline
    \omit\vrule\hfill\vrule & & \omit\hspace{64mm}\hfill\vrule\\
    \centering \texttt{FE80}& 
    \centering\texttt{0000 .   .   .   .   . 0000} &
    Interface Id \\
    \omit\vrule\hfill\vrule & &\\
    \hline
\end{tabular}
\caption{Formato di un indirizzo \textit{link-local}.}
\label{tab:IP_ipv6_linklocal}
\end{table}

Ci sono due tipi di indirizzi, \textit{link-local} e \textit{site-local}. Il
primo è usato per un singolo link; la struttura è mostrata in
tab.~\ref{tab:IP_ipv6_linklocal}, questi indirizzi iniziano sempre con un
valore nell'intervallo \texttt{FE80}--\texttt{FEBF} e vengono in genere usati
per la configurazione automatica dell'indirizzo al bootstrap e per la ricerca
dei vicini (vedi \ref{sec:IP_ipv6_autoconf}); un pacchetto che abbia tale
indirizzo come sorgente o destinazione non deve venire ritrasmesso dai router,
sono gli indirizzi che identificano la macchina sulla rete locale, per questo
sono chiamati in questo modo, in quanto sono usati solo su di essa.

Un indirizzo \textit{site-local} invece è usato per l'indirizzamento
all'interno di un sito che non necessita di un prefisso globale; la struttura
è mostrata in tab.~\ref{tab:IP_ipv6_sitelocal}, questi indirizzi iniziano
sempre con un valore nell'intervallo \texttt{FEC0}--\texttt{FEFF} e non devono
venire ritrasmessi dai router all'esterno del sito stesso; sono in sostanza
gli equivalenti degli indirizzi riservati per reti private definiti su IPv4.
Per entrambi gli indirizzi il campo \textit{Interface Id} è un identificatore
che deve essere unico nel dominio in cui viene usato, un modo immediato per
costruirlo è quello di usare il MAC-address delle schede di rete.
 
\begin{table}[!h]
  \centering
  \footnotesize
  \begin{tabular} {@{\vrule}p{10mm}@{\vrule}p{38mm}@{\vrule}p{16mm}
      @{\vrule}c@{\vrule}}
    \multicolumn{1}{c}{10} &\multicolumn{1}{c}{38 bit} & 
    \multicolumn{1}{c}{16 bit} &\multicolumn{1}{c}{64 bit} \\
    \hline
    \omit\vrule\hfill\vrule& & & \omit\hspace{64mm}\hfill\vrule\\
    \centering \texttt{FEC0}& 
    \centering \texttt{0000 .   .   . 0000}& 
    \centering Subnet Id &
    Interface Id\\
    \omit\vrule\hfill\vrule& & &\\
    \hline
\end{tabular}
\caption{Formato di un indirizzo \textit{site-local}.}
\label{tab:IP_ipv6_sitelocal}
\end{table}

Gli indirizzi di uso locale consentono ad una organizzazione che non è
(ancora) connessa ad Internet di operare senza richiedere un prefisso globale,
una volta che in seguito l'organizzazione venisse connessa a Internet
potrebbe continuare a usare la stessa suddivisione effettuata con gli
indirizzi \textit{site-local} utilizzando un prefisso globale e la
rinumerazione degli indirizzi delle singole macchine sarebbe automatica.

\subsection{Indirizzi riservati}
\label{sec:IP_ipv6_reserved}

Alcuni indirizzi sono riservati per scopi speciali, in particolare per scopi
di compatibilità.

Un primo tipo sono gli indirizzi \textit{IPv4 mappati su IPv6} (mostrati in
tab.~\ref{tab:IP_ipv6_map}), questo sono indirizzi unicast che vengono usati
per consentire ad applicazioni IPv6 di comunicare con host capaci solo di
IPv4; questi sono ad esempio gli indirizzi generati da un DNS quando l'host
richiesto supporta solo IPv4; l'uso di un tale indirizzo in un socket IPv6
comporta la generazione di un pacchetto IPv4 (ovviamente occorre che sia IPv4
che IPv6 siano supportati sull'host di origine).

\begin{table}[!htb]
  \centering
  \footnotesize
  \begin{tabular} {@{\vrule}p{80mm}@{\vrule}p{16mm}@{\vrule}c@{\vrule}}
    \multicolumn{1}{c}{80 bit} &\multicolumn{1}{c}{16 bit} & 
    \multicolumn{1}{c}{32 bit} \\
    \hline
    \omit\vrule\hfill\vrule& &\omit\hspace{32mm}\hfill\vrule\\ 
    \centering
    \texttt{0000 .   .   .   .   .   .   .   .   .   .   .   . 0000} & 
    \centering\texttt{FFFF} &
    IPv4 address \\
    \omit\vrule\hfill\vrule& &\\ 
    \hline
\end{tabular}
\caption{Formato di un indirizzo IPV4 mappato su IPv6.}
\label{tab:IP_ipv6_map}
\end{table}

Un secondo tipo di indirizzi di compatibilità sono gli \textsl{IPv4
  compatibili IPv6} (vedi tab.~\ref{tab:IP_ipv6_comp}) usati nella transizione
da IPv4 a IPv6: quando un nodo che supporta sia IPv6 che IPv4 non ha un router
IPv6 deve usare nel DNS un indirizzo di questo tipo, ogni pacchetto IPv6
inviato a un tale indirizzo verrà automaticamente incapsulato in IPv4.

\begin{table}[htb]
  \centering
  \footnotesize
  \begin{tabular} {@{\vrule}p{80mm}@{\vrule}p{16mm}@{\vrule}p{32mm}@{\vrule}}
    \multicolumn{1}{c}{80 bit} &\multicolumn{1}{c}{16 bit} & 
    \multicolumn{1}{c}{32 bit} \\
    \hline
    \omit\vrule\hfill\vrule& &\omit\hspace{32mm}\hfill\vrule\\ 
    \centering
    \texttt{0000 .   .   .   .   .   .   .   .   .   .   .   . 0000} & 
    \centering\texttt{0000} &
    \parbox{32mm}{\centering IPv4 address} \\
    \omit\vrule\hfill\vrule& &\\ 
    \hline
\end{tabular}
\caption{Formato di un indirizzo IPV4 mappato su IPv6.}
\label{tab:IP_ipv6_comp}
\end{table}

Altri indirizzi speciali sono il \textit{loopback address}, costituito da 127
zeri ed un uno (cioè \texttt{::1}) e l'\textsl{indirizzo generico}
costituito da tutti zeri (scritto come \texttt{0::0} o ancora più
semplicemente come \texttt{:}) usato in genere quando si vuole indicare
l'accettazione di una connessione da qualunque host.

\subsection{Multicasting}
\label{sec:IP_ipv6_multicast}

\itindbeg{multicast}

Gli indirizzi \textit{multicast} sono usati per inviare un pacchetto a un
gruppo di interfacce; l'indirizzo identifica uno specifico gruppo di
\textit{multicast} e il pacchetto viene inviato a tutte le interfacce di detto
gruppo.  Un'interfaccia può appartenere ad un numero qualunque numero di
gruppi di \textit{multicast}. Il formato degli indirizzi \textit{multicast} è
riportato in tab.~\ref{tab:IP_ipv6_multicast}:

\begin{table}[htb]
  \centering
  \footnotesize
  \begin{tabular} {@{\vrule}p{12mm}
      @{\vrule}p{6mm}@{\vrule}p{6mm}@{\vrule}c@{\vrule}}
    \multicolumn{1}{c}{8}&\multicolumn{1}{c}{4}&
    \multicolumn{1}{c}{4}&\multicolumn{1}{c}{112 bit} \\
    \hline
    \omit\vrule\hfill\vrule& & & \omit\hspace{104mm}\hfill\vrule\\
    \centering\texttt{FF}& 
    \centering flag &
    \centering scop& 
    Group Id\\
    \omit\vrule\hfill\vrule &&&\\ 
    \hline
  \end{tabular}
\caption{Formato di un indirizzo \textit{multicast}.}
\label{tab:IP_ipv6_multicast}
\end{table}

Il prefisso di formato per tutti gli indirizzi \textit{multicast} è
\texttt{FF}, ad esso seguono i due campi il cui significato è il seguente:

\begin{itemize}
\item \textsl{flag}: un insieme di 4 bit, di cui i primi tre sono riservati e
  posti a zero, l'ultimo è zero se l'indirizzo è permanente (cioè un
  indirizzo noto, assegnato dalla IANA), ed è uno se invece l'indirizzo è
  transitorio.
\item \textsl{scop} è un numero di quattro bit che indica il raggio di
  validità dell'indirizzo, i valori assegnati per ora sono riportati in
  tab.~\ref{tab:IP_ipv6_multiscope}.
\end{itemize}



\begin{table}[!htb]
  \centering 
  \footnotesize
  \begin{tabular}[c]{|c|l|c|l|}
    \hline
    \multicolumn{4}{|c|}{\bf Gruppi di multicast} \\
    \hline
    \hline
    0 & riservato & 8 & organizzazione locale \\
    1 & nodo locale & 9 & non assegnato \\
    2 & collegamento locale & A & non assegnato \\
    3 & non assegnato & B & non assegnato \\
    4 & non assegnato & C & non assegnato \\ 
    5 & sito locale & D & non assegnato \\
    6 & non assegnato & E & globale \\
    7 & non assegnato & F & riservato \\
    \hline
  \end{tabular}
\caption{Possibili valori del campo \textsl{scop} di un indirizzo
  \textit{multicast}.} 
\label{tab:IP_ipv6_multiscope}
\end{table}

Infine l'ultimo campo identifica il gruppo di \textit{multicast}, sia
permanente che transitorio, all'interno del raggio di validità del medesimo.
Alcuni indirizzi \textit{multicast}, riportati in tab.~\ref{tab:multiadd} sono
già riservati per il funzionamento della rete.

\begin{table}[!htb]
  \centering 
  \footnotesize
  \begin{tabular}[c]{|l|l|l|}
    \hline
    \textbf{Uso}& \textbf{Indirizzi riservati} & \textbf{Definizione}\\
    \hline 
    \hline 
    all-nodes       & \texttt{FFxx:0:0:0:0:0:0:1}  & 
                      \href{http://www.ietf.org/rfc/rfc1970.txt}{RFC~1970} \\
    all-routers     & \texttt{FFxx:0:0:0:0:0:0:2}  & 
                      \href{http://www.ietf.org/rfc/rfc1970.txt}{RFC~1970} \\
    all-rip-routers & \texttt{FFxx:0:0:0:0:0:0:9}  & 
                      \href{http://www.ietf.org/rfc/rfc2080.txt}{RFC~2080} \\
    all-cbt-routers & \texttt{FFxx:0:0:0:0:0:0:10} & \\
    reserved        & \texttt{FFxx:0:0:0:0:0:1:0}  & IANA \\
    link-name       & \texttt{FFxx:0:0:0:0:0:1:1}  &  \\
    all-dhcp-agents & \texttt{FFxx:0:0:0:0:0:1:2}  & \\
    all-dhcp-servers& \texttt{FFxx:0:0:0:0:0:1:3}  & \\
    all-dhcp-relays & \texttt{FFxx:0:0:0:0:0:1:4}  & \\
    solicited-nodes & \texttt{FFxx:0:0:0:0:1:0:0}  & 
                      \href{http://www.ietf.org/rfc/rfc1970.txt}{RFC~1970} \\
    \hline
  \end{tabular}
\caption{Gruppi di \textit{multicast} predefiniti.}
\label{tab:multiadd}
\end{table}

L'utilizzo del campo di \textit{scope} e di questi indirizzi predefiniti serve
a recuperare le funzionalità del \textit{broadcasting} (ad esempio inviando un
pacchetto all'indirizzo \texttt{FF02:0:0:0:0:0:0:1} si raggiungono tutti i
nodi locali).

\itindend{multicast}

\subsection{Indirizzi \textit{anycast}}
\label{sec:IP_anycast}

Gli indirizzi \textit{anycast} sono indirizzi che vengono assegnati ad un
gruppo di interfacce: un pacchetto indirizzato a questo tipo di indirizzo
viene inviato al componente del gruppo più ``\textsl{vicino}'' secondo la
distanza di instradamento calcolata dai router.

Questi indirizzi sono allocati nello stesso spazio degli indirizzi unicast,
usando uno dei formati disponibili, e per questo, sono da essi assolutamente
indistinguibili. Quando un indirizzo unicast viene assegnato a più interfacce
(trasformandolo in un anycast) il computer su cui è l'interfaccia deve essere
configurato per tener conto del fatto.

Gli indirizzi anycast consentono a un nodo sorgente di inviare pacchetti a una
destinazione su un gruppo di possibili interfacce selezionate. La sorgente non
deve curarsi di come scegliere l'interfaccia più vicina, compito che tocca al
sistema di instradamento (in sostanza la sorgente non ha nessun controllo
sulla selezione).

Gli indirizzi anycast, quando vengono usati come parte di una sequenza di
instradamento, consentono ad esempio ad un nodo di scegliere quale fornitore
vuole usare (configurando gli indirizzi anycast per identificare i router di
uno stesso provider).

Questi indirizzi pertanto possono essere usati come indirizzi intermedi in una
intestazione di instradamento o per identificare insiemi di router connessi a
una particolare sottorete, o che forniscono l'accesso a un certo sotto
dominio.

L'idea alla base degli indirizzi anycast è perciò quella di utilizzarli per
poter raggiungere il fornitore di servizio più vicino; ma restano aperte tutta
una serie di problematiche, visto che una connessione con uno di questi
indirizzi non è possibile, dato che per una variazione delle distanze di
routing non è detto che due pacchetti successivi finiscano alla stessa
interfaccia.

La materia è pertanto ancora controversa e in via di definizione.


\subsection{Le estensioni}
\label{sec:IP_ipv6_extens}

Come già detto in precedenza IPv6 ha completamente cambiato il trattamento
delle opzioni; queste ultime infatti sono state tolte dall'intestazione del
pacchetto, e poste in apposite \textsl{intestazioni di estensione} (o
\textit{extension header}) poste fra l'intestazione di IPv6 e l'intestazione
del protocollo di trasporto.

Per aumentare la velocità di elaborazione, sia dei dati del livello seguente
che di ulteriori opzioni, ciascuna estensione deve avere una lunghezza
multipla di 8 byte per mantenere l'allineamento a 64~bit di tutti le
intestazioni seguenti.

Dato che la maggior parte di queste estensioni non sono esaminate dai router
durante l'instradamento e la trasmissione dei pacchetti, ma solo all'arrivo
alla destinazione finale, questa scelta ha consentito un miglioramento delle
prestazioni rispetto a IPv4 dove la presenza di un'opzione comportava l'esame
di tutte quante.

Un secondo miglioramento è che rispetto a IPv4 le opzioni possono essere di
lunghezza arbitraria e non limitate a 40 byte; questo, insieme al modo in cui
vengono trattate, consente di utilizzarle per scopi come l'autenticazione e la
sicurezza, improponibili con IPv4.

Le estensioni definite al momento sono le seguenti:
\begin{itemize}
\item \textbf{Hop by hop} devono seguire immediatamente l'intestazione
  principale; indicano le opzioni che devono venire processate ad ogni
  passaggio da un router, fra di esse è da menzionare la \textit{jumbo
    payload} che segnala la presenza di un pacchetto di dati di dimensione
  superiore a 65535 byte.
\item \textbf{Destination options} opzioni che devono venire esaminate al nodo
  di ricevimento, nessuna di esse è tuttora definita.
\item \textbf{Routing} definisce una \textit{source route} (come la analoga
  opzione di IPv4) cioè una lista di indirizzi IP di nodi per i quali il
  pacchetto deve passare. 
\item \textbf{Fragmentation} viene generato automaticamente quando un host
  vuole frammentare un pacchetto, ed è riprocessato automaticamente alla
  destinazione che riassembla i frammenti.
\item \textbf{Authentication} gestisce l'autenticazione e il controllo di
  integrità dei pacchetti; è documentato
  dall'\href{http://www.ietf.org/rfc/rfc1826.txt}{RFC~1826}.
\item \textbf{Encapsulation} serve a gestire la segretezza del contenuto
  trasmesso; è documentato
  dall'\href{http://www.ietf.org/rfc/rfc1827.txt}{RFC~1827}.
\end{itemize}

La presenza di opzioni è rilevata dal valore del campo \textit{next header}
che indica qual è l'intestazione successiva a quella di IPv6; in assenza di
opzioni questa sarà l'intestazione di un protocollo di trasporto del livello
superiore, per cui il campo assumerà lo stesso valore del campo
\textit{protocol} di IPv4, altrimenti assumerà il valore dell'opzione
presente; i valori possibili sono riportati in
tab.~\ref{tab:IP_ipv6_nexthead}.

\begin{table}[htb]
  \begin{center}
    \footnotesize
    \begin{tabular}{|c|l|l|}
      \hline
      \textbf{Valore} & \textbf{Keyword} & \textbf{Tipo di protocollo} \\
      \hline
      \hline
      0  &      & Riservato.\\
         & HBH  & Hop by Hop.\\
      1  & ICMP & Internet Control Message (IPv4 o IPv6).\\
      2  & IGMP & Internet Group Management (IPv4).\\
      3  & GGP  & Gateway-to-Gateway.\\
      4  & IP   & IP in IP (IPv4 encapsulation).\\
      5  & ST   & Stream.\\
      6  & TCP  & Trasmission Control.\\
      17 & UDP  & User Datagram.\\
      43 & RH   & Routing Header (IPv6).\\
      44 & FH   & Fragment Header (IPv6).\\
      45 & IDRP & Inter Domain Routing.\\
      51 & AH   & Authentication Header (IPv6).\\
      52 & ESP  & Encrypted Security Payload (IPv6).\\
      59 & Null & No next header (IPv6).\\
      88 & IGRP & Internet Group Routing.\\
      89 & OSPF & Open Short Path First.\\
      255&      & Riservato.\\
    \hline
    \end{tabular}
    \caption{Tipi di protocolli e intestazioni di estensione}
    \label{tab:IP_ipv6_nexthead}
  \end{center}
\end{table}

Questo meccanismo permette la presenza di più opzioni in successione prima
del pacchetto del protocollo di trasporto; l'ordine raccomandato per le
estensioni è quello riportato nell'elenco precedente con la sola differenza
che le opzioni di destinazione sono inserite nella posizione ivi indicata solo
se, come per il tunnelling, devono essere esaminate dai router, quelle che
devono essere esaminate solo alla destinazione finale vanno in coda.


\subsection{Qualità di servizio}
\label{sec:IP_ipv6_qos}

Una delle caratteristiche innovative di IPv6 è quella di avere introdotto un
supporto per la qualità di servizio che è importante per applicazioni come
quelle multimediali o ``real-time'' che richiedono un qualche grado di
controllo sulla stabilità della banda di trasmissione, sui ritardi o la
dispersione dei temporale del flusso dei pacchetti.


\subsection{Etichette di flusso}
\label{sec:IP_ipv6_flow}
L'introduzione del campo \textit{flow label} può essere usata dall'origine
della comunicazione per etichettare quei pacchetti per i quali si vuole un
trattamento speciale da parte dei router come un una garanzia di banda minima
assicurata o un tempo minimo di instradamento/trasmissione garantito.

Questo aspetto di IPv6 è ancora sperimentale per cui i router che non
supportino queste funzioni devono porre a zero il \textit{flow label} per i
pacchetti da loro originanti e lasciare invariato il campo per quelli in
transito.

Un flusso è una sequenza di pacchetti da una particolare origine a una
particolare destinazione per il quale l'origine desidera un trattamento
speciale da parte dei router che lo manipolano; la natura di questo
trattamento può essere comunicata ai router in vari modi (come un protocollo
di controllo o con opzioni del tipo \textit{hop-by-hop}). 

Ci possono essere più flussi attivi fra un'origine e una destinazione, come
del traffico non assegnato a nessun flusso, un flusso viene identificato
univocamente dagli indirizzi di origine e destinazione e da una etichetta di
flusso diversa da zero, il traffico normale deve avere l'etichetta di flusso
posta a zero.

L'etichetta di flusso è assegnata dal nodo di origine, i valori devono
essere scelti in maniera (pseudo)casuale nel range fra 1 e FFFFFF in modo da
rendere utilizzabile un qualunque sottoinsieme dei bit come chiavi di hash per
i router.

\subsection{Priorità}
\label{sec:IPv6_prio}

Il campo di priorità consente di indicare il livello di priorità dei
pacchetti relativamente agli altri pacchetti provenienti dalla stessa
sorgente. I valori sono divisi in due intervalli, i valori da 0 a 7 sono usati
per specificare la priorità del traffico per il quale la sorgente provvede
un controllo di congestione cioè per il traffico che può essere ``tirato
indietro'' in caso di congestione come quello di TCP, i valori da 8 a 15 sono
usati per i pacchetti che non hanno questa caratteristica, come i pacchetti
``real-time'' inviati a ritmo costante.

Per il traffico con controllo di congestione sono raccomandati i seguenti
valori di priorità a seconda del tipo di applicazione:

\begin{table}[htb]
  \centering
  \footnotesize
  \begin{tabular}{|c|l|}
    \hline
    \textbf{Valore} & \textbf{Tipo di traffico} \\
    \hline
    \hline
    0 & Traffico generico.\\
    1 & Traffico di riempimento (es. news).\\
    2 & Trasferimento dati non interattivo (es. e-mail).\\
    3 & Riservato.\\
    4 & Trasferimento dati interattivo (es. FTP, HTTP, NFS).\\
    5 & Riservato.\\
    \hline
\end{tabular}
\caption{Formato di un indirizzo \textit{site-local}.}
\label{tab:priority}
\end{table}

Per il traffico senza controllo di congestione la priorità più bassa
dovrebbe essere usata per quei pacchetti che si preferisce siano scartati
più facilmente in caso di congestione.


\subsection{Sicurezza a livello IP}
\label{sec:security}

La attuale implementazione di Internet presenta numerosi problemi di
sicurezza, in particolare i dati presenti nelle intestazioni dei vari
protocolli sono assunti essere corretti, il che da adito alla possibilità di
varie tipologie di attacco forgiando pacchetti false, inoltre tutti questi
dati passano in chiaro sulla rete e sono esposti all'osservazione di chiunque
si trovi in mezzo.

Con IPv4 non è possibile realizzare un meccanismo di autenticazione e
riservatezza a un livello inferiore al primo (quello di applicazione), con
IPv6 è stato progettata la possibilità di intervenire al livello di rete (il
terzo) prevedendo due apposite estensioni che possono essere usate per fornire
livelli di sicurezza a seconda degli utenti. La codifica generale di questa
architettura è riportata
nell'\href{http://www.ietf.org/rfc/rfc2401.txt}{RFC~2401}.

Il meccanismo in sostanza si basa su due estensioni:
\begin{itemize}
\item una intestazione di sicurezza (\textit{authentication header}) che
  garantisce al destinatario l'autenticità del pacchetto
\item un carico di sicurezza (\textit{Encrypted Security Payload}) che
  assicura che solo il legittimo ricevente può leggere il pacchetto.
\end{itemize}

Perché tutto questo funzioni le stazioni sorgente e destinazione devono
usare una stessa chiave crittografica e gli stessi algoritmi, l'insieme degli
accordi fra le due stazioni per concordare chiavi e algoritmi usati va sotto
il nome di associazione di sicurezza.

I pacchetti autenticati e crittografati portano un indice dei parametri di
sicurezza (SPI, \textit{Security Parameter Index}) che viene negoziato prima
di ogni comunicazione ed è definito dalla stazione sorgente. Nel caso di
\textit{multicast} dovrà essere lo stesso per tutte le stazioni del gruppo.

\subsection{Autenticazione}
\label{sec:auth} 

Il primo meccanismo di sicurezza è quello dell'intestazione di autenticazione
(\textit{authentication header}) che fornisce l'autenticazione e il controllo
di integrità (ma senza riservatezza) dei pacchetti IP.

L'intestazione di autenticazione ha il formato descritto in
fig.~\ref{fig:autent_estens}: il campo \textit{Next Header} indica
l'intestazione successiva, con gli stessi valori del campo omonimo
nell'intestazione principale di IPv6, il campo \textit{Length} indica la
lunghezza dell'intestazione di autenticazione in numero di parole a 32 bit, il
campo riservato deve essere posto a zero, seguono poi l'indice di sicurezza,
stabilito nella associazione di sicurezza, e un numero di sequenza che la
stazione sorgente deve incrementare di pacchetto in pacchetto.

Completano l'intestazione i dati di autenticazione che contengono un valore di
controllo di integrità (ICV, \textit{Integrity Check Value}), che deve essere
di dimensione pari a un multiplo intero di 32 bit e può contenere un padding
per allineare l'intestazione a 64 bit. Tutti gli algoritmi di autenticazione
devono provvedere questa capacità.

\begin{figure}[!htb]
  \centering
  \includegraphics[width=10cm]{img/ah_option}
    \caption{Formato dell'intestazione dell'estensione di autenticazione.}
    \label{fig:autent_estens}
\end{figure}

L'intestazione di autenticazione può essere impiegata in due modi diverse
modalità: modalità trasporto e modalità tunnel.

La modalità trasporto è utilizzabile solo per comunicazioni fra stazioni
singole che supportino l'autenticazione. In questo caso l'intestazione di
autenticazione è inserita dopo tutte le altre intestazioni di estensione
eccezion fatta per la \textit{Destination Option} che può comparire sia
prima che dopo. 

\begin{figure}[!htb]
  \centering
  \includegraphics[width=10cm]{img/IP_AH_packet}
  \caption{Formato di un pacchetto IPv6 che usa l'opzione di autenticazione.}
  \label{fig:AH_autent_head}
\end{figure}

La modalità tunnel può essere utilizzata sia per comunicazioni fra stazioni
singole che con un gateway di sicurezza; in questa modalità ...


L'intestazione di autenticazione è una intestazione di estensione inserita
dopo l'intestazione principale e prima del carico dei dati. La sua presenza
non ha perciò alcuna influenza sui livelli superiori dei protocolli di
trasmissione come il TCP.


La procedura di autenticazione cerca di garantire l'autenticità del pacchetto
nella massima estensione possibile, ma dato che alcuni campi dell'intestazione
di IP possono variare in maniera impredicibile alla sorgente, il loro valore
non può essere protetto dall'autenticazione.

Il calcolo dei dati di autenticazione viene effettuato alla sorgente su una
versione speciale del pacchetto in cui il numero di salti nell'intestazione
principale è impostato a zero, così come le opzioni che possono essere
modificate nella trasmissione, e l'intestazione di routing (se usata) è posta
ai valori che deve avere all'arrivo.

L'estensione è indipendente dall'algoritmo particolare, e il protocollo è
ancora in fase di definizione; attualmente è stato suggerito l'uso di una
modifica dell'MD5 chiamata \textit{keyed MD5} che combina alla codifica anche
una chiave che viene inserita all'inizio e alla fine degli altri campi.


\subsection{Riservatezza}
\label{sec:ecry}

Per garantire una trasmissione riservata dei dati, è stata previsto la
possibilità di trasmettere pacchetti con i dati criptati: il cosiddetto ESP,
\textit{Encripted Security Payload}. Questo viene realizzato usando con una
apposita opzione che deve essere sempre l'ultima delle intestazioni di
estensione; ad essa segue il carico del pacchetto che viene criptato.

Un pacchetto crittografato pertanto viene ad avere una struttura del tipo di
quella mostrata in fig.~\ref{fig:ESP_criptopack}, tutti i campi sono in chiaro
fino al vettore di inizializzazione, il resto è crittografato.

\begin{figure}[!htb]
  \centering
  \includegraphics[width=10cm]{img/esp_option}
  \caption{Schema di pacchetto crittografato.}
  \label{fig:ESP_criptopack}
\end{figure}



\subsection{Auto-configurazione}
\label{sec:IP_ipv6_autoconf}

Una delle caratteristiche salienti di IPv6 è quella dell'auto-configurazione,
il protocollo infatti fornisce la possibilità ad un nodo di scoprire
automaticamente il suo indirizzo acquisendo i parametri necessari per potersi
connettere a internet.

L'auto-configurazione sfrutta gli indirizzi \textit{link-local}; qualora sul nodo sia
presente una scheda di rete che supporta lo standard IEEE802 (ethernet) questo
garantisce la presenza di un indirizzo fisico a 48 bit unico; pertanto il nodo
può assumere automaticamente senza pericoli di collisione l'indirizzo
\textit{link-local} \texttt{FE80::xxxx:xxxx:xxxx} dove \texttt{xxxx:xxxx:xxxx} è
l'indirizzo hardware della scheda di rete. 

Nel caso in cui non sia presente una scheda che supporta lo standard IEEE802
allora il nodo assumerà ugualmente un indirizzo \textit{link-local} della forma
precedente, ma il valore di \texttt{xxxx:xxxx:xxxx} sarà generato
casualmente; in questo caso la probabilità di collisione è di 1 su 300
milioni. In ogni caso per prevenire questo rischio il nodo invierà un
messaggio ICMP \textit{Solicitation} all'indirizzo scelto attendendo un certo
lasso di tempo; in caso di risposta l'indirizzo è duplicato e il
procedimento dovrà essere ripetuto con un nuovo indirizzo (o interrotto
richiedendo assistenza).

Una volta ottenuto un indirizzo locale valido diventa possibile per il nodo
comunicare con la rete locale; sono pertanto previste due modalità di
auto-configurazione, descritte nelle seguenti sezioni. In ogni caso
l'indirizzo \textit{link-local} resta valido.

\subsection{Auto-configurazione stateless}
\label{sec:stateless}

Questa è la forma più semplice di auto-configurazione, possibile quando
l'indirizzo globale può essere ricavato dall'indirizzo \textit{link-local} cambiando
semplicemente il prefisso a quello assegnato dal provider per ottenere un
indirizzo globale.

La procedura di configurazione è la seguente: all'avvio tutti i nodi IPv6
iniziano si devono aggregare al gruppo di \textit{multicast}
\textit{all-nodes} programmando la propria interfaccia per ricevere i messaggi
dall'indirizzo \textit{multicast} \texttt{FF02::1} (vedi
sez.~\ref{sec:IP_ipv6_multicast}); a questo punto devono inviare un messaggio
ICMP \textit{Router solicitation} a tutti i router locali usando l'indirizzo
\textit{multicast} \texttt{FF02::2} usando come sorgente il proprio indirizzo
\textit{link-local}.

Il router risponderà con un messaggio ICMP \textit{Router Advertisement} che
fornisce il prefisso e la validità nel tempo del medesimo, questo tipo di
messaggio può essere trasmesso anche a intervalli regolari. Il messaggio
contiene anche l'informazione che autorizza un nodo a autocostruire
l'indirizzo, nel qual caso, se il prefisso unito all'indirizzo \textit{link-local} non
supera i 128 bit, la stazione ottiene automaticamente il suo indirizzo
globale.

\subsection{Auto-configurazione stateful}
\label{sec:stateful}

Benché estremamente semplice l'auto-configurazione stateless presenta alcuni
problemi; il primo è che l'uso degli indirizzi delle schede di rete è
molto inefficiente; nel caso in cui ci siano esigenze di creare una gerarchia
strutturata su parecchi livelli possono non restare 48~bit per l'indirizzo
della singola stazione; il secondo problema è di sicurezza, dato che basta
introdurre in una rete una stazione autoconfigurante per ottenere un accesso
legale.

Per questi motivi è previsto anche un protocollo stateful basato su un server
che offra una versione IPv6 del DHCP; un apposito gruppo di \textit{multicast}
\texttt{FF02::1:0} è stato riservato per questi server; in questo caso il nodo
interrogherà il server su questo indirizzo di \textit{multicast} con
l'indirizzo \textit{link-local} e riceverà un indirizzo unicast globale.


\section{Il protocollo ICMP}
\label{sec:ICMP_protocol}

Come già accennato nelle sezioni precedenti, l'\textit{Internet Control
  Message Protocol} è un protocollo di servizio fondamentale per il
funzionamento del livello di rete. Il protocollo ICMP viene trasportato
direttamente su IP, ma proprio per questa sua caratteristica di protocollo di
servizio è da considerarsi a tutti gli effetti appartenente al livello di
rete.

\subsection{L'intestazione di ICMP}
\label{sec:ICMP_header}

Il protocollo ICMP è estremamente semplice, ed il suo unico scopo è quello di
inviare messaggi di controllo; in fig.~\ref{fig:ICMP_header} si è riportata la
struttura dell'intestazione di un pacchetto ICMP generico. 

\begin{figure}[!htb]
  \centering \includegraphics[width=12cm]{img/icmp_head}
  \caption{L'intestazione del protocollo ICMP.}
  \label{fig:ICMP_header}
\end{figure}

Ciascun pacchetto ICMP è contraddistinto dal valore del primo campo, il tipo,
che indica appunto che tipo di messaggio di controllo viene veicolato dal
pacchetto in questione; i valori possibili per questo campo, insieme al
relativo significato, sono riportati in tab.~\ref{tab:ICMP_type}.

\begin{table}[!htb]
  \centering
  \footnotesize
  \begin{tabular}{|l|l|p{9.5cm}|}
    \hline
    \textbf{Valore}&\textbf{Tipo}&\textbf{Significato}\\
    \hline
    \hline
    \texttt{any} & -- & Seleziona tutti i possibili valori \\
    \hline
    \textit{echo-reply}             &0& Inviato in risposta ad un ICMP
                                        \textit{echo-request}.\\ 
    \textit{destination-unreachable}&3& Segnala una destinazione 
                                        irraggiungibile, viene
                                        inviato all'IP sorgente di un
                                        pacchetto quando un router realizza
                                        che questo non può essere inviato a
                                        destinazione.\\
    \textit{source-quench}          &4& Inviato in caso di congestione della
                                        rete per indicare all'IP sorgente di
                                        diminuire il traffico inviato.\\
    \textit{redirect}               &5& Inviato per segnalare un errore di
                                        routing, richiede che la macchina
                                        sorgente rediriga il traffico ad un
                                        altro router da esso specificato.\\
    \textit{echo-request}           &8& Richiede l'invio in risposta di un
                                        \textit{echo-reply}.\\
%    \textit{router-advertisement}   & & \\
%    \textit{router-solicitation}    & & \\
    \textit{time-exceeded}          &11& Inviato quando il TTL di un pacchetto
                                         viene azzerato.\\
    \textit{parameter-problem}      &12& Inviato da un router che rileva dei
                                         problemi con l'intestazione di un
                                         pacchetto.\\
    \textit{timestamp-request}      &13& Richiede l'invio in risposta di un
                                         \textit{timestamp-reply}.\\
    \textit{timestamp-reply}        &14& Inviato in risposta di un
                                         \textit{timestamp-request}.\\
    \textit{info-request}           &15& Richiede l'invio in risposta di un
                                         \textit{info-reply}.\\
    \textit{info-reply}             &16& Inviato in risposta di un
                                         \textit{info-request}.\\
    \textit{address-mask-request}   &17& Richiede l'invio in risposta di un
                                         \textit{address-mask-reply}.\\
    \textit{address-mask-reply}     &18& Inviato in risposta di un
                                         \textit{address-mask-request}.\\
    \hline
  \end{tabular}
  \caption{I valori del \textsl{tipo} per i pacchetti ICMP.}
\label{tab:ICMP_type}
\end{table}

Per alcuni tipi di messaggi ICMP, esiste un secondo campo, detto codice, che
specifica ulteriormente la natura del messaggio; i soli messaggi che
utilizzano un valore per questo campo sono quelli di tipo
\textit{destination-unreachable}, \textit{redirect}, \textit{time-exceeded} e
\textit{parameter-problem}. I possibili valori del codice relativi a ciascuno
di essi sono stati riportati nelle quattro sezioni in cui si è suddivisa
tab.~\ref{tab:ICMP_code}, rispettivamente nell'ordine con cui sono appena
elencati i tipi a cui essi fanno riferimento. 

\begin{table}[!htb]
  \centering
  \footnotesize
  \begin{tabular}{|l|l|}
    \hline
    \textbf{Valore}&\textbf{Codice}\\
    \hline
    \hline
    \textit{network-unreachable}      &0\\
    \textit{host-unreachable}         &1\\
    \textit{protocol-unreachable}     &2\\
    \textit{port-unreachable}         &3 \\
    \textit{fragmentation-needed}     &4\\
    \textit{source-route-failed}      &5\\
    \textit{network-unknown}          &6\\
    \textit{host-unknown}             &7\\
    \textit{host-isolated}            &8\\
    \textit{network-prohibited}       &9\\
    \textit{host-prohibited}          &10 \\
    \textit{TOS-network-unreachable}  &11 \\
    \textit{TOS-host-unreachable}     &12 \\
    \textit{communication-prohibited} &13 \\
    \textit{host-precedence-violation}&14 \\
    \textit{precedence-cutoff}        &15 \\
    \hline
    \textit{network-redirect}         &0  \\
    \textit{host-redirect}            &1  \\
    \textit{TOS-network-redirect}     &2  \\
    \textit{TOS-host-redirect}        &3  \\
    \hline
    \textit{ttl-zero-during-transit}  &0 \\
    \textit{ttl-zero-during-reassembly}&1 \\
    \hline
    \textit{ip-header-bad}            &0 \\
    \textit{required-option-missing}  &1 \\
    \hline
  \end{tabular}
  \caption{Valori del campo \textsl{codice} per il protocollo ICMP.}
\label{tab:ICMP_code}
\end{table}


% LocalWords:  sez Protocol IPv dall' RFC Ethernet Token FDDI Universal host of
% LocalWords:  addressing Best effort l'host router IANA Assigned Number tab to
% LocalWords:  Authority quest'ultime multicast group reserved for CIDR Domain
% LocalWords:  Classless Routing TOS Type Service IPTOS LOWDELAY THROUGHPUT QoS
% LocalWords:  RELIABILITY MINCOST optval anycast unicast fig header version FE
% LocalWords:  priority flow label payload length next hop limit live source FF
% LocalWords:  destination identification fragment checksum TCP UDP ICMPv type
% LocalWords:  service head total fragmentation protocol broadcast broadcasting
% LocalWords:  multicasting path MTU discovery NSAP IPX based geografic local
% LocalWords:  routing format prefix Registry Subscriber Intra Regional SSH
% LocalWords:  Register INTERNIC NCC APNIC subscriber Interface MAC address Reg
% LocalWords:  Subnet Naz Prov Subscr FEBF bootstrap FEC FEFF DNS socket FFFF
% LocalWords:  sull'host loopback scop all nodes routers rip cbt name dhcp HBH
% LocalWords:  agents servers relays solicited extension options route Keyword
% LocalWords:  Authentication Encapsulation ICMP Control Message GGP Gateway ST
% LocalWords:  encapsulation Stream Trasmission Datagram RH FH IDRP ESP Null ip
% LocalWords:  Encrypted Security IGRP OSPF Short First tunnelling FFFFFF hash
% LocalWords:  news FTP NFS authentication Parameter Index ICV Integrity Value
% LocalWords:  padding Option gateway dell'MD keyed Encripted IEEE ethernet any
% LocalWords:  Solicitation netfilter iptables TFTP streaming Normal IGMP
% LocalWords:  stateless solicitation Advertisement stateful Transfer Unit echo
% LocalWords:  l'autoconfigurazione reply request unreachable all'IP quench TTL
% LocalWords:  redirect exceeded parameter problem timestamp info mask port ttl
% LocalWords:  needed failed unknown isolated prohibited communication cutoff
% LocalWords:  precedence violation during reassembly bad required option
% LocalWords:  missing



%%% Local Variables: 
%%% mode: latex
%%% TeX-master: "gapil"
%%% End: 

%% tcpprot.tex
%%
%% Copyright (C) 2002-2018 Simone Piccardi.  Permission is granted to copy,
%% distribute and/or modify this document under the terms of the GNU Free
%% Documentation License, Version 1.1 or any later version published by the
%% Free Software Foundation; with the Invariant Sections being "Un preambolo",
%% with no Front-Cover Texts, and with no Back-Cover Texts.  A copy of the
%% license is included in the section entitled "GNU Free Documentation
%% License".
%%

\chapter{Il livello di trasporto}
\label{cha:transport_layer}

In questa appendice tratteremo i vari protocolli relativi al livello di
trasporto.\footnote{al solito per la definizione dei livelli si faccia
  riferimento alle spiegazioni fornite in sez.~\ref{sec:net_protocols}.} In
particolare gran parte del capitolo sarà dedicato al più importante di questi,
il TCP, che è pure il più complesso ed utilizzato su internet.


\section{Il protocollo TCP}
\label{sec:tcp_protocol}

In questa sezione prenderemo in esame i vari aspetti del protocollo TCP, il
protocollo più comunemente usato dalle applicazioni di rete.


\subsection{Gli stati del TCP}
\label{sec:TCP_states}

In sez.~\ref{sec:TCP_connession} abbiamo descritto in dettaglio le modalità con
cui il protocollo TCP avvia e conclude una connessione, ed abbiamo accennato
alla presenza dei vari stati del protocollo. In generale infatti il
funzionamento del protocollo segue una serie di regole, che possono essere
riassunte nel comportamento di una macchina a stati, il cui diagramma di
transizione è riportato in fig.~\ref{fig:TCP_state_diag}.

\begin{figure}[!htb]
  \centering \includegraphics[width=10cm]{img/tcp_state_diagram}  
  \caption{Il diagramma degli stati del TCP.}
  \label{fig:TCP_state_diag}
\end{figure}

Il protocollo prevede l'esistenza di 11 diversi stati per una connessione ed
un insieme di regole per le transizioni da uno stato all'altro basate sullo
stato corrente, sull'operazione effettuata dall'applicazione o sul tipo di
segmento ricevuto; i nomi degli stati mostrati in
fig.~\ref{fig:TCP_state_diag} sono gli stessi che vengono riportati del
comando \cmd{netstat} nel campo \textit{State}.

\begin{figure}[!htb]
  \centering \includegraphics[width=10cm]{img/tcp_head}  
  \caption{L'intestazione del protocollo TCP.}
  \label{fig:TCP_header}
\end{figure}


\itindbeg{Maximum~Segment~Size~(MSS)}
% TODO trattare la MSS
\itindend{Maximum~Segment~Size~(MSS)}

\itindbeg{advertised~window}
% TODO trattare la advertised window
\itindend{advertised~window}

\index{algoritmo~di~Nagle|(}
% TODO trattare l'algoritmo di Nagle
\index{algoritmo~di~Nagle|)}


\section{Il protocollo UDP}
\label{sec:udp_protocol}

In questa sezione prenderemo in esame i vari aspetti del protocollo UDP, che
dopo il TCP è il protocollo più usato dalle applicazioni di rete. 


\begin{figure}[!htb]
  \centering \includegraphics[width=10cm]{img/udp_head}  
  \caption{L'intestazione del protocollo UDP.}
  \label{fig:UDP_header}
\end{figure}



% LocalWords:  sez TCP fig netstat UDP Segment Size advertised window


%%% Local Variables: 
%%% mode: latex
%%% TeX-master: "gapil"
%%% End: 

%% build.tex
%%
%% Copyright (C) 1999-2018 Simone Piccardi.  Permission is granted to copy,
%% distribute and/or modify this document under the terms of the GNU Free
%% Documentation License, Version 1.1 or any later version published by the
%% Free Software Foundation; with the Invariant Sections being "Un preambolo",
%% with no Front-Cover Texts, and with no Back-Cover Texts.  A copy of the
%% license is included in the section entitled "GNU Free Documentation
%% License".
%%

\chapter{Gli strumenti di ausilio per la programmazione}
\label{cha:build_manage}

Tratteremo in questa appendice in maniera superficiale i principali strumenti
che vengono utilizzati per programmare in ambito Linux, ed in particolare gli
strumenti per la compilazione e la costruzione di programmi e librerie, e gli
strumenti di gestione dei sorgenti e di controllo di versione.

Questo materiale è ripreso da un vecchio articolo, ed al momento è molto
obsoleto.


\section{L'uso di \texttt{make} per l'automazione della compilazione}
\label{sec:build_make}

Il comando \texttt{make} serve per automatizzare il processo di costruzione di
un programma ed effettuare una compilazione intelligente di tutti i file
relativi ad un progetto software, ricompilando solo i file necessari ed
eseguendo automaticamente tutte le operazioni che possono essere necessarie
alla produzione del risultato finale.\footnote{in realtà \texttt{make} non si
  applica solo ai programmi, ma in generale alla automazione di processi di
  costruzione, ad esempio anche la creazione dei file di questa guida viene
  fatta con \texttt{make}.}


\subsection {Introduzione a \texttt{make}}
\label{sec:make_intro}

Con \texttt{make} si possono definire i simboli del preprocessore C che
consentono la compilazione condizionale dei programmi (anche in Fortran); è
pertanto possibile gestire la ricompilazione dei programmi con diverse
configurazioni con la modifica di un unico file.

La sintassi normale del comando (quella che si usa quasi sempre, per le
opzioni vedere la pagina di manuale) è semplicemente \texttt{make}. Questo
comando esegue le istruzioni contenute in un file standard (usualmente
\texttt{Makefile}, o \texttt{makefile} nella directory corrente).

Il formato normale dei comandi contenuti in un \texttt{Makefile} è:
\begin{verbatim}
bersaglio: dipendenza1 dipendenza2 ...
        regola1
        regola2
        ...
\end{verbatim}
dove lo spazio all'inizio deve essere un tabulatore (metterci degli spazi è un
errore comune, fortunatamente ben segnalato dalle ultime versioni del
programma), il bersaglio e le dipendenze nomi di file e le regole comandi di
shell.

Il concetto di base è che se uno dei file di dipendenza è più recente (nel
senso di tempo di ultima modifica) del file bersaglio quest'ultimo viene
ricostruito di nuovo usando le regole elencate nelle righe successive.

Il comando \texttt{make} ricostruisce di default il primo bersaglio che viene
trovato nella scansione del \texttt{Makefile}, se in un \texttt{Makefile} sono
contenuti più bersagli indipendenti, si può farne ricostruire un altro che non
sia il primo passandolo esplicitamente al comando come argomento, con qualcosa
del tipo di: \texttt{make altrobersaglio}.

Si tenga presente che le dipendenze stesse possono essere dichiarate come 
bersagli dipendenti da altri file; in questo modo è possibile creare una 
catena di ricostruzioni. 

In esempio comune di quello che si fa è mettere come primo bersaglio il
programma principale che si vuole usare, e come dipendenze tutte gli oggetti
delle funzioni subordinate che utilizza, con i quali deve essere collegato; a
loro volta questi oggetti sono bersagli che hanno come dipendenza i relativi
sorgenti. In questo modo il cambiamento di una delle funzioni subordinate
comporta solo la ricompilazione della medesima e del programma finale.


\subsection{Utilizzo di \texttt{make}}

Il comando \texttt{make} mette a disposizione una serie molto complesse di
opzioni e di regole standard predefinite e sottintese, che permettono una
gestione estremamente rapida e concisa di progetti anche molto complessi; per
questo piuttosto che fare una replica del manuale preferisco commentare un
esempio di \texttt{makefile}, quello usato per ricompilare i programmi di
analisi dei dati dei test su fascio del tracciatore di Pamela.

\small
\begin{verbatim}
#----------------------------------------------------------------------
#
# Makefile for a Linux System:
# use GNU FORTRAN compiler g77
# Makefile done for tracker test data
#
#----------------------------------------------------------------------
# Fortran flags
FC=g77
FFLAGS= -fvxt -fno-automatic -Wall -O6 -DPC # -DDEBUG
CC=gcc
CFLAGS= -Wall -O6
CFLADJ=-c #-DDEBUG
#
# FC            Fortran compiler for standard rules
# FFLAGS        Fortran flags for standard rules
# CC            C Compiler for standard rules
# CFLAGS        C compiler flags for standard rules 
LIBS=  -L/cern/pro/lib -lkernlib -lpacklib -lgraflib -lmathlib
OBJ=cnoise.o fit2.o pedsig.o loop.o badstrp.o cutcn.o readevnt.o \
erasepedvar.o readinit.o dumpval.o writeinit.o

riduzione: riduzione.F $(OBJ) commondef.f readfile.o
        $(FC) $(FFLAGS) -o riduzione riduzione.F readfile.o $(OBJ) $(LIBS) 

readfile.o: readfile.c
        $(CC) $(CFLAGS) -o readfile.o readfile.c

$(OBJ): commondef.f

.PHONY : clean
clean:
        rm -f *.o 
        rm -f *~
        rm -f riduzione
        rm -f *.rz
        rm -f output
\end{verbatim}
\normalsize

Anzitutto i commenti, ogni linea che inizia con un \texttt{\#} è un
commento e non viene presa in considerazione.

Con \texttt{make} possono essere definite delle variabili, da potersi riusare
a piacimento, per leggibilità si tende a definirle tutte maiuscole,
nell'esempio ne sono definite varie: \small
\begin{verbatim}
FC=g77
FFLAGS= -fvxt -fno-automatic -Wall -O6 -DPC # -DDEBUG
CC=gcc
CFLAGS= -Wall -O6
CFLADJ=-c #-DDEBUG
...
LIBS=  -L/cern/pro/lib -lkernlib -lpacklib -lgraflib -lmathlib
OBJ=cnoise.o fit2.o pedsig.o loop.o badstrp.o cutcn.o readevnt.o \
\end{verbatim}
\normalsize

La sintassi è \texttt{NOME=}, alcuni nomi però hanno un significato speciale
(nel caso \texttt{FC}, \texttt{FLAGS}, \texttt{CC}, \texttt{CFLAGS}) in quanto
sono usati da \texttt{make} nelle cosiddette \textsl{regole implicite} (su cui
torneremo dopo).

Nel caso specifico, vedi anche i commenti, abbiamo definito i comandi di 
compilazione da usare per il C e il Fortran, e i rispettivi flag, una variabile
che contiene il pathname e la lista delle librerie del CERN e una variabile con una
lista di file oggetto.

Per richiamare una variabile si usa la sintassi \texttt{\$(NOME)}, ad esempio
nel makefile abbiamo usato:
\begin{verbatim}
        $(FC) $(FFLAGS) -o riduzione riduzione.F readfile.o $(OBJ) $(LIBS) 
\end{verbatim}
e questo significa che la regola verrà trattata come se avessimo scritto 
esplicitamente i valori delle variabili.

Veniamo ora alla parte principale del makefile che esegue la costruzione del
programma:
\begin{verbatim}
riduzione: riduzione.F $(OBJ) commondef.f readfile.o
        $(FC) $(FFLAGS) -o riduzione riduzione.F readfile.o $(OBJ) $(LIBS) 
readfile.o: readfile.c
        $(CC) $(CFLAGS) -o readfile.o readfile.c
$(OBJ): commondef.f
\end{verbatim}

Il primo bersaglio del makefile, che definisce il bersaglio di default, è il
programma di riduzione dei dati; esso dipende dal suo sorgente da tutti gli
oggetti definiti dalla variabile \texttt{OBJ}, dal file di definizioni
\texttt{commondef.f} e dalla routine C \texttt{readfile.o}; si noti il
\texttt{.F} del sorgente, che significa che il file prima di essere compilato
viene fatto passare attraverso il preprocessore C (cosa che non avviene per i
\texttt{.f}) che permette di usare i comandi di compilazione condizionale del
preprocessore C con la relativa sintassi. Sotto segue il comando di
compilazione che sfrutta le variabili definite in precedenza per specificare
quale compilatore e opzioni usare e specifica di nuovo gli oggetti e le
librerie.

Il secondo bersaglio definisce le regole per la compilazione della routine
in C; essa dipende solo dal suo sorgente. Si noti che per la compilazione
vengono usate le variabili relative al compilatore C. Si noti anche che se
questa regola viene usata, allora lo sarà anche la precedente, dato che
\texttt{riduzione} dipende da \texttt{readfile.o}.

Il terzo bersaglio è apparentemente incomprensibile dato che vi compare
solo il riferimento alla variabile \texttt{OBJ} con una sola dipendenza e
nessuna regola, essa però mostra le possibilità (oltre che la
complessità) di \texttt{make} connesse alla presenza di quelle regole
implicite a cui avevamo accennato.

Anzitutto una peculiarità di \texttt{make} è che si possono anche usare più
bersagli per una stessa regola (nell'esempio quelli contenuti nella variabile
\texttt{OBJ} che viene espansa in una lista); in questo caso la regola di
costruzione sarà applicata a ciascuno che si potrà citare nella regola stessa
facendo riferimento con la variabile automatica: \texttt{\$@}.  L'esempio
usato per la nostra costruzione però sembra non avere neanche la regola di
costruzione.

Questa mancanza sia di regola che di dipendenze (ad esempio dai vari sorgenti) 
illustra le capacità di funzionamento automatico di \texttt{make}.
Infatti è facile immaginarsi che un oggetto dipenda da un sorgente, e che 
per ottenere l'oggetto si debba compilare quest'ultimo.

Il comando \texttt{make} sa tutto questo per cui quando un bersaglio è un
oggetto (cioè ha un nome tipo \texttt{qualcosa.o}) non è necessario
specificare il sorgente, ma il programma lo va a cercare nella directory
corrente (ma è possibile pure dirgli di cercarlo altrove, il caso è trattato
nel manuale).  Nel caso specifico allora si è messo come dipendenza solo il
file delle definizioni che viene incluso in ogni subroutine.

Inoltre come dicevamo in genere per costruire un oggetto si deve compilarne il
sorgente; \texttt{make} sa anche questo e sulla base dell'estensione del
sorgente trovato (che nel caso sarà un \texttt{qualcosa.f}) applica la regola
implicita.  In questo caso la regola è quella di chiamare il compilatore
fortran applicato al file oggetto e al relativo sorgente, questo viene fatto
usando la variabile \texttt{FC} che è una delle variabili standard usata dalle
regole implicite (come \texttt{CC} nel caso di file \texttt{.c}); per una
maggiore flessibilità poi la regola standard usa anche la variabile
\texttt{FFLAGS} per specificare, a scelta dell'utente che non ha che da
definirla, quali flag di compilazione usare (nella documentazione sono
riportate tutte le regole implicite e le relative variabili usate).

In questo modo è stato possibile usare una sola riga per indicare la serie di
dipendenze e relative compilazioni delle singole subroutine; inoltre con l'uso
della variabile \texttt{OBJ} l'aggiunta di una nuova eventuale routine
\texttt{nuova.f} comporta solo l'aggiunta di \texttt{nuova.o} alla definizione
di \texttt{OBJ}.


\section{Source Control Management}
\label{sec:build_scm}

Uno dei problemi più comuni che si hanno nella programmazione è quella di
poter disporre di un sistema che consenta di tenere conto del lavoro
effettuato, di tracciare l'evoluzione del codice, e, soprattutto nel caso di
progetti portati avanti da più persone, consentire un accesso opportunamente
coordinato fra i vari partecipanti alla base comune dei sorgenti dello
sviluppo.

I programmi che servono a questo scopo vanno sotto il nome comune di SCM
(\textit{Source Control Manager}), e ne esistono di diversi tipi con diverse
filosofie progettuali, in particolare nelle modalità con cui gestiscono
l'accesso alla base di codice comune da parte dei singoli programmatori che vi
accedono.

Fra questi uno dei più usati, nonostante la sua architettura sia considerata
superata, è Subversion, un sistema di archiviazione centralizzata del codice
che consente di tenere traccia di tutte le modifiche e di condividere un
archivio comune per progetti portati avanti da diverse persone.

\subsection{Introduzione a Subversion}
\label{sec:build_subversion}

Subversion è basato sul concetto di \textit{repository}, un archivio
centralizzato in cui vengono riposti e da cui vengono presi i sorgenti dei
programmi. L'archivio tiene traccia delle diverse versioni registrate; i
programmatori inviano le modifiche usando una copia locale che hanno nella
loro directory di lavoro.

Subversion può gestire più di un progetto all'interno di un singolo server,
ciascuno dei quali viene associato ad un \textit{repository} distinto, ma si
possono anche creare sotto-progetti suddividendo un \textit{repository} in
diverse directory; ma ciascun progetto avrà meccanismi di controllo (ad 
esempio quelli che consentono di inviare email all'inserimento di nuovo
codice) comuni.

Una delle caratteristiche che contraddistinguono Subversion dal suo
predecessore CVS è quella di essere gestibile in maniera molto flessibile
l'accesso al repository, che può avvenire sia in maniera diretta facendo
riferimento alla directory in cui questo è stato installato che via rete,
tramite diversi protocolli. L'accesso più comune è fatto direttamente via
HTTP, utilizzando opportune estensioni del protocollo DAV, ma è possibile
passare attraverso SSH o fornire un servizio di rete dedicato.\footnote{esiste
  all'uopo il programma \texttt{svnserve}, ma il suo uso è sconsigliato per le
  scarse prestazioni e le difficoltà riscontrate a gestire accessi di utenti
  diversi; la modalità di accesso preferita resta quella tramite le estensioni
  al protocollo DAV.}

In generale è comunque necessario preoccuparsi delle modalità di accesso al
codice soltanto in fase di primo accesso al \textit{repository}, che occorrerà
identificare o con il pathname alla directory dove questo si trova o con una
opportuna URL (con il comune accesso via web del tutto analoga a quella che si
usa in un browser), dopo di che detto indirizzo sarà salvato nella propria
copia locale dei dati ed il riferimento diventerà implicito.

Il programma prevede infatti che in ogni directory che si è ottenuta come
copia locale sia presente una directory \texttt{.svn} contenente tutti i dati
necessari al programma. Inoltre il programma usa la directory
\texttt{.subversion} nella home dell'utente per mantenere le configurazioni
generali del client e le eventuali informazioni di autenticazione. 

Tutte le operazioni di lavoro sul \textit{repository} vengono effettuate lato
client tramite il comando \texttt{svn} che vedremo in
sez.~\ref{sec:build_subversion} ma la creazione e la inizializzazione dello
stesso (così come la gestione lato server) devono essere fatte tramite il
comando \texttt{svnadmin} eseguito sulla macchina che lo ospita. In generale
infatti il comando \texttt{svn} richiede che si faccia riferimento ad un
\textit{repository} (al limite anche vuoto) esistente e questo deve essere
opportunamente creato.

Il comando \texttt{svnadmin} utilizza una sintassi che richiede sempre
l'ulteriore specificazione di un sotto-comando, seguito da eventuali altri
argomenti. L'inizializzazione di un \textit{repository} (che sarà creato
sempre vuoto) viene eseguita con il comando:
\begin{verbatim}
svnadmin create /path/to/repository
\end{verbatim}
dove \texttt{/path/to/repository} è la directory dove verranno creati e
mantenuti tutti i file, una volta creato il \textit{repository} si potrà
iniziare ad utilizzarlo ed inserirvi i contenuti con il comando \texttt{svn}. 

Non essendo questo un testo di amministrazione di sistema non tratteremo qui i
dettagli della configurazione del server per l'accesso via rete al
\textit{repository}, per i quali si rimanda alla documentazione del progetto
ed alla documentazione sistemistica scritta per Truelite
Srl.\footnote{rispettivamente disponibili su \url{svn.tigris.org} e
  \url{labs.truelite.it/truedoc}.}

\subsection{Utilizzo di \texttt{svn}}
\label{sec:build_svn_use}

Una volta che si abbia a disposizione un \textit{repository} si potrà creare
un nuovo progetto sottoposto a controllo di versione importando al suo interno
i dati disponibili. In genere è pratica comune suddividere il contenuto di un
repository in tre directory secondo il seguente schema:
\begin{basedescript}{\desclabelwidth{2.7cm}\desclabelstyle{\nextlinelabel}}
\item[\texttt{trunk}] contiene la versione corrente i sviluppo, su cui vengono
  effettuate normalmente le modifiche e gli aggiornamenti;
\item[\texttt{tags}] contiene le diverse versioni \textsl{fotografate} ad un
  certo istante del processo di sviluppo, ad esempio in occasione del rilascio
  di una versione stabile, così che sia possibile identificarle facilmente;
\item[\texttt{branches}] contiene \textsl{rami} alternativi di sviluppo, ad
  esempio quello delle correzioni eseguite ad una versione stabile, che
  vengono portati avanti in maniera indipendente dalla versione principale.
\end{basedescript}

Questa suddivisione consente di sfruttare la capacità di Subversion di creare
senza spesa copie diverse del proprio contenuto, pertanto in genere si pone il
proprio progetto di sviluppo sotto \texttt{trunk}, e si copia quest'ultima in
occasione delle varie versioni di rilascio in altrettante sottocartelle di
\texttt{tags} e qualora si voglia aprire un ramo alternativo di sviluppo
basterà copiarsi il punto di partenza del ramo sotto \texttt{branches} e
iniziare ad eseguire le modifiche su di esso.

Le operazioni di gestione di un progetto con Subversion vengono eseguite con
il comando \texttt{svn}, che analogamente al precedente \texttt{svnadmin}
utilizza una sintassi basata sulla specificazione degli opportuni
sotto-comandi. Si sono riportati quelli più importanti in
tab.~\ref{tab:svn_mainsubcommands}. 

\begin{table}[htb]
  \centering
  \footnotesize
  \begin{tabular}[c]{|l|l|p{8cm}|}
    \hline
     \multicolumn{2}{|c|}{\textbf{Sotto-comando}}& \textbf{Significato} \\
    \hline
    \hline
    \texttt{import}  &  --       & Importa i file della directory corrente sul
                                   \textit{repository}.\\ 
    \texttt{checkout}&\texttt{co}& Scarica una versione del progetto dal
                                   \textit{repository}.\\
    \texttt{commit}  &\texttt{ci}& Invia le modifiche effettuate localmente al 
                                   \textit{repository}.\\
    \texttt{add}     &  --       & Richiede l'aggiunta un file o una directory
                                   al \textit{repository}.\\
    \texttt{remove}  &\texttt{rm}& Richiede la rimozione un file o una 
                                   directory dal \textit{repository}.\\
    \texttt{copy}    &\texttt{cp}& Richiede la copia un file o una cartella del
                                   progetto (mantenendone la storia).\\
    \texttt{move}    &\texttt{mv}& Richiede lo spostamento un file o una 
                                   directory (equivalente ad un \texttt{cp}
                                   seguito da un \texttt{rm}).\\
    \texttt{update}  &  --       & Aggiorna la copia locale.\\
    \texttt{resolved}&  --       & Rimuove una situazione di conflitto presente
                                   su un file.\\ 
    \hline
  \end{tabular}
  \caption{Tabella riassuntiva dei principali sotto-comandi di \texttt{svn}.}
  \label{tab:svn_mainsubcommands}
\end{table}

In genere però è piuttosto raro iniziare un progetto totalmente da zero, è
molto più comune avere una qualche versione iniziale dei propri file
all'interno di una cartella. In questo caso il primo passo è quello di
eseguire una inizializzazione del \textit{repository} importando al suo
interno quanto già esistente.  Per far questo occorre eseguire il comando:
\begin{verbatim}
svn import [/pathname] URL
\end{verbatim}
questo può essere eseguito direttamente nella directory contenente la versione
iniziale dei propri sorgenti nel qual caso il comando richiede come ulteriore
argomento la directory o la URL con la quale indicare il \textit{repository}
da usare. Alternativamente si può passare come primo argomento il pathname
della directory da importare, seguito dall'indicazione della URL del
\textit{repository}.

Si tenga presente che l'operazione di importazione inserisce sul
\textit{repository} il contenuto completo della directory indicata, compresi
eventuali file nascosti e sotto-directory. È anche possibile eseguire
l'importazione di più directory da inserire in diverse sezioni del
\textit{repository}, ma un tal caso ciascuna importazione sarà vista con una
diversa \textit{release}. Ad ogni operazione di modifica del
\textit{repository} viene infatti assegnato un numero progressivo che consente
di identificarne la storia delle modifiche e riportarsi ad un dato punto della
stessa in ogni momento successivo.\footnote{a differenza di CVS Subversion non
  assegna un numero di versione progressivo distinto ad ogni file, ma un
  numero di \textit{release} progressivo ad ogni cambiamento globale del
  repository, pertanto non esiste il concetto di versione di un singolo file,
  quanto di stato di tutto il \textit{repository} ad un dato momento, è
  comunque possibile richiedere in maniera indipendente la versione di ogni
  singolo file a qualunque \textit{release} si desideri.}

Una volta eseguita l'importazione di una versione iniziale è d'uopo cancellare
la directory originale e ripartire dal progetto appena creato. L'operazione di
recuperare ex-novo di tutti i file che fanno parte di un progetto, chiamata
usualmente \textit{checkout}, viene eseguita con il
comando:\footnote{alternativamente si può usare l'abbreviazione \texttt{svn
    co}.}
\begin{verbatim}
svn checkout URL [/pathname]
\end{verbatim}
che creerà nella directory corrente una directory corrispondente al nome
specificato in coda alla URL passata come argomento, scaricando l'ultima
versione dei file archiviati sul repository; alternativamente si può
specificare come ulteriore argomento la directory su cui scaricare i file.

Sia in caso di \texttt{import} che di \texttt{checkout} è sempre possibile
operare su una qualunque sotto cartella contenuta all'interno di un
\textit{repository}, ignorando totalmente quello che sta al di sopra, basterà
indicare in sede di importazione o di estrazione iniziale un pathname o una
URL che identifichi quella parte del progetto.

Se quando si effettua lo scaricamento non si vuole usare la versione più
aggiornata, ma una versione precedente si può usare l'opzione \texttt{-r}
seguita da un numero che scaricherà esattamente quella \textit{release},
alternativamente al posto del numero si può indicare una data, e verrà presa
la \textit{release} più prossima a quella data. 

A differenza di CVS Subversion non supporta l'uso di \textsl{etichette}
associate ad una certa versione del proprio progetto, per questo è invalso
l'uso di strutturare il repository secondo lo schema illustrato inizialmente;
è infatti molto semplice (e non comporta nessun tipo di aggravio) creare delle
copie complete di una qualunque parte del \textit{repository} su un'altra
parte dello stesso, per cui se si è eseguito lo sviluppo sulla cartella
\texttt{trunk} sarà possibile creare banalmente una versione con etichetta
\texttt{label} (o quel che si preferisce) semplicemente con una copia eseguita
con:
\begin{verbatim}
svn cp trunk tags/label
\end{verbatim}

Il risultato di questo comando è la creazione della nuova cartella
\texttt{label} sotto \texttt{tags}, che sarà assolutamente identica, nel
contenuto (e nella sua storia) a quanto presente in \texttt{trunk} al momento
dell'esecuzione del comando. In questo modo, una volta salvate le
modifiche,\footnote{la copia viene eseguita localmente verrà creata anche sul
  \textit{repository} solo dopo un \textit{commit}.} si potrà ottenere la
versione \texttt{label} del proprio progetto semplicemente eseguendo un
\textit{checkout} di \texttt{tags/label} in un'altra
directory.\footnote{ovviamente una volta presa la suddetta versione si deve
  aver cura di non eseguire nessuna modifica a partire dalla stessa, per
  questo se si deve modificare una versione etichettata si usa
  \texttt{branches}.}

Una volta creata la propria copia locale dei programmi, è possibile lavorare
su di essi ponendosi nella relativa directory, e apportare tutte le modifiche
che si vogliono ai file ivi presenti; due comandi permettono inoltre di
schedulare la rimozione o l'aggiunta di file al
\textit{repository}:\footnote{a differenza di CVS si possono aggiungere e
  rimuovere, ed anche spostare con \texttt{svn mv}, sia file che directory.}
\begin{verbatim}
svn add file1.c
svn remove file2.c
\end{verbatim}
ma niente viene modificato sul \textit{repository} fintanto che non viene
eseguito il cosiddetto \textit{commit} delle modifiche, vale a dire fintanto
che non viene dato il comando:\footnote{in genere anche questo viene
  abbreviato, con \texttt{svn ci}.}
\begin{verbatim}
svn commit [file]
\end{verbatim}
ed è possibile eseguire il \textit{commit} delle modifiche per un singolo
file, indicandolo come ulteriore argomento, mentre se non si indica nulla
verranno inviate tutte le modifiche presenti.

Si tenga presente però che il \textit{commit} non verrà eseguito se nel
frattempo i file del \textit{repository} sono stati modificati; in questo caso
\texttt{svn} rileverà la presenza di differenze fra la propria
\textit{release} e quella del \textit{repository} e chiederà che si effettui
preventivamente un aggiornamento. Questa è una delle operazioni di base di
Subversion, che in genere si compie tutte le volte che si inizia a lavorare,
il comando che la esegue è:
\begin{verbatim}
svn update
\end{verbatim}

Questo comando opera a partire dalla directory in cui viene eseguito e ne
aggiorna il contenuto (compreso quello di eventuali sotto-directory) alla
versione presente, scaricando le ultime versioni dei file esistenti o nuovi
file o directory aggiunte, cancellando eventuali file e directory rimossi dal
\textit{repository}.  Esso inoltre esso cerca, in caso di presenza di
modifiche eseguite in maniera indipendente sulla propria copia locale, di
eseguire un \textsl{raccordo} (il cosiddetto \textit{merging}) delle stesse
con quelle presenti sulla versione del \textit{repository}.

Fintanto che sono state modificate parti indipendenti di un file di testo in
genere il processo di \textit{merging} ha successo e le modifiche vengono
incorporate automaticamente in conseguenza dell'aggiornamento, ma quando le
modifiche attengono alla stessa parte di un file nel ci si troverà di fronte
ad un conflitto ed a quel punto sarà richiesto al ``committente'' di
intervenire manualmente sui file per i quali sono stati rilevati i conflitti
per risolverli.

Per aiutare il committente nel suo compito quando l'operazione di
aggiornamento fallisce nel raccordo delle modifiche lascia sezioni di codice
in conflitto opportunamente marcate e separate fra loro come nell'esempio
seguente:
\begin{verbatim}
<<<<<<< .mine
        $(CC) $(CFLAGS) -o pamacq pamacq.c -lm
=======
        $(CC) $(CFLAGS) -o pamacq pamacq.c

>>>>>>> r.122
\end{verbatim}

In questo caso si c'è stata una modifica sul file (mostrata nella parte
superiore) incompatibile con quella fatta nel \textit{repository} (mostrata
nella parte inferiore). Prima di eseguire un \textit{commit} occorrerà
pertanto integrare le modifiche e salvare nuovamente il file rimuovendo i
marcatori, inoltre prima che il \textit{commit} ritorni possibile si dovrà
esplicitare la risoluzione del conflitto con il comando:
\begin{verbatim}
svn resolved file
\end{verbatim}

\begin{table}[htb]
  \centering
  \footnotesize
  \begin{tabular}[c]{|c|l|}
    \hline
    \textbf{Flag} & \textbf{Significato} \\
    \hline
    \hline
    \texttt{?}& File sconosciuto.\\
    \texttt{M}& File modificato localmente.\\
    \texttt{A}& File aggiunto.\\
    \texttt{C}& File con conflitto.\\
%    \const{G}& File con modifiche integrate.\\
%    \const{}& .\\
    \hline
  \end{tabular}
  \caption{Caratteri associati ai vari stati dei file.}
  \label{tab:svn_status}
\end{table}


Infine per capire la situazione della propria copia locale si può utilizzare
il comando \texttt{svn status} che confronta i file presenti nella directory
locale rispetto alla ultima versione scaricata dal \textit{repository} e per
tutti quelli che non corrispondono stampa a schermo delle informazioni di
stato nella forma di un carattere seguito dal nome del file, secondo quanto
illustrato in tab.~\ref{tab:svn_status}.




% LocalWords:  make Makefile makefile shell FC FLAGS CFLAGS pathname CERN SCM
% LocalWords:  OBJ commondef readfile FFLAGS Cuncurrent Version System CVS home
% LocalWords:  adidsp geometry muonacq andaranno CVSROOT startup cvs ssh rsh to
% LocalWords:  project tag Root commit update merging Locally Modified Added co
% LocalWords:  Needs Patch Merge log rtag checkout module altrobersaglio Source
% LocalWords:   Control Subversion repository DAV svnserve URL svn subversion
% LocalWords:  client sez svnadmin Truelite Srl trunk tags branches tab add rm
% LocalWords:  remove copy cp move resolved label



%%% Local Variables: 
%%% mode: latex
%%% TeX-master: "gapil"
%%% End: 

%% ringraziamenti.tex
%%
%% Copyright (C) 2000-2018 Simone Piccardi.  Permission is granted to
%% copy, distribute and/or modify this document under the terms of the GNU Free
%% Documentation License, Version 1.1 or any later version published by the
%% Free Software Foundation; with the Invariant Sections being "Un preambolo",
%% with no Front-Cover Texts, and with no Back-Cover Texts.  A copy of the
%% license is included in the section entitled "GNU Free Documentation
%% License".
%%
\chapter{Ringraziamenti}
\label{cha:ringraziamenti}

Desidero ringraziare tutti coloro che a vario titolo e a più riprese mi hanno
aiutato ed han contribuito a migliorare in molteplici aspetti la qualità di
GaPiL. In ordine rigorosamente alfabetico desidero citare:
\begin{description}
\item[\textbf{Alessio Frusciante}] per l'apprezzamento, le innumerevoli
  correzioni ed i suggerimenti per rendere più chiara l'esposizione.
\item[\textbf{Daniele Masini}] per la rilettura puntuale, le innumerevoli
  correzioni, i consigli sull'esposizione ed i contributi relativi alle
  calling convention dei linguaggi e al confronto delle diverse tecniche di
  gestione della memoria.
\item[\textbf{Mirko Maischberger}] per la rilettura, le numerose correzioni,
  la segnalazione dei passi poco chiari e soprattutto per il grande lavoro
  svolto per produrre una versione della guida in un HTML piacevole ed
  accurato.
\item[\textbf{Fabio Rossi}] per la rilettura, le innumerevoli correzioni, ed i
  vari consigli stilistici ed i suggerimenti per il miglioramento della
  comprensione di vari passaggi.
\end{description}

Infine, vorrei ringraziare il Firenze Linux User Group (FLUG), di cui mi
pregio di fare parte, che ha messo a disposizione il repository CVS su cui era
presente la prima versione della Guida, ed il relativo spazio web, e
\href{https://www.truelite.it}{Truelite Srl}, l'azienda che ho fondato e di
cui sono amministratore, che fornisce le risorse per il nuovo repository Git
ed per il sito del progetto: \url{https://gapil.gnulinux.it}.


% LocalWords:  GaPiL Masini calling convention Maischberger HTML Group FLUG CVS
% LocalWords:  repository Truelite Srl SVN


%%% Local Variables: 
%%% mode: latex
%%% TeX-master: "gapil"
%%% End: 


\include{fdl-1.3}

% at the end put the bibliography
\backmatter
\printindex
%\bibliographystyle{phaip}
%\bibliographystyle{ieeetr}
\bibliographystyle{abstract}
\bibliography{biblio}

\end{document}

%%% Local Variables:
%%% mode: latex
%%% TeX-master: t
%%% End:
